\IfFileExists{stacks-project-book.cls}{%
\documentclass{stacks-project-book}
}{%
\documentclass{amsbook}
}

\usepackage{verbatim}
\newenvironment{reference}{\comment}{\endcomment}
\newenvironment{slogan}{\comment}{\endcomment}
\newenvironment{history}{\comment}{\endcomment}

\usepackage[all]{xy}

\xyoption{2cell}
\UseAllTwocells

\usepackage{multicol}

\usepackage{lmodern}
\usepackage[T1]{fontenc}
\usepackage[french]{babel}
\AtBeginDocument{\shorthandoff{?}}


\usepackage{hyperref}


\theoremstyle{plain}
\newtheorem{theorem}[subsection]{Théorème}
\newtheorem{proposition}[subsection]{Proposition}
\newtheorem{lemma}[subsection]{Lemme}

\theoremstyle{definition}
\newtheorem{definition}[subsection]{Définition}
\newtheorem{example}[subsection]{Exemple}
\newtheorem{exercise}[subsection]{Exercice}
\newtheorem{situation}[subsection]{Situation}

\theoremstyle{remark}
\newtheorem{remark}[subsection]{Remarque}
\newtheorem{remarks}[subsection]{Remarques}

\renewcommand{\proofname}{Démonstration}

\numberwithin{equation}{subsection}

\def\lim{\mathop{\mathrm{lim}}\nolimits}
\def\colim{\mathop{\mathrm{colim}}\nolimits}
\def\Spec{\mathop{\mathrm{Spec}}}
\def\Hom{\mathop{\mathrm{Hom}}\nolimits}
\def\Ext{\mathop{\mathrm{Ext}}\nolimits}
\def\SheafHom{\mathop{\mathcal{H}\!\mathit{om}}\nolimits}
\def\SheafExt{\mathop{\mathcal{E}\!\mathit{xt}}\nolimits}
\def\Sch{\mathit{Sch}}
\def\Mor{\mathop{\mathrm{Mor}}\nolimits}
\def\Ob{\mathop{\mathrm{Ob}}\nolimits}
\def\Sh{\mathop{\mathit{Sh}}\nolimits}
\def\NL{\mathop{N\!L}\nolimits}
\def\CH{\mathop{\mathrm{CH}}\nolimits}
\def\proetale{{pro\text{-}\acute{e}tale}}
\def\etale{{\acute{e}tale}}
\def\QCoh{\mathit{QCoh}}
\def\Ker{\mathop{\mathrm{Ker}}}
\def\Im{\mathop{\mathrm{Im}}}
\def\Coker{\mathop{\mathrm{Coker}}}
\def\Coim{\mathop{\mathrm{Coim}}}

\DeclareMathSymbol{\boxtimes}{\mathbin}{AMSa}{"02}

\def\QCohstack{\mathcal{QC}\!\mathit{oh}}
\def\Cohstack{\mathcal{C}\!\mathit{oh}}
\def\Spacesstack{\mathcal{S}\!\mathit{paces}}
\def\Quotfunctor{\mathrm{Quot}}
\def\Hilbfunctor{\mathrm{Hilb}}
\def\Curvesstack{\mathcal{C}\!\mathit{urves}}
\def\Polarizedstack{\mathcal{P}\!\mathit{olarized}}
\def\Complexesstack{\mathcal{C}\!\mathit{omplexes}}
\def\Pic{\mathop{\mathrm{Pic}}\nolimits}
\def\Picardstack{\mathcal{P}\!\mathit{ic}}
\def\Picardfunctor{\mathrm{Pic}}
\def\Deformationcategory{\mathcal{D}\!\mathit{ef}}
\begin{document}
\title{Le projet Stacks --- édition française, partie 6}
\maketitle
\tableofcontents

% OK, start here.
%

\chapter{Schémas en groupoïdes}



\phantomsection
\label{groupoids-section-phantom}


\section{Introduction}
\label{groupoids-section-introduction}

\noindent
Ce chapitre est consacré aux généralités sur les schémas en groupoïdes.
Voir par exemple le bel article \cite{K-M} de Keel et Mori.





\section{Notations}
\label{groupoids-section-notation}

\noindent
Soit $S$ un schéma. Si $U$ et $T$ sont des schémas sur $S$, nous notons
$U(T)$ l'ensemble des points à valeurs dans $T$ du schéma $U$ {\it au-dessus de} $S$. En formule :
$U(T) = \Mor_S(T, U)$. Nous nous efforçons de réserver la lettre $T$ pour désigner
un ``schéma test'' sur $S$, comme dans la discussion qui suit.
Soient $X$ et $Y$ des schémas sur
$S$, et soit $f : X \to Y$ un morphisme de schémas sur $S$.
Pour tout schéma $T$ sur $S$, on obtient une application induite entre ensembles
$$
f : X(T) \longrightarrow Y(T)
$$
que, comme indiqué, nous notons aussi $f$. En fait, cette construction
est fonctorielle en $T/S$. Le lemme de Yoneda, voir Catégories,
lemme \ref{categories-lemma-yoneda}, affirme que $f$ détermine et est
déterminé par cette transformation de foncteurs $f : h_X \to h_Y$.
Plus généralement, nous employons la même notation pour les applications entre produits
fibrés. Par exemple, si
$X$, $Y$ et $Z$ sont des schémas sur $S$, et si
$m : X \times_S Y \to Z \times_S Z$ est
un morphisme de schémas sur $S$, alors nous considérons $m$ comme correspondant
à une famille d'applications entre les ensembles de points à valeurs dans $T$
$$
X(T) \times Y(T) \longrightarrow Z(T) \times Z(T).
$$
Et ainsi de suite.

\medskip\noindent
Nous conservons notre convention consistant à indexer les morphismes de projection à partir de
l'indice $0$; on a donc $\text{pr}_0 : X \times_S Y \to X$ et
$\text{pr}_1 : X \times_S Y \to Y$.






\section{Relations d'équivalence}
\label{groupoids-section-equivalence-relations}

\noindent
Rappelons qu'une {\it relation} $R$ sur un ensemble $A$ est simplement un sous-ensemble
$R \subset A \times A$. On écrit habituellement $a R b$ pour indiquer que
$(a, b) \in R$. On dit que la relation est {\it transitive} si
$a R b, b R c \Rightarrow a R c$. On dit que la relation est
{\it réflexive} si $a R a$ pour tout $a \in A$. On dit que la relation est
{\it symétrique} si $a R b \Rightarrow b R a$.
Une relation est appelée {\it relation d'équivalence} si
elle est transitive, réflexive et symétrique.

\medskip\noindent
Dans le cadre des schémas, nous allons assouplir quelque peu la notion de
relation et exiger seulement que $R \to A \times A$ soit
un morphisme. Voici la définition.

\begin{definition}
\label{groupoids-definition-equivalence-relation}
Soit $S$ un schéma. Soit $U$ un schéma sur $S$.
\begin{enumerate}
\item Une {\it prérelation} sur $U$ relativement à $S$ est un morphisme quelconque
de schémas $j : R \to U \times_S U$. Dans ce cas, on pose
$t = \text{pr}_0 \circ j$ et $s = \text{pr}_1 \circ j$, de sorte
que $j = (t, s)$.
\item Une {\it relation} sur $U$ relativement à $S$ est un monomorphisme
de schémas $j : R \to U \times_S U$.
\item Une {\it prérelation d'équivalence} est une prérelation
$j : R \to U \times_S U$ telle que l'image de
$j : R(T) \to U(T) \times U(T)$ est une relation d'équivalence pour
tout $T/S$.
\item On dit qu'un morphisme de schémas $R \to U \times_S U$ est
une {\it relation d'équivalence sur $U$ relativement à $S$}
si et seulement si, pour tout schéma $T$ sur $S$, les points à valeurs dans $T$
du schéma $R$ définissent une relation d'équivalence
sur l'ensemble des points à valeurs dans $T$ du schéma $U$.
\end{enumerate}
\end{definition}

\noindent
Autrement dit, une relation d'équivalence est une prérelation d'équivalence
telle que $j$ soit une relation.

\begin{lemma}
\label{groupoids-lemma-restrict-relation}
Soit $S$ un schéma.
Soit $U$ un schéma sur $S$.
Soit $j : R \to U \times_S U$ une prérelation.
Soit $g : U' \to U$ un morphisme de schémas.
Enfin, posons
$$
R' = (U' \times_S U')\times_{U \times_S U} R
\xrightarrow{j'}
U' \times_S U'
$$
Alors $j'$ est une prérelation sur $U'$ relativement à $S$.
Si $j$ est une relation, alors $j'$ est une relation.
Si $j$ est une prérelation d'équivalence, alors $j'$ est une prérelation d'équivalence.
Si $j$ est une relation d'équivalence, alors $j'$ est une relation d'équivalence.
\end{lemma}

\begin{proof}
Démonstration omise.
\end{proof}

\begin{definition}
\label{groupoids-definition-restrict-relation}
Soit $S$ un schéma.
Soit $U$ un schéma sur $S$.
Soit $j : R \to U \times_S U$ une prérelation.
Soit $g : U' \to U$ un morphisme de schémas.
La prérelation $j' : R' \to U' \times_S U'$ est appelée
la {\it restriction}, ou l'{\it image réciproque}, de la prérelation $j$ sur $U'$.
Dans cette situation, nous écrivons parfois $R' = R|_{U'}$.
\end{definition}

\begin{lemma}
\label{groupoids-lemma-pre-equivalence-equivalence-relation-points}
Soit $j : R \to U \times_S U$ une prérelation.
Considérons la relation sur les points du schéma $U$ définie par
la règle
$$
x \sim y
\Leftrightarrow
\exists\ r \in R :
t(r) = x,
s(r) = y.
$$
Si $j$ est une prérelation d'équivalence, alors la relation ainsi définie est une
relation d'équivalence.
\end{lemma}

\begin{proof}
Supposons que $x \sim y$ et que $y \sim z$.
Choisissons $r \in R$ tel que $t(r) = x$, $s(r) = y$ et
choisissons $r' \in R$ tel que $t(r') = y$, $s(r') = z$.
Choisissons un corps $K$ s'insérant dans le diagramme commutatif
suivant
$$
\xymatrix{
\kappa(r) \ar[r] & K \\
\kappa(y) \ar[u] \ar[r] & \kappa(r') \ar[u]
}
$$
Désignons par $x_K, y_K, z_K : \Spec(K) \to U$
les morphismes
$$
\begin{matrix}
\Spec(K) \to \Spec(\kappa(r))
\to
\Spec(\kappa(x)) \to U \\
\Spec(K) \to \Spec(\kappa(r))
\to
\Spec(\kappa(y)) \to U \\
\Spec(K) \to \Spec(\kappa(r'))
\to
\Spec(\kappa(z)) \to U
\end{matrix}
$$
Par construction, $(x_K, y_K) \in j(R(K))$ et
$(y_K, z_K) \in j(R(K))$. Comme $j$ est une prérelation d'équivalence,
on a aussi $(x_K, z_K) \in j(R(K))$.
Cela implique clairement que $x \sim z$.

\medskip\noindent
La démonstration du fait que $\sim$ est réflexive et symétrique est omise.
\end{proof}

\begin{lemma}
\label{groupoids-lemma-etale-equivalence-relation}
Soit $j : R \to U \times_S U$ une prérelation. Supposons que
\begin{enumerate}
\item $s, t$ soient non ramifiés,
\item pour tout corps algébriquement clos $k$ sur $S$
l'application $R(k) \to U(k) \times U(k)$ définisse une relation d'équivalence,
\item il existe des morphismes $e : U \to R$, $i : R \to R$,
$c : R \times_{s, U, t} R \to R$ tels que
$$
\xymatrix@C-1pc{
U \ar[r]_e \ar[d]_\Delta &
R \ar[d]_j &
R \ar[d]^j \ar[r]_i &
R \ar[d]^j &
R \times_{s, U, t} R \ar[d]^{j \times j} \ar[r]_c &
R \ar[d]^j \\
U \times_S U \ar[r] &
U \times_S U &
U \times_S U \ar[r]^{flip} &
U \times_S U &
U \times_S U \times_S U \ar[r]^{\text{pr}_{02}} &
U \times_S U
}
$$
sont commutatifs.
\end{enumerate}
Alors $j$ est une relation d'équivalence.
\end{lemma}

\begin{proof}
D'après la condition (1) et le
lemme \ref{morphisms-lemma-unramified-permanence} de Morphismes,
on voit que $j$ est non ramifié. Ainsi,
$\Delta_j : R \to R \times_{U \times_S U} R$ est une immersion ouverte d'après
le lemme \ref{morphisms-lemma-diagonal-unramified-morphism} de Morphismes.
Or la condition (2) affirme que $\Delta_j$ est bijective sur les
points à valeurs dans $k$, donc $\Delta_j$ est un isomorphisme, donc $j$
est un monomorphisme. Il découle alors aisément des diagrammes
commutatifs que $R(T) \subset U(T) \times U(T)$ est une relation
d'équivalence pour tout schéma $T$ sur $S$.
\end{proof}













\section{Schémas en groupes}
\label{groupoids-section-group-schemes}

\noindent
Rappelons qu'un {\it groupe} est un couple
$(G, m)$, où $G$ est un ensemble et $m : G \times G \to G$ est
une application d'ensembles possédant les propriétés suivantes :
\begin{enumerate}
\item (associativité) $m(g, m(g', g'')) = m(m(g, g'), g'')$
pour tous $g, g', g'' \in G$,
\item (élément neutre) il existe un unique élément $e \in G$
(appelé {\it élément neutre}, {\it unité}, ou $1$ de $G$) tel que
$m(g, e) = m(e, g) = g$ pour tout $g \in G$, et
\item (inverse) pour tout $g \in G$ il existe un élément $i(g) \in G$
tel que $m(g, i(g)) = m(i(g), g) = e$, où $e$ est
l'élément neutre.
\end{enumerate}
Nous obtenons ainsi une application $e : \{*\} \to G$ et une application
$i : G \to G$ de sorte que le quadruplet $(G, m, e, i)$
satisfait les axiomes énumérés ci-dessus.

\medskip\noindent
Un {\it homomorphisme de groupes} $\psi : (G, m) \to (G', m')$
est une application d'ensembles $\psi : G \to G'$ telle que
$m'(\psi(g), \psi(g')) = \psi(m(g, g'))$. Il s'ensuit automatiquement
que $\psi(e) = e'$ et $i'(\psi(g)) = \psi(i(g))$.
(Notations évidentes.) Nous nous en servirons ci-dessous.

\begin{definition}
\label{groupoids-definition-group-scheme}
Soit $S$ un schéma.
\begin{enumerate}
\item Un {\it schéma en groupes sur $S$} est un couple $(G, m)$, où
$G$ est un schéma sur $S$ et $m : G \times_S G \to G$ est
un morphisme de schémas sur $S$ vérifiant la propriété suivante :
Pour tout schéma $T$ sur $S$, le couple $(G(T), m)$
est un groupe.
\item Un {\it morphisme $\psi : (G, m) \to (G', m')$ de schémas en groupes sur $S$}
est un morphisme $\psi : G \to G'$ de schémas sur $S$ tel que, pour
tout $T/S$, l'application induite $\psi : G(T) \to G'(T)$ soit un homomorphisme
de groupes.
\end{enumerate}
\end{definition}

\noindent
Soit $(G, m)$ un schéma en groupes sur le schéma $S$.
D'après la discussion ci-dessus (et celle de la section \ref{groupoids-section-notation})
nous obtenons des morphismes de schémas sur $S$ :
(élément neutre) $e : S \to G$ et (inverse) $i : G \to G$ tels que
pour tout $T$, le quadruplet $(G(T), m, e, i)$ satisfait les
axiomes d'un groupe énumérés ci-dessus.

\medskip\noindent
Soient $(G, m)$ et $(G', m')$ des schémas en groupes sur $S$.
Soit $f : G \to G'$ un morphisme de schémas sur $S$.
Il résulte de la définition que $f$ est un morphisme
de schémas en groupes sur $S$ si et seulement si le diagramme suivant
est commutatif :
$$
\xymatrix{
G \times_S G \ar[r]_-{f \times f} \ar[d]_m &
G' \times_S G' \ar[d]^{m'} \\
G \ar[r]^f & G'
}
$$

\begin{lemma}
\label{groupoids-lemma-base-change-group-scheme}
Soit $(G, m)$ un schéma en groupes sur $S$.
Soit $S' \to S$ un morphisme de schémas.
Le changement de base $(G_{S'}, m_{S'})$ est un schéma en groupes sur $S'$.
\end{lemma}

\begin{proof}
La démonstration est omise.
\end{proof}

\begin{definition}
\label{groupoids-definition-closed-subgroup-scheme}
Soit $S$ un schéma. Soit $(G, m)$ un schéma en groupes sur $S$.
\begin{enumerate}
\item Un {\it sous-schéma en groupes fermé} de $G$ est un sous-schéma fermé
$H \subset G$ tel que $m|_{H \times_S H}$ se factorise par $H$ et induise
sur $H$ une structure de schéma en groupes sur $S$.
\item Un {\it sous-schéma en groupes ouvert} de $G$ est un sous-schéma ouvert
$G' \subset G$ tel que $m|_{G' \times_S G'}$ se factorise par $G'$
et induise sur $G'$ une structure de schéma en groupes sur $S$.
\end{enumerate}
\end{definition}

\noindent
On pourrait également dire que $H$ est un sous-schéma en groupes fermé de $G$
si c'est un schéma en groupes sur $S$ muni d'un morphisme de schémas en groupes
$i : H \to G$ sur $S$ qui identifie $H$ à un sous-schéma fermé de $G$.

\begin{lemma}
\label{groupoids-lemma-closed-subgroup-scheme}
Soit $S$ un schéma. Soit $(G, m, e, i)$ un schéma en groupes sur $S$.
\begin{enumerate}
\item Un sous-schéma fermé $H \subset G$ est un sous-schéma en groupes fermé
si et seulement si $e : S \to G$, $m|_{H \times_S H} : H \times_S H \to G$,
et $i|_H : H \to G$ se factorisent par $H$.
\item Un sous-schéma ouvert $H \subset G$ est un sous-schéma en groupes ouvert
si et seulement si $e : S \to G$, $m|_{H \times_S H} : H \times_S H \to G$,
et $i|_H : H \to G$ se factorisent par $H$.
\end{enumerate}
\end{lemma}

\begin{proof}
En considérant les points à valeurs dans $T$, cela se traduit par les conditions bien connues
caractérisant les parties d'un groupe qui en sont des sous-groupes.
\end{proof}

\begin{definition}
\label{groupoids-definition-smooth-group-scheme}
Soit $S$ un schéma. Soit $(G, m)$ un schéma en groupes sur $S$.
\begin{enumerate}
\item On dit que $G$ est un {\it schéma en groupes lisse} si son morphisme
structural $G \to S$ est lisse.
\item On dit que $G$ est un {\it schéma en groupes plat} si son morphisme
structural $G \to S$ est plat.
\item On dit que $G$ est un {\it schéma en groupes séparé} si son morphisme
structural $G \to S$ est séparé.
\end{enumerate}
On pourra compléter cette liste au besoin.
\end{definition}






\section{Exemples de schémas en groupes}
\label{groupoids-section-examples-group-schemes}

\begin{example}[Schéma en groupes multiplicatif]
\label{groupoids-example-multiplicative-group}
Considérons le foncteur qui associe
à tout schéma $T$ le groupe $\Gamma(T, \mathcal{O}_T^*)$
des unités de l'anneau des sections globales du faisceau structural.
Ce foncteur est représentable par le schéma
$$
\mathbf{G}_m = \Spec(\mathbf{Z}[x, x^{-1}])
$$
Le morphisme définissant la structure de groupe est le morphisme
\begin{eqnarray*}
\mathbf{G}_m \times \mathbf{G}_m & \to & \mathbf{G}_m \\
\Spec(\mathbf{Z}[x, x^{-1}] \otimes_{\mathbf{Z}} \mathbf{Z}[x, x^{-1}])
& \to &
\Spec(\mathbf{Z}[x, x^{-1}]) \\
\mathbf{Z}[x, x^{-1}] \otimes_{\mathbf{Z}} \mathbf{Z}[x, x^{-1}]
& \leftarrow &
\mathbf{Z}[x, x^{-1}] \\
x \otimes x & \leftarrow & x
\end{eqnarray*}
On voit donc que $\mathbf{G}_m$ est un schéma en groupes sur $\mathbf{Z}$.
Pour tout schéma $S$, le changement de base $\mathbf{G}_{m, S}$ est un
schéma en groupes sur $S$ dont le foncteur des points est
$$
T/S
\longmapsto
\mathbf{G}_{m, S}(T) = \mathbf{G}_m(T) = \Gamma(T, \mathcal{O}_T^*)
$$
comme précédemment.
\end{example}

\begin{example}[Racines de l'unité]
\label{groupoids-example-roots-of-unity}
Soit $n \in \mathbf{N}$.
Considérons le foncteur qui associe
à tout schéma $T$ le sous-groupe de $\Gamma(T, \mathcal{O}_T^*)$
formé des racines $n$-ièmes de l'unité.
Ce foncteur est représentable par le schéma
$$
\mu_n = \Spec(\mathbf{Z}[x]/(x^n - 1)).
$$
Le morphisme définissant la structure de groupe est le morphisme
\begin{eqnarray*}
\mu_n \times \mu_n & \to & \mu_n \\
\Spec(
\mathbf{Z}[x]/(x^n - 1)
\otimes_{\mathbf{Z}}
\mathbf{Z}[x]/(x^n - 1))
& \to &
\Spec(\mathbf{Z}[x]/(x^n - 1)) \\
\mathbf{Z}[x]/(x^n - 1) \otimes_{\mathbf{Z}} \mathbf{Z}[x]/(x^n - 1)
& \leftarrow &
\mathbf{Z}[x]/(x^n - 1) \\
x \otimes x & \leftarrow & x
\end{eqnarray*}
On voit donc que $\mu_n$ est un schéma en groupes sur $\mathbf{Z}$.
Pour tout schéma $S$, le changement de base $\mu_{n, S}$ est un
schéma en groupes sur $S$ dont le foncteur des points est
$$
T/S
\longmapsto
\mu_{n, S}(T) = \mu_n(T) = \{f \in \Gamma(T, \mathcal{O}_T^*) \mid f^n = 1\}
$$
comme précédemment.
\end{example}


\begin{example}[Schéma en groupes additif]
\label{groupoids-example-additive-group}
Considérons le foncteur qui associe
à tout schéma $T$ le groupe $\Gamma(T, \mathcal{O}_T)$
des sections globales du faisceau structural.
Ce foncteur est représentable par le schéma
$$
\mathbf{G}_a = \Spec(\mathbf{Z}[x])
$$
Le morphisme définissant la structure de groupe est le morphisme
\begin{eqnarray*}
\mathbf{G}_a \times \mathbf{G}_a & \to & \mathbf{G}_a \\
\Spec(\mathbf{Z}[x] \otimes_{\mathbf{Z}} \mathbf{Z}[x])
& \to &
\Spec(\mathbf{Z}[x]) \\
\mathbf{Z}[x] \otimes_{\mathbf{Z}} \mathbf{Z}[x]
& \leftarrow &
\mathbf{Z}[x] \\
x \otimes 1 + 1 \otimes x & \leftarrow & x
\end{eqnarray*}
On voit donc que $\mathbf{G}_a$ est un schéma en groupes sur $\mathbf{Z}$.
Pour tout schéma $S$, le changement de base $\mathbf{G}_{a, S}$ est un
schéma en groupes sur $S$ dont le foncteur des points est
$$
T/S
\longmapsto
\mathbf{G}_{a, S}(T) = \mathbf{G}_a(T) = \Gamma(T, \mathcal{O}_T)
$$
comme précédemment.
\end{example}

\begin{example}[Schéma en groupes linéaire général]
\label{groupoids-example-general-linear-group}
Soit $n \geq 1$.
Considérons le foncteur qui associe
à tout schéma $T$ le groupe
$$
\text{GL}_n(\Gamma(T, \mathcal{O}_T))
$$
des matrices inversibles $n \times n$ à coefficients dans
l'anneau des sections globales du faisceau structural.
Ce foncteur est représentable par le schéma
$$
\text{GL}_n = \Spec(\mathbf{Z}[\{x_{ij}\}_{1 \leq i, j \leq n}][1/d])
$$
où $d = \det((x_{ij}))$ et où $(x_{ij})$ est la matrice $n \times n$
ayant pour coefficient $x_{ij}$ en position $(i, j)$.
Le morphisme définissant la structure de groupe est le morphisme
\begin{eqnarray*}
\text{GL}_n \times \text{GL}_n & \to & \text{GL}_n \\
\Spec(\mathbf{Z}[x_{ij}, 1/d] \otimes_{\mathbf{Z}}
\mathbf{Z}[x_{ij}, 1/d])
& \to &
\Spec(\mathbf{Z}[x_{ij}, 1/d]) \\
\mathbf{Z}[x_{ij}, 1/d] \otimes_{\mathbf{Z}} \mathbf{Z}[x_{ij}, 1/d]
& \leftarrow &
\mathbf{Z}[x_{ij}, 1/d] \\
\sum x_{ik} \otimes x_{kj} & \leftarrow & x_{ij}
\end{eqnarray*}
Ainsi, on voit que $\text{GL}_n$ est un schéma en groupes sur $\mathbf{Z}$.
Pour tout schéma $S$, le changement de base $\text{GL}_{n, S}$ est un
schéma en groupes sur $S$ dont le foncteur des points est donné par
$$
T/S
\longmapsto
\text{GL}_{n, S}(T) = \text{GL}_n(T) =\text{GL}_n(\Gamma(T, \mathcal{O}_T))
$$
comme précédemment.
\end{example}

\begin{example}
\label{groupoids-example-determinant}
Le déterminant définit un morphisme de schémas en groupes
$$
\det : \text{GL}_n \longrightarrow \mathbf{G}_m
$$
sur $\mathbf{Z}$. Par changement de base, il donne un morphisme
de schémas en groupes $\text{GL}_{n, S} \to \mathbf{G}_{m, S}$
sur tout schéma de base $S$.
\end{example}

\begin{example}[Groupe constant]
\label{groupoids-example-constant-group}
Soit $G$ un groupe abstrait. Considérons le foncteur
qui associe à tout schéma $T$ le groupe
des applications localement constantes $T \to G$ (où $T$ est muni de la topologie de Zariski
et $G$ de la topologie discrète). Ce foncteur est représentable par le schéma
$$
G_{\Spec(\mathbf{Z})} =
\coprod\nolimits_{g \in G} \Spec(\mathbf{Z}).
$$
Le morphisme définissant la structure de groupe est le morphisme
$$
G_{\Spec(\mathbf{Z})}
\times_{\Spec(\mathbf{Z})}
G_{\Spec(\mathbf{Z})}
\longrightarrow
G_{\Spec(\mathbf{Z})}
$$
qui envoie la composante correspondant au couple $(g, g')$ sur la
composante correspondant à $gg'$. Pour tout schéma $S$, le changement de base
$G_S$ est un schéma en groupes sur $S$ dont le foncteur des points est donné par
$$
T/S
\longmapsto
G_S(T) = \{f : T \to G \text{ localement constante}\}
$$
comme précédemment.
\end{example}





\section{Propriétés des schémas en groupes}
\label{groupoids-section-properties-group-schemes}

\noindent
Dans cette section, nous rassemblons quelques propriétés élémentaires des schémas en groupes qui
sont valables sur toute base.

\begin{lemma}
\label{groupoids-lemma-group-scheme-separated}
Soit $S$ un schéma.
Soit $G$ un schéma en groupes sur $S$.
Alors $G \to S$ est séparé (resp.\ quasi-séparé) si et seulement si
le morphisme unité $e : S \to G$ est une immersion fermée
(resp.\ quasi-compact).
\end{lemma}

\begin{proof}
Rappelons que, d'après le
lemme \ref{schemes-lemma-section-immersion} de Schémas,
$e$ est une immersion, laquelle est fermée
(resp.\ quasi-compacte) si $G \to S$ est séparé (resp.\ quasi-séparé).
Pour la réciproque, considérons le diagramme
$$
\xymatrix{
G \ar[r]_-{\Delta_{G/S}} \ar[d] &
G \times_S G \ar[d]^{(g, g') \mapsto m(i(g), g')} \\
S \ar[r]^e & G
}
$$
Montrer que ce diagramme est cartésien est un exercice d'application du point de vue
fonctoriel en géométrie algébrique. Autrement dit, on voit que
$\Delta_{G/S}$ s'obtient à partir de $e$ par changement de base. Ainsi, si $e$ est une
immersion fermée (resp.\ quasi-compact), il en est de même de $\Delta_{G/S}$, voir
Schémas, lemme \ref{schemes-lemma-base-change-immersion}
(resp.\ Schémas, Lemme
\ref{schemes-lemma-quasi-compact-preserved-base-change}).
\end{proof}

\begin{lemma}
\label{groupoids-lemma-flat-action-on-group-scheme}
Soit $S$ un schéma.
Soit $G$ un schéma en groupes sur $S$.
Soit $T$ un schéma sur $S$ et soit $\psi : T \to G$ un morphisme sur $S$.
Si $T$ est plat sur $S$, alors le morphisme
$$
T \times_S G \longrightarrow G, \quad
(t, g) \longmapsto m(\psi(t), g)
$$
est plat. En particulier, si $G$ est plat sur $S$, alors
$m : G \times_S G \to G$ est plat.
\end{lemma}

\begin{proof}
Considérons le diagramme
$$
\xymatrix{
T \times_S G \ar[rrr]_{(t, g) \mapsto (t, m(\psi(t), g))} & & &
T \times_S G \ar[r]_{\text{pr}} \ar[d] &
G \ar[d] \\
& & &
T \ar[r] &
S
}
$$
La flèche horizontale supérieure gauche est un isomorphisme et le
carré est cartésien. Le lemme résulte donc de
Morphismes, Lemme \ref{morphisms-lemma-base-change-flat}.
\end{proof}

\begin{lemma}
\label{groupoids-lemma-group-scheme-module-differentials}
\begin{reference}
\cite[Proposition 3.15]{BookAV}
\end{reference}
Soit $(G, m, e, i)$ un schéma en groupes sur le schéma $S$.
Notons $f : G \to S$ le morphisme structural.
Il existe alors des isomorphismes canoniques
$$
\Omega_{G/S} \cong f^*\mathcal{C}_{S/G} \cong f^*e^*\Omega_{G/S}
$$
où $\mathcal{C}_{S/G}$ désigne le faisceau conormal de
l'immersion $e$. En particulier, si $S$ est le spectre d'un corps, alors
$\Omega_{G/S}$ est un $\mathcal{O}_G$-module libre.
\end{lemma}

\begin{proof}
D'après Morphismes, Lemme \ref{morphisms-lemma-base-change-differentials}, nous avons
$$
\Omega_{G \times_S G/G} = \text{pr}_0^*\Omega_{G/S}
$$
où, dans le membre de gauche, nous considérons $G \times_S G$ comme un schéma sur $G$
au moyen de $\text{pr}_1$.
Soit $\tau : G \times_S G \to G \times_S G$ l'application dite ``de cisaillement''
donnée sur les points par $(g, h) \mapsto (m(g, h), h)$. Cette application est un automorphisme
de $G \times_S G$ considéré comme un schéma sur $G$ via la projection $\text{pr}_1$.
En combinant ces deux remarques, nous obtenons un isomorphisme
$$
\tau^*\text{pr}_0^*\Omega_{G/S} \to \text{pr}_0^*\Omega_{G/S}
$$
Puisque $\text{pr}_0 \circ \tau = m$, celui-ci peut se réécrire sous la forme
$$
m^*\Omega_{G/S} \to \text{pr}_0^*\Omega_{G/S}
$$
En tirant cet isomorphisme en arrière par
$(e \circ f, \text{id}_G) : G \to G \times_S G$
et en utilisant que $m \circ (e \circ f, \text{id}_G) = \text{id}_G$
et que $\text{pr}_0 \circ (e \circ f, \text{id}_G) = e \circ f$,
nous obtenons un isomorphisme
$$
\Omega_{G/S} \to f^*e^*\Omega_{G/S}
$$
comme souhaité. D'après
Morphismes, Lemme \ref{morphisms-lemma-differentials-relative-immersion-section},
nous avons $\mathcal{C}_{S/G} \cong e^*\Omega_{G/S}$.
Si $S$ est le spectre d'un corps, alors
tout $\mathcal{O}_S$-module sur $S$ est libre,
d'où la dernière assertion.
\end{proof}

\begin{lemma}
\label{groupoids-lemma-group-scheme-addition-tangent-vectors}
Soit $S$ un schéma. Soit $G$ un schéma en groupes sur $S$.
Soit $s \in S$. Alors la composée
$$
T_{G/S, e(s)} \oplus T_{G/S, e(s)} = T_{G \times_S G/S, (e(s), e(s))}
\rightarrow T_{G/S, e(s)}
$$
est l'addition des vecteurs tangents. Ici, le signe $=$ provient de
Variétés, Lemme \ref{varieties-lemma-tangent-space-product},
et la flèche de droite est induite par $m : G \times_S G \to G$ via
Variétés, Lemme \ref{varieties-lemma-map-tangent-spaces}.
\end{lemma}

\begin{proof}
Nous utiliserons Variétés, Équation (\ref{varieties-equation-tangent-space-fibre})
et travaillerons avec les vecteurs tangents dans les fibres.
Un élément $\theta$ du premier facteur $T_{G_s/s, e(s)}$
est l'image de $\theta$ par l'application
$T_{G_s/s, e(s)} \to T_{G_s \times G_s/s, (e(s), e(s))}$
provenant de $(1, e) : G_s \to G_s \times G_s$.
Puisque $m \circ (1, e) = 1$, nous voyons que $\theta$ s'envoie sur $\theta$
par fonctorialité. Comme l'application est linéaire, nous voyons que
$(\theta_1, \theta_2)$ s'envoie sur $\theta_1 + \theta_2$.
\end{proof}





\section{Propriétés des schémas en groupes sur un corps}
\label{groupoids-section-properties-group-schemes-field}

\noindent
Dans cette section, nous rassemblons quelques propriétés des schémas en groupes sur un
corps. Dans le cas des schémas en groupes qui sont (localement) algébriques
sur un corps, nous pouvons en dire beaucoup plus, voir
section \ref{groupoids-section-algebraic-group-schemes}.

\begin{lemma}
\label{groupoids-lemma-group-scheme-over-field-open-multiplication}
Si $(G, m)$ est un schéma en groupes sur un corps $k$, alors le
morphisme de multiplication $m : G \times_k G \to G$ est ouvert.
\end{lemma}

\begin{proof}
Le morphisme de multiplication est isomorphe au morphisme de projection
$\text{pr}_0 : G \times_k G \to G$
car le diagramme
$$
\xymatrix{
G \times_k G \ar[d]^m \ar[rrr]_{(g, g') \mapsto (m(g, g'), g')} & & &
G \times_k G \ar[d]^{(g, g') \mapsto g} \\
G \ar[rrr]^{\text{id}} & & & G
}
$$
est commutatif, les flèches horizontales étant des isomorphismes. La projection
est ouverte d'après
Morphismes, Lemme \ref{morphisms-lemma-scheme-over-field-universally-open}.
\end{proof}

\begin{lemma}
\label{groupoids-lemma-group-scheme-over-field-translate-open}
Soit $(G, m)$ un schéma en groupes sur un corps $k$. Soit $U \subset G$
un ouvert et soit $T \to G$ un morphisme de schémas. Alors l'image de
la composée $T \times_k U \to G \times_k G \to G$ est ouverte.
\end{lemma}

\begin{proof}
Pour toute extension de corps $K/k$, le morphisme $G_K \to G$ est ouvert
(Morphismes, Lemme \ref{morphisms-lemma-scheme-over-field-universally-open}).
Tout point $\xi$ de $T \times_k U$ est l'image d'un morphisme
$(t, u) : \Spec(K) \to T \times_k U$ pour un certain $K$. Alors l'image de
$T_K \times_K U_K = (T \times_k U)_K \to G_K$ contient le translaté
$t \cdot U_K$, qui est ouvert. En combinant ces faits, nous voyons que
l'image de $T \times_k U \to G$ contient un voisinage ouvert de
l'image de $\xi$. Comme $\xi$ était arbitraire, le résultat s'ensuit.
\end{proof}

\begin{lemma}
\label{groupoids-lemma-group-scheme-over-field-separated}
Soit $G$ un schéma en groupes sur un corps.
Alors $G$ est un schéma séparé.
\end{lemma}

\begin{proof}
Posons $S = \Spec(k)$, où $k$ est un corps, et soit $G$ un schéma en groupes
sur $S$. D'après
le Lemme \ref{groupoids-lemma-group-scheme-separated},
nous devons montrer que $e : S \to G$ est une immersion fermée. D'après
Morphismes, Lemme
\ref{morphisms-lemma-algebraic-residue-field-extension-closed-point-fibre},
l'image de $e : S \to G$ est un point fermé de $G$.
Il est clair que $\mathcal{O}_G \to e_*\mathcal{O}_S$ est surjectif,
puisque $e_*\mathcal{O}_S$ est un faisceau gratte-ciel supporté en l'élément
neutre de $G$, de valeur $k$. Nous concluons que $e$ est une immersion fermée d'après
Schémas, Lemme \ref{schemes-lemma-characterize-closed-immersions}.
\end{proof}

\begin{lemma}
\label{groupoids-lemma-group-scheme-field-geometrically-irreducible}
Soit $G$ un schéma en groupes sur un corps $k$.
Alors
\begin{enumerate}
\item tout anneau local $\mathcal{O}_{G, g}$ de $G$ possède un unique idéal
premier minimal,
\item il existe exactement une composante irréductible $Z$ de $G$
passant par $e$, et
\item $Z$ est géométriquement irréductible sur $k$.
\end{enumerate}
\end{lemma}

\begin{proof}
Pour tout point $g \in G$, il existe une extension de corps
$K/k$ et un point $K$-rationnel $g' \in G(K)$ dont l'image est
$g$. Si l'on considère $g'$ comme un point $K$-rationnel du
schéma en groupes $G_K$, nous voyons alors que
$\mathcal{O}_{G, g} \to \mathcal{O}_{G_K, g'}$ est un homomorphisme local
d'anneaux fidèlement plat (puisque $G_K \to G$ est plat et qu'un homomorphisme
local plat est fidèlement plat, voir
Algèbre, Lemme \ref{algebra-lemma-local-flat-ff}).
Le résultat pour $\mathcal{O}_{G_K, g'}$ entraîne le
résultat pour $\mathcal{O}_{G, g}$, voir
Algèbre, Lemme \ref{algebra-lemma-injective-minimal-primes-in-image}.
Ainsi, pour démontrer (1), il suffit de
le démontrer pour les points $k$-rationnels $g$ de $G$. Dans ce cas,
la translation par $g$ définit un automorphisme $G \to G$
qui envoie $e$ sur $g$. Ainsi $\mathcal{O}_{G, g} \cong \mathcal{O}_{G, e}$.
On voit de cette façon que (2) implique (1), puisque les composantes irréductibles
passant par $e$ correspondent bijectivement aux idéaux premiers minimaux
de $\mathcal{O}_{G, e}$.

\medskip\noindent
Pour démontrer (2) et (3), il suffit de démontrer (2) lorsque $k$
est algébriquement clos. Dans ce cas, soient $Z_1$, $Z_2$ deux
composantes irréductibles de $G$ passant par $e$.
Puisque $k$ est algébriquement clos, le sous-schéma fermé
$Z_1 \times_k Z_2 \subset G \times_k G$ est lui aussi irréductible, voir
Variétés, Lemme \ref{varieties-lemma-bijection-irreducible-components}.
Par conséquent, $m(Z_1 \times_k Z_2)$ est contenu dans une composante irréductible
de $G$. D'autre part, il contient
$Z_1$ et $Z_2$ puisque $m|_{e \times G} = \text{id}_G$ et
$m|_{G \times e} = \text{id}_G$. Nous concluons que $Z_1 = Z_2$, comme souhaité.
\end{proof}

\begin{remark}
\label{groupoids-remark-warning-group-scheme-geometrically-irreducible}
Attention : le résultat du
Lemme \ref{groupoids-lemma-group-scheme-field-geometrically-irreducible}
ne signifie pas que toute composante irréductible de $G/k$ soit
géométriquement irréductible. Par exemple, le schéma en groupes
$\mu_{3, \mathbf{Q}} = \Spec(\mathbf{Q}[x]/(x^3 - 1))$
sur $\mathbf{Q}$ possède deux composantes irréductibles correspondant
à la factorisation $x^3 - 1 = (x - 1)(x^2 + x + 1)$.
Le premier facteur correspond à la composante irréductible
passant par l'élément neutre, tandis que la seconde composante irréductible
n'est pas géométriquement irréductible sur $\Spec(\mathbf{Q})$.
\end{remark}

\begin{lemma}
\label{groupoids-lemma-reduced-subgroup-scheme-perfect}
Soit $G$ un schéma en groupes sur un corps parfait $k$.
Alors le réduit $G_{red}$ de $G$ est un sous-schéma en groupes fermé de $G$.
\end{lemma}

\begin{proof}
Démonstration omise. Indication : utiliser le fait que $G_{red} \times_k G_{red}$ est réduit d'après
Variétés, Lemmes \ref{varieties-lemma-perfect-reduced} et
\ref{varieties-lemma-geometrically-reduced-any-base-change}.
\end{proof}

\begin{lemma}
\label{groupoids-lemma-open-subgroup-closed-over-field}
Soit $k$ un corps. Soit $\psi : G' \to G$ un morphisme de schémas en groupes
sur $k$. Si $\psi(G')$ est ouvert dans $G$, alors $\psi(G')$ est fermé dans $G$.
\end{lemma}

\begin{proof}
Posons $U = \psi(G') \subset G$. Soit $Z = G \setminus \psi(G') = G \setminus U$,
muni de la structure réduite induite de sous-schéma fermé. D'après
le Lemme \ref{groupoids-lemma-group-scheme-over-field-translate-open},
l'image de
$$
Z \times_k G' \longrightarrow
Z \times_k U \longrightarrow G
$$
est ouverte (la première flèche est surjective). D'autre part, puisque $\psi$
est un homomorphisme de schémas en groupes, l'image de $Z \times_k G' \to G$
est contenue dans $Z$ (car la translation par $\psi(g')$ préserve
$U$ pour tout point $g'$ de $G'$ ; un petit détail est omis).
Par conséquent, $Z \subset G$ est un ouvert (mais pas
nécessairement un sous-schéma ouvert). Ainsi $U = \psi(G')$ est fermé.
\end{proof}

\begin{lemma}
\label{groupoids-lemma-immersion-group-schemes-closed-over-field}
Soit $i : G' \to G$ une immersion de schémas en groupes sur un corps $k$.
Alors $i$ est une immersion fermée, c'est-à-dire que $i(G')$ est un sous-schéma en groupes fermé
de $G$.
\end{lemma}

\begin{proof}
Pour montrer que $i$ est une immersion fermée, il suffit de montrer que
$i(G')$ est un sous-ensemble fermé de $G$. Soit $k \subset k'$ une inclusion dans un corps parfait
qui est une extension de $k$. Si $i(G'_{k'}) \subset G_{k'}$ est fermé, alors
$i(G') \subset G$ est fermé d'après
Morphismes, Lemme \ref{morphisms-lemma-fpqc-quotient-topology}
(puisque $G_{k'} \to G$ est plat, quasi-compact et surjectif).
Nous pouvons donc supposer, et supposons, que $k$ est parfait. Nous utiliserons sans autre
mention le fait que les produits de schémas réduits sur $k$ sont réduits.
Nous pouvons remplacer $G'$ et $G$ par leurs réduits, voir
le Lemme \ref{groupoids-lemma-reduced-subgroup-scheme-perfect}.
Soit $\overline{G'} \subset G$ l'adhérence de $i(G')$, considérée
comme un sous-schéma fermé réduit. D'après
Variétés, Lemme \ref{varieties-lemma-closure-of-product},
nous concluons que $\overline{G'} \times_k \overline{G'}$
est l'adhérence de l'image de $G' \times_k G' \to G \times_k G$. Par conséquent
$$
m\Big(\overline{G'} \times_k \overline{G'}\Big)
\subset \overline{G'}
$$
puisque $m$ est continu. Il s'ensuit que $\overline{G'} \subset G$
est un sous-schéma en groupes fermé (réduit). D'après
le Lemme \ref{groupoids-lemma-open-subgroup-closed-over-field},
nous voyons que $i(G') \subset \overline{G'}$ est également fermé,
ce qui implique que $i(G') = \overline{G'}$, comme souhaité.
\end{proof}

\begin{lemma}
\label{groupoids-lemma-irreducible-group-scheme-over-field-qc}
Soit $G$ un schéma en groupes sur un corps $k$. Si $G$ est irréductible,
alors $G$ est quasi-compact.
\end{lemma}

\begin{proof}
Supposons que $K/k$ soit une extension de corps. Si $G_K$
est quasi-compact, alors $G$ l'est aussi puisque $G_K \to G$ est surjectif.
D'après le Lemme \ref{groupoids-lemma-group-scheme-field-geometrically-irreducible},
nous voyons que $G_K$ est irréductible. Il suffit donc de démontrer le lemme
après avoir remplacé $k$ par une extension. Choisissons comme $K$ un corps algébriquement
clos, de très grand cardinal. Alors, d'après
Variétés, Lemme \ref{varieties-lemma-make-Jacobson},
nous voyons que $G_K$ est un schéma de Jacobson dont tous les points fermés ont
un corps résiduel égal à $K$. Autrement dit, nous pouvons supposer que $G$ est un
schéma de Jacobson dont tous les points fermés ont pour corps résiduel $k$.

\medskip\noindent
Soit $U \subset G$ un ouvert affine non vide. Soit $g \in G(k)$. Alors
$gU \cap U \not = \emptyset$. Nous voyons donc que $g$ appartient à l'image
du morphisme
$$
U \times_{\Spec(k)} U \longrightarrow G, \quad
(u_1, u_2) \longmapsto u_1u_2^{-1}
$$
Puisque l'image de ce morphisme est ouverte
(Lemme \ref{groupoids-lemma-group-scheme-over-field-open-multiplication}),
nous voyons que cette image est $G$ tout entier (car $G$ est un schéma de Jacobson
dont tous les points fermés sont $k$-rationnels).
Puisque $U$ est affine, $U \times_{\Spec(k)} U$ l'est aussi. Ainsi $G$ est
l'image d'un schéma quasi-compact, et est donc quasi-compact.
\end{proof}

\begin{lemma}
\label{groupoids-lemma-connected-group-scheme-over-field-irreducible}
Soit $G$ un schéma en groupes sur un corps $k$. Si $G$ est connexe,
alors $G$ est irréductible.
\end{lemma}

\begin{proof}
D'après Variétés, Lemme
\ref{varieties-lemma-geometrically-connected-if-connected-and-point},
nous voyons que $G$ est géométriquement connexe. Si nous montrons que $G_K$
est irréductible pour une extension de corps $K/k$, alors
le lemme en découle. Nous pouvons donc appliquer
Variétés, Lemme \ref{varieties-lemma-make-Jacobson},
pour nous ramener au cas où $k$ est algébriquement clos,
$G$ est un schéma de Jacobson et tous les points fermés sont $k$-rationnels.

\medskip\noindent
Soit $Z \subset G$ l'unique composante irréductible de $G$ qui contient
l'élément neutre, voir
le Lemme \ref{groupoids-lemma-group-scheme-field-geometrically-irreducible}.
Munissons $Z$ de la structure réduite induite de sous-schéma fermé ;
nous voyons que $Z \times_k Z$ est réduit et irréductible
(Variétés, Lemmes
\ref{varieties-lemma-geometrically-reduced-any-base-change} et
\ref{varieties-lemma-bijection-irreducible-components}).
Nous en concluons que $m|_{Z \times_k Z} : Z \times_k Z \to G$ se factorise
par $Z$. Ainsi, $Z$ devient un sous-schéma en groupes fermé de $G$.

\medskip\noindent
Pour obtenir une contradiction, supposons qu'il existe une autre composante
irréductible $Z' \subset G$. Alors $Z \cap Z' = \emptyset$ d'après
le Lemme \ref{groupoids-lemma-group-scheme-field-geometrically-irreducible}.
D'après le Lemme \ref{groupoids-lemma-irreducible-group-scheme-over-field-qc},
nous voyons que $Z$ est quasi-compact. Nous pouvons donc choisir un ouvert quasi-compact
$U \subset G$ tel que $Z \subset U$ et $U \cap Z' = \emptyset$.
L'image $W$ de $Z \times_k U \to G$ est ouverte dans $G$ d'après
le Lemme \ref{groupoids-lemma-group-scheme-over-field-translate-open}.
D'autre part, $W$ est quasi-compact en tant qu'image d'un
espace quasi-compact. Nous affirmons que $W$ est fermé.

\medskip\noindent
Preuve de l'affirmation. Puisque $W$ est quasi-compact et $G$ est séparé
(Lemme \ref{groupoids-lemma-group-scheme-over-field-separated}), l'inclusion
$W \to G$ est quasi-compacte d'après Schémas, Lemme
\ref{schemes-lemma-quasi-compact-permanence}. Nous en concluons que
les points de l'adhérence de $W$ sont des spécialisations de points de $W$
(Morphismes, Lemme \ref{morphisms-lemma-reach-points-scheme-theoretic-image}).
Ainsi, nous devons montrer que toute composante irréductible
$Z'' \subset G$ de $G$ qui rencontre $W$ est contenue dans $W$.
Comme $G$ est un schéma de Jacobson et que ses points fermés sont rationnels,
$Z'' \cap W$ contient un point rationnel
$g \in Z''(k) \cap W(k)$, et donc $Z'' = Zg$. Mais $W = m(Z \times_k W)$
par construction, donc $Z'' \cap W \not = \emptyset$ implique
$Z'' \subset W$.

\medskip\noindent
D'après l'affirmation, $W \subset G$ est un sous-ensemble ouvert et fermé de $G$.
Or $W \cap Z' = \emptyset$, car sinon, d'après l'argument donné dans
le paragraphe précédent, nous aurions $Z' = Zg$ pour un certain $g \in W(k)$.
Alors, puisque $Z$ est un sous-groupe, nous pourrions même choisir $g \in U(k)$, ce qui
contredirait $Z' \cap U = \emptyset$. Ainsi, $W \subset G$ est un sous-ensemble propre, ouvert
et fermé, ce qui contredit l'hypothèse que $G$ est connexe.
\end{proof}

\begin{proposition}
\label{groupoids-proposition-connected-component}
Soit $G$ un schéma en groupes sur un corps $k$. Il existe un sous-schéma
en groupes fermé canonique $G^0 \subset G$ ayant les propriétés suivantes
\begin{enumerate}
\item $G^0 \to G$ est une immersion fermée plate,
\item $G^0 \subset G$ est la composante connexe neutre,
\item $G^0$ est géométriquement irréductible, et
\item $G^0$ est quasi-compact.
\end{enumerate}
\end{proposition}

\begin{proof}
Soit $G^0$ la composante connexe neutre, munie de sa structure canonique
de schéma (Morphismes, Définition
\ref{morphisms-definition-scheme-structure-connected-component}).
Pour montrer que $G^0$ est un sous-schéma en groupes fermé, nous utiliserons le
critère du Lemme \ref{groupoids-lemma-closed-subgroup-scheme}.
Le morphisme $e : \Spec(k) \to G$ se factorise par $G^0$ puisque nous avons choisi
$G^0$ comme la composante connexe de $G$ qui contient $e$.
Puisque $i : G \to G$ est un automorphisme fixant $e$, nous voyons que
$i$ envoie $G^0$ dans lui-même.
D'après Variétés, Lemme \ref{varieties-lemma-geometrically-connected-criterion},
le schéma $G^0$ est géométriquement connexe sur $k$.
Ainsi, $G^0 \times_k G^0$ est connexe
(Variétés, Lemme \ref{varieties-lemma-bijection-connected-components}).
Ainsi, $m(G^0 \times_k G^0) \subset G^0$ ensemblistement.
Ainsi, $m|_{G^0 \times_k G^0} : G^0 \times_k G^0 \to G$
se factorise par $G^0$ d'après
Morphismes, Lemme \ref{morphisms-lemma-characterize-flat-closed-immersions}.
Par conséquent, $G^0$ est un sous-schéma en groupes fermé de $G$.
D'après le Lemme \ref{groupoids-lemma-connected-group-scheme-over-field-irreducible},
nous voyons que $G^0$ est irréductible. D'après
le Lemme \ref{groupoids-lemma-group-scheme-field-geometrically-irreducible},
nous voyons que $G^0$ est géométriquement irréductible. D'après
le Lemme \ref{groupoids-lemma-irreducible-group-scheme-over-field-qc},
nous voyons que $G^0$ est quasi-compact.
\end{proof}

\begin{lemma}
\label{groupoids-lemma-profinite-product-over-field}
Soit $k$ un corps. Soit $T = \Spec(A)$, où $A$ est une colimite filtrante
d'algèbres qui sont des produits finis de copies de $k$. Pour tout schéma $X$
sur $k$, on a $|T \times_k X| = |T| \times |X|$ comme espaces topologiques.
\end{lemma}

\begin{proof}
En prenant un recouvrement par des ouverts affines, on se ramène au cas où $X$ est affine.
Écrivons $X = \Spec(B)$.
Écrivons $A = \colim A_i$ avec $A_i = \prod_{t \in T_i} k$ et $T_i$ fini.
Alors $T_i = |\Spec(A_i)|$ muni de la topologie discrète, et les morphismes
de transition $A_i \to A_{i'}$ sont donnés par des applications $T_{i'} \to T_i$. Ainsi,
$|T| = \lim T_i$ comme espace topologique, voir
Limites, Lemme \ref{limits-lemma-topology-limit}. De même, nous avons
\begin{align*}
|T \times_k X| & =
|\Spec(A \otimes_k B)| \\
& =
|\Spec(\colim A_i \otimes_k B)| \\
& =
\lim |\Spec(A_i \otimes_k B)| \\
& =
\lim |\Spec(\prod\nolimits_{t \in T_i} B)| \\
& =
\lim T_i \times |X| \\
& =
(\lim T_i) \times |X| \\
& =
|T| \times |X|
\end{align*}
d'après le lemme précédent et le fait que les limites commutent entre elles.
\end{proof}

\noindent
Le lemme suivant affirme qu'en fait on peut placer une
``famille profinie algébrique de points'' dans un ouvert affine.
Nous invitons le lecteur à lire d'abord le Lemme \ref{groupoids-lemma-points-in-affine}.

\begin{lemma}
\label{groupoids-lemma-compact-set-in-affine}
Soit $k$ un corps algébriquement clos. Soit $G$ un schéma en groupes sur $k$.
Supposons que $G$ soit un schéma de Jacobson et que tous ses points fermés soient $k$-rationnels.
Soit $T = \Spec(A)$, où $A$ est une colimite filtrante d'algèbres qui
sont des produits finis de copies de $k$. Pour tout morphisme $f : T \to G$,
il existe un ouvert affine $U \subset G$ contenant $f(T)$.
\end{lemma}

\begin{proof}
Soit $G^0 \subset G$ le sous-schéma en groupes fermé obtenu dans la
Proposition \ref{groupoids-proposition-connected-component}. Les deux premiers paragraphes
servent à se ramener au cas $G = G^0$.

\medskip\noindent
Remarquons que $T$ est une limite projective filtrante d'espaces topologiques finis
(Limites, Lemme \ref{limits-lemma-topology-limit}), donc profini comme
espace topologique (Topologie, Définition \ref{topology-definition-profinite}).
Soit $W \subset G$ un ouvert quasi-compact contenant l'image de $T \to G$.
Après avoir remplacé $W$ par l'image de $G^0 \times W \to G \times G \to G$, on peut
supposer que $W$ est invariant sous l'action de $G^0$ par translations à gauche, voir
le Lemme \ref{groupoids-lemma-group-scheme-over-field-translate-open}.
Considérons la composée
$$
\psi = \pi \circ f : T \xrightarrow{f} W \xrightarrow{\pi} \pi_0(W)
$$
L'espace $\pi_0(W)$ est profini
(Topologie, Lemme \ref{topology-lemma-spectral-pi0} et
Propriétés, Lemme
\ref{properties-lemma-quasi-compact-quasi-separated-spectral}).
Soit $F_\xi \subset T$ la fibre de $T \to \pi_0(W)$ au-dessus de $\xi \in \pi_0(W)$.
Supposons que, pour tout $\xi \in \pi_0(W)$, nous ayons trouvé un ouvert affine
$U_\xi \subset W$ tel que $f(F_\xi) \subset U_\xi$. Comme $\psi$ est fermée d'après
Topologie, Lemme \ref{topology-lemma-closed-map},
$\psi(T \setminus f^{-1}(U_\xi))$ est un sous-ensemble fermé de $\pi_0(W)$
qui ne contient pas $\xi$. Son complémentaire $V_\xi$ est donc un
voisinage ouvert de $\xi$, et $\psi^{-1}(V_\xi) \subseteq f^{-1}(U_\xi)$.
Le recouvrement ouvert de $\pi_0(W)$ par les $V_\xi$ peut être raffiné en
un recouvrement fini $\pi_0(W) = V_1 \amalg \ldots \amalg V_n$
par des sous-ensembles ouverts et fermés deux à deux disjoints (Topologie,
Lemme \ref{topology-lemma-profinite-refine-open-covering}).
Pour tout $i = 1, \ldots, n$, choisissons $\xi_i \in \pi_0(W)$
tel que $V_i \subset V_{\xi_i}$. Comme $V_i$ est fermé et
$U_{\xi_i}$ affine, $U_{\xi_i} \cap \pi^{-1}(V_i)$ est également affine.
Ainsi, en remplaçant $U_{\xi_i}$ par $U_{\xi_i} \cap \pi^{-1}(V_i)$, nous
avons $\psi^{-1}(V_i) = f^{-1}(U_{\xi_i})$. Comme les $V_i$ recouvrent
$\pi_0(W)$, les $f^{-1}(U_{\xi_i})$ recouvrent $T$, et donc
$f(T) \subseteq U = \bigcup_i U_{\xi_i}$. Alors $U$ est la réunion disjointe
d'un nombre fini d'ouverts affines et est donc affine, comme souhaité.

\medskip\noindent
Soit $Z$ une composante connexe de $G$ qui rencontre $f(T)$. Alors $Z$
possède un point $k$-rationnel $z$ (car tous les corps résiduels du schéma $T$
sont isomorphes à $k$). Ainsi, $Z = G^0 z$. Par notre choix de $W$, nous voyons
que $Z \subset W$. L'argument du paragraphe précédent nous ramène
au problème de trouver un voisinage ouvert affine de $f(T) \cap Z$ dans $W$.
Après translation par un point rationnel, nous pouvons supposer que $Z = G^0$
(détails omis). Remarquons que l'image réciproque schématique
$T' = f^{-1}(G^0) \subset T$ est un sous-schéma fermé, qui est du même type.
Après avoir remplacé $T$ par $T'$, nous pouvons supposer que $f(T) \subset G^0$.
Choisissons un voisinage ouvert affine $U \subset G$
de $e \in G$, de sorte qu'en particulier $U \cap G^0$ soit non vide. Nous allons montrer
qu'il existe $g \in G^0(k)$ tel que $f(T) \subset g^{-1}U$.
Cela achèvera la preuve, car $g^{-1}U \subset W$ en vertu de l'invariance de $W$
sous les translations à gauche par $G^0$.

\medskip\noindent
Les arguments des deux paragraphes précédents permettent de passer à $G^0$
et de ramener le problème au suivant :
Supposons que $G$ soit irréductible et que $U \subset G$ soit un
voisinage ouvert affine de $e$. Montrons que $f(T) \subset g^{-1}U$
pour un certain $g \in G(k)$. Considérons le morphisme
$$
U \times_k T \longrightarrow G \times_k T,\quad
(t, u) \longrightarrow (uf(t)^{-1}, t)
$$
qui est une immersion ouverte (car le prolongement de ce morphisme à
$G \times_k T \to G \times_k T$ est un isomorphisme).
D'après notre hypothèse sur $T$, nous avons $|U \times_k T| = |U| \times |T|$
et de même pour $G \times_k T$, voir
le Lemme \ref{groupoids-lemma-profinite-product-over-field}.
L'image de l'immersion ouverte affichée est donc une réunion finie
de rectangles $\bigcup_{i = 1, \ldots, n} U_i \times V_i$, avec
$V_i \subset T$ et $U_i \subset G$ ouverts quasi-compacts. Cela signifie que
les ouverts possibles $Uf(t)^{-1}$, pour $t \in T$, sont en nombre fini, disons
$Uf(t_1)^{-1}, \ldots, Uf(t_r)^{-1}$. Puisque $G$ est irréductible,
l'intersection
$$
Uf(t_1)^{-1} \cap \ldots \cap Uf(t_r)^{-1}
$$
est non vide et, puisque $G$ est un schéma de Jacobson à points fermés $k$-rationnels,
nous pouvons choisir dans cette intersection un point $k$-rationnel $g \in G(k)$.
Nous voyons alors que $g \in Uf(t)^{-1}$ pour tout $t \in T$, ce qui
signifie que $f(t) \in g^{-1}U$, comme souhaité.
\end{proof}

\begin{remark}
\label{groupoids-remark-easy}
Si $G$ est un schéma en groupes sur un corps, existe-t-il toujours un sous-schéma en groupes quasi-compact,
ouvert et fermé ? D'après la
Proposition \ref{groupoids-proposition-connected-component},
cette question n'est intéressante que si $G$ possède une infinité de composantes connexes
(géométriquement).
\end{remark}

\begin{lemma}
\label{groupoids-lemma-group-scheme-field-countable-affine}
Soit $G$ un schéma en groupes sur un corps.
Il existe un sous-schéma ouvert et fermé $G' \subset G$
qui est une réunion dénombrable d'ouverts affines.
\end{lemma}

\begin{proof}
Soit $e \in U(k)$ un voisinage ouvert quasi-compact de l'élément
neutre. En remplaçant $U$ par $U \cap i(U)$, nous pouvons supposer que
$U$ est invariant par l'application d'inversion. Comme $G$ est séparé, c'est
encore un ensemble quasi-compact. Posons
$$
G' = \bigcup\nolimits_{n \geq 1} m_n(U \times_k \ldots \times_k U)
$$
où $m_n : G \times_k \ldots \times_k G \to G$ est l'application de multiplication à $n$ facteurs
suivante
$(g_1, \ldots, g_n) \mapsto m(m(\ldots (m(g_1, g_2), g_3), \ldots ), g_n)$.
Chacune de ces applications est ouverte (voir
Lemme \ref{groupoids-lemma-group-scheme-over-field-open-multiplication}) ;
ainsi, $G'$ est un sous-schéma en groupes ouvert. D'après le
Lemme \ref{groupoids-lemma-open-subgroup-closed-over-field},
c'est aussi un sous-schéma en groupes fermé.
\end{proof}






\section{Propriétés des schémas en groupes algébriques}
\label{groupoids-section-algebraic-group-schemes}

\noindent
Rappelons qu'un schéma sur un corps $k$ est (localement) algébrique s'il est
(localement) de type fini sur $\Spec(k)$, voir
Variétés, Définition \ref{varieties-definition-algebraic-scheme}.
C'est dans ce sens que nous employons le terme algébrique dans le titre de cette section.

\begin{lemma}
\label{groupoids-lemma-group-scheme-finite-type-field}
Soit $k$ un corps. Soit $G$ un schéma en groupes localement algébrique sur $k$.
Alors $G$ est équidimensionnel et $\dim(G) = \dim_g(G)$ pour tout $g \in G$.
Pour tout point fermé $g \in G$, on a $\dim(G) = \dim(\mathcal{O}_{G, g})$.
\end{lemma}

\begin{proof}
Montrons d'abord que $\dim_g(G) = \dim_{g'}(G)$ pour tout
couple de points $g, g' \in G$. D'après
Morphismes, Lemme \ref{morphisms-lemma-dimension-fibre-after-base-change},
nous pouvons étendre à volonté le corps de base. Nous pouvons donc supposer que
$g$ et $g'$ sont tous deux définis sur $k$. Il existe alors un
automorphisme de $G$ envoyant $g$ sur $g'$, d'où l'égalité.
D'après
Morphismes, Lemme \ref{morphisms-lemma-dimension-fibre-at-a-point},
on a
$\dim_g(G) = \dim(\mathcal{O}_{G, g}) +
\text{trdeg}_k(\kappa(g))$.
D'autre part, la dimension de $G$ (ou de tout ouvert de $G$)
est la borne supérieure des dimensions des anneaux locaux de $G$, voir
Propriétés, Lemme \ref{properties-lemma-codimension-local-ring}.
Il est clair que cette borne supérieure est atteinte aux points fermés $g$, auquel cas
$\text{trdeg}_k(\kappa(g)) = 0$ (par le théorème des zéros de Hilbert, voir
Morphismes, section \ref{morphisms-section-points-finite-type}).
Le lemme en résulte.
\end{proof}

\noindent
Le résultat suivant est parfois appelé théorème de Cartier.

\begin{lemma}
\label{groupoids-lemma-group-scheme-characteristic-zero-smooth}
Soit $k$ un corps de caractéristique $0$. Soit $G$ un
schéma en groupes localement algébrique sur $k$. Alors le morphisme
structural $G \to \Spec(k)$ est lisse, c'est-à-dire que $G$ est un
schéma en groupes lisse.
\end{lemma}

\begin{proof}
D'après le
Lemme \ref{groupoids-lemma-group-scheme-module-differentials},
le module des différentielles de $G$ sur $k$ est libre.
La lissité résulte donc de
Variétés, Lemme \ref{varieties-lemma-char-zero-differentials-free-smooth}.
\end{proof}

\begin{remark}
\label{groupoids-remark-when-reduced}
Tout schéma en groupes sur un corps de caractéristique $0$ est réduit, voir
\cite[I, théorème 1.1 et I, corollaire 3.9, et II, théorème 2.4]{Perrin-thesis}
ainsi que
\cite[Proposition 4.2.8]{Perrin}.
Cette question avait été posée dans
\cite[page 80]{Oort}.
Nous avons vu dans le
Lemme \ref{groupoids-lemma-group-scheme-characteristic-zero-smooth}
que cela est vrai lorsque le schéma en groupes est localement de type fini.
\end{remark}

\begin{lemma}
\label{groupoids-lemma-reduced-group-scheme-prefect-field-characteristic-p-smooth}
Soit $k$ un corps parfait de caractéristique $p > 0$ (voir le
Lemme \ref{groupoids-lemma-group-scheme-characteristic-zero-smooth}
pour le cas de caractéristique nulle).
Soit $G$ un schéma en groupes localement algébrique sur $k$.
Si $G$ est réduit, alors le morphisme
structural $G \to \Spec(k)$ est lisse, c'est-à-dire que $G$ est un
schéma en groupes lisse.
\end{lemma}

\begin{proof}
D'après le
Lemme \ref{groupoids-lemma-group-scheme-module-differentials},
le faisceau $\Omega_{G/k}$ est libre. Le lemme résulte donc de
Variétés, Lemme \ref{varieties-lemma-char-p-differentials-free-smooth}.
\end{proof}

\begin{remark}
\label{groupoids-remark-reduced-smooth-not-true-general}
Soit $k$ un corps de caractéristique $p > 0$.
Soit $\alpha \in k$ un élément qui n'est pas une puissance $p$-ième.
Le sous-schéma en groupes fermé
$$
G = V(x^p + \alpha y^p) \subset \mathbf{G}_{a, k}^2
$$
est réduit et irréductible, mais non lisse (pas même normal).
\end{remark}

\noindent
Le lemme suivant est un cas particulier du
Lemme \ref{groupoids-lemma-compact-set-in-affine},
avec une démonstration un peu plus simple.

\begin{lemma}
\label{groupoids-lemma-points-in-affine}
Soit $k$ un corps algébriquement clos.
Soit $G$ un schéma en groupes localement algébrique sur $k$.
Soient $g_1, \ldots, g_n \in G(k)$ des points $k$-rationnels.
Alors il existe un ouvert affine $U \subset G$ contenant $g_1, \ldots, g_n$.
\end{lemma}

\begin{proof}
Nous raisonnons d'abord par récurrence sur $n$ pour nous ramener au cas où tous les $g_i$
appartiennent à la même composante connexe de $G$. En effet, sinon, nous pouvons
trouver une décomposition $G = W_1 \amalg W_2$, où les $W_i$ sont ouverts dans $G$, et
(après une éventuelle renumérotation) $g_1, \ldots, g_r \in W_1$ et
$g_{r + 1}, \ldots, g_n \in W_2$ pour un certain $0 < r < n$. Par
récurrence, nous pouvons trouver des ouverts affines $U_1$ et $U_2$ de $G$ tels que
$g_1, \ldots, g_r \in U_1$ et $g_{r + 1}, \ldots, g_n \in U_2$.
Alors
$$
g_1, \ldots, g_n \in (U_1 \cap W_1) \cup (U_2 \cap W_2)
$$
fournit une solution au problème. Nous pouvons donc supposer que $g_1, \ldots, g_n$
appartiennent tous à la même composante connexe de $G$. En translatant par $g_1^{-1}$,
nous pouvons supposer que $g_1, \ldots, g_n \in G^0$, où $G^0 \subset G$ est comme dans la
Proposition \ref{groupoids-proposition-connected-component}. Choisissons un voisinage
ouvert affine $U$ de $e$ ; en particulier, $U \cap G^0$ est non vide.
Comme $G^0$ est irréductible, on voit que
$$
G^0 \cap (Ug_1^{-1} \cap \ldots \cap Ug_n^{-1})
$$
est non vide. Comme $G \to \Spec(k)$ est localement de type fini, il en va de même de
$G^0 \to \Spec(k)$ ; par conséquent, tout ouvert non vide
possède un point $k$-rationnel. Nous pouvons donc choisir $g \in G^0(k)$ tel que
$g \in Ug_i^{-1}$ pour tout $i$. Alors $g_i \in g^{-1}U$ pour tout $i$,
et $g^{-1}U$ est l'ouvert affine recherché.
\end{proof}

\begin{lemma}
\label{groupoids-lemma-algebraic-quasi-projective}
Soit $k$ un corps. Soit $G$ un schéma en groupes algébrique sur $k$.
Alors $G$ est quasi-projectif sur $k$.
\end{lemma}

\begin{proof}
D'après Variétés, Lemme \ref{varieties-lemma-ample-after-field-extension},
nous pouvons supposer que $k$ est algébriquement clos. Soit $G^0 \subset G$
la composante connexe neutre de $G$ comme dans la
proposition \ref{groupoids-proposition-connected-component}.
Alors toute autre composante connexe de $G$ possède un point $k$-rationnel
et est donc isomorphe à $G^0$ comme schéma.
Comme $G$ est quasi-compact et noethérien, il ne possède qu'un nombre fini de
composantes connexes. Nous nous ramenons ainsi au cas considéré au
paragraphe suivant.

\medskip\noindent
Soit $G$ un schéma en groupes algébrique connexe sur un corps algébriquement clos
$k$. Si $k$ est de caractéristique nulle, alors $G$ est lisse sur
$k$ d'après le Lemme \ref{groupoids-lemma-group-scheme-characteristic-zero-smooth}.
Si $k$ est de caractéristique $p > 0$, soit $H = G_{red}$
la réduction de $G$. D'après
Diviseurs, proposition \ref{divisors-proposition-push-down-ample},
il suffit de montrer que $H$ possède un faisceau inversible ample.
(Pour un schéma algébrique sur $k$, posséder un faisceau inversible
ample équivaut à être quasi-projectif sur $k$ ; voir
par exemple le résultat très général
Plus sur les morphismes, Lemme \ref{more-morphisms-lemma-quasi-projective}.)
D'après le Lemme \ref{groupoids-lemma-reduced-subgroup-scheme-perfect},
on voit que $H$ est un schéma en groupes sur $k$.
D'après le Lemme \ref{groupoids-lemma-reduced-group-scheme-prefect-field-characteristic-p-smooth},
on voit que $H$ est lisse sur $k$.
Cela nous ramène à la situation considérée au
paragraphe suivant.

\medskip\noindent
Soit $G$ un schéma en groupes lisse, irréductible et quasi-compact sur un
corps algébriquement clos $k$. Observons que les anneaux locaux de $G$
sont réguliers et donc factoriels
(Variétés, Lemme \ref{varieties-lemma-smooth-regular} et
Plus sur l'algèbre, Lemme \ref{more-algebra-lemma-regular-local-UFD}).
Le complémentaire d'un ouvert affine non vide de $G$
est le support d'un diviseur de Cartier effectif $D$.
Cela résulte de Diviseurs, Lemme
\ref{divisors-lemma-complement-open-affine-effective-cartier-divisor}.
(Observons que $G$ est séparé d'après le
Lemme \ref{groupoids-lemma-group-scheme-over-field-separated}.)
Nous en concluons qu'il existe un diviseur de Cartier effectif $D \subset G$
tel que $G \setminus D$ soit affine. Nous utiliserons ci-dessous que,
pour tout $n \geq 1$ et tous $g_1, \ldots, g_n \in G(k)$, le complémentaire
$G \setminus \bigcup D g_i$ est affine. En effet, c'est l'intersection
des ouverts affines $G \setminus Dg_i \cong G \setminus D$
dans le schéma séparé $G$.

\medskip\noindent
Nous pouvons choisir la ligne supérieure du diagramme
$$
\xymatrix{
G & U \ar[l]_j \ar[r]^\pi & \mathbf{A}^d_k \\
& W \ar[r]^{\pi'} \ar[u] & V \ar[u]
}
$$
de telle sorte que $U \not = \emptyset$, que $j : U \to G$ soit une immersion ouverte et que
$\pi$ soit \'etale, voir
Morphismes, Lemme \ref{morphisms-lemma-smooth-etale-over-affine-space}.
Il existe un ouvert affine non vide $V \subset \mathbf{A}^d_k$ tel que,
si $W = \pi^{-1}(V)$, le morphisme $\pi' = \pi|_W : W \to V$ soit fini \'etale.
En particulier, $\pi'$ est fini localement libre, disons de degré $n$.
Considérons le diviseur de Cartier effectif
$$
\mathcal{D} = \{(g, w) \mid m(g, j(w)) \in D\} \subset G \times W
$$
(Il s'agit de la restriction à $G \times W$ de l'image réciproque de $D \subset G$
par le morphisme plat $m : G \times G \to G$.)
Considérons le fermé\footnote{En utilisant les résultats
de Diviseurs, section \ref{divisors-section-norms},
nous pourrions prendre pour diviseur de Cartier effectif
$E$ la norme du diviseur de Cartier effectif $\mathcal{D}$
le long du morphisme fini localement libre $1 \times \pi'$, ce qui permettrait d'éviter
une partie des arguments.}
$T = (1 \times \pi')(\mathcal{D}) \subset G \times V$.
Comme $\pi'$ est fini localement libre, toute composante irréductible
de $T$ est de codimension $1$ dans $G \times V$. Comme $G \times V$
est lisse sur $k$, nous en concluons que ces composantes sont des diviseurs de Cartier
effectifs (Diviseurs, Lemme \ref{divisors-lemma-weil-divisor-is-cartier-UFD}
et les lemmes cités ci-dessus),
et donc que $T$ est le support d'un diviseur de Cartier effectif
$E$ dans $G \times V$. Si $v \in V(k)$, alors
$(\pi')^{-1}(v) = \{w_1, \ldots, w_n\} \subset W(k)$ et l'on voit que
$$
E_v = \bigcup\nolimits_{i = 1, \ldots, n} D j(w_i)^{-1}
$$
dans $G$, du point de vue ensembliste. En particulier, on voit que $G \setminus E_v$
est un ouvert affine (voir ci-dessus).
De plus, si $g \in G(k)$, il existe $v \in V$
tel que $g \not \in E_v$. En effet, l'ensemble $W'$ des $w \in W$ tels que
$g \not \in Dj(w)^{-1}$ est un ouvert non vide, et il suffit de choisir $v$
tel que la fibre de $W' \to V$ au-dessus de $v$ possède $n$ éléments.

\medskip\noindent
Considérons le faisceau inversible
$\mathcal{M} = \mathcal{O}_{G \times V}(E)$ sur $G \times V$.
D'après Variétés, Lemme \ref{varieties-lemma-rational-equivalence-for-Pic},
la classe d'isomorphisme $\mathcal{L}$ de la restriction
$\mathcal{M}_v = \mathcal{O}_G(E_v)$ est indépendante de $v \in V(k)$.
D'autre part, pour tout $g \in G(k)$, nous pouvons trouver $v$
tel que $g \not \in E_v$ et que $G \setminus E_v$
soit affine. Ainsi, la section canonique
(Diviseurs, Définition
\ref{divisors-definition-invertible-sheaf-effective-Cartier-divisor})
de $\mathcal{O}_G(E_v)$
correspond à une section $s_v$ de $\mathcal{L}$ qui ne
s'annule pas en $g$ et telle que $G_{s_v}$ soit affine.
Cela signifie que $\mathcal{L}$ est ample par définition
(Propriétés, Définition \ref{properties-definition-ample}).
\end{proof}

\begin{lemma}
\label{groupoids-lemma-algebraic-center}
Soit $k$ un corps. Soit $G$ un schéma en groupes localement algébrique sur $k$.
Alors le centre de $G$ est un sous-schéma en groupes fermé de $G$.
\end{lemma}

\begin{proof}
Notons $\text{Aut}(G)$ le foncteur contravariant sur la catégorie des
schémas sur $k$ qui associe à $S/k$ l'ensemble des automorphismes
du schéma $G_S$ obtenu par changement de base, en tant que schéma en groupes sur $S$. Il existe une transformation
naturelle
$$
G \longrightarrow \text{Aut}(G),\quad
g \longmapsto \text{inn}_g
$$
qui, pour tout $S$, associe à un point $g$ de $G$ à valeurs dans ce schéma l'automorphisme intérieur de $G$
déterminé par $g$. Le centre $C$ de $G$ est, par définition, le noyau de
cette transformation, c'est-à-dire le foncteur qui associe à $S$ les
$g \in G(S)$ dont l'automorphisme intérieur associé est trivial. L'énoncé
du lemme affirme que ce foncteur est représentable par un sous-schéma en groupes
fermé de $G$.

\medskip\noindent
Choisissons un entier $n \geq 1$. Soit $G_n \subset G$ le $n$-ième voisinage
infinitésimal de l'élément neutre $e$ de $G$. Pour tout schéma $S/k$,
le changement de base $G_{n, S}$ est le $n$-ième voisinage infinitésimal de
$e_S : S \to G_S$. Nous voyons ainsi qu'il existe une transformation naturelle
$\text{Aut}(G) \to \text{Aut}(G_n)$, où le membre de droite est le
foncteur des automorphismes de $G_n$ comme schéma ($G_n$ n'est pas, en général,
un schéma en groupes). Observons que $G_n$ est le spectre d'un anneau
local artinien $A_n$, de corps résiduel $k$ et de dimension finie
comme espace vectoriel sur $k$
(Variétés, Lemme \ref{varieties-lemma-algebraic-scheme-dim-0}).
Comme tout automorphisme de $G_n$ induit en particulier une application
linéaire inversible $A_n \to A_n$, nous obtenons des transformations de foncteurs
$$
G \to \text{Aut}(G) \to \text{Aut}(G_n) \to \text{GL}(A_n)
$$
Le dernier foncteur à valeurs dans les groupes est représentable, voir
Exemple \ref{groupoids-example-general-linear-group}, et la
dernière flèche est manifestement injective.
Ainsi, pour tout $n$, nous obtenons un sous-schéma en groupes fermé
$$
H_n = \Ker(G \to \text{Aut}(G_n)) = \Ker(G \to \text{GL}(A_n)).
$$
En première approximation, posons $H = \bigcap_{n \geq 1} H_n$
(intersection schématique). C'est un sous-schéma en groupes fermé
qui contient le centre $C$.

\medskip\noindent
Soit $h$ un $S$-point de $H$, où $S$ est localement noethérien.
Alors l'automorphisme $\text{inn}_h$ induit l'identité sur tous
les sous-schémas fermés $G_{n, S}$. Considérons le noyau
$K = \Ker(\text{inn}_h : G_S \to G_S)$.
C'est un sous-schéma en groupes fermé de $G_S$ sur $S$ contenant les
sous-schémas fermés $G_{n, S}$ pour $n \geq 1$.
Cela implique que $K$ contient un voisinage ouvert de
$e(S) \subset G_S$, voir
Algèbre, Remarque \ref{algebra-remark-intersection-powers-ideal}.
Soit $G^0 \subset G$ comme dans la proposition \ref{groupoids-proposition-connected-component}.
Comme $G^0$ est géométriquement irréductible, nous en concluons que $K$ contient
$G^0_S$ (pour tout ouvert non vide $U \subset G^0_{k'}$ et toute extension de corps
$k'/k$, on a $U \cdot U^{-1} = G^0_{k'}$ ; voir la démonstration du
Lemme \ref{groupoids-lemma-irreducible-group-scheme-over-field-qc}).
En appliquant ceci avec $S = H$, nous constatons que $G^0$ et $H$
sont des sous-schémas en groupes de $G$ dont les points commutent : pour tout schéma $S$
et tous les points à valeurs dans $S$, $g \in G^0(S)$ et $h \in H(S)$, on a
$gh = hg$ dans $G(S)$.

\medskip\noindent
Supposons que $k$ soit algébriquement clos. Nous pouvons alors choisir un point à valeurs dans $k$,
noté $g_i$, dans chaque composante irréductible $G_i$ de $G$. Observons que,
dans ce cas, les composantes connexes de $G$ sont les composantes irréductibles
de $G$, lesquelles sont les translatées de $G^0$ par nos $g_i$. Nous affirmons que
$$
C = H \cap \bigcap\nolimits_i \Ker(\text{inn}_{g_i} : G \to G)
\quad (\text{intersection schématique})
$$
En effet, $C$ est contenu dans le membre de droite. Réciproquement, tout
point à valeurs dans $S$, noté $h$, du membre de droite commute avec $G^0$
et avec les $g_i$, donc avec tout élément de $G = \bigcup G^0g_i$.

\medskip\noindent
Le cas d'un corps de base quelconque $k$ découle du résultat pour la
clôture algébrique $\overline{k}$ par descente. En effet, soit
$A \subset G_{\overline{k}}$ le sous-schéma en groupes fermé représentant
le centre de $G_{\overline{k}}$. Alors on a
$$
A \times_{\Spec(k)} \Spec(\overline{k}) =
\Spec(\overline{k}) \times_{\Spec(k)} A
$$
comme sous-schémas fermés de $G_{\overline{k} \otimes_k \overline{k}}$, par
fonctorialité du centre. On voit donc que $A$ descend
en un sous-schéma en groupes fermé $Z \subset G$ d'après
Descente, Lemme \ref{descent-lemma-closed-immersion}
(et Descente, Lemme \ref{descent-lemma-descending-property-closed-immersion}).
Alors $Z$ représente $C$ (nous omettons un petit argument), ce qui achève la démonstration.
\end{proof}







\section{Variétés abéliennes}
\label{groupoids-section-abelian-varieties}

\noindent
Une excellente référence pour ce sujet est le livre de Mumford sur les
variétés abéliennes, voir \cite{AVar}. Nous encourageons le lecteur à
le consulter. Il existe de nombreuses définitions équivalentes ; en voici une.

\begin{definition}
\label{groupoids-definition-abelian-variety}
Soit $k$ un corps. Une {\it variété abélienne} est un schéma en groupes sur
$k$ qui est aussi une variété propre et géométriquement intègre sur
$k$\footnote{Pour des définitions équivalentes, voir la Remarque \ref{groupoids-remark-16-equivalent}.}.
\end{definition}

\noindent
Nous démontrons quelques lemmes sur cette notion, puis nous rassemblons
tous les résultats dans la
proposition \ref{groupoids-proposition-review-abelian-varieties}.

\begin{lemma}
\label{groupoids-lemma-abelian-variety-projective}
Soit $k$ un corps. Soit $A$ une variété abélienne sur $k$.
Alors $A$ est projective.
\end{lemma}

\begin{proof}
Cela résulte du
Lemme \ref{groupoids-lemma-algebraic-quasi-projective} et de
Plus sur les morphismes, Lemme \ref{more-morphisms-lemma-projective}.
\end{proof}

\begin{lemma}
\label{groupoids-lemma-abelian-variety-change-field}
Soit $k$ un corps. Soit $A$ une variété abélienne sur $k$.
Pour toute extension de corps $K/k$, le changement de base $A_K$ est une
variété abélienne sur $K$.
\end{lemma}

\begin{proof}
Démonstration omise. Notons que c'est la raison pour laquelle nous avons exigé que $A$ soit
géométriquement intègre ; sans cette condition, ce lemme
(et beaucoup d'autres ci-dessous) serait faux.
\end{proof}

\begin{lemma}
\label{groupoids-lemma-abelian-variety-smooth}
Soit $k$ un corps. Soit $A$ une variété abélienne sur $k$.
Alors $A$ est lisse sur $k$.
\end{lemma}

\begin{proof}
Si $k$ est parfait, cela résulte du
Lemme \ref{groupoids-lemma-group-scheme-characteristic-zero-smooth}
(caractéristique nulle) et du
Lemme \ref{groupoids-lemma-reduced-group-scheme-prefect-field-characteristic-p-smooth}
(caractéristique positive).
Nous pouvons nous ramener à ce cas dans la situation générale par descente de la lissité
(Descente, Lemme \ref{descent-lemma-descending-property-smooth})
et en passant à la clôture parfaite grâce au
Lemme \ref{groupoids-lemma-abelian-variety-change-field}.
\end{proof}

\begin{lemma}
\label{groupoids-lemma-abelian-variety-abelian}
Une variété abélienne est un schéma en groupes commutatif, c'est-à-dire que la loi
de groupe est commutative.
\end{lemma}

\begin{proof}
Soit $k$ un corps. Soit $A$ une variété abélienne sur $k$.
D'après le Lemme \ref{groupoids-lemma-abelian-variety-change-field}, nous pouvons remplacer
$k$ par sa clôture algébrique. Considérons le morphisme
$$
h : A \times_k A \longrightarrow A \times_k A,\quad
(x, y) \longmapsto (x, xyx^{-1}y^{-1})
$$
C'est un morphisme au-dessus de $A$ pour la première projection de chaque côté.
Soit $e \in A(k)$ l'élément neutre. Alors on voit que $h|_{e \times A}$ est
constant de valeur $(e, e)$. D'après Plus sur les morphismes, Lemme
\ref{more-morphisms-lemma-flat-proper-family-cannot-collapse-fibre},
il existe un voisinage ouvert $U \subset A$ de $e$
tel que $h|_{U \times A}$ se factorise par un certain $Z \subset U \times A$
fini sur $U$. Cela signifie que, pour $x \in U(k)$, le morphisme
$A \to A$, $y \mapsto xyx^{-1}y^{-1}$, ne prend qu'un nombre fini de valeurs.
Bien entendu, cela signifie qu'il est constant de valeur $e$. Ainsi,
$(x, y) \mapsto xyx^{-1}y^{-1}$ est
constant de valeur $e$ sur $U \times A$, ce qui implique
que la loi de groupe sur $A$ est commutative.
\end{proof}

\begin{lemma}
\label{groupoids-lemma-apply-cube}
Soit $k$ un corps. Soit $A$ une variété abélienne sur $k$.
Soit $\mathcal{L}$ un $\mathcal{O}_A$-module inversible.
Il existe alors un isomorphisme
$$
m_{1, 2, 3}^*\mathcal{L} \otimes
m_1^*\mathcal{L} \otimes
m_2^*\mathcal{L} \otimes
m_3^*\mathcal{L} \cong
m_{1, 2}^*\mathcal{L} \otimes
m_{1, 3}^*\mathcal{L} \otimes
m_{2, 3}^*\mathcal{L}
$$
de modules inversibles sur $A \times_k A \times_k A$,
où $m_{i_1, \ldots, i_t} : A \times_k A \times_k A \to A$
est le morphisme $(x_1, x_2, x_3) \mapsto \sum x_{i_j}$.
\end{lemma}

\begin{proof}
Appliquons le théorème du cube
(Plus sur les morphismes, Théorème \ref{more-morphisms-theorem-of-the-cube})
à la différence
$$
\mathcal{M} =
m_{1, 2, 3}^*\mathcal{L} \otimes
m_1^*\mathcal{L} \otimes
m_2^*\mathcal{L} \otimes
m_3^*\mathcal{L} \otimes
m_{1, 2}^*\mathcal{L}^{\otimes -1} \otimes
m_{1, 3}^*\mathcal{L}^{\otimes -1} \otimes
m_{2, 3}^*\mathcal{L}^{\otimes -1}
$$
Cela s'applique car la restriction de $\mathcal{M}$
à $A \times A \times e = A \times A$ est égale à
$$
n_{1, 2}^*\mathcal{L} \otimes
n_1^*\mathcal{L} \otimes
n_2^*\mathcal{L} \otimes
n_{1, 2}^*\mathcal{L}^{\otimes -1} \otimes
n_1^*\mathcal{L}^{\otimes -1} \otimes
n_2^*\mathcal{L}^{\otimes -1} \cong \mathcal{O}_{A \times_k A}
$$
où $n_{i_1, \ldots, i_t} : A \times_k A \to A$
est le morphisme $(x_1, x_2) \mapsto \sum x_{i_j}$.
Il en va de même pour $A \times e \times A$ et $e \times A \times A$.
\end{proof}

\begin{lemma}
\label{groupoids-lemma-pullbacks-by-n}
Soit $k$ un corps. Soit $A$ une variété abélienne sur $k$.
Soit $\mathcal{L}$ un $\mathcal{O}_A$-module inversible.
Alors
$$
[n]^*\mathcal{L} \cong
\mathcal{L}^{\otimes n(n + 1)/2} \otimes
([-1]^*\mathcal{L})^{\otimes n(n - 1)/2}
$$
où $[n] : A \to A$ envoie $x$ sur $x + x + \ldots + x$, somme de $n$ termes,
et où $[-1] : A \to A$ est l'application d'inversion de $A$.
\end{lemma}

\begin{proof}
Considérons le morphisme $A \to A \times_k A \times_k A$,
$x \mapsto (x, x, -x)$, où $-x = [-1](x)$. En prenant l'image réciproque
de la relation du Lemme \ref{groupoids-lemma-apply-cube}, nous obtenons
$$
\mathcal{L} \otimes
\mathcal{L} \otimes
\mathcal{L} \otimes
[-1]^*\mathcal{L} \cong
[2]^*\mathcal{L}
$$
ce qui démontre le résultat pour $n = 2$. Supposons par récurrence que le résultat vaille
pour $1, 2, \ldots, n$. Considérons alors le morphisme
$A \to A \times_k A \times_k A$, $x \mapsto (x, x, [n - 1]x)$.
En prenant l'image réciproque
de la relation du Lemme \ref{groupoids-lemma-apply-cube}, nous obtenons
$$
[n + 1]^*\mathcal{L} \otimes
\mathcal{L} \otimes
\mathcal{L} \otimes
[n - 1]^*\mathcal{L} \cong
[2]^*\mathcal{L} \otimes
[n]^*\mathcal{L} \otimes
[n]^*\mathcal{L}
$$
et le résultat découle d'un calcul élémentaire.
\end{proof}

\begin{lemma}
\label{groupoids-lemma-degree-multiplication-by-d}
Soit $k$ un corps. Soit $A$ une variété abélienne sur $k$.
Soit $[d] : A \to A$ la multiplication par $d$.
Alors $[d]$ est fini localement libre de degré $d^{2\dim(A)}$.
\end{lemma}

\begin{proof}
{\emergencystretch=1em D'après le Lemme \ref{groupoids-lemma-abelian-variety-projective}
(et Plus sur les morphismes, Lemme \ref{more-morphisms-lemma-projective}),
on voit que $A$ possède un module inversible ample $\mathcal{L}$.
Comme $[-1] : A \to A$ est un automorphisme, on voit que
$[-1]^*\mathcal{L}$ est lui aussi un $\mathcal{O}_A$-module inversible ample.
Ainsi, $\mathcal{N} = \mathcal{L} \otimes [-1]^*\mathcal{L}$
est ample, voir
Propriétés, Lemme \ref{properties-lemma-ample-tensor-globally-generated}.
Comme $\mathcal{N} \cong [-1]^*\mathcal{N}$, on voit que
$[d]^*\mathcal{N} \cong \mathcal{N}^{\otimes d^2}$ d'après le
Lemme \ref{groupoids-lemma-pullbacks-by-n}.\par}

\medskip\noindent
Pour obtenir une contradiction, soit $C$ une courbe propre contenue dans une
fibre de $[d]$. Alors $\mathcal{N}^{\otimes d^2}|_C \cong \mathcal{O}_C$
est un $\mathcal{O}_C$-module inversible ample de degré $0$, ce qui
contredit par exemple Variétés, Lemme \ref{varieties-lemma-ample-curve}.
(On peut aussi utiliser Variétés, Lemme \ref{varieties-lemma-ample-positive}.)
Ainsi, toute fibre de $[d]$ est de dimension $0$, et $[d]$ est donc fini,
par exemple d'après Cohomologie des schémas, Lemme
\ref{coherent-lemma-characterize-finite}.
De plus, comme $A$ est lisse sur $k$ d'après le
Lemme \ref{groupoids-lemma-abelian-variety-smooth},
on voit que $[d] : A \to A$ est plat d'après
Algèbre, Lemme \ref{algebra-lemma-CM-over-regular-flat}
(on utilise aussi que les schémas lisses sur un corps sont réguliers et que
les anneaux réguliers sont des anneaux de Cohen-Macaulay, voir
Variétés, lemme \ref{varieties-lemma-smooth-regular} et
Algèbre, lemme \ref{algebra-lemma-regular-ring-CM}).
Ainsi, $[d]$ est fini et plat, donc fini localement libre d'après
Morphismes, lemme \ref{morphisms-lemma-finite-flat}.

\medskip\noindent
Enfin, venons-en à la formule du degré. D'après
Variétés, lemme \ref{varieties-lemma-degree-finite-morphism-in-terms-degrees},
on voit que
$$
\deg_{\mathcal{N}^{\otimes d^2}}(A) = \deg([d]) \deg_\mathcal{N}(A)
$$
Puisque le degré de $A$ relativement à
$\mathcal{N}^{\otimes d^2}$, respectivement à $\mathcal{N}$,
est le coefficient de $n^{\dim(A)}$ dans le polynôme
$$
n \longmapsto \chi(A, \mathcal{N}^{\otimes nd^2}),\quad
\text{respectivement}\quad n \longmapsto \chi(A, \mathcal{N}^{\otimes n})
$$
on voit que $\deg([d]) = d^{2 \dim(A)}$.
\end{proof}

\begin{lemma}
\label{groupoids-lemma-abelian-variety-multiplication-by-d-etale}
\begin{slogan}
La multiplication par un entier sur une variété abélienne est un morphisme étale
si et seulement si cet entier est inversible dans le corps de base.
\end{slogan}
Soit $k$ un corps. Soit $A$ une variété abélienne non nulle sur $k$.
Alors $[d] : A \to A$ est \'etale si et seulement si $d$ est inversible dans $k$.
\end{lemma}

\begin{proof}
Observons que $[d](x + y) = [d](x) + [d](y)$. Comme la translation par un
point est un automorphisme de $A$, on voit que le lieu où
$[d] : A \to A$ est \'etale est soit vide, soit égal à $A$ (certains détails
sont omis). Il suffit donc de déterminer si $[d]$ est \'etale en
l'élément neutre $e \in A(k)$. Puisque l'on sait que $[d]$ est fini localement libre
(lemme \ref{groupoids-lemma-degree-multiplication-by-d}),
dire qu'il est \'etale en $e$ équivaut à
montrer que $\text{d}[d] : T_{A/k, e} \to T_{A/k, e}$ est injective. Voir
Variétés, lemme \ref{varieties-lemma-injective-tangent-spaces-unramified} et
Morphismes, lemme \ref{morphisms-lemma-flat-unramified-etale}.
Le lemme \ref{groupoids-lemma-group-scheme-addition-tangent-vectors} montre que
$\text{d}[d]$ est donnée par la multiplication par $d$ sur $T_{A/k, e}$.
\end{proof}

\begin{lemma}
\label{groupoids-lemma-abelian-variety-multiplication-by-p}
Soit $k$ un corps de caractéristique $p > 0$. Soit $A$ une variété abélienne
de dimension $g$
sur $k$. La fibre de $[p] : A \to A$ au-dessus de $0$ possède au plus
$p^g$ points distincts.
\end{lemma}

\begin{proof}
Pour le démontrer, on peut remplacer $k$ par sa clôture algébrique, ce que l'on fait.
D'après le lemme \ref{groupoids-lemma-group-scheme-addition-tangent-vectors},
la différentielle de $[p]$ est la multiplication par $p$ en tant qu'application
$T_{A/k, e} \to T_{A/k, e}$ et est donc nulle (comparer
avec la démonstration du lemme \ref{groupoids-lemma-abelian-variety-multiplication-by-d-etale}).
Comme $[p]$ commute aux translations, on en conclut que la différentielle de $[p]$
est nulle partout, c'est-à-dire que l'application induite
$[p]^*\Omega_{A/k} \to \Omega_{A/k}$ est nulle.
En passant aux points génériques, on voit que
l'homomorphisme correspondant $[p]^* : k(A) \to k(A)$
entre corps de fonctions induit l'application nulle sur $\Omega_{k(A)/k}$.
Soit $t_1, \ldots, t_g$ une p-base de $k(A)$ sur $k$
(Compléments d'algèbre, définition \ref{more-algebra-definition-p-basis} et
lemme \ref{more-algebra-lemma-p-basis}). Alors $[p]^*(t_i)$
possède une racine $p$-ième d'après
Algèbre, lemme \ref{algebra-lemma-derivative-zero-pth-power}.
On en conclut que
$k(A)[x_1, \ldots, x_g]/(x_1^p - t_1, \ldots, x_g^p - t_g)$ est une sous-extension
de celle définie par $[p]^* : k(A) \to k(A)$.
On peut donc trouver un ouvert affine $U \subset A$ tel que
$t_i \in \mathcal{O}_A(U)$ et $x_i \in \mathcal{O}_A([p]^{-1}(U))$.
On obtient une factorisation
$$
[p]^{-1}(U)
\xrightarrow{\pi_1}
\Spec(\mathcal{O}(U)[x_1, \ldots, x_g]/(x_1^p - t_1, \ldots, x_g^p - t_g))
\xrightarrow{\pi_2}
U
$$
de $[p]$ au-dessus de $U$. Quitte à rétrécir $U$, on peut supposer que $\pi_1$
est fini localement libre (par exemple par platitude générique -- en fait, il est
déjà fini localement libre dans notre cas).
D'après le lemme \ref{groupoids-lemma-degree-multiplication-by-d}, on voit que
$[p]$ est de degré $p^{2g}$. Comme $\pi_2$
est de degré $p^g$, on voit que $\pi_1$ est également de degré $p^g$.
Le morphisme $\pi_2$ est un homéomorphisme universel ; ses fibres sont donc
réduites à un point. On en conclut que les fibres ensemblistes de $[p]^{-1}(U) \to U$
sont les fibres de $\pi_1$. Elles
ont donc au plus $p^g$ éléments. Comme $[p]$ est un homomorphisme de schémas
en groupes sur $k$, la fibre de $[p] : A(k) \to A(k)$ possède le
même cardinal pour tout $a \in A(k)$, ce qui achève la démonstration.
\end{proof}

\begin{proposition}
\label{groupoids-proposition-review-abelian-varieties}
\begin{reference}
Exposé de manière remarquable dans \cite{AVar}.
\end{reference}
Soit $A$ une variété abélienne sur un corps $k$. Alors
\begin{enumerate}
\item $A$ est projective sur $k$,
\item $A$ est un schéma en groupes commutatif,
\item le morphisme $[n] : A \to A$ est surjectif pour tout $n \geq 1$,
\item si $k$ est algébriquement clos, alors $A(k)$ est un groupe abélien divisible,
\item $A[n] = \Ker([n] : A \to A)$ est un schéma en groupes fini de degré
$n^{2\dim A}$ sur $k$,
\item $A[n]$ est \'etale sur $k$ si et seulement si $n \in k^*$,
\item si $n \in k^*$ et si $k$ est algébriquement clos,
alors $A(k)[n] \cong (\mathbf{Z}/n\mathbf{Z})^{\oplus 2\dim(A)}$,
\item si $k$ est algébriquement clos de caractéristique $p > 0$, alors
il existe un entier $0 \leq f \leq \dim(A)$ tel que
$A(k)[p^m] \cong (\mathbf{Z}/p^m\mathbf{Z})^{\oplus f}$
pour tout $m \geq 1$.
\end{enumerate}
\end{proposition}

\begin{proof}
L'assertion (1) découle du lemme \ref{groupoids-lemma-abelian-variety-projective}.
L'assertion (2) découle du lemme \ref{groupoids-lemma-abelian-variety-abelian}.
L'assertion (3) découle du lemme \ref{groupoids-lemma-degree-multiplication-by-d}.
Si $k$ est algébriquement clos, les morphismes surjectifs de variétés
sur $k$ induisent des applications surjectives sur les points $k$-rationnels, donc
(4) découle de (3).
L'assertion (5) découle du lemme \ref{groupoids-lemma-degree-multiplication-by-d}
et du fait qu'un changement de base d'un morphisme fini localement libre
de degré $N$ est un morphisme fini localement libre de degré $N$.
L'assertion (6) découle du 
lemme \ref{groupoids-lemma-abelian-variety-multiplication-by-d-etale}.
En effet, si $n$ est inversible dans $k$, alors $[n]$ est \'etale
et donc $A[n]$ est \'etale sur $k$.
D'autre part, si $n$ n'est pas inversible dans $k$, alors
$[n]$ n'est pas \'etale en $e$ et il en résulte que $A[n]$
n'est pas \'etale sur $k$ en $e$ (utiliser
Morphismes, lemmes \ref{morphisms-lemma-flat-unramified-etale} et
\ref{morphisms-lemma-set-points-where-fibres-unramified}).

\medskip\noindent
Supposons $k$ algébriquement clos. Posons $g = \dim(A)$. Démonstration de (7).
Soit $\ell$ un nombre premier inversible dans $k$. On voit alors que
$$
A[\ell](k) = A(k)[\ell]
$$
est un groupe abélien fini, annulé par $\ell$, d'ordre $\ell^{2g}$.
Il s'ensuit qu'il est isomorphe à $(\mathbf{Z}/\ell\mathbf{Z})^{2g}$
d'après le théorème de structure des groupes abéliens finis. Considérons ensuite
la suite exacte courte
$$
0 \to A(k)[\ell] \to A(k)[\ell^2] \xrightarrow{\ell} A(k)[\ell] \to 0
$$
En raisonnant de la même manière que ci-dessus, nous concluons que 
$A(k)[\ell^2] \cong (\mathbf{Z}/\ell^2\mathbf{Z})^{2g}$.
Par récurrence sur l'exposant, nous obtenons
$A(k)[\ell^m] \cong (\mathbf{Z}/\ell^m\mathbf{Z})^{2g}$.
Pour les entiers composés $n$ premiers à la caractéristique de $k$,
nous considérons les composantes primaires et obtenons la forme correcte de la
$n$-torsion de $A(k)$.
La démonstration de (8) procède exactement de la même manière, en utilisant que
le lemme \ref{groupoids-lemma-abelian-variety-multiplication-by-p} donne
$A(k)[p] \cong (\mathbf{Z}/p\mathbf{Z})^{\oplus f}$ pour un certain $0 \leq f \leq g$.
\end{proof}

\begin{remark}
\label{groupoids-remark-16-equivalent}
Soit $k$ un corps. Il existe $2 \times 4 \times 2 = 16$ définitions
équivalentes des variétés abéliennes. Soient
\begin{itemize}
\item projectif, propre,
\item {\emergencystretch=2em géométriquement irréductible, irréductible, géométriquement connexe,
connexe,\par}
\item lisse, géométriquement réduit
\end{itemize}
trois ensembles de propriétés ; choisissons-en une dans chacun d'eux, et soit $A$
un schéma en groupes sur $k$ possédant les propriétés choisies sur $k$.
Alors $A$ est une variété abélienne.
Si nous choisissons « propre, géométriquement irréductible, géométriquement
réduit », nous retrouvons la définition \ref{groupoids-definition-abelian-variety}
(utiliser Variétés, lemme \ref{varieties-lemma-geometrically-integral}).
Les choix les plus faibles possibles seraient « propre, connexe
et géométriquement réduit », voir par exemple
Morphismes, lemme \ref{morphisms-lemma-locally-projective-proper} et
Variétés, lemme \ref{varieties-lemma-smooth-geometrically-normal}.
Supposons donc que $A$ soit un schéma en groupes propre, connexe
et géométriquement réduit sur $k$.
Alors $A$ est géométriquement irréductible d'après les
lemmes \ref{groupoids-lemma-connected-group-scheme-over-field-irreducible} et
\ref{groupoids-lemma-group-scheme-field-geometrically-irreducible},
et est donc une variété abélienne.
Enfin, si $A/k$ est une variété abélienne, alors elle est projective
et lisse sur $k$ (proposition
\ref{groupoids-proposition-review-abelian-varieties}) ; elle satisfait donc
les choix les plus forts possibles « projectif, géométriquement irréductible, lisse ».
\end{remark}







\section{Actions de schémas en groupes}
\label{groupoids-section-action-group-scheme}

\noindent
Soient $(G, m)$ un groupe et $V$ un ensemble.
Rappelons qu'une {\it action (à gauche)} de $G$ sur $V$ est donnée
par une application $a : G \times V \to V$ telle que
\begin{enumerate}
\item (associativité) $a(m(g, g'), v) = a(g, a(g', v))$ pour tous
$g, g' \in G$ et $v \in V$, et
\item (élément neutre) $a(e, v) = v$ pour tout $v \in V$.
\end{enumerate}
On dit également que $V$ est un {\it $G$-ensemble} (cela signifie en général que nous
omettons $a$ de la notation --- ce qui constitue un abus de notation).
Un {\it morphisme de $G$-ensembles} $\psi : V \to V'$ est une application
telle que $\psi(a(g, v)) = a(g, \psi(v))$ pour tout $v \in V$.

\begin{definition}
\label{groupoids-definition-action-group-scheme}
Soit $S$ un schéma. Soit $(G, m)$ un schéma en groupes sur $S$.
\begin{enumerate}
\item Une {\it action de $G$ sur le schéma $X/S$} est
un morphisme $a : G \times_S X \to X$ sur $S$ tel que
pour tout $T/S$, l'application $a : G(T) \times X(T) \to X(T)$
définit la structure d'un $G(T)$-ensemble sur $X(T)$.
\item Supposons que $X$, $Y$ soient des schémas sur $S$, chacun muni
d'une action de $G$. Un morphisme {\it équivariant}, ou plus précisément
un morphisme {\it $G$-équivariant} $\psi : X \to Y$,
est un morphisme de schémas sur $S$
tel que pour tout $T/S$, l'application $\psi : X(T) \to Y(T)$ soit
un morphisme de $G(T)$-ensembles.
\end{enumerate}
\end{definition}

\noindent
Dans la situation (1), cela signifie que les diagrammes
\begin{equation}
\label{groupoids-equation-action}
\vcenter{
\xymatrix{
G \times_S G \times_S X \ar[r]_-{1_G \times a} \ar[d]_{m \times 1_X} &
G \times_S X \ar[d]^a \\
G \times_S X \ar[r]^a & X
}
}
\quad\quad
\vcenter{
\xymatrix{
G \times_S X \ar[r]_-a & X \\
X\ar[u]^{e \times 1_X} \ar[ru]_{1_X}
}
}
\end{equation}
sont commutatifs. Dans la situation (2), cela signifie simplement que le diagramme
$$
\xymatrix{
G \times_S X \ar[r]_-{\text{id} \times \psi} \ar[d]_a &
G \times_S Y \ar[d]^a \\
X \ar[r]^\psi & Y
}
$$
est commutatif.

\begin{definition}
\label{groupoids-definition-free-action}
Soient $S$, $G \to S$ et $X \to S$ comme dans la
définition \ref{groupoids-definition-action-group-scheme}.
Soit $a : G \times_S X \to X$ une action de $G$ sur $X/S$.
Nous disons que l'action est {\it libre} si, pour tout schéma $T$ sur $S$,
l'action $a : G(T) \times X(T) \to X(T)$ est une action libre
du groupe $G(T)$ sur l'ensemble $X(T)$.
\end{definition}

\begin{lemma}
\label{groupoids-lemma-free-action}
Dans la situation de la définition \ref{groupoids-definition-free-action},
l'action $a$ est libre si et seulement si
$$
G \times_S X \to X \times_S X, \quad (g, x) \mapsto (a(g, x), x)
$$
est un monomorphisme.
\end{lemma}

\begin{proof}
Cela résulte immédiatement des définitions.
\end{proof}






\section{Torseurs et espaces homogènes principaux}
\label{groupoids-section-principal-homogeneous}

\noindent
Dans
Cohomologie sur les sites, définition \ref{sites-cohomology-definition-torsor},
nous avons défini un torseur sous un faisceau de groupes sur un site.
Supposons que $\tau \in\allowbreak \{Zariski,\allowbreak \etale,\allowbreak smooth,\allowbreak syntomic,\allowbreak fppf\}$ soit une
topologie et que $(G, m)$ soit un schéma en groupes sur $S$. Puisque $\tau$ est plus fine que
la topologie canonique (voir
Descente, lemme \ref{descent-lemma-fpqc-universal-effective-epimorphisms}),
nous voyons que $\underline{G}$ (voir
Sites, définition \ref{sites-definition-representable-sheaf})
est un faisceau de groupes sur $(\Sch/S)_\tau$.
Nous savons donc déjà ce que signifie être un
torseur sous $\underline{G}$ sur $(\Sch/S)_\tau$. Une situation particulière
se présente lorsque ce faisceau est représentable. Dans les définitions suivantes,
nous définissons directement ce que signifie, pour le schéma qui le représente, être un
torseur sous $G$.

\begin{definition}
\label{groupoids-definition-pseudo-torsor}
Soit $S$ un schéma.
Soit $(G, m)$ un schéma en groupes sur $S$.
Soit $X$ un schéma sur $S$, et soit
$a : G \times_S X \to X$ une action de $G$ sur $X$.
\begin{enumerate}
\item Nous disons que $X$ est un {\it pseudo-torseur sous $G$} ou que $X$ est
{\it formellement homogène principal sous $G$} si le
morphisme induit de schémas $G \times_S X \to X \times_S X$,
$(g, x) \mapsto (a(g, x), x)$ est un isomorphisme de schémas sur $S$.
\item {\emergencystretch=1em Un pseudo-torseur $X$ sous $G$ est dit {\it trivial} s'il existe
un isomorphisme $G$-équivariant $G \to X$ sur $S$, où $G$ agit sur
$G$ par multiplication à gauche.\par}
\end{enumerate}
\end{definition}

\noindent
Il est clair que, si $S' \to S$ est un morphisme de schémas, alors
le changement de base $X_{S'}$ d'un pseudo-torseur sous $G$ sur $S$ est un
pseudo-torseur sous $G_{S'}$ sur $S'$.

\begin{lemma}
\label{groupoids-lemma-characterize-trivial-pseudo-torsors}
Dans la situation de la
définition \ref{groupoids-definition-pseudo-torsor}.
\begin{enumerate}
\item Le schéma $X$ est un pseudo-torseur sous $G$ si et seulement si, pour tout schéma
$T$ sur $S$, l'ensemble $X(T)$ est vide ou l'action du groupe $G(T)$
sur $X(T)$ est simplement transitive.
\item Un pseudo-torseur $X$ sous $G$ est trivial si et seulement si le morphisme
$X \to S$ possède une section.
\end{enumerate}
\end{lemma}

\begin{proof}
Démonstration omise.
\end{proof}

\begin{definition}
\label{groupoids-definition-principal-homogeneous-space}
Soit $S$ un schéma.
Soit $(G, m)$ un schéma en groupes sur $S$.
Soit $X$ un pseudo-torseur sous $G$ sur $S$.
\begin{enumerate}
\item Nous disons que $X$ est un {\it espace homogène principal}
ou un {\it torseur sous $G$} s'il existe un recouvrement fpqc\footnote{Cela signifie
que, par défaut, un torseur est un pseudo-torseur trivial sur un
recouvrement fpqc. C'est la définition de \cite[Expos\'e IV, 6.5]{SGA3}.
Elle est quelque peu malcommode pour nous, puisque nous travaillons le plus souvent pour la
topologie fppf.}
$\{S_i \to S\}_{i \in I}$ tel que chaque
$X_{S_i} \to S_i$ possède une section (c'est-à-dire qu'il est un pseudo-torseur trivial sous $G_{S_i}$).
\item Soit $\tau \in\allowbreak \{Zariski,\allowbreak \etale,\allowbreak smooth,\allowbreak syntomic,\allowbreak fppf\}$.
Nous disons que $X$ est un {\it torseur sous $G$ pour la topologie $\tau$}, ou un
{\it $\tau$-torseur sous $G$}, ou simplement un {\it $\tau$-torseur},
s'il existe un recouvrement $\tau$ $\{S_i \to S\}_{i \in I}$
tel que chaque $X_{S_i} \to S_i$ possède une section.
\item Si $X$ est un torseur sous $G$, nous disons qu'il est
{\it quasi-isotrivial} s'il est un torseur pour la topologie \'etale.
\item Si $X$ est un torseur sous $G$, nous disons qu'il est
{\it localement trivial} s'il est un torseur pour la topologie de Zariski.
\end{enumerate}
\end{definition}

\noindent
Nous disons parfois « soit $X$ un torseur sous $G$ sur $S$ » pour indiquer que
$X$ est un schéma sur $S$ muni d'une action de $G$ qui en fait
un espace homogène principal sur $S$.
Montrons maintenant que cela concorde avec la terminologie introduite plus haut
lorsque les deux s'appliquent.

\begin{lemma}
\label{groupoids-lemma-torsor}
Soit $S$ un schéma.
Soit $(G, m)$ un schéma en groupes sur $S$.
Soit $X$ un schéma sur $S$, et soit
$a : G \times_S X \to X$ une action de $G$ sur $X$.
Soit $\tau \in\allowbreak \{Zariski,\allowbreak \etale,\allowbreak smooth,\allowbreak syntomic,\allowbreak fppf\}$.
Alors $X$ est un torseur sous $G$ pour la topologie $\tau$ si et seulement si
$\underline{X}$ est un torseur sous $\underline{G}$ sur $(\Sch/S)_\tau$.
\end{lemma}

\begin{proof}
Démonstration omise.
\end{proof}

\begin{remark}
\label{groupoids-remark-fun-with-torsors}
Soit $(G, m)$ un schéma en groupes sur le schéma $S$.
Dans cette situation, les types naturels de questions suivants se posent :
\begin{enumerate}
\item Si $X \to S$ est un pseudo-torseur sous $G$ et si $X \to S$ est surjectif,
alors $X$ est-il nécessairement un torseur sous $G$ ?
\item Tout torseur sous $\underline{G}$ sur $(\Sch/S)_{fppf}$ est-il
représentable ? Autrement dit, tout torseur sous $\underline{G}$
provient-il d'un torseur fppf sous $G$ ?
\item Tout torseur sous $G$ est-il un
torseur fppf (resp.\ lisse, resp.\ \'etale, resp.\ de Zariski) ?
\end{enumerate}
En général, les réponses à ces questions sont négatives. Pour obtenir une réponse positive,
nous devons imposer des conditions supplémentaires à $G \to S$.
Par exemple :
si $S$ est le spectre d'un corps, alors la réponse à (1) est affirmative
car $\{X \to S\}$ est alors un recouvrement fpqc qui trivialise $X$.
Si $G \to S$ est affine, alors la réponse à (2) est affirmative
(cela résulte de Descente, lemme \ref{descent-lemma-affine}).
Si $G = \text{GL}_{n, S}$, alors la réponse à (3) est affirmative
et, en fait, tout torseur sous $\text{GL}_{n, S}$ est localement trivial
(cela résulte de
Descente, lemme \ref{descent-lemma-finite-locally-free-descends}).
\end{remark}



\section{Faisceaux quasi-cohérents équivariants}
\label{groupoids-section-equivariant}

\noindent
Nous considérons les « fonctions » comme duales de l'« espace ». Ainsi, pour un morphisme d'espaces,
l'application induite sur les fonctions va dans le sens opposé. De plus, nous considérons les
sections d'un faisceau de modules comme des « fonctions ». Cela nous conduit naturellement
au sens des flèches choisi dans la définition suivante.

\begin{definition}
\label{groupoids-definition-equivariant-module}
Soit $S$ un schéma, soit $(G, m)$ un schéma en groupes sur $S$, et
soit $a : G \times_S X \to X$ une action du schéma en groupes $G$
sur $X/S$. Un {\it module $G$-équivariant quasi-cohérent sur $\mathcal{O}_X$},
ou simplement un {\it module quasi-cohérent équivariant sur $\mathcal{O}_X$},
est une paire $(\mathcal{F}, \alpha)$, où $\mathcal{F}$ est un module quasi-cohérent
sur $\mathcal{O}_X$, et $\alpha$ est un morphisme de $\mathcal{O}_{G \times_S X}$-modules
donné par
$$
\alpha : a^*\mathcal{F} \longrightarrow \text{pr}_1^*\mathcal{F}
$$
où $\text{pr}_1 : G \times_S X \to X$ est la projection,
tel que
\begin{enumerate}
\item le diagramme
$$
\xymatrix{
(1_G \times a)^*\text{pr}_1^*\mathcal{F} \ar[r]_-{\text{pr}_{12}^*\alpha} &
\text{pr}_2^*\mathcal{F} \\
(1_G \times a)^*a^*\mathcal{F} \ar[u]^{(1_G \times a)^*\alpha} \ar@{=}[r] &
(m \times 1_X)^*a^*\mathcal{F} \ar[u]_{(m \times 1_X)^*\alpha}
}
$$
soit commutatif dans la catégorie des
$\mathcal{O}_{G \times_S G \times_S X}$-modules, et
\item le tiré en arrière
$$
(e \times 1_X)^*\alpha : \mathcal{F} \longrightarrow \mathcal{F}
$$
soit l'identité.
\end{enumerate}
Pour une explication, comparer avec les diagrammes correspondants de
l'équation (\ref{groupoids-equation-action}).
\end{definition}

\noindent
Notons que la commutativité du premier diagramme garantit que
$(e \times 1_X)^*\alpha$ est un endomorphisme idempotent de $\mathcal{F}$,
de sorte que la condition (2) signifie simplement que c'est un isomorphisme.

\begin{lemma}
\label{groupoids-lemma-pullback-equivariant}
Soit $S$ un schéma. Soit $G$ un schéma en groupes sur $S$.
Soit $f : Y \to X$ un morphisme $G$-équivariant entre des $S$-schémas
munis d'actions de $G$. La règle
$(\mathcal{F}, \alpha) \mapsto (f^*\mathcal{F}, (1_G \times f)^*\alpha)$
définit un foncteur de la catégorie des modules quasi-cohérents $G$-équivariants
sur $\mathcal{O}_X$ vers la catégorie des
modules quasi-cohérents $G$-équivariants sur $\mathcal{O}_Y$.
\end{lemma}

\begin{proof}
Démonstration omise.
\end{proof}

\noindent
Donnons un exemple.

\begin{example}
\label{groupoids-example-Gm-on-affine}
Soit $A$ un anneau $\mathbf{Z}$-gradué, c'est-à-dire que $A$ est muni d'une décomposition en
somme directe $A = \bigoplus_{n \in \mathbf{Z}} A_n$ et que
$A_n \cdot A_m \subset A_{n + m}$.
Posons $X = \Spec(A)$. Nous obtenons alors une action de $\mathbf{G}_m$
$$
a : \mathbf{G}_m \times X \longrightarrow X
$$
à partir de l'homomorphisme d'anneaux $\mu : A \to A \otimes \mathbf{Z}[x, x^{-1}]$,
$f \mapsto f \otimes x^{\deg(f)}$.
En effet, pour le vérifier, nous devons vérifier que
$$
\xymatrix{
A \ar[r]_\mu \ar[d]_\mu &
A \otimes \mathbf{Z}[x, x^{-1}] \ar[d]^{\mu \otimes 1} \\
A \otimes \mathbf{Z}[x, x^{-1}] \ar[r]^-{1 \otimes m} &
A \otimes \mathbf{Z}[x, x^{-1}] \otimes \mathbf{Z}[x, x^{-1}]
}
$$
où $m(x) = x \otimes x$, voir l'exemple \ref{groupoids-example-multiplicative-group}.
Cela est immédiatement clair en évaluant sur un élément homogène.
Supposons que $M$ soit un $A$-module gradué. Nous obtenons alors un
module quasi-cohérent $\mathbf{G}_m$-équivariant sur $\mathcal{O}_X$
$\mathcal{F} = \widetilde{M}$ en prenant $\alpha$ comme dans la
définition \ref{groupoids-definition-equivariant-module}, correspondant au
morphisme de $A \otimes \mathbf{Z}[x, x^{-1}]$-modules
$$
M \otimes_{A, \mu} (A \otimes \mathbf{Z}[x, x^{-1}])
\longrightarrow
M \otimes_{A, \text{id}_A \otimes 1} (A \otimes \mathbf{Z}[x, x^{-1}])
$$
qui envoie $m \otimes 1 \otimes 1$ sur $m \otimes 1 \otimes x^{\deg(m)}$
pour $m \in M$ homogène.
\end{example}

\begin{lemma}
\label{groupoids-lemma-complete-reducibility-Gm}
Soit $a : \mathbf{G}_m \times X \to X$ une action sur un schéma affine.
Alors $X$ est le spectre d'un anneau $\mathbf{Z}$-gradué
et l'action est celle de l'exemple \ref{groupoids-example-Gm-on-affine}.
\end{lemma}

\begin{proof}
Soit $f \in A = \Gamma(X, \mathcal{O}_X)$. Alors nous pouvons écrire
$$
a^\sharp(f) = \sum\nolimits_{n \in \mathbf{Z}} f_n \otimes x^n
\quad\text{dans}\quad
A \otimes \mathbf{Z}[x, x^{-1}] =
\Gamma(\mathbf{G}_m \times X, \mathcal{O}_{\mathbf{G}_m \times X})
$$
sous la forme d'une somme finie, les $f_n$ de $A$ étant déterminés de façon unique.
Nous obtenons ainsi des applications $A \to A$, $f \mapsto f_n$.
Puisque $a$ est une action, en évaluant en $x = 1$,
nous obtenons $f = \sum f_n$. Puisque $a$ est une action,
nous trouvons que
$$
\sum (f_n)_m \otimes x^m \otimes x^n = \sum f_n x^n \otimes x^n
$$
(comparer avec le calcul de l'exemple \ref{groupoids-example-Gm-on-affine}).
Ainsi $(f_n)_m = 0$ si $n \not = m$ et $(f_n)_n = f_n$.
Ainsi, si nous posons
$$
A_n = \{f \in A \mid f_n = f\}
$$
alors nous obtenons $A = \sum A_n$. En revanche, la somme est nécessairement
directe puisque $f = 0$ implique $f_n = 0$ dans la situation ci-dessus.
\end{proof}

\begin{lemma}
\label{groupoids-lemma-Gm-equivariant-module}
Soit $A$ un anneau gradué. Soit $X = \Spec(A)$ muni de l'action
$a : \mathbf{G}_m \times X \to X$ de l'exemple \ref{groupoids-example-Gm-on-affine}.
Soit $\mathcal{F}$ un module quasi-cohérent $\mathbf{G}_m$-équivariant
sur $\mathcal{O}_X$. Alors $M = \Gamma(X, \mathcal{F})$
possède une graduation canonique qui en fait un $A$-module gradué
et telle que l'isomorphisme $\widetilde{M} \to \mathcal{F}$
(Schémas, lemme \ref{schemes-lemma-quasi-coherent-affine})
soit un isomorphisme de modules $\mathbf{G}_m$-équivariants, où
la structure $\mathbf{G}_m$-équivariante sur $\widetilde{M}$
est celle de l'exemple \ref{groupoids-example-Gm-on-affine}.
\end{lemma}

\begin{proof}
On peut soit démontrer ceci en répétant les arguments du
lemme \ref{groupoids-lemma-complete-reducibility-Gm} pour le module $M$.
On peut aussi considérer le schéma
$(X', \mathcal{O}_{X'}) = (X, \mathcal{O}_X \oplus \mathcal{F})$,
où $\mathcal{F}$ est considéré comme un idéal de carré nul.
Il existe une action naturelle $a' : \mathbf{G}_m \times X' \to X'$
définie à l'aide de l'action sur $X$ et sur $\mathcal{F}$. Appliquons alors
le lemme \ref{groupoids-lemma-complete-reducibility-Gm} à $X'$ et concluons.
(L'intérêt de cet argument est qu'il montre immédiatement
que les graduations de $A$ et de $M$ sont compatibles, c'est-à-dire que $M$
est un $A$-module gradué.)
Détails omis.
\end{proof}



\section{Groupoïdes}
\label{groupoids-section-groupoids}

\noindent
Rappelons qu'un groupoïde est une catégorie dans laquelle tout morphisme
est un isomorphisme ; voir
Catégories, définition \ref{categories-definition-groupoid}.
Ainsi, un groupoïde possède un ensemble d'objets $\text{Ob}$,
un ensemble de flèches $\text{Arrows}$, des applications {\it source} et {\it but}
$s, t : \text{Arrows} \to \text{Ob}$, ainsi qu'une {\it loi de composition}
$c : \text{Arrows} \times_{s, \text{Ob}, t} \text{Arrows}
\to \text{Arrows}$.
Ces applications satisfont exactement aux axiomes suivants :
\begin{enumerate}
\item {\emergencystretch=2em (associativité) $c \circ (1, c) = c \circ (c, 1)$ comme applications
$\text{Arrows} \times_{s, \text{Ob}, t}
\text{Arrows} \times_{s, \text{Ob}, t}
\text{Arrows} \to \text{Arrows}$,\par}
\item (identité) il existe une application $e : \text{Ob} \to \text{Arrows}$
telle que
\begin{enumerate}
\item $s \circ e = t \circ e = \text{id}$ comme applications $\text{Ob} \to \text{Ob}$,
\item $c \circ (1, e \circ s) = c \circ (e \circ t, 1) = 1$
comme applications $\text{Arrows} \to \text{Arrows}$,
\end{enumerate}
\item (inverse) il existe une application $i : \text{Arrows} \to \text{Arrows}$
telle que
\begin{enumerate}
\item $s \circ i = t$, $t \circ i = s$ comme applications $\text{Arrows} \to \text{Ob}$,
et
\item $c \circ (1, i) = e \circ t$ et $c \circ (i, 1) = e \circ s$
comme applications $\text{Arrows} \to \text{Arrows}$.
\end{enumerate}
\end{enumerate}
Dans ce cas, les applications $e$ et $i$ sont déterminées de manière unique et
$i$ est une bijection. Notons que si $(\text{Ob}', \text{Arrows}', s', t', c')$
est une seconde catégorie groupoïde, alors un foncteur
$f : (\text{Ob}, \text{Arrows}, s, t, c) \to
(\text{Ob}', \text{Arrows}', s', t', c')$
est défini par un couple d'applications d'ensembles $f : \text{Ob} \to \text{Ob}'$ et
$f : \text{Arrows} \to \text{Arrows}'$ telles que
$s' \circ f = f \circ s$, $t' \circ f = f \circ t$ et
$c' \circ (f, f) = f \circ c$. La compatibilité avec l'identité et
l'inverse est automatique. Nous utiliserons ce fait ci-dessous.
(Attention : la compatibilité avec l'identité
doit être imposée dans le cas des catégories générales.)

\begin{definition}
\label{groupoids-definition-groupoid}
Soit $S$ un schéma.
\begin{enumerate}
\item Un {\it schéma en groupoïdes sur $S$}, ou simplement un
{\it groupoïde sur $S$}, est un
quintuplet $(U, R, s, t, c)$, où
$U$ et $R$ sont des schémas sur $S$ et
$s, t : R \to U$ et $c : R \times_{s, U, t} R \to R$
sont des morphismes de schémas sur $S$ ayant la
propriété suivante : pour tout schéma
$T$ sur $S$, le quintuplet
$$
(U(T), R(T), s, t, c)
$$
est une catégorie groupoïde au sens décrit ci-dessus.
\item Un {\it morphisme
$f : (U, R, s, t, c) \to (U', R', s', t', c')$
de schémas en groupoïdes sur $S$} est défini par des morphismes
de schémas $f : U \to U'$ et $f : R \to R'$ ayant la
propriété suivante : pour tout schéma
$T$ sur $S$, les applications $f$ définissent un foncteur de la
catégorie groupoïde $(U(T), R(T), s, t, c)$ vers la
catégorie groupoïde $(U'(T), R'(T), s', t', c')$.
\end{enumerate}
\end{definition}

\noindent
Soit $(U, R, s, t, c)$ un groupoïde sur $S$.
Notons que, d'après les remarques précédant la définition et le lemme de Yoneda,
il existe d'uniques morphismes de schémas
$e : U \to R$ et
$i : R \to R$ sur $S$ tels que, pour tout schéma $T$ sur $S$,
l'application induite $e : U(T) \to R(T)$ soit l'application des identités et que
$i : R(T) \to R(T)$ soit l'application des inverses de
la catégorie groupoïde. Le septuplet $(U, R, s, t, c, e, i)$
satisfait aux diagrammes commutatifs correspondant à chacun des
axiomes (1), (2)(a), (2)(b), (3)(a) et (3)(b) ci-dessus et, réciproquement,
étant donné un septuplet ayant cette propriété, le quintuplet $(U, R, s, t, c)$
est un schéma en groupoïdes. Notons que $i$ est un isomorphisme
et que $e$ est une section à la fois de $s$ et de $t$.
En outre, étant donné un schéma en groupoïdes sur $S$, on note
$$
j = (t, s) : R \longrightarrow U \times_S U
$$
ce qui est compatible avec nos conventions de la section
\ref{groupoids-section-equivalence-relations} ci-dessus.
On dit parfois « soit $(U, R, s, t, c, e, i)$ un
groupoïde sur $S$ » pour souligner l'existence des identités et des
inverses.

\begin{lemma}
\label{groupoids-lemma-groupoid-pre-equivalence}
Étant donné un schéma en groupoïdes $(U, R, s, t, c)$ sur $S$,
le morphisme $j : R \to U \times_S U$ est une prérelation
d'équivalence.
\end{lemma}

\begin{proof}
La preuve est omise.
C'est un joli exercice sur les définitions.
\end{proof}

\begin{lemma}
\label{groupoids-lemma-equivalence-groupoid}
Étant donné une relation d'équivalence $j : R \to U \times_S U$ sur $S$,
il existe une unique manière de l'étendre en un groupoïde
$(U, R, s, t, c)$ sur $S$.
\end{lemma}

\begin{proof}
La preuve est omise.
C'est un joli exercice sur les définitions.
\end{proof}

\begin{lemma}
\label{groupoids-lemma-diagram}
Soit $S$ un schéma.
Soit $(U, R, s, t, c)$ un groupoïde sur $S$.
Dans le diagramme commutatif
$$
\xymatrix{
& U & \\
R \ar[d]_s \ar[ru]^t &
R \times_{s, U, t} R
\ar[l]^-{\text{pr}_0} \ar[d]^{\text{pr}_1} \ar[r]_-c &
R \ar[d]^s \ar[lu]_t \\
U & R \ar[l]_t \ar[r]^s & U
}
$$
les deux carrés inférieurs sont des carrés cartésiens.
De plus, le triangle supérieur (qui est en réalité un carré)
est lui aussi cartésien.
\end{lemma}

\begin{proof}
La preuve est omise.
C'est un exercice sur les définitions et le point de
vue fonctoriel en géométrie algébrique.
\end{proof}

\begin{lemma}
\label{groupoids-lemma-diagram-pull}
Soit $S$ un schéma.
Soit $(U, R, s, t, c, e, i)$ un groupoïde sur $S$.
Le diagramme
\begin{equation}
\label{groupoids-equation-pull}
\xymatrix{
R \times_{t, U, t} R
\ar@<1ex>[r]^-{\text{pr}_1} \ar@<-1ex>[r]_-{\text{pr}_0}
\ar[d]_{(\text{pr}_0, c \circ (i, 1))} &
R \ar[r]^t \ar[d]^{\text{id}_R} &
U \ar[d]^{\text{id}_U} \\
R \times_{s, U, t} R
\ar@<1ex>[r]^-c \ar@<-1ex>[r]_-{\text{pr}_0} \ar[d]_{\text{pr}_1} &
R \ar[r]^t \ar[d]^s &
U \\
R \ar@<1ex>[r]^s \ar@<-1ex>[r]_t &
U
}
\end{equation}
est commutatif. Les deux lignes supérieures sont isomorphes au moyen des applications verticales indiquées.
Les deux carrés inférieurs de gauche sont cartésiens.
\end{lemma}

\begin{proof}
La commutativité du diagramme résulte des axiomes d'un groupoïde.
Notons qu'en termes de groupoïdes, la flèche verticale supérieure gauche associe
à une paire de morphismes $(\alpha, \beta)$ de même but la paire
de morphismes $(\alpha, \alpha^{-1} \circ \beta)$. Dans tout groupoïde,
ceci définit une bijection entre
$\text{Arrows} \times_{t, \text{Ob}, t} \text{Arrows}$
et
$\text{Arrows} \times_{s, \text{Ob}, t} \text{Arrows}$. D'où la deuxième
assertion du lemme.
La dernière assertion résulte du lemme \ref{groupoids-lemma-diagram}.
\end{proof}

\begin{lemma}
\label{groupoids-lemma-base-change-groupoid}
Soit $(U, R, s, t, c)$ un groupoïde sur le schéma $S$.
Soit $S' \to S$ un morphisme. Alors les changements de base $U' = S' \times_S U$,
$R' = S' \times_S R$, munis des changements de base $s'$, $t'$, $c'$
des morphismes $s, t, c$, forment un schéma en groupoïdes
$(U', R', s', t', c')$ sur $S'$, et les projections
déterminent un morphisme
$(U', R', s', t', c') \to (U, R, s, t, c)$
de schémas en groupoïdes sur $S$.
\end{lemma}

\begin{proof}
La preuve est omise. Indication :
$R' \times_{s', U', t'} R' = S' \times_S (R \times_{s, U, t} R)$.
\end{proof}









\section{Faisceaux quasi-cohérents sur les groupoïdes}
\label{groupoids-section-groupoids-quasi-coherent}

\noindent
Voir l'introduction de la section \ref{groupoids-section-equivariant} pour nos
choix concernant le sens des flèches.

\begin{definition}
\label{groupoids-definition-groupoid-module}
Soit $S$ un schéma et soit $(U, R, s, t, c)$ un schéma en groupoïdes sur $S$.
Un {\it module quasi-cohérent sur $(U, R, s, t, c)$}
est une paire $(\mathcal{F}, \alpha)$, où $\mathcal{F}$ est un
$\mathcal{O}_U$-module quasi-cohérent, et $\alpha$ est un homomorphisme de $\mathcal{O}_R$-modules
donné par
$$
\alpha : t^*\mathcal{F} \longrightarrow s^*\mathcal{F}
$$
tel que
\begin{enumerate}
\item le diagramme
$$
\xymatrix{
& \text{pr}_1^*t^*\mathcal{F} \ar[r]_-{\text{pr}_1^*\alpha} &
\text{pr}_1^*s^*\mathcal{F} \ar@{=}[rd] & \\
\text{pr}_0^*s^*\mathcal{F} \ar@{=}[ru] & & & c^*s^*\mathcal{F} \\
& \text{pr}_0^*t^*\mathcal{F} \ar[lu]^{\text{pr}_0^*\alpha} \ar@{=}[r] &
c^*t^*\mathcal{F} \ar[ru]_{c^*\alpha}
}
$$
soit commutatif dans la catégorie des
$\mathcal{O}_{R \times_{s, U, t} R}$-modules, et
\item l'image inverse
$$
e^*\alpha : \mathcal{F} \longrightarrow \mathcal{F}
$$
soit l'identité.
\end{enumerate}
Comparer avec les diagrammes commutatifs du lemme \ref{groupoids-lemma-diagram}.
\end{definition}

\noindent
La commutativité du premier diagramme force l'opérateur $e^*\alpha$
à être idempotent. La seconde condition peut donc être reformulée en disant
que $e^*\alpha$ est un isomorphisme. En fait, cette condition implique que
$\alpha$ est un isomorphisme.

\begin{lemma}
\label{groupoids-lemma-isomorphism}
Soit $S$ un schéma et soit $(U, R, s, t, c)$ un schéma en groupoïdes sur $S$.
Si $(\mathcal{F}, \alpha)$ est un module quasi-cohérent sur $(U, R, s, t, c)$,
alors $\alpha$ est un isomorphisme.
\end{lemma}

\begin{proof}
Prenons l'image inverse du diagramme commutatif de la
définition \ref{groupoids-definition-groupoid-module}
par le morphisme $(i, 1) : R \to R \times_{s, U, t} R$.
On voit alors que $i^*\alpha \circ \alpha = s^*e^*\alpha$.
En prenant l'image inverse par le morphisme $(1, i)$, on obtient la relation
$\alpha \circ i^*\alpha = t^*e^*\alpha$. D'après la seconde hypothèse,
ces morphismes sont l'identité. Ainsi, $i^*\alpha$ est un inverse de
$\alpha$.
\end{proof}

\begin{lemma}
\label{groupoids-lemma-pullback}
Soit $S$ un schéma. Considérons un morphisme
$f : (U, R, s, t, c) \to (U', R', s', t', c')$
de schémas en groupoïdes sur $S$. Alors l'image inverse $f^*$ donnée par
$$
(\mathcal{F}, \alpha) \mapsto (f^*\mathcal{F}, f^*\alpha)
$$
définit un foncteur de la catégorie des faisceaux quasi-cohérents sur
$(U', R', s', t', c')$ vers la catégorie des faisceaux quasi-cohérents sur
$(U, R, s, t, c)$.
\end{lemma}

\begin{proof}
La preuve est omise.
\end{proof}

\begin{lemma}
\label{groupoids-lemma-pushforward}
Soit $S$ un schéma. Considérons un morphisme
$f : (U, R, s, t, c) \to (U', R', s', t', c')$
de schémas en groupoïdes sur $S$. Supposons que
\begin{enumerate}
\item $f : U \to U'$ soit quasi-compact et quasi-séparé,
\item le carré
$$
\xymatrix{
R \ar[d]_t \ar[r]_f & R' \ar[d]^{t'} \\
U \ar[r]^f & U'
}
$$
soit cartésien, et
\item $s'$ et $t'$ soient plats.
\end{enumerate}
Alors l'image directe $f_*$ donnée par
$$
(\mathcal{F}, \alpha) \mapsto (f_*\mathcal{F}, f_*\alpha)
$$
définit un foncteur de la catégorie des faisceaux quasi-cohérents sur
$(U, R, s, t, c)$ vers la catégorie des faisceaux quasi-cohérents sur
$(U', R', s', t', c')$, adjoint à droite de l'image inverse définie dans le
lemme \ref{groupoids-lemma-pullback}.
\end{lemma}

\begin{proof}
Puisque $U \to U'$ est quasi-compact et quasi-séparé, on voit que
$f_*$ transforme les faisceaux quasi-cohérents en faisceaux quasi-cohérents
(Schémas, lemme \ref{schemes-lemma-push-forward-quasi-coherent}).
De plus, puisque les carrés
$$
\vcenter{
\xymatrix{
R \ar[d]_t \ar[r]_f & R' \ar[d]^{t'} \\
U \ar[r]^f & U'
}
}
\quad\text{et}\quad
\vcenter{
\xymatrix{
R \ar[d]_s \ar[r]_f & R' \ar[d]^{s'} \\
U \ar[r]^f & U'
}
}
$$
sont cartésiens, on obtient $(t')^*f_*\mathcal{F} = f_*t^*\mathcal{F}$
et $(s')^*f_*\mathcal{F} = f_*s^*\mathcal{F}$, voir
Cohomologie des schémas, lemme
\ref{coherent-lemma-flat-base-change-cohomology}.
Il est donc légitime de considérer $f_*\alpha$ comme un homomorphisme
$(t')^*f_*\mathcal{F} \to (s')^*f_*\mathcal{F}$. Un argument semblable
montre que $f_*\alpha$ satisfait à la condition de cocycle.
Ce foncteur est adjoint au foncteur image inverse, car les foncteurs image inverse
et image directe sur les modules d'espaces annelés sont adjoints.
Certains détails sont omis.
\end{proof}

\begin{lemma}
\label{groupoids-lemma-colimits}
Soit $S$ un schéma. Soit $(U, R, s, t, c)$ un schéma en groupoïdes sur $S$.
La catégorie des modules quasi-cohérents sur $(U, R, s, t, c)$ admet des colimites.
\end{lemma}

\begin{proof}
Soit $i \mapsto (\mathcal{F}_i, \alpha_i)$ un diagramme sur la catégorie
d'indices $\mathcal{I}$. On peut former la colimite
$\mathcal{F} = \colim \mathcal{F}_i$,
qui est un faisceau quasi-cohérent sur $U$, voir
Schémas, section \ref{schemes-section-quasi-coherent}.
Puisque l'image inverse commute aux colimites, on a
$s^*\mathcal{F} = \colim s^*\mathcal{F}_i$ et, de même,
$t^*\mathcal{F} = \colim t^*\mathcal{F}_i$. On peut donc poser
$\alpha = \colim \alpha_i$. Nous omettons la preuve que $(\mathcal{F}, \alpha)$
est la colimite du diagramme dans la catégorie des modules quasi-cohérents
sur $(U, R, s, t, c)$.
\end{proof}

\begin{lemma}
\label{groupoids-lemma-abelian}
Soit $S$ un schéma.
Soit $(U, R, s, t, c)$ un schéma en groupoïdes sur $S$.
Si $s$ et $t$ sont plats, alors la catégorie des modules quasi-cohérents sur
$(U, R, s, t, c)$ est abélienne.
\end{lemma}

\begin{proof}
Soit $\varphi : (\mathcal{F}, \alpha) \to (\mathcal{G}, \beta)$ un
homomorphisme de modules quasi-cohérents sur $(U, R, s, t, c)$. Puisque
$s$ est plat, on voit que
$$
0 \to s^*\Ker(\varphi)
\to s^*\mathcal{F} \to s^*\mathcal{G} \to s^*\Coker(\varphi) \to 0
$$
est exacte, de même que son analogue obtenu après image inverse par $t$. Ainsi, $\alpha$ et $\beta$
induisent des isomorphismes
$\kappa : t^*\Ker(\varphi) \to s^*\Ker(\varphi)$ et
$\lambda : t^*\Coker(\varphi) \to s^*\Coker(\varphi)$
qui satisfont à la condition de cocycle. Il est alors immédiat de
vérifier que $(\Ker(\varphi), \kappa)$ et
$(\Coker(\varphi), \lambda)$ sont respectivement un noyau et un conoyau dans la
catégorie des modules quasi-cohérents sur $(U, R, s, t, c)$. En outre,
la relation $\Coim(\varphi) = \Im(\varphi)$ résulte
du fait qu'elle est satisfaite sur $U$.
\end{proof}













\section{Colimites de modules quasi-cohérents}
\label{groupoids-section-colimits}

\noindent
Dans cette section, nous démontrons quelques résultats techniques affirmant que, sous
des hypothèses convenables, tout module quasi-cohérent sur un groupoïde est
une colimite filtrante de « petits » modules quasi-cohérents.

\begin{lemma}
\label{groupoids-lemma-construct-quasi-coherent}
Soit $(U, R, s, t, c)$ un schéma en groupoïdes sur $S$.
Supposons que $s, t$ soient plats, quasi-compacts et quasi-séparés.
Pour tout module quasi-cohérent $\mathcal{G}$ sur $U$, il existe
un isomorphisme canonique
$\alpha : t^*s_*t^*\mathcal{G} \to s^*s_*t^*\mathcal{G}$
qui fait de $(s_*t^*\mathcal{G}, \alpha)$ un module quasi-cohérent
sur $(U, R, s, t, c)$. Cette construction définit un foncteur
$$
\QCoh(\mathcal{O}_U) \longrightarrow \QCoh(U, R, s, t, c)
$$
qui est adjoint à droite du foncteur d'oubli
$(\mathcal{F}, \beta) \mapsto \mathcal{F}$.
\end{lemma}

\begin{proof}
L'image directe d'un module quasi-cohérent par un morphisme quasi-compact et
quasi-séparé est quasi-cohérente, voir Schémas, lemme
\ref{schemes-lemma-push-forward-quasi-coherent}. Ainsi, $s_*t^*\mathcal{G}$
est quasi-cohérent. Avec les notations du lemme \ref{groupoids-lemma-diagram}, on a
$$
t^*s_*t^*\mathcal{G} =
\text{pr}_{1, *}\text{pr}_0^*t^*\mathcal{G} =
\text{pr}_{1, *}c^*t^*\mathcal{G} =
s^*s_*t^*\mathcal{G}
$$
L'égalité du milieu vient de ce que $t \circ c = t \circ \text{pr}_0$ comme
morphismes $R \times_{s, U, t} R \to U$, tandis que les première et dernière
égalités viennent de ce que l'on sait que le changement de base et l'image directe commutent
à ces étapes, d'après Cohomologie des schémas, lemme
\ref{coherent-lemma-flat-base-change-cohomology}.

\medskip\noindent
Pour vérifier la condition de cocycle de la définition \ref{groupoids-definition-groupoid-module}
pour $\alpha$, ainsi que la propriété d'adjonction, décrivons la construction
$\mathcal{G} \mapsto (s_*t^*\mathcal{G}, \alpha)$ d'une autre manière.
Considérons le schéma en groupoïdes
$(R, R \times_{t, U, t} R, \text{pr}_0, \text{pr}_1, \text{pr}_{02})$
associé à la relation d'équivalence $R \times_{t, U, t} R$
sur $R$, voir le lemme \ref{groupoids-lemma-equivalence-groupoid}.
Il existe un morphisme
$$
f :
(R, R \times_{t, U, t} R, \text{pr}_1, \text{pr}_0, \text{pr}_{02})
\longrightarrow
(U, R, s, t, c)
$$
de schémas en groupoïdes donné par $s : R \to U$ et par l'application $R \times_{t, U, t} R \to R$
qui est définie par $(r_0, r_1) \mapsto r_0^{-1} \circ r_1$ ; nous omettons la vérification
de la commutativité des diagrammes voulus. Puisque
$t, s : R \to U$ sont quasi-compacts, quasi-séparés et plats,
et puisque l'on a un carré cartésien
$$
\xymatrix{
R \times_{t, U, t} R \ar[d]_{\text{pr}_0}
\ar[rr]_-{(r_0, r_1) \mapsto r_0^{-1} \circ r_1} & & R \ar[d]^t \\
R \ar[rr]^s & & U
}
$$
Le lemme \ref{groupoids-lemma-diagram-pull} montre que ce carré est cartésien ; par conséquent, le
lemme \ref{groupoids-lemma-pushforward} s'applique à $f$.
Ainsi, les foncteurs image directe et image inverse des modules quasi-cohérents associés à
$f$ sont adjoints. Pour terminer la preuve, nous allons identifier ces foncteurs
aux foncteurs décrits ci-dessus. Pour cela, notons que
$$
t^* :
\QCoh(\mathcal{O}_U)
\longrightarrow
\QCoh(R, R \times_{t, U, t} R, \text{pr}_1, \text{pr}_0, \text{pr}_{02})
$$
est une équivalence d'après la théorie de la descente des faisceaux quasi-cohérents, car
$\{t : R \to U\}$ est un recouvrement fpqc, voir
Descente, proposition \ref{descent-proposition-fpqc-descent-quasi-coherent}.

\medskip\noindent
L'image directe par $f$, précomposée par l'équivalence $t^*$, envoie
$\mathcal{G}$ sur $(s_*t^*\mathcal{G}, \alpha)$ ;
nous omettons la vérification que l'isomorphisme $\alpha$
obtenu de cette manière est le même que celui construit ci-dessus.

\medskip\noindent
L'image inverse par $f$, postcomposée par l'inverse de l'équivalence $t^*$,
envoie $(\mathcal{F}, \beta)$ sur la descente relativement à $\{t : R \to U\}$
du module $s^*\mathcal{F}$ muni de la donnée de descente $\gamma$ sur
$R \times_{t, U, t} R$, cette donnée étant l'image inverse de $\beta$ par
$(r_0, r_1) \mapsto r_0^{-1} \circ r_1$.
Considérons l'isomorphisme $\beta : t^*\mathcal{F} \to s^*\mathcal{F}$.
La donnée de descente canonique (Descente, définition
\ref{descent-definition-descent-datum-effective-quasi-coherent})
sur $t^*\mathcal{F}$ relativement à $\{t : R \to U\}$
se traduit, via $\beta$, par l'application
$$
\text{pr}_0^*s^*\mathcal{F}
\xrightarrow{\text{pr}_0^*\beta^{-1}}
\text{pr}_0^*t^*\mathcal{F}
\xrightarrow{can}
\text{pr}_1^*t^*\mathcal{F}
\xrightarrow{\text{pr}_1^*\beta}
\text{pr}_1^*s^*\mathcal{F}
$$
Puisque $\beta$ satisfait à la condition de cocycle, cette application est égale à l'image inverse
de $\beta$ par $(r_0, r_1) \mapsto r_0^{-1} \circ r_1$. Pour le voir,
prenons la relation de cocycle proprement dite dans la définition \ref{groupoids-definition-groupoid-module}
et prenons son image inverse par le morphisme
$(\text{pr}_0, c \circ (i, 1)) : R \times_{t, U, t} R \to R \times_{s, U, t} R$
qui intervient également dans le diagramme commutatif du
lemme \ref{groupoids-lemma-diagram-pull}. Il en résulte que
$(s^*\mathcal{F}, \gamma)$ est isomorphe à $(t^*\mathcal{F}, can)$.
En définitive, nous concluons que l'image inverse par $f$, postcomposée par
l'inverse de l'équivalence $t^*$, est isomorphe au foncteur d'oubli
$(\mathcal{F}, \beta) \mapsto \mathcal{F}$.
\end{proof}

\begin{remark}
\label{groupoids-remark-adjunction-map}
Dans la situation du lemme \ref{groupoids-lemma-construct-quasi-coherent}, notons
$$
F : \QCoh(U, R, s, t, c) \to \QCoh(\mathcal{O}_U),\quad
(\mathcal{F}, \beta) \mapsto \mathcal{F}
$$
le foncteur d'oubli, et notons
$$
G : \QCoh(\mathcal{O}_U) \to \QCoh(U, R, s, t, c),\quad
\mathcal{G} \mapsto (s_*t^*\mathcal{G}, \alpha)
$$
l'adjoint à droite construit dans le lemme. Alors l'unité
$\eta : \text{id} \to G \circ F$ de l'adjonction, évaluée
en $(\mathcal{F}, \beta)$, est donnée par l'application
$$
\mathcal{F} \to s_*s^*\mathcal{F} \xrightarrow{\beta^{-1}} s_*t^*\mathcal{F}
$$
Nous omettons la vérification.
\end{remark}

\begin{lemma}
\label{groupoids-lemma-push-pull}
Soit $f : Y \to X$ un morphisme de schémas. Soit $\mathcal{F}$
un $\mathcal{O}_X$-module quasi-cohérent, soit $\mathcal{G}$
un $\mathcal{O}_Y$-module quasi-cohérent et soit
$\varphi : \mathcal{G} \to f^*\mathcal{F}$ un homomorphisme de modules. Supposons que
\begin{enumerate}
\item $\varphi$ soit injectif,
\item $f$ soit quasi-compact, quasi-séparé, plat et surjectif,
\item $X$ et $Y$ soient localement noethériens, et
\item $\mathcal{G}$ soit un $\mathcal{O}_Y$-module cohérent.
\end{enumerate}
Alors $\mathcal{F} \cap f_*\mathcal{G}$, défini comme le produit fibré
$$
\xymatrix{
\mathcal{F} \ar[r] & f_*f^*\mathcal{F} \\
\mathcal{F} \cap f_*\mathcal{G} \ar[u] \ar[r] &
f_*\mathcal{G} \ar[u]
}
$$
est un $\mathcal{O}_X$-module cohérent.
\end{lemma}

\begin{proof}
Nous utiliserons librement la caractérisation des modules cohérents de
Cohomologie des schémas, lemme \ref{coherent-lemma-coherent-Noetherian},
ainsi que le fait que les modules cohérents forment une sous-catégorie de Serre
de $\QCoh(\mathcal{O}_X)$, voir
Cohomologie des schémas,
lemme \ref{coherent-lemma-coherent-Noetherian-quasi-coherent-sub-quotient}.
Si $f$ possède une section $\sigma$, on voit que
$\mathcal{F} \cap f_*\mathcal{G}$ est contenu dans l'image de
$\sigma^*\mathcal{G} \to \sigma^*f^*\mathcal{F} = \mathcal{F}$,
donc est cohérent. En général, pour montrer que $\mathcal{F} \cap f_*\mathcal{G}$
est cohérent, il suffit de montrer que
$f^*(\mathcal{F} \cap f_*\mathcal{G})$ est cohérent (voir
Descente, lemme \ref{descent-lemma-finite-type-descends}).
Puisque $f$ est plat, ce module est égal à $f^*\mathcal{F} \cap f^*f_*\mathcal{G}$.
Puisque $f$ est plat, quasi-compact et quasi-séparé, on a
$f^*f_*\mathcal{G} = p_*q^*\mathcal{G}$, où $p, q : Y \times_X Y \to Y$
sont les projections, voir
Cohomologie des schémas, lemme \ref{coherent-lemma-flat-base-change-cohomology}.
Puisque $p$ possède une section, on conclut.
\end{proof}

\noindent
Soit $S$ un schéma. Soit $(U, R, s, t, c)$ un schéma en groupoïdes sur $S$.
Supposons que $U$ soit localement noethérien. Dans le lemme ci-dessous, un
faisceau quasi-cohérent $(\mathcal{F}, \alpha)$ sur $(U, R, s, t, c)$ est dit
{\it cohérent} si $\mathcal{F}$ est un $\mathcal{O}_U$-module cohérent.

\begin{lemma}
\label{groupoids-lemma-colimit-coherent}
Soit $(U, R, s, t, c)$ un schéma en groupoïdes sur $S$.
Supposons que
\begin{enumerate}
\item $U$, $R$ soient noethériens,
\item $s, t$ soient plats, quasi-compacts et quasi-séparés.
\end{enumerate}
Alors tout module quasi-cohérent $(\mathcal{F}, \beta)$ sur $(U, R, s, t, c)$
est une colimite filtrante de modules cohérents.
\end{lemma}

\begin{proof}
Nous utiliserons sans plus ample mention la caractérisation donnée dans
Cohomologie des schémas, lemme \ref{coherent-lemma-coherent-Noetherian}, des
modules cohérents sur un schéma localement noethérien. On peut écrire
$\mathcal{F} = \colim \mathcal{H}_i$ comme colimite filtrante
de sous-modules cohérents $\mathcal{H}_i \subset \mathcal{F}$, voir
Cohomologie des schémas, lemme \ref{coherent-lemma-directed-colimit-coherent}.
Étant donné un faisceau quasi-cohérent $\mathcal{H}$ sur $U$, on note
$(s_*t^*\mathcal{H}, \alpha)$ le faisceau quasi-cohérent sur $(U, R, s, t, c)$ du
lemme \ref{groupoids-lemma-construct-quasi-coherent}. Considérons le morphisme d'adjonction
$(\mathcal{F}, \beta) \to (s_*t^*\mathcal{F}, \alpha)$ dans
$\QCoh(U, R, s, t, c)$, voir remarque \ref{groupoids-remark-adjunction-map}.
Posons
$$
(\mathcal{F}_i, \beta_i) =
(\mathcal{F}, \beta)
\times_{(s_*t^*\mathcal{F}, \alpha)}
(s_*t^*\mathcal{H}_i, \alpha)
$$
{\emergencystretch=1em dans $\QCoh(U,\allowbreak R,\allowbreak s,\allowbreak t,\allowbreak c)$. Puisque la restriction à $U$ est un foncteur exact
sur $\QCoh(U,\allowbreak R,\allowbreak s,\allowbreak t,\allowbreak c)$ d'après la démonstration du lemme \ref{groupoids-lemma-abelian},
on obtient un diagramme cartésien\par}
$$
\xymatrix{
\mathcal{F} \ar[r] & s_*t^*\mathcal{F} \\
\mathcal{F}_i \ar[r] \ar[u] & s_*t^*\mathcal{H}_i \ar[u]
}
$$
autrement dit $\mathcal{F}_i = \mathcal{F} \cap s_*t^*\mathcal{H}_i$.
D'après la description du morphisme d'adjonction dans la remarque \ref{groupoids-remark-adjunction-map},
ce diagramme est isomorphe au diagramme
$$
\xymatrix{
\mathcal{F} \ar[r] & s_*s^*\mathcal{F} \\
\mathcal{F}_i \ar[r] \ar[u] & s_*t^*\mathcal{H}_i \ar[u]
}
$$
où la flèche verticale de droite s'obtient en appliquant $s_*$ au morphisme
$$
t^*\mathcal{H}_i \to t^*\mathcal{F} \xrightarrow{\beta} s^*\mathcal{F}
$$
Cette flèche est injective puisque $t$ est un morphisme plat. Il s'ensuit que
$\mathcal{F}_i$ est cohérent par le lemme \ref{groupoids-lemma-push-pull}.
Enfin, puisque $s$ est quasi-compact et quasi-séparé, on voit que
$s_*$ commute aux colimites (voir
Cohomologie des schémas, lemme \ref{coherent-lemma-colimit-cohomology}).
Ainsi $s_*t^*\mathcal{F} = \colim s_*t^*\mathcal{H}_i$ et donc
$(\mathcal{F}, \beta) = \colim (\mathcal{F}_i, \beta_i)$, comme voulu.
\end{proof}

\noindent
Voici un lemme curieux, utile lorsqu'on travaille avec des groupoïdes
sur des corps. C'est en fait l'argument usuel pour démontrer que toute
représentation d'un groupe algébrique est une colimite de représentations
de dimension finie.

\begin{lemma}
\label{groupoids-lemma-colimit-finite-type}
Soit $(U, R, s, t, c)$ un schéma en groupoïdes sur $S$.
Supposons que
\begin{enumerate}
\item $U$, $R$ soient affines,
\item il existe des $e_i \in \mathcal{O}_R(R)$ tels que
tout élément $g \in \mathcal{O}_R(R)$ s'écrive de manière unique sous la forme
$\sum s^*(f_i)e_i$ avec des $f_i \in \mathcal{O}_U(U)$.
\end{enumerate}
Alors tout module quasi-cohérent $(\mathcal{F}, \alpha)$ sur $(U, R, s, t, c)$
est une colimite filtrante de modules quasi-cohérents de type fini.
\end{lemma}

\begin{proof}
{\emergencystretch=1em L'hypothèse signifie que $\mathcal{O}_R(R)$ est un
$\mathcal{O}_U(U)$-module libre via $s$, de base $e_i$. Ainsi,
pour tout $\mathcal{O}_U$-module quasi-cohérent $\mathcal{G}$,
on a $s^*\mathcal{G}(R) = \bigoplus_i \mathcal{G}(U)e_i$.
On écrira $s(-)$ pour désigner l'image réciproque des sections par $s$, et
de même pour les autres morphismes.
Soit $(\mathcal{F}, \alpha)$ un module quasi-cohérent sur
$(U, R, s, t, c)$. Soit $\sigma \in \mathcal{F}(U)$. D'après ce qui précède,
on peut écrire\par}
$$
\alpha(t(\sigma)) = \sum s(\sigma_i) e_i
$$
pour des $\sigma_i \in \mathcal{F}(U)$ uniques (tous sauf un nombre fini
sont bien entendu nuls). On peut aussi écrire
$$
c(e_i) = \sum \text{pr}_1(f_{ij}) \text{pr}_0(e_j)
$$
en tant que fonctions sur $R \times_{s, U, t}R$. Alors la commutativité du diagramme
de la définition \ref{groupoids-definition-groupoid-module} signifie que
$$
\sum \text{pr}_1(\alpha(t(\sigma_i))) \text{pr}_0(e_i)
=
\sum \text{pr}_1(s(\sigma_i)f_{ij}) \text{pr}_0(e_j)
$$
(calcul omis). En identifiant les coefficients de $\text{pr}_0(e_l)$,
on voit que $\alpha(t(\sigma_l)) = \sum s(\sigma_i)f_{il}$. Ainsi,
le sous-module $\mathcal{G} \subset \mathcal{F}$ engendré par les
éléments $\sigma_i$ définit un module quasi-cohérent de type fini
préservé par $\alpha$. C'est donc un sous-objet de $\mathcal{F}$ dans
$\QCoh(U, R, s, t, c)$. Ce sous-module contient $\sigma$
(comme on le voit en prenant l'image réciproque de la première relation par $e$). D'où le résultat.
\end{proof}

\noindent
Nous suggérons au lecteur de sauter le reste de cette section. Soit $S$ un schéma.
Soit $(U, R, s, t, c)$ un schéma en groupoïdes sur $S$. Soit $\kappa$ un
cardinal. Dans la suite, on dira qu'un faisceau quasi-cohérent
$(\mathcal{F}, \alpha)$ sur $(U, R, s, t, c)$ est $\kappa$-engendré si
$\mathcal{F}$ est $\kappa$-engendré comme $\mathcal{O}_U$-module, voir
Propriétés, définition \ref{properties-definition-kappa-generated}.

\begin{lemma}
\label{groupoids-lemma-set-of-iso-classes}
Soit $(U, R, s, t, c)$ un schéma en groupoïdes sur $S$.
Soit $\kappa$ un cardinal.
Il existe un ensemble $T$ et une famille $(\mathcal{F}_t, \alpha_t)_{t \in T}$ de
modules quasi-cohérents $\kappa$-engendrés sur $(U, R, s, t, c)$
telle que tout module quasi-cohérent $\kappa$-engendré sur
$(U, R, s, t, c)$ soit isomorphe à l'un des $(\mathcal{F}_t, \alpha_t)$.
\end{lemma}

\begin{proof}
Pour tout module quasi-cohérent $\mathcal{F}$ sur $U$, il existe un
ensemble (éventuellement vide) de morphismes $\alpha : t^*\mathcal{F} \to s^*\mathcal{F}$
tels que $(\mathcal{F}, \alpha)$ soit un module quasi-cohérent sur
$(U, R, s, t, c)$. D'après
Propriétés, lemme \ref{properties-lemma-set-of-iso-classes},
il existe un ensemble de classes d'isomorphisme d'objets $\kappa$-engendrés qui sont des
$\mathcal{O}_U$-modules quasi-cohérents.
\end{proof}

\begin{lemma}
\label{groupoids-lemma-colimit-kappa}
Soit $(U, R, s, t, c)$ un schéma en groupoïdes sur $S$.
Supposons que $s, t$ soient plats. Il existe un
cardinal $\kappa$ tel que tout module quasi-cohérent
$(\mathcal{F}, \alpha)$ sur $(U, R, s, t, c)$
soit la colimite filtrante de ses sous-modules $\kappa$-engendrés
quasi-cohérents.
\end{lemma}

\begin{proof}
Dans l'énoncé du lemme et dans cette démonstration,
un {\it sous-module} d'un module quasi-cohérent $(\mathcal{F}, \alpha)$
est un sous-module quasi-cohérent $\mathcal{G} \subset \mathcal{F}$
tel que $\alpha(t^*\mathcal{G}) = s^*\mathcal{G}$ comme sous-faisceaux de
$s^*\mathcal{F}$. Cela a un sens car, $s, t$ étant plats, les
images réciproques $s^*$ et $t^*$ sont exactes, c'est-à-dire préservent les sous-faisceaux.
La démonstration reprendra celle de
Propriétés, lemme \ref{properties-lemma-colimit-kappa}.
Nous invitons vivement le lecteur à lire d'abord cette démonstration.

\medskip\noindent
Choisissons un recouvrement ouvert affine $U = \bigcup_{i \in I} U_i$.
Pour chaque paire $i, j$, choisissons des recouvrements ouverts affines
$$
U_i \cap U_j = \bigcup\nolimits_{k \in I_{ij}} U_{ijk}
\quad\text{et}\quad
s^{-1}(U_i) \cap t^{-1}(U_j) = \bigcup\nolimits_{k \in J_{ij}} W_{ijk}.
$$
Écrivons $U_i = \Spec(A_i)$, $U_{ijk} = \Spec(A_{ijk})$,
$W_{ijk} = \Spec(B_{ijk})$.
Soit $\kappa$ un cardinal infini qui soit $\geq$ la cardinalité
de chacun des ensembles $I$, $I_{ij}$, $J_{ij}$.

\medskip\noindent
Soit $(\mathcal{F}, \alpha)$ un module quasi-cohérent sur $(U, R, s, t, c)$.
Posons $M_i = \mathcal{F}(U_i)$, $M_{ijk} = \mathcal{F}(U_{ijk})$.
Notons que
$$
M_i \otimes_{A_i} A_{ijk} = M_{ijk} = M_j \otimes_{A_j} A_{ijk}
$$
et que $\alpha$ donne des isomorphismes
$$
\alpha|_{W_{ijk}} :
M_i \otimes_{A_i, t} B_{ijk}
\longrightarrow
M_j \otimes_{A_j, s} B_{ijk}
$$
voir
Schémas, lemme \ref{schemes-lemma-widetilde-pullback}.
Au moyen de l'axiome du choix, choisissons une application
$$
(i, j, k, m) \mapsto S(i, j, k, m)
$$
qui associe à tous $i, j \in I$, $k \in I_{ij}$ ou $k \in J_{ij}$
et $m \in M_i$ une partie finie $S(i, j, k, m) \subset M_j$
telle que l'on ait
$$
m \otimes 1 = \sum\nolimits_{m' \in S(i, j, k, m)} m' \otimes a_{m'}
\quad\text{ou}\quad
\alpha(m \otimes 1) = \sum\nolimits_{m' \in S(i, j, k, m)} m' \otimes b_{m'}
$$
dans $M_{ijk}$ pour certains $a_{m'} \in A_{ijk}$ ou $b_{m'} \in B_{ijk}$.
Convenons en outre que $S(i, i, k, m) = \{m\}$ pour tous
$i, j = i, k, m$ lorsque $k \in I_{ij}$. Fixons une telle famille $S(i, j, k, m)$.

\medskip\noindent
Étant donnée une famille $\mathcal{S} = (S_i)_{i \in I}$ de parties
$S_i \subset M_i$ de cardinalité au plus $\kappa$, posons
$\mathcal{S}' = (S'_i)$, où
$$
S'_j = \bigcup\nolimits_{(i, j, k, m)\text{ tel que }m \in S_i}
S(i, j, k, m)
$$
Notons que $S_i \subset S'_i$. Notons que $S'_i$ est de cardinalité au plus
$\kappa$, car c'est une réunion, indexée par un ensemble de cardinalité au plus $\kappa$,
de parties finies. Posons $\mathcal{S}^{(0)} = \mathcal{S}$,
$\mathcal{S}^{(1)} = \mathcal{S}'$ et, par récurrence,
$\mathcal{S}^{(n + 1)} = (\mathcal{S}^{(n)})'$. Posons ensuite
$\mathcal{S}^{(\infty)} = \bigcup_{n \geq 0} \mathcal{S}^{(n)}$.
En écrivant $\mathcal{S}^{(\infty)} = (S^{(\infty)}_i)$, on voit que,
pour tout élément $m \in S^{(\infty)}_i$, l'image de $m$ dans
$M_{ijk}$ peut s'écrire comme une somme finie $\sum m' \otimes a_{m'}$
avec $m' \in S_j^{(\infty)}$. Ainsi, en posant
$$
N_i = A_i\text{-sous-module de }M_i\text{ engendré par }S^{(\infty)}_i
$$
on a
$$
N_i \otimes_{A_i} A_{ijk} = N_j \otimes_{A_j} A_{ijk}
\quad\text{et}\quad
\alpha(N_i \otimes_{A_i, t} B_{ijk}) = N_j \otimes_{A_j, s} B_{ijk}
$$
comme sous-modules de $M_{ijk}$ ou de $M_j \otimes_{A_j, s} B_{ijk}$.
Il existe donc un sous-module quasi-cohérent
$\mathcal{G} \subset \mathcal{F}$ avec $\mathcal{G}(U_i) = N_i$
tel que $\alpha(t^*\mathcal{G}) = s^*\mathcal{G}$ comme sous-modules
de $s^*\mathcal{F}$. Autrement dit,
$(\mathcal{G}, \alpha|_{t^*\mathcal{G}})$ est un sous-module de
$(\mathcal{F}, \alpha)$.
De plus, par construction, $\mathcal{G}$ est $\kappa$-engendré.

\medskip\noindent
{\emergencystretch=1em Soit $\{(\mathcal{G}_t, \alpha_t)\}_{t \in T}$ l'ensemble des
sous-modules quasi-cohérents $\kappa$-engendrés de $(\mathcal{F}, \alpha)$.
Si $t, t' \in T$, alors $\mathcal{G}_t + \mathcal{G}_{t'}$ est encore un
sous-module quasi-cohérent $\kappa$-engendré, car il est l'image du morphisme
$\mathcal{G}_t \oplus \mathcal{G}_{t'} \to \mathcal{F}$.
Le système, ordonné par inclusion, est donc filtrant.
Les arguments précédents montrent que toute section de $\mathcal{F}$ sur $U_i$
appartient à l'un des $\mathcal{G}_t$ (car on peut partir de $\mathcal{S}$
tel que la section donnée soit un élément de $S_i$). Par conséquent,
$\colim_t \mathcal{G}_t \to \mathcal{F}$ est à la fois injectif et surjectif,
comme voulu.\par}
\end{proof}





\section{Groupoïdes et schémas en groupes}
\label{groupoids-section-groupoids-group-schemes}

\noindent
Il existe de nombreuses façons de construire un groupoïde à partir d'une action $a$
d'un groupe $G$ sur un ensemble $V$. Nous choisissons celle où l'on considère
un élément $g \in G$ comme une flèche de source $v$ et de but $a(g, v)$.
Cela conduit à la construction suivante pour les actions de schémas en
groupes.

\begin{lemma}
\label{groupoids-lemma-groupoid-from-action}
Soit $S$ un schéma.
Soit $Y$ un schéma sur $S$.
Soit $(G, m)$ un schéma en groupes sur $Y$, de
section unité $e_G$ et de morphisme d'inversion $i_G$.
Soit $X/Y$ un schéma sur $Y$ et soit $a : G \times_Y X \to X$
une action de $G$ sur $X/Y$.
On obtient alors un schéma en groupoïdes $(U, R, s, t, c, e, i)$ sur $S$
de la manière suivante :
\begin{enumerate}
\item On pose $U = X$, et $R = G \times_Y X$.
\item On prend $s : R \to U$ égal à $(g, x) \mapsto x$.
\item On prend $t : R \to U$ égal à $(g, x) \mapsto a(g, x)$.
\item On prend $c : R \times_{s, U, t} R \to R$ égal à
$((g, x), (g', x')) \mapsto (m(g, g'), x')$.
\item On prend $e : U \to R$ égal à $x \mapsto (e_G(x), x)$.
\item On prend $i : R \to R$ égal à $(g, x) \mapsto (i_G(g), a(g, x))$.
\end{enumerate}
\end{lemma}

\begin{proof}
Omis. Indication : il suffit de vérifier que cela fonctionne au niveau
des ensembles. Pour cela, utiliser la description précédant le lemme, qui voit
$g$ comme une flèche de $v$ vers $a(g, v)$.
\end{proof}

\begin{lemma}
\label{groupoids-lemma-action-groupoid-modules}
Soit $S$ un schéma.
Soit $Y$ un schéma sur $S$.
Soit $(G, m)$ un schéma en groupes sur $Y$.
Soit $X$ un schéma sur $Y$ et soit $a : G \times_Y X \to X$
une action de $G$ sur $X$ au-dessus de $Y$. Soit $(U, R, s, t, c)$
le schéma en groupoïdes construit dans le lemme \ref{groupoids-lemma-groupoid-from-action}.
La règle
$(\mathcal{F}, \alpha) \mapsto (\mathcal{F}, \alpha)$ définit
une équivalence de catégories entre les objets $G$-équivariants qui sont des
$\mathcal{O}_X$-modules et la catégorie des modules quasi-cohérents
sur $(U, R, s, t, c)$.
\end{lemma}

\begin{proof}
L'assertion a un sens puisque $t = a$ et $s = \text{pr}_1$
comme morphismes $R = G \times_Y X \to X$, voir les
définitions \ref{groupoids-definition-equivariant-module} et
\ref{groupoids-definition-groupoid-module}.
En utilisant la traduction donnée dans le lemme \ref{groupoids-lemma-groupoid-from-action},
les conditions de commutativité
des deux définitions coïncident exactement.
\end{proof}





\section{Le schéma en groupes stabilisateur}
\label{groupoids-section-stabilizer}

\noindent
À partir d'un schéma en groupoïdes, on obtient un schéma en groupes comme suit.

\begin{lemma}
\label{groupoids-lemma-groupoid-stabilizer}
Soit $S$ un schéma.
Soit $(U, R, s, t, c)$ un groupoïde sur $S$.
Le schéma $G$ défini par le carré cartésien
$$
\xymatrix{
G \ar[r] \ar[d] & R \ar[d]^{j = (t, s)} \\
U \ar[r]^-{\Delta} & U \times_S U
}
$$
est un schéma en groupes sur $U$, dont la loi de composition
$m$ est induite par la loi de composition $c$.
\end{lemma}

\begin{proof}
Cela résulte du fait que, dans une catégorie groupoïde, l'ensemble
des endomorphismes de tout objet forme un groupe.
\end{proof}

\noindent
Puisque $\Delta$ est une immersion, on voit que $G = j^{-1}(\Delta_{U/S})$
est un sous-schéma localement fermé de $R$. De ce point de vue,
le morphisme structural $j^{-1}(\Delta_{U/S}) \to U$ est induit par
$s$ ou par $t$ (c'est le même), et $m$ est induit par $c$.

\begin{definition}
\label{groupoids-definition-stabilizer-groupoid}
Soit $S$ un schéma.
Soit $(U, R, s, t, c)$ un groupoïde sur $S$.
Le schéma en groupes $j^{-1}(\Delta_{U/S})\to U$
est appelé le {\it stabilisateur du schéma en groupoïdes
$(U, R, s, t, c)$}.
\end{definition}

\noindent
Dans la littérature, le schéma en groupes stabilisateur est souvent noté $S$
(sans doute parce que le mot stabilisateur commence par un ``s'');
nous ne pouvons le faire, puisque nous avons déjà employé $S$ pour le schéma de base.

\begin{lemma}
\label{groupoids-lemma-groupoid-action-stabilizer}
Soit $S$ un schéma.
Soit $(U, R, s, t, c)$ un groupoïde sur $S$, et soit $G/U$ son stabilisateur.
Notons $R_t/U$ le schéma $R$ considéré comme schéma sur $U$ via le
morphisme $t : R \to U$.
Il existe une action à gauche canonique
$$
a : G \times_U R_t \longrightarrow R_t
$$
induite par la loi de composition $c$.
\end{lemma}

\begin{proof}
En termes de points au-dessus de $T/S$, on définit $a(g, r) = c(g, r)$.
\end{proof}

\begin{lemma}
\label{groupoids-lemma-groupoid-action-stabilizer-pseudo-torsor}
Soit $S$ un schéma. Soit $(U, R, s, t, c)$ un schéma en groupoïdes
sur $S$. Soit $G$ le schéma en groupes stabilisateur de $R$.
Soit
$$
G_0 = G \times_{U, \text{pr}_0} (U \times_S U) = G \times_S U
$$
comme schéma en groupes sur $U \times_S U$. L'action de $G$ sur $R$ du
lemme \ref{groupoids-lemma-groupoid-action-stabilizer}
induit une action de $G_0$ sur $R$ au-dessus de $U \times_S U$
qui fait de $R$ un pseudo-torseur sous $G_0$ sur $U \times_S U$.
\end{lemma}

\begin{proof}
Cela résulte du fait que, dans une catégorie groupoïde $\mathcal{C}$, l'ensemble
$\Mor_\mathcal{C}(x, y)$ est un espace homogène principal
sous le groupe $\Mor_\mathcal{C}(y, y)$.
\end{proof}

\begin{lemma}
\label{groupoids-lemma-fibres-j}
Soit $S$ un schéma. Soit $(U, R, s, t, c)$ un schéma en groupoïdes sur $S$.
Soit $p \in U \times_S U$ un point. Notons
$R_p$ la fibre schématique de $j = (t, s) : R \to U \times_S U$.
Si $R_p \not = \emptyset$, alors l'action
$$
G_{0, \kappa(p)} \times_{\kappa(p)} R_p \longrightarrow R_p
$$
(voir
lemme \ref{groupoids-lemma-groupoid-action-stabilizer-pseudo-torsor})
fait de $R_p$ un torseur sous $G_{\kappa(p)}$ sur $\kappa(p)$.
\end{lemma}

\begin{proof}
On a un pseudo-torseur d'après le lemme cité dans l'énoncé.
Et si $R_p$ n'est pas le schéma vide, alors $\{R_p \to p\}$
est un recouvrement fpqc qui trivialise le pseudo-torseur.
\end{proof}







\section{Restriction des groupoïdes}
\label{groupoids-section-restrict-groupoid}

\noindent
Considérons un groupoïde (ordinaire)
$\mathcal{C} = (\text{Ob}, \text{Arrows}, s, t, c)$.
Supposons donnée une application d'ensembles $g : \text{Ob}' \to \text{Ob}$.
On peut alors construire un groupoïde
$\mathcal{C}' = (\text{Ob}', \text{Arrows}', s', t', c')$
en considérant un morphisme entre des éléments $x', y'$ de $\text{Ob}'$
comme un morphisme dans $\mathcal{C}$ entre $g(x'), g(y')$.
Autrement dit, on pose
$$
\text{Arrows}' =
\text{Ob}'
\times_{g, \text{Ob}, t}
\text{Arrows}
\times_{s, \text{Ob}, g}
\text{Ob}'.
$$
On fait les choix évidents pour $s'$, $t'$ et $c'$. Il existe un foncteur canonique
$\mathcal{C}' \to \mathcal{C}$ qui est pleinement fidèle,
mais pas nécessairement essentiellement surjectif. Ce groupoïde $\mathcal{C}'$,
muni du foncteur $\mathcal{C}' \to \mathcal{C}$,
est appelé la {\it restriction} du groupoïde
$\mathcal{C}$ à $\text{Ob}'$.

\begin{lemma}
\label{groupoids-lemma-restrict-groupoid}
Soit $S$ un schéma.
Soit $(U, R, s, t, c)$ un schéma en groupoïdes sur $S$.
Soit $g : U' \to U$ un morphisme de schémas.
Considérons le diagramme suivant
$$
\xymatrix{
R' \ar[d] \ar[r] \ar@/_3pc/[dd]_{t'} \ar@/^1pc/[rr]^{s'}&
R \times_{s, U} U' \ar[r] \ar[d] &
U' \ar[d]^g \\
U' \times_{U, t} R \ar[d] \ar[r] &
R \ar[r]^s \ar[d]_t &
U \\
U' \ar[r]^g &
U
}
$$
où tous les carrés sont cartésiens. Il existe alors une
loi de composition canonique $c' : R' \times_{s', U', t'} R' \to R'$
telle que $(U', R', s', t', c')$ soit un schéma en groupoïdes sur
$S$ et que $U' \to U$, $R' \to R$ définissent un morphisme
$(U', R', s', t', c') \to (U, R, s, t, c)$ de groupoïdes en schémas sur $S$.
De plus, pour tout schéma $T$ sur $S$, le foncteur de groupoïdes
$$
(U'(T), R'(T), s', t', c') \to (U(T), R(T), s, t, c)
$$
est la restriction (voir ci-dessus) de $(U(T), R(T), s, t, c)$ par l'application
$U'(T) \to U(T)$.
\end{lemma}

\begin{proof}
Omis.
\end{proof}

\begin{definition}
\label{groupoids-definition-restrict-groupoid}
Soit $S$ un schéma.
Soit $(U, R, s, t, c)$ un schéma en groupoïdes sur $S$.
Soit $g : U' \to U$ un morphisme de schémas.
Le morphisme de groupoïdes
$(U', R', s', t', c') \to (U, R, s, t, c)$
construit dans le lemme \ref{groupoids-lemma-restrict-groupoid} est appelé
la {\it restriction de $(U, R, s, t, c)$ à $U'$}.
On utilise parfois la notation $R' = R|_{U'}$ dans ce cas.
\end{definition}

\begin{lemma}
\label{groupoids-lemma-restrict-groupoid-relation}
Les notions de restriction des groupoïdes et des
relations de (pré-)équivalence définies dans les définitions
\ref{groupoids-definition-restrict-groupoid} et \ref{groupoids-definition-restrict-relation}
coïncident via les constructions des
lemmes \ref{groupoids-lemma-groupoid-pre-equivalence} et
\ref{groupoids-lemma-equivalence-groupoid}.
\end{lemma}

\begin{proof}
Ce que nous affirmons ici est que le $R'$ du
lemme \ref{groupoids-lemma-restrict-groupoid} est également
égal à
$$
R' = (U' \times_S U')\times_{U \times_S U} R
\longrightarrow
U' \times_S U'
$$
En fait, c'eût peut-être été une manière plus claire d'énoncer ce lemme.
\end{proof}

\begin{lemma}
\label{groupoids-lemma-restrict-stabilizer}
Soit $S$ un schéma.
Soit $(U, R, s, t, c)$ un schéma en groupoïdes sur $S$.
Soit $g : U' \to U$ un morphisme de schémas.
Soit $(U', R', s', t', c')$ la restriction de $(U, R, s, t, c)$ par $g$.
Soit $G$ le stabilisateur de $(U, R, s, t, c)$ et soit
$G'$ le stabilisateur de $(U', R', s', t', c')$.
Alors $G'$ est le changement de base de $G$ par $g$, c'est-à-dire
qu'il existe une identification canonique $G' = U' \times_{g, U} G$.
\end{lemma}

\begin{proof}
Omis.
\end{proof}






\section{Sous-schémas invariants}
\label{groupoids-section-invariant}

\noindent
Dans cette section, nous examinons brièvement la notion de sous-schéma invariant.

\begin{definition}
\label{groupoids-definition-invariant-open}
Soit $(U, R, s, t, c)$ un schéma en groupoïdes sur le schéma de base $S$.
\begin{enumerate}
\item Une partie $W \subset U$ est {\it $R$-invariante au sens ensembliste}
si $t(s^{-1}(W)) \subset W$.
\item Un ouvert $W \subset U$ est {\it $R$-invariant} si
$t(s^{-1}(W)) \subset W$.
\item Un sous-schéma fermé $Z \subset U$ est dit {\it $R$-invariant}
si $t^{-1}(Z) = s^{-1}(Z)$. On utilise ici l'image réciproque schématique, voir
Schémas, définition \ref{schemes-definition-inverse-image-closed-subscheme}.
\item Un monomorphisme de schémas $T \to U$ est {\it $R$-invariant} si
$T \times_{U, t} R = R \times_{s, U} T$ comme schémas sur $R$.
\end{enumerate}
\end{definition}

\noindent
Pour les parties et les sous-schémas ouverts $W \subset U$, l'$R$-invariance
équivaut aussi à demander que $s^{-1}(W) = t^{-1}(W)$
comme parties de $R$. Si $W \subset U$ est un sous-schéma ouvert $R$-équivariant,
alors la restriction de $R$ à $W$ est simplement $R_W = s^{-1}(W) = t^{-1}(W)$.
De même, si $Z \subset U$ est un sous-schéma fermé $R$-invariant, alors
la restriction de $R$ à $Z$ est simplement $R_Z = s^{-1}(Z) = t^{-1}(Z)$.

\begin{lemma}
\label{groupoids-lemma-constructing-invariant-opens}
Soit $S$ un schéma.
Soit $(U, R, s, t, c)$ un schéma en groupoïdes sur $S$.
\begin{enumerate}
\item Pour toute partie $W \subset U$, la partie $t(s^{-1}(W))$
est $R$-invariante au sens ensembliste.
\item Si $s$ et $t$ sont ouverts, alors, pour tout ouvert $W \subset U$,
l'ouvert $t(s^{-1}(W))$ est un sous-schéma ouvert $R$-invariant.
\item Si $s$ et $t$ sont ouverts et quasi-compacts, alors $U$ possède un
recouvrement ouvert formé de sous-schémas ouverts quasi-compacts $R$-invariants.
\end{enumerate}
\end{lemma}

\begin{proof}
{\emergencystretch=2em La partie (1) résulte des
Lemmes \ref{groupoids-lemma-pre-equivalence-equivalence-relation-points} et
\ref{groupoids-lemma-groupoid-pre-equivalence}, à savoir que $t(s^{-1}(W))$
est l'ensemble des points de $U$ équivalents à un point de $W$.
Supposons ensuite $s$ et $t$ ouverts, et $W \subset U$ ouvert.
Comme $s$ est ouvert, l'ensemble $W' = t(s^{-1}(W))$ est un ouvert de $U$.
Supposons enfin que $s$, $t$ soient tous deux ouverts et quasi-compacts.
Alors, si $W \subset U$ est un ouvert quasi-compact,
$W' = t(s^{-1}(W))$ est lui aussi un ouvert quasi-compact et il est invariant d'après la
discussion précédente. En faisant parcourir à $W$ tous les ouverts affines de $U$,
on obtient (3).\par}
\end{proof}

\begin{lemma}
\label{groupoids-lemma-first-observation}
Soit $S$ un schéma. Soit $(U, R, s, t, c)$ un schéma en groupoïdes sur $S$.
Supposons $s$ et $t$ quasi-compacts et plats, et $U$ quasi-séparé.
Soit $W \subset U$ un ouvert quasi-compact. Alors $t(s^{-1}(W))$
est l'intersection d'une famille non vide d'ouverts quasi-compacts de $U$.
\end{lemma}

\begin{proof}
Notons que $s^{-1}(W)$ est un ouvert quasi-compact de $R$.
En tant qu'application continue, $t$ envoie le sous-ensemble quasi-compact
$s^{-1}(W)$ sur un sous-ensemble quasi-compact $t(s^{-1}(W))$.
Comme $t$ est plat et que $s^{-1}(W)$ est stable par générisation,
il en va de même de $t(s^{-1}(W))$, voir
(Morphismes, Lemme \ref{morphisms-lemma-generalizations-lift-flat} et
Topologie, Lemme \ref{topology-lemma-lift-specializations-images}).
Choisissons un ouvert quasi-compact $W' \subset U$ contenant $t(s^{-1}(W))$. D'après
Propriétés, Lemme \ref{properties-lemma-quasi-compact-quasi-separated-spectral},
on voit que $W'$ est un espace spectral (c'est ici que l'on utilise que $U$ est quasi-séparé).
Le lemme résulte alors de
Topologie, Lemme \ref{topology-lemma-make-spectral-space},
appliqué à $t(s^{-1}(W)) \subset W'$.
\end{proof}

\begin{lemma}
\label{groupoids-lemma-second-observation}
Hypothèses et notations comme dans le Lemme \ref{groupoids-lemma-first-observation}.
Il existe un ouvert $R$-invariant $V \subset U$ et un ouvert quasi-compact
$W'$ tels que $W \subset V \subset W' \subset U$.
\end{lemma}

\begin{proof}
Posons $E = t(s^{-1}(W))$. Rappelons que $E$ est $R$-invariant au sens ensembliste
(Lemme \ref{groupoids-lemma-constructing-invariant-opens}).
D'après le Lemme \ref{groupoids-lemma-first-observation}, il existe un ouvert quasi-compact
$W'$ contenant $E$. Posons $Z = U \setminus W'$ et considérons
$T = t(s^{-1}(Z))$. Observons que $Z \subset T$ et que
$E \cap T = \emptyset$, car $s^{-1}(E) = t^{-1}(E)$ est disjoint
de $s^{-1}(Z)$. Puisque $T$ est l'image du sous-ensemble fermé
$s^{-1}(Z) \subset R$ par le morphisme quasi-compact $t : R \to U$,
on voit que tout point $\xi$ de l'adhérence $\overline{T}$
est une spécialisation d'un point de $T$, voir
Morphismes, Lemme \ref{morphisms-lemma-reach-points-scheme-theoretic-image} (et
Morphismes, Lemme \ref{morphisms-lemma-quasi-compact-scheme-theoretic-image}
pour voir que l'image schématique est l'adhérence de l'image).
Écrivons $\xi' \leadsto \xi$ avec $\xi' \in T$. Supposons que $r \in R$ et
$s(r) = \xi$. Comme $s$ est plat, on peut trouver une spécialisation $r' \leadsto r$
dans $R$ telle que $s(r') = \xi'$
(Morphismes, Lemme \ref{morphisms-lemma-generalizations-lift-flat}).
Alors $t(r') \leadsto t(r)$. On conclut que $t(r') \in T$, puisque $T$
est invariant au sens ensembliste d'après
le Lemme \ref{groupoids-lemma-constructing-invariant-opens}.
Ainsi, $\overline{T}$ est un sous-ensemble fermé $R$-invariant au sens ensembliste
et $V = U \setminus \overline{T}$ est l'ouvert que nous
cherchions. Il est contenu dans $W'$, ce qui achève la démonstration.
\end{proof}




\section{Faisceaux quotients}
\label{groupoids-section-quotient-sheaves}

\noindent
Soit $\tau \in\allowbreak \{Zariski,\allowbreak \etale,\allowbreak fppf,\allowbreak smooth,\allowbreak syntomic\}$.
Soit $S$ un schéma.
Soit $j : R \to U \times_S U$ une prérelation sur $S$.
Supposons que $U, R, S$ soient des objets d'un $\tau$-site $\Sch_\tau$
(voir Topologies, section \ref{topologies-section-procedure}).
On peut alors considérer les foncteurs
$$
h_U, h_R :
(\Sch/S)_\tau^{opp}
\longrightarrow
\textit{Ensembles}.
$$
Ce sont des faisceaux, voir
Descente, Lemme \ref{descent-lemma-fpqc-universal-effective-epimorphisms}.
Le morphisme $j$ induit une application $j : h_R \to h_U \times h_U$.
Pour tout objet $T \in \Ob((\Sch/S)_\tau)$,
on peut prendre la relation d'équivalence $\sim_T$ engendrée par
$j(T) : R(T) \to U(T) \times U(T)$ et considérer le quotient.
On obtient ainsi un préfaisceau
\begin{equation}
\label{groupoids-equation-quotient-presheaf}
(\Sch/S)_\tau^{opp}
\longrightarrow
\textit{Ensembles}, \quad
T \longmapsto U(T)/\sim_T
\end{equation}

\begin{definition}
\label{groupoids-definition-quotient-sheaf}
Soient $\tau$, $S$ et la prérelation $j : R \to U \times_S U$ comme ci-dessus.
Dans cette situation, le {\it faisceau quotient $U/R$} associé
à $j$ est le faisceau associé au préfaisceau
(\ref{groupoids-equation-quotient-presheaf}) pour la topologie $\tau$.
Si $j : R \to U \times_S U$ provient de l'action d'un schéma en groupes
$G/S$ sur $U$ comme dans le Lemme \ref{groupoids-lemma-groupoid-from-action}, on
note parfois le faisceau quotient $U/G$.
\end{definition}

\noindent
Cela signifie exactement que le diagramme
$$
\xymatrix{
h_R \ar@<1ex>[r] \ar@<-1ex>[r] &
h_U \ar[r] &
U/R
}
$$
est un diagramme coégalisateur dans la catégorie des faisceaux d'ensembles
sur $(\Sch/S)_\tau$. Au moyen du plongement de Yoneda, on
peut considérer $(\Sch/S)_\tau$ comme une sous-catégorie pleine des
faisceaux sur $(\Sch/S)_\tau$ et ainsi identifier les schémas
aux foncteurs représentables. Par cet abus de notation,
nous représenterons souvent le diagramme ci-dessus simplement par
$$
\xymatrix{
R \ar@<1ex>[r]^s \ar@<-1ex>[r]_t &
U \ar[r] &
U/R
}
$$
Nous travaillerons le plus souvent avec la topologie fppf lorsque nous considérerons
les faisceaux quotients de groupoïdes ou de relations d'équivalence.

\begin{definition}
\label{groupoids-definition-representable-quotient}
Dans la situation de la Définition \ref{groupoids-definition-quotient-sheaf}.
On dit que la prérelation $j$ admet un
{\it quotient représentable} si le faisceau $U/R$ est représentable.
On dira qu'un groupoïde $(U, R, s, t, c)$ admet un
{\it quotient représentable}
si le quotient $U/R$ correspondant à $j = (t, s)$ est représentable.
\end{definition}

\noindent
Le lemme suivant caractérise les schémas $M$ qui représentent le quotient.
Il s'applique par exemple si $\tau = fppf$, si $U \to M$ est plat,
de présentation finie et surjectif, et si $R \cong U \times_M U$.

\begin{lemma}
\label{groupoids-lemma-criterion-quotient-representable}
Dans la situation de la Définition \ref{groupoids-definition-quotient-sheaf}.
Supposons qu'il existe un schéma $M$ et un morphisme $U \to M$ tels que
\begin{enumerate}
\item le morphisme $U \to M$ égalise $s, t$,
\item le morphisme $U \to M$ induit une surjection de faisceaux
$h_U \to h_M$ pour la topologie $\tau$, et
\item l'application induite $(t, s) : R \to U \times_M U$ induit une
surjection de faisceaux $h_R \to h_{U \times_M U}$ pour la topologie $\tau$.
\end{enumerate}
Alors $M$ représente le faisceau quotient $U/R$.
\end{lemma}

\begin{proof}
La condition (1) dit que $h_U \to h_M$ se factorise par $U/R$.
La condition (2) dit que $U/R \to h_M$ est surjectif comme morphisme de faisceaux.
La condition (3) dit que $U/R \to h_M$ est injectif comme morphisme de faisceaux.
Le lemme en résulte.
\end{proof}

\noindent
Le lemme suivant est faux si l'on n'exige pas que $j$ soit une
prérelation d'équivalence (mais seulement, disons, une prérelation).

\begin{lemma}\emergencystretch=1em
\label{groupoids-lemma-quotient-pre-equivalence}
Soit $\tau \in\allowbreak \{Zariski,\allowbreak \etale,\allowbreak fppf,\allowbreak smooth,\allowbreak syntomic\}$.
Soit $S$ un schéma.
Soit $j : R \to U \times_S U$ une prérelation d'équivalence sur $S$.
Supposons que $U, R, S$ soient des objets d'un $\tau$-site $\Sch_\tau$.
Pour $T \in \Ob((\Sch/S)_\tau)$ et
$a, b \in U(T)$, les conditions suivantes sont équivalentes :
\begin{enumerate}
\item $a$ et $b$ ont la même image dans $(U/R)(T)$, et
\item il existe un $\tau$-recouvrement $\{f_i : T_i \to T\}$ de $T$
et des morphismes $r_i : T_i \to R$ tels que
$a \circ f_i = s \circ r_i$ et $b \circ f_i = t \circ r_i$.
\end{enumerate}
Autrement dit, dans ce cas, le morphisme de $\tau$-faisceaux
$$
h_R \longrightarrow h_U \times_{U/R} h_U
$$
est surjectif.
\end{lemma}

\begin{proof}
Omis. Indication : cela fonctionne parce que, dans ce cas, le préfaisceau
(\ref{groupoids-equation-quotient-presheaf}) est en réalité donné
par $T \mapsto U(T)/j(R(T))$, puisque $j(R(T)) \subset U(T) \times U(T)$
est une relation d'équivalence, voir
Définition \ref{groupoids-definition-equivalence-relation}.
\end{proof}

\begin{lemma}\emergencystretch=1em
\label{groupoids-lemma-quotient-pre-equivalence-relation-restrict}
Soit $\tau \in\allowbreak \{Zariski,\allowbreak \etale,\allowbreak fppf,\allowbreak smooth,\allowbreak syntomic\}$.
Soit $S$ un schéma.
Soit $j : R \to U \times_S U$ une prérelation d'équivalence sur $S$
et $g : U' \to U$ un morphisme de schémas sur $S$.
Soit $j' : R' \to U' \times_S U'$ la restriction de $j$ à $U'$.
Supposons que $U, U', R, S$ soient des objets d'un $\tau$-site $\Sch_\tau$.
Le morphisme de faisceaux quotients
$$
U'/R' \longrightarrow U/R
$$
est injectif. Si $g$ définit une surjection $h_{U'} \to h_U$ de faisceaux
pour la topologie $\tau$ (par exemple, si $\{g : U' \to U\}$ est un
$\tau$-recouvrement), alors $U'/R' \to U/R$ est un isomorphisme.
\end{lemma}

\begin{proof}
Supposons que $\xi, \xi' \in (U'/R')(T)$ soient des sections dont les
images dans $U/R$ sont égales.
On peut alors trouver un $\tau$-recouvrement $\mathcal{T} = \{T_i \to T\}$ de $T$
tel que $\xi|_{T_i}, \xi'|_{T_i}$ soient données par $a_i, a_i' \in U'(T_i)$. D'après le
Lemme \ref{groupoids-lemma-quotient-pre-equivalence}
et les axiomes d'un site, on peut, après avoir raffiné
$\mathcal{T}$, supposer qu'il existe des morphismes $r_i : T_i \to R$
tels que $g \circ a_i = s \circ r_i$, $g \circ a_i' = t \circ r_i$.
Puisque, par construction,
$R' = R \times_{U \times_S U} (U' \times_S U')$,
on voit que $(r_i, (a_i, a_i')) \in R'(T_i)$, ce qui
montre que $a_i$ et $a_i'$ définissent la même section
de $U'/R'$ au-dessus de $T_i$. La condition de faisceau implique alors
$\xi = \xi'$.

\medskip\noindent
Si $h_{U'} \to h_U$ est une surjection
de faisceaux, alors, bien sûr, $U'/R' \to U/R$ est également surjectif.
Si $\{g : U' \to U\}$ est un $\tau$-recouvrement, alors
le morphisme de faisceaux $h_{U'} \to h_U$ est surjectif, voir
Sites, Lemme \ref{sites-lemma-covering-surjective-after-sheafification}.
Ainsi, $U'/R' \to U/R$ est également surjectif dans ce cas.
\end{proof}

\begin{lemma}\emergencystretch=1em
\label{groupoids-lemma-quotient-groupoid-restrict}
Soit $\tau \in\allowbreak \{Zariski,\allowbreak \etale,\allowbreak fppf,\allowbreak smooth,\allowbreak syntomic\}$.
Soit $S$ un schéma.
Soit $(U, R, s, t, c)$ un schéma en groupoïdes sur $S$.
Soit $g : U' \to U$ un morphisme de schémas sur $S$.
Soit $(U', R', s', t', c')$ la restriction de $(U, R, s, t, c)$ à $U'$.
Supposons que $U, U', R, S$ soient des objets d'un $\tau$-site $\Sch_\tau$.
Le morphisme de faisceaux quotients
$$
U'/R' \longrightarrow U/R
$$
est injectif. Si la composée
$$
\xymatrix{
U' \times_{g, U, t} R \ar[r]_-{\text{pr}_1} \ar@/^3ex/[rr]^h
& R \ar[r]_s & U
}
$$
définit une surjection de faisceaux pour la topologie $\tau$, alors
le morphisme est bijectif. Cela vaut par exemple si
$\{h : U' \times_{g, U, t} R \to U\}$ est un $\tau$-recouvrement, ou
si $U' \to U$ définit une surjection de faisceaux pour la topologie $\tau$, ou si
$\{g : U' \to U\}$ est un recouvrement pour la topologie $\tau$.
\end{lemma}

\begin{proof}
L'injectivité résulte de la combinaison des
Lemmes \ref{groupoids-lemma-groupoid-pre-equivalence} et
\ref{groupoids-lemma-quotient-pre-equivalence-relation-restrict}.
Pour établir la surjectivité (voir
Sites, section \ref{sites-section-sheaves-injective},
pour une caractérisation des morphismes surjectifs de faisceaux), raisonnons comme suit.
Supposons que $T$ soit un schéma et que $\sigma \in U/R(T)$.
Il existe un recouvrement $\{T_i \to T\}$ tel que $\sigma|_{T_i}$
soit l'image d'un élément $f_i \in U(T_i)$. On peut donc
supposer que $\sigma$ est l'image de $f \in U(T)$.
Comme, par hypothèse, $h$ est une surjection de faisceaux, on
peut trouver un $\tau$-recouvrement $\{\varphi_i : T_i \to T\}$ et des morphismes
$f_i : T_i \to U' \times_{g, U, t} R$ tels que
$f \circ \varphi_i = h \circ f_i$. Notons
$f'_i = \text{pr}_0 \circ f_i : T_i \to U'$. Alors on voit que
$f'_i \in U'(T_i)$ s'envoie sur $g \circ f'_i \in U(T_i)$ et
que $g \circ f'_i \sim_{T_i} h \circ f_i = f \circ \varphi_i$,
avec les notations de (\ref{groupoids-equation-quotient-presheaf}). Plus précisément, un
élément de $R(T_i)$ donnant la relation est $\text{pr}_1 \circ f_i$.
Cela signifie que la restriction
de $\sigma$ à $T_i$ appartient à l'image de $U'/R'(T_i) \to U/R(T_i)$,
comme voulu.

\medskip\noindent
Si $\{h\}$ est un $\tau$-recouvrement, il induit une surjection de faisceaux, voir
Sites, Lemme \ref{sites-lemma-covering-surjective-after-sheafification}.
Si $U' \to U$ est surjectif, alors $h$ l'est aussi, car $s$ possède une section
(à savoir l'élément neutre $e$ du schéma en groupoïdes).
\end{proof}

\begin{lemma}\emergencystretch=1em
\label{groupoids-lemma-criterion-fibre-product}
Soit $S$ un schéma. Soit $f : (U, R, j) \to (U', R', j')$ un morphisme
entre des relations d'équivalence sur $S$. Supposons que
$$
\xymatrix{
R \ar[d]_s \ar[r]_f & R' \ar[d]^{s'} \\
U \ar[r]^f & U'
}
$$
soit cartésien. Pour tout
$\tau \in\allowbreak \{Zariski,\allowbreak \etale,\allowbreak fppf,\allowbreak smooth,\allowbreak syntomic\}$,
le diagramme
$$
\xymatrix{
U \ar[d] \ar[r] & U/R \ar[d]^f \\
U' \ar[r] & U'/R'
}
$$
est un carré cartésien de $\tau$-faisceaux.
\end{lemma}

\begin{proof}
D'après le Lemme \ref{groupoids-lemma-quotient-pre-equivalence}, les faisceaux quotients
admettent une description simple, que nous utiliserons ci-dessous sans autre mention.
Montrons d'abord que
$$
U \longrightarrow U' \times_{U'/R'} U/R
$$
est injectif. En effet, supposons que $a, b \in U(T)$ aient la même image
dans le membre de droite. Alors $f(a) = f(b)$. Après avoir remplacé $T$ par les
membres d'un $\tau$-recouvrement, on peut supposer qu'il existe
$r \in R(T)$ tel que $a = s(r)$ et $b = t(r)$. Alors $r' = f(r)$
est un $T$-point de $R'$ tel que $s'(r') = t'(r')$. Ainsi,
$r' = e'(f(a))$ (où $e'$ est l'identité du groupoïde
associé à $j'$, voir le Lemme \ref{groupoids-lemma-equivalence-groupoid}).
Comme le premier diagramme du lemme est cartésien, cela implique
nécessairement que $r$ est égal à $e(a)$. Ainsi, $a = b$.

\medskip\noindent
Enfin, montrons que la flèche affichée est surjective. Soit
$T$ un schéma sur $S$ et soit $(a', \overline{b})$ une section
du faisceau $U' \times_{U'/R'} U/R$ au-dessus de $T$. Après avoir remplacé $T$
par les membres d'un $\tau$-recouvrement, on peut supposer que $\overline{b}$
est la classe d'un élément $b \in U(T)$. Après avoir remplacé $T$
par les membres d'un $\tau$-recouvrement, on peut supposer qu'il existe
$r' \in R'(T)$ tel que $a' = t(r')$ et $s'(r') = f(b)$.
Comme le premier diagramme du lemme est cartésien, on peut trouver
$r \in R(T)$ tel que $s(r) = b$ et $f(r) = r'$. Il est alors clair
que $a = t(r) \in U(T)$ est une section qui s'envoie sur
$(a', \overline{b})$.
\end{proof}








\section{Descente en termes de groupoïdes}
\label{groupoids-section-groupoids-descent}

\noindent
Les morphismes cartésiens sont définis comme suit.

\begin{definition}
\label{groupoids-definition-cartesian-morphism}
Soit $S$ un schéma. Soit $f : (U', R', s', t', c') \to (U, R, s, t, c)$
un morphisme de schémas en groupoïdes sur $S$. On dit que $f$ est {\it cartésien}, ou
que {\it $(U', R', s', t', c')$ est cartésien au-dessus de $(U, R, s, t, c)$},
si le diagramme
$$
\xymatrix{
R' \ar[r]_f \ar[d]_{s'} & R \ar[d]^s \\
U' \ar[r]^f & U
}
$$
est un carré cartésien dans la catégorie des schémas. Un {\it morphisme de schémas en groupoïdes
cartésien au-dessus de $(U, R, s, t, c)$} est un morphisme de schémas en groupoïdes
compatible avec les morphismes structuraux vers $(U, R, s, t, c)$.
\end{definition}

\noindent
Les morphismes cartésiens sont reliés aux données de descente. Nous démontrons d'abord un
lemme général décrivant la catégorie des schémas en groupoïdes cartésiens au-dessus d'un
schéma en groupoïdes fixé.

\begin{lemma}
\label{groupoids-lemma-characterize-cartesian-schemes}
Soit $S$ un schéma. Soit $(U, R, s, t, c)$ un schéma en groupoïdes sur $S$.
La catégorie des schémas en groupoïdes cartésiens au-dessus de $(U, R, s, t, c)$
est équivalente à la catégorie des couples $(V, \varphi)$, où $V$ est un
schéma sur $U$ et où
$$
\varphi :
V \times_{U, t} R
\longrightarrow
R \times_{s, U} V
$$
est un isomorphisme au-dessus de $R$ tel que $e^*\varphi = \text{id}_V$ et tel que
$$
c^*\varphi = \text{pr}_1^*\varphi \circ \text{pr}_0^*\varphi
$$
comme morphismes de schémas au-dessus de $R \times_{s, U, t} R$.
\end{lemma}

\begin{proof}
Dans le lemme, la notation de tiré en arrière désigne le changement de base. La formule
affichée a un sens parce que
$$
(R \times_{s, U, t} R) \times_{\text{pr}_1, R, \text{pr}_1} (V \times_{U, t} R)
=
(R \times_{s, U, t} R) \times_{\text{pr}_0, R, \text{pr}_0} (R \times_{s, U} V)
$$
comme schémas au-dessus de $R \times_{s, U, t} R$.

\medskip\noindent
Étant donné $(V, \varphi)$, posons $U' = V$ et $R' = V \times_{U, t} R$.
On prend pour $t' : R' \to U'$ la projection $V \times_{U, t} R \to V$.
On prend pour $s'$ le composé de $\varphi$ et de la projection
$R \times_{s, U} V \to V$. On prend pour $c'$ le composé
\begin{align*}
R' \times_{s', U', t'} R'
& \xrightarrow{\varphi, 1}
(R \times_{s, U} V) \times_V (V \times_{U, t} R) \\
& \xrightarrow{}
R \times_{s, U} V \times_{U, t} R \\
& \xrightarrow{\varphi^{-1}, 1}
V \times_{U, t} (R \times_{s, U, t} R) \\
& \xrightarrow{1, c}
V \times_{U, t} R = R'
\end{align*}
Un calcul, que nous omettons, montre que l'on obtient un schéma en groupoïdes
au-dessus de $(U, R, s, t, c)$. Il est clair que ce schéma en groupoïdes est
cartésien au-dessus de $(U, R, s, t, c)$.

\medskip\noindent
Réciproquement, étant donné $f : (U', R', s', t', c') \to (U, R, s, t, c)$
cartésien, les morphismes
$$
U' \times_{U, t} R \xleftarrow{t', f} R' \xrightarrow{f, s'} R \times_{s, U} U'
$$
sont des isomorphismes; on peut donc poser $V = U'$ et prendre pour $\varphi$ le
composé $(f, s') \circ (t', f)^{-1}$. Nous omettons la démonstration du fait que
$\varphi$ satisfait aux conditions du lemme. Nous omettons la démonstration du fait que
ces constructions sont inverses l'une de l'autre.
\end{proof}

\noindent
Soit $S$ un schéma. Soit $f : X \to Y$ un morphisme de schémas sur $S$. On
obtient alors un schéma en groupoïdes $(X, X \times_Y X, \text{pr}_1, \text{pr}_0, c)$
sur $S$. En effet, $j : X \times_Y X \to X \times_S X$ est une relation
d'équivalence et l'on peut prendre le groupoïde associé, voir le
Lemme \ref{groupoids-lemma-equivalence-groupoid}.

\begin{lemma}
\label{groupoids-lemma-cartesian-equivalent-descent-datum}
Soit $S$ un schéma. Soit $f : X \to Y$ un morphisme de schémas sur $S$.
La construction du Lemme \ref{groupoids-lemma-characterize-cartesian-schemes}
détermine une équivalence
$$
\begin{matrix}
\text{catégorie des schémas en groupoïdes} \\
\text{cartésiens au-dessus de } (X, X \times_Y X, \ldots)
\end{matrix}
\longrightarrow
\begin{matrix}
\text{ catégorie des données de descente} \\
\text{ relatives à } X/Y
\end{matrix}
$$
\end{lemma}

\begin{proof}
Cela résulte immédiatement du
Lemme \ref{groupoids-lemma-characterize-cartesian-schemes}
et de la définition des données de descente sur les schémas dans
Descente, Définition \ref{descent-definition-descent-datum}.
\end{proof}







\section{Conditions de séparation}
\label{groupoids-section-separation}

\noindent
Il s'agit en réalité de conditions sur le morphisme $j : R \to U \times_S U$
lorsque l'on se donne un groupoïde $(U, R, s, t, c)$ sur $S$. Comme dans la section
précédente, nous formulons d'abord le diagramme correspondant.

\begin{lemma}
\label{groupoids-lemma-diagram-diagonal}
Soit $S$ un schéma.
Soit $(U, R, s, t, c)$ un groupoïde sur $S$.
Soit $G \to U$ le schéma en groupes stabilisateur.
Le diagramme commutatif
$$
\xymatrix{
R \ar[d]^{\Delta_{R/U \times_S U}} \ar[rrr]_{f \mapsto (f, s(f))} & & &
R \times_{s, U} U \ar[d] \ar[r] & U \ar[d] \\
R \times_{(U \times_S U)} R \ar[rrr]^{(f, g) \mapsto (f, f^{-1} \circ g)} & & &
R \times_{s, U} G \ar[r] & G
}
$$
possède deux flèches horizontales de gauche qui sont des isomorphismes,
et son carré de droite est cartésien.
\end{lemma}

\begin{proof}
Omis.
Exercice sur les définitions et le point de vue fonctoriel en
géométrie algébrique.
\end{proof}

\begin{lemma}
\label{groupoids-lemma-diagonal}
Soit $S$ un schéma.
Soit $(U, R, s, t, c)$ un groupoïde sur $S$.
Soit $G \to U$ le schéma en groupes stabilisateur.
\begin{enumerate}
\item Les conditions suivantes sont équivalentes :
\begin{enumerate}
\item $j : R \to U \times_S U$ est séparé,
\item $G \to U$ est séparé, et
\item $e : U \to G$ est une immersion fermée.
\end{enumerate}
\item Les conditions suivantes sont équivalentes :
\begin{enumerate}
\item $j : R \to U \times_S U$ est quasi-séparé,
\item $G \to U$ est quasi-séparé, et
\item $e : U \to G$ est quasi-compact.
\end{enumerate}
\end{enumerate}
\end{lemma}

\begin{proof}
Le schéma en groupes $G \to U$ est le changement de base de $R \to U \times_S U$
par le morphisme diagonal $U \to U \times_S U$, voir le
Lemme \ref{groupoids-lemma-groupoid-stabilizer}. Par conséquent, si
$j$ est séparé (resp.\ quasi-séparé),
alors $G \to U$ est séparé (resp.\ quasi-séparé).
(Voir Schémas, Lemme
\ref{schemes-lemma-separated-permanence}).
Ainsi, (a) $\Rightarrow$ (b) dans (1) comme dans (2).

\medskip\noindent
Si $G \to U$ est séparé (resp.\ quasi-séparé), alors le morphisme
$U \to G$, en tant que section du morphisme structural $G \to U$, est une immersion
fermée (resp.\ quasi-compact), voir
Schémas, Lemme \ref{schemes-lemma-section-immersion}.
Ainsi, (b) $\Rightarrow$ (a) dans (1) comme dans (2).

\medskip\noindent
D'après le résultat du
Lemme \ref{groupoids-lemma-diagram-diagonal}
(et Schémas, Lemmes \ref{schemes-lemma-base-change-immersion}
et \ref{schemes-lemma-quasi-compact-preserved-base-change}),
on voit que, si $e$ est une immersion fermée (resp.\ quasi-compact),
$\Delta_{R/U \times_S U}$ est une immersion
fermée (resp.\ quasi-compact).
Ainsi, (c) $\Rightarrow$ (a) dans (1) comme dans (2).
\end{proof}









\section{Groupoïdes finis plats, cas affine}
\label{groupoids-section-finite-flat}

\noindent
Soit $S$ un schéma.
Soit $(U, R, s, t, c)$ un schéma en groupoïdes sur $S$.
Supposons que $U = \Spec(A)$ et $R = \Spec(B)$ soient affines.
Dans ce cas, on obtient deux homomorphismes d'anneaux
$s^\sharp, t^\sharp : A \longrightarrow B$.
Soit $C$ l'égalisateur de $s^\sharp$ et $t^\sharp$. En formule,
\begin{equation}
\label{groupoids-equation-invariants}
C = \{a \in A \mid t^\sharp(a) = s^\sharp(a) \}.
\end{equation}
Nous l'appellerons parfois {\it anneau des fonctions $R$-invariantes sur $U$}.
Quelles propriétés $M = \Spec(C)$ possède-t-il ? La première observation est
que le diagramme
$$
\xymatrix{
R \ar[r]_s \ar[d]_t & U \ar[d] \\
U \ar[r] & M
}
$$
est commutatif, c'est-à-dire que le morphisme $U \to M$ égalise $s, t$.
De plus, si $T$ est un schéma affine quelconque, et si $U \to T$ est
un morphisme qui égalise $s, t$, alors $U \to T$ se factorise par $U \to M$.
Autrement dit, $U \to M$ est un coégalisateur dans la catégorie des schémas affines.

\medskip\noindent
Nous aimerions trouver des conditions garantissant que le morphisme $U \to M$ est
véritablement un ``quotient'' dans la catégorie des schémas. Nous en discuterons
longuement ailleurs (insérer ici une référence ultérieure) ; ici, nous examinons seulement
quelques cas particuliers. Plus précisément, nous nous concentrerons sur le cas où $s, t$ sont
finis localement libres.

\begin{example}
\label{groupoids-example-quotient-projective-line}
Soit $k$ un corps. Soit $U = \text{GL}_{2, k}$. Soit $B \subset \text{GL}_2$
le sous-schéma en groupes fermé des matrices triangulaires supérieures.
Alors le faisceau quotient $\text{GL}_{2, k}/B$ (pour la topologie de Zariski, \'etale ou
fppf, voir la Définition \ref{groupoids-definition-quotient-sheaf}) est
représentable par la droite projective : $\mathbf{P}^1 = \text{GL}_{2, k}/B$.
(Détails omis.)
En revanche, l'anneau des fonctions invariantes dans ce cas est simplement $k$.
Notons que, dans ce cas, les morphismes
$s, t : R = \text{GL}_{2, k} \times_k B \to \text{GL}_{2, k} = U$ sont lisses
de dimension relative $3$.
\end{example}

\noindent
Rappelons que, dans Exercices, Exercices \ref{exercises-exercise-trace-det} et
\ref{exercises-exercise-trace-det-rings}, nous avons défini le déterminant
et la norme pour les modules finis localement libres et les extensions d'anneaux
finies localement libres. Si $\varphi : A \to B$ est un homomorphisme d'anneaux fini localement libre, alors
nous noterons $\text{Norm}_\varphi(b) \in A$ la norme de $b \in B$. Dans le
cas d'un morphisme de schémas fini localement libre, la norme a été construite dans
Diviseurs, Lemme \ref{divisors-lemma-finite-locally-free-has-norm}.

\begin{lemma}
\label{groupoids-lemma-determinant-trick}
Soit $S$ un schéma. Soit $(U, R, s, t, c)$ un schéma en groupoïdes sur $S$.
Supposons que $U = \Spec(A)$ et $R = \Spec(B)$ soient affines et que
$s, t : R \to U$ soient finis localement libres.
Soit $C$ comme dans (\ref{groupoids-equation-invariants}).
Soit $f \in A$. Alors $\text{Norm}_{s^\sharp}(t^\sharp(f)) \in C$.
\end{lemma}

\begin{proof}
Considérons le diagramme commutatif
$$
\xymatrix{
& U & \\
R \ar[d]_s \ar[ru]^t &
R \times_{s, U, t} R
\ar[l]^-{\text{pr}_0} \ar[d]^{\text{pr}_1} \ar[r]_-c &
R \ar[d]^s \ar[lu]_t \\
U & R \ar[l]_t \ar[r]^s & U
}
$$
{\emergencystretch=2em Il s'agit du diagramme du lemme \ref{groupoids-lemma-diagram}.
Considérons $f \in \Gamma(U, \mathcal{O}_U)$. La commutativité de la
partie supérieure du diagramme montre que
$\text{pr}_0^\sharp(t^\sharp(f)) = c^\sharp(t^\sharp(f))$ comme éléments de
$\Gamma(R \times_{S, U, t} R, \mathcal{O})$.
En examinant le carré cartésien inférieur droit,
la compatibilité de la construction de la norme avec le changement de base montre que
$s^\sharp(\text{Norm}_{s^\sharp}(t^\sharp(f))) =
\text{Norm}_{\text{pr}_1^\sharp}(c^\sharp(t^\sharp(f)))$.
De même, nous obtenons
$t^\sharp(\text{Norm}_{s^\sharp}(t^\sharp(f))) =
\text{Norm}_{\text{pr}_1^\sharp}(\text{pr}_0^\sharp(t^\sharp(f)))$.
Par conséquent, la première égalité de cette preuve donne
$s^\sharp(\text{Norm}_{s^\sharp}(t^\sharp(f))) =
t^\sharp(\text{Norm}_{s^\sharp}(t^\sharp(f)))$, comme souhaité.\par}
\end{proof}

\begin{lemma}
\label{groupoids-lemma-finite-locally-free-disjoint-free}
Soit $S$ un schéma. Soit $(U, R, s, t, c)$ un schéma en groupoïdes sur $S$.
Supposons que $s, t : R \to U$ soient finis localement libres.
Alors
$$
U = \coprod\nolimits_{r \geq 1} U_r
$$
est une réunion disjointe d'ouverts invariants par $R$ telle que la restriction $R_r$ de
$R$ à $U_r$ soit telle que $s, t : R_r \to U_r$ soient finis localement
libres de rang $r$.
\end{lemma}

\begin{proof}
D'après
Morphismes, Lemme \ref{morphisms-lemma-finite-locally-free},
il existe une décomposition
$U = \coprod\nolimits_{r \geq 0} U_r$
telle que $s : s^{-1}(U_r) \to U_r$ soit fini localement libre de rang $r$.
Comme $s$ est surjectif, on voit que $U_0 = \emptyset$.
Notons que $u \in U_r \Leftrightarrow$ la fibre schématique
$s^{-1}(u)$ est de degré $r$ sur $\kappa(u)$. Maintenant, si $z \in R$ avec $s(z) = u$
et $t(z) = u'$, alors, avec les notations du Lemme \ref{groupoids-lemma-diagram},
$$
\text{pr}_1^{-1}(z) \to \Spec(\kappa(z))
$$
est à la fois le changement de base de
$s^{-1}(u) \to \Spec(\kappa(u))$ et de $s^{-1}(u') \to \Spec(\kappa(u'))$
d'après le lemme cité. Ainsi $u \in U_r \Leftrightarrow u' \in U_r$ ;
autrement dit, les ouverts $U_r$ sont invariants par $R$.
En particulier, la restriction de $R$ à $U_r$ est simplement
$s^{-1}(U_r)$ et $s : R_r \to U_r$ est fini localement libre de rang $r$.
Comme $t : R_r \to U_r$ est isomorphe à $s$ par l'inversion de $R_r$,
on voit qu'il est également de rang $r$.
\end{proof}

\begin{lemma}
\label{groupoids-lemma-integral-over-invariants}
Soit $S$ un schéma. Soit $(U, R, s, t, c)$ un schéma en groupoïdes sur $S$.
Supposons que $U = \Spec(A)$ et $R = \Spec(B)$ soient affines et que
$s, t : R \to U$ soient finis localement libres.
Soit $C \subset A$ comme dans (\ref{groupoids-equation-invariants}).
Alors $A$ est entier sur $C$.
\end{lemma}

\begin{proof}
Tout d'abord, le Lemme \ref{groupoids-lemma-finite-locally-free-disjoint-free}
nous apprend que $(U, R, s, t, c)$ est une réunion disjointe
de schémas en groupoïdes $(U_r, R_r, s, t, c)$ tels que chaque $s, t : R_r \to U_r$
soit de rang constant $r$. Comme $U$ est quasi-compact, on a $U_r = \emptyset$ pour
presque tout $r$. Il suffit de démontrer le lemme pour chaque $(U_r, R_r, s, t, c)$ ;
nous pouvons donc supposer que $s, t$ sont finis localement libres de rang $r$.

\medskip\noindent
Supposons que $s, t$ soient finis localement libres de rang $r$.
Soit $f \in A$. Considérons l'élément $x - f \in A[x]$, où nous regardons
$x$ comme la coordonnée sur $\mathbf{A}^1$.
Puisque
$$
(U \times \mathbf{A}^1, R \times \mathbf{A}^1,
s \times \text{id}_{\mathbf{A}^1},
t \times \text{id}_{\mathbf{A}^1},
c \times \text{id}_{\mathbf{A}^1})
$$
est lui aussi un schéma en groupoïdes dont la source et le but sont finis, nous pouvons lui appliquer
le Lemme \ref{groupoids-lemma-determinant-trick} et nous voyons que
$P(x) = \text{Norm}_{s^\sharp}(t^\sharp(x - f))$
est un élément de $C[x]$. Puisque $s^\sharp : A \to B$ est fini localement
libre de rang $r$, on voit que $P$ est unitaire de degré $r$.
De plus, $P(f) = 0$ d'après Cayley-Hamilton
(Algèbre, Lemme \ref{algebra-lemma-charpoly}).
Plus précisément, Cayley-Hamilton affirme que le polynôme
$s^\sharp(P) \in B[x]$ annule l'élément $t^\sharp(f)$.
Puisque les coefficients de $P$ appartiennent à $C$, nous obtenons
$t^\sharp(P(f)) = 0$ dans $B$. Comme $t$ est fidèlement plat,
nous obtenons $P(f) = 0$.
\end{proof}

\begin{lemma}
\label{groupoids-lemma-invariants-base-change}
Soit $S$ un schéma. Soit $(U, R, s, t, c)$ un schéma en groupoïdes sur $S$.
Supposons que $U = \Spec(A)$ et $R = \Spec(B)$ soient affines et que
$s, t : R \to U$ soient finis localement libres. Soit $C \subset A$ comme dans
(\ref{groupoids-equation-invariants}). Soit $C \to C'$ un homomorphisme d'anneaux, et posons
$U' = \Spec(A \otimes_C C')$,
$R' = \Spec(B \otimes_C C')$.
Alors
\begin{enumerate}
\item Les applications $s, t, c$ donnent des applications $s', t', c'$.
Celles-ci font de $(U', R', s', t', c')$ un schéma en groupoïdes. Soit $C^1 \subset A'$
l'anneau des fonctions $R'$-invariantes sur $U'$.
\item L'application canonique $\varphi : C' \to C^1$ satisfait aux propriétés suivantes :
\begin{enumerate}
\item pour tout $f \in C^1$, il existe un $n > 0$ et un
polynôme $P \in C'[x]$ dont l'image dans $C^1[x]$ est
$(x - f)^n$, et
\item pour tout $f \in \Ker(\varphi)$, il existe
un $n > 0$ tel que $f^n = 0$.
\end{enumerate}
\item Si $C \to C'$ est plat, alors $\varphi$ est un isomorphisme.
\end{enumerate}
\end{lemma}

\begin{proof}
La preuve de la partie (1) est omise. Notons $A' = A \otimes_C C'$ et
$B' = B \otimes_C C'$. Alors nous avons
$$
C^1
= \{a \in A' \mid (t')^\sharp(a) = (s')^\sharp(a) \}
= \{a \in A \otimes_C C' \mid t^\sharp \otimes 1(a) = s^\sharp \otimes 1(a) \}.
$$
Autrement dit, $C^1$ est le noyau de l'application différence
$(t^\sharp - s^\sharp) \otimes 1$, qui est précisément le changement de base
de l'application $C$-linéaire $t^\sharp - s^\sharp : A \to B$ par $C \to C'$.
Ainsi, (3) en résulte.

\medskip\noindent
Preuve de la partie (2)(b). Puisque $C \to A$ est entier
(Lemme \ref{groupoids-lemma-integral-over-invariants}) et injectif, on voit que
$\Spec(A) \to \Spec(C)$ est surjectif, voir
Algèbre, Lemme \ref{algebra-lemma-integral-overring-surjective}.
Ainsi, $\Spec(A') \to \Spec(C')$ est également surjectif
comme changement de base d'un morphisme surjectif
(Morphismes, Lemme \ref{morphisms-lemma-base-change-surjective}).
Par conséquent, $\Spec(C^1) \to \Spec(C')$ est lui aussi surjectif.
Cela implique (2)(b), par exemple d'après
Algèbre, Lemme \ref{algebra-lemma-image-dense-generic-points}.

\medskip\noindent
Preuve de la partie (2)(a). D'après le Lemme \ref{groupoids-lemma-finite-locally-free-disjoint-free},
notre schéma en groupoïdes $(U, R, s, t, c)$ se décompose en une réunion disjointe finie
de schémas en groupoïdes $(U_r, R_r, s, t, c)$ tels que $s, t : R_r \to U_r$
soient finis localement libres de rang $r$. En tirant en arrière par $U' = \Spec(C') \to U$,
nous obtenons une décomposition analogue de $U'$ et de $U^1 = \Spec(C^1)$.
Nous montrerons dans le paragraphe suivant que (2)(a) vaut pour le système correspondant
d'anneaux $A_r, B_r, C_r, C'_r, C^1_r$, avec $n = r$.
Étant donné $f \in C^1$, soit alors $P_r \in C_r[x]$ le polynôme
dont l'image dans $C^1_r[x]$ est l'image de $(x - f)^r$.
En choisissant un entier $n$ suffisamment divisible, on voit qu'il
existe un polynôme $P \in C'[x]$ dont l'image dans $C^1[x]$ est
$(x - f)^n$ ; à savoir, nous prenons pour $P$ l'unique élément de
$C'[x]$ dont l'image dans $C'_r[x]$ est $P_r^{n/r}$.

\medskip\noindent
Dans ce paragraphe, nous démontrons (2)(a) lorsque les homomorphismes d'anneaux
$s^\sharp, t^\sharp : A \to B$ sont finis localement libres d'un rang fixé $r$.
Soit $f \in C^1 \subset A' = A \otimes_C C'$. Choisissons une
$C$-algèbre plate $D$ et une surjection $D \to C'$. Choisissons un relèvement
$g \in A \otimes_C D$ de $f$.
Considérons le polynôme
$$
P = \text{Norm}_{s^\sharp \otimes 1}(t^\sharp \otimes 1(x - g))
$$
dans $(A \otimes_C D)[x]$. D'après le Lemme \ref{groupoids-lemma-determinant-trick}
et la partie (3) du présent lemme, les coefficients de $P$ appartiennent à $D$
(comparer avec la preuve du Lemme \ref{groupoids-lemma-integral-over-invariants}).
D'autre part, l'image de $P$ dans $(A \otimes_C C')[x]$ est
$(x - f)^r$, car $t^\sharp \otimes 1(x - f) = s^\sharp(x - f)$
et $s^\sharp$ est fini localement libre de rang $r$.
Cela démontre l'assertion avec $P$ comme dans l'énoncé (2)(a),
donné par l'image de notre $P$ par l'application $D[x] \to C'[x]$.
\end{proof}

\begin{lemma}
\label{groupoids-lemma-points}
Soit $S$ un schéma. Soit $(U, R, s, t, c)$ un schéma en groupoïdes sur $S$.
Supposons que $U = \Spec(A)$ et $R = \Spec(B)$ soient affines et que
$s, t : R \to U$ soient finis localement libres. Soit $C \subset A$ comme dans
(\ref{groupoids-equation-invariants}). Alors $U \to M = \Spec(C)$ possède
les propriétés suivantes :
\begin{enumerate}
\item l'application sur les points $|U| \to |M|$ est surjective et
$u_0, u_1 \in |U|$ ont la même image si et seulement s'il
existe un $r \in |R|$ tel que $t(r) = u_0$ et $s(r) = u_1$, ce qui donne
la formule
$$
|M| = |U|/|R|
$$
\item pour tout corps algébriquement clos $k$, on a
$$
M(k) = U(k)/R(k)
$$
\end{enumerate}
\end{lemma}

\begin{proof}
Puisque $C \to A$ est entier (Lemme \ref{groupoids-lemma-integral-over-invariants})
et injectif, on voit que $\Spec(A) \to \Spec(C)$ est surjectif, voir
Algèbre, Lemme \ref{algebra-lemma-integral-overring-surjective}.
Ainsi, $|U| \to |M|$ est surjectif.

\medskip\noindent
Soit $k$ un corps algébriquement clos et soit $C \to k$ un homomorphisme d'anneaux.
Puisque les morphismes surjectifs sont préservés par changement de base
(Morphismes, Lemme \ref{morphisms-lemma-base-change-surjective}),
on voit que $A \otimes_C k$ n'est pas nul. Or $k \subset A \otimes_C k$ est une
extension entière non nulle. Ainsi, tout corps résiduel de $A \otimes_C k$
est une extension algébrique de $k$, donc est égal à $k$. Nous voyons donc que
$U(k) \to M(k)$ est surjectif.

\medskip\noindent
Soient $a_0, a_1 : A \to k$ deux homomorphismes d'anneaux. S'il existe un homomorphisme d'anneaux
$b : B \to k$ telle que $a_0 = b \circ t^\sharp$ et $a_1 = b \circ s^\sharp$,
alors on voit par définition que $a_0|_C = a_1|_C$. Ainsi, l'application
$U(k) \to M(k)$ égalise les deux applications $R(k) \to U(k)$.
Réciproquement, supposons que $a_0|_C = a_1|_C$. Notons cette application
d'algèbres $c : C \to k$. Considérons le diagramme
$$
\xymatrix{
& &
B \ar@{-->}[lld] \\
k & &
A
\ar@<0.5ex>[ll]^{a_0}
\ar@<-0.5ex>[ll]_{a_1}
\ar@<1ex>[u]
\ar@<-1ex>[u] \\
& &
C \ar[u] \ar[llu]^c
}
$$
Si nous pouvons construire une flèche pointillée rendant le diagramme commutatif, alors
la preuve de la partie (2) du lemme est achevée. Puisque $s : A \to B$ est fini,
il existe un nombre fini d'homomorphismes d'anneaux
$b_1, \ldots, b_n : B \to k$ telles que $b_i \circ s^\sharp = a_1$.
Si la flèche pointillée n'existe pas, nous voyons qu'aucune des
$a'_i = b_i \circ t^\sharp$, $i = 1, \ldots, n$, n'est égale à $a_0$.
Par conséquent, les idéaux maximaux
$$
\mathfrak m'_i = \Ker(a_i' \otimes 1 : A \otimes_C k \to k)
$$
de $A \otimes_C k$ sont distincts de
$\mathfrak m = \Ker(a_0 \otimes 1 : A \otimes_C k \to k)$.
D'après Algèbre, Lemme \ref{algebra-lemma-silly}, nous obtiendrions un élément
$f \in A \otimes_C k$ tel que $f \in \mathfrak m$, mais que
$f \not \in \mathfrak m_i'$ pour $i = 1, \ldots, n$.
Considérons la norme
$$
g = \text{Norm}_{s^\sharp \otimes 1}(t^\sharp \otimes 1(f))
\in
A \otimes_C k
$$
D'après le Lemme \ref{groupoids-lemma-determinant-trick}, celle-ci appartient à l'anneau des invariants
$C^1 \subset A \otimes_C k$ du schéma en groupoïdes obtenu par
changement de base (au moyen de l'application $c : C \to k$). D'une part,
$a_1(g) \in k^*$, puisque
la valeur de $t^\sharp(f)$ en tous les points (qui correspondent à
$b_1, \ldots, b_n$) situés au-dessus de $a_1$ est
inversible (insérer ici une référence ultérieure à cette propriété du déterminant).
D'autre part, puisque $f \in \mathfrak m$, nous voyons que
$f$ n'est pas une unité, et donc que $t^\sharp(f)$ n'est pas une unité
(car $t^\sharp \otimes 1$ est fidèlement plat),
et donc que sa norme n'est pas une unité (insérer ici une référence ultérieure
à cette propriété du déterminant). Nous concluons que $C^1$ contient
un élément qui n'est ni nilpotent
ni une unité. Nous allons maintenant montrer que cela conduit à une contradiction.
En effet, appliquons le Lemme \ref{groupoids-lemma-invariants-base-change}
à l'application $c : C \to C' = k$ ; alors
nous voyons que l'application de $k$ dans l'anneau des invariants $C^1$ est injective
et que, de plus, pour tout élément $x \in C^1$, il existe un entier
$n > 0$ tel que $x^n \in k$. Ainsi, tout élément de $C^1$ est
soit une unité, soit nilpotent.

\medskip\noindent
Il nous reste à achever la preuve de (1). Nous savons déjà que
$|U| \to |M|$ est surjectif. Il est clair que $|U| \to |M|$ est
invariante par $|R|$. Enfin, supposons que $u_0, u_1 \in U$ aient la même
image $m \in M$. Alors les extensions de corps induites $\kappa(u_0)/\kappa(m)$
et $\kappa(u_1)/\kappa(m)$ sont algébriques (car $A$ est entier sur $C$,
comme utilisé ci-dessus). Par conséquent, si $k$ est une clôture algébrique de $\kappa(m)$,
nous pouvons trouver des $\kappa(m)$-plongements $\overline{u}_0 : \kappa(u_0) \to k$
et $\overline{u}_1 : \kappa(u_1) \to k$. Ceux-ci déterminent des points à valeurs dans $k$,
$\overline{u}_0, \overline{u}_1 \in U(k)$, ayant la même
image dans $M(k)$. D'après la partie (2), il existe un point
$\overline{r} \in R(k)$ tel que $s(\overline{r}) = \overline{u}_0$ et
$t(\overline{r}) = \overline{u}_1$. L'image $r \in R$ de $\overline{r}$
est un point tel que $s(r) = u_0$ et $t(r) = u_1$, comme souhaité.
\end{proof}

\begin{lemma}
\label{groupoids-lemma-etale}
Soit $S$ un schéma. Soit $f : (U', R', s', t') \to (U, R, s, t, c)$ un
morphisme de schémas en groupoïdes sur $S$. Supposons que
\begin{enumerate}
\item $U$, $R$, $U'$, $R'$ soient affines,
\item $s, t, s', t'$ soient finis localement libres,
\item les diagrammes
$$
\xymatrix{
R' \ar[d]_{s'} \ar[r]_f & R \ar[d]^s \\
U' \ar[r]^f & U
}
\quad
\quad
\xymatrix{
R' \ar[d]_{t'} \ar[r]_f & R \ar[d]^t \\
U' \ar[r]^f & U
}
\quad
\quad
\xymatrix{
G' \ar[d] \ar[r]_f & G \ar[d] \\
U' \ar[r]^f & U
}
$$
soient cartésiens, où $G$ et $G'$ sont les schémas en groupes stabilisateurs, et
\item $f : U' \to U$ soit \'etale.
\end{enumerate}
Alors l'application $C \to C'$ de l'anneau des fonctions $R$-invariantes sur $U$
vers l'anneau des fonctions $R'$-invariantes sur $U'$ est \'etale et
$U' = \Spec(C') \times_{\Spec(C)} U$.
\end{lemma}

\begin{proof}
Posons $M = \Spec(C)$ et $M' = \Spec(C')$.
Écrivons $U = \Spec(A)$, $U' = \Spec(A')$, $R = \Spec(B)$ et
$R' = \Spec(B')$. Nous utiliserons les résultats des
Lemmes \ref{groupoids-lemma-integral-over-invariants},
\ref{groupoids-lemma-invariants-base-change} et
\ref{groupoids-lemma-points}
sans autre mention.

\medskip\noindent
Supposons que $C$ soit un anneau local strictement hensélien. Soit $p \in M$
le point fermé et soit $p' \in M'$ un point au-dessus de $p$.
Affirmation : dans ce cas, il existe une décomposition en réunion disjointe
$(U', R', s', t', c') = (U, R, s, t, c) \amalg (U'', R'', s'', t'', c'')$
au-dessus de $(U, R, s, t, c)$ telle que, pour la décomposition en
réunion disjointe correspondante $M' = M \amalg M''$ au-dessus de $M$,
le point $p'$ corresponde à $p \in M$.

\medskip\noindent
L'affirmation implique le lemme. Supposons que $M_1 \to M$ soit un morphisme plat
de schémas affines. Nous pouvons alors tout changer de base vers $M_1$
sans modifier les hypothèses (1) -- (4).
Le Lemme \ref{groupoids-lemma-invariants-base-change} montre que
$M_1$, resp.\ $M_1'$, est le spectre de l'anneau des
fonctions $R_1$-invariantes sur $U_1$,
resp.\ des fonctions $R'_1$-invariantes sur $U'_1$.
Supposons que $p' \in M'$ ait pour image $p \in M$.
Soit $M_1$ le spectre de l'hensélisé strict de
$\mathcal{O}_{M, p}$, de point fermé $p_1 \in M_1$.
Choisissons un point $p'_1 \in M'_1$ au-dessus de $p_1$ et de $p'$.
L'affirmation donne
$$
(U'_1, R'_1, s'_1, t'_1, c'_1) =
(U_1, R_1, s_1, t_1, c_1) \amalg
(U''_1, R''_1, s''_1, t''_1, c''_1)
$$
et, de manière correspondante, $M'_1 = M_1 \amalg M''_1$ comme schéma au-dessus de $M_1$.
Écrivons $M_1 = \Spec(C_1)$ et $C_1 = \colim C_i$ comme une limite inductive filtrante
de $C$-algèbres \'etales. Posons $M_i = \Spec(C_i)$.
Alors $M_1 = \lim M_i$, et de même pour les autres schémas.
D'après Limites, Lemmes \ref{limits-lemma-descend-opens} et
\ref{limits-lemma-descend-isomorphism},
nous pouvons trouver un $i$ tel que
$$
(U'_i, R'_i, s'_i, t'_i, c'_i) =
(U_i, R_i, s_i, t_i, c_i) \amalg
(U''_i, R''_i, s''_i, t''_i, c''_i)
$$
Nous concluons que $M'_i = M_i \amalg M''_i$. En particulier,
$M' \to M$ devient \'etale en un point au-dessus de $p'$ après un
changement de base \'etale. Cela implique que $M' \to M$ est \'etale en $p'$
(par exemple d'après Morphismes, Lemme
\ref{morphisms-lemma-set-points-where-fibres-etale}).
Nous démontrerons $U' \cong M' \times_M U$ après avoir prouvé l'affirmation.

\medskip\noindent
Preuve de l'affirmation. Observons que $U_p$ et $U'_{p'}$ ont un nombre fini de points.
Pour $u \in U_p$, l'extension $\kappa(u)/\kappa(p)$ est algébrique ;
ainsi, $\kappa(u)$ est séparablement clos.
Comme $U' \to U$ est \'etale, nous concluons que le morphisme $U'_{p'} \to U_p$
induit des isomorphismes de corps résiduels.
Soit $u' \in U'_{p'}$ d'image $u \in U_p$.
Par l'hypothèse (3), les morphismes de fibres schématiques
$(s')^{-1}(u') \to s^{-1}(u)$,
$(t')^{-1}(u') \to t^{-1}(u)$ et
$G'_{u'} \to G_u$ sont des isomorphismes. Comme $U_p = t(s^{-1}(u))$
(ensemblistement), nous concluons que l'application de l'ensemble des points de $U'_{p'}$
dans l'ensemble des points de $U_p$ est surjective.
Supposons que $u'_1$ et $u'_2$ soient des points de $U'_{p'}$ ayant
la même image $u$ dans $U_p$. Il existe alors un point
$r' \in R'_{p'}$ tel que $s'(r') = u'_1$ et $t'(r') = u'_2$.
Considérons les deux tours de corps
$$
\kappa(r')/\kappa(u'_1)/\kappa(u)/\kappa(p) \quad
\kappa(r')/\kappa(u'_2)/\kappa(u)/\kappa(p)
$$
dont les ``extrémités'' sont les mêmes que les deux ``extrémités'' des deux tours
$$
\kappa(r')/\kappa(u'_1)/\kappa(p')/\kappa(p) \quad
\kappa(r')/\kappa(u'_2)/\kappa(p')/\kappa(p)
$$
{\emergencystretch=3em Ces deux tours induisent les mêmes applications $\kappa(p') \to \kappa(r')$, puisque
$(U'_{p'}, R'_{p'}, s', t', c')$ est un schéma en groupoïdes sur $p'$.
Puisque $\kappa(u)/\kappa(p)$ est purement inséparable,
nous concluons que les deux applications induites
$\kappa(u) \to \kappa(r')$ sont les mêmes.
Par conséquent, $r'$ s'envoie sur un point de la fibre $G_u$.
Par l'hypothèse (3), nous concluons que $r' \in (G')_{u'_1}$.
En effet, nous pouvons considérer $G$ comme un sous-schéma fermé de $R$,
vu comme un schéma sur $U$ au moyen de $s$, et utiliser le fait que
son changement de base à $U'$ donne $G' \subset R'$.
En particulier, nous avons $u'_1 = u'_2$.
Nous concluons que $U'_{p'} \to U_p$ est une application bijective
sur les points et qu'elle induit des isomorphismes sur les corps résiduels.
Il s'ensuit que $U'_{p'}$ est un ensemble fini de points fermés
(Algèbre, Lemme \ref{algebra-lemma-finite-residue-extension-closed})
et donc que $U'_{p'}$ est fermé dans $U'$.
Soit $J' \subset A'$ l'idéal radical définissant $U'_{p'}$
ensemblistement.\par}

\medskip\noindent
Seconde partie de la preuve de l'affirmation.
Soit $\mathfrak m \subset C$ l'idéal maximal.
Observons que $(A, \mathfrak m A)$ est un couple hensélien d'après
Compléments d'algèbre, Lemme \ref{more-algebra-lemma-integral-over-henselian-pair}.
Soit $J = \sqrt{\mathfrak m A}$.
Alors $(A, J)$ est un couple hensélien
(Compléments d'algèbre, Lemme \ref{more-algebra-lemma-change-ideal-henselian-pair})
et l'homomorphisme d'anneaux \'etale
$A \to A'$ induit un isomorphisme $A/J \to A'/J'$
d'après nos considérations ci-dessus.
Nous concluons que $A' = A \times A''$ d'après
Compléments d'algèbre, Lemme \ref{more-algebra-lemma-characterize-henselian-pair}.
Considérons la décomposition en réunion disjointe
correspondante $U' = U \amalg U''$. L'ouvert $(s')^{-1}(U)$ est l'ensemble
des points de $R'$ qui se spécialisent en un point de $R'_{p'}$.
Il en va de même pour $(t')^{-1}(U)$. De même, nous avons
$(s')^{-1}(U'') = (t')^{-1}(U'')$, car il s'agit de l'ensemble des
points qui ne se spécialisent pas dans $R'_{p'}$.
Nous obtenons donc une décomposition en réunion disjointe
$$
(U', R', s', t', c') =
(U, R, s, t, c) \amalg
(U'', R'', s'', t'', c'')
$$
Celle-ci donne immédiatement $M' = M \amalg M''$, et la preuve de l'affirmation
est achevée.

\medskip\noindent
Il nous reste à montrer que l'application canonique $U' \to M' \times_M U$
est un isomorphisme. C'est un morphisme \'etale
(Morphismes, Lemme \ref{morphisms-lemma-etale-permanence}).
D'autre part, en effectuant des changements de base vers des anneaux locaux strictement henséliens
(comme dans le troisième paragraphe de la preuve) et en utilisant la bijectivité
$U'_{p'} \to U_p$ établie au cours de la preuve de l'affirmation,
nous voyons que $U' \to M' \times_M U$ est universellement bijectif
(certains détails sont omis). Or un morphisme universellement bijectif
et \'etale est un isomorphisme
(Descente, Lemme \ref{descent-lemma-universally-injective-etale-open-immersion}),
ce qui achève la preuve.
\end{proof}

\begin{lemma}
\label{groupoids-lemma-basis}
Soit $S$ un schéma.
Soit $(U, R, s, t, c)$ un schéma en groupoïdes sur $S$.
Supposons que
\begin{enumerate}
\item $U = \Spec(A)$ et $R = \Spec(B)$ soient affines, et qu'
\item il existe des éléments $x_i \in A$, $i \in I$, tels que
$B = \bigoplus_{i \in I} s^\sharp(A)t^\sharp(x_i)$.
\end{enumerate}
Alors $A = \bigoplus_{i\in I} Cx_i$ et $B \cong A \otimes_C A$,
où $C \subset A$ est l'anneau des fonctions $R$-invariantes
sur $U$, comme dans (\ref{groupoids-equation-invariants}).
\end{lemma}

\begin{proof}
Dans cette preuve, nous écrirons $s, t : A \to B$ au lieu de
$s^\sharp, t^\sharp$, et de même $c : B \to B \otimes_{s, A, t} B$.
Nous écrivons $p_0 : B \to B \otimes_{s, A, t} B$, $b \mapsto b \otimes 1$ et
$p_1 : B \to B \otimes_{s, A, t} B$, $b \mapsto 1 \otimes b$. D'après le
Lemme \ref{groupoids-lemma-diagram-pull}
et la définition de $C$, nous avons le diagramme
commutatif suivant :
$$
\xymatrix{
B \otimes_{s, A, t} B &
B \ar@<-1ex>[l]_-c \ar@<1ex>[l]^-{p_0} &
A \ar[l]^t \\
B \ar[u]^{p_1} &
A \ar@<-1ex>[l]_s \ar@<1ex>[l]^t \ar[u]_s &
C \ar[u] \ar[l]
}
$$
De plus, les deux carrés de gauche sont cocartésiens dans la catégorie des anneaux, et
la ligne supérieure est isomorphe au diagramme
$$
\xymatrix{
B \otimes_{t, A, t} B &
B \ar@<-1ex>[l]_-{p_1} \ar@<1ex>[l]^-{p_0} &
A \ar[l]^t
}
$$
qui est un diagramme égalisateur d'après
Descente, Lemme \ref{descent-lemma-ff-exact}, car la condition (2) implique
en particulier que $s$ (et donc aussi la flèche $t$ qui lui est isomorphe)
est fidèlement plat.
La ligne inférieure est un diagramme égalisateur par définition de $C$.
Nous pouvons utiliser les $x_i$ et obtenir un diagramme commutatif
$$
\xymatrix{
B \otimes_{s, A, t} B &
B \ar@<-1ex>[l]_-c \ar@<1ex>[l]^-{p_0} &
A \ar[l]^t \\
\bigoplus_{i \in I} B x_i \ar[u]^{p_1} &
\bigoplus_{i \in I} A x_i \ar@<-1ex>[l]_s \ar@<1ex>[l]^t \ar[u]_s &
\bigoplus_{i \in I} C x_i \ar[u] \ar[l]
}
$$
où la flèche verticale de droite envoie $x_i$ sur $x_i$,
la flèche verticale du milieu envoie $x_i$ sur $t(x_i)$ et
la flèche verticale de gauche envoie $x_i$ sur
$c(t(x_i)) = t(x_i) \otimes 1 = p_0(t(x_i))$ (égalité résultant de la commutativité
de la partie supérieure du diagramme du lemme \ref{groupoids-lemma-diagram}). Le diagramme
est alors commutatif. De plus, la flèche verticale du milieu est un isomorphisme
par hypothèse. Les deux carrés de gauche étant cocartésiens, on
en conclut que la flèche verticale de gauche est elle aussi un isomorphisme.
D'autre part, les lignes horizontales sont exactes (c'est-à-dire que ce sont des
égalisateurs). On en conclut donc que la flèche verticale de droite
est elle aussi un isomorphisme.
\end{proof}

\begin{proposition}
\label{groupoids-proposition-finite-flat-equivalence}
Soit $S$ un schéma.
Soit $(U, R, s, t, c)$ un schéma en groupoïdes sur $S$.
Supposons que
\begin{enumerate}
\item $U = \Spec(A)$ et $R = \Spec(B)$ soient affines,
\item $s, t : R \to U$ soient finis localement libres, et
\item $j = (t, s)$ soit une relation d'équivalence.
\end{enumerate}
Dans ce cas, soit $C \subset A$ comme dans
(\ref{groupoids-equation-invariants}). Alors $U \to M = \Spec(C)$
est fini localement libre et $R = U \times_M U$.
De plus, $M$ représente le faisceau quotient $U/R$
pour la topologie fppf (voir la définition \ref{groupoids-definition-quotient-sheaf}).
\end{proposition}

\begin{proof}
Dans cette démonstration, on utilise la notation $s, t : A \to B$
au lieu de la notation $s^\sharp, t^\sharp$.
D'après le lemme \ref{groupoids-lemma-criterion-quotient-representable},
il suffit de montrer que $C \to A$ est fini localement libre
et que l'application
$$
t \otimes s : A \otimes_C A \longrightarrow B
$$
est un isomorphisme. Notons d'abord que $j$ est un monomorphisme, et
qu'il est aussi fini (puisque $s$ et $t$ sont déjà finis). On voit donc
que $j$ est une immersion fermée d'après
Morphismes, lemme \ref{morphisms-lemma-finite-monomorphism-closed}.
Ainsi, $A \otimes_C A \to B$ est surjective.

\medskip\noindent
On effectuera des changements de base par des homomorphismes plats d'anneaux $C \to C'$, comme dans
le lemme \ref{groupoids-lemma-invariants-base-change}, et l'on utilisera le fait que
la formation des invariants commute au changement de base plat, voir
la partie (3) du lemme cité.
On montrera ci-dessous que, pour tout idéal premier $\mathfrak p \subset C$, il existe
un homomorphisme local plat d'anneaux $C_{\mathfrak p} \to C_{\mathfrak p}'$
tel que le résultat soit vrai après changement de base à $C_{\mathfrak p}'$.
Cela implique immédiatement
que $A \otimes_C A \to B$ est injective (utiliser
Algèbre, lemme \ref{algebra-lemma-characterize-zero-local}).
Cela implique aussi que $C \to A$ est plat, en combinant
Algèbre, lemmes \ref{algebra-lemma-local-flat-ff},
\ref{algebra-lemma-flat-localization} et
\ref{algebra-lemma-flatness-descends}. Puisque $U \to \Spec(C)$
est également surjectif (lemme \ref{groupoids-lemma-points}), on en conclut que $C \to A$
est fidèlement plat. L'isomorphisme $B \cong A \otimes_C A$
implique alors que $A$ est un $C$-module de présentation finie, voir
Algèbre, lemme \ref{algebra-lemma-descend-properties-modules}.
Ainsi, $A$ est fini localement libre sur $C$, voir
Algèbre, lemme \ref{algebra-lemma-finite-projective}.

\medskip\noindent
D'après le lemme \ref{groupoids-lemma-finite-locally-free-disjoint-free},
on sait que $A$ est un produit fini
d'anneaux $A_r$ et que $B$ est un produit fini d'anneaux $B_r$,
de sorte que le schéma en groupoïdes se décompose de façon correspondante (voir la démonstration
du lemme \ref{groupoids-lemma-integral-over-invariants}).
Alors $C$ est lui aussi un produit d'anneaux $C_r$, et
$C'$ se décompose de façon correspondante en un produit. On peut donc, et l'on va,
supposer que les homomorphismes d'anneaux $s, t : A \to B$ sont finis
localement libres d'un rang fixé $r$.

\medskip\noindent
Les homomorphismes locaux d'anneaux $C_{\mathfrak p} \to C_{\mathfrak p}'$ que l'on va
utiliser sont des morphismes locaux plats quelconques tels que le corps résiduel de
$C_{\mathfrak p}'$ soit infini.
D'après Algèbre, lemme \ref{algebra-lemma-flat-local-given-residue-field},
de tels morphismes locaux existent.

\medskip\noindent
Supposons que $C$ soit un anneau local d'idéal maximal $\mathfrak m$ et
de corps résiduel infini, et que $s, t : A \to B$ soient
finis localement libres de rang constant $r > 0$.
Comme l'extension $C \subset A$ est entière (lemme \ref{groupoids-lemma-integral-over-invariants}),
tous les idéaux premiers au-dessus de $\mathfrak m$ sont maximaux, et tous les idéaux maximaux
de $A$ sont au-dessus de $\mathfrak m$. Il en va de même pour $C \subset B$.
Choisissons un idéal maximal $\mathfrak m'$
de $A$ au-dessus de $\mathfrak m$ (il en existe d'après le lemme \ref{groupoids-lemma-points}).
Comme $t : A \to B$ est fini localement libre, il n'existe qu'un nombre fini
d'idéaux maximaux de $B$ au-dessus de $\mathfrak m'$. On en conclut
(de nouveau par le lemme \ref{groupoids-lemma-points})
que $A$ a un nombre fini d'idéaux maximaux, c'est-à-dire que
$A$ est semi-local. Cela implique à son tour que $B$ est également
semi-local. Or, puisque $t \otimes s : A \otimes_C A \to B$ est surjective,
on peut appliquer
Algèbre, lemme \ref{algebra-lemma-semi-local-module-basis-in-submodule}
à l'homomorphisme d'anneaux $C \to A$, au $A$-module $M = B$ (considéré comme un $A$-module
via $t$) et au $C$-sous-module $s(A) \subset B$. Ce lemme implique qu'il
existe $x_1, \ldots, x_r \in A$ tels que $M$ soit libre sur $A$
de base $s(x_1), \ldots, s(x_r)$. On en conclut que $C \to A$ est fini libre
et que $B \cong A \otimes_C A$, en appliquant
le lemme \ref{groupoids-lemma-basis}.
\end{proof}



\section{Groupoïdes finis plats}
\label{groupoids-section-finite-flat-general}

\noindent
Dans cette section, on démontre un lemme qui aidera à montrer que le quotient
d'un schéma par une relation d'équivalence finie plate est un schéma, pourvu que
chaque classe d'équivalence soit contenue dans un ouvert affine. Voir
Propriétés des espaces,
proposition \ref{spaces-properties-proposition-finite-flat-equivalence-global}.

\begin{lemma}
\label{groupoids-lemma-find-invariant-affine}
Soit $S$ un schéma.
Soit $(U, R, s, t, c)$ un schéma en groupoïdes sur $S$.
Supposons que $s$, $t$ soient finis localement libres.
Soit $u \in U$ un point tel que $t(s^{-1}(\{u\}))$
soit contenu dans un ouvert affine de $U$.
Alors il existe un voisinage ouvert affine $R$-invariant
de $u$ dans $U$.
\end{lemma}

\begin{proof}
Puisque $s$ est fini localement libre, ses fibres sont finies. Ainsi,
$t(s^{-1}(\{u\})) = \{u_1, \ldots, u_n\}$ est un ensemble fini.
Notons que $u \in \{u_1, \ldots, u_n\}$.
Soit $W \subset U$ un ouvert affine contenant $\{u_1, \ldots, u_n\}$ ;
en particulier, $u \in W$. Considérons
$Z = R \setminus s^{-1}(W) \cap t^{-1}(W)$. C'est un fermé
de $R$. L'image $t(Z)$ est un fermé de $U$ que l'on peut décrire
informellement comme l'ensemble des points de $U$ qui sont $R$-équivalents à un point
de $U \setminus W$. Ainsi, $W' = U \setminus t(Z)$ est un sous-schéma ouvert $R$-invariant
de $U$ contenu dans $W$, et $\{u_1, \ldots, u_n\} \subset W'$.
La situation est la suivante :
$$
\{u_1, \ldots, u_n\} \subset W' \subset W \subset U.
$$
Soit $f \in \Gamma(W, \mathcal{O}_W)$ un élément tel que
$\{u_1, \ldots, u_n\} \subset D(f) \subset W'$. Un tel $f$ existe d'après
Algèbre, lemme \ref{algebra-lemma-silly}. Par notre choix de $W'$,
on a $s^{-1}(W') \subset t^{-1}(W)$, et l'on obtient donc un diagramme
$$
\xymatrix{
s^{-1}(W') \ar[d]_s \ar[r]_-t & W \\
W'
}
$$
La flèche verticale est finie localement libre par hypothèse. Posons
$$
g = \text{Norm}_s(t^\sharp f) \in \Gamma(W', \mathcal{O}_{W'})
$$
Par construction, $g$ est une fonction sur $W'$ qui est
non nulle en $u$, car $t^\sharp(f)$ est non nulle en chacun des points de
$R$ au-dessus de $u$, puisque $f$ est non nulle en $u_1, \ldots, u_n$.
De même, $D(g) \subset W'$ est égal à
l'ensemble des points $w$ tels que $f$ ne s'annule en aucun des points
équivalents à $w$. Cela signifie que $D(g)$ est un
ouvert affine $R$-invariant de $W'$. La situation finale est
$$
\{u_1, \ldots, u_n\} \subset D(g) \subset D(f) \subset W' \subset W \subset U
$$
et l'on obtient donc le résultat souhaité.
\end{proof}







\section{Descente des schémas quasi-projectifs}
\label{groupoids-section-quasi-projective}

\noindent
On peut utiliser le lemme \ref{groupoids-lemma-find-invariant-affine}
pour montrer qu'un certain type de donnée de descente est effectif.

\begin{lemma}
\label{groupoids-lemma-descend-along-finite-with-bound}
Soit $X \to Y$ un morphisme surjectif fini localement libre.
Soit $d$ un entier strictement positif. Supposons que, pour tout point géométrique
$\bar{y} : \mathrm{Spec}(k) \to Y$, la fibre $X_{\bar{y}}$
ait au plus $d$ points.
Soit $V$ un schéma sur $X$ tel que, pour tout
$(y, v_1, \ldots, v_d)$ où $y \in Y$ et
$v_1, \ldots, v_d \in V_y$, il existe un ouvert affine
$U \subset V$ tel que $v_1, \ldots, v_d \in U$.
Alors toute donnée de descente sur $V/X/Y$ est effective.
\end{lemma}

\begin{proof}
Soit $\varphi$ une donnée de descente au sens de
Descente, définition \ref{descent-definition-descent-datum}.
Rappelons que le foncteur des schémas sur $Y$ vers les données de descente
relatives à $\{X \to Y\}$ est pleinement fidèle, voir
Descente, lemme \ref{descent-lemma-refine-coverings-fully-faithful}.
Ainsi, en utilisant Constructions, lemme \ref{constructions-lemma-relative-glueing},
il suffit de démontrer le lemme dans le cas où $Y$ est affine.
Certains détails sont omis (on peut éviter cet argument si $Y$ est
séparé ou possède une diagonale affine, car tout morphisme d'un
schéma affine vers $X$ est alors affine).

\medskip\noindent
Supposons que $Y$ soit affine. Si $V$ est lui aussi affine, on a l'effectivité d'après
Descente, lemme \ref{descent-lemma-affine}. Ainsi, d'après
Descente, lemme \ref{descent-lemma-effective-for-fpqc-is-local-upstairs},
il suffit de démontrer que tout point $v$ de $V$ possède un voisinage ouvert affine $\varphi$-invariant.
Considérons le groupoïde
$(X, X \times_Y X, \text{pr}_1, \text{pr}_0, \text{pr}_{02})$.
D'après le lemme \ref{groupoids-lemma-cartesian-equivalent-descent-datum},
la donnée de descente $\varphi$ détermine, et est déterminée par,
un morphisme cartésien de schémas en groupoïdes
$$
(V, R, s, t, c)
\longrightarrow
(X, X \times_Y X, \text{pr}_1, \text{pr}_0, \text{pr}_{02})
$$
sur $\Spec(\mathbf{Z})$.
Puisque $X \to Y$ est fini localement libre, on voit que
$\text{pr}_i : X \times_Y X \to X$, et donc $s$ et $t$,
sont finis localement libres, et leurs fibres géométriques ont au plus $d$
points. En particulier, la $R$-orbite $t(s^{-1}(\{v\}))$ de notre point
$v \in V$ a au plus $d$ points. En utilisant encore une fois l'équivalence de catégories du
lemme \ref{groupoids-lemma-cartesian-equivalent-descent-datum},
on voit que les ouverts $\varphi$-invariants de $V$
sont exactement les ouverts $R$-invariants de $V$.
Notre hypothèse montre qu'il existe un ouvert affine de $V$
contenant l'orbite $t(s^{-1}(\{v\}))$, car tous les points
de cette orbite ont la même image dans $Y$.
Le lemme \ref{groupoids-lemma-find-invariant-affine}
fournit donc un ouvert affine $R$-invariant contenant $v$.
\end{proof}

\begin{lemma}
\label{groupoids-lemma-descend-along-finite}
Soit $X \to Y$ un morphisme surjectif fini localement libre.
Soit $V$ un schéma sur $X$ tel que, pour tout
$(y, v_1, \ldots, v_d)$ où $y \in Y$ et
$v_1, \ldots, v_d \in V_y$, il existe un ouvert affine
$U \subset V$ tel que $v_1, \ldots, v_d \in U$.
Alors toute donnée de descente sur $V/X/Y$ est effective.
\end{lemma}

\begin{proof}
Comme dans la démonstration du lemme \ref{groupoids-lemma-descend-along-finite-with-bound}, on
peut supposer que $Y$ est affine, donc en particulier quasi-compact. Il
existe alors un entier $d$ tel que toutes les fibres géométriques de $X \to Y$
aient au plus $d$ points, voir la discussion dans Morphismes, section
\ref{morphisms-section-universally-bounded}.
Notre hypothèse sur $V$ implique trivialement que l'on peut appliquer
le lemme \ref{groupoids-lemma-descend-along-finite-with-bound}, ce qui termine la démonstration.
\end{proof}

\begin{lemma}
\label{groupoids-lemma-descend-along-finite-quasi-projective}
Soit $X \to Y$ un morphisme surjectif fini localement libre.
Soit $V$ un schéma sur $X$ tel que l'une des conditions suivantes soit satisfaite :
\begin{enumerate}
\item $V \to X$ est projectif,
\item $V \to X$ est quasi-projectif,
\item il existe un faisceau inversible ample sur $V$,
\item il existe un faisceau inversible $X$-ample sur $V$,
\item il existe un faisceau inversible $X$-très ample sur $V$.
\end{enumerate}
Alors toute donnée de descente sur $V/X/Y$ est effective.
\end{lemma}

\begin{proof}
Vérifions la condition du lemme \ref{groupoids-lemma-descend-along-finite}.
Soient $y \in Y$ et $v_1, \ldots, v_d \in V$ des points au-dessus de $y$.
Le cas (1) est un cas particulier de (2), voir
Morphismes, lemme \ref{morphisms-lemma-projective-quasi-projective}.
Le cas (2) est un cas particulier de (4), voir
Morphismes, définition \ref{morphisms-definition-quasi-projective}.
S'il existe un faisceau inversible ample sur $V$, alors
il existe un ouvert affine contenant $v_1, \ldots, v_d$ d'après
Propriétés, lemme \ref{properties-lemma-ample-finite-set-in-affine}.
Ainsi, la condition (3) est satisfaite.
Dans les cas (4) et (5), on peut sans inconvénient remplacer $Y$ par un
voisinage ouvert affine de $y$.
Alors $X$ est lui aussi affine.
Dans le cas (4), on voit que $V$ possède un faisceau inversible ample
d'après Morphismes, définition \ref{morphisms-definition-relatively-ample},
et le résultat découle du cas (3).
Dans le cas (5), on peut remplacer $V$ par un ouvert quasi-compact contenant
$v_1, \ldots, v_d$, et l'on se ramène au cas (4) par
Morphismes, lemme \ref{morphisms-lemma-ample-very-ample}.
\end{proof}

\begin{lemma}
\label{groupoids-lemma-descend-along-finite-radicial}
Soit $X \to Y$ un morphisme surjectif fini localement libre qui est radiciel.
Soit $V$ un schéma sur $X$. Alors toute donnée de descente sur $V/X/Y$ est effective.
\end{lemma}

\begin{proof}
Appliquer le lemme \ref{groupoids-lemma-descend-along-finite-with-bound} avec $d = 1$.
\end{proof}












% OK, start here.
%

\chapter{Compléments sur les schémas en groupoïdes}



\phantomsection
\label{more-groupoids-section-phantom}


\section{Introduction}
\label{more-groupoids-section-introduction}

\noindent
Ce chapitre est consacré à des sujets avancés sur les schémas en groupoïdes.
Bien que les résultats soient énoncés en termes de schémas en groupoïdes, le
lecteur doit garder à l'esprit le diagramme $2$-cartésien
\begin{equation}
\label{more-groupoids-equation-quotient-stack}
\vcenter{
\xymatrix{
R \ar[r] \ar[d] & U \ar[d] \\
U \ar[r] & [U/R]
}
}
\end{equation}
où $[U/R]$ est le champ quotient, voir
Groupoïdes dans les espaces, Remarque \ref{spaces-groupoids-remark-fundamental-square}.
Nombre de résultats sont motivés par la considération de ce diagramme.
Voir par exemple le très bel article \cite{K-M} de Keel et Mori.





\section{Notations}
\label{more-groupoids-section-notation}

\noindent
Nous continuons à suivre les conventions et les notations introduites dans
Groupoïdes, section \ref{groupoids-section-notation}.






\section{Diagrammes utiles}
\label{more-groupoids-section-diagrams}

\noindent
Nous rappelons brièvement les résultats des
Lemmes \ref{groupoids-lemma-diagram} et
\ref{groupoids-lemma-diagram-pull} de Groupoïdes,
afin de pouvoir nous y référer aisément dans ce chapitre.
Soit $S$ un schéma.
Soit $(U, R, s, t, c)$ un schéma en groupoïdes sur $S$.
Dans le diagramme commutatif
\begin{equation}
\label{more-groupoids-equation-diagram}
\vcenter{
\xymatrix{
& U & \\
R \ar[d]_s \ar[ru]^t &
R \times_{s, U, t} R
\ar[l]^-{\text{pr}_0} \ar[d]^{\text{pr}_1} \ar[r]_-c &
R \ar[d]^s \ar[lu]_t \\
U & R \ar[l]_t \ar[r]^s & U
}
}
\end{equation}
les deux carrés inférieurs sont des carrés de produit fibré.
De plus, le triangle supérieur (qui est en réalité un carré)
est lui aussi cartésien.

\medskip\noindent
Le diagramme
\begin{equation}
\label{more-groupoids-equation-pull}
\vcenter{
\xymatrix{
R \times_{t, U, t} R
\ar@<1ex>[r]^-{\text{pr}_1} \ar@<-1ex>[r]_-{\text{pr}_0}
\ar[d]_{\text{pr}_0 \times c \circ (i, 1)} &
R \ar[r]^t \ar[d]^{\text{id}_R} &
U \ar[d]^{\text{id}_U} \\
R \times_{s, U, t} R
\ar@<1ex>[r]^-c \ar@<-1ex>[r]_-{\text{pr}_0} \ar[d]_{\text{pr}_1} &
R \ar[r]^t \ar[d]^s &
U \\
R \ar@<1ex>[r]^s \ar@<-1ex>[r]_t &
U
}
}
\end{equation}
est commutatif. Les deux lignes supérieures sont isomorphes par les applications verticales indiquées.
Les deux carrés inférieurs de gauche sont cartésiens.





\section{Faisceau des différentielles}
\label{more-groupoids-section-differentials}

\noindent
Le lemme suivant est l'analogue du
Lemme \ref{groupoids-lemma-group-scheme-module-differentials} de Groupoïdes.

\begin{lemma}
\label{more-groupoids-lemma-sheaf-differentials}
Soit $S$ un schéma.
Soit $(U, R, s, t, c)$ un schéma en groupoïdes sur $S$.
Le faisceau des différentielles de $R$, considéré comme un schéma sur
$U$ par $t$, est un quotient de l'image réciproque par $t$ du faisceau conormal de
l'immersion $e : U \to R$. En formule : il existe une surjection canonique
$t^*\mathcal{C}_{U/R} \to \Omega_{R/U}$. Si $s$ est plat, alors
cette application est un isomorphisme.
\end{lemma}

\begin{proof}
Notons que $e : U \to R$ est une immersion, puisqu'il s'agit d'une section
du morphisme $s$, voir
Schémas, Lemme \ref{schemes-lemma-section-immersion}.
Considérons le diagramme suivant
$$
\xymatrix{
R \ar[r]_-{(1, i)} \ar[d]_t &
R \times_{s, U, t} R \ar[d]^c \ar[rr]_{(\text{pr}_0, i \circ \text{pr}_1)} & &
R \times_{t, U, t} R \\
U \ar[r]^e &
R
}
$$
Le carré de gauche est cartésien, car si $a \circ b = e$, alors
$b = i(a)$. La composée des applications horizontales est le
morphisme diagonal de $t : R \to U$. La flèche horizontale supérieure droite est un
isomorphisme. Puisque $\Omega_{R/U}$ est le faisceau conormal de la
composée, il est isomorphe au faisceau conormal de
$(1, i)$. Par le
Lemme \ref{morphisms-lemma-conormal-functorial-flat} de Morphismes,
nous obtenons la surjection $t^*\mathcal{C}_{U/R} \to \Omega_{R/U}$,
et, si $c$ est plat, il s'agit d'un isomorphisme. Puisque $c$ est un changement de base
de $s$ d'après les propriétés du Diagramme (\ref{more-groupoids-equation-pull}),
nous concluons que, si $s$ est plat, alors $c$ est plat, voir
Morphismes, Lemme \ref{morphisms-lemma-base-change-flat}.
\end{proof}






\section{Structure locale}
\label{more-groupoids-section-local}

\noindent
Soit $S$ un schéma.
Soit $(U, R, s, t, c, e, i)$ un schéma en groupoïdes sur $S$.
Soit $u \in U$ un point. Dans cette section, nous expliquons quelle
structure nous obtenons sur les anneaux locaux
$$
A = \mathcal{O}_{U, u}
\quad\text{et}\quad
B = \mathcal{O}_{R, e(u)}
$$
Nous conviendrons de désigner les homomorphismes d'anneaux locaux
induits par les morphismes $s, t, c, e, i$ par les lettres correspondantes.
En particulier, nous avons un diagramme commutatif
$$
\xymatrix{
A \ar[rd]_t \ar[rrd]^1 \\
& B \ar[r]^e & A \\
A \ar[ru]^s \ar[rru]_1
}
$$
d'anneaux locaux. Ainsi, si $I \subset B$ désigne le noyau de $e : B \to A$,
alors $B = s(A) \oplus I = t(A) \oplus I$. Notons
$$
C = \mathcal{O}_{R \times_{s, U, t} R, (e(u), e(u))}
$$
Nous avons alors
$$
C = (B \otimes_{s, A, t} B)_{\mathfrak m_B \otimes B + B \otimes \mathfrak m_B}
$$
Soit $J \subset C$ l'idéal de $C$ engendré par $I \otimes B + B \otimes I$.
Alors $J$ est également le noyau de l'homomorphisme d'anneaux locaux
$$
(e, e) : C \longrightarrow A
$$
La loi de composition $c : R \times_{s, U, t} R \to R$ correspond à un
homomorphisme d'anneaux
$$
c : B \longrightarrow C
$$
envoyant $I$ dans $J$.

\begin{lemma}
\label{more-groupoids-lemma-first-order-structure-c}
L'application $I/I^2 \to J/J^2$ induite par $c$ est la composée
$$
I/I^2 \xrightarrow{(1, 1)} I/I^2 \oplus I/I^2 \to J/J^2
$$
où la seconde flèche provient de l'égalité
$J = (I \otimes B + B \otimes I)C$.
L'application $i : B \to B$ induit l'application $-1 : I/I^2 \to I/I^2$.
\end{lemma}

\begin{proof}
Pour décrire un homomorphisme local de $C$ vers un autre anneau local,
il suffit d'indiquer ce qu'il advient des éléments de la forme
$b_1 \otimes b_2$. En gardant cela à l'esprit, nous avons les deux applications canoniques
$$
e_2 : C \to B,\ b_1 \otimes b_2 \mapsto b_1s(e(b_2)),\quad
e_1 : C \to B,\ b_1 \otimes b_2 \mapsto t(e(b_1))b_2
$$
correspondant aux applications
$R \to R \times_{s, U, t} R$ données par
$r \mapsto (r, e(s(r)))$ et $r \mapsto (e(t(r)), r)$.
Ces applications définissent des applications $J/J^2 \to I/I^2$ qui, ensemble,
donnent un inverse de l'application $I/I^2 \oplus I/I^2 \to J/J^2$
du lemme. Ainsi, pour démontrer l'assertion, il nous suffit de montrer
que $e_1 \circ c : B \to B$ et $e_2 \circ c : B \to B$
sont les applications identiques. Cela résulte du fait que les deux
composées $R \to R \times_{s, U, t} R \to R$ sont des identités.

\medskip\noindent
L'assertion concernant $i$ résulte de celle concernant $c$ et du
fait que $c \circ (1, i) = e \circ t$. Nous omettons quelques détails.
\end{proof}






\section{Propriétés des groupoïdes}
\label{more-groupoids-section-technical-lemma}

\noindent
Soit $(U, R, s, t, c)$ un schéma en groupoïdes.
L'idée qui sous-tend les résultats de cette section est que $s: R \to U$
est un changement de base du morphisme $U \to [U/R]$ (voir
le Diagramme (\ref{more-groupoids-equation-quotient-stack})).
Les propriétés locales de $s : R \to U$ devraient donc refléter les
propriétés locales du morphisme $U \to [U/R]$.
Cela ne fonctionne pas, car $[U/R]$ n'est pas toujours un champ algébrique et
nous ne pouvons donc pas parler des propriétés géométriques ou algébriques de
$U \to [U/R]$.
Il s'avère toutefois que nous pouvons en récupérer une partie sans même
faire référence au champ quotient.

\medskip\noindent
Voici un premier exemple d'un tel résultat. Grosso modo, l'ouvert $W \subset U'$
trouvé dans le lemme est le lieu où le morphisme
$U' \to [U/R]$ possède la propriété $\mathcal{P}$.

\begin{lemma}
\label{more-groupoids-lemma-local-source}
Soit $S$ un schéma.
Soit $(U, R, s, t, c, e, i)$ un groupoïde sur $S$.
Soit $g : U' \to U$ un morphisme de schémas.
Notons $h$ la composée
$$
\xymatrix{
h : U' \times_{g, U, t} R \ar[r]_-{\text{pr}_1} & R \ar[r]_s & U.
}
$$
Soient $\mathcal{P}, \mathcal{Q}, \mathcal{R}$ des propriétés de morphismes
de schémas. Supposons que
\begin{enumerate}
\item $\mathcal{R} \Rightarrow \mathcal{Q}$,
\item $\mathcal{Q}$ est préservée par changement de base et par composition,
\item pour tout morphisme $f : X \to Y$ qui possède $\mathcal{Q}$, il existe un
plus grand ouvert $W(\mathcal{P}, f) \subset X$ tel que $f|_{W(\mathcal{P}, f)}$
possède $\mathcal{P}$, et
\item pour tout morphisme $f : X \to Y$ qui possède $\mathcal{Q}$,
et tout morphisme $Y' \to Y$ qui possède $\mathcal{R}$, nous avons
$Y' \times_Y W(\mathcal{P}, f) = W(\mathcal{P}, f')$, où
$f' : X_{Y'} \to Y'$ est le changement de base de $f$.
\end{enumerate}
Si $s, t$ possèdent $\mathcal{R}$ et si $g$ possède $\mathcal{Q}$, alors
il existe un sous-schéma ouvert $W \subset U'$ tel que
$W \times_{g, U, t} R = W(\mathcal{P}, h)$.
\end{lemma}

\begin{proof}
Notons que le diagramme suivant est commutatif
$$
\xymatrix{
U' \times_{g, U, t} R \times_{t, U, t} R
\ar[rr]_-{\text{pr}_{12}}
\ar@<1ex>[d]^-{\text{pr}_{02}} \ar@<-1ex>[d]_-{\text{pr}_{01}} & &
R \times_{t, U, t} R
\ar@<1ex>[d]^-{\text{pr}_1} \ar@<-1ex>[d]_-{\text{pr}_0}
\\
U' \times_{g, U, t} R \ar[rr]^{\text{pr}_1} & & R
}
$$
dont les deux carrés sont cartésiens (on utilise ici que les deux applications
$t \circ \text{pr}_i : R \times_{t, U, t} R \to U$ sont égales).
En combinant ceci avec les propriétés du diagramme (\ref{more-groupoids-equation-pull}),
nous obtenons un diagramme commutatif
$$
\xymatrix{
U' \times_{g, U, t} R \times_{t, U, t} R
\ar[rr]_-{c \circ (i, 1)}
\ar@<1ex>[d]^-{\text{pr}_{02}} \ar@<-1ex>[d]_-{\text{pr}_{01}} & &
R
\ar@<1ex>[d]^-{s} \ar@<-1ex>[d]_-{t}
\\
U' \times_{g, U, t} R \ar[rr]^h & & U
}
$$
dont les deux carrés sont cartésiens.

\medskip\noindent
Supposons que $s, t$ possèdent $\mathcal{R}$ et que $g$ possède $\mathcal{Q}$.
Alors $h$ possède $\mathcal{Q}$, car il est la composée de $s$ (qui possède
$\mathcal{R}$, donc $\mathcal{Q}$) et d'un changement de base de $g$ (qui
possède $\mathcal{Q}$). Ainsi, $W(\mathcal{P}, h) \subset U' \times_{g, U, t} R$
existe. D'après nos hypothèses, nous avons
$\text{pr}_{01}^{-1}(W(\mathcal{P}, h)) =
\text{pr}_{02}^{-1}(W(\mathcal{P}, h))$
car ce sont tous deux le plus grand ouvert sur lequel $c \circ (i, 1)$ possède $\mathcal{P}$.
Notons que la projection $U' \times_{g, U, t} R \to U'$ possède une section, à savoir
$\sigma : U' \to U' \times_{g, U, t} R$, $u' \mapsto (u', e(g(u')))$.
De plus, par l'isomorphisme
$$
(U' \times_{g, U, t} R) \times_{U'} (U' \times_{g, U, t} R)
=
U' \times_{g, U, t} R \times_{t, U, t} R
$$
les deux projections du membre de gauche
vers $U' \times_{g, U, t} R$ coïncident avec les morphismes $\text{pr}_{01}$
et $\text{pr}_{02}$ du membre de droite. Puisque
$\text{pr}_{01}^{-1}(W(\mathcal{P}, h)) =
\text{pr}_{02}^{-1}(W(\mathcal{P}, h))$
nous concluons que $W(\mathcal{P}, h)$ est l'image réciproque d'un sous-ensemble de $U$,
qui est nécessairement l'ouvert
$W = \sigma^{-1}(W(\mathcal{P}, h))$.
\end{proof}

\begin{remark}
\label{more-groupoids-remark-local-source-warning}
Avertissement :
le Lemme \ref{more-groupoids-lemma-local-source}
doit être utilisé avec précaution.
Il s'applique par exemple à $\mathcal{P}=$``plat'', $\mathcal{Q}=$``sans condition'',
et $\mathcal{R}=$``plat et localement de présentation finie''. Toutefois, pour un
morphisme de schémas $f : X \to Y$, le plus grand ouvert $W \subset X$ tel que
$f|_W$ soit plat n'est {\it pas} l'ensemble des points où $f$ est plat !
\end{remark}

\begin{remark}
\label{more-groupoids-remark-local-source-apply}
Malgré l'avertissement de la
Remarque \ref{more-groupoids-remark-local-source-warning},
il existe certains cas où le
Lemme \ref{more-groupoids-lemma-local-source}
peut être utilisé sans introduire trop d'ambiguïté.
Nous en donnons une liste. Dans chaque cas, nous omettons la vérification des
hypothèses (1) et (2) et donnons des références qui impliquent
(3) et (4). Voici la liste :
\begin{enumerate}
\item $\mathcal{Q} = \mathcal{R} =$``localement de type fini'', et
$\mathcal{P} =$``de dimension relative $\leq d$''.
Voir
Morphismes, Définition \ref{morphisms-definition-relative-dimension-d}
et
Morphismes, Lemmes \ref{morphisms-lemma-openness-bounded-dimension-fibres} et
\ref{morphisms-lemma-dimension-fibre-after-base-change}.
\item $\mathcal{Q} = \mathcal{R} =$``localement de type fini'', et
$\mathcal{P} =$``localement quasi-fini''.
C'est le cas $d = 0$ du point précédent, voir
Morphismes, Lemme \ref{morphisms-lemma-locally-quasi-finite-rel-dimension-0}.
\item $\mathcal{Q} = \mathcal{R} =$``localement de type fini'', et
$\mathcal{P} =$``non ramifié''.
Voir
Morphismes, Lemmes \ref{morphisms-lemma-unramified-characterize} et
\ref{morphisms-lemma-set-points-where-fibres-unramified}.
\end{enumerate}
L'intérêt des cas énumérés ci-dessus est qu'il n'est pas nécessaire de
supposer que $s, t$ sont plats pour obtenir une conclusion sur le lieu
où le morphisme $h$ possède la propriété $\mathcal{P}$. Poursuivons la
liste :
\begin{enumerate}
\item[(4)] $\mathcal{Q} =$``localement de présentation finie'',
$\mathcal{R} =$``plat et localement de présentation finie'', et
$\mathcal{P} =$``plat''. Voir
Compléments sur les morphismes, Théorème
\ref{more-morphisms-theorem-openness-flatness} et
Lemme \ref{more-morphisms-lemma-flat-locus-base-change}.
\item[(5)] $\mathcal{Q} =$``localement de présentation finie'',
$\mathcal{R} =$``plat et localement de présentation finie'', et
$\mathcal{P}=$``de Cohen-Macaulay''. Voir
Compléments sur les morphismes, Définition \ref{more-morphisms-definition-CM}
et
Compléments sur les morphismes, Lemmes \ref{more-morphisms-lemma-base-change-CM} et
\ref{more-morphisms-lemma-flat-finite-presentation-CM-open}.
\item[(6)] $\mathcal{Q} =$``localement de présentation finie'',
$\mathcal{R} =$``plat et localement de présentation finie'', et
$\mathcal{P}=$``syntomique'' ; voir
Morphismes, Lemme \ref{morphisms-lemma-set-points-where-fibres-lci}
(le lieu est automatiquement ouvert).
\item[(7)] $\mathcal{Q} =$``localement de présentation finie'',
$\mathcal{R} =$``plat et localement de présentation finie'', et
$\mathcal{P}=$``lisse''. Voir
Morphismes, Lemme \ref{morphisms-lemma-set-points-where-fibres-smooth}
(le lieu est automatiquement ouvert).
\item[(8)] $\mathcal{Q} =$``localement de présentation finie'',
$\mathcal{R} =$``plat et localement de présentation finie'', et
$\mathcal{P}=$``\'etale''. Voir
Morphismes, Lemme \ref{morphisms-lemma-set-points-where-fibres-etale}
(le lieu est automatiquement ouvert).
\end{enumerate}
\end{remark}

\noindent
Voici le second résultat. Il faut penser à l'ouvert $W \subset U$ invariant par $R$ comme à
l'image réciproque du plus grand ouvert de $[U/R]$ au-dessus duquel
le morphisme $U \to [U/R]$ possède la propriété $\mathcal{P}$.

\begin{lemma}
\label{more-groupoids-lemma-property-invariant}
Soit $S$ un schéma.
Soit $(U, R, s, t, c)$ un groupoïde sur $S$.
Soit $\tau \in \{Zariski, \linebreak[0] fppf,
\linebreak[0] \etale, \linebreak[0]
smooth, \linebreak[0] syntomic\}$\footnote{Le fait que $fpqc$ soit absent
n'est pas une coquille.}. Soit $\mathcal{P}$ une propriété de morphismes de schémas
qui est $\tau$-locale sur la base
(Descente, Définition \ref{descent-definition-property-morphisms-local}).
Supposons que $\{s : R \to U\}$ et $\{t : R \to U\}$ soient des recouvrements pour la
topologie $\tau$. Soit $W \subset U$ le sous-schéma ouvert maximal tel que
$s|_{s^{-1}(W)} : s^{-1}(W) \to W$ possède la propriété $\mathcal{P}$.
Alors $W$ est invariant par $R$, voir
Groupoïdes, Définition \ref{groupoids-definition-invariant-open}.
\end{lemma}

\begin{proof}
L'existence et les propriétés de l'ouvert $W \subset U$ sont décrites dans
Descente, Lemme \ref{descent-lemma-largest-open-of-the-base}.
Dans le
Diagramme (\ref{more-groupoids-equation-diagram}),
soit $W_1 \subset R$ le sous-schéma ouvert maximal au-dessus duquel le morphisme
$\text{pr}_1 : R \times_{s, U, t} R \to R$ possède la propriété $\mathcal{P}$.
Il résulte du
Lemme \ref{descent-lemma-largest-open-of-the-base} de Descente déjà cité
et de l'hypothèse que $\{s : R \to U\}$ et $\{t : R \to U\}$ sont des recouvrements
pour la topologie $\tau$ que $t^{-1}(W) = W_1 = s^{-1}(W)$, comme souhaité.
\end{proof}

\begin{lemma}
\label{more-groupoids-lemma-property-G-invariant}
Soit $S$ un schéma.
Soit $(U, R, s, t, c)$ un groupoïde sur $S$.
Soit $G \to U$ son schéma en groupes stabilisateur.
Soit $\tau \in \{fppf, \linebreak[0] \etale, \linebreak[0]
smooth, \linebreak[0] syntomic\}$.
Soit $\mathcal{P}$ une propriété de morphismes qui est $\tau$-locale
sur la base. Supposons que $\{s : R \to U\}$ et $\{t : R \to U\}$ soient des recouvrements
pour la topologie $\tau$. Soit $W \subset U$ le sous-schéma ouvert maximal
tel que $G_W \to W$ possède la propriété $\mathcal{P}$. Alors $W$ est invariant par $R$
(voir
Groupoïdes, Définition
\ref{groupoids-definition-invariant-open}).
\end{lemma}

\begin{proof}
L'existence et les propriétés de l'ouvert $W \subset U$ sont décrites dans
Descente, Lemme \ref{descent-lemma-largest-open-of-the-base}.
Le morphisme
$$
G \times_{U, t} R \longrightarrow R \times_{s, U} G, \quad
(g, r) \longmapsto (r, r^{-1} \circ g \circ r)
$$
est un isomorphisme au-dessus de $R$ (où $\circ$ désigne
la composition dans le groupoïde). Ainsi, $s^{-1}(W) = t^{-1}(W)$ par les
propriétés de $W$ démontrées dans le
Lemme \ref{descent-lemma-largest-open-of-the-base} de Descente déjà cité.
\end{proof}



\section{Comparaison des fibres}
\label{more-groupoids-section-fibres}

\noindent
Soit $(U, R, s, t, c, e, i)$ un schéma en groupoïdes sur $S$.
Le Diagramme (\ref{more-groupoids-equation-diagram})
nous fournit un moyen de comparer les fibres de l'application $s : R \to U$ dans un groupoïde.
Pour un point $u \in U$, nous noterons $F_u = s^{-1}(u)$ la fibre
schématique de $s : R \to U$ au-dessus de $u$. Par exemple, le diagramme
implique que, si $u, u' \in U$ sont des points tels
que $s(r) = u$ et $t(r) = u'$, alors
$(F_u)_{\kappa(r)} \cong (F_{u'})_{\kappa(r)}$.
C'est un cas particulier du résultat plus général et plus précis du
Lemme \ref{more-groupoids-lemma-two-fibres}
ci-dessous. Pour le voir, prenons $r' = i(r)$.

\medskip\noindent
Un couple $(X, x)$ formé d'un schéma $X$ et d'un point $x \in X$ est parfois
appelé le {\it germe de $X$ en $x$}.
Un {\it morphisme de germes} $f : (X, x) \to (S, s)$
est un morphisme $f : U \to S$ défini sur un voisinage ouvert
de $x$ et tel que $f(x) = s$. Deux tels
$f$, $f'$ sont dits définir le même morphisme de germes
si et seulement si $f$ et $f'$ coïncident sur un voisinage ouvert de $x$.
Soit $\tau \in \{Zariski, \etale, smooth, syntomic, fppf\}$.
Introduisons provisoirement la notion suivante : nous disons que deux morphismes
de germes $f : (X, x) \to (S, s)$ et $f' : (X', x') \to (S', s')$
sont {\it localement isomorphes sur la base pour la topologie $\tau$},
s'il existe un schéma pointé $(S'', s'')$ et des morphismes de germes
$g : (S'', s'') \to (S, s)$ et $g' : (S'', s'') \to (S', s')$
tels que
\begin{enumerate}
\item $g$ et $g'$ sont des immersions ouvertes (resp.\ \'etales, lisses, syntomiques,
plates et localement de présentation finie) en $s''$,
\item il existe un isomorphisme
$$
(S'' \times_{g, S, f} X, \tilde x)
\cong
(S'' \times_{g', S', f'} X', \tilde  x')
$$
de germes au-dessus du germe $(S'', s'')$, pour un certain choix de points
$\tilde x$ et $\tilde x'$ situés au-dessus de $(s'', x)$ et $(s'', x')$.
\end{enumerate}
Enfin, nous disons simplement que les morphismes de germes
$f : (X, x) \to (S, s)$ et $f' : (X', x') \to (S', s')$
sont {\it isomorphes après des changements de base plats} s'il existe
$S'', s'', g, g'$ comme ci-dessus, mais en remplaçant (1) par
la condition que $g$ et $g'$ soient plats en $s''$ (celle-ci est
bien plus faible que toutes les conditions $\tau$ ci-dessus,
car un morphisme plat n'est pas nécessairement ouvert).

\begin{lemma}
\label{more-groupoids-lemma-two-fibres}
Soit $S$ un schéma.
Soit $(U, R, s, t, c)$ un groupoïde sur $S$.
Soient $r, r' \in R$ tels que $t(r) = t(r')$ dans $U$.
Posons $u = s(r)$, $u' = s(r')$.
Notons $F_u = s^{-1}(u)$ et $F_{u'} = s^{-1}(u')$ les fibres
schématiques.
\begin{enumerate}
\item Il existe une extension de corps commune
$\kappa(u) \subset k$, $\kappa(u') \subset k$ et
un isomorphisme $(F_u)_k \cong (F_{u'})_k$.
\item On peut choisir l'isomorphisme de (1) de telle sorte qu'un point
situé au-dessus de $r$ s'envoie sur un point situé au-dessus de $r'$.
\item Si les morphismes $s$, $t$ sont plats, alors les morphismes de germes
$s : (R, r) \to (U, u)$ et $s : (R, r') \to (U, u')$ sont
localement isomorphes sur la base pour la topologie plate.
\item Si les morphismes $s$, $t$ sont \'etales
(resp.\ lisses, syntomiques, ou plats et localement de présentation finie),
alors les morphismes de germes $s : (R, r) \to (U, u)$ et
$s : (R, r') \to (U, u')$ sont localement isomorphes sur la base
pour la topologie \'etale (resp.\ lisse, syntomique, ou fppf).
\end{enumerate}
\end{lemma}

\begin{proof}
Nous utilisons à plusieurs reprises les propriétés et l'existence du
diagramme (\ref{more-groupoids-equation-diagram}).
D'après les propriétés du diagramme (et
Schémas, lemme \ref{schemes-lemma-points-fibre-product}),
il existe un point $\xi$ de $R \times_{s, U, t} R$
tel que $\text{pr}_0(\xi) = r$ et $c(\xi) = r'$.
Posons $\tilde r = \text{pr}_1(\xi) \in R$.

\medskip\noindent
Démonstration de (1). Posons $k = \kappa(\tilde r)$. Puisque $t(\tilde r) = u$
et $s(\tilde r) = u'$, nous voyons que $k$ est une extension commune
de $\kappa(u)$ et de $\kappa(u')$ et, en fait, que
$(F_u)_k$ et $(F_{u'})_k$ sont tous deux isomorphes à la fibre de
$\text{pr}_1 : R \times_{s, U, t} R \to R$ au-dessus de $\tilde r$.
Ainsi, (1) est démontré.

\medskip\noindent
L'assertion (2) en résulte, puisque le point $\xi$ s'envoie sur $r$, resp.\ $r'$.

\medskip\noindent
L'assertion (3) découle clairement de ce qui précède (en utilisant le point $\xi$ pour
$\tilde u$ et $\tilde u'$) et des définitions.

\medskip\noindent
Si $s$ et $t$ sont plats et de présentation finie, alors
ce sont des morphismes ouverts (Morphismes, lemme \ref{morphisms-lemma-fppf-open}).
L'image d'un certain voisinage ouvert affine $V''$ de $\tilde r$ recouvrira donc
un voisinage ouvert $V$ de $u$, resp.\ $V'$ de $u'$.
On peut les utiliser pour montrer que les propriétés (1) et (2) de la
définition de ``localement isomorphes sur la base pour la
topologie $\tau$''.
\end{proof}







\section{Présentations de Cohen-Macaulay}
\label{more-groupoids-section-CM}

\noindent
Étant donné un groupoïde $(U, R, s, t, c)$ avec $s, t$ plats et
localement de présentation finie, il existe un groupoïde ``équivalent''
$(U', R', s', t', c')$ tel que $s'$ et $t'$ soient des morphismes
de Cohen-Macaulay (et localement de présentation finie). Voir
Compléments sur les morphismes, section \ref{more-morphisms-section-CM},
pour plus de renseignements sur les morphismes de Cohen-Macaulay.
Ici, ``équivalent'' peut signifier que les champs quotients
$[U/R]$ et $[U'/R']$ sont des champs équivalents, voir
Groupoïdes dans les espaces, section \ref{spaces-groupoids-section-stacks}
et la section \ref{spaces-groupoids-section-quotient-stack-restrict}.

\begin{lemma}
\label{more-groupoids-lemma-make-CM}
Soit $S$ un schéma.
Soit $(U, R, s, t, c)$ un groupoïde au-dessus de $S$.
Supposons $s$ et $t$ plats et localement de présentation finie.
Alors il existe un ouvert $U' \subset U$ tel que
\begin{enumerate}
\item $t^{-1}(U') \subset R$ est le plus grand sous-schéma ouvert de
$R$ sur lequel le morphisme $s$ est de Cohen-Macaulay,
\item $s^{-1}(U') \subset R$ est le plus grand sous-schéma ouvert de
$R$ sur lequel le morphisme $t$ est de Cohen-Macaulay,
\item le morphisme $t|_{s^{-1}(U')} : s^{-1}(U') \to U$ est
surjectif,
\item le morphisme $s|_{t^{-1}(U')} : t^{-1}(U') \to U$ est
surjectif, et
\item la restriction $R' = s^{-1}(U') \cap t^{-1}(U')$
de $R$ à $U'$ définit un groupoïde $(U', R', s', t', c')$ qui a la propriété
que les morphismes $s'$ et $t'$ sont de Cohen-Macaulay et localement de
présentation finie.
\end{enumerate}
\end{lemma}

\begin{proof}
Appliquons le
lemme \ref{more-groupoids-lemma-local-source}
avec
$g = \text{id}$ et
$\mathcal{Q} =$``localement de présentation finie'',
$\mathcal{R} =$``plat et localement de présentation finie'', et
$\mathcal{P}=$``de Cohen-Macaulay'', voir
la remarque \ref{more-groupoids-remark-local-source-apply}.
Cela nous donne un ouvert $U' \subset U$ tel que
Alors $t^{-1}(U') \subset R$ est le plus grand sous-schéma ouvert de $R$
sur lequel le morphisme $s$ est de Cohen-Macaulay.
Ceci prouve (1).
Soit $i : R \to R$ l'inverse du groupoïde.
Puisque $i$ est un isomorphisme, et $s \circ i = t$ et $t \circ i = s$,
nous voyons que $s^{-1}(U')$ est aussi le plus grand ouvert de $R$ sur lequel $t$ est
de Cohen-Macaulay. Ceci prouve (2).
D'après
Compléments sur les morphismes,
lemme \ref{more-morphisms-lemma-flat-finite-presentation-CM-open},
l'ouvert $t^{-1}(U')$ est dense dans chaque fibre de $s : R \to U$.
Ceci prouve (3). Même argument pour (4).
L'assertion (5) est une conséquence formelle de (1) et (2) et de la discussion
sur les restrictions dans
Groupoïdes, section \ref{groupoids-section-restrict-groupoid}.
\end{proof}








\section{Restriction des groupoïdes}
\label{more-groupoids-section-restricting-groupoids}

\noindent
Dans cette section, nous rassemblons plusieurs lemmes sur les
propriétés des groupoïdes qui se transmettent aux restrictions.
La plupart de ces lemmes se démontrent en considérant le
diagramme qui définit
\begin{equation}
\label{more-groupoids-equation-restriction}
\vcenter{
\xymatrix{
R' \ar[d] \ar[r] \ar@/_3pc/[dd]_{t'} \ar@/^1pc/[rr]^{s'}&
R \times_{s, U} U' \ar[r] \ar[d] &
U' \ar[d]^g \\
U' \times_{U, t} R \ar[d] \ar[r] &
R \ar[r]^s \ar[d]_t &
U \\
U' \ar[r]^g &
U
}
}
\end{equation}
une restriction. Voir
Groupoïdes, lemme \ref{groupoids-lemma-restrict-groupoid}.

\begin{lemma}
\label{more-groupoids-lemma-restrict-preserves-type}
Soit $S$ un schéma.
Soit $(U, R, s, t, c)$ un schéma en groupoïdes au-dessus de $S$.
Soit $g : U' \to U$ un morphisme de schémas.
Soit $(U', R', s', t', c')$ la restriction de
$(U, R, s, t, c)$ par $g$.
\begin{enumerate}
\item Si $s, t$ sont localement de type fini et $g$ est localement de
type fini, alors $s', t'$ sont localement de type fini.
\item Si $s, t$ sont localement de présentation finie et $g$ est localement de
présentation finie, alors $s', t'$ sont localement de présentation finie.
\item Si $s, t$ sont plats et $g$ est plat, alors $s', t'$ sont plats.
\item Ajouter ici.
\end{enumerate}
\end{lemma}

\begin{proof}
La propriété d'être localement de type fini est stable par composition
et par changement de base arbitraire, voir
Morphismes, lemmes \ref{morphisms-lemma-composition-finite-type} et
\ref{morphisms-lemma-base-change-finite-type}.
Par conséquent, (1) résulte clairement du diagramme (\ref{more-groupoids-equation-restriction}).
Pour les autres cas, voir
Morphismes, lemmes \ref{morphisms-lemma-composition-finite-presentation},
\ref{morphisms-lemma-base-change-finite-presentation},
\ref{morphisms-lemma-composition-flat}, et
\ref{morphisms-lemma-base-change-flat}.
\end{proof}

\noindent
Le lemme suivant aurait pu servir à démontrer les résultats du lemme
précédent de manière plus uniforme.

\begin{lemma}
\label{more-groupoids-lemma-restrict-property}
Soit $S$ un schéma.
Soit $(U, R, s, t, c)$ un schéma en groupoïdes au-dessus de $S$.
Soit $g : U' \to U$ un morphisme de schémas.
Soit $(U', R', s', t', c')$ la restriction de
$(U, R, s, t, c)$ par $g$, et soit
$h = s \circ \text{pr}_1 : U' \times_{g, U, t} R \to U$. Si
$\mathcal{P}$ est une propriété des morphismes de schémas telle que
\begin{enumerate}
\item $h$ possède la propriété $\mathcal{P}$, et
\item $\mathcal{P}$ est préservée par changement de base,
\end{enumerate}
alors $s', t'$ possèdent la propriété $\mathcal{P}$.
\end{lemma}

\begin{proof}
Cela est clair, puisque $s'$ est le changement de base de $h$ d'après le
diagramme (\ref{more-groupoids-equation-restriction})
et que $t'$ est isomorphe à $s'$ en tant que morphisme de schémas.
\end{proof}

\begin{lemma}
\label{more-groupoids-lemma-double-restrict}
Soit $S$ un schéma.
Soit $(U, R, s, t, c)$ un schéma en groupoïdes au-dessus de $S$.
Soient $g : U' \to U$ et $g' : U'' \to U'$ des morphismes de schémas.
Posons $g'' = g \circ g'$.
Soit $(U', R', s', t', c')$ la restriction de $R$ à $U'$.
Soit $h = s \circ \text{pr}_1 : U' \times_{g, U, t} R \to U$,
soit $h' = s' \circ \text{pr}_1 : U'' \times_{g', U', t} R \to U'$, et
soit $h'' = s \circ \text{pr}_1 : U'' \times_{g'', U, t} R \to U$.
Le diagramme suivant est commutatif
$$
\xymatrix{
U'' \times_{g', U', t} R' \ar[d]^{h'} &
(U' \times_{g, U, t} R) \times_U (U'' \times_{g'', U, t} R)
\ar[l] \ar[r] \ar[d] &
U'' \times_{g'', U, t} R \ar[d]_{h''} \\
U' &
U' \times_{g, U, t} R \ar[l]_{\text{pr}_0} \ar[r]^h &
U
}
$$
ses deux carrés étant cartésiens, et la flèche horizontale supérieure gauche
étant donnée par la règle
$$
\begin{matrix}
(U' \times_{g, U, t} R) \times_U (U'' \times_{g'', U, t} R) &
\longrightarrow &
U'' \times_{g', U', t} R' \\
((u', r_0), (u'', r_1)) &
\longmapsto &
(u'', (c(r_1, i(r_0)), (g'(u''), u')))
\end{matrix}
$$
avec les notations expliquées dans la démonstration.
\end{lemma}

\begin{proof}
Nous allons expliciter ceci en exploitant le point de vue fonctoriel
et en ramenant le lemme à un énoncé sur les flèches dans les restrictions
d'une catégorie groupoïde. Dans la dernière formule du lemme, la
notation $((u', r_0), (u'', r_1))$ désigne un point à valeurs dans $T$ de
$(U' \times_{g, U, t} R) \times_U (U'' \times_{g'', U, t} R)$.
Cela signifie que $u', u'', r_0, r_1$ sont des points à valeurs dans $T$ de $U', U'', R, R$
et que $g(u') = t(r_0)$, $g(g'(u'')) = g''(u'') = t(r_1)$, et
$s(r_0) = s(r_1)$. Il serait plus correct d'écrire ici
$g \circ u' = t \circ r_0$ et ainsi de suite, mais cela rend la notation
encore plus illisible. Si nous considérons $r_1$ et $r_0$ comme des flèches dans
une catégorie groupoïde, nous pouvons représenter ceci par le diagramme
$$
\xymatrix{
t(r_0) = g(u') &
s(r_0) = s(r_1) \ar[l]_{r_0} \ar[r]^-{r_1} &
t(r_1) = g(g'(u''))
}
$$
Ce diagramme montre notamment que la composition
$c(r_1, i(r_0))$ a un sens. Rappelons que
$$
R' = R \times_{(t, s), U \times_S U, g \times g} U' \times_S U'
$$
ainsi, un point à valeurs dans $T$ de $R'$ s'écrit $(r, (u'_0, u'_1))$
avec $t(r) = g(u'_0)$ et $s(r) = g(u'_1)$. En particulier, étant donné
$((u', r_0), (u'', r_1))$ comme ci-dessus, on obtient le point à valeurs dans $T$
$(c(r_1, i(r_0)), (g'(u''), u'))$ de $R'$, car
$t(c(r_1, i(r_0))) = t(r_1) = g(g'(u''))$ et
$s(c(r_1, i(r_0))) = s(i(r_0)) = t(r_0) = g(u')$.
Nous laissons au lecteur le soin de montrer que le carré gauche commute
avec cette définition.

\medskip\noindent
Pour montrer que le carré gauche est cartésien,
supposons donnés $(v'', p')$ et $(v', p)$ qui sont des points à valeurs dans $T$ de
$U'' \times_{g', U', t} R'$ et $U' \times_{g, U, t} R$ avec
$v' = s'(p')$. Cela signifie aussi que $g'(v'') = t'(p')$ et
$g(v') = t(p)$. D'après la discussion ci-dessus, nous savons que nous pouvons écrire
$p' = (r, (u_0', u_1'))$ avec $t(r) = g(u'_0)$ et
$s(r) = g(u'_1)$. Avec cette notation, nous voyons que
$v' = s'(p') = u_1'$ et
$g'(v'') = t'(p') = u_0'$. Voici un diagramme
$$
\xymatrix{
s(p) \ar[r]^-p &
g(v') = g(u'_1) \ar[r]^-r &
g(u'_0) = g(g'(v''))
}
$$
Il s'agit de montrer qu'il existe un unique point à valeurs dans $T$
$((u', r_0), (u'', r_1))$ comme ci-dessus tel que
$v' = u'$, $p = r_0$, $v'' = u''$ et $p' = (c(r_1, i(r_0)), (g'(u''), u'))$.
En comparant les deux diagrammes ci-dessus, il est clair que nous n'avons pas d'autre choix
que de poser
$$
((u', r_0), (u'', r_1)) = ((v', p), (v'', c(r, p))
$$
Quelques détails sont omis.
\end{proof}

\begin{lemma}
\label{more-groupoids-lemma-double-restrict-property}
Soit $S$ un schéma.
Soit $(U, R, s, t, c)$ un schéma en groupoïdes au-dessus de $S$.
Soient $g : U' \to U$ et $g' : U'' \to U'$ des morphismes de schémas.
Posons $g'' = g \circ g'$.
Soit $(U', R', s', t', c')$ la restriction de $R$ à $U'$.
Soit $h = s \circ \text{pr}_1 : U' \times_{g, U, t} R \to U$,
soit $h' = s' \circ \text{pr}_1 : U'' \times_{g', U', t} R \to U'$, et
soit $h'' = s \circ \text{pr}_1 : U'' \times_{g'', U, t} R \to U$.
Soit $\tau \in \{Zariski, \linebreak[0] \etale, \linebreak[0]
smooth, \linebreak[0] syntomic, \linebreak[0] fppf, \linebreak[0] fpqc\}$. Soit
$\mathcal{P}$ une propriété des morphismes de schémas
qui est préservée par changement de base et qui
est locale sur la base pour la topologie $\tau$. Si
\begin{enumerate}
\item $h(U' \times_U R)$ est ouvert dans $U$,
\item $\{h : U' \times_U R \to h(U' \times_U R)\}$ est un $\tau$-recouvrement,
\item $h'$ possède la propriété $\mathcal{P}$,
\end{enumerate}
alors $h''$ possède la propriété $\mathcal{P}$. Réciproquement, si
\begin{enumerate}
\item[(a)] $\{t : R \to U\}$ est un $\tau$-recouvrement,
\item[(d)] $h''$ possède la propriété $\mathcal{P}$,
\end{enumerate}
alors $h'$ possède la propriété $\mathcal{P}$.
\end{lemma}

\begin{proof}
Cela résulte formellement des propriétés du diagramme du
lemme \ref{more-groupoids-lemma-double-restrict}.
Dans le premier cas, remarquons que l'image du morphisme
$h''$ est contenue dans l'image de $h$, puisque $g'' = g \circ g'$.
On peut donc remplacer le $U$ situé dans le coin inférieur droit du
diagramme par $h(U' \times_U R)$. Cela explique l'importance des
conditions (1) et (2) du lemme. Dans le second cas, remarquons que
$\{\text{pr}_0 : U' \times_{g, U, t} R \to U'\}$ est un $\tau$-recouvrement
en tant que changement de base de $\tau$ et de la condition (a).
\end{proof}






\section{Propriétés des groupoïdes sur les corps}
\label{more-groupoids-section-properties-groupoids-on-fields}

\noindent
Un ``groupoïde sur un corps'' désigne un schéma en groupoïdes $(U, R, s, t, c)$
où $U$ est le spectre d'un corps. Cela ne signifie {\bf pas} que
$(U, R, s, t, c)$ est défini sur un corps; plus précisément, cela ne
signifie {\bf pas} que les morphismes $s, t : R \to U$ sont égaux.
Étant donné un corps $k$, un groupe abstrait $G$ et un homomorphisme de groupes
$\varphi : G \to \text{Aut}(k)$, on obtient un schéma en groupoïdes
$(U, R, s, t, c)$ au-dessus de $\mathbf{Z}$ en posant
\begin{align*}
U & = \Spec(k) \\
R & = \coprod\nolimits_{g \in G} \Spec(k) \\
s & = \coprod\nolimits_{g \in G} \Spec(\text{id}_k) \\
t & = \coprod\nolimits_{g \in G} \Spec(\varphi(g)) \\
c & = \text{composition dans }G
\end{align*}
Cet exemple est encore un schéma en groupoïdes au-dessus de $\Spec(k^G)$.
Par conséquent, si $G$ est fini, alors $U = \Spec(k)$ est fini sur
$\Spec(k^G)$.
En un certain sens, notre but dans cette section est de montrer que des
conditions de finitude convenables sur $s, t$ imposent à tout groupoïde sur un corps
d'être défini sur un sous-corps d'indice fini $k' \subset k$.

\medskip\noindent
Si $k$ est un corps et $(G, m)$ un schéma en groupes sur $k$, de morphisme
structural $p : G \to \Spec(k)$, alors $(\Spec(k), G, p, p, m)$
est un exemple de groupoïde sur un corps (et, dans ce cas, bien entendu, toute la
structure est définie sur un corps). Ainsi, cette section peut être considérée comme
l'analogue de
Groupoïdes, section \ref{groupoids-section-properties-group-schemes-field}.

\begin{lemma}
\label{more-groupoids-lemma-groupoid-on-field-open-multiplication}
Soit $S$ un schéma. Soit $(U, R, s, t, c)$ un schéma en groupoïdes
au-dessus de $S$. Si $U$ est le spectre d'un corps, alors le morphisme de composition
$c : R \times_{s, U, t} R \to R$ est ouvert.
\end{lemma}

\begin{proof}
La composition est isomorphe à la projection
$\text{pr}_1 : R \times_{t, U, t} R \to R$ d'après le
diagramme (\ref{more-groupoids-equation-pull}).
La projection est ouverte d'après
Morphismes, lemme \ref{morphisms-lemma-scheme-over-field-universally-open}.
\end{proof}

\begin{lemma}
\label{more-groupoids-lemma-groupoid-on-field-separated}
Soit $S$ un schéma. Soit $(U, R, s, t, c)$ un schéma en groupoïdes
au-dessus de $S$. Si $U$ est le spectre d'un corps,
alors $R$ est un schéma séparé.
\end{lemma}

\begin{proof}
D'après
Groupoïdes, lemme \ref{groupoids-lemma-group-scheme-over-field-separated},
le schéma en groupes stabilisateur $G \to U$ est séparé. D'après
Groupoïdes, lemme \ref{groupoids-lemma-diagonal},
le morphisme $j = (t, s) : R \to U \times_S U$ est séparé.
Comme $U$ est le spectre d'un corps, le schéma
$U \times_S U$ est affine (par la construction des produits fibrés dans
Schémas, section \ref{schemes-section-fibre-products}).
Par conséquent, $R$ est un schéma séparé, voir
Schémas, lemme \ref{schemes-lemma-separated-permanence}.
\end{proof}

\begin{lemma}
\label{more-groupoids-lemma-groupoid-on-field-homogeneous}
Soit $S$ un schéma. Soit $(U, R, s, t, c)$ un schéma en groupoïdes
au-dessus de $S$. Supposons $U = \Spec(k)$, où $k$ est un corps.
Pour tous points $r, r' \in R$, il existe une extension de corps
$k'/k$ et des points
$r_1, r_2 \in R \times_{s, \Spec(k)} \Spec(k')$
et un diagramme
$$
\xymatrix{
R &
R \times_{s, \Spec(k)} \Spec(k')
\ar[l]_-{\text{pr}_0} \ar[r]^\varphi &
R \times_{s, \Spec(k)} \Spec(k')
\ar[r]^-{\text{pr}_0} &
R
}
$$
tel que $\varphi$ soit un isomorphisme de schémas au-dessus de $\Spec(k')$, et que
l'on ait $\varphi(r_1) = r_2$, $\text{pr}_0(r_1) = r$, et
$\text{pr}_0(r_2) = r'$.
\end{lemma}

\begin{proof}
C'est un cas particulier du
lemme \ref{more-groupoids-lemma-two-fibres},
parties (1) et (2).
\end{proof}

\begin{lemma}
\label{more-groupoids-lemma-restrict-groupoid-on-field}
Soit $S$ un schéma. Soit $(U, R, s, t, c)$ un schéma en groupoïdes
au-dessus de $S$. Supposons $U = \Spec(k)$, où $k$ est un corps.
Soit $k'/k$ une extension de corps, $U' = \Spec(k')$,
et soit $(U', R', s', t', c')$ la restriction de
$(U, R, s, t, c)$ par $U' \to U$. Dans le diagramme qui la définit
$$
\xymatrix{
R' \ar[d] \ar[r] \ar@/_3pc/[dd]_{t'} \ar@/^1pc/[rr]^{s'} \ar@{..>}[rd] &
R \times_{s, U} U' \ar[r] \ar[d] &
U' \ar[d] \\
U' \times_{U, t} R \ar[d] \ar[r] &
R \ar[r]^s \ar[d]_t &
U \\
U' \ar[r] &
U
}
$$
tous les morphismes sont surjectifs, plats et universellement ouverts.
La flèche pointillée $R' \to R$ est en outre affine.
\end{lemma}

\begin{proof}
Le morphisme $U' \to U$ est égal à $\Spec(k') \to \Spec(k)$;
il est donc affine, surjectif et plat. Les morphismes $s, t : R \to U$
ainsi que le morphisme $U' \to U$ sont universellement ouverts d'après
Morphismes, lemme \ref{morphisms-lemma-scheme-over-field-universally-open}.
Comme $R$ n'est pas vide et que $U$ est le spectre d'un corps, les morphismes
$s, t : R \to U$ sont surjectifs et plats. On conclut alors à l'aide de
Morphismes, lemmes \ref{morphisms-lemma-base-change-surjective},
\ref{morphisms-lemma-composition-surjective},
\ref{morphisms-lemma-composition-open},
\ref{morphisms-lemma-base-change-affine},
\ref{morphisms-lemma-composition-affine},
\ref{morphisms-lemma-base-change-flat}, et
\ref{morphisms-lemma-composition-flat}.
\end{proof}

\begin{lemma}
\label{more-groupoids-lemma-groupoid-on-field-explain-points}
Soit $S$ un schéma. Soit $(U, R, s, t, c)$ un schéma en groupoïdes
au-dessus de $S$. Supposons $U = \Spec(k)$, où $k$ est un corps.
Pour tout point $r \in R$, il existe
\begin{enumerate}
\item une extension de corps $k'/k$ avec $k'$ algébriquement clos,
\item un point $r' \in R'$ où $(U', R', s', t', c')$ est la
restriction de $(U, R, s, t, c)$ par $\Spec(k') \to \Spec(k)$
\end{enumerate}
de telle sorte que
\begin{enumerate}
\item le point $r'$ a pour image $r$ par le morphisme $R' \to R$, et
\item les applications $s', t' : R' \to \Spec(k')$ induisent des isomorphismes
$k' \to \kappa(r')$.
\end{enumerate}
\end{lemma}

\begin{proof}
En traduisant l'énoncé géométrique en un énoncé sur les corps,
cela signifie que nous pouvons trouver un diagramme
$$
\xymatrix{
k' & k' \ar[l]^1 & \\
k' \ar[u]^\tau & \kappa(r) \ar[lu]^\sigma & k \ar[l]^-s \ar[lu]_i \\
& k \ar[lu]^i \ar[u]_t
}
$$
où $i : k \to k'$ est le plongement de $k$ dans $k'$,
les applications $s, t : k \to \kappa(r)$ sont induites par $s, t : R \to U$, et
l'application $\tau : k' \to k'$ est un automorphisme. Pour construire un tel
diagramme, on peut procéder comme suit :
\begin{enumerate}
\item Choisir un homomorphisme de corps $i : k \to k'$, avec $k'$ algébriquement clos et de
très grand degré de transcendance sur $k$.
\item Choisir un plongement $\sigma : \kappa(r) \to k'$ tel que
$\sigma \circ s = i$. Un tel $\sigma$ existe, car il suffit de
choisir une base de transcendance $\{x_\alpha\}_{\alpha \in A}$ de $\kappa(r)$
sur $k$ et des $y_\alpha \in k'$, $\alpha \in A$, qui soient algébriquement
indépendants sur $i(k)$, puis d'envoyer $s(k)(\{x_\alpha\})$ dans $k'$
selon les règles $s(\lambda) \mapsto i(\lambda)$ pour $\lambda \in k$
et $x_\alpha \mapsto y_\alpha$ pour $\alpha \in A$.
Prolonger ensuite en $\tau : \kappa(\alpha) \to k'$ en utilisant le fait que $k'$ est
algébriquement clos.
\item Choisir un automorphisme $\tau : k' \to k'$ tel que
$\tau \circ i = \sigma \circ t$. Pour cela, choisir une base de transcendance
$\{x_\alpha\}_{\alpha \in A}$ de $k$ sur son sous-corps premier.
D'une part, compléter $\{i(x_\alpha)\}$ en une base de transcendance de
$k'$ en ajoutant $\{y_\beta\}_{\beta \in B}$, et compléter
$\{\sigma(t(x_\alpha))\}$ en une base de transcendance de $k'$ en ajoutant
$\{z_\gamma\}_{\gamma \in C}$.
Comme $k'$ est algébriquement clos, nous pouvons prolonger l'isomorphisme
$\sigma \circ t \circ i^{-1} : i(k) \to \sigma(t(k))$
en un isomorphisme $\tau' : \overline{i(k)} \to \overline{\sigma(t(k))}$
entre leurs clôtures algébriques dans $k'$.
Comme $k'$ a un grand degré de transcendance,
nous voyons que les ensembles $B$ et $C$ ont le même cardinal.
Ainsi, nous pouvons utiliser une bijection
$B \to C$ pour prolonger $\tau'$ en un isomorphisme
$$
\overline{i(k)}(\{y_\beta\})
\longrightarrow
\overline{\sigma(t(k))}(\{z_\gamma\})
$$
et alors, puisque $k'$ est la clôture algébrique des deux membres, nous
voyons que cela se prolonge en un automorphisme $\tau : k' \to k'$
comme souhaité.
\end{enumerate}
Cela démontre le lemme.
\end{proof}

\begin{lemma}
\label{more-groupoids-lemma-groupoid-on-field-move-point}
Soit $S$ un schéma. Soit $(U, R, s, t, c)$ un schéma en groupoïdes
sur $S$. Supposons $U = \Spec(k)$, où $k$ est un corps.
Si $r \in R$ est un point tel que $s, t$ induisent des
isomorphismes $k \to \kappa(r)$, alors l'application
$$
R \longrightarrow R, \quad
x \longmapsto c(r, x)
$$
(voir la démonstration pour une notation précise) est un automorphisme $R \to R$
qui envoie $e$ sur $r$.
\end{lemma}

\begin{proof}
Cela est tout à fait évident si l'on envisage les groupoïdes de manière
fonctorielle. Mais nous allons aussi l'expliciter entièrement.
Notons $a : U \to R$ le morphisme d'image $r$ tel que
$s \circ a = \text{id}_U$, qui existe par l'hypothèse
selon laquelle $s : k \to \kappa(r)$ est un isomorphisme. De même, notons
$b : U \to R$ le morphisme d'image $r$ tel que
$t \circ b = \text{id}_U$. Remarquons que
$b = a \circ (t \circ a)^{-1}$, et, en particulier,
$a \circ s \circ b = b$.

\medskip\noindent
Considérons le morphisme $\Psi : R \to R$ défini sur les points à valeurs dans $T$
par
$$
(f : T \to R) \longmapsto (c(a \circ t \circ f, f) : T \to R)
$$
Pour voir qu'il est défini, nous devons vérifier que
$s \circ a \circ t \circ f = t \circ f$, ce qui est évident puisque $s \circ a = 1$.
Remarquons que $\Phi(e) = a$ ; ainsi, pour démontrer le lemme, il
suffit de montrer que $\Phi$ est un automorphisme de $R$.
Soit $\Phi : R \to R$ le morphisme défini sur les points à valeurs dans $T$ par
$$
(g : T \to R) \longmapsto (c(i \circ b \circ t \circ g, g) : T \to R).
$$
Il est défini car
$s \circ i \circ b \circ t \circ g = t \circ b \circ t \circ g =
t \circ g$. Nous affirmons que $\Phi$ et $\Psi$ sont inverses
l'un de l'autre. Pour le voir, calculons
\begin{align*}
& c(a \circ t \circ c(i \circ b \circ t \circ g, g),
c(i \circ b \circ t \circ g, g)) \\
& =
c(a \circ t \circ i \circ b \circ t \circ g,
c(i \circ b \circ t \circ g, g)) \\
& =
c(a \circ s \circ b \circ t \circ g,
c(i \circ b \circ t \circ g, g)) \\
& =
c(b \circ t \circ g, c(i \circ b \circ t \circ g, g)) \\
& =
c(c(b \circ t \circ g, i \circ b \circ t \circ g), g)) \\
& =
c(e, g) \\
& = g
\end{align*}
où nous avons utilisé la relation $a \circ s \circ b = b$ établie ci-dessus.
Dans l'autre sens, nous avons
\begin{align*}
& c(i \circ b \circ t \circ c(a \circ t \circ f, f), c(a \circ t \circ f, f)) \\
& =
c(i \circ b \circ t \circ a \circ t \circ f, c(a \circ t \circ f, f)) \\
& =
c(i \circ a \circ (t \circ a)^{-1} \circ t \circ a \circ t \circ f,
c(a \circ t \circ f, f)) \\
& =
c(i \circ a \circ t \circ f, c(a \circ t \circ f, f)) \\
& =
c(c(i \circ a \circ t \circ f, a \circ t \circ f), f) \\
& =
c(e, f) \\
& = f
\end{align*}
Le lemme est démontré.
\end{proof}

\begin{lemma}
\label{more-groupoids-lemma-groupoid-on-field-translate-open}
Soit $S$ un schéma. Soit $(U, R, s, t, c)$ un schéma en groupoïdes
sur $S$. Si $U$ est le spectre d'un corps, si $W \subset R$ est ouvert
et si $Z \to R$ est un morphisme de schémas, alors l'image de la
composée $Z \times_{s, U, t} W \to R \times_{s, U, t} R \to R$ est ouverte.
\end{lemma}

\begin{proof}
Écrivons $U = \Spec(k)$. Considérons une extension de corps $k'/k$. Notons
$U' = \Spec(k')$. Soit $R'$ la restriction de $R$ par $U' \to U$.
Posons $Z' = Z \times_R R'$ et $W' = R' \times_R W$.
Considérons un point $\xi = (z, w)$ de $Z \times_{s, U, t} W$.
Soit $r \in R$ l'image de $z$ par $Z \to R$.
Choisissons $k' \supset k$ et $r' \in R'$ comme dans le
lemme \ref{more-groupoids-lemma-groupoid-on-field-explain-points}.
Nous pouvons choisir $z' \in Z'$ se projetant sur $z$ et sur $r'$.
Nous pouvons alors trouver $\xi' \in Z' \times_{s', U', t'} W'$
se projetant sur $z'$ et sur $\xi$. L'ouvert $c(r', W')$
(lemme \ref{more-groupoids-lemma-groupoid-on-field-move-point}) est
contenu dans l'image de $Z' \times_{s', U', t'} W' \to R'$.
Observons que $Z' \times_{s', U', t'} W' = (Z \times_{s, U, t} W)
\times_{R \times_{s, U, t} R} (R' \times_{s', U', t'} R')$.
Par conséquent, l'image de $Z' \times_{s', U', t'} W' \to R' \to R$
est contenue dans l'image de $Z \times_{s, U, t} W \to R$.
Comme le morphisme $R' \to R$ est ouvert (lemme \ref{more-groupoids-lemma-restrict-groupoid-on-field}),
nous concluons que l'image contient un voisinage ouvert de
l'image de $\xi$, comme souhaité.
\end{proof}

\begin{lemma}
\label{more-groupoids-lemma-groupoid-on-field-geometrically-irreducible}
Soit $S$ un schéma. Soit $(U, R, s, t, c)$ un schéma en groupoïdes
sur $S$. Supposons $U = \Spec(k)$, où $k$ est un corps.
Par abus de notation, notons $e \in R$ l'image du morphisme identité
$e : U \to R$. Alors :
\begin{enumerate}
\item tout anneau local $\mathcal{O}_{R, r}$ de $R$ possède un unique
idéal premier minimal ;
\item il existe exactement une composante irréductible $Z$ de $R$
passant par $e$ ;
\item $Z$ est géométriquement irréductible sur $k$ pour l'un ou l'autre
des morphismes $s$ et $t$.
\end{enumerate}
\end{lemma}

\begin{proof}
Soit $r \in R$ un point.
Dans cette démonstration, nous utiliserons sans autre mention la correspondance entre les composantes irréductibles
de $R$ passant par un point $r$ et les idéaux premiers minimaux de l'anneau
local $\mathcal{O}_{R, r}$.
Choisissons $k \subset k'$ et $r' \in R'$ comme dans le
lemme \ref{more-groupoids-lemma-groupoid-on-field-explain-points}.
Remarquons que $\mathcal{O}_{R, r} \to \mathcal{O}_{R', r'}$
est fidèlement plat et local, voir le
lemme \ref{more-groupoids-lemma-restrict-groupoid-on-field}.
Ainsi, le résultat pour $r' \in R'$ entraîne le résultat pour $r \in R$.
Autrement dit, nous pouvons supposer que $s, t : k \to \kappa(r)$
sont des isomorphismes. D'après le
lemme \ref{more-groupoids-lemma-groupoid-on-field-move-point},
il existe un automorphisme envoyant $e$ sur $r$.
Nous pouvons donc supposer $r = e$, c'est-à-dire que (1) découle de (2).

\medskip\noindent
Démontrons d'abord (2) dans le cas où $k$ est séparablement clos.
Soient en effet $X, Y \subset R$ des composantes irréductibles
passant par $e$. Alors, d'après
Variétés, lemme \ref{varieties-lemma-bijection-irreducible-components} et
\ref{varieties-lemma-separably-closed-irreducible},
le schéma $X \times_{s, U, t} Y$ est lui aussi irréductible.
Ainsi, $c(X \times_{s, U, t} Y) \subset R$ est un sous-ensemble irréductible.
Nous affirmons qu'il contient à la fois $X$ et $Y$ (comme sous-ensembles de $R$).
En effet, soit $T$ le spectre d'un corps. Si $x : T \to X$ est un $T$-point
de $X$, alors $c(x, e \circ s \circ x) = x$ et $e \circ s \circ x$
se factorise par $Y$ puisque $e \in Y$. Il en va de même pour les points de $Y$.
Cela implique manifestement que $X = Y$, c'est-à-dire qu'il existe une unique composante irréductible
de $R$ passant par $e$.

\medskip\noindent
Démonstration de (2) et (3) dans le cas général.
Soit $k \subset k'$ une clôture séparable, et
soit $(U', R', s', t', c')$ la restriction de
$(U, R, s, t, c)$ par $\Spec(k') \to \Spec(k)$.
D'après le paragraphe précédent, il existe exactement une composante irréductible
$Z'$ de $R'$ passant par $e'$.
Notons $e'' \in R \times_{s, U} U'$ le changement de base de $e$.
Comme $R' \to R \times_{s, U} U'$ est fidèlement plat, voir le
lemme \ref{more-groupoids-lemma-restrict-groupoid-on-field},
et que $e' \mapsto e''$, il existe exactement une composante
irréductible $Z''$ de $R \times_{s, k} k'$ passant
par $e''$. Comme $R \times_k k' \to R$ est fidèlement
plat, cela implique qu'il existe exactement une composante irréductible $Z$ de $R$
passant par $e$. Cela démontre (2).

\medskip\noindent
Pour démontrer (3), soit $Z''' \subset R \times_k k'$ une composante irréductible quelconque
de $Z \times_k k'$. D'après
Variétés, lemme \ref{varieties-lemma-orbit-irreducible-components},
nous avons $Z''' = \sigma(Z'')$ pour un certain $\sigma \in \text{Gal}(k'/k)$.
Comme $\sigma(e'') = e''$, nous avons $e'' \in Z'''$, et donc
$Z''' = Z''$. Cela signifie que $Z$ est géométriquement irréductible
sur $\Spec(k)$ pour le morphisme $s$.
Le même argument montre que $Z$ est géométriquement irréductible
sur $\Spec(k)$ pour le morphisme $t$.
\end{proof}

\begin{lemma}
\label{more-groupoids-lemma-groupoid-on-field-locally-finite-type-dimension}
Soit $S$ un schéma. Soit $(U, R, s, t, c)$ un schéma en groupoïdes
sur $S$. Supposons $U = \Spec(k)$, où $k$ est un corps.
Supposons que $s, t$ soient localement de type fini.
Alors :
\begin{enumerate}
\item $R$ est équidimensionnel ;
\item $\dim(R) = \dim_r(R)$ pour tout $r \in R$ ;
\item pour tout $r \in R$, nous avons
$\text{trdeg}_{s(k)}(\kappa(r)) = \text{trdeg}_{t(k)}(\kappa(r))$ ; et
\item pour tout point fermé $r \in R$, nous avons
$\dim(R) = \dim(\mathcal{O}_{R, r})$.
\end{enumerate}
\end{lemma}

\begin{proof}
Soient $r, r' \in R$.
Alors $\dim_r(R) = \dim_{r'}(R)$ d'après le
lemme \ref{more-groupoids-lemma-groupoid-on-field-homogeneous}
et
Morphismes, lemme \ref{morphisms-lemma-dimension-fibre-after-base-change}.
D'après
Morphismes, lemme \ref{morphisms-lemma-dimension-fibre-at-a-point},
nous avons
$$
\dim_r(R) =
\dim(\mathcal{O}_{R, r}) + \text{trdeg}_{s(k)}(\kappa(r)) =
\dim(\mathcal{O}_{R, r}) + \text{trdeg}_{t(k)}(\kappa(r)).
$$
D'autre part, la dimension de $R$ (ou de tout ouvert de $R$)
est la borne supérieure des dimensions des anneaux locaux de $R$, voir
Propriétés, lemme \ref{properties-lemma-codimension-local-ring}.
La dimension de l'anneau local est manifestement maximale aux points fermés $r$, auquel cas
$\text{trdeg}_k(\kappa(r)) = 0$ (par le théorème des zéros de Hilbert, voir
Morphismes, section \ref{morphisms-section-points-finite-type}).
Le lemme en résulte.
\end{proof}

\begin{lemma}
\label{more-groupoids-lemma-groupoid-on-field-dimension-equal-stabilizer}
Soit $S$ un schéma. Soit $(U, R, s, t, c)$ un schéma en groupoïdes
sur $S$. Supposons $U = \Spec(k)$, où $k$ est un corps.
Supposons que $s, t$ soient localement de type fini.
Alors $\dim(R) = \dim(G)$, où $G$ est le schéma en groupes stabilisateur de $R$.
\end{lemma}

\begin{proof}
Soit $Z \subset R$ la composante irréductible passant par $e$ (voir le
lemme \ref{more-groupoids-lemma-groupoid-on-field-geometrically-irreducible}),
considérée comme un sous-schéma fermé intègre de $R$.
Soit $k'_s$, resp.\ $k'_t$, la clôture intégrale de
$s(k)$, resp.\ $t(k)$, dans $\Gamma(Z, \mathcal{O}_Z)$.
Rappelons que $k'_s$ et $k'_t$ sont des corps, voir
Variétés, lemme \ref{varieties-lemma-integral-closure-ground-field}.
D'après
Variétés, proposition \ref{varieties-proposition-unique-base-field},
nous avons $k'_s = k'_t$ comme sous-anneaux de $\Gamma(Z, \mathcal{O}_Z)$.
Comme $e$ se factorise par $Z$, nous obtenons un diagramme commutatif
$$
\xymatrix{
k \ar[rd]_t \ar[rrd]^1 \\
& \Gamma(Z, \mathcal{O}_Z) \ar[r]^e & k \\
k \ar[ru]^s \ar[rru]_1
}
$$
Cela montre d'une part que $k'_s = s(k)$ et $k'_t = t(k)$, donc que
$s(k) = t(k)$, ce qui, combiné au diagramme ci-dessus, implique
que $s = t$ ! Autrement dit, nous concluons que $Z$ est un sous-schéma fermé
de $G = R \times_{(t, s), U \times_S U, \Delta} U$.
Le lemme en résulte, car $G$ et $R$ sont tous deux équidimensionnels, voir
lemme \ref{more-groupoids-lemma-groupoid-on-field-locally-finite-type-dimension} et
Groupoïdes, lemme \ref{groupoids-lemma-group-scheme-finite-type-field}.
\end{proof}

\begin{remark}
\label{more-groupoids-remark-warn-dimension-groupoid-on-field}
Avertissement :
Le lemme \ref{more-groupoids-lemma-groupoid-on-field-dimension-equal-stabilizer}
est faux sans la condition que $s$ et $t$ soient localement de
type fini.
Un exemple simple consiste à partir de l'action
$$
\mathbf{G}_{m, \mathbf{Q}} \times_{\mathbf{Q}} \mathbf{A}^1_{\mathbf{Q}}
\to \mathbf{A}^1_{\mathbf{Q}}
$$
et à restreindre le schéma en groupoïdes correspondant au point générique de
$\mathbf{A}^1_{\mathbf{Q}}$. Autrement dit, effectuons la restriction suivant le morphisme
$\Spec(\mathbf{Q}(x)) \to
\Spec(\mathbf{Q}[x]) = \mathbf{A}^1_{\mathbf{Q}}$.
On obtient alors un schéma en groupoïdes
$(U, R, s, t, c)$ avec
$U = \Spec(\mathbf{Q}(x))$
et
$$
R = \Spec\left(
\mathbf{Q}(x)[y]\left[
\frac{1}{P(xy)}, P \in \mathbf{Q}[T], P \not = 0
\right]
\right)
$$
Dans ce cas, $\dim(R) = 1$ et $\dim(G) = 0$.
\end{remark}

\begin{lemma}
\label{more-groupoids-lemma-groupoid-characteristic-zero-smooth}
Soit $S$ un schéma. Soit $(U, R, s, t, c)$ un schéma en groupoïdes
sur $S$. Supposons que les conditions suivantes soient satisfaites :
\begin{enumerate}
\item $U = \Spec(k)$, où $k$ est un corps ;
\item $s, t$ sont localement de type fini ; et
\item la caractéristique de $k$ est nulle.
\end{enumerate}
Alors $s, t : R \to U$ sont lisses.
\end{lemma}

\begin{proof}
D'après le
lemme \ref{more-groupoids-lemma-sheaf-differentials},
le faisceau des différentielles de $R \to U$ est libre.
La lissité résulte donc de
Variétés, lemme \ref{varieties-lemma-char-zero-differentials-free-smooth}.
\end{proof}

\begin{lemma}
\label{more-groupoids-lemma-reduced-group-scheme-perfect-field-characteristic-p-smooth}
Soit $S$ un schéma. Soit $(U, R, s, t, c)$ un schéma en groupoïdes
sur $S$. Supposons que les conditions suivantes soient satisfaites :
\begin{enumerate}
\item $U = \Spec(k)$, où $k$ est un corps ;
\item $s, t$ sont localement de type fini ;
\item $R$ est réduit ; et
\item $k$ est parfait.
\end{enumerate}
Alors $s, t : R \to U$ sont lisses.
\end{lemma}

\begin{proof}
D'après le
lemme \ref{more-groupoids-lemma-sheaf-differentials},
le faisceau $\Omega_{R/U}$ est libre. Le lemme résulte donc de
Variétés, lemme \ref{varieties-lemma-char-p-differentials-free-smooth}.
\end{proof}










\section{Morphismes de groupoïdes sur des corps}
\label{more-groupoids-section-morphisms-groupoids-on-fields}

\noindent
Cette section étudie les morphismes entre groupoïdes sur des corps.
C'est légèrement plus général, mais très proche de l'étude des
morphismes de schémas en groupes sur un corps.

\begin{situation}
\label{more-groupoids-situation-morphism-groupoids-on-field}
Soit $S$ un schéma.
Soit $U = \Spec(k)$ un schéma sur $S$, où $k$ est un corps.
Soient $(U, R_1, s_1, t_1, c_1)$ et $(U, R_2, s_2, t_2, c_2)$ des schémas en groupoïdes
sur $S$ ayant la même première composante. Soit $a : R_1 \to R_2$ un morphisme
tel que $(\text{id}_U, a)$ définisse un morphisme de schémas en
groupoïdes sur $S$, voir
Groupoïdes, définition \ref{groupoids-definition-groupoid}.
En particulier, les diagrammes suivants sont commutatifs
$$
\vcenter{
\xymatrix{
R_1 \ar[rrd]^{t_1} \ar[rdd]_{s_1} \ar[rd]_a \\
& R_2 \ar[d]^{t_2} \ar[r]_{s_2} & U \\
& U
}
}
\quad\quad
\vcenter{
\xymatrix{
R_1 \times_{s_1, U, t_1} R_1 \ar[r]_-{c_1} \ar[d]_{a \times a} &
R_1 \ar[d]^a \\
R_2 \times_{s_2, U, t_2} R_2 \ar[r]^-{c_2} &
R_2
}
}
$$
\end{situation}

\noindent
Le lemme suivant généralise
Groupoïdes, lemme \ref{groupoids-lemma-open-subgroup-closed-over-field}.

\begin{lemma}
\label{more-groupoids-lemma-open-image-is-closed}
Notations et hypothèses comme dans la
situation \ref{more-groupoids-situation-morphism-groupoids-on-field}.
Si $a(R_1)$ est ouvert dans $R_2$, alors $a(R_1)$ est fermé dans $R_2$.
\end{lemma}

\begin{proof}
Soit $r_2 \in R_2$ un point appartenant à l'adhérence de $a(R_1)$.
Nous voulons montrer que $r_2 \in a(R_1)$. Choisissons $k \subset k'$ et
$r_2' \in R'_2$ adapté à $(U, R_2, s_2, t_2, c_2)$ et à $r_2$ comme dans le
lemme \ref{more-groupoids-lemma-groupoid-on-field-explain-points}.
Soit $R_i'$ la restriction de $R_i$ par le morphisme
$U' = \Spec(k') \to U = \Spec(k)$.
Soit $a' : R'_1 \to R_2'$ le changement de base de $a$. Le diagramme
$$
\xymatrix{
R'_1 \ar[r]_{a'} \ar[d]_{p_1} & R'_2 \ar[d]^{p_2} \\
R_1 \ar[r]^a & R_2
}
$$
est un carré cartésien. L'image de $a'$ est donc l'image réciproque de
l'image de $a$ par le morphisme $p_2 : R'_2 \to R_2$. D'après le
lemme \ref{more-groupoids-lemma-restrict-groupoid-on-field},
l'application $p_2$ est surjective et ouverte. Ainsi, d'après
Topologie, lemme \ref{topology-lemma-open-morphism-quotient-topology},
nous voyons que $r_2'$ appartient à l'adhérence de $a'(R'_1)$.
Cela signifie que nous pouvons supposer que $r_2 \in R_2$ vérifie
que les applications $k \to \kappa(r_2)$ induites
par $s_2$ et $t_2$ sont des isomorphismes.

\medskip\noindent
Dans ce cas, nous pouvons utiliser le
lemme \ref{more-groupoids-lemma-groupoid-on-field-move-point}.
Ce lemme implique que $c(r_2, a(R_1))$ est un voisinage ouvert de $r_2$.
Ainsi, $a(R_1) \cap c(r_2, a(R_1)) \not = \emptyset$, puisque nous avons supposé
que $r_2$ appartenait à l'adhérence de $a(R_1)$.
En utilisant les morphismes d'inversion de $R_2$ et de $R_1$, nous voyons que cela signifie que
$c_2(a(R_1), a(R_1))$ contient $r_2$.
Comme $c_2(a(R_1), a(R_1)) \subset a(c_1(R_1, R_1)) = a(R_1)$,
nous concluons que $r_2 \in a(R_1)$, comme souhaité.
\end{proof}

\begin{lemma}
\label{more-groupoids-lemma-map-groupoids-on-field-image}
Notations et hypothèses comme dans la
situation \ref{more-groupoids-situation-morphism-groupoids-on-field}.
Soit $Z \subset R_2$ le sous-schéma fermé réduit (voir
Schémas, définition \ref{schemes-definition-reduced-induced-scheme})
dont l'espace topologique sous-jacent est l'adhérence de l'image de
$a : R_1 \to R_2$. Alors
$c_2(Z \times_{s_2, U, t_2} Z) \subset Z$
ensemblistement.
\end{lemma}

\begin{proof}
Considérons le diagramme commutatif
$$
\xymatrix{
R_1 \times_{s_1, U, t_1} R_1 \ar[r] \ar[d] & R_1 \ar[d] \\
R_2 \times_{s_2, U, t_2} R_2 \ar[r] & R_2
}
$$
D'après
Variétés, lemme \ref{varieties-lemma-closure-image-product-map},
l'adhérence de l'image de la flèche verticale de gauche est (ensemblistement)
$Z \times_{s_2, U, t_2} Z$.
Le résultat en découle.
\end{proof}

\begin{lemma}
\label{more-groupoids-lemma-map-groupoids-on-perfect-field-image}
Notations et hypothèses comme dans la
situation \ref{more-groupoids-situation-morphism-groupoids-on-field}.
Supposons que $k$ soit parfait.
Soit $Z \subset R_2$ le sous-schéma fermé réduit (voir
Schémas, définition \ref{schemes-definition-reduced-induced-scheme})
dont l'espace topologique sous-jacent est l'adhérence de l'image de
$a : R_1 \to R_2$. Alors
$$
(U, Z, s_2|_Z, t_2|_Z, c_2|_Z)
$$
est un schéma en groupoïdes sur $S$.
\end{lemma}

\begin{proof}
Expliquons d'abord pourquoi l'énoncé a un sens. Puisque $U$ est le spectre
d'un corps parfait $k$, le schéma $Z$ est géométriquement réduit
sur $k$ (pour l'une ou l'autre projection), voir
Variétés, lemme \ref{varieties-lemma-perfect-reduced}.
Ainsi, le schéma $Z \times_{s_2, U, t_2} Z \subset Z$
est réduit, voir
Variétés, lemme \ref{varieties-lemma-geometrically-reduced-any-base-change}.
Par conséquent, d'après le
lemme \ref{more-groupoids-lemma-map-groupoids-on-field-image},
nous voyons que $c$ induit un morphisme
$Z \times_{s_2, U, t_2} Z \to Z$.
Enfin, il est clair que $e_2$ se factorise par $Z$
et que le morphisme $i_2 : R_2 \to R_2$ préserve $Z$. Puisque les morphismes
du septuple
$(U, R_2, s_2, t_2, c_2, e_2, i_2)$
satisfont aux axiomes d'un groupoïde, il s'ensuit qu'après restriction
à $Z$, ils satisfont encore à ces axiomes.
\end{proof}

\begin{lemma}
\label{more-groupoids-lemma-locally-closed-image-is-closed}
Notations et hypothèses comme dans la
situation \ref{more-groupoids-situation-morphism-groupoids-on-field}.
Si l'image $a(R_1)$ est un sous-ensemble localement fermé de $R_2$,
alors c'est un sous-ensemble fermé.
\end{lemma}

\begin{proof}
Soit $k \subset k'$ une clôture parfaite du corps $k$.
Soit $R_i'$ la restriction de $R_i$ par le morphisme
$U' = \Spec(k') \to \Spec(k)$. Remarquons que les
morphismes $R_i' \to R_i$ sont des homéomorphismes universels, car ce sont des
composés de changements de base de l'homéomorphisme universel
$U' \to U$ (voir le diagramme dans l'énoncé du
lemme \ref{more-groupoids-lemma-restrict-groupoid-on-field}).
Il suffit donc de démontrer que l'image $a'(R_1')$ est fermée
dans $R_2'$. Autrement dit, nous pouvons supposer que $k$ est parfait.

\medskip\noindent
Si $k$ est parfait, l'adhérence de l'image est
un schéma en groupoïdes $Z \subset R_2$, d'après le
lemme \ref{more-groupoids-lemma-map-groupoids-on-perfect-field-image}.
En appliquant le même lemme à
$\text{id}_{R_1} : R_1 \to R_1$,
nous voyons que $(R_2)_{red}$ est un schéma en groupoïdes.
Nous pouvons donc appliquer le
lemme \ref{more-groupoids-lemma-open-image-is-closed}
au morphisme
$a|_{(R_2)_{red}} : (R_2)_{red} \to Z$
pour conclure que $Z$ est égal à l'image de $a$.
\end{proof}

\begin{lemma}
\label{more-groupoids-lemma-quasi-compact-map-groupoids-on-field-image}
Notations et hypothèses comme dans la
situation \ref{more-groupoids-situation-morphism-groupoids-on-field}.
Supposons que $a : R_1 \to R_2$ soit un morphisme quasi-compact.
Soit $Z \subset R_2$ l'image schématique (voir
Morphismes, définition \ref{morphisms-definition-scheme-theoretic-image})
de $a : R_1 \to R_2$. Alors
$$
(U, Z, s_2|_Z, t_2|_Z, c_2|_Z)
$$
est un schéma en groupoïdes sur $S$.
\end{lemma}

\begin{proof}
La principale difficulté consiste à montrer que $c_2|_{Z \times_{s_2, U, t_2} Z}$
a son image dans $Z$. Considérons le diagramme commutatif
$$
\xymatrix{
R_1 \times_{s_1, U, t_1} R_1 \ar[r] \ar[d]^{a \times a} & R_1 \ar[d] \\
R_2 \times_{s_2, U, t_2} R_2 \ar[r] & R_2
}
$$
D'après le
lemme \ref{varieties-lemma-scheme-theoretic-image-product-map} du chapitre Variétés,
nous voyons que l'image schématique de $a \times a$ est
$Z \times_{s_2, U, t_2} Z$. Par commutativité du diagramme, nous
concluons que $Z \times_{s_2, U, t_2} Z$ a son image dans $Z$ par la flèche
horizontale inférieure. Comme dans la démonstration du
lemme \ref{more-groupoids-lemma-map-groupoids-on-perfect-field-image},
on a aussi $i_2(Z) \subset Z$ et
$e_2$ se factorise par $Z$. Nous concluons donc comme dans la
démonstration de ce lemme.
\end{proof}

\begin{lemma}
\label{more-groupoids-lemma-groupoid-on-field-image}
Soit $S$ un schéma. Soit $(U, R, s, t, c)$ un schéma en groupoïdes
sur $S$. Supposons que $U$ soit le spectre d'un corps.
Soit $Z \subset U \times_S U$ le sous-schéma fermé réduit (voir
Schémas, définition \ref{schemes-definition-reduced-induced-scheme})
dont l'espace topologique sous-jacent est l'adhérence de l'image de
$j = (t, s) : R \to U \times_S U$. Alors
$\text{pr}_{02}(Z \times_{\text{pr}_1, U, \text{pr}_0} Z) \subset Z$
ensemblistement.
\end{lemma}

\begin{proof}
Puisque $(U, U \times_S U, \text{pr}_1, \text{pr}_0, \text{pr}_{02})$
est un schéma en groupoïdes sur $S$, c'est un cas particulier du
lemme \ref{more-groupoids-lemma-map-groupoids-on-field-image}.
Mais nous pouvons aussi le démontrer directement comme suit.

\medskip\noindent
Écrivons $U = \Spec(k)$. Notons
$R_s$ (resp.\ $Z_s$, resp.\ $U^2_s$) le schéma
$R$ (resp.\ $Z$, resp.\ $U \times_S U$) considéré comme un schéma sur $k$ au moyen de
$s$ (resp.\ $\text{pr}_1|_Z$, resp.\ $\text{pr}_1$).
De même, notons
${}_tR$ (resp.\ ${}_tZ$, resp.\ ${}_tU^2$) le schéma
$R$ (resp.\ $Z$, resp.\ $U \times_S U$) considéré comme un schéma sur $k$ au moyen de
$t$ (resp.\ $\text{pr}_0|_Z$, resp.\ $\text{pr}_0$).
Le morphisme $j$ induit des morphismes de schémas
$j_s : R_s \to U^2_s$ et ${}_tj : {}_tR \to {}_tU^2$ sur $k$.
Considérons le diagramme commutatif
$$
\xymatrix{
R_s \times_k {}_tR \ar[r]^c \ar[d]_{j_s \times {}_tj} &  R \ar[d]^j \\
U^2_s \times_k {}_tU^2 \ar[r] & U \times_S U
}
$$
D'après le
lemme \ref{varieties-lemma-closure-image-product-map} du chapitre Variétés,
nous voyons que l'adhérence de l'image de $j_s \times {}_tj$ est
$Z_s \times_k {}_tZ$. Par commutativité du diagramme, nous
concluons que $Z_s \times_k {}_tZ$ a son image dans $Z$ par la flèche
horizontale inférieure.
\end{proof}

\begin{lemma}
\label{more-groupoids-lemma-groupoid-on-perfect-field-image}
Soit $S$ un schéma. Soit $(U, R, s, t, c)$ un schéma en groupoïdes
sur $S$. Supposons que $U$ soit le spectre d'un corps parfait.
Soit $Z \subset U \times_S U$ le sous-schéma fermé réduit (voir
Schémas, définition \ref{schemes-definition-reduced-induced-scheme})
dont l'espace topologique sous-jacent est l'adhérence de l'image de
$j = (t, s) : R \to U \times_S U$.
Alors
$$
(U, Z, \text{pr}_0|_Z, \text{pr}_1|_Z,
\text{pr}_{02}|_{Z \times_{\text{pr}_1, U, \text{pr}_0} Z})
$$
est un schéma en groupoïdes sur $S$.
\end{lemma}

\begin{proof}
Puisque $(U, U \times_S U, \text{pr}_1, \text{pr}_0, \text{pr}_{02})$
est un schéma en groupoïdes sur $S$, c'est un cas particulier du
lemme \ref{more-groupoids-lemma-map-groupoids-on-perfect-field-image}.
Mais nous pouvons aussi le démontrer directement comme suit.

\medskip\noindent
Expliquons d'abord pourquoi l'énoncé a un sens. Puisque $U$ est le spectre
d'un corps parfait $k$, le schéma $Z$ est géométriquement réduit
sur $k$ (au moyen de l'une ou l'autre projection), voir
Variétés, lemme \ref{varieties-lemma-perfect-reduced}.
Le schéma $Z \times_{\text{pr}_1, U, \text{pr}_0} Z \subset Z$ est donc
réduit, voir
Variétés, lemme \ref{varieties-lemma-geometrically-reduced-any-base-change}.
Ainsi, d'après le
lemme \ref{more-groupoids-lemma-groupoid-on-field-image},
nous voyons que $\text{pr}_{02}$ induit un morphisme
$Z \times_{\text{pr}_1, U, \text{pr}_0} Z \to Z$.
Enfin, il est clair que $\Delta_{U/S}$ se factorise par $Z$
et que l'application
$\sigma : U \times_S U \to U \times_S U$, $(x, y) \mapsto (y, x)$
préserve $Z$. Puisque
$(U, U \times_S U, \text{pr}_0, \text{pr}_1, \text{pr}_{02},
\Delta_{U/S}, \sigma)$
satisfait aux axiomes d'un groupoïde, il s'ensuit qu'après restriction
à $Z$, ces morphismes satisfont aux axiomes.
\end{proof}

\begin{lemma}
\label{more-groupoids-lemma-quasi-compact-groupoid-on-field-image}
Soit $S$ un schéma. Soit $(U, R, s, t, c)$ un schéma en groupoïdes
sur $S$. Supposons que $U$ soit le spectre d'un corps et
supposons que $R$ soit quasi-compact (de manière équivalente, $s, t$ sont quasi-compacts).
Soit $Z \subset U \times_S U$ l'image schématique (voir
Morphismes, définition \ref{morphisms-definition-scheme-theoretic-image})
de $j = (t, s) : R \to U \times_S U$.
Alors
$$
(U, Z, \text{pr}_0|_Z, \text{pr}_1|_Z,
\text{pr}_{02}|_{Z \times_{\text{pr}_1, U, \text{pr}_0} Z})
$$
est un schéma en groupoïdes sur $S$.
\end{lemma}

\begin{proof}
Puisque $(U, U \times_S U, \text{pr}_1, \text{pr}_0, \text{pr}_{02})$
est un schéma en groupoïdes sur $S$, c'est un cas particulier du
lemme \ref{more-groupoids-lemma-quasi-compact-map-groupoids-on-field-image}.
Mais nous pouvons aussi le démontrer directement comme suit.

\medskip\noindent
La difficulté principale consiste à montrer que
$\text{pr}_{02}|_{Z \times_{\text{pr}_1, U, \text{pr}_0} Z}$
a son image dans $Z$.
Écrivons $U = \Spec(k)$. Notons
$R_s$ (resp.\ $Z_s$, resp.\ $U^2_s$) le schéma
$R$ (resp.\ $Z$, resp.\ $U \times_S U$) considéré comme un schéma sur $k$ au moyen de
$s$ (resp.\ $\text{pr}_1|_Z$, resp.\ $\text{pr}_1$).
De même, notons
${}_tR$ (resp.\ ${}_tZ$, resp.\ ${}_tU^2$) le schéma
$R$ (resp.\ $Z$, resp.\ $U \times_S U$) considéré comme un schéma sur $k$ au moyen de
$t$ (resp.\ $\text{pr}_0|_Z$, resp.\ $\text{pr}_0$).
Le morphisme $j$ induit des morphismes de schémas
$j_s : R_s \to U^2_s$ et ${}_tj : {}_tR \to {}_tU^2$ sur $k$.
Considérons le diagramme commutatif
$$
\xymatrix{
R_s \times_k {}_tR \ar[r]^c \ar[d]_{j_s \times {}_tj} &  R \ar[d]^j \\
U^2_s \times_k {}_tU^2 \ar[r] & U \times_S U
}
$$
D'après le
lemme \ref{varieties-lemma-scheme-theoretic-image-product-map} du chapitre Variétés,
nous voyons que l'image schématique de $j_s \times {}_tj$ est
$Z_s \times_k {}_tZ$. Par commutativité du diagramme, nous
concluons que $Z_s \times_k {}_tZ$ a son image dans $Z$ par la flèche
horizontale inférieure. Comme dans la démonstration du
lemme \ref{more-groupoids-lemma-groupoid-on-perfect-field-image},
on a aussi $\sigma(Z) \subset Z$ et
$\Delta_{U/S}$ se factorise par $Z$. Nous concluons donc comme dans la
démonstration de ce lemme.
\end{proof}







\section{Tranches de schémas en groupoïdes}
\label{more-groupoids-section-slicing}

\noindent
Le lemme suivant montre que nous pouvons prendre une tranche d'un schéma en groupoïdes de Cohen-Macaulay
afin de réduire la dimension des fibres, pourvu que la
dimension du stabilisateur soit petite. C'est une étape essentielle du
processus d'amélioration d'une présentation donnée d'un champ quotient.

\begin{situation}
\label{more-groupoids-situation-slice}
Soit $S$ un schéma.
Soit $(U, R, s, t, c)$ un schéma en groupoïdes sur $S$.
Soit $g : U' \to U$ un morphisme de schémas.
Soit $u \in U$ un point, et soit $u' \in U'$ un point tel que
$g(u') = u$. Ces données étant fixées, notons $(U', R', s', t', c')$
la restriction de $(U, R, s, t, c)$ par le morphisme $g$.
Notons $G \to U$ le schéma en groupes stabilisateur de $R$, qui
est un sous-schéma localement fermé de $R$.
Notons $h$ la composée
$$
h = s \circ \text{pr}_1 : U' \times_{g, U, t} R \longrightarrow U.
$$
Notons $F_u = s^{-1}(u)$ (fibre schématique), et $G_u$ la
fibre schématique de $G$ au-dessus de $u$.
De même, pour $R'$, notons $F'_{u'} = (s')^{-1}(u')$.
Comme $g(u') = u$, on a
$$
F'_{u'} = h^{-1}(u) \times_{\Spec(\kappa(u))} \Spec(\kappa(u')).
$$
Le point $e(u) \in R$ peut être considéré comme un point de $G_u$ et de $F_u$, et
$e'(u')$ est un point de $R'$ (resp.\ $G'_{u'}$, resp.\ $F'_{u'}$) qui s'envoie
sur $e(u)$ dans $R$ (resp.\ $G_u$, resp.\ $F_u$).
\end{situation}

\begin{lemma}
\label{more-groupoids-lemma-slice}
Soit $S$ un schéma.
Soit $(U, R, s, t, c, e, i)$ un schéma en groupoïdes sur $S$.
Soit $G \to U$ le schéma en groupes stabilisateur.
Supposons que $s$ et $t$ soient de Cohen-Macaulay et localement de présentation finie.
Soit $u \in U$ un point de type fini du schéma $U$, voir
Morphismes, définition \ref{morphisms-definition-finite-type-point}.
Avec les notations de la
situation \ref{more-groupoids-situation-slice},
posons
$$
d_1 = \dim(G_u), \quad
d_2 = \dim_{e(u)}(F_u).
$$
Si $d_2 > d_1$, alors il existe un schéma affine $U'$
et un morphisme $g : U' \to U$ tels que (avec les notations de la
situation \ref{more-groupoids-situation-slice})
\begin{enumerate}
\item $g$ est une immersion,
\item $u \in U'$,
\item $g$ est localement de présentation finie,
\item le morphisme $h : U' \times_{g, U, t} R \longrightarrow U$
est de Cohen-Macaulay au point $(u, e(u))$, et
\item on a $\dim_{e'(u)}(F'_u) = d_2 - 1$.
\end{enumerate}
\end{lemma}

\begin{proof}
Soit $\Spec(A) \subset U$ un voisinage affine de $u$
tel que $u$ corresponde à un point fermé de $U$, voir
Morphismes, lemme \ref{morphisms-lemma-identify-finite-type-points}.
Soit $\Spec(B) \subset R$ un voisinage affine de $e(u)$
dont l'image par $j$ est contenue dans l'ouvert
$\Spec(A) \times_S \Spec(A) \subset U \times_S U$.
Soit $\mathfrak m \subset A$ l'idéal maximal correspondant à $u$.
Soit $\mathfrak q \subset B$ l'idéal premier correspondant à $e(u)$.
Diagrammes :
$$
\vcenter{
\xymatrix{
B & A \ar[l]^s \\
A \ar[u]^t
}
}
\quad\text{et}\quad
\vcenter{
\xymatrix{
B_{\mathfrak q} & A_{\mathfrak m} \ar[l]^s \\
A_{\mathfrak m} \ar[u]^t
}
}
$$
Notons que les deux applications induites
$s, t : \kappa(\mathfrak m) \to \kappa(\mathfrak q)$
sont égales et sont des isomorphismes puisque $s \circ e = t \circ e = \text{id}_U$.
En particulier, nous voyons que $\mathfrak q$
est aussi un idéal maximal. Les homomorphismes d'anneaux $s, t : A \to B$ sont
de présentation finie et plats. Par hypothèse, l'anneau
$$
\mathcal{O}_{F_u, e(u)} = B_{\mathfrak q}/s(\mathfrak m)B_{\mathfrak q}
$$
est de Cohen-Macaulay de dimension $d_2$. L'égalité des dimensions résulte de
Morphismes, lemme \ref{morphisms-lemma-dimension-fibre-at-a-point}.

\medskip\noindent
Soit $R''$ la restriction de $R$ à $u = \Spec(\kappa(u))$
par le morphisme $\Spec(\kappa(u)) \to U$.
Comme $u \to U$ est localement de type fini,
nous voyons que $(\Spec(\kappa(u)), R'', s'', t'', c'')$
est un schéma en groupoïdes avec $s'', t''$ localement de type fini, voir
le lemme \ref{more-groupoids-lemma-restrict-preserves-type}.
D'après le
lemme \ref{more-groupoids-lemma-groupoid-on-field-dimension-equal-stabilizer},
cela implique que $\dim(G'') = \dim(R'')$. On a aussi
$\dim(R'') = \dim_{e''}(R'') = \dim(\mathcal{O}_{R'', e''})$, voir
le lemme \ref{more-groupoids-lemma-groupoid-on-field-locally-finite-type-dimension}.
D'après
Groupoïdes, lemme \ref{groupoids-lemma-restrict-stabilizer},
on a $G'' = G_u$. Nous concluons donc que
$\dim(\mathcal{O}_{R'', e''}) = d_1$.

\medskip\noindent
En tant que schéma, $R''$ est
$$
R'' =
R \times_{(U \times_S U)}
\Big(
\Spec(\kappa(\mathfrak m)) \times_S \Spec(\kappa(\mathfrak m))
\Big)
$$
Par conséquent, un voisinage ouvert affine de $e''$ est le spectre de l'anneau
$$
B \otimes_{(A \otimes A)} (\kappa(\mathfrak m) \otimes \kappa(\mathfrak m))
=
B/s(\mathfrak m)B + t(\mathfrak m)B
$$
Nous en concluons que
$$
\mathcal{O}_{R'', e''} =
B_{\mathfrak q}/s(\mathfrak m)B_{\mathfrak q} + t(\mathfrak m)B_{\mathfrak q}
$$
et nous savons donc maintenant que cet anneau est de dimension $d_1$.

\medskip\noindent
Nous affirmons que cela implique que l'on peut trouver
un élément $f \in \mathfrak m$ tel que
$$
\dim(B_{\mathfrak q}/(s(\mathfrak m)B_{\mathfrak q} + fB_{\mathfrak q}) < d_2
$$
En effet, supposons que $\mathfrak n_j \supset s(\mathfrak m)B_{\mathfrak q}$,
$j = 1, \ldots, m$, correspondent aux idéaux premiers minimaux de l'anneau local
$B_{\mathfrak q}/s(\mathfrak m)B_{\mathfrak q}$.
Ils sont en nombre fini puisque cet anneau est noethérien (car il est essentiellement
de type fini sur un corps -- mais aussi parce qu'un anneau de Cohen-Macaulay est
noethérien). Par la propriété de Cohen-Macaulay, on a
$\dim(B_{\mathfrak q}/\mathfrak n_j) = d_2$, par exemple d'après
Algèbre, lemme \ref{algebra-lemma-CM-dim-formula}.
Notons que
$\dim(B_{\mathfrak q}/(\mathfrak n_j + t(\mathfrak m)B_{\mathfrak q}))
\leq d_1$
car c'est un quotient de l'anneau
$\mathcal{O}_{R'', e''} =
B_{\mathfrak q}/s(\mathfrak m)B_{\mathfrak q} + t(\mathfrak m)B_{\mathfrak q}$
qui est de dimension $d_1$. Comme $d_1 < d_2$, cela implique que
$\mathfrak m \not \subset t^{-1}(\mathfrak n_i)$.
Par le lemme d'évitement des idéaux premiers, voir
Algèbre, lemme \ref{algebra-lemma-silly},
nous pouvons trouver $f \in \mathfrak m$ tel que $t(f) \not \in \mathfrak n_j$ pour
$j = 1, \ldots, m$. Pour ce choix de $f$, on a
l'inégalité affichée ci-dessus, voir
Algèbre, lemme \ref{algebra-lemma-one-equation}.

\medskip\noindent
Posons $A' = A/fA$ et $U' = \Spec(A')$. Il est alors clair que
$U' \to U$ est une immersion, localement de présentation finie,
et que $u \in U'$. Ainsi, les assertions (1), (2) et (3) du lemme sont vérifiées.
Le morphisme
$$
U' \times_{g, U, t} R \longrightarrow U
$$
se factorise par $\Spec(A)$ et correspond à l'homomorphisme d'anneaux
$$
\xymatrix{
B/t(f)B \ar@{=}[r] & A/(f) \otimes_{A, t} B & A \ar[l]_-s
}
$$
Nous voyons maintenant que $t(f)$ n'est pas un diviseur de zéro dans
$B_{\mathfrak q}/s(\mathfrak m)B_{\mathfrak q}$, car c'est un anneau de
Cohen-Macaulay de dimension positive et $f$ n'appartient à
aucun idéal premier minimal, voir par exemple
Algèbre, lemme \ref{algebra-lemma-reformulate-CM}.
Par conséquent, d'après
Algèbre, lemme \ref{algebra-lemma-grothendieck-general},
nous concluons que $s : A_{\mathfrak m} \to B_{\mathfrak q}/t(f)B_{\mathfrak q}$
est plat et a pour anneau de la fibre
$B_{\mathfrak q}/(s(\mathfrak m)B_{\mathfrak q} + t(f)B_{\mathfrak q})$,
qui est de Cohen-Macaulay, encore d'après
Algèbre, lemme \ref{algebra-lemma-reformulate-CM}.
Cela implique l'assertion (4) du lemme.
Pour voir l'assertion (5), notons que, d'après le diagramme (\ref{more-groupoids-equation-restriction}),
la fibre $F'_u$ est égale à la fibre de $h$ au-dessus de $u$.
Par conséquent,
$\dim_{e'(u)}(F'_u) =
\dim(B_{\mathfrak q}/(s(\mathfrak m)B_{\mathfrak q} + t(f)B_{\mathfrak q}))$
d'après
Morphismes, lemme \ref{morphisms-lemma-dimension-fibre-at-a-point},
et la dimension de cet anneau est $d_2 - 1$ d'après
Algèbre, lemme \ref{algebra-lemma-reformulate-CM},
une fois encore. Cela démontre la dernière assertion du lemme et achève la démonstration.
\end{proof}

\noindent
Maintenant que nous savons prendre une tranche, nous pouvons combiner ce résultat avec ce qui précède
pour obtenir le résultat ``optimal'' suivant. Il est optimal au sens où,
puisque $G_u$ est un sous-schéma localement fermé de $F_u$, on a toujours
l'inégalité $\dim(G_u) = \dim_{e(u)}(G_u) \leq \dim_{e(u)}(F_u)$ ; il
n'est donc pas possible de pousser plus loin la prise d'une tranche que dans le lemme.

\begin{lemma}
\label{more-groupoids-lemma-max-slice}
Soit $S$ un schéma.
Soit $(U, R, s, t, c, e, i)$ un schéma en groupoïdes sur $S$.
Soit $G \to U$ le schéma en groupes stabilisateur.
Supposons que $s$ et $t$ soient de Cohen-Macaulay et localement de présentation finie.
Soit $u \in U$ un point de type fini du schéma $U$, voir
Morphismes, définition \ref{morphisms-definition-finite-type-point}.
Avec les notations de la
situation \ref{more-groupoids-situation-slice},
il existe un schéma affine $U'$ et un morphisme $g : U' \to U$ tels que
\begin{enumerate}
\item $g$ est une immersion,
\item $u \in U'$,
\item $g$ est localement de présentation finie,
\item le morphisme $h : U' \times_{g, U, t} R \longrightarrow U$
est de Cohen-Macaulay et localement de présentation finie,
\item les morphismes $s', t' : R' \to U'$ sont de Cohen-Macaulay et
localement de présentation finie, et
\item $\dim_{e(u)}(F'_u) = \dim(G'_u)$.
\end{enumerate}
\end{lemma}

\begin{proof}
Comme $s$ est localement de présentation finie, le schéma $F_u$ est
localement de type fini sur $\kappa(u)$. Par conséquent,
$\dim_{e(u)}(F_u) < \infty$ et nous pouvons raisonner par récurrence sur
$\dim_{e(u)}(F_u)$.

\medskip\noindent
Si $\dim_{e(u)}(F_u) = \dim(G_u)$, il n'y a rien à démontrer.
Supposons $\dim_{e(u)}(F_u) > \dim(G_u)$. Cela signifie que le
lemme \ref{more-groupoids-lemma-slice}
s'applique, et nous trouvons un morphisme $g : U' \to U$ qui possède
les propriétés (1), (2), (3), tandis qu'au lieu de (6), nous avons
$\dim_{e(u)}(F'_u) < \dim_{e(u)}(F_u)$,
et qu'au lieu de (4) et (5), nous savons que la composée
$$
h = s \circ \text{pr}_1 : U' \times_{g, U, t} R \longrightarrow U
$$
est de Cohen-Macaulay au point $(u, e(u))$. Nous appliquons la
remarque \ref{more-groupoids-remark-local-source-apply}
et obtenons un sous-schéma ouvert $U'' \subset U'$ tel que
$U'' \times_{g, U, t} R \subset U' \times_{g, U, t} R$
soit le plus grand sous-schéma ouvert sur lequel $h$ est de Cohen-Macaulay.
Puisque $(u, e(u)) \in U'' \times_{g, U, t} R$, nous voyons que $u \in U''$.
Nous pouvons donc remplacer $U'$ par $U''$ et supposer qu'en fait $h$ est
de Cohen-Macaulay partout ! D'après le
lemme \ref{more-groupoids-lemma-restrict-property},
nous concluons que $s', t'$ sont localement de
présentation finie et de Cohen-Macaulay (utiliser
Morphismes, lemme \ref{morphisms-lemma-base-change-finite-presentation},
et
Compléments sur les morphismes, lemme \ref{more-morphisms-lemma-base-change-CM}).

\medskip\noindent
Par construction, $\dim_{e'(u)}(F'_u) < \dim_{e(u)}(F_u)$,
de sorte que nous pouvons appliquer l'hypothèse de récurrence à $(U', R', s', t', c')$
et au point $u \in U'$. Notons que $u$ est également un point de type fini
de $U'$ (on peut par exemple le voir en utilisant la caractérisation des
points de type fini donnée dans
Morphismes, lemme \ref{morphisms-lemma-identify-finite-type-points}).
Soient $g' : U'' \to U'$ et $(U'', R'', s'', t'', c'')$ la solution
du problème correspondant partant de $(U', R', s', t', c')$
et du point $u \in U'$. Nous affirmons que la composée
$$
g'' = g \circ g' : U'' \longrightarrow U
$$
est une solution du problème initial. Les propriétés (1), (2), (3), (5)
et (6) sont immédiates. Pour voir (4), notons que le morphisme
$$
h'' = s \circ \text{pr}_1 : U'' \times_{g'', U, t} R \longrightarrow U
$$
est localement de présentation finie et de Cohen-Macaulay par application du
lemme \ref{more-groupoids-lemma-double-restrict-property}
(utiliser
Compléments sur les morphismes, lemme \ref{more-morphisms-lemma-CM-local-source-and-target},
pour voir que la propriété d'être de Cohen-Macaulay est locale pour la topologie fppf sur la base).
\end{proof}

\noindent
Lorsque les fibres du schéma en groupes stabilisateur sont de dimension 0,
cela donne le lemme de tranche suivant.

\begin{lemma}
\label{more-groupoids-lemma-max-slice-quasi-finite}
Soit $S$ un schéma.
Soit $(U, R, s, t, c, e, i)$ un schéma en groupoïdes sur $S$.
Soit $G \to U$ le schéma en groupes stabilisateur.
Supposons que $s$ et $t$ soient de Cohen-Macaulay et localement de présentation finie.
Soit $u \in U$ un point de type fini du schéma $U$, voir
Morphismes, définition \ref{morphisms-definition-finite-type-point}.
Supposons que $G \to U$ soit localement quasi-fini.
Avec les notations de la
situation \ref{more-groupoids-situation-slice},
il existe un schéma affine $U'$ et un morphisme $g : U' \to U$ tels que
\begin{enumerate}
\item $g$ est une immersion,
\item $u \in U'$,
\item $g$ est localement de présentation finie,
\item le morphisme $h : U' \times_{g, U, t} R \longrightarrow U$
est plat, localement de présentation finie et localement quasi-fini, et
\item les morphismes $s', t' : R' \to U'$ sont plats,
localement de présentation finie et localement quasi-finis.
\end{enumerate}
\end{lemma}

\begin{proof}
Prenons $g : U' \to U$ comme dans le
lemme \ref{more-groupoids-lemma-max-slice}.
Puisque $h^{-1}(u) = F'_u$, nous voyons que $h$ est de dimension relative
$\leq 0$ en $(u, e(u))$. Par conséquent, d'après la
remarque \ref{more-groupoids-remark-local-source-apply},
nous obtenons un sous-schéma ouvert $U'' \subset U'$ tel que
$u \in U''$ et que $U'' \times_{g, U, t} R$ soit le plus grand sous-schéma ouvert
de $U' \times_{g, U, t} R$ sur lequel $h$ est de dimension relative $\leq 0$.
Après avoir remplacé $U'$ par $U''$, nous voyons que $h$ est de dimension relative $\leq 0$.
Cela implique que $h$ est localement quasi-fini d'après
Morphismes, lemme \ref{morphisms-lemma-locally-quasi-finite-rel-dimension-0}.
Comme il est encore localement de présentation finie et de Cohen-Macaulay, nous voyons
qu'il est plat, localement de présentation finie et localement quasi-fini,
c'est-à-dire que (4) ci-dessus est vérifiée. Cela implique que $s'$ est plat, localement de
présentation finie et localement quasi-fini en tant que changement de base de $h$, voir
le lemme \ref{more-groupoids-lemma-restrict-property}.
\end{proof}







\section{Localisation \'etale des schémas en groupoïdes}
\label{more-groupoids-section-etale-localize}

\noindent
Dans cette section, nous commençons à appliquer aux schémas en groupoïdes les techniques de localisation \'etale de
Compléments sur les morphismes, section \ref{more-morphisms-section-etale-localization}.
On trouvera des résultats plus avancés de ce type dans
Compléments sur les groupoïdes dans les espaces,
section \ref{spaces-more-groupoids-section-etale-localize}.
Le lemme \ref{more-groupoids-lemma-quasi-finite-over-base-j-proper}
sera utilisé pour démontrer des résultats sur les espaces algébriques
séparés et quasi-finis sur un schéma, à savoir
Morphismes d'espaces, proposition
\ref{spaces-morphisms-proposition-locally-quasi-finite-separated-over-scheme}
et son corollaire
Morphismes d'espaces, lemme
\ref{spaces-morphisms-lemma-locally-quasi-finite-separated-representable}.

\begin{lemma}
\label{more-groupoids-lemma-quasi-finite-over-base}
Soit $S$ un schéma.
Soit $(U, R, s, t, c)$ un schéma en groupoïdes sur $S$.
Soit $p \in S$ un point, et soit $u \in U$ un point situé au-dessus de $p$.
Supposons que
\begin{enumerate}
\item $U \to S$ soit localement de type fini,
\item $U \to S$ soit quasi-fini en $u$,
\item $U \to S$ soit séparé,
\item $R \to S$ soit séparé,
\item $s$, $t$ soient plats et localement de présentation finie, et
\item $s^{-1}(\{u\})$ soit fini.
\end{enumerate}
Alors il existe un voisinage \'etale $(S', p') \to (S, p)$ tel que
$\kappa(p) = \kappa(p')$ et un diagramme de changement de base
$$
\xymatrix{
R' \amalg W'
\ar@{=}[r] &
S' \times_S R
\ar[r] \ar@<2ex>[d]^{s'} \ar@<-2ex>[d]_{t'} &
R \ar@<1ex>[d]^s \ar@<-1ex>[d]_t \\
U' \amalg W
\ar@{=}[r] &
S' \times_S U
\ar[r] \ar[d] &
U \ar[d] \\
 &
S' \ar[r] &
S
}
$$
où les signes d'égalité sont des décompositions en sous-schémas ouverts et
fermés telles que
\begin{enumerate}
\item[(a)] il existe un point $u'$ de $U'$ qui s'envoie sur $u$ dans $U$,
\item[(b)] la fibre $(U')_{p'}$ est égale à $t'\big((s')^{-1}(\{u'\})\big)$
ensemblistement,
\item[(c)] la fibre $(R')_{p'}$ est égale à $(s')^{-1}\big((U')_{p'}\big)$
ensemblistement,
\item[(d)] les schémas $U'$ et $R'$ sont finis sur $S'$,
\item[(e)] on a $s'(R') \subset U'$ et $t'(R') \subset U'$,
\item[(f)] on a
$c'(R' \times_{s', U', t'} R') \subset R'$
où $c'$ est le changement de base de $c$, et
\item[(g)] les morphismes $s', t', c'$ déterminent une structure de groupoïde
en prenant le système
$(U', R', s'|_{R'}, t'|_{R'}, c'|_{R' \times_{s', U', t'} R'})$.
\end{enumerate}
\end{lemma}

\begin{proof}
Notons $f : U \to S$ le morphisme structural de $U$.
D'après l'hypothèse (6), on peut écrire $s^{-1}(\{u\}) = \{r_1, \ldots, r_n\}$.
Comme cet ensemble est fini, on voit que $s$ est quasi-fini en chacun de
ces antécédents, en nombre fini, voir
Morphismes, Lemme \ref{morphisms-lemma-finite-fibre}.
On voit donc que $f \circ s : R \to S$ est quasi-fini en chaque $r_i$
(Morphismes, Lemme \ref{morphisms-lemma-composition-quasi-finite}).
Ainsi, $r_i$ est isolé dans la fibre $R_p$, voir
Morphismes, Lemme \ref{morphisms-lemma-quasi-finite-at-point-characterize}.
Écrivons $t(\{r_1, \ldots, r_n\}) = \{u_1, \ldots, u_m\}$.
Notons qu'il peut arriver que $m < n$ et que
$u \in \{u_1, \ldots, u_m\}$.
Comme $t$ est plat et localement de présentation finie,
le morphisme de fibres $t_p : R_p \to U_p$ est plat et localement de
présentation finie (Morphismes,
Lemmes \ref{morphisms-lemma-base-change-flat} et
\ref{morphisms-lemma-base-change-finite-presentation}),
donc ouvert (Morphismes,
Lemme \ref{morphisms-lemma-fppf-open}).
Le fait que chaque $r_i$ soit isolé dans $R_p$ implique que
chaque $u_j = t(r_i)$ est isolé dans $U_p$. En utilisant de nouveau
Morphismes, Lemme \ref{morphisms-lemma-quasi-finite-at-point-characterize},
on voit que $f$ est quasi-fini en $u_1, \ldots, u_m$.

\medskip\noindent
Notons $F_u = s^{-1}(u)$ et $F_{u_j} = s^{-1}(u_j)$ les fibres schématiques.
Notons que $F_u$ est fini sur $\kappa(u)$, car il est localement de type fini
sur $\kappa(u)$ et possède un nombre fini de points (cela résulte par exemple
du résultat beaucoup plus général
Morphismes, Lemme
\ref{morphisms-lemma-locally-quasi-finite-qc-source-universally-bounded}).
D'après
le Lemme \ref{more-groupoids-lemma-two-fibres},
on voit que $F_u$ et $F_{u_j}$ deviennent isomorphes après extension à un corps
qui est une extension commune de $\kappa(u)$ et $\kappa(u_j)$. On voit donc
que $F_{u_j}$ est fini sur $\kappa(u_j)$. En particulier, on voit que
$s^{-1}(\{u_j\})$ est un ensemble fini pour chaque $j = 1, \ldots, m$.
Ainsi, les hypothèses (2) et (6) sont aussi vérifiées pour chaque $u_j$
(on a vu plus haut que $U \to S$ est quasi-fini en $u_j$).
L'argument du premier paragraphe s'applique donc à chaque $u_j$,
et l'on voit que $R \to U$ est quasi-fini en chacun des points de
$$
\{r_1, \ldots, r_N\} = s^{-1}(\{u_1, \ldots, u_m\})
$$
Notons que $t(\{r_1, \ldots, r_N\}) = \{u_1, \ldots, u_m\}$ et
$t^{-1}(\{u_1, \ldots, u_m\}) = \{r_1, \ldots, r_N\}$
car $R$ est un groupoïde\footnote{Explication dans le langage des groupoïdes :
L'ensemble initial $\{r_1, \ldots, r_n\}$ était l'ensemble des flèches de
source $u$. L'ensemble $\{u_1, \ldots, u_m\}$ était l'ensemble des objets
isomorphes à $u$. Et $\{r_1, \ldots, r_N\}$ est l'ensemble de toutes les flèches
entre tous les objets équivalents à $u$.}. De plus, on a
$\text{pr}_0(c^{-1}(\{r_1, \ldots, r_N\})) = \{r_1, \ldots, r_N\}$
et
$\text{pr}_1(c^{-1}(\{r_1, \ldots, r_N\})) = \{r_1, \ldots, r_N\}$.
De même, on obtient $e(\{u_1, \ldots, u_m\}) \subset \{r_1, \ldots, r_N\}$
et $i(\{r_1, \ldots, r_N\}) = \{r_1, \ldots, r_N\}$.

\medskip\noindent
On peut appliquer
Compléments sur les morphismes,
Lemme \ref{more-morphisms-lemma-etale-splits-off-quasi-finite-part-technical}
aux couples
$(U \to S, \{u_1, \ldots, u_m\})$ et
$(R \to S, \{r_1, \ldots, r_N\})$
pour obtenir un voisinage \'etale $(S', p') \to (S, p)$
qui induit une identification $\kappa(p) = \kappa(p')$
tel que $S' \times_S U$ et $S' \times_S R$ se décomposent sous la forme
$$
S' \times_S U = U' \amalg W, \quad
S' \times_S R = R' \amalg W'
$$
où $U' \to S'$ est fini et $(U')_{p'}$ s'envoie bijectivement sur
$\{u_1, \ldots, u_m\}$, et où $R' \to S'$ est fini et
$(R')_{p'}$ s'envoie bijectivement sur $\{r_1, \ldots, r_N\}$.
De plus, aucun point de $W_{p'}$ (resp.\ de $(W')_{p'}$) n'a pour image
aucun des points $u_j$ (resp.\ $r_i$). À ce stade, (a), (b), (c) et (d)
du lemme sont satisfaites. En outre, les inclusions de (e) et (f) valent
sur les fibres au-dessus de $p'$, c'est-à-dire $s'((R')_{p'}) \subset (U')_{p'}$,
$t'((R')_{p'}) \subset (U')_{p'}$, et
$c'((R' \times_{s', U', t'} R')_{p'}) \subset (R')_{p'}$.

\medskip\noindent
Nous affirmons que l'on peut remplacer $S'$ par un voisinage ouvert de Zariski
de $p'$ de sorte que les inclusions de (e) et (f) soient vérifiées.
Considérons par exemple l'ensemble $E = (s'|_{R'})^{-1}(W)$.
Il est ouvert et fermé dans $R'$ et ne contient aucun point
de $R'$ situé au-dessus de $p'$. Comme $R' \to S'$ est fermé,
après avoir remplacé $S'$ par $S' \setminus (R' \to S')(E)$, on arrive à une
situation où $E$ est vide. Autrement dit, $s'$ envoie $R'$ dans $U'$.
Notons que cette propriété est conservée lorsque l'on rétrécit encore $S'$.
De même, on peut faire en sorte que $t'$ envoie $R'$ dans $U'$.
À ce stade, (e) est vérifiée. De la même manière, considérons l'ensemble
$E = (c'|_{R' \times_{s', U', t'} R'})^{-1}(W')$.
Il est ouvert et fermé dans le schéma $R' \times_{s', U', t'} R'$,
qui est fini sur $S'$, et ne contient aucun point situé
au-dessus de $p'$. Après avoir remplacé $S'$ par
$S' \setminus (R' \times_{s', U', t'} R' \to S')(E)$,
on arrive donc à une situation où $E$ est vide. Autrement dit, on obtient
l'inclusion de (f). On peut répéter l'argument avec l'identité
$e' : S' \times_S U \to S' \times_S R$ et l'inverse
$i' : S' \times_S R \to S' \times_S R$, afin de pouvoir supposer
(après avoir encore rétréci $S'$) que $(e'|_{U'})^{-1}(W') = \emptyset$
et $(i'|_{R'})^{-1}(W') = \emptyset$.

\medskip\noindent
À ce stade, on voit que l'on peut considérer la structure
$$
(U', R', s'|_{R'}, t'|_{R'}, c'|_{R' \times_{t', U', s'} R'},
e'|_{U'}, i'|_{R'}).
$$
Les axiomes d'un schéma en groupoïdes sur $S'$ sont vérifiés
parce qu'ils le sont pour le schéma en groupoïdes
$(S' \times_S U, S' \times_S R, s', t', c', e', i')$.
\end{proof}

\begin{lemma}
\label{more-groupoids-lemma-quasi-finite-over-base-j-proper}
Soit $S$ un schéma.
Soit $(U, R, s, t, c)$ un schéma en groupoïdes sur $S$.
Soient $p \in S$ un point et $u \in U$ un point situé au-dessus de $p$.
Supposons que les hypothèses (1) -- (6) du
Lemme \ref{more-groupoids-lemma-quasi-finite-over-base}
soient vérifiées, ainsi que
\begin{enumerate}
\item[(7)] $j : R \to U \times_S U$ est universellement fermé\footnote{Compte tenu
des autres conditions, cela équivaut à demander que $j$ soit propre.}.
\end{enumerate}
Alors on peut choisir $(S', p') \to (S, p)$, des décompositions
$S' \times_S U = U' \amalg W$ et $S' \times_S R = R' \amalg W'$
et $u' \in U'$ tels que (a) -- (g) du
Lemme \ref{more-groupoids-lemma-quasi-finite-over-base}
soient vérifiées, ainsi que
\begin{enumerate}
\item[(h)] $R'$ est la restriction de $S' \times_S R$ à $U'$.
\end{enumerate}
\end{lemma}

\begin{proof}
On applique le Lemme \ref{more-groupoids-lemma-quasi-finite-over-base} au
groupoïde $(U, R, s, t, c)$ sur le schéma $S$, avec les points $p$ et $u$.
On obtient donc un voisinage \'etale
$(S', p') \to (S, p)$ et des décompositions en unions disjointes
$$
S' \times_S U = U' \amalg W, \quad
S' \times_S R = R' \amalg W'
$$
ainsi que $u' \in U'$ satisfaisant les conclusions (a), (b), (c), (d), (e), (f) et (g).
On peut rétrécir $S'$ à un voisinage plus petit de $p'$ sans
affecter les conclusions (a) -- (g). Nous allons montrer que, pour un
rétrécissement convenable, la conclusion (h) est également vérifiée.
Notons $j'$ le changement de base de $j$ à $S'$.
D'après la conclusion (e), il est clair que
$$
j'^{-1}(U' \times_{S'} U') = R' \amalg Rest
$$
pour une certaine partie $Rest$ ouverte et fermée. Comme $U' \to S'$ est fini
d'après la conclusion (d), on voit que $U' \times_{S'} U'$ est fini sur $S'$.
Comme $j$ est universellement fermé, $j'$ l'est aussi, et
donc $j'|_{Rest}$ est lui aussi universellement fermé. D'après les conclusions
(b) et (c), on voit que la fibre de
$$
(U' \times_{S'} U' \to S') \circ j'|_{Rest} :
Rest
\longrightarrow
S'
$$
au-dessus de $p'$ est vide. Ainsi, puisque $Rest \to S'$ est fermé comme composé
de morphismes fermés, après avoir remplacé $S'$ par
$S' \setminus \Im(Rest \to S')$, on peut supposer que
$Rest = \emptyset$. C'est exactement la condition disant que $R'$ est
la restriction de $S' \times_S R$ au sous-schéma ouvert
$U' \subset S' \times_S U$, voir
Groupoïdes, Lemme \ref{groupoids-lemma-restrict-groupoid-relation}
et sa démonstration.
\end{proof}






\section{Groupoïdes finis}
\label{more-groupoids-section-finite-groupoids}

\noindent
Un schéma en groupoïdes $(U, R, s, t, c)$ est parfois dit {\it fini} si les
morphismes $s$ et $t$ sont finis. Cette terminologie peut prêter à confusion, car elle
n'implique pas que $U$, $R$ ou le faisceau quotient $U/R$ soient finis sur quoi que ce soit.

\begin{lemma}
\label{more-groupoids-lemma-finite-stratify}
Soit $(U, R, s, t, c)$ un schéma en groupoïdes sur un schéma $S$. Supposons $s, t$
finis. Il existe une suite de sous-schémas fermés invariants par $R$
$$
U = Z_0 \supset Z_1 \supset Z_2 \supset \ldots
$$
telle que $\bigcap Z_r = \emptyset$ et que
$s^{-1}(Z_{r - 1}) \setminus s^{-1}(Z_r) \to Z_{r - 1} \setminus Z_r$
soit fini localement libre de rang $r$.
\end{lemma}

\begin{proof}
Soit $\{Z_r\}$ la stratification de $U$ donnée par les idéaux de Fitting
du module quasi-cohérent de type fini $s_*\mathcal{O}_R$. Voir
Diviseurs, Lemme \ref{divisors-lemma-locally-free-rank-r-pullback}.
Comme l'identité $e : U \to R$ est une section de $s$, on voit que
$s_*\mathcal{O}_R$ contient $\mathcal{O}_S$ comme facteur direct.
Ainsi $U = Z_{-1} = Z_0$ (détails omis).
Comme la formation des idéaux de Fitting commute au changement de base
(Compléments sur l'algèbre, Lemme \ref{more-algebra-lemma-fitting-ideal-basics}),
on trouve que $s^{-1}(Z_r)$ correspond au $r$-ième idéal de Fitting
de $\text{pr}_{1, *}\mathcal{O}_{R \times_{s, U, t} R}$, car
le carré inférieur droit du diagramme (\ref{more-groupoids-equation-diagram}) est cartésien.
En utilisant le fait que le carré inférieur gauche est lui aussi cartésien, on conclut
que $s^{-1}(Z_r) = t^{-1}(Z_r)$ ; autrement dit, $Z_r$ est invariant par $R$.
Le morphisme
$s^{-1}(Z_{r - 1}) \setminus s^{-1}(Z_r) \to Z_{r - 1} \setminus Z_r$
est fini localement libre de rang $r$, car le module
$s_*\mathcal{O}_R$ se restreint en un module fini localement libre de rang $r$
sur $Z_{r - 1} \setminus Z_r$, d'après
Diviseurs, Lemme \ref{divisors-lemma-locally-free-rank-r-pullback}.
\end{proof}

\begin{lemma}
\label{more-groupoids-lemma-finite-flat-over-almost-dense-subscheme}
Soit $(U, R, s, t, c)$ un schéma en groupoïdes sur un schéma $S$. Supposons $s, t$
finis. Il existe un sous-schéma ouvert $W \subset U$ et un sous-schéma fermé
$W' \subset W$ tels que
\begin{enumerate}
\item $W$ et $W'$ soient invariants par $R$,
\item $U = t(s^{-1}(\overline{W}))$ ensemblistement,
\item $W$ soit un épaississement de $W'$, et
\item les morphismes $s'$, $t'$ de la restriction $(W', R', s', t', c')$
soient finis localement libres.
\end{enumerate}
\end{lemma}

\begin{proof}
Considérons la stratification $U = Z_0 \supset Z_1 \supset Z_2 \supset \ldots$
du Lemme \ref{more-groupoids-lemma-finite-stratify}.

\medskip\noindent
Nous allons construire des unions disjointes $W = \coprod_{r \geq 1} W_r$ et
$W' = \coprod_{r \geq 1} W'_r$, chaque $W'_r \to W_r$ étant un épaississement
entre des sous-schémas invariants par $R$ de $U$, de sorte que les morphismes
$s_r', t_r'$ des restrictions $(W_r', R_r', s_r', t_r', c_r')$
soient finis localement libres de rang $r$. Pour commencer, posons
$W_1 = W'_1 = U \setminus Z_1$. C'est un sous-schéma ouvert invariant par $R$
de $U$ ; il est vrai que $W_0$ est un épaississement de $W'_0$,
et les morphismes $s_1'$, $t_1'$ de la
restriction $(W_1', R_1', s_1', t_1', c_1')$ sont des isomorphismes, c'est-à-dire
finis localement libres de rang $1$.
De plus, tout point de $U \setminus Z_1$ appartient à $t(s^{-1}(\overline{W_1}))$.

\medskip\noindent
Supposons avoir trouvé des sous-schémas $W'_r \subset W_r \subset U$ pour $r \leq n$
tels que
\begin{enumerate}
\item $W_1, \ldots, W_n$ soient disjoints,
\item $W_r$ et $W_r'$ soient invariants par $R$,
\item $U \setminus Z_n \subset \bigcup_{r \leq n} t(s^{-1}(\overline{W_r}))$
ensemblistement,
\item $W_r$ soit un épaississement de $W'_r$,
\item les morphismes $s_r'$, $t_r'$ de la restriction
$(W_r', R_r', s_r', t_r', c_r')$ soient finis localement libres de rang $r$.
\end{enumerate}
Posons alors
$$
W_{n + 1} = Z_n \setminus
\left(
Z_{n + 1} \cup \bigcup\nolimits_{r \leq n} t(s^{-1}(\overline{W_r}))
\right)
$$
ensemblistement et
$$
W'_{n + 1} = Z_n \setminus
\left(
Z_{n + 1} \cup \bigcup\nolimits_{r \leq n} t(s^{-1}(\overline{W_r}))
\right)
$$
schématiquement. Alors $W_{n + 1}$ est un sous-schéma ouvert invariant par $R$
de $U$, car $Z_{n + 1} \setminus \overline{U \setminus Z_{n + 1}}$
est ouvert dans $U$ et $\overline{U \setminus Z_{n + 1}}$ est contenu
dans le fermé $\bigcup\nolimits_{r \leq n} t(s^{-1}(\overline{W_r}))$
que l'on retire, par la propriété (3) et le fait que $t$ est un morphisme fermé.
Il est clair que $W'_{n + 1}$ est un sous-schéma fermé
de $W_{n + 1}$ ayant le même espace topologique sous-jacent.
Enfin, les propriétés (1), (2) et (3) sont claires, et la propriété (5) résulte du
Lemme \ref{more-groupoids-lemma-finite-stratify}.

\medskip\noindent
D'après le Lemme \ref{more-groupoids-lemma-finite-stratify}, on a $\bigcap Z_r = \emptyset$.
Tout point de $U$ est donc contenu dans $U \setminus Z_n$
pour un certain $n$. On voit ainsi que
$U = \bigcup_{r \geq 1} t(s^{-1}(\overline{W_r}))$
ensemblistement, et l'on voit que (2) est vérifiée.
Ainsi, les sous-schémas $W' \subset W$ satisfont (1), (2), (3) et (4).
\end{proof}

\noindent
Soit $(U, R, s, t, c)$ un schéma en groupoïdes. Étant donné un point $u \in U$,
l'{\it orbite sous $R$} de $u$ est le sous-ensemble $t(s^{-1}(\{u\}))$ de $U$.

\begin{lemma}
\label{more-groupoids-lemma-finite-flat-over-almost-dense-subscheme-addendum}
Dans le Lemme \ref{more-groupoids-lemma-finite-flat-over-almost-dense-subscheme},
supposons en outre que $s$ et $t$ soient de présentation finie.
Alors
\begin{enumerate}
\item le morphisme $W' \to W$ est de présentation finie, et
\item si $u \in U$ est un point dont l'orbite sous $R$ est constituée de
points génériques de composantes irréductibles de $U$, alors $u \in W$.
\end{enumerate}
\end{lemma}

\begin{proof}
Dans ce cas, la stratification
$U = Z_0 \supset Z_1 \supset Z_2 \supset \ldots$ du
Lemme \ref{more-groupoids-lemma-finite-stratify} est donnée par des immersions fermées $Z_k \to U$
de présentation finie, voir
Diviseurs, Lemme \ref{divisors-lemma-locally-free-rank-r-pullback}.
La partie (1) en résulte immédiatement, puisque $W' \to W$ est localement obtenu
par l'intersection de l'ouvert $W$ avec $Z_r$. Pour voir la partie (2),
soit $\{u_1, \ldots, u_n\}$ l'orbite de $u$.
Comme les sous-schémas fermés $Z_k$ sont invariants par $R$ et que
$\bigcap Z_k = \emptyset$, on trouve un $k$ tel que $u_i \in Z_k$
et $u_i \not \in Z_{k + 1}$ pour tout $i$.
L'image de $Z_k \to U$ et celle de $Z_{k + 1} \to U$ sont localement constructibles
(Morphismes, Théorème \ref{morphisms-theorem-chevalley}).
Comme $u_i \in U$ est un point générique d'une composante irréductible
de $U$, il existe un voisinage ouvert $U_i$ de $u_i$ qui
est contenu dans $Z_k \setminus Z_{k + 1}$ ensemblistement
(Propriétés, Lemme \ref{properties-lemma-generic-point-in-constructible}).
Dans la démonstration du Lemme \ref{more-groupoids-lemma-finite-flat-over-almost-dense-subscheme},
nous avons construit $W$ comme une union disjointe $\coprod W_r$,
avec $W_r \subset Z_{r - 1} \setminus Z_r$, de sorte que
$U = \bigcup t(s^{-1}(\overline{W_r}))$. Comme $\{u_1, \ldots, u_n\}$
est une orbite sous $R$, on voit que $u \in t(s^{-1}(\overline{W_r}))$
implique $u_i \in \overline{W_r}$ pour un certain $i$, ce qui implique
$U_i \cap W_r \not = \emptyset$, et donc $r = k$.
On en conclut que $u$ appartient à
$$
W_{k + 1} = Z_k \setminus
\left(
Z_{k + 1} \cup \bigcup\nolimits_{r \leq k} t(s^{-1}(\overline{W_r}))
\right)
$$
comme voulu.
\end{proof}

\begin{lemma}
\label{more-groupoids-lemma-invariant-affine-open-around-generic-point}
Soit $(U, R, s, t, c)$ un schéma en groupoïdes sur un schéma $S$. Supposons $s, t$
finis et de présentation finie, et $U$ quasi-séparé. Soient
$u_1, \ldots, u_m \in U$ des points dont les orbites sont constituées de points génériques
de composantes irréductibles de $U$. Il existe alors des sous-schémas invariants par $R$
$V' \subset V \subset U$ tels que
\begin{enumerate}
\item $u_1, \ldots, u_m \in V'$,
\item $V$ soit ouvert dans $U$,
\item $V'$ et $V$ soient affines,
\item $V' \subset V$ soit un épaississement de présentation finie,
\item les morphismes $s', t'$ de la restriction $(V', R', s', t', c')$
soient finis localement libres.
\end{enumerate}
\end{lemma}

\begin{proof}
Soient $W' \subset W \subset U$ comme dans le
Lemme \ref{more-groupoids-lemma-finite-flat-over-almost-dense-subscheme}.
D'après le Lemme \ref{more-groupoids-lemma-finite-flat-over-almost-dense-subscheme-addendum},
on a $u_j \in W$ pour chaque indice, et $W' \to W$ est un épaississement de présentation finie.
D'après Limites, Lemme \ref{limits-lemma-affines-glued-in-closed-affine},
il suffit de trouver un sous-schéma ouvert affine invariant par $R$
$V'$ de $W'$ contenant $u_j$ (car on peut alors prendre pour $V \subset W$
le sous-schéma ouvert correspondant, qui sera affine).
On peut donc remplacer $(U, R, s, t, c)$ par la restriction
$(W', R', s', t', c')$ à $W'$.
Autrement dit, on peut supposer que l'on a un schéma en groupoïdes $(U, R, s, t, c)$
dont les morphismes $s$ et $t$ sont finis localement libres.
D'après Propriétés, Lemme \ref{properties-lemma-maximal-points-affine},
on peut trouver un ouvert affine contenant la réunion des orbites de
$u_1, \ldots, u_m$. Enfin, on peut appliquer
Groupoïdes, Lemme \ref{groupoids-lemma-find-invariant-affine}
pour conclure.
\end{proof}

\noindent
Le lemme suivant est un cas particulier du
Lemme \ref{more-groupoids-lemma-invariant-affine-open-around-generic-point},
mais nous reprenons l'argument, car il est légèrement plus simple dans ce cas
(il évite d'utiliser le
Lemme \ref{more-groupoids-lemma-finite-flat-over-almost-dense-subscheme-addendum}).

\begin{lemma}
\label{more-groupoids-lemma-invariant-affine-open-around-generic-point-Noetherian}
Soit $(U, R, s, t, c)$ un schéma en groupoïdes sur un schéma $S$.
Supposons $s, t$ finis et $U$ localement noethérien, et soient $u_1, \ldots, u_m \in U$
des points dont les orbites sont constituées de points génériques de
composantes irréductibles de $U$. Il existe alors des sous-schémas invariants par $R$
$V' \subset V \subset U$ tels que
\begin{enumerate}
\item $u_1, \ldots, u_m \in V'$,
\item $V$ soit ouvert dans $U$,
\item $V'$ et $V$ soient affines,
\item $V' \subset V$ soit un épaississement,
\item les morphismes $s', t'$ de la restriction $(V', R', s', t', c')$
soient finis localement libres.
\end{enumerate}
\end{lemma}

\begin{proof}
Soit $\{u_{j1}, \ldots, u_{jn_j}\}$ l'orbite de $u_j$.
Soient $W' \subset W \subset U$ comme dans le
Lemme \ref{more-groupoids-lemma-finite-flat-over-almost-dense-subscheme}.
Comme $U = t(s^{-1}(\overline{W}))$, on voit qu'il existe au moins
un $u_{ji} \in \overline{W}$. Comme $u_{ji}$ est un point générique
d'une composante irréductible et que $U$ est localement noethérien,
cela implique que $u_{ji} \in W$. Comme $W$ est invariant par $R$, on
conclut que $u_j \in W$ et, en fait, que l'orbite entière est contenue dans $W$.
D'après Cohomologie des schémas, Lemme
\ref{coherent-lemma-image-affine-finite-morphism-affine-Noetherian},
il suffit de trouver un sous-schéma ouvert affine $V'$ invariant par $R$
de $W'$ contenant $u_1, \ldots, u_m$ (car on peut alors prendre pour $V \subset W$
le sous-schéma ouvert correspondant, qui sera affine).
On peut donc remplacer $(U, R, s, t, c)$
par la restriction $(W', R', s', t', c')$ à $W'$.
Autrement dit, on peut supposer que l'on a un schéma en groupoïdes $(U, R, s, t, c)$
dont les morphismes $s$ et $t$ sont finis localement libres.
D'après Propriétés, Lemme \ref{properties-lemma-maximal-points-affine},
on peut trouver un ouvert affine contenant $\{u_{ij}\}$
(un schéma localement noethérien est quasi-séparé d'après
Propriétés, Lemme \ref{properties-lemma-locally-Noetherian-quasi-separated}).
Enfin, on peut appliquer
Groupoïdes, Lemme \ref{groupoids-lemma-find-invariant-affine}
pour conclure.
\end{proof}

\begin{lemma}
\label{more-groupoids-lemma-find-affine-integral}
Soit $(U, R, s, t, c)$ un schéma en groupoïdes sur un schéma $S$
avec $s, t$ intégraux. Soit $g : U' \to U$ un morphisme intégral
tel que toute orbite sous $R$ dans $U$ rencontre $g(U')$. Soit $(U', R', s', t', c')$
la restriction de $R$ à $U'$. Si $u' \in U'$ appartient à un
ouvert affine invariant par $R'$, alors l'image $u \in U$ appartient
à un ouvert affine de $U$ invariant par $R$.
\end{lemma}

\begin{proof}
Soit $W' \subset U'$ un ouvert affine invariant par $R'$.
Posons $\tilde R = U' \times_{g, U, t} R$, muni des applications
$\text{pr}_0 : \tilde R \to U'$ et $h = s \circ \text{pr}_1 : \tilde R \to U$.
Observons que $\text{pr}_0$ et $h$ sont intégraux.
Il s'ensuit que $\tilde W = \text{pr}_0^{-1}(W')$ est affine.
Comme $W'$ est invariant par $R'$, l'image
$W = h(\tilde W)$ est ensemblistement invariante par $R$ et
$\tilde W = h^{-1}(W)$ ensemblistement (détails omis).
Ainsi, si l'on peut montrer que $W$ est ouvert, alors $W$ est un schéma
et le morphisme $\tilde W \to W$ est intégral surjectif,
ce qui implique que $W$ est affine d'après
Limites, Proposition \ref{limits-proposition-affine}.
Cependant, l'hypothèse selon laquelle les orbites rencontrent $U'$ implique
que $h : \tilde R \to U$ est surjectif. Comme un
morphisme intégral surjectif est submersif
(Topologie, Lemme \ref{topology-lemma-closed-morphism-quotient-topology}
et Morphismes, Lemme \ref{morphisms-lemma-integral-universally-closed}),
il s'ensuit que $W$ est ouvert.
\end{proof}

\noindent
Le lemme technique suivant produit des fonctions « presque » invariantes
dans la situation d'un groupoïde fini sur un
schéma quasi-affine.

\begin{lemma}
\label{more-groupoids-lemma-find-almost-invariant-function}
Soit $(U, R, s, t, c)$ un schéma en groupoïdes avec $s, t$ finis et de
présentation finie. Soient $u_1, \ldots, u_m \in U$ des points dont les orbites sous $R$
sont constituées de points génériques de composantes irréductibles de $U$.
Soit $j : U \to \Spec(A)$ une immersion.
Soit $I \subset A$ un idéal tel que $j(U) \cap V(I) = \emptyset$
et que $V(I) \cup j(U)$ soit fermé dans $\Spec(A)$.
Il existe alors un élément $h \in I$ tel que $j^{-1}D(h)$
soit un sous-schéma ouvert affine de $U$ invariant par $R$ et contenant
$u_1, \ldots, u_m$.
\end{lemma}

\begin{proof}
Soient $u_1, \ldots, u_m \in V' \subset V \subset U$ comme dans le
Lemme \ref{more-groupoids-lemma-invariant-affine-open-around-generic-point}.
Comme $U \setminus V$ est fermé dans $U$, que $j$ est une immersion et que $V(I) \cup j(U)$
est fermé dans $\Spec(A)$, on peut trouver un idéal
$J \subset I$ tel que $V(J) = V(I) \cup j(U \setminus V)$.
On peut par exemple prendre l'idéal des éléments de $I$ qui
s'annulent sur $j(U \setminus V)$. On peut donc remplacer
$(U, R, s, t, c)$, $j : U \to \Spec(A)$ et $I$ par
$(V', R', s', t', c')$, $j|_{V'} : V' \to \Spec(A)$ et $J$.
Autrement dit, on peut supposer que $U$ est affine et que
$s$ et $t$ sont finis localement libres.
Prenons un quelconque $f \in I$ qui ne s'annule en aucun des
points des orbites sous $R$ de $u_1, \ldots, u_m$
(Algèbre, Lemme \ref{algebra-lemma-silly}). Considérons
$$
g = \text{Norm}_s(t^\sharp(j^\sharp(f))) \in \Gamma(U, \mathcal{O}_U)
$$
Comme $f \in I$ et que $V(I) \cup j(U)$ est fermé, on voit que
$U \cap D(f) \to D(f)$ est une immersion fermée. Ainsi, $f^ng$ est
l'image d'un élément $h \in I$ pour un certain $n > 0$. Nous affirmons que $h$ convient.
En effet, nous avons vu dans
Groupoïdes, Lemme \ref{groupoids-lemma-determinant-trick},
que $g$ est une fonction invariante par $R$ ; ainsi, $D(g) \subset U$
est invariant par $R$. Comme $f$ ne s'annule pas sur l'orbite de $u_j$,
la fonction $g$ ne s'annule pas en $u_j$. De plus, on a
$V(g) \supset V(j^\sharp(f))$ et donc $j^{-1}D(h) = D(g)$.
\end{proof}

\begin{lemma}
\label{more-groupoids-lemma-no-specializations-map-to-same-point}
Soit $(U, R, s, t, c)$ un schéma en groupoïdes. Si $s, t$ sont finis,
et si $u, u' \in R$ sont des points distincts de la même orbite,
alors $u'$ n'est pas une spécialisation de $u$.
\end{lemma}

\begin{proof}
Soit $r \in R$ tel que $s(r) = u$ et $t(r) = u'$. Si $u \leadsto u'$
alors nous pouvons trouver une spécialisation non triviale $r \leadsto r'$ telle que
$s(r') = u'$, voir 
Schémas, lemme \ref{schemes-lemma-quasi-compact-closed}.
Posons $u'' = t(r')$. Remarquons que $u'' \not = u'$ puisqu'il n'existe pas de
spécialisations dans les fibres d'un morphisme fini.
Nous pouvons donc continuer et trouver une spécialisation non triviale
$r' \leadsto r''$ telle que $s(r'') = u''$, et ainsi de suite. Cela montre que
l'orbite de $u$ contient une suite infinie
$u \leadsto u' \leadsto u'' \leadsto \ldots$
de spécialisations, ce qui est absurde puisque l'orbite
$t(s^{-1}(\{u\}))$ est finie.
\end{proof}

\begin{lemma}
\label{more-groupoids-lemma-get-affine}
Soit $j : V \to \Spec(A)$ une immersion quasi-compacte de schémas.
Soit $f \in A$ tel que $j^{-1}D(f)$ soit affine et que $j(V) \cap V(f)$
soit fermé. Alors $V$ est affine.
\end{lemma}

\begin{proof}
Cela résulte de Morphismes, lemme \ref{morphisms-lemma-get-affine},
mais nous allons aussi en donner une démonstration directe.
Posons $A' = \Gamma(V, \mathcal{O}_V)$. Alors $j' : V \to \Spec(A')$ est une
immersion ouverte quasi-compacte, voir
Propriétés, lemme \ref{properties-lemma-quasi-affine}.
Soit $f' \in A'$ l'image de $f$. Alors $(j')^{-1}D(f') = j^{-1}D(f)$
est affine. D'autre part, $j'(V) \cap V(f')$ est un sous-schéma de
$\Spec(A')$ qui s'envoie isomorphiquement sur le sous-schéma fermé
$j(V) \cap V(f)$ de $\Spec(A)$. Il est donc fermé dans $\Spec(A')$,
par exemple d'après Schémas, lemme \ref{schemes-lemma-section-immersion}.
Ainsi, nous pouvons remplacer $A$ par $A'$ et supposer que $j$ est une immersion ouverte
et que $A = \Gamma(V, \mathcal{O}_V)$.

\medskip\noindent
Dans ce cas, nous affirmons que $j(V) = \Spec(A)$, ce qui achève la démonstration.
Sinon, nous pouvons trouver un ouvert affine principal $D(g) \subset \Spec(A)$
qui rencontre le complémentaire et évite la partie fermée $j(V) \cap V(f)$.
Remarquons que $j$ envoie $j^{-1}D(f)$ isomorphiquement sur $D(f)$, voir
Propriétés, lemme \ref{properties-lemma-invert-f-affine}.
Ainsi $D(g)$ rencontre $V(f)$. D'autre part, $j^{-1}D(g)$
est un ouvert principal de l'ouvert affine $j^{-1}D(f)$, donc est affine.
Par conséquent, en appliquant de nouveau
Propriétés, lemme \ref{properties-lemma-invert-f-affine},
nous voyons que $D(g)$ est isomorphe à $j^{-1}D(g) \subset j^{-1}D(f)$,
ce qui implique que $D(g) \subset D(f)$. Cette contradiction achève
la démonstration.
\end{proof}

\begin{lemma}
\label{more-groupoids-lemma-find-affine-codimension-1}
Soit $(U, R, s, t, c)$ un schéma en groupoïdes. Soit $u \in U$. Supposons que
\begin{enumerate}
\item $s, t$ soient des morphismes finis,
\item $U$ soit séparé et localement noethérien,
\item $\dim(\mathcal{O}_{U, u'}) \leq 1$ pour tout point $u'$
de l'orbite de $u$.
\end{enumerate}
Alors $u$ est contenu dans un ouvert affine de $U$ invariant par $R$.
\end{lemma}

\begin{proof}
L'orbite sous $R$ de $u$ est finie. D'après les conditions (2) et (3), elle est contenue
dans un ouvert affine $U'$ de $U$, voir
Variétés, proposition
\ref{varieties-proposition-finite-set-of-points-of-codim-1-in-affine}.
Alors $t(s^{-1}(U \setminus U'))$ est une partie fermée et invariante par $R$
de $U$ qui ne contient pas $u$. Ainsi
$U \setminus t(s^{-1}(U \setminus U'))$ est un ouvert invariant par $R$
de $U'$ contenant $u$.
En remplaçant $U$ par cet ouvert, nous pouvons supposer que $U$ est quasi-affine.

\medskip\noindent
D'après le lemme \ref{more-groupoids-lemma-find-affine-integral}, nous pouvons remplacer $U$ par son réduit
et supposer que $U$ est réduit. Cela signifie que les sous-schémas invariants par $R$
$W' \subset W \subset U$ du
lemme \ref{more-groupoids-lemma-finite-flat-over-almost-dense-subscheme}
sont égaux : $W' = W$. Comme $U = t(s^{-1}(\overline{W}))$, un point
$u'$ appartenant à l'orbite sous $R$ de $u$ est contenu dans $\overline{W}$
et, d'après le lemme \ref{more-groupoids-lemma-find-affine-integral},
nous pouvons remplacer $U$ par $\overline{W}$ et $u$ par $u'$.
Nous pouvons donc supposer qu'il existe
un sous-schéma ouvert dense $W \subset U$ invariant par $R$ tel que
les morphismes $s_W, t_W$ de la restriction $(W, R_W, s_W, t_W, c_W)$ soient
finis localement libres.

\medskip\noindent
Si $u \in W$, nous avons terminé d'après
Groupoïdes, lemme \ref{groupoids-lemma-find-invariant-affine}
(car $W$ est quasi-affine, donc tout ensemble fini de points
de $W$ est contenu dans un ouvert affine, voir
Propriétés, lemme \ref{properties-lemma-ample-finite-set-in-affine}).
Nous supposons donc que $u \not \in W$ et, par suite, qu'aucun des points de
l'orbite de $u$ n'appartient à $W$. Soit $\xi \in U$
un point qui se spécialise non trivialement en un point $u'$ de l'orbite
de $u$. Puisqu'il n'y a pas de spécialisations entre les points de
l'orbite de $u$ (lemme \ref{more-groupoids-lemma-no-specializations-map-to-same-point}),
nous voyons que $\xi$ n'appartient pas à l'orbite.
D'après l'hypothèse (3), nous voyons que $\xi$ est un point générique de $U$
et donc que $\xi \in W$.
Comme $U$ est noethérien, il n'y a qu'un nombre fini de ces
points $\xi_1, \ldots, \xi_m \in W$. Puisque $s_W, t_W$ sont plats, l'orbite
de chaque $\xi_j$ est formée de points génériques de composantes irréductibles
de $W$ (et donc de $U$).

\medskip\noindent
Soit $j : U \to \Spec(A)$ une immersion de $U$ dans un schéma affine
(cela est possible puisque $U$ est quasi-affine). Soit $J \subset A$
un idéal tel que $V(J) \cap j(W) = \emptyset$ et que $V(J) \cup j(W)$
soit fermé. Appliquons le lemme \ref{more-groupoids-lemma-find-almost-invariant-function}
au schéma en groupoïdes $(W, R_W, s_W, t_W, c_W)$, au morphisme
$j|_W : W \to \Spec(A)$, aux points $\xi_j$ et à l'idéal $J$
pour trouver $f \in J$ tel que $(j|_W)^{-1}D(f)$ soit un ouvert invariant par $R_W$
affine contenant $\xi_j$ pour tout $j$. Puisque $f \in J$,
nous voyons que $j^{-1}D(f) \subset W$, c'est-à-dire que $j^{-1}D(f)$ est
un ouvert affine de $U$ invariant par $R$ et contenu dans $W$
et contenant tous les $\xi_j$.

\medskip\noindent
Soit $Z$ le sous-schéma fermé réduit induit sur
$$
U \setminus j^{-1}D(f) = j^{-1}V(f).
$$
Alors $Z$ est, du point de vue ensembliste,
invariant par $R$ (mais il peut ne pas être invariant par $R$ au sens schématique).
Soit $(Z, R_Z, s_Z, t_Z, c_Z)$ la restriction de $R$ à $Z$.
Puisque $Z \to U$ est fini, il s'ensuit que $s_Z$ et $t_Z$ sont finis.
Puisque $u \in Z$, l'orbite de $u$ est contenue dans $Z$ et coïncide avec
l'orbite sous $R_Z$ de $u$, considéré comme un point de $Z$. Puisque
$\dim(\mathcal{O}_{U, u'}) \leq 1$ et puisque $\xi_j \not \in Z$
pour tout $j$, nous voyons que $\dim(\mathcal{O}_{Z, u'}) \leq 0$ pour
tout $u'$ de l'orbite de $u$. Autrement dit, l'orbite sous $R_Z$ de $u$
est formée de points génériques de composantes irréductibles de $Z$.

\medskip\noindent
Soit $I \subset A$ un idéal tel que $V(I) \cap j(U) =\emptyset$
et que $V(I) \cup j(U)$ soit fermé. Appliquons
le lemme \ref{more-groupoids-lemma-find-almost-invariant-function} au
schéma en groupoïdes $(Z, R_Z, s_Z, t_Z, c_Z)$, à la restriction $j|_Z$,
à l'idéal $I$ et au point $u \in Z$, pour obtenir $h \in I$ tel que
$j^{-1}D(h) \cap Z$ soit un ouvert affine invariant par $R_Z$ et contenant $u$.

\medskip\noindent
Considérons la fonction invariante par $R_W$ (Groupoïdes, lemme
\ref{groupoids-lemma-determinant-trick})
$$
g = 
\text{Norme}_{s_W}(t_W^\sharp(j^\sharp(h)|_W)) \in \Gamma(W, \mathcal{O}_W)
$$
(Dans ce qui suit, nous n'avons besoin que de la restriction de $g$ à $j^{-1}D(f)$ et,
dans ce cas, la norme est prise le long d'un morphisme fini localement libre entre schémas affines.)
Nous affirmons que
$$
V = (W_g \cap j^{-1}D(f)) \cup (j^{-1}D(h) \cap Z)
$$
est un ouvert affine de $U$ invariant par $R$, ce qui achève la démonstration du lemme.
Il est invariant par $R$ du point de vue ensembliste par construction. Comme $V$ est une
partie constructible, pour voir qu'il est ouvert il suffit de montrer qu'il est
stable par générisation dans $U$ (Topologie, lemme
\ref{topology-lemma-characterize-closed-Noetherian}
ou, plus généralement,
Topologie, lemme
\ref{topology-lemma-constructible-stable-specialization-closed}).
Puisque $W_g \cap j^{-1}D(f)$ est ouvert dans $U$, il suffit de considérer
une spécialisation $u_1 \leadsto u_2$ de $U$ telle que
$u_2 \in j^{-1}D(h) \cap Z$.
Cela signifie que $h$ est non nul en $j(u_2)$ et que $u_2 \in Z$.
Si $u_1 \in Z$, alors $j(u_1) \leadsto j(u_2)$ et, puisque
$h$ est non nul en $j(u_2)$, il est non nul en $j(u_1)$, ce qui
implique que $u_1 \in V$. Si $u_1 \not \in Z$ et n'appartient pas non plus
à $W_g \cap j^{-1}D(f)$, alors $u_1 \in W$, $u_1 \not \in W_g$
car le complémentaire de $Z = j^{-1}V(f)$ est contenu dans $W \cap j^{-1}D(f)$.
Il existe donc un point $r_1 \in R$ tel que $s(r_1) = u_1$
et tel que $h$ soit nul en $t(r_1)$. Puisque $s$ est fini, nous
pouvons trouver une spécialisation $r_1 \leadsto r_2$ telle que $s(r_2) = u_2$.
Mais nous en déduisons alors que $h$ est nul en $u'_2 = t(r_2)$,
ce qui contredit le fait que $j^{-1}D(h) \cap Z$ est
invariant par $R$ et que $u_2$ lui appartient. Ainsi $V$ est ouvert.

\medskip\noindent
Observons que $V \subset j^{-1}D(h)$ pour notre fonction $h \in I$.
Nous obtenons donc une immersion
$$
j' : V \longrightarrow \Spec(A_h)
$$
Soit $f' \in A_h$ l'image de $f$. Alors $(j')^{-1}D(f')$
est l'ouvert principal déterminé par $g$ dans l'ouvert
affine $j^{-1}D(f)$ de $U$.
Ainsi $(j')^{-1}D(f)$ est affine. Enfin,
$j'(V) \cap V(f') = j'(j^{-1}D(h) \cap Z)$
est fermé dans $\Spec(A_h/(f')) = \Spec((A/f)_h) = D(h) \cap V(f)$
en vertu du choix de $h \in I$ et de l'idéal $I$. Nous pouvons donc appliquer le
lemme \ref{more-groupoids-lemma-get-affine}
pour conclure que $V$ est affine, comme affirmé ci-dessus.
\end{proof}







\section{Descente des morphismes ind-quasi-affines}
\label{more-groupoids-section-ind-quasi-affine}

\noindent
Les morphismes ind-quasi-affines ont été définis dans
Compléments sur les morphismes, section \ref{more-morphisms-section-ind-quasi-affine}.
Cette section est l'analogue de
Descente, section \ref{descent-section-quasi-affine},
pour les morphismes ind-quasi-affines.

\medskip\noindent
Soit $X$ un schéma quasi-séparé. Soit $E \subset X$ une partie
qui est l'intersection d'une famille non vide d'ouverts quasi-compacts de $X$.
Écrivons $E = \bigcap_{i \in I} U_i$, où $U_i \subset X$ est un ouvert quasi-compact
et où $I$ est non vide.
En ajoutant les intersections finies, nous pouvons supposer que, pour $i, j \in I$,
il existe $k \in I$ tel que $U_k \subset U_i \cap U_j$.
Dans cette situation, nous avons
\begin{equation}
\label{more-groupoids-equation-sections-of-intersection}
\Gamma(E, \mathcal{F}|_E) = \colim \Gamma(U_i, \mathcal{F}|_{U_i})
\end{equation}
pour tout faisceau $\mathcal{F}$ défini sur $X$. En effet, fixons $i_0 \in I$
et remplaçons $X$ par $U_{i_0}$ et $I$ par
$\{i \in I \mid U_i \subset U_{i_0}\}$. Alors $X$ est quasi-compact
et quasi-séparé, donc est un espace spectral, voir
Propriétés, lemme
\ref{properties-lemma-quasi-compact-quasi-separated-spectral}.
Nous voyons alors que l'égalité résulte de
Topologie, lemme \ref{topology-lemma-make-spectral-space}, et de
Faisceaux, lemme \ref{sheaves-lemma-descend-opens}.
(En fait, la formule vaut aussi pour les groupes de cohomologie supérieurs
lorsque $\mathcal{F}$ est abélien, voir
Cohomologie, lemme \ref{cohomology-lemma-colimit}.)

\begin{lemma}
\label{more-groupoids-lemma-sits-in-functions}
Soit $X$ un schéma ind-quasi-affine. Soit $E \subset X$ une
intersection d'une famille non vide d'ouverts quasi-compacts de $X$.
Posons $A = \Gamma(E, \mathcal{O}_X|_E)$ et $Y = \Spec(A)$.
Alors le morphisme canonique
$$
j : (E, \mathcal{O}_X|_E) \longrightarrow (Y, \mathcal{O}_Y)
$$
fourni par Schémas, lemme \ref{schemes-lemma-morphism-into-affine},
détermine un isomorphisme
$(E, \mathcal{O}_X|_E) \to (E', \mathcal{O}_Y|_{E'})$
où $E' \subset Y$ est une intersection d'ouverts quasi-compacts.
Si $W \subset E$ est ouvert dans $X$, alors $j(W)$ est ouvert dans $Y$.
\end{lemma}

\begin{proof}
Remarquons que $(E, \mathcal{O}_X|_E)$ est un espace localement annelé, de sorte que
Schémas, lemme \ref{schemes-lemma-morphism-into-affine}, s'applique
à $A \to \Gamma(E, \mathcal{O}_X|_E)$. Écrivons $E = \bigcap_{i \in I} U_i$,
avec $I \not = \emptyset$ et $U_i \subset X$ ouvert quasi-compact.
Nous pouvons supposer, et nous le faisons, que pour $i, j \in I$ il existe $k \in I$ tel que
$U_k \subset U_i \cap U_j$. Posons $A_i = \Gamma(U_i, \mathcal{O}_{U_i})$.
Nous obtenons les diagrammes commutatifs
$$
\xymatrix{
(E, \mathcal{O}_X|_E) \ar[r] \ar[d] &
(\Spec(A), \mathcal{O}_{\Spec(A)}) \ar[d] \\
(U_i, \mathcal{O}_{U_i}) \ar[r] &
(\Spec(A_i), \mathcal{O}_{\Spec(A_i)})
}
$$
Puisque $U_i$ est quasi-affine, nous voyons que $U_i \to \Spec(A_i)$
est une immersion ouverte quasi-compacte. D'autre part,
$A = \colim A_i$. Ainsi $\Spec(A) = \lim \Spec(A_i)$ comme espaces
topologiques (Limites, lemme \ref{limits-lemma-topology-limit}). Puisque
$E = \lim U_i$ (d'après Topologie, lemme \ref{topology-lemma-make-spectral-space}),
nous voyons que $E \to \Spec(A)$ est un homéomorphisme sur son
image $E'$ et que $E'$ est l'intersection des images réciproques
des ouverts $U_i \subset \Spec(A_i)$ dans $\Spec(A)$. Pour tout
$e \in E$, l'anneau local $\mathcal{O}_{X, e}$ est
$\mathcal{O}_{U_i, e}$, qui coïncide avec l'anneau local correspondant sur $\Spec(A)$.

\medskip\noindent
Pour démontrer la dernière assertion du lemme, raisonnons comme suit.
Choisissons $i, j \in I$ tels que $U_i \subset U_j$.
Considérons les diagrammes commutatifs suivants
$$
\xymatrix{
U_i \ar[r] \ar[d] & \Spec(A_i) \ar[d] \\
U_i \ar[r] & \Spec(A_j)
}
\quad\quad
\xymatrix{
W \ar[r] \ar[d] & \Spec(A_i) \ar[d] \\
W \ar[r] & \Spec(A_j)
}
\quad\quad
\xymatrix{
W \ar[r] \ar[d] & \Spec(A) \ar[d] \\
W \ar[r] & \Spec(A_j)
}
$$
D'après Propriétés, lemme
\ref{properties-lemma-cartesian-diagram-quasi-affine},
le premier diagramme est cartésien. Le deuxième l'est donc aussi.
En passant à la limite, nous trouvons que le troisième diagramme
est cartésien, si bien que la flèche horizontale supérieure de ce diagramme est une immersion ouverte.
\end{proof}

\begin{lemma}
\label{more-groupoids-lemma-affine-base-change}
Supposons donné un diagramme cartésien
$$
\xymatrix{
X \ar[d]_f \ar[r] & \Spec(B) \ar[d] \\
Y \ar[r] & \Spec(A)
}
$$
de schémas. Soit $E \subset Y$ une intersection d'une famille non vide
d'ouverts quasi-compacts de $Y$. Alors
$$
\Gamma(f^{-1}(E), \mathcal{O}_X|_{f^{-1}(E)}) =
\Gamma(E, \mathcal{O}_Y|_E) \otimes_A B
$$
pourvu que $Y$ soit quasi-séparé et que $A \to B$ soit plat.
\end{lemma}

\begin{proof}
Écrivons $E = \bigcap_{i \in I} V_i$, où $V_i \subset Y$ est un ouvert quasi-compact.
Nous pouvons supposer, et nous le faisons, que pour $i, j \in I$ il existe $k \in I$ tel que
$V_k \subset V_i \cap V_j$. Nous avons alors de même
$f^{-1}(E) = \bigcap_{i \in I} f^{-1}(V_i)$ dans $X$.
Le résultat découle donc de l'équation (\ref{more-groupoids-equation-sections-of-intersection})
et du résultat correspondant pour $V_i$ et $f^{-1}(V_i)$, qui est
Cohomologie des schémas, lemme \ref{coherent-lemma-flat-base-change-cohomology}.
\end{proof}

\begin{lemma}[Gabber]
\label{more-groupoids-lemma-ind-quasi-affine}
Soit $S$ un schéma. Soit $\{X_i \to S\}_{i\in I}$ un recouvrement fpqc.
Soit $(V_i/X_i, \varphi_{ij})$ une donnée de descente relativement à
$\{X_i \to S\}$, voir Descente, définition
\ref{descent-definition-descent-datum-for-family-of-morphisms}. 
Si chaque morphisme $V_i \to X_i$ est ind-quasi-affine, alors la donnée de descente
est effective.
\end{lemma}

\begin{proof}
Être ind-quasi-affine est une propriété des morphismes de schémas
qui est préservée par tout changement de base, voir
Compléments sur les morphismes, lemme
\ref{more-morphisms-lemma-base-change-ind-quasi-affine}.
Ainsi Descente, lemme \ref{descent-lemma-descending-types-morphisms}, s'applique,
et il suffit de démontrer l'énoncé du lemme
dans le cas où le recouvrement fpqc est constitué du singleton
$\{X \to S\}$, formé d'un morphisme plat surjectif de schémas affines.
Écrivons $X = \Spec(A)$ et $S = \Spec(R)$, de sorte
que $R \to A$ est un homomorphisme d'anneaux fidèlement plat.
Soit $(V, \varphi)$ une donnée de descente relative à $X$ au-dessus de $S$
et supposons que $V \to X$ soit ind-quasi-affine, autrement dit que
$V$ soit ind-quasi-affine.

\medskip\noindent
Soit $(U, R, s, t, c)$ le schéma en groupoïdes au-dessus de $S$ où
$U = X$ et $R = X \times_S X$, et où $s$, $t$, $c$ sont les morphismes usuels.
D'après Groupoïdes, lemme \ref{groupoids-lemma-cartesian-equivalent-descent-datum},
la paire $(V, \varphi)$ correspond à un morphisme cartésien
$(U', R', s', t', c') \to (U, R, s, t, c)$ de schémas en groupoïdes.
Soit $u' \in U'$ un point quelconque. D'après
Groupoïdes, lemmes \ref{groupoids-lemma-constructing-invariant-opens},
\ref{groupoids-lemma-first-observation} et
\ref{groupoids-lemma-second-observation},
nous pouvons choisir $u' \in W \subset E \subset U'$
où $W$ est ouvert et invariant par $R'$, et où
$E$ est invariant par $R'$ du point de vue ensembliste et
est une intersection d'une famille non vide d'ouverts quasi-compacts.

\medskip\noindent
En revenant à $(V, \varphi)$, pour tout $v \in V$ nous pouvons choisir
$v \in W \subset E \subset V$ de sorte que les propriétés suivantes soient vérifiées :
(a) $W$ est ouvert et $\varphi(W \times_S X) = X \times_S W$, et
(b) $E$ est une intersection d'ouverts quasi-compacts et
$\varphi(E \times_S X) = X \times_S E$ du point de vue ensembliste.
Ici, la notation $E \times_S X$ désigne
l'image réciproque de $E$ dans $V \times_S X$ par le morphisme de projection, et
de même pour $X \times_S E$. D'après le lemme \ref{more-groupoids-lemma-affine-base-change},
cela implique que $\varphi$ définit un isomorphisme
\begin{align*}
\Gamma(E, \mathcal{O}_V|_E) \otimes_R A
& =
\Gamma(E \times_S X, \mathcal{O}_{V \times_S X}|_{E \times_S X}) \\
& \to
\Gamma(X \times_S E, \mathcal{O}_{X \times_S V}|_{X \times_S E}) \\
& =
A \otimes_R \Gamma(E, \mathcal{O}_V|_E)
\end{align*}
d'algèbres sur $A \otimes_R A$, que nous appellerons $\psi$.
La condition de cocycle pour $\varphi$
se traduit par la condition de cocycle pour $\psi$ comme dans
Descente, définition \ref{descent-definition-descent-datum-modules}
(détails omis). D'après Descente, proposition
\ref{descent-proposition-descent-module},
nous trouvons une $R$-algèbre $R'$ et un isomorphisme
$\chi : R' \otimes_R A \to \Gamma(E, \mathcal{O}_V|_E)$
d'$A$-algèbres, compatible avec $\psi$ et avec la
donnée de descente canonique sur $R' \otimes_R A$.

\medskip\noindent
D'après le lemme \ref{more-groupoids-lemma-sits-in-functions}, nous obtenons un ``plongement'' canonique
$$
j : (E, \mathcal{O}_V|_E) \longrightarrow
\Spec(\Gamma(E, \mathcal{O}_V|_E)) = \Spec(R' \otimes_R A)
$$
d'espaces localement annelés. La construction de cette application est canonique
et nous obtenons un diagramme commutatif
$$
\xymatrix{
& E \times_S X \ar[rr]_\varphi \ar[ld] \ar[rd]^{j'} & &
X \times_S E \ar[rd] \ar[ld]_{j''} \\
E \ar[rd]^j  & &
\Spec(R' \otimes_R A \otimes_R A) \ar[ld] \ar[rd] & &
E \ar[ld]_j \\
& \Spec(R' \otimes_R A) \ar[rd] && \Spec(R' \otimes_R A) \ar[ld] \\
& & \Spec(R')
}
$$
où $j'$ et $j''$ proviennent de la même construction appliquée à
$E \times_S X \subset V \times_S X$ et à $X \times_S E \subset X \times_S V$
au moyen de $\chi$ et des identifications utilisées pour construire $\psi$.
Il s'ensuit que $j(W)$ est un sous-schéma ouvert de $\Spec(R' \otimes_R A)$
dont les images réciproques par les deux projections
$\Spec(R' \otimes_R A \otimes_R A) \to \Spec(R' \otimes_R A)$
sont égales. D'après Descente, lemme \ref{descent-lemma-open-fpqc-covering},
nous trouvons un ouvert $W_0 \subset \Spec(R')$ dont le changement de base
à $\Spec(A)$ est $j(W)$. En examinant le diagramme ci-dessus,
nous voyons que la donnée de descente $(W, \varphi|_{W \times_S X})$
est effective. D'après Descente, lemme
\ref{descent-lemma-effective-for-fpqc-is-local-upstairs},
nous voyons que notre donnée de descente est effective.
\end{proof}










% OK, start here.
%

\chapter{Morphismes \'etales de schémas}


\phantomsection
\label{etale-section-phantom}




\section{Introduction}
\label{etale-section-introduction}

\noindent
Dans ce chapitre, nous étudions les morphismes \'etales de schémas. Nous illustrons
certains des concepts les plus importants en travaillant dans le cas noethérien.
Notre objectif principal est de réunir pour le lecteur suffisamment de résultats
d'algèbre commutative pour commencer la lecture d'un traité de cohomologie \'etale. Un
objectif auxiliaire est de fournir assez d'éléments pour que le lecteur puisse
cesser de qualifier l'expression « la topologie \'etale des schémas »
d'exercice de pur formalisme abstrait, s'il se livre à un tel blasphème.

\medskip\noindent
Nous renverrons aux autres
chapitres du projet Champs pour les résultats classiques de géométrie algébrique
(sur les schémas et l'algèbre commutative). Nous donnerons des démonstrations
détaillées des nouveaux résultats que nous énonçons ici.




\section{Conventions}
\label{etale-section-conventions}

\noindent
Dans ce chapitre, les schémas seront souvent supposés localement noethériens
et les anneaux seront souvent supposés noethériens. Mais, dans tous les énoncés,
nous le répéterons chaque fois que cela sera nécessaire et veillerons à énumérer
toutes les hypothèses ! Voici par ailleurs quelques faits généraux que nous utiliserons
souvent et qu'il est utile de garder à l'esprit :
\begin{enumerate}
\item Un homomorphisme d'anneaux $A \to B$ de type fini, avec $A$ noethérien,
est de présentation finie. Voir Algèbre,
Lemme \ref{algebra-lemma-Noetherian-finite-type-is-finite-presentation}.
\item Un morphisme (localement) de type fini entre schémas localement noethériens
est automatiquement (localement) de présentation finie.
Voir Morphismes,
Lemme \ref{morphisms-lemma-noetherian-finite-type-finite-presentation}.
\item Ajouter ici d'autres faits de ce genre.
\end{enumerate}




\section{Morphismes non ramifiés}
\label{etale-section-unramified-definition}

\noindent
Nous définissons d'abord les « homomorphismes non ramifiés d'anneaux locaux » pour les anneaux
locaux noethériens. Nous ne pouvons employer simplement le terme « non ramifié », car il existe déjà
une notion
d'homomorphisme d'anneaux non ramifié (Algèbre, section \ref{algebra-section-unramified})
et elle est différente. Après avoir quelque peu étudié cette notion, nous
la globaliserons pour décrire les morphismes non ramifiés de schémas localement noethériens.

\begin{definition}
\label{etale-definition-unramified-rings}
Soient $A$ et $B$ des anneaux locaux noethériens. Un homomorphisme local $A \to B$
est appelé un {\it homomorphisme non ramifié d'anneaux locaux} si
\begin{enumerate}
\item $\mathfrak m_AB = \mathfrak m_B$,
\item $\kappa(\mathfrak m_B)$ est une extension finie séparable de
$\kappa(\mathfrak m_A)$, et
\item $B$ est essentiellement de type fini sur $A$ (cela signifie
que $B$ est le localisé d'une $A$-algèbre de type fini en un idéal premier).
\end{enumerate}
\end{definition}

\noindent
C'est la version locale de la
définition donnée dans Algèbre, section \ref{algebra-section-unramified}.
Dans cette section, un homomorphisme d'anneaux $R \to S$ est dit non ramifié si et
seulement s'il est de type fini et si $\Omega_{S/R} = 0$.
Nous disons que $R \to S$ est non ramifié en un idéal premier $\mathfrak q \subset S$
s'il existe $g \in S$, $g \not \in \mathfrak q$, tel que
$R \to S_g$ soit un homomorphisme d'anneaux non ramifié. Il est démontré dans
Algèbre, Lemmes \ref{algebra-lemma-unramified-at-prime} et
\ref{algebra-lemma-characterize-unramified}, qu'étant donnés un homomorphisme
d'anneaux $R \to S$ de type fini et un idéal premier $\mathfrak q$ de $S$
au-dessus de $\mathfrak p \subset R$, on a
$$
R \to S\text{ est non ramifié en }\mathfrak q
\Leftrightarrow
\mathfrak pS_{\mathfrak q} = \mathfrak q S_{\mathfrak q}
\text{ et }
\kappa(\mathfrak p) \subset \kappa(\mathfrak q)\text{ est finie séparable}
$$
Nous voyons donc que, pour un homomorphisme local d'anneaux locaux, les propriétés
de notre définition ci-dessus sont étroitement liées au fait
d'être non ramifié. Nous avons en fait démontré le lemme suivant.

\begin{lemma}
\label{etale-lemma-characterize-unramified-Noetherian}
\begin{slogan}
La non-ramification est une condition locale sur les germes.
\end{slogan}
Soit $A \to B$ un homomorphisme de type fini, avec $A$ noethérien.
Soit $\mathfrak q$ un idéal premier de $B$ au-dessus de $\mathfrak p \subset A$.
Alors $A \to B$ est non ramifié en $\mathfrak q$ si et seulement si
$A_{\mathfrak p} \to B_{\mathfrak q}$ est un homomorphisme non ramifié
d'anneaux locaux.
\end{lemma}

\begin{proof}
Voir la discussion précédente.
\end{proof}

\noindent
Nous allons caractériser la propriété d'être non ramifié au moyen
des complétés. Pour un anneau local noethérien $A$,
nous notons $A^\wedge$ le complété de $A$ pour
l'idéal maximal. C'est encore un anneau local noethérien ; voir
Algèbre, Lemme \ref{algebra-lemma-completion-Noetherian-Noetherian}.

\begin{lemma}
\label{etale-lemma-unramified-completions}
Soient $A$ et $B$ des anneaux locaux noethériens.
Soit $A \to B$ un homomorphisme local.
\begin{enumerate}
\item si $A \to B$ est un homomorphisme non ramifié d'anneaux locaux,
alors $B^\wedge$ est un $A^\wedge$-module de type fini,
\item si $A \to B$ est un homomorphisme non ramifié d'anneaux locaux et si
$\kappa(\mathfrak m_A) = \kappa(\mathfrak m_B)$,
alors $A^\wedge \to B^\wedge$ est surjectif,
\item si $A \to B$ est un homomorphisme non ramifié d'anneaux locaux
et si $\kappa(\mathfrak m_A)$
est séparablement clos, alors $A^\wedge \to B^\wedge$ est surjectif,
\item si $A$ et $B$ sont des anneaux de valuation discrète complets, alors
$A \to B$ est un homomorphisme non ramifié d'anneaux locaux
si et seulement si une uniformisante de $A$ s'envoie sur une uniformisante de $B$,
et si l'extension des corps résiduels est finie séparable (et si $B$ est
essentiellement de type fini sur $A$).
\end{enumerate}
\end{lemma}

\begin{proof}
L'assertion (1) est un cas particulier du résultat
Algèbre, Lemme \ref{algebra-lemma-finite-after-completion}.
Pour l'assertion (2), remarquons que l'espace vectoriel sur $\kappa(\mathfrak m_A)$
$B^\wedge/\mathfrak m_{A^\wedge}B^\wedge$
est engendré par $1$. Le lemme de Nakayama
(Algèbre, Lemme \ref{algebra-lemma-NAK}) montre donc que l'application
$A^\wedge \to B^\wedge$ est surjective.
L'assertion (3) est un cas particulier de l'assertion (2).
L'assertion (4) découle immédiatement des définitions.
\end{proof}

\begin{lemma}
\label{etale-lemma-characterize-unramified-completions}
Soient $A$ et $B$ des anneaux locaux noethériens.
Soit $A \to B$ un homomorphisme local tel que $B$ soit
essentiellement de type fini sur $A$.
Les conditions suivantes sont équivalentes :
\begin{enumerate}
\item $A \to B$ est un homomorphisme non ramifié d'anneaux locaux,
\item $A^\wedge \to B^\wedge$ est un homomorphisme non ramifié d'anneaux locaux, et
\item $A^\wedge \to B^\wedge$ est non ramifié.
\end{enumerate}
\end{lemma}

\begin{proof}
L'équivalence de (1) et (2) résulte du fait que
$\mathfrak m_AA^\wedge$ est l'idéal maximal de $A^\wedge$
(et de même pour $B$), ainsi que de la fidèle platitude de $B \to B^\wedge$.
Par exemple, si $A^\wedge \to B^\wedge$ est non ramifié, alors
$\mathfrak m_AB^\wedge = (\mathfrak m_AB)B^\wedge = \mathfrak m_BB^\wedge$
et donc $\mathfrak m_AB = \mathfrak m_B$.

\medskip\noindent
Supposons satisfaites les conditions équivalentes (1) et (2).
Le Lemme \ref{etale-lemma-unramified-completions}
montre que $A^\wedge \to B^\wedge$ est
fini. Ainsi, $A^\wedge \to B^\wedge$ est de présentation finie et, par
Algèbre, Lemme \ref{algebra-lemma-characterize-unramified},
nous concluons que $A^\wedge \to B^\wedge$ est non ramifié en
$\mathfrak m_{B^\wedge}$. Comme $B^\wedge$ est local, nous concluons
que $A^\wedge \to B^\wedge$ est non ramifié.

\medskip\noindent
Supposons (3). Par Algèbre, Lemme \ref{algebra-lemma-unramified-at-prime},
nous concluons que $A^\wedge \to B^\wedge$ est un homomorphisme non ramifié
d'anneaux locaux, c'est-à-dire que (2) est satisfaite.
\end{proof}

\begin{definition}
\label{etale-definition-unramified-schemes}
(Voir Morphismes, Définition \ref{morphisms-definition-unramified},
pour la définition dans le cas général.)
Soit $Y$ un schéma localement noethérien.
Soit $f : X \to Y$ un morphisme localement de type fini.
Soit $x \in X$.
\begin{enumerate}
\item Nous disons que $f$ est {\it non ramifié en $x$} si
$\mathcal{O}_{Y, f(x)} \to \mathcal{O}_{X, x}$
est un homomorphisme non ramifié d'anneaux locaux.
\item Le morphisme $f : X \to Y$ est dit {\it non ramifié}
s'il est non ramifié en tout point de $X$.
\end{enumerate}
\end{definition}

\noindent
Montrons que cette définition coïncide avec celle du
chapitre sur les morphismes de schémas. Cela garantit en particulier que
l'ensemble des points où un morphisme est non ramifié est ouvert.

\begin{lemma}
\label{etale-lemma-unramified-definition}
Soit $Y$ un schéma localement noethérien.
Soit $f : X \to Y$ un morphisme localement de type fini.
Soit $x \in X$. Le morphisme $f$ est non ramifié en $x$ au
sens de la Définition \ref{etale-definition-unramified-schemes}
si et seulement s'il est non ramifié au
sens de Morphismes, Définition \ref{morphisms-definition-unramified}.
\end{lemma}

\begin{proof}
Cela résulte du Lemme \ref{etale-lemma-characterize-unramified-Noetherian}
et des définitions.
\end{proof}

\noindent
Voici quelques résultats sur les morphismes non ramifiés.
Sous la forme donnée dans cette liste, ils ne s'appliquent qu'aux
morphismes localement de type fini entre schémas localement noethériens.
Dans chaque cas, nous renvoyons au résultat général
démontré plus tôt dans le projet, même si, dans certains cas, on peut
démontrer le résultat plus facilement dans le cas noethérien.
En voici la liste :
\begin{enumerate}
\item La non-ramification est locale sur la source et sur le but pour la
topologie de Zariski.
\item Les morphismes non ramifiés sont stables par changement de base et par composition.
Voir Morphismes, Lemmes \ref{morphisms-lemma-base-change-unramified}
et \ref{morphisms-lemma-composition-unramified}.
\item Les morphismes non ramifiés de schémas sont localement quasi-finis,
et les morphismes non ramifiés quasi-compacts sont quasi-finis.
Voir Morphismes, Lemme \ref{morphisms-lemma-unramified-quasi-finite}.
\item Les morphismes non ramifiés sont de dimension relative $0$. Voir
Morphismes, Définition \ref{morphisms-definition-relative-dimension-d}
et
Morphismes, Lemme \ref{morphisms-lemma-locally-quasi-finite-rel-dimension-0}.
\item Un morphisme est non ramifié si et seulement si toutes ses fibres sont non ramifiées.
Autrement dit, la non-ramification se vérifie sur les fibres schématiques. Voir
Morphismes, Lemme \ref{morphisms-lemma-unramified-etale-fibres}.
\item Soient $X$ et $Y$ non ramifiés sur un schéma de base $S$.
Tout $S$-morphisme de $X$ vers $Y$ est non ramifié.
Voir Morphismes, Lemme \ref{morphisms-lemma-unramified-permanence}.
\end{enumerate}

\section{Trois autres caractérisations des morphismes non ramifiés}
\label{etale-section-three-other}

\noindent
Le théorème suivant donne trois notions équivalentes de
non-ramification en un point. Voir
Morphismes, Lemme \ref{morphisms-lemma-unramified-at-point},
pour une partie de l'énoncé concernant les schémas généraux.

\begin{theorem}
\label{etale-theorem-unramified-equivalence}
Soit $Y$ un schéma localement noethérien.
Soit $f : X \to Y$ un morphisme de schémas localement de type fini.
Soit $x$ un point de $X$. Les conditions suivantes sont équivalentes :
\begin{enumerate}
\item $f$ est non ramifié en $x$,
\item le germe $\Omega_{X/Y, x}$ du module des différentielles relatives
en $x$ est nul,
\item il existe des voisinages ouverts $U$ de $x$ et $V$ de $f(x)$, ainsi qu'un
diagramme commutatif
$$
\xymatrix{
U \ar[rr]_i \ar[rd] & & \mathbf{A}^n_V \ar[ld] \\
& V
}
$$
où $i$ est une immersion fermée définie par un
faisceau quasi-cohérent d'idéaux $\mathcal{I}$ tel que les différentielles
$\text{d}g$, pour $g \in \mathcal{I}_{i(x)}$, engendrent
$\Omega_{\mathbf{A}^n_V/V, i(x)}$, et
\item la diagonale $\Delta_{X/Y} : X \to X \times_Y X$
est un isomorphisme local en $x$.
\end{enumerate}
\end{theorem}

\begin{proof}
L'équivalence de (1) et (2) est démontrée dans
Morphismes, Lemme \ref{morphisms-lemma-unramified-at-point}.

\medskip\noindent
Si $f$ est non ramifié en $x$, alors $f$ est non ramifié sur un voisinage
ouvert de $x$ ; cela ne résulte pas immédiatement
de la Définition \ref{etale-definition-unramified-schemes} de ce chapitre,
mais cela découle de
Morphismes, Définition \ref{morphisms-definition-unramified}, dont nous
avons démontré l'équivalence au
Lemme \ref{etale-lemma-unramified-definition}.
Choisissons des ouverts affines $V \subset Y$ et $U \subset X$
tels que $f(U) \subset V$ et $x \in U$, et tels que $f$ soit
non ramifié sur $U$, c'est-à-dire que $f|_U : U \to V$ soit non ramifié.
Par Morphismes, Lemme \ref{morphisms-lemma-diagonal-unramified-morphism},
le morphisme $U \to U \times_V U$
est une immersion ouverte. Cela démontre que (1) implique (4).

\medskip\noindent
Si $\Delta_{X/Y}$ est un isomorphisme local en $x$, alors
$\Omega_{X/Y, x} = 0$ par
Morphismes, Lemme \ref{morphisms-lemma-differentials-diagonal}.
Nous voyons donc que (4) implique (2).
À ce stade, nous savons que (1), (2) et (4) sont toutes équivalentes.

\medskip\noindent
Supposons (3). L'hypothèse sur le diagramme, combinée au résultat
Morphismes, Lemme \ref{morphisms-lemma-differentials-relative-immersion},
montre que $\Omega_{U/V, x} = 0$. Comme $\Omega_{U/V, x} = \Omega_{X/Y, x}$,
nous concluons que (2) est satisfaite.

\medskip\noindent
Supposons enfin que (2) soit satisfaite. Pour démontrer (3), nous pouvons localiser sur
$X$ et $Y$ et supposer que $X$ et $Y$ sont affines.
Écrivons $X = \Spec(B)$ et $Y = \Spec(A)$.
Le point $x \in X$ correspond à un idéal premier $\mathfrak q \subset B$.
Notre hypothèse est que $\Omega_{B/A, \mathfrak q} = 0$
(voir Morphismes, Lemme \ref{morphisms-lemma-differentials-affine}, pour le
lien entre les différentielles sur les schémas et les modules
de différentielles en algèbre commutative).
Comme $Y$ est localement noethérien et $f$ localement de type fini,
nous voyons que $A$ est noethérien et que
$B \cong A[x_1, \ldots, x_n]/(f_1, \ldots, f_m)$ ; voir
Propriétés, Lemme \ref{properties-lemma-locally-Noetherian}, et
Morphismes, Lemme \ref{morphisms-lemma-locally-finite-type-characterize}.
En particulier, $\Omega_{B/A}$ est un $B$-module de type fini. Nous
pouvons donc trouver un seul $g \in B$, $g \not \in \mathfrak q$, tel que
le localisé principal $(\Omega_{B/A})_g$ soit nul. Ainsi, après
avoir remplacé $B$ par $B_g$, nous voyons que $\Omega_{B/A} = 0$ (la formation
des modules de différentielles commute à la localisation ; voir
Algèbre, Lemme \ref{algebra-lemma-differentials-localize}). Cela signifie que
les $\text{d}(f_j)$ engendrent le noyau de l'application canonique
$\Omega_{A[x_1, \ldots, x_n]/A} \otimes_A B \to \Omega_{B/A}$.
Ainsi, la surjection $A[x_1, \ldots, x_n] \to B$ de $A$-algèbres fournit le
diagramme commutatif de (3), et le théorème est démontré.
\end{proof}

\noindent
Comment utiliser ce théorème ? Voici quelques remarques :
\begin{enumerate}
\item Supposons que
$f : X \to Y$ et $g : Y \to Z$ soient deux morphismes localement de type
fini entre schémas localement noethériens. Il existe une
suite exacte courte canonique
$$
f^*(\Omega_{Y/Z}) \to \Omega_{X/Z} \to \Omega_{X/Y} \to 0
$$
voir Morphismes, Lemme \ref{morphisms-lemma-triangle-differentials}.
Le théorème implique donc que, si $g \circ f$ est non ramifié,
alors $f$ l'est aussi. C'est
Morphismes, Lemme \ref{morphisms-lemma-unramified-permanence}.
\item Puisque $\Omega_{X/Y}$ est isomorphe au faisceau conormal
du morphisme diagonal
(Morphismes, Lemme \ref{morphisms-lemma-differentials-diagonal}),
nous voyons que, si $X \to Y$ est un monomorphisme de
schémas localement noethériens et est localement de type fini,
alors $X \to Y$ est non ramifié.
En particulier, les immersions ouvertes et fermées de schémas localement noethériens
sont non ramifiées. Voir
Morphismes, Lemmes
\ref{morphisms-lemma-open-immersion-unramified} et
\ref{morphisms-lemma-closed-immersion-unramified}.
\item Le théorème implique aussi que l'ensemble des points
où un morphisme $f : X \to Y$ (localement de type fini entre schémas localement
noethériens) n'est pas non ramifié est
le support du faisceau cohérent $\Omega_{X/Y}$.
Cela permet de donner une définition schématique du
« lieu de ramification ».
\end{enumerate}

\section{La caractérisation fonctorielle des morphismes non ramifiés}
\label{etale-section-functorial-unramified}

\noindent
En géométrie algébrique élémentaire, on apprend que certaines classes de morphismes peuvent être
caractérisées fonctoriellement et que de telles descriptions sont fort utiles.
Les morphismes non ramifiés admettent eux aussi une telle caractérisation.

\begin{theorem}
\label{etale-theorem-formally-unramified}
Soit $f : X \to S$ un morphisme de schémas.
Supposons que $S$ soit un schéma localement noethérien et que $f$ soit localement de type fini.
Les conditions suivantes sont alors équivalentes :
\begin{enumerate}
\item $f$ est non ramifié,
\item le morphisme $f$ est formellement non ramifié :
pour tout $S$-schéma affine $T$ et tout sous-schéma $T_0$ de $T$
défini par un idéal de carré nul,
l'application naturelle
$$
\Hom_S(T, X) \longrightarrow \Hom_S(T_0, X)
$$
est injective.
\end{enumerate}
\end{theorem}

\begin{proof}
Voir Compléments sur les morphismes,
Lemme \ref{more-morphisms-lemma-unramified-formally-unramified},
pour un énoncé et une démonstration plus généraux.
Voici une esquisse de la démonstration dans le cas présent.

\medskip\noindent
On vérifie d'abord que les deux propriétés sont locales sur la source et sur le but.
Nous pouvons donc supposer que $S$ et $X$ sont affines.
Écrivons $X = \Spec(B)$ et $S = \Spec(R)$.
Écrivons $T = \Spec(C)$. Soit $J$ l'idéal de carré nul de $C$
tel que $T_0 = \Spec(C/J)$. Supposons donné le diagramme
$$
\xymatrix{
& B \ar[d]^\phi \ar[rd]^{\bar{\phi}}
& \\
R \ar[r] \ar[ur] & C \ar[r]
& C/J
}
$$
On vérifie ensuite que l'association $\phi' \mapsto \phi' - \phi$
définit une bijection entre l'ensemble des relèvements de $\bar{\phi}$ et le module
$\text{Der}_R(B, J)$. Ainsi, nous obtenons l'implication (1) $\Rightarrow$ (2)
par la description des morphismes non ramifiés dont le module
de différentielles est nul ; voir le Théorème \ref{etale-theorem-unramified-equivalence}.

\medskip\noindent
Pour obtenir l'implication réciproque, considérons la surjection
$q : C = (B \otimes_R B)/I^2 \to B = C/J$ définie par l'idéal de carré nul
$J = I/I^2$, où $I$ est le noyau de l'application de multiplication
$B \otimes_R B \to B$. Nous disposons déjà d'un relèvement $B \to C$ défini, par exemple,
par $b \mapsto b \otimes 1$. Ainsi, par le même raisonnement que ci-dessus, nous obtenons une
correspondance bijective entre les relèvements de $\text{id} : B \to C/J$ et
$\text{Der}_R(B, J)$. L'hypothèse implique donc que ce dernier module est
nul. Or nous savons que $J \cong \Omega_{B/R}$. Ainsi, $B/R$ est non ramifié.
\end{proof}



\section{Propriétés topologiques des morphismes non ramifiés}
\label{etale-section-topological-unramified}

\noindent
Le premier résultat topologique qui nous sera utile affirme
que les morphismes non ramifiés et séparés ont de « bonnes » sections.
Les résultats de cette section ne nécessitent aucune hypothèse noethérienne.

\begin{proposition}
\label{etale-proposition-properties-sections}
sections des morphismes non ramifiés.
\begin{enumerate}
\item Toute section d'un morphisme non ramifié est une immersion ouverte.
\item Toute section d'un morphisme séparé est une immersion fermée.
\item Toute section d'un morphisme non ramifié et séparé est ouverte et fermée.
\end{enumerate}
\end{proposition}

\begin{proof}
Fixons un schéma de base $S$.
Si $f : X' \to X$ est un $S$-morphisme quelconque, alors le graphe
$\Gamma_f : X' \to X' \times_S X$
s'obtient par changement de base de la diagonale
$\Delta_{X/S} : X \to X \times_S X$ le long de la projection
$X' \times_S X \to X \times_S X$.
Si $g : X \to S$ est séparé (resp. non ramifié),
alors la diagonale est une immersion fermée (resp. une immersion ouverte)
par Schémas, Définition \ref{schemes-definition-separated}
(resp.\ Morphismes, Lemme \ref{morphisms-lemma-diagonal-unramified-morphism}).
Il en va donc de même du graphe par changement de base (voir
Schémas, Lemme \ref{schemes-lemma-base-change-immersion}).
Dans le cas particulier $X' = S$, nous obtenons (1), resp.\ (2).
L'assertion (3) résulte de la combinaison de (1) et (2).
\end{proof}

\noindent
Nous pouvons maintenant décrire explicitement les sections des morphismes non ramifiés.

\begin{theorem}
\label{etale-theorem-sections-unramified-maps}
Soit $Y$ un schéma connexe.
Soit $f : X \to Y$ un morphisme non ramifié et séparé.
Toute section de $f$ est un isomorphisme sur une composante connexe.
Il existe une correspondance bijective
$$
\text{sections de }f
\leftrightarrow
\left\{
\begin{matrix}
\text{composantes connexes }X'\text{ de }X\text{ telles que}\\
\text{l'application induite }X' \to Y\text{ soit un isomorphisme}
\end{matrix}
\right\}
$$
En particulier, pour $x \in X$, il existe au plus une
section passant par $x$.
\end{theorem}

\begin{proof}
Cela découle directement de la Proposition \ref{etale-proposition-properties-sections}, assertion (3).
\end{proof}

\noindent
Le théorème précédent donne une idée de la « rigidité » des morphismes non ramifiés.
La proposition suivante en donne une autre manifestation ;
outre son intérêt propre, elle est également utile dans la
théorie du groupe fondamental algébrique (voir \cite[Expos\'e V]{SGA1}).
Voir aussi le résultat plus général
Morphismes, Lemme \ref{morphisms-lemma-value-at-one-point}.

\begin{proposition}
\label{etale-proposition-equality}
Soit $S$ un schéma.
Soit $\pi : X \to S$ un morphisme non ramifié et séparé.
Soit $Y$ un $S$-schéma et soit $y \in Y$ un point.
Soient $f, g : Y \to X$ deux $S$-morphismes. Supposons que
\begin{enumerate}
\item $Y$ soit connexe,
\item $x = f(y) = g(y)$, et
\item les applications induites $f^\sharp, g^\sharp : \kappa(x) \to \kappa(y)$
sur les corps résiduels soient égales.
\end{enumerate}
Alors $f = g$.
\end{proposition}

\begin{proof}
Les applications $f, g : Y \to X$ définissent des applications $f', g' : Y \to X_Y = Y \times_S X$
qui sont des sections du morphisme structural $X_Y \to Y$.
Remarquons que $f = g$ si et seulement si $f' = g'$.
Le morphisme structural $X_Y \to Y$ est le changement de base de $\pi$ et est donc
lui aussi non ramifié et séparé (voir
Morphismes, Lemmes \ref{morphisms-lemma-base-change-unramified} et
Schémas, Lemme \ref{schemes-lemma-separated-permanence}).
Ainsi, d'après le Théorème \ref{etale-theorem-sections-unramified-maps},
il suffit de démontrer que $f'$ et $g'$ passent par le même
point de $X_Y$. C'est précisément ce que les hypothèses (2) et (3)
garantissent, à savoir $f'(y) = g'(y) \in X_Y$.
\end{proof}

\begin{lemma}
\label{etale-lemma-finitely-many-maps-to-unramified}
Soit $S$ un schéma noethérien. Soit $X \to S$ un morphisme non ramifié quasi-compact.
Soit $Y \to S$ un morphisme, où $Y$ est noethérien. Alors
$\Mor_S(Y, X)$ est un ensemble fini.
\end{lemma}

\begin{proof}
Supposons d'abord que $X \to S$ soit séparé (ce qui est souvent le cas en pratique).
Puisque $Y$ est noethérien, il possède un nombre fini de composantes connexes. Nous
pouvons donc supposer $Y$ connexe. Choisissons un point $y \in Y$, d'image $s \in S$.
Comme $X \to S$ est non ramifié et quasi-compact,
la fibre $X_s$ est finie, disons $X_s = \{x_1, \ldots, x_n\}$,
et $\kappa(x_i)/\kappa(s)$ est une extension finie de corps.
Voir Morphismes, Lemmes \ref{morphisms-lemma-unramified-quasi-finite},
\ref{morphisms-lemma-residue-field-quasi-finite} et
\ref{morphisms-lemma-quasi-finite}.
Pour chaque $i$, il n'existe qu'un nombre fini d'homomorphismes de $\kappa(s)$-algèbres
$\kappa(x_i) \to \kappa(y)$ (par la théorie élémentaire des corps).
Ainsi, $\Mor_S(Y, X)$ est fini par
la Proposition \ref{etale-proposition-equality}.

\medskip\noindent
Cas général. Il existe un ouvert non vide $U \subset S$ tel
que $X_U \to U$ soit fini (donc, en particulier, séparé) ; voir
Morphismes, Lemme \ref{morphisms-lemma-generically-finite}
(ce lemme s'applique puisque nous avons déjà vu ci-dessus qu'un morphisme non ramifié
quasi-compact est quasi-fini et puisque $X \to S$ est quasi-séparé par
Morphismes, Lemme \ref{morphisms-lemma-finite-type-Noetherian-quasi-separated}).
Soit $Z \subset S$ le sous-schéma fermé réduit dont le support est
le complément de $U$. Par récurrence noethérienne, nous voyons que
$\Mor_Z(Y_Z, X_Z)$ est fini (détails omis).
D'après le résultat du premier paragraphe, l'ensemble
$\Mor_U(Y_U, X_U)$ est fini. Il suffit donc de montrer que
$$
\Mor_S(Y, X) \longrightarrow \Mor_Z(Y_Z, X_Z) \times \Mor_U(Y_U, X_U)
$$
est injective. Cela résulte du fait que le lieu où
deux morphismes $a, b : Y \to X$ coïncident est ouvert dans $Y$, puisque
la diagonale $\Delta : X \to X \times_S X$ est ouverte ; voir
Morphismes, Lemme \ref{morphisms-lemma-diagonal-unramified-morphism}.
\end{proof}







\section{Morphismes universellement injectifs et non ramifiés}
\label{etale-section-universally-injective-unramified}

\noindent
Rappelons qu'un morphisme de schémas $f : X \to Y$ est universellement
injectif si tout changement de base de $f$ est injectif (sur les espaces
topologiques sous-jacents) ; voir
Morphismes, Définition \ref{morphisms-definition-universally-injective}.
Les morphismes universellement injectifs et non ramifiés peuvent être
caractérisés comme suit.

\begin{lemma}
\label{etale-lemma-universally-injective-unramified}
Soit $f : X \to S$ un morphisme de schémas.
Les conditions suivantes sont équivalentes :
\begin{enumerate}
\item $f$ est non ramifié et est un monomorphisme,
\item $f$ est non ramifié et universellement injectif,
\item $f$ est localement de type fini et est un monomorphisme,
\item $f$ est universellement injectif, localement de type fini et
formellement non ramifié,
\item $f$ est localement de type fini et $X_s$ est vide
ou bien $X_s \to s$ est un isomorphisme, pour tout $s \in S$.
\end{enumerate}
\end{lemma}

\begin{proof}
Nous avons vu dans
Compléments sur les morphismes, Lemme
\ref{more-morphisms-lemma-unramified-formally-unramified},
qu'être formellement non ramifié et localement de type fini équivaut
à être non ramifié. Ainsi, (4) équivaut à (2).
Un monomorphisme est certainement universellement injectif et
formellement non ramifié ; (3) implique donc (4).
Il est clair que (1) implique (3). Enfin, si (2) est satisfaite, alors
$\Delta : X \to X \times_S X$ est à la fois une immersion ouverte
(Morphismes, Lemme \ref{morphisms-lemma-diagonal-unramified-morphism})
et surjective
(Morphismes, Lemme \ref{morphisms-lemma-universally-injective}) ;
c'est donc un isomorphisme, autrement dit $f$ est un monomorphisme. Ainsi,
(2) implique (1).

\medskip\noindent
La condition (3) implique (5), car les monomorphismes sont préservés par
changement de base
(Schémas, Lemme \ref{schemes-lemma-base-change-monomorphism})
et en vertu de la description des monomorphismes à valeurs dans les spectres de corps
donnée dans
Schémas, Lemme \ref{schemes-lemma-mono-towards-spec-field}.
La condition (5) implique (4) par
Morphismes, Lemmes \ref{morphisms-lemma-universally-injective} et
\ref{morphisms-lemma-unramified-etale-fibres}.
\end{proof}

\noindent
On en déduit la caractérisation utile suivante des immersions fermées.

\begin{lemma}
\label{etale-lemma-characterize-closed-immersion}
Soit $f : X \to S$ un morphisme de schémas.
Les conditions suivantes sont équivalentes :
\begin{enumerate}
\item $f$ est une immersion fermée,
\item $f$ est un monomorphisme propre,
\item $f$ est propre, non ramifié et universellement injectif,
\item $f$ est universellement fermé, non ramifié et est un monomorphisme,
\item $f$ est universellement fermé, non ramifié et universellement injectif,
\item $f$ est universellement fermé, localement de type fini et est un monomorphisme,
\item $f$ est universellement fermé, universellement injectif, localement de
type fini et formellement non ramifié.
\end{enumerate}
\end{lemma}

\begin{proof}
L'équivalence des conditions (4) à (7) résulte immédiatement du
Lemme \ref{etale-lemma-universally-injective-unramified}.

\medskip\noindent
Soit $f : X \to S$ un morphisme qui satisfait à (6). Alors $f$ est séparé ; voir
Schémas, Lemme \ref{schemes-lemma-monomorphism-separated},
et ses fibres sont finies. Ainsi,
Compléments sur les morphismes, Lemme \ref{more-morphisms-lemma-characterize-finite},
montre que $f$ est fini. Puis
Morphismes, Lemme \ref{morphisms-lemma-finite-monomorphism-closed},
implique que $f$ est une immersion fermée, c'est-à-dire que (1) est satisfaite.

\medskip\noindent
Remarquons que (1) $\Rightarrow$ (2), car une immersion fermée est
propre et est un monomorphisme
(Morphismes, Lemme \ref{morphisms-lemma-closed-immersion-proper},
et
Schémas, Lemme \ref{schemes-lemma-immersions-monomorphisms}).
Par le
Lemme \ref{etale-lemma-universally-injective-unramified},
nous voyons que (2) implique (3). Il est clair que (3) implique (5).
\end{proof}

\noindent
Voici un autre résultat du même genre.

\begin{lemma}
\label{etale-lemma-finite-unramified-one-point}
Soit $\pi : X \to S$ un morphisme de schémas. Soit $s \in S$.
Supposons que
\begin{enumerate}
\item $\pi$ soit fini,
\item $\pi$ soit non ramifié,
\item $\pi^{-1}(\{s\}) = \{x\}$, et
\item $\kappa(s) \subset \kappa(x)$ soit purement
inséparable\footnote{Compte tenu de la condition (2),
cela équivaut à $\kappa(s) = \kappa(x)$.}.
\end{enumerate}
Alors il existe un voisinage ouvert $U$ de $s$ tel que
$\pi|_{\pi^{-1}(U)} : \pi^{-1}(U) \to U$ soit une immersion fermée.
\end{lemma}

\begin{proof}
La question est locale sur $S$. Nous pouvons donc supposer que $S = \Spec(A)$.
Par définition d'un morphisme fini, cela implique que $X = \Spec(B)$.
Notons que l'homomorphisme d'anneaux $\varphi : A \to B$ définissant $\pi$
est un homomorphisme d'anneaux fini non ramifié.
Soit $\mathfrak p \subset A$ l'idéal premier correspondant à $s$.
Soit $\mathfrak q \subset B$ l'idéal premier correspondant à $x$.
Les conditions (2), (3) et (4) impliquent que
$B_{\mathfrak q}/\mathfrak pB_{\mathfrak q} = \kappa(\mathfrak p)$.
D'après Algèbre, Lemme \ref{algebra-lemma-unique-prime-over-localize-below},
nous avons $B_{\mathfrak q} = B_{\mathfrak p}$
(notons qu'un homomorphisme d'anneaux fini vérifie la propriété de montée, voir
Algèbre, section \ref{algebra-section-going-up}.)
Nous voyons donc que
$B_{\mathfrak p}/\mathfrak pB_{\mathfrak p} = \kappa(\mathfrak p)$.
Comme $B$ est un $A$-module fini, le lemme de Nakayama (voir
Algèbre, Lemme \ref{algebra-lemma-NAK})
montre que $B_{\mathfrak p} = \varphi(A_{\mathfrak p})$. Ainsi (en utilisant de
nouveau la finitude de $B$ comme $A$-module), il existe
$f \in A$, $f \not \in \mathfrak p$ tel que $B_f = \varphi(A_f)$,
ce qui est le résultat voulu.
\end{proof}

\noindent
Les résultats topologiques présentés ci-dessus serviront à donner une
caractérisation fonctorielle des morphismes \'etales analogue au Théorème
\ref{etale-theorem-formally-unramified}.




\section{Exemples de morphismes non ramifiés}
\label{etale-section-examples}

\noindent
Voici quelques exemples.

\begin{example}
\label{etale-example-etale-field-extensions}
Soit $k$ un corps.
Les morphismes quasi-compacts non ramifiés $X \to \Spec(k)$ sont affines.
En effet, $X$ est de dimension $0$ et noethérien,
donc c'est un ensemble discret fini, et chaque point fournit un ouvert affine ;
ainsi $X$ est une réunion disjointe finie de schémas affines, donc est affine.
Le lemme de normalisation de Noether entraîne que $X$ est le spectre d'une
$k$-algèbre finie $A$.
Cette algèbre est un produit d'extensions finies séparables de $k$.
Ainsi, un morphisme quasi-compact non ramifié vers $\Spec(k)$
correspond à un nombre fini d'extensions finies séparables de $k$.
En particulier, un morphisme non ramifié dont la source est connexe et la cible
réduite à un point correspond nécessairement à une extension finie séparable de corps.
Comme nous le verrons plus loin, $X \to \Spec(k)$ est \'etale si et
seulement s'il est non ramifié. Ainsi, dans ce cas au moins, nous obtenons une
description très simple de la topologie \'etale d'un schéma. Bien entendu, la
cohomologie de cette topologie est une autre affaire.
\end{example}

\begin{example}
\label{etale-example-standard-etale}
La propriété (3) du
Théorème \ref{etale-theorem-unramified-equivalence}
nous fournit une source canonique d'exemples de morphismes non ramifiés.
Fixons un anneau $R$ et un entier $n$. Soit $I = (g_1, \ldots, g_m)$ un
idéal de $R[x_1, \ldots, x_n]$. Soit $\mathfrak q \subset R[x_1, \ldots, x_n]$
un idéal premier. Supposons $I \subset \mathfrak q$ et que la matrice
$$
\left(\frac{\partial g_i}{\partial x_j}\right) \bmod \mathfrak q
\quad\in\quad
\text{Mat}(n \times m, \kappa(\mathfrak q))
$$
soit de rang $n$. Alors le morphisme
$f : Z = \Spec(R[x_1, \ldots, x_n]/I) \to \Spec(R)$
est non ramifié au point $x \in Z \subset \mathbf{A}^n_R$ correspondant
à $\mathfrak q$. Il est clair que nous devons avoir $m \geq n$.
Dans le cas extrême $m = n$, c'est-à-dire lorsque la différentielle de l'application
$\mathbf{A}^n_R \to \mathbf{A}^n_R$ définie par les $g_i$
est un isomorphisme des espaces tangents, alors $f$ est aussi plat
en $x$ et est, par conséquent, un morphisme \'etale (voir Algèbre,
Définition \ref{algebra-definition-standard-smooth},
Lemme \ref{algebra-lemma-standard-smooth} et
Exemple \ref{algebra-example-make-standard-smooth}).
\end{example}

\begin{example}
\label{etale-example-number-theory-etale}
Fixons une extension de corps de nombres $L/K$, dont les anneaux d'entiers sont
$\mathcal{O}_L$ et $\mathcal{O}_K$. L'injection $K \to L$ définit un
morphisme $f : \Spec(\mathcal{O}_L) \to \Spec(\mathcal{O}_K)$.
Comme expliqué ci-dessus, les points où $f$ est non ramifié au sens adopté ici
sont exactement ceux où $f$ est non ramifié au sens
usuel. Au sens usuel, le lieu de ramification dans
$\Spec(\mathcal{O}_L)$ est le lieu des zéros de la
différente, qui est un idéal de $\mathcal{O}_L$. En fait, la différente
n'est autre que l'annulateur du module
$\Omega_{\mathcal{O}_L/\mathcal{O}_K}$. De même, le
discriminant est un idéal de $\mathcal{O}_K$, à savoir la
norme de la différente.
Le lieu des zéros du discriminant est précisément l'ensemble
des points de $K$ qui se ramifient dans $L$.
Ainsi, en désignant par $X$ le complémentaire du sous-ensemble fermé
défini par la différente dans $\Spec(\mathcal{O}_L)$,
nous obtenons un morphisme $X \to \Spec(\mathcal{O}_K)$ qui est non ramifié.
De plus, ce morphisme est aussi plat, puisque tout homomorphisme local
d'anneaux de valuation discrète est plat, et ce morphisme est donc
en fait \'etale. Si $L/K$ est une extension galoisienne finie, en désignant par
$Y$ le complémentaire du sous-ensemble fermé défini par le discriminant dans
$\Spec(\mathcal{O}_K)$, nous voyons que nous obtenons même un
morphisme fini \'etale $X \to Y$.
C'est donc un exemple de revêtement \'etale fini.
\end{example}





\section{Morphismes plats}
\label{etale-section-flat-morphisms}

\noindent
Cette section a simplement pour objet de résumer les propriétés de la platitude qui
nous seront utiles. Nous nous contenterons donc d'énoncer précisément les théorèmes
et de donner des références pour leurs démonstrations.

\medskip\noindent
Après avoir brièvement rappelé les faits nécessaires concernant les modules plats sur les anneaux
noethériens, nous énonçons un théorème de Grothendieck qui donne des conditions suffisantes
pour que les « sections hyperplanes » de certains modules soient plates.

\begin{definition}
\label{etale-definition-flat-rings}
Platitude des modules et des anneaux.
\begin{enumerate}
\item Un module $N$ sur un anneau $A$ est dit {\it plat}
si le foncteur $M \mapsto M \otimes_A N$ est exact.
\item Si ce foncteur est également fidèle, nous disons que
$N$ est {\it fidèlement plat} sur $A$.
\item Un morphisme d'anneaux $f : A \to B$ est dit
{\it plat (resp. fidèlement plat)}
si le foncteur $M \mapsto M \otimes_A B$ est exact
(resp. fidèle et exact).
\end{enumerate}
\end{definition}

\noindent
Voici une liste de faits, accompagnés de références au chapitre d'algèbre.
\begin{enumerate}
\item Les modules libres et projectifs sont plats. Cela est clair pour les modules libres
et s'en déduit pour les modules projectifs, car ceux-ci sont facteurs directs de modules
libres et $\otimes$ commute aux sommes directes.
\item La platitude est une propriété locale, c'est-à-dire que $M$ est plat sur $A$
si et seulement si $M_{\mathfrak p}$ est plat sur $A_{\mathfrak p}$ pour tout
$\mathfrak p \in \Spec(A)$.
Voir Algèbre, Lemme \ref{algebra-lemma-flat-localization}.
\item Si $M$ est un $A$-module plat et $A \to B$ un homomorphisme d'anneaux,
alors $M \otimes_A B$ est un $B$-module plat. Voir
Algèbre, Lemme \ref{algebra-lemma-flat-base-change}.
\item Les modules finis et plats sur les anneaux locaux sont libres.
Voir Algèbre, Lemme \ref{algebra-lemma-finite-flat-local}.
\item Si $f : A \to B$ est un morphisme d'anneaux quelconques,
$f$ est plat si et seulement si les homomorphismes induits
$A_{f^{-1}(\mathfrak q)} \to B_{\mathfrak q}$ sont plats pour tout
$\mathfrak q \in \Spec(B)$.
Voir Algèbre, Lemme \ref{algebra-lemma-flat-localization}.
\item Si $f : A \to B$ est un homomorphisme local d'anneaux locaux,
$f$ est plat si et seulement s'il est fidèlement plat.
Voir Algèbre, Lemme \ref{algebra-lemma-local-flat-ff}.
\item Un homomorphisme d'anneaux $A \to B$ est fidèlement plat si et seulement s'il est
plat et si l'application induite sur les spectres est surjective.
Voir Algèbre, Lemme \ref{algebra-lemma-ff-rings}.
\item Si $A$ est un anneau local noethérien, le complété
$A^\wedge$ est fidèlement plat sur $A$.
Voir Algèbre, Lemme \ref{algebra-lemma-completion-faithfully-flat}.
\item Soient $A$ un anneau local noethérien et $M$ un $A$-module.
Alors $M$ est plat sur $A$ si et seulement si $M \otimes_A A^\wedge$
est plat sur $A^\wedge$. (Combiner l'assertion précédente avec
Algèbre, Lemme \ref{algebra-lemma-flatness-descends}.)
\end{enumerate}
Avant de passer à la catégorie géométrique, nous présentons le théorème de
Grothendieck, qui fournit un procédé commode pour construire des modules
plats.

\begin{theorem}
\label{etale-theorem-flatness-grothendieck}
Soient $A$, $B$ des anneaux locaux noethériens.
Soit $f : A \to B$ un homomorphisme local.
Si $M$ est un $B$-module fini qui est plat comme $A$-module,
et si $t \in \mathfrak m_B$ est un élément tel que la multiplication
par $t$ soit injective sur $M/\mathfrak m_AM$, alors $M/tM$ est également plat sur $A$.
\end{theorem}

\begin{proof}
Voir Algèbre, Lemme \ref{algebra-lemma-mod-injective}.
Voir aussi \cite[section 20]{MatCA}.
\end{proof}

\begin{definition}
\label{etale-definition-flat-schemes}
(Voir Morphismes, Définition \ref{morphisms-definition-flat}.)
Soit $f : X \to Y$ un morphisme de schémas.
Soit $\mathcal{F}$ un $\mathcal{O}_X$-module quasi-cohérent.
\begin{enumerate}
\item Soit $x \in X$. Nous disons que $\mathcal{F}$ est
{\it plat sur $Y$ en $x \in X$} si $\mathcal{F}_x$
est un $\mathcal{O}_{Y, f(x)}$-module plat.
On utilise l'homomorphisme $\mathcal{O}_{Y, f(x)} \to \mathcal{O}_{X, x}$ pour
considérer $\mathcal{F}_x$ comme un $\mathcal{O}_{Y, f(x)}$-module.
\item Soit $x \in X$. Nous disons que $f$ est {\it plat en $x \in X$}
si $\mathcal{O}_{Y, f(x)} \to \mathcal{O}_{X, x}$ est plat.
\item Nous disons que $f$ est {\it plat} s'il est plat en tout point de $X$.
\item Un morphisme $f : X \to Y$ plat et surjectif est parfois
dit {\it fidèlement plat}.
\end{enumerate}
\end{definition}

\noindent
Voici, une fois encore, une liste de résultats :
\begin{enumerate}
\item La propriété (pour un morphisme) d'être plat est, par définition,
locale pour la topologie de Zariski sur la source et la cible.
\item Les immersions ouvertes sont plates. (C'est clair, puisqu'elles induisent des isomorphismes
sur les anneaux locaux.)
\item Les morphismes plats sont stables par changement de base et par composition.
Morphismes, Lemmes \ref{morphisms-lemma-base-change-flat} et
\ref{morphisms-lemma-composition-flat}.
\item Si $f : X \to Y$ est plat, alors le foncteur image inverse
$\QCoh(\mathcal{O}_Y) \to \QCoh(\mathcal{O}_X)$ est exact.
Cela résulte immédiatement de l'examen des germes.
\item Soit $f : X \to Y$ un morphisme de schémas, et supposons $Y$
quasi-compact et quasi-séparé. Dans ce cas,
si le foncteur $f^*$ est exact, alors $f$ est plat.
(Démonstration omise. Indication : utiliser
Propriétés, Lemme \ref{properties-lemma-extend-trivial} pour voir que
$Y$ possède « assez » de faisceaux d'idéaux et utiliser la caractérisation de la
platitude dans Algèbre, Lemme \ref{algebra-lemma-flat}.)
\end{enumerate}



\section{Propriétés topologiques des morphismes plats}
\label{etale-section-topological-flat}

\noindent
Nous « rappelons » ci-dessous quelques propriétés d'ouverture dont jouissent les morphismes plats.

\begin{theorem}
\label{etale-theorem-flat-open}
Soit $Y$ un schéma localement noethérien.
Soit $f : X \to Y$ un morphisme localement de type fini.
Soit $\mathcal{F}$ un $\mathcal{O}_X$-module cohérent.
L'ensemble des points de $X$ où $\mathcal{F}$ est plat sur $Y$ est ouvert.
En particulier, l'ensemble des points où $f$ est plat est ouvert dans $X$.
\end{theorem}

\begin{proof}
Voir Compléments sur les morphismes, Théorème \ref{more-morphisms-theorem-openness-flatness}.
\end{proof}

\begin{theorem}
\label{etale-theorem-flat-map-open}
Soit $Y$ un schéma localement noethérien.
Soit $f : X \to Y$ un morphisme plat et localement de type fini.
Alors $f$ est (universellement) ouvert.
\end{theorem}

\begin{proof}
Voir Morphismes, Lemme \ref{morphisms-lemma-fppf-open}.
\end{proof}

\begin{theorem}
\label{etale-theorem-flat-is-quotient}
Un morphisme fidèlement plat quasi-compact est une application quotient pour
la topologie de Zariski.
\end{theorem}

\begin{proof}
Voir Morphismes, Lemme \ref{morphisms-lemma-fpqc-quotient-topology}.
\end{proof}

\noindent
Une raison importante d'étudier les morphismes plats est qu'ils fournissent le cadre
adéquat pour saisir la notion de famille de schémas paramétrée par les
points d'un autre schéma. Naïvement, on pourrait penser que tout morphisme $f : X \to S$
doit être considéré comme une famille paramétrée par les points de $S$. Cependant,
sans hypothèse de platitude sur $f$, des phénomènes véritablement bizarres peuvent se produire dans
cette prétendue famille. Par exemple, rien ne garantit que la dimension relative
(la dimension des fibres) soit constante dans une famille. D'autres invariants numériques,
comme le polynôme de Hilbert, peuvent eux aussi varier d'une fibre à
l'autre. La platitude empêche de tels phénomènes et confère donc
une certaine « continuité » aux fibres.


\section{Morphismes \'etales}
\label{etale-section-etale-morphisms}

\noindent
Dans cette section, nous allons définir les morphismes \'etales et démontrer plusieurs
propriétés importantes à leur sujet. La plus importante est sans aucun doute la
caractérisation fonctorielle présentée dans le Théorème \ref{etale-theorem-formally-etale}.
Ensuite, nous examinerons aussi quelques propriétés des anneaux qui sont
insensibles à une extension \'etale (propriétés qui valent pour un anneau
si et seulement si elles valent pour toutes ses extensions \'etales), afin de motiver le
principe fondamental de la cohomologie \'etale --- les morphismes \'etales sont l'analogue algébrique des
isomorphismes locaux.

\medskip\noindent
Comme le titre le suggère, nous allons définir la classe des morphismes \'etales --- la
classe des morphismes (dont nous considérerons les familles surjectives comme des recouvrements)
dans la catégorie des schémas au-dessus d'un schéma de base $S$, afin de définir le
site \'etale $S_\etale$. Intuitivement, un morphisme \'etale est censé
exprimer l'idée d'un espace de revêtement et devrait donc être proche d'un
isomorphisme local. Si nous travaillons avec des variétés sur des corps algébriquement clos,
cette dernière assertion peut servir de définition, à condition de remplacer
« isomorphisme local » par « isomorphisme local formel » (isomorphisme après
complétion). On peut alors donner une définition sur un corps de base quelconque en demandant
que le changement de base à la clôture algébrique soit \'etale (au
sens précité). Mais, plutôt que de procéder par de telles constructions peu
esthétiques, nous adopterons une approche algébrique plus nette, quoique légèrement plus
abstraite.

\medskip\noindent
Nous définissons d'abord les « homomorphismes \'etales d'anneaux locaux » pour les anneaux
locaux noethériens. Nous ne pouvons pas employer simplement le terme « \'etale », puisqu'il
existe déjà une notion d'homomorphisme d'anneaux \'etale
(Algèbre, section \ref{algebra-section-etale})
et qu'elle est différente.

\begin{definition}
\label{etale-definition-etale-ring}
Soient $A$, $B$ des anneaux locaux noethériens.
Un homomorphisme local $f : A \to B$ est appelé
{\it homomorphisme \'etale d'anneaux locaux}
s'il est plat et non ramifié en tant qu'homomorphisme d'anneaux locaux
(voir Définition \ref{etale-definition-unramified-rings}).
\end{definition}

\noindent
Il s'agit de la version locale de la définition d'un homomorphisme d'anneaux \'etale dans
Algèbre, section \ref{algebra-section-etale}.
La définition précise
donnée dans cette section est celle d'un homomorphisme d'anneaux lisse de dimension
relative $0$. Il est démontré (dans
Algèbre, Lemme \ref{algebra-lemma-etale-standard-smooth})
qu'une $R$-algèbre \'etale $S$ possède toujours une présentation
$$
S = R[x_1, \ldots, x_n]/(f_1, \ldots, f_n)
$$
telle que
$$
g =
\det
\left(
\begin{matrix}
\partial f_1/\partial x_1 &
\partial f_2/\partial x_1 &
\ldots &
\partial f_n/\partial x_1 \\
\partial f_1/\partial x_2 &
\partial f_2/\partial x_2 &
\ldots &
\partial f_n/\partial x_2 \\
\ldots & \ldots & \ldots & \ldots \\
\partial f_1/\partial x_n &
\partial f_2/\partial x_n &
\ldots &
\partial f_n/\partial x_n
\end{matrix}
\right)
$$
ait pour image un élément inversible de $S$. Les deux lemmes suivants relient les deux
notions.

\begin{lemma}
\label{etale-lemma-characterize-etale-Noetherian}
Soit $A \to B$ un homomorphisme de type fini, où $A$ est un anneau noethérien.
Soit $\mathfrak q$ un idéal premier de $B$ au-dessus de $\mathfrak p \subset A$.
Alors $A \to B$ est \'etale en $\mathfrak q$ si et seulement si
$A_{\mathfrak p} \to B_{\mathfrak q}$ est un homomorphisme \'etale
d'anneaux locaux.
\end{lemma}

\begin{proof}
Voir Algèbre, Lemmes \ref{algebra-lemma-etale} (platitude des homomorphismes \'etales),
\ref{algebra-lemma-etale-at-prime} (les homomorphismes \'etales sont non ramifiés)
et \ref{algebra-lemma-characterize-etale} (les homomorphismes plats et non ramifiés
sont \'etales).
\end{proof}

\begin{lemma}
\label{etale-lemma-characterize-etale-completions}
Soient $A$, $B$ des anneaux locaux noethériens.
Soit $A \to B$ un homomorphisme local tel que $B$ soit essentiellement de
type fini sur $A$.
Les assertions suivantes sont équivalentes :
\begin{enumerate}
\item $A \to B$ est un homomorphisme \'etale d'anneaux locaux
\item $A^\wedge \to B^\wedge$ est un homomorphisme \'etale d'anneaux locaux, et
\item $A^\wedge \to B^\wedge$ est \'etale.
\end{enumerate}
De plus, dans ce cas $B^\wedge \cong (A^\wedge)^{\oplus n}$ comme
$A^\wedge$-modules pour un certain $n \geq 1$.
\end{lemma}

\begin{proof}
Pour établir les équivalences de (1), (2) et (3), puisque nous disposons des résultats
correspondants pour les homomorphismes d'anneaux non ramifiés
(Lemme \ref{etale-lemma-characterize-unramified-completions}),
il suffit de démontrer que
$A \to B$ est plat si et seulement si $A^\wedge \to B^\wedge$ est plat.
Cela ressort de nos listes de propriétés des homomorphismes plats, puisque
les homomorphismes d'anneaux $A \to A^\wedge$ et $B \to B^\wedge$ sont fidèlement plats.
Pour la dernière assertion, Lemme \ref{etale-lemma-unramified-completions}
montre que $B^\wedge$ est un $A^\wedge$-module fini et plat.
Il est donc fini et libre d'après notre liste
des propriétés des modules plats de la section \ref{etale-section-flat-morphisms}.
\end{proof}

\noindent
L'entier $n$ qui intervient dans le lemme ci-dessus
n'est autre que le degré
$[\kappa(\mathfrak m_B) : \kappa(\mathfrak m_A)]$ de l'extension des corps
résiduels. En particulier, si $\kappa(\mathfrak m_A)$
est séparablement clos, nous voyons que $A^\wedge \to B^\wedge$
est un isomorphisme, ce qui confirme nos assertions précédentes.

\begin{definition}
\label{etale-definition-etale-schemes-1}
(Voir Morphismes, Définition \ref{morphisms-definition-etale}.)
Soit $Y$ un schéma localement noethérien.
Soit $f : X \to Y$ un morphisme de schémas localement de type fini.
\begin{enumerate}
\item Soit $x \in X$. Nous disons que $f$ est {\it \'etale en $x \in X$} si
$\mathcal{O}_{Y, f(x)} \to \mathcal{O}_{X, x}$ est un
homomorphisme \'etale d'anneaux locaux.
\item Le morphisme est dit {\it \'etale} s'il est \'etale en chacun de ses
points.
\end{enumerate}
\end{definition}

\noindent
Démontrons que cette définition coïncide avec celle du
chapitre sur les morphismes de schémas. Cela garantit en particulier que
l'ensemble des points où un morphisme est \'etale est ouvert.

\begin{lemma}
\label{etale-lemma-etale-definition}
Soit $Y$ un schéma localement noethérien.
Soit $f : X \to Y$ localement de type fini.
Soit $x \in X$. Le morphisme $f$ est \'etale en $x$ au
sens de la Définition \ref{etale-definition-etale-schemes-1}
si et seulement s'il est \'etale en $x$ au
sens de Morphismes, Définition \ref{morphisms-definition-etale}.
\end{lemma}

\begin{proof}
Cela résulte du Lemme \ref{etale-lemma-characterize-etale-Noetherian}
et des définitions.
\end{proof}

\noindent
Voici quelques résultats sur les morphismes \'etales.
Les formulations données dans cette liste ne s'appliquent qu'aux
morphismes localement de type fini entre schémas localement noethériens.
Dans chaque cas, nous donnons une référence au résultat général tel qu'il a été
démontré plus tôt dans le projet, mais on peut, dans certains cas,
démontrer plus facilement le résultat dans le cas noethérien.
Voici la liste :
\begin{enumerate}
\item Un morphisme \'etale est non ramifié. (Cela ressort de nos définitions.)
\item La propriété d'être \'etale est locale sur la source et la cible pour la
topologie de Zariski.
\item Les morphismes \'etales sont stables par changement de base et par composition.
Voir Morphismes, Lemmes \ref{morphisms-lemma-base-change-etale}
et \ref{morphisms-lemma-composition-etale}.
\item Les morphismes \'etales de schémas sont localement quasi-finis,
et les morphismes \'etales quasi-compacts sont quasi-finis. (Cela est
vrai parce que c'est le cas pour les morphismes non ramifiés, comme vu précédemment.)
\item Les morphismes \'etales sont de dimension relative $0$. Voir
Morphismes, Définition \ref{morphisms-definition-relative-dimension-d}
et
Morphismes, Lemme \ref{morphisms-lemma-locally-quasi-finite-rel-dimension-0}.
\item Un morphisme est \'etale si et seulement s'il est plat et si
toutes ses fibres sont \'etales. Voir
Morphismes, Lemme \ref{morphisms-lemma-etale-flat-etale-fibres}.
\item Les morphismes \'etales sont ouverts. C'est vrai parce qu'un morphisme \'etale
est plat, d'après le Théorème \ref{etale-theorem-flat-map-open}.
\item Soient $X$ et $Y$ \'etales sur un schéma de base $S$.
Tout $S$-morphisme de $X$ vers $Y$ est \'etale.
Voir Morphismes, Lemme \ref{morphisms-lemma-etale-permanence}.
\end{enumerate}






\section{Le théorème de structure}
\label{etale-section-structure-etale-map}

\noindent
Nous présentons un théorème qui décrit la structure locale des morphismes \'etales
et non ramifiés. Outre son importance propre évidente,
ce théorème nous permet aussi de passer à une autre
définition des morphismes \'etales qui rend mieux compte de l'intuition géométrique
que celle que nous avons employée jusqu'ici.

\medskip\noindent
Pour l'énoncer, nous avons besoin de la notion d'{\it homomorphisme d'anneaux \'etale standard}, voir
Algèbre, Définition \ref{algebra-definition-standard-etale}.
Plus précisément, supposons que $R$ soit un anneau et que $f, g \in R[t]$ soient des polynômes
tels que
\begin{enumerate}
\item[(a)] $f$ soit un polynôme unitaire, et
\item[(b)] $f' = \text{d}f/\text{d}t$ soit inversible dans le localisé
$R[t]_g/(f)$.
\end{enumerate}
Alors l'application
$$
R \longrightarrow R[t]_g/(f) = R[t, 1/g]/(f)
$$
est une algèbre \'etale standard, et toute algèbre \'etale standard est isomorphe
à l'une d'elles. C'est un exercice agréable de démontrer qu'un tel homomorphisme d'anneaux
est plat et non ramifié, donc \'etale (comme prévu, bien entendu).
Un cas particulier d'homomorphisme d'anneaux \'etale standard est tout homomorphisme d'anneaux
$$
R \longrightarrow R[t]_{f'}/(f) = R[t, 1/f']/(f)
$$
où $f$ est un polynôme unitaire, et toute algèbre \'etale standard est (isomorphe à)
un localisé principal de l'une d'elles.

\begin{theorem}
\label{etale-theorem-structure-etale}
Soit $f : A \to B$ un homomorphisme \'etale d'anneaux locaux.
Alors il existe $f, g \in A[t]$ tels que
\begin{enumerate}
\item $B' = A[t]_g/(f)$ soit une algèbre \'etale standard --- voir (a) et (b) ci-dessus, et
\item $B$ soit isomorphe au localisé de $B'$ en un idéal premier.
\end{enumerate}
\end{theorem}

\begin{proof}
Écrivons $B = B'_{\mathfrak q}$ pour une certaine $A$-algèbre de type fini $B'$
(ce qui est possible puisque $B$ est essentiellement de type fini sur $A$).
D'après le Lemme \ref{etale-lemma-characterize-etale-Noetherian},
nous voyons que $A \to B'$ est \'etale en $\mathfrak q$.
Nous pouvons donc appliquer
Algèbre, Proposition \ref{algebra-proposition-etale-locally-standard}
pour voir qu'un localisé principal de $B'$ est \'etale standard.
\end{proof}

\noindent
Voici la version pour les homomorphismes non ramifiés d'anneaux locaux.

\begin{theorem}
\label{etale-theorem-structure-unramified}
Soit $f : A \to B$ un homomorphisme non ramifié d'anneaux locaux.
Alors il existe $f, g \in A[t]$ tels que
\begin{enumerate}
\item $B' = A[t]_g/(f)$ soit une algèbre \'etale standard --- voir (a) et (b) ci-dessus, et
\item $B$ soit isomorphe à un quotient d'un localisé de $B'$ en un idéal premier.
\end{enumerate}
\end{theorem}

\begin{proof}
Écrivons $B = B'_{\mathfrak q}$ pour une certaine $A$-algèbre de type fini $B'$
(ce qui est possible puisque $B$ est essentiellement de type fini sur $A$).
D'après le Lemme \ref{etale-lemma-characterize-unramified-Noetherian},
nous voyons que $A \to B'$ est non ramifié en $\mathfrak q$.
Nous pouvons donc appliquer
Algèbre, Proposition \ref{algebra-proposition-unramified-locally-standard}
pour voir qu'un localisé principal de $B'$ est un quotient d'une
$A$-algèbre \'etale standard.
\end{proof}

\noindent
Des arguments usuels de relèvement donnent alors l'énoncé géométrique
suivant, qui nous sera d'une utilité essentielle.

\begin{theorem}
\label{etale-theorem-geometric-structure}
Soit $\varphi : X \to Y$ un morphisme de schémas. Soit $x \in X$.
Soit $V \subset Y$ un voisinage ouvert affine de $\varphi(x)$.
Si $\varphi$ est \'etale en $x$, alors il existe un ouvert affine
$U \subset X$ avec $x \in U$ et $\varphi(U) \subset V$
tel que l'on ait le diagramme suivant :
$$
\xymatrix{
X \ar[d] & U \ar[l] \ar[d] \ar[r]_-j & \Spec(R[t]_{f'}/(f)) \ar[d] \\
Y & V \ar[l] \ar@{=}[r] & \Spec(R)
}
$$
où $j$ est une immersion ouverte et $f \in R[t]$ est un polynôme unitaire.
\end{theorem}

\begin{proof}
Cet énoncé équivaut à celui de
Morphismes, Lemme \ref{morphisms-lemma-etale-locally-standard-etale},
bien que les énoncés diffèrent légèrement.
Voir aussi Variétés, Lemme \ref{varieties-lemma-geometric-structure-unramified}
pour une variante relative aux morphismes non ramifiés.
\end{proof}


\section{Morphismes \'etales et lisses}
\label{etale-section-etale-smooth}

\noindent
Un morphisme \'etale est lisse de dimension relative zéro.
La projection $\mathbf{A}^n_S \to S$ est un exemple standard
de morphisme lisse de dimension relative $n$.
Il se trouve que, localement pour la topologie \'etale, tout morphisme lisse
est de cette forme. Voici l'énoncé précis.

\begin{theorem}
\label{etale-theorem-smooth-etale-over-n-space}
Soit $\varphi : X \to Y$ un morphisme de schémas.
Soit $x \in X$.
Si $\varphi$ est lisse en $x$, alors
il existe un entier $n \geq 0$ et des ouverts affines
$V \subset Y$ et $U \subset X$ avec $x \in U$ et $\varphi(U) \subset V$
de sorte qu'il existe un diagramme commutatif
$$
\xymatrix{
X \ar[d] & U \ar[l] \ar[d] \ar[r]_-\pi &
\mathbf{A}^n_R \ar[d] \ar@{=}[r] &  \Spec(R[x_1, \ldots, x_n]) \ar[dl] \\
Y & V \ar[l] \ar@{=}[r] & \Spec(R)
}
$$
où $\pi$ est \'etale.
\end{theorem}

\begin{proof}
Voir
Morphismes, Lemme \ref{morphisms-lemma-smooth-etale-over-affine-space}.
\end{proof}




\section{Propriétés topologiques des morphismes \'etales}
\label{etale-section-topological-etale}

\noindent
Nous présentons quelques propriétés topologiques des morphismes \'etales et
non ramifiés. Tout d'abord, nous donnons ce que Grothendieck
appelle la {\it propriété fondamentale des morphismes \'etales}, voir
\cite[Expos\'e I.5]{SGA1}.

\begin{theorem}
\label{etale-theorem-etale-radicial-open}
Soit $f : X \to Y$ un morphisme de schémas.
Les assertions suivantes sont équivalentes :
\begin{enumerate}
\item $f$ est une immersion ouverte,
\item $f$ est universellement injectif et \'etale, et
\item $f$ est un monomorphisme plat, localement de présentation finie.
\end{enumerate}
\end{theorem}

\begin{proof}
Une immersion ouverte est universellement injective,
car tout changement de base d'une immersion ouverte
est une immersion ouverte. De plus, elle est \'etale d'après
Morphismes, Lemme \ref{morphisms-lemma-open-immersion-etale}.
Ainsi (1) implique (2).

\medskip\noindent
Supposons que $f$ soit universellement injectif et \'etale.
Comme $f$ est \'etale, il est plat et localement de présentation finie, voir
Morphismes, Lemmes \ref{morphisms-lemma-etale-flat} et
\ref{morphisms-lemma-etale-locally-finite-presentation}.
D'après
Lemme \ref{etale-lemma-universally-injective-unramified},
nous voyons que $f$ est un monomorphisme. Ainsi (2) implique (3).

\medskip\noindent
Supposons que $f$ soit plat, localement de présentation finie et un monomorphisme.
Alors $f$ est ouvert, voir
Morphismes, Lemme \ref{morphisms-lemma-fppf-open}.
Nous pouvons donc remplacer $Y$ par $f(X)$ et supposer que $f$ est
surjectif. Alors $f$ est ouvert et bijectif, donc c'est un homéomorphisme.
Par conséquent, $f$ est quasi-compact. Dès lors,
Descente, lemme
\ref{descent-lemma-flat-surjective-quasi-compact-monomorphism-isomorphism}
montre que $f$ est un isomorphisme, ce qui conclut.
\end{proof}

\noindent
Voici un autre résultat de même nature.

\begin{lemma}
\label{etale-lemma-finite-etale-one-point}
Soit $\pi : X \to S$ un morphisme de schémas. Soit $s \in S$.
Supposons que
\begin{enumerate}
\item $\pi$ soit fini ;
\item $\pi$ soit \'etale ;
\item $\pi^{-1}(\{s\}) = \{x\}$ ; et
\item $\kappa(s) \subset \kappa(x)$ soit purement
inséparable\footnote{Compte tenu de la condition (2)
cela équivaut à $\kappa(s) = \kappa(x)$.}.
\end{enumerate}
Alors il existe un voisinage ouvert $U$ de $s$ tel que
$\pi|_{\pi^{-1}(U)} : \pi^{-1}(U) \to U$ soit un isomorphisme.
\end{lemma}

\begin{proof}
D'après le
lemme \ref{etale-lemma-finite-unramified-one-point},
il existe un voisinage ouvert $U$ de $s$ tel que
$\pi|_{\pi^{-1}(U)} : \pi^{-1}(U) \to U$ soit une immersion fermée.
Mais un morphisme qui est à la fois \'etale et une immersion fermée est une
immersion ouverte (par exemple d'après le
théorème \ref{etale-theorem-etale-radicial-open}).
En rétrécissant $U$, on obtient donc un isomorphisme.
\end{proof}

\begin{lemma}
\label{etale-lemma-relative-frobenius-etale}
Soit $U \to X$ un morphisme \'etale de schémas
où $X$ est un schéma de caractéristique $p$.
Alors le Frobenius relatif $F_{U/X} : U \to U \times_{X, F_X} X$
est un isomorphisme.
\end{lemma}

\begin{proof}
Le morphisme $F_{U/X}$ est un homéomorphisme universel d'après
Variétés, lemme \ref{varieties-lemma-relative-frobenius}.
Le morphisme $F_{U/X}$ est \'etale en tant que
morphisme entre des schémas \'etales sur $X$
(Morphismes, lemme \ref{morphisms-lemma-etale-permanence}).
Par conséquent, $F_{U/X}$ est un isomorphisme d'après le
théorème \ref{etale-theorem-etale-radicial-open}.
\end{proof}






\section{Invariance topologique de la topologie \'etale}
\label{etale-section-topological-invariance}

\noindent
Nous présentons maintenant un théorème absolument crucial qui affirme, en gros,
que le caractère \'etale est une propriété topologique.

\begin{theorem}
\label{etale-theorem-etale-topological}
Soient $X$ et $Y$ deux schémas sur un schéma de base $S$. Soit $S_0$ un sous-schéma fermé
de $S$ ayant le même espace topologique sous-jacent
(par exemple si le faisceau d'idéaux de $S_0$ dans $S$ est de carré nul).
Notons $X_0$ (resp.\ $Y_0$) le changement de base $S_0 \times_S X$
(resp.\ $S_0 \times_S Y$).
Si $X$ est \'etale sur $S$, alors l'application
$$
\Mor_S(Y, X) \longrightarrow \Mor_{S_0}(Y_0, X_0)
$$
est bijective.
\end{theorem}

\begin{proof}
Après le changement de base par $Y \to S$, nous pouvons supposer que $Y = S$.
Dans ce cas, le théorème affirme que tout $S$-morphisme $\sigma_0 : S_0 \to X$
se factorise en fait de manière unique par une section $S \to X$ du
morphisme structural \'etale $f : X \to S$.

\medskip\noindent
Unicité. Supposons que nous ayons deux sections $\sigma, \sigma'$
par lesquelles $\sigma_0$ se factorise. Comme $X \to S$ est \'etale,
nous voyons que $\Delta : X \to X \times_S X$ est une immersion ouverte
(Morphismes, lemme \ref{morphisms-lemma-diagonal-unramified-morphism}).
Le morphisme $(\sigma, \sigma') : S \to X \times_S X$ se factorise
par cet ouvert, car, pour tout $s \in S$, nous avons
$(\sigma, \sigma')(s) = (\sigma_0(s), \sigma_0(s))$. Ainsi,
$\sigma = \sigma'$.

\medskip\noindent
Pour démontrer l'existence, nous nous ramenons d'abord au cas affine
(nous suggérons au lecteur de sauter cette étape).
Soit $X = \bigcup X_i$ un recouvrement par des ouverts affines tel
que chaque $X_i$ s'envoie dans un ouvert affine $S_i$ de $S$.
Pour tout $s \in S$, nous pouvons choisir $i$ tel que
$\sigma_0(s) \in X_i$.
Choisissons un voisinage ouvert affine $U \subset S_i$ de $s$
tel que $\sigma_0(U_0) \subset X_{i, 0}$. Remarquons que
$X' = X_i \times_S U = X_i \times_{S_i} U$ est affine.
Si l'on peut relever $\sigma_0|_{U_0} : U_0 \to X'_0$ en
un morphisme $U \to X'$, alors, par unicité, ces relèvements locaux se recollent
en un morphisme global $S \to X$. Nous pouvons donc supposer $S$ et
$X$ affines.

\medskip\noindent
Existence lorsque $S$ et $X$ sont affines. Écrivons $S = \Spec(A)$
et $X = \Spec(B)$. Alors $A \to B$ est \'etale et, en particulier,
lisse (de dimension relative $0$). Puisque $|S_0| = |S|$, nous voyons
que $S_0 = \Spec(A/I)$, où $I \subset A$ est localement nilpotent.
L'existence résulte donc de
Algèbre, lemme \ref{algebra-lemma-smooth-strong-lift}.
\end{proof}

\noindent
La démonstration du théorème précédent nous donne aussi une implication de la
caractérisation fonctorielle promise des morphismes \'etales. Le
théorème suivant sera renforcé dans
Cohomologie \'etale,
théorème \ref{etale-cohomology-theorem-topological-invariance}.

\begin{theorem}[Une équivalence remarquable de cat\'egories]
\label{etale-theorem-remarkable-equivalence}
\begin{reference}
\cite[IV, Theorem 18.1.2]{EGA}
\end{reference}
Soit $S$ un schéma.
Soit $S_0 \subset S$ un sous-schéma fermé ayant le même espace
topologique sous-jacent (par exemple si le faisceau d'idéaux de $S_0$ dans $S$
est de carré nul). Le foncteur
$$
X \longmapsto X_0 = S_0 \times_S X
$$
définit une équivalence de catégories
$$
\{
\text{schémas }X\text{ \'etales sur }S
\}
\leftrightarrow
\{
\text{schémas }X_0\text{ \'etales sur }S_0
\}
$$
\end{theorem}

\begin{proof}
Le théorème \ref{etale-theorem-etale-topological}
montre que ce foncteur est pleinement fidèle.
Il reste à démontrer qu'il est essentiellement surjectif.
Soit $Y \to S_0$ un morphisme \'etale de schémas.

\medskip\noindent
Supposons le résultat vrai lorsque $S$ et $Y$ sont affines.
Dans ce cas, choisissons un recouvrement par des ouverts affines
$Y = \bigcup V_j$ tel que chaque $V_j$ s'envoie
dans un ouvert affine de $S$. Par hypothèse (cas affine), nous pouvons
trouver des morphismes \'etales $W_j \to S$ tels que $W_{j, 0} \cong V_j$
(comme schémas sur $S_0$). Soit $W_{j, j'} \subset W_j$
le sous-schéma ouvert dont l'espace topologique sous-jacent
correspond à $V_j \cap V_{j'}$. Puisque nous avons des isomorphismes
$$
W_{j, j', 0} \cong V_j \cap V_{j'} \cong W_{j', j, 0}
$$
de schémas sur $S_0$, la pleine fidélité montre que nous
obtenons des isomorphismes
$\theta_{j, j'} : W_{j, j'} \to W_{j', j}$ de schémas sur $S$.
Nous omettons de vérifier que ces isomorphismes satisfont à la
condition de cocycle de Schémas, section \ref{schemes-section-glueing-schemes}.
En appliquant Schémas, lemme \ref{schemes-lemma-glue-schemes}
nous obtenons un schéma $X \to S$ en
recollant les schémas $W_j$ au moyen des identifications $\theta_{j, j'}$.
Il est clair, par construction, que $X \to S$ est \'etale et que $X_0 \cong Y$.

\medskip\noindent
Il suffit donc de démontrer le lemme lorsque $S$ et $Y$ sont affines.
Écrivons $S = \Spec(R)$ et $S_0 = \Spec(R/I)$, où $I$ est localement nilpotent.
D'après Algèbre, lemme \ref{algebra-lemma-etale-standard-smooth}, nous savons que
$Y$ est le spectre d'un anneau $\overline{A}$ avec
$$
\overline{A} = (R/I)[x_1, \ldots, x_n]/(\overline{f}_1, \ldots, \overline{f}_n)
$$
tel que
$$
\overline{g} =
\det
\left(
\begin{matrix}
\partial \overline{f}_1/\partial x_1 &
\partial \overline{f}_2/\partial x_1 &
\ldots &
\partial \overline{f}_n/\partial x_1 \\
\partial \overline{f}_1/\partial x_2 &
\partial \overline{f}_2/\partial x_2 &
\ldots &
\partial \overline{f}_n/\partial x_2 \\
\ldots & \ldots & \ldots & \ldots \\
\partial \overline{f}_1/\partial x_n &
\partial \overline{f}_2/\partial x_n &
\ldots &
\partial \overline{f}_n/\partial x_n
\end{matrix}
\right)
$$
a pour image un élément inversible dans $\overline{A}$. Choisissons des relèvements quelconques
$f_i \in R[x_1, \ldots, x_n]$. Posons
$$
A = R[x_1, \ldots, x_n]/(f_1, \ldots, f_n)
$$
Puisque $I$ est localement nilpotent, l'idéal $IA$ est localement nilpotent
(Algèbre, lemme \ref{algebra-lemma-locally-nilpotent}).
Remarquons que $\overline{A} = A/IA$.
Il s'ensuit que le déterminant de la matrice des dérivées partielles des
$f_i$ est inversible dans l'algèbre $A$, d'après
Algèbre, lemme \ref{algebra-lemma-locally-nilpotent-unit}.
Par conséquent, $R \to A$ est \'etale, ce qui achève la démonstration.
\end{proof}



\section{La caractérisation fonctorielle}
\label{etale-section-functorial-etale}

\noindent
Nous présentons enfin la caractérisation fonctorielle annoncée.
Il y a donc quatre façons d'envisager les morphismes \'etales de schémas :
\begin{enumerate}
\item comme des morphismes lisses de dimension relative $0$,
\item comme des morphismes localement de présentation finie, plats et non ramifiés,
\item au moyen du théorème de structure, et
\item au moyen de la caractérisation fonctorielle.
\end{enumerate}

\begin{theorem}
\label{etale-theorem-formally-etale}
Soit $f : X \to S$ un morphisme localement de présentation finie.
Les assertions suivantes sont équivalentes
\begin{enumerate}
\item $f$ est \'etale,
\item pour tout $S$-schéma affine $Y$ et tout sous-schéma fermé $Y_0 \subset Y$
défini par un idéal de carré nul, l'application naturelle
$$
\Mor_S(Y, X) \longrightarrow \Mor_S(Y_0, X)
$$
est bijective.
\end{enumerate}
\end{theorem}

\begin{proof}
C'est
Compléments sur les morphismes, lemme \ref{more-morphisms-lemma-etale-formally-etale}.
\end{proof}

\noindent
Cette caractérisation affirme que les solutions des équations définissant $X$ se
relèvent de façon unique à travers les épaississements nilpotents.



\section{Structure locale \'etale des morphismes non ramifiés}
\label{etale-section-unramified-etale-local}

\noindent
Au chapitre
Compléments sur les morphismes, section \ref{more-morphisms-section-etale-localization},
le lecteur trouvera quelques résultats sur la structure locale \'etale des
morphismes quasi-finis. Dans cette section, nous voulons les combiner
avec les propriétés topologiques des morphismes non ramifiés rencontrées
dans ce chapitre. La configuration générale à garder à l'esprit est
$$
\xymatrix{
V \ar[r] \ar[dr] & X_U \ar[d] \ar[r] & X \ar[d]^f \\
& U \ar[r] & S
}
$$
voir
Compléments sur les morphismes, équation (\ref{more-morphisms-equation-basic-diagram}).
Commençons par un cas très général.

\begin{lemma}
\label{etale-lemma-unramified-etale-local}
Soit $f : X \to S$ un morphisme de schémas.
Soient $x_1, \ldots, x_n \in X$ des points ayant même image $s$ dans $S$.
Supposons que $f$ soit non ramifié en chacun des $x_i$.
Il existe alors un voisinage \'etale $(U, u) \to (S, s)$
et des ouverts $V_{i, j} \subset X_U$, $i = 1, \ldots, n$, $j = 1, \ldots, m_i$,
tels que
\begin{enumerate}
\item $V_{i, j} \to U$ est une immersion fermée dont l'image contient $u$,
\item $u$ n'appartient pas à l'image de $V_{i, j} \cap V_{i', j'}$, sauf si
$i = i'$ et $j = j'$, et
\item tout point de $(X_U)_u$ s'envoyant sur $x_i$ appartient à l'un des $V_{i, j}$.
\end{enumerate}
\end{lemma}

\begin{proof}
D'après
Morphismes, définition \ref{morphisms-definition-unramified},
il existe, pour chaque $x_i$, un voisinage ouvert qui est localement de type
fini sur $S$. En remplaçant $X$ par un voisinage ouvert de $\{x_1, \ldots, x_n\}$,
nous pouvons supposer que $f$ est localement de type fini. Appliquons
Compléments sur les morphismes, lemme
\ref{more-morphisms-lemma-etale-makes-quasi-finite-finite-multiple-points-var}
pour obtenir le voisinage \'etale $(U, u)$ et les ouverts $V_{i, j}$ finis sur
$U$. D'après
le lemme \ref{etale-lemma-finite-unramified-one-point},
quitte à rétrécir $U$, le morphisme $V_{i, j} \to U$ est une immersion
fermée.
\end{proof}

\begin{lemma}
\label{etale-lemma-unramified-etale-local-technical}
Soit $f : X \to S$ un morphisme de schémas.
Soient $x_1, \ldots, x_n \in X$ des points ayant même image $s$ dans $S$.
Supposons que $f$ soit séparé et que $f$ soit non ramifié en chacun des $x_i$.
Il existe alors un voisinage \'etale $(U, u) \to (S, s)$
et une décomposition en somme disjointe
$$
X_U =
W \amalg \coprod\nolimits_{i, j} V_{i, j}
$$
telle que
\begin{enumerate}
\item $V_{i, j} \to U$ est une immersion fermée dont l'image contient $u$,
\item la fibre $W_u$ ne contient aucun point s'envoyant sur l'un des $x_i$.
\end{enumerate}
En particulier, si $f^{-1}(\{s\}) = \{x_1, \ldots, x_n\}$, alors
la fibre $W_u$ est vide.
\end{lemma}

\begin{proof}
Appliquons
le lemme \ref{etale-lemma-unramified-etale-local}.
Nous pouvons supposer que $U$ est affine, de sorte que $X_U$ est séparé.
Alors $V_{i, j} \to X_U$ est une application fermée ; voir
Morphismes, lemme \ref{morphisms-lemma-image-proper-scheme-closed}.
Supposons que $(i, j) \not = (i', j')$.
Alors $V_{i, j} \cap V_{i', j'}$ est fermé dans $V_{i, j}$ et
son image dans $U$ ne contient pas $u$.
Par conséquent, quitte à rétrécir $U$, nous pouvons supposer que
$V_{i, j} \cap V_{i', j'} = \emptyset$. De plus, $\bigcup V_{i, j}$ est
un sous-schéma ouvert et fermé de $X_U$ et admet donc un
complémentaire ouvert et fermé $W$. Cela achève la démonstration.
\end{proof}

\noindent
Le lemme suivant est en un certain sens beaucoup plus faible que le précédent,
mais il peut être utile de l'énoncer explicitement ici. Il affirme qu'un morphisme fini
non ramifié est, localement pour la topologie \'etale sur la base, une immersion fermée.

\begin{lemma}
\label{etale-lemma-finite-unramified-etale-local}
Soit $f : X \to S$ un morphisme fini non ramifié de schémas.
Soit $s \in S$.
Il existe un voisinage \'etale $(U, u) \to (S, s)$
et une décomposition en somme disjointe finie
$$
X_U = \coprod\nolimits_j V_j
$$
telle que chaque $V_j \to U$ soit une immersion fermée.
\end{lemma}

\begin{proof}
Puisque $X \to S$ est fini, la fibre au-dessus de $s$ est un ensemble fini
$\{x_1, \ldots, x_n\}$ de points de $X$. Appliquons le
lemme \ref{etale-lemma-unramified-etale-local-technical}
à cet ensemble (un morphisme fini est séparé ; voir
Morphismes, section \ref{morphisms-section-integral}).
L'image de $W$ dans $U$ est un sous-ensemble fermé
(car $X_U \to U$ est fini, donc propre) qui ne
contient pas $u$. Après l'avoir retiré de $U$, nous voyons que $W = \emptyset$
ce qui est le résultat voulu.
\end{proof}




\section{Structure locale \'etale des morphismes \'etales}
\label{etale-section-etale-local-etale}

\noindent
C'est un peu artificiel, mais peut-être utile pour se forger une intuition des morphismes \'etales.
Nous recopions simplement les résultats de la
section \ref{etale-section-unramified-etale-local}
et remplaçons « immersion fermée » par « isomorphisme ».

\begin{lemma}
\label{etale-lemma-etale-etale-local}
Soit $f : X \to S$ un morphisme de schémas.
Soient $x_1, \ldots, x_n \in X$ des points ayant la même image $s$ dans $S$.
Supposons que $f$ soit \'etale en chacun des $x_i$.
Alors il existe un voisinage \'etale $(U, u) \to (S, s)$
et des ouverts $V_{i, j} \subset X_U$, $i = 1, \ldots, n$, $j = 1, \ldots, m_i$
tels que
\begin{enumerate}
\item $V_{i, j} \to U$ soit un isomorphisme ;
\item $u$ n'appartienne pas à l'image de $V_{i, j} \cap V_{i', j'}$, sauf si
$i = i'$ et $j = j'$ ; et
\item tout point de $(X_U)_u$ qui s'envoie sur $x_i$ appartienne à un certain $V_{i, j}$.
\end{enumerate}
\end{lemma}

\begin{proof}
Un morphisme \'etale est non ramifié ; nous pouvons donc appliquer le
lemme \ref{etale-lemma-unramified-etale-local}.
Or $V_{i, j} \to U$ est une immersion fermée et est \'etale.
C'est donc une immersion ouverte, par exemple d'après le
théorème \ref{etale-theorem-etale-radicial-open}.
Remplaçons $U$ par l'intersection des images des morphismes $V_{i, j} \to U$
pour obtenir le lemme.
\end{proof}

\begin{lemma}
\label{etale-lemma-etale-etale-local-technical}
Soit $f : X \to S$ un morphisme de schémas.
Soient $x_1, \ldots, x_n \in X$ des points ayant la même image $s$ dans $S$.
Supposons $f$ séparé et $f$ \'etale en chacun des $x_i$.
Alors il existe un voisinage \'etale $(U, u) \to (S, s)$
et une décomposition en somme disjointe finie
$$
X_U =
W \amalg \coprod\nolimits_{i, j} V_{i, j}
$$
de schémas tels que
\begin{enumerate}
\item $V_{i, j} \to U$ soit un isomorphisme ;
\item la fibre $W_u$ ne contienne aucun point s'envoyant sur l'un des $x_i$.
\end{enumerate}
En particulier, si $f^{-1}(\{s\}) = \{x_1, \ldots, x_n\}$, alors
la fibre $W_u$ est vide.
\end{lemma}

\begin{proof}
Un morphisme \'etale est non ramifié ; nous pouvons donc appliquer le
lemme \ref{etale-lemma-unramified-etale-local-technical}.
Comme dans la démonstration du
lemme \ref{etale-lemma-etale-etale-local},
les morphismes $V_{i, j} \to U$ sont des immersions ouvertes et
l'on conclut en remplaçant $U$ par l'intersection de leurs
images.
\end{proof}

\noindent
Le lemme suivant est en un certain sens beaucoup plus faible que le précédent,
mais il peut être utile de l'énoncer explicitement ici. Il affirme qu'un morphisme fini
\'etale est, localement pour la topologie \'etale sur la base, un
« espace de revêtement topologique », c'est-à-dire un produit fini de copies de la base.

\begin{lemma}
\label{etale-lemma-finite-etale-etale-local}
Soit $f : X \to S$ un morphisme fini \'etale de schémas.
Soit $s \in S$. Il existe un voisinage \'etale $(U, u) \to (S, s)$
et une décomposition en somme disjointe finie
$$
X_U = \coprod\nolimits_j V_j
$$
de schémas telle que chaque $V_j \to U$ soit un isomorphisme.
\end{lemma}

\begin{proof}
Un morphisme \'etale est non ramifié ; nous pouvons donc appliquer le
lemme \ref{etale-lemma-finite-unramified-etale-local}.
Comme dans la démonstration du
lemme \ref{etale-lemma-etale-etale-local},
nous voyons que $V_{i, j} \to U$ est une immersion ouverte et l'on conclut
en remplaçant $U$ par l'intersection de leurs images.
\end{proof}




\section{Propriétés de permanence}
\label{etale-section-properties-permanence}

\noindent
Dans ce qui suit, nous présentons quelques propriétés de « permanence »
des homomorphismes \'etales d'anneaux locaux noethériens
(au sens de la définition \ref{etale-definition-etale-ring}). Voir
Compléments d'algèbre, sections \ref{more-algebra-section-permanence-completion} et
\ref{more-algebra-section-permanence-henselization}
pour les résultats analogues concernant la complétion et
l'hensélisation d'un anneau local noethérien.

\begin{lemma}
\label{etale-lemma-etale-dimension}
Soient $A$ et $B$ des anneaux locaux noethériens.
Soit $A \to B$ un homomorphisme local \'etale.
Alors $\dim(A) = \dim(B)$.
\end{lemma}

\begin{proof}
Voir par exemple
Algèbre, lemme \ref{algebra-lemma-dimension-base-fibre-equals-total}.
\end{proof}

\begin{proposition}
\label{etale-proposition-etale-depth}
Soient $A$ et $B$ des anneaux locaux noethériens.
Soit $f : A \to B$ un homomorphisme local \'etale.
Alors $\text{profondeur}(A) = \text{profondeur}(B)$
\end{proposition}

\begin{proof}
Voir Algèbre, lemme \ref{algebra-lemma-apply-grothendieck}.
\end{proof}

\begin{proposition}
\label{etale-proposition-etale-CM}
\begin{slogan}
La propriété d'être de Cohen-Macaulay monte et descend le long des homomorphismes \'etales.
\end{slogan}
Soient $A$ et $B$ des anneaux locaux noethériens.
Soit $f : A \to B$ un homomorphisme local \'etale.
Alors $A$ est de Cohen-Macaulay si et seulement si $B$ l'est.
\end{proposition}

\begin{proof}
Un anneau local $A$ est de Cohen-Macaulay si et seulement si $\dim(A) = \text{profondeur}(A)$.
Puisque ces deux invariants sont préservés par une extension \'etale,
l'assertion en résulte.
\end{proof}

\begin{proposition}
\label{etale-proposition-etale-regular}
Soient $A$, $B$ des anneaux locaux noethériens.
Soit $f : A \to B$ un homomorphisme local \'etale.
Alors $A$ est régulier si et seulement si $B$ l'est.
\end{proposition}

\begin{proof}
Si $B$ est régulier, alors $A$ est régulier d'après
Algèbre, lemme \ref{algebra-lemma-flat-under-regular}.
Supposons $A$ régulier. Soit $\mathfrak m$ l'idéal maximal
de $A$. Alors $\dim_{\kappa(\mathfrak m)} \mathfrak m/\mathfrak m^2 =
\dim(A) = \dim(B)$ (voir le lemme \ref{etale-lemma-etale-dimension}).
D'autre part, $\mathfrak mB$ est l'idéal maximal de
$B$ et par conséquent $\mathfrak m_B/\mathfrak m_B = \mathfrak mB/\mathfrak m^2B$
est engendré par au plus $\dim(B)$ éléments. Ainsi $B$ est régulier.
(On peut aussi utiliser le résultat légèrement plus général
Algèbre, lemme \ref{algebra-lemma-flat-over-regular-with-regular-fibre}.)
\end{proof}


\begin{proposition}
\label{etale-proposition-etale-reduced}
Soient $A$, $B$ des anneaux locaux noethériens.
Soit $f : A \to B$ un homomorphisme local \'etale.
Alors $A$ est réduit si et seulement si $B$ l'est.
\end{proposition}

\begin{proof}
La fidèle platitude de $A \to B$ montre clairement que, si $B$ est réduit,
$A$ l'est aussi. Voir également Algèbre, lemme \ref{algebra-lemma-descent-reduced}.
Réciproquement, supposons $A$ réduit. Par hypothèse, $B$ est un localisé
d'une $A$-algèbre de type fini $B'$ en un idéal premier $\mathfrak q$.
Quitte à remplacer $B'$ par un localisé, nous pouvons supposer que $B'$
est \'etale sur $A$, voir le lemme \ref{etale-lemma-characterize-etale-Noetherian}.
Alors Algèbre, lemme \ref{algebra-lemma-reduced-goes-up}, s'applique à
$A \to B'$ et $B'$ est réduit. Par conséquent, $B$ est réduit.
\end{proof}

\begin{remark}
\label{etale-remark-technicality-needed}
Le résultat sur la « propriété d'être réduit » ne vaut pas avec une définition plus
faible des homomorphismes locaux \'etales $A \to B$, dans laquelle on
supprime l'hypothèse que $B$ est essentiellement de type fini sur $A$.
En effet, il peut arriver qu'un anneau local noethérien intègre $A$ ait un
complété non réduit $A^\wedge$, voir
Exemples, section \ref{examples-section-local-completion-nonreduced}.
Mais l'homomorphisme d'anneaux $A \to A^\wedge$ est plat, et $\mathfrak m_AA^\wedge$
est l'idéal maximal de $A^\wedge$ et, bien sûr, $A$ et $A^\wedge$ ont
le même corps résiduel. C'est pourquoi il importe de ne considérer
cette notion que pour des extensions d'anneaux essentiellement de type fini
(ou essentiellement de présentation finie si $A$ n'est pas noethérien).
\end{remark}

\begin{proposition}
\label{etale-proposition-etale-normal}
\begin{reference}
\cite[Expose I, Theorem 9.5 part (i)]{SGA1}
\end{reference}
Soient $A$, $B$ des anneaux locaux noethériens.
Soit $f : A \to B$ un homomorphisme local \'etale.
Alors $A$ est un anneau intègre normal si et seulement si $B$ l'est.
\end{proposition}

\begin{proof}
Voir
Algèbre, lemme \ref{algebra-lemma-descent-normal},
pour la descente de la normalité. Réciproquement, supposons $A$ normal.
Par hypothèse, $B$ est un localisé d'une $A$-algèbre de type fini
$B'$ en un idéal premier $\mathfrak q$. Après avoir remplacé $B'$ par un localisé,
nous pouvons supposer que $B'$ est \'etale sur $A$, voir
le lemme \ref{etale-lemma-characterize-etale-Noetherian}.
Alors
Algèbre, lemme \ref{algebra-lemma-normal-goes-up},
s'applique à $A \to B'$ et nous concluons que $B'$ est normal.
Par conséquent, $B$ est un anneau intègre normal.
\end{proof}

\noindent
Les propositions précédentes donnent une idée de la raison pour laquelle nous voulons considérer
les morphismes \'etales comme des « isomorphismes locaux ». Une autre propriété qui indique
très clairement que nous avons la « bonne » définition est le fait que,
pour les $\mathbf{C}$-schémas de type fini, un morphisme est \'etale si et seulement si
le morphisme associé entre espaces analytiques (les points à valeurs dans $\mathbf{C}$ munis
de la topologie complexe) est un isomorphisme local au sens analytique (une immersion
ouverte localement sur la source). Ce fait peut être démontré à l'aide du
théorème de structure et du fait que l'analytification commute à la
formation des anneaux locaux complétés -- les détails sont laissés au lecteur.








\section{Descente des morphismes \'etales}
\label{etale-section-descending-etale}

\noindent
Pour comprendre le langage employé dans cette section, nous invitons
le lecteur à consulter
Descente, section \ref{descent-section-descent-datum}.
Soit $f : X \to S$ un morphisme de schémas. Considérons le
foncteur de changement de base
\begin{equation}
\label{etale-equation-descent-etale}
\text{schémas }U\text{ \'etales sur }S \longrightarrow
\begin{matrix}
\text{données de descente }(V, \varphi)\text{ relativement à }X/S \\
\text{ avec }V\text{ \'etale sur }X
\end{matrix}
\end{equation}
qui envoie $U$ sur la donnée de descente canonique $(X \times_S U, can)$.

\begin{lemma}
\label{etale-lemma-faithful}
Si $f : X \to S$ est surjectif, alors le foncteur
(\ref{etale-equation-descent-etale}) est fidèle.
\end{lemma}

\begin{proof}
Soient $a, b : U_1 \to U_2$ deux morphismes entre des schémas \'etales sur $S$.
Supposons que les changements de base de $a$ et $b$ par $X$ coïncident.
Nous devons montrer que $a = b$.
D'après la proposition \ref{etale-proposition-equality}, il suffit de
montrer que $a$ et $b$ coïncident sur les points et les corps résiduels.
C'est clair, car, pour tout $u \in U_1$, nous pouvons trouver un point
$v \in X \times_S U_1$ s'envoyant sur $u$.
\end{proof}

\begin{lemma}
\label{etale-lemma-fully-faithful}
Supposons que $f : X \to S$ soit submersif et que tout changement de base \'etale
de $f$ soit submersif. Alors le foncteur
(\ref{etale-equation-descent-etale}) est pleinement fidèle.
\end{lemma}

\begin{proof}
D'après le lemme \ref{etale-lemma-faithful}, le foncteur est fidèle.
Soient $U_1 \to S$ et $U_2 \to S$ des morphismes \'etales
et soit $a : X \times_S U_1 \to X \times_S U_2$ un
morphisme compatible avec les données de descente canoniques.
Nous allons montrer que $a$ est le changement de base d'un morphisme $U_1 \to U_2$.

\medskip\noindent
Soit $U'_2 \subset U_2$ un sous-schéma ouvert. Considérons
$W = a^{-1}(X \times_S U'_2)$. C'est un sous-schéma ouvert
de $X \times_S U_1$ compatible avec la donnée de descente
canonique sur $V_1 = X \times_S U_1$. Cela signifie que les
deux images inverses de $W$ par les projections
$V_1 \times_{U_1} V_1 \to V_1$ coïncident. Puisque $V_1 \to U_1$
est surjectif (étant le changement de base de $X \to S$), nous en déduisons
que $W$ est l'image inverse d'une partie $U'_1 \subset U_1$.
Puisque $W$ est ouvert, notre hypothèse sur $f$ implique que $U'_1 \subset U_1$
est ouvert.

\medskip\noindent
Soit $U_2 = \bigcup U_{2, i}$ un recouvrement par des ouverts affines.
D'après le résultat du paragraphe précédent, nous obtenons un recouvrement ouvert
$U_1 = \bigcup U_{1, i}$ tel que
$X \times_S U_{1, i} = a^{-1}(X \times_S U_{2, i})$.
Si nous pouvons montrer qu'il existe un morphisme $U_{1, i} \to U_{2, i}$
dont le changement de base est le morphisme
$a_i : X \times_S U_{1, i} \to X \times_S U_{2, i}$,
alors nous pouvons recoller ces morphismes en un morphisme $U_1 \to U_2$
(en utilisant la fidélité). Nous nous ramenons ainsi au cas où
$U_2$ est affine. En particulier, $U_2 \to S$ est séparé
(Schémas, lemme \ref{schemes-lemma-compose-after-separated}).

\medskip\noindent
Supposons que $U_2 \to S$ soit séparé. Alors le graphe $\Gamma_a$ de $a$
est un sous-schéma fermé de
$$
V = (X \times_S U_1) \times_X (X \times_S U_2) = X \times_S U_1 \times_S U_2
$$
d'après Schémas, lemme \ref{schemes-lemma-semi-diagonal}.
D'autre part, le graphe est ouvert, par exemple
parce que c'est une section d'un morphisme \'etale
(proposition \ref{etale-proposition-properties-sections}).
Comme $a$ est un morphisme de données de descente, les deux images inverses de
$\Gamma_a \subset V$ par les projections
$V \times_{U_1 \times_S U_2} V \to V$ sont les mêmes.
Ainsi, en raisonnant comme au deuxième alinéa de la démonstration, nous
trouvons un sous-schéma ouvert et fermé $\Gamma \subset U_1 \times_S U_2$
dont le changement de base à $X$ donne $\Gamma_a$. Alors
$\Gamma \to U_1$ est un morphisme \'etale dont le changement de base
à $X$ est un isomorphisme. Cela signifie que $\Gamma \to U_1$
est universellement bijectif, donc est un isomorphisme
d'après le théorème \ref{etale-theorem-etale-radicial-open}.
Ainsi, $\Gamma$ est le graphe d'un morphisme $U_1 \to U_2$
et le changement de base de ce morphisme est $a$, ce qui est le résultat voulu.
\end{proof}

\begin{lemma}
\label{etale-lemma-fully-faithful-cases}
Soit $f : X \to S$ un morphisme de schémas. Dans les cas suivants,
le foncteur (\ref{etale-equation-descent-etale}) est pleinement fidèle :
\begin{enumerate}
\item $f$ est surjectif et universellement fermé
(par exemple fini, entier ou propre),
\item $f$ est surjectif et universellement ouvert
(par exemple localement de présentation finie et plat, lisse ou étale),
\item $f$ est surjectif, quasi-compact et plat.
\end{enumerate}
\end{lemma}

\begin{proof}
Cela résulte du lemme \ref{etale-lemma-fully-faithful}.
Par exemple, une application fermée surjective d'espaces topologiques
est submersive (Topologie, lemme
\ref{topology-lemma-closed-morphism-quotient-topology}).
Les morphismes finis, entiers et propres sont universellement fermés, voir
Morphismes, lemmes \ref{morphisms-lemma-integral-universally-closed} et
\ref{morphisms-lemma-finite-proper}, et
la définition \ref{morphisms-definition-proper}.
D'autre part, une application ouverte surjective d'espaces topologiques
est submersive (Topologie, lemme
\ref{topology-lemma-open-morphism-quotient-topology}).
Les morphismes plats localement de présentation finie, lisses et \'etales sont
universellement ouverts, voir
Morphismes, lemmes \ref{morphisms-lemma-fppf-open},
\ref{morphisms-lemma-smooth-open} et
\ref{morphisms-lemma-etale-open}.
Le cas des morphismes surjectifs, quasi-compacts et plats résulte
de Morphismes, lemme \ref{morphisms-lemma-fpqc-quotient-topology}.
\end{proof}

\begin{lemma}
\label{etale-lemma-reduce-to-affine}
Soit $f : X \to S$ un morphisme de schémas.
Soit $(V, \varphi)$ une donnée de descente relativement à $X/S$
avec $V \to X$ \'etale. Soit $S = \bigcup S_i$ un
recouvrement ouvert. Supposons que
\begin{enumerate}
\item la donnée de descente $(V, \varphi)$ obtenue par changement de base
à $X \times_S S_i/S_i$ soit effective,
\item le foncteur (\ref{etale-equation-descent-etale})
pour $X \times_S (S_i \cap S_j) \to (S_i \cap S_j)$ soit pleinement fidèle, et
\item le foncteur (\ref{etale-equation-descent-etale})
pour $X \times_S (S_i \cap S_j \cap S_k) \to (S_i \cap S_j \cap S_k)$
soit fidèle.
\end{enumerate}
Alors $(V, \varphi)$ est effective.
\end{lemma}

\begin{proof}
(Rappelons que les changements de base des données de descente sont définis dans
Descente, définition \ref{descent-definition-pullback-functor}.)
Posons $X_i = X \times_S S_i$. Notons $(V_i, \varphi_i)$ la donnée obtenue
par changement de base à partir de $(V, \varphi)$ sur $X_i/S_i$.
D'après l'hypothèse (1), nous pouvons trouver un morphisme \'etale $U_i \to S_i$
muni d'un isomorphisme $X_i \times_{S_i} U_i \to V_i$ compatible avec
$can$ et $\varphi_i$. D'après l'hypothèse (2), nous obtenons des isomorphismes
$\psi_{ij} : U_i \times_{S_i} (S_i \cap S_j) \to
U_j \times_{S_j} (S_i \cap S_j)$.
D'après l'hypothèse (3), ces isomorphismes satisfont à la condition de cocycle,
de sorte que $(U_i, \psi_{ij})$ est une donnée de descente pour le
recouvrement de Zariski $\{S_i \to S\}$. Alors Descente, lemme
\ref{descent-lemma-Zariski-refinement-coverings-equivalence}
(qui n'est essentiellement qu'une reformulation de
Schémas, section \ref{schemes-section-glueing-schemes})
affirme qu'il existe un morphisme de schémas $U \to S$
et des isomorphismes $U \times_S S_i \to U_i$ compatibles
avec $\psi_{ij}$. Les isomorphismes $U \times_S S_i \to U_i$
déterminent des isomorphismes correspondants $X_i \times_S U \to V_i$
qui se recollent en un morphisme $X \times_S U \to V$ compatible
avec la donnée de descente canonique et $\varphi$.
\end{proof}

\begin{lemma}
\label{etale-lemma-split-henselian}
Soit $(A, I)$ un couple hensélien. Soit $U \to \Spec(A)$ un morphisme
quasi-compact, séparé et \'etale tel que
$U \times_{\Spec(A)} \Spec(A/I) \to \Spec(A/I)$ soit fini.
Alors
$$
U = U_{fin} \amalg U_{away}
$$
où $U_{fin} \to \Spec(A)$ est fini et $U_{away}$ ne possède
aucun point au-dessus de $Z$.
\end{lemma}

\begin{proof}
D'après le théorème principal de Zariski, le schéma $U$ est quasi-affine.
En fait, nous pouvons trouver une immersion ouverte $U \to T$, avec $T$ affine et
$T \to \Spec(A)$ fini, voir Compléments sur les morphismes, lemme
\ref{more-morphisms-lemma-quasi-finite-separated-pass-through-finite}.
Écrivons $Z = \Spec(A/I)$ et notons $U_Z \to T_Z$ le changement de base.
Comme $U_Z \to Z$ est fini, nous voyons que $U_Z \to T_Z$ est fermé
aussi bien qu'ouvert. Par conséquent, d'après
Compléments sur l'algèbre, lemme \ref{more-algebra-lemma-characterize-henselian-pair},
nous obtenons une unique décomposition $T = T' \amalg T''$ telle que $T'_Z = U_Z$.
Posons $U_{fin} = U \cap T'$ et $U_{away} = U \cap T''$. Puisque
$T'_Z \subset U_Z$, nous voyons que tous les points fermés de $T'$ sont dans $U$,
donc que $T' \subset U$, puis que $U_{fin} = T'$, et donc que $U_{fin} \to \Spec(A)$
est fini. Nous omettons la démonstration
de l'unicité de la décomposition.
\end{proof}

\begin{proposition}
\label{etale-proposition-effective}
Soit $f : X \to S$ un morphisme entier surjectif.
Le foncteur (\ref{etale-equation-descent-etale}) induit une équivalence
$$
\begin{matrix}
\text{schémas quasi-compacts,}\\
\text{séparés et \'etales sur }S
\end{matrix}
\longrightarrow
\begin{matrix}
\text{données de descente }(V, \varphi)\text{ relativement à }X/S\text{ avec}\\
V\text{ quasi-compact, séparé et \'etale sur }X
\end{matrix}
$$
\end{proposition}

\begin{proof}
D'après le lemme \ref{etale-lemma-fully-faithful-cases}, le
foncteur (\ref{etale-equation-descent-etale})
est pleinement fidèle, et cela reste vrai après tout
changement de base $S \to S'$. Soit $(V, \varphi)$ une donnée de descente
relativement à $X/S$, avec $V \to X$ quasi-compact, séparé et \'etale.
Nous pouvons utiliser le lemme \ref{etale-lemma-reduce-to-affine}
pour voir qu'il suffit de prouver l'effectivité
localement pour la topologie de Zariski sur $S$. En particulier, nous pouvons et allons
supposer que $S$ est affine.

\medskip\noindent
Si $S$ est affine, nous pouvons trouver un ensemble ordonné filtrant $\Lambda$ et
un système projectif $X_\lambda \to S_\lambda$
de morphismes finis de schémas affines de type fini sur
$\Spec(\mathbf{Z})$ tel que $(X \to S) = \lim (X_\lambda \to S_\lambda)$.
Voir Algèbre, lemme \ref{algebra-lemma-limit-integral}.
Puisque les limites commutent entre elles, nous en déduisons que
$X \times_S X = \lim X_\lambda \times_{S_\lambda} X_\lambda$
et que
$X \times_S X \times_S X = \lim
X_\lambda \times_{S_\lambda} X_\lambda \times_{S_\lambda} X_\lambda$.
Remarquons que $V \to X$ est un morphisme de présentation finie.
En utilisant Limites, lemme \ref{limits-lemma-descend-finite-presentation},
nous pouvons trouver un $\lambda$ et une donnée de descente $(V_\lambda, \varphi_\lambda)$
relativement à $X_\lambda/S_\lambda$ dont le changement de base à $X/S$ est
$(V, \varphi)$. Il suffit bien sûr de montrer que
$(V_\lambda, \varphi_\lambda)$ est effective. Notons que $V_\lambda$
est quasi-compact par construction.
Quitte à augmenter $\lambda$, nous pouvons supposer
que $V_\lambda \to X_\lambda$ est séparé et \'etale, voir
Limites, lemme \ref{limits-lemma-descend-separated-finite-presentation} et
\ref{limits-lemma-descend-etale}.
Ainsi, nous pouvons supposer que $f$ est fini surjectif et que
$S$ est affine de type fini sur $\mathbf{Z}$.

\medskip\noindent
Considérons un ouvert $S' \subset S$ tel que la donnée $(V', \varphi')$ obtenue
par changement de base à partir de $(V, \varphi)$ sur $X' = X \times_S S'$ soit effective. Nous prouverons ci-dessous
que $S' \not = S$ implique l'existence d'un ouvert strictement plus grand sur
lequel la donnée de descente est effective. Comme $S$ est noethérien (et que son
espace topologique sous-jacent est donc noethérien), cela achèvera la démonstration.
Soit $\xi \in S$ un point générique d'une composante irréductible du
fermé $Z = S \setminus S'$.
Si $\xi \in S'' \subset S$ est un ouvert sur lequel la donnée de descente est
effective, alors celle-ci est effective sur
$S' \cup S''$ par l'argument de recollement du premier alinéa. Ainsi,
dans la suite de la démonstration, nous pouvons remplacer $S$ par un voisinage ouvert affine
de $\xi$.

\medskip\noindent
Après un premier remplacement de ce type, nous pouvons supposer que $Z$ est irréductible
de point générique $Z$. Munissons $Z$ de la structure réduite induite
de sous-schéma fermé. Après une nouvelle restriction, nous pouvons supposer que
$X_Z = X \times_S Z = f^{-1}(Z) \to Z$ est plat, voir
Morphismes, proposition \ref{morphisms-proposition-generic-flatness}.
Soit $(V_Z, \varphi_Z)$ la donnée de descente obtenue par changement de base à $X_Z/Z$.
D'après Compléments sur les morphismes, lemme
\ref{more-morphisms-lemma-separated-locally-quasi-finite-morphisms-fppf-descend},
cette donnée de descente est effective, et nous obtenons un morphisme \'etale
$U_Z \to Z$ dont le changement de base est isomorphe à $V_Z$ d'une manière
compatible avec les données de descente.
Bien sûr, $U_Z \to Z$ est quasi-compact et séparé
(Descente, lemmes \ref{descent-lemma-descending-property-quasi-compact} et
\ref{descent-lemma-descending-property-separated}).
Ainsi, après une nouvelle restriction, nous pouvons supposer
que $U_Z \to Z$ est fini, voir
Morphismes, lemme \ref{morphisms-lemma-generically-finite}.

\medskip\noindent
Soit $S = \Spec(A)$ et soit $I \subset A$ l'idéal premier correspondant
à $Z \subset S$. Soit $(A^h, IA^h)$ l'hensélisation du couple
$(A, I)$. Notons $S^h = \Spec(A^h)$ et $Z^h = V(IA^h) \cong Z$.
Nous affirmons qu'il suffit de montrer l'effectivité après changement de base à
$S^h$. En effet, $\{S^h \to S, S' \to S\}$ est un recouvrement fpqc
($A \to A^h$ est plat d'après Compléments sur l'algèbre, lemme
\ref{more-algebra-lemma-henselization-flat}) et,
d'après Compléments sur les morphismes, lemme
\ref{more-morphisms-lemma-separated-locally-quasi-finite-morphisms-fppf-descend},
nous disposons de la descente fpqc pour les morphismes \'etales séparés.
En effet, si $U^h \to S^h$ et $U' \to S'$ sont les objets
correspondant aux données $(V^h, \varphi^h)$ et
$(V', \varphi')$, alors les isomorphismes requis
$$
U^h \times_S S^h \to S^h \times_S V^h
\quad\text{et}\quad
U^h \times_S S' \to S^h \times_S U'
$$
sont obtenus grâce à la pleine fidélité signalée au premier
alinéa. Nous nous ramenons ainsi à la situation décrite dans
l'alinéa suivant.

\medskip\noindent
Ici, $S = \Spec(A)$, $Z = V(I)$, $S' = S \setminus Z$, où
$(A, I)$ est un couple hensélien ; nous avons $U' \to S'$ correspondant
à la donnée de descente $(V', \varphi')$ et un morphisme \'etale fini
$U_Z \to Z$ correspondant à la donnée de descente
$(V_Z, \varphi_Z)$. Nous ne disposons plus de l'hypothèse que $A$ est de type fini
sur $\mathbf{Z}$, mais le reste de l'argument n'utilisera même pas
que $A$ est noethérien.
D'après Compléments sur l'algèbre, lemme \ref{more-algebra-lemma-finite-etale-equivalence},
nous pouvons trouver un morphisme \'etale fini $U_{fin} \to S$ dont la
restriction à $Z$ est isomorphe à $U_Z \to Z$.
Écrivons $X = \Spec(B)$ et $Y = V(IB)$. Comme $(B, IB)$ est un couple hensélien
(Compléments sur l'algèbre, lemme \ref{more-algebra-lemma-integral-over-henselian-pair})
et comme la restriction de $V \to X$ à $Y$
est finie (car elle est le changement de base de $U_Z \to Z$), nous voyons
qu'il existe une décomposition canonique en union disjointe
$$
V = V_{fin} \amalg V_{away}
$$
où $V_{fin} \to X$ est fini et où $V_{away}$ ne possède aucun
point au-dessus de $Y$. Voir le lemme \ref{etale-lemma-split-henselian}.
En utilisant l'unicité de cette décomposition sur $X \times_S X$,
nous voyons que $\varphi$ la respecte, et nous obtenons
$$
(V, \varphi) = (V_{fin}, \varphi_{fin}) \amalg (V_{away}, \varphi_{away})
$$
dans la catégorie des données de descente.
D'après Compléments sur l'algèbre, lemme \ref{more-algebra-lemma-finite-etale-equivalence},
il existe un unique isomorphisme
$$
X \times_S U_{fin} \longrightarrow V_{fin}
$$
compatible avec l'isomorphisme donné $Y \times_Z U_Z \to V \times_X Y$
sur $Y$.
Par unicité, nous voyons que cet isomorphisme est compatible
avec les données de descente, c'est-à-dire que
$(X \times_S U_{fin}, can) \cong (V_{fin}, \varphi_{fin})$.
Notons $U'_{fin} = U_{fin} \times_S S'$. Par pleine fidélité,
nous obtenons un morphisme $U'_{fin} \to U'$ qui est
l'inclusion d'un sous-schéma ouvert (et fermé).
Posons alors $U = U_{fin} \amalg_{U'_{fin}} U'$ (recollement de schémas comme
dans Schémas, section \ref{schemes-section-glueing-schemes}).
Les morphismes $X \times_S U_{fin} \to V$ et
$X \times_S U' \to V$ se recollent en un morphisme $X \times_S U \to V$
qui est l'isomorphisme recherché.
\end{proof}





\section{Diviseurs à croisements normaux}
\label{etale-section-normal-crossings}

\noindent
Voici la définition.

\begin{definition}
\label{etale-definition-strict-normal-crossings}
Soit $X$ un schéma localement noethérien. Un
{\it diviseur strictement à croisements normaux}
sur $X$ est un diviseur de Cartier effectif $D \subset X$ tel que,
pour tout $p \in D$, l'anneau local $\mathcal{O}_{X, p}$ soit régulier
et qu'il existe un système régulier de paramètres
$x_1, \ldots, x_d \in \mathfrak m_p$ et $1 \leq r \leq d$
tel que $D$ soit défini par $x_1 \ldots x_r$ dans $\mathcal{O}_{X, p}$.
\end{definition}

\noindent
Nous rencontrons souvent des diviseurs de Cartier effectifs $E$ sur des schémas
localement noethériens $X$ tels qu'il existe un diviseur strictement à croisements normaux $D$
avec $E \subset D$ ensemblistement.
Dans ce cas, nous avons
$E = \sum a_i D_i$, avec $a_i \geq 0$, où $D = \bigcup_{i \in I} D_i$
est la décomposition de $D$ en ses composantes irréductibles.
Remarquons que $D' = \bigcup_{a_i > 0} D_i$ est un diviseur à croisements
normaux stricts tel que $E = D'$ ensemblistement.
Lorsque ce qui précède se produit, nous dirons que
$E$ est {\it à support contenu dans un diviseur strictement à croisements normaux}.

\begin{lemma}
\label{etale-lemma-strict-normal-crossings}
Soit $X$ un schéma localement noethérien. Soit $D \subset X$ un
diviseur de Cartier effectif. Soient $D_i \subset D$, $i \in I$, ses
composantes irréductibles considérées comme des sous-schémas fermés réduits de $X$.
Les conditions suivantes sont équivalentes :
\begin{enumerate}
\item $D$ est un diviseur strictement à croisements normaux, et
\item $D$ est réduit, chaque $D_i$ est un diviseur de Cartier effectif, et
pour $J \subset I$ fini, l'intersection schématique
$D_J = \bigcap_{j \in J} D_j$ est un schéma
régulier dont chaque composante irréductible est de
codimension $|J|$ dans $X$.
\end{enumerate}
\end{lemma}

\begin{proof}
Supposons que $D$ soit un diviseur strictement à croisements normaux. Choisissons $p \in D$
et un système régulier de paramètres $x_1, \ldots, x_d \in \mathfrak m_p$
ainsi que $1 \leq r \leq d$ comme dans la
définition \ref{etale-definition-strict-normal-crossings}.
Comme $\mathcal{O}_{X, p}/(x_i)$ est un anneau local régulier
(et en particulier un anneau intègre), nous voyons que les composantes irréductibles
$D_1, \ldots, D_r$ de $D$ passant par $p$ sont en correspondance bijective ($1$ à $1$)
aux idéaux premiers de hauteur un $(x_1), \ldots, (x_r)$ de $\mathcal{O}_{X, p}$.
D'après Algèbre, lemme \ref{algebra-lemma-regular-ring-CM},
nous trouvons que l'intersection $D_{i_1} \cap \ldots \cap D_{i_s}$
est de codimension $s$ dans un voisinage ouvert de $p$
et que son anneau local en $p$ est régulier.
Comme cela vaut pour tout $p \in D$, nous concluons que (2) est satisfaite.

\medskip\noindent
Supposons (2). Soit $p \in D$. Comme $\mathcal{O}_{X, p}$ est de dimension
finie, nous voyons que $p$ peut appartenir à au plus
$\dim(\mathcal{O}_{X, p})$ des composantes $D_i$.
Disons que $p \in D_1, \ldots, D_r$ pour un certain $r \geq 1$.
Soient $x_1, \ldots, x_r \in \mathfrak m_p$ des équations locales
de $D_1, \ldots, D_r$. Alors $x_1$ est un non diviseur de zéro dans $\mathcal{O}_{X, p}$
et $\mathcal{O}_{X, p}/(x_1) = \mathcal{O}_{D_1, p}$ est régulier.
Ainsi, $\mathcal{O}_{X, p}$ est régulier, voir
Algèbre, lemme \ref{algebra-lemma-regular-mod-x}.
Comme $D_1 \cap \ldots \cap D_r$ est un schéma régulier (donc normal),
c'est une union disjointe de ses composantes irréductibles
(Propriétés, lemme \ref{properties-lemma-normal-Noetherian}).
Soit $Z \subset D_1 \cap \ldots \cap D_r$
la composante irréductible contenant $p$.
Alors $\mathcal{O}_{Z, p} = \mathcal{O}_{X, p}/(x_1, \ldots, x_r)$
est régulier de codimension $r$ (notons que, puisque nous savons déjà
que $\mathcal{O}_{X, p}$ est régulier et donc de Cohen-Macaulay,
il n'y a aucune ambiguïté sur la codimension, car l'anneau est caténaire, voir
Algèbre, lemmes \ref{algebra-lemma-regular-ring-CM} et
\ref{algebra-lemma-CM-dim-formula}).
Ainsi, $\dim(\mathcal{O}_{Z, p}) = \dim(\mathcal{O}_{X, p}) - r$.
Choisissons des éléments supplémentaires $x_{r + 1}, \ldots, x_n \in \mathfrak m_p$
dont les images forment un système minimal de générateurs de $\mathfrak m_{Z, p}$.
Alors $\mathfrak m_p = (x_1, \ldots, x_n)$ d'après le lemme de Nakayama,
et nous voyons que $D$ est un diviseur à croisements normaux.
\end{proof}

\begin{lemma}
\label{etale-lemma-smooth-pullback-strict-normal-crossings}
\begin{slogan}
L'image inverse d'un diviseur strictement à croisements normaux par un morphisme
lisse est un diviseur strictement à croisements normaux.
\end{slogan}
Soit $X$ un schéma localement noethérien. Soit $D \subset X$ un
diviseur strictement à croisements normaux. Si $f : Y \to X$ est un morphisme
lisse de schémas, alors l'image inverse $f^*D$ est un
diviseur strictement à croisements normaux sur $Y$.
\end{lemma}

\begin{proof}
Comme $f$ est plat, l'image inverse est définie par
Diviseurs, lemme \ref{divisors-lemma-pullback-effective-Cartier-defined},
donc l'énoncé a un sens.
Soit $q \in f^*D$ d'image $p \in D$. Choisissons un système régulier
de paramètres $x_1, \ldots, x_d \in \mathfrak m_p$
et $1 \leq r \leq d$ comme dans la
définition \ref{etale-definition-strict-normal-crossings}.
Comme $f$ est lisse, l'homomorphisme d'anneaux locaux
$\mathcal{O}_{X, p} \to \mathcal{O}_{Y, q}$ est plat
et l'anneau de la fibre
$$
\mathcal{O}_{Y, q}/\mathfrak m_p \mathcal{O}_{Y, q} =
\mathcal{O}_{Y_p, q}
$$
est un anneau local régulier (voir par exemple
Algèbre, lemme \ref{algebra-lemma-characterize-smooth-over-field}).
Choisissons $y_1, \ldots, y_n \in \mathfrak m_q$ dont les images forment un système
régulier de paramètres dans $\mathcal{O}_{Y_p, q}$.
Alors $x_1, \ldots, x_d, y_1, \ldots, y_n$ engendrent l'idéal
maximal $\mathfrak m_q$. Ainsi, $\mathcal{O}_{Y, q}$
est un anneau local régulier de dimension
$d + n$ d'après Algèbre, lemme \ref{algebra-lemma-dimension-base-fibre-equals-total},
et $x_1, \ldots, x_d, y_1, \ldots, y_n$
est un système régulier de paramètres. Comme $f^*D$ est défini
par $x_1 \ldots x_r$ dans $\mathcal{O}_{Y, q}$, nous concluons
que le lemme est démontré.
\end{proof}

\noindent
Voici la définition d'un diviseur à croisements normaux.

\begin{definition}
\label{etale-definition-normal-crossings}
Soit $X$ un schéma localement noethérien. Un {\it diviseur à croisements normaux}
sur $X$ est un diviseur de Cartier effectif $D \subset X$ tel que, pour
tout $p \in D$, il existe un morphisme \'etale $U \to X$ dont l'image
contient $p$ et tel que $D \times_X U$ soit un
diviseur strictement à croisements normaux sur $U$.
\end{definition}

\noindent
Par exemple, $D = V(x^2 + y^2)$ est un diviseur à croisements normaux
(mais non strict) sur
$\Spec(\mathbf{R}[x, y])$, car, après changement de base au
recouvrement \'etale $\Spec(\mathbf{C}[x, y])$, nous obtenons $(x - iy)(x + iy) = 0$.

\begin{lemma}
\label{etale-lemma-smooth-pullback-normal-crossings}
\begin{slogan}
L'image inverse d'un diviseur à croisements normaux par un morphisme
lisse est un diviseur à croisements normaux.
\end{slogan}
Soit $X$ un schéma localement noethérien. Soit $D \subset X$ un
diviseur à croisements normaux. Si $f : Y \to X$ est un morphisme
lisse de schémas, alors l'image inverse $f^*D$ est un
diviseur à croisements normaux sur $Y$.
\end{lemma}

\begin{proof}
Comme $f$ est plat, l'image inverse est définie par
Diviseurs, lemme \ref{divisors-lemma-pullback-effective-Cartier-defined},
donc l'énoncé a un sens.
Soit $q \in f^*D$ d'image $p \in D$.
Choisissons un morphisme \'etale $U \to X$ dont l'image contient $p$
tel que $D \times_X U \subset U$ soit un diviseur à croisements normaux
stricts, comme dans la définition \ref{etale-definition-normal-crossings}.
Posons $V = Y \times_X U$. Alors $V \to Y$ est \'etale comme changement
de base de $U \to X$
(Morphismes, lemme \ref{morphisms-lemma-base-change-etale}),
et l'image inverse $D \times_X V$ est un diviseur strictement à croisements
normaux sur $V$ d'après le lemme \ref{etale-lemma-smooth-pullback-strict-normal-crossings}.
Nous avons donc vérifié la condition de la
définition \ref{etale-definition-normal-crossings}
pour $q \in f^*D$, ce qui permet de conclure.
\end{proof}

\begin{lemma}
\label{etale-lemma-characterize-normal-crossings-normalization}
Soit $X$ un schéma localement noethérien. Soit $D \subset X$ un sous-schéma
fermé. Les conditions suivantes sont équivalentes :
\begin{enumerate}
\item $D$ est un diviseur à croisements normaux dans $X$,
\item $D$ est réduit, la normalisation $\nu : D^\nu \to D$ est non ramifiée,
et, pour tout $n \geq 1$, le schéma
$$
Z_n = D^\nu \times_D \ldots \times_D D^\nu
\setminus \{(p_1, \ldots, p_n) \mid p_i = p_j\text{ pour certains }i\not = j\}
$$
est régulier ; le morphisme $Z_n \to X$ est localement d'intersection complète
et son faisceau conormal est localement libre de rang $n$.
\end{enumerate}
\end{lemma}

\begin{proof}
Expliquons d'abord comment interpréter la condition (2).
La diagonale d'un morphisme non ramifié est ouverte
(Morphismes, lemme \ref{morphisms-lemma-diagonal-unramified-morphism}).
D'autre part, $D^\nu \to D$ est séparé, donc la
diagonale $D^\nu \to D^\nu \times_D D^\nu$ est fermée.
Ainsi, $Z_n$ est un sous-schéma ouvert et fermé de
$D^\nu \times_D \ldots \times_D D^\nu$. D'autre part,
$Z_n \to X$ est non ramifié, car c'est la composée
$$
Z_n \to D^\nu \times_D \ldots \times_D D^\nu \to \ldots \to
D^\nu \times_D D^\nu \to D^\nu \to D \to X
$$
et chacune des flèches est non ramifiée.
Comme un morphisme non ramifié est formellement non ramifié
(Compléments sur les morphismes, lemme
\ref{more-morphisms-lemma-unramified-formally-unramified}),
nous avons un faisceau conormal
$\mathcal{C}_n = \mathcal{C}_{Z_n/X}$ du morphisme $Z_n \to X$, voir
Compléments sur les morphismes, définition
\ref{more-morphisms-definition-universal-thickening}.

\medskip\noindent
La formation de la normalisation commute à la localisation \'etale d'après
Compléments sur les morphismes, lemme \ref{more-morphisms-lemma-normalization-and-smooth}.
La régularité des anneaux locaux, le fait
qu'un morphisme soit non ramifié ou localement
d'intersection complète, ou le fait qu'un morphisme soit
non ramifié et possède un faisceau conormal qui soit
localement libre d'un rang donné, se vérifient localement pour la topologie \'etale (voir
Compléments sur l'algèbre, lemme \ref{more-algebra-lemma-regular-etale-extension},
Descente, lemme \ref{descent-lemma-descending-property-unramified},
Compléments sur les morphismes, lemme \ref{more-morphisms-lemma-descending-property-lci}
et
Descente, lemme \ref{descent-lemma-finite-locally-free-descends}).

\medskip\noindent
D'après la remarque de l'alinéa précédent et la définition
des diviseurs à croisements normaux, il suffit de prouver qu'un
diviseur strictement à croisements normaux $D = \bigcup_{i \in I} D_i$
satisfait à (2). Dans ce cas, $D^\nu = \coprod D_i$
et $D^\nu \to D$ est non ramifié (la propriété d'être non ramifié
est locale sur la source, et $D_i \to D$ est une immersion
fermée qui est non ramifiée). De même, $Z_1 = D^\nu \to X$
est un morphisme localement d'intersection complète, car nous pouvons
le vérifier localement sur la source et chaque morphisme $D_i \to X$
est une immersion régulière, puisque c'est l'inclusion d'un diviseur de Cartier
(voir le lemme \ref{etale-lemma-strict-normal-crossings} et
Compléments sur les morphismes, lemme \ref{more-morphisms-lemma-regular-immersion-lci}).
Comme un diviseur de Cartier effectif possède un faisceau
conormal inversible, nous concluons que la condition portant sur le
faisceau conormal est satisfaite.
De même, le schéma $Z_n$, pour $n \geq 2$, est l'union disjointe
des schémas $D_J = \bigcap_{j \in J} D_j$, où $J \subset I$
parcourt les sous-ensembles de cardinal $n$. Comme $D_J \to X$ est
une immersion régulière de codimension $n$
(par la définition des diviseurs strictement à croisements normaux et le
fait que nous pouvons le vérifier sur les germes d'après
Diviseurs, lemme \ref{divisors-lemma-Noetherian-scheme-regular-ideal}),
il s'ensuit de la même manière que $Z_n \to X$ possède les propriétés
requises. Nous omettons quelques détails.

\medskip\noindent
Supposons (2). Soit $p \in D$. Comme $D^\nu \to D$ est non ramifié, il est
fini (d'après Morphismes, lemme \ref{morphisms-lemma-finite-integral}).
Ainsi, $D^\nu \to X$ est fini non ramifié.
D'après le lemme \ref{etale-lemma-finite-unramified-etale-local}
et par localisation \'etale (permise par la discussion
du deuxième alinéa et par la définition des diviseurs à
croisements normaux), nous nous ramenons au cas où
$D^\nu = \coprod_{i \in I} D_i$,
avec $I$ fini et $D_i \to U$ une immersion fermée.
Quitte à restreindre $X$, nous pouvons supposer
que $p \in D_i$ pour tout $i \in I$. La condition selon laquelle $Z_1 =  D^\nu \to X$ est un
morphisme non ramifié localement d'intersection complète
dont le faisceau conormal est localement libre de rang $1$
implique que $D_i \subset X$ est un diviseur de Cartier effectif, voir
Compléments sur les morphismes, lemme \ref{more-morphisms-lemma-lci} et
Diviseurs, lemme \ref{divisors-lemma-regular-immersion-noetherian}.
Pour achever la démonstration, nous pouvons supposer que $X = \Spec(A)$ est affine
et que $D_i = V(f_i)$, avec $f_i \in A$ un non diviseur de zéro.
Si $I = \{1, \ldots, r\}$, alors $p \in Z_r = V(f_1, \ldots, f_r)$.
La même référence que ci-dessus implique que
$(f_1, \ldots, f_r)$ est un idéal Koszul-régulier dans $A$.
Comme le faisceau conormal est de rang $r$, nous voyons que
$f_1, \ldots, f_r$ est un système minimal de générateurs de
l'idéal définissant $Z_r$ dans $\mathcal{O}_{X, p}$.
Cela implique que $f_1, \ldots, f_r$ est une suite régulière
dans $\mathcal{O}_{X, p}$ telle que $\mathcal{O}_{X, p}/(f_1, \ldots, f_r)$
soit régulier. Nous concluons ainsi du
lemme \ref{algebra-lemma-regular-mod-x} d'Algèbre
que la suite $f_1, \ldots, f_r$ peut être complétée en un système régulier de paramètres
dans $\mathcal{O}_{X, p}$, ce qui achève la démonstration.
\end{proof}

\begin{lemma}
\label{etale-lemma-characterize-normal-crossings}
Soit $X$ un schéma localement noethérien. Soit $D \subset X$ un sous-schéma
fermé. Si $X$ est J-2 ou de Nagata, les conditions suivantes sont équivalentes :
\begin{enumerate}
\item $D$ est un diviseur à croisements normaux dans $X$,
\item pour tout $p \in D$, l'image inverse de $D$ sur le spectre de
l'hensélisé strict $\mathcal{O}_{X, p}^{sh}$
est un diviseur strictement à croisements normaux.
\end{enumerate}
\end{lemma}

\begin{proof}
L'implication (1) $\Rightarrow$ (2) est immédiate et
ne nécessite pas l'hypothèse que $X$ est J-2 ou de Nagata.
En effet, soit $p \in D$ et choisissons un voisinage \'etale
$(U, u) \to (X, p)$ tel que l'image inverse de $D$ soit
un diviseur strictement à croisements normaux sur $U$.
Alors $\mathcal{O}_{X, p}^{sh} = \mathcal{O}_{U, u}^{sh}$,
et nous voyons que la trace de $D$ sur $\Spec(\mathcal{O}_{U, u}^{sh})$
est définie par une partie d'un système régulier de paramètres,
puisque c'est déjà le cas dans $\mathcal{O}_{U, u}$.

\medskip\noindent
Pour prouver l'implication dans l'autre sens,
nous utiliserons le critère du
lemme \ref{etale-lemma-characterize-normal-crossings-normalization}.
Remarquons que la formation de la normalisation $D^\nu \to D$
commute à l'hensélisation stricte, voir
Compléments sur les morphismes, lemme
\ref{more-morphisms-lemma-normalization-and-henselization}.
Si nous pouvons montrer que $D^\nu \to D$ est fini,
alors nous voyons que $D^\nu \to D$ et les schémas
$Z_n$ satisfont à toutes les propriétés voulues, car celles-ci
peuvent toutes se vérifier au niveau des anneaux locaux
(mais la finitude du morphisme $D^\nu \to D$
ne peut pas se vérifier sur les anneaux locaux).
Nous omettons les vérifications détaillées.

\medskip\noindent
Si $X$ est de Nagata, alors $D^\nu \to D$ est fini d'après
Morphismes, lemme \ref{morphisms-lemma-nagata-normalization}.

\medskip\noindent
Supposons que $X$ soit J-2. Choisissons un point $p \in D$. Nous montrerons
que $D^\nu \to D$ est fini au-dessus d'un voisinage de $p$.
Par hypothèse, il existe un système régulier de
paramètres $f_1, \ldots, f_d$ de $\mathcal{O}_{X, p}^{sh}$
et $1 \leq r \leq d$ tels que la trace de $D$ sur
$\Spec(\mathcal{O}_{X, p}^{sh})$ soit définie par $f_1 \ldots f_r$.
Alors
$$
D^\nu \times_X \Spec(\mathcal{O}_{X, p}^{sh}) = 
\coprod\nolimits_{i = 1, \ldots, r} V(f_i)
$$
Choisissons un voisinage \'etale affine
$(U, u) \to (X, p)$ tel que $f_i$ provienne d'un élément
$f_i \in \mathcal{O}_U(U)$. Posons $D_i = V(f_i) \subset U$.
L'hensélisé strict de $\mathcal{O}_{D_i, u}$
est $\mathcal{O}_{X, p}^{sh}/(f_i)$, qui est régulier.
Ainsi, $\mathcal{O}_{D_i, u}$ est régulier (par exemple d'après
Compléments sur l'algèbre, lemme \ref{more-algebra-lemma-henselization-regular}).
Comme $X$ est J-2, le lieu régulier est ouvert dans $D_i$.
Ainsi, quitte à remplacer $U$ par un ouvert de Zariski, nous pouvons supposer
que $D_i$ est régulier pour tout $i$. Il s'ensuit que
$$
\coprod\nolimits_{i = 1, \ldots, r} D_i = D^\nu \times_X U
\longrightarrow D \times_X U
$$
est le morphisme de normalisation, et il est clairement fini.
Autrement dit, nous avons trouvé
un voisinage \'etale $(U, u)$ de $(X, p)$ tel que
le changement de base de $D^\nu \to D$ à ce voisinage soit fini.
Cela implique que $D^\nu \to D$ est fini par descente
(Descente, lemme \ref{descent-lemma-descending-property-finite}),
et la démonstration est achevée.
\end{proof}











% OK, start here.
%

\chapter{Homologie de Chow et classes de Chern}


\phantomsection
\label{chow-section-phantom}




\section{Introduction}
\label{chow-section-introduction}

\noindent
Dans ce chapitre, nous étudions les groupes d'homologie de Chow et la construction
des classes de Chern des fibrés vectoriels comme éléments des groupes
de cohomologie de Chow opérationnelle (le tout à coefficients dans $\mathbf{Z}$).

\medskip\noindent
Nous commençons ce chapitre en donnant la démonstration algébrique la plus courte
possible du lemme clé \ref{chow-lemma-milnor-gersten-low-degree}.
Nous définissons d'abord le quotient de Herbrand
(section \ref{chow-section-periodic-complexes})
et le calculons dans certains cas
(section \ref{chow-section-calculation}).
Ensuite, nous démontrons quelques lemmes algébriques simples sur
l'existence de factorisations convenables après des modifications
(section \ref{chow-section-preparation-tame-symbol}).
À l'aide de ceux-ci, nous construisons et définissons le symbole modéré à la
section \ref{chow-section-tame-symbol}.
Seules les propriétés les plus élémentaires du symbole modéré
sont nécessaires pour démontrer le lemme clé, ce que nous faisons
à la section \ref{chow-section-key-lemma}.

\medskip\noindent
Ensuite, nous introduisons le cadre de base dans lequel nous travaillons dans la suite de ce
chapitre à la section \ref{chow-section-setup}. Afin de rendre la matière
un peu plus stimulante, nous avons décidé de traiter un cas un peu plus général
que celui habituellement considéré. Plus précisément, nous supposons que nos schémas $X$ sont localement
de type fini sur un schéma de base fixé, localement noethérien, qui est universellement
caténaire et muni d'une fonction de dimension. Ces hypothèses suffisent
pour pouvoir définir les groupes d'homologie de Chow $\CH_*(X)$ et l'action sur
ceux-ci du produit cap avec les classes de Chern. Cela indique que nous devrions
également pouvoir les définir pour les champs algébriques localement de type fini
sur une telle base.

\medskip\noindent
Ensuite, nous suivons les premiers chapitres de \cite{F} afin de définir
les cycles, l'image réciproque plate, l'image directe propre et l'équivalence rationnelle,
si ce n'est que nous avons été moins précis quant aux supports des cycles
considérés.

\medskip\noindent
Nous nous écartons de la présentation donnée dans \cite{F} en utilisant le
lemme clé mentionné ci-dessus pour démontrer une relation de commutativité fondamentale à la
section \ref{chow-section-key}. Nous en déduisons que l'opération
d'intersection avec un faisceau inversible passe au quotient par l'équivalence
rationnelle et est commutative, voir la section \ref{chow-section-commutativity}.
Une nouvelle application du lemme
clé démontre que l'application de Gysin d'un diviseur de Cartier effectif
passe au quotient par l'équivalence rationnelle, voir la section \ref{chow-section-gysin}.
Cela étant démontré, il est immédiat de définir les classes de Chern
des fibrés vectoriels, d'établir l'additivité, de démontrer le principe de scindage,
d'introduire les caractères de Chern et les classes de Todd, puis d'énoncer le
théorème de Grothendieck--Riemann--Roch.

\medskip\noindent
Il y a deux appendices. Dans l'appendice A (section \ref{chow-section-appendix-A}),
nous étudions une autre construction, plus longue, du
symbole modéré et la démonstration correspondante du lemme clé.
Enfin, dans l'appendice B (section \ref{chow-section-appendix-chow}),
nous examinons brièvement le rapport avec la $K$-théorie des faisceaux
cohérents et étudions quelques lemmes sur les éclatements.
Nous suggérons au lecteur de consulter leurs introductions pour
de plus amples informations.

\medskip\noindent
Nous reviendrons, au chapitre suivant, sur les groupes de Chow $\CH_*(X)$ des variétés projectives lisses
sur des corps algébriquement clos. À l'aide d'un lemme de déplacement
comme dans \cite{Samuel}, \cite{ChevalleyI} et \cite{ChevalleyII},
ainsi que de la formule de Serre pour les Tor
(voir \cite{Serre_local_algebra} ou \cite{Serre_algebre_locale}),
nous définirons une structure d'anneau sur $\CH_*(X)$. Voir
Théorie de l'intersection, section \ref{intersection-section-introduction} et suivantes.








\section{Complexes périodiques et quotients de Herbrand}
\label{chow-section-periodic-complexes}

\noindent
Il existe bien entendu une notion très générale de complexe périodique.
On peut imposer la périodicité des applications ou celle des objets.
Nous ajouterons ici ces conditions selon les besoins. Pour le moment, seuls les
cas suivants nous sont nécessaires.

\begin{definition}
\label{chow-definition-periodic-complex}
Soit $R$ un anneau.
\begin{enumerate}
\item Un {\it complexe $2$-périodique} sur $R$ est donné
par un quadruplet $(M, N, \varphi, \psi)$ constitué de
$R$-modules $M$, $N$ et d'applications de $R$-modules $\varphi : M \to N$,
$\psi : N \to M$ telles que
$$
\xymatrix{
\ldots \ar[r] &
M \ar[r]^\varphi &
N \ar[r]^\psi &
M \ar[r]^\varphi &
N \ar[r] & \ldots
}
$$
soit un complexe. Dans ce cadre, nous définissons les {\it modules de cohomologie}
du complexe comme étant les $R$-modules
$$
H^0(M, N, \varphi, \psi) = \Ker(\varphi)/\Im(\psi)
\quad\text{et}\quad
H^1(M, N, \varphi, \psi) = \Ker(\psi)/\Im(\varphi).
$$
Nous disons que le complexe $2$-périodique est {\it exact} si ses modules de cohomologie
sont nuls.
\item Un {\it complexe $(2, 1)$-périodique} sur $R$ est donné
par un triplet $(M, \varphi, \psi)$ constitué d'un $R$-module $M$ et
d'applications de $R$-modules $\varphi : M \to M$, $\psi : M \to M$
telles que
$$
\xymatrix{
\ldots \ar[r] &
M \ar[r]^\varphi &
M \ar[r]^\psi &
M \ar[r]^\varphi &
M \ar[r] & \ldots
}
$$
soit un complexe. Comme il s'agit d'un cas particulier d'un complexe $2$-périodique,
nous disposons de ses {\it modules de cohomologie} $H^0(M, \varphi, \psi)$,
$H^1(M, \varphi, \psi)$ et d'une notion d'exactitude.
\end{enumerate}
\end{definition}

\noindent
Dans la suite, nous utiliserons tout résultat démontré pour les complexes $2$-périodiques
sans autre mention dans le cas des complexes $(2, 1)$-périodiques.
Il est clair que la collection des complexes $2$-périodiques forme une
catégorie, dont les morphismes
$(f, g) : (M, N, \varphi, \psi) \to (M', N', \varphi', \psi')$
sont les couples d'applications $f : M \to M'$ et $g : N \to N'$ tels
que $\varphi' \circ f = g \circ \varphi$ et $\psi' \circ g = f \circ \psi$.
On obtient une catégorie abélienne, avec noyaux et conoyaux comme dans
Homologie, Lemme \ref{homology-lemma-cat-chain-abelian}.

\begin{definition}
\label{chow-definition-periodic-length}
Soit $(M, N, \varphi, \psi)$ un complexe $2$-périodique
sur un anneau $R$ dont les modules de cohomologie sont de longueur finie.
Dans ce cas, nous définissons la {\it multiplicité} de $(M, N, \varphi, \psi)$
comme étant l'entier
$$
e_R(M, N, \varphi, \psi) =
\text{longueur}_R(H^0(M, N, \varphi, \psi))
-
\text{longueur}_R(H^1(M, N, \varphi, \psi))
$$
Dans le cas d'un complexe $(2, 1)$-périodique $(M, \varphi, \psi)$,
nous la notons $e_R(M, \varphi, \psi)$ et l'appellerons parfois
le {\it quotient de Herbrand (additif)}.
\end{definition}

\noindent
Si les groupes de cohomologie de $(M, \varphi, \psi)$
sont des groupes abéliens finis, il est d'usage d'appeler
{\it quotient de Herbrand (multiplicatif)}
$$
q(M, \varphi, \psi) =
\frac{\# H^0(M, \varphi, \psi)}{\# H^1(M, \varphi, \psi)}
$$
Autrement dit, le quotient de Herbrand multiplicatif est le nombre d'éléments de
$H^0$ divisé par le nombre d'éléments de $H^1$. Si $R$ est local et si
le corps résiduel de $R$ est fini à $q$ éléments, alors
$$
q(M, \varphi, \psi) = q^{e_R(M, \varphi, \psi)}
$$

\medskip\noindent
Un exemple de complexe $(2, 1)$-périodique sur un anneau $R$ est tout triplet
de la forme $(M, 0, \psi)$, où $M$ est un $R$-module et $\psi$ une application
$R$-linéaire. Si le noyau et le conoyau de $\psi$ sont de longueur finie,
alors on obtient
\begin{equation}
\label{chow-equation-multiplicity-coker-ker}
e_R(M, 0, \psi) = \text{longueur}_R(\Coker(\psi)) - \text{longueur}_R(\Ker(\psi))
\end{equation}
Énonçons et démontrons les lemmes usuels relatifs à ces notations.

\begin{lemma}
\label{chow-lemma-additivity-periodic-length}
Soit $R$ un anneau. Supposons donnée une suite exacte courte de
complexes $2$-périodiques
$$
0 \to (M_1, N_1, \varphi_1, \psi_1)
\to (M_2, N_2, \varphi_2, \psi_2)
\to (M_3, N_3, \varphi_3, \psi_3)
\to 0
$$
Si deux d'entre eux ont des modules de cohomologie de longueur finie, il en va de même
du troisième, et l'on a
$$
e_R(M_2, N_2, \varphi_2, \psi_2) =
e_R(M_1, N_1, \varphi_1, \psi_1) +
e_R(M_3, N_3, \varphi_3, \psi_3).
$$
\end{lemma}

\begin{proof}
Abrégeons $A = (M_1, N_1, \varphi_1, \psi_1)$,
$B = (M_2, N_2, \varphi_2, \psi_2)$ et $C = (M_3, N_3, \varphi_3, \psi_3)$.
Nous avons une suite exacte longue de cohomologie
$$
\ldots
\to H^1(C)
\to H^0(A)
\to H^0(B)
\to H^0(C)
\to H^1(A)
\to H^1(B)
\to H^1(C)
\to \ldots
$$
Elle donne une suite exacte finie
$$
0 \to I
\to H^0(A)
\to H^0(B)
\to H^0(C)
\to H^1(A)
\to H^1(B)
\to K \to 0
$$
avec $0 \to K \to H^1(C) \to I \to 0$ comme filtration. Par additivité de
la fonction longueur (Algèbre, Lemme \ref{algebra-lemma-length-additive}),
on obtient le résultat.
\end{proof}

\begin{lemma}
\label{chow-lemma-finite-periodic-length}
Soit $R$ un anneau. Si $(M, N, \varphi, \psi)$ est un complexe $2$-périodique
tel que $M$ et $N$ soient de longueur finie, alors
$e_R(M, N, \varphi, \psi) = \text{longueur}_R(M) - \text{longueur}_R(N)$.
En particulier, si $(M, \varphi, \psi)$ est un complexe $(2, 1)$-périodique
tel que $M$ soit de longueur finie, alors
$e_R(M, \varphi, \psi) = 0$.
\end{lemma}

\begin{proof}
Observons que, dans la catégorie des complexes $2$-périodiques où $M$ et $N$
sont de longueur finie, la quantité « $\text{longueur}_R(M) - \text{longueur}_R(N)$ »
est additive pour les suites exactes courtes (l'énoncé précis est laissé au lecteur).
Considérons la suite exacte courte
$$
0 \to (M, \Im(\varphi), \varphi, 0) \to
(M, N, \varphi, \psi) \to (0, N/\Im(\varphi), 0, 0) \to 0
$$
La remarque initiale, combinée à l'additivité du
Lemme \ref{chow-lemma-additivity-periodic-length},
nous ramène aux cas (a) $M = 0$ et (b) $\varphi$ est surjective.
Nous les laissons au lecteur.
\end{proof}

\begin{lemma}
\label{chow-lemma-compare-periodic-lengths}
Soit $R$ un anneau. Soit $f : (M, \varphi, \psi) \to (M', \varphi', \psi')$
une application de complexes $(2, 1)$-périodiques dont les modules de cohomologie
sont de longueur finie. Si $\Ker(f)$ et $\Coker(f)$ sont de longueur finie,
alors $e_R(M, \varphi, \psi) = e_R(M', \varphi', \psi')$.
\end{lemma}

\begin{proof}
Appliquons l'additivité du Lemme \ref{chow-lemma-additivity-periodic-length}
et observons que $(\Ker(f), \varphi, \psi)$ et
$(\Coker(f), \varphi', \psi')$ ont une multiplicité nulle d'après le
Lemme \ref{chow-lemma-finite-periodic-length}.
\end{proof}




\section{Calcul de quelques multiplicités}
\label{chow-section-calculation}

\noindent
Pour démontrer plus loin l'égalité de certains cycles, nous devons
calculer quelques multiplicités. Notre principal outil, outre les
lemmes élémentaires sur les multiplicités donnés à la section précédente,
sera Algèbre, Lemme \ref{algebra-lemma-order-vanishing-determinant}.

\begin{lemma}
\label{chow-lemma-length-multiplication}
Soit $R$ un anneau local noethérien.
Soit $M$ un $R$-module fini. Soit $x \in R$. Supposons que
\begin{enumerate}
\item $\dim(\text{Supp}(M)) \leq 1$, et
\item $\dim(\text{Supp}(M/xM)) \leq 0$.
\end{enumerate}
Écrivons
$\text{Supp}(M) = \{\mathfrak m, \mathfrak q_1, \ldots, \mathfrak q_t\}$.
Alors
$$
e_R(M, 0, x) =
\sum\nolimits_{i = 1, \ldots, t}
\text{ord}_{R/\mathfrak q_i}(x)
\text{longueur}_{R_{\mathfrak q_i}}(M_{\mathfrak q_i}).
$$
\end{lemma}

\begin{proof}
Commençons par quelques remarques préparatoires.
Le résultat du lemme vaut si $M$ est de longueur finie, c'est-à-dire si $t = 0$,
car le membre de gauche et celui de droite sont alors tous deux nuls,
voir le Lemme \ref{chow-lemma-finite-periodic-length}.
De même, si l'on a une suite exacte courte $0 \to M \to M' \to M'' \to 0$
de modules satisfaisant (1) et (2), le lemme pour deux de ces trois
modules implique le lemme pour le troisième, par
additivité des longueurs (Algèbre, Lemme \ref{algebra-lemma-length-additive}) et
additivité des multiplicités (Lemme \ref{chow-lemma-additivity-periodic-length}).

\medskip\noindent
Notons $M_i$ l'image de $M$ dans $M_{\mathfrak q_i}$, de sorte que
$\text{Supp}(M_i) = \{\mathfrak m, \mathfrak q_i\}$.
Le noyau et le conoyau de l'application $M \to \bigoplus M_i$
ont pour support $\{\mathfrak m\}$ et sont donc de longueur finie.
Il résulte de nos remarques préparatoires qu'il suffit de
démontrer le lemme pour chacun des $M_i$. Nous pouvons donc supposer que
$\text{Supp}(M) = \{\mathfrak m, \mathfrak q\}$.
Dans ce cas, nous avons une filtration finie
$M \supset \mathfrak qM \supset \mathfrak q^2M \supset \ldots \supset
\mathfrak q^nM = 0$ d'après Algèbre, Lemme
\ref{algebra-lemma-Noetherian-power-ideal-kills-module}.
Une nouvelle application de l'additivité montre qu'il suffit de démontrer le lemme
lorsque $M$ est annulé par $\mathfrak q$.
Nous pouvons alors considérer $M$ comme un $R/\mathfrak q$-module,
c'est-à-dire supposer que $R$ est un anneau local noethérien intègre
de dimension $1$, de corps des fractions $K$.
En quotientant par le sous-module de torsion, c'est-à-dire par le
noyau de $M \to M \otimes_R K = V$ (la torsion étant de
longueur finie, nos remarques préliminaires s'y appliquent),
nous pouvons supposer que $M \subset V$ est un réseau
(Algèbre, Définition \ref{algebra-definition-lattice}).
Alors $x : M \to M$ est injective et
$\text{longueur}_R(M/xM) = d(M, xM)$
(Algèbre, Définition \ref{algebra-definition-distance}). Comme
$\text{longueur}_K(V) = \dim_K(V)$,
on voit que $\det(x : V \to V) = x^{\dim_K(V)}$ et
$\text{ord}_R(\det(x : V \to V)) = \dim_K(V) \text{ord}_R(x)$.
L'égalité cherchée résulte donc de
Algèbre, Lemme \ref{algebra-lemma-order-vanishing-determinant}
dans ce cas.
\end{proof}

\begin{lemma}
\label{chow-lemma-additivity-divisors-restricted}
Soit $R$ un anneau local noethérien.
Soit $x \in R$. Si $M$ est un module de Cohen--Macaulay fini sur $R$
avec $\dim(\text{Supp}(M)) = 1$ et $\dim(\text{Supp}(M/xM)) = 0$, alors
$$
\text{longueur}_R(M/xM)
=
\sum\nolimits_i \text{longueur}_R(R/(x, \mathfrak q_i))
\text{longueur}_{R_{\mathfrak q_i}}(M_{\mathfrak q_i}).
$$
où $\mathfrak q_1, \ldots, \mathfrak q_t$ sont les idéaux premiers
minimaux du support de $M$. Si $I \subset R$ est un idéal
tel que $x$ soit un non-diviseur de zéro dans $R/I$ et $\dim(R/I) = 1$, alors
$$
\text{longueur}_R(R/(x, I))
=
\sum\nolimits_i \text{longueur}_R(R/(x, \mathfrak q_i))
\text{longueur}_{R_{\mathfrak q_i}}((R/I)_{\mathfrak q_i})
$$
où $\mathfrak q_1, \ldots, \mathfrak q_n$ sont les idéaux premiers minimaux
contenant $I$.
\end{lemma}

\begin{proof}
Ce sont des cas particuliers du Lemme \ref{chow-lemma-length-multiplication}.
\end{proof}

\noindent
Voici un autre cas où l'on peut déterminer la valeur d'une multiplicité.

\begin{lemma}
\label{chow-lemma-powers-period-length-zero}
Soit $R$ un anneau. Soit $M$ un $R$-module.
Soit $\varphi : M \to M$ un endomorphisme et soit $n > 0$
tel que $\varphi^n = 0$ et que $\Ker(\varphi)/\Im(\varphi^{n - 1})$
soit de longueur finie comme $R$-module.
Alors
$$
e_R(M, \varphi^i, \varphi^{n - i}) = 0
$$
pour $i = 0, \ldots, n$.
\end{lemma}

\begin{proof}
Les cas $i = 0, n$ sont triviaux puisque, par convention, $\varphi^0 = \text{id}_M$.
Considérons $M$ comme un $R[t]$-module où la multiplication par $t$
est donnée par $\varphi$. Écrivons
$K_i = \Ker(t^i : M \to M)$ et
$$
a_i = \text{longueur}_R(K_i/t^{n - i}M),\quad
b_i = \text{longueur}_R(K_i/tK_{i + 1}),\quad
c_i = \text{longueur}_R(K_1/t^iK_{i + 1})
$$
Les valeurs au bord sont $a_0 = a_n = b_0 = c_0 = 0$.
Les $c_i$ sont des entiers pour $i < n$, car $K_1/t^iK_{i + 1}$
est un quotient de $K_1/t^{n - 1}M$, qui est supposé de longueur finie.
Nous utiliserons souvent l'égalité $K_i \cap t^jM = t^jK_{i + j}$.
Pour $0 < i < n - 1$, nous avons une suite exacte
$$
0 \to
K_1/t^{n - i - 1}K_{n - i} \to
K_{i + 1}/t^{n - i - 1}M \xrightarrow{t} K_i/t^{n - i}M
\to K_i/tK_{i + 1} \to 0
$$
Par récurrence sur $i$, on en déduit que $a_i$ et $b_i$ sont
des entiers pour $i < n$ et que
$$
c_{n - i - 1} - a_{i + 1} + a_i - b_i = 0
$$
Pour $0 < i < n - 1$, il existe une suite exacte
$$
0 \to
K_i/tK_{i + 1} \to
K_{i + 1}/tK_{i + 2} \xrightarrow{t^i}
K_1/t^{i + 1}K_{i + 2} \to
K_1/t^iK_{i + 1} \to 0
$$
qui donne
$$
b_i - b_{i + 1} + c_{i + 1} - c_i = 0
$$
Comme $b_0 = c_0$, on en conclut que $b_i = c_i$ pour $i < n$.
On voit alors que
$$
a_2 = a_1 + b_{n - 2} - b_1,\quad
a_3 = a_2 + b_{n - 3} - b_2,\quad \ldots
$$
On vérifie immédiatement que cela implique $a_i = a_{n - i}$, comme voulu.
\end{proof}

\begin{lemma}
\label{chow-lemma-multiply-period-length}
Soit $(R, \mathfrak m)$ un anneau local noethérien. Soit
$(M, \varphi, \psi)$ un complexe $(2, 1)$-périodique sur $R$,
avec $M$ fini et dont les groupes de cohomologie sont de longueur finie sur $R$.
Soit $x \in R$ tel que $\dim(\text{Supp}(M/xM)) \leq 0$. Alors
$$
e_R(M, x\varphi, \psi) = e_R(M, \varphi, \psi) - e_R(\Im(\varphi), 0, x)
$$
et
$$
e_R(M, \varphi, x\psi) = e_R(M, \varphi, \psi) + e_R(\Im(\psi), 0, x)
$$
\end{lemma}

\begin{proof}
Nous ne démontrerons que la première formule, la seconde se démontrant
exactement de la même manière.
Soit $M' = M[x^\infty]$ le sous-module de $x$-torsion de $M$.
Considérons la suite exacte courte $0 \to M' \to M \to M'' \to 0$.
Alors $M''$ est sans $x$-torsion (Compléments sur l'algèbre, Lemme
\ref{more-algebra-lemma-divide-by-torsion}).
Comme $\varphi$ et $\psi$ envoient $M'$ dans $M'$,
nous obtenons une suite exacte courte
$$
0 \to (M', \varphi', \psi') \to (M, \varphi, \psi) \to
(M'', \varphi'', \psi'') \to 0
$$
de complexes $(2, 1)$-périodiques. De plus, on obtient une suite exacte courte
$0 \to M' \cap \Im(\varphi) \to \Im(\varphi) \to \Im(\varphi'') \to 0$.
Nous avons
$e_R(M', \varphi, \psi) = e_R(M', x\varphi, \psi) =
e_R(M' \cap \Im(\varphi), 0, x) = 0$
d'après le Lemme \ref{chow-lemma-compare-periodic-lengths}.
Par additivité (Lemme \ref{chow-lemma-additivity-periodic-length}),
on voit qu'il suffit de démontrer le lemme pour $(M'', \varphi'', \psi'')$.
Cela nous ramène au cas étudié dans le paragraphe suivant.

\medskip\noindent
Supposons que $x : M \to M$ soit injective.
Dans ce cas, $\Ker(x\varphi) = \Ker(\varphi)$.
D'autre part, nous avons une suite exacte courte
$$
0 \to \Im(\varphi)/x\Im(\varphi) \to
\Ker(\psi)/\Im(x\varphi) \to \Ker(\psi)/\Im(\varphi) \to 0
$$
Ceci, joint à (\ref{chow-equation-multiplicity-coker-ker}), démontre la formule.
\end{proof}







\section{Préparation aux symboles modérés}
\label{chow-section-preparation-tame-symbol}

\noindent
Dans cette section, nous donnons quelques lemmes qui nous aideront à définir le
symbole modéré à la section suivante.

\begin{lemma}
\label{chow-lemma-glue-at-max}
Soit $A$ un anneau noethérien. Soient $\mathfrak m_1, \ldots, \mathfrak m_r$
des idéaux maximaux deux à deux distincts de $A$. Pour $i = 1, \ldots, r$, soit
$\varphi_i : A_{\mathfrak m_i} \to B_i$ un morphisme d'anneaux dont le
noyau et le conoyau sont annulés par une puissance
de $\mathfrak m_i$. Alors il existe un morphisme d'anneaux $\varphi : A \to B$ tel
que
\begin{enumerate}
\item le localisé de $\varphi$ en $\mathfrak m_i$ soit
isomorphe à $\varphi_i$, et
\item $\Ker(\varphi)$ et $\Coker(\varphi)$ soient annulés
par une puissance de $\mathfrak m_1 \cap \ldots \cap \mathfrak m_r$.
\end{enumerate}
De plus, si chaque $\varphi_i$ est fini, injectif ou
surjectif, il en va de même de $\varphi$.
\end{lemma}

\begin{proof}
Posons $I = \mathfrak m_1 \cap \ldots \cap \mathfrak m_r$. Posons
$A_i = A_{\mathfrak m_i}$ et $A' = \prod A_i$.
Alors $IA' = \prod \mathfrak m_i A_i$ et $A \to A'$
est un morphisme d'anneaux plat tel que $A/I \cong A'/IA'$.
Nous pouvons donc appliquer Compléments sur l'algèbre, Lemme
\ref{more-algebra-lemma-application-formal-glueing}
pour voir qu'il existe une application de $A$-modules $\varphi : A \to B$
tel que $\varphi_i$ soit isomorphe à la localisation de $\varphi$
en $\mathfrak m_i$. On peut ensuite utiliser la discussion de
Compléments sur l'algèbre, Remarque \ref{more-algebra-remark-formal-glueing-algebras}
pour munir $B$ d'une structure de $A$-algèbre
qui induit la structure de $A$-algèbre donnée sur $B_i$.
La dernière assertion du lemme résulte aisément du
fait que $\Ker(\varphi)_{\mathfrak m_i} \cong \Ker(\varphi_i)$
et $\Coker(\varphi)_{\mathfrak m_i} \cong \Coker(\varphi_i)$.
\end{proof}

\noindent
Le lemme suivant est très semblable à
Algèbre, Lemme \ref{algebra-lemma-nonregular-dimension-one}.

\begin{lemma}
\label{chow-lemma-Noetherian-domain-dim-1-two-elements}
Soit $(R, \mathfrak m)$ un anneau local noethérien de dimension $1$.
Soient $a, b \in R$ des non-diviseurs de zéro.
Il existe une extension finie d'anneaux $R \subset R'$
telle que $R'/R$ soit annulé par une puissance de $\mathfrak m$,
et des non-diviseurs de zéro $t, a', b' \in R'$ tels que
$a = ta'$, $b = tb'$ et $R' = a'R' + b'R'$.
\end{lemma}

\begin{proof}
Si $a$ ou $b$ est une unité, le lemme est vrai avec $R = R'$.
Nous pouvons donc supposer $a, b \in \mathfrak m$.
Posons $I = (a, b)$. L'idée consiste à éclater $R$ le long de $I$.
Au lieu d'un raisonnement algébrique, raisonnons géométriquement.
Soit $X = \text{Proj}(\bigoplus_{d \geq 0} I^d)$.
D'après Diviseurs, Lemme
\ref{divisors-lemma-blowing-up-gives-effective-Cartier-divisor},
le morphisme $X \to \Spec(R)$ est un isomorphisme au-dessus du
spectre épointé $U = \Spec(R) \setminus \{\mathfrak m\}$.
Nous pouvons et allons donc considérer $U$ comme un sous-schéma ouvert de $X$.
Le morphisme $X \to \Spec(R)$ est projectif d'après
Diviseurs, Lemme \ref{divisors-lemma-blowing-up-projective}.
De plus, tout point générique de $X$ appartient à $U$, par exemple
d'après Diviseurs, Lemme \ref{divisors-lemma-blow-up-and-irreducible-components}.
Il résulte de Variétés, Lemme \ref{varieties-lemma-finite-in-codim-1}
que $X \to \Spec(R)$ est fini. Ainsi $X = \Spec(R')$ est
affine et $R \to R'$ est fini. Nous avons $R_a \cong R'_a$ puisque $U = D(a)$.
Par conséquent, une puissance de $a$ annule le $R$-module fini $R'/R$.
Comme $\mathfrak m = \sqrt{(a)}$, le module $R'/R$ est annulé
par une puissance de $\mathfrak m$. D'après
Diviseurs, Lemme \ref{divisors-lemma-blowing-up-gives-effective-Cartier-divisor},
on voit que $IR'$ est un idéal localement principal.
Comme $R'$ est semi-local, $IR'$ est principal,
voir Algèbre, Lemme \ref{algebra-lemma-locally-free-semi-local-free},
disons $IR' = (t)$. Alors $a = a't$ et $b = b't$, et tout est
clair.
\end{proof}

\begin{lemma}
\label{chow-lemma-not-infinitely-divisible}
Soit $(R, \mathfrak m)$ un anneau local noethérien de dimension $1$.
Soient $a, b \in R$ des non-diviseurs de zéro avec $a \in \mathfrak m$.
Il existe un entier $n = n(R, a, b)$ tel que, pour toute extension finie
d'anneaux $R \subset R'$, si $b = a^m c$ pour un certain $c \in R'$, alors $m \leq n$.
\end{lemma}

\begin{proof}
Choisissons un idéal premier minimal $\mathfrak q \subset R$. Observons que
$\dim(R/\mathfrak q) = 1$ ; en particulier, $R/\mathfrak q$ n'est pas un corps.
On peut choisir un anneau de valuation discrète $A$ dominant $R/\mathfrak q$
et ayant le même corps des fractions, voir
Algèbre, Lemme \ref{algebra-lemma-dominate-by-dimension-1}. Observons que
$a$ et $b$ s'envoient sur des éléments non nuls de $A$, puisque les non-diviseurs de zéro de $R$
n'appartiennent pas à $\mathfrak q$. Soit $v$ la valuation discrète sur $A$.
Alors $v(a) > 0$ puisque $a \in \mathfrak m$.
Nous affirmons que $n = v(b)/v(a)$ convient.

\medskip\noindent
Soit donnée une extension $R \subset R'$. Posons $A' = A \otimes_R R'$.
Comme $\Spec(R') \to \Spec(R)$ est surjectif
(Algèbre, Lemme \ref{algebra-lemma-integral-overring-surjective}),
$\Spec(A') \to \Spec(A)$ est aussi surjectif
(Algèbre, Lemme \ref{algebra-lemma-surjective-spec-radical-ideal}).
Choisissons un idéal premier $\mathfrak q' \subset A'$ au-dessus de $(0) \subset A$.
Alors $A \subset A'' = A'/\mathfrak q'$ est une extension finie d'anneaux
(induisant de nouveau une surjection sur les spectres).
Choisissons un idéal maximal $\mathfrak m'' \subset A''$
au-dessus de l'idéal maximal de $A$, ainsi qu'un anneau de valuation discrète
$A'''$ dominant $A''_{\mathfrak m''}$ (voir le lemme cité ci-dessus).
Alors $A \to A'''$ est une extension d'anneaux de valuation discrète
et nous avons $b = a^m c$ dans $A'''$. Ainsi $v'''(b) \geq mv'''(a)$.
Comme $v''' = ev$, où $e$ est l'indice de ramification
de $A'''/A$, on obtient $m \leq n$, comme voulu.
\end{proof}

\begin{lemma}
\label{chow-lemma-prepare-tame-symbol}
Soit $(A, \mathfrak m)$ un anneau local noethérien de dimension $1$.
Soit $r \geq 2$ et soient $a_1, \ldots, a_r \in A$ des non-diviseurs de zéro
qui ne sont pas tous des unités.
Alors il existe
\begin{enumerate}
\item une extension finie d'anneaux $A \subset B$ telle que
$B/A$ soit annulé par une puissance de $\mathfrak m$,
\item pour tout idéal maximal $\mathfrak m_j \subset B$,
un non-diviseur de zéro $\pi_j \in B_j = B_{\mathfrak m_j}$, et
\item des factorisations $a_i = u_{i, j} \pi_j^{e_{i, j}}$ dans $B_j$,
où $u_{i, j} \in B_j$ est une unité et $e_{i, j} \geq 0$.
\end{enumerate}
\end{lemma}

\begin{proof}
Puisqu'au moins un $a_i$ n'est pas une unité, $\mathfrak m$
n'est pas un idéal premier associé de $A$. De plus, pour toute extension $A \subset B$
comme dans l'énoncé, $\mathfrak m$ n'est pas un idéal premier associé de $B$
et $\mathfrak m_j$ n'est pas un idéal premier associé de $B_j$.
Garder cela à l'esprit aidera à vérifier les raisonnements ci-dessous.

\medskip\noindent
Tout d'abord, nous affirmons qu'il suffit de démontrer le lemme pour $r = 2$.
Nous allons le montrer par récurrence sur $r$ ; nous suggérons au lecteur
de sauter la démonstration. Supposons donnés $A \subset B$ et des $\pi_j$ dans
$B_j = B_{\mathfrak m_j}$, ainsi que des factorisations
$a_i = u_{i, j} \pi_j^{e_{i, j}}$ pour $i = 1, \ldots, r - 1$ dans $B_j$,
où $u_{i, j} \in B_j$ est une unité et $e_{i, j} \geq 0$.
En appliquant le cas $r = 2$ à $\pi_j$ et $a_r$ dans $B_j$,
on peut trouver des extensions $B_j \subset C_j$ et, pour tout idéal maximal
$\mathfrak m_{j, k}$ de $C_j$, un non-diviseur de zéro
$\pi_{j, k} \in C_{j, k} = (C_j)_{\mathfrak m_{j, k}}$
et des factorisations
$$
\pi_j = v_{j, k} \pi_{j, k}^{f_{j, k}}
\quad\text{et}\quad
a_r = w_{j, k} \pi_{j, k}^{g_{j, k}}
$$
comme dans le lemme. Il existe une unique extension finie $B \subset C$
telle que $C/B$ soit annulé par une puissance de $\mathfrak m$ et
que $C_j \cong C_{\mathfrak m_j}$ pour tout $j$, voir
le Lemme \ref{chow-lemma-glue-at-max}.
Les idéaux maximaux de $C$ correspondent $1$ à $1$
aux idéaux maximaux $\mathfrak m_{j, k}$ dans les localisés,
et dans ces localisés nous avons
$$
a_i = u_{i, j} \pi_j^{e_{i, j}} =
u_{i, j} v_{j, k}^{e_{i, j}} \pi_{j, k}^{e_{i, j}f_{j, k}}
$$
pour $i \leq r - 1$. Comme $a_r$ se factorise lui aussi correctement, la
démonstration de l'étape de récurrence est achevée.

\medskip\noindent
Démonstration du cas $r = 2$. Nous raisonnons par récurrence sur
$$
\ell = \min(\text{longueur}_A(A/a_1A),\ \text{longueur}_A(A/a_2A)).
$$
Si $\ell = 0$, alors $a_1$ ou $a_2$ est une unité et
le lemme vaut avec $A = B$. Nous pouvons donc supposer $\ell > 0$.

\medskip\noindent
Supposons donnée une extension finie d'anneaux $A \subset A'$ telle que
$A'/A$ soit annulé par une puissance de $\mathfrak m$ et que
$\mathfrak m$ ne soit pas un idéal premier associé de $A'$.
Soient $\mathfrak m_1, \ldots, \mathfrak m_r \subset A'$
les idéaux maximaux, et posons $A'_i = A'_{\mathfrak m_i}$.
Si l'on peut résoudre le problème pour $a_1, a_2$ dans chaque $A'_i$,
alors on peut appliquer le Lemme \ref{chow-lemma-glue-at-max}
pour produire une solution pour $a_1, a_2$ dans $A$.
Choisissons $x \in \{a_1, a_2\}$ tel que $\ell = \text{longueur}_A(A/xA)$.
D'après le Lemme \ref{chow-lemma-compare-periodic-lengths} 
et (\ref{chow-equation-multiplicity-coker-ker}),
nous avons $\text{longueur}_A(A/xA) = \text{longueur}_A(A'/xA')$.
D'autre part,
$$
\text{longueur}_A(A'/xA') =
\sum [\kappa(\mathfrak m_i) : \kappa(\mathfrak m)]
\text{longueur}_{A'_i}(A'_i/xA'_i)
$$
d'après Algèbre, Lemme \ref{algebra-lemma-pushdown-module}.
Comme $x \in \mathfrak m$, chaque terme du membre de droite
est positif. On en conclut que l'hypothèse de récurrence s'applique
à $a_1, a_2$ dans chaque $A'_i$ si $r > 1$, ou si $r = 1$ et
$[\kappa(\mathfrak m_1) : \kappa(\mathfrak m)] > 1$.
On en conclut que l'on peut supposer que tout $A'$ comme ci-dessus est local et a
le même corps résiduel que $A$.

\medskip\noindent
En appliquant la discussion du paragraphe précédent,
nous pouvons remplacer $A$ par l'anneau construit dans le
Lemme \ref{chow-lemma-Noetherian-domain-dim-1-two-elements}
pour $a_1, a_2 \in A$. Puisque $A$ est local, nous obtenons alors,
quitte à échanger $a_1$ et $a_2$, que $a_2 \in (a_1)$.
Écrivons $a_2 = a_1^m c$, avec $m > 0$ maximal. En fait, d'après le
Lemme \ref{chow-lemma-not-infinitely-divisible},
nous pouvons supposer que $m$ demeure maximal même après avoir remplacé $A$
par toute extension finie $A \subset A'$ comme au paragraphe précédent.
Si $c$ est une unité, nous avons terminé. Sinon, nous remplaçons
$A$ par l'anneau construit dans le
Lemme \ref{chow-lemma-Noetherian-domain-dim-1-two-elements}
pour $a_1, c \in A$. Alors ou bien (1) $c = a_1 c'$, ou bien
(2) $a_1 = c a'_1$. Le premier cas est impossible, puisqu'il
donnerait $a_2 = a_1^{m + 1} c'$, ce qui contredirait la
maximalité de $m$. Dans le second cas, on obtient
$a_1 = c a'_1$ et $a_2 = c^{m + 1} (a'_1)^m$.
Il suffit alors de démontrer le lemme pour $A$ et $c, a'_1$.
Si $a'_1$ est une unité, nous avons terminé ; sinon,
$\text{longueur}_A(A/cA) < \ell$, car $cA$ est un idéal strictement
plus grand que $a_1A$. Nous concluons donc par l'hypothèse de récurrence.
\end{proof}






\section{Symboles modérés}
\label{chow-section-tame-symbol}

\noindent
Considérons un anneau local noethérien $(A, \mathfrak m)$ de dimension $1$.
Notons $Q(A)$ l'anneau total des fractions de $A$, voir
Algèbre, Exemple \ref{algebra-example-localize-at-prime}.
Le {\it symbole modéré} sera une application
$$
\partial_A(-, -) : Q(A)^* \times Q(A)^* \longrightarrow \kappa(\mathfrak m)^*
$$
qui satisfait les propriétés suivantes :
\begin{enumerate}
\item $\partial_A(f, gh) = \partial_A(f, g) \partial_A(f, h)$
\label{chow-item-bilinear}
pour $f, g, h \in Q(A)^*$,
\item $\partial_A(f, g) \partial_A(g, f) = 1$
\label{chow-item-skew}
pour $f, g \in Q(A)^*$,
\item $\partial_A(f, 1 - f) = 1$
\label{chow-item-1-x}
pour $f \in Q(A)^*$ tel que $1 - f \in Q(A)^*$,
\item $\partial_A(aa', b) = \partial_A(a, b)\partial_A(a', b)$
\label{chow-item-bilinear-better}
et $\partial_A(a, bb') = \partial_A(a, b)\partial_A(a, b')$
pour des non-diviseurs de zéro $a, a', b, b' \in A$,
\item $\partial_A(b, b) = (-1)^m$
\label{chow-item-skew-better}
avec $m = \text{longueur}_A(A/bA)$,
pour tout non-diviseur de zéro $b \in A$,
\item $\partial_A(u, b) = u^m \bmod \mathfrak m$
\label{chow-item-normalization}
avec $m = \text{longueur}_A(A/bA)$, pour toute unité $u \in A$ et tout
non-diviseur de zéro $b \in A$, et
\item
\label{chow-item-1-x-better}
$\partial_A(a, b - a)\partial_A(b, b) = \partial_A(b, b - a)\partial_A(a, b)$
pour $a, b \in A$ tels que $a, b, b - a$ soient des non-diviseurs de zéro.
\end{enumerate}
Comme il est plus facile de travailler avec des éléments de $A$, nous
considérerons souvent $\partial_A$ comme une application définie sur les couples de
non-diviseurs de zéro de $A$ qui satisfait (\ref{chow-item-bilinear-better}),
(\ref{chow-item-skew-better}), (\ref{chow-item-normalization}),
(\ref{chow-item-1-x-better}). C'est un exercice de vérifier qu'en
posant
$$
\partial_A(\frac{a}{b}, \frac{c}{d}) =
\partial_A(a, c) \partial_A(a, d)^{-1} \partial_A(b, c)^{-1} \partial_A(b, d)
$$
on obtient une application bien définie $Q(A)^* \times Q(A)^* \to \kappa(\mathfrak m)^*$
qui satisfait (\ref{chow-item-bilinear}), (\ref{chow-item-skew}), (\ref{chow-item-1-x}),
ainsi que les autres propriétés.

\medskip\noindent
Nous ne prétendons pas qu'il existe une unique application ayant ces propriétés.
Nous allons plutôt donner une recette pour construire une telle application.
Plus précisément, étant donnés des non-diviseurs de zéro $a_1, a_2 \in A$, nous choisissons
une extension d'anneaux $A \subset B$ et des factorisations locales
comme dans le Lemme \ref{chow-lemma-prepare-tame-symbol}.
Puis nous définissons
\begin{equation}
\label{chow-equation-tame-symbol}
\partial_A(a_1, a_2) = \prod\nolimits_j
\text{Norme}_{\kappa(\mathfrak m_j)/\kappa(\mathfrak m)}
((-1)^{e_{1, j}e_{2, j}}u_{1, j}^{e_{2, j}}u_{2, j}^{-e_{1, j}}
\bmod \mathfrak m_j)^{m_j}
\end{equation}
où $m_j = \text{longueur}_{B_j}(B_j/\pi_j B_j)$ et le produit
est pris sur les idéaux maximaux $\mathfrak m_1, \ldots, \mathfrak m_r$ de $B$.

\begin{lemma}
\label{chow-lemma-well-defined-tame-symbol}
La formule (\ref{chow-equation-tame-symbol}) détermine un
élément bien défini de $\kappa(\mathfrak m)^*$. Autrement dit, le
membre de droite ne dépend ni du choix des
factorisations locales, ni du choix de $B$.
\end{lemma}

\begin{proof}
Indépendance du choix des factorisations. Supposons donnés
un anneau local noethérien de dimension $1$, noté $B$, des éléments $a_1, a_2 \in B$,
et des non-diviseurs de zéro $\pi, \theta$ tels que l'on puisse écrire
$$
a_1 = u_1 \pi^{e_1} = v_1 \theta^{f_1},\quad
a_2 = u_2 \pi^{e_2} = v_2 \theta^{f_2}
$$
avec des entiers $e_i, f_i \geq 0$ et des unités $u_i, v_i$ de $B$.
Observons que cela implique
$$
a_1^{e_2} = u_1^{e_2}u_2^{-e_1}a_2^{e_1},\quad
a_1^{f_2} = v_1^{f_2}v_2^{-f_1}a_2^{f_1}
$$
D'autre part, en posant
$m = \text{longueur}_B(B/\pi B)$ et $k = \text{longueur}_B(B/\theta B)$,
on obtient $e_2 m = \text{longueur}_B(B/a_2 B) = f_2 k$.
En développant $a_1^{e_2m} = a_1^{f_2 k}$ à l'aide des égalités ci-dessus, on trouve
$$
(u_1^{e_2}u_2^{-e_1})^m =  (v_1^{f_2}v_2^{-f_1})^k
$$
Cela démontre l'égalité voulue au signe près. Pour vérifier que les signes
conviennent, il faut montrer que $me_1e_2$ est pair si et seulement si
$kf_1f_2$ est pair. Cela résulte des deux égalités $me_2 = kf_2$ et
$me_1 = kf_1$ (par le même raisonnement que ci-dessus).

\medskip\noindent
Indépendance du choix de $B$. Supposons données deux extensions
$A \subset B$ et $A \subset B'$ comme dans le Lemme \ref{chow-lemma-prepare-tame-symbol}.
Alors
$$
C = (B \otimes_A B')/(\mathfrak m\text{-torsion})
$$
en fournira une troisième. Nous pouvons donc supposer que nous avons
$A \subset B \subset C$, ainsi que des factorisations sur les
anneaux locaux de $B$, et il faut montrer que l'emploi
des mêmes factorisations sur les anneaux locaux de $C$
donne le même élément de $\kappa(\mathfrak m)$.
Par transitivité des normes
(Corps, Lemme \ref{fields-lemma-trace-and-norm-tower}),
cela se ramène au problème suivant :
si $B$ est un anneau local noethérien de dimension $1$
et $\pi \in B$ un non-diviseur de zéro, alors
$$
\lambda^m = \prod \text{Norme}_{\kappa_k/\kappa}(\lambda)^{m_k}
$$
Nous avons employé ici les notations suivantes :
(1) $\kappa$ est le corps résiduel de $B$,
(2) $\lambda$ est un élément de $\kappa$,
(3) les $\mathfrak m_k \subset C$ sont les idéaux maximaux de $C$,
(4) $\kappa_k = \kappa(\mathfrak m_k)$ est le corps résiduel de
$C_k = C_{\mathfrak m_k}$,
(5) $m = \text{longueur}_B(B/\pi B)$, et
(6) $m_k = \text{longueur}_{C_k}(C_k/\pi C_k)$.
L'égalité affichée vaut parce que
$\text{Norme}_{\kappa_k/\kappa}(\lambda) = \lambda^{[\kappa_k : \kappa]}$,
puisque $\lambda \in \kappa$, et parce que $m = \sum m_k[\kappa_k:\kappa]$.
Premièrement, nous avons $m = \text{longueur}_B(B/xB) = \text{longueur}_B(C/\pi C)$
d'après le Lemme \ref{chow-lemma-compare-periodic-lengths} 
et (\ref{chow-equation-multiplicity-coker-ker}).
Enfin, $\text{longueur}_B(C/\pi C) = \sum m_k[\kappa_k:\kappa]$
d'après Algèbre, Lemme \ref{algebra-lemma-pushdown-module}.
\end{proof}

\begin{lemma}
\label{chow-lemma-tame-symbol}
Le symbole modéré (\ref{chow-equation-tame-symbol}) satisfait
(\ref{chow-item-bilinear-better}), (\ref{chow-item-skew-better}),
(\ref{chow-item-normalization}), (\ref{chow-item-1-x-better}) et fournit donc
une application $\partial_A : Q(A)^* \times Q(A)^* \to \kappa(\mathfrak m)^*$
qui satisfait (\ref{chow-item-bilinear}), (\ref{chow-item-skew}), (\ref{chow-item-1-x}).
\end{lemma}

\begin{proof}
Démontrons (\ref{chow-item-bilinear-better}).
Soient $a_1, a_2, a_3 \in A$ des non-diviseurs de zéro.
Choisissons $A \subset B$ comme dans le Lemme \ref{chow-lemma-prepare-tame-symbol}
pour $a_1, a_2, a_3$. Alors l'égalité
$$
\partial_A(a_1a_2, a_3) = \partial_A(a_1, a_3) \partial_A(a_2, a_3)
$$
résulte de l'égalité
$$
(-1)^{(e_{1, j} + e_{2, j})e_{3, j}}
(u_{1, j}u_{2, j})^{e_{3, j}}u_{3, j}^{-e_{1, j} - e_{2, j}} =
(-1)^{e_{1, j}e_{3, j}}
u_{1, j}^{e_{3, j}}u_{3, j}^{-e_{1, j}}
(-1)^{e_{2, j}e_{3, j}}
u_{2, j}^{e_{3, j}}u_{3, j}^{-e_{2, j}}
$$
dans $B_j$. Les propriétés (\ref{chow-item-skew-better}) et
(\ref{chow-item-normalization}) sont tout aussi immédiates.

\medskip\noindent
Démontrons (\ref{chow-item-1-x-better}). Soient $a_1, a_2, a_1 - a_2 \in A$
des non-diviseurs de zéro et posons $a_3 = a_1 - a_2$.
Choisissons $A \subset B$ comme dans le Lemme \ref{chow-lemma-prepare-tame-symbol}
pour $a_1, a_2, a_3$. Il suffit alors de montrer que
$$
(-1)^{e_{1, j}e_{2, j} + e_{1, j}e_{3, j} + e_{2, j}e_{3, j} + e_{2, j}}
u_{1, j}^{e_{2, j} - e_{3, j}}
u_{2, j}^{e_{3, j} - e_{1, j}}
u_{3, j}^{e_{1, j} - e_{2, j}} \bmod \mathfrak m_j = 1
$$
C'est clair si $e_{1, j} = e_{2, j} = e_{3, j}$.
Supposons $e_{1, j} > e_{2, j}$. Alors $e_{3, j} = e_{2, j}$,
car $a_3 = a_1 - a_2$, et $u_{3, j}$
a la même classe résiduelle que $-u_{2, j}$. Par conséquent,
la formule est vraie ; les signes conviennent eux aussi,
et cette vérification explique le choix des signes
dans (\ref{chow-equation-tame-symbol}).
Les autres cas se traitent exactement de la même manière.
\end{proof}

\begin{lemma}
\label{chow-lemma-norm-down-tame-symbol}
Soit $(A, \mathfrak m)$ un anneau local noethérien de dimension $1$.
Soit $A \subset B$ une extension finie d'anneaux telle que $B/A$
soit annulé par une puissance de $\mathfrak m$ et que $\mathfrak m$ ne soit pas
un idéal premier associé de $B$.
Pour des non-diviseurs de zéro $a, b \in A$, nous avons
$$
\partial_A(a, b) = \prod
\text{Norme}_{\kappa(\mathfrak m_j)/\kappa(\mathfrak m)}(\partial_{B_j}(a, b))
$$
où le produit est pris sur les idéaux maximaux $\mathfrak m_j$ de $B$
et $B_j = B_{\mathfrak m_j}$.
\end{lemma}

\begin{proof}
Choisissons $B_j \subset C_j$ comme dans le
Lemme \ref{chow-lemma-prepare-tame-symbol} pour $a, b$.
D'après le Lemme \ref{chow-lemma-glue-at-max}, on peut choisir une extension finie
d'anneaux $B \subset C$ telle que $C_j \cong C_{\mathfrak m_j}$ pour tout $j$.
Soient $\mathfrak m_{j, k} \subset C$ les idéaux maximaux de $C$
au-dessus de $\mathfrak m_j$. Soient
$$
a = u_{j, k}\pi_{j, k}^{f_{j, k}},\quad
b = v_{j, k}\pi_{j, k}^{g_{j, k}}
$$
les factorisations locales dont l'existence résulte de notre choix de
$C_j \cong C_{\mathfrak m_j}$. Par définition, nous avons
$$
\partial_A(a, b) = 
\prod\nolimits_{j, k}
\text{Norme}_{\kappa(\mathfrak m_{j, k})/\kappa(\mathfrak m)}
((-1)^{f_{j, k}g_{j, k}}u_{j, k}^{g_{j, k}}v_{j, k}^{-f_{j, k}}
\bmod \mathfrak m_{j, k})^{m_{j, k}}
$$
et
$$
\partial_{B_j}(a, b) = 
\prod\nolimits_k
\text{Norme}_{\kappa(\mathfrak m_{j, k})/\kappa(\mathfrak m_j)}
((-1)^{f_{j, k}g_{j, k}}u_{j, k}^{g_{j, k}}v_{j, k}^{-f_{j, k}}
\bmod \mathfrak m_{j, k})^{m_{j, k}}
$$
Le résultat découle de la transitivité des normes
pour $\kappa(\mathfrak m_{j, k})/\kappa(\mathfrak m_j)/\kappa(\mathfrak m)$, voir
Corps, Lemme \ref{fields-lemma-trace-and-norm-tower}.
\end{proof}

\begin{lemma}
\label{chow-lemma-tame-symbol-formally-smooth}
Soit $(A, \mathfrak m, \kappa) \to (A', \mathfrak m', \kappa')$
un homomorphisme local d'anneaux locaux noethériens. Supposons $A \to A'$
plat et $\dim(A) = \dim(A') = 1$. Posons
$m = \text{longueur}_{A'}(A'/\mathfrak mA')$.
Pour des non-diviseurs de zéro $a_1, a_2 \in A$,
$\partial_A(a_1, a_2)^m$ s'envoie sur $\partial_{A'}(a_1, a_2)$
par $\kappa \to \kappa'$.
\end{lemma}

\begin{proof}
Si $a_1, a_2$ sont tous deux des unités, alors $\partial_A(a_1, a_2) = 1$,
$\partial_{A'}(a_1, a_2) = 1$, et le résultat est vrai.
Sinon, on peut choisir une extension d'anneaux $A \subset B$ et
des factorisations locales comme dans le Lemme \ref{chow-lemma-prepare-tame-symbol}.
Notons $\mathfrak m_1, \ldots, \mathfrak m_m$
les idéaux maximaux de $B$. Soient $\mathfrak m_1, \ldots, \mathfrak m_m$
les idéaux maximaux de $B$, de corps résiduels $\kappa_1, \ldots, \kappa_m$.
Pour chaque $j \in \{1, \ldots, m\}$, notons $\pi_j \in B_j = B_{\mathfrak m_j}$
un non-diviseur de zéro tel que nous ayons des factorisations
$a_i = u_{i, j}\pi_j^{e_{i, j}}$ comme dans le lemme.
Par définition, nous avons
$$
\partial_A(a_1, a_2) = \prod\nolimits_j
\text{Norme}_{\kappa_j/\kappa}
((-1)^{e_{1, j}e_{2, j}}u_{1, j}^{e_{2, j}}u_{2, j}^{-e_{1, j}}
\bmod \mathfrak m_j)^{m_j}
$$
où $m_j = \text{longueur}_{B_j}(B_j/\pi_j B_j)$.

\medskip\noindent
Posons $B' = A' \otimes_A B$. Comme $A'$ est plat sur $A$, on voit
que $A' \subset B'$ est une extension d'anneaux telle que $B'/A'$ soit annulé
par une puissance de $\mathfrak m'$. Soient
$$
\mathfrak m'_{j, l},\quad l = 1, \ldots, n_j
$$
les idéaux maximaux de $B'$ au-dessus de $\mathfrak m_j$. Notons
$\kappa'_{j, l}$ le corps résiduel de $\mathfrak m'_{j, l}$. Notons
$B'_{j, l}$ le localisé de $B'$ en $\mathfrak m'_{j, l}$.
Comme factorisations de $a_1$ et $a_2$ dans $B'_{j, l}$,
nous utilisons l'image des factorisations
$a_i = u_{i, j} \pi_j^{e_{i, j}}$ qui nous sont données dans $B_j$.
Par définition, nous avons
$$
\partial_{A'}(a_1, a_2) = \prod\nolimits_{j, l}
\text{Norme}_{\kappa'_{j, l}/\kappa'}
((-1)^{e_{1, j}e_{2, j}}u_{1, j}^{e_{2, j}}u_{2, j}^{-e_{1, j}}
\bmod \mathfrak m'_{j, l})^{m'_{j, l}}
$$
où $m'_{j, l} = \text{longueur}_{B'_{j, l}}(B'_{j, l}/\pi_j B'_{j, l})$.

\medskip\noindent
En comparant les formules ci-dessus, on voit qu'il suffit de montrer que
pour tout $j$ et toute unité $u \in B_j$, on a
\begin{equation}
\label{chow-equation-to-prove}
\left(\text{Norme}_{\kappa_j/\kappa}(u \bmod \mathfrak m_j)^{m_j}\right)^m
=
\prod\nolimits_l
\text{Norme}_{\kappa'_{j, l}/\kappa'}(u \bmod \mathfrak m'_{j, l})^{m'_{j, l}}
\end{equation}
dans $\kappa'$. Pour le démontrer, nous allons utiliser la construction des déterminants
d'endomorphismes de modules de longueur finie donnée dans
Compléments sur l'algèbre, section \ref{more-algebra-section-determinants-finite-length}.
Posons $M = B_j/\pi_j B_j$. D'après
Compléments sur l'algèbre, Lemme \ref{more-algebra-lemma-multiplication}, nous avons
$$
\text{Norme}_{\kappa_j/\kappa}(u \bmod \mathfrak m_j)^{m_j} =
\det\nolimits_\kappa(u : M \to M)
$$
Ainsi, d'après
Compléments sur l'algèbre, Lemme \ref{more-algebra-lemma-flat-base-change-det},
le membre de gauche de (\ref{chow-equation-to-prove}) est égal à
$\det_{\kappa'}(u : M \otimes_A A' \to M \otimes_A A')$.
Nous avons un isomorphisme
$$
M \otimes_A A' = (B_j/\pi_j B_j) \otimes_A A' =
\bigoplus\nolimits_l B'_{j, l}/\pi_j B'_{j, l}
$$
de $A'$-modules. En posant $M'_l = B'_{j, l}/\pi_j B'_{j, l}$, on voit que
$\text{Norme}_{\kappa'_{j, l}/\kappa'}(u \bmod \mathfrak m'_{j, l})^{m'_{j, l}}
= \det_{\kappa'}(u_j : M'_l \to M'_l)$ d'après, de nouveau,
Compléments sur l'algèbre, Lemme \ref{more-algebra-lemma-multiplication}.
Ainsi (\ref{chow-equation-to-prove}) résulte de la multiplicativité de la construction
du déterminant, voir Compléments sur l'algèbre, Lemme \ref{more-algebra-lemma-ses}.
\end{proof}






\section{Un lemme clé}
\label{chow-section-key-lemma}

\noindent
Dans cette section, nous appliquons les résultats précédents pour démontrer le
Lemme \ref{chow-lemma-milnor-gersten-low-degree}.
Ce lemme est le cas en bas degré de l'assertion selon laquelle
il existe, pour la K-théorie de Milnor, un complexe analogue
au complexe de Gersten--Quillen en K-théorie de Quillen.
Voir la Remarque \ref{chow-remark-gersten-complex-milnor}.

\begin{lemma}
\label{chow-lemma-perpare-key}
Soit $(A, \mathfrak m)$ un anneau local noethérien de dimension $2$.
Soit $t \in \mathfrak m$ un non-diviseur de zéro. Écrivons
$V(t) = \{\mathfrak m, \mathfrak q_1, \ldots, \mathfrak q_r\}$.
Soit $A_{\mathfrak q_i} \subset B_i$ une extension finie d'anneaux
telle que $B_i/A_{\mathfrak q_i}$ soit annulé par une puissance de
$t$. Alors il existe une extension finie $A \subset B$
d'anneaux locaux, identifiant les corps résiduels,
telle que $B_i \cong B_{\mathfrak q_i}$ et que $B/A$ soit annulé
par une puissance de $t$.
\end{lemma}

\begin{proof}
Choisissons $n > 0$ tel que $B_i \subset t^{-n}A_{\mathfrak q_i}$.
Soient $M \subset t^{-n}A$, respectivement $M' \subset t^{-2n}A$, les
sous-$A$-modules formés des éléments dont les images appartiennent à $B_i$
dans $t^{-n}A_{\mathfrak q_i}$ et $t^{-2n}A_{\mathfrak q_i}$, respectivement.
L'inclusion $M \subset M'$ est alors une inclusion de $A$-modules finis, puisque $A$ est noethérien,
et $M_{\mathfrak q_i} = M'_{\mathfrak q_i} = B_i$, car la localisation
est exacte. Ainsi $M'/M$ est annulé par $\mathfrak m^c$ pour un certain
$c > 0$. Observons que $M \cdot M \subset M'$ pour la multiplication
$t^{-n}A \times t^{-n}A \to t^{-2n}A$. Par conséquent,
$B = A + \mathfrak m^{c + 1}M$ est une $A$-algèbre finie ayant les bons
localisés. Nous omettons de vérifier que $B$ est local, d'idéal maximal
$\mathfrak m + \mathfrak m^{c + 1}M$.
\end{proof}

\begin{lemma}
\label{chow-lemma-key-nonzerodivisors}
Soit $(A, \mathfrak m)$ un anneau local noethérien de dimension $2$.
Soient $a, b \in A$ des non-diviseurs de zéro.
Alors
$$
\sum
\text{ord}_{A/\mathfrak q}(\partial_{A_{\mathfrak q}}(a, b))
=
0
$$
où la somme porte sur les idéaux premiers de hauteur $1$, notés $\mathfrak q$, de $A$.
\end{lemma}

\begin{proof}
Si $\mathfrak q$ est un idéal premier de hauteur $1$ de $A$ tel que $a, b$
s'envoient sur une unité de $A_\mathfrak q$, alors $\partial_{A_\mathfrak q}(a, b) = 1$.
La somme est donc finie. En fait, si
$V(ab) = \{\mathfrak m, \mathfrak q_1, \ldots, \mathfrak q_r\}$,
la somme porte sur $i = 1, \ldots, r$.
Pour chaque $i$, choisissons une extension $A_{\mathfrak q_i} \subset B_i$
comme dans le Lemme \ref{chow-lemma-prepare-tame-symbol} pour $a, b$.
D'après le Lemme \ref{chow-lemma-perpare-key}, appliqué avec $t = ab$ et la liste d'idéaux premiers donnée,
nous pouvons supposer donnée une extension locale finie $A \subset B$
telle que $B/A$ soit annulé par une puissance de $ab$ et que,
pour chaque $i$, on ait $B_{\mathfrak q_i} \cong B_i$.
Observons que, si les $\mathfrak q_{i, j}$ sont les idéaux premiers de $B$
au-dessus de $\mathfrak q_i$, alors
$$
\text{ord}_{A/\mathfrak q_i}(\partial_{A_{\mathfrak q_i}}(a, b))
=
\sum\nolimits_j
\text{ord}_{B/\mathfrak q_{i, j}}(\partial_{B_{\mathfrak q_{i, j}}}(a, b))
$$
d'après le Lemme \ref{chow-lemma-norm-down-tame-symbol} et
Algèbre, Lemme \ref{algebra-lemma-finite-extension-dim-1}.
Nous pouvons donc remplacer $A$ par $B$ et
nous ramener au cas traité dans le paragraphe suivant.

\medskip\noindent
Supposons que, pour chaque $i$, il existe un non-diviseur de zéro
$\pi_i \in A_{\mathfrak q_i}$ et des unités $u_i, v_i \in A_{\mathfrak q_i}$
tels que, pour certains entiers $e_i, f_i \geq 0$, nous ayons
$$
a = u_i \pi_i^{e_i},\quad b = v_i \pi_i^{f_i}
$$
dans $A_{\mathfrak q_i}$. En posant
$m_i = \text{longueur}_{A_{\mathfrak q_i}}(A_{\mathfrak q_i}/\pi_i)$,
nous avons
$\partial_{A_{\mathfrak q_i}}(a, b) =
((-1)^{e_if_i}u_i^{f_i}v_i^{-e_i})^{m_i}$ par définition.
Comme $a, b$ sont des non-diviseurs de zéro, le
complexe $(2, 1)$-périodique $(A/(ab), a, b)$ a une cohomologie nulle.
Notons $M_i$ l'image de $A/(ab)$ dans $A_{\mathfrak q_i}/(ab)$.
Nous avons alors une application
$$
A/(ab) \longrightarrow \bigoplus M_i
$$
dont le noyau et le conoyau ont leur support dans $\{\mathfrak m\}$
et sont donc de longueur finie. On voit ainsi que
$$
\sum e_A(M_i, a, b) = 0
$$
d'après le Lemme \ref{chow-lemma-compare-periodic-lengths}. Il suffit donc de montrer
$e_A(M_i, a, b) =
- \text{ord}_{A/\mathfrak q_i}(\partial_{A_{\mathfrak q_i}}(a, b))$.

\medskip\noindent
Démontrons d'abord ceci dans le cas où $\pi_i, u_i, v_i$ sont les images
d'éléments $\pi_i, u_i, v_i \in A$ (l'emploi des mêmes symboles
ne devrait prêter à aucune confusion). Dans ce cas, on obtient
\begin{align*}
e_A(M_i, a, b) & =
e_A(M_i, u_i\pi_i^{e_i}, v_i\pi_i^{f_i}) \\
& =
e_A(M_i, \pi_i^{e_i}, \pi_i^{f_i}) -
e_A(\pi_i^{e_i}M_i, 0, u_i) +
e_A(\pi_i^{f_i}M_i, 0, v_i) \\
& =
0 -
f_im_i\text{ord}_{A/\mathfrak q_i}(u_i) +
e_im_i\text{ord}_{A/\mathfrak q_i}(v_i) \\
& =
-m_i\text{ord}_{A/\mathfrak q_i}(u_i^{f_i}v_i^{-e_i}) =
-\text{ord}_{A/\mathfrak q_i}(\partial_{A_{\mathfrak q_i}}(a, b))
\end{align*}
La deuxième égalité résulte du Lemme \ref{chow-lemma-multiply-period-length}.
Observons que
$M_i \subset (M_i)_{\mathfrak q_i} = A_{\mathfrak q_i}/(\pi_i^{e_i + f_i})$
et que
$(\pi_i^{e_i}M_i)_{\mathfrak q_i} \cong A_{\mathfrak q_i}/\pi_i^{f_i}$ et
$(\pi_i^{f_i}M_i)_{\mathfrak q_i} \cong A_{\mathfrak q_i}/\pi_i^{e_i}$.
Le $0$ de la troisième égalité provient du
Lemme \ref{chow-lemma-powers-period-length-zero},
et les deux autres termes proviennent du
Lemme \ref{chow-lemma-length-multiplication}.
Les deux dernières égalités résultent de la multiplicativité de
la fonction ordre et de la définition de notre symbole modéré.

\medskip\noindent
En général, on peut d'abord choisir $c \in A$, $c \not \in \mathfrak q_i$
tel que $c\pi_i \in A$. En remplaçant $\pi_i$ par $c\pi_i$,
$u_i$ par $c^{-e_i}u_i$ et $v_i$ par $c^{-f_i}v_i$,
nous pouvons et allons supposer que $\pi_i$ appartient à $A$.
Ensuite, choisissons un $c \in A$, $c \not \in \mathfrak q_i$,
tel que $cu_i, cv_i \in A$. On observe alors que
$$
e_A(M_i, ca, cb) = e_A(M_i, a, b) - e_A(aM_i, 0, c) + e_A(bM_i, 0, c)
$$
d'après le Lemme \ref{chow-lemma-length-multiplication}.
D'autre part, nous avons
$$
\partial_{A_{\mathfrak q_i}}(ca, cb) =
c^{m_i(f_i - e_i)}\partial_{A_{\mathfrak q_i}}(a, b)
$$
dans $\kappa(\mathfrak q_i)^*$, puisque $c$ est une unité de $A_{\mathfrak q_i}$.
Les raisonnements du paragraphe précédent montrent que
$e_A(M_i, ca, cb) = -
\text{ord}_{A/\mathfrak q_i}(\partial_{A_{\mathfrak q_i}}(ca, cb))$.
Il suffit donc de démontrer
$$
e_A(aM_i, 0, c) = \text{ord}_{A/\mathfrak q_i}(c^{m_if_i})
\quad\text{et}\quad
e_A(bM_i, 0, c) = \text{ord}_{A/\mathfrak q_i}(c^{m_ie_i})
$$
ce qui résulte du Lemme \ref{chow-lemma-length-multiplication}
par la description donnée ci-dessus de ce qui se produit
après localisation en $\mathfrak q_i$.
\end{proof}

\begin{lemma}[Lemme clé]
\label{chow-lemma-milnor-gersten-low-degree}
\begin{reference}
Lorsque $A$ est un anneau excellent, il s'agit de \cite[Proposition 1]{Kato-Milnor-K}.
\end{reference}
Soit $A$ un anneau local noethérien intègre de dimension $2$, de corps des fractions $K$.
Soient $f, g \in K^*$.
Soient $\mathfrak q_1, \ldots, \mathfrak q_t$ les idéaux premiers de hauteur
$1$, notés $\mathfrak q$, de $A$ tels que $f$ ou $g$ n'appartienne pas à
$A^*_{\mathfrak q}$.
Alors
$$
\sum\nolimits_{i = 1, \ldots, t}
\text{ord}_{A/\mathfrak q_i}(\partial_{A_{\mathfrak q_i}}(f, g))
=
0
$$
On peut aussi l'écrire
$$
\sum\nolimits_{\text{hauteur}(\mathfrak q) = 1}
\text{ord}_{A/\mathfrak q}(\partial_{A_{\mathfrak q}}(f, g))
=
0
$$
car, en tout idéal premier de hauteur $1$, noté $\mathfrak q$,
de $A$ tel que $f, g \in A^*_{\mathfrak q}$,
nous avons $\partial_{A_{\mathfrak q}}(f, g) = 1$.
\end{lemma}

\begin{proof}
Comme les symboles modérés $\partial_{A_{\mathfrak q}}(f, g)$ sont
bilinéaires et les fonctions ordre $\text{ord}_{A/\mathfrak q}$
additives, il suffit de démontrer la formule lorsque
$f$ et $g$ sont des éléments de $A$. Ce cas est démontré au
Lemme \ref{chow-lemma-key-nonzerodivisors}.
\end{proof}

\begin{remark}[K-théorie de Milnor]
\label{chow-remark-gersten-complex-milnor}
Pour un corps $k$, notons $K^M_*(k)$ le quotient de
l'algèbre tensorielle sur $k^*$ par l'idéal bilatère
engendré par les éléments $x \otimes 1 - x$ pour $x \in k \setminus \{0, 1\}$.
Ainsi $K^M_0(k) = \mathbf{Z}$, $K_1^M(k) = k^*$, et
$$
K^M_2(k) = k^* \otimes_\mathbf{Z} k^* / \langle x \otimes 1 - x \rangle
$$
Si $A$ est un anneau de valuation discrète de corps des fractions $F = \text{Frac}(A)$
et de corps résiduel $\kappa$, il existe un symbole modéré
$$
\partial_A : K_{i + 1}^M(F) \to K_i^M(\kappa)
$$
défini comme à la section \ref{chow-section-tame-symbol} ; voir \cite{Kato-Milnor-K}.
Plus généralement, cette application peut être étendue au cas où $A$ est un
anneau local intègre excellent de dimension $1$, à l'aide de la normalisation et des applications
normes sur $K_i^M$, voir \cite{Kato-Milnor-K} ; vraisemblablement, la méthode de la
section \ref{chow-section-tame-symbol} permet d'étendre la construction
du symbole modéré $\partial_A$ à tout anneau local noethérien intègre $A$
de dimension $1$. Ensuite, soit $X$ un schéma noethérien muni d'une
fonction de dimension $\delta$. On peut alors utiliser ces symboles modérés pour obtenir
les flèches suivantes :
$$
\bigoplus\nolimits_{\delta(x) = j + 1} K^M_{i + 1}(\kappa(x))
\longrightarrow
\bigoplus\nolimits_{\delta(x) = j} K^M_i(\kappa(x))
\longrightarrow
\bigoplus\nolimits_{\delta(x) = j - 1} K^M_{i - 1}(\kappa(x))
$$
Cependant, il n'est pas clair que la composée soit nulle, c'est-à-dire que
l'on obtienne un complexe de groupes abéliens.
Pour $X$ excellent, cela est démontré dans \cite{Kato-Milnor-K}.
Lorsque $i = 1$ et $j$ est quelconque, cela résulte du
Lemme \ref{chow-lemma-milnor-gersten-low-degree}.
\end{remark}





















\section{Cadre}
\label{chow-section-setup}

\noindent
Nous travaillerons partout sur une base localement noethérienne et universellement
caténaire $S$, munie d'une fonction de dimension $\delta$.
Bien qu'il soit probablement possible de généraliser certaines parties de la
discussion de ce chapitre, cela semble constituer une bonne première
approximation. C'est exactement le degré de généralité étudié dans \cite{Thorup}.
Nous ne supposons généralement pas nos schémas
séparés ni quasi-compacts. De nombreux champs algébriques intéressants
ne sont ni séparés ni quasi-compacts, et c'est un bon
cas d'étude pour voir comment élaborer aussi pour eux une théorie raisonnable.
Afin de pouvoir faire référence à ces hypothèses, nous les numérotons.

\begin{situation}
\label{chow-situation-setup}
Ici, $S$ est un schéma localement noethérien et universellement caténaire.
De plus, nous supposons $S$ muni d'une fonction de dimension
$\delta : S \longrightarrow \mathbf{Z}$.
\end{situation}

\noindent
Voir Morphismes, Définition \ref{morphisms-definition-universally-catenary},
pour la notion de schéma universellement caténaire, et
Topologie, Définition \ref{topology-definition-dimension-function},
pour la notion de fonction de dimension. Rappelons que tout schéma localement
noethérien et caténaire possède localement une fonction de dimension, voir
Propriétés, Lemme \ref{properties-lemma-catenary-dimension-function}.
De plus, beaucoup de schémas sont universellement caténaires,
voir Morphismes, Lemme \ref{morphisms-lemma-ubiquity-uc}.

\medskip\noindent
Soit $(S, \delta)$ comme dans la Situation \ref{chow-situation-setup}.
Tout schéma $X$ localement de type fini sur $S$ est localement noethérien
et caténaire. En fait, $X$ possède une fonction de dimension canonique
$$
\delta = \delta_{X/S} : X \longrightarrow \mathbf{Z}
$$
associée à $(f : X \to S, \delta)$ et donnée par la règle
$\delta_{X/S}(x) = \delta(f(x)) + \text{trdeg}_{\kappa(f(x))}\kappa(x)$.
Voir Morphismes, Lemme \ref{morphisms-lemma-dimension-function-propagates}.
De plus, si $h : X \to Y$ est un morphisme de schémas localement de type
fini sur $S$, et si $x \in X$, $y = h(x)$,
alors, de manière évidente,
$\delta_{X/S}(x) = \delta_{Y/S}(y) + \text{trdeg}_{\kappa(y)}\kappa(x)$.
Nous utiliserons librement cette fonction et ses propriétés dans la suite.

\medskip\noindent
Voici les exemples fondamentaux de cadres comme ci-dessus.
En fait, le principal intérêt réside dans le cas où la base
est le spectre d'un corps, ou dans celui où la base
est le spectre d'un anneau de Dedekind (par exemple $\mathbf{Z}$
ou un anneau de valuation discrète).

\begin{example}
\label{chow-example-field}
Ici, $S = \Spec(k)$ et $k$ est un corps.
On pose $\delta(pt) = 0$, où $pt$ désigne l'unique point de $S$.
Le couple $(S, \delta)$ est un exemple de situation comme dans la
Situation \ref{chow-situation-setup}, d'après
Morphismes, Lemme \ref{morphisms-lemma-ubiquity-uc}.
\end{example}

\begin{example}
\label{chow-example-domain-dimension-1}
Ici, $S = \Spec(A)$, où $A$ est un anneau noethérien intègre
de dimension $1$.
On peut par exemple considérer $A = \mathbf{Z}$.
On pose $\delta(\mathfrak p) = 0$ si
$\mathfrak p$ est un idéal maximal, et $\delta(\mathfrak p) = 1$
si $\mathfrak p = (0)$ correspond au point générique.
C'est un exemple de Situation \ref{chow-situation-setup}, d'après
Morphismes, Lemme \ref{morphisms-lemma-ubiquity-uc}.
\end{example}

\begin{example}
\label{chow-example-CM-irreducible}
Ici, $S$ est un schéma de Cohen--Macaulay. Alors $S$ est universellement caténaire d'après
Morphismes, Lemme \ref{morphisms-lemma-ubiquity-uc}.
On pose $\delta(s) = -\dim(\mathcal{O}_{S, s})$.
Si $s' \leadsto s$ est une spécialisation non triviale de points de $S$,
alors $\mathcal{O}_{S, s'}$ est le localisé de $\mathcal{O}_{S, s}$
en un idéal premier non maximal $\mathfrak p \subset \mathcal{O}_{S, s}$, voir
Schémas, Lemme \ref{schemes-lemma-specialize-points}.
Ainsi $\dim(\mathcal{O}_{S, s}) =
\dim(\mathcal{O}_{S, s'}) + \dim(\mathcal{O}_{S, s}/\mathfrak p) >
\dim(\mathcal{O}_{S, s'})$ d'après
Algèbre, Lemme \ref{algebra-lemma-CM-dim-formula}.
Par conséquent, $\delta(s') > \delta(s)$. Si $s' \leadsto s$ est
une spécialisation immédiate, alors il n'existe aucun idéal
premier strictement compris entre $\mathfrak p$ et $\mathfrak m_s$,
et l'on obtient $\delta(s') = \delta(s) + 1$. Ainsi $\delta$
est une fonction de dimension. Autrement dit, le couple $(S, \delta)$
est un exemple de Situation \ref{chow-situation-setup}.
\end{example}

\noindent
Si $S$ est jacobsonien et si $\delta$ envoie les points fermés sur zéro, alors $\delta$
est la fonction qui associe à un point la dimension de son adhérence.

\begin{lemma}
\label{chow-lemma-delta-is-dimension}
Soit $(S, \delta)$ comme dans la Situation \ref{chow-situation-setup}.
Supposons en outre que $S$ soit un schéma jacobsonien et que $\delta(s) = 0$ pour tout
point fermé $s$ de $S$. Soit $X$ localement de type fini sur $S$.
Soit $Z \subset X$ un sous-schéma fermé intègre et soit
$\xi \in Z$ son point générique. Les entiers suivants sont égaux :
\begin{enumerate}
\item $\delta_{X/S}(\xi)$,
\item $\dim(Z)$, et
\item $\dim(\mathcal{O}_{Z, z})$, où $z$ est un point fermé de $Z$.
\end{enumerate}
\end{lemma}

\begin{proof}
Soient $X \to S$ et $\xi \in Z \subset X$ comme dans le lemme.
Comme $X$ est localement de type fini sur $S$, on voit que
$X$ est jacobsonien, voir
Morphismes, Lemme \ref{morphisms-lemma-Jacobson-universally-Jacobson}.
Par conséquent, les points fermés de $X$ sont denses dans tout fermé de $Z$
et s'envoient sur des points fermés de $S$. Ainsi, toute chaîne
de fermés irréductibles de $Z$ peut être terminée par un point fermé de $Z$.
Il s'ensuit que $\dim(Z) = \sup_z(\dim(\mathcal{O}_{Z, z})$
(voir Propriétés, Lemme \ref{properties-lemma-codimension-local-ring}),
où $z \in Z$ parcourt les points fermés de $Z$.
Notons que $\dim(\mathcal{O}_{Z, z}) = \delta(\xi) - \delta(z)$
d'après les propriétés d'une fonction de dimension.
Pour tout point fermé $z \in Z$, l'extension de corps
$\kappa(z)/\kappa(f(z))$ est finie, voir Morphismes,
Lemme \ref{morphisms-lemma-jacobson-finite-type-points}.
Ainsi $\delta_{X/S}(z) = \delta(f(z)) = 0$ pour $z \in Z$ fermé.
Il s'ensuit que les trois entiers sont égaux.
\end{proof}

\noindent
Dans la situation du lemme précédent, la
valeur de $\delta$ au point générique d'un fermé irréductible
est la dimension de ce fermé irréductible.
Cependant, on ne peut s'attendre à ce que cette égalité vaille en général.
Par exemple, si $S = \Spec(\mathbf{C}[[t]])$ et
$X = \Spec(\mathbf{C}((t)))$, on obtiendrait alors
$\delta(x) = 1$ pour l'unique point de $X$, mais $\dim(X) = 0$.
Nous voulons néanmoins considérer $\delta_{X/S}$ comme donnant la
dimension des sous-schémas fermés irréductibles. Nous introduisons donc
la terminologie suivante.

\begin{definition}
\label{chow-definition-delta-dimension}
Soit $(S, \delta)$ comme dans la Situation \ref{chow-situation-setup}.
Pour tout schéma $X$ localement de type fini sur $S$
et tout fermé irréductible $Z \subset X$, nous définissons
$$
\dim_\delta(Z) = \delta(\xi)
$$
où $\xi \in Z$ est le point générique de $Z$.
Nous appellerons cet entier la {\it $\delta$-dimension de $Z$}.
Si $Z$ est un sous-schéma fermé de $X$, nous définissons
$\dim_\delta(Z)$ comme la borne supérieure des $\delta$-dimensions
de ses composantes irréductibles.
\end{definition}







\section{Cycles}
\label{chow-section-cycles}

\noindent
Comme nous ne supposons pas nos schémas quasi-compacts, il faut
être un peu prudent dans la définition des cycles. Nous devons autoriser
les sommes infinies, car une fonction rationnelle peut, par exemple, avoir une infinité de
pôles. Quoi qu'il en soit, si $X$ est quasi-compact, un
cycle est une somme finie, comme d'habitude.

\begin{definition}
\label{chow-definition-cycles}
Soit $(S, \delta)$ comme dans la Situation \ref{chow-situation-setup}.
Soit $X$ localement de type fini sur $S$.
Soit $k \in \mathbf{Z}$.
\begin{enumerate}
\item Un {\it cycle sur $X$} est une somme formelle
$$
\alpha = \sum n_Z [Z]
$$
où la somme porte sur les sous-schémas fermés intègres $Z \subset X$,
où chaque $n_Z \in \mathbf{Z}$, et où la collection
$\{Z; n_Z \not = 0\}$ est localement finie
(Topologie, Définition \ref{topology-definition-locally-finite}).
\item Un {\it $k$-cycle} sur $X$ est un cycle
$$
\alpha = \sum n_Z [Z]
$$
tel que $n_Z \not = 0 \Rightarrow \dim_\delta(Z) = k$.
\item Le groupe abélien de tous les $k$-cycles sur $X$ est noté $Z_k(X)$.
\end{enumerate}
\end{definition}

\noindent
Autrement dit, un $k$-cycle sur $X$
est une combinaison $\mathbf{Z}$-linéaire formelle et localement finie
de sous-schémas fermés intègres de $\delta$-dimension $k$.
L'addition des $k$-cycles $\alpha = \sum n_Z[Z]$ et
$\beta = \sum m_Z[Z]$ est donnée par
$$
\alpha + \beta = \sum (n_Z + m_Z)[Z],
$$
c'est-à-dire par addition des coefficients.

\begin{remark}
\label{chow-remark-cycles-pointwise}
Soit $(S, \delta)$ comme dans la Situation \ref{chow-situation-setup}.
Soit $X$ localement de type fini sur $S$. Soit $k \in \mathbf{Z}$.
Alors on peut écrire
$$
Z_k(X) = \bigoplus\nolimits_{\delta(x) = k}' K_0^M(\kappa(x))
\quad\subset\quad
\bigoplus\nolimits_{\delta(x) = k} K_0^M(\kappa(x))
$$
avec les notations et conventions suivantes :
\begin{enumerate}
\item $K_0^M(\kappa(x)) = \mathbf{Z}$ est la partie de degré $0$ de
la K-théorie de Milnor du corps résiduel $\kappa(x)$ du point
$x \in X$ (voir la Remarque \ref{chow-remark-gersten-complex-milnor}), et
\item la somme directe de droite porte sur tous les points $x \in X$
tels que $\delta(x) = k$,
\item la notation $\bigoplus'_x$ signifie que nous considérons le
sous-groupe constitué des éléments localement finis, à savoir des éléments
$\sum_x n_x$ tels que, pour tout ouvert quasi-compact $U \subset X$,
l'ensemble des $x \in U$ tels que $n_x \not = 0$ soit fini.
\end{enumerate}
\end{remark}

\begin{definition}
\label{chow-definition-support-cycle}
Soit $(S, \delta)$ comme dans la Situation \ref{chow-situation-setup}.
Soit $X$ localement de type fini sur $S$.
Le {\it support} d'un cycle $\alpha = \sum n_Z [Z]$ sur $X$
est
$$
\text{Supp}(\alpha) = \bigcup\nolimits_{n_Z \not = 0} Z \subset X
$$
\end{definition}

\noindent
Comme la collection $\{Z; n_Z \not = 0\}$ est localement finie,
on voit que $\text{Supp}(\alpha)$ est un fermé de $X$.
Si $\alpha$ est un $k$-cycle, toute composante irréductible $Z$
de $\text{Supp}(\alpha)$ est de $\delta$-dimension $k$.

\begin{definition}
\label{chow-definition-effective-cycle}
Soit $(S, \delta)$ comme dans la Situation \ref{chow-situation-setup}.
Soit $X$ localement de type fini sur $S$.
Un cycle $\alpha$ sur $X$ est {\it effectif} s'il
peut s'écrire $\alpha =\sum n_Z [Z]$ avec $n_Z \geq 0$ pour tout $Z$.
\end{definition}

\noindent
L'ensemble de tous les cycles effectifs est un monoïde, car la somme de
deux cycles effectifs est effective, mais ce n'est pas un groupe
(sauf si $X = \emptyset$).




\section{Cycle associé à un sous-schéma fermé}
\label{chow-section-cycle-of-closed-subscheme}

\begin{lemma}
\label{chow-lemma-multiplicity-finite}
Soit $(S, \delta)$ comme dans la Situation \ref{chow-situation-setup}.
Soit $X$ localement de type fini sur $S$.
Soit $Z \subset X$ un sous-schéma fermé.
\begin{enumerate}
\item Soit $Z' \subset Z$ une composante irréductible et
soit $\xi \in Z'$ son point générique.
Alors
$$
\text{longueur}_{\mathcal{O}_{X, \xi}} \mathcal{O}_{Z, \xi} < \infty
$$
\item Si $\dim_\delta(Z) \leq k$ et si $\xi \in Z$ vérifie
$\delta(\xi) = k$, alors $\xi$ est le point générique d'une
composante irréductible de $Z$.
\end{enumerate}
\end{lemma}

\begin{proof}
Soient $Z' \subset Z$ et $\xi \in Z'$ comme dans (1).
Alors $\dim(\mathcal{O}_{Z, \xi}) = 0$ (par exemple d'après
Propriétés, Lemme \ref{properties-lemma-codimension-local-ring}).
Ainsi $\mathcal{O}_{Z, \xi}$ est un anneau noethérien
local de dimension zéro, donc de longueur finie sur
lui-même (voir
Algèbre, Proposition \ref{algebra-proposition-dimension-zero-ring}).
Il est donc aussi de longueur finie sur $\mathcal{O}_{X, \xi}$, voir
Algèbre, Lemme \ref{algebra-lemma-length-independent}.

\medskip\noindent
Supposons $\xi \in Z$ et $\delta(\xi) = k$.
Considérons l'adhérence $Z' = \overline{\{\xi\}}$. C'est un sous-schéma
fermé irréductible tel que $\dim_\delta(Z') = k$ par définition.
Comme $\dim_\delta(Z) = k$, ce doit être une composante irréductible
de $Z$. Cela démontre (2).
\end{proof}

\begin{definition}
\label{chow-definition-cycle-associated-to-closed-subscheme}
Soit $(S, \delta)$ comme dans la Situation \ref{chow-situation-setup}.
Soit $X$ localement de type fini sur $S$.
Soit $Z \subset X$ un sous-schéma fermé.
\begin{enumerate}
\item Pour toute composante irréductible $Z' \subset Z$ de point générique $\xi$,
l'entier
$m_{Z', Z} = \text{longueur}_{\mathcal{O}_{X, \xi}} \mathcal{O}_{Z, \xi}$
(Lemme \ref{chow-lemma-multiplicity-finite})
est appelé la {\it multiplicité de $Z'$ dans $Z$}.
\item Supposons $\dim_\delta(Z) \leq k$.
Le {\it $k$-cycle associé à $Z$} est
$$
[Z]_k
=
\sum m_{Z', Z}[Z']
$$
où la somme porte sur les composantes irréductibles de $Z$
de $\delta$-dimension $k$. (C'est un $k$-cycle d'après
Diviseurs, Lemme \ref{divisors-lemma-components-locally-finite}.)
\end{enumerate}
\end{definition}

\noindent
Il importe de noter que nous ne définissons $[Z]_k$ que si la $\delta$-dimension
de $Z$ ne dépasse pas $k$. Autrement dit, par convention, écrire
$[Z]_k$ implique que $\dim_\delta(Z) \leq k$.



\section{Cycle associé à un faisceau cohérent}
\label{chow-section-cycle-of-coherent-sheaf}



\begin{lemma}
\label{chow-lemma-length-finite}
Soit $(S, \delta)$ comme dans la Situation \ref{chow-situation-setup}.
Soit $X$ localement de type fini sur $S$.
Soit $\mathcal{F}$ un $\mathcal{O}_X$-module cohérent.
\begin{enumerate}
\item La collection des composantes irréductibles du support de
$\mathcal{F}$ est localement finie.
\item Soit $Z' \subset \text{Supp}(\mathcal{F})$
une composante irréductible, et
soit $\xi \in Z'$ son point générique.
Alors
$$
\text{longueur}_{\mathcal{O}_{X, \xi}} \mathcal{F}_\xi < \infty
$$
\item Si $\dim_\delta(\text{Supp}(\mathcal{F})) \leq k$
et si $\xi \in \text{Supp}(\mathcal{F})$ vérifie $\delta(\xi) = k$, alors $\xi$ est un
point générique d'une composante irréductible de $\text{Supp}(\mathcal{F})$.
\end{enumerate}
\end{lemma}

\begin{proof}
D'après Cohomologie des schémas, Lemme \ref{coherent-lemma-coherent-support-closed},
le support $Z$ de $\mathcal{F}$ est un fermé de $X$.
Nous pouvons considérer $Z$ comme un sous-schéma fermé réduit de $X$
(Schémas, Lemme \ref{schemes-lemma-reduced-closed-subscheme}).
Ainsi (1) résulte de
Diviseurs, Lemme \ref{divisors-lemma-components-locally-finite}, appliqué à $Z$,
et (3) résulte du
Lemme \ref{chow-lemma-multiplicity-finite}, appliqué à $Z$.

\medskip\noindent
Soit $\xi \in Z'$ comme dans (2). Dans ce cas, pour toute spécialisation
$\xi' \leadsto \xi$ dans $X$, nous avons $\mathcal{F}_{\xi'} = 0$.
Rappelons que les idéaux premiers non maximaux de $\mathcal{O}_{X, \xi}$ correspondent
aux points de $X$ qui se spécialisent en $\xi$
(Schémas, Lemme \ref{schemes-lemma-specialize-points}).
Ainsi $\mathcal{F}_\xi$ est un $\mathcal{O}_{X, \xi}$-module fini
de support $\{\mathfrak m_\xi\}$. Il est donc de longueur finie
d'après Algèbre, Lemme \ref{algebra-lemma-support-point}.
\end{proof}

\begin{definition}
\label{chow-definition-cycle-associated-to-coherent-sheaf}
Soit $(S, \delta)$ comme dans la Situation \ref{chow-situation-setup}.
Soit $X$ localement de type fini sur $S$.
Soit $\mathcal{F}$ un $\mathcal{O}_X$-module cohérent.
\begin{enumerate}
\item Pour toute composante irréductible $Z' \subset \text{Supp}(\mathcal{F})$
de point générique $\xi$, l'entier
$m_{Z', \mathcal{F}} = \text{longueur}_{\mathcal{O}_{X, \xi}} \mathcal{F}_\xi$
(Lemme \ref{chow-lemma-length-finite})
est appelé la {\it multiplicité de $Z'$ dans $\mathcal{F}$}.
\item Supposons $\dim_\delta(\text{Supp}(\mathcal{F})) \leq k$.
Le {\it $k$-cycle associé à $\mathcal{F}$} est
$$
[\mathcal{F}]_k
=
\sum m_{Z', \mathcal{F}}[Z']
$$
où la somme porte sur les composantes irréductibles de
$\text{Supp}(\mathcal{F})$ de $\delta$-dimension $k$.
(C'est un $k$-cycle d'après le Lemme \ref{chow-lemma-length-finite}.)
\end{enumerate}
\end{definition}

\noindent
Il importe de noter que nous ne définissons $[\mathcal{F}]_k$
que si $\mathcal{F}$ est cohérent et si la $\delta$-dimension
de $\text{Supp}(\mathcal{F})$ ne dépasse pas $k$. Autrement dit,
par convention, écrire $[\mathcal{F}]_k$ implique que
$\mathcal{F}$ est cohérent sur $X$ et que
$\dim_\delta(\text{Supp}(\mathcal{F})) \leq k$.

\begin{lemma}
\label{chow-lemma-cycle-closed-coherent}
Soit $(S, \delta)$ comme dans la Situation \ref{chow-situation-setup}.
Soit $X$ localement de type fini sur $S$.
Soit $Z \subset X$ un sous-schéma fermé.
Si $\dim_\delta(Z) \leq k$, alors $[Z]_k = [{\mathcal O}_Z]_k$.
\end{lemma}

\begin{proof}
Cela vient de ce que, dans ce cas, les multiplicités $m_{Z', Z}$ et
$m_{Z', \mathcal{O}_Z}$ coïncident par définition.
\end{proof}

\begin{lemma}
\label{chow-lemma-additivity-sheaf-cycle}
Soit $(S, \delta)$ comme dans la Situation \ref{chow-situation-setup}.
Soit $X$ localement de type fini sur $S$.
Soit $0 \to \mathcal{F} \to \mathcal{G} \to \mathcal{H} \to 0$
une suite exacte courte de faisceaux cohérents sur $X$.
Supposons que la $\delta$-dimension des supports
de $\mathcal{F}$, $\mathcal{G}$ et $\mathcal{H}$ soit $\leq k$.
Alors $[\mathcal{G}]_k = [\mathcal{F}]_k + [\mathcal{H}]_k$.
\end{lemma}

\begin{proof}
Cela résulte immédiatement de l'additivité des longueurs, voir
Algèbre, Lemme \ref{algebra-lemma-length-additive}.
\end{proof}









\section{Préparation à l'image directe propre}
\label{chow-section-preparation-pushforward}

\begin{lemma}
\label{chow-lemma-equal-dimension}
Soit $(S, \delta)$ comme dans la Situation \ref{chow-situation-setup}.
Soient $X$ et $Y$ localement de type fini sur $S$.
Soit $f : X \to Y$ un morphisme.
Supposons $X$ et $Y$ intègres et $\dim_\delta(X) = \dim_\delta(Y)$.
Alors ou bien $f(X)$ est contenu dans un sous-schéma fermé strict
de $Y$, ou bien $f$ est dominant et l'extension de corps de fonctions
$R(X)/R(Y)$ est finie.
\end{lemma}

\begin{proof}
L'adhérence $\overline{f(X)} \subset Y$ est irréductible, puisque $X$
est irréductible (Topologie, Lemmes
\ref{topology-lemma-image-irreducible-space} et
\ref{topology-lemma-irreducible}).
Si $\overline{f(X)} \not = Y$, nous avons terminé.
Si $\overline{f(X)} = Y$, alors $f$ est dominant et, d'après
Morphismes,
Lemme \ref{morphisms-lemma-dominant-finite-number-irreducible-components},
le point générique $\eta_Y$ de $Y$ appartient à l'image de $f$.
Cela implique bien sûr que $f(\eta_X) = \eta_Y$, où $\eta_X \in X$
est le point générique de $X$. Comme $\delta(\eta_X) = \delta(\eta_Y)$,
on voit que $R(Y) = \kappa(\eta_Y) \subset \kappa(\eta_X) = R(X)$
est une extension de degré de transcendance $0$.
Par conséquent, $R(Y) \subset R(X)$ est une extension finie d'après
Morphismes, Lemme \ref{morphisms-lemma-finite-degree}
(qui s'applique en vertu de
Morphismes, Lemme \ref{morphisms-lemma-permanence-finite-type}).
\end{proof}

\begin{lemma}
\label{chow-lemma-quasi-compact-locally-finite}
Soit $(S, \delta)$ comme dans la Situation \ref{chow-situation-setup}.
Soient $X$ et $Y$ localement de type fini sur $S$.
Soit $f : X \to Y$ un morphisme.
Supposons $f$ quasi-compact et $\{Z_i\}_{i \in I}$ une collection localement
finie de fermés de $X$.
Alors $\{\overline{f(Z_i)}\}_{i \in I}$ est une collection localement finie
de fermés de $Y$.
\end{lemma}

\begin{proof}
Soit $V \subset Y$ un ouvert quasi-compact.
Puisque $f$ est quasi-compact, l'ouvert $f^{-1}(V)$ est
quasi-compact. Ainsi l'ensemble
$\{i \in I \mid Z_i \cap f^{-1}(V) \not = \emptyset \}$
est fini par un argument topologique simple que nous omettons.
Comme il s'agit du même ensemble que
$$
\{i \in I \mid f(Z_i) \cap V \not = \emptyset \} =
\{i \in I \mid \overline{f(Z_i)} \cap V \not = \emptyset \}
$$
le lemme est démontré.
\end{proof}









\section{Image directe propre}
\label{chow-section-proper-pushforward}

\begin{definition}
\label{chow-definition-proper-pushforward}
Soit $(S, \delta)$ comme dans la Situation \ref{chow-situation-setup}.
Soient $X$ et $Y$ localement de type fini sur $S$.
Soit $f : X \to Y$ un morphisme.
Supposons $f$ propre.
\begin{enumerate}
\item Soit $Z \subset X$ un sous-schéma fermé intègre
tel que $\dim_\delta(Z) = k$. Nous définissons
$$
f_*[Z] =
\left\{
\begin{matrix}
0 & \text{si} & \dim_\delta(f(Z))< k, \\
\deg(Z/f(Z)) [f(Z)] & \text{si} & \dim_\delta(f(Z)) = k.
\end{matrix}
\right.
$$
Ici, nous considérons $f(Z) \subset Y$ comme un sous-schéma fermé intègre.
Le degré de $Z$ sur $f(Z)$ est fini si
$\dim_\delta(f(Z)) = \dim_\delta(Z)$,
d'après le Lemme \ref{chow-lemma-equal-dimension}.
\item Soit $\alpha = \sum n_Z [Z]$ un $k$-cycle sur $X$. Nous appelons
{\it image directe} de $\alpha$ la somme
$$
f_* \alpha = \sum n_Z f_*[Z]
$$
où chaque $f_*[Z]$ est défini comme ci-dessus. La somme est localement finie
d'après le Lemme \ref{chow-lemma-quasi-compact-locally-finite} ci-dessus.
\end{enumerate}
\end{definition}

\noindent
Par définition, l'image directe propre des cycles
$$
f_* : Z_k(X) \longrightarrow Z_k(Y)
$$
est un homomorphisme de groupes abéliens. Elle fait de $X \mapsto Z_k(X)$
un foncteur covariant sur la catégorie des schémas localement
de type fini sur $S$, dont les morphismes sont les morphismes propres.

\begin{lemma}
\label{chow-lemma-compose-pushforward}
Soit $(S, \delta)$ comme dans la Situation \ref{chow-situation-setup}.
Soient $X$, $Y$ et $Z$ localement de type fini sur $S$.
Soient $f : X \to Y$ et $g : Y \to Z$ des morphismes propres.
Alors $g_* \circ f_* = (g \circ f)_*$ comme applications $Z_k(X) \to Z_k(Z)$.
\end{lemma}

\begin{proof}
Soit $W \subset X$ un sous-schéma fermé intègre de dimension $k$.
Considérons $W' = f(W) \subset Y$ et $W'' = g(f(W)) \subset Z$.
Comme $f$ et $g$ sont propres, $W'$ (respectivement $W''$) est
un sous-schéma fermé intègre de $Y$ (respectivement de $Z$).
Il faut montrer que $g_*(f_*[W]) = (g \circ f)_*[W]$.
Si $\dim_\delta(W'') < k$, les deux membres sont nuls.
Si $\dim_\delta(W'') = k$, les morphismes induits
$$
W \longrightarrow
W' \longrightarrow
W''
$$
satisfont tous deux les hypothèses du Lemme \ref{chow-lemma-equal-dimension}. Ainsi
$$
g_*(f_*[W]) = \deg(W/W')\deg(W'/W'')[W''],
\quad
(g \circ f)_*[W] = \deg(W/W'')[W''].
$$
On peut alors appliquer
Morphismes, Lemme \ref{morphisms-lemma-degree-composition},
pour conclure.
\end{proof}

\noindent
Une immersion fermée est propre. Si $i : Z \to X$ est une immersion fermée,
alors les applications
$$
i_* : Z_k(Z) \longrightarrow Z_k(X)
$$
sont toutes {\it injectives}.

\begin{lemma}
\label{chow-lemma-exact-sequence-closed}
Soit $(S, \delta)$ comme dans la Situation \ref{chow-situation-setup}.
Soit $X$ localement de type fini sur $S$. Soient $X_1, X_2 \subset X$
des sous-schémas fermés tels que $X = X_1 \cup X_2$ ensemblistement.
Pour tout $k \in \mathbf{Z}$, la suite de groupes abéliens
$$
\xymatrix{
Z_k(X_1 \cap X_2) \ar[r] &
Z_k(X_1) \oplus Z_k(X_2) \ar[r] &
Z_k(X) \ar[r] &
0
}
$$
est exacte. Ici, $X_1 \cap X_2$ est l'intersection schématique et
les applications sont les applications d'image directe, l'une étant multipliée par $-1$.
\end{lemma}

\begin{proof}
Supposons d'abord $X$ quasi-compact. Alors $Z_k(X)$ est un $\mathbf{Z}$-module libre
ayant pour base les éléments $[Z]$, où $Z \subset X$ est fermé intègre
de $\delta$-dimension $k$. Les groupes
$Z_k(X_1)$, $Z_k(X_2)$ et $Z_k(X_1 \cap X_2)$ sont libres sur le sous-ensemble des
$Z$ tels que $Z \subset X_1$, $Z \subset X_2$, $Z \subset X_1 \cap X_2$.
Cela démontre immédiatement le lemme dans ce cas. Le cas général est analogue
et la démonstration est omise.
\end{proof}

\begin{lemma}
\label{chow-lemma-cycle-push-sheaf}
Soit $(S, \delta)$ comme dans la Situation \ref{chow-situation-setup}.
Soit $f : X \to Y$ un morphisme propre de schémas qui sont
localement de type fini sur $S$.
\begin{enumerate}
\item Soit $Z \subset X$ un sous-schéma fermé tel que $\dim_\delta(Z) \leq k$.
Alors
$$
f_*[Z]_k = [f_*{\mathcal O}_Z]_k.
$$
\item Soit $\mathcal{F}$ un faisceau cohérent sur $X$ tel que
$\dim_\delta(\text{Supp}(\mathcal{F})) \leq k$. Alors
$$
f_*[\mathcal{F}]_k = [f_*{\mathcal F}]_k.
$$
\end{enumerate}
Notons que l'énoncé a un sens puisque $f_*\mathcal{F}$ et
$f_*\mathcal{O}_Z$ sont des $\mathcal{O}_Y$-modules cohérents d'après
Cohomologie des schémas, Proposition
\ref{coherent-proposition-proper-pushforward-coherent}.
\end{lemma}

\begin{proof}
La partie (1) résulte de (2) et du Lemme \ref{chow-lemma-cycle-closed-coherent}.
Soit $\mathcal{F}$ un faisceau cohérent sur $X$.
Supposons que $\dim_\delta(\text{Supp}(\mathcal{F})) \leq k$.
D'après Cohomologie des schémas, Lemme \ref{coherent-lemma-coherent-support-closed},
il existe un sous-schéma fermé $i : Z \to X$ et un
$\mathcal{O}_Z$-module cohérent $\mathcal{G}$ tels que
$i_*\mathcal{G} \cong \mathcal{F}$ et que le support
de $\mathcal{F}$ soit $Z$. Soit $Z' \subset Y$ l'image schématique
de $f|_Z : Z \to Y$. Considérons le diagramme commutatif de schémas
$$
\xymatrix{
Z \ar[r]_i \ar[d]_{f|_Z} &
X \ar[d]^f \\
Z' \ar[r]^{i'} & Y
}
$$
Nous avons $f_*\mathcal{F} = f_*i_*\mathcal{G} = i'_*(f|_Z)_*\mathcal{G}$
en parcourant le diagramme des deux façons. Supposons le résultat établi
pour les immersions fermées et pour $f|_Z$. Alors
$$
f_*[\mathcal{F}]_k = f_*i_*[\mathcal{G}]_k
= (i')_*(f|_Z)_*[\mathcal{G}]_k =
(i')_*[(f|_Z)_*\mathcal{G}]_k =
[(i')_*(f|_Z)_*\mathcal{G}]_k = [f_*\mathcal{F}]_k
$$
comme voulu. Le cas d'une immersion fermée est immédiat (et omis).
Notons que $f|_Z : Z \to Z'$ est un morphisme dominant (voir
Morphismes, Lemme \ref{morphisms-lemma-quasi-compact-scheme-theoretic-image}).
Nous nous sommes donc ramenés au cas où
$\dim_\delta(X) \leq k$ et $f : X \to Y$ est propre et dominant.

\medskip\noindent
Supposons $\dim_\delta(X) \leq k$ et $f : X \to Y$ propre et dominant.
Puisque $f$ est dominant, pour toute composante irréductible $Z \subset Y$
de point générique $\eta$, il existe un point $\xi \in X$ tel
que $f(\xi) = \eta$. Ainsi $\delta(\eta) \leq \delta(\xi) \leq k$.
On voit donc que, dans les expressions
$$
f_*[\mathcal{F}]_k = \sum n_Z[Z],
\quad
\text{et}
\quad
[f_*\mathcal{F}]_k = \sum m_Z[Z].
$$
dès que $n_Z \not = 0$ ou que $m_Z \not = 0$, le sous-schéma fermé
intègre $Z$ est en fait une composante irréductible de $Y$ de
$\delta$-dimension $k$. Choisissons un tel sous-schéma fermé intègre
$Z \subset Y$ et notons $\eta$ son point générique. Notons que, pour
tout $\xi \in X$ tel que $f(\xi) = \eta$, nous avons $\delta(\xi) \geq k$,
et que $\xi$ est donc lui aussi le point générique d'une composante irréductible
de $X$ de $\delta$-dimension $k$
(voir le Lemme \ref{chow-lemma-multiplicity-finite}). Comme $f$ est quasi-compact
et $X$ localement noethérien, il ne peut y en avoir qu'un nombre fini,
et $f^{-1}(\{\eta\})$ est donc fini.
D'après Morphismes, Lemme \ref{morphisms-lemma-generically-finite}, il existe
un voisinage ouvert $\eta \in V \subset Y$ tel que $f^{-1}(V) \to V$
soit fini. En remplaçant $Y$ par $V$ et $X$ par $f^{-1}(V)$, on se ramène au
cas où $Y$ est affine et où $f$ est fini.

\medskip\noindent
Écrivons $Y = \Spec(R)$ et $X = \Spec(A)$ (ce qui est possible car
un morphisme fini est affine).
Alors $R$ et $A$ sont des anneaux noethériens, et $A$ est fini sur $R$.
De plus, $\mathcal{F} = \widetilde{M}$ pour un certain $A$-module fini
$M$. Notons que $f_*\mathcal{F}$ correspond à $M$ considéré comme $R$-module.
Soit $\mathfrak p \subset R$ l'idéal premier minimal correspondant
à $\eta \in Y$. Le coefficient de $Z$ dans $[f_*\mathcal{F}]_k$
est manifestement $\text{longueur}_{R_{\mathfrak p}}(M_{\mathfrak p})$.
Soient $\mathfrak q_i$, $i = 1, \ldots, t$, les idéaux premiers de $A$
au-dessus de $\mathfrak p$. Alors $A_{\mathfrak p} = \prod A_{\mathfrak q_i}$,
puisque $A_{\mathfrak p}$ est un anneau artinien, étant fini sur l'anneau
local noethérien de dimension zéro $R_{\mathfrak p}$.
Le coefficient de $Z$ dans $f_*[\mathcal{F}]_k$ est manifestement
$$
\sum\nolimits_{i = 1, \ldots, t}
[\kappa(\mathfrak q_i) : \kappa(\mathfrak p)]
\text{longueur}_{A_{\mathfrak q_i}}(M_{\mathfrak q_i})
$$
L'égalité voulue résulte donc de
Algèbre, Lemme \ref{algebra-lemma-pushdown-module}.
\end{proof}




















\section{Préparation à l'image réciproque plate}
\label{chow-section-preparation-flat-pullback}


\noindent
Rappelons qu'un morphisme $f : X \to Y$ qui est localement de type fini
est dit de dimension relative $r$ si toute fibre non vide
est équidimensionnelle de dimension $r$. Voir
Morphismes, Définition \ref{morphisms-definition-relative-dimension-d}.

\begin{lemma}
\label{chow-lemma-flat-inverse-image-dimension}
Soit $(S, \delta)$ comme dans la Situation \ref{chow-situation-setup}.
Soient $X$ et $Y$ localement de type fini sur $S$.
Soit $f : X \to Y$ un morphisme.
Supposons $f$ plat de dimension relative $r$.
Pour tout fermé $Z \subset Y$, nous avons
$$
\dim_\delta(f^{-1}(Z)) = \dim_\delta(Z) + r.
$$
pourvu que $f^{-1}(Z)$ soit non vide.
Si $Z$ est irréductible et si $Z' \subset f^{-1}(Z)$ est une composante irréductible,
alors $Z'$ domine $Z$ et
$\dim_\delta(Z') = \dim_\delta(Z) + r$.
\end{lemma}

\begin{proof}
Il suffit de démontrer la dernière assertion.
Nous pouvons remplacer $Y$ par le sous-schéma fermé intègre $Z$ et
$X$ par l'image réciproque schématique $f^{-1}(Z) = Z \times_Y X$.
Nous pouvons donc supposer $Z = Y$ intègre et $f$ un morphisme plat
de dimension relative $r$. Comme $Y$ est localement noethérien, le
morphisme $f$, qui est localement de type fini,
est en fait localement de présentation finie. Par conséquent,
Morphismes, Lemme \ref{morphisms-lemma-fppf-open},
s'applique et montre que $f$ est ouvert.
Soit $\xi \in X$ le point générique d'une composante irréductible
de $X$. Comme $f$ est ouvert, $f(\xi)$ est le
point générique $\eta$ de $Z = Y$. Notons que $\dim_\xi(X_\eta) = r$,
car $f$ est, par hypothèse, de dimension relative $r$. D'autre
part, comme $\xi$ est un point générique de $X$, on voit que
$\mathcal{O}_{X, \xi} = \mathcal{O}_{X_\eta, \xi}$ ne possède qu'un
idéal premier et est donc de dimension $0$. Ainsi, d'après
Morphismes, Lemme \ref{morphisms-lemma-dimension-fibre-at-a-point},
on conclut que le degré de transcendance
de $\kappa(\xi)$ sur $\kappa(\eta)$ est $r$.
Autrement dit, $\delta(\xi) = \delta(\eta) + r$, comme voulu.
\end{proof}

\noindent
Voici le lemme que nous utiliserons pour démontrer que l'image réciproque plate
d'une collection localement finie de sous-schémas fermés est localement finie.

\begin{lemma}
\label{chow-lemma-inverse-image-locally-finite}
Soit $(S, \delta)$ comme dans la Situation \ref{chow-situation-setup}.
Soient $X$ et $Y$ localement de type fini sur $S$.
Soit $f : X \to Y$ un morphisme.
Supposons que $\{Z_i\}_{i \in I}$ soit une collection localement
finie de fermés de $Y$.
Alors $\{f^{-1}(Z_i)\}_{i \in I}$ est une collection localement finie
de fermés de $X$.
\end{lemma}

\begin{proof}
Soit $U \subset X$ un ouvert quasi-compact.
Comme l'image $f(U) \subset Y$ est un sous-ensemble quasi-compact,
il existe un ouvert quasi-compact $V \subset Y$ tel que
$f(U) \subset V$. Notons que
$$
\{i \in I \mid f^{-1}(Z_i) \cap U \not = \emptyset \}
\subset
\{i \in I \mid Z_i \cap V \not = \emptyset \}.
$$
Le membre de droite étant fini par hypothèse, nous avons terminé.
\end{proof}



\section{Image réciproque plate}
\label{chow-section-flat-pullback}

\noindent
Dans la suite, nous appelons $f^{-1}(Z)$
l'{\it image réciproque schématique} d'un sous-schéma fermé
$Z \subset Y$ par un morphisme de schémas $f : X \to Y$.
Rappelons que l'image réciproque schématique est le produit fibré
$$
\xymatrix{
f^{-1}(Z) \ar[r] \ar[d] & X \ar[d] \\
Z \ar[r] & Y
}
$$
et qu'elle est aussi le sous-schéma fermé de $X$ défini par le
faisceau quasi-cohérent d'idéaux $f^{-1}(\mathcal{I})\mathcal{O}_X$, si
$\mathcal{I} \subset \mathcal{O}_Y$ est le faisceau quasi-cohérent d'idéaux
qui correspond à $Z$ dans $Y$.
(Ceci est étudié dans
Schémas, section \ref{schemes-section-closed-immersion} et
Lemme \ref{schemes-lemma-fibre-product-immersion},
ainsi que dans la Définition \ref{schemes-definition-inverse-image-closed-subscheme}.)

\begin{definition}
\label{chow-definition-flat-pullback}
Soit $(S, \delta)$ comme dans la Situation \ref{chow-situation-setup}.
Soient $X$ et $Y$ localement de type fini sur $S$.
Soit $f : X \to Y$ un morphisme.
Supposons $f$ plat de dimension relative $r$.
\begin{enumerate}
\item Soit $Z \subset Y$ un sous-schéma fermé intègre de
$\delta$-dimension $k$. Nous définissons $f^*[Z]$ comme le
$(k+r)$-cycle sur $X$ associé à l'image réciproque schématique
$$
f^*[Z] = [f^{-1}(Z)]_{k+r}.
$$
Cela a un sens puisque $\dim_\delta(f^{-1}(Z)) = k + r$
d'après le Lemme \ref{chow-lemma-flat-inverse-image-dimension}.
\item Soit $\alpha = \sum n_i [Z_i]$ un
$k$-cycle sur $Y$. L'{\it image réciproque plate de $\alpha$ par $f$}
est la somme
$$
f^* \alpha = \sum n_i f^*[Z_i]
$$
où chaque $f^*[Z_i]$ est défini comme ci-dessus.
La somme est localement finie d'après le Lemme \ref{chow-lemma-inverse-image-locally-finite}.
\item Nous notons $f^* : Z_k(Y) \to Z_{k + r}(X)$ l'application de groupes
abéliens ainsi obtenue.
\end{enumerate}
\end{definition}

\noindent
Une immersion ouverte est plate. C'est un cas particulier important, quoique trivial,
d'un morphisme plat. Si $U \subset X$ est ouvert, on appelle parfois
l'image réciproque par $j : U \to X$ d'un cycle la {\it restriction} du
cycle à $U$. Notons que, dans ce cas, les applications
$$
j^* : Z_k(X) \longrightarrow Z_k(U)
$$
sont toutes {\it surjectives}. En effet, étant donné tout sous-schéma fermé
intègre $Z' \subset U$, on peut prendre l'adhérence $Z$ de $Z'$ dans $X$
et la considérer comme un sous-schéma fermé réduit de $X$ (voir
Schémas, Lemme \ref{schemes-lemma-reduced-closed-subscheme}).
Et, manifestement, $Z \cap U = Z'$ ; autrement dit,
$j^*[Z] = [Z']$, d'où la surjectivité. En fait, un peu plus
est vrai.

\begin{lemma}
\label{chow-lemma-exact-sequence-open}
Soit $(S, \delta)$ comme dans la Situation \ref{chow-situation-setup}.
Soit $X$ localement de type fini sur $S$.
Soit $U \subset X$ un sous-schéma ouvert et notons
$i : Y = X \setminus U \to X$ le sous-schéma fermé réduit de $X$.
Pour tout $k \in \mathbf{Z}$, la suite
$$
\xymatrix{
Z_k(Y) \ar[r]^{i_*} & Z_k(X) \ar[r]^{j^*} & Z_k(U) \ar[r] & 0
}
$$
est un complexe exact de groupes abéliens.
\end{lemma}

\begin{proof}
Supposons d'abord $X$ quasi-compact. Alors $Z_k(X)$ est un $\mathbf{Z}$-module libre
ayant pour base les éléments $[Z]$, où $Z \subset X$ est fermé intègre
de $\delta$-dimension $k$. Un tel élément de base s'envoie
soit sur l'élément de base $[Z \cap U]$, soit sur zéro si $Z \subset Y$.
Le lemme est donc clair dans ce cas. Le cas général est analogue
et la démonstration est omise.
\end{proof}

\begin{lemma}
\label{chow-lemma-compose-flat-pullback}
Soit $(S, \delta)$ comme dans la Situation \ref{chow-situation-setup}.
Soient $X, Y, Z$ localement de type fini sur $S$.
Soient $f : X \to Y$ et $g : Y \to Z$ des morphismes plats de dimensions relatives
$r$ et $s$. Alors $g \circ f$ est plat de dimension relative
$r + s$ et
$$
f^* \circ g^* = (g \circ f)^*
$$
comme applications $Z_k(Z) \to Z_{k + r + s}(X)$.
\end{lemma}

\begin{proof}
La composée est plate de dimension relative $r + s$ d'après
Morphismes, Lemme \ref{morphisms-lemma-composition-relative-dimension-d}.
Supposons que
\begin{enumerate}
\item $W \subset Z$ soit un sous-schéma fermé intègre de $\delta$-dimension $k$,
\item $W' \subset Y$ soit un sous-schéma fermé intègre de $\delta$-dimension
$k + s$ tel que $W' \subset g^{-1}(W)$, et
\item $W'' \subset Y$ soit un sous-schéma fermé intègre de $\delta$-dimension
$k + s + r$ tel que $W'' \subset f^{-1}(W')$.
\end{enumerate}
Il faut montrer que le coefficient $n$ de $[W'']$ dans
$(g \circ f)^*[W]$ coïncide avec le coefficient $m$ de
$[W'']$ dans $f^*(g^*[W])$. Le fait qu'il suffise de vérifier le lemme dans ces
cas résulte du Lemme \ref{chow-lemma-flat-inverse-image-dimension}.
Soient $\xi'' \in W''$, $\xi' \in W'$
et $\xi \in W$ les points génériques. Considérons les anneaux locaux
$A = \mathcal{O}_{Z, \xi}$, $B = \mathcal{O}_{Y, \xi'}$
et $C = \mathcal{O}_{X, \xi''}$. Nous avons alors des morphismes locaux plats d'anneaux
$A \to B$, $B \to C$ et, de plus,
$$
n = \text{longueur}_C(C/\mathfrak m_AC),
\quad
\text{et}
\quad
m = \text{longueur}_C(C/\mathfrak m_BC) \text{longueur}_B(B/\mathfrak m_AB)
$$
L'égalité résulte donc de
Algèbre, Lemme \ref{algebra-lemma-pullback-transitive}.
\end{proof}

\begin{lemma}
\label{chow-lemma-pullback-coherent}
Soit $(S, \delta)$ comme dans la Situation \ref{chow-situation-setup}.
Soient $X, Y$ localement de type fini sur $S$.
Soit $f : X \to Y$ un morphisme plat de dimension relative $r$.
\begin{enumerate}
\item Soit $Z \subset Y$ un sous-schéma fermé tel que
$\dim_\delta(Z) \leq k$. Alors
$\dim_\delta(f^{-1}(Z)) \leq k + r$
et $[f^{-1}(Z)]_{k + r} = f^*[Z]_k$ dans $Z_{k + r}(X)$.
\item Soit $\mathcal{F}$ un faisceau cohérent sur $Y$ tel que
$\dim_\delta(\text{Supp}(\mathcal{F})) \leq k$.
Alors $\dim_\delta(\text{Supp}(f^*\mathcal{F})) \leq k + r$
et
$$
f^*[{\mathcal F}]_k = [f^*{\mathcal F}]_{k+r}
$$
dans $Z_{k + r}(X)$.
\end{enumerate}
\end{lemma}

\begin{proof}
Les assertions sur les dimensions résultent immédiatement du
Lemme \ref{chow-lemma-flat-inverse-image-dimension}.
La partie (1) résulte de la partie (2), du Lemme \ref{chow-lemma-cycle-closed-coherent}
et du fait que $f^*\mathcal{O}_Z = \mathcal{O}_{f^{-1}(Z)}$.

\medskip\noindent
Démonstration de (2).
Comme $X$ et $Y$ sont localement noethériens, nous pouvons appliquer
Cohomologie des schémas, Lemme \ref{coherent-lemma-coherent-Noetherian}, pour voir
que $\mathcal{F}$ est de type fini, donc que $f^*\mathcal{F}$ est
de type fini (Modules, Lemme \ref{modules-lemma-pullback-finite-type}),
et donc que $f^*\mathcal{F}$ est cohérent
(de nouveau Cohomologie des schémas, Lemme \ref{coherent-lemma-coherent-Noetherian}).
Le lemme a donc un sens. Soit $W \subset Y$ un sous-schéma fermé intègre
de $\delta$-dimension $k$, et soit $W' \subset X$ un sous-schéma fermé
intègre de dimension $k + r$ qui s'envoie dans $W$
par $f$. Il faut montrer que le coefficient $n$ de
$[W']$ dans $f^*[{\mathcal F}]_k$ coïncide avec le coefficient
$m$ de $[W']$ dans $[f^*{\mathcal F}]_{k+r}$. Soient $\xi \in W$ et
$\xi' \in W'$ soient les points génériques. Posons
$A = \mathcal{O}_{Y, \xi}$, $B = \mathcal{O}_{X, \xi'}$
et considérons $M = \mathcal{F}_\xi$ comme un $A$-module. (Notons que
$M$ est de longueur finie d'après nos hypothèses sur les dimensions, mais nous
n'avons en fait pas besoin de le vérifier. Voir le
Lemme \ref{chow-lemma-length-finite}.)
Nous avons $f^*\mathcal{F}_{\xi'} = B \otimes_A M$.
Nous voyons ainsi que
$$
n = \text{longueur}_B(B \otimes_A M)
\quad
\text{et}
\quad
m = \text{longueur}_A(M) \text{longueur}_B(B/\mathfrak m_AB)
$$
L'égalité résulte donc du
Lemme \ref{algebra-lemma-pullback-module} d'Algèbre.
\end{proof}



\section{Images directes et images réciproques}
\label{chow-section-push-pull}

\noindent
Dans cette section, nous vérifions que l'image directe propre et l'image réciproque plate
sont compatibles lorsque cela a un sens. D'après le travail effectué ci-dessus, ceci
est une conséquence de la cohomologie et du changement de base.

\begin{lemma}
\label{chow-lemma-flat-pullback-proper-pushforward}
Soit $(S, \delta)$ comme dans la Situation \ref{chow-situation-setup}.
Soit
$$
\xymatrix{
X' \ar[r]_{g'} \ar[d]_{f'} & X \ar[d]^f \\
Y' \ar[r]^g & Y
}
$$
un diagramme cartésien de schémas localement de type fini sur $S$.
Supposons $f : X \to Y$ propre et $g : Y' \to Y$ plat de dimension relative $r$.
Alors $f'$ est aussi propre et $g'$ est aussi plat de dimension relative $r$.
Pour tout $k$-cycle $\alpha$ sur $X$, nous avons
$$
g^*f_*\alpha = f'_*(g')^*\alpha
$$
dans $Z_{k + r}(Y')$.
\end{lemma}

\begin{proof}
Le fait que $f'$ soit propre résulte du
Lemme \ref{morphisms-lemma-base-change-proper} de Morphismes.
Le fait que $g'$ soit plat de dimension relative $r$ résulte des
Lemmes \ref{morphisms-lemma-base-change-relative-dimension-d}
et \ref{morphisms-lemma-base-change-flat} de Morphismes.
Il suffit de démontrer l'égalité de cycles lorsque $\alpha = [W]$
pour un sous-schéma fermé intègre $W \subset X$ de $\delta$-dimension $k$.
Notons que, dans ce cas, nous avons $\alpha = [\mathcal{O}_W]_k$ ; voir le
Lemme \ref{chow-lemma-cycle-closed-coherent}.
D'après les Lemmes \ref{chow-lemma-cycle-push-sheaf} et
\ref{chow-lemma-pullback-coherent}, il suffit donc
de montrer que $f'_*(g')^*\mathcal{O}_W$ est isomorphe à
$g^*f_*\mathcal{O}_W$. Ceci résulte de la cohomologie et du
changement de base ; voir le
Lemme \ref{coherent-lemma-flat-base-change-cohomology} de Cohomologie des schémas.
\end{proof}

\begin{lemma}
\label{chow-lemma-finite-flat}
Soit $(S, \delta)$ comme dans la Situation \ref{chow-situation-setup}.
Soient $X$, $Y$ localement de type fini sur $S$.
Soit $f : X \to Y$ un morphisme fini localement libre
de degré $d$ (voir la
Définition \ref{morphisms-definition-finite-locally-free} de Morphismes).
Alors $f$ est à la fois propre et plat de dimension relative $0$, et
$$
f_*f^*\alpha = d\alpha
$$
pour tout $\alpha \in Z_k(Y)$.
\end{lemma}

\begin{proof}
Un morphisme fini localement libre est plat et fini d'après le
Lemme \ref{morphisms-lemma-finite-flat} de Morphismes,
et un morphisme fini est propre
d'après le Lemme \ref{morphisms-lemma-finite-proper} de Morphismes.
Nous omettons de montrer qu'un morphisme fini
est de dimension relative $0$. La formule a donc un sens.
Pour la démontrer, soit $Z \subset Y$ un sous-schéma fermé intègre
de $\delta$-dimension $k$. Il suffit de démontrer la formule
pour $\alpha = [Z]$. Puisque le changement de base d'un morphisme fini localement libre
est fini localement libre
(Lemme \ref{morphisms-lemma-base-change-finite-locally-free} de Morphismes),
nous voyons que $f_*f^*\mathcal{O}_Z$ est un faisceau fini localement libre de
rang $d$ sur $Z$. Ainsi
$$
f_*f^*[Z] = f_*f^*[\mathcal{O}_Z]_k =
[f_*f^*\mathcal{O}_Z]_k = d[Z]
$$
où nous avons utilisé les Lemmes \ref{chow-lemma-pullback-coherent} et
\ref{chow-lemma-cycle-push-sheaf}.
\end{proof}








\section{Préparation aux diviseurs principaux}
\label{chow-section-preparation-principal-divisors}

\noindent
Une partie du contenu de cette section recoupe partiellement la
discussion de Diviseurs, section \ref{divisors-section-Weil-divisors}.

\begin{lemma}
\label{chow-lemma-divisor-delta-dimension}
Soit $(S, \delta)$ comme dans la Situation \ref{chow-situation-setup}.
Soit $X$ localement de type fini sur $S$. Supposons $X$
intègre.
\begin{enumerate}
\item Si $Z \subset X$ est un sous-schéma fermé intègre, alors
les conditions suivantes sont équivalentes :
\begin{enumerate}
\item $Z$ est un diviseur premier,
\item $Z$ est de codimension $1$ dans $X$, et
\item $\dim_\delta(Z) = \dim_\delta(X) - 1$.
\end{enumerate}
\item Si $Z$ est une composante irréductible d'un diviseur de Cartier
effectif sur $X$, alors $\dim_\delta(Z) = \dim_\delta(X) - 1$.
\end{enumerate}
\end{lemma}

\begin{proof}
La partie (1) résulte de la définition d'un diviseur premier
(Définition \ref{divisors-definition-Weil-divisor} de Diviseurs)
et de la définition d'une fonction de dimension
(Définition \ref{topology-definition-dimension-function} de Topologie).
Soit $\xi \in Z$ le point générique d'une composante irréductible $Z$ d'un
diviseur de Cartier effectif $D \subset X$.
Alors $\dim(\mathcal{O}_{D, \xi}) = 0$ et
$\mathcal{O}_{D, \xi} = \mathcal{O}_{X, \xi}/(f)$ pour un certain
élément non diviseur de zéro $f \in \mathcal{O}_{X, \xi}$ (Lemme
\ref{divisors-lemma-effective-Cartier-in-points} de Diviseurs).
Alors $\dim(\mathcal{O}_{X, \xi}) = 1$ d'après le
Lemme \ref{algebra-lemma-one-equation} d'Algèbre. Ainsi $Z$ est comme dans (1) d'après le
Lemme \ref{properties-lemma-codimension-local-ring} de Propriétés,
ce qui achève la démonstration.
\end{proof}

\begin{lemma}
\label{chow-lemma-finite-in-codimension-one}
Soit $f : X \to Y$ un morphisme de schémas.
Soit $\xi \in Y$ un point.
Supposons que
\begin{enumerate}
\item $X$, $Y$ sont intègres,
\item $Y$ est localement noethérien,
\item $f$ est propre, dominant et l'extension $R(Y) \subset R(X)$ est finie, et
\item $\dim(\mathcal{O}_{Y, \xi}) = 1$.
\end{enumerate}
Alors il existe un voisinage ouvert $V \subset Y$ de $\xi$
tel que $f|_{f^{-1}(V)} : f^{-1}(V) \to V$ soit fini.
\end{lemma}

\begin{proof}
Ce lemme est un cas particulier du
Lemme \ref{varieties-lemma-finite-in-codim-1} de Variétés.
Voici un argument direct dans ce cas.
D'après le Lemme
\ref{coherent-lemma-proper-finite-fibre-finite-in-neighbourhood} de Cohomologie des schémas,
il suffit de démontrer que $f^{-1}(\{\xi\})$ est fini.
Nous remplaçons $Y$ par un ouvert affine, disons $Y = \Spec(R)$.
Notons que $R$ est noethérien, puisque $Y$ est supposé localement noethérien.
Puisque $f$ est propre, il est quasi-compact. Nous pouvons donc trouver un recouvrement
ouvert affine fini $X = U_1 \cup \ldots \cup U_n$, avec
chaque $U_i = \Spec(A_i)$. Notons que $R \to A_i$ est un
homomorphisme injectif de type fini de domaines tel que
l'extension induite des corps des fractions soit finie.
Le lemme résulte donc
du Lemme \ref{algebra-lemma-finite-in-codim-1} d'Algèbre.
\end{proof}


\section{Diviseurs principaux}
\label{chow-section-principal-divisors}

\noindent
La définition suivante est l'analogue de la
Définition \ref{divisors-definition-principal-divisor} de Diviseurs
dans notre cadre actuel.

\begin{definition}
\label{chow-definition-principal-divisor}
Soit $(S, \delta)$ comme dans la Situation \ref{chow-situation-setup}.
Soit $X$ localement de type fini sur $S$. Supposons $X$
intègre avec $\dim_\delta(X) = n$.
Soit $f \in R(X)^*$. Le {\it diviseur principal
associé à $f$} est le $(n - 1)$-cycle
$$
\text{div}(f) = \text{div}_X(f) = \sum \text{ord}_Z(f) [Z]
$$
défini dans Diviseurs, Définition \ref{divisors-definition-principal-divisor}.
Ceci a un sens parce que les diviseurs premiers sont de $\delta$-dimension $n - 1$ d'après le
Lemme \ref{chow-lemma-divisor-delta-dimension}.
\end{definition}

\noindent
Dans la situation de la définition, pour $f, g \in R(X)^*$, nous avons
$$
\text{div}_X(fg) = \text{div}_X(f) + \text{div}_X(g)
$$
dans $Z_{n - 1}(X)$. Voir Diviseurs, Lemme \ref{divisors-lemma-div-additive}.
Le lemme suivant sera supplanté par le Lemme plus général
\ref{chow-lemma-flat-pullback-rational-equivalence}.

\begin{lemma}
\label{chow-lemma-flat-pullback-principal-divisor}
Soit $(S, \delta)$ comme dans la Situation \ref{chow-situation-setup}.
Soient $X$, $Y$ localement de type fini sur $S$. Supposons $X$, $Y$
intègres et $n = \dim_\delta(Y)$.
Soit $f : X \to Y$ un morphisme plat de dimension relative $r$.
Soit $g \in R(Y)^*$. Alors
$$
f^*(\text{div}_Y(g)) = \text{div}_X(g)
$$
dans $Z_{n + r - 1}(X)$.
\end{lemma}

\begin{proof}
Notons que, puisque $f$ est plat, il est dominant, de sorte que
$f$ induit un plongement $R(Y) \subset R(X)$ ; nous pouvons donc
considérer $g$ comme un élément de $R(X)^*$.
Soit $Z \subset X$ un sous-schéma fermé intègre de
$\delta$-dimension $n + r - 1$. Soit $\xi \in Z$
son point générique. Si $\dim_\delta(f(Z)) > n - 1$,
alors nous voyons que le coefficient de $[Z]$ dans les membres de gauche et
de droite de l'équation est nul.
Nous pouvons donc supposer que $Z' = \overline{f(Z)}$ est un
sous-schéma fermé intègre de $Y$ de $\delta$-dimension $n - 1$.
Soit $\xi' = f(\xi)$. C'est le point générique de $Z'$.
Posons $A = \mathcal{O}_{Y, \xi'}$, $B = \mathcal{O}_{X, \xi}$.
L'homomorphisme d'anneaux $A \to B$ est un homomorphisme local plat de
domaines locaux noethériens de dimension $1$.
Nous avons $g$ dans le corps des fractions de $A$. Ce que nous devons montrer est que
$$
\text{ord}_A(g) \text{longueur}_B(B/\mathfrak m_AB)
=
\text{ord}_B(g).
$$
Ceci résulte du Lemme \ref{algebra-lemma-pullback-module} d'Algèbre
(détails omis).
\end{proof}











\section{Diviseurs principaux et image directe}
\label{chow-section-two-fun}

\noindent
Le premier lemme implique que l'image directe d'un diviseur
principal par un morphisme génériquement fini est un diviseur principal.

\begin{lemma}
\label{chow-lemma-proper-pushforward-alteration}
Soit $(S, \delta)$ comme dans la Situation \ref{chow-situation-setup}.
Soient $X$, $Y$ localement de type fini sur $S$. Supposons $X$, $Y$
intègres et $n = \dim_\delta(X) = \dim_\delta(Y)$.
Soit $p : X \to Y$ un morphisme propre dominant.
Soit $f \in R(X)^*$. Posons
$$
g = \text{Nm}_{R(X)/R(Y)}(f).
$$
Alors nous avons
$p_*\text{div}(f) = \text{div}(g)$.
\end{lemma}

\begin{proof}
Soit $Z \subset Y$ un sous-schéma fermé intègre de $\delta$-dimension
$n - 1$. Nous voulons montrer que les coefficients de $[Z]$ dans
$p_*\text{div}(f)$ et $\text{div}(g)$ sont égaux. Nous pouvons appliquer le
Lemme \ref{chow-lemma-finite-in-codimension-one}
au morphisme $p : X \to Y$ et au point générique $\xi \in Z$.
Nous pouvons donc remplacer $Y$ par un
voisinage ouvert affine de $\xi$ et supposer que $p : X \to Y$ est fini.
Écrivons $Y = \Spec(R)$ et $X = \Spec(A)$, où $p$ est induit
par un homomorphisme fini $R \to A$ de domaines noethériens qui induit
une extension finie $L/K$ des corps des fractions.
Nous avons alors $f \in L$, $g = \text{Nm}(f) \in K$,
et un idéal premier $\mathfrak p \subset R$ avec $\dim(R_{\mathfrak p}) = 1$.
Le coefficient de $[Z]$ dans $\text{div}_Y(g)$ est
$\text{ord}_{R_\mathfrak p}(g)$.
Le coefficient de $[Z]$ dans $p_*\text{div}_X(f)$ est
$$
\sum\nolimits_{\mathfrak q\text{ au-dessus de }\mathfrak p}
[\kappa(\mathfrak q) : \kappa(\mathfrak p)]
\text{ord}_{A_{\mathfrak q}}(f)
$$
L'égalité souhaitée résulte donc du
Lemme \ref{algebra-lemma-finite-extension-dim-1} d'Algèbre.
\end{proof}

\noindent
Dans l'étude des diviseurs principaux, un rôle important
est joué par le diviseur principal « universel » $[0] - [\infty]$
sur $\mathbf{P}^1_S$. Pour préciser cela, notons
\begin{equation}
\label{chow-equation-zero-infty}
D_0, D_\infty \subset
\mathbf{P}^1_S = \underline{\text{Proj}}_S(\mathcal{O}_S[T_0, T_1])
\end{equation}
les sous-schémas fermés définis par la section $T_1$, resp.\ $T_0$
de $\mathcal{O}(1)$. Ce sont des diviseurs de Cartier effectifs ; voir la
Définition \ref{divisors-definition-effective-Cartier-divisor}
et le Lemme \ref{divisors-lemma-characterize-OD} de Diviseurs.
Le lemme suivant dit que, de façon informelle, nous avons
« $\text{div}(T_1/T_0) = [D_0] - [D_1]$ » et que c'est le
diviseur principal universel.

\begin{lemma}
\label{chow-lemma-rational-function}
Soit $(S, \delta)$ comme dans la Situation \ref{chow-situation-setup}.
Soit $X$ localement de type fini sur $S$. Supposons $X$
intègre et $n = \dim_\delta(X)$. Soit $f \in R(X)^*$.
Soit $U \subset X$ un ouvert non vide tel que $f$
corresponde à une section $f \in \Gamma(U, \mathcal{O}_X^*)$.
Soit $Y \subset X \times_S \mathbf{P}^1_S$ l'adhérence du
graphe de $f : U \to \mathbf{P}^1_S$.
Alors
\begin{enumerate}
\item le morphisme de projection $p : Y \to X$ est propre,
\item $p|_{p^{-1}(U)} : p^{-1}(U) \to U$ est un isomorphisme,
\item les images réciproques $Y_0 = q^{-1}D_0$ et $Y_\infty = q^{-1}D_\infty$
par le morphisme $q : Y \to \mathbf{P}^1_S$ sont définies
(Diviseurs, Définition
\ref{divisors-definition-pullback-effective-Cartier-divisor}),
\item nous avons
$$
\text{div}_Y(f) = [Y_0]_{n - 1} - [Y_\infty]_{n - 1}
$$
\item nous avons
$$
\text{div}_X(f) = p_*\text{div}_Y(f)
$$
\item si nous considérons $Y_0$ et $Y_\infty$ comme des sous-schémas fermés de $X$
par le morphisme $p$, alors nous avons
$$
\text{div}_X(f) = [Y_0]_{n - 1} - [Y_\infty]_{n - 1}
$$
\end{enumerate}
\end{lemma}

\begin{proof}
Puisque $X$ est intègre, nous voyons que $U$ est intègre.
Ainsi $Y$ est intègre, et $(1, f)(U) \subset Y$ est un sous-schéma ouvert dense.
Notons aussi que le sous-schéma fermé $Y \subset X \times_S \mathbf{P}^1_S$
ne dépend pas du choix de l'ouvert $U$, puisqu'il s'agit en définitive de
l'adhérence de l'ensemble à un point $\{\eta'\} = \{(1, f)(\eta)\}$,
où $\eta \in X$ est le point générique. Cela étant dit, démontrons les
assertions du lemme.

\medskip\noindent
Pour (1), notons que $p$ est la composée de l'immersion fermée
$Y \to X \times_S \mathbf{P}^1_S = \mathbf{P}^1_X$ et du morphisme propre
$\mathbf{P}^1_X \to X$. Comme une composée de morphismes propres
est propre (Morphismes, Lemme \ref{morphisms-lemma-composition-proper}),
nous concluons.

\medskip\noindent
Il est clair que $Y \cap U \times_S \mathbf{P}^1_S = (1, f)(U)$.
Ainsi (2) en résulte. Il en résulte aussi que $\dim_\delta(Y) = n$.

\medskip\noindent
Notons que $q(\eta') = f(\eta)$ n'est contenu ni dans $D_0$ ni dans $D_\infty$,
puisque $f \in R(X)^*$. D'où (3) par le
Lemme \ref{divisors-lemma-pullback-effective-Cartier-defined} de Diviseurs.
Nous obtenons $\dim_\delta(Y_0) = n - 1$
et $\dim_\delta(Y_\infty) = n - 1$ par le
Lemme \ref{chow-lemma-divisor-delta-dimension}.

\medskip\noindent
Considérons le diviseur de Cartier effectif $Y_0$.
En tout point $\xi \in Y_0$, nous avons $f \in \mathcal{O}_{Y, \xi}$ et
une équation locale de $Y_0$ est donnée par $f$.
En particulier, si $\delta(\xi) = n - 1$, de sorte que $\xi$ est le point générique
d'un sous-schéma fermé intègre $Z$ de $\delta$-dimension $n - 1$,
alors nous voyons que le coefficient de $[Z]$ dans $\text{div}_Y(f)$ est
$$
\text{ord}_Z(f) =
\text{longueur}_{\mathcal{O}_{Y, \xi}}
(\mathcal{O}_{Y, \xi}/f\mathcal{O}_{Y, \xi}) =
\text{longueur}_{\mathcal{O}_{Y, \xi}}
(\mathcal{O}_{Y_0, \xi})
$$
qui est le coefficient de $[Z]$ dans $[Y_0]_{n - 1}$. Un argument semblable
utilisant la fonction rationnelle $1/f$ montre que
$-[Y_\infty]$ coïncide avec les termes de coefficients négatifs dans
l'expression de $\text{div}_Y(f)$. D'où (4).

\medskip\noindent
Notons que $D_0 \to S$ est un isomorphisme. Nous voyons donc que
$X \times_S D_0 \to X$ est également un isomorphisme. Il est clair que
nous avons $Y_0 = Y \cap X \times_S D_0$ (intersection schématique)
dans $X \times_S \mathbf{P}^1_S$. Il s'agit donc bien du cas où
$Y_0 \to X$ est une immersion fermée. Il en résulte que
$$
p_*\mathcal{O}_{Y_0} = \mathcal{O}_{Y'_0}
$$
où $Y'_0 \subset X$ est l'image de $Y_0 \to X$.
D'après le Lemme \ref{chow-lemma-cycle-push-sheaf}, nous
avons $p_*[Y_0]_{n - 1} = [Y'_0]_{n - 1}$. Il en va de même
pour $D_\infty$ et $Y_\infty$. Ainsi (6) est une conséquence de (5).
Enfin, (5) résulte immédiatement du
Lemme \ref{chow-lemma-proper-pushforward-alteration}.
\end{proof}

\noindent
Le lemme suivant dit que le degré d'un diviseur principal sur
une courbe propre est nul.

\begin{lemma}
\label{chow-lemma-curve-principal-divisor}
Soit $K$ un corps quelconque. Soit $X$ un schéma intègre de dimension $1$
muni d'un morphisme propre $c : X \to \Spec(K)$.
Soit $f \in K(X)^*$ une fonction rationnelle inversible.
Alors
$$
\sum\nolimits_{x \in X \text{ fermé}}
[\kappa(x) : K] \text{ord}_{\mathcal{O}_{X, x}}(f)
=
0
$$
où $\text{ord}$ est comme dans la
Définition \ref{algebra-definition-ord} d'Algèbre.
Autrement dit, $c_*\text{div}(f) = 0$.
\end{lemma}

\begin{proof}
Considérons le diagramme
$$
\xymatrix{
Y \ar[r]_p \ar[d]_q & X \ar[d]^c \\
\mathbf{P}^1_K \ar[r]^-{c'} & \Spec(K)
}
$$
que nous avons construit dans le Lemme \ref{chow-lemma-rational-function}
à partir de $X$ et de la fonction rationnelle $f$ sur $S = \Spec(K)$.
Nous utiliserons sans autre mention tous les résultats de ce lemme.
Nous devons montrer que $c_*\text{div}_X(f) = c_*p_*\text{div}_Y(f) = 0$.
Cela revient à démontrer que $c'_*q_*\text{div}_Y(f) = 0$.
Si $q(Y)$ est un point fermé de $\mathbf{P}^1_K$, alors nous
voyons que $\text{div}_X(f) = 0$ et le lemme est démontré.
Nous pouvons donc supposer que $q$ est dominant.
Supposons que nous puissions montrer que $q : Y \to \mathbf{P}^1_K$ est fini
localement libre de degré $d$ (voir la
Définition \ref{morphisms-definition-finite-locally-free} de Morphismes).
Puisque $\text{div}_Y(f) = [q^{-1}D_0]_0 - [q^{-1}D_\infty]_0$,
nous voyons (par définition de l'image réciproque plate) que
$\text{div}_Y(f) = q^*([D_0]_0 - [D_\infty]_0)$.
Alors le Lemme \ref{chow-lemma-finite-flat} donne
$q_*\text{div}_Y(f) = d([D_0]_0 - [D_\infty]_0)$.
Comme il est clair que $c'_*[D_0]_0 = c'_*[D_\infty]_0$, nous avons le résultat.

\medskip\noindent
Il reste à montrer que $q$ est fini localement libre.
(Il aura automatiquement un certain degré donné puisque $\mathbf{P}^1_K$
est connexe.)
Comme $\dim(\mathbf{P}^1_K) = 1$, nous voyons que $q$ est fini, par exemple
d'après le Lemme \ref{chow-lemma-finite-in-codimension-one}.
Tous les anneaux locaux de $\mathbf{P}^1_K$ aux
points fermés sont des anneaux locaux réguliers de dimension $1$
(autrement dit, des anneaux de valuation discrète), puisqu'ils sont des
localisés de $K[T]$ (voir le
Lemme \ref{algebra-lemma-dim-affine-space} d'Algèbre).
Ainsi, pour $y\in Y$ fermé, l'anneau local $\mathcal{O}_{Y, y}$
sera plat sur $\mathcal{O}_{\mathbf{P}^1_K, q(y)}$ dès qu'il est
sans torsion (Plus sur l'algèbre, Lemme
\ref{more-algebra-lemma-dedekind-torsion-free-flat}).
C'est manifestement le cas puisque
$\mathcal{O}_{Y, y}$ est un domaine et $q$ est dominant.
Ainsi $q$ est plat. Par conséquent, $q$ est fini localement libre d'après le
Lemme \ref{morphisms-lemma-finite-flat} de Morphismes.
\end{proof}





\section{Équivalence rationnelle}
\label{chow-section-rational-equivalence}

\noindent
Dans cette section, nous définissons l'{\it équivalence rationnelle} sur les $k$-cycles.
Nous autoriserons les sommes localement finies d'images de
diviseurs principaux (par des immersions fermées). Cela conduit à des
phénomènes assez étranges ; voir l'Exemple \ref{chow-example-weird}.
Cependant, si nous ne les autorisons pas, nous ne savons pas démontrer que
le produit avec les classes de Chern de fibrés inversibles se factorise par l'équivalence
rationnelle.

\begin{definition}
\label{chow-definition-rational-equivalence}
Soit $(S, \delta)$ comme dans la Situation \ref{chow-situation-setup}.
Soit $X$ un schéma localement de type fini sur $S$.
Soit $k \in \mathbf{Z}$.
\begin{enumerate}
\item Étant donnée toute famille localement finie $\{W_j \subset X\}$
de sous-schémas fermés intègres avec $\dim_\delta(W_j) = k + 1$,
et des $f_j \in R(W_j)^*$ quelconques, nous pouvons considérer
$$
\sum (i_j)_*\text{div}(f_j) \in Z_k(X)
$$
où $i_j : W_j \to X$ est le morphisme d'inclusion.
Ceci a un sens car le morphisme
$\coprod i_j : \coprod W_j \to X$ est propre.
\item Nous disons que $\alpha \in Z_k(X)$ est {\it rationnellement équivalent à zéro}
si $\alpha$ est un cycle de la forme affichée ci-dessus.
\item Nous disons que $\alpha, \beta \in Z_k(X)$ sont
{\it rationnellement équivalents} et nous écrivons $\alpha \sim_{rat} \beta$
si $\alpha - \beta$ est rationnellement équivalent à zéro.
\item Nous définissons
$$
\CH_k(X) = Z_k(X) / \sim_{rat}
$$
comme étant le {\it groupe de Chow des $k$-cycles sur $X$}. On l'appelle parfois
le {\it groupe de Chow des $k$-cycles modulo équivalence rationnelle sur $X$}.
\end{enumerate}
\end{definition}

\noindent
Il existe beaucoup d'autres relations d'équivalence (adéquates) intéressantes.
L'équivalence rationnelle est la moins fine de toutes.

\begin{remark}
\label{chow-remark-chow-group-pointwise}
Soit $(S, \delta)$ comme dans la Situation \ref{chow-situation-setup}.
Soit $X$ localement de type fini sur $S$. Soit $k \in \mathbf{Z}$.
Montrons que nous avons une présentation
$$
\bigoplus\nolimits_{\delta(x) = k + 1}' K_1^M(\kappa(x))
\xrightarrow{\partial}
\bigoplus\nolimits_{\delta(x) = k}' K_0^M(\kappa(x)) \to
\CH_k(X) \to 0
$$
Nous utilisons ici les notations et conventions introduites dans la
Remarque \ref{chow-remark-cycles-pointwise} et, de plus,
\begin{enumerate}
\item $K_1^M(\kappa(x)) = \kappa(x)^*$ est la partie de degré $1$ de
la K-théorie de Milnor du corps résiduel $\kappa(x)$ du point
$x \in X$ (voir la Remarque \ref{chow-remark-gersten-complex-milnor}), et
\item la différentielle $\partial$ est définie comme suit :
étant donné un élément $\xi = \sum_x f_x$, nous notons $W_x = \overline{x}$
le sous-schéma fermé intègre de $X$ de point générique $x$ et nous posons
$$
\partial(\xi) = \sum (W_x \to X)_*\text{div}(f_x)
$$
dans $Z_k(X)$, ce qui a un sens puisque nous avons vu que le deuxième
terme du complexe est égal à $Z_k(X)$ d'après la
Remarque \ref{chow-remark-cycles-pointwise}.
\end{enumerate}
Le fait que nous obtenions une présentation de $\CH_k(X)$ résulte
immédiatement de la comparaison avec la Définition \ref{chow-definition-rational-equivalence}.
\end{remark}

\noindent
Voici un lemme très simple mais important.

\begin{lemma}
\label{chow-lemma-restrict-to-open}
Soit $(S, \delta)$ comme dans la Situation \ref{chow-situation-setup}.
Soit $X$ un schéma localement de type fini sur $S$.
Soit $U \subset X$ un sous-schéma ouvert, et notons
$i : Y = X \setminus U \to X$ le sous-schéma fermé réduit de $X$.
Soit $k \in \mathbf{Z}$.
Supposons $\alpha, \beta \in Z_k(X)$.
Si $\alpha|_U \sim_{rat} \beta|_U$, alors il existe un cycle
$\gamma \in Z_k(Y)$ tel que
$$
\alpha \sim_{rat} \beta + i_*\gamma.
$$
Autrement dit, la suite
$$
\xymatrix{
\CH_k(Y) \ar[r]^{i_*} & \CH_k(X) \ar[r]^{j^*} & \CH_k(U) \ar[r] & 0
}
$$
est un complexe exact de groupes abéliens.
\end{lemma}

\begin{proof}
Soit $\{W_j\}_{j \in J}$ une famille localement finie de sous-schémas fermés
intègres de $U$, de $\delta$-dimension $k + 1$, et soient $f_j \in R(W_j)^*$
des éléments tels que $(\alpha - \beta)|_U = \sum (i_j)_*\text{div}(f_j)$
comme dans la définition. Posons $W_j' \subset X$ égal
à l'adhérence de $W_j$. Supposons que $V \subset X$ soit un ouvert
quasi-compact. Alors $V \cap U$ est aussi un ouvert quasi-compact de $U$, car
$V$ est noethérien. Par conséquent, l'ensemble
$\{j \in J \mid W_j \cap V \not = \emptyset\}
= \{j \in J \mid W'_j \cap V \not = \emptyset\}$
est fini puisque $\{W_j\}$ est localement finie. Autrement dit, nous voyons que
$\{W'_j\}$ est également localement finie. Comme $R(W_j) = R(W'_j)$, nous voyons
que
$$
\alpha - \beta - \sum (i'_j)_*\text{div}(f_j)
$$
est un cycle supporté sur $Y$, et le lemme en résulte (voir le
Lemme \ref{chow-lemma-exact-sequence-open}).
\end{proof}

\begin{lemma}
\label{chow-lemma-exact-sequence-closed-chow}
Soit $(S, \delta)$ comme dans la Situation \ref{chow-situation-setup}.
Soit $X$ localement de type fini sur $S$. Soient $X_1, X_2 \subset X$
des sous-schémas fermés tels que $X = X_1 \cup X_2$ ensemblistement.
Pour tout $k \in \mathbf{Z}$, la suite de groupes abéliens
$$
\xymatrix{
\CH_k(X_1 \cap X_2) \ar[r] &
\CH_k(X_1) \oplus \CH_k(X_2) \ar[r] &
\CH_k(X) \ar[r] &
0
}
$$
est exacte. Ici $X_1 \cap X_2$ est l'intersection schématique, et les
applications sont les applications d'image directe, l'une étant multipliée par $-1$.
\end{lemma}

\begin{proof}
D'après le Lemme \ref{chow-lemma-exact-sequence-closed}, la flèche
$\CH_k(X_1) \oplus \CH_k(X_2) \to \CH_k(X)$ est surjective.
Supposons que $(\alpha_1, \alpha_2)$ soit envoyé sur zéro par cette application.
Écrivons $\alpha_1 = \sum n_{1, i}[W_{1, i}]$ et
$\alpha_2 = \sum n_{2, i}[W_{2, i}]$. Nous obtenons alors une famille
localement finie $\{W_j\}_{j \in J}$ de sous-schémas fermés intègres
de $X$ de $\delta$-dimension $k + 1$ et des $f_j \in R(W_j)^*$
tels que
$$
\sum n_{1, i}[W_{1, i}] + \sum n_{2, i}[W_{2, i}] = \sum (i_j)_*\text{div}(f_j)
$$
comme cycles sur $X$, où $i_j : W_j \to X$ est le morphisme d'inclusion.
Choisissons une décomposition en union disjointe $J = J_1 \amalg J_2$ telle que
$W_j \subset X_1$ si $j \in J_1$ et $W_j \subset X_2$ si $j \in J_2$.
(C'est possible parce que les $W_j$ sont intègres.) Nous pouvons alors écrire
l'équation ci-dessus sous la forme
$$
\sum n_{1, i}[W_{1, i}] - \sum\nolimits_{j \in J_1} (i_j)_*\text{div}(f_j) =
- \sum n_{2, i}[W_{2, i}] + \sum\nolimits_{j \in J_2} (i_j)_*\text{div}(f_j)
$$
Cette expression est donc un cycle (!) sur $X_1 \cap X_2$. Autrement dit,
l'élément $(\alpha_1, \alpha_2)$ appartient à l'image de la première flèche,
ce qui achève la démonstration.
\end{proof}

\begin{example}
\label{chow-example-weird}
Voici un exemple « étrange ».
Supposons que $S$ soit le spectre d'un corps $k$,
avec $\delta$ comme dans l'Exemple \ref{chow-example-field}.
Supposons que $X = C_1 \cup C_2 \cup \ldots$ soit une union infinie
de courbes $C_j \cong \mathbf{P}^1_k$ recollées
de la manière suivante : le point $\infty \in C_j$ est recollé
transversalement au point $0 \in C_{j + 1}$ pour $j = 1, 2, 3, \ldots$.
Prenons le point $0 \in C_1$. Cela donne un zéro-cycle
$[0] \in Z_0(X)$. L'« étrangeté » de cette situation est
qu'en fait $[0] \sim_{rat} 0$ ! En effet, nous pouvons choisir
la fonction rationnelle $f_j \in R(C_j)$ comme étant la fonction
qui a un zéro simple en $0$, un pôle simple en $\infty$
et aucun autre zéro ni pôle. Nous voyons alors que la somme
$\sum (i_j)_*\text{div}(f_j)$ est exactement le $0$-cycle
$[0]$. En fait, il s'avère que $\CH_0(X) = 0$ dans cet exemple.
Si vous trouvez cela trop bizarre, il vous suffit de
veiller à ce que vos espaces soient toujours quasi-compacts
(ainsi $X$ n'existe même pas pour vous).
\end{example}

\begin{remark}
\label{chow-remark-infinite-sums-rational-equivalences}
Soit $(S, \delta)$ comme dans la Situation \ref{chow-situation-setup}.
Soit $X$ un schéma localement de type fini sur $S$.
Supposons données des familles infinies $\alpha_i, \beta_i \in Z_k(X)$,
$i \in I$, de $k$-cycles sur $X$. Supposons que les supports
des $\alpha_i$ et des $\beta_i$ forment des familles localement finies
de fermés de $X$, de sorte que $\sum \alpha_i$
et $\sum \beta_i$ soient définies comme cycles. Supposons de plus que
$\alpha_i \sim_{rat} \beta_i$ pour tout $i$. Il n'est alors pas
évident que $\sum \alpha_i \sim_{rat} \sum \beta_i$. En effet,
le problème est que les équivalences rationnelles peuvent être
données par des
familles localement finies $\{W_{i, j}, f_{i, j} \in R(W_{i, j})^*\}_{j \in J_i}$,
mais que l'union $\{W_{i, j}\}_{i \in I, j\in J_i}$ peut ne pas
être localement finie.

\medskip\noindent
Dans de nombreux cas pratiques, on dispose d'une famille localement finie de fermés
$\{T_i\}_{i \in I}$ telle que $\alpha_i, \beta_i$
soient supportés sur $T_i$ et telle que $\alpha_i = \beta_i$
dans $\CH_k(T_i)$ ; autrement dit, les familles
$\{W_{i, j}, f_{i, j} \in R(W_{i, j})^*\}_{j \in J_i}$
sont constituées de sous-schémas $W_{i, j} \subset T_i$. Dans ce cas, il est vrai que
$\sum \alpha_i \sim_{rat} \sum \beta_i$ sur $X$, simplement parce que
la famille $\{W_{i, j}\}_{i \in I, j\in J_i}$ est alors automatiquement
localement finie.
\end{remark}






\section{Équivalence rationnelle et images directe et réciproque}
\label{chow-section-properties-rational-equivalence}

\noindent
Dans cette section, nous montrons que l'image réciproque plate et l'image directe propre
commutent avec l'équivalence rationnelle.

\begin{lemma}
\label{chow-lemma-prepare-flat-pullback-rational-equivalence}
Soit $(S, \delta)$ comme dans la Situation \ref{chow-situation-setup}.
Soient $X$, $Y$ des schémas localement de type fini sur $S$.
Supposons $Y$ intègre avec $\dim_\delta(Y) = k$.
Soit $f : X \to Y$ un morphisme plat de
dimension relative $r$. Alors, pour $g \in R(Y)^*$, nous avons
$$
f^*\text{div}_Y(g) =
\sum n_j i_{j, *}\text{div}_{X_j}(g \circ f|_{X_j})
$$
comme $(k + r - 1)$-cycles sur $X$, où la somme porte sur les composantes
irréductibles $X_j$ de $X$ et $n_j$ est la multiplicité de $X_j$ dans $X$.
\end{lemma}

\begin{proof}
Soit $Z \subset X$ un sous-schéma fermé intègre de $\delta$-dimension
$k + r - 1$. Nous devons montrer que le coefficient $n$ de $[Z]$ dans
$f^*\text{div}(g)$ est égal au coefficient
$m$ de $[Z]$ dans $\sum i_{j, *} \text{div}(g \circ f|_{X_j})$.
Soit $Z'$ l'adhérence de $f(Z)$, qui est un sous-schéma fermé
intègre de $Y$. D'après le Lemme \ref{chow-lemma-flat-inverse-image-dimension},
nous avons $\dim_\delta(Z') \geq k - 1$. Ainsi, ou bien $Z' = Y$,
ou bien $Z'$ est un diviseur premier sur $Y$. Si $Z' = Y$, alors les coefficients
$n$ et $m$ sont tous deux nuls : c'est clair pour $n$ par définition
de $f^*$, et cela résulte pour $m$ du fait que $g \circ f|_{X_j}$ est
une unité en tout point de $X_j$ envoyé sur le point générique de $Y$.
Nous pouvons donc supposer que $Z' \subset Y$ est un diviseur premier.

\medskip\noindent
Nous allons traduire en algèbre l'égalité de $n$ et $m$.
En effet, soient $\xi' \in Z'$ et $\xi \in Z$ les points génériques.
Posons $A = \mathcal{O}_{Y, \xi'}$ et $B = \mathcal{O}_{X, \xi}$.
Notons que $A$, $B$ sont noethériens, que $A \to B$ est plat et local,
que $A$ est un domaine et que $\mathfrak m_AB$ est un idéal de définition
de l'anneau local $B$. La fonction rationnelle $g$ est un élément
du corps des fractions $Q(A)$ de $A$.
Par construction, les sous-schémas fermés $X_j$
qui rencontrent $\xi$ correspondent $1$-à-$1$ aux idéaux premiers minimaux
$$
\mathfrak q_1, \ldots, \mathfrak q_s \subset B
$$
Les entiers $n_j$ sont les longueurs correspondantes
$$
n_i = \text{longueur}_{B_{\mathfrak q_i}}(B_{\mathfrak q_i})
$$
Les fonctions rationnelles $g \circ f|_{X_j}$ correspondent à l'image
$g_i \in \kappa(\mathfrak q_i)^*$ de $g \in Q(A)$.
En rassemblant tout cela, nous voyons que
$$
n = \text{ord}_A(g) \text{longueur}_B(B/\mathfrak m_AB)
$$
et que
$$
m = \sum \text{ord}_{B/\mathfrak q_i}(g_i)
\text{longueur}_{B_{\mathfrak q_i}}(B_{\mathfrak q_i})
$$
En écrivant $g = x/y$ pour des éléments non nuls $x, y \in A$, nous voyons qu'il suffit
de démontrer
$$
\text{longueur}_A(A/(x)) \text{longueur}_B(B/\mathfrak m_AB) =
\text{longueur}_B(B/xB)
$$
(l'égalité utilise Algèbre, Lemme \ref{algebra-lemma-pullback-module})
égale à
$$
\sum\nolimits_{i = 1, \ldots, s}
\text{longueur}_{B/\mathfrak q_i}(B/(x, \mathfrak q_i))
\text{longueur}_{B_{\mathfrak q_i}}(B_{\mathfrak q_i})
$$
et de même pour $y$. Comme $A \to B$ est plat, il s'ensuit que $x$
est un non-diviseur de zéro dans $B$. L'égalité souhaitée résulte alors du
Lemme \ref{chow-lemma-additivity-divisors-restricted}.
\end{proof}

\begin{lemma}
\label{chow-lemma-flat-pullback-rational-equivalence}
Soit $(S, \delta)$ comme dans la Situation \ref{chow-situation-setup}.
Soient $X$, $Y$ des schémas localement de type fini sur $S$.
Soit $f : X \to Y$ un morphisme plat de dimension relative $r$.
Soient $\alpha \sim_{rat} \beta$ des $k$-cycles rationnellement équivalents sur $Y$.
Alors $f^*\alpha \sim_{rat} f^*\beta$ comme $(k + r)$-cycles sur $X$.
\end{lemma}

\begin{proof}
Que devons-nous montrer ? Supposons que l'on nous donne une famille
$$
i_j : W_j \longrightarrow Y
$$
d'immersions fermées, avec chaque $W_j$ intègre de $\delta$-dimension $k + 1$,
et des fonctions rationnelles $g_j \in R(W_j)^*$. Supposons de plus que
la famille $\{i_j(W_j)\}_{j \in J}$ soit localement finie sur $Y$.
Nous devons alors montrer que
$$
f^*(\sum i_{j, *}\text{div}(g_j)) = \sum f^*i_{j, *}\text{div}(g_j)
$$
est rationnellement équivalent à zéro sur $X$. La somme de droite
a un sens car $\{W_j\}$ est localement finie dans $X$ d'après le
Lemme \ref{chow-lemma-inverse-image-locally-finite}.

\medskip\noindent
Considérons les produits fibrés
$$
i'_j : W'_j = W_j \times_Y X \longrightarrow X.
$$
et notons $f_j : W'_j \to W_j$ la première projection.
D'après le Lemme \ref{chow-lemma-flat-pullback-proper-pushforward},
nous pouvons écrire la somme ci-dessus sous la forme
$$
\sum i'_{j, *}(f_j^*\text{div}(g_j))
$$
D'après le Lemme \ref{chow-lemma-prepare-flat-pullback-rational-equivalence},
nous voyons que chaque $f_j^*\text{div}(g_j)$ est rationnellement équivalent
à zéro sur $W'_j$. Ainsi, chaque $i'_{j, *}(f_j^*\text{div}(g_j))$
est rationnellement équivalent à zéro. Il en va alors de même pour
la somme affichée d'après la discussion de la
Remarque \ref{chow-remark-infinite-sums-rational-equivalences}.
\end{proof}

\begin{lemma}
\label{chow-lemma-proper-pushforward-rational-equivalence}
Soit $(S, \delta)$ comme dans la Situation \ref{chow-situation-setup}.
Soient $X$, $Y$ des schémas localement de type fini sur $S$.
Soit $p : X \to Y$ un morphisme propre.
Supposons que $\alpha, \beta \in Z_k(X)$ soient rationnellement équivalents.
Alors $p_*\alpha$ est rationnellement équivalent à $p_*\beta$.
\end{lemma}

\begin{proof}
Que devons-nous montrer ? Supposons que l'on nous donne une famille
$$
i_j : W_j \longrightarrow X
$$
d'immersions fermées, avec chaque $W_j$ intègre de $\delta$-dimension $k + 1$,
et des fonctions rationnelles $f_j \in R(W_j)^*$.
Supposons de plus que
la famille $\{i_j(W_j)\}_{j \in J}$ soit localement finie sur $X$.
Nous devons alors montrer que
$$
p_*\left(\sum i_{j, *}\text{div}(f_j)\right)
$$
est rationnellement équivalent à zéro sur $X$.

\medskip\noindent
Notons que la somme est égale à
$$
\sum p_*i_{j, *}\text{div}(f_j).
$$
Soit $W'_j \subset Y$ le sous-schéma fermé intègre qui est l'
image de $p \circ i_j$. La famille $\{W'_j\}$ est localement finie
dans $Y$ d'après le Lemme \ref{chow-lemma-quasi-compact-locally-finite}.
Il suffit donc de montrer, pour un $j$ donné, que soit
$p_*i_{j, *}\text{div}(f_j) = 0$, soit il
est égal à $i'_{j, *}\text{div}(g_j)$ pour un certain $g_j \in R(W'_j)^*$.

\medskip\noindent
Les arguments ci-dessus nous ramènent donc au cas d'un seul
sous-schéma fermé intègre $W \subset X$ de $\delta$-dimension $k + 1$.
Soit $f \in R(W)^*$. Soit $W' = p(W)$ comme ci-dessus.
Nous obtenons un diagramme commutatif de morphismes
$$
\xymatrix{
W \ar[r]_i \ar[d]_{p'} & X \ar[d]^p \\
W' \ar[r]^{i'} & Y
}
$$
Notons que $p_*i_*\text{div}(f) = i'_*(p')_*\text{div}(f)$ d'après le
Lemme \ref{chow-lemma-compose-pushforward}. Comme expliqué ci-dessus,
nous devons montrer que $(p')_*\text{div}(f)$
est le diviseur d'une fonction rationnelle sur $W'$, ou zéro.
Il y a trois cas à distinguer.

\medskip\noindent
Le cas $\dim_\delta(W') < k$. Dans ce cas, automatiquement,
$(p')_*\text{div}(f) = 0$ et il n'y a rien à démontrer.

\medskip\noindent
Le cas $\dim_\delta(W') = k$. Montrons que $(p')_*\text{div}(f) = 0$
dans ce cas. Soit $\eta \in W'$ le point générique.
Notons que $c : W_\eta \to \Spec(K)$
est une courbe intègre propre sur $K = \kappa(\eta)$
dont le corps des fonctions $K(W_\eta)$ s'identifie à $R(W)$.
Voici un diagramme
$$
\xymatrix{
W_\eta \ar[r] \ar[d]_c & W \ar[d]^{p'} \\
\Spec(K) \ar[r] & W'
}
$$
Notons $f_\eta \in K(W_\eta)^*$ la fonction rationnelle
correspondant à $f \in R(W)^*$.
De plus, les points fermés $\xi$ de $W_\eta$ correspondent $1 - 1$ aux
sous-schémas intègres fermés $Z = Z_\xi \subset W$ de $\delta$-dimension $k$
avec $p'(Z) = W'$. Notons que la multiplicité
de $Z_\xi$ dans $\text{div}(f)$ est égale à
$\text{ord}_{\mathcal{O}_{W_\eta, \xi}}(f_\eta)$, simplement parce que les
anneaux locaux $\mathcal{O}_{W_\eta, \xi}$ et $\mathcal{O}_{W, \xi}$
sont identifiés (comme sous-anneaux de leurs corps des fractions).
Nous voyons donc que la multiplicité de $[W']$ dans
$(p')_*\text{div}(f)$ est égale à la multiplicité de
$[\Spec(K)]$ dans $c_*\text{div}(f_\eta)$.
D'après le Lemme \ref{chow-lemma-curve-principal-divisor}, celle-ci est nulle.

\medskip\noindent
Le cas $\dim_\delta(W') = k + 1$. Dans ce cas, le
Lemme \ref{chow-lemma-proper-pushforward-alteration} s'applique,
et nous voyons qu'en effet $p'_*\text{div}(f) = \text{div}(g)$
pour un certain $g \in R(W')^*$, comme souhaité.
\end{proof}















\section{Équivalence rationnelle et droite projective}
\label{chow-section-different-rational-equivalence}

\noindent
Soit $(S, \delta)$ comme dans la Situation \ref{chow-situation-setup}.
Soit $X$ un schéma localement de type fini sur $S$.
Étant donné un sous-schéma fermé quelconque
$Z \subset X \times_S \mathbf{P}^1_S = X \times \mathbf{P}^1$,
nous notons $Z_0$, resp.\ $Z_\infty$, le sous-schéma fermé schématique
$Z_0 = \text{pr}_2^{-1}(D_0)$,
resp.\ $Z_\infty = \text{pr}_2^{-1}(D_\infty)$.
Ici $D_0$, $D_\infty$ sont comme dans (\ref{chow-equation-zero-infty}).

\begin{lemma}
\label{chow-lemma-rational-equivalence-family}
Soit $(S, \delta)$ comme dans la Situation \ref{chow-situation-setup}.
Soit $X$ un schéma localement de type fini sur $S$.
Soit $W \subset X \times_S \mathbf{P}^1_S$ un sous-schéma fermé
intègre de $\delta$-dimension $k + 1$.
Supposons $W \not = W_0$ et $W \not = W_\infty$. Alors
\begin{enumerate}
\item $W_0$, $W_\infty$ sont des diviseurs de Cartier effectifs de $W$,
\item $W_0$, $W_\infty$ peuvent être considérés comme des sous-schémas fermés
de $X$ et
$$
[W_0]_k \sim_{rat} [W_\infty]_k,
$$
\item pour toute famille localement finie de
sous-schémas fermés intègres
$W_i \subset X \times_S \mathbf{P}^1_S$
de $\delta$-dimension $k + 1$ avec $W_i \not = (W_i)_0$ et
$W_i \not = (W_i)_\infty$, nous avons
$\sum ([(W_i)_0]_k - [(W_i)_\infty]_k) \sim_{rat} 0$
sur $X$, et
\item pour tout $\alpha \in Z_k(X)$ avec $\alpha \sim_{rat} 0$,
il existe une famille localement finie de
sous-schémas fermés intègres $W_i \subset X \times_S \mathbf{P}^1_S$
comme ci-dessus telle que $\alpha = \sum ([(W_i)_0]_k - [(W_i)_\infty]_k)$.
\end{enumerate}
\end{lemma}

\begin{proof}
La partie (1) résulte du
Lemme \ref{divisors-lemma-pullback-effective-Cartier-defined} de Diviseurs,
puisque, par hypothèse, le point générique
de $W$ n'est envoyé ni dans $D_0$ ni dans $D_\infty$ par la projection
$X \times_S \mathbf{P}^1_S \to \mathbf{P}^1_S$.

\medskip\noindent
Comme $X \times_S D_0 \to X$ est une immersion fermée, nous voyons que $W_0$
est isomorphe à un sous-schéma fermé de $X$. Il en va de même pour $W_\infty$.
Le morphisme $p : W \to X$ est propre comme composée de
l'immersion fermée $W \to X \times_S \mathbf{P}^1_S$ et du
morphisme propre $X \times_S \mathbf{P}^1_S \to X$. D'après le
Lemme \ref{chow-lemma-rational-function}, nous avons
$[W_0]_k \sim_{rat} [W_\infty]_k$ comme cycles sur $W$. Ainsi la partie (2) résulte du
Lemme \ref{chow-lemma-proper-pushforward-rational-equivalence}, puisque manifestement
$p_*[W_0]_k = [W_0]_k$ et de même pour $W_\infty$.

\medskip\noindent
Le seul contenu de l'assertion (3), compte tenu des parties (1) et (2), est que
la famille $\{(W_i)_0, (W_i)_\infty\}$ est une famille localement finie
de sous-schémas fermés de $X$. C'est clair.

\medskip\noindent
Supposons que $\alpha \sim_{rat} 0$.
Par définition, cela signifie qu'il existe des sous-schémas fermés intègres
$V_i \subset X$ de $\delta$-dimension $k + 1$ et des fonctions rationnelles
$f_i \in R(V_i)^*$ tels que la famille
$\{V_i\}_{i \in I}$ soit localement finie dans $X$ et tels que
$\alpha = \sum (V_i \to X)_*\text{div}(f_i)$.
Soit
$$
W_i \subset V_i \times_S \mathbf{P}^1_S \subset X \times_S \mathbf{P}^1_S
$$
l'adhérence du graphe de l'application rationnelle $f_i$, comme dans le
Lemme \ref{chow-lemma-rational-function}.
Alors $(V_i \to X)_*\text{div}(f_i)$
est égal à $[(W_i)_0]_k - [(W_i)_\infty]_k$ d'après ce même lemme.
Le résultat est donc clair.
\end{proof}

\begin{lemma}
\label{chow-lemma-closed-subscheme-cross-p1}
Soit $(S, \delta)$ comme dans la Situation \ref{chow-situation-setup}.
Soit $X$ un schéma localement de type fini sur $S$.
Soit $Z$ un sous-schéma fermé de $X \times \mathbf{P}^1$.
Supposons
\begin{enumerate}
\item $\dim_\delta(Z) \leq k + 1$,
\item $\dim_\delta(Z_0) \leq k$, $\dim_\delta(Z_\infty) \leq k$, et
\item pour tout point immergé $\xi$ (Diviseurs, Définition
\ref{divisors-definition-embedded}) de $Z$, soit
$\xi \not \in Z_0 \cup Z_\infty$, soit $\delta(\xi) < k$.
\end{enumerate}
Alors $[Z_0]_k \sim_{rat} [Z_\infty]_k$ comme $k$-cycles sur $X$.
\end{lemma}

\begin{proof}
Soit $\{W_i\}_{i \in I}$ la famille des composantes
irréductibles de $Z$ qui sont de $\delta$-dimension $k + 1$.
Écrivons
$$
[Z]_{k + 1} = \sum n_i[W_i]
$$
avec $n_i > 0$, conformément à la définition. Notons que $\{W_i\}$
est une famille localement finie de fermés de
$X \times_S \mathbf{P}^1_S$ d'après le
Lemme \ref{divisors-lemma-components-locally-finite} de Diviseurs.
Nous affirmons que
$$
[Z_0]_k = \sum n_i[(W_i)_0]_k
$$
et de même pour $[Z_\infty]_k$. Si nous le démontrons, alors le lemme
résulte du Lemme \ref{chow-lemma-rational-equivalence-family}.

\medskip\noindent
Soit $Z' \subset X$ un sous-schéma fermé intègre de $\delta$-dimension $k$.
Pour démontrer l'égalité ci-dessus, il suffit de montrer que le coefficient $n$
de $[Z']$ dans $[Z_0]_k$ est égal au coefficient $m$ de
$[Z']$ dans $\sum n_i[(W_i)_0]_k$. Soit $\xi' \in Z'$ le point générique.
Posons $\xi = (\xi', 0) \in  X \times_S \mathbf{P}^1_S$.
Considérons l'anneau local $A = \mathcal{O}_{X \times_S \mathbf{P}^1_S, \xi}$.
Soit $I \subset A$ l'idéal définissant $Z$, autrement dit tel que
$A/I = \mathcal{O}_{Z, \xi}$. Soit $t \in A$ l'élément définissant
$X \times_S D_0$ (c'est-à-dire le relevé de la coordonnée de $\mathbf{P}^1$ en zéro).
Par notre choix de $\xi' \in Z'$, nous avons $\delta(\xi) = k$,
et donc $\dim(A/I) = 1$. Puisque $\xi$ n'est pas un point immergé d'après
l'hypothèse (3), nous voyons que $A/I$ est de Cohen-Macaulay. Puisque $\dim_\delta(Z_0)
= k$, nous voyons que $\dim(A/(t, I)) = 0$, ce qui implique que $t$
est un non-diviseur de zéro dans $A/I$. Enfin, les sous-schémas fermés irréductibles
$W_i$ passant par $\xi$ correspondent aux idéaux premiers minimaux
$I \subset \mathfrak q_i$ au-dessus de $I$. Les multiplicités $n_i$ correspondent
aux longueurs $\text{longueur}_{A_{\mathfrak q_i}}(A/I)_{\mathfrak q_i}$.
Nous voyons donc que
$$
n = \text{longueur}_A(A/(t, I))
$$
et
$$
m = \sum
\text{longueur}_A(A/(t, \mathfrak q_i))
\text{longueur}_{A_{\mathfrak q_i}}(A/I)_{\mathfrak q_i}
$$
Le résultat résulte donc du
Lemme \ref{chow-lemma-additivity-divisors-restricted}.
\end{proof}

\begin{lemma}
\label{chow-lemma-coherent-sheaf-cross-p1}
Soit $(S, \delta)$ comme dans la Situation \ref{chow-situation-setup}.
Soit $X$ un schéma localement de type fini sur $S$.
Soit $\mathcal{F}$ un faisceau cohérent sur $X \times \mathbf{P}^1$.
Soient $i_0, i_\infty : X \to X \times \mathbf{P}^1$ les immersions fermées
telles que $i_t(x) = (x, t)$. Notons $\mathcal{F}_0 = i_0^*\mathcal{F}$ et
$\mathcal{F}_\infty = i_\infty^*\mathcal{F}$.
Supposons
\begin{enumerate}
\item $\dim_\delta(\text{Supp}(\mathcal{F})) \leq k + 1$,
\item $\dim_\delta(\text{Supp}(\mathcal{F}_0)) \leq k$,
$\dim_\delta(\text{Supp}(\mathcal{F}_\infty)) \leq k$, et
\item pour tout point associé immergé $\xi$ de $\mathcal{F}$, soit
$\xi \not \in (X \times \mathbf{P}^1)_0 \cup (X \times \mathbf{P}^1)_\infty$,
soit $\delta(\xi) < k$.
\end{enumerate}
Alors $[\mathcal{F}_0]_k \sim_{rat} [\mathcal{F}_\infty]_k$ comme $k$-cycles sur $X$.
\end{lemma}

\begin{proof}
Soit $\{W_i\}_{i \in I}$ la famille des composantes
irréductibles de $\text{Supp}(\mathcal{F})$
qui sont de $\delta$-dimension $k + 1$.
Écrivons
$$
[\mathcal{F}]_{k + 1} = \sum n_i[W_i]
$$
avec $n_i > 0$, conformément à la définition. Notons que $\{W_i\}$
est une famille localement finie de fermés de
$X \times_S \mathbf{P}^1_S$ d'après le Lemme \ref{chow-lemma-length-finite}.
Nous affirmons que
$$
[\mathcal{F}_0]_k = \sum n_i[(W_i)_0]_k
$$
et de même pour $[\mathcal{F}_\infty]_k$. Si nous le démontrons, alors le lemme
résulte du Lemme \ref{chow-lemma-rational-equivalence-family}.

\medskip\noindent
Soit $Z' \subset X$ un sous-schéma fermé intègre de $\delta$-dimension $k$.
Pour démontrer l'égalité ci-dessus, il suffit de montrer que le coefficient $n$
de $[Z']$ dans $[\mathcal{F}_0]_k$ est égal au coefficient $m$ de
$[Z']$ dans $\sum n_i[(W_i)_0]_k$. Soit $\xi' \in Z'$ le point générique.
Posons $\xi = (\xi', 0) \in  X \times_S \mathbf{P}^1_S$.
Considérons l'anneau local $A = \mathcal{O}_{X \times_S \mathbf{P}^1_S, \xi}$.
Considérons $M = \mathcal{F}_\xi$ comme un $A$-module.
Soit $t \in A$ l'élément définissant $X \times_S D_0$
(c'est-à-dire le relevé de la coordonnée de $\mathbf{P}^1$ en zéro).
Par notre choix de $\xi' \in Z'$, nous avons $\delta(\xi) = k$,
et donc $\dim(\text{Supp}(M)) = 1$. Puisque $\xi$ n'est pas un point associé
de $\mathcal{F}$ d'après l'hypothèse (3), nous voyons que $M$ est un module de Cohen-Macaulay.
Puisque $\dim_\delta(\text{Supp}(\mathcal{F}_0)) = k$,
nous voyons que $\dim(\text{Supp}(M/tM)) = 0$, ce qui implique que $t$
est un non-diviseur de zéro dans $M$. Enfin, les sous-schémas fermés irréductibles
$W_i$ passant par $\xi$ correspondent aux idéaux premiers minimaux
$\mathfrak q_i$ de $\text{Ass}(M)$. Les multiplicités $n_i$ correspondent
aux longueurs $\text{longueur}_{A_{\mathfrak q_i}}M_{\mathfrak q_i}$.
Nous voyons donc que
$$
n = \text{longueur}_A(M/tM)
$$
et
$$
m = \sum
\text{longueur}_A(A/(t, \mathfrak q_i)A)
\text{longueur}_{A_{\mathfrak q_i}}M_{\mathfrak q_i}
$$
Le résultat résulte donc du
Lemme \ref{chow-lemma-additivity-divisors-restricted}.
\end{proof}







\section{Groupes de Chow et enveloppes}
\label{chow-section-envelopes}

\noindent
Voici la définition.

\begin{definition}
\label{chow-definition-envelope}
\begin{reference}
\cite[Definition 18.3]{F}
\end{reference}
Soit $X$ un schéma. Une {\it enveloppe} est un morphisme propre $f : Y \to X$
qui est complètement décomposé
(Plus sur les morphismes, Définition \ref{more-morphisms-definition-cd-morphism}).
\end{definition}

\noindent
La suite exacte du Lemme \ref{chow-lemma-envelope}
est la motivation principale de la définition.

\begin{lemma}
\label{chow-lemma-composition-envelope}
Soit $(S, \delta)$ comme dans la Situation \ref{chow-situation-setup}.
Soit $X$ un schéma localement de type fini sur $S$.
Si $f : Y \to X$ et $g : Z \to Y$ sont des enveloppes, alors
$f \circ g$ est une enveloppe.
\end{lemma}

\begin{proof}
Ceci résulte du Lemme \ref{morphisms-lemma-composition-proper} de Morphismes
et du Lemme \ref{more-morphisms-lemma-composition-cd} de Plus sur les morphismes.
\end{proof}

\begin{lemma}
\label{chow-lemma-base-change-envelope}
Soit $(S, \delta)$ comme dans la Situation \ref{chow-situation-setup}.
Soit $X' \to X$ un morphisme de schémas localement de type fini sur $S$.
Si $f : Y \to X$ est une enveloppe, alors le changement de base $f' : Y' \to X'$
de $f$ est aussi une enveloppe.
\end{lemma}

\begin{proof}
Ceci résulte du Lemme \ref{morphisms-lemma-base-change-proper} de Morphismes
et du Lemme \ref{more-morphisms-lemma-base-change-cd} de Plus sur les morphismes.
\end{proof}

\begin{lemma}
\label{chow-lemma-envelope}
Soit $(S, \delta)$ comme dans la Situation \ref{chow-situation-setup}.
Soit $X$ un schéma localement de type fini sur $S$.
Soit $f : Y \to X$ une enveloppe. Alors
nous avons une suite exacte
$$
\CH_k(Y \times_X Y) \xrightarrow{p_* - q_*}
\CH_k(Y) \xrightarrow{f_*}
\CH_k(X) \to 0
$$
pour tout $k \in \mathbf{Z}$. Ici $p, q : Y \times_X Y \to Y$ sont
les projections.
\end{lemma}

\begin{proof}
Puisque $f$ est une enveloppe, $f$ est propre, et l'image directe des
cycles et des classes de cycles est donc définie ; voir les
sections \ref{chow-section-proper-pushforward} et \ref{chow-section-push-pull}.
De même, les morphismes $p$ et $q$ sont propres comme changements de base de $f$.
La composée des flèches est nulle, car
$f_* \circ p_* = (p \circ f)_* = (q \circ f)_* = f_* \circ q_*$ ; voir le
Lemme \ref{chow-lemma-compose-pushforward}.

\medskip\noindent
Montrons que $f_* : Z_k(Y) \to Z_k(X)$ est surjectif.
En effet, supposons que nous ayons $\alpha = \sum n_i[Z_i] \in Z_k(X)$,
où les $Z_i \subset X$ forment une famille localement finie de sous-schémas
fermés intègres. Soit $x_i \in Z_i$ le point générique.
Puisque $f$ est une enveloppe et donc complètement décomposé,
il existe un point $y_i \in Y$ avec $f(y_i) = x_i$
et avec $\kappa(y_i)/\kappa(x_i)$ triviale. Soit $W_i \subset Y$
le sous-schéma fermé intègre de point générique $y_i$.
Comme $f$ est fermé, nous voyons que $f(W_i) = Z_i$.
Il en résulte que la famille de sous-schémas fermés $W_i$ est localement finie
sur $Y$. Comme $\kappa(y_i)/\kappa(x_i)$ est triviale, nous voyons
que $\dim_\delta(W_i) = \dim_\delta(Z_i) = k$. Ainsi
$\beta = \sum n_i[W_i]$ appartient à $Z_k(Y)$. Enfin, puisque
$\kappa(y_i)/\kappa(x_i)$ est triviale, le degré du morphisme dominant
$f|_{W_i} : W_i \to Z_i$ est $1$, et nous concluons
que $f_*\beta = \alpha$.

\medskip\noindent
Puisque $f_* : Z_k(Y) \to Z_k(X)$ est surjectif, l'application
$f_* : \CH_k(Y) \to \CH_k(X)$ est a fortiori surjective.

\medskip\noindent
Soit $\beta \in Z_k(Y)$ un élément tel que $f_*\beta$ soit nul dans
$\CH_k(X)$. Cela signifie que nous pouvons trouver une famille localement finie de
sous-schémas fermés intègres $Z_j \subset X$ avec $\dim_\delta(Z_j) = k + 1$
et des $f_j \in R(Z_j)^*$ tels que
$$
f_*\beta = \sum (Z_j \to X)_*\text{div}(f_j)
$$
comme cycles, où $i_j : Z_j \to X$ est l'immersion fermée donnée.
En raisonnant exactement comme ci-dessus, nous pouvons trouver une famille
localement finie de sous-schémas fermés intègres $W_j \subset Y$
avec $f(W_j) = Z_j$ et telle que $W_j \to Z_j$ soit birationnel, c'est-à-dire
induise un isomorphisme $R(Z_j) = R(W_j)$. Notons $g_j \in R(W_j)^*$
l'élément correspondant à $f_j$. Observons que $W_j \to Z_j$
est propre et que $(W_j \to Z_j)_*\text{div}(g_j) = \text{div}(f_j)$
comme cycles sur $Z_j$. Il s'ensuit que, si nous remplaçons
$\beta$ par le cycle rationnellement équivalent
$$
\beta' = \beta - \sum (W_j \to Y)_*\text{div}(g_j)
$$
nous obtenons $f_*\beta' = 0$.
(On utilise ici le Lemme \ref{chow-lemma-compose-pushforward}.)
Ainsi, pour achever la démonstration
du lemme, il suffit de démontrer l'affirmation du paragraphe suivant.

\medskip\noindent
Affirmation : si $\beta \in Z_k(Y)$ et $f_*\beta = 0$, alors
$\beta = \delta + p_*\gamma - q_*\gamma$ dans $Z_k(Y)$ pour un certain
$\gamma \in Z_k(Y \times_X Y)$. En effet, écrivons $\beta = \sum_{j \in J} n_j[W_j]$,
où $\{W_j\}_{j \in J}$ est une famille localement finie de sous-schémas fermés
intègres de $Y$ avec $\dim_\delta(W_j) = k$.
Fixons un sous-schéma fermé intègre $Z \subset X$. Considérons le sous-ensemble
$J_Z = \{j \in J : f(W_j) = Z\}$. C'est un ensemble fini. Il y a trois
cas :
\begin{enumerate}
\item $J_Z = \emptyset$. Dans ce cas, nous posons $\gamma_Z = 0$.
\item $J_Z \not = \emptyset$ et $\dim_\delta(Z) = k$.
La condition $f_*\beta = 0$ implique, en considérant le
coefficient de $Z$, que $\sum_{j \in J_Z} n_j\deg(W_j/Z) = 0$.
Dans ce cas, nous choisissons un sous-schéma fermé intègre $W \subset Y$
qui s'envoie birationnellement sur $Z$ (voir ci-dessus). En considérant les points
génériques, nous voyons que $W_j \times_Z W$ possède une unique composante
irréductible $W'_j \subset W_j \times_Z W \subset Y \times_X Y$
qui s'envoie birationnellement sur $W_j$. Alors $W'_j \to W$ est dominant
et $\deg(W'_j/W) = \deg(W_j/W)$. Ainsi, si nous posons
$\gamma_Z = \sum_{j \in J_Z} n_j[W'_j]$,
alors nous voyons que
$p_*\gamma_Z = \sum_{j \in J_Z} n_j[W_j]$ et
$q_*\gamma_Z = \sum_{j \in J_Z} n_j\deg(W'_j/W)[W] = 0$.
\item $J_Z \not = \emptyset$ et $\dim_\delta(Z) < k$.
Dans ce cas, nous choisissons un sous-schéma fermé intègre $W \subset Y$
qui s'envoie birationnellement sur $Z$ (voir ci-dessus). En considérant les points
génériques, nous voyons que $W_j \times_Z W$ possède une unique composante
irréductible $W'_j \subset W_j \times_Z W \subset Y \times_X Y$
qui s'envoie birationnellement sur $W_j$. Alors $W'_j \to W$ est dominant
et $k = \dim_\delta(W'_j) > \dim_\delta(W) = \dim_\delta(Z)$.
Ainsi, si nous posons $\gamma_Z = \sum_{j \in J_Z} n_j[W'_j]$,
alors nous voyons que
$p_*\gamma_Z = \sum_{j \in J_Z} n_j[W_j]$ et
$q_*\gamma_Z = 0$.
\end{enumerate}
Puisque la famille de sous-schémas fermés intègres $\{f(W_j)\}$
est localement finie sur $X$
(Lemme \ref{chow-lemma-quasi-compact-locally-finite}),
nous voyons que le $k$-cycle
$$
\gamma = \sum\nolimits_{Z \subset X\text{ fermé intègre}} \gamma_Z
$$
sur $Y \times_X Y$ est bien défini. Nos calculs ci-dessus montrent que
$p_*\gamma_Z = \beta$ et $q_*\gamma_Z = 0$, ce qui implique
ce que nous voulions démontrer.
\end{proof}













\section{Groupes de Chow et groupes K}
\label{chow-section-chow-and-K}

\noindent
Dans cette section, nous allons comparer le $K_0$ de la
catégorie des faisceaux cohérents aux groupes de Chow.

\medskip\noindent
Soit $(S, \delta)$ comme dans la Situation \ref{chow-situation-setup}.
Soit $X$ un schéma localement de type fini sur $S$.
Nous notons $\textit{Coh}(X) = \textit{Coh}(\mathcal{O}_X)$
la catégorie des faisceaux cohérents sur $X$.
C'est une catégorie abélienne ; voir le
Lemme \ref{coherent-lemma-coherent-abelian-Noetherian} de Cohomologie des schémas.
Pour tout $k \in \mathbf{Z}$, nous notons $\textit{Coh}_{\leq k}(X)$
la sous-catégorie pleine de $\textit{Coh}(X)$
formée des faisceaux cohérents $\mathcal{F}$
tels que $\dim_\delta(\text{Supp}(\mathcal{F})) \leq k$.

\begin{lemma}
\label{chow-lemma-Serre-subcategories}
Soit $(S, \delta)$ comme dans la Situation \ref{chow-situation-setup}.
Soit $X$ un schéma localement de type fini sur $S$.
Les catégories $\textit{Coh}_{\leq k}(X)$ sont des sous-catégories de Serre
de la catégorie abélienne $\textit{Coh}(X)$.
\end{lemma}

\begin{proof}
La définition d'une sous-catégorie de Serre est la
Définition \ref{homology-definition-serre-subcategory} d'Homologie.
La démonstration du lemme est immédiate et omise.
\end{proof}

\begin{lemma}
\label{chow-lemma-cycles-k-group}
Soit $(S, \delta)$ comme dans la Situation \ref{chow-situation-setup}.
Soit $X$ un schéma localement de type fini sur $S$.
Les applications
$$
Z_k(X)
\longrightarrow
K_0(\textit{Coh}_{\leq k}(X)/\textit{Coh}_{\leq k - 1}(X)),
\quad
\sum n_Z[Z] \mapsto
\left[\bigoplus\nolimits_{n_Z > 0} \mathcal{O}_Z^{\oplus n_Z}\right]
-
\left[\bigoplus\nolimits_{n_Z < 0} \mathcal{O}_Z^{\oplus -n_Z}\right]
$$
et
$$
K_0(\textit{Coh}_{\leq k}(X)/\textit{Coh}_{\leq k - 1}(X))
\longrightarrow
Z_k(X),\quad
\mathcal{F} \longmapsto [\mathcal{F}]_k
$$
sont des isomorphismes réciproques l'un de l'autre.
\end{lemma}

\begin{proof}
Notons que, si $\sum n_Z[Z]$ appartient à $Z_k(X)$, alors
les sommes directes
$\bigoplus\nolimits_{n_Z > 0} \mathcal{O}_Z^{\oplus n_Z}$ et
$\bigoplus\nolimits_{n_Z < 0} \mathcal{O}_Z^{\oplus -n_Z}$
sont des faisceaux cohérents sur $X$, puisque la famille $\{Z \mid n_Z > 0\}$
est localement finie sur $X$.
L'application $\mathcal{F} \to [\mathcal{F}]_k$ est additive
sur $\textit{Coh}_{\leq k}(X)$ ; voir le
Lemme \ref{chow-lemma-additivity-sheaf-cycle}. Et $[\mathcal{F}]_k = 0$
si $\mathcal{F} \in \textit{Coh}_{\leq k - 1}(X)$. D'après la partie (1)
du Lemme \ref{homology-lemma-serre-subcategory-K-groups} d'Homologie,
cela implique que la deuxième application est elle aussi bien définie.
Il est clair que la composée de la première application avec la deuxième
est l'identité.

\medskip\noindent
Réciproquement, supposons que nous partions d'un faisceau cohérent $\mathcal{F}$
sur $X$. Écrivons $[\mathcal{F}]_k = \sum_{i \in I} n_i[Z_i]$
avec $n_i > 0$ et $Z_i \subset X$, $i \in I$,
des sous-schémas fermés intègres deux à deux distincts de $\delta$-dimension $k$.
Nous devons montrer que
$$
[\mathcal{F}] = [\bigoplus\nolimits_{i \in I} \mathcal{O}_{Z_i}^{\oplus n_i}]
$$
dans $K_0(\textit{Coh}_{\leq k}(X)/\textit{Coh}_{\leq k - 1}(X))$.
Notons $\xi_i \in Z_i$ le point générique.
Si nous posons
$$
\mathcal{F}' = \Ker(\mathcal{F} \to \bigoplus \xi_{i, *}\mathcal{F}_{\xi_i})
$$
alors $\mathcal{F}'$ est le sous-module cohérent maximal de $\mathcal{F}$
dont le support est de dimension $\leq k - 1$. En particulier, $\mathcal{F}$
et $\mathcal{F}/\mathcal{F}'$ ont la même classe dans
$K_0(\textit{Coh}_{\leq k}(X)/\textit{Coh}_{\leq k - 1}(X))$.
Ainsi, après avoir remplacé $\mathcal{F}$ par $\mathcal{F}/\mathcal{F}'$,
nous pouvons supposer, et supposons, que le noyau $\mathcal{F}'$ affiché
ci-dessus est nul.

\medskip\noindent
Pour tout $i \in I$, choisissons une filtration
$$
\mathcal{F}_{\xi_i} = \mathcal{F}_i^0 \supset \mathcal{F}_i^1 \supset
\ldots \supset \mathcal{F}_i^{n_i} = 0
$$
telle que les quotients successifs soient de dimension $1$ sur le corps
résiduel en $\xi_i$. C'est possible puisque la longueur de $\mathcal{F}_{\xi_i}$
sur $\mathcal{O}_{X, \xi_i}$ est $n_i$.
Pour $p > n_i$, posons $\mathcal{F}_i^p = 0$. Pour $p \geq 0$, nous notons
$$
\mathcal{F}^p =
\Ker\left(\mathcal{F} \longrightarrow \bigoplus
\xi_{i, *}(\mathcal{F}_{\xi_i}/\mathcal{F}_i^p)\right)
$$
Alors $\mathcal{F}^p$ est cohérent, $\mathcal{F}^0 = \mathcal{F}$, et
$\mathcal{F}^p/\mathcal{F}^{p + 1}$ est isomorphe à un
$\mathcal{O}_{Z_i}$-module libre de rang $1$ (si $n_i > p$), ou à $0$
(si $n_i \leq p$), dans un voisinage ouvert de $\xi_i$. De plus,
$\mathcal{F}' = \bigcap \mathcal{F}^p = 0$. Puisque tout ouvert quasi-compact
$U \subset X$ ne contient qu'un nombre fini de $\xi_i$,
nous concluons que $\mathcal{F}^p|_U$ est nul pour $p \gg 0$.
Ainsi $\bigoplus_{p \geq 0} \mathcal{F}^p$ est un
$\mathcal{O}_X$-module cohérent. Considérons les suites exactes courtes
$$
0 \to
\bigoplus\nolimits_{p > 0} \mathcal{F}^p \to
\bigoplus\nolimits_{p \geq 0} \mathcal{F}^p \to
\bigoplus\nolimits_{p > 0} \mathcal{F}^p/\mathcal{F}^{p + 1} \to 0
$$
et
$$
0 \to
\bigoplus\nolimits_{p > 0} \mathcal{F}^p \to
\bigoplus\nolimits_{p \geq 0} \mathcal{F}^p \to
\mathcal{F} \to 0
$$
de $\mathcal{O}_X$-modules cohérents. Cela montre déjà que
$$
[\mathcal{F}] = [\bigoplus \mathcal{F}^p/\mathcal{F}^{p + 1}]
$$
dans $K_0(\textit{Coh}_{\leq k}(X)/\textit{Coh}_{\leq k - 1}(X))$.
Ensuite, pour tout $p \geq 0$ et tout $i \in I$ tel que $n_i > p$,
nous choisissons un faisceau d'idéaux non nul $\mathcal{I}_{i, p} \subset \mathcal{O}_{Z_i}$
et une application $\mathcal{I}_{i, p} \to \mathcal{F}^p/\mathcal{F}^{p + 1}$ sur $X$
qui est un isomorphisme sur le voisinage ouvert de $\xi_i$
mentionné ci-dessus. C'est possible d'après le
Lemme \ref{coherent-lemma-extend-coherent} de Cohomologie des schémas.
Nous considérons alors la suite exacte courte
$$
0 \to
\bigoplus\nolimits_{p \geq 0, i \in I, n_i > p} \mathcal{I}_{i, p}
\to
\bigoplus \mathcal{F}^p/\mathcal{F}^{p + 1} \to
\mathcal{Q} \to 0
$$
et la suite exacte courte
$$
0 \to
\bigoplus\nolimits_{p \geq 0, i \in I, n_i > p} \mathcal{I}_{i, p}
\to
\bigoplus\nolimits_{p \geq 0, i \in I, n_i > p} \mathcal{O}_{Z_i}
\to
\mathcal{Q}' \to 0
$$
Observons que $\mathcal{Q}$ et $\mathcal{Q}'$ sont tous deux nuls dans un voisinage
des points $\xi_i$ et qu'ils sont supportés sur $\bigcup Z_i$.
Ainsi $\mathcal{Q}$ et $\mathcal{Q}'$ appartiennent à
$\textit{Coh}_{\leq k - 1}(X)$.
Puisque
$$
\bigoplus\nolimits_{i \in I} \mathcal{O}_{Z_i}^{\oplus n_i} \cong
\bigoplus\nolimits_{p \geq 0, i \in I, n_i > p} \mathcal{O}_{Z_i}
$$
ceci achève la démonstration.
\end{proof}

\begin{lemma}
\label{chow-lemma-finite-cycles-k-group}
Soit $\pi : X \to Y$ un morphisme fini de schémas localement de type fini
sur $(S, \delta)$ comme dans la Situation \ref{chow-situation-setup}. Alors
$\pi_* : \textit{Coh}(X) \to \textit{Coh}(Y)$ est un foncteur exact
qui envoie $\textit{Coh}_{\leq k}(X)$ dans $\textit{Coh}_{\leq k}(Y)$
et induit des homomorphismes sur les $K_0$ de ces catégories et de
leurs quotients. Les applications du Lemme \ref{chow-lemma-cycles-k-group}
s'insèrent dans un diagramme commutatif
$$
\xymatrix{
Z_k(X) \ar[d]^{\pi_*} \ar[r] &
K_0(\textit{Coh}_{\leq k}(X)/\textit{Coh}_{\leq k - 1}(X))
\ar[d]^{\pi_*} \ar[r] &
Z_k(X) \ar[d]^{\pi_*} \\
Z_k(Y) \ar[r] &
K_0(\textit{Coh}_{\leq k}(Y)/\textit{Coh}_{\leq k - 1}(Y)) \ar[r] &
Z_k(Y)
}
$$
\end{lemma}

\begin{proof}
Un morphisme fini est affine ; l'image directe des modules quasi-cohérents
par $\pi$ est donc un foncteur exact d'après le
Lemme \ref{coherent-lemma-relative-affine-vanishing} de Cohomologie des schémas.
Un morphisme fini est propre ; $\pi_*$ envoie donc les faisceaux cohérents
sur des faisceaux cohérents ; voir la Proposition
\ref{coherent-proposition-proper-pushforward-coherent} de Cohomologie des schémas.
L'assertion sur les dimensions des supports est claire.
La commutativité à droite résulte immédiatement du
Lemme \ref{chow-lemma-cycle-push-sheaf}.
Puisque les flèches horizontales sont des bijections, nous constatons que
nous avons aussi la commutativité à gauche.
\end{proof}

\begin{lemma}
\label{chow-lemma-from-chow-to-K}
Soit $X$ un schéma localement de type fini sur $(S, \delta)$
comme dans la Situation \ref{chow-situation-setup}. Il existe une application canonique
$$
\CH_k(X)
\longrightarrow
K_0(\textit{Coh}_{\leq k + 1}(X)/\textit{Coh}_{\leq k - 1}(X))
$$
induite par l'application
$Z_k(X) \to K_0(\textit{Coh}_{\leq k}(X)/\textit{Coh}_{\leq k - 1}(X))$
du Lemme \ref{chow-lemma-cycles-k-group}.
\end{lemma}

\begin{proof}
Nous devons montrer qu'un élément $\alpha$ de $Z_k(X)$ qui est rationnellement
équivalent à zéro est envoyé sur zéro dans
$K_0(\textit{Coh}_{\leq k + 1}(X)/\textit{Coh}_{\leq k - 1}(X))$.
Écrivons $\alpha = \sum (i_j)_*\text{div}(f_j)$ comme dans la
Définition \ref{chow-definition-rational-equivalence}.
Observons que
$$
\pi = \coprod i_j : W = \coprod W_j \longrightarrow X
$$
est un morphisme fini, car chaque $i_j : W_j \to X$ est une immersion fermée
et la famille des $W_j$ est localement finie dans $X$. Nous pouvons donc utiliser le
Lemme \ref{chow-lemma-finite-cycles-k-group} pour nous ramener au cas de $W$.
Puisque $W$ est une union disjointe de schémas intègres, nous nous ramenons
au cas étudié dans le paragraphe suivant.

\medskip\noindent
Supposons $X$ intègre de $\delta$-dimension $k + 1$.
Soit $f$ une fonction rationnelle non nulle sur $X$.
Soit $\alpha = \text{div}(f)$. Nous devons montrer que
$\alpha$ est envoyé sur zéro dans
$K_0(\textit{Coh}_{\leq k + 1}(X)/\textit{Coh}_{\leq k - 1}(X))$.
Soit $\mathcal{I} \subset \mathcal{O}_X$ l'idéal des dénominateurs
de $f$ ; voir Diviseurs, Définition
\ref{divisors-definition-regular-meromorphic-ideal-denominators}.
Nous avons alors des suites exactes courtes
$$
0 \to \mathcal{I} \to \mathcal{O}_X \to \mathcal{O}_X/\mathcal{I} \to 0
$$
et
$$
0 \to \mathcal{I} \xrightarrow{f} \mathcal{O}_X \to
\mathcal{O}_X/f\mathcal{I} \to 0
$$
Voir Diviseurs, Lemme
\ref{divisors-lemma-regular-meromorphic-ideal-denominators}.
Nous affirmons que
$$
[\mathcal{O}_X/\mathcal{I}]_k - [\mathcal{O}_X/f\mathcal{I}]_k =
\text{div}(f)
$$
Cette affirmation implique que l'élément $\alpha = \text{div}(f)$ est représenté par
$[\mathcal{O}_X/\mathcal{I}] - [\mathcal{O}_X/f\mathcal{I}]$
dans $K_0(\textit{Coh}_{\leq k}(X)/\textit{Coh}_{\leq k - 1}(X))$.
Les suites exactes courtes montrent alors que cet élément s'envoie sur
zéro dans $K_0(\textit{Coh}_{\leq k + 1}(X)/\textit{Coh}_{\leq k - 1}(X))$.

\medskip\noindent
Pour démontrer l'affirmation, soit $Z \subset X$ un sous-schéma fermé intègre
de $\delta$-dimension $k$, et soit $\xi \in Z$ son point générique.
Alors $I = \mathcal{I}_\xi \subset A = \mathcal{O}_{X, \xi}$
est un idéal tel que $fI \subset A$. Or le coefficient de
$[Z]$ dans $\text{div}(f)$ est $\text{ord}_A(f)$. (Bien sûr, comme d'habitude,
nous identifions le corps des fonctions de $X$ au corps des fractions de $A$.)
D'autre part, le coefficient de $[Z]$ dans
$[\mathcal{O}_X/\mathcal{I}] - [\mathcal{O}_X/f\mathcal{I}]$
est
$$
\text{longueur}_A(A/I) - \text{longueur}_A(A/fI)
$$
À l'aide de la fonction distance de la
Définition \ref{algebra-definition-distance} d'Algèbre,
nous pouvons récrire ceci sous la forme
$$
d(A, I) - d(A, fI) = d(I, fI) = \text{ord}_A(f)
$$
Les égalités sont vraies d'après les Lemmes
\ref{algebra-lemma-properties-distance-function} et
\ref{algebra-lemma-order-vanishing-determinant} d'Algèbre.
(L'utilisation de ces lemmes n'est pas nécessaire, mais commode.)
\end{proof}

\begin{remark}
\label{chow-remark-good-cases-K-A}
Soit $(S, \delta)$ comme dans la Situation \ref{chow-situation-setup}.
Soit $X$ un schéma localement de type fini sur $S$.
Nous verrons plus loin (dans le Lemme \ref{chow-lemma-cycles-rational-equivalence-K-group})
que l'application
$$
\CH_k(X)
\longrightarrow
K_0(\textit{Coh}_{k + 1}(X)/\textit{Coh}_{\leq k - 1}(X))
$$
du Lemme \ref{chow-lemma-from-chow-to-K} est injective.
En la composant avec l'application canonique
$$
K_0(\textit{Coh}_{k + 1}(X)/\textit{Coh}_{\leq k - 1}(X))
\longrightarrow
K_0(\textit{Coh}(X)/\textit{Coh}_{\leq k - 1}(X))
$$
nous obtenons une application canonique
$$
\CH_k(X)
\longrightarrow
K_0(\textit{Coh}(X)/\textit{Coh}_{\leq k - 1}(X)).
$$
Nous n'avons trouvé dans la littérature aucun énoncé ni aucune conjecture
indiquant si cette application devrait être injective ou non.
Il semble raisonnable de s'attendre à ce que le noyau de cette application soit de torsion.
Nous reviendrons sur cette question (insérer une référence ultérieure).
\end{remark}

\begin{lemma}
\label{chow-lemma-K-coherent-supported-on-closed}
Soit $X$ un schéma localement noethérien. Soit $Z \subset X$ un sous-schéma
fermé. Notons $\textit{Coh}_Z(X) \subset \textit{Coh}(X)$
la sous-catégorie de Serre des $\mathcal{O}_X$-modules cohérents dont le
support ensembliste est contenu dans $Z$. Alors le foncteur exact d'inclusion
$\textit{Coh}(Z) \to \textit{Coh}_Z(X)$ induit
un isomorphisme
$$
K'_0(Z) = K_0(\textit{Coh}(Z)) \longrightarrow K_0(\textit{Coh}_Z(X))
$$
\end{lemma}

\begin{proof}
Soit $\mathcal{F}$ un objet de $\textit{Coh}_Z(X)$.
Soit $\mathcal{I} \subset \mathcal{O}_X$ le faisceau quasi-cohérent
d'idéaux de $Z$. Considérons la filtration décroissante
$$
\ldots \subset
\mathcal{F}^p = \mathcal{I}^p \mathcal{F} \subset
\mathcal{F}^{p - 1} \subset \ldots \subset \mathcal{F}^0 = \mathcal{F}
$$
Exactement comme dans la démonstration du Lemme \ref{chow-lemma-from-chow-to-K}, cette filtration
est localement finie, et donc
$\bigoplus_{p \geq 0} \mathcal{F}^p$,
$\bigoplus_{p \geq 1} \mathcal{F}^p$ et
$\bigoplus_{p \geq 0} \mathcal{F}^p/\mathcal{F}^{p + 1}$
sont des $\mathcal{O}_X$-modules cohérents supportés sur $Z$.
Nous obtenons ainsi
$$
[\mathcal{F}] =
[\bigoplus\nolimits_{p \geq 0} \mathcal{F}^p/\mathcal{F}^{p + 1}]
$$
dans $K_0(\textit{Coh}_Z(X))$, exactement comme dans la démonstration du
Lemme \ref{chow-lemma-from-chow-to-K}. Puisque le module cohérent
$\bigoplus_{p \geq 0} \mathcal{F}^p/\mathcal{F}^{p + 1}$
est annulé par $\mathcal{I}$, nous concluons que
$[\mathcal{F}]$ appartient à l'image. En fait, nous affirmons que l'application
$$
\mathcal{F} \longmapsto 
c(\mathcal{F}) =
[\bigoplus\nolimits_{p \geq 0} \mathcal{F}^p/\mathcal{F}^{p + 1}]
$$
se factorise par $K_0(\textit{Coh}_Z(X))$ et est réciproque de
l'application de l'énoncé du lemme. Pour le voir, il nous
suffit de montrer que, si
$$
0 \to \mathcal{F} \to \mathcal{G} \to \mathcal{H} \to 0
$$
est une suite exacte courte dans $\textit{Coh}_Z(X)$, alors nous
obtenons $c(\mathcal{G}) = c(\mathcal{F}) + c(\mathcal{H})$.
Observons que, pour tout $q \geq 0$, nous avons une suite exacte courte
$$
0 \to
(\mathcal{F} \cap \mathcal{I}^q\mathcal{G})/
(\mathcal{F} \cap \mathcal{I}^{q + 1}\mathcal{G}) \to
\mathcal{G}^q/\mathcal{G}^{q + 1} \to
\mathcal{H}^q/\mathcal{H}^{q + 1} \to 0
$$
Pour $p, q \geq 0$, considérons le sous-module cohérent
$$
\mathcal{F}^{p, q} = \mathcal{I}^p\mathcal{F} \cap \mathcal{I}^q\mathcal{G}
$$
En raisonnant exactement comme ci-dessus et en utilisant que les filtrations
$\mathcal{F}^p = \mathcal{I}^p\mathcal{F}$ et
$\mathcal{F} \cap \mathcal{I}^q\mathcal{G}$ sont localement finies,
nous obtenons
$$
[\bigoplus\nolimits_{p \geq 0} \mathcal{F}^p/\mathcal{F}^{p + 1}] =
[\bigoplus\nolimits_{p, q \geq 0}
\mathcal{F}^{p, q}/(\mathcal{F}^{p + 1, q} + \mathcal{F}^{p, q + 1})] =
[\bigoplus\nolimits_{q \geq 0}
(\mathcal{F} \cap \mathcal{I}^q\mathcal{G})/
(\mathcal{F} \cap \mathcal{I}^{q + 1}\mathcal{G})]
$$
dans $K_0(\textit{Coh}(Z))$. Avec les suites exactes ci-dessus, nous obtenons
le résultat souhaité. Certains détails sont omis.
\end{proof}













\section{Le diviseur associé à un faisceau inversible}
\label{chow-section-divisor-invertible-sheaf}

\noindent
La définition suivante est l'analogue de la
Définition \ref{divisors-definition-divisor-invertible-sheaf} de Diviseurs
dans notre cadre actuel.

\begin{definition}
\label{chow-definition-divisor-invertible-sheaf}
Soit $(S, \delta)$ comme dans la Situation \ref{chow-situation-setup}.
Soit $X$ localement de type fini sur $S$. Supposons $X$
intègre et $n = \dim_\delta(X)$.
Soit $\mathcal{L}$ un $\mathcal{O}_X$-module inversible.
\begin{enumerate}
\item Pour toute section méromorphe non nulle $s$ de $\mathcal{L}$,
nous définissons le {\it diviseur de Weil associé à $s$} comme le
$(n - 1)$-cycle
$$
\text{div}_\mathcal{L}(s) =
\sum \text{ord}_{Z, \mathcal{L}}(s) [Z]
$$
défini dans Diviseurs, Définition
\ref{divisors-definition-divisor-invertible-sheaf}.
Ceci a un sens parce que les diviseurs de Weil sont de $\delta$-dimension $n - 1$
d'après le Lemme \ref{chow-lemma-divisor-delta-dimension}.
\item Nous définissons le {\it diviseur de Weil associé à $\mathcal{L}$} par
$$
c_1(\mathcal{L}) \cap [X] =
\text{classe de }\text{div}_\mathcal{L}(s) \in \CH_{n - 1}(X)
$$
où $s$ est une section méromorphe non nulle quelconque de $\mathcal{L}$ sur
$X$. Ceci est bien défini d'après le
Lemme \ref{divisors-lemma-divisor-meromorphic-well-defined} de Diviseurs.
\end{enumerate}
\end{definition}

\noindent
Soient $X$ et $S$ comme dans la Définition \ref{chow-definition-divisor-invertible-sheaf}
ci-dessus. Posons $n = \dim_\delta(X)$. Il est clair d'après les définitions que
$Cl(X) = \CH_{n - 1}(X)$, où $Cl(X)$ est le groupe des classes de diviseurs de Weil de $X$
défini dans Diviseurs, Définition \ref{divisors-definition-class-group}.
L'application
$$
\Pic(X) \longrightarrow \CH_{n - 1}(X), \quad
\mathcal{L} \longmapsto c_1(\mathcal{L}) \cap [X]
$$
est la même que l'application $\Pic(X) \to Cl(X)$ construite dans
Diviseurs, Équation (\ref{divisors-equation-c1}), pour les schémas intègres
localement noethériens quelconques. En particulier, cette application
est un homomorphisme de groupes abéliens ; elle est injective si $X$ est
un schéma normal, et un isomorphisme si tous les anneaux locaux de $X$
sont factoriels. Voir Diviseurs, Lemmes \ref{divisors-lemma-normal-c1-injective} et
\ref{divisors-lemma-local-rings-UFD-c1-bijective}.
Dans certains cas, il est facile de calculer le
diviseur de Weil associé à un faisceau inversible.

\begin{lemma}
\label{chow-lemma-compute-c1}
Soit $(S, \delta)$ comme dans la Situation \ref{chow-situation-setup}.
Soit $X$ localement de type fini sur $S$. Supposons $X$
intègre et $n = \dim_\delta(X)$.
Soit $\mathcal{L}$ un $\mathcal{O}_X$-module inversible.
Soit $s \in \Gamma(X, \mathcal{L})$ une section globale non nulle.
Alors
$$
\text{div}_\mathcal{L}(s) = [Z(s)]_{n - 1}
$$
dans $Z_{n - 1}(X)$ et
$$
c_1(\mathcal{L}) \cap [X] = [Z(s)]_{n - 1}
$$
dans $\CH_{n - 1}(X)$.
\end{lemma}

\begin{proof}
Soit $Z \subset X$ un sous-schéma fermé intègre de
$\delta$-dimension $n - 1$. Soit $\xi \in Z$ son point
générique. Choisissons un générateur $s_\xi \in \mathcal{L}_\xi$.
Écrivons $s = fs_\xi$ pour un certain $f \in \mathcal{O}_{X, \xi}$.
Par définition de $Z(s)$, voir la
Définition \ref{divisors-definition-zero-scheme-s} de Diviseurs,
nous voyons que $Z(s)$ est défini par un faisceau quasi-cohérent
d'idéaux $\mathcal{I} \subset \mathcal{O}_X$ tel
que $\mathcal{I}_\xi = (f)$. Ainsi
$\text{longueur}_{\mathcal{O}_{X, \xi}}(\mathcal{O}_{Z(s), \xi})
=
\text{longueur}_{\mathcal{O}_{X, \xi}}(\mathcal{O}_{X, \xi}/(f))
=
\text{ord}_{\mathcal{O}_{X, \xi}}(f)$, comme souhaité.
\end{proof}

\noindent
Le lemme suivant sera supplanté par le Lemme plus général
\ref{chow-lemma-flat-pullback-cap-c1}.

\begin{lemma}
\label{chow-lemma-flat-pullback-divisor-invertible-sheaf}
Soit $(S, \delta)$ comme dans la Situation \ref{chow-situation-setup}.
Soient $X$, $Y$ localement de type fini sur $S$. Supposons $X$, $Y$
intègres et $n = \dim_\delta(Y)$.
Soit $\mathcal{L}$ un $\mathcal{O}_Y$-module inversible.
Soit $f : X \to Y$ un morphisme plat de dimension relative $r$. Alors
$$
f^*(c_1(\mathcal{L}) \cap [Y]) = c_1(f^*\mathcal{L}) \cap [X]
$$
dans $\CH_{n + r - 1}(X)$.
\end{lemma}

\begin{proof}
Soit $s$ une section méromorphe non nulle de $\mathcal{L}$.
Nous allons montrer qu'en fait
$f^*\text{div}_\mathcal{L}(s) = \text{div}_{f^*\mathcal{L}}(f^*s)$,
et donc que le lemme est vrai.
Pour le voir, soit $\xi \in Y$ un point, et soit $s_\xi \in \mathcal{L}_\xi$
un générateur. Écrivons $s = gs_\xi$ avec $g \in R(Y)^*$.
Il existe alors un voisinage ouvert $V \subset Y$ de $\xi$
tel que $s_\xi \in \mathcal{L}(V)$ et que $s_\xi$ engendre
$\mathcal{L}|_V$. Nous voyons donc que
$$
\text{div}_\mathcal{L}(s)|_V = \text{div}_Y(g)|_V.
$$
De la même manière exactement, puisque $f^*s_\xi$ engendre $f^*\mathcal{L}$
sur $f^{-1}(V)$ et puisque $f^*s = g f^*s_\xi$, nous
avons aussi
$$
\text{div}_\mathcal{L}(f^*s)|_{f^{-1}(V)}
=
\text{div}_X(g)|_{f^{-1}(V)}.
$$
L'égalité souhaitée de cycles sur $f^{-1}(V)$ résulte donc du
résultat correspondant pour les images réciproques de diviseurs principaux ; voir le
Lemme \ref{chow-lemma-flat-pullback-principal-divisor}.
\end{proof}



\section{Intersection par un faisceau inversible}
\label{chow-section-intersecting-with-divisors}

\noindent
Dans cette section, nous étudions la construction suivante.

\begin{definition}
\label{chow-definition-cap-c1}
Soit $(S, \delta)$ comme dans la Situation \ref{chow-situation-setup}.
Soit $X$ localement de type fini sur $S$.
Soit $\mathcal{L}$ un $\mathcal{O}_X$-module inversible.
Nous définissons, pour tout entier $k$, une opération
$$
c_1(\mathcal{L}) \cap - :
Z_{k + 1}(X) \to \CH_k(X)
$$
appelée {\it intersection avec la première classe de Chern de $\mathcal{L}$}.
\begin{enumerate}
\item Étant donné un sous-schéma fermé intègre $i : W \to X$ avec
$\dim_\delta(W) = k + 1$, nous définissons
$$
c_1(\mathcal{L}) \cap [W] = i_*(c_1({i^*\mathcal{L}}) \cap [W])
$$
où le membre de droite est défini dans la
Définition \ref{chow-definition-divisor-invertible-sheaf}.
\item Pour un $(k + 1)$-cycle général $\alpha = \sum n_i [W_i]$, nous posons
$$
c_1(\mathcal{L}) \cap \alpha = \sum n_i c_1(\mathcal{L}) \cap [W_i]
$$
\end{enumerate}
\end{definition}

\noindent
Écrivons chaque $c_1(\mathcal{L}) \cap W_i = \sum_j n_{i, j} [Z_{i, j}]$,
où $\{Z_{i, j}\}_j$ est une somme localement finie
de sous-schémas fermés intègres de $W_i$. Puisque $\{W_i\}$ est une famille
localement finie de sous-schémas fermés intègres de $X$, il en résulte
facilement que $\{Z_{i, j}\}_{i, j}$ est une famille localement finie
de sous-schémas fermés de $X$. Ainsi
$c_1(\mathcal{L}) \cap \alpha = \sum n_in_{i, j}[Z_{i, j}]$
est un cycle. Une autre manière, plus commode, de voir cela
est d'observer que le morphisme $\coprod W_i \to X$ est
propre. Ainsi $c_1(\mathcal{L}) \cap \alpha$ peut être considéré
comme l'image directe d'une classe de $\CH_k(\coprod W_i) = \prod \CH_k(W_i)$.
Cela explique aussi pourquoi le résultat est bien défini modulo équivalence
rationnelle sur $X$.

\medskip\noindent
L'objectif principal des quelques sections suivantes est de montrer que l'intersection avec
$c_1(\mathcal{L})$ se factorise par l'équivalence rationnelle.
Ce n'est pas une trivialité.

\begin{lemma}
\label{chow-lemma-c1-cap-additive}
Soit $(S, \delta)$ comme dans la Situation \ref{chow-situation-setup}.
Soit $X$ localement de type fini sur $S$.
Soient $\mathcal{L}$, $\mathcal{N}$ des faisceaux inversibles sur $X$.
Alors
$$
c_1(\mathcal{L}) \cap \alpha  + c_1(\mathcal{N}) \cap \alpha =
c_1(\mathcal{L} \otimes_{\mathcal{O}_X} \mathcal{N}) \cap \alpha
$$
dans $\CH_k(X)$ pour tout $\alpha \in Z_{k + 1}(X)$. De plus,
$c_1(\mathcal{O}_X) \cap \alpha = 0$ pour tout $\alpha$.
\end{lemma}

\begin{proof}
L'additivité résulte directement du
Lemme \ref{divisors-lemma-c1-additive} de Diviseurs
et des définitions. Pour voir que $c_1(\mathcal{O}_X) \cap \alpha = 0$,
considérons la section $1 \in \Gamma(X, \mathcal{O}_X)$. Sa restriction
est une section partout non nulle sur tout sous-schéma fermé intègre
$W \subset X$. Ainsi $c_1(\mathcal{O}_X) \cap [W] = 0$, comme souhaité.
\end{proof}

\noindent
Rappelons que $Z(s) \subset X$ désigne le schéma des zéros d'une section globale
$s$ d'un faisceau inversible sur un schéma $X$ ; voir la
Définition \ref{divisors-definition-zero-scheme-s} de Diviseurs.

\begin{lemma}
\label{chow-lemma-prepare-geometric-cap}
Soit $(S, \delta)$ comme dans la Situation \ref{chow-situation-setup}.
Soit $Y$ localement de type fini sur $S$.
Soit $\mathcal{L}$ un $\mathcal{O}_Y$-module inversible.
Soit $s \in \Gamma(Y, \mathcal{L})$.
Supposons
\begin{enumerate}
\item $\dim_\delta(Y) \leq k + 1$,
\item $\dim_\delta(Z(s)) \leq k$, et
\item pour tout point générique $\xi$ d'une composante irréductible de
$Z(s)$ de $\delta$-dimension $k$, la multiplication par $s$
induit une injection $\mathcal{O}_{Y, \xi} \to \mathcal{L}_\xi$.
\end{enumerate}
Écrivons $[Y]_{k + 1} = \sum n_i[Y_i]$, où les $Y_i \subset Y$ sont les
composantes irréductibles de $Y$ de $\delta$-dimension $k + 1$.
Posons $s_i = s|_{Y_i} \in \Gamma(Y_i, \mathcal{L}|_{Y_i})$. Alors
\begin{equation}
\label{chow-equation-equal-as-cycles}
[Z(s)]_k =  \sum n_i[Z(s_i)]_k
\end{equation}
comme $k$-cycles sur $Y$.
\end{lemma}

\begin{proof}
Soit $Z \subset Y$ un sous-schéma fermé intègre de
$\delta$-dimension $k$. Soit $\xi \in Z$ son point générique.
Nous voulons comparer le coefficient $n$ de $[Z]$ dans l'expression
$\sum n_i[Z(s_i)]_k$ au coefficient $m$ de $[Z]$ dans
l'expression $[Z(s)]_k$. Choisissons un générateur $s_\xi \in \mathcal{L}_\xi$.
Écrivons $A = \mathcal{O}_{Y, \xi}$, $L = \mathcal{L}_\xi$.
Alors $L = As_\xi$. Écrivons $s = f s_\xi$ pour un certain (unique) $f \in A$.
L'hypothèse (3) signifie que $f : A \to A$ est injectif.
Puisque $\dim_\delta(Y) \leq k + 1$ et $\dim_\delta(Z) = k$,
nous avons $\dim(A) = 0$ ou $1$. Nous avons
$$
m = \text{longueur}_A(A/(f))
$$
qui est finie dans les deux cas.

\medskip\noindent
Si $\dim(A) = 0$, alors l'injectivité de $f : A \to A$
implique que $f \in A^*$. Ainsi, dans ce cas, $m$ est nul.
De plus, la condition $\dim(A) = 0$ signifie que $\xi$
n'appartient à aucune composante irréductible de $\delta$-dimension
$k + 1$, c'est-à-dire que $n = 0$ également.

\medskip\noindent
Supposons maintenant $\dim(A) = 1$.
Puisque $A$ est un anneau local noethérien, il possède un nombre fini
d'idéaux premiers minimaux $\mathfrak q_1, \ldots, \mathfrak q_t$.
Ceux-ci correspondent bijectivement aux $Y_i$ passant par $\xi'$.
De plus, $n_i = \text{longueur}_{A_{\mathfrak q_i}}(A_{\mathfrak q_i})$.
La multiplicité de $[Z]$ dans $[Z(s_i)]_k$ est aussi
$\text{longueur}_A(A/(f, \mathfrak q_i))$.
L'équation à démontrer dans ce cas est donc
$$
\text{longueur}_A(A/(f))
=
\sum \text{longueur}_{A_{\mathfrak q_i}}(A_{\mathfrak q_i})
\text{longueur}_A(A/(f, \mathfrak q_i))
$$
qui résulte du
Lemme \ref{chow-lemma-additivity-divisors-restricted}.
\end{proof}

\noindent
Le lemme suivant est un résultat utile pour calculer le produit d'intersection
de $c_1$ d'un faisceau inversible avec le cycle associé
à un sous-schéma fermé.
Rappelons que $Z(s) \subset X$ désigne le schéma des zéros d'une section globale
$s$ d'un faisceau inversible sur un schéma $X$ ; voir la
Définition \ref{divisors-definition-zero-scheme-s} de Diviseurs.

\begin{lemma}
\label{chow-lemma-geometric-cap}
Soit $(S, \delta)$ comme dans la Situation \ref{chow-situation-setup}.
Soit $X$ localement de type fini sur $S$.
Soit $\mathcal{L}$ un $\mathcal{O}_X$-module inversible.
Soit $Y \subset X$ un sous-schéma fermé.
Soit $s \in \Gamma(Y, \mathcal{L}|_Y)$.
Supposons
\begin{enumerate}
\item $\dim_\delta(Y) \leq k + 1$,
\item $\dim_\delta(Z(s)) \leq k$, et
\item pour tout point générique $\xi$ d'une composante irréductible de
$Z(s)$ de $\delta$-dimension $k$, la multiplication par $s$
induit une injection
$\mathcal{O}_{Y, \xi} \to (\mathcal{L}|_Y)_\xi$\footnote{Par exemple,
ceci est vrai si $s$ est une section régulière de $\mathcal{L}|_Y$.}.
\end{enumerate}
Alors
$$
c_1(\mathcal{L}) \cap [Y]_{k + 1} = [Z(s)]_k
$$
dans $\CH_k(X)$.
\end{lemma}

\begin{proof}
Écrivons
$$
[Y]_{k + 1} = \sum n_i[Y_i]
$$
où les $Y_i \subset Y$ sont les composantes irréductibles de
$Y$ de $\delta$-dimension $k + 1$ et $n_i > 0$.
Par hypothèse, la restriction
$s_i = s|_{Y_i} \in \Gamma(Y_i, \mathcal{L}|_{Y_i})$ n'est pas
nulle, et est donc une section régulière. D'après le Lemme \ref{chow-lemma-compute-c1},
nous voyons que $[Z(s_i)]_k$ représente $c_1(\mathcal{L}|_{Y_i})$.
Ainsi, par définition,
$$
c_1(\mathcal{L}) \cap [Y]_{k + 1} = \sum n_i[Z(s_i)]_k
$$
Le résultat résulte donc du Lemme \ref{chow-lemma-prepare-geometric-cap}.
\end{proof}




\section{Intersection par un faisceau inversible et images directe et réciproque}
\label{chow-section-intersecting-with-divisors-push-pull}

\noindent
Dans cette section, nous démontrons que l'opération $c_1(\mathcal{L}) \cap -$
commute avec l'image réciproque plate et l'image directe propre.

\begin{lemma}
\label{chow-lemma-prepare-flat-pullback-cap-c1}
Soit $(S, \delta)$ comme dans la Situation \ref{chow-situation-setup}.
Soient $X$, $Y$ localement de type fini sur $S$.
Soit $f : X \to Y$ un morphisme plat de dimension relative $r$.
Soit $\mathcal{L}$ un faisceau inversible sur $Y$.
Supposons $Y$ intègre et $n = \dim_\delta(Y)$.
Soit $s$ une section méromorphe non nulle de $\mathcal{L}$.
Alors nous avons
$$
f^*\text{div}_\mathcal{L}(s) = \sum n_i\text{div}_{f^*\mathcal{L}|_{X_i}}(s_i)
$$
dans $Z_{n + r - 1}(X)$. Ici, la somme porte sur les composantes
irréductibles $X_i \subset X$ de $\delta$-dimension $n + r$,
la section $s_i = f|_{X_i}^*(s)$ est l'image réciproque de $s$, et
$n_i = m_{X_i, X}$ est la multiplicité de $X_i$ dans $X$.
\end{lemma}

\begin{proof}
Pour démontrer cette égalité de cycles, nous pouvons travailler localement sur $Y$.
Nous pouvons donc supposer $Y$ affine et $s = p/q$ pour des sections non nulles
$p \in \Gamma(Y, \mathcal{L})$ et $q \in \Gamma(Y, \mathcal{O})$.
Si nous pouvons démontrer à la fois
$$
f^*\text{div}_\mathcal{L}(p) =
\sum n_i\text{div}_{f^*\mathcal{L}|_{X_i}}(p_i)
\quad\text{et}\quad
f^*\text{div}_\mathcal{O}(q) =
\sum n_i\text{div}_{\mathcal{O}_{X_i}}(q_i)
$$
(avec les notations évidentes), alors nous avons le résultat par
additivité ; voir le Lemme \ref{divisors-lemma-c1-additive} de Diviseurs.
Nous pouvons donc supposer que $s \in \Gamma(Y, \mathcal{L})$.
Dans ce cas, nous pouvons appliquer l'égalité
(\ref{chow-equation-equal-as-cycles}) pour voir que
$$
[Z(f^*(s))]_{k + r - 1} =
\sum n_i\text{div}_{f^*\mathcal{L}|_{X_i}}(s_i)
$$
où $f^*(s) \in f^*\mathcal{L}$ désigne l'image réciproque de $s$ sur $X$.
D'autre part, nous avons
$$
f^*\text{div}_\mathcal{L}(s) = f^*[Z(s)]_{k - 1}
= [f^{-1}(Z(s))]_{k + r - 1},
$$
d'après les Lemmes \ref{chow-lemma-compute-c1} et \ref{chow-lemma-pullback-coherent}.
Puisque $Z(f^*(s)) = f^{-1}(Z(s))$, nous avons le résultat.
\end{proof}

\begin{lemma}
\label{chow-lemma-flat-pullback-cap-c1}
Soit $(S, \delta)$ comme dans la Situation \ref{chow-situation-setup}.
Soient $X$, $Y$ localement de type fini sur $S$.
Soit $f : X \to Y$ un morphisme plat de dimension relative $r$.
Soit $\mathcal{L}$ un faisceau inversible sur $Y$.
Soit $\alpha$ un $k$-cycle sur $Y$.
Alors
$$
f^*(c_1(\mathcal{L}) \cap \alpha) = c_1(f^*\mathcal{L}) \cap f^*\alpha
$$
dans $\CH_{k + r - 1}(X)$.
\end{lemma}

\begin{proof}
Écrivons $\alpha = \sum n_i[W_i]$. Nous allons montrer que
$$
f^*(c_1(\mathcal{L}) \cap [W_i]) = c_1(f^*\mathcal{L}) \cap f^*[W_i]
$$
dans $\CH_{k + r - 1}(X)$ en produisant une équivalence rationnelle
sur le sous-schéma fermé $f^{-1}(W_i)$ de $X$.
D'après la discussion de la
Remarque \ref{chow-remark-infinite-sums-rational-equivalences},
ceci démontrera l'égalité du lemme.

\medskip\noindent
Soit $W \subset Y$ un sous-schéma fermé intègre de $\delta$-dimension $k$.
Considérons le sous-schéma fermé $W' = f^{-1}(W) = W \times_Y X$,
de sorte que nous avons le diagramme cartésien
$$
\xymatrix{
W' \ar[r] \ar[d]_h & X \ar[d]^f \\
W \ar[r] & Y
}
$$
Nous devons montrer que
$f^*(c_1(\mathcal{L}) \cap [W]) = c_1(f^*\mathcal{L}) \cap f^*[W]$.
Choisissons une section méromorphe non nulle $s$ de $\mathcal{L}|_W$.
Soient $W'_i \subset W'$ les composantes irréductibles de
$\delta$-dimension $k + r$. Écrivons $[W']_{k + r} = \sum n_i[W'_i]$,
où $n_i$ est la multiplicité de $W'_i$ dans $W'$, conformément à la définition.
Ainsi $f^*[W] = \sum n_i[W'_i]$ dans $Z_{k + r}(X)$.
Puisque chaque $W'_i \to W$ est dominant, nous
voyons que $s_i = s|_{W'_i}$ est une section méromorphe non nulle pour
tout $i$. D'après le Lemme \ref{chow-lemma-prepare-flat-pullback-cap-c1},
nous avons l'égalité de cycles suivante
$$
h^*\text{div}_{\mathcal{L}|_W}(s) =
\sum n_i\text{div}_{f^*\mathcal{L}|_{W'_i}}(s_i)
$$
dans $Z_{k + r - 1}(W')$. Ceci achève la démonstration, puisque
le membre de gauche est un cycle sur $W'$ dont l'image directe est
$f^*(c_1(\mathcal{L}) \cap [W])$ dans $\CH_{k + r - 1}(X)$,
et le membre de droite est un cycle sur $W'$ dont l'image directe est
$c_1(f^*\mathcal{L}) \cap f^*[W]$ dans $\CH_{k + r - 1}(X)$.
\end{proof}

\begin{lemma}
\label{chow-lemma-equal-c1-as-cycles}
Soit $(S, \delta)$ comme dans la Situation \ref{chow-situation-setup}.
Soient $X$, $Y$ localement de type fini sur $S$.
Soit $f : X \to Y$ un morphisme propre.
Soit $\mathcal{L}$ un faisceau inversible sur $Y$.
Soit $s$ une section méromorphe non nulle, encore notée $s$, de $\mathcal{L}$ sur $Y$.
Supposons $X$, $Y$ intègres, $f$ dominant et $\dim_\delta(X) = \dim_\delta(Y)$.
Alors
$$
f_*\left(\text{div}_{f^*\mathcal{L}}(f^*s)\right) =
[R(X) : R(Y)]\text{div}_\mathcal{L}(s).
$$
comme cycles sur $Y$. En particulier,
$$
f_*(c_1(f^*\mathcal{L}) \cap [X]) =
[R(X) : R(Y)] c_1(\mathcal{L}) \cap [Y] =
c_1(\mathcal{L}) \cap f_*[X]
$$
\end{lemma}

\begin{proof}
La dernière égalité résulte de la première puisque $f_*[X] = [R(X) : R(Y)][Y]$
par définition. Il se trouve que nous pouvons réutiliser le
Lemme \ref{chow-lemma-proper-pushforward-alteration}
pour le démontrer. En effet, puisque nous cherchons à démontrer une égalité
de cycles, nous pouvons travailler localement sur $Y$. Nous pouvons donc supposer
que $\mathcal{L} = \mathcal{O}_Y$. Dans ce cas, $s$
correspond à une fonction rationnelle $g \in R(Y)$, et
nous cherchons simplement à démontrer
$$
f_*\left(\text{div}_X(g)\right) =
[R(X) : R(Y)]\text{div}_Y(g).
$$
En comparant avec le résultat du
Lemme \ref{chow-lemma-proper-pushforward-alteration} précité,
nous voyons que c'est vrai puisque
$\text{Nm}_{R(X)/R(Y)}(g) = g^{[R(X) : R(Y)]}$,
car $g \in R(Y)^*$.
\end{proof}

\begin{lemma}
\label{chow-lemma-pushforward-cap-c1}
Soit $(S, \delta)$ comme dans la Situation \ref{chow-situation-setup}.
Soient $X$, $Y$ localement de type fini sur $S$.
Soit $p : X \to Y$ un morphisme propre.
Soit $\alpha \in Z_{k + 1}(X)$.
Soit $\mathcal{L}$ un faisceau inversible sur $Y$.
Alors
$$
p_*(c_1(p^*\mathcal{L}) \cap \alpha) = c_1(\mathcal{L}) \cap p_*\alpha
$$
dans $\CH_k(Y)$.
\end{lemma}

\begin{proof}
Supposons que $p$ ait la propriété que, pour tout sous-schéma fermé
intègre $W \subset X$, l'application $p|_W : W \to Y$
soit une immersion fermée. Alors, par définition du produit
avec $c_1(\mathcal{L})$, le lemme est vrai.

\medskip\noindent
Nous allons utiliser cette remarque pour nous ramener à un cas particulier. En effet,
écrivons $\alpha = \sum n_i[W_i]$ avec $n_i \not = 0$ et les $W_i$ deux à deux
distincts. Soit $W'_i \subset Y$ l'image de $W_i$ (comme sous-schéma
fermé intègre). Considérons le diagramme
$$
\xymatrix{
X' = \coprod W_i \ar[r]_-q \ar[d]_{p'} & X \ar[d]^p \\
Y' = \coprod W'_i \ar[r]^-{q'} & Y.
}
$$
Puisque $\{W_i\}$ est localement finie sur $X$ et que $p$ est propre,
nous voyons que $\{W'_i\}$ est localement finie sur $Y$ et que
$q, q', p'$ sont aussi des morphismes propres.
Nous pouvons aussi considérer $\sum n_i[W_i]$ comme un $k$-cycle
$\alpha' \in Z_k(X')$. Il est clair que $q_*\alpha' = \alpha$.
Nous avons
$q_*(c_1(q^*p^*\mathcal{L}) \cap \alpha')
= c_1(p^*\mathcal{L}) \cap q_*\alpha'$
et
$(q')_*(c_1((q')^*\mathcal{L}) \cap p'_*\alpha') =
c_1(\mathcal{L}) \cap q'_*p'_*\alpha'$ d'après la remarque initiale
de la démonstration. Il suffit donc de démontrer le lemme
pour le morphisme $p'$ et le cycle $\sum n_i[W_i]$.
Il est clair que cela signifie que nous pouvons supposer $X$, $Y$ intègres,
$f : X \to Y$ dominant et $\alpha = [X]$.
Dans ce cas, le résultat résulte du
Lemme \ref{chow-lemma-equal-c1-as-cycles}.
\end{proof}






\section{La formule clé}
\label{chow-section-key}

\noindent
Soit $(S, \delta)$ comme dans la Situation \ref{chow-situation-setup}.
Soit $X$ localement de type fini sur $S$. Supposons
$X$ intègre et $\dim_\delta(X) = n$.
Soient $\mathcal{L}$ et $\mathcal{N}$ des faisceaux inversibles sur $X$.
Soit $s$ une section méromorphe non nulle de $\mathcal{L}$ et
soit $t$ une section méromorphe non nulle de $\mathcal{N}$.
Soit $Z_i \subset X$, $i \in I$, une famille localement finie de fermés
irréductibles de codimension $1$ ayant la propriété suivante :
si $Z \not \in \{Z_i\}$ a pour point générique $\xi$, alors $s$ est un générateur
de $\mathcal{L}_\xi$ et $t$ est un générateur de $\mathcal{N}_\xi$.
Une telle famille existe d'après le
Lemme \ref{divisors-lemma-divisor-meromorphic-locally-finite} de Diviseurs.
Alors
$$
\text{div}_\mathcal{L}(s) = \sum \text{ord}_{Z_i, \mathcal{L}}(s) [Z_i]
$$
et, de même,
$$
\text{div}_\mathcal{N}(t) = \sum \text{ord}_{Z_i, \mathcal{N}}(t) [Z_i]
$$
En développant davantage les définitions, choisissons, pour chaque $i$, des générateurs
$s_i \in \mathcal{L}_{\xi_i}$ et $t_i \in \mathcal{N}_{\xi_i}$,
où $\xi_i$ est le point générique de $Z_i$. Nous pouvons alors écrire
$$
s = f_i s_i
\quad\text{et}\quad
t = g_i t_i
$$
Posons $B_i = \mathcal{O}_{X, \xi_i}$. Alors, par définition,
$$
\text{ord}_{Z_i, \mathcal{L}}(s) = \text{ord}_{B_i}(f_i)
\quad\text{et}\quad
\text{ord}_{Z_i, \mathcal{N}}(t) = \text{ord}_{B_i}(g_i)
$$
Puisque $t_i$ est un générateur de $\mathcal{N}_{\xi_i}$, nous voyons que
son image dans la fibre $\mathcal{N}_{\xi_i} \otimes \kappa(\xi_i)$
est une section méromorphe non nulle de $\mathcal{N}|_{Z_i}$. Nous noterons
cette image $t_i|_{Z_i}$. Il résulte de nos définitions que
$$
c_1(\mathcal{N}) \cap \text{div}_\mathcal{L}(s) =
\sum \text{ord}_{B_i}(f_i)
(Z_i \to X)_*\text{div}_{\mathcal{N}|_{Z_i}}(t_i|_{Z_i})
$$
et, de même,
$$
c_1(\mathcal{L}) \cap \text{div}_\mathcal{N}(t) =
\sum \text{ord}_{B_i}(g_i)
(Z_i \to X)_*\text{div}_{\mathcal{L}|_{Z_i}}(s_i|_{Z_i})
$$
dans $\CH_{n - 2}(X)$. Nous allons trouver une équivalence rationnelle entre
ces deux cycles. Pour cela, nous considérons le symbole modéré
$$
\partial_{B_i}(f_i, g_i) \in \kappa(\xi_i)^*
$$
voir la section \ref{chow-section-tame-symbol}.

\begin{lemma}[Formule fondamentale]
\label{chow-lemma-key-formula}
Dans la situation ci-dessus, le cycle
$$
\sum
(Z_i \to X)_*\left(
\text{ord}_{B_i}(f_i) \text{div}_{\mathcal{N}|_{Z_i}}(t_i|_{Z_i}) -
\text{ord}_{B_i}(g_i) \text{div}_{\mathcal{L}|_{Z_i}}(s_i|_{Z_i}) \right)
$$
est égal au cycle
$$
\sum (Z_i \to X)_*\text{div}(\partial_{B_i}(f_i, g_i))
$$
\end{lemma}

\begin{proof}
Examinons d'abord ce qui se passe si l'on remplace $s_i$ par $us_i$,
où $u$ est une unité de $B_i$. Alors $f_i$ est remplacé par $u^{-1} f_i$.
La première partie de la première expression du lemme reste donc inchangée,
et l'on ajoute à la seconde partie
$$
-\text{ord}_{B_i}(g_i)\text{div}(u|_{Z_i})
$$
(où $u|_{Z_i}$ est l'image de $u$ dans le corps résiduel), d'après
Diviseurs, lemme \ref{divisors-lemma-divisor-meromorphic-well-defined},
tandis que l'on ajoute à la seconde expression
$$
\text{div}(\partial_{B_i}(u^{-1}, g_i))
$$
par bilinéarité du symbole modéré. Ces termes coïncident d'après la propriété
(\ref{chow-item-normalization}) du symbole modéré.

\medskip\noindent
Soit $Z \subset X$ une partie fermée irréductible telle que $\dim_\delta(Z) = n - 2$.
Pour montrer que les coefficients de $Z$ dans les deux cycles du lemme
coïncident, on peut effectuer un remplacement $s_i \mapsto us_i$ comme au
paragraphe précédent. De la même manière, on montre que l'on peut effectuer un remplacement
$t_i \mapsto vt_i$, où $v$ est une unité de $B_i$.

\medskip\noindent
Puisqu'il s'agit de démontrer une égalité de cycles, on peut raisonner coefficient
par coefficient. Choisissons donc une partie fermée irréductible $Z \subset X$
telle que $\dim_\delta(Z) = n - 2$, et comparons les coefficients. Soit $\xi \in Z$
son point générique, et posons $A = \mathcal{O}_{X, \xi}$. C'est un anneau local
noethérien intègre de dimension $2$. Choisissons des générateurs $\sigma$ et $\tau$
de $\mathcal{L}_\xi$ et $\mathcal{N}_\xi$. Après avoir restreint $X$, on peut
supposer, et l'on suppose, que $\sigma$ et $\tau$ définissent des trivialisations
des faisceaux inversibles $\mathcal{L}$ et $\mathcal{N}$ sur tout $X$.
La famille $Z_i$ étant localement
finie, on peut, après avoir restreint $X$, supposer $Z \subset Z_i$ pour tout $i \in I$
et $I$ fini. Alors $\xi_i$ correspond à un idéal premier
$\mathfrak q_i \subset A$ de hauteur $1$.
On peut écrire $s_i = a_i \sigma$ et $t_i = b_i \tau$,
où $a_i$ et $b_i$ sont des unités de $A_{\mathfrak q_i}$.
D'après les remarques précédentes, il suffit de démontrer le lemme lorsque
$a_i = b_i = 1$ pour tout $i$.

\medskip\noindent
Supposons $a_i = b_i = 1$ pour tout $i$. Alors la première expression du
lemme est nulle, car nous avons choisi $\sigma$ et $\tau$ comme sections
trivialisantes. Écrivons $s = f\sigma$ et $t = g \tau$, avec $f$ et $g$ dans le
corps des fractions de $A$. Le paragraphe précédent nous a ramenés au cas où
$f_i = f$ et $g_i = g$ pour tout $i$. De plus, pour tout idéal premier de hauteur $1$
$\mathfrak q$ de $A$ qui n'appartient pas à $\{\mathfrak q_i\}$,
$f$ et $g$ sont des unités de $A_\mathfrak q$ (d'après le choix de
la famille $\{Z_i\}$ dans la discussion précédant le lemme). Ainsi,
le coefficient de $Z$ dans la seconde expression du lemme est
$$
\sum\nolimits_i \text{ord}_{A/\mathfrak q_i}(\partial_{B_i}(f, g))
$$
qui est nul d'après le lemme fondamental \ref{chow-lemma-milnor-gersten-low-degree}.
\end{proof}

\begin{remark}
\label{chow-remark-higher-chow-pointwise}
Soit $(S, \delta)$ comme dans la situation \ref{chow-situation-setup}.
Soit $X$ localement de type fini sur $S$. Soit $k \in \mathbf{Z}$.
Nous affirmons qu'il existe un complexe
$$
\bigoplus\nolimits_{\delta(x) = k + 2}' K_2^M(\kappa(x))
\xrightarrow{\partial}
\bigoplus\nolimits_{\delta(x) = k + 1}' K_1^M(\kappa(x))
\xrightarrow{\partial}
\bigoplus\nolimits_{\delta(x) = k}' K_0^M(\kappa(x))
$$
Nous utilisons ici les notations et conventions introduites dans la
remarque \ref{chow-remark-chow-group-pointwise}, avec en outre les précisions suivantes :
\begin{enumerate}
\item $K_2^M(\kappa(x))$ est la composante de degré $2$ de
la K-théorie de Milnor du corps résiduel $\kappa(x)$ du point
$x \in X$ (voir la remarque \ref{chow-remark-gersten-complex-milnor}) ; c'est
le quotient de $\kappa(x)^* \otimes_\mathbf{Z} \kappa(x)^*$
par le sous-groupe engendré par les éléments de la forme
$\lambda \otimes (1 - \lambda)$, pour
$\lambda \in \kappa(x) \setminus \{0, 1\}$, et
\item la première différentielle $\partial$ est définie comme suit :
étant donné un élément $\xi = \sum_x \alpha_x$ du premier terme,
on pose
$$
\partial(\xi) = \sum\nolimits_{x \leadsto x',\ \delta(x') = k + 1}
\partial_{\mathcal{O}_{W_x, x'}}(\alpha_x)
$$
où
$\partial_{\mathcal{O}_{W_x, x'}} : K_2^M(\kappa(x)) \to K_1^M(\kappa(x))$
est le symbole modéré construit dans la section \ref{chow-section-tame-symbol}.
\end{enumerate}
Nous affirmons que l'on obtient un complexe, c'est-à-dire que $\partial \circ \partial = 0$.
Pour le voir, il suffit de prendre un élément $\xi$ comme ci-dessus et un point
$x'' \in X$ tel que $\delta(x'') = k$, et de vérifier que le coefficient de
$x''$ dans l'élément $\partial(\partial(\xi))$ est nul.
Comme $\xi = \sum \alpha_x$ est une somme localement finie, on
peut en fait supposer, par additivité, que $\xi = \alpha_x$ pour
un certain $x \in X$ tel que $\delta(x) = k + 2$ et $\alpha_x \in K_2^M(\kappa(x))$.
Par linéarité, on peut encore supposer que $\alpha_x = f \otimes g$ pour
certains $f, g \in \kappa(x)^*$. Notons $W \subset X$ le sous-schéma fermé
intègre de point générique $x$. Si $x'' \not \in W$, il est
immédiatement clair que le coefficient de $x$ dans $\partial(\partial(\xi))$
est nul. Si $x'' \in W$, on voit que le coefficient de $x''$
dans $\partial(\partial(x))$ est égal à
$$
\sum\nolimits_{x \leadsto x' \leadsto x'',\ \delta(x') = k + 1}
\text{ord}_{\mathcal{O}_{\overline{\{x'\}}, x''}}(
\partial_{\mathcal{O}_{W, x'}}(f, g))
$$
Le lemme algébrique fondamental \ref{chow-lemma-milnor-gersten-low-degree}
affirme précisément que cette somme est nulle.
\end{remark}

\begin{remark}
\label{chow-remark-higher-chow}
Soit $(S, \delta)$ comme dans la situation \ref{chow-situation-setup}.
Soit $X$ localement de type fini sur $S$. Soit $k \in \mathbf{Z}$.
Le complexe de la remarque \ref{chow-remark-higher-chow-pointwise} et la
présentation de $\CH_k(X)$ dans la remarque \ref{chow-remark-chow-group-pointwise}
suggèrent de définir un premier groupe de Chow supérieur par
$$
\CH^M_k(X, 1) =
H_1(\text{le complexe de la remarque \ref{chow-remark-higher-chow-pointwise}})
$$
Nous utilisons l'exposant ${}^M$ pour distinguer notre notation de celle des
groupes de Chow supérieurs définis dans la littérature, par exemple dans les articles
de Spencer Bloch (\cite{Bloch} et \cite{Bloch-moving}).
Soit $U \subset X$ un ouvert dont le complémentaire est $Y \subset X$ (muni de sa structure de
sous-schéma fermé réduit). On obtient alors une suite exacte courte scindée
$$
0 \to
\bigoplus\nolimits_{y \in Y, \delta(y) = k + i}' K_i^M(\kappa(y)) \to
\bigoplus\nolimits_{x \in X, \delta(x) = k + i}' K_i^M(\kappa(x)) \to
\bigoplus\nolimits_{u \in U, \delta(u) = k + i}' K_i^M(\kappa(u)) \to 0
$$
pour $i = 2, 1, 0$, compatible avec les opérateurs de bord des complexes
de la remarque \ref{chow-remark-higher-chow-pointwise}. En appliquant le lemme du serpent
(voir Homologie, lemme \ref{homology-lemma-long-exact-sequence-chain}),
on obtient une suite exacte à six termes
$$
\CH^M_k(Y, 1) \to \CH^M_k(X, 1) \to \CH^M_k(U, 1) \to
\CH_k(Y) \to \CH_k(X) \to \CH_k(U) \to 0
$$
prolongeant la suite exacte canonique du lemme \ref{chow-lemma-restrict-to-open}.
Au prix de quelques efforts, on peut aussi définir l'image réciproque plate et l'image directe propre
pour le premier groupe de Chow supérieur $\CH^M_k(X, 1)$. Nous y reviendrons
plus loin (insérer ici un renvoi ultérieur).
\end{remark}






\section{Intersection avec un faisceau inversible et équivalence rationnelle}
\label{chow-section-commutativity}

\noindent
L'application du lemme fondamental donne les propriétés essentielles de l'intersection
avec les faisceaux inversibles. Nous verrons en particulier que
$c_1(\mathcal{L}) \cap -$ passe au quotient par l'équivalence rationnelle et
que ces opérations commutent pour différents faisceaux inversibles.

\begin{lemma}
\label{chow-lemma-commutativity-on-integral}
Soit $(S, \delta)$ comme dans la situation \ref{chow-situation-setup}.
Soit $X$ localement de type fini sur $S$.
Supposons $X$ intègre et $\dim_\delta(X) = n$.
Soient $\mathcal{L}$, $\mathcal{N}$ des faisceaux inversibles sur $X$.
Choisissons une section méromorphe non nulle $s$ de $\mathcal{L}$
et une section méromorphe non nulle $t$ de $\mathcal{N}$.
Posons $\alpha = \text{div}_\mathcal{L}(s)$ et
$\beta = \text{div}_\mathcal{N}(t)$.
Alors
$$
c_1(\mathcal{N}) \cap \alpha
=
c_1(\mathcal{L}) \cap \beta
$$
dans $\CH_{n - 2}(X)$.
\end{lemma}

\begin{proof}
Cela découle immédiatement du lemme fondamental \ref{chow-lemma-key-formula}
et de la discussion qui le précède.
\end{proof}

\begin{lemma}
\label{chow-lemma-factors}
Soit $(S, \delta)$ comme dans la situation \ref{chow-situation-setup}.
Soit $X$ localement de type fini sur $S$.
Soit $\mathcal{L}$ un faisceau inversible sur $X$.
L'opération $\alpha \mapsto c_1(\mathcal{L}) \cap \alpha$
passe au quotient par l'équivalence rationnelle et donne une opération
$$
c_1(\mathcal{L}) \cap - : \CH_{k + 1}(X) \to \CH_k(X)
$$
\end{lemma}

\begin{proof}
Soit $\alpha \in Z_{k + 1}(X)$, avec $\alpha \sim_{rat} 0$.
Il faut montrer que $c_1(\mathcal{L}) \cap \alpha$,
au sens de la définition \ref{chow-definition-cap-c1}, est nul.
D'après la définition \ref{chow-definition-rational-equivalence}, il
existe une famille localement finie $\{W_j\}$ de sous-schémas fermés
intègres tels que $\dim_\delta(W_j) = k + 2$, et des fonctions rationnelles
$f_j \in R(W_j)^*$ telles que
$$
\alpha = \sum (i_j)_*\text{div}_{W_j}(f_j)
$$
Remarquons que $p : \coprod W_j \to X$ est un morphisme propre ;
on a donc $\alpha = p_*\alpha'$, où $\alpha' \in Z_{k + 1}(\coprod W_j)$
est la somme des diviseurs principaux $\text{div}_{W_j}(f_j)$.
D'après le lemme \ref{chow-lemma-pushforward-cap-c1}, on a
$c_1(\mathcal{L}) \cap \alpha = p_*(c_1(p^*\mathcal{L}) \cap \alpha')$.
Il suffit donc de montrer que chacun des
$c_1(\mathcal{L}|_{W_j}) \cap \text{div}_{W_j}(f_j)$ est nul.
Autrement dit, on peut supposer $X$ intègre et
$\alpha = \text{div}_X(f)$ pour un certain $f \in R(X)^*$.

\medskip\noindent
Supposons $X$ intègre et $\alpha = \text{div}_X(f)$ pour un certain $f \in R(X)^*$.
On peut considérer $f$ comme une section méromorphe régulière du faisceau
inversible $\mathcal{N} = \mathcal{O}_X$. Choisissons une section méromorphe
$s$ de $\mathcal{L}$ et notons $\beta = \text{div}_\mathcal{L}(s)$.
D'après le lemme \ref{chow-lemma-commutativity-on-integral},
on en conclut que
$$
c_1(\mathcal{L}) \cap \alpha = c_1(\mathcal{O}_X) \cap \beta.
$$
Or le lemme \ref{chow-lemma-c1-cap-additive} montre que le membre de droite
est nul dans $\CH_k(X)$, comme voulu.
\end{proof}

\noindent
Soit $(S, \delta)$ comme dans la situation \ref{chow-situation-setup}.
Soit $X$ localement de type fini sur $S$.
Soit $\mathcal{L}$ un faisceau inversible sur $X$.
Nous noterons
$$
c_1(\mathcal{L}) \cap - : \CH_{k + 1}(X) \to \CH_k(X)
$$
l'opération $c_1(\mathcal{L}) \cap - $. Cette notation a un sens d'après le
lemme \ref{chow-lemma-factors}. Nous noterons $c_1(\mathcal{L})^s \cap -$
la $s$-ième itérée de cette opération pour tout $s \geq 0$.

\begin{lemma}
\label{chow-lemma-cap-commutative}
Soit $(S, \delta)$ comme dans la situation \ref{chow-situation-setup}.
Soit $X$ localement de type fini sur $S$.
Soient $\mathcal{L}$, $\mathcal{N}$ des faisceaux inversibles sur $X$.
Pour tout $\alpha \in \CH_{k + 2}(X)$, on a
$$
c_1(\mathcal{L}) \cap c_1(\mathcal{N}) \cap \alpha
=
c_1(\mathcal{N}) \cap c_1(\mathcal{L}) \cap \alpha
$$
comme éléments de $\CH_k(X)$.
\end{lemma}

\begin{proof}
Écrivons $\alpha = \sum m_j[Z_j]$ pour une famille localement finie
de sous-schémas fermés intègres $Z_j \subset X$
tels que $\dim_\delta(Z_j) = k + 2$.
Considérons le morphisme propre $p : \coprod Z_j \to X$.
Posons $\alpha' = \sum m_j[Z_j]$, considéré comme un $(k + 2)$-cycle sur
$\coprod Z_j$. Plusieurs applications du
lemme \ref{chow-lemma-pushforward-cap-c1} donnent
$c_1(\mathcal{L}) \cap c_1(\mathcal{N}) \cap \alpha
= p_*(c_1(p^*\mathcal{L}) \cap c_1(p^*\mathcal{N}) \cap \alpha')$
et
$c_1(\mathcal{N}) \cap c_1(\mathcal{L}) \cap \alpha
= p_*(c_1(p^*\mathcal{N}) \cap c_1(p^*\mathcal{L}) \cap \alpha')$.
Il suffit donc de démontrer la formule lorsque $X$ est intègre
et $\alpha = [X]$. Dans ce cas, le résultat découle
du lemme \ref{chow-lemma-commutativity-on-integral} et des définitions.
\end{proof}






\section{Homomorphismes de Gysin}
\label{chow-section-intersecting-effective-Cartier}

\noindent
Dans cette section, nous définissons l'application de Gysin pour le lieu des zéros $D$ d'une
section d'un faisceau inversible. Le cas où $D$
est un diviseur de Cartier effectif est intéressant, mais la généralisation à un $D$ quelconque
donne la souplesse nécessaire pour formuler diverses compatibilités ; voir la
remarque \ref{chow-remark-pullback-pairs} et les
lemmes \ref{chow-lemma-closed-in-X-gysin}, \ref{chow-lemma-gysin-flat-pullback} et
\ref{chow-lemma-gysin-commutes-gysin}.
Ces résultats se généralisent aux sous-schémas fermés localement principaux
munis d'un fibré normal virtuel
(remarque \ref{chow-remark-generalize-to-virtual}) ou aux
pseudo-diviseurs (remarque \ref{chow-remark-generalize-to-pseudo-divisor}).

\medskip\noindent
Rappelons que les diviseurs de Cartier effectifs correspondent $1$ à $1$ aux
classes d'isomorphisme de couples $(\mathcal{L}, s)$, où $\mathcal{L}$
est un faisceau inversible et $s$ une section globale régulière ; voir
Diviseurs, lemme \ref{divisors-lemma-characterize-OD}.
Si $D$ correspond à $(\mathcal{L}, s)$, alors
$\mathcal{L} = \mathcal{O}_X(D)$. On gardera ce fait à l'esprit
dans la lecture de cette section.

\begin{definition}
\label{chow-definition-gysin-homomorphism}
Soit $(S, \delta)$ comme dans la situation \ref{chow-situation-setup}.
Soit $X$ localement de type fini sur $S$.
Soit $(\mathcal{L}, s)$ un couple formé d'un faisceau
inversible et d'une section globale $s \in \Gamma(X, \mathcal{L})$.
Soit $D = Z(s)$ le schéma des zéros de $s$, et
notons $i : D \to X$ l'immersion fermée.
Nous définissons, pour tout entier $k$, un {\it homomorphisme de Gysin}
$$
i^* : Z_{k + 1}(X) \to \CH_k(D).
$$
par les règles suivantes :
\begin{enumerate}
\item Étant donné un sous-schéma fermé intègre $W \subset X$ tel que
$\dim_\delta(W) = k + 1$, on définit
\begin{enumerate}
\item si $W \not \subset D$, alors $i^*[W] = [D \cap W]_k$, considéré comme
$k$-cycle sur $D$, et
\item si $W \subset D$, alors
$i^*[W] = i'_*(c_1(\mathcal{L}|_W) \cap [W])$,
où $i' : W \to D$ est l'immersion fermée induite.
\end{enumerate}
\item Pour un $(k + 1)$-cycle quelconque $\alpha = \sum n_j[W_j]$,
on pose
$$
i^*\alpha = \sum n_j i^*[W_j]
$$
\item Si $D$ est un diviseur de Cartier effectif, on note
$D \cdot \alpha = i_*i^*\alpha$ l'image directe de la classe $i^*\alpha$,
considérée comme classe sur $X$.
\end{enumerate}
\end{definition}

\noindent
En fait, comme nous le verrons plus loin, cet homomorphisme de Gysin $i^*$ peut être considéré
comme un exemple d'image réciproque non plate. Nous appellerons donc parfois, de manière informelle,
la classe $i^*\alpha$ l'{\it image réciproque} de la classe $\alpha$.

\begin{remark}
\label{chow-remark-generalize-to-virtual}
Soit $X$ un schéma localement de type fini sur $S$ comme dans la
situation \ref{chow-situation-setup}. Soit $(D, \mathcal{N}, \sigma)$
un triplet formé d'un sous-schéma fermé localement principal (Diviseurs, définition
\ref{divisors-definition-effective-Cartier-divisor})
$i : D \to X$, d'un $\mathcal{O}_D$-module inversible $\mathcal{N}$ et
d'une surjection $\sigma : \mathcal{N}^{\otimes -1} \to i^*\mathcal{I}_D$
de $\mathcal{O}_D$-modules\footnote{Cette condition assure que, si
$D$ est un diviseur de Cartier effectif, alors $\mathcal{N} = \mathcal{O}_X(D)|_D$.}.
Il faut ici considérer $\mathcal{N}$ comme
un {\it fibré normal virtuel de $D$ dans $X$}. La construction de
$i^* : Z_{k + 1}(X) \to \CH_k(D)$ dans la
définition \ref{chow-definition-gysin-homomorphism}
se généralise à de tels triplets ; voir la
section \ref{chow-section-gysin-higher-codimension}.
\end{remark}

\begin{remark}
\label{chow-remark-generalize-to-pseudo-divisor}
Soit $X$ un schéma localement de type fini sur $S$ comme dans la
situation \ref{chow-situation-setup}. Dans \cite{F}, un {\it pseudo-diviseur} sur $X$
est défini comme un triplet $D = (\mathcal{L}, Z, s)$, où $\mathcal{L}$
est un $\mathcal{O}_X$-module inversible, $Z \subset X$ une partie fermée,
et $s \in \Gamma(X \setminus Z, \mathcal{L})$ une section qui ne s'annule
nulle part. Comme ci-dessus, on peut définir pour tout $\alpha$
dans $\CH_{k + 1}(X)$ un produit $D \cdot \alpha$ dans $\CH_k(Z \cap |\alpha|)$,
où $|\alpha|$ est le support de $\alpha$.
\end{remark}

\begin{lemma}
\label{chow-lemma-support-cap-effective-Cartier}
Soit $(S, \delta)$ comme dans la situation \ref{chow-situation-setup}. Soit $X$ localement
de type fini sur $S$. Soit $(\mathcal{L}, s, i : D \to X)$ comme dans la
définition \ref{chow-definition-gysin-homomorphism}. Soit $\alpha$ un
$(k + 1)$-cycle sur $X$. Alors $i_*i^*\alpha = c_1(\mathcal{L}) \cap \alpha$
dans $\CH_k(X)$. En particulier, si $D$ est un diviseur de Cartier effectif, alors
$D \cdot \alpha = c_1(\mathcal{O}_X(D)) \cap \alpha$.
\end{lemma}

\begin{proof}
Écrivons $\alpha = \sum n_j[W_j]$, où $i_j : W_j \to X$ sont des sous-schémas fermés
intègres tels que $\dim_\delta(W_j) = k$.
Puisque $D$ est le schéma des zéros de $s$, on voit que $D \cap W_j$ est le schéma des zéros
de la restriction $s|_{W_j}$. Ainsi, pour tout $j$ tel que
$W_j \not \subset D$, on a
$c_1(\mathcal{L}) \cap [W_j] = [D \cap W_j]_k$
d'après le lemme \ref{chow-lemma-geometric-cap}. On a donc
$$
c_1(\mathcal{L}) \cap \alpha
=
\sum\nolimits_{W_j \not \subset D} n_j[D \cap W_j]_k
+
\sum\nolimits_{W_j \subset D}
n_j i_{j, *}(c_1(\mathcal{L})|_{W_j}) \cap [W_j])
$$
dans $\CH_k(X)$, d'après la définition \ref{chow-definition-cap-c1}.
Le membre de droite coïncide terme à terme avec l'image directe de la classe
$i^*\alpha$ sur $D$ de la définition \ref{chow-definition-gysin-homomorphism}.
On obtient donc le résultat.
\end{proof}

\begin{lemma}
\label{chow-lemma-easy-gysin}
Soit $(S, \delta)$ comme dans la situation \ref{chow-situation-setup}.
Soit $X$ localement de type fini sur $S$.
Soit $(\mathcal{L}, s, i : D \to X)$ comme dans la
définition \ref{chow-definition-gysin-homomorphism}.
\begin{enumerate}
\item Soit $Z \subset X$ un sous-schéma fermé tel
que $\dim_\delta(Z) \leq k + 1$ et tel que
$D \cap Z$ soit un diviseur de Cartier effectif sur $Z$. Alors
$i^*[Z]_{k + 1} = [D \cap Z]_k$.
\item Soit $\mathcal{F}$ un faisceau cohérent sur $X$
tel que $\dim_\delta(\text{Supp}(\mathcal{F})) \leq k + 1$ et que
$s : \mathcal{F} \to \mathcal{F} \otimes_{\mathcal{O}_X} \mathcal{L}$
soit injectif. Alors
$$
i^*[\mathcal{F}]_{k + 1} = [i^*\mathcal{F}]_k
$$
dans $\CH_k(D)$.
\end{enumerate}
\end{lemma}

\begin{proof}
Soit $Z \subset X$ comme dans (1). Posons $\mathcal{F} = \mathcal{O}_Z$.
L'hypothèse selon laquelle $D \cap Z$ est un diviseur de Cartier effectif
équivaut à l'injectivité de
$s : \mathcal{F} \to \mathcal{F} \otimes_{\mathcal{O}_X} \mathcal{L}$.
De plus, $[Z]_{k + 1} = [\mathcal{F}]_{k + 1}]$
et $[D \cap Z]_k = [\mathcal{O}_{D \cap Z}]_k = [i^*\mathcal{F}]_k$.
Voir le lemme \ref{chow-lemma-cycle-closed-coherent}.
L'assertion (1) découle donc de (2).

\medskip\noindent
Écrivons $[\mathcal{F}]_{k + 1} = \sum m_j[W_j]$, avec $m_j > 0$
et des sous-schémas fermés intègres deux à deux distincts $W_j \subset X$
de $\delta$-dimension $k + 1$. L'hypothèse d'injectivité de
$s : \mathcal{F} \to \mathcal{F} \otimes_{\mathcal{O}_X} \mathcal{L}$
entraîne que $W_j \not \subset D$ pour tout $j$.
Par définition, on a
$$
i^*[\mathcal{F}]_{k + 1} = \sum m_j [D \cap W_j]_k.
$$
Nous affirmons que
$$
\sum [D \cap W_j]_k = [i^*\mathcal{F}]_k
$$
comme cycles.
Soit $Z \subset D$ un sous-schéma fermé intègre de $\delta$-dimension
$k$. Soit $\xi \in Z$ son point générique. Soit $A = \mathcal{O}_{X, \xi}$.
Soit $M = \mathcal{F}_\xi$. Soit $f \in A$ un élément engendrant
l'idéal de $D$, c'est-à-dire tel que $\mathcal{O}_{D, \xi} = A/fA$.
Par hypothèse, $\dim(\text{Supp}(M)) = 1$,
l'application $f : M \to M$ est injective, et
$\text{longueur}_A(M/fM) < \infty$. De plus, $\text{longueur}_A(M/fM)$
est le coefficient de $[Z]$ dans $[i^*\mathcal{F}]_k$. D'autre
part, soient $\mathfrak q_1, \ldots, \mathfrak q_t$ les idéaux premiers minimaux
du support de $M$. Alors
$$
\sum
\text{longueur}_{A_{\mathfrak q_i}}(M_{\mathfrak q_i})
\text{ord}_{A/\mathfrak q_i}(f)
$$
est le coefficient de $[Z]$ dans $\sum [D \cap W_j]_k$.
L'égalité résulte donc du
lemme \ref{chow-lemma-additivity-divisors-restricted}.
\end{proof}

\begin{remark}
\label{chow-remark-gysin-on-cycles}
Soient $X \to S$, $\mathcal{L}$, $s$, $i : D \to X$ comme dans la
définition \ref{chow-definition-gysin-homomorphism}, et supposons
que $\mathcal{L}|_D \cong \mathcal{O}_D$. Dans ce cas, on
peut définir une application canonique $i^* : Z_{k + 1}(X) \to Z_k(D)$
sur les cycles, en imposant $i^*[W] = 0$ lorsque $W \subset D$
est un sous-schéma fermé intègre.
Cette possibilité nous sera utile par la suite.
\end{remark}

\begin{remark}
\label{chow-remark-pullback-pairs}
Soit $f : X' \to X$ un morphisme de schémas localement de type fini sur $S$
comme dans la situation \ref{chow-situation-setup}. Soit $(\mathcal{L}, s, i : D \to X)$
un triplet comme dans la définition \ref{chow-definition-gysin-homomorphism}.
On peut alors poser $\mathcal{L}' = f^*\mathcal{L}$, $s' = f^*s$ et
$D' = X' \times_X D = Z(s')$. Cela donne un diagramme commutatif
$$
\xymatrix{
D' \ar[d]_g \ar[r]_{i'} & X' \ar[d]^f \\
D \ar[r]^i & X
}
$$
et l'on peut étudier diverses compatibilités entre $i^*$ et $(i')^*$.
\end{remark}

\begin{lemma}
\label{chow-lemma-closed-in-X-gysin}
Soit $(S, \delta)$ comme dans la situation \ref{chow-situation-setup}.
Soit $f : X' \to X$ un morphisme propre de schémas
localement de type fini sur $S$.
Soit $(\mathcal{L}, s, i : D \to X)$ comme dans la
définition \ref{chow-definition-gysin-homomorphism}.
Formons le diagramme
$$
\xymatrix{
D' \ar[d]_g \ar[r]_{i'} & X' \ar[d]^f \\
D \ar[r]^i & X
}
$$
comme dans la remarque \ref{chow-remark-pullback-pairs}.
Pour tout $(k + 1)$-cycle $\alpha'$ sur $X'$, on a
$i^*f_*\alpha' = g_*(i')^*\alpha'$ dans $\CH_k(D)$
(ce qui a un sens, puisque $f_*$ est défini au niveau des cycles).
\end{lemma}

\begin{proof}
Supposons $\alpha = [W']$ pour un sous-schéma fermé intègre
$W' \subset X'$. Soit $W = f(W') \subset X$. Si $W' \not \subset D'$,
alors $W \not \subset D$, et l'on voit que
$$
[W' \cap D']_k = \text{div}_{\mathcal{L}'|_{W'}}({s'|_{W'}})
\quad\text{et}\quad
[W \cap D]_k = \text{div}_{\mathcal{L}|_W}(s|_W)
$$
et donc que l'image du premier cycle par $f_*$ est égale au second cycle, d'après le
lemme \ref{chow-lemma-equal-c1-as-cycles}. L'égalité
est donc vraie comme égalité de cycles. Si $W' \subset D'$, alors
$W \subset D$ et $f_*(c_1(\mathcal{L}|_{W'}) \cap [W'])$
est égal à $c_1(\mathcal{L}|_W) \cap [W]$ dans $\CH_k(W)$, d'après la seconde
assertion du lemme \ref{chow-lemma-equal-c1-as-cycles}.
D'après la remarque \ref{chow-remark-infinite-sums-rational-equivalences},
le résultat s'ensuit pour un $\alpha'$ quelconque.
\end{proof}

\begin{lemma}
\label{chow-lemma-gysin-flat-pullback}
Soit $(S, \delta)$ comme dans la situation \ref{chow-situation-setup}. Soit $f : X' \to X$
un morphisme plat de dimension relative $r$ entre schémas localement de type fini
sur $S$. Soit $(\mathcal{L}, s, i : D \to X)$ comme dans la
définition \ref{chow-definition-gysin-homomorphism}. Formons le diagramme
$$
\xymatrix{
D' \ar[d]_g \ar[r]_{i'} & X' \ar[d]^f \\
D \ar[r]^i & X
}
$$
comme dans la remarque \ref{chow-remark-pullback-pairs}.
Pour tout $(k + 1)$-cycle $\alpha$ sur $X$, on a
$(i')^*f^*\alpha = g^*i^*\alpha$ dans $\CH_{k + r}(D')$
(ce qui a un sens, puisque $f^*$ est défini au niveau des cycles).
\end{lemma}

\begin{proof}
Supposons $\alpha = [W]$ pour un sous-schéma fermé intègre
$W \subset X$. Soit $W' = f^{-1}(W) \subset X'$. Si $W \not \subset D$,
alors $W' \not \subset D'$, et l'on voit que
$$
W' \cap D' = g^{-1}(W \cap D)
$$
comme sous-schémas fermés de $D'$. L'égalité
est donc vraie comme égalité de cycles ; voir le lemme \ref{chow-lemma-pullback-coherent}.
Si $W \subset D$, alors $W' \subset D'$ et $W' = g^{-1}(W)$,
avec $[W']_{k + 1 + r} = g^*[W]$, et l'égalité est vraie dans
$\CH_{k + r}(D')$ d'après le lemme \ref{chow-lemma-flat-pullback-cap-c1}.
D'après la remarque \ref{chow-remark-infinite-sums-rational-equivalences},
le résultat s'ensuit pour un $\alpha'$ quelconque.
\end{proof}








\section{Homomorphismes de Gysin et équivalence rationnelle}
\label{chow-section-gysin}

\noindent
Dans cette section, nous utilisons la formule fondamentale pour montrer que l'homomorphisme
de Gysin passe au quotient par l'équivalence rationnelle. Nous démontrons aussi une importante
propriété de commutativité.

\begin{lemma}
\label{chow-lemma-gysin-factors-general}
Soit $(S, \delta)$ comme dans la situation \ref{chow-situation-setup}.
Soit $X$ localement de type fini sur $S$.
Supposons $X$ intègre et $n = \dim_\delta(X)$.
Soit $i : D \to X$ un diviseur de Cartier effectif.
Soit $\mathcal{N}$ un $\mathcal{O}_X$-module inversible,
et soit $t$ une section méromorphe non nulle de $\mathcal{N}$.
Alors $i^*\text{div}_\mathcal{N}(t) = c_1(\mathcal{N}|_D) \cap [D]_{n - 1}$
dans $\CH_{n - 2}(D)$.
\end{lemma}

\begin{proof}
Écrivons $\text{div}_\mathcal{N}(t) = \sum \text{ord}_{Z_i, \mathcal{N}}(t)[Z_i]$
pour des sous-schémas fermés intègres $Z_i \subset X$ de $\delta$-dimension
$n - 1$. On peut supposer que la famille $\{Z_i\}$ est localement
finie, que $t \in \Gamma(U, \mathcal{N}|_U)$ est un générateur,
où $U = X \setminus \bigcup Z_i$, et que toute composante irréductible
de $D$ est l'un des $Z_i$ ; voir
Diviseurs, lemmes \ref{divisors-lemma-components-locally-finite},
\ref{divisors-lemma-divisor-locally-finite} et
\ref{divisors-lemma-divisor-meromorphic-locally-finite}.

\medskip\noindent
Posons $\mathcal{L} = \mathcal{O}_X(D)$. Notons
$s \in \Gamma(X, \mathcal{O}_X(D)) = \Gamma(X, \mathcal{L})$
la section canonique. Nous allons appliquer la discussion de la
section \ref{chow-section-key} à la situation présente.
Pour chaque $i$, soit $\xi_i \in Z_i$ son point générique. Soit
$B_i = \mathcal{O}_{X, \xi_i}$. Pour chaque $i$, choisissons des générateurs
$s_i \in \mathcal{L}_{\xi_i}$ et $t_i \in \mathcal{N}_{\xi_i}$
sur $B_i$, en imposant toutefois $s_i = s$ si $Z_i \not \subset D$.
Écrivons $s = f_i s_i$ et $t = g_i t_i$, avec $f_i, g_i \in B_i$.
Alors $\text{ord}_{Z_i, \mathcal{N}}(t) = \text{ord}_{B_i}(g_i)$.
D'autre part, on a $f_i \in B_i$ et
$$
[D]_{n - 1} = \sum \text{ord}_{B_i}(f_i)[Z_i]
$$
d'après le choix des $s_i$. Nous affirmons que
$$
i^*\text{div}_\mathcal{N}(t) =
\sum \text{ord}_{B_i}(g_i) \text{div}_{\mathcal{L}|_{Z_i}}(s_i|_{Z_i})
$$
comme cycles. Plus précisément, le membre de droite est un cycle
représentant le membre de gauche. Cela résulte en effet de notre
formule pour $\text{div}_\mathcal{N}(t)$ et du fait que
$\text{div}_{\mathcal{L}|_{Z_i}}(s_i|_{Z_i}) = [Z(s_i|_{Z_i})]_{n - 2} =
[Z_i \cap D]_{n - 2}$ lorsque $Z_i \not \subset D$, car dans
ce cas $s_i|_{Z_i} = s|_{Z_i}$ est une section régulière ; voir le
lemme \ref{chow-lemma-compute-c1}. De même,
$$
c_1(\mathcal{N}) \cap [D]_{n - 1} =
\sum \text{ord}_{B_i}(f_i) \text{div}_{\mathcal{N}|_{Z_i}}(t_i|_{Z_i})
$$
La formule fondamentale (lemme \ref{chow-lemma-key-formula}) donne l'égalité
$$
\sum \left(
\text{ord}_{B_i}(f_i) \text{div}_{\mathcal{N}|_{Z_i}}(t_i|_{Z_i}) -
\text{ord}_{B_i}(g_i) \text{div}_{\mathcal{L}|_{Z_i}}(s_i|_{Z_i}) \right) =
\sum \text{div}_{Z_i}(\partial_{B_i}(f_i, g_i))
$$
de cycles. Si $Z_i \not \subset D$, alors $f_i = 1$, et donc
$\text{div}_{Z_i}(\partial_{B_i}(f_i, g_i)) = 0$. On obtient ainsi une équivalence
rationnelle entre les cycles particuliers représentant
$i^*\text{div}_\mathcal{N}(t)$ et $c_1(\mathcal{N}) \cap [D]_{n - 1}$
sur $D$. Cela achève la démonstration.
\end{proof}

\begin{lemma}
\label{chow-lemma-gysin-factors}
Soit $(S, \delta)$ comme dans la situation \ref{chow-situation-setup}.
Soit $X$ localement de type fini sur $S$.
Soit $(\mathcal{L}, s, i : D \to X)$ comme dans la
définition \ref{chow-definition-gysin-homomorphism}.
L'homomorphisme de Gysin passe au quotient par l'équivalence rationnelle et
donne une application $i^* : \CH_{k + 1}(X) \to \CH_k(D)$.
\end{lemma}

\begin{proof}
Soit $\alpha \in Z_{k + 1}(X)$, et supposons $\alpha \sim_{rat} 0$.
Cela signifie qu'il existe une famille localement finie de sous-schémas
fermés intègres $W_j \subset X$ de $\delta$-dimension $k + 2$
et des $f_j \in R(W_j)^*$ tels que
$\alpha = \sum i_{j, *}\text{div}_{W_j}(f_j)$.
Posons $X' = \coprod W_i$ et considérons le diagramme
$$
\xymatrix{
D' \ar[d]_q \ar[r]_{i'} & X' \ar[d]^p \\
D \ar[r]^i & X
}
$$
de la remarque \ref{chow-remark-pullback-pairs}. Puisque $X' \to X$ est propre,
on a $i^*p_* = q_*(i')^*$ d'après le lemme \ref{chow-lemma-closed-in-X-gysin}.
Comme on sait que $q_*$ passe au quotient par l'équivalence rationnelle
(lemme \ref{chow-lemma-proper-pushforward-rational-equivalence}), il suffit
de démontrer le résultat pour $\alpha' = \sum \text{div}_{W_j}(f_j)$
sur $X'$. Cela nous ramène clairement au cas où $X$ est intègre
et $\alpha = \text{div}(f)$ pour un certain $f \in R(X)^*$.

\medskip\noindent
Supposons $X$ intègre et $\alpha = \text{div}(f)$ pour un certain $f \in R(X)^*$.
Si $X = D$, on voit que $i^*\alpha$ est égal
à $c_1(\mathcal{L}) \cap \alpha$.
Cette classe est rationnellement équivalente à zéro d'après le lemme \ref{chow-lemma-factors}.
Si $D \not = X$, on voit que $i^*\text{div}_X(f)$ est égal à
$c_1(\mathcal{O}_D) \cap [D]_{n - 1}$ dans $\CH_{n - 2}(D)$ d'après le
lemme \ref{chow-lemma-gysin-factors-general}. Bien entendu,
le produit cap avec $c_1(\mathcal{O}_D)$ est l'application nulle
(lemme \ref{chow-lemma-c1-cap-additive}).
\end{proof}

\begin{lemma}
\label{chow-lemma-gysin-back}
Soit $(S, \delta)$ comme dans la situation \ref{chow-situation-setup}. Soit $X$ localement
de type fini sur $S$. Soit $(\mathcal{L}, s, i : D \to X)$ comme dans la
définition \ref{chow-definition-gysin-homomorphism}. Alors
$i^*i_* : \CH_k(D) \to \CH_{k - 1}(D)$ envoie $\alpha$ sur
$c_1(\mathcal{L}|_D) \cap \alpha$.
\end{lemma}

\begin{proof}
Cela découle immédiatement de la définition de $i_*$ sur les cycles
et de celle de $i^*$ donnée dans la
définition \ref{chow-definition-gysin-homomorphism}.
\end{proof}

\begin{lemma}
\label{chow-lemma-gysin-commutes-cap-c1}
Soit $(S, \delta)$ comme dans la situation \ref{chow-situation-setup}. Soit $X$
localement de type fini sur $S$. Soit $(\mathcal{L}, s, i : D \to X)$
un triplet comme dans la définition \ref{chow-definition-gysin-homomorphism}.
Soit $\mathcal{N}$ un $\mathcal{O}_X$-module inversible.
Alors $i^*(c_1(\mathcal{N}) \cap \alpha) = c_1(i^*\mathcal{N}) \cap i^*\alpha$
dans $\CH_{k - 2}(D)$ pour tout $\alpha \in \CH_k(X)$.
\end{lemma}

\begin{proof}
Par une démonstration identique à celle du lemme \ref{chow-lemma-gysin-factors},
cela résulte des lemmes
\ref{chow-lemma-pushforward-cap-c1},
\ref{chow-lemma-cap-commutative} et
\ref{chow-lemma-gysin-factors-general}.
\end{proof}

\begin{lemma}
\label{chow-lemma-gysin-commutes-gysin}
Soit $(S, \delta)$ comme dans la situation \ref{chow-situation-setup}. Soit $X$ localement
de type fini sur $S$. Soient $(\mathcal{L}, s, i : D \to X)$ et
$(\mathcal{L}', s', i' : D' \to X)$ deux triplets comme dans la
définition \ref{chow-definition-gysin-homomorphism}. Alors le diagramme
$$
\xymatrix{
\CH_k(X) \ar[r]_{i^*} \ar[d]_{(i')^*} & \CH_{k - 1}(D) \ar[d]^{j^*} \\
\CH_{k - 1}(D') \ar[r]^{(j')^*} & \CH_{k - 2}(D \cap D')
}
$$
est commutatif, chacune des applications étant une application de Gysin.
\end{lemma}

\begin{proof}
Notons $j : D \cap D' \to D$ et $j' : D \cap D' \to D'$ les immersions
fermées correspondant à $(\mathcal{L}|_{D'}, s|_{D'}$ et
$(\mathcal{L}'_D, s|_D)$. Il faut montrer que
$(j')^*i^*\alpha = j^* (i')^*\alpha$ pour tout $\alpha \in \CH_k(X)$.
Soit $W \subset X$ un sous-schéma fermé intègre de dimension $k$.
Démontrons l'égalité lorsque $\alpha = [W]$. Nous la déduirons
de la formule fondamentale.

\medskip\noindent
Soit $\sigma$ une section méromorphe non nulle de $\mathcal{L}|_W$,
que l'on impose égale à $s|_W$ si $W \not \subset D$.
Soit $\sigma'$ une section méromorphe non nulle de $\mathcal{L}'|_W$,
que l'on impose égale à $s'|_W$ si $W \not \subset D'$.
Écrivons
$$
\text{div}_{\mathcal{L}|_W}(\sigma) =
\sum \text{ord}_{Z_i, \mathcal{L}|_W}(\sigma)[Z_i] = \sum n_i[Z_i]
$$
et, de même,
$$
\text{div}_{\mathcal{L}'|_W}(\sigma') =
\sum \text{ord}_{Z_i, \mathcal{L}'|_W}(\sigma')[Z_i] = \sum n'_i[Z_i]
$$
comme dans la discussion de la section \ref{chow-section-key}.
On voit alors que $Z_i \subset D$ si $n_i \not = 0$, et
$Z'_i \subset D'$ si $n'_i \not = 0$. Pour chaque $i$, soit $\xi_i \in Z_i$
le point générique. Comme dans la section \ref{chow-section-key}, choisissons
pour chaque $i$ un élément
$\sigma_i \in \mathcal{L}_{\xi_i}$, resp.\ $\sigma'_i \in \mathcal{L}'_{\xi_i}$,
qui soit un générateur sur $B_i = \mathcal{O}_{W, \xi_i}$
et qui soit égal à l'image de
$s$, resp.\ $s'$ si $Z_i \not \subset D$, resp.\ $Z_i \not \subset D'$.
Écrivons $\sigma = f_i \sigma_i$ et $\sigma' = f'_i\sigma'_i$, de sorte que
$n_i = \text{ord}_{B_i}(f_i)$ et
$n'_i = \text{ord}_{B_i}(f'_i)$.
Il résulte de nos définitions que
$$
(j')^*i^*[W] =
\sum \text{ord}_{B_i}(f_i) \text{div}_{\mathcal{L}'|_{Z_i}}(\sigma'_i|_{Z_i})
$$
comme cycles, et que
$$
j^*(i')^*[W] =
\sum \text{ord}_{B_i}(f'_i) \text{div}_{\mathcal{L}|_{Z_i}}(\sigma_i|_{Z_i})
$$
La formule fondamentale (lemme \ref{chow-lemma-key-formula}) donne maintenant l'égalité
$$
\sum \left(
\text{ord}_{B_i}(f_i) \text{div}_{\mathcal{L}'|_{Z_i}}(\sigma'_i|_{Z_i}) -
\text{ord}_{B_i}(f'_i) \text{div}_{\mathcal{L}|_{Z_i}}(\sigma_i|_{Z_i})
\right) =
\sum \text{div}_{Z_i}(\partial_{B_i}(f_i, f'_i))
$$
de cycles. Remarquons que $\text{div}_{Z_i}(\partial_{B_i}(f_i, f'_i)) = 0$ si
$Z_i \not \subset D \cap D'$, car, dans ce cas, soit $f_i = 1$,
soit $f'_i = 1$. On obtient ainsi une équivalence rationnelle entre nos cycles
particuliers représentant $(j')^*i^*[W]$ et $j^*(i')^*[W]$ sur $D \cap D' \cap W$.
D'après la remarque \ref{chow-remark-infinite-sums-rational-equivalences},
le résultat s'ensuit pour un $\alpha$ quelconque.
\end{proof}
















\section{Diviseurs de Cartier effectifs relatifs}
\label{chow-section-relative-effective-cartier}

\noindent
Les diviseurs de Cartier effectifs relatifs sont définis et étudiés
dans Diviseurs, section \ref{divisors-section-effective-Cartier-morphisms}.
Pour établir les résultats fondamentaux sur les classes de Chern des fibrés vectoriels,
nous n'avons besoin que du cas où le schéma ambiant et le diviseur de Cartier
effectif sont tous deux plats sur la base.

\begin{lemma}
\label{chow-lemma-relative-effective-cartier}
Soit $(S, \delta)$ comme dans la situation \ref{chow-situation-setup}.
Soient $X$, $Y$ localement de type fini sur $S$.
Soit $p : X \to Y$ un morphisme plat de dimension relative $r$.
Soit $i : D \to X$ un diviseur de Cartier effectif relatif
(Diviseurs, définition
\ref{divisors-definition-relative-effective-Cartier-divisor}).
Soit $\mathcal{L} = \mathcal{O}_X(D)$.
Pour tout $\alpha \in \CH_{k + 1}(Y)$, on a
$$
i^*p^*\alpha = (p|_D)^*\alpha
$$
dans $\CH_{k + r}(D)$ et
$$
c_1(\mathcal{L}) \cap p^*\alpha = i_* ((p|_D)^*\alpha)
$$
dans $\CH_{k + r}(X)$.
\end{lemma}

\begin{proof}
Soit $W \subset Y$ un sous-schéma fermé intègre de $\delta$-dimension
$k + 1$. D'après Diviseurs, lemme \ref{divisors-lemma-relative-Cartier},
on voit que $D \cap p^{-1}W$ est un diviseur de Cartier
effectif sur $p^{-1}W$. Le lemme \ref{chow-lemma-easy-gysin}
donne la première égalité dans
$$
i^*[p^{-1}W]_{k + r + 1} =
[D \cap p^{-1}W]_{k + r} =
[(p|_D)^{-1}(W)]_{k + r}.
$$
et la seconde découle de l'égalité $D \cap p^{-1}(W) = (p|_D)^{-1}(W)$ de schémas.
Puisque, par définition, $p^*[W] = [p^{-1}W]_{k + r + 1}$, on voit que
$i^*p^*[W] = (p|_D)^*[W]$ comme cycles. Si $\alpha = \sum m_j[W_j]$ est un
$k + 1$-cycle quelconque, on obtient
$i^*\alpha = \sum m_j i^*p^*[W_j] = \sum m_j(p|_D)^*[W_j]$ comme cycles.
Cela démontre la première égalité. Pour en déduire la seconde,
on applique le lemme \ref{chow-lemma-support-cap-effective-Cartier}.
\end{proof}








\section{Fibrés affines}
\label{chow-section-affine-vector}

\noindent
Pour un fibré affine, l'application d'image réciproque est surjective sur les groupes de Chow.

\begin{lemma}
\label{chow-lemma-pullback-affine-fibres-surjective}
Soit $(S, \delta)$ comme dans la situation \ref{chow-situation-setup}.
Soient $X$, $Y$ localement de type fini sur $S$.
Soit $f : X \to Y$ un morphisme plat de dimension relative $r$.
Supposons que, pour tout $y \in Y$, il existe un voisinage ouvert
$U \subset Y$ tel que $f|_{f^{-1}(U)} : f^{-1}(U) \to U$
s'identifie au morphisme $U \times \mathbf{A}^r \to U$.
Alors $f^* : \CH_k(Y) \to \CH_{k + r}(X)$ est surjectif pour tout
$k \in \mathbf{Z}$.
\end{lemma}

\begin{proof}
Soit $\alpha \in \CH_{k + r}(X)$.
Écrivons $\alpha = \sum m_j[W_j]$, avec $m_j \not = 0$ et
des sous-schémas fermés intègres $W_j$ deux à deux distincts de
$\delta$-dimension $k + r$. Alors la famille $\{W_j\}$
est localement finie dans $X$. Pour tout ouvert quasi-compact
$V \subset Y$, on voit que $f^{-1}(V) \cap W_j$
n'est non vide que pour un nombre fini de $j$. La
famille $Z_j = \overline{f(W_j)}$ des adhérences
des images est donc une famille localement finie de sous-schémas
fermés intègres de $Y$.

\medskip\noindent
Considérons les diagrammes de produit fibré
$$
\xymatrix{
f^{-1}(Z_j) \ar[r] \ar[d]_{f_j} & X \ar[d]^f \\
Z_j \ar[r] & Y
}
$$
Supposons que $[W_j] \in Z_{k + r}(f^{-1}(Z_j))$
soit rationnellement équivalent à $f_j^*\beta_j$ pour un
$k$-cycle $\beta_j \in \CH_k(Z_j)$. Alors
$\beta = \sum m_j \beta_j$ sera un $k$-cycle sur $Y$,
et $f^*\beta = \sum m_j f_j^*\beta_j$ sera rationnellement
équivalent à $\alpha$ (voir la
remarque \ref{chow-remark-infinite-sums-rational-equivalences}).
Cela nous ramène au cas où $Y$ est intègre et
$\alpha = [W]$ pour un sous-schéma fermé intègre
de $X$ dominant $Y$. En particulier, on peut
supposer que $d = \dim_\delta(Y) < \infty$.

\medskip\noindent
On peut donc raisonner par récurrence sur $d = \dim_\delta(Y)$.
Si $d < k$, alors $\CH_{k + r}(X) = 0$ et le lemme est vrai.
Par hypothèse, il existe un ouvert dense $V \subset Y$ tel
que $f^{-1}(V) \cong V \times \mathbf{A}^r$ comme schémas sur $V$.
Supposons que l'on puisse montrer que $\alpha|_{f^{-1}(V)} = f^*\beta$
pour un certain $\beta \in Z_k(V)$. D'après le lemme \ref{chow-lemma-exact-sequence-open},
on voit que
$\beta = \beta'|_V$ pour un certain $\beta' \in Z_k(Y)$.
D'après la suite exacte
$\CH_k(f^{-1}(Y \setminus V)) \to \CH_k(X) \to \CH_k(f^{-1}(V))$
du lemme \ref{chow-lemma-restrict-to-open},
on voit que $\alpha - f^*\beta'$ provient
d'un cycle $\alpha' \in \CH_{k + r}(f^{-1}(Y \setminus V))$.
Comme $\dim_\delta(Y \setminus V) < d$, la récurrence sur
$d$ donne le résultat.

\medskip\noindent
On peut donc supposer que $X = Y \times \mathbf{A}^r$.
Dans ce cas, on peut factoriser $f$ sous la forme
$$
X = Y \times \mathbf{A}^r \to
Y \times \mathbf{A}^{r - 1} \to \ldots \to
Y \times \mathbf{A}^1 \to Y.
$$
Il suffit donc de traiter le cas $r = 1$. L'argument du
deuxième paragraphe de la démonstration nous ramène au cas où
$\alpha = [W]$, $Y$ est intègre et $W \to Y$ dominant.
On peut à nouveau raisonner par récurrence sur $d = \dim_\delta(Y)$.
Si $W = Y \times \mathbf{A}^1$, alors $[W] = f^*[Y]$.
Enfin, si $W \subset Y \times \mathbf{A}^1$ est une inclusion stricte,
alors $W \to Y$ induit une extension finie de corps $R(W)/R(Y)$.
Soit $P(T) \in R(Y)[T]$ le polynôme unitaire irréductible tel
que la fibre générique de $W \to Y$ soit définie par $P$ dans
$\mathbf{A}^1_{R(Y)}$. Soit $V \subset Y$ un ouvert non vide tel
que $P \in \Gamma(V, \mathcal{O}_Y)[T]$ et tel que
$W \cap f^{-1}(V)$ soit encore défini par $P$. On voit alors que
$\alpha|_{f^{-1}(V)} \sim_{rat} 0$, et donc que $\alpha \sim_{rat} \alpha'$
pour un cycle $\alpha'$ sur $(Y \setminus V) \times \mathbf{A}^1$.
La récurrence sur la dimension donne le résultat.
\end{proof}

\begin{lemma}
\label{chow-lemma-linebundle}
Soit $(S, \delta)$ comme dans la situation \ref{chow-situation-setup}.
Soit $X$ localement de type fini sur $S$.
Soit $\mathcal{L}$ un $\mathcal{O}_X$-module inversible.
Soit
$$
p :
L = \underline{\Spec}(\text{Sym}^*(\mathcal{L}))
\longrightarrow
X
$$
le fibré vectoriel associé sur $X$.
Alors $p^* : \CH_k(X) \to \CH_{k + 1}(L)$ est un isomorphisme pour tout $k$.
\end{lemma}

\begin{proof}
Pour la surjectivité, voir le lemme \ref{chow-lemma-pullback-affine-fibres-surjective}.
Soit $o : X \to L$ la section nulle de $L \to X$, c'est-à-dire le morphisme
correspondant à la surjection $\text{Sym}^*(\mathcal{L}) \to \mathcal{O}_X$
qui envoie $\mathcal{L}^{\otimes n}$ sur zéro pour tout $n > 0$.
Alors $p \circ o = \text{id}_X$, et $o(X)$ est un diviseur de Cartier
effectif sur $L$. Le lemme \ref{chow-lemma-relative-effective-cartier}
donne donc $o^* \circ p^* = \text{id}$, et l'on en conclut que $p^*$ est
également injectif.
\end{proof}

\begin{remark}
\label{chow-remark-when-isomorphism}
Nous verrons plus loin (lemme \ref{chow-lemma-vectorbundle}) que, si $X$ est un
fibré vectoriel de rang $r$ sur $Y$, alors l'application d'image réciproque
$\CH_k(Y) \to \CH_{k + r}(X)$
est un isomorphisme. Cela est vrai dès que $X \to Y$ satisfait
aux hypothèses du lemme \ref{chow-lemma-pullback-affine-fibres-surjective} ; voir
\cite[Lemme 2.2]{Totaro-group}. Nous esquisserons une démonstration dans la
remarque \ref{chow-remark-higher-chow-isomorphism} à l'aide des groupes de Chow supérieurs.
\end{remark}

\begin{lemma}
\label{chow-lemma-linebundle-formulae}
Dans la situation du lemme \ref{chow-lemma-linebundle}, notons $o : X \to L$
la section nulle (voir la démonstration du lemme). On a alors
\begin{enumerate}
\item $o(X)$ est le schéma des zéros d'une section globale régulière de
$p^*\mathcal{L}^{\otimes -1}$,
\item $o_* : \CH_k(X) \to \CH_k(L)$, puisque $o$ est une immersion fermée,
\item $o^* : \CH_{k + 1}(L) \to \CH_k(X)$, puisque $o(X)$
est un diviseur de Cartier effectif,
\item $o^* p^* : \CH_k(X) \to \CH_k(X)$ est l'application identique,
\item $o_*\alpha = - p^*(c_1(\mathcal{L}) \cap \alpha)$ pour tout
$\alpha \in \CH_k(X)$, et
\item $o^* o_* : \CH_k(X) \to \CH_{k - 1}(X)$ est égal à l'application
$\alpha \mapsto - c_1(\mathcal{L}) \cap \alpha$.
\end{enumerate}
\end{lemma}

\begin{proof}
Puisque $p_*\mathcal{O}_L = \text{Sym}^*(\mathcal{L})$, on a
$p_*(p^*\mathcal{L}^{\otimes -1}) =
\text{Sym}^*(\mathcal{L}) \otimes_{\mathcal{O}_X} \mathcal{L}^{\otimes -1}$
d'après la formule de projection
(Cohomologie, lemme \ref{cohomology-lemma-projection-formula}),
et la section mentionnée dans (1) est
la trivialisation canonique $\mathcal{O}_X \to
\mathcal{L} \otimes_{\mathcal{O}_X} \mathcal{L}^{\otimes -1}$.
Nous omettons la démonstration du fait que le lieu des zéros de
cette section est précisément $o(X)$. Cela démontre (1).

\medskip\noindent
Nous avons établi les assertions (2), (3) et (4) au cours de la démonstration du
lemme \ref{chow-lemma-linebundle}. Bien entendu, (4) est la première
formule du lemme \ref{chow-lemma-relative-effective-cartier}.

\medskip\noindent
L'assertion (5) découle de la seconde formule du
lemme \ref{chow-lemma-relative-effective-cartier},
de l'additivité du produit cap avec $c_1$ (lemme \ref{chow-lemma-c1-cap-additive})
et du fait que ce produit cap avec $c_1$ commute à l'image réciproque plate
(lemme \ref{chow-lemma-flat-pullback-cap-c1}).

\medskip\noindent
L'assertion (6) découle du lemme \ref{chow-lemma-gysin-back}
et du fait que $o^*p^*\mathcal{L} = \mathcal{L}$.
\end{proof}

\begin{lemma}
\label{chow-lemma-decompose-section}
Soit $Y$ un schéma. Soient $\mathcal{L}_i$, $i = 1, 2$ des
$\mathcal{O}_Y$-modules inversibles. Soit $s$ une section globale de
$\mathcal{L}_1 \otimes_{\mathcal{O}_Y} \mathcal{L}_2$.
Notons $i : D \to Y$ le schéma des zéros de $s$.
Il existe alors un diagramme commutatif
$$
\xymatrix{
D_1 \ar[r]_{i_1} \ar[d]_{p_1} &
L \ar[d]^p &
D_2 \ar[l]^{i_2} \ar[d]^{p_2} \\
D \ar[r]^i &
Y &
D \ar[l]_i
}
$$
et des sections $s_i$ de $p^*\mathcal{L}_i$ telles que
les assertions suivantes soient vérifiées :
\begin{enumerate}
\item $p^*s = s_1 \otimes s_2$,
\item $p$ est de type fini et plat de dimension relative $1$,
\item $D_i$ est le schéma des zéros de $s_i$,
\item $D_i \cong
\underline{\Spec}(\text{Sym}^*(\mathcal{L}_{3 - i}^{\otimes -1})|_D))$
sur $D$ pour $i = 1, 2$,
\item $p^{-1}D = D_1 \cup D_2$ (réunion schématique),
\item $D_1 \cap D_2$ (intersection schématique) s'envoie
isomorphiquement sur $D$, et
\item $D_1 \cap D_2 \to D_i$
est la section nulle du fibré en droites $D_i \to D$ pour $i = 1, 2$.
\end{enumerate}
De plus, la formation de ce diagramme et des sections $s_i$
commute à tout changement de base.
\end{lemma}

\begin{proof}
Soit $p : X \to Y$ le spectre relatif du faisceau quasi-cohérent
de $\mathcal{O}_Y$-algèbres
$$
\mathcal{A} =
\left(\bigoplus\nolimits_{a_1, a_2 \geq 0}
\mathcal{L}_1^{\otimes -a_1} \otimes_{\mathcal{O}_Y}
\mathcal{L}_2^{\otimes -a_2}\right)/\mathcal{J}
$$
où $\mathcal{J}$ est l'idéal engendré par les sections locales de
la forme $st - t$, pour $t$ section locale de l'un des facteurs directs
$\mathcal{L}_1^{\otimes -a_1} \otimes \mathcal{L}_2^{\otimes -a_2}$
avec $a_1, a_2 > 0$. Les sections $s_i$, considérées comme applications
$p^*\mathcal{L}_i^{\otimes -1} \to \mathcal{O}_X$, sont définies comme les adjointes
des applications $\mathcal{L}_i^{\otimes -1} \to \mathcal{A} = p_*\mathcal{O}_X$.
Pour tout $y \in Y$, on peut choisir un
ouvert affine $V \subset Y$, disons $V = \Spec(A)$, contenant $y$, et des
trivialisations $z_i : \mathcal{O}_V \to \mathcal{L}_i^{\otimes -1}|_V$.
Remarquons que $f = s(z_1z_2) \in A$ définit le sous-schéma fermé $D$.
Alors, clairement,
$$
p^{-1}(V) = \Spec(B[z_1, z_2]/(z_1 z_2 - f))
$$
Puisque $D_i$ est défini par $z_i$, tout est clair.
\end{proof}

\begin{lemma}
\label{chow-lemma-decompose-section-formulae}
Dans la situation du lemme \ref{chow-lemma-decompose-section},
supposons $Y$ localement de type fini sur $(S, \delta)$ comme dans la
situation \ref{chow-situation-setup}. Alors on a
$i_1^*p^*\alpha = p_1^*i^*\alpha$
dans $\CH_k(D_1)$ pour tout $\alpha \in \CH_k(Y)$.
\end{lemma}

\begin{proof}
Soit $W \subset Y$ un sous-schéma fermé intègre de $\delta$-dimension $k$.
Distinguons deux cas.

\medskip\noindent
Supposons $W \subset D$. Alors
$i^*[W] = c_1(\mathcal{L}_1) \cap [W] + c_1(\mathcal{L}_2) \cap [W]$
dans $\CH_{k - 1}(D)$, d'après notre définition des homomorphismes de Gysin et
l'additivité du lemme \ref{chow-lemma-c1-cap-additive}.
Ainsi $p_1^*i^*[W] =
p_1^*(c_1(\mathcal{L}_1) \cap [W]) + p_1^*(c_1(\mathcal{L}_2) \cap [W])$.
D'autre part, on a
$p^*[W] = [p^{-1}(W)]_{k + 1}$ par construction de l'image réciproque plate.
De plus, $p^{-1}(W) = W_1 \cup W_2$ (au sens schématique),
où $W_i = p_i^{-1}(W)$ est un fibré en droites sur $W$
d'après le lemme (puisque la formation du diagramme commute au changement de base).
Alors $[p^{-1}(W)]_{k + 1} = [W_1] + [W_2]$, car les $W_i$ sont des sous-schémas fermés
intègres de $L$ de $\delta$-dimension $k + 1$. Ainsi
\begin{align*}
i_1^*p^*[W]
& =
i_1^*[p^{-1}(W)]_{k + 1} \\
& =
i_1^*([W_1] + [W_2]) \\
& =
c_1(p_1^*\mathcal{L}_1) \cap [W_1] + [W_1 \cap W_2]_k \\
& =
c_1(p_1^*\mathcal{L}_1) \cap p_1^*[W] + [W_1 \cap W_2]_k \\
& =
p_1^*(c_1(\mathcal{L}_1) \cap [W]) + [W_1 \cap W_2]_k
\end{align*}
par construction des homomorphismes de Gysin, par définition de l'image réciproque plate
(pour la deuxième égalité) et par compatibilité de $c_1 \cap -$
avec l'image réciproque plate (lemme \ref{chow-lemma-flat-pullback-cap-c1}).
Puisque $W_1 \cap W_2$ est la section nulle du fibré en droites
$W_1 \to W$, le lemme \ref{chow-lemma-linebundle-formulae} donne
$[W_1 \cap W_2]_k = p_1^*(c_1(\mathcal{L}_2) \cap [W])$.
On utilise ici le fait que $D_1$ est le fibré en droites
qui est le spectre relatif de l'inverse de $\mathcal{L}_2$.
On obtient donc la même expression que précédemment.

\medskip\noindent
Supposons $W \not \subset D$. Dans ce cas, $i_1^*p^*[W]$
et $p_1^*i^*[W]$ sont tous deux représentés par le $k - 1$-cycle associé
à l'image réciproque schématique de $W$ dans $D_1$.
\end{proof}

\begin{lemma}
\label{chow-lemma-normal-cone-effective-Cartier}
Dans la situation \ref{chow-situation-setup}, soit $X$ un schéma localement
de type fini sur $S$. Soit $(\mathcal{L}, s, i : D \to X)$
un triplet comme dans la définition \ref{chow-definition-gysin-homomorphism}.
Il existe un diagramme commutatif
$$
\xymatrix{
D' \ar[r]_{i'} \ar[d]_p & X' \ar[d]^g \\
D \ar[r]^i & X
}
$$
tel que
\begin{enumerate}
\item $p$ et $g$ soient de type fini et plats de dimension relative $1$,
\item $p^* : \CH_k(D) \to \CH_{k + 1}(D')$ soit injectif pour tout $k$,
\item $D' \subset X'$ soit le schéma des zéros d'une section globale
$s' \in \Gamma(X', \mathcal{O}_{X'})$,
\item $p^*i^* = (i')^*g^*$ comme applications $\CH_k(X) \to \CH_k(D')$.
\end{enumerate}
De plus, ces propriétés restent vraies après tout changement de base
par un morphisme $Y \to X$ localement de type fini.
\end{lemma}

\begin{proof}
Remarquons que $(i')^*$ est défini, car on dispose du triplet
$(\mathcal{O}_{X'}, s', i' : D' \to X')$ comme dans la
définition \ref{chow-definition-gysin-homomorphism}. L'énoncé a donc un sens.

\medskip\noindent
Posons $\mathcal{L}_1 = \mathcal{O}_X$, $\mathcal{L}_2 = \mathcal{L}$,
et appliquons le lemme \ref{chow-lemma-decompose-section} à la section $s$ de
$\mathcal{L} = \mathcal{L}_1 \otimes_{\mathcal{O}_X} \mathcal{L}_2$.
Prenons $D' = D_1$. Les résultats découlent alors de ce lemme, du
lemme \ref{chow-lemma-decompose-section-formulae},
et l'injectivité du
lemme \ref{chow-lemma-linebundle}.
\end{proof}

\begin{remark}
\label{chow-remark-higher-chow-isomorphism}
Soit $(S, \delta)$ comme dans la situation \ref{chow-situation-setup}.
Soit $Y$ localement de type fini sur $S$. Soit $r \geq 0$. Soit
$f : X \to Y$ un morphisme de schémas. Supposons que tout $y \in Y$
soit contenu dans un ouvert $V \subset Y$ tel que
$f^{-1}(V) \cong V \times \mathbf{A}^r$ comme schémas sur $V$.
Dans cette remarque, nous esquissons une démonstration du fait que
$f^* : \CH_k(Y) \to \CH_{k + r}(X)$ est un isomorphisme.
D'abord, d'après le lemme \ref{chow-lemma-pullback-affine-fibres-surjective},
l'application est surjective. Soit $\alpha \in \CH_k(Y)$ tel que $f^*\alpha = 0$.
Nous allons démontrer que $\alpha = 0$.

\medskip\noindent
Étape 1. On peut supposer que $\dim_\delta(Y) < \infty$. (En pratique,
cela est immédiat dans tous les cas, et nous suggérons donc de sauter cette étape.)
En effet, toute équivalence rationnelle témoignant de l'égalité $f^*\alpha = 0$ sur $X$
fait intervenir une famille localement finie de sous-schémas fermés intègres
de dimension $k + r + 1$. En prenant la réunion des adhérences de leurs images
dans $Y$, on obtient un sous-schéma fermé $Y' \subset Y$
tel que $\dim_\delta(Y') \leq k + r + 1$, que
$\alpha$ soit l'image d'un certain $\alpha' \in \CH_k(Y')$,
et que $(f')^*\alpha = 0$, où $f'$ est le changement de base de $f$ à $Y'$.

\medskip\noindent
Étape 2. Supposons $d = \dim_\delta(Y) < \infty$. On peut alors raisonner par récurrence sur
$d$. Si $d < k$, alors $\alpha = 0$ et le résultat est acquis ; c'est le cas initial
de la récurrence.
En général, notre hypothèse sur $f$ montre que l'on peut choisir un ouvert dense
$V \subset Y$ tel que $U = f^{-1}(V) = \mathbf{A}^r_V$.
Notons $Y' \subset Y$ le complémentaire de $V$, muni de sa structure de sous-schéma fermé
réduit, et posons $X' = f^{-1}(Y')$. Considérons
$$
\xymatrix{
\CH^M_{k + r}(U, 1) \ar[r] &
\CH_{k + r}(X') \ar[r] &
\CH_{k + r}(X) \ar[r] &
\CH_{k + r}(U) \ar[r] &
0 \\
\CH^M_k(V, 1) \ar[r] \ar[u] &
\CH_k(Y') \ar[r] \ar[u] &
\CH_k(Y) \ar[r] \ar[u] &
\CH_k(V) \ar[r] \ar[u] &
0
}
$$
Nous utilisons ici les premiers groupes de Chow supérieurs de $V$ et de $U$ et les
suites exactes à six termes construites dans la remarque \ref{chow-remark-higher-chow},
ainsi que l'image réciproque plate pour ces groupes de Chow supérieurs et
sa compatibilité avec ces suites exactes à six termes.
Puisque $U = \mathbf{A}^r_V$, l'application verticale de droite est un isomorphisme.
L'application $\CH_k(Y') \to \CH_{k + r}(X')$ est bijective par récurrence sur $d$.
Pour achever le raisonnement, il suffit donc de montrer que
$$
\CH^M_k(V, 1) \longrightarrow \CH^M_{k + r}(U, 1)
$$
est surjective. En raisonnant comme dans la démonstration du
lemme \ref{chow-lemma-pullback-affine-fibres-surjective},
on se ramène à l'étape 3 ci-dessous.

\medskip\noindent
Étape 3. Soit $F$ un corps. Alors $\CH^M_0(\mathbf{A}^1_F, 1) = 0$.
(Dans la démonstration du lemme cité ci-dessus, nous avons montré de manière analogue que
$\CH_0(\mathbf{A}^1_F) = 0$.) On a
$$
\CH^M_0(\mathbf{A}^1_F, 1) = \Coker\left(
\partial : K^M_2(F(T)) \longrightarrow
\bigoplus\nolimits_{\mathfrak p \subset F[T]\text{ maximal}}
\kappa(\mathfrak p)^*\right)
$$
L'argument classique pour établir la nullité du conoyau consiste à
montrer, par récurrence sur le degré de $\kappa(\mathfrak p)/F$,
que le facteur direct correspondant à $\mathfrak p$ est contenu dans l'image.
Si $\mathfrak p$ est engendré par le
polynôme unitaire irréductible $P(T) \in F[T]$ et si
$u \in \kappa(x)^*$ est la classe résiduelle d'un
$Q(T) \in F[T]$ tel que $\deg(Q) < \deg(P)$, on montre que
$\partial(Q, P)$ donne l'élément $u$ en $\mathfrak p$,
et éventuellement d'autres unités en des idéaux premiers divisant $Q$, qui
ont un degré inférieur. Cela achève l'esquisse de démonstration.
\end{remark}











\section{Théorie bivariante de l'intersection}
\label{chow-section-bivariant}

\noindent
Afin de traiter convenablement les classes de Chern supérieures des fibrés
vectoriels, nous introduisons les classes de Chow bivariantes comme dans \cite{F}.
Notre définition diffère de celle de \cite{F} sur deux points :
(1) nous travaillons dans un cadre différent, et (2) nous exigeons seulement
que nos classes bivariantes commutent aux homomorphismes de Gysin pour
les schémas des zéros de sections de modules inversibles
(section \ref{chow-section-intersecting-effective-Cartier}).
Nous verrons plus loin, dans le lemme \ref{chow-lemma-gysin-commutes}, que nos
classes bivariantes commutent à tous les homomorphismes de Gysin en codimension supérieure,
et satisfont donc à toutes les propriétés exigées dans \cite{F} ; voir
aussi \cite[Théorème 17.1]{F}.

\begin{definition}
\label{chow-definition-bivariant-class}
\begin{reference}
Analogue à \cite[Définition 17.1]{F}
\end{reference}
Soit $(S, \delta)$ comme dans la situation \ref{chow-situation-setup}.
Soit $f : X \to Y$ un morphisme de schémas localement de type fini sur $S$.
Soit $p \in \mathbf{Z}$.
Une {\it classe bivariante $c$ de degré $p$ pour $f$} est donnée par une règle
qui associe à tout morphisme localement de type fini $Y' \to Y$
et à tout $k$ une application
$$
c \cap - : \CH_k(Y') \longrightarrow \CH_{k - p}(X')
$$
où $X' = Y' \times_Y X$, satisfaisant aux conditions suivantes :
\begin{enumerate}
\item si $Y'' \to Y'$ est propre, alors
$c \cap (Y'' \to Y')_*\alpha'' = (X'' \to X')_*(c \cap \alpha'')$
pour tout $\alpha''$ sur $Y''$, où $X'' = Y'' \times_Y X$,
\item si $Y'' \to Y'$ est plat, localement de type fini et de
dimension relative fixée, alors
$c \cap (Y'' \to Y')^*\alpha' = (X'' \to X')^*(c \cap \alpha')$
pour tout $\alpha'$ sur $Y'$, et
\item si $(\mathcal{L}', s', i' : D' \to Y')$ est comme dans la
définition \ref{chow-definition-gysin-homomorphism},
avec pour image réciproque $(\mathcal{N}', t', j' : E' \to X')$ sur $X'$,
alors on a $c \cap (i')^*\alpha' = (j')^*(c \cap \alpha')$
pour tout $\alpha'$ sur $Y'$.
\end{enumerate}
L'ensemble des classes bivariantes de degré $p$ pour $f$ est
noté $A^p(X \to Y)$.
\end{definition}

\noindent
Soit $(S, \delta)$ comme dans la situation \ref{chow-situation-setup}. Soient $X \to Y$
et $Y \to Z$
des morphismes de schémas localement de type fini sur $S$. Soit
$p \in \mathbf{Z}$. Il est clair que $A^p(X \to Y)$ est un groupe abélien.
De plus, on dispose clairement d'une composition bilinéaire
$$
A^p(X \to Y) \times A^q(Y \to Z) \to A^{p + q}(X \to Z)
$$
qui est associative.

\begin{lemma}
\label{chow-lemma-flat-pullback-bivariant}
Soit $(S, \delta)$ comme dans la situation \ref{chow-situation-setup}.
Soit $f : X \to Y$ un morphisme plat de dimension relative $r$
entre schémas localement de type fini sur $S$.
Alors la règle qui associe à $Y' \to Y$
$(f')^* : \CH_k(Y') \to \CH_{k + r}(X')$, où $X' = X \times_Y Y'$,
est une classe bivariante de degré $-r$.
\end{lemma}

\begin{proof}
Cela découle des
lemmes \ref{chow-lemma-flat-pullback-rational-equivalence},
\ref{chow-lemma-compose-flat-pullback},
\ref{chow-lemma-flat-pullback-proper-pushforward} et
\ref{chow-lemma-gysin-flat-pullback}.
\end{proof}

\begin{lemma}
\label{chow-lemma-gysin-bivariant}
Soit $(S, \delta)$ comme dans la situation \ref{chow-situation-setup}.
Soit $X$ localement de type fini sur $S$.
Soit $(\mathcal{L}, s, i : D \to X)$ un triplet comme dans la
définition \ref{chow-definition-gysin-homomorphism}.
Alors la règle qui associe à $f : X' \to X$
$(i')^* : \CH_k(X') \to \CH_{k - 1}(D')$, où $D' = D \times_X X'$,
est une classe bivariante de degré $1$.
\end{lemma}

\begin{proof}
Cela découle des lemmes \ref{chow-lemma-gysin-factors},
\ref{chow-lemma-closed-in-X-gysin},
\ref{chow-lemma-gysin-flat-pullback} et
\ref{chow-lemma-gysin-commutes-gysin}.
\end{proof}

\begin{lemma}
\label{chow-lemma-push-proper-bivariant}
Soit $(S, \delta)$ comme dans la situation \ref{chow-situation-setup}.
Soient $f : X \to Y$ et $g : Y \to Z$ des morphismes de
schémas localement de type fini sur $S$.
Soit $c \in A^p(X \to Z)$, et supposons $f$ propre.
Alors la règle qui associe à $Z' \to Z$
$\alpha \longmapsto f'_*(c \cap \alpha)$
est une classe bivariante notée $f_* \circ c \in A^p(Y \to Z)$.
\end{lemma}

\begin{proof}
Cela découle des lemmes \ref{chow-lemma-compose-pushforward},
\ref{chow-lemma-flat-pullback-proper-pushforward} et
\ref{chow-lemma-closed-in-X-gysin}.
\end{proof}

\begin{remark}
\label{chow-remark-restriction-bivariant}
Soit $(S, \delta)$ comme dans la situation \ref{chow-situation-setup}. Soient $X \to Y$
et $Y' \to Y$ des morphismes de schémas localement de type fini sur $S$.
Soit $X' = Y' \times_Y X$. On dispose alors d'une application de restriction évidente
$$
A^p(X \to Y) \longrightarrow A^p(X' \to Y'),\quad
c \longmapsto res(c)
$$
obtenue en considérant un schéma $Y''$ localement de type fini sur $Y'$
comme un schéma localement de type fini sur $Y$ et en posant
$res(c) \cap \alpha'' = c \cap \alpha''$ pour tout $\alpha'' \in \CH_k(Y'')$.
Cette opération de restriction est compatible avec les compositions de manière
évidente.
\end{remark}

\begin{remark}
\label{chow-remark-bivariant-commute}
Soit $(S, \delta)$ comme dans la situation \ref{chow-situation-setup}. Soit $X$
localement de type fini sur $S$. Pour $i = 1, 2$, soit $Z_i \to X$
un morphisme de schémas localement de type fini. Soient
$c_i \in A^{p_i}(Z_i \to X)$, $i = 1, 2$ des classes bivariantes.
Pour tout $\alpha \in \CH_k(X)$, on peut se demander si
$$
c_1 \cap c_2 \cap \alpha = c_2 \cap c_1 \cap \alpha
$$
dans $\CH_{k - p_1 - p_2}(Z_1 \times_X Z_2)$. Si cette égalité est vraie et le demeure
après tout changement de base par un morphisme $X' \to X$ localement de type fini, on dit
que $c_1$ et $c_2$ {\it commutent}. Bien entendu, cela revient à dire que
$$
res(c_1) \circ c_2 = res(c_2) \circ c_1
$$
dans $A^{p_1 + p_2}(Z_1 \times_X Z_2 \to X)$. Ici,
$res(c_1) \in A^{p_1}(Z_1 \times_X Z_2 \to Z_2)$ est la restriction de $c_1$
au sens de la remarque \ref{chow-remark-restriction-bivariant} ; de même pour $res(c_2)$.
\end{remark}

\begin{example}
\label{chow-example-gysin-commute}
Soit $(S, \delta)$ comme dans la situation \ref{chow-situation-setup}. Soit $X$ localement
de type fini sur $S$. Soit $(\mathcal{L}, s, i : D \to X)$ un triplet comme dans la
définition \ref{chow-definition-gysin-homomorphism}. Soit $Z \to X$ un morphisme
de schémas localement de type fini, et soit $c \in A^p(Z \to X)$ une
classe bivariante. Alors la classe de Gysin bivariante $c' \in A^1(D \to X)$ du
lemme \ref{chow-lemma-gysin-bivariant} commute à $c$ au sens de la
remarque \ref{chow-remark-bivariant-commute}. Il s'agit en effet d'une reformulation de la
condition (3) de la définition \ref{chow-definition-bivariant-class}.
\end{example}

\begin{remark}
\label{chow-remark-more-general-bivariant}
Il existe un type plus général de classe bivariante qui ne semble pas être
considéré dans la littérature. Supposons en effet donné un diagramme
$$
X \longrightarrow Z \longleftarrow Y
$$
de schémas localement de type fini sur $(S, \delta)$ comme dans la
situation \ref{chow-situation-setup}. Soit $p \in \mathbf{Z}$.
On peut alors considérer une règle $c$ qui associe à tout $Z' \to Z$
localement de type fini des applications
$$
c \cap - : \CH_k(Y') \longrightarrow \CH_{k - p}(X')
$$
pour tout $k \in \mathbf{Z}$,
où $X' = X \times_Z Z'$ et $Y' = Z' \times_Z Y$, compatibles aux
\begin{enumerate}
\item images directes propres pour $Z'' \to Z'$ propre,
\item images réciproques plates pour $Z'' \to Z'$ plat
de dimension relative fixée, et
\item applications de Gysin pour $D' \subset Z'$ comme dans la
définition \ref{chow-definition-gysin-homomorphism}.
\end{enumerate}
Nous omettons les formulations détaillées. Notons l'ensemble
de toutes ces opérations $A^p(X \to Z \leftarrow Y)$. Un exemple simple
de l'utilité de cette notion est fourni par un morphisme propre
$f : X_2 \to X_1$. Alors $f_*$ n'est pas une opération bivariante au sens de la
définition \ref{chow-definition-bivariant-class}, mais l'est au sens
généralisé ci-dessus ; on a $f_* \in A^0(X_1 \to X_1 \leftarrow X_2)$.
\end{remark}





\section{Cohomologie de Chow et première classe de Chern}
\label{chow-section-chow-cohomology}

\noindent
Nous nous intéresserons surtout à $A^p(X) = A^p(X \to X)$, qui désignera toujours
les classes de cohomologie bivariante pour $\text{id}_X$. C'est en effet là que
se trouveront les classes de Chern.

\begin{definition}
\label{chow-definition-chow-cohomology}
Soit $(S, \delta)$ comme dans la situation \ref{chow-situation-setup}.
Soit $X$ localement de type fini sur $S$. La {\it cohomologie de Chow}
de $X$ est la $\mathbf{Z}$-algèbre graduée $A^*(X)$ dont la composante de degré
$p$ est $A^p(X \to X)$.
\end{definition}

\noindent
Attention : il n'est pas clair que la structure de $\mathbf{Z}$-algèbre
sur $A^*(X)$ soit commutative, mais nous verrons que les classes de Chern appartiennent
à son centre.

\begin{remark}
\label{chow-remark-pullback-cohomology}
Soit $(S, \delta)$ comme dans la situation \ref{chow-situation-setup}.
Soit $f : Y' \to Y$ un morphisme de schémas localement de type fini sur $S$.
Comme cas particulier de la remarque \ref{chow-remark-restriction-bivariant},
on dispose d'un homomorphisme canonique de $\mathbf{Z}$-algèbres $res : A^*(Y) \to A^*(Y')$.
Cet homomorphisme est souvent noté $f^*$ dans la littérature.
\end{remark}

\begin{lemma}
\label{chow-lemma-cap-c1-bivariant}
Soit $(S, \delta)$ comme dans la situation \ref{chow-situation-setup}.
Soit $X$ localement de type fini sur $S$.
Soit $\mathcal{L}$ un $\mathcal{O}_X$-module inversible.
Alors la règle qui associe à $f : X' \to X$
$c_1(f^*\mathcal{L}) \cap - : \CH_k(X') \to \CH_{k - 1}(X')$
est une classe bivariante de degré $1$.
\end{lemma}

\begin{proof}
Cela découle des lemmes \ref{chow-lemma-factors},
\ref{chow-lemma-pushforward-cap-c1},
\ref{chow-lemma-flat-pullback-cap-c1} et
\ref{chow-lemma-gysin-commutes-cap-c1}.
\end{proof}

\noindent
Le lemme précédent nous permet enfin de donner la définition suivante.

\begin{definition}
\label{chow-definition-first-chern-class}
Soit $(S, \delta)$ comme dans la situation \ref{chow-situation-setup}. Soit $X$
localement de type fini sur $S$. Soit $\mathcal{L}$ un
$\mathcal{O}_X$-module inversible. La {\it première classe de Chern}
$c_1(\mathcal{L}) \in A^1(X)$ de $\mathcal{L}$
est la classe bivariante du lemme \ref{chow-lemma-cap-c1-bivariant}.
\end{definition}

\noindent
Pour les modules localement libres de rang fini, nous construisons les classes de Chern dans la
section \ref{chow-section-intersecting-chern-classes}.
Montrons que $c_1(\mathcal{L})$ appartient au centre de $A^*(X)$.

\begin{lemma}
\label{chow-lemma-c1-center}
Soit $(S, \delta)$ comme dans la situation \ref{chow-situation-setup}.
Soit $X$ localement de type fini sur $S$.
Soit $\mathcal{L}$ un $\mathcal{O}_X$-module inversible.
Alors
\begin{enumerate}
\item $c_1(\mathcal{L}) \in A^1(X)$ appartient au centre de $A^*(X)$, et
\item si $f : X' \to X$ est localement de type fini et $c \in A^*(X' \to X)$,
alors $c \circ c_1(\mathcal{L}) = c_1(f^*\mathcal{L}) \circ c$.
\end{enumerate}
\end{lemma}

\begin{proof}
Bien entendu, (2) entraîne (1).
Soit $p : L \to X$ comme dans le lemme \ref{chow-lemma-linebundle}, et soit $o : X \to L$
la section nulle. Notons $p' : L' \to X'$ et $o' : X' \to L'$
leurs changements de base. D'après le lemme \ref{chow-lemma-linebundle-formulae}, on a
$$
p^*(c_1(\mathcal{L}) \cap \alpha) = - o_* \alpha
\quad\text{et}\quad
(p')^*(c_1(f^*\mathcal{L}) \cap \alpha') = - o'_* \alpha'
$$
Puisque $c$ est une classe bivariante, on a
\begin{align*}
(p')^*(c \cap c_1(\mathcal{L}) \cap \alpha)
& =
c \cap p^*(c_1(\mathcal{L}) \cap \alpha) \\
& =
- c \cap o_* \alpha \\
& =
- o'_*(c \cap \alpha) \\
& =
(p')^*(c_1(f^*\mathcal{L}) \cap c \cap \alpha)
\end{align*}
Puisque $(p')^*$ est injectif d'après l'un des lemmes cités ci-dessus, on obtient
$c \cap c_1(\mathcal{L}) \cap \alpha =
c_1(f^*\mathcal{L}) \cap c \cap \alpha$.
Il en est de même après tout changement de base par un morphisme $Y \to X$ localement de type fini,
d'où l'égalité de classes bivariantes énoncée dans (2).
\end{proof}

\begin{lemma}
\label{chow-lemma-vanish-above-dimension}
Soit $(S, \delta)$ comme dans la situation \ref{chow-situation-setup}. Soit $X$ un
schéma de type fini sur $S$ possédant un faisceau inversible ample.
Supposons $d = \dim(X) < \infty$ (il s'agit bien ici de la dimension et
non de la $\delta$-dimension).
Alors, pour tous faisceaux inversibles $\mathcal{L}_1, \ldots, \mathcal{L}_{d + 1}$
sur $X$, on a
$c_1(\mathcal{L}_1) \circ \ldots \circ c_1(\mathcal{L}_{d + 1}) = 0$
dans $A^{d + 1}(X)$.
\end{lemma}

\begin{proof}
Raisonnons par récurrence sur $d$. Le cas initial $d = 0$ est vrai, car
dans ce cas $X$ est un ensemble fini de points fermés, et tout module
inversible est donc trivial. Supposons $d > 0$. D'après Diviseurs, lemme
\ref{divisors-lemma-quasi-projective-Noetherian-pic-effective-Cartier},
on peut écrire $\mathcal{L}_{d + 1} \cong \mathcal{O}_X(D) \otimes
\mathcal{O}_X(D')^{\otimes -1}$ pour des diviseurs de Cartier effectifs
$D, D' \subset X$. Alors $c_1(\mathcal{L}_{d + 1})$ est la différence
de $c_1(\mathcal{O}_X(D))$ et de $c_1(\mathcal{O}_X(D'))$ ; on peut donc
supposer $\mathcal{L}_{d + 1} = \mathcal{O}_X(D)$ pour un
diviseur de Cartier effectif.

\medskip\noindent
Notons $i : D \to X$ le morphisme d'inclusion, et notons
$i^* \in A^1(D \to X)$ la classe bivariante donnée par
l'homomorphisme de Gysin comme dans le lemme \ref{chow-lemma-gysin-bivariant}.
On a $i_* \circ i^* = c_1(\mathcal{L}_{d + 1})$
dans $A^1(X)$ d'après le lemme \ref{chow-lemma-support-cap-effective-Cartier}
(et le lemme \ref{chow-lemma-push-proper-bivariant}
pour donner un sens au membre de gauche).
Puisque $c_1(\mathcal{L}_i)$ commute
à $i_*$ et à $i^*$ (par définition des classes bivariantes),
on en conclut que
$$
c_1(\mathcal{L}_1) \circ \ldots \circ c_1(\mathcal{L}_{d + 1}) =
i_* \circ c_1(\mathcal{L}_1) \circ \ldots \circ c_1(\mathcal{L}_d) \circ i^* =
i_* \circ c_1(\mathcal{L}_1|_D) \circ \ldots \circ c_1(\mathcal{L}_d|_D)
\circ i^*
$$
On conclut donc par récurrence sur $d$. En effet, on a $\dim(D) < d$,
car aucun des points génériques de $X$ n'appartient à $D$.
\end{proof}

\begin{remark}
\label{chow-remark-ring-loc-classes}
Soit $(S, \delta)$ comme dans la situation \ref{chow-situation-setup}.
Soit $Z \to X$ une immersion fermée de schémas localement de
type fini sur $S$, et soit $p \geq 0$. Dans ce cadre, on définit
$$
A^{(p)}(Z \to X) =
\prod\nolimits_{i \leq p - 1} A^i(X) \times
\prod\nolimits_{i \geq p} A^i(Z \to X).
$$
Alors $A^{(p)}(Z \to X)$ est canoniquement muni d'une structure
d'algèbre graduée. En fait, plus généralement, on dispose d'une multiplication
$$
A^{(p)}(Z \to X) \times A^{(q)}(Z \to X)
\longrightarrow A^{(\max(p, q))}(Z \to X)
$$
Pour les définir, on définit des applications
\begin{align*}
A^i(Z \to X) \times A^j(X) & \to A^{i + j}(Z \to X) \\
A^i(X) \times A^j(Z \to X) & \to A^{i + j}(Z \to X) \\
A^i(Z \to X) \times A^j(Z \to X) & \to A^{i + j}(Z \to X)
\end{align*}
Pour la première, on utilise la composition des classes bivariantes.
Pour la deuxième, on utilise la restriction
$A^i(X) \to A^i(Z)$ (remarque \ref{chow-remark-restriction-bivariant}) et la
composition $A^i(Z) \times A^j(Z \to X) \to A^{i + j}(Z \to X)$.
Pour la troisième, on envoie $(c, c')$ sur $res(c) \circ c'$,
où $res : A^i(Z \to X) \to A^i(Z)$ est l'application de restriction (voir la
remarque \ref{chow-remark-restriction-bivariant}). Nous omettons la
vérification de l'associativité de ces multiplications en un sens approprié.
\end{remark}

\begin{remark}
\label{chow-remark-res-push}
Soit $(S, \delta)$ comme dans la situation \ref{chow-situation-setup}.
Soit $Z \to X$ une immersion fermée de schémas localement de
type fini sur $S$. Notons $res : A^p(Z \to X) \to A^p(Z)$
l'application de restriction de la remarque \ref{chow-remark-restriction-bivariant}.
Pour $c \in A^p(Z \to X)$, on a
$res(c) \cap \alpha = c \cap i_*\alpha$ pour $\alpha \in \CH_*(Z)$.
En effet, $res(c) \cap \alpha = c \cap \alpha$,
et la compatibilité de $c$ avec l'image directe propre
donne $(Z \to Z)_*(c \cap \alpha) = c \cap (Z \to X)_*\alpha$.
\end{remark}






\section{Lemmes sur les classes bivariantes}
\label{chow-section-bivariant-II}

\noindent
Dans cette section, nous démontrons quelques résultats élémentaires sur les
classes bivariantes. Voici un critère permettant de voir qu'une opération
passe au quotient par équivalence rationnelle.

\begin{lemma}
\label{chow-lemma-factors-through-rational-equivalence}
\begin{reference}
Forme très faible de \cite[Théorème 17.1]{F}
\end{reference}
Soit $(S, \delta)$ comme dans la situation \ref{chow-situation-setup}.
Soit $f : X \to Y$ un morphisme de schémas localement de type fini sur $S$.
Soit $p \in \mathbf{Z}$. Supposons donnée une règle
qui associe à tout morphisme localement de type fini $Y' \to Y$
et à tout $k$ une application
$$
c \cap - : Z_k(Y') \longrightarrow \CH_{k - p}(X')
$$
où $Y' = X' \times_X Y$, satisfaisant à la condition (3) de la
définition \ref{chow-definition-bivariant-class}
dès que $\mathcal{L}'|_{D'} \cong \mathcal{O}_{D'}$. Alors
$c \cap -$ se factorise par l'équivalence rationnelle.
\end{lemma}

\begin{proof}
L'énoncé a un sens car, étant donné un triplet
$(\mathcal{L}, s, i : D \to X)$ comme dans la
définition \ref{chow-definition-gysin-homomorphism}
tel que $\mathcal{L}|_D \cong \mathcal{O}_D$, l'opération
$i^*$ est définie au niveau des cycles, voir la
remarque \ref{chow-remark-gysin-on-cycles}.
Soit $\alpha \in Z_k(X')$ un cycle rationnellement équivalent à zéro.
Nous devons montrer que $c \cap \alpha = 0$. D'après le
lemme \ref{chow-lemma-rational-equivalence-family},
il existe un cycle $\beta \in Z_{k + 1}(X' \times \mathbf{P}^1)$
tel que $\alpha = i_0^*\beta - i_\infty^*\beta$,
où $i_0, i_\infty : X' \to X' \times \mathbf{P}^1$ sont les
immersions fermées de $X'$ au-dessus de $0, \infty$. Puisque ce sont
des exemples de diviseurs de Cartier effectifs à fibrés normaux
triviaux, on voit que $c \cap i_0^*\beta = j_0^*(c \cap \beta)$
et $c \cap i_\infty^*\beta = j_\infty^*(c \cap \beta)$,
où $j_0, j_\infty : Y' \to Y' \times \mathbf{P}^1$ sont les
immersions fermées comme ci-dessus. Comme
$j_0^*(c \cap \beta) \sim_{rat} j_\infty^*(c \cap \beta)$
(cela résulte du lemme \ref{chow-lemma-rational-equivalence-family}), on conclut.
\end{proof}

\begin{lemma}
\label{chow-lemma-bivariant-weaker}
\begin{reference}
Forme faible de \cite[Théorème 17.1]{F}
\end{reference}
Soit $(S, \delta)$ comme dans la situation \ref{chow-situation-setup}.
Soit $f : X \to Y$ un morphisme de schémas localement de type fini sur $S$.
Soit $p \in \mathbf{Z}$. Supposons donnée une règle
qui associe à tout morphisme localement de type fini $Y' \to Y$
et à tout $k$ une application
$$
c \cap - : \CH_k(Y') \longrightarrow \CH_{k - p}(X')
$$
où $Y' = X' \times_X Y$, satisfaisant aux conditions (1), (2) de la
définition \ref{chow-definition-bivariant-class}
et à la condition (3) dès que $\mathcal{L}'|_{D'} \cong \mathcal{O}_{D'}$. Alors
$c \cap -$ est une classe bivariante.
\end{lemma}

\begin{proof}
Soit $Y' \to Y$ un morphisme de schémas localement de type fini.
Soit $(\mathcal{L}', s', i' : D' \to Y')$ comme dans la
définition \ref{chow-definition-gysin-homomorphism},
et soit $(\mathcal{N}', t', j' : E' \to X')$ son image réciproque sur $X'$.
Nous devons montrer que $c \cap (i')^*\alpha' = (j')^*(c \cap \alpha')$
pour tout $\alpha' \in \CH_k(Y')$.

\medskip\noindent
Notons $g : Y'' \to Y'$ le morphisme lisse de dimension relative
$1$, avec $i'' : D'' \to Y''$ et $p : D'' \to D'$,
construit dans le lemme \ref{chow-lemma-normal-cone-effective-Cartier}.
(Attention : $D''$ n'est pas l'image réciproque entière de $D'$.)
Notons $f : X'' \to X'$ et $E'' \subset X''$
leurs changements de base par $X' \to Y'$. On a le diagramme
$$
\xymatrix{
& X'' \ar[rr] \ar'[d][dd]_h & & Y'' \ar[dd]^g \\
E'' \ar[rr] \ar[dd]_q \ar[ru]^{j''} & & D'' \ar[dd]^p \ar[ru]^{i''} & \\
& X' \ar'[r][rr] & & Y' \\
E' \ar[rr] \ar[ru]^{j'} & & D' \ar[ru]^{i'}
}
$$
D'après les propriétés données dans le lemme, nous savons que $\beta' = (i')^*\alpha'$
est l'unique élément de $\CH_{k - 1}(D')$ tel que
$p^*\beta' = (i'')^*g^*\alpha'$. De même, nous savons que
$\gamma' = (j')^*(c \cap \alpha')$ est l'unique élément de
$\CH_{k - 1 - p}(E')$ tel que $q^*\gamma' = (j'')^*h^*(c \cap \alpha')$.
Or nous savons que
$$
(j'')^*h^*(c \cap \alpha') =
(j'')^*(c \cap g^*\alpha') =
c \cap (i'')^*g^*\alpha'
$$
par nos hypothèses sur $c$ ; remarquons que la version modifiée de (3)
supposée dans l'énoncé du lemme s'applique à $i''$
et à son changement de base $j''$. Nous savons de même que
$$
q^*(c \cap \beta') = c \cap p^*\beta'
$$
On conclut que $\gamma' = c \cap \beta'$ par l'unicité signalée
ci-dessus.
\end{proof}

\noindent
Voici un critère pour qu'une classe bivariante soit nulle.

\begin{lemma}
\label{chow-lemma-bivariant-zero}
Soit $(S, \delta)$ comme dans la situation \ref{chow-situation-setup}.
Soit $f : X \to Y$ un morphisme de schémas localement de type fini sur $S$.
Soit $c \in A^p(X \to Y)$. Pour $Y'' \to Y' \to Y$, posons
$X'' = Y'' \times_Y X$ et $X' = Y' \times_Y X$.
Les assertions suivantes sont équivalentes :
\begin{enumerate}
\item $c$ est nulle,
\item $c \cap [Y'] = 0$ dans $\CH_*(X')$ pour tout schéma intègre $Y'$
localement de type fini sur $Y$, et
\item pour tout schéma intègre $Y'$ localement de type fini sur $Y$,
il existe un morphisme propre birationnel $Y'' \to Y'$ tel que
$c \cap [Y''] = 0$ dans $\CH_*(X'')$.
\end{enumerate}
\end{lemma}

\begin{proof}
Les implications (1) $\Rightarrow$ (2) $\Rightarrow$ (3) sont claires.
L'hypothèse (3) implique (2) car $(Y'' \to Y')_*[Y''] = [Y']$
et donc $c \cap [Y'] = (X'' \to X')_*(c \cap [Y''])$, puisque $c$
est une classe bivariante. Supposons (2).
Soit $Y' \to Y$ localement de type fini. Soit $\alpha \in \CH_k(Y')$.
Écrivons $\alpha = \sum n_i [Y'_i]$, où les $Y'_i \subset Y'$ forment une
famille localement finie de sous-schémas fermés intègres de $\delta$-dimension $k$.
On voit alors que $\alpha$ est l'image directe du cycle
$\alpha' = \sum n_i[Y'_i]$ sur $Y'' = \coprod Y'_i$ par le
morphisme propre $Y'' \to Y'$. Par les propriétés des classes bivariantes,
il suffit de prouver que $c \cap \alpha' = 0$ dans $\CH_{k - p}(X'')$.
On a $\CH_{k - p}(X'') = \prod \CH_{k - p}(X'_i)$, où
$X'_i = Y'_i \times_Y X$. Cela résulte immédiatement
des définitions. Les applications de projection
$\CH_{k - p}(X'') \to \CH_{k - p}(X'_i)$ sont données par l'image réciproque
plate. Puisque le produit cap avec $c$ commute à l'image réciproque
plate, on voit qu'il suffit de montrer que $c \cap [Y'_i]$
est nul dans $\CH_{k - p}(X'_i)$, ce qui est vrai par hypothèse.
\end{proof}

\begin{lemma}
\label{chow-lemma-disjoint-decomposition-bivariant}
Soit $(S, \delta)$ comme dans la situation \ref{chow-situation-setup}.
Soit $f : X \to Y$ un morphisme de schémas localement de type fini sur $S$.
Supposons données des décompositions en réunions disjointes
$X = \coprod_{i \in I} X_i$ et $Y = \coprod_{j \in J} Y_j$
par des sous-schémas ouverts et fermés,
ainsi qu'une application d'ensembles $a : I \to J$ telle que $f(X_i) \subset Y_{a(i)}$.
Alors
$$
A^p(X \to Y) = \prod\nolimits_{i \in I} A^p(X_i \to Y_{a(i)})
$$
\end{lemma}

\begin{proof}
Supposons donné un élément $(c_i) \in \prod_i A^p(X_i \to Y_{a(i)})$.
Étant donné $\beta \in \CH_k(Y)$, on peut alors l'envoyer sur l'élément de
$\CH_{k - p}(X)$ dont la restriction à $X_i$ est $c_i \cap \beta|_{Y_{a(i)}}$.
Cela fonctionne car $\CH_{k - p}(X) = \prod_i \CH_{k - p}(X_i)$.
La même construction fonctionne après tout changement de base $Y' \to Y$
localement de type fini, et l'on obtient $c \in A^p(X \to Y)$.
On obtient ainsi une application $\Psi$ du membre droit de la formule
vers son membre gauche.
Réciproquement, étant donnés $c \in A^p(X \to Y)$ et un élément
$\beta_i \in \CH_k(Y_{a(i)})$, on peut considérer l'élément
$(c \cap (Y_{a(i)} \to Y)_*\beta_i)|_{X_i}$ de $\CH_{k - p}(X_i)$.
La même chose fonctionne après tout changement de base $Y' \to Y$
localement de type fini, et l'on obtient $c_i \in A^p(X_i \to Y_{a(i)})$.
On obtient ainsi une application $\Phi$ du membre gauche
de la formule vers son membre droit.
Il est immédiat que $\Phi \circ \Psi = \text{id}$.
Réciproquement, supposons que $c \in A^p(X \to Y)$ et
$\beta \in \CH_k(Y)$. Écrivons $\Phi(c) = (c_i)$. Soit $j \in J$.
Comme $c$ commute à l'image réciproque plate, on a
$$
(c \cap \beta)|_{\coprod_{a(i) = j} X_i} =
c \cap \beta|_{Y_j}
$$
Comme $c$ commute à l'image directe propre, on a
$$
(\coprod\nolimits_{a(i) = j} X_i \to X)_*
((c \cap \beta)|_{\coprod_{a(i) = j} X_i})
=
c \cap (Y_j \to Y)_*\beta|_{Y_j}
$$
Le membre gauche est le cycle sur $X$ qui se restreint en $(c \cap \beta)|_{X_i}$
sur $X_i$ pour $i \in I$ tel que $a(i) = j$, et qui vaut $0$ sinon.
Le membre droit est un cycle sur $X$ dont la restriction à $X_i$
est $c_i \cap \beta|_{Y_j}$ pour $i \in I$ tel que $a(i) = j$.
Ainsi, $c \cap \beta = \Psi((c_i))$, comme voulu.
\end{proof}

\begin{remark}
\label{chow-remark-completion-bivariant}
Soit $(S, \delta)$ comme dans la situation \ref{chow-situation-setup}.
Soit $f : X \to Y$ un morphisme de schémas localement de type fini sur $S$.
Soient $X = \coprod_{i \in I} X_i$ et $Y = \coprod_{j \in J} Y_j$
les décompositions de $X$ et de $Y$ en leurs composantes connexes
(les composantes connexes sont ouvertes puisque $X$ et $Y$ sont localement noethériens ; voir
Topologie, Lemme \ref{topology-lemma-locally-Noetherian-locally-connected}, et
Propriétés, Lemme \ref{properties-lemma-Noetherian-topology}).
Soit $a(i) \in J$ l'indice tel que $f(X_i) \subset Y_{a(i)}$.
Alors $A^p(X \to Y) = \prod A^p(X_i \to Y_{a(i)})$ d'après le
Lemme \ref{chow-lemma-disjoint-decomposition-bivariant}.
Dans ce cadre, il est commode de poser
$$
A^*(X \to Y)^\wedge = \prod\nolimits_i A^*(X_i \to Y_{a(i)})
$$
Ce groupe bivariant « complété » est le sous-ensemble
$$
A^*(X \to Y)^\wedge \quad\subset\quad \prod\nolimits_{p \geq 0} A^p(X \to Y)
$$
formé des éléments $c = (c_0, c_1, c_2, \ldots)$ tels que,
pour toute composante connexe $X_i$, l'image de $c_p$ dans
$A^p(X_i \to Y_{a(i)})$ soit nulle pour presque tout $p$.
Si $Y \to Z$ est un second morphisme, alors la
composition $A^*(X \to Y) \times A^*(Y \to Z) \to A^*(X \to Z)$
s'étend en une composition
$A^*(X \to Y)^\wedge \times A^*(Y \to Z)^\wedge \to A^*(X \to Z)^\wedge$
des complétés.
On appelle parfois $A^*(X)^\wedge = A^*(X \to X)^\wedge$
l'{\it anneau de cohomologie bivariante complété} de $X$.
\end{remark}

\begin{lemma}
\label{chow-lemma-envelope-bivariant}
Soit $(S, \delta)$ comme dans la situation \ref{chow-situation-setup}.
Soit $f : X \to Y$ un morphisme de schémas localement de type fini sur $S$.
Soit $g : Y' \to Y$ une enveloppe (Définition \ref{chow-definition-envelope})
et notons $X' = Y' \times_Y X$. Soit $p \in \mathbf{Z}$ et soit
$c' \in A^p(X' \to Y')$. Si les deux restrictions
$$
res_1(c') = res_2(c') \in A^p(X' \times_X X' \to Y' \times_Y Y')
$$
sont égales (voir la démonstration), alors il existe un unique $c \in A^p(X \to Y)$
dont la restriction $res(c) = c'$ dans $A^p(X' \to Y')$.
\end{lemma}

\begin{proof}
On a un diagramme commutatif
$$
\xymatrix{
X' \times_X X' \ar[d]^{f''} \ar@<1ex>[r]^-a \ar@<-1ex>[r]_-b &
X' \ar[d]^{f'} \ar[r]_h &
X \ar[d]^f \\
Y' \times_Y Y' \ar@<1ex>[r]^-p \ar@<-1ex>[r]_-q &
Y' \ar[r]^g &
Y
}
$$
L'élément $res_1(c')$ est la restriction (voir
Remarque \ref{chow-remark-restriction-bivariant}) de $c'$ pour le carré cartésien
de morphismes $a, f', p, f''$, et l'élément $res_2(c')$ est la restriction
de $c'$ pour le carré cartésien de morphismes $b, f', q, f''$.
Supposons $res_1(c') = res_2(c')$ et soit $\beta \in \CH_k(Y)$.
D'après le Lemme \ref{chow-lemma-envelope}, on peut trouver un $\beta' \in \CH_k(Y')$
tel que $g_*\beta' = \beta$. On pose alors
$$
c \cap \beta = h_*(c' \cap \beta')
$$
Pour voir que ceci ne dépend pas du choix de
$\beta'$, il suffit de montrer que
$h_*(c' \cap (p_*\gamma - q_*\gamma))$ est nul
pour $\gamma \in \CH_k(Y' \times_Y Y')$.
Puisque $c'$ est une classe bivariante, on a
$$
h_*(c' \cap (p_*\gamma - q_*\gamma)) =
h_*(a_*(c' \cap \gamma) - b_*(c' \cap \gamma)) = 0
$$
la dernière égalité résultant de $h_* \circ a_* = h_* \circ b_*$,
car $h \circ a = h \circ b$.

\medskip\noindent
Remarquons que notre choix de $c \cap \beta$ est imposé
par la condition $res(c) = c'$ et par la compatibilité
des classes bivariantes avec l'image directe propre.

\medskip\noindent
Bien entendu, pour définir la classe bivariante $c$, il nous faut
construire des applications $c \cap -: \CH_k(Y_1) \to \CH_{k + p}(Y_1 \times_Y X)$
pour tout morphisme $Y_1 \to Y$ localement de type fini satisfaisant aux
conditions énumérées dans la Définition \ref{chow-definition-bivariant-class}.
Notons $Y'_1 = Y' \times_Y Y_1$, $X_1 = X \times_Y Y_1$.
Le morphisme $Y'_1 \to Y_1$ est une enveloppe d'après le
Lemme \ref{chow-lemma-base-change-envelope}. Nous pouvons donc utiliser le
diagramme obtenu par changement de base
$$
\xymatrix{
X'_1 \times_{X_1} X'_1 \ar[d]^{f''_1} \ar@<1ex>[r]^-{a_1} \ar@<-1ex>[r]_-{b_1} &
X'_1 \ar[d]^{f'_1} \ar[r]_{h_1} &
X_1 \ar[d]^{f_1} \\
Y'_1 \times_{Y_1} Y'_1 \ar@<1ex>[r]^-{p_1} \ar@<-1ex>[r]_-{q_1} &
Y'_1 \ar[r]^{g_1} &
Y_1
}
$$
et les mêmes arguments pour obtenir une application bien définie
$c \cap - : \CH_k(Y_1) \to \CH_{k + p}(X_1)$ comme ci-dessus.

\medskip\noindent
Il nous faut ensuite vérifier les conditions (1), (2) et (3) de la
Définition \ref{chow-definition-bivariant-class} pour $c$.
Par exemple, supposons que $t : Y_2 \to Y_1$ soit un morphisme propre
de schémas localement de type fini sur $Y$. Notons comme ci-dessus
les changements de base du premier diagramme à $Y_1$, resp.\ à $Y_2$,
par les indices ${}_1$, resp.\ ${}_2$. Notons $t' : Y'_2 \to Y'_1$,
$s : X_2 \to X_1$ et $s' : X'_2 \to X'_1$ les changements de base de $t$
à $Y'$, $X$ et $X'$. Nous devons montrer que
$$
s_*(c \cap \beta_2) = c \cap t_*\beta_2
$$
pour $\beta_2 \in \CH_k(Y_2)$. Choisissons $\beta'_2 \in \CH_k(Y'_2)$
tel que $g_{2, *}\beta'_2 = \beta_2$. Puisque $c'$ est une classe bivariante
et que les diagrammes
$$
\vcenter{
\xymatrix{
X'_2 \ar[d]_{s'} \ar[r]_{h_2} & X_2 \ar[d]^s \\
X'_1 \ar[r]^{h_1} & X_1
}
}
\quad\text{et}\quad
\vcenter{
\xymatrix{
X'_2 \ar[d]_{s'} \ar[r]_{f'_2} & Y'_2 \ar[d]^{t'} \\
X'_2 \ar[r]^{f'_1} & Y'_1
}
}
$$
sont cartésiens, on a
$$
s_*(c \cap \beta_2) =
s_*(h_{2, *}(c' \cap \beta'_2)) =
h_{1, *}s'_*(c' \cap \beta'_2) =
h_{1, *}(c' \cap (t'_*\beta'_2))
$$
et l'expression finale calcule $c \cap t_*\beta_2$ par construction :
$t'_*\beta'_2 \in \CH_k(Y'_1)$ est une classe dont l'image par $g_{1, *}$ est
$t_*\beta_2$. Ceci démontre la condition (1).
Les autres conditions se démontrent de la même manière,
et nous omettons les arguments détaillés.
\end{proof}















\section{Formule du fibré projectif}
\label{chow-section-projective-space-bundle-formula}

\noindent
Soit $(S, \delta)$ comme dans la situation \ref{chow-situation-setup}.
Soit $X$ localement de type fini sur $S$.
Considérons un $\mathcal{O}_X$-module localement libre de type fini
$\mathcal{E}$ de rang $r$.
Notre convention est que le {\it fibré projectif associé à
$\mathcal{E}$} est le morphisme
$$
\xymatrix{
\mathbf{P}(\mathcal{E}) =
\underline{\text{Proj}}_X(\text{Sym}^*(\mathcal{E}))
\ar[r]^-\pi
& X
}
$$
au-dessus de $X$, avec
$\mathcal{O}_{\mathbf{P}(\mathcal{E})}(1)$ normalisé de sorte que
$\pi_*(\mathcal{O}_{\mathbf{P}(\mathcal{E})}(1)) = \mathcal{E}$.
En particulier, il existe une surjection
$\pi^*\mathcal{E} \to \mathcal{O}_{\mathbf{P}(\mathcal{E})}(1)$.
Nous dirons informellement « soit $(\pi : P \to X, \mathcal{O}_P(1))$
le fibré projectif associé à $\mathcal{E}$ » pour désigner
la situation où $P = \mathbf{P}(\mathcal{E})$ et
$\mathcal{O}_P(1) = \mathcal{O}_{\mathbf{P}(\mathcal{E})}(1)$.

\begin{lemma}
\label{chow-lemma-cap-projective-bundle}
Soit $(S, \delta)$ comme dans la situation \ref{chow-situation-setup}.
Soit $X$ localement de type fini sur $S$.
Soit $\mathcal{E}$ un $\mathcal{O}_X$-module localement libre de type fini
$\mathcal{E}$ de rang $r$. Soit $(\pi : P \to X, \mathcal{O}_P(1))$
le fibré projectif associé à $\mathcal{E}$.
Pour tout $\alpha \in \CH_k(X)$, l'élément
$$
\pi_*\left(
c_1(\mathcal{O}_P(1))^s \cap \pi^*\alpha
\right)
\in
\CH_{k + r - 1 - s}(X)
$$
est $0$ si $s < r - 1$ et est égal à $\alpha$ lorsque $s = r - 1$.
\end{lemma}

\begin{proof}
Soit $Z \subset X$ un sous-schéma fermé intègre de $\delta$-dimension $k$.
Remarquons que $\pi^*[Z] = [\pi^{-1}(Z)]$, car $\pi^{-1}(Z)$ est intègre de
$\delta$-dimension $r - 1$.
Si $s < r - 1$, alors, par construction,
$c_1(\mathcal{O}_P(1))^s \cap \pi^*[Z]$
est représenté par un $(k + r - 1 - s)$-cycle supporté sur
$\pi^{-1}(Z)$. L'image directe de ce cycle
est donc nulle pour des raisons de dimension.

\medskip\noindent
Soit $s = r - 1$. L'argument donné ci-dessus montre que
$\pi_*(c_1(\mathcal{O}_P(1))^s \cap \pi^*\alpha) = n [Z]$
pour un certain $n \in \mathbf{Z}$. Nous voulons montrer que $n = 1$.
Pour les mêmes raisons de dimension
que ci-dessus, il suffit de démontrer ce résultat après avoir remplacé $X$ par
$X \setminus T$, où $T \subset Z$ est une partie fermée propre.
Soit $\xi$ le point générique de $Z$.
On peut choisir des éléments $e_1, \ldots, e_{r - 1} \in \mathcal{E}_\xi$
qui font partie d'une base de $\mathcal{E}_\xi$.
Ils donnent des sections rationnelles $s_1, \ldots, s_{r - 1}$
de $\mathcal{O}_P(1)|_{\pi^{-1}(Z)}$ dont le lieu commun des zéros
est l'adhérence de l'image d'une section rationnelle de
$\mathbf{P}(\mathcal{E}|_Z) \to Z$, réunie à une partie fermée dont le
support s'envoie dans une partie fermée propre $T$ de $Z$.
Après avoir retiré $T$ de $X$ (et, de façon correspondante, $\pi^{-1}(T)$
de $P$), on voit que $s_1, \ldots, s_n$ forment une suite
de sections globales
$s_i \in \Gamma(\pi^{-1}(Z), \mathcal{O}_{\pi^{-1}(Z)}(1))$
dont le lieu commun des zéros est l'image d'une section $Z \to \pi^{-1}(Z)$.
On voit donc successivement que
\begin{eqnarray*}
\pi^*[Z] & = & [\pi^{-1}(Z)] \\
c_1(\mathcal{O}_P(1)) \cap \pi^*[Z] & = & [Z(s_1)] \\
c_1(\mathcal{O}_P(1))^2 \cap \pi^*[Z] & = & [Z(s_1) \cap Z(s_2)] \\
\ldots & = & \ldots \\
c_1(\mathcal{O}_P(1))^{r - 1} \cap \pi^*[Z] & = &
[Z(s_1) \cap \ldots \cap Z(s_{r - 1})]
\end{eqnarray*}
par applications répétées du Lemme \ref{chow-lemma-geometric-cap}.
Puisque l'image directe par $\pi$ de l'image d'une
section de $\pi$ au-dessus de $Z$ est clairement $[Z]$, on obtient le résultat
lorsque $\alpha = [Z]$. Nous omettons la vérification que ces
arguments donnent le résultat pour un cycle général $\alpha = \sum n_j [Z_j]$.
\end{proof}

\begin{lemma}[Formule du fibré projectif]
\label{chow-lemma-chow-ring-projective-bundle}
Soit $(S, \delta)$ comme dans la situation \ref{chow-situation-setup}.
Soit $X$ localement de type fini sur $S$.
Soit $\mathcal{E}$ un $\mathcal{O}_X$-module localement libre de type fini
$\mathcal{E}$ de rang $r$. Soit $(\pi : P \to X, \mathcal{O}_P(1))$
le fibré projectif associé à $\mathcal{E}$.
L'application
$$
\bigoplus\nolimits_{i = 0}^{r - 1}
\CH_{k + i}(X)
\longrightarrow
\CH_{k + r - 1}(P),
$$
$$
(\alpha_0, \ldots, \alpha_{r-1})
\longmapsto
\pi^*\alpha_0 +
c_1(\mathcal{O}_P(1)) \cap \pi^*\alpha_1
+ \ldots +
c_1(\mathcal{O}_P(1))^{r - 1} \cap \pi^*\alpha_{r-1}
$$
est un isomorphisme.
\end{lemma}

\begin{proof}
Fixons $k \in \mathbf{Z}$. Montrons d'abord que l'application est injective.
Supposons que $(\alpha_0, \ldots, \alpha_{r - 1})$ soit un élément
du membre gauche qui s'envoie sur zéro.
D'après le Lemme \ref{chow-lemma-cap-projective-bundle}, on voit que
$$
0 = \pi_*(\pi^*\alpha_0 +
c_1(\mathcal{O}_P(1)) \cap \pi^*\alpha_1
+ \ldots +
c_1(\mathcal{O}_P(1))^{r - 1} \cap \pi^*\alpha_{r-1})
= \alpha_{r - 1}
$$
Ensuite, on voit que
$$
0 = \pi_*(c_1(\mathcal{O}_P(1)) \cap (\pi^*\alpha_0 +
c_1(\mathcal{O}_P(1)) \cap \pi^*\alpha_1
+ \ldots +
c_1(\mathcal{O}_P(1))^{r - 2} \cap \pi^*\alpha_{r - 2}))
= \alpha_{r - 2}
$$
et ainsi de suite. L'application est donc injective.

\medskip\noindent
Il reste à montrer que l'application est surjective.
Soient $X_i$, $i \in I$, les composantes irréductibles de $X$.
Alors $P_i = \mathbf{P}(\mathcal{E}|_{X_i})$, $i \in I$,
sont les composantes irréductibles de $P$. Considérons le
diagramme commutatif
$$
\xymatrix{
\coprod P_i \ar[d]_{\coprod \pi_i} \ar[r]_p & P \ar[d]^\pi \\
\coprod X_i \ar[r]^q & X
}
$$
Remarquons que $p_*$ est surjective. Si $\beta \in \CH_k(\coprod X_i)$,
alors $\pi^* q_* \beta = p_*(\coprod \pi_i)^* \beta$ ; voir le
Lemme \ref{chow-lemma-flat-pullback-proper-pushforward}. Il en va de même pour
le produit cap avec $c_1(\mathcal{O}(1))$, d'après le
Lemme \ref{chow-lemma-pushforward-cap-c1}.
Par conséquent, si l'application du lemme est surjective pour chacun
des morphismes $\pi_i : P_i \to X_i$, alors elle est
surjective pour $\pi : P \to X$. On peut donc supposer $X$ irréductible.
Ainsi, $\dim_\delta(X) < \infty$, et l'on peut en particulier raisonner
par récurrence sur $\dim_\delta(X)$.

\medskip\noindent
Le résultat est clair si $\dim_\delta(X) < k$.
Soit $\alpha \in \CH_{k + r - 1}(P)$.
Pour tout sous-schéma localement fermé $T \subset X$, notons
$\gamma_T : \bigoplus \CH_{k + i}(T) \to \CH_{k + r - 1}(\pi^{-1}(T))$
l'application
$$
\gamma_T(\alpha_0, \ldots, \alpha_{r - 1})
= \pi^*\alpha_0 + \ldots +
c_1(\mathcal{O}_{\pi^{-1}(T)}(1))^{r - 1} \cap \pi^*\alpha_{r - 1}.
$$
Supposons que, pour un certain ouvert non vide $U \subset X$, on ait
$\alpha|_{\pi^{-1}(U)} = \gamma_U(\alpha_0, \ldots, \alpha_{r - 1})$.
On peut alors choisir des relèvements $\alpha'_i \in \CH_{k + i}(X)$, et l'on
voit que $\alpha - \gamma_X(\alpha'_0, \ldots, \alpha'_{r - 1})$
est, d'après le Lemme \ref{chow-lemma-restrict-to-open},
rationnellement équivalent à un $k$-cycle sur $P_Y = \mathbf{P}(\mathcal{E}|_Y)$,
où $Y = X \setminus U$ est muni de sa structure de sous-schéma fermé réduit.
Remarquons que $\dim_\delta(Y) < \dim_\delta(X)$.
Par récurrence, le résultat vaut
pour $P_Y \to Y$, et il vaut donc pour $\alpha$.
On peut par conséquent remplacer $X$ par n'importe quel ouvert non vide de $X$.

\medskip\noindent
En particulier, on peut supposer que $\mathcal{E} \cong \mathcal{O}_X^{\oplus r}$.
Dans ce cas, $\mathbf{P}(\mathcal{E}) = X \times \mathbf{P}^{r - 1}$.
Utilisons la stratification
$$
\mathbf{P}^{r - 1} = \mathbf{A}^{r - 1}
\amalg \mathbf{A}^{r - 2}
\amalg \ldots
\amalg \mathbf{A}^0
$$
L'adhérence de chaque strate est un $\mathbf{P}^{r - 1 - i}$ qui représente
$c_1(\mathcal{O}(1))^i \cap [\mathbf{P}^{r - 1}]$.
Ainsi, $P$ possède une stratification analogue
$$
P = U^{r - 1} \amalg U^{r - 2} \amalg \ldots \amalg U^0
$$
Soit $P^i$ l'adhérence de $U^i$. Soit $\pi^i : P^i \to X$
la restriction de $\pi$ à $P^i$.
Soit $\alpha \in \CH_{k + r - 1}(P)$. D'après le
Lemme \ref{chow-lemma-pullback-affine-fibres-surjective},
on peut écrire $\alpha|_{U^{r - 1}} = \pi^*\alpha_0|_{U^{r - 1}}$
pour un certain $\alpha_0 \in \CH_k(X)$. La différence
$\alpha - \pi^*\alpha_0$ est donc l'image d'un certain
$\alpha' \in \CH_{k + r - 1}(P^{r - 2})$.
En appliquant de nouveau le Lemme \ref{chow-lemma-pullback-affine-fibres-surjective},
on peut écrire
$\alpha'|_{U^{r - 2}} = (\pi^{r - 2})^*\alpha_1|_{U^{r - 2}}$
pour un certain $\alpha_1 \in \CH_{k + 1}(X)$.
D'après le Lemme \ref{chow-lemma-relative-effective-cartier},
on voit que l'image de $(\pi^{r - 2})^*\alpha_1$
représente $c_1(\mathcal{O}_P(1)) \cap \pi^*\alpha_1$.
On voit également que
$\alpha - \pi^*\alpha_0 - c_1(\mathcal{O}_P(1)) \cap \pi^*\alpha_1$
est l'image d'un certain $\alpha'' \in \CH_{k + r - 1}(P^{r - 3})$.
Et ainsi de suite.
\end{proof}

\begin{lemma}
\label{chow-lemma-vectorbundle}
Soit $(S, \delta)$ comme dans la situation \ref{chow-situation-setup}.
Soit $X$ localement de type fini sur $S$.
Soit $\mathcal{E}$ un faisceau localement libre de type fini de rang $r$ sur $X$.
Soit
$$
p :
E = \underline{\Spec}(\text{Sym}^*(\mathcal{E}))
\longrightarrow
X
$$
le fibré vectoriel associé au-dessus de $X$.
Alors $p^* : \CH_k(X) \to \CH_{k + r}(E)$ est un isomorphisme pour tout $k$.
\end{lemma}

\begin{proof}
(Pour le cas des fibrés en droites, voir le Lemme \ref{chow-lemma-linebundle}.)
Pour la surjectivité, voir le Lemme \ref{chow-lemma-pullback-affine-fibres-surjective}.
Soit $(\pi  : P \to X, \mathcal{O}_P(1))$
le fibré projectif associé
au faisceau localement libre de type fini $\mathcal{E} \oplus \mathcal{O}_X$.
Soit $s \in \Gamma(P, \mathcal{O}_P(1))$ correspondant à la section globale
$(0, 1) \in \Gamma(X, \mathcal{E} \oplus \mathcal{O}_X)$.
Soit $D = Z(s) \subset P$. Remarquons que
$(\pi|_D : D \to X , \mathcal{O}_P(1)|_D)$
est le fibré projectif associé
à $\mathcal{E}$. On note $\pi_D = \pi|_D$ et
$\mathcal{O}_D(1) = \mathcal{O}_P(1)|_D$.
De plus, $D$ est un diviseur de Cartier effectif
sur $P$. Ainsi, $\mathcal{O}_P(D) = \mathcal{O}_P(1)$
(voir Diviseurs, Lemme \ref{divisors-lemma-characterize-OD}).
Il existe aussi un isomorphisme
$E \cong P \setminus D$. Notons $j : E \to P$
l'immersion ouverte correspondante.
Pour l'injectivité, on utilise que le noyau de
$$
j^* :
\CH_{k + r}(P)
\longrightarrow
\CH_{k + r}(E)
$$
est formé des cycles supportés dans le diviseur de Cartier effectif $D$ ;
voir le Lemme \ref{chow-lemma-restrict-to-open}. Ainsi, si $p^*\alpha = 0$, alors
$\pi^*\alpha = i_*\beta$ pour un certain $\beta \in \CH_{k + r}(D)$.
D'après le Lemme \ref{chow-lemma-chow-ring-projective-bundle}, on peut écrire
$$
\beta = \pi_D^*\beta_0 +
\ldots + c_1(\mathcal{O}_D(1))^{r - 1} \cap \pi_D^* \beta_{r - 1}.
$$
pour certains $\beta_i \in \CH_{k + i}(X)$.
D'après les Lemmes \ref{chow-lemma-relative-effective-cartier}
et \ref{chow-lemma-pushforward-cap-c1},
cela implique
$$
\pi^*\alpha = i_*\beta =
c_1(\mathcal{O}_P(1)) \cap \pi^*\beta_0 +
\ldots +
c_1(\mathcal{O}_D(1))^r \cap \pi^*\beta_{r - 1}.
$$
Puisque le rang de $\mathcal{E} \oplus \mathcal{O}_X$ est $r + 1$,
cela contredit le Lemme \ref{chow-lemma-pushforward-cap-c1}, à moins que tous les
$\alpha$ et tous les $\beta_i$ ne soient nuls.
\end{proof}








\section{Les classes de Chern d'un fibré vectoriel}
\label{chow-section-chern-classes-vector-bundles}

\noindent
Nous pouvons utiliser la formule du fibré en espaces projectifs pour définir les
classes de Chern d'un fibré vectoriel de rang $r$ au moyen du développement
de $c_1(\mathcal{O}(1))^r$ suivant les puissances inférieures, voir
la formule (\ref{chow-equation-chern-classes}).
La raison des signes sera expliquée plus loin.

\begin{definition}
\label{chow-definition-chern-classes}
Soit $(S, \delta)$ comme dans la Situation \ref{chow-situation-setup}.
Soit $X$ localement de type fini sur $S$.
Supposons $X$ intègre et $n = \dim_\delta(X)$.
Soit $\mathcal{E}$ un faisceau localement libre de type fini de rang $r$
sur $X$. Soit $(\pi : P \to X, \mathcal{O}_P(1))$ le fibré en espaces
projectifs associé à $\mathcal{E}$.
\begin{enumerate}
\item D'après le Lemme \ref{chow-lemma-chow-ring-projective-bundle}, il existe des
éléments $c_i \in \CH_{n - i}(X)$, $i = 0, \ldots, r$
tels que $c_0 = [X]$, et
\begin{equation}
\label{chow-equation-chern-classes}
\sum\nolimits_{i = 0}^r
(-1)^i c_1(\mathcal{O}_P(1))^i \cap \pi^*c_{r - i}
= 0.
\end{equation}
\item Avec les notations ci-dessus, nous posons
$c_i(\mathcal{E}) \cap [X] = c_i$
comme élément de $\CH_{n - i}(X)$.
Nous les appelons les {\it classes de Chern de $\mathcal{E}$ sur $X$}.
\item La {\it classe de Chern totale de $\mathcal{E}$ sur $X$}
est la somme
$$
c({\mathcal E}) \cap [X] =
c_0({\mathcal E}) \cap [X]
+ c_1({\mathcal E}) \cap [X] + \ldots
+ c_r({\mathcal E}) \cap [X]
$$
qui est un élément de
$\CH_*(X) = \bigoplus_{k \in \mathbf{Z}} \CH_k(X)$.
\end{enumerate}
\end{definition}

\noindent
Vérifions que cela ne donne pas une notion nouvelle lorsque le
fibré vectoriel est de rang $1$.

\begin{lemma}
\label{chow-lemma-first-chern-class}
Soit $(S, \delta)$ comme dans la Situation \ref{chow-situation-setup}.
Soit $X$ localement de type fini sur $S$.
Supposons $X$ intègre et $n = \dim_\delta(X)$.
Soit $\mathcal{L}$ un $\mathcal{O}_X$-module inversible.
La première classe de Chern de $\mathcal{L}$ sur $X$ de la
Définition \ref{chow-definition-chern-classes}
est égale au diviseur de Weil associé à $\mathcal{L}$
par la Définition \ref{chow-definition-divisor-invertible-sheaf}.
\end{lemma}

\begin{proof}
Dans cette démonstration, nous employons $c_1(\mathcal{L}) \cap [X]$ pour désigner la
construction de la Définition \ref{chow-definition-divisor-invertible-sheaf}.
Comme $\mathcal{L}$ est de rang $1$, nous avons
$\mathbf{P}(\mathcal{L}) = X$ et
$\mathcal{O}_{\mathbf{P}(\mathcal{L})}(1) = \mathcal{L}$
d'après nos normalisations. Ainsi, (\ref{chow-equation-chern-classes})
s'écrit
$$
(-1)^1 c_1(\mathcal{L}) \cap c_0 + (-1)^0 c_1 = 0
$$
Puisque $c_0 = [X]$, nous concluons que $c_1 = c_1(\mathcal{L}) \cap [X]$,
comme voulu.
\end{proof}

\begin{remark}
\label{chow-remark-equation-signs}
Nous pourrions aussi récrire l'équation \ref{chow-equation-chern-classes} sous la forme
\begin{equation}
\label{chow-equation-signs}
\sum\nolimits_{i = 0}^r
c_1(\mathcal{O}_P(-1))^i \cap \pi^*c_{r - i}
= 0.
\end{equation}
mais il nous paraît plus facile de travailler avec le faisceau quotient
tautologique $\mathcal{O}_P(1)$ plutôt qu'avec
son dual.
\end{remark}




\section{Intersection avec les classes de Chern}
\label{chow-section-intersecting-chern-classes}

\noindent
Dans cette section, nous définissons les classes de Chern des fibrés vectoriels sur $X$ comme
classes bivariantes sur $X$, voir le Lemme \ref{chow-lemma-cap-cp-bivariant}
et la discussion qui le suit. Notre construction suit le procédé familier
consistant à définir d'abord l'opération sur les cycles premiers, puis à
sommer. Dans le Lemme \ref{chow-lemma-determine-intersections}, nous montrons
que le résultat est déterminé par la formule usuelle sur le fibré
projectif associé. Ensuite, nous montrons que l'intersection avec les classes de Chern
passe à l'équivalence rationnelle, commute à l'image directe propre,
commute à l'image inverse par un morphisme plat et commute aux homomorphismes de Gysin pour les
inclusions de diviseurs de Cartier effectifs. On aurait pu éviter ces lemmes
en utilisant directement la caractérisation du
Lemme \ref{chow-lemma-determine-intersections} et le
Lemme \ref{chow-lemma-push-proper-bivariant} ; le lecteur qui souhaite
en voir les détails pourra consulter
Groupes de Chow des espaces, Lemme \ref{spaces-chow-lemma-segre-classes}.

\begin{definition}
\label{chow-definition-cap-chern-classes}
Soit $(S, \delta)$ comme dans la Situation \ref{chow-situation-setup}.
Soit $X$ localement de type fini sur $S$.
Soit $\mathcal{E}$ un faisceau localement libre de type fini de rang $r$ sur $X$.
Nous définissons, pour tout entier $k$ et tout $0 \leq j \leq r$,
une opération
$$
c_j(\mathcal{E}) \cap - : Z_k(X) \to \CH_{k - j}(X)
$$
appelée {\it intersection avec la $j$-ième classe de Chern de $\mathcal{E}$}.
\begin{enumerate}
\item Étant donné un sous-schéma fermé intègre $i : W \to X$ de $\delta$-dimension
$k$, nous définissons
$$
c_j(\mathcal{E}) \cap [W] = i_*(c_j({i^*\mathcal{E}}) \cap [W])
\in
\CH_{k - j}(X)
$$
où $c_j({i^*\mathcal{E}}) \cap [W]$ est défini comme dans la
Définition \ref{chow-definition-chern-classes}.
\item Pour un $k$-cycle quelconque $\alpha = \sum n_i [W_i]$, nous posons
$$
c_j(\mathcal{E}) \cap \alpha = \sum n_i c_j(\mathcal{E}) \cap [W_i]
$$
\end{enumerate}
\end{definition}

\noindent
Si $\mathcal{E}$ est de rang $1$, cela coïncide avec notre
définition précédente (Définition \ref{chow-definition-cap-c1})
d'après le Lemme \ref{chow-lemma-first-chern-class}.

\begin{lemma}
\label{chow-lemma-determine-intersections}
Soit $(S, \delta)$ comme dans la Situation \ref{chow-situation-setup}.
Soit $X$ localement de type fini sur $S$.
Soit $\mathcal{E}$ un faisceau localement libre de type fini de rang $r$ sur $X$.
Soit $(\pi : P \to X, \mathcal{O}_P(1))$ le fibré projectif
associé à $\mathcal{E}$.
Pour $\alpha \in Z_k(X)$, les éléments
$c_j(\mathcal{E}) \cap \alpha$ sont les uniques éléments
$\alpha_j$ de $\CH_{k - j}(X)$
tels que $\alpha_0 = \alpha$ et que
$$
\sum\nolimits_{i = 0}^r
(-1)^i c_1(\mathcal{O}_P(1))^i \cap
\pi^*(\alpha_{r - i}) = 0
$$
dans le groupe de Chow de $P$.
\end{lemma}

\begin{proof}
L'unicité de $\alpha_0, \ldots, \alpha_r$ tels que
$\alpha_0 = \alpha$ et que
l'équation affichée soit satisfaite résulte de
la formule du fibré en espaces projectifs,
Lemme \ref{chow-lemma-chow-ring-projective-bundle}.
Par définition, l'identité est satisfaite pour $\alpha = [W]$, où $W$
est un sous-schéma fermé intègre de $X$.
Pour un $k$-cycle quelconque $\alpha$ sur $X$, écrivons
$\alpha = \sum n_a[W_a]$ avec $n_a \not = 0$, et
$i_a : W_a \to X$ des sous-schémas fermés intègres deux à deux distincts.
Alors la famille $\{W_a\}$ est localement finie sur $X$.
Posons $P_a = \pi^{-1}(W_a) = \mathbf{P}(\mathcal{E}|_{W_a})$.
Notons $i'_a : P_a \to P$ les immersions fermées correspondantes.
Considérons le diagramme de produit fibré
$$
\xymatrix{
P' \ar@{=}[r] \ar[d]_{\pi'} &
\coprod P_a \ar[d]_{\coprod \pi_a} \ar[r]_{\coprod i'_a} &
P \ar[d]^\pi \\
X' \ar@{=}[r] &
\coprod W_a \ar[r]^{\coprod i_a} &
X
}
$$
Le morphisme $p : X' \to X$ est propre. De plus,
$\pi' : P' \to X'$, muni du faisceau inversible
$\mathcal{O}_{P'}(1) = \coprod \mathcal{O}_{P_a}(1)$,
qui est également l'image inverse de $\mathcal{O}_P(1)$,
est le fibré projectif associé à
$\mathcal{E}' = p^*\mathcal{E}$. Par définition,
$$
c_j(\mathcal{E}) \cap [\alpha]
=
\sum i_{a, *}(c_j(\mathcal{E}|_{W_a}) \cap [W_a]).
$$
Écrivons $\beta_{a, j} = c_j(\mathcal{E}|_{W_a}) \cap [W_a]$,
qui est un élément de $\CH_{k - j}(W_a)$. Nous avons
$$
\sum\nolimits_{i = 0}^r
(-1)^i c_1(\mathcal{O}_{P_a}(1))^i \cap \pi_a^*(\beta_{a, r - i}) = 0
$$
pour tout $a$, par définition. Il est donc clair que nous avons
$$
\sum\nolimits_{i = 0}^r
(-1)^i c_1(\mathcal{O}_{P'}(1))^i \cap (\pi')^*(\beta_{r - i}) = 0
$$
avec $\beta_j = \sum n_a\beta_{a, j} \in \CH_{k - j}(X')$. Notons
$p' : P' \to P$ le morphisme $\coprod i'_a$.
Nous avons $\pi^*p_*\beta_j = p'_*(\pi')^*\beta_j$
d'après le Lemme \ref{chow-lemma-flat-pullback-proper-pushforward}.
Par la formule de projection du Lemme \ref{chow-lemma-pushforward-cap-c1},
nous concluons que
$$
\sum\nolimits_{i = 0}^r
(-1)^i c_1(\mathcal{O}_P(1))^i \cap \pi^*(p_*\beta_j) = 0
$$
Puisque $p_*\beta_j$ est un représentant de $c_j(\mathcal{E}) \cap \alpha$,
le résultat s'ensuit.
\end{proof}

\noindent
Nous utiliserons systématiquement cette caractérisation des classes de Chern
pour en démontrer bien d'autres propriétés.

\begin{lemma}
\label{chow-lemma-cap-chern-class-factors-rational-equivalence}
Soit $(S, \delta)$ comme dans la Situation \ref{chow-situation-setup}.
Soit $X$ localement de type fini sur $S$.
Soit $\mathcal{E}$ un faisceau localement libre de type fini de rang $r$ sur $X$.
Si $\alpha \sim_{rat} \beta$ sont des $k$-cycles rationnellement équivalents
sur $X$, alors $c_j(\mathcal{E}) \cap \alpha = c_j(\mathcal{E}) \cap \beta$
dans $\CH_{k - j}(X)$.
\end{lemma}

\begin{proof}
D'après le Lemme \ref{chow-lemma-determine-intersections}, les éléments
$\alpha_j = c_j(\mathcal{E}) \cap \alpha$, $j \geq 1$, et
$\beta_j = c_j(\mathcal{E}) \cap \beta$, $j \geq 1$, sont uniquement déterminés
par la {\it même} équation dans le groupe de Chow du fibré
projectif associé à $\mathcal{E}$. (Cela repose bien entendu sur le fait que
l'image inverse par un morphisme plat est compatible avec l'équivalence rationnelle, voir le
Lemme \ref{chow-lemma-flat-pullback-rational-equivalence}.) Ils sont donc égaux.
\end{proof}

\noindent
Autrement dit, l'intersection avec les classes de Chern des
faisceaux localement libres de type fini passe à l'équivalence rationnelle
et donne des applications
$$
c_j(\mathcal{E}) \cap - : \CH_k(X) \to \CH_{k - j}(X).
$$
Notre tâche suivante consiste à montrer que les classes de Chern sont des classes bivariantes, voir
la Définition \ref{chow-definition-bivariant-class}.

\begin{lemma}
\label{chow-lemma-pushforward-cap-cj}
Soit $(S, \delta)$ comme dans la Situation \ref{chow-situation-setup}.
Soient $X$, $Y$ localement de type fini sur $S$.
Soit $\mathcal{E}$ un faisceau localement libre de type fini de rang $r$ sur $X$.
Soit $p : X \to Y$ un morphisme propre.
Soit $\alpha$ un $k$-cycle sur $X$.
Soit $\mathcal{E}$ un faisceau localement libre de type fini sur $Y$.
Alors
$$
p_*(c_j(p^*\mathcal{E}) \cap \alpha) = c_j(\mathcal{E}) \cap p_*\alpha
$$
\end{lemma}

\begin{proof}
Soit $(\pi : P \to Y, \mathcal{O}_P(1))$ le fibré projectif associé
à $\mathcal{E}$. Alors $P_X = X \times_Y P$ est le fibré projectif associé
à $p^*\mathcal{E}$ et $\mathcal{O}_{P_X}(1)$ est l'image inverse de
$\mathcal{O}_P(1)$. Écrivons $\alpha_j = c_j(p^*\mathcal{E}) \cap \alpha$, de sorte que
$\alpha_0 = \alpha$. D'après le Lemme \ref{chow-lemma-determine-intersections}, nous avons
$$
\sum\nolimits_{i = 0}^r
(-1)^i c_1(\mathcal{O}_P(1))^i \cap
\pi_X^*(\alpha_{r - i}) = 0
$$
dans le groupe de Chow de $P_X$. Considérons le diagramme de produit fibré
$$
\xymatrix{
P_X \ar[r]_-{p'} \ar[d]_{\pi_X} & P \ar[d]^\pi \\
X \ar[r]^p & Y
}
$$
Appliquons l'image directe propre $p'_*$,
(Lemme \ref{chow-lemma-proper-pushforward-rational-equivalence})
à l'égalité affichée ci-dessus. En utilisant les
Lemmes \ref{chow-lemma-pushforward-cap-c1} et
\ref{chow-lemma-flat-pullback-proper-pushforward}, nous obtenons
$$
\sum\nolimits_{i = 0}^r
(-1)^i c_1(\mathcal{O}_P(1))^i \cap
\pi^*(p_*\alpha_{r - i}) = 0
$$
dans le groupe de Chow de $P$. Par la caractérisation du
Lemme \ref{chow-lemma-determine-intersections}, nous concluons.
\end{proof}

\begin{lemma}
\label{chow-lemma-flat-pullback-cap-cj}
Soit $(S, \delta)$ comme dans la Situation \ref{chow-situation-setup}.
Soient $X$, $Y$ localement de type fini sur $S$.
Soit $\mathcal{E}$ un faisceau localement libre de type fini de rang $r$ sur $Y$.
Soit $f : X \to Y$ un morphisme plat de dimension relative $r$.
Soit $\alpha$ un $k$-cycle sur $Y$.
Alors
$$
f^*(c_j(\mathcal{E}) \cap \alpha) = c_j(f^*\mathcal{E}) \cap f^*\alpha
$$
\end{lemma}

\begin{proof}
Écrivons $\alpha_j = c_j(\mathcal{E}) \cap \alpha$, de sorte que $\alpha_0 = \alpha$.
D'après le Lemme \ref{chow-lemma-determine-intersections}, nous avons
$$
\sum\nolimits_{i = 0}^r
(-1)^i c_1(\mathcal{O}_P(1))^i \cap
\pi^*(\alpha_{r - i}) = 0
$$
dans le groupe de Chow du fibré projectif
$(\pi : P \to Y, \mathcal{O}_P(1))$
associé à $\mathcal{E}$. Considérons le diagramme de produit fibré
$$
\xymatrix{
P_X = \mathbf{P}(f^*\mathcal{E}) \ar[r]_-{f'} \ar[d]_{\pi_X} &
P \ar[d]^\pi \\
X \ar[r]^f & Y
}
$$
Notons que $\mathcal{O}_{P_X}(1)$ est l'image inverse de $\mathcal{O}_P(1)$.
Appliquons l'image inverse plate $(f')^*$,
(Lemme \ref{chow-lemma-flat-pullback-rational-equivalence}) à l'équation
affichée ci-dessus. D'après les Lemmes \ref{chow-lemma-flat-pullback-cap-c1} et
\ref{chow-lemma-compose-flat-pullback}, nous voyons que
$$
\sum\nolimits_{i = 0}^r
(-1)^i c_1(\mathcal{O}_{P_X}(1))^i \cap
\pi_X^*(f^*\alpha_{r - i}) = 0
$$
dans le groupe de Chow de $P_X$. Par la caractérisation du
Lemme \ref{chow-lemma-determine-intersections}, nous concluons.
\end{proof}

\begin{lemma}
\label{chow-lemma-cap-chern-class-commutes-with-gysin}
Soit $(S, \delta)$ comme dans la Situation \ref{chow-situation-setup}.
Soit $X$ localement de type fini sur $S$.
Soit $\mathcal{E}$ un faisceau localement libre de type fini de rang $r$ sur $X$.
Soit $(\mathcal{L}, s, i : D \to X)$ comme dans la
Définition \ref{chow-definition-gysin-homomorphism}.
Alors $c_j(\mathcal{E}|_D) \cap i^*\alpha = i^*(c_j(\mathcal{E}) \cap \alpha)$
pour tout $\alpha \in \CH_k(X)$.
\end{lemma}

\begin{proof}
Écrivons $\alpha_j = c_j(\mathcal{E}) \cap \alpha$, de sorte que $\alpha_0 = \alpha$.
D'après le Lemme \ref{chow-lemma-determine-intersections}, nous avons
$$
\sum\nolimits_{i = 0}^r
(-1)^i c_1(\mathcal{O}_P(1))^i \cap
\pi^*(\alpha_{r - i}) = 0
$$
dans le groupe de Chow du fibré projectif
$(\pi : P \to X, \mathcal{O}_P(1))$
associé à $\mathcal{E}$. Considérons le diagramme de produit fibré
$$
\xymatrix{
P_D = \mathbf{P}(\mathcal{E}|_D) \ar[r]_-{i'} \ar[d]_{\pi_D} &
P \ar[d]^\pi \\
D \ar[r]^i & X
}
$$
Notons que $\mathcal{O}_{P_D}(1)$ est l'image inverse de $\mathcal{O}_P(1)$.
Appliquons l'homomorphisme de Gysin $(i')^*$ (Lemme \ref{chow-lemma-gysin-factors}) à
l'équation affichée ci-dessus.
En appliquant les Lemmes \ref{chow-lemma-gysin-commutes-cap-c1} et
\ref{chow-lemma-gysin-flat-pullback}, nous obtenons
$$
\sum\nolimits_{i = 0}^r
(-1)^i c_1(\mathcal{O}_{P_D}(1))^i \cap
\pi_D^*(i^*\alpha_{r - i}) = 0
$$
dans le groupe de Chow de $P_D$.
Par la caractérisation du Lemme \ref{chow-lemma-determine-intersections},
nous concluons.
\end{proof}

\noindent
Nous disposons maintenant de suffisamment d'éléments pour démontrer que
l'intersection avec les classes de Chern définit une classe bivariante.

\begin{lemma}
\label{chow-lemma-cap-cp-bivariant}
Soit $(S, \delta)$ comme dans la Situation \ref{chow-situation-setup}.
Soit $X$ localement de type fini sur $S$.
Soit $\mathcal{E}$ un $\mathcal{O}_X$-module localement libre
de rang $r$. Soit $0 \leq p \leq r$.
Alors la règle qui, à $f : X' \to X$, associe
$c_p(f^*\mathcal{E}) \cap - : \CH_k(X') \to \CH_{k - p}(X')$
est une classe bivariante de degré $p$.
\end{lemma}

\begin{proof}
Cela résulte immédiatement des Lemmes
\ref{chow-lemma-cap-chern-class-factors-rational-equivalence},
\ref{chow-lemma-pushforward-cap-cj},
\ref{chow-lemma-flat-pullback-cap-cj}, et
\ref{chow-lemma-cap-chern-class-commutes-with-gysin},
ainsi que de la Définition \ref{chow-definition-bivariant-class}.
\end{proof}

\noindent
Ce lemme nous permet de définir comme suit les classes de Chern d'un module
localement libre de type fini.

\begin{definition}
\label{chow-definition-chern-classes-final}
Soit $(S, \delta)$ comme dans la Situation \ref{chow-situation-setup}.
Soit $X$ localement de type fini sur $S$.
Soit $\mathcal{E}$ un $\mathcal{O}_X$-module localement libre
de rang $r$. Pour $i = 0, \ldots, r$, la {\it $i$-ième classe de Chern}
de $\mathcal{E}$ est la classe bivariante
$c_i(\mathcal{E}) \in A^i(X)$ de degré $i$
construite dans le Lemme \ref{chow-lemma-cap-cp-bivariant}. La
{\it classe de Chern totale} de $\mathcal{E}$ est la somme formelle
$$
c(\mathcal{E}) = 
c_0(\mathcal{E}) + c_1(\mathcal{E}) + \ldots + c_r(\mathcal{E})
$$
considérée comme une classe bivariante non homogène sur $X$.
\end{definition}

\noindent
D'après la remarque qui suit la Définition \ref{chow-definition-cap-chern-classes},
si $\mathcal{E}$ est inversible, cette définition coïncide avec la
Définition \ref{chow-definition-first-chern-class}.
Nous voyons ensuite que les classes de Chern appartiennent au centre de l'anneau de
cohomologie bivariante de Chow $A^*(X)$.

\begin{lemma}
\label{chow-lemma-cap-commutative-chern}
Soit $(S, \delta)$ comme dans la Situation \ref{chow-situation-setup}.
Soit $X$ localement de type fini sur $S$.
Soit $\mathcal{E}$ un $\mathcal{O}_X$-module localement libre de rang $r$.
Alors
\begin{enumerate}
\item $c_j(\mathcal{E}) \in A^j(X)$ appartient au centre de $A^*(X)$, et
\item si $f : X' \to X$ est localement de type fini et $c \in A^*(X' \to X)$,
alors $c \circ c_j(\mathcal{E}) = c_j(f^*\mathcal{E}) \circ c$.
\end{enumerate}
En particulier, si $\mathcal{F}$ est un second
$\mathcal{O}_X$-module localement libre sur $X$ de rang $s$, alors
$$
c_i(\mathcal{E}) \cap c_j(\mathcal{F}) \cap \alpha
=
c_j(\mathcal{F}) \cap c_i(\mathcal{E}) \cap \alpha
$$
comme éléments de $\CH_{k - i - j}(X)$ pour tout $\alpha \in \CH_k(X)$.
\end{lemma}

\begin{proof}
Il est immédiat que (2) implique (1).
Soit $\alpha \in \CH_k(X)$. Écrivons $\alpha_j = c_j(\mathcal{E}) \cap \alpha$, de sorte que
$\alpha_0 = \alpha$. D'après le Lemme \ref{chow-lemma-determine-intersections}, nous avons
$$
\sum\nolimits_{i = 0}^r
(-1)^i c_1(\mathcal{O}_P(1))^i \cap
\pi^*(\alpha_{r - i}) = 0
$$
dans le groupe de Chow du fibré projectif
$(\pi : P \to Y, \mathcal{O}_P(1))$
associé à $\mathcal{E}$. Notons $\pi' : P' \to X'$ le changement de base
de $\pi$ par $f$. En utilisant le Lemme \ref{chow-lemma-c1-center} et
les propriétés des classes bivariantes, nous obtenons
\begin{align*}
0 & = c \cap \left(\sum\nolimits_{i = 0}^r
(-1)^i c_1(\mathcal{O}_P(1))^i \cap
\pi^*(\alpha_{r - i})\right) \\
& =
\sum\nolimits_{i = 0}^r
(-1)^i c_1(\mathcal{O}_{P'}(1))^i \cap
(\pi')^*(c \cap \alpha_{r - i})
\end{align*}
dans le groupe de Chow de $P'$ (calcul omis).
Nous voyons donc que $c \cap \alpha_j$ est
égal à $c_j(f^*\mathcal{E}) \cap (c \cap \alpha)$ par la caractérisation
du Lemme \ref{chow-lemma-determine-intersections}.
Cela démontre le lemme.
\end{proof}

\begin{remark}
\label{chow-remark-extend-to-finite-locally-free}
Soit $(S, \delta)$ comme dans la Situation \ref{chow-situation-setup}.
Soit $X$ localement de type fini sur $S$.
Soit $\mathcal{E}$ un $\mathcal{O}_X$-module localement libre de type fini.
Si le rang de $\mathcal{E}$ n'est pas constant, nous pouvons
encore définir les classes de Chern de $\mathcal{E}$. En effet, dans ce
cas, nous pouvons écrire
$$
X = X_0 \amalg X_1 \amalg X_2 \amalg \ldots
$$
où $X_r \subset X$ est le sous-espace ouvert et fermé sur lequel
le rang de $\mathcal{E}$ est $r$. D'après le
Lemme \ref{chow-lemma-disjoint-decomposition-bivariant},
nous avons $A^p(X) = \prod A^p(X_r)$.
Nous pouvons donc définir $c_p(\mathcal{E})$ comme le
produit des classes $c_p(\mathcal{E}|_{X_r})$ dans $A^p(X_r)$.
Explicitement, si $X' \to X$ est un morphisme localement de type fini,
nous obtenons par image inverse une décomposition correspondante de $X'$,
et nous trouvons que
$$
\CH_*(X') = \prod\nolimits_{r \geq 0} \CH_*(X'_r)
$$
d'après nos définitions. Alors $c_p(\mathcal{E}) \in A^p(X)$
est la classe bivariante qui préserve ces décompositions
en produits directs et agit par les opérations déjà définies
$c_i(\mathcal{E}|_{X_r}) \cap -$
sur les facteurs. Observons que, dans ce cadre, il peut arriver
que $c_p(\mathcal{E})$ soit non nul pour une infinité de $p$.
Il s'ensuit que la classe de Chern totale est un élément
$$
c(\mathcal{E}) =
c_0(\mathcal{E}) + c_1(\mathcal{E}) + c_2(\mathcal{E}) + \ldots
\in A^*(X)^\wedge
$$
de l'anneau de cohomologie bivariante complété, voir la
Remarque \ref{chow-remark-completion-bivariant}.
Dans ce cadre, nous définissons le « rang » de $\mathcal{E}$
comme l'élément $r(\mathcal{E}) \in A^0(X)$,
à savoir l'opération bivariante qui envoie $(\alpha_r) \in \prod \CH_*(X'_r)$
sur $(r\alpha_r) \in \prod \CH_*(X'_r)$.
Notons qu'il reste vrai que $c_p(\mathcal{E})$ et $r(\mathcal{E})$
appartiennent au centre de $A^*(X)$.
\end{remark}

\begin{remark}
\label{chow-remark-top-chern-class}
Soit $(S, \delta)$ comme dans la Situation \ref{chow-situation-setup}. Soit $X$
localement de type fini sur $S$. Soit $\mathcal{E}$ un
$\mathcal{O}_X$-module localement libre de type fini. En général,
nous écrivons $X = \coprod X_r$ comme dans la
Remarque \ref{chow-remark-extend-to-finite-locally-free}.
Si seul un nombre fini des $X_r$ sont non vides, alors
nous pouvons poser
$$
c_{top}(\mathcal{E}) = \sum\nolimits_r c_r(\mathcal{E}|_{X_r})
\in A^*(X) = \bigoplus A^*(X_r)
$$
où l'égalité est le Lemme \ref{chow-lemma-disjoint-decomposition-bivariant}.
Si une infinité des $X_r$ sont non vides, nous emploierons la même
notation pour désigner
$$
c_{top}(\mathcal{E}) = \prod c_r(\mathcal{E}|_{X_r})
\in \prod A^r(X_r) \subset A^*(X)^\wedge
$$
voir la Remarque \ref{chow-remark-completion-bivariant} pour les notations.
\end{remark}











\section{Relations polynomiales entre classes de Chern}
\label{chow-section-relations-chern-classes}

\noindent
Soit $(S, \delta)$ comme dans la Situation \ref{chow-situation-setup}. Soit $X$ localement
de type fini sur $S$. Les $\mathcal{E}_i$ forment une famille finie de faisceaux
localement libres de type fini sur $X$. D'après le Lemme \ref{chow-lemma-cap-commutative-chern},
nous voyons que les classes de Chern
$$
c_j(\mathcal{E}_i) \in A^*(X)
$$
engendrent une $\mathbf{Z}$-sous-algèbre commutative (et même centrale) de
l'algèbre de cohomologie de Chow $A^*(X)$.
Nous pouvons donc préciser ce que signifie qu'un polynôme en ces classes de Chern
est nul, ou que deux polynômes sont égaux. Par exemple, dire que
$c_1(\mathcal{E}_1)^5 + c_2(\mathcal{E}_2)c_3(\mathcal{E}_3) = 0$
signifie que les opérations
$$
\CH_k(Y) \longrightarrow \CH_{k - 5}(Y), \quad
\alpha \longmapsto
c_1(\mathcal{E}_1)^5 \cap \alpha +
c_2(\mathcal{E}_2) \cap c_3(\mathcal{E}_3) \cap \alpha
$$
sont nulles pour tout morphisme $f : Y \to X$ localement de type fini.
D'après le Lemme \ref{chow-lemma-bivariant-zero},
cela équivaut à exiger que, pour tout morphisme
$f : Y \to X$  où $Y$ est un schéma intègre
localement de type fini sur $S$, le cycle
$$
c_1(\mathcal{E}_1)^5 \cap [Y] +
c_2(\mathcal{E}_2) \cap c_3(\mathcal{E}_3) \cap [Y]
$$
soit nul dans $\CH_{\dim(Y) - 5}(Y)$.

\medskip\noindent
Un exemple particulier est la relation
$$
c_1(\mathcal{L} \otimes_{\mathcal{O}_X} \mathcal{N})
=
c_1(\mathcal{L}) + c_1(\mathcal{N})
$$
démontrée dans le Lemme \ref{chow-lemma-c1-cap-additive}.
Plus généralement, voici ce qui se passe lorsque l'on tensorise un
faisceau localement libre quelconque par un faisceau inversible.

\begin{lemma}
\label{chow-lemma-chern-classes-E-tensor-L}
Soit $(S, \delta)$ comme dans la Situation \ref{chow-situation-setup}.
Soit $X$ localement de type fini sur $S$.
Soit $\mathcal{E}$ un faisceau localement libre de type fini de
rang $r$ sur $X$. Soit $\mathcal{L}$ un faisceau inversible
sur $X$. Alors nous avons
\begin{equation}
\label{chow-equation-twist}
c_i({\mathcal E} \otimes {\mathcal L})
=
\sum\nolimits_{j = 0}^i
\binom{r - i + j}{j} c_{i - j}({\mathcal E}) c_1({\mathcal L})^j
\end{equation}
dans $A^*(X)$.
\end{lemma}

\begin{proof}
Cela doit être vrai pour tout triplet $(X, \mathcal{E}, \mathcal{L})$.
En particulier, cela doit être vrai lorsque $X$ est intègre et, d'après le
Lemme \ref{chow-lemma-bivariant-zero},
il suffit de démontrer
que cela vaut après intersection avec $[X]$ pour un tel $X$. Supposons donc
que $X$ est intègre. Soit $(\pi : P \to X, \mathcal{O}_P(1))$,
resp.\ $(\pi' : P' \to X, \mathcal{O}_{P'}(1))$, le
fibré en espaces projectifs associé à $\mathcal{E}$,
resp.\ à $\mathcal{E} \otimes \mathcal{L}$. Considérons le morphisme canonique
$$
\xymatrix{
P \ar[rd]_\pi \ar[rr]_g & & P' \ar[ld]^{\pi'} \\
& X &
}
$$
voir Constructions, Lemme \ref{constructions-lemma-twisting-and-proj}.
Il possède la propriété
$g^*\mathcal{O}_{P'}(1)
= \mathcal{O}_P(1) \otimes \pi^* {\mathcal L}$.
Cela signifie que nous avons
$$
\sum\nolimits_{i = 0}^r
(-1)^i
(\xi + x)^i \cap \pi^*(c_{r - i}(\mathcal{E} \otimes \mathcal{L}) \cap [X])
=
0
$$
dans $\CH_*(P)$, où $\xi$ représente
$c_1(\mathcal{O}_P(1))$ et $x$
représente $c_1(\pi^*\mathcal{L})$. Un calcul algébrique élémentaire montre que cela
équivaut à
$$
\sum\nolimits_{i = 0}^r
(-1)^i \xi^i \left(
\sum\nolimits_{j = i}^r
(-1)^{j - i}
\binom{j}{i}
x^{j - i} \cap
\pi^*(c_{r - j}(\mathcal{E} \otimes \mathcal{L}) \cap [X])
\right)
=
0
$$
En comparant avec
l'Équation (\ref{chow-equation-chern-classes}), il en résulte que
$$
c_{r - i}(\mathcal{E}) \cap [X] =
\sum\nolimits_{j = i}^r
\binom{j}{i}
(-c_1(\mathcal{L}))^{j - i} \cap
c_{r - j}(\mathcal{E} \otimes \mathcal{L}) \cap [X]
$$
En remaniant ceci (en éliminant les signes moins et en renumérotant), nous obtenons
la relation voulue.
\end{proof}

\noindent
Quelques cas particuliers de (\ref{chow-equation-twist}) sont
\begin{align*}
c_1(\mathcal{E} \otimes \mathcal{L})
& =
c_1(\mathcal{E}) +
r c_1(\mathcal{L}) \\
c_2(\mathcal{E} \otimes \mathcal{L})
& =
c_2(\mathcal{E}) +
(r - 1) c_1(\mathcal{E}) c_1(\mathcal{L}) +
\binom{r}{2} c_1(\mathcal{L})^2 \\
c_3(\mathcal{E} \otimes \mathcal{L})
& =
c_3(\mathcal{E}) +
(r - 2) c_2(\mathcal{E})c_1(\mathcal{L}) +
\binom{r - 1}{2} c_1(\mathcal{E})c_1(\mathcal{L})^2 +
\binom{r}{3} c_1(\mathcal{L})^3
\end{align*}








\section{Additivité des classes de Chern}
\label{chow-section-additivity-chern-classes}

\noindent
Tous les lemmes préliminaires découlent trivialement du
résultat final.

\begin{lemma}
\label{chow-lemma-get-rid-of-trivial-subbundle}
Soit $(S, \delta)$ comme dans la Situation \ref{chow-situation-setup}.
Soit $X$ localement de type fini sur $S$.
Soient $\mathcal{E}$, $\mathcal{F}$ des faisceaux localement libres de type fini
sur $X$, de rangs $r$, $r - 1$, qui s'insèrent dans une suite
exacte courte
$$
0 \to \mathcal{O}_X \to \mathcal{E} \to \mathcal{F} \to 0
$$
Alors nous avons
$$
c_r(\mathcal{E}) = 0, \quad
c_j(\mathcal{E}) = c_j(\mathcal{F}), \quad j = 0, \ldots, r - 1
$$
dans $A^*(X)$.
\end{lemma}

\begin{proof}
D'après le Lemme \ref{chow-lemma-bivariant-zero},
il suffit de montrer que, si $X$ est intègre,
alors $c_j(\mathcal{E}) \cap [X] = c_j(\mathcal{F}) \cap [X]$.
Soit $(\pi : P \to X, \mathcal{O}_P(1))$,
resp.\ $(\pi' : P' \to X, \mathcal{O}_{P'}(1))$, le
fibré en espaces projectifs associé à $\mathcal{E}$, resp.\ à $\mathcal{F}$.
La surjection $\mathcal{E} \to \mathcal{F}$ donne lieu
à une immersion fermée
$$
i : P' \longrightarrow P
$$
sur $X$. De plus, l'élément
$1 \in \Gamma(X, \mathcal{O}_X) \subset \Gamma(X, \mathcal{E})$
donne lieu à une section globale $s \in \Gamma(P, \mathcal{O}_P(1))$
dont le lieu des zéros est exactement $P'$. Ainsi, $P'$ est un diviseur de Cartier effectif
sur $P$ tel que $\mathcal{O}_P(P') \cong \mathcal{O}_P(1)$.
Nous voyons donc que
$$
c_1(\mathcal{O}_P(1)) \cap \pi^*\alpha = i_*((\pi')^*\alpha)
$$
pour toute classe de cycle $\alpha$ sur $X$, d'après le
Lemme \ref{chow-lemma-relative-effective-cartier}.
D'après le Lemme \ref{chow-lemma-determine-intersections}, nous voyons que
$\alpha_j = c_j(\mathcal{F}) \cap [X]$, $j = 0, \ldots, r - 1$,
satisfont
$$
\sum\nolimits_{j = 0}^{r - 1} (-1)^jc_1(\mathcal{O}_{P'}(1))^j
\cap (\pi')^*\alpha_j = 0
$$
En prenant l'image directe dans $P$ et en utilisant la remarque ci-dessus ainsi que le
Lemme \ref{chow-lemma-pushforward-cap-c1}, nous obtenons
$$
\sum\nolimits_{j = 0}^{r - 1}
(-1)^j c_1(\mathcal{O}_P(1))^{j + 1}
\cap \pi^*\alpha_j = 0
$$
Par l'unicité du Lemme \ref{chow-lemma-determine-intersections},
nous concluons que
$c_r(\mathcal{E}) \cap [X] = 0$ et
$c_j(\mathcal{E}) \cap [X] = \alpha_j = c_j(\mathcal{F}) \cap [X]$
pour $j = 0, \ldots, r - 1$. Le lemme est donc vrai.
\end{proof}

\begin{lemma}
\label{chow-lemma-additivity-invertible-subsheaf}
Soit $(S, \delta)$ comme dans la Situation \ref{chow-situation-setup}.
Soit $X$ localement de type fini sur $S$.
Soient $\mathcal{E}$, $\mathcal{F}$ des faisceaux localement libres de type fini
sur $X$, de rangs $r$, $r - 1$, qui s'insèrent dans une suite
exacte courte
$$
0 \to \mathcal{L} \to \mathcal{E} \to \mathcal{F} \to 0
$$
où $\mathcal{L}$ est un faisceau inversible.
Alors
$$
c(\mathcal{E}) = c(\mathcal{L}) c(\mathcal{F})
$$
dans $A^*(X)$.
\end{lemma}

\begin{proof}
Cette relation dit simplement que
$c_i(\mathcal{E}) = c_i(\mathcal{F}) + c_1(\mathcal{L})c_{i - 1}(\mathcal{F})$.
D'après le Lemme \ref{chow-lemma-get-rid-of-trivial-subbundle},
nous avons $c_j(\mathcal{E} \otimes \mathcal{L}^{\otimes -1})
= c_j(\mathcal{F} \otimes \mathcal{L}^{\otimes -1})$ pour
$j = 0, \ldots, r$ étaient nous posons
$c_r(\mathcal{F} \otimes \mathcal{L}^{-1}) = 0$ par convention.
En appliquant le Lemme \ref{chow-lemma-chern-classes-E-tensor-L}, nous déduisons
$$
\sum_{j = 0}^i
\binom{r - i + j}{j} (-1)^j c_{i - j}({\mathcal E}) c_1({\mathcal L})^j
=
\sum_{j = 0}^i
\binom{r - 1 - i + j}{j} (-1)^j c_{i - j}({\mathcal F}) c_1({\mathcal L})^j
$$
En posant
$c_i(\mathcal{E}) = c_i(\mathcal{F}) + c_1(\mathcal{L})c_{i - 1}(\mathcal{F})$,
on obtient une « solution » de cette équation. Le lemme s'ensuit si nous montrons
que c'est la seule solution possible. Nous omettons la vérification.
\end{proof}

\begin{lemma}
\label{chow-lemma-additivity-chern-classes}
Soit $(S, \delta)$ comme dans la Situation \ref{chow-situation-setup}.
Soit $X$ un schéma localement de type fini sur $S$.
Supposons que ${\mathcal E}$ s'insère dans une
suite exacte
$$
0
\to
{\mathcal E}_1
\to
{\mathcal E}
\to
{\mathcal E}_2
\to
0
$$
de faisceaux localement libres de type fini $\mathcal{E}_i$ de rang $r_i$.
Les classes de Chern totales satisfont
$$
c({\mathcal E}) = c({\mathcal E}_1) c({\mathcal E}_2)
$$
dans $A^*(X)$.
\end{lemma}

\begin{proof}
D'après le Lemme \ref{chow-lemma-bivariant-zero}, nous pouvons supposer $X$ intègre
et nous devons montrer l'identité après intersection avec $[X]$.
Raisonnons par récurrence sur $r_1$. Le cas $r_1 = 1$ est le
Lemme \ref{chow-lemma-additivity-invertible-subsheaf}.
Supposons $r_1 > 1$. Soit $(\pi : P \to X, \mathcal{O}_P(1))$
le fibré en espaces projectifs associé à $\mathcal{E}_1$. Notons que
\begin{enumerate}
\item $\pi^* : \CH_*(X) \to \CH_*(P)$ est injective, et
\item $\pi^*\mathcal{E}_1$ s'insère dans une suite exacte courte
$0 \to \mathcal{F} \to \pi^*\mathcal{E}_1 \to \mathcal{L} \to 0$,
où $\mathcal{L}$ est inversible.
\end{enumerate}
La première assertion résulte de la formule du fibré en espaces projectifs
et la seconde de la définition d'un fibré en espaces projectifs.
(En fait, $\mathcal{L} = \mathcal{O}_P(1)$.)
Soit $Q = \pi^*\mathcal{E}/\mathcal{F}$, qui s'insère dans une
suite exacte $0 \to \mathcal{L} \to Q \to \pi^*\mathcal{E}_2 \to 0$.
Par récurrence, nous avons
\begin{eqnarray*}
c(\pi^*\mathcal{E}) \cap [P]
& = &
c(\mathcal{F}) \cap c(\pi^*\mathcal{E}/\mathcal{F}) \cap [P] \\
& = &
c(\mathcal{F}) \cap c(\mathcal{L}) \cap c(\pi^*\mathcal{E}_2) \cap [P] \\
& = &
c(\pi^*\mathcal{E}_1) \cap c(\pi^*\mathcal{E}_2) \cap [P]
\end{eqnarray*}
Puisque $[P] = \pi^*[X]$, nous
concluons par le Lemme \ref{chow-lemma-flat-pullback-cap-cj}.
\end{proof}

\begin{lemma}
\label{chow-lemma-chern-filter-by-linebundles}
Soit $(S, \delta)$ comme dans la Situation \ref{chow-situation-setup}.
Soit $X$ localement de type fini sur $S$.
Soient ${\mathcal L}_i$, $i = 1, \ldots, r$, des
$\mathcal{O}_X$-modules inversibles sur $X$.
Soit $\mathcal{E}$ un module localement libre de rang
$\mathcal{O}_X$-module muni d'une filtration
$$
0 = \mathcal{E}_0 \subset \mathcal{E}_1 \subset \mathcal{E}_2
\subset \ldots \subset \mathcal{E}_r = \mathcal{E}
$$
telle que $\mathcal{E}_i/\mathcal{E}_{i - 1} \cong \mathcal{L}_i$.
Posons $c_1({\mathcal L}_i) = x_i$. Alors
$$
c(\mathcal{E})
=
\prod\nolimits_{i = 1}^r (1 + x_i)
$$
dans $A^*(X)$.
\end{lemma}

\begin{proof}
Appliquer le Lemme \ref{chow-lemma-additivity-invertible-subsheaf} et procéder par récurrence.
\end{proof}




\section{Degrés des zéro-cycles}
\label{chow-section-degree-zero-cycles}

\noindent
Nous commençons par définir le degré d'un zéro-cycle sur un schéma propre sur un corps.
Une méthode consiste à le définir directement comme dans le
Lemme \ref{chow-lemma-spell-out-degree-zero-cycle}, puis à montrer
qu'il est bien défini à l'aide du
Lemme \ref{chow-lemma-curve-principal-divisor}.
Nous le définissons plutôt comme suit.

\begin{definition}
\label{chow-definition-degree-zero-cycle}
Soit $k$ un corps (Exemple \ref{chow-example-field}). Soit $p : X \to \Spec(k)$
propre. Le {\it degré d'un zéro-cycle} sur $X$ est donné par l'image
directe propre
$$
p_* : \CH_0(X) \to \CH_0(\Spec(k))
$$
(Lemme \ref{chow-lemma-proper-pushforward-rational-equivalence})
composée avec l'isomorphisme naturel $\CH_0(\Spec(k)) = \mathbf{Z}$
qui envoie $[\Spec(k)]$ sur $1$. Notation : $\deg(\alpha)$.
\end{definition}

\noindent
Explicitons davantage cette définition.

\begin{lemma}
\label{chow-lemma-spell-out-degree-zero-cycle}
Soit $k$ un corps. Soit $X$ propre sur $k$. Soit $\alpha = \sum n_i[Z_i]$
un élément de $Z_0(X)$. Alors
$$
\deg(\alpha) = \sum n_i\deg(Z_i)
$$
où $\deg(Z_i)$ est le degré de $Z_i \to \Spec(k)$, c'est-à-dire
$\deg(Z_i) = \dim_k \Gamma(Z_i, \mathcal{O}_{Z_i})$.
\end{lemma}

\begin{proof}
C'est la définition de l'image directe propre
(Définition \ref{chow-definition-proper-pushforward}).
\end{proof}

\noindent
Nous établissons ensuite le lien avec les degrés des fibrés vectoriels
sur les schémas propres de dimension $1$ sur des corps, tels qu'ils sont définis dans
Variétés, section \ref{varieties-section-divisors-curves}.

\begin{lemma}
\label{chow-lemma-degree-vector-bundle}
Soit $k$ un corps. Soit $X$ un schéma propre sur $k$ de dimension $\leq 1$.
Soit $\mathcal{E}$ un $\mathcal{O}_X$-module localement libre de type fini de
rang constant. Alors
$$
\deg(\mathcal{E}) = \deg(c_1(\mathcal{E}) \cap [X]_1)
$$
où le membre de gauche est défini dans
Variétés, Définition \ref{varieties-definition-degree-invertible-sheaf}.
\end{lemma}

\begin{proof}
Soient $C_i \subset X$, $i = 1, \ldots, t$, les composantes irréductibles
de dimension $1$, munies de leur structure de schéma induite réduite, et soit $m_i$ la
multiplicité de $C_i$ dans $X$. Alors $[X]_1 = \sum m_i[C_i]$ et
$c_1(\mathcal{E}) \cap [X]_1$ est la somme des images directes des cycles
$m_i c_1(\mathcal{E}|_{C_i}) \cap [C_i]$. Puisque nous avons une décomposition analogue
du degré de $\mathcal{E}$ par
Variétés, Lemme \ref{varieties-lemma-degree-in-terms-of-components},
il suffit de démontrer le lemme lorsque $X$ est une courbe propre sur $k$.

\medskip\noindent
Supposons que $X$ soit une courbe propre sur $k$.
Par Diviseurs, Lemme \ref{divisors-lemma-filter-after-modification},
il existe une modification $f : X' \to X$ telle que $f^*\mathcal{E}$
possède une filtration dont les quotients successifs sont des
$\mathcal{O}_{X'}$-modules inversibles. Puisque $f_*[X']_1 = [X]_1$, nous déduisons
du Lemme \ref{chow-lemma-pushforward-cap-cj} que
$$
\deg(c_1(\mathcal{E}) \cap [X]_1) = \deg(c_1(f^*\mathcal{E}) \cap [X']_1)
$$
Puisque nous avons une relation analogue pour le degré par
Variétés, Lemme \ref{varieties-lemma-degree-birational-pullback},
nous nous ramenons au cas où $\mathcal{E}$ possède une filtration dont les
quotients successifs sont des $\mathcal{O}_X$-modules inversibles.
Dans ce cas, nous pouvons utiliser l'additivité du degré
(Variétés, Lemme \ref{varieties-lemma-degree-additive})
et celle des premières classes de Chern (Lemme \ref{chow-lemma-additivity-chern-classes})
pour nous ramener au cas traité dans le paragraphe suivant.

\medskip\noindent
Supposons que $X$ soit une courbe propre sur $k$ et que $\mathcal{E}$ soit un
$\mathcal{O}_X$-module inversible. Par
Diviseurs, Lemme
\ref{divisors-lemma-quasi-projective-Noetherian-pic-effective-Cartier},
nous voyons que $\mathcal{E}$ est isomorphe à
$\mathcal{O}_X(D) \otimes \mathcal{O}_X(D')^{\otimes -1}$
pour certains diviseurs de Cartier effectifs $D, D'$ sur $X$ (cela utilise aussi
que $X$ est projectif, voir par exemple
Variétés, Lemme \ref{varieties-lemma-dim-1-proper-projective}).
Par additivité du degré relativement au produit tensoriel de faisceaux inversibles
(Variétés, Lemme \ref{varieties-lemma-degree-tensor-product})
et additivité de $c_1$ relativement au produit tensoriel de faisceaux inversibles
(Lemme \ref{chow-lemma-c1-cap-additive} ou \ref{chow-lemma-chern-classes-E-tensor-L}),
nous nous ramenons au cas $\mathcal{E} = \mathcal{O}_X(D)$.
Dans ce cas, le membre de gauche donne $\deg(D)$
(Variétés, Lemme \ref{varieties-lemma-degree-effective-Cartier-divisor})
et le membre de droite donne $\deg([D]_0)$ par le
Lemme \ref{chow-lemma-geometric-cap}.
Puisque
$$
[D]_0 = \sum\nolimits_{x \in D}
\text{longueur}_{\mathcal{O}_{X, x}}(\mathcal{O}_{D, x}) [x] =
\sum\nolimits_{x \in D}
\text{longueur}_{\mathcal{O}_{D, x}}(\mathcal{O}_{D, x}) [x]
$$
par définition, nous voyons que
$$
\deg([D]_0) = \sum\nolimits_{x \in D}
\text{longueur}_{\mathcal{O}_{D, x}}(\mathcal{O}_{D, x}) [\kappa(x) : k] =
\dim_k \Gamma(D, \mathcal{O}_D) = \deg(D)
$$
L'avant-dernière égalité résulte de
Algèbre, Lemme \ref{algebra-lemma-pushdown-module},
en utilisant que $D$ est affine.
\end{proof}

\noindent
Enfin, nous pouvons tout relier aux intersections numériques
définies dans Variétés, section \ref{varieties-section-num}.

\begin{lemma}
\label{chow-lemma-degrees-and-numerical-intersections}
Soit $k$ un corps. Soit $X$ un schéma propre sur $k$.
Soit $Z \subset X$ un sous-schéma fermé de dimension $d$.
Soient $\mathcal{L}_1, \ldots, \mathcal{L}_d$ des
$\mathcal{O}_X$-modules inversibles. Alors
$$
(\mathcal{L}_1 \cdots \mathcal{L}_d \cdot Z) =
\deg(
c_1(\mathcal{L}_1) \cap \ldots \cap c_1(\mathcal{L}_d) \cap [Z]_d)
$$
où le membre de gauche est défini dans
Variétés, Définition \ref{varieties-definition-intersection-number}.
En particulier,
$$
\deg_\mathcal{L}(Z) = \deg(c_1(\mathcal{L})^d \cap [Z]_d)
$$
si $\mathcal{L}$ est un $\mathcal{O}_X$-module inversible ample.
\end{lemma}

\begin{proof}
Nous allons le démontrer par récurrence sur $d$. Si $d = 0$, le résultat est
vrai par Variétés, Lemme \ref{varieties-lemma-chi-tensor-finite}.
Supposons $d > 0$.

\medskip\noindent
Soient $Z_i \subset Z$, $i = 1, \ldots, t$, les composantes irréductibles
de dimension $d$, munies de leur structure de schéma induite réduite, et soit $m_i$ la
multiplicité de $Z_i$ dans $Z$. Alors $[Z]_d = \sum m_i[Z_i]$ et
$c_1(\mathcal{L}_1) \cap \ldots \cap c_1(\mathcal{L}_d) \cap [Z]_d$
est la somme des cycles
$m_i c_1(\mathcal{L}_1) \cap \ldots \cap c_1(\mathcal{L}_d) \cap [Z_i]$.
Puisque nous avons une décomposition analogue pour
$(\mathcal{L}_1 \cdots \mathcal{L}_d \cdot Z)$ par
Variétés, Lemme \ref{varieties-lemma-numerical-polynomial-leading-term},
il suffit de démontrer le lemme dans le cas où $Z = X$
est une variété propre de dimension $d$ sur $k$.

\medskip\noindent
Par le lemme de Chow, il existe un morphisme propre birationnel $f : Y \to X$
tel que $Y$ soit H-projectif sur $k$. Voir Cohomologie des schémas, Lemme
\ref{coherent-lemma-chow-Noetherian} et Remarque
\ref{coherent-remark-chow-Noetherian}. Alors
$$
(f^*\mathcal{L}_1 \cdots f^*\mathcal{L}_d \cdot Y) =
(\mathcal{L}_1 \cdots \mathcal{L}_d \cdot X)
$$
par Variétés, Lemme \ref{varieties-lemma-intersection-number-and-pullback},
et nous avons
$$
f_*(c_1(f^*\mathcal{L}_1) \cap \ldots \cap c_1(f^*\mathcal{L}_d) \cap [Y]) =
c_1(\mathcal{L}_1) \cap \ldots \cap c_1(\mathcal{L}_d) \cap [X]
$$
par le Lemme \ref{chow-lemma-pushforward-cap-c1}. Ainsi, nous pouvons remplacer $X$ par $Y$
et supposer que $X$ est projectif sur $k$.

\medskip\noindent
Si $X$ est une variété projective propre de dimension $d$, nous pouvons
écrire $\mathcal{L}_1 = \mathcal{O}_X(D) \otimes \mathcal{O}_X(D')^{\otimes -1}$
pour certains diviseurs de Cartier effectifs $D, D' \subset X$
par Diviseurs, Lemme
\ref{divisors-lemma-quasi-projective-Noetherian-pic-effective-Cartier}.
Par additivité des deux membres de l'équation
(Variétés, Lemme \ref{varieties-lemma-intersection-number-additive} et
Lemme \ref{chow-lemma-c1-cap-additive}),
nous nous ramenons au cas $\mathcal{L}_1 = \mathcal{O}_X(D)$ pour un
diviseur de Cartier effectif $D$.
Par Variétés, Lemme
\ref{varieties-lemma-numerical-intersection-effective-Cartier-divisor},
nous avons
$$
(\mathcal{L}_1 \cdots \mathcal{L}_d \cdot X) =
(\mathcal{L}_2 \cdots \mathcal{L}_d \cdot D)
$$
et, par le Lemme \ref{chow-lemma-geometric-cap}, nous avons
$$
c_1(\mathcal{L}_1) \cap \ldots \cap c_1(\mathcal{L}_d) \cap [X] =
c_1(\mathcal{L}_2) \cap \ldots \cap c_1(\mathcal{L}_d) \cap [D]_{d - 1}
$$
Nous obtenons ainsi le résultat par notre hypothèse de récurrence.
\end{proof}








\section{Cycles de codimension donnée}
\label{chow-section-cycles-codimension}

\noindent
Dans certains cas, le groupe abélien de tous les cycles possède une seconde
graduation donnée par la codimension.

\begin{lemma}
\label{chow-lemma-locally-equidimensional}
Soit $(S, \delta)$ comme dans la Situation \ref{chow-situation-setup}.
Soit $X$ localement de type fini sur $S$. Écrivons
$\delta = \delta_{X/S}$ comme dans la section \ref{chow-section-setup}.
Les assertions suivantes sont équivalentes :
\begin{enumerate}
\item Il existe une décomposition $X = \coprod_{n \in \mathbf{Z}} X_n$
en sous-schémas ouverts et fermés telle que $\delta(\xi) = n$ dès que
$\xi \in X_n$ est un point générique d'une composante irréductible de $X_n$.
\item Pour tout $x \in X$, il existe un voisinage ouvert $U \subset X$
de $x$ et un entier $n$ tels que $\delta(\xi) = n$ dès que
$\xi \in U$ est un point générique d'une composante irréductible de $U$.
\item Pour tout $x \in X$, il existe un entier $n_x$ tel que
$\delta(\xi) = n_x$ pour tout point générique $\xi$ d'une composante
irréductible de $X$ contenant $x$.
\end{enumerate}
Les conditions sont satisfaites si $X$ est soit
normal, soit de Cohen--Macaulay\footnote{En fait, il suffit que
$X$ soit $(S_2)$. Comparer avec Cohomologie locale, Lemme
\ref{local-cohomology-lemma-catenary-S2-equidimensional}.}.
\end{lemma}

\begin{proof}
Il est clair que (1) $\Rightarrow$ (2) $\Rightarrow$ (3).
Réciproquement, si (3) est vérifié, alors nous posons $X_n = \{x \in X \mid n_x = n\}$
et nous obtenons une décomposition comme dans (1). En effet, $X_n$ est ouvert car,
pour $x$ donné, l'union des composantes irréductibles de $X$ passant par $x$
privée de l'union des composantes irréductibles de $X$ ne passant pas par $x$
est un voisinage ouvert de $x$. Si $X$ est normal, alors $X$ est une
réunion disjointe de schémas intègres
(Propriétés, Lemme \ref{properties-lemma-normal-locally-Noetherian})
et les propriétés sont donc vérifiées.
Si $X$ est de Cohen--Macaulay, alors
$\delta' : X \to \mathbf{Z}$, $x \mapsto -\dim(\mathcal{O}_{X, x})$
est une fonction de dimension sur $X$ (voir Exemple \ref{chow-example-CM-irreducible}).
Puisque $\delta - \delta'$ est localement constante
(Topologie, Lemme \ref{topology-lemma-dimension-function-unique})
et puisque $\delta'(\xi) = 0$ pour tout point générique $\xi$ de $X$,
nous voyons que (2) est vérifié.
\end{proof}

\noindent
Soit $(S, \delta)$ comme dans la Situation \ref{chow-situation-setup}. Soit $X$
localement de type fini sur $S$ et satisfaisant aux conditions équivalentes
du Lemme \ref{chow-lemma-locally-equidimensional}. Pour un sous-schéma fermé
intègre $Z \subset X$, nous avons la codimension $\text{codim}(Z, X)$
de $Z$ dans $X$, voir Topologie, Définition \ref{topology-definition-codimension}.
Nous définissons un {\it cycle de codimension $p$} comme un cycle $\alpha = \sum n_Z[Z]$
sur $X$ tel que $n_Z \not = 0 \Rightarrow \text{codim}(Z, X) = p$.
Le groupe abélien de tous les cycles de codimension $p$ est noté $Z^p(X)$.
Soit $X = \coprod X_n$ la décomposition donnée dans la
partie (1) du Lemme \ref{chow-lemma-locally-equidimensional}.
Rappelant que nos cycles sont définis comme des sommes localement finies, il est clair que
$$
Z^p(X) = \prod\nolimits_n Z_{n - p}(X_n)
$$
De plus, nous voyons que $\prod_p Z^p(X) = \prod_k Z_k(X)$. Nous pourrions maintenant définir
l'équivalence rationnelle des cycles de codimension $p$ sur $X$ exactement de la même
manière que précédemment et, en fait, redévelopper toute la théorie à partir de zéro
pour les cycles de codimension donnée lorsque $X$ est comme dans le
Lemme \ref{chow-lemma-locally-equidimensional}. Nous nous contentons plutôt de
définir le {\it groupe de Chow des cycles de codimension $p$} par
$$
\CH^p(X) = \prod\nolimits_n \CH_{n - p}(X_n)
$$
Comme précédemment, nous avons $\prod_p \CH^p(X) = \prod_k \CH_k(X)$.
Si $X$ est quasi-compact, alors le produit dans la formule est fini
(et c'est donc une somme directe) et nous avons
$\bigoplus_p \CH^p(X) = \bigoplus_k \CH_k(X)$. Si $X$ est quasi-compact
et de dimension finie, alors seul un nombre fini de ces groupes
est non nul.

\medskip\noindent
De nombreuses constructions et de nombreux résultats pour les groupes de Chow démontrés ci-dessus ont
des analogues naturels pour les groupes de Chow $\CH^*(X)$. Chacun s'obtient
en décomposant les schémas pertinents en morceaux « équidimensionnels »
comme dans le Lemme \ref{chow-lemma-locally-equidimensional},
puis en appliquant les résultats déjà démontrés aux
facteurs de la décomposition en produit donnée ci-dessus.
Énumérons-en quelques-uns.
\begin{enumerate}
\item Si $f : X \to Y$ est un morphisme plat de schémas localement de type fini
sur $S$ et si $X$ et $Y$ satisfont aux conditions équivalentes du
Lemme \ref{chow-lemma-locally-equidimensional}, alors l'image inverse plate détermine une application
$$
f^* : \CH^p(Y) \to \CH^p(X)
$$
\item Si $f : X \to Y$ est un morphisme de schémas localement de type fini
sur $S$ et si $X$ et $Y$ satisfont aux conditions équivalentes du
Lemme \ref{chow-lemma-locally-equidimensional}, disons que $f$ est de
{\it codimension} $r \in \mathbf{Z}$ si, pour toute paire de composantes irréductibles
$Z \subset X$, $W \subset Y$ telles que $f(Z) \subset W$, nous avons
$\dim_\delta(W) - \dim_\delta(Z) = r$.
\item Si $f : X \to Y$ est un morphisme propre de schémas localement de type fini
sur $S$, si $X$ et $Y$ satisfont aux conditions équivalentes du
Lemme \ref{chow-lemma-locally-equidimensional} et si $f$ est de codimension $r$,
alors l'image directe propre est une application
$$
f_* : \CH^p(X) \to \CH^{p + r}(Y)
$$
\item Si $f : X \to Y$ est un morphisme de schémas localement de type fini sur $S$,
si $X$ et $Y$ satisfont aux conditions équivalentes du
Lemme \ref{chow-lemma-locally-equidimensional}, si $f$ est de codimension $r$
et si $c \in A^q(X \to Y)$, alors $c$ induit des applications
$$
c \cap -  : \CH^p(Y) \to \CH^{p + q - r}(X)
$$
\item Si $X$ est un schéma localement de type fini sur $S$
satisfaisant aux conditions équivalentes du
Lemme \ref{chow-lemma-locally-equidimensional} et si
$\mathcal{L}$ est un $\mathcal{O}_X$-module inversible,
alors
$$
c_1(\mathcal{L}) \cap - : \CH^p(X) \to \CH^{p + 1}(X)
$$
\item Si $X$ est un schéma localement de type fini sur $S$
satisfaisant aux conditions équivalentes du
Lemme \ref{chow-lemma-locally-equidimensional} et si
$\mathcal{E}$ est un $\mathcal{O}_X$-module localement libre de type fini,
alors
$$
c_i(\mathcal{E}) \cap - : \CH^p(X) \to \CH^{p + i}(X)
$$
\end{enumerate}
Attention : la propriété, pour un morphisme, d'être de codimension $r$
n'est pas préservée par changement de base.

\begin{remark}
\label{chow-remark-fundamental-class}
Soit $(S, \delta)$ comme dans la Situation \ref{chow-situation-setup}.
Soit $X$ localement de type fini sur $S$ et satisfaisant aux
conditions équivalentes du Lemme \ref{chow-lemma-locally-equidimensional}.
Soit $X = \coprod X_n$ la décomposition en sous-schémas ouverts et fermés
telle que toute composante irréductible de $X_n$ soit de
dimension $\delta$ égale à $n$. Dans cette situation, nous posons parfois
$$
[X] = \sum\nolimits_n [X_n]_n \in \CH^0(X)
$$
Cette classe est une sorte de « classe fondamentale » de $X$ dans la théorie de Chow.
\end{remark}




\section{Le principe de scindage}
\label{chow-section-splitting-principle}

\noindent
Dans notre cadre, il n'est pas si facile de dire ce que le principe de scindage
dit/est exactement. En voici une formulation possible.

\begin{lemma}
\label{chow-lemma-splitting-principle}
Soit $(S, \delta)$ comme dans la Situation \ref{chow-situation-setup}. Soit $X$ localement
de type fini sur $S$. Les $\mathcal{E}_i$ forment une collection finie de
$\mathcal{O}_X$-modules localement libres de rang $r_i$. Il existe un morphisme
projectif plat $\pi : P \to X$ de dimension relative $d$ tel que
\begin{enumerate}
\item pour tout morphisme $f : Y \to X$, l'application
$\pi_Y^* : \CH_*(Y) \to \CH_{* + d}(Y \times_X P)$ est injective, et
\item chaque $\pi^*\mathcal{E}_i$ possède une filtration
dont les quotients successifs $\mathcal{L}_{i, 1}, \ldots, \mathcal{L}_{i, r_i}$
sont des ${\mathcal O}_P$-modules inversibles.
\end{enumerate}
De plus, lorsque (1) est vérifié, l'application de restriction $A^*(X) \to A^*(P)$
(Remarque \ref{chow-remark-pullback-cohomology}) est injective.
\end{lemma}

\begin{proof}
Nous pouvons supposer $r_i \geq 1$ pour tout $i$. Nous démontrerons le lemme par récurrence
sur $\sum (r_i - 1)$. Si cet entier vaut $0$, alors $\mathcal{E}_i$
est inversible pour tout $i$ et nous concluons en prenant $\pi = \text{id}_X$.
Sinon, nous pouvons choisir un $i$ tel que $r_i > 1$ et considérer le
morphisme $\pi_i : P_i = \mathbf{P}(\mathcal{E}_i) \to X$.
Nous avons une suite exacte courte
$$
0 \to \mathcal{F} \to \pi_i^*\mathcal{E}_i \to \mathcal{O}_{P_i}(1) \to 0
$$
de $\mathcal{O}_{P_i}$-modules localement libres de type fini de rangs $r_i - 1$,
$r_i$ et $1$. Observons que $\pi_i^*$ est injective sur les groupes de Chow
après tout changement de base, par la formule du fibré projectif
(Lemme \ref{chow-lemma-chow-ring-projective-bundle}).
Par l'hypothèse de récurrence appliquée aux
$\mathcal{O}_{P_i}$-modules localement libres de type fini $\mathcal{F}$ et $\pi_{i'}^*\mathcal{E}_{i'}$
pour $i' \not = i$, nous trouvons un morphisme $\pi : P \to P_i$ ayant
les propriétés énoncées dans le lemme. Alors la composée
$\pi_i \circ \pi : P \to X$ convient. Certains détails sont omis.
\end{proof}

\begin{remark}
\label{chow-remark-the-proof-shows-more}
La démonstration du Lemme \ref{chow-lemma-splitting-principle}
montre que le morphisme $\pi : P \to X$ possède les propriétés supplémentaires
suivantes :
\begin{enumerate}
\item $\pi$ est une composée finie de fibrés en espaces projectifs
associés à des modules localement libres de rang fini constant, et
\item pour tout $\alpha \in \CH_k(X)$, nous avons
$\alpha = \pi_*(\xi_1 \cap \ldots \cap \xi_d \cap \pi^*\alpha)$,
où $\xi_i$ est la première classe de Chern d'un
$\mathcal{O}_P$-module inversible.
\end{enumerate}
La seconde observation résulte de la première et du
Lemme \ref{chow-lemma-cap-projective-bundle}.
Nous ajouterons ici d'autres observations selon les besoins.
\end{remark}

\noindent
Soient $(S, \delta)$, $X$ et $\mathcal{E}_i$ comme dans le
Lemme \ref{chow-lemma-splitting-principle}.
Le {\it principe de scindage} désigne l'usage consistant à écrire symboliquement
$$
c(\mathcal{E}_i) = \prod (1 + x_{i, j})
$$
Les symboles $x_{i, 1}, \ldots, x_{i, r_i}$ sont appelés les {\it racines de Chern}
de $\mathcal{E}_i$. Autrement dit, la $p$-ième classe de Chern de $\mathcal{E}_i$
est la $p$-ième fonction symétrique élémentaire des racines de Chern.
L'utilité du principe de scindage vient de l'assertion selon laquelle,
pour démontrer une relation polynomiale entre les classes de Chern des
$\mathcal{E}_i$, il suffit de démontrer la relation correspondante entre les
racines de Chern.

\medskip\noindent
En effet, soit $\pi : P \to X$ comme dans le Lemme \ref{chow-lemma-splitting-principle}.
Rappelons qu'il existe un morphisme canonique de $\mathbf{Z}$-algèbres
$\pi^* : A^*(X) \to A^*(P)$, voir
Remarque \ref{chow-remark-pullback-cohomology}. L'injectivité de $\pi_Y^*$
sur les groupes de Chow pour tout $Y$ sur $X$ implique que l'application
$\pi^* : A^*(X) \to A^*(P)$ est injective (détails omis).
Nous avons
$$
\pi^*c(\mathcal{E}_i) = \prod (1 + c_1(\mathcal{L}_{i, j}))
$$
par le Lemme \ref{chow-lemma-chern-filter-by-linebundles}. Nous pouvons donc considérer les
racines de Chern $x_{i, j}$ comme les éléments $c_1(\mathcal{L}_{i, j}) \in A^*(P)$
et l'équation affichée comme ayant lieu dans $A^*(P)$ après
application du morphisme injectif $\pi^* : A^*(X) \to A^*(P)$ au
membre de gauche de l'équation.

\medskip\noindent
Pour voir comment cela fonctionne, il est préférable de donner quelques exemples.

\begin{lemma}
\label{chow-lemma-chern-classes-dual}
Dans la Situation \ref{chow-situation-setup}, soit $X$ localement de type fini sur $S$.
Soit $\mathcal{E}$ un $\mathcal{O}_X$-module localement libre de type fini
de dual $\mathcal{E}^\vee$. Alors
$$
c_i(\mathcal{E}^\vee) = (-1)^i c_i(\mathcal{E})
$$
dans $A^i(X)$.
\end{lemma}

\begin{proof}
Choisissons un morphisme $\pi : P \to X$ comme dans le
Lemme \ref{chow-lemma-splitting-principle}.
Par l'injectivité de $\pi^*$ (après tout changement de base),
il suffit de démontrer la relation entre
les classes de Chern de $\mathcal{E}$ et de $\mathcal{E}^\vee$
après image inverse sur $P$. Nous pouvons donc supposer qu'il
existe des $\mathcal{O}_X$-modules inversibles
${\mathcal L}_i$, $i = 1, \ldots, r$,
et une filtration
$$
0 = \mathcal{E}_0 \subset \mathcal{E}_1 \subset \mathcal{E}_2
\subset \ldots \subset \mathcal{E}_r = \mathcal{E}
$$
telle que $\mathcal{E}_i/\mathcal{E}_{i - 1} \cong \mathcal{L}_i$.
Nous obtenons alors la filtration duale
$$
0 = \mathcal{E}_r^\perp \subset \mathcal{E}_1^\perp \subset \mathcal{E}_2^\perp
\subset \ldots \subset \mathcal{E}_0^\perp = \mathcal{E}^\vee
$$
telle que $\mathcal{E}_{i - 1}^\perp/\mathcal{E}_i^\perp \cong
\mathcal{L}_i^{\otimes -1}$.
Posons $x_i = c_1(\mathcal{L}_i)$.
Alors $c_1(\mathcal{L}_i^{\otimes -1}) = - x_i$
par le Lemme \ref{chow-lemma-c1-cap-additive}.
Par le Lemme \ref{chow-lemma-chern-filter-by-linebundles},
nous avons
$$
c(\mathcal{E}) = \prod\nolimits_{i = 1}^r (1 + x_i)
\quad\text{et}\quad
c(\mathcal{E}^\vee) = \prod\nolimits_{i = 1}^r (1 - x_i)
$$
dans $A^*(X)$. Le résultat résulte d'un calcul formel
que nous omettons.
\end{proof}

\begin{lemma}
\label{chow-lemma-chern-classes-tensor-product}
Dans la Situation \ref{chow-situation-setup}, soit $X$ localement de type fini sur $S$.
Soient $\mathcal{E}$ et $\mathcal{F}$ deux
$\mathcal{O}_X$-modules localement libres de type fini, de rangs $r$ et $s$. Alors nous avons
$$
c_1(\mathcal{E} \otimes \mathcal{F})
=
r c_1(\mathcal{F}) + s c_1(\mathcal{E})
$$
$$
c_2(\mathcal{E} \otimes \mathcal{F})
=
r c_2(\mathcal{F}) + s c_2(\mathcal{E}) +
{r \choose 2} c_1(\mathcal{F})^2 +
(rs - 1) c_1(\mathcal{F})c_1(\mathcal{E}) +
{s \choose 2} c_1(\mathcal{E})^2
$$
et ainsi de suite dans $A^*(X)$.
\end{lemma}

\begin{proof}
En raisonnant exactement comme dans la démonstration du Lemme \ref{chow-lemma-chern-classes-dual},
nous pouvons supposer que nous avons
des $\mathcal{O}_X$-modules inversibles
${\mathcal L}_i$, $i = 1, \ldots, r$,
${\mathcal N}_i$, $i = 1, \ldots, s$,
des filtrations
$$
0 = \mathcal{E}_0 \subset \mathcal{E}_1 \subset \mathcal{E}_2
\subset \ldots \subset \mathcal{E}_r = \mathcal{E}
\quad\text{et}\quad
0 = \mathcal{F}_0 \subset \mathcal{F}_1 \subset \mathcal{F}_2
\subset \ldots \subset \mathcal{F}_s = \mathcal{F}
$$
telles que $\mathcal{E}_i/\mathcal{E}_{i - 1} \cong \mathcal{L}_i$
et telles que $\mathcal{F}_j/\mathcal{F}_{j - 1} \cong \mathcal{N}_j$.
En ordonnant les paires $(i, j)$ lexicographiquement,
nous obtenons une filtration
$$
0 \subset \ldots \subset
\mathcal{E}_i \otimes \mathcal{F}_j
+
\mathcal{E}_{i - 1} \otimes \mathcal{F}
\subset \ldots \subset \mathcal{E} \otimes \mathcal{F}
$$
de quotients successifs
$$
\mathcal{L}_1 \otimes \mathcal{N}_1,
\mathcal{L}_1 \otimes \mathcal{N}_2,
\ldots,
\mathcal{L}_1 \otimes \mathcal{N}_s,
\mathcal{L}_2 \otimes \mathcal{N}_1,
\ldots,
\mathcal{L}_r \otimes \mathcal{N}_s
$$
Par le Lemme \ref{chow-lemma-chern-filter-by-linebundles},
nous avons
$$
c(\mathcal{E}) = \prod (1 + x_i),
\quad
c(\mathcal{F}) = \prod (1 + y_j),
\quad\text{et}\quad
c(\mathcal{E} \otimes \mathcal{F}) = \prod (1 + x_i + y_j),
$$
dans $A^*(X)$. Le résultat résulte d'un calcul formel
que nous omettons.
\end{proof}

\begin{remark}
\label{chow-remark-equalities-nonconstant-rank}
Les égalités démontrées ci-dessus restent vraies même lorsque nous travaillons avec
des modules localement libres de type fini
sur $\mathcal{O}_X$ dont le rang peut ne pas être constant.
En fait, nous pouvons travailler avec des polynômes en le rang et les
classes de Chern comme suit. Considérons l'anneau gradué de polynômes
$\mathbf{Z}[r, c_1, c_2, c_3, \ldots]$,
où $r$ est de degré $0$ et $c_i$ de degré $i$. Soit
$$
P \in \mathbf{Z}[r, c_1, c_2, c_3, \ldots]
$$
un polynôme homogène de degré $p$. Alors, pour tout
$\mathcal{O}_X$-module localement libre de type fini $\mathcal{E}$ sur $X$, nous pouvons considérer
$$
P(\mathcal{E}) =
P(r(\mathcal{E}), c_1(\mathcal{E}), c_2(\mathcal{E}), c_3(\mathcal{E}), \ldots)
\in A^p(X)
$$
voir la Remarque \ref{chow-remark-extend-to-finite-locally-free} pour les notations et
conventions. Pour démontrer des relations entre ces polynômes (pour plusieurs
modules localement libres de type fini), nous pouvons travailler localement sur $X$ et utiliser le principe de
scindage comme ci-dessus. Par exemple, nous affirmons que
$$
c_2(\SheafHom_{\mathcal{O}_X}(\mathcal{E}, \mathcal{E})) =
P(\mathcal{E})
$$
où $P = 2rc_2 - (r - 1)c_1^2$.
En effet, puisque $\SheafHom_{\mathcal{O}_X}(\mathcal{E}, \mathcal{E}) =
\mathcal{E} \otimes \mathcal{E}^\vee$, cela résulte facilement des
Lemmes \ref{chow-lemma-chern-classes-dual} et
\ref{chow-lemma-chern-classes-tensor-product}
ci-dessus, en décomposant $X$ en morceaux où le rang
de $\mathcal{E}$ est constant comme dans la
Remarque \ref{chow-remark-extend-to-finite-locally-free}.
\end{remark}

\begin{example}
\label{chow-example-power-sum}
Pour tout $p \geq 1$, il existe un unique polynôme homogène
$P_p \in \mathbf{Z}[c_1, c_2, c_3, \ldots]$ de degré $p$
tel que, pour tout $n \geq p$, nous ayons
$$
P_p(s_1, s_2, \ldots, s_p) = \sum x_i^p
$$
dans $\mathbf{Z}[x_1, \ldots, x_n]$, où $s_1, \ldots, s_p$ sont les
polynômes symétriques élémentaires en $x_1, \ldots, x_n$, donc
$$
s_i = \sum\nolimits_{1 \leq j_1 < \ldots < j_i \leq n}
x_{j_1}x_{j_2} \ldots x_{j_i}
$$
L'existence de $P_p$ résulte du fait bien connu que
les fonctions symétriques élémentaires engendrent l'anneau de
toutes les fonctions symétriques sur les entiers. Une autre manière de
caractériser $P_p \in \mathbf{Z}[c_1, c_2, c_3, \ldots]$ est que nous avons
$$
\log(1 + c_1 + c_2 + c_3 + \ldots) =
\sum\nolimits_{p \geq 1} (-1)^{p - 1}\frac{P_p}{p}
$$
comme série formelle. Cela est clair en écrivant
$1 + c_1 + c_2 + \ldots = \prod (1 + x_i)$ et en appliquant
la série de la fonction logarithme. En développant le membre de
gauche, nous obtenons
\begin{align*}
& (c_1 + c_2 + \ldots) - (1/2)(c_1 + c_2 + \ldots)^2 +
(1/3)(c_1 + c_2 + \ldots)^3 - \ldots \\
& =
c_1 + (c_2 - (1/2)c_1^2) + (c_3 - c_1c_2 + (1/3)c_1^3) + \ldots
\end{align*}
De cette manière, nous trouvons
\begin{align*}
P_1 & = c_1, \\
P_2 & = c_1^2 - 2c_2, \\
P_3 & = c_1^3 - 3c_1c_2 + 3c_3, \\
P_4 & = c_1^4 - 4c_1^2c_2 + 4c_1c_3 + 2c_2^2 - 4c_4,
\end{align*}
et ainsi de suite. Puisque les classes de Chern d'un
$\mathcal{O}_X$-module localement libre de type fini $\mathcal{E}$ sont les polynômes
symétriques élémentaires en les racines de Chern $x_i$, nous voyons que
$$
P_p(\mathcal{E}) = \sum x_i^p
$$
Pour plus de commodité, nous posons $P_0 = r$ dans $\mathbf{Z}[r, c_1, c_2, c_3, \ldots]$
de sorte que $P_0(\mathcal{E}) = r(\mathcal{E})$ comme classe bivariante
(comme dans les Remarques \ref{chow-remark-extend-to-finite-locally-free} et
\ref{chow-remark-equalities-nonconstant-rank}).
\end{example}






\section{Classes de Chern et sections}
\label{chow-section-top-chern-class}

\noindent
Une brève section dont le résultat principal est que nous pouvons calculer la
classe de Chern de degré maximal d'un module localement libre de type fini au moyen du
lieu des zéros d'une « section régulière ».

\medskip\noindent
Soit $(S, \delta)$ comme dans la Situation \ref{chow-situation-setup}. Soit $X$ un schéma
localement de type fini sur $S$. Soit $\mathcal{E}$ un
$\mathcal{O}_X$-module localement libre de type fini. Soit $f : X' \to X$ localement de type fini. Soit
$$
s \in \Gamma(X', f^*\mathcal{E})
$$
une section globale de l'image inverse de $\mathcal{E}$ sur $X'$. Soit
$Z(s) \subset X'$ le schéma des zéros de $s$. Plus précisément, nous définissons
$Z(s)$ comme le sous-schéma fermé dont le faisceau quasi-cohérent
d'idéaux est l'image de l'application
$s : f^*\mathcal{E}^\vee \to \mathcal{O}_{X'}$.

\begin{lemma}
\label{chow-lemma-top-chern-class}
Dans la situation décrite juste ci-dessus, supposons $\dim_\delta(X') = n$,
que $f^*\mathcal{E}$ soit de rang constant $r$, que
$\dim_\delta(Z(s)) \leq n - r$, et que, pour tout point générique
$\xi \in Z(s)$ avec $\delta(\xi) = n - r$, l'idéal de $Z(s)$
dans $\mathcal{O}_{X', \xi}$ soit engendré par une suite régulière
de longueur $r$. Alors
$$
c_r(\mathcal{E}) \cap [X']_n = [Z(s)]_{n - r}
$$
dans $\CH_*(X')$.
\end{lemma}

\begin{proof}
Puisque $c_r(\mathcal{E})$ est une classe bivariante
(Lemme \ref{chow-lemma-cap-cp-bivariant}),
nous pouvons supposer $X = X'$ et nous devons montrer que
$c_r(\mathcal{E}) \cap [X]_n = [Z(s)]_{n - r}$ dans $\CH_{n - r}(X)$.
Nous démontrerons le lemme par récurrence sur $r \geq 0$. (Le cas
$r = 0$ est trivial.) Le cas $r = 1$
est traité par le Lemme \ref{chow-lemma-geometric-cap}. Supposons $r > 1$.

\medskip\noindent
Soit $\pi : P \to X$ le fibré en espaces projectifs associé à
$\mathcal{E}$ et considérons la suite exacte courte
$$
0 \to \mathcal{E}' \to \pi^*\mathcal{E} \to \mathcal{O}_P(1) \to 0
$$
Par la formule du fibré en espaces projectifs
(Lemme \ref{chow-lemma-chow-ring-projective-bundle}),
il suffit de démontrer l'égalité après image inverse par $\pi$.
Observons que $\pi^{-1}Z(s) = Z(\pi^*s)$ est de dimension $\delta$
$\leq n - 1$ et que l'hypothèse sur les suites régulières aux
points génériques de dimension $\delta$ égale à $n - 1$ est vérifiée par
image inverse plate, voir
Algèbre, Lemme \ref{algebra-lemma-flat-increases-depth}.
Soit $t \in \Gamma(P, \mathcal{O}_P(1))$ l'image de $\pi^*s$.
Nous affirmons que
$$
[Z(t)]_{n + r - 2} = c_1(\mathcal{O}_P(1)) \cap [P]_{n + r - 1}
$$
En admettant cette assertion, nous terminons la démonstration comme suit.
La restriction $\pi^*s|_{Z(t)}$ s'envoie sur zéro dans
$\mathcal{O}_P(1)|_{Z(t)}$ et provient donc d'un unique
élément $s' \in \Gamma(Z(t), \mathcal{E}'|_{Z(t)})$.
Notons que $Z(s') = Z(\pi^*s)$ comme sous-schémas fermés de $P$.
Si $\xi \in Z(s')$ est un point générique avec $\delta(\xi) = n - 1$,
alors l'idéal de $Z(s')$ dans $\mathcal{O}_{Z(t), \xi}$
peut être engendré par une suite régulière de longueur $r - 1$ : il est engendré par
$r - 1$ éléments qui sont les images de $r - 1$ éléments de
$\mathcal{O}_{P, \xi}$ qui, avec un générateur de
l'idéal de $Z(t)$ dans $\mathcal{O}_{P, \xi}$, forment une suite régulière
de longueur $r$ dans $\mathcal{O}_{P, \xi}$. Nous pouvons donc appliquer
l'hypothèse de récurrence à $s'$ sur $Z(t)$ pour obtenir
$c_{r - 1}(\mathcal{E}') \cap [Z(t)]_{n + r - 2} = [Z(s')]_{n - 1}$.
En combinant tout ce qui précède, nous obtenons
\begin{align*}
c_r(\pi^*\mathcal{E}) \cap [P]_{n + r - 1}
& =
c_{r - 1}(\mathcal{E}') \cap c_1(\mathcal{O}_P(1)) \cap [P]_{n + r - 1} \\
& =
c_{r - 1}(\mathcal{E}') \cap [Z(t)]_{n + r - 2} \\
& =
[Z(s')]_{n - 1} \\
& = [Z(\pi^*s)]_{n - 1}
\end{align*}
ce qui est précisément ce qu'il fallait montrer.

\medskip\noindent
Démonstration de l'assertion. Elle résultera d'une application du
Lemme \ref{chow-lemma-geometric-cap}, déjà utilisé.
Nous avons $\pi^{-1}(Z(s)) = Z(\pi^*s) \subset Z(t)$.
D'autre part, pour $x \in X$, si $P_x \subset Z(t)$, alors
$t|_{P_x} = 0$, ce qui implique que $s$ est nul dans la fibre
$\mathcal{E} \otimes \kappa(x)$, ce qui implique $x \in Z(s)$.
Il s'ensuit que $\dim_\delta(Z(t)) \leq n + (r - 1) - 1$.
Enfin, soit $\xi \in Z(t)$ un point générique tel que
$\delta(\xi) = n + r - 2$. Si $\xi$ n'est pas le point générique
de la fibre de $P \to X$, il est immédiat que
toute équation locale de $Z(t)$ est un non-diviseur de zéro dans $\mathcal{O}_{P, \xi}$
(car nous pouvons le vérifier sur la fibre par
Algèbre, Lemme \ref{algebra-lemma-grothendieck}).
Si $\xi$ est le point générique d'une fibre, alors $x = \pi(\xi) \in Z(s)$
et $\delta(x) = n + r - 2 - (r - 1) = n - 1$. Cela contredit
$\dim_\delta(Z(s)) \leq n - r$ puisque $r > 1$,
de sorte que ce cas ne se produit pas.
\end{proof}

\begin{lemma}
\label{chow-lemma-easy-virtual-class}
Soit $(S, \delta)$ comme dans la Situation \ref{chow-situation-setup}. Soit $X$
un schéma localement de type fini sur $S$. Soit
$$
0 \to \mathcal{N}' \to \mathcal{N} \to \mathcal{E} \to 0
$$
une suite exacte courte de $\mathcal{O}_X$-modules localement libres de type fini.
Considérons l'immersion fermée
$$
i :
N' = \underline{\Spec}_X(\text{Sym}((\mathcal{N}')^\vee))
\longrightarrow
N = \underline{\Spec}_X(\text{Sym}(\mathcal{N}^\vee))
$$
Pour $\alpha \in \CH_k(X)$, nous avons
$$
i_*(p')^*\alpha = p^*(c_{top}(\mathcal{E}) \cap \alpha)
$$
où $p' : N' \to X$ et $p : N \to X$ sont les morphismes structuraux.
\end{lemma}

\begin{proof}
Ici, $c_{top}(\mathcal{E})$ est la classe bivariante définie dans la
Remarque \ref{chow-remark-top-chern-class}. Par sa définition même, pour
vérifier la formule, nous pouvons supposer que $\mathcal{E}$
est de rang constant. Nous pouvons de même supposer que $\mathcal{N}'$ et
$\mathcal{N}$ sont de rangs constants, disons $r'$ et $r$, de sorte que
$\mathcal{E}$ est de rang $r - r'$ et
$c_{top}(\mathcal{E}) = c_{r - r'}(\mathcal{E})$.
Observons que $p^*\mathcal{E}$ possède une section canonique
$$
s \in \Gamma(N, p^*\mathcal{E}) = \Gamma(X, p_*p^*\mathcal{E}) =
\Gamma(X, \mathcal{E} \otimes_{\mathcal{O}_X} \text{Sym}(\mathcal{N}^\vee)
\supset \Gamma(X, \SheafHom(\mathcal{N}, \mathcal{E}))
$$
correspondant à la surjection $\mathcal{N} \to \mathcal{E}$ donnée
dans l'énoncé du lemme. Le schéma d'annulation de cette section
est exactement $N' \subset N$. Soit $Y \subset X$ un sous-schéma fermé
intègre de dimension $\delta$ égale à $n$. Alors nous avons
\begin{enumerate}
\item $p^*[Y] = [p^{-1}(Y)]$ puisque $p^{-1}(Y)$ est intègre de
dimension $\delta$ égale à $n + r$,
\item $(p')^*[Y] = [(p')^{-1}(Y)]$ puisque $(p')^{-1}(Y)$ est intègre de
dimension $\delta$ égale à $n + r'$,
\item la restriction de $s$ à $p^{-1}Y$ a pour schéma d'annulation
$(p')^{-1}Y$ et l'immersion fermée $(p')^{-1}Y \to p^{-1}Y$
est une immersion régulière (localement découpée par une suite régulière).
\end{enumerate}
Nous concluons que
$$
(p')^*[Y] = c_{r - r'}(p^*\mathcal{E}) \cap p^*[Y]
\quad\text{dans}\quad \CH_*(N)
$$
par le Lemme \ref{chow-lemma-top-chern-class}. Cela démontre le lemme.
\end{proof}










\section{Le caractère de Chern et les produits tensoriels}
\label{chow-section-chern-classes-tensor}

\noindent
Soit $(S, \delta)$ comme dans la Situation \ref{chow-situation-setup}.
Soit $X$ localement de type fini sur $S$.
Nous définissons le {\it caractère de Chern} d'un
$\mathcal{O}_X$-module localement libre de type fini comme l'expression formelle
$$
ch({\mathcal E}) = \sum\nolimits_{i=1}^r e^{x_i}
$$
si les $x_i$ sont les racines de Chern de ${\mathcal E}$. En écrivant cette
expression comme un polynôme en les classes de Chern, nous obtenons
\begin{align*}
ch(\mathcal{E})
& =
r(\mathcal{E})
+
c_1(\mathcal{E}) +
\frac{1}{2}(c_1(\mathcal{E})^2 - 2c_2(\mathcal{E}))
+
\frac{1}{6}(c_1(\mathcal{E})^3 - 3c_1(\mathcal{E})c_2(\mathcal{E}) + 3c_3(\mathcal{E})) \\
& \quad\quad +
\frac{1}{24}(c_1(\mathcal{E})^4 - 4c_1(\mathcal{E})^2c_2(\mathcal{E}) + 4c_1(\mathcal{E})c_3(\mathcal{E}) + 2c_2(\mathcal{E})^2 - 4c_4(\mathcal{E}))
+
\ldots \\
& =
\sum\nolimits_{p = 0, 1, 2, \ldots} \frac{P_p(\mathcal{E})}{p!}
\end{align*}
où $P_p$ sont les polynômes en les classes de Chern de
l'Exemple \ref{chow-example-power-sum}. La composante de degré $p$ de
ce qui précède est
$$
ch_p(\mathcal{E}) = \frac{P_p(\mathcal{E})}{p!} \in A^p(X) \otimes \mathbf{Q}
$$
Que signifie le fait que les coefficients soient des nombres rationnels ?
Eh bien, cela signifie simplement que nous considérons
$ch_p(\mathcal{E})$ comme un élément de $A^p(X) \otimes \mathbf{Q}$.

\begin{remark}
\label{chow-remark-extend-chern-character-to-finite-locally-free}
Dans la discussion ci-dessus, nous avons défini les composantes du caractère de Chern
$ch_p(\mathcal{E}) \in A^p(X) \otimes \mathbf{Q}$
de $\mathcal{E}$ même si le rang de $\mathcal{E}$
n'est pas constant. Voir les Remarques \ref{chow-remark-extend-to-finite-locally-free} et
\ref{chow-remark-equalities-nonconstant-rank}. Ainsi, le caractère de Chern complet
de $\mathcal{E}$ est
un élément de $\prod_{p \geq 0} (A^p(X) \otimes \mathbf{Q})$. Si $X$
est quasi-compact et $\dim(X) < \infty$ (dimension ordinaire), alors on peut montrer,
en utilisant le Lemme \ref{chow-lemma-vanish-above-dimension} et le principe de scindage,
que $ch(\mathcal{E}) \in A^*(X) \otimes \mathbf{Q}$.
\end{remark}

\begin{lemma}
\label{chow-lemma-chern-character-additive}
Soit $(S, \delta)$ comme dans la Situation \ref{chow-situation-setup}. Soit $X$ localement
de type fini sur $S$. Soit
$
0 \to \mathcal{E}_1 \to \mathcal{E} \to \mathcal{E}_2 \to 0
$
une suite exacte courte de $\mathcal{O}_X$-modules localement libres de type fini.
Alors nous avons l'égalité
$$
ch(\mathcal{E}) = ch(\mathcal{E}_1) + ch(\mathcal{E}_2)
$$
Plus précisément, nous avons
$P_p(\mathcal{E}) = P_p(\mathcal{E}_1) + P_p(\mathcal{E}_2)$
dans $A^p(X)$, où $P_p$ est comme dans l'Exemple \ref{chow-example-power-sum}.
\end{lemma}

\begin{proof}
Il suffit de démontrer l'énoncé plus précis. Par la
section \ref{chow-section-splitting-principle},
cela résulte du fait que si $x_{1, i}$, $i = 1, \ldots, r_1$,
et $x_{2, i}$, $i = 1, \ldots, r_2$, sont les
racines de Chern de $\mathcal{E}_1$ et $\mathcal{E}_2$, alors
$x_{1, 1}, \ldots, x_{1, r_1}, x_{2, 1}, \ldots, x_{2, r_2}$
sont les racines de Chern de $\mathcal{E}$. Le résultat découle donc
de notre choix de $P_p$ dans l'Exemple \ref{chow-example-power-sum}.
\end{proof}

\begin{lemma}
\label{chow-lemma-chern-character-multiplicative}
Soit $(S, \delta)$ comme dans la Situation \ref{chow-situation-setup}. Soit $X$ localement
de type fini sur $S$. Soient $\mathcal{E}_1$ et $\mathcal{E}_2$
des $\mathcal{O}_X$-modules localement libres de type fini.
Alors nous avons l'égalité
$$
ch(\mathcal{E}_1 \otimes_{\mathcal{O}_X} \mathcal{E}_2) =
ch(\mathcal{E}_1) ch(\mathcal{E}_2)
$$
Plus précisément, nous avons
$$
P_p(\mathcal{E}_1 \otimes_{\mathcal{O}_X} \mathcal{E}_2) =
\sum\nolimits_{p_1 + p_2 = p}
{p \choose p_1} P_{p_1}(\mathcal{E}_1) P_{p_2}(\mathcal{E}_2)
$$
dans $A^p(X)$, où $P_p$ est comme dans l'Exemple \ref{chow-example-power-sum}.
\end{lemma}

\begin{proof}
Il suffit de démontrer l'énoncé plus précis. Par la
section \ref{chow-section-splitting-principle},
cela résulte du fait que si $x_{1, i}$, $i = 1, \ldots, r_1$,
et $x_{2, i}$, $i = 1, \ldots, r_2$, sont les
racines de Chern de $\mathcal{E}_1$ et $\mathcal{E}_2$, alors
$x_{1, i} + x_{2, j}$, $1 \leq i \leq r_1$, $1 \leq j \leq r_2$,
sont les racines de Chern de $\mathcal{E}_1 \otimes \mathcal{E}_2$.
Le résultat découle donc de la formule du binôme pour
$(x_{1, i} + x_{2, j})^p$ et de la
forme de nos polynômes $P_p$ dans l'Exemple \ref{chow-example-power-sum}.
\end{proof}

\begin{lemma}
\label{chow-lemma-chern-character-dual}
Dans la Situation \ref{chow-situation-setup}, soit $X$ localement de type fini sur $S$.
Soit $\mathcal{E}$ un $\mathcal{O}_X$-module localement libre de type fini
de dual $\mathcal{E}^\vee$. Alors
$ch_i(\mathcal{E}^\vee) = (-1)^i ch_i(\mathcal{E})$ dans
$A^i(X) \otimes \mathbf{Q}$.
\end{lemma}

\begin{proof}
Cela résulte du résultat correspondant pour les classes de Chern
(Lemme \ref{chow-lemma-chern-classes-dual}).
\end{proof}











\section{Classes de Chern et catégorie dérivée}
\label{chow-section-pre-derived}

\noindent
Dans cette section, nous définissons la classe de Chern totale d'un objet parfait $E$
de la catégorie dérivée d'un schéma $X$, sous l'hypothèse que $E$
peut être représenté par un complexe fini de modules localement libres de type fini
sur une enveloppe de $X$.

\medskip\noindent
Soit $(S, \delta)$ comme dans la Situation \ref{chow-situation-setup}.
Soit $X$ localement de type fini sur $S$. Soit
$$
\mathcal{E}^a \to \mathcal{E}^{a + 1} \to \ldots \to \mathcal{E}^b
$$
un complexe borné de $\mathcal{O}_X$-modules localement libres de type fini
de rang constant.
Nous définissons alors la {\it classe de Chern totale du complexe} par la formule
$$
c(\mathcal{E}^\bullet) = \prod\nolimits_{n = a, \ldots, b}
c(\mathcal{E}^n)^{(-1)^n} \in \prod\nolimits_{p \geq 0} A^p(X)
$$
Ici, l'inverse est l'inverse formel, donc
$$
(1 + c_1 + c_2 + c_3 + \ldots)^{-1} =
1 - c_1 + c_1^2 - c_2 - c_1^3 + 2c_1 c_2 - c_3 + \ldots
$$
Nous noterons $c_p(\mathcal{E}^\bullet) \in A^p(X)$
la partie de degré $p$ de $c(\mathcal{E}^\bullet)$.
Nous définissons de même le {\it caractère de Chern du complexe} par
la formule
$$
ch(\mathcal{E}^\bullet) = \sum\nolimits_{n = a, \ldots, b}
(-1)^n ch(\mathcal{E}^n) \in
\prod\nolimits_{p \geq 0} (A^p(X) \otimes \mathbf{Q})
$$
Nous noterons $ch_p(\mathcal{E}^\bullet) \in A^p(X) \otimes \mathbf{Q}$
la partie de degré $p$ de $ch(\mathcal{E}^\bullet)$.
Enfin, pour $P_p \in \mathbf{Z}[r, c_1, c_2, c_3, \ldots]$
comme dans l'Exemple \ref{chow-example-power-sum}, nous définissons
$$
P_p(\mathcal{E}^\bullet) = \sum\nolimits_{n = a, \ldots, b}
(-1)^n P_p(\mathcal{E}^n)
$$
dans $A^p(X)$. Alors nous avons
$ch_p(\mathcal{E}^\bullet) = (1/p!)P_p(\mathcal{E}^\bullet)$
comme d'habitude. Le lemme suivant montre que ces constructions ne dépend
que de l'image du complexe dans la catégorie dérivée.

\begin{lemma}
\label{chow-lemma-pre-derived-chern-class}
Soit $(S, \delta)$ comme dans la Situation \ref{chow-situation-setup}.
Soit $X$ localement de type fini sur $S$. Soit $E \in D(\mathcal{O}_X)$
un objet tel qu'il existe un complexe localement borné
$\mathcal{E}^\bullet$ de $\mathcal{O}_X$-modules localement libres de type fini
représentant $E$. Alors une légère généralisation des constructions ci-dessus
$$
c(\mathcal{E}^\bullet) \in \prod\nolimits_{p \geq 0} A^p(X),\quad
ch(\mathcal{E}^\bullet) \in
\prod\nolimits_{p \geq 0} A^p(X) \otimes \mathbf{Q},\quad
P_p(\mathcal{E}^\bullet) \in A^p(X)
$$
sont indépendantes du choix du complexe $\mathcal{E}^\bullet$.
\end{lemma}

\begin{proof}
Nous le démontrons pour la classe de Chern totale ; les deux autres cas résultent
des mêmes arguments en utilisant le
Lemme \ref{chow-lemma-chern-character-additive}
à la place du
Lemme \ref{chow-lemma-additivity-chern-classes}.

\medskip\noindent
Comme dans la Remarque \ref{chow-remark-extend-to-finite-locally-free}, afin
de définir la classe de Chern totale $c(\mathcal{E}^\bullet)$,
nous décomposons $X$ en sous-schémas ouverts et fermés
$$
X = \coprod\nolimits_{i \in I} X_i
$$
tels que le rang de $\mathcal{E}^n$ soit constant sur $X_i$ pour
tous $n$ et $i$. (Comme ces rangs sont des fonctions localement constantes
sur $X$, nous pouvons le faire.) Puisque $\mathcal{E}^\bullet$ est localement
borné, nous voyons que seul un nombre fini des faisceaux
$\mathcal{E}^n|_{X_i}$ sont non nuls pour $i$ fixé. Nous pouvons donc
définir
$$
c(\mathcal{E}^\bullet|_{X_i}) =
\prod\nolimits_n c(\mathcal{E}^n|_{X_i})^{(-1)^n}
\in \prod\nolimits_{p \geq 0} A^p(X_i)
$$
comme ci-dessus. Par le Lemme \ref{chow-lemma-disjoint-decomposition-bivariant},
nous avons $A^p(X) = \prod_i A^p(X_i)$. Ainsi, pour tout $p \in \mathbf{Z}$,
nous avons un unique élément $c_p(\mathcal{E}^\bullet) \in A^p(X)$ dont la restriction
est $c_p(\mathcal{E}^\bullet|_{X_i})$ sur $X_i$ pour tout $i$.

\medskip\noindent
Supposons que nous ayons un second complexe localement borné
$\mathcal{F}^\bullet$ de $\mathcal{O}_X$-modules localement libres de type fini
représentant $E$.
Soit $g : Y \to X$ un morphisme localement de type fini avec $Y$ intègre.
Par le Lemme \ref{chow-lemma-bivariant-zero}, il suffit de montrer que
avec $c(g^*\mathcal{E}^\bullet) \cap [Y]$ est le même que
$c(g^*\mathcal{F}^\bullet) \cap [Y]$, et il suffit même de le démontrer
après avoir remplacé $Y$ par un schéma intègre propre et birationnel
sur $Y$. Nous concluons alors d'abord que $g^*\mathcal{E}^\bullet$
et $g^*\mathcal{F}^\bullet$ sont des complexes bornés de
$\mathcal{O}_Y$-modules localement libres de type fini de rang constant. Ensuite, par
Plus sur la platitude, Lemme \ref{flat-lemma-blowup-complex-integral},
nous pouvons supposer que $H^i(Lg^*E)$ est parfait de dimension de Tor $\leq 1$
pour tout $i \in \mathbf{Z}$.
Cela nous ramène au cas traité dans le paragraphe suivant.

\medskip\noindent
Supposons que $X$ soit intègre, que $\mathcal{E}^\bullet$ et $\mathcal{F}^\bullet$
soient des complexes bornés de modules localement libres de type fini de rang constant, et que
$H^i(E)$ soit un $\mathcal{O}_X$-module parfait de dimension de Tor $\leq 1$
pour tout $i \in \mathbf{Z}$. Nous devons
montrer que $c(\mathcal{E}^\bullet) \cap [X]$ est le même que
$c(\mathcal{F}^\bullet) \cap [X]$. Notons
$d_\mathcal{E}^i : \mathcal{E}^i \to \mathcal{E}^{i + 1}$ et
$d_\mathcal{F}^i : \mathcal{F}^i \to \mathcal{F}^{i + 1}$
les différentielles de nos complexes. Par
Plus sur la platitude, Remarque \ref{flat-remark-when-you-have-a-complex},
nous savons que $\Im(d_\mathcal{E}^i)$, $\Ker(d_\mathcal{E}^i)$,
$\Im(d_\mathcal{F}^i)$ et $\Ker(d_\mathcal{F}^i)$
sont des $\mathcal{O}_X$-modules localement libres de type fini pour tout $i$.
Par additivité (Lemme \ref{chow-lemma-additivity-chern-classes}), nous voyons que
$$
c(\mathcal{E}^\bullet) = \prod\nolimits_i
c(\Ker(d_\mathcal{E}^i))^{(-1)^i} c(\Im(d_\mathcal{E}^i))^{(-1)^i}
$$
et de même pour $\mathcal{F}^\bullet$. Puisque nous avons les
suites exactes courtes
$$
0 \to \Im(d_\mathcal{E}^i) \to \Ker(d_\mathcal{E}^i) \to H^i(E) \to 0
\quad\text{et}\quad
0 \to \Im(d_\mathcal{F}^i) \to \Ker(d_\mathcal{F}^i) \to H^i(E) \to 0
$$
nous nous ramenons au problème énoncé et résolu dans le paragraphe suivant.

\medskip\noindent
Supposons que $X$ soit intègre et que nous ayons deux suites exactes courtes
$$
0 \to \mathcal{E}' \to \mathcal{E} \to \mathcal{Q} \to 0
\quad\text{et}\quad
0 \to \mathcal{F}' \to \mathcal{F} \to \mathcal{Q} \to 0
$$
où $\mathcal{E}$, $\mathcal{E}'$, $\mathcal{F}$, $\mathcal{F}'$ sont
localement libres de type fini. Problème : montrer que
$c(\mathcal{E})c(\mathcal{E}')^{-1} \cap [X] =
c(\mathcal{F})c(\mathcal{F}')^{-1} \cap [X]$.
Pour ce faire, considérons la suite exacte courte
$$
0 \to \mathcal{G} \to \mathcal{E} \oplus \mathcal{F} \to \mathcal{Q} \to 0
$$
qui définit $\mathcal{G}$. Puisque $\mathcal{Q}$ est de dimension de Tor $\leq 1$,
nous voyons que $\mathcal{G}$ est localement libre de type fini. Une chasse au diagramme
montre que le noyau de la surjection $\mathcal{G} \to \mathcal{F}$
s'envoie isomorphiquement sur $\mathcal{E}'$ dans $\mathcal{E}$ et
que le noyau de la surjection $\mathcal{G} \to \mathcal{E}$ s'envoie
isomorphiquement sur $\mathcal{F}'$ dans $\mathcal{F}$. (En travaillant localement
sur des ouverts affines, cela résulte du lemme de Schanuel, ou lui est équivalent, voir
Algèbre, Lemme \ref{algebra-lemma-Schanuel}.)
Nous concluons que
$$
c(\mathcal{E})c(\mathcal{F}') = c(\mathcal{G}) =
c(\mathcal{F})c(\mathcal{E}')
$$
comme souhaité.
\end{proof}

\begin{lemma}
\label{chow-lemma-defined-by-envelope}
Soit $(S, \delta)$ comme dans la Situation \ref{chow-situation-setup}.
Soit $X$ localement de type fini sur $S$. Soit $E \in D(\mathcal{O}_X)$
un objet parfait. Supposons qu'il existe une enveloppe
$f : Y \to X$ (Définition \ref{chow-definition-envelope})
telle que $Lf^*E$ soit isomorphe dans $D(\mathcal{O}_Y)$
à un complexe localement borné $\mathcal{E}^\bullet$ de
$\mathcal{O}_Y$-modules localement libres de type fini. Alors il existe uniques des classes bivariantes
$c(E) \in \prod_{p \geq 0} A^p(X)$,
$ch(E) \in \prod_{p \geq 0} A^p(X) \otimes \mathbf{Q}$ et
$P_p(E) \in A^p(X)$, indépendantes du choix de $f : Y \to X$
et de $\mathcal{E}^\bullet$, telles que les restrictions de ces classes
à $Y$ soient égales à $c(\mathcal{E}^\bullet)$,
$ch(\mathcal{E}^\bullet)$ et $P_p(\mathcal{E}^\bullet)$.
\end{lemma}

\begin{proof}
Fixons $p \in \mathbf{Z}$. Nous démontrerons le lemme pour la classe
de Chern $c_p(E) \in A^p(X)$ et omettrons les arguments pour les autres cas.

\medskip\noindent
Soit $g : T \to X$ un morphisme localement de type fini pour lequel
il existe un complexe localement borné $\mathcal{E}^\bullet$ de
$\mathcal{O}_T$-modules localement libres de type fini représentant $Lg^*E$ dans $D(\mathcal{O}_T)$.
La classe bivariante $c_p(\mathcal{E}^\bullet) \in A^p(T)$
est indépendante du choix de $\mathcal{E}^\bullet$ par le
Lemme \ref{chow-lemma-pre-derived-chern-class}.
Notons $c_p(Lg^*E) \in A^p(T)$ cette classe.
Pour tout morphisme ultérieur $h : T' \to T$ qui est localement
de type fini, en posant $g' = g \circ h$, nous voyons que
$L(g')^*E = L(g \circ h)^*E = Lh^*Lg^*E$ est représenté par
$h^*\mathcal{E}^\bullet$ dans $D(\mathcal{O}_{T'})$.
Nous concluons que $c_p(L(g')^*E)$ a un sens et est égal à la
restriction (Remarque \ref{chow-remark-restriction-bivariant})
de $c_p(Lg^*E)$ à $T'$ (strictement parlant, cela requiert une application du
Lemme \ref{chow-lemma-cap-cp-bivariant}).

\medskip\noindent
Soient $f : Y \to X$ et $\mathcal{E}^\bullet$ comme dans l'énoncé
du lemme. Nous obtenons une classe bivariante $c_p(E) \in A^p(X)$ par
une application du Lemme \ref{chow-lemma-envelope-bivariant}
à $f : Y \to X$ et à la classe $c' = c_p(Lf^*E)$ que nous avons construite
dans le paragraphe précédent. L'hypothèse du lemme est
satisfaite car, d'après la discussion du paragraphe précédent, nous avons
$res_1(c') = c_p(Lg^*E) = res_2(c')$, où
$g = f \circ p = f \circ q : Y \times_X Y \to X$.

\medskip\noindent
Enfin, supposons que $f' : Y' \to X$ soit une seconde enveloppe
telle que $L(f')^*E$ soit représenté par un complexe borné de
$\mathcal{O}_{Y'}$-modules localement libres de type fini. Il s'ensuit
que les restrictions de $c_p(Lf^*E)$ et $c_p(L(f')^*E)$
à $Y \times_X Y'$ sont égales. Puisque $Y \times_X Y' \to X$
est une enveloppe (Lemmes \ref{chow-lemma-base-change-envelope} et
\ref{chow-lemma-composition-envelope}), nous voyons que nos deux candidats pour $c_p(E)$
coïncident par l'unicité dans le Lemme \ref{chow-lemma-envelope-bivariant}.
\end{proof}

\begin{definition}
\label{chow-definition-defined-on-perfect}
Soit $(S, \delta)$ comme dans la Situation \ref{chow-situation-setup}.
Soit $X$ localement de type fini sur $S$. Soit $E \in D(\mathcal{O}_X)$
un objet parfait.
\begin{enumerate}
\item Nous disons que les {\it classes de Chern de $E$ sont définies}\footnote{Voir le
Lemme \ref{chow-lemma-chern-classes-defined} pour quelques critères.} s'il existe
une enveloppe $f : Y \to X$ telle que $Lf^*E$ soit isomorphe dans
$D(\mathcal{O}_Y)$ à un complexe localement borné de
$\mathcal{O}_Y$-modules localement libres de type fini.
\item Si les classes de Chern de $E$ sont définies, alors nous définissons
$$
c(E) \in \prod\nolimits_{p \geq 0} A^p(X),\quad
ch(E) \in
\prod\nolimits_{p \geq 0} A^p(X) \otimes \mathbf{Q},\quad
P_p(E) \in A^p(X)
$$
par une application du Lemme \ref{chow-lemma-defined-by-envelope}.
\end{enumerate}
\end{definition}

\noindent
Cette définition s'applique dans beaucoup, mais non dans toutes, les situations envisagées
dans ce chapitre, voir le Lemme \ref{chow-lemma-chern-classes-defined}.
Peut-être existe-t-il une construction élémentaire de ces classes bivariantes pour
$E/X/(S,\delta)$ général comme dans la définition ; nous l'ignorons.

\begin{lemma}
\label{chow-lemma-chern-classes-defined}
Soit $(S, \delta)$ comme dans la Situation \ref{chow-situation-setup}.
Soit $X$ localement de type fini sur $S$. Soit $E \in D(\mathcal{O}_X)$
un objet parfait. Si l'une des conditions suivantes est satisfaite, alors
les classes de Chern de $E$ sont définies :
\begin{enumerate}
\item il existe une enveloppe $f : Y \to X$ telle que $Lf^*E$
soit isomorphe dans $D(\mathcal{O}_Y)$ à un complexe localement borné de
$\mathcal{O}_Y$-modules localement libres de type fini,
\item $E$ peut être représenté par un complexe borné de
$\mathcal{O}_X$-modules localement libres de type fini,
\item les composantes irréductibles de $X$ sont quasi-compactes,
\item $X$ est quasi-compact,
\item il existe un morphisme $X \to X'$ de schémas localement de type fini
sur $S$ tel que $E$ soit l'image inverse d'un objet parfait $E'$ sur $X'$
dont les classes de Chern sont définies, ou
\item ajouter davantage ici.
\end{enumerate}
\end{lemma}

\begin{proof}
La condition (1) est exactement la partie (1) de la Définition \ref{chow-definition-defined-on-perfect}.
La condition (2) implique (1).

\medskip\noindent
Comme dans (3), supposons que les composantes irréductibles $X_i$ de $X$ soient quasi-compactes.
Nous considérons $X_i$ comme un sous-schéma fermé intègre réduit de $X$.
Le morphisme $\coprod X_i \to X$ est une enveloppe. Pour tout $i$, il
existe une enveloppe $X'_i \to X_i$ telle que $X'_i$ possède une famille ample
de modules inversibles, voir Plus sur les morphismes, Proposition
\ref{more-morphisms-proposition-envelope-with-resolution-property}.
Observons que $f : Y = \coprod X'_i \to X$ est une enveloppe ; un petit détail est omis.
Par Catégories dérivées de schémas, Lemme
\ref{perfect-lemma-resolution-property-ample-family},
chaque $X'_i$ possède la propriété de résolution.
Ainsi, l'objet parfait $L(f|_{X'_i})^*E$ de $D(\mathcal{O}_{X'_i})$
peut être représenté par un complexe borné
de $\mathcal{O}_{X'_i}$-modules localement libres de type fini, voir
Catégories dérivées de schémas, Lemme
\ref{perfect-lemma-resolution-property-perfect-complex}.
Cela démontre que (3) implique (1).

\medskip\noindent
La partie (4) implique (3).

\medskip\noindent
Soient $g : X \to X'$ et $E'$ comme dans la partie (5). Alors il existe une
enveloppe $f' : Y' \to X'$ telle que $L(f')^*E'$ soit représenté par un
complexe localement borné $(\mathcal{E}')^\bullet$ de $\mathcal{O}_{Y'}$-modules.
Alors le changement de base $f : Y \to X$ est une enveloppe par le
Lemme \ref{chow-lemma-base-change-envelope}. De plus, la tiragee en arrière
$\mathcal{E}^\bullet = g^*(\mathcal{E}')^\bullet$ représente $Lf^*E$
et nous voyons que les classes de Chern de $E$ sont définies.
\end{proof}

\begin{lemma}
\label{chow-lemma-chern-classes-computed}
Soit $(S, \delta)$ comme dans la Situation \ref{chow-situation-setup}.
Soit $X$ localement de type fini sur $S$. Soit $E \in D(\mathcal{O}_X)$
un objet parfait. Supposons que les classes de Chern de $E$ soient définies.
Pour $g : W \to X$ localement de type fini avec $W$ intègre, il existe
un diagramme commutatif
$$
\xymatrix{
W' \ar[rd]_{g'} \ar[rr]_b & & W \ar[ld]^g \\
& X
}
$$
avec $W'$ intègre et $b : W' \to W$ propre birationnel tel que $L(g')^*E$
soit représenté par un complexe borné $\mathcal{E}^\bullet$ de
$\mathcal{O}_{W'}$-modules localement libres de rang constant et que nous ayons
$res(c_p(E)) = c_p(\mathcal{E}^\bullet)$ dans $A^p(W')$.
\end{lemma}

\begin{proof}
Choisissons une enveloppe $f : Y \to X$ telle que $Lf^*E$ soit isomorphe dans
$D(\mathcal{O}_Y)$ à un complexe localement borné $\mathcal{E}^\bullet$
de $\mathcal{O}_Y$-modules localement libres de type fini. Le changement de base
$Y \times_X W \to W$ de $f$ est une enveloppe par le
Lemme \ref{chow-lemma-base-change-envelope}. Choisissons un point
$\xi \in Y \times_X W$ s'envoyant sur le point générique de $W$
et ayant le même corps résiduel. Considérons le sous-schéma fermé intègre
$W' \subset Y \times_X W$ de point générique $\xi$. La restriction
de la projection $Y \times_X W \to W$ à $W'$ est un morphisme propre
birationnel $b : W' \to W$. Posons $g' = g \circ b$. Enfin, considérons le complexe
image inverse $(W' \to Y)^*\mathcal{E}^\bullet$. C'est un complexe localement borné
de modules localement libres de type fini sur $W'$. Puisque $W'$ est intègre,
il s'ensuit qu'il est borné et que ses termes sont de rang constant.
Enfin, par construction, $(W' \to Y)^*\mathcal{E}^\bullet$
représente $L(g')^*E$ et, par construction, sa $p$-ième classe de Chern
donne la restriction de $c_p(E)$ par $W' \to X$. Cela achève la démonstration.
\end{proof}

\begin{lemma}
\label{chow-lemma-commutative-chern-perfect}
Soit $(S, \delta)$ comme dans la Situation \ref{chow-situation-setup}.
Soit $X$ localement de type fini sur $S$.
Soit $E \in D(\mathcal{O}_X)$ parfait. Si les classes de Chern
de $E$ sont définies, alors
\begin{enumerate}
\item $c_p(E)$ appartient au centre de l'algèbre $A^*(X)$, et
\item si $g : X' \to X$ est localement de type fini et $c \in A^*(X' \to X)$,
alors $c \circ c_p(E) = c_p(Lg^*E) \circ c$.
\end{enumerate}
\end{lemma}

\begin{proof}
La partie (1) résulte immédiatement de la partie (2). Soient $g : X' \to X$ et
$c \in A^*(X' \to X)$ comme dans (2). Pour montrer que
$c \circ c_p(E) - c_p(Lg^*E) \circ c = 0$, nous utilisons le critère du
Lemme \ref{chow-lemma-bivariant-zero}. Nous pouvons donc supposer que $X$
est intègre et, par le Lemme \ref{chow-lemma-chern-classes-computed},
nous pouvons même supposer que $E$ est représenté
par un complexe borné $\mathcal{E}^\bullet$
de $\mathcal{O}_X$-modules localement libres de type fini de rang constant.
Nous devons alors montrer que
$$
c \cap c_p(\mathcal{E}^\bullet) \cap [X] =
c_p(\mathcal{E}^\bullet) \cap c \cap [X] 
$$
dans $\CH_*(X')$. Cela résulte immédiatement du
Lemme \ref{chow-lemma-cap-commutative-chern} et de la construction
de $c_p(\mathcal{E}^\bullet)$ comme polynôme en les
classes de Chern des modules localement libres $\mathcal{E}^n$.
\end{proof}

\begin{lemma}
\label{chow-lemma-additivity-on-perfect}
Soit $(S, \delta)$ comme dans la Situation \ref{chow-situation-setup}.
Soit $X$ localement de type fini sur $S$. Soit
$$
E_1 \to E_2 \to E_3 \to E_1[1]
$$
un triangle distingué d'objets parfaits dans $D(\mathcal{O}_X)$.
Si l'une des conditions suivantes est satisfaite :
\begin{enumerate}
\item il existe une enveloppe $f : Y \to X$ telle que
$Lf^*E_1 \to Lf^*E_2$ puisse être représenté par une application de complexes
localement bornés de $\mathcal{O}_Y$-modules localement libres de type fini,
\item $E_1 \to E_2$ peut être représenté par une application de complexes localement bornés
de $\mathcal{O}_X$-modules localement libres de type fini,
\item les composantes irréductibles de $X$ sont quasi-compactes,
\item $X$ est quasi-compact, ou
\item ajouter davantage ici,
\end{enumerate}
alors les classes de Chern de $E_1$, $E_2$, $E_3$ sont définies et nous avons
$c(E_2) = c(E_1) c(E_3)$, $ch(E_2) = ch(E_1) + ch(E_3)$ et
$P_p(E_2) = P_p(E_1) + P_p(E_3)$.
\end{lemma}

\begin{proof}
Soit $f : Y \to X$ une enveloppe et soit
$\alpha^\bullet : \mathcal{E}_1^\bullet \to \mathcal{E}_2^\bullet$
une application de complexes localement bornés de
$\mathcal{O}_Y$-modules localement libres de type fini représentant $Lf^*E_1 \to Lf^*E_2$.
Alors le cône $C(\alpha)^\bullet$ représente $Lf^*E_3$. Puisque
$C(\alpha)^n = \mathcal{E}_2^n \oplus \mathcal{E}_1^{n + 1}$,
nous voyons que $C(\alpha)^\bullet$ est un complexe localement borné
de $\mathcal{O}_Y$-modules localement libres de type fini.
Nous concluons que les classes de Chern de $E_1$, $E_2$, $E_3$
sont définies. De plus, rappelons que $c_p(E_1)$ est défini
comme l'unique élément de $A^p(X)$ dont la restriction est
$c_p(\mathcal{E}_1^\bullet)$ dans $A^p(Y)$. De même pour
$E_2$ et $E_3$. Il suffit donc
de démontrer $c(\mathcal{E}_2^\bullet) =
c(\mathcal{E}_1^\bullet) c(C(\alpha)^\bullet)$
dans $\prod_{p \geq 0} A^p(Y)$.
À son tour, il suffit de le démontrer après restriction
à une composante connexe de $Y$. Nous pouvons donc supposer que
les complexes $\mathcal{E}_1^\bullet$ et $\mathcal{E}_2^\bullet$
sont des complexes bornés de $\mathcal{O}_Y$-modules localement libres de type fini
de rang fixé. Dans ce cas, l'égalité souhaitée résulte de la
multiplicativité du Lemme \ref{chow-lemma-additivity-chern-classes}.
Dans le cas de $ch$ ou $P_p$, nous utilisons les
Lemmes \ref{chow-lemma-chern-character-additive}.

\medskip\noindent
Dans le paragraphe précédent, nous avons vu que le lemme est vrai si
la condition (1) est satisfaite. Puisque (2) implique (1), cela traite
le deuxième cas. Supposons (3). En raisonnant exactement comme dans la démonstration du
Lemme \ref{chow-lemma-chern-classes-defined}, nous trouvons une enveloppe
$f : Y \to X$ tel que $Y$ soit une réunion disjointe $Y = \coprod Y_i$
de schémas quasi-compacts (et quasi-séparés), chacun possédant la
propriété de résolution. Nous pouvons alors représenter la restriction de
$Lf^*E_1 \to Lf^*E_2$ à $Y_i$ par un morphisme de complexes bornés
de modules localement libres de rang fini ; voir
Catégories dérivées de schémas, proposition
\ref{perfect-proposition-perfect-resolution-property}.
Nous voyons ainsi que la condition (3) implique la condition (1).
Bien entendu, la condition (4) implique la condition (3), ce qui achève
la démonstration.
\end{proof}

\begin{remark}
\label{chow-remark-splitting-principle-perfect}
Les classes de Chern d'un complexe parfait, lorsqu'elles sont définies, satisfont une sorte de
principe de scindage. Plus précisément, supposons que $(S, \delta), X, E$ soient comme dans la
définition \ref{chow-definition-defined-on-perfect}
et que les classes de Chern de $E$ soient définies.
Supposons que nous voulions démontrer une relation entre les classes bivariantes
$c_p(E)$, $P_p(E)$ et $ch_p(E)$. Pour cela, nous pouvons choisir une
enveloppe $f : Y \to X$ et un complexe localement borné
$\mathcal{E}^\bullet$ de $\mathcal{O}_X$-modules localement libres de rang fini
représentant $E$. Par l'unicité dans le lemme \ref{chow-lemma-defined-by-envelope},
il suffit de démontrer la relation voulue entre les classes bivariantes
$c_p(\mathcal{E}^\bullet)$, $P_p(\mathcal{E}^\bullet)$ et
$ch_p(\mathcal{E}^\bullet)$. Nous pouvons donc remplacer $X$ par une composante
connexe de $Y$ et supposer que $E$ est représenté par un complexe borné
$\mathcal{E}^\bullet$ de modules localement libres de rang constant et fini.
En utilisant le principe de scindage
(lemme \ref{chow-lemma-splitting-principle}), nous pouvons supposer que chaque
$\mathcal{E}^i$ possède une filtration dont les quotients
successifs $\mathcal{L}_{i, j}$ sont des modules inversibles.
En posant $x_{i, j} = c_1(\mathcal{L}_{i, j})$, nous voyons que
$$
c(E) =
\prod\nolimits_{i\text{ pair}} (1 + x_{i, j})
\prod\nolimits_{i\text{ impair}} (1 + x_{i, j})^{-1}
$$
et
$$
P_p(E) =  \sum\nolimits_{i\text{ pair}} (x_{i, j})^p -
\sum\nolimits_{i\text{ impair}} (x_{i, j})^p
$$
En prenant formellement le logarithme de l'expression de $c(E)$ ci-dessus,
nous obtenons
$$
\log(c(E)) = \sum (-1)^{p - 1}\frac{P_p(E)}{p}
$$
La construction des polynômes $P_p$ dans
l'exemple \ref{chow-example-power-sum} montre que $P_p(E)$
est exactement la même expression dans les classes de Chern de $E$
que dans le cas des fibrés vectoriels ; autrement dit, nous avons
\begin{align*}
P_1(E) & = c_1(E), \\
P_2(E) & = c_1(E)^2 - 2c_2(E), \\
P_3(E) & = c_1(E)^3 - 3c_1(E)c_2(E) + 3c_3(E), \\
P_4(E) & = c_1(E)^4 - 4c_1(E)^2c_2(E) + 4c_1(E)c_3(E) + 2c_2(E)^2 - 4c_4(E),
\end{align*}
et ainsi de suite. D'autre part, la classe bivariante $P_0(E) = r(E) = ch_0(E)$
ne peut être retrouvée à partir de la classe de Chern $c(E)$ de $E$ ; la classe de Chern
ne contient pas l'information sur le rang du complexe.
\end{remark}

\begin{lemma}
\label{chow-lemma-chern-classes-perfect-dual}
Dans la situation \ref{chow-situation-setup}, soit $X$ localement de type fini sur $S$.
Soit $E \in D(\mathcal{O}_X)$ un objet parfait dont les classes de Chern sont
définies. Alors $c_i(E^\vee) = (-1)^i c_i(E)$, $P_i(E^\vee) = (-1)^iP_i(E)$
et $ch_i(E^\vee) = (-1)^ich_i(E)$ dans $A^i(X)$.
\end{lemma}

\begin{proof}
Première démonstration : raisonnons comme dans la démonstration du
lemme \ref{chow-lemma-commutative-chern-perfect}
pour nous ramener au cas où $E$ est représenté
par un complexe borné de modules localement libres de rang
constant et fini, puis appliquons le lemme \ref{chow-lemma-chern-classes-dual}.
Seconde démonstration : utilisons le principe de scindage exposé
dans la remarque \ref{chow-remark-splitting-principle-perfect}
et le fait que les racines de Chern de $E^\vee$ sont les opposées
des racines de Chern de $E$.
\end{proof}

\begin{lemma}
\label{chow-lemma-chern-class-perfect-tensor-invertible}
Dans la situation \ref{chow-situation-setup}, soit $X$ localement de type fini sur $S$.
Soit $E$ un objet parfait de $D(\mathcal{O}_X)$ dont les classes de Chern
sont définies.
Soit $\mathcal{L}$ un $\mathcal{O}_X$-module inversible. Alors
$$
c_i(E \otimes \mathcal{L}) =
\sum\nolimits_{j = 0}^i
\binom{r - i + j}{j} c_{i - j}(E) c_1(\mathcal{L})^j
$$
pourvu que $E$ soit de rang constant $r \in \mathbf{Z}$.
\end{lemma}

\begin{proof}
Lorsque $E$ est localement libre de rang $r$, c'est le
lemme \ref{chow-lemma-chern-classes-E-tensor-L}. Le lecteur peut déduire
le lemme de ce cas particulier par un calcul formel.
On peut aussi utiliser le principe de scindage de la
remarque \ref{chow-remark-splitting-principle-perfect}.
Dans ce cas, on est ramené à démontrer le fait
d'algèbre suivant : si nous écrivons formellement
$$
\frac{\prod_{a = 1, \ldots, n} (1 + x_a)}{\prod_{n = 1, \ldots, m} (1 + y_b)}
= 1 + c_1 + c_2 + c_3 + \ldots
$$
où $c_i$ est homogène de degré $i$
dans $\mathbf{Z}[x_i, y_j]$, alors nous avons
$$
\frac{\prod_{a = 1, \ldots, n} (1 + x_a + t)}{\prod_{b = 1, \ldots, m} (1 + y_b + t)}
= \sum\nolimits_{i \geq 0} \sum\nolimits_{j = 0}^i
\binom{r - i + j}{j} c_{i - j} t^j
$$
où $r = n - m$. Nous omettons les détails.
\end{proof}

\begin{lemma}
\label{chow-lemma-chern-classes-perfect-tensor-product}
Dans la situation \ref{chow-situation-setup}, soit $X$ localement de type fini sur $S$.
Soient $E$ et $F$ des objets parfaits de $D(\mathcal{O}_X)$ dont les classes de Chern
sont définies. Alors nous avons
$$
c_1(E \otimes_{\mathcal{O}_X}^\mathbf{L} F) =
r(E) c_1(\mathcal{F}) + r(F) c_1(\mathcal{E})
$$
et, pour $c_2(E \otimes_{\mathcal{O}_X}^\mathbf{L} F)$, l'expression est
$$
r(E) c_2(F) + r(F) c_2(E) + {r(E) \choose 2} c_1(F)^2 +
(r(E)r(F) - 1) c_1(F)c_1(E) + {r(F) \choose 2} c_1(E)^2
$$
et ainsi de suite pour les classes de Chern supérieures dans $A^*(X)$. De même, nous avons
$ch(E \otimes_{\mathcal{O}_X}^\mathbf{L} F) = ch(E) ch(F)$
dans $A^*(X) \otimes \mathbf{Q}$. Plus précisément, nous avons
$$
P_p(E \otimes_{\mathcal{O}_X}^\mathbf{L} F) = \sum\nolimits_{p_1 + p_2 = p}
{p \choose p_1} P_{p_1}(E) P_{p_2}(F)
$$
dans $A^p(X)$.
\end{lemma}

\begin{proof}
Après avoir choisi une enveloppe $f : Y \to X$ telle que $Lf^*E$ et $Lf^*F$
puissent être représentés par des complexes localement bornés de
$\mathcal{O}_X$-modules localement libres de rang fini, ceci résulte, par calcul, du
résultat correspondant pour les fibrés vectoriels dans les
lemmes \ref{chow-lemma-chern-classes-tensor-product} et
\ref{chow-lemma-chern-character-multiplicative}.
Une meilleure démonstration consiste probablement à utiliser le principe de scindage comme dans la
remarque \ref{chow-remark-splitting-principle-perfect}
et à réduire le lemme à des calculs dans des anneaux de polynômes 
que nous décrivons au paragraphe suivant.

\medskip\noindent
Soit $A$ un anneau commutatif (pour nous, ce sera le sous-anneau de l'anneau de Chow
bivariant de $X$ engendré par les classes de Chern).
Soit $S$ un ensemble fini muni d'applications $\epsilon : S \to \{\pm 1\}$
et $f : S \to A$. Définissons
$$
P_p(S, f , \epsilon) = \sum\nolimits_{s \in S} \epsilon(s) f(s)^p
$$
dans $A$. Étant donné un second triplet $(S', \epsilon', f')$,
l'égalité qu'il faut démontrer pour $P_p$ est
$$
P_p(S \times S', f + f' , \epsilon \epsilon') = 
\sum\nolimits_{p_1 + p_2 = p}
{p \choose p_1} P_{p_1}(S, f, \epsilon) P_{p_2}(S', f', \epsilon')
$$
Pour le voir, on se ramène à l'anneau de polynômes en les variables
$S \amalg S'$ et l'on montre que chaque terme $f(s)^if'(s')^j$ apparaît
dans les deux membres avec le même coefficient.
Pour vérifier les formules de $c_1(E \otimes_{\mathcal{O}_X}^\mathbf{L} F)$
et $c_2(E \otimes_{\mathcal{O}_X}^\mathbf{L} F)$, nous utilisons le principe de scindage
pour nous ramener à vérifier ces formules dans un anneau sans torsion.
Nous utilisons ensuite la relation entre $P_j(E)$ et $c_i(E)$ démontrée
dans la remarque \ref{chow-remark-splitting-principle-perfect}. Par exemple,
$$
c_1(E \otimes F) = P_1(E \otimes F) = r(F)P_1(E) + r(E)P_1(F) =
r(F)c_1(E) + r(E)c_1(F)
$$
l'égalité du milieu résultant de $r(E) = P_0(E)$ par définition. De même, nous avons
\begin{align*}
& 2c_2(E \otimes F) \\
& = c_1(E \otimes F)^2 - P_2(E \otimes F) \\
& =
(r(F)c_1(E) + r(E)c_1(F))^2 -
r(F)P_2(E) - P_1(E)P_1(F) - r(E)P_2(F) \\
& =
(r(F)c_1(E) + r(E)c_1(F))^2 -
r(F)(c_1(E)^2 - 2c_2(E)) - c_1(E)c_1(F) - \\
& \quad r(E)(c_1(F)^2 - 2c_2(F))
\end{align*}
ce que le lecteur peut vérifier être égal à la formule de l'énoncé
du lemme à un facteur $2$ près.
\end{proof}







\section{Un cas élémentaire de classes de Chern localisées}
\label{chow-section-preparation-localized-chern}

\noindent
Dans cette section, nous étudions quelques propriétés des classes bivariantes
construites dans le lemme suivant ; la plupart de ces propriétés résultent
immédiatement de la caractérisation donnée dans le lemme. Nous recommandons vivement au
lecteur de sauter le reste de la section.

\begin{lemma}
\label{chow-lemma-silly}
Soit $(S, \delta)$ comme dans la situation \ref{chow-situation-setup}. Soit $X$
localement de type fini sur $S$. Soient $i_j : X_j \to X$, $j = 1, 2$,
des immersions fermées telles que $X = X_1 \cup X_2$ ensemblistement. Soit
$E_2 \in D(\mathcal{O}_{X_2})$ un objet parfait. Supposons que
\begin{enumerate}
\item les classes de Chern de $E_2$ soient définies ;
\item la restriction $E_2|_{X_1 \cap X_2}$ soit nulle,
resp.\ isomorphe à un $\mathcal{O}_{X_1 \cap X_2}$-module localement libre de rang fini
de rang $< p$ placé en degré cohomologique $0$.
\end{enumerate}
Il existe alors une classe bivariante canonique
$$
P'_p(E_2),\text{ resp. }c'_p(E_2) \in A^p(X_2 \to X)
$$
caractérisée par la propriété
$$
P'_p(E_2) \cap i_{2, *} \alpha_2 = P_p(E_2) \cap \alpha_2
\quad\text{et}\quad
P'_p(E_2) \cap i_{1, *} \alpha_1 = 0,
$$
respectivement
$$
c'_p(E_2) \cap i_{2, *} \alpha_2 = c_p(E_2) \cap \alpha_2
\quad\text{et}\quad
c'_p(E_2) \cap i_{1, *} \alpha_1 = 0
$$
pour $\alpha_i \in \CH_k(X_i)$, et de même après tout changement de base
$X' \to X$ localement de type fini.
\end{lemma}

\begin{proof}
Nous allons utiliser les résultats de la section \ref{chow-section-pre-derived}
sans autre mention.

\medskip\noindent
Supposons que $E_2|_{X_1 \cap X_2}$ soit nul.
Considérons un morphisme de schémas $X' \to X$
localement de type fini et notons $i'_j : X'_j \to X'$ le
changement de base de $i_j$. D'après le lemme \ref{chow-lemma-exact-sequence-closed-chow},
nous pouvons écrire tout élément $\alpha' \in \CH_k(X')$ sous la forme
$i'_{1, *}\alpha'_1 + i'_{2, *}\alpha'_2$, où
$\alpha'_2 \in \CH_k(X'_2)$
est bien défini à un élément de l'image de l'image directe
par $X'_1 \cap X'_2 \to X'_2$ près. Nous pouvons alors poser
$P'_p(E_2) \cap \alpha' = P_p(E_2) \cap \alpha'_2 \in \CH_{k - p}(X'_2)$. Ceci
est bien défini puisque, par hypothèse, la restriction de $E_2$
à $X_1 \cap X_2$ est nulle.

\medskip\noindent
Si $E_2|_{X_1 \cap X_2}$ est isomorphe à un module localement libre de rang fini
sur $\mathcal{O}_{X_1 \cap X_2}$, de rang $< p$ et placé en
degré cohomologique $0$, alors $c_p(E_2|_{X_1 \cap X_2}) = 0$
pour des raisons de rang et nous pouvons raisonner exactement de la même manière.
\end{proof}

\begin{lemma}
\label{chow-lemma-silly-independent}
Dans le lemme \ref{chow-lemma-silly}, la classe bivariante
$P'_p(E_2)$, resp.\ $c'_p(E_2)$ dans $A^p(X_2 \to X)$
ne dépend pas du choix de $X_1$.
\end{lemma}

\begin{proof}
Supposons que $X_1' \subset X$ soit un autre sous-schéma fermé tel que
$X = X'_1 \cup X_2$ ensemblistement et que la restriction
$E_2|_{X'_1 \cap X_2}$ soit nulle, resp.\ isomorphe à un
$\mathcal{O}_{X'_1 \cap X_2}$-module localement libre de rang fini
de rang $< p$ placé en degré cohomologique $0$.
Alors $X = (X_1 \cap X'_1) \cup X_2$. Nous pouvons donc écrire
tout élément $\alpha \in \CH_k(X)$ sous la forme $i_*\beta + i_{2, *}\alpha_2$, avec
$\alpha_2 \in \CH_k(X'_2)$ et $\beta \in \CH_k(X_1 \cap X'_1)$.
Il est donc clair que
$P'_p(E_2) \cap \alpha = P_p(E_2) \cap \alpha_2 \in \CH_{k - p}(X_2)$,
resp.\  $c'_p(E_2) \cap \alpha = c_p(E_2) \cap \alpha_2 \in \CH_{k - p}(X_2)$,
ne dépend pas de l'emploi de $X_1$ ou de $X'_1$. Il en va de même
après tout changement de base.
\end{proof}

\begin{lemma}
\label{chow-lemma-base-change-silly}
Dans le lemme \ref{chow-lemma-silly}, soit $X' \to X$ un morphisme
localement de type fini. Notons $X' = X'_1 \cup X'_2$
et $E'_2 \in D(\mathcal{O}_{X'_2})$ les images inverses sur $X'$.
Alors la classe $P'_p(E_2')$, resp.\ $c'_p(E_2')$ dans
$A^p(X_2' \to X')$, construite dans le lemme \ref{chow-lemma-silly} à partir de
$X' = X'_1 \cup X'_2$ et de $E_2'$, est la restriction
(remarque \ref{chow-remark-restriction-bivariant})
de la classe $P'_p(E_2)$, resp.\ $c'_p(E_2)$ dans $A^p(X_2 \to X)$.
\end{lemma}

\begin{proof}
Cela résulte immédiatement de la caractérisation de ces classes dans le
lemme \ref{chow-lemma-silly}.
\end{proof}

\begin{lemma}
\label{chow-lemma-silly-silly}
Dans le lemme \ref{chow-lemma-silly}, supposons que $E_2$ soit la restriction d'un objet
parfait $E \in D(\mathcal{O}_X)$ tel que $E|_{X_1}$ soit nul,
resp.\ isomorphe à un $\mathcal{O}_{X_1}$-module localement libre de rang fini
de rang $< p$ placé en degré cohomologique $0$.
Si les classes de Chern de $E$ sont définies, alors
$i_{2, *} \circ P'_p(E_2) = P_p(E)$,
resp.\ $i_{2, *} \circ c'_p(E_2) = c_p(E)$
(où $\circ$ est comme dans le lemme \ref{chow-lemma-push-proper-bivariant}).
\end{lemma}

\begin{proof}
Supposons d'abord que $E|_{X_1}$ soit nul.
Avec les notations de la démonstration du lemme \ref{chow-lemma-silly},
le lemme résulte dans ce cas de
\begin{align*}
P_p(E) \cap \alpha'
& =
i'_{1, *}(P_p(E) \cap \alpha'_1) +
i'_{2, *}(P_p(E) \cap \alpha'_2) \\
& =
i'_{1, *}(P_p(E|_{X_1}) \cap \alpha'_1) +
i'_{2, *}(P'_p(E_2) \cap \alpha') \\
& =
i'_{2, *}(P'_p(E_2) \cap \alpha')
\end{align*}
Le cas où $E|_{X_1}$ est isomorphe à un
$\mathcal{O}_{X_1}$-module localement libre de rang fini, de rang $< p$ et placé en degré cohomologique $0$,
est analogue.
\end{proof}

\begin{lemma}
\label{chow-lemma-silly-shrink}
Dans le lemme \ref{chow-lemma-silly}, supposons donnés des sous-schémas fermés
$X'_2 \subset X_2$ et $X_1 \subset X'_1 \subset X$ tels que
$X = X'_1 \cup X'_2$ ensemblistement. Supposons que $E_2|_{X'_1 \cap X_2}$
soit nul, resp.\ isomorphe à un module localement libre de rang fini
de rang $< p$ placé en degré $0$. Alors nous avons
$(X'_2 \to X_2)_* \circ P'_p(E_2|_{X'_2}) = P'_p(E_2)$,
resp.\  $(X'_2 \to X_2)_* \circ c'_p(E_2|_{X'_2}) = c_p(E_2)$
(où $\circ$ est comme dans le lemme \ref{chow-lemma-push-proper-bivariant}).
\end{lemma}

\begin{proof}
Cela résulte immédiatement de la caractérisation de ces classes
dans le lemme \ref{chow-lemma-silly}.
\end{proof}

\begin{lemma}
\label{chow-lemma-silly-commutes}
Dans le lemme \ref{chow-lemma-silly}, soit $f : Y \to X$ localement de type fini
et soit $c \in A^*(Y \to X)$. Alors
$$
c \circ P'_p(E_2) = P'_p(Lf_2^*E_2) \circ c
\quad\text{resp.}\quad
c \circ c'_p(E_2) = c'_p(Lf_2^*E_2) \circ c
$$
dans $A^*(Y_2 \to Y)$, où $f_2 : Y_2 \to X_2$ est le changement de base de $f$.
\end{lemma}

\begin{proof}
Soit $\alpha \in \CH_k(X)$. Nous pouvons écrire
$$
\alpha = \alpha_1 + \alpha_2
$$
avec $\alpha_i \in \CH_k(X_i)$ ; nous omettons les images directes
par les immersions fermées $X_i \to X$. Le lecteur vérifie alors que
$c'_p(E_2) \cap \alpha = c_p(E_2) \cap \alpha_2$,
$c \cap c'_p(E_2) \cap \alpha = c \cap c_p(E_2) \cap \alpha_2$,
$c \cap \alpha = c \cap \alpha_1 + c \cap \alpha_2$, et
$c'_p(Lf_2^*E_2) \cap c \cap \alpha = c_p(Lf_2^*E_2) \cap c \cap \alpha_2$.
Nous concluons par le lemme \ref{chow-lemma-commutative-chern-perfect}.
\end{proof}

\begin{lemma}
\label{chow-lemma-silly-compose}
Dans le lemme \ref{chow-lemma-silly}, supposons que $E_2|_{X_1 \cap X_2}$ soit nul. Alors
\begin{align*}
P'_1(E_2) & = c'_1(E_2), \\
P'_2(E_2) & = c'_1(E_2)^2 - 2c'_2(E_2), \\
P'_3(E_2) & = c'_1(E_2)^3 - 3c'_1(E_2)c'_2(E_2) + 3c'_3(E_2), \\
P'_4(E_2) & = c'_1(E_2)^4 - 4c'_1(E_2)^2c'_2(E_2) +
4c'_1(E_2)c'_3(E_2) + 2c'_2(E_2)^2 - 4c'_4(E_2),
\end{align*}
et ainsi de suite, la multiplication étant celle de la remarque \ref{chow-remark-ring-loc-classes}.
\end{lemma}

\begin{proof}
L'énoncé a un sens car le faisceau nul est de rang $< 1$ et
les classes $c'_p(E_2)$ sont donc définies pour tout $p \geq 1$. Les égalités
résultent immédiatement de la caractérisation des classes construites
par le lemme \ref{chow-lemma-silly} et du résultat correspondant pour
l'intersection avec les classes de Chern de $E_2$ donné dans la
remarque \ref{chow-remark-splitting-principle-perfect}.
\end{proof}

\begin{lemma}
\label{chow-lemma-silly-sum-c}
Soit $(S, \delta)$ comme dans la situation \ref{chow-situation-setup}. Soit $X$
localement de type fini sur $S$. Soient $i_j : X_j \to X$, $j = 1, 2$,
des immersions fermées telles que $X = X_1 \cup X_2$ ensemblistement. Soient
$E, F \in D(\mathcal{O}_X)$ des objets parfaits. Supposons que
\begin{enumerate}
\item les classes de Chern de $E$ et de $F$ soient définies ;
\item les restrictions $E|_{X_1 \cap X_2}$ et $F|_{X_1 \cap X_2}$
soient isomorphes à deux $\mathcal{O}_{X_1}$-modules localement libres,
de rangs respectivement $< p$ et $< q$, placés en degré cohomologique $0$.
\end{enumerate}
Avec la notation de la remarque \ref{chow-remark-ring-loc-classes}, posons
$$
c^{(p)}(E) = 1 + c_1(E) + \ldots + c_{p - 1}(E) +
c'_p(E|_{X_2}) + c'_{p + 1}(E|_{X_2}) + \ldots \in A^{(p)}(X_2 \to X)
$$
où $c'_p(E|_{X_2})$ est comme dans le lemme \ref{chow-lemma-silly}. De même
pour $c^{(q)}(F)$ et $c^{(p + q)}(E \oplus F)$.
Alors $c^{(p + q)}(E \oplus F) = c^{(p)}(E)c^{(q)}(F)$
dans $A^{(p + q)}(X_2 \to X)$.
\end{lemma}

\begin{proof}
Cela résulte immédiatement de la caractérisation des classes dans le
lemme \ref{chow-lemma-silly} et de l'additivité dans le
lemme \ref{chow-lemma-additivity-on-perfect}.
\end{proof}

\begin{lemma}
\label{chow-lemma-silly-sum-P}
Soit $(S, \delta)$ comme dans la situation \ref{chow-situation-setup}. Soit $X$
localement de type fini sur $S$. Soient $i_j : X_j \to X$, $j = 1, 2$,
des immersions fermées telles que $X = X_1 \cup X_2$ ensemblistement. Soient
$E, F \in D(\mathcal{O}_{X_2})$ des objets parfaits. Supposons que
\begin{enumerate}
\item les classes de Chern de $E$ et de $F$ soient définies ;
\item les restrictions $E|_{X_1 \cap X_2}$ et $F|_{X_1 \cap X_2}$ soient nulles,
\end{enumerate}
Notons $P'_p(E), P'_p(F), P'_p(E \oplus F) \in A^p(X_2 \to X)$, pour $p \geq 0$,
les classes construites dans le lemme \ref{chow-lemma-silly}. Alors
$P'_p(E \oplus F) = P'_p(E) + P'_p(F)$.
\end{lemma}

\begin{proof}
Cela résulte immédiatement de la caractérisation des classes dans le
lemme \ref{chow-lemma-silly} et de l'additivité dans le
lemme \ref{chow-lemma-additivity-on-perfect}.
\end{proof}

\begin{lemma}
\label{chow-lemma-silly-tensor-invertible}
Dans le lemme \ref{chow-lemma-silly}, supposons que $E_2$ soit de rang constant $0$.
Soit $\mathcal{L}$ un $\mathcal{O}_X$-module inversible. Alors
$$
c'_i(E_2 \otimes \mathcal{L}) =
\sum\nolimits_{j = 0}^i
\binom{- i + j}{j} c'_{i - j}(E_2) c_1(\mathcal{L})^j
$$
\end{lemma}

\begin{proof}
L'hypothèse sur le rang implique que $E_2|_{X_1 \cap X_2}$ est nul.
Ainsi $c'_i(E_2)$ est défini pour tout $i \geq 1$ et l'énoncé
a un sens. L'égalité elle-même résulte
immédiatement du lemme \ref{chow-lemma-chern-class-perfect-tensor-invertible}
et de la caractérisation de $c'_i$ dans le lemme \ref{chow-lemma-silly}.
\end{proof}

\begin{lemma}
\label{chow-lemma-silly-tensor-product}
Dans la situation \ref{chow-situation-setup}, soit $X$ localement de type fini sur $S$.
Soit
$$
X = X_1 \cup X_2 = X'_1 \cup X'_2
$$
deux écritures de $X$ comme réunion ensembliste de sous-schémas fermés.
Soient $E$, $E'$ des objets parfaits de $D(\mathcal{O}_X)$
dont les classes de Chern sont définies.
Supposons que $E|_{X_1}$ et $E'|_{X'_1}$ soient nuls\footnote{Il existe vraisemblablement
une variante de ce lemme où l'on suppose seulement que ces restrictions sont
isomorphes à deux modules localement libres,
de rangs respectivement $< p$ et $< p'$.} pour $i = 1, 2$. Notons
\begin{enumerate}
\item $r = P'_0(E) \in A^0(X_2 \to X)$ et
$r' = P'_0(E') \in A^0(X'_2 \to X)$,
\item $\gamma_p = c'_p(E|_{X_2}) \in A^p(X_2 \to X)$ et
$\gamma'_p = c'_p(E'|_{X'_2}) \in A^p(X'_2 \to X)$,
\item $\chi_p = P'_p(E|_{X_2}) \in A^p(X_2 \to X)$ et
$\chi'_p = P'_p(E'|_{X'_2}) \in A^p(X'_2 \to X)$
\end{enumerate}
les classes construites dans le lemme \ref{chow-lemma-silly}. Alors nous avons
$$
c'_1((E \otimes_{\mathcal{O}_X}^\mathbf{L} E')|_{X_2 \cap X'_2}) =
r \gamma'_1 + r' \gamma_1
$$
dans $A^1(X_2 \cap X'_2 \to X)$ et
$$
c'_2((E \otimes_{\mathcal{O}_X}^\mathbf{L} E')|_{X_2 \cap X'_2}) =
r \gamma'_2 + r' \gamma_2 + {r \choose 2} (\gamma'_1)^2 +
(rr' - 1) \gamma'_1\gamma_1 + {r' \choose 2} \gamma_1^2
$$
dans $A^2(X_2 \cap X'_2 \to X)$, et ainsi de suite pour les classes de Chern supérieures.
De même, nous avons
$$
P'_p((E \otimes_{\mathcal{O}_X}^\mathbf{L} E')|_{X_2 \cap X'_2}) =
\sum\nolimits_{p_1 + p_2 = p}
{p \choose p_1} \chi_{p_1} \chi'_{p_2}
$$
dans $A^p(X_2 \cap X'_2 \to X)$.
\end{lemma}

\begin{proof}
Observons d'abord que l'énoncé a un sens. En effet, nous avons
$X = (X_2 \cap X'_2) \cup Y$, où
$Y = (X_1 \cap X'_1) \cup (X_1 \cap X'_2) \cup (X_2 \cap X'_1)$,
et la restriction de l'objet $E \otimes_{\mathcal{O}_X}^\mathbf{L} E'$
à $Y$ est nulle.
Les égalités elles-mêmes résultent de la caractérisation
de nos classes dans le lemme \ref{chow-lemma-silly}
et des égalités du lemme \ref{chow-lemma-chern-classes-perfect-tensor-product}.
Nous omettons les détails.
\end{proof}







\section{Gysin à l'infini}
\label{chow-section-gysin-at-infty}

\noindent
Cette section porte sur la classe bivariante construite dans le prochain
lemme. Nous recommandons vivement au lecteur de sauter le reste de la section.

\begin{lemma}
\label{chow-lemma-gysin-at-infty}
Soit $(S, \delta)$ comme dans la situation \ref{chow-situation-setup}.
Soit $X$ localement de type fini sur $S$. Soit
$b : W \to \mathbf{P}^1_X$ un morphisme propre de schémas
qui est un isomorphisme au-dessus de $\mathbf{A}^1_X$.
Notons $i_\infty : W_\infty \to W$ l'image inverse du diviseur
$D_\infty \subset \mathbf{P}^1_X$ dont le complémentaire est $\mathbf{A}^1_X$.
Il existe alors une classe bivariante canonique
$$
C \in A^0(W_\infty \to X)
$$
ayant la propriété suivante :
$i_{\infty, *}(C \cap \alpha) = i_{0, *}\alpha$
pour $\alpha \in \CH_k(X)$, et de même après tout changement de base par
$X' \to X$ localement de type fini.
\end{lemma}

\begin{proof}
Étant donné $\alpha \in \CH_k(X)$, il existe $\beta \in \CH_{k + 1}(W)$
dont la restriction est l'image inverse plate de $\alpha$ sur $b^{-1}(\mathbf{A}^1_X)$ ; voir
le lemme \ref{chow-lemma-exact-sequence-open}.
Un second choix, encore noté $\beta$, diffère du choix initial $\beta$ par un cycle
à support dans $W_\infty$ ; voir
le lemme \ref{chow-lemma-restrict-to-open}. Puisque le fibré normal du diviseur de Cartier
effectif $D_\infty \subset \mathbf{P}^1_X$ de
(\ref{chow-equation-zero-infty}) est trivial,
l'homomorphisme de Gysin $i_\infty^*$ annule les classes de cycles
à support dans $W_\infty$ ; voir la remarque \ref{chow-remark-gysin-on-cycles}.
Ainsi, poser $C \cap \alpha = i_\infty^*\beta$ donne une définition indépendante des choix.

\medskip\noindent
Puisque $W_\infty$ et $W_0 = X \times \{0\}$
sont les images inverses des diviseurs de Cartier effectifs rationnellement équivalents
$D_0, D_\infty$ dans $\mathbf{P}^1_X$, nous voyons que $i_\infty^*\beta$ et
$i_0^*\beta$ ont la même image comme classe de cycle sur $W$ ; en effet, tous deux
représentent la classe $c_1(\mathcal{O}_{\mathbf{P}^1_X}(1)) \cap \beta$ d'après le
lemme \ref{chow-lemma-support-cap-effective-Cartier}. Par notre choix de
$\beta$, nous avons $i_0^*\beta = \alpha$ comme cycles sur
$W_0 = X \times \{0\}$ ; voir par exemple le
lemme \ref{chow-lemma-relative-effective-cartier}.
Nous voyons donc que $i_{\infty, *}(C \cap \alpha) = i_{0, *}\alpha$,
comme l'affirme le lemme.

\medskip\noindent
Remarquons que les hypothèses sur $b$ sont préservées par tout changement de base
par $X' \to X$ localement de type fini. Nous obtenons donc une opération
$C \cap - : \CH_k(X') \to \CH_k(W'_\infty)$ par la même construction que ci-dessus.
Pour voir que cette famille d'opérations définit une classe bivariante,
considérons le diagramme
$$
\xymatrix{
& & & \CH_*(X) \ar[d]^{\text{image inverse plate}} \\
\CH_{* + 1}(W_\infty) \ar[r] \ar[rd]^0 &
\CH_{* + 1}(W) \ar[d]^{i_\infty^*} \ar[rr]^{\text{image inverse plate}} & &
\CH_{* + 1}(\mathbf{A}^1_X) \ar[r] \ar@{..>}[lld]^{C \cap -} &
0 \\
& \CH_*(W_\infty)
}
$$
pour $X$ comme indiqué, ainsi que le changement de base de ce diagramme pour tout $X' \to X$.
Nous savons que l'image inverse plate et $i_\infty^*$ sont des opérations bivariantes ; voir les
lemmes \ref{chow-lemma-flat-pullback-bivariant} et \ref{chow-lemma-gysin-bivariant}.
Un argument formel (faisant intervenir d'immenses diagrammes de schémas et de leurs
groupes de Chow) montre alors que la flèche pointillée est une opération bivariante.
\end{proof}

\begin{lemma}
\label{chow-lemma-base-change-gysin-at-infty}
Dans le lemme \ref{chow-lemma-gysin-at-infty}, soit $X' \to X$ un morphisme
localement de type fini. Notons $b' : W' \to \mathbf{P}^1_{X'}$
et $i'_\infty : W'_\infty \to W'$ les changements de base de $b$ et de $i_\infty$.
Alors la classe $C' \in A^0(W'_\infty \to X')$ construite comme dans le
lemme \ref{chow-lemma-gysin-at-infty} à l'aide de $b'$ est la restriction
(remarque \ref{chow-remark-restriction-bivariant}) de $C$.
\end{lemma}

\begin{proof}
Cela résulte immédiatement de la construction et du fait qu'un énoncé analogue
vaut pour l'image inverse plate et $i_\infty^*$.
\end{proof}

\begin{lemma}
\label{chow-lemma-gysin-at-infty-independent}
Dans le lemme \ref{chow-lemma-gysin-at-infty}, soit $g : W' \to W$ un morphisme propre
qui est un isomorphisme au-dessus de $\mathbf{A}^1_X$. Soient
$C' \in A^0(W'_\infty \to X)$ et $C \in A^0(W_\infty \to X)$
les classes construites dans le lemme \ref{chow-lemma-gysin-at-infty}.
Alors $g_{\infty, *} \circ C' = C$ dans $A^0(W_\infty \to X)$.
\end{lemma}

\begin{proof}
Posons $b' = b \circ g : W' \to \mathbf{P}^1_X$. Notons
$i'_\infty : W'_\infty \to W'$ le morphisme d'inclusion.
Notons $g_\infty : W'_\infty \to W_\infty$ la restriction de $g$.
Étant donné $\alpha \in \CH_k(X)$, choisissons $\beta' \in \CH_{k + 1}(W')$
dont la restriction est l'image inverse plate de $\alpha$ sur $(b')^{-1}\mathbf{A}^1_X$.
Alors $\beta = g_*\beta' \in \CH_{k + 1}(W)$ se restreint en
l'image inverse plate de $\alpha$ sur $b^{-1}\mathbf{A}^1_X$.
Nous avons alors $i_\infty^*\beta = g_{\infty, *}(i'_\infty)^*\beta'$
d'après le lemme \ref{chow-lemma-closed-in-X-gysin}.
Ceci, ainsi que le fait correspondant après changement de base par des
morphismes $X' \to X$ localement de type fini, équivaut
à l'assertion du lemme.
\end{proof}

\begin{lemma}
\label{chow-lemma-homomorphism-pre}
Dans le lemme \ref{chow-lemma-gysin-at-infty}, nous avons
$C \circ (W_\infty \to X)_* \circ i_\infty^* = i_\infty^*$.
\end{lemma}

\begin{proof}
Soit $\beta \in \CH_{k + 1}(W)$. Notons $i_0 : X = X \times \{0\} \to W$
l'immersion fermée de la fibre au-dessus de $0$ dans $\mathbf{P}^1$. Alors
$(W_\infty \to X)_* i_\infty^* \beta = i_0^*\beta$ dans $\CH_k(X)$, car
$i_{\infty, *}i_\infty^*\beta$ et $i_{0, *}i_0^*\beta$
représentent la même classe sur $W$ (par exemple d'après le
lemme \ref{chow-lemma-support-cap-effective-Cartier})
et ont donc la même image directe sur $X$.
La restriction de $\beta$ à $b^{-1}(\mathbf{A}^1_X)$
se restreint en l'image inverse plate de
$i_0^*\beta = (W_\infty \to X)_* i_\infty^* \beta$, car nous pouvons le vérifier
après image inverse par $i_0$ ; voir les
lemmes \ref{chow-lemma-linebundle} et \ref{chow-lemma-linebundle-formulae}.
Nous pouvons donc utiliser $\beta$ pour calculer l'image de
$(W_\infty \to X)_*i_\infty^*\beta$ par $C$,
et nous obtenons le résultat voulu.
\end{proof}

\begin{lemma}
\label{chow-lemma-gysin-at-infty-commutes}
Dans le lemme \ref{chow-lemma-gysin-at-infty}, soit $f : Y \to X$ un morphisme
localement de type fini et soit $c \in A^*(Y \to X)$. Alors $C \circ c = c \circ C$
dans $A^*(W_\infty \times_X Y \to X)$.
\end{lemma}

\begin{proof}
Considérons le diagramme commutatif
$$
\xymatrix{
W_\infty \times_X Y \ar@{=}[r] &
W_{Y, \infty} \ar[r]_{i_{Y, \infty}} \ar[d] &
W_Y \ar[r]_{b_Y} \ar[d] &
\mathbf{P}^1_Y \ar[r]_{p_Y} \ar[d] &
Y \ar[d]^f \\
& W_\infty \ar[r]^{i_\infty} &
W \ar[r]^b &
\mathbf{P}^1_X \ar[r]^p &
X
}
$$
dont les carrés sont cartésiens. Pour un élément $\alpha \in \CH_k(X)$,
choisissons $\beta \in \CH_{k + 1}(W)$ dont la restriction à $b^{-1}(\mathbf{A}^1_X)$
est l'image inverse plate de $\alpha$. Alors $c \cap \beta$ est une classe
dans $\CH_*(W_Y)$ dont la restriction à $b_Y^{-1}(\mathbf{A}^1_Y)$
est l'image inverse plate de $c \cap \alpha$. Ensuite, nous avons
$$
i_{Y, \infty}^*(c \cap \beta) = c \cap i_\infty^*\beta
$$
car $c$ est une classe bivariante. Cela signifie exactement que
$C \cap c \cap \alpha = c \cap C \cap \alpha$. Le même argument
vaut après tout changement de base par $X' \to X$ localement de type fini.
Cela démontre le lemme.
\end{proof}





\section{Préparation aux classes de Chern localisées}
\label{chow-section-preparation-localized-chern-II}

\noindent
Dans cette section, nous étudions quelques propriétés des classes bivariantes
construites dans le lemme suivant. Nous recommandons vivement au
lecteur de sauter le reste de la section.

\begin{lemma}
\label{chow-lemma-localized-chern-pre}
Soit $(S, \delta)$ comme dans la situation \ref{chow-situation-setup}. Soit $X$
localement de type fini sur $S$. Soit $Z \subset X$ un sous-schéma fermé.
Soit
$$
b : W \longrightarrow \mathbf{P}^1_X
$$
un morphisme propre de schémas. Soit $Q \in D(\mathcal{O}_W)$ un
objet parfait. Notons $W_\infty \subset W$ l'image inverse du diviseur
$D_\infty \subset \mathbf{P}^1_X$ dont le complémentaire est $\mathbf{A}^1_X$.
Nous supposons que
\begin{enumerate}
\item[(A0)] les classes de Chern de $Q$ sont définies
(section \ref{chow-section-pre-derived}) ;
\item[(A1)] $b$ est un isomorphisme au-dessus de $\mathbf{A}^1_X$ ;
\item[(A2)] il existe un sous-schéma fermé $T \subset W_\infty$
contenant tous les points de $W_\infty$ situés au-dessus de $X \setminus Z$ tel que
$Q|_T$ soit nul, resp.\ isomorphe à un
$\mathcal{O}_T$-module localement libre de rang fini, de rang $< p$ et placé en degré cohomologique $0$.
\end{enumerate}
Il existe alors une classe bivariante canonique
$$
P'_p(Q),\text{ resp. }c'_p(Q) \in A^p(Z \to X)
$$
telle que
$(Z \to X)_* \circ P'_p(Q) = P_p(Q|_{X \times \{0\}})$,
resp.\ $(Z \to X)_* \circ c'_p(Q) = c_p(Q|_{X \times \{0\}})$.
\end{lemma}

\begin{proof}
Notons $E \subset W_\infty$ l'image inverse de $Z$. Alors
$W_\infty = T \cup E$ et $b$ induit un morphisme propre $E \to Z$.
Notons $C \in A^0(W_\infty \to X)$ la classe bivariante construite
dans le lemme \ref{chow-lemma-gysin-at-infty}. Notons $P'_p(Q|_E)$, resp.\ $c'_p(Q|_E)$
dans $A^p(E \to W_\infty)$ la classe bivariante construite
dans le lemme \ref{chow-lemma-silly}. Cela a un sens car
$(Q|_E)|_{T \cap E}$ est nul, resp.\ isomorphe à un
$\mathcal{O}_{E \cap T}$-module localement libre de rang fini, de rang $< p$ et placé en
degré cohomologique $0$ par l'hypothèse (A2). Nous définissons alors
$$
P'_p(Q) = (E \to Z)_* \circ P'_p(Q|_E) \circ C,\text{ resp. }
c'_p(Q) = (E \to Z)_* \circ c'_p(Q|_E) \circ C
$$
C'est une classe bivariante ; voir le lemme \ref{chow-lemma-push-proper-bivariant}.
Puisque $E \to Z \to X$ est égal à $E \to W_\infty \to W \to X$, nous voyons que
\begin{align*}
(Z \to X)_* \circ c'_p(Q)
& =
(W \to X)_* \circ i_{\infty, *} \circ (E \to W_\infty)_*
\circ c'_p(Q|_E) \circ C \\
& =
(W \to X)_* \circ i_{\infty, *} \circ c_p(Q|_{W_\infty}) \circ C \\
& =
(W \to X)_* \circ c_p(Q) \circ i_{\infty, *} \circ C \\
& =
(W \to X)_*\circ c_p(Q) \circ i_{0, *} \\
& =
(W \to X)_* \circ i_{0, *} \circ c_p(Q|_{X \times \{0\}}) \\
& =
c_p(Q|_{X \times \{0\}})
\end{align*}
La deuxième égalité résulte du lemme \ref{chow-lemma-silly-silly}.
La troisième égalité vient du fait que $c_p(Q)$ est une classe bivariante.
La quatrième égalité résulte du lemme \ref{chow-lemma-gysin-at-infty}.
La cinquième égalité vient du fait que $c_p(Q)$ est une classe bivariante.
La dernière égalité vient de ce que $(W_0 \to W) \circ (W \to X)$
est l'identité de $X$ si l'on identifie $W_0$ à $X$ comme nous l'avons
fait ci-dessus. La même suite d'égalités permet de
démontrer la propriété pour $P'_p(Q)$.
\end{proof}

\begin{lemma}
\label{chow-lemma-base-change-localized-chern-pre}
Dans le lemme \ref{chow-lemma-localized-chern-pre}, soit $X' \to X$ un morphisme
localement de type fini. Notons
$Z'$, $b' : W' \to \mathbf{P}^1_{X'}$ et $T' \subset W'_\infty$
les changements de base de $Z$, $b : W \to \mathbf{P}^1_X$ et $T \subset W_\infty$.
Posons $Q' = (W' \to W)^*Q$. Alors la classe
$P'_p(Q')$, resp.\ $c'_p(Q')$ dans $A^p(Z' \to X')$, construite comme dans le
lemme \ref{chow-lemma-localized-chern-pre} à l'aide de $b'$, $Q'$ et $T'$,
est la restriction (remarque \ref{chow-remark-restriction-bivariant})
de la classe $P'_p(Q)$, resp.\ $c'_p(Q)$ dans $A^p(Z \to X)$.
\end{lemma}

\begin{proof}
Rappelons que la construction est la suivante :
$$
P'_p(Q) = (E \to Z)_* \circ P'_p(Q|_E) \circ C,\text{ resp. }
c'_p(Q) = (E \to Z)_* \circ c'_p(Q|_E) \circ C
$$
Le lemme résulte donc de la propriété de changement de base correspondante
pour $C$ (lemme \ref{chow-lemma-base-change-gysin-at-infty})
et du fait que la même propriété de changement de base vaut pour les classes
construites dans le lemme \ref{chow-lemma-silly} (nous omettons un petit détail).
\end{proof}

\begin{lemma}
\label{chow-lemma-localized-chern-pre-independent}
Dans le lemme \ref{chow-lemma-localized-chern-pre}, la classe bivariante
$P'_p(Q)$, resp.\ $c'_p(Q)$,
est indépendante du choix du sous-schéma fermé $T$.
De plus, étant donné un morphisme propre $g : W' \to W$ qui est un
isomorphisme au-dessus de $\mathbf{A}^1_X$, en posant $Q' = g^*Q$,
nous avons $P'_p(Q) = P'_p(Q')$, resp.\ $c'_p(Q) = c'_p(Q')$.
\end{lemma}

\begin{proof}
L'indépendance à l'égard de $T$ résulte immédiatement du
lemme \ref{chow-lemma-silly-independent}.

\medskip\noindent
Soit $g : W' \to W$ un morphisme propre qui est un isomorphisme au-dessus de
$\mathbf{A}^1_X$. Remarquons que $T' = g^{-1}(T) \subset W'_\infty$
est un sous-schéma fermé satisfaisant à (A2), de sorte que l'opérateur
$P'_p(Q')$, resp.\ $c'_p(Q')$ dans $A^p(Z \to X)$,
correspondant à $b' = b \circ g : W' \to \mathbf{P}^1_X$
et à $Q'$, est défini. Notons $E' \subset W'_\infty$
l'image inverse de $Z$ dans $W'_\infty$. Rappelons que
$$
c'_p(Q') = (E' \to Z)_* \circ c'_p(Q'|_{E'}) \circ C'
$$
avec $C' \in A^0(W'_\infty \to X)$ et
$c'_p(Q'|_{E'}) \in A^p(E' \to W'_\infty)$.
D'après le lemme \ref{chow-lemma-gysin-at-infty-independent}, nous avons
$g_{\infty, *} \circ C' = C$. Remarquons que $E'$ est aussi
l'image inverse de $E$ dans $W'_\infty$ par $g_\infty$.
De plus, puisque $Q' = g^*Q$, nous constatons que $c'_p(Q'|_{E'})$ est simplement la
restriction de $c'_p(Q|_E)$ aux schémas au-dessus de $W'_\infty$ ; voir la
remarque \ref{chow-remark-restriction-bivariant}. Nous obtenons donc
\begin{align*}
c'_p(Q')
& = 
(E' \to Z)_* \circ c'_p(Q'|_{E'}) \circ C' \\
& =
(E \to Z)_* \circ (E' \to E)_* \circ c'_p(Q|_E) \circ C' \\
& =
(E \to Z)_* \circ c'_p(Q|_E) \circ g_{\infty, *} \circ C' \\
& =
(E \to Z)_* \circ c'_p(Q|_E) \circ C \\
& =
c'_p(Q)
\end{align*}
Dans la troisième égalité, nous avons utilisé que $c'_p(Q|_E)$
commute à l'image directe propre puisqu'il s'agit d'une
classe bivariante. L'égalité $P'_p(Q) = P'_p(Q')$
se démontre exactement de la même manière.
\end{proof}

\begin{lemma}
\label{chow-lemma-homomorphism}
Dans le lemme \ref{chow-lemma-localized-chern-pre}, supposons que $Q|_T$ soit isomorphe
à un $\mathcal{O}_T$-module localement libre de rang fini, de rang $< p$.
Notons $C \in A^0(W_\infty \to X)$ la classe du
lemme \ref{chow-lemma-gysin-at-infty}. Alors
$$
C \circ c_p(Q|_{X \times \{0\}}) =
C \circ (Z \to X)_* \circ c'_p(Q) = c_p(Q|_{W_\infty}) \circ C
$$
\end{lemma}

\begin{proof}
La première égalité vaut car $c_p(Q|_{X \times \{0\}}) =
(Z \to X)_* \circ c'_p(Q)$ d'après le lemme \ref{chow-lemma-localized-chern-pre}.
Nous pouvons démontrer la seconde égalité classe de cycle par classe de cycle
(voir le lemme \ref{chow-lemma-bivariant-zero}). Comme la construction des
classes bivariantes du lemme est compatible au changement de base,
nous pouvons supposer disposer d'un $\alpha \in \CH_k(X)$ et devoir montrer que
$C \cap (Z \to X)_*(c'_p(Q) \cap \alpha) =
c_p(Q|_{W_\infty}) \cap C \cap \alpha$. Remarquons que
\begin{align*}
C \cap (Z \to X)_*(c'_p(Q) \cap \alpha)
& =
C \cap (Z \to X)_* (E \to Z)_*(c'_p(Q|_E) \cap C \cap \alpha) \\
& =
C \cap (W_\infty \to X)_*(E \to W_\infty)_*(c'_p(Q|_E) \cap C \cap \alpha) \\
& =
C \cap (W_\infty \to X)_*(E \to W_\infty)_*(c'_p(Q|_E) \cap i_\infty^*\beta) \\
& =
C \cap (W_\infty \to X)_*(c_p(Q|_{W_\infty}) \cap i_\infty^*\beta) \\
& =
C \cap (W_\infty \to X)_*i_\infty^*(c_p(Q) \cap \beta) \\
& =
i_\infty^*(c_p(Q) \cap \beta) \\
& =
c_p(Q|_{W_\infty}) \cap i_\infty^*\beta \\
& =
c_p(Q|_{W_\infty}) \cap C \cap \alpha
\end{align*}
comme souhaité. Pour la première égalité, nous avons utilisé
$c'_p(Q) = (E \to Z)_* \circ c'_p(Q|_E) \circ C$, où $E \subset W_\infty$
est l'image inverse de $Z$ et $c'_p(Q|_E)$ la classe construite
dans le lemme \ref{chow-lemma-silly}. La deuxième égalité exprime simplement
que $E \to Z \to X$ est égal à $E \to W_\infty \to X$.
Pour la troisième égalité, choisissons $\beta \in \CH_{k + 1}(W)$ dont la restriction à
$b^{-1}(\mathbf{A}^1_X)$ est l'image inverse plate de $\alpha$, de sorte que
$C \cap \alpha = i_\infty^*\beta$ par construction. La quatrième égalité est le
lemme \ref{chow-lemma-silly-silly}. La cinquième vient du fait que
$c_p(Q)$ est une classe bivariante et commute donc à $i_\infty^*$.
La sixième égalité est le lemme \ref{chow-lemma-homomorphism-pre}.
La septième utilise de nouveau le fait que $c_p(Q)$ est une classe bivariante.
La dernière vaut puisque $C \cap \alpha = i_\infty^*\beta$.
\end{proof}

\begin{lemma}
\label{chow-lemma-homomorphism-commute}
Dans le lemme \ref{chow-lemma-localized-chern-pre}, soit $Y \to X$ un morphisme
localement de type fini et soit $c \in A^*(Y \to X)$ une classe bivariante.
Alors
$$
P'_p(Q) \circ c = c \circ P'_p(Q)
\quad\text{resp.}\quad
c'_p(Q) \circ c = c \circ c'_p(Q)
$$
dans $A^*(Y \times_X Z \to X)$.
\end{lemma}

\begin{proof}
Soit $E \subset W_\infty$ l'image inverse de $Z$.
Rappelons que $P'_p(Q) = (E \to Z)_* \circ P'_p(Q|_E) \circ C$,
resp.\ $c'_p(Q) = (E \to Z)_* \circ c'_p(Q|_E) \circ C$,
où $C$ est comme dans le lemme \ref{chow-lemma-gysin-at-infty} et
$P'_p(Q|_E)$, resp.\ $c'_p(Q|_E)$ sont comme dans le
lemme \ref{chow-lemma-silly}.
D'après le lemme \ref{chow-lemma-gysin-at-infty-commutes},
nous voyons que $C$ commute à $c$,
et le lemme \ref{chow-lemma-silly-commutes} montre que
$P'_p(Q|_E)$, resp.\ $c'_p(Q|_E)$ commute à $c$.
Comme $c$ est une classe bivariante, elle commute par définition à l'image
directe propre par $E \to Z$. Cela achève la démonstration.
\end{proof}

\begin{lemma}
\label{chow-lemma-localized-chern-pre-compose}
Dans le lemme \ref{chow-lemma-localized-chern-pre}, supposons que $Q|_T$ soit nul. Dans
$A^*(Z \to X)$, nous avons
\begin{align*}
P'_1(Q) & = c'_1(Q), \\
P'_2(Q) & = c'_1(Q)^2 - 2c'_2(Q), \\
P'_3(Q) & = c'_1(Q)^3 - 3c'_1(Q)c'_2(Q) + 3c'_3(Q), \\
P'_4(Q) & = c'_1(Q)^4 - 4c'_1(Q)^2c'_2(Q) +
4c'_1(Q)c'_3(Q) + 2c'_2(Q)^2 - 4c'_4(Q),
\end{align*}
et ainsi de suite, la multiplication étant celle de la remarque \ref{chow-remark-ring-loc-classes}.
\end{lemma}

\begin{proof}
L'énoncé a un sens car le faisceau nul est de rang $< 1$ et
les classes $c'_p(Q)$ sont donc définies pour tout $p \geq 1$.
Dans la démonstration du lemme \ref{chow-lemma-localized-chern-pre}, nous avons construit
les classes $P'_p(Q)$ et $c'_p(Q)$ en utilisant la classe bivariante
$C \in A^0(W_\infty \to X)$ du lemme \ref{chow-lemma-gysin-at-infty}
et les classes bivariantes
$P'_p(Q|_E)$ et $c'_p(Q|_E)$ du lemme \ref{chow-lemma-silly} pour la restriction
$Q|_E$ de $Q$ à l'image inverse $E$ de $Z$ dans $W_\infty$.
Remarquons que le lemme \ref{chow-lemma-silly-compose} donne la relation voulue
entre $P'_p(Q|_E)$ et $c'_p(Q|_E)$. Rappelons que
$$
P'_p(Q) = (E \to Z)_* \circ P'_p(Q|_E) \circ C
\quad\text{et}\quad
c'_p(Q) = (E \to Z)_* \circ c'_p(Q|_E) \circ C
$$
Pour achever la démonstration, il suffit de montrer que les multiplications définies
dans la remarque \ref{chow-remark-ring-loc-classes} sur les classes $a_p = c'_p(Q)$
et sur les classes $b_p = c'_p(Q|_E)$ concordent :
$$
a_{p_1}a_{p_2} \ldots a_{p_r} =
(E \to Z)_* \circ b_{p_1}b_{p_2} \ldots b_{p_r} \circ C
$$
Nous omettons certains détails. Si $r = 1$, c'est vrai.
Pour $r > 1$, remarquons que, d'après la remarque \ref{chow-remark-res-push}, la multiplication de la
remarque \ref{chow-remark-ring-loc-classes} s'effectue
en insérant $(Z \to X)_*$, resp.\ $(E \to W_\infty)_*$ entre
les facteurs du produit
$a_{p_1}a_{p_2} \ldots a_{p_r}$, resp.\ $b_{p_1}b_{p_2} \ldots b_{p_r}$,
et en prenant les compositions comme classes bivariantes.
Or, d'après le lemme \ref{chow-lemma-silly}, nous avons
$$
(E \to W_\infty)_* \circ b_{p_i} = c_{p_i}(Q|_{W_\infty})
$$
et, d'après le lemme \ref{chow-lemma-homomorphism}, nous avons
$$
C \circ (Z \to X)_* \circ a_{p_i} = c_{p_i}(Q|_{W_\infty}) \circ C
$$
pour $i = 2, \ldots, r$. Un calcul
montre que le membre de gauche et le membre de droite de l'égalité voulue
se simplifient tous deux en
$$
(E \to Z)_* \circ c'_{p_1}(Q|_E) \circ
c_{p_2}(Q|_{W_\infty}) \circ \ldots \circ
c_{p_r}(Q|_{W_\infty}) \circ C
$$
ce qui achève la démonstration.
\end{proof}

\begin{lemma}
\label{chow-lemma-localized-chern-pre-sum-c}
Dans le lemme \ref{chow-lemma-localized-chern-pre}, supposons que $Q|_T$ soit isomorphe
à un $\mathcal{O}_T$-module localement libre de rang fini, de rang $< p$.
Supposons donné un autre objet parfait $Q' \in D(\mathcal{O}_W)$
dont les classes de Chern sont définies et tel que $Q'|_T$ soit isomorphe à un
$\mathcal{O}_T$-module localement libre de rang fini, de rang $< p'$, placé
en degré cohomologique $0$. Avec la notation de la
remarque \ref{chow-remark-ring-loc-classes}, posons
$$
c^{(p)}(Q) = 1 + c_1(Q|_{X \times \{0\}}) + \ldots +
c_{p - 1}(Q|_{X \times \{0\}}) +
c'_{p}(Q) + c'_{p + 1}(Q) + \ldots
$$
dans $A^{(p)}(Z \to X)$, où $c'_i(Q)$, pour $i \geq p$, est comme dans le
lemme \ref{chow-lemma-localized-chern-pre}. De même pour $c^{(p')}(Q')$ et
$c^{(p + p')}(Q \oplus Q')$.
Alors $c^{(p + p')}(Q \oplus Q') = c^{(p)}(Q)c^{(p')}(Q')$
dans $A^{(p + p')}(Z \to X)$.
\end{lemma}

\begin{proof}
Rappelons que l'image de $c'_i(Q)$ dans $A^p(X)$ est égale à
$c_i(Q|_{X \times \{0\}})$ pour $i \geq p$, et de même pour
$Q'$ et $Q \oplus Q'$ ; voir le lemme \ref{chow-lemma-localized-chern-pre}.
L'égalité en degrés $< p + p'$ résulte donc de
l'additivité du lemme \ref{chow-lemma-additivity-on-perfect}.

\medskip\noindent
Prenons $n \geq p + p'$.
Comme dans la démonstration du lemme \ref{chow-lemma-localized-chern-pre},
notons $E \subset W_\infty$ l'image inverse de $Z$.
Remarquons que nous avons l'égalité
$$
c^{(p + p')}(Q|_E \oplus Q'|_E) =
c^{(p)}(Q|_E)c^{(p')}(Q'|_E)
$$
dans $A^{(p + p')}(E \to W_\infty)$ d'après le lemme \ref{chow-lemma-silly-sum-c}.
Puisque, par construction,
$$
c'_p(Q \oplus Q') = (E \to Z)_* \circ c'_p(Q|_E \oplus Q'|_E) \circ C
$$
nous concluons qu'il suffit de montrer, pour tout $i + j = n$, que
$$
(E \to Z)_* \circ c^{(p)}_i(Q|_E)c^{(p')}_j(Q'|_E) \circ C
=
c^{(p)}_i(Q)c^{(p')}_j(Q')
$$
dans $A^n(Z \to X)$, où la multiplication est celle de la
remarque \ref{chow-remark-ring-loc-classes} dans les deux membres. Il y a
trois cas, selon que $i \geq p$, $j \geq p'$, ou les deux.

\medskip\noindent
Supposons $i \geq p$ et $j \geq p'$. Dans ce cas, les produits sont
définis en insérant $(E \to W_\infty)_*$, resp.\ $(Z \to X)_*$ entre
les deux facteurs et en prenant les compositions comme classes bivariantes ; voir la
remarque \ref{chow-remark-res-push}.
Autrement dit, nous devons montrer
$$
(E \to Z)_* \circ c'_i(Q|_E) \circ
(E \to W_\infty)_* \circ c'_j(Q'|_E) \circ C =
c'_i(Q) \circ (Z \to X)_* \circ c'_j(Q')
$$
D'après le lemme \ref{chow-lemma-silly}, le membre de gauche est égal à
$$
(E \to Z)_* \circ c'_i(Q|_E) \circ c_j(Q'|_{W_\infty}) \circ C
$$
Puisque $c'_i(Q) = (E \to Z)_* \circ c'_i(Q|_E) \circ C$,
le membre de droite est égal à
$$
(E \to Z)_* \circ c'_i(Q|_E) \circ C \circ (Z \to X)_* \circ c'_j(Q')
$$
ce qui est immédiatement égal à l'expression ci-dessus
d'après le lemme \ref{chow-lemma-homomorphism}.

\medskip\noindent
Supposons $i \geq p$ et $j < p$. En développant les produits
dans ce cas, nous devons montrer
$$
(E \to Z)_* \circ c'_i(Q|_E) \circ c_j(Q'|_{W_\infty}) \circ C =
c'_i(Q) \circ c_j(Q'|_{X \times \{0\}})
$$
En utilisant de nouveau $c'_i(Q) = (E \to Z)_* \circ c'_i(Q|_E) \circ C$,
nous voyons qu'il suffit de montrer $c_j(Q'|_{W_\infty}) \circ C =
C \circ c_j(Q'|_{X \times \{0\}})$, ce qui fait partie du
lemme \ref{chow-lemma-homomorphism}.

\medskip\noindent
Supposons $i < p$ et $j \geq p'$. En développant les produits
dans ce cas, nous devons montrer
$$
(E \to Z)_* \circ c_i(Q|_E) \circ c'_j(Q'|_E) \circ C =
c_i(Q|_{Z \times \{0\}}) \circ c'_j(Q')
$$
Cependant, puisque $c'_j(Q|_E)$ et $c'_j(Q')$ sont des
classes bivariantes, elles commutent à l'intersection avec les classes de Chern
(lemme \ref{chow-lemma-cap-commutative-chern}). Il suffit donc de démontrer
$$
(E \to Z)_* \circ c'_j(Q'|_E) \circ c_i(Q|_{W_\infty}) \circ C =
c'_j(Q') \circ c_i(Q|_{X \times \{0\}})
$$
ce qui nous ramène au cas traité au paragraphe précédent.
\end{proof}

\begin{lemma}
\label{chow-lemma-localized-chern-pre-sum-P}
Dans le lemme \ref{chow-lemma-localized-chern-pre}, supposons que $Q|_T$ soit nul.
Supposons donné un autre objet parfait $Q' \in D(\mathcal{O}_W)$
dont les classes de Chern sont définies et dont la restriction $Q'|_T$ est nulle.
Dans ce cas, les classes
$P'_p(Q), P'_p(Q'), P'_p(Q \oplus Q') \in A^p(Z \to X)$
construites dans le lemme \ref{chow-lemma-localized-chern-pre}
satisfont $P'_p(Q \oplus Q') = P'_p(Q) + P'_p(Q')$.
\end{lemma}

\begin{proof}
Cela résulte immédiatement de la construction de ces
classes et du lemme \ref{chow-lemma-silly-sum-P}.
\end{proof}









 
\section{Classes de Chern localisées}
\label{chow-section-localized-chern}

\noindent
Aperçu de la construction. Soit $F$ un corps,
soit $X$ une variété sur $F$, soit $E$ un objet parfait de
$D(\mathcal{O}_X)$, et soit $Z \subset X$ un sous-schéma fermé
tel que $E|_{X \setminus Z} = 0$. Nous voulons alors construire des éléments
$$
c_p(Z \to X, E) \in A^p(Z \to X)
$$
Nous le ferons en construisant un diagramme
$$
\xymatrix{
W \ar[d]_f \ar[r]_q & X \\
\mathbf{P}^1_F
}
$$
et un objet parfait $Q$ de $D(\mathcal{O}_W)$ tels que
\begin{enumerate}
\item $f$ est plat, et $f$, $q$ sont propres ; pour $t \in \mathbf{P}^1_F$,
notons $W_t$ la fibre de $f$,
$q_t : W_t \to X$ la restriction de $q$, et $Q_t = Q|_{W_t}$ ;
\item $q_t : W_t \to X$ est un isomorphisme et $Q_t = q_t^*E$
pour $t \in \mathbf{A}^1_F$ ;
\item $q_\infty : W_\infty \to X$ est un isomorphisme au-dessus de $X \setminus Z$ ;
\item si $T \subset W_\infty$ est l'adhérence de
$q_\infty^{-1}(X \setminus Z)$, alors $Q_\infty|_T$ est nul.
\end{enumerate}
L'idée est de considérer ceci comme une famille $\{(W_t, Q_t)\}$
paramétrée par $t \in \mathbf{P}^1$.
Pour $t \not = \infty$, nous voyons que $c_p(Q_t)$ est simplement $c_p(E)$
sur les groupes de Chow de $Q_t = X$. Mais pour $t = \infty$, nous voyons que
$c_p(Q_\infty)$ envoie les classes sur $Q_\infty$ sur des classes à support dans
$E = q_\infty^{-1}(Z)$ puisque $Q_\infty|_T = 0$.
Nous considérons $E$ comme le lieu exceptionnel de $q_\infty : W_\infty \to X$.
Puisque tout $\alpha \in \CH_*(X)$ donne naissance à une « famille »
de cycles $\alpha_t \in \CH_*(W_t)$, il est naturel de définir
$c_p(Z \to X, E) \cap \alpha$ comme
l'image directe $(E \to Z)_*(c_p(Q_\infty) \cap \alpha_\infty)$.

\medskip\noindent
Pour réaliser cette construction, il y a deux ingrédients principaux : (1) la construction de
$W$ et de $Q$ est une sorte de construction algébrique du graphe de Macpherson ; elle
est effectuée dans Compléments sur la platitude, section \ref{flat-section-blowup-complexes-III}.
(2) La construction de la classe elle-même à partir de $W$ et $Q$ est faite dans la
section \ref{chow-section-preparation-localized-chern-II}, en s'appuyant sur les
sections \ref{chow-section-gysin-at-infty} et
\ref{chow-section-preparation-localized-chern}.

\begin{situation}
\label{chow-situation-loc-chern}
Soit $(S, \delta)$ comme dans la situation \ref{chow-situation-setup}. Soit $X$ un schéma
localement de type fini sur $S$. Soit $i : Z \to X$ une immersion fermée.
Soit $E \in D(\mathcal{O}_X)$ un objet. Soit $p \geq 0$. Supposons que
\begin{enumerate}
\item $E$ soit un objet parfait de $D(\mathcal{O}_X)$ ;
\item la restriction $E|_{X \setminus Z}$ soit nulle, resp.\ isomorphe à un
$\mathcal{O}_{X \setminus Z}$-module localement libre de rang fini, de rang $< p$,
placé en degré cohomologique $0$ ; et
\item au moins une\footnote{Veuillez ignorer cette condition technique en
première lecture ; voir la discussion de la remarque \ref{chow-remark-loc-chern}.}
des assertions suivantes soit vraie :
(a) $X$ est quasi-compact ;
(b) les composantes irréductibles de $X$ sont quasi-compactes ;
(c) il existe un complexe localement borné de
$\mathcal{O}_X$-modules localement libres de rang fini représentant $E$ ; ou
(d) il existe un morphisme $X \to X'$ de schémas localement de type fini
sur $S$ tel que $E$ soit l'image inverse d'un objet parfait sur $X'$ et que
les composantes irréductibles de $X'$ soient quasi-compactes.
\end{enumerate}
\end{situation}

\begin{lemma}
\label{chow-lemma-independent-loc-chern}
Dans la situation \ref{chow-situation-loc-chern}, il existe une classe bivariante canonique
$$
P_p(Z \to X, E) \in A^p(Z \to X),
\quad\text{resp.}\quad
c_p(Z \to X, E) \in A^p(Z \to X)
$$
ayant la propriété suivante :
\begin{equation}
\label{chow-equation-defining-property-localized-classes}
i_* \circ P_p(Z \to X, E) = P_p(E),
\quad\text{resp.}\quad
i_* \circ c_p(Z \to X, E) = c_p(E)
\end{equation}
comme classes bivariantes sur $X$ (où $\circ$ est comme dans le
lemme \ref{chow-lemma-push-proper-bivariant}).
\end{lemma}

\begin{proof}
La construction de ces classes bivariantes est la suivante. Soient
$$
b : W \longrightarrow \mathbf{P}^1_X
\quad\text{et}\quad
T \longrightarrow W_\infty
\quad\text{et}\quad
Q
$$
l'éclatement, l'objet parfait $Q$ de $D(\mathcal{O}_W)$
et l'immersion fermée construits dans
Compléments sur la platitude, section \ref{flat-section-blowup-complexes-III}
et le lemme \ref{flat-lemma-graph-construction}.
Soit $T' \subset T$ le sous-schéma ouvert et fermé tel que
$Q|_{T'}$ soit nul, resp.\ isomorphe à un
$\mathcal{O}_{T'}$-module localement libre de rang fini, de rang $< p$,
placé en degré cohomologique $0$. D'après la condition (2) de la
situation \ref{chow-situation-loc-chern}, les morphismes
$$
T' \to T \to W_\infty \to X
$$
sont tous des isomorphismes de schémas au-dessus du sous-schéma ouvert $X \setminus Z$ de $X$.
Nous vérifions ci-dessous que les classes de Chern de $Q$ sont définies.
Rappelant que $Q|_{X \times \{0\}} \cong E$ par construction, nous concluons
que la classe bivariante construite dans le lemme \ref{chow-lemma-localized-chern-pre}
à partir de $W, b, Q, T'$ nous donne les classes
$$
P_p(Z \to X, E) = P'_p(Q) \in A^p(Z \to X)
$$
et
$$
c_p(Z \to X, E) = c'_p(Q) \in A^p(Z \to X)
$$
satisfaisant à (\ref{chow-equation-defining-property-localized-classes}).

\medskip\noindent
Dans ce paragraphe, nous démontrons que les classes de Chern de $Q$ sont définies
(définition \ref{chow-definition-defined-on-perfect}) ; nous suggérons au lecteur de sauter
ce passage. Si l'hypothèse (3)(a) ou (3)(b) de la situation \ref{chow-situation-loc-chern}
est satisfaite, c'est-à-dire si les composantes irréductibles de $X$ sont quasi-compactes, alors il en
va de même pour $W$ (car $W \to X$ est propre). Nous concluons donc que
les classes de Chern de tout objet parfait de $D(\mathcal{O}_W)$ sont définies d'après le
lemme \ref{chow-lemma-chern-classes-defined}. Si (3)(c) est satisfaite, c'est-à-dire si
$E$ peut être représenté par un complexe localement borné de modules localement
libres de rang fini, alors l'objet $Q$ peut être représenté par un complexe localement borné
de $\mathcal{O}_W$-modules localement libres de rang fini d'après la partie (5) de
Compléments sur la platitude, lemme \ref{flat-lemma-graph-construction}. Les
classes de Chern de $Q$ sont donc définies. Enfin, supposons (3)(d) satisfaite, c'est-à-dire
qu'il existe un morphisme $X \to X'$ de schémas localement de type fini
sur $S$ tel que $E$ soit l'image inverse d'un objet parfait $E'$ sur $X'$ et que
les composantes irréductibles de $X'$ soient quasi-compactes.
Soient $b' : W' \to \mathbf{P}^1_{X'}$ et $Q' \in D(\mathcal{O}_{W'})$
le morphisme et l'objet parfait construits comme
dans Compléments sur la platitude, section \ref{flat-section-blowup-complexes-III},
à partir du triplet
$(\mathbf{P}^1_{X'}, (\mathbf{P}^1_{X'})_\infty, L(p')^*E')$.
La discussion ci-dessus montre que les classes de Chern de $Q'$
sont définies. Puisque $b$ et $b'$ ont été construits par une application de
Compléments sur la platitude, lemme \ref{flat-lemma-complex-and-divisor-blowup-pre},
il résulte de
Compléments sur la platitude, lemme \ref{flat-lemma-complex-and-divisor-blowup-base-change},
qu'il existe un morphisme $W \to W'$ tel que
$Q = L(W \to W')^*Q'$. Le
lemme \ref{chow-lemma-chern-classes-defined}
montre alors que les classes de Chern de $Q$ sont définies.
\end{proof}

\begin{definition}
\label{chow-definition-localized-chern}
Avec $(S, \delta)$, $X$, $E \in D(\mathcal{O}_X)$ et $i : Z \to X$ comme dans la
situation \ref{chow-situation-loc-chern}.
\begin{enumerate}
\item Si la restriction $E|_{X \setminus Z}$ est nulle, alors, pour tout
$p \geq 0$, nous définissons
$$
P_p(Z \to X, E) \in A^p(Z \to X)
$$
par la construction du lemme \ref{chow-lemma-independent-loc-chern},
et nous définissons le {\it caractère de Chern localisé} par la formule
$$
ch(Z \to X, E) =
\sum\nolimits_{p = 0, 1, 2, \ldots} \frac{P_p(Z \to X, E)}{p!}
\quad\text{dans}\quad \prod\nolimits_{p \geq 0} A^p(Z \to X) \otimes \mathbf{Q}
$$
\item Si la restriction $E|_{X \setminus Z}$ est isomorphe à un
$\mathcal{O}_{X \setminus Z}$-module localement libre de rang fini, de rang $< p$,
placé en degré cohomologique $0$, nous définissons alors la
{\it $p$-ième classe de Chern localisée} $c_p(Z \to X, E)$ par la construction
du lemme \ref{chow-lemma-independent-loc-chern}.
\end{enumerate}
\end{definition}

\noindent
Dans la situation de la définition, supposons que $E|_{X \setminus Z}$ soit nul.
Pour être précis, nous avons alors l'égalité
$$
i_* \circ ch(Z \to X, E) = ch(E)
$$
dans $A^*(X) \otimes \mathbf{Q}$, car nous avons démontré ci-dessus
l'égalité (\ref{chow-equation-defining-property-localized-classes}).

\medskip\noindent
Voici une vérification de cohérence importante.

\begin{lemma}
\label{chow-lemma-base-change-loc-chern}
Dans la situation \ref{chow-situation-loc-chern},
soit $f : X' \to X$ un morphisme de schémas localement de type fini.
Notons $E' = f^*E$ et $Z' = f^{-1}(Z)$. Alors la classe bivariante
de la définition \ref{chow-definition-localized-chern}
$$
P_p(Z' \to X', E') \in A^p(Z' \to X'),
\quad\text{resp.}\quad
c_p(Z' \to X', E') \in A^p(Z' \to X')
$$
construite comme dans le lemme \ref{chow-lemma-independent-loc-chern}
à partir de $X', Z', E'$ est la restriction
(remarque \ref{chow-remark-restriction-bivariant}) de la
classe bivariante $P_p(Z \to X, E) \in A^p(Z \to X)$,
resp.\ $c_p(Z \to X, E) \in A^p(Z \to X)$.
\end{lemma}

\begin{proof}
Notons $p : \mathbf{P}^1_X \to X$ et $p' : \mathbf{P}^1_{X'} \to X'$
les morphismes structuraux.
Rappelons que $b : W \to \mathbf{P}^1_X$ et $b' : W' \to \mathbf{P}^1_{X'}$
sont les morphismes construits à partir des triplets
$(\mathbf{P}^1_X, (\mathbf{P}^1_X)\infty, p^*E)$ et
$(\mathbf{P}^1_{X'}, (\mathbf{P}^1_{X'})\infty, (p')^*E')$
dans Compléments sur la platitude, lemme \ref{flat-lemma-complex-and-divisor-blowup-pre}.
En outre, $Q = L\eta_{\mathcal{I}_\infty}p^*E$ et
$Q = L\eta_{\mathcal{I}'_\infty}(p')^*E'$
où
$\mathcal{I}_\infty \subset \mathcal{O}_W$ est le
faisceau d'idéaux de $W_\infty$ et
$\mathcal{I}'_\infty \subset \mathcal{O}_{W'}$ est le
faisceau d'idéaux de $W'_\infty$.
Ensuite, $h : \mathbf{P}^1_{X'} \to \mathbf{P}^1_X$ est un morphisme de
schémas tel que l'image inverse du diviseur de Cartier effectif
$(\mathbf{P}^1_X)_\infty$ soit le diviseur de Cartier effectif
$(\mathbf{P}^1_{X'})_\infty$ et que $h^*p^*E = (p')^*E'$.
D'après Compléments sur la platitude, lemme
\ref{flat-lemma-complex-and-divisor-blowup-base-change},
nous obtenons un diagramme commutatif
$$
\xymatrix{
W' \ar[rd]_{b'} \ar[r]_-g &
\mathbf{P}^1_{X'} \times_{\mathbf{P}^1_X} W \ar[d]_r \ar[r]_-q &
W \ar[d]^b \\
&
\mathbf{P}^1_{X'} \ar[r] &
\mathbf{P}^1_X
}
$$
tel que $W'$ soit le « transformé strict » de $\mathbf{P}^1_{X'}$
par rapport à $b$ et que $Q' = (q \circ g)^*Q$.
Rappelons maintenant que $P_p(Z \to X, E) = P'_p(Q)$,
resp.\ $c_p(Z \to X, E) = c'_p(Q)$, où $P'_p(Q)$, resp.\ $c'_p(Q)$,
sont construits dans le lemme \ref{chow-lemma-localized-chern-pre}
à l'aide de $b, Q, T'$, où $T'$ est un sous-schéma fermé $T' \subset W_\infty$
possédant les deux propriétés suivantes :
(a) $T'$ contient tous les points de $W_\infty$ situés au-dessus de $X \setminus Z$,
et (b) $Q|_{T'}$ est nul, resp.\ isomorphe à un module localement libre de rang fini,
de rang $< p$ et placé en degré $0$.
Dans la construction du lemme \ref{chow-lemma-localized-chern-pre},
nous avons choisi un sous-schéma fermé particulier $T'$ possédant les propriétés (a) et (b),
mais le choix précis de $T'$ est sans importance ; voir le
lemme \ref{chow-lemma-localized-chern-pre-independent}.

\medskip\noindent
Ensuite, d'après le lemme \ref{chow-lemma-base-change-localized-chern-pre},
la restriction de la classe bivariante $P_p(Z \to X, E) = P'_p(Q)$,
resp.\ $c_p(Z \to X, E) = c_p(Q')$,
à $X'$ correspond à la classe $P'_p(q^*Q)$, resp.\ $c'_p(q^*Q)$,
construite comme dans le lemme \ref{chow-lemma-localized-chern-pre} à l'aide de
$r : \mathbf{P}^1_{X'} \times_{\mathbf{P}^1_X} W \to \mathbf{P}^1_{X'}$,
du complexe $q^*Q$ et de l'image inverse $q^{-1}(T')$.

\medskip\noindent
Maintenant, d'après la seconde assertion du
lemme \ref{chow-lemma-localized-chern-pre-independent},
nous avons $P'_p(Q') = P'_p(q^*Q)$, resp.\ $c'_p(q^*Q) = c'_p(Q')$.
Puisque $P_p(Z' \to X', E') = P'_p(Q')$, resp.\ $c_p(Z' \to X', E') = c'_p(Q')$,
nous concluons que le lemme est vrai.
\end{proof}

\begin{remark}
\label{chow-remark-loc-chern}
Dans la situation \ref{chow-situation-loc-chern}, il aurait été plus naturel
de remplacer l'hypothèse (3) par l'hypothèse : « les classes de Chern
de $E$ sont définies ». En fait, en combinant les
lemmes \ref{chow-lemma-independent-loc-chern} et
\ref{chow-lemma-base-change-loc-chern}
avec le lemme \ref{chow-lemma-envelope-bivariant},
il est facile d'étendre la définition à ce cas légèrement plus général.
Si cela nous est un jour nécessaire, nous le ferons ici.
\end{remark}

\begin{lemma}
\label{chow-lemma-loc-chern-after-pushforward}
Dans la situation \ref{chow-situation-loc-chern}, nous avons
$$
P_p(Z \to X, E) \cap i_*\alpha = P_p(E|_Z) \cap \alpha,
\quad\text{resp.}\quad
c_p(Z \to X, E) \cap i_*\alpha = c_p(E|_Z) \cap \alpha
$$
dans $\CH_*(Z)$ pour tout $\alpha \in \CH_*(Z)$.
\end{lemma}

\begin{proof}
Nous démontrons seulement la seconde égalité et omettons la démonstration de la première.
Puisque $c_p(Z \to X, E)$ est une classe bivariante et que le changement de
base de $Z \to X$ par $Z \to X$ est $\text{id} : Z \to Z$, nous avons
$c_p(Z \to X, E) \cap i_*\alpha = c_p(Z \to X, E) \cap \alpha$.
D'après le lemme \ref{chow-lemma-base-change-loc-chern}, la restriction de
$c_p(Z \to X, E)$ à $Z$ (!) est la classe de Chern localisée pour
$\text{id} : Z \to Z$ et $E|_Z$. Le résultat résulte donc de
(\ref{chow-equation-defining-property-localized-classes}) avec $X = Z$.
\end{proof}

\begin{lemma}
\label{chow-lemma-loc-chern-disjoint}
Dans la situation \ref{chow-situation-loc-chern},
si $\alpha \in \CH_k(X)$ a un support disjoint de $Z$, alors
$P_p(Z \to X, E) \cap \alpha = 0$, resp.\ $c_p(Z \to X, E) \cap \alpha = 0$.
\end{lemma}

\begin{proof}
Cela résulte immédiatement de la construction des classes de Chern localisées.
Cela résulte aussi du fait que nous pouvons calculer
$c_p(Z \to X, E) \cap \alpha$ en restreignant d'abord $c_p(Z \to X, E)$ au
support de $\alpha$, puis en utilisant
le lemme \ref{chow-lemma-base-change-loc-chern}
pour voir que cette restriction est nulle.
\end{proof}

\begin{lemma}
\label{chow-lemma-loc-chern-shrink-Z}
Dans la situation \ref{chow-situation-loc-chern},
supposons que $Z \subset Z' \subset X$, où $Z'$ est un sous-schéma fermé de $X$.
Alors
$P_p(Z' \to X, E) = (Z \to Z')_* \circ P_p(Z \to X, E)$,
resp.\ $c_p(Z' \to X, E) = (Z \to Z')_* \circ c_p(Z \to X, E)$
(avec $\circ$ comme dans le lemme \ref{chow-lemma-push-proper-bivariant}).
\end{lemma}

\begin{proof}
La construction de $P_p(Z' \to X, E)$,
resp.\ $c_p(Z' \to X, E)$ dans le lemme \ref{chow-lemma-independent-loc-chern}
utilise exactement le même morphisme
$b : W \to \mathbf{P}^1_X$ et le même objet parfait $Q$ de $D(\mathcal{O}_W)$.
Nous pouvons alors utiliser le lemme \ref{chow-lemma-silly-shrink} pour conclure.
Certains détails sont omis.
\end{proof}

\begin{lemma}
\label{chow-lemma-loc-chern-agree}
Dans le lemme \ref{chow-lemma-silly}, supposons que $E_2$ soit la restriction d'un objet parfait
$E \in D(\mathcal{O}_X)$ dont la restriction à $X_1$ est nulle,
resp.\ est isomorphe à un $\mathcal{O}_{X_1}$-module localement libre de rang fini
de rang $< p$ placé en degré cohomologique $0$. Alors la classe
$P'_p(E_2)$, resp.\ $c'_p(E_2)$ du lemme \ref{chow-lemma-silly} coïncide avec
$P_p(X_2 \to X, E)$, resp.\ $c_p(X_2 \to X, E)$ de la
définition \ref{chow-definition-localized-chern}, pourvu que $E$ satisfasse
l'hypothèse (3) de la situation \ref{chow-situation-loc-chern}.
\end{lemma}

\begin{proof}
Les hypothèses sur $E$ impliquent qu'il existe un ouvert $U \subset X$
contenant $X_1$ tel que $E|_U$ soit nul, resp.\ soit isomorphe à un
$\mathcal{O}_U$-module localement libre de rang $< p$. Voir Compléments d'algèbre, lemme
\ref{more-algebra-lemma-lift-perfect-from-residue-field}.
Soit $Z \subset X$ le complémentaire de $U$ dans $X$, muni de
la structure de sous-schéma fermé induite réduite. Alors
$P_p(X_2 \to X, E) = (Z \to X_2)_* \circ P_p(Z \to X, E)$,
resp.\ $c_p(X_2 \to X, E) = (Z \to X_2)_* \circ c_p(Z \to X, E)$
par le lemme \ref{chow-lemma-loc-chern-shrink-Z}.
Nous pouvons maintenant démontrer que $P_p(X_2 \to X, E)$, resp.\ $c_p(X_2 \to X, E)$
satisfait la caractérisation de $P'_p(E_2)$, resp.\ $c'_p(E_2)$
donnée dans le lemme \ref{chow-lemma-silly}. En effet, par la relation
$P_p(X_2 \to X, E) = (Z \to X_2)_* \circ P_p(Z \to X, E)$,
resp.\ $c_p(X_2 \to X, E) = (Z \to X_2)_* \circ c_p(Z \to X, E)$
que nous venons de démontrer et le fait que $X_1 \cap Z = \emptyset$,
la composée $P_p(X_2 \to X, E) \circ i_{1, *}$,
resp.\ $c_p(X_2 \to X, E) \circ i_{1, *}$ est nulle
par le lemme \ref{chow-lemma-loc-chern-disjoint}.
D'autre part,
$P_p(X_2 \to X, E) \circ i_{2, *} = P_p(E_2)$,
resp.\ $c_p(X_2 \to X, E) \circ i_{2, *} = c_p(E_2)$
par le lemme \ref{chow-lemma-loc-chern-after-pushforward}.
\end{proof}





\section{Deux lemmes techniques}
\label{chow-section-tools-loc-chern}

\noindent
Dans cette section, nous développons quelques outils supplémentaires afin de pouvoir travailler
plus commodément avec les classes de Chern localisées. Le lemme suivant
est une version plus précise d'un fait que nous avons déjà rencontré dans
les démonstrations des lemmes \ref{chow-lemma-localized-chern-pre-compose} et
\ref{chow-lemma-localized-chern-pre-sum-c}.

\begin{lemma}
\label{chow-lemma-homomorphism-final}
Soit $(S, \delta)$ comme dans la situation \ref{chow-situation-setup}. Soit $X$
localement de type fini sur $S$. Soit $b : W \longrightarrow \mathbf{P}^1_X$
un morphisme propre de schémas. Soit $n \geq 1$. Pour $i = 1, \ldots, n$,
soit $Z_i \subset X$ un sous-schéma fermé, soit $Q_i \in D(\mathcal{O}_W)$
un objet parfait, soit $p_i \geq 0$ un entier et soient
$T_i \subset W_\infty$, $i = 1, \ldots, n$, des sous-schémas fermés.
Notons $W_i = b^{-1}(\mathbf{P}^1_{Z_i})$. Supposons que
\begin{enumerate}
\item pour $i = 1, \ldots, n$, l'hypothèse du
lemme \ref{chow-lemma-localized-chern-pre} soit satisfaite pour
$b, Z_i, Q_i, T_i, p_i$,
\item $Q_i|_{W \setminus W_i}$ soit nul, resp.\ soit isomorphe à un module
localement libre de rang $< p_i$ placé en degré cohomologique $0$,
\item $Q_i$ sur $W$ satisfasse
l'hypothèse (3) de la situation \ref{chow-situation-loc-chern}.
\end{enumerate}
Alors $P'_{p_n}(Q_n) \circ \ldots \circ P'_{p_1}(Q_1)$ est égal à
$$
(W_{n, \infty} \cap \ldots \cap W_{1, \infty} \to
Z_n \cap \ldots \cap Z_1)_* \circ
P'_{p_n}(Q_n|_{W_{n, \infty}}) \circ \ldots \circ P'_{p_1}(Q_1|_{W_{1, \infty}})
\circ C
$$
dans $A^{p_n + \ldots + p_1}(Z_n \cap \ldots \cap Z_1 \to X)$,
resp.\ $c'_{p_n}(Q_n) \circ \ldots \circ c'_{p_1}(Q_1)$ est égal à
$$
(W_{n, \infty} \cap \ldots \cap W_{1, \infty} \to
Z_n \cap \ldots \cap Z_1)_* \circ
c'_{p_n}(Q_n|_{W_{n, \infty}}) \circ \ldots \circ c'_{p_1}(Q_1|_{W_{1, \infty}})
\circ C
$$
dans $A^{p_n + \ldots + p_1}(Z_n \cap \ldots \cap Z_1 \to X)$.
\end{lemma}

\begin{proof}
Démontrons l'assertion sur les classes de Chern par récurrence sur $n$ ;
l'assertion sur $P_p(-)$ se démontre exactement de la même manière.
Le cas $n = 1$ est la construction de $c'_{p_1}(Q_1)$, car
$W_{1, \infty}$ est l'image inverse de $Z_1$ dans $W_\infty$.
Pour $n > 1$, nous avons par récurrence que
$c'_{p_n}(Q_n) \circ \ldots \circ c'_{p_1}(Q_1)$ est égal à
$$
c'_{p_n}(Q_n) \circ
(W_{n - 1, \infty} \cap \ldots \cap W_{1, \infty} \to
Z_{n - 1} \cap \ldots \cap Z_1)_* \circ
c'_{p_{n - 1}}(Q_{n - 1}|_{W_{n - 1}, \infty}) \circ \ldots \circ
c'_{p_1}(Q_1|_{W_{1, \infty}})
\circ C
$$
Par le lemme \ref{chow-lemma-base-change-localized-chern-pre}, la restriction de
$c'_{p_n}(Q_n)$ à $Z_{n - 1} \cap \ldots \cap Z_1$ est calculée par
le sous-ensemble fermé $Z_n \cap \ldots \cap Z_1$, le morphisme
$b' : W_{n - 1} \cap \ldots \cap W_1 \to
\mathbf{P}^1_{Z_{n - 1} \cap \ldots \cap Z_1}$
et la restriction de $Q_n$ à $W_{n - 1} \cap \ldots \cap W_1$.
Observons que $(b')^{-1}(Z_n) = W_n \cap \ldots \cap W_1$
et que $(W_n \cap \ldots \cap W_1)_\infty =
W_{n, \infty} \cap \ldots \cap W_{1, \infty}$.
Notons $C_{n - 1} \in A^0(W_{n - 1, \infty} \cap \ldots \cap W_{1, \infty} \to
Z_{n - 1} \cap \ldots \cap Z_1)$ la classe du lemme \ref{chow-lemma-gysin-at-infty}.
Nous concluons que la restriction de $c'_{p_n}(Q_n)$ à
$Z_{n - 1} \cap \ldots \cap Z_1$ est
\begin{align*}
&
(W_{n, \infty} \cap \ldots \cap W_{1, \infty} \to Z_n \cap \ldots \cap Z_1)_*
\circ
c'_{p_n}(Q_n|_{(W_n \cap \ldots \cap W_1)_\infty})
\circ
C_{n - 1} \\
& =
(W_{n, \infty} \cap \ldots \cap W_{1, \infty} \to Z_n \cap \ldots \cap Z_1)_*
\circ
c'_{p_n}(Q_n|_{W_{n, \infty}})
\circ
C_{n - 1}
\end{align*}
où l'égalité résulte du lemme \ref{chow-lemma-base-change-silly}
(nous omettons d'écrire la restriction à droite). Ainsi, l'expression ci-dessus devient
\begin{align*}
(W_{n, \infty} \cap \ldots \cap W_{1, \infty} \to
Z_n \cap \ldots \cap Z_1)_* \circ
c'_{p_n}(Q_n|_{W_n, \infty}) \circ \\
C_{n - 1} \circ
(W_{n - 1, \infty} \cap \ldots \cap W_{1, \infty} \to
Z_{n - 1} \cap \ldots \cap Z_1)_* \\
\circ
c'_{p_{n - 1}}(Q_{n - 1}|_{W_{n - 1}, \infty}) \circ \ldots \circ
c'_{p_1}(Q_1|_{W_{1, \infty}})
\circ C
\end{align*}
Par le lemme \ref{chow-lemma-homomorphism-pre},
nous savons que la composée
$C_{n - 1} \circ (W_{n - 1, \infty} \cap \ldots \cap W_{1, \infty} \to
Z_{n - 1} \cap \ldots \cap Z_1)_*$
est l'identité sur les éléments de l'image de l'homomorphisme de Gysin
$$
(W_{n - 1, \infty} \cap \ldots \cap W_{1, \infty} \to
W_{n - 1} \cap \ldots \cap W_1)^*
$$
Il suffit donc de montrer que tout élément de l'image de
$c'_{p_{n - 1}}(Q_{n - 1}|_{W_{n - 1}, \infty}) \circ \ldots \circ
c'_{p_1}(Q_1|_{W_{1, \infty}}) \circ C$
est dans l'image de l'homomorphisme de Gysin. Nous pouvons écrire
$$
c'_{p_i}(Q_i|_{W_{i, \infty}}) = \text{restriction de } c_{p_i}(W_i \to W, Q_i)
\text{ à } W_{i, \infty}
$$
par le lemme \ref{chow-lemma-loc-chern-agree} et les hypothèses (2) et (3) sur $Q_i$
dans l'énoncé du lemme. Ainsi, si $\beta \in \CH_{k + 1}(W)$
se restreint à l'image inverse plate de $\alpha$ sur $b^{-1}(\mathbf{A}^1_X)$,
alors
\begin{align*}
& c'_{p_{n - 1}}(Q_{n - 1}|_{W_{n - 1}, \infty}) \cap \ldots \cap
c'_{p_1}(Q_1|_{W_{1, \infty}})
\cap C \cap \alpha \\
& =
c'_{p_{n - 1}}(Q_{n - 1}|_{W_{n - 1}, \infty}) \cap \ldots \cap
c'_{p_1}(Q_1|_{W_{1, \infty}})
\cap i_\infty^* \beta \\
& =
c_{p_{n - 1}}(W_{n - 1} \to W, Q_{n - 1}) \cap \ldots \cap
c_{p_{n - 1}}(W_1 \to W, Q_1) \cap i_\infty^* \beta \\
& =
(W_{n - 1, \infty} \cap \ldots \cap W_{1, \infty} \to
W_{n - 1} \cap \ldots \cap W_1)^*
\left(c_{p_{n - 1}}(W_{n - 1} \to W, Q_{n - 1}) \cap \ldots \cap
c_{p_1}(W_1 \to W, Q_1) \cap \beta\right)
\end{align*}
comme désiré. En effet, pour la dernière égalité, nous utilisons le fait que
$c_{p_i}(W_i \to W, Q_i)$ est une classe bivariante et donc
commute à $i_\infty^*$ par définition.
\end{proof}

\noindent
Le lemme suivant nous donne une flexibilité considérable
si nous voulons calculer les classes de Chern localisées d'un complexe.

\begin{lemma}
\label{chow-lemma-independent-loc-chern-bQ}
Supposons que $(S, \delta), X, Z, b : W \to \mathbf{P}^1_X, Q, T, p$
satisfassent les hypothèses du lemme \ref{chow-lemma-localized-chern-pre}.
Soit $F \in D(\mathcal{O}_X)$ un objet parfait tel que
\begin{enumerate}
\item la restriction de $Q$ à $b^{-1}(\mathbf{A}^1_X)$ soit
isomorphe à l'image inverse de $F$,
\item $F|_{X \setminus Z}$ soit nul, resp.\ soit isomorphe à un
$\mathcal{O}_{X \setminus Z}$-module localement libre de rang $< p$
placé en degré cohomologique $0$, et
\item $Q$ sur $W$ et $F$ sur $X$ satisfassent l'hypothèse (3) de la
situation \ref{chow-situation-loc-chern}.
\end{enumerate}
Alors la classe $P'_p(Q)$, resp.\ $c'_p(Q)$ dans $A^p(Z \to X)$ construite
dans le lemme \ref{chow-lemma-localized-chern-pre}
est égale à $P_p(Z \to X, F)$, resp.\ $c_p(Z \to X, F)$
de la définition \ref{chow-definition-localized-chern}.
\end{lemma}

\begin{proof}
Les hypothèses sont préservées par changement de base par un morphisme
$X' \to X$ localement de type fini. Il suffit donc de montrer que
$P_p(Z \to X, F) \cap \alpha = P'_p(Q) \cap \alpha$,
resp.\ $c_p(Z \to X, F) \cap \alpha = c'_p(Q) \cap \alpha$
pour tout $\alpha \in \CH_k(X)$. Choisissons $\beta \in \CH_{k + 1}(W)$
dont la restriction à $b^{-1}(\mathbf{A}^1_X)$ est égale à
l'image inverse plate de $\alpha$, comme dans la construction de
$C$ dans le lemme \ref{chow-lemma-gysin-at-infty}.
Notons $W' = b^{-1}(Z)$ et notons
$E = W'_\infty \subset W_\infty$ l'image inverse de $Z$
par $W_\infty \to X$.
Le lemme résulte de
la suite d'égalités suivante (le cas de $P_p$ est analogue)
\begin{align*}
c'_p(Q) \cap \alpha
& =
(E \to Z)_*(c'_p(Q|_E) \cap i_\infty^*\beta) \\
& =
(E \to Z)_*(c_p(E \to W_\infty, Q|_{W_\infty}) \cap i_\infty^*\beta) \\
& =
(W'_\infty \to Z)_*(c_p(W' \to W, Q) \cap i_\infty^*\beta) \\
& =
(W'_\infty \to Z)_*((i'_\infty)^*(c_p(W' \to W, Q) \cap \beta)) \\
& =
(W'_\infty \to Z)_*((i'_\infty)^*(c_p(Z' \to X, F) \cap \beta)) \\
& =
(W'_0 \to Z)_*((i'_0)^*(c_p(Z' \to X, F) \cap \beta)) \\
& =
(W'_0 \to Z)_*(c_p(Z' \to X, F) \cap i_0^*\beta)) \\
& =
c_p(Z \to X, F) \cap \alpha
\end{align*}
La première égalité est la construction de $c'_p(Q)$ dans le
lemme \ref{chow-lemma-localized-chern-pre}.
La deuxième est le lemme \ref{chow-lemma-loc-chern-agree}.
Le changement de base de $W' \to W$ par $W_\infty \to W$ est le
morphisme $E = W'_\infty \to W_\infty$. La troisième égalité résulte donc
du lemme \ref{chow-lemma-base-change-loc-chern}. La quatrième
égalité, dans laquelle $i'_\infty : W'_\infty \to W'$ est le
morphisme d'inclusion, résulte du fait que $c_p(W' \to W, Q)$
est une classe bivariante. Pour la cinquième égalité, observons que
$c_p(W' \to W, Q)$ et $c_p(Z' \to X, F)$
se restreignent à la même classe bivariante dans
$A^p((b')^{-1} \to b^{-1}(\mathbf{A}^1_X))$ par
l'hypothèse (1) du lemme, qui dit que $Q$ et $F$ se restreignent
au même objet de $D(\mathcal{O}_{b^{-1}(\mathbf{A}^1_X)})$ ;
utilisons le lemme \ref{chow-lemma-base-change-loc-chern}.
Puisque $(i'_\infty)^*$ annule les cycles à support dans $W'_\infty$
(voir la remarque \ref{chow-remark-gysin-on-cycles}), nous concluons que la cinquième égalité
est vraie. La sixième égalité est satisfaite parce que $W'_\infty$ et $W'_0$
sont les images inverses des diviseurs de Cartier effectifs rationnellement équivalents
$D_0, D_\infty$ dans $\mathbf{P}^1_Z$ et donc $i_\infty^*\beta$ et
$i_0^*\beta$ s'envoient sur la même classe de cycle sur $W'$; plus précisément, tous deux
représentent la classe
$c_1(\mathcal{O}_{\mathbf{P}^1_Z}(1)) \cap c_p(Z \to X, F_) \cap \beta$ par le
lemme \ref{chow-lemma-support-cap-effective-Cartier}.
La septième égalité est satisfaite parce que $c_p(Z \to X, F)$ est
une classe bivariante. Par construction, $W'_0 = Z$ et $i_0^*\beta = \alpha$,
ce qui explique pourquoi la dernière égalité est satisfaite.
\end{proof}







\section{Propriétés des classes de Chern localisées}
\label{chow-section-properties-loc-chern}

\noindent
Les principaux résultats de cette section sont l'additivité et la multiplicativité
des classes de Chern localisées.

\begin{lemma}
\label{chow-lemma-loc-chern-character}
Dans la situation \ref{chow-situation-loc-chern}, supposons que $E|_{X \setminus Z}$ soit nul.
Alors
\begin{align*}
P_1(Z \to X, E) & = c_1(Z \to X, E), \\
P_2(Z \to X, E) & = c_1(Z \to X, E)^2 - 2c_2(Z \to X, E), \\
P_3(Z \to X, E) & = c_1(Z \to X, E)^3 - 3c_1(Z \to X, E)c_2(Z \to X, E)
+ 3c_3(Z \to X, E),
\end{align*}
et ainsi de suite, où les produits sont pris dans l'algèbre $A^{(1)}(Z \to X)$
de la remarque \ref{chow-remark-ring-loc-classes}.
\end{lemma}

\begin{proof}
L'énoncé a un sens parce que le faisceau nul est de rang $< 1$ et
donc les classes $c_p(Z \to X, E)$ sont définies pour tout $p \geq 1$.
Le résultat lui-même découle immédiatement du plus général
lemme \ref{chow-lemma-localized-chern-pre-compose}, puisque les classes de Chern localisées
ont été définies au moyen du procédé du
lemme \ref{chow-lemma-localized-chern-pre}
dans la section \ref{chow-section-localized-chern}.
\end{proof}

\begin{lemma}
\label{chow-lemma-loc-chern-classes-commute}
Dans la situation \ref{chow-situation-loc-chern},
soit $Y \to X$ localement de type fini et soit $c \in A^*(Y \to X)$.
Alors
$$
P_p(Z \to X, E) \circ c = c \circ P_p(Z \to X, E),
$$
respectivement
$$
c_p(Z \to X, E) \circ c = c \circ c_p(Z \to X, E)
$$
dans $A^*(Y \times_X Z \to X)$.
\end{lemma}

\begin{proof}
Cela résulte du lemme \ref{chow-lemma-homomorphism-commute}.
Plus précisément, soient
$$
b : W \to \mathbf{P}^1_X
\quad\text{et}\quad
Q
\quad\text{et}\quad
T' \subset T \subset W_\infty
$$
comme dans la démonstration du lemme \ref{chow-lemma-independent-loc-chern}.
Par définition, $c_p(Z \to X, E) = c'_p(Q)$
comme opérations bivariantes,
où le membre de droite est la classe bivariante construite dans le
lemme \ref{chow-lemma-localized-chern-pre} à l'aide de $W, b, Q, T'$.
Par le lemme \ref{chow-lemma-homomorphism-commute}, nous avons
$P'_p(Q) \circ c = c \circ P'_p(Q)$, resp.\ $c'_p(Q) \circ c = c \circ c'_p(Q)$
dans $A^*(Y \times_X Z \to X)$, et nous concluons.
\end{proof}

\begin{remark}
\label{chow-remark-loc-chern-classes}
Dans la situation \ref{chow-situation-loc-chern}, il est commode de définir
$$
c^{(p)}(Z \to X, E) = 1 + c_1(E) + \ldots + c_{p - 1}(E) +
c_p(Z \to X, E) + c_{p + 1}(Z \to X, E) + \ldots
$$
comme un élément de l'algèbre $A^{(p)}(Z \to X)$ considérée dans la
remarque \ref{chow-remark-ring-loc-classes}.
\end{remark}

\begin{lemma}
\label{chow-lemma-additivity-loc-chern-c}
Soit $(S, \delta)$ comme dans la situation \ref{chow-situation-setup}.
Soit $X$ localement de type fini sur $S$. Soit $Z \to X$
une immersion fermée. Soit
$$
E_1 \to E_2 \to E_3 \to E_1[1]
$$
un triangle distingué d'objets parfaits de $D(\mathcal{O}_X)$.
Supposons que
\begin{enumerate}
\item les restrictions $E_1|_{X \setminus Z}$ et $E_3|_{X \setminus Z}$
soient isomorphes à des $\mathcal{O}_{X \setminus Z}$-modules localement libres de rang fini
de rangs $< p_1$ et $< p_3$ placés en degré $0$, et
\item au moins l'une des conditions suivantes soit satisfaite :
(a) $X$ est quasi-compact,
(b) $X$ a des composantes irréductibles quasi-compactes,
(c) $E_3 \to E_1[1]$ peut être représenté par un morphisme de complexes
localement bornés de $\mathcal{O}_X$-modules localement libres de rang fini, ou
(d) il existe une enveloppe $f : Y \to X$ telle que $Lf^*E_3 \to Lf^*E_1[1]$
puisse être représenté par un morphisme de complexes localement bornés de
$\mathcal{O}_Y$-modules localement libres de rang fini.
\end{enumerate}
Avec la notation de la remarque \ref{chow-remark-loc-chern-classes}, nous avons
$$
c^{(p_1 + p_3)}(Z \to X, E_2) = c^{(p_1)}(Z \to X, E_1)c^{(p_3)}(Z \to X, E_3)
$$
dans $A^{(p_1 + p_3)}(Z \to X)$.
\end{lemma}

\begin{proof}
Observons que les hypothèses impliquent que $E_2|_{X \setminus Z}$ est nul,
resp.\ est isomorphe à un $\mathcal{O}_{X \setminus Z}$-module localement libre de rang fini
de rang $< p_1 + p_3$. L'énoncé a donc un sens.

\medskip\noindent
Soit $f : Y \to X$ une enveloppe. En développant les membres de gauche et de droite
de la formule dans l'énoncé du lemme, nous voyons que nous devons démontrer
certaines égalités de classes dans $A^*(X)$ et dans $A^*(Z \to X)$. Grâce à
l'unicité du lemme \ref{chow-lemma-envelope-bivariant}, il suffit de démontrer les
relations correspondantes dans $A^*(Y)$ et $A^*(Z \to Y)$. Puisque, de plus,
la construction des classes en jeu est compatible au changement de base
(lemme \ref{chow-lemma-base-change-loc-chern}), nous pouvons remplacer $X$ par $Y$
et le triangle distingué par son image inverse.

\medskip\noindent
Dans la démonstration du lemme \ref{chow-lemma-additivity-on-perfect}, nous avons
vu que les conditions (2)(a), (2)(b) et (2)(c) impliquent la condition
(2)(d). Combiné à la discussion du paragraphe précédent, cela nous
ramène au cas examiné dans le paragraphe suivant.

\medskip\noindent
Soit $\varphi^\bullet : \mathcal{E}_3^\bullet[-1] \to \mathcal{E}_1^\bullet$
un morphisme de complexes localement bornés de modules localement libres de rang fini
sur $\mathcal{O}_X$, représentant le morphisme $E_3[-1] \to E_1$
dans la catégorie dérivée. Considérons le schéma
$X' = \mathbf{A}^1 \times X$, de projection
$g : X' \to X$. Posons $Z' = g^{-1}(Z) = \mathbf{A}^1 \times Z$.
Notons $t$ la coordonnée sur $\mathbf{A}^1$. Considérons le cône
$\mathcal{C}^\bullet$ du morphisme de complexes
$$
t g^*\varphi^\bullet :
g^*\mathcal{E}_3^\bullet[-1]
\longrightarrow
g^*\mathcal{E}_1^\bullet
$$
sur $X'$. Nous obtenons un triangle distingué
$$
g^*\mathcal{E}_1^\bullet \to \mathcal{C}^\bullet \to
g^*\mathcal{E}_3^\bullet \to g^*\mathcal{E}_1^\bullet[1]
$$
où les trois premiers termes forment une suite exacte courte scindée terme à terme
de complexes. Manifestement, $\mathcal{C}^\bullet$ est un
complexe borné de $\mathcal{O}_{X'}$-modules localement libres de rang fini
dont la restriction à $X' \setminus Z'$ est isomorphe à un
module localement libre de rang fini
sur $\mathcal{O}_{X' \setminus Z'}$, de rang $< p_1 + p_3$,
placé en degré $0$. Nous avons donc les classes de Chern localisées
$$
c_p(Z' \to X', \mathcal{C}^\bullet) \in A^p(Z' \to X')
$$
pour $p \geq p_1 + p_3$. Pour tout $\alpha \in \CH_k(X)$, considérons
$$
c_p(Z' \to X', \mathcal{C}^\bullet) \cap g^*\alpha
\in \CH_{k + 1 - p}(\mathbf{A}^1 \times X)
$$
Si nous nous restreignons à $t = 0$, alors le morphisme $t g^*\varphi^\bullet$
se restreint à zéro et $\mathcal{C}^\bullet|_{t = 0}$
est la somme directe de $\mathcal{E}_1^\bullet$ et de $\mathcal{E}_3^\bullet$.
Par compatibilité des classes de Chern localisées au changement de base
(lemme \ref{chow-lemma-base-change-loc-chern}), nous concluons que
$$
i_0^* \circ c^{(p_1 + p_3)}(Z' \to X', \mathcal{C}^\bullet) \circ g^* =
c^{(p_1 + p_2)}(Z \to X, E_1 \oplus E_3)
$$
dans $A^{(p_1 + p_3)}(Z \to X)$. D'autre part, si nous nous restreignons à $t = 1$,
alors le morphisme $t g^*\varphi^\bullet$
se restreint à $\varphi$ et $\mathcal{C}^\bullet|_{t = 1}$
est un complexe borné de modules localement libres de rang fini représentant $E_2$.
Nous concluons que
$$
i_1^* \circ c^{(p_1 + p_3)}(Z' \to X', \mathcal{C}^\bullet) \circ g^* =
c^{(p_1 + p_2)}(Z \to X, E_2)
$$
dans $A^{(p_1 + p_3)}(Z \to X)$. Puisque $i_0^* = i_1^*$ par définition de
l'équivalence rationnelle (plus précisément, cela résulte des formules du
lemme \ref{chow-lemma-linebundle-formulae}), nous concluons que
$$
c^{(p_1 + p_2)}(Z \to X, E_2) = c^{(p_1 + p_2)}(Z \to X, E_1 \oplus E_3)
$$
Cela nous ramène au cas examiné dans le paragraphe suivant.

\medskip\noindent
Supposons que $E_2 = E_1 \oplus E_3$ et que les triplets $(X, Z, E_i)$
soient comme dans la situation \ref{chow-situation-loc-chern}.
Pour $i = 1, 3$, soient
$$
b_i : W_i \to \mathbf{P}^1_X
\quad\text{et}\quad
Q_i
\quad\text{et}\quad
T'_i \subset T_i \subset W_{i, \infty}
$$
comme dans la démonstration du lemme \ref{chow-lemma-independent-loc-chern}.
Par définition,
$$
c_p(Z \to X, E_i) = c'_p(Q_i)
$$
où le membre de droite est la classe bivariante construite dans le
lemme \ref{chow-lemma-localized-chern-pre} à l'aide de $W_i, b_i, Q_i, T'_i$.
Posons $W = W_1 \times_{b_1, \mathbf{P}^1_X, b_2} W_2$ et considérons
le diagramme cartésien
$$
\xymatrix{
W \ar[d]_{g_1} \ar[rd]^b \ar[r]_{g_3} & W_3 \ar[d]^{b_3} \\
W_1 \ar[r]^{b_1} & \mathbf{P}^1_X
}
$$
Bien sûr, $b^{-1}(\mathbf{A}^1)$ s'envoie isomorphiquement sur $\mathbf{A}^1_X$.
Observons que $T' = g_1^{-1}(T'_1) \cap g_2^{-1}(T'_2)$ contient encore
tous les points de $W_\infty$ situés au-dessus de $X \setminus Z$.
Par le lemme \ref{chow-lemma-localized-chern-pre-independent}, nous pouvons utiliser
$W$, $b$, $g_i^*\mathcal{Q}_i$ et
$T'$ pour construire $c_p(Z \to X, E_i)$ pour $i = 1, 3$.
De plus, par l'indépendance plus forte donnée dans le
lemme \ref{chow-lemma-independent-loc-chern-bQ}, nous pouvons utiliser
$W$, $b$, $g_1^*Q_1 \oplus g_3^*Q_3$ et $T'$
pour calculer les classes $c_p(Z \to X, E_2)$.
L'égalité désirée résulte donc du
lemme \ref{chow-lemma-localized-chern-pre-sum-c}.
\end{proof}

\begin{lemma}
\label{chow-lemma-additivity-loc-chern-P}
Soit $(S, \delta)$ comme dans la situation \ref{chow-situation-setup}.
Soit $X$ localement de type fini sur $S$. Soit $Z \to X$
une immersion fermée. Soit
$$
E_1 \to E_2 \to E_3 \to E_1[1]
$$
un triangle distingué d'objets parfaits de $D(\mathcal{O}_X)$.
Supposons que
\begin{enumerate}
\item les restrictions $E_1|_{X \setminus Z}$ et $E_3|_{X \setminus Z}$
soient nulles, et
\item au moins l'une des conditions suivantes soit satisfaite :
(a) $X$ est quasi-compact,
(b) $X$ a des composantes irréductibles quasi-compactes,
(c) $E_3 \to E_1[1]$ peut être représenté par un morphisme de complexes
localement bornés de $\mathcal{O}_X$-modules localement libres de rang fini, ou
(d) il existe une enveloppe $f : Y \to X$ telle que $Lf^*E_3 \to Lf^*E_1[1]$
puisse être représenté par un morphisme de complexes localement bornés de
$\mathcal{O}_Y$-modules localement libres de rang fini.
\end{enumerate}
Alors nous avons
$$
P_p(Z \to X, E_2) = P_p(Z \to X, E_1) + P_p(Z \to X, E_3)
$$
pour tout $p \in \mathbf{Z}$ et, par conséquent,
$ch(Z \to X, E_2) = ch(Z \to X, E_1) + ch(Z \to X, E_3)$.
\end{lemma}

\begin{proof}
La démonstration est exactement la même que celle du
lemme \ref{chow-lemma-additivity-loc-chern-c},
si ce n'est qu'elle utilise le
lemme \ref{chow-lemma-localized-chern-pre-sum-P}
à la toute fin. Pour $p > 0$, nous pouvons déduire ce lemme
du lemme \ref{chow-lemma-additivity-loc-chern-c} avec $p_1 = p_3 = 1$
et de la relation entre $P_p(Z \to X, E)$ et $c_p(Z \to X, E)$ donnée dans le
lemme \ref{chow-lemma-loc-chern-character}. Le cas $p = 0$ peut se démontrer
directement (il n'est intéressant que si $X$ a une composante connexe
entièrement contenue dans $Z$).
\end{proof}

\begin{lemma}
\label{chow-lemma-loc-chern-tensor-product}
Dans la situation \ref{chow-situation-setup}, soit $X$ localement de type fini sur $S$.
Soient $Z_i \subset X$, $i = 1, 2$, des sous-schémas fermés. Soient $F_i$, $i = 1, 2$,
des objets parfaits de $D(\mathcal{O}_X)$. Supposons, pour $i = 1, 2$, que
$F_i|_{X \setminus Z_i}$ soit nul\footnote{Il existe vraisemblablement une
variante de ce lemme où nous supposons seulement que $F_i|_{X \setminus Z_i}$
est isomorphe à un $\mathcal{O}_{X \setminus Z_i}$-module localement libre de rang fini
de rang $< p_i$.} et que $F_i$ sur $X$ satisfasse l'hypothèse
(3) de la situation \ref{chow-situation-loc-chern}. Notons
$r_i = P_0(Z_i \to X, F_i) \in A^0(Z_i \to X)$.
Alors nous avons
$$
c_1(Z_1 \cap Z_2 \to X, F_1 \otimes_{\mathcal{O}_X}^\mathbf{L} F_2) =
r_1 c_1(Z_2 \to X, F_2) + r_2 c_1(Z_1 \to X, F_1)
$$
dans $A^1(Z_1 \cap Z_2 \to X)$ et
\begin{align*}
c_2(Z_1 \cap Z_2 \to X, F_1 \otimes_{\mathcal{O}_X}^\mathbf{L} F_2)
& =
r_1 c_2(Z_2 \to X, F_2) +
r_2 c_2(Z_1 \to X, F_1) + \\
& {r_1 \choose 2} c_1(Z_2 \to X, F_2)^2 + \\
& (r_1r_2 - 1) c_1(Z_2 \to X, F_2)c_1(Z_1 \to X, F_1) + \\
& {r_2 \choose 2} c_1(Z_1 \to X, F_1)^2
\end{align*}
dans $A^2(Z_1 \cap Z_2 \to X)$, et ainsi de suite pour les classes de Chern supérieures.
De même, nous avons
$$
ch(Z_1 \cap Z_2 \to X, F_1 \otimes_{\mathcal{O}_X}^\mathbf{L} F_2) =
ch(Z_1 \to X, F_1) ch(Z_2 \to X, F_2)
$$
dans $\prod_{p \geq 0} A^p(Z_1 \cap Z_2 \to X) \otimes \mathbf{Q}$.
Plus précisément, nous avons
$$
P_p(Z_1 \cap Z_2 \to X, F_1 \otimes_{\mathcal{O}_X}^\mathbf{L} F_2) =
\sum\nolimits_{p_1 + p_2 = p}
{p \choose p_1} P_{p_1}(Z_1 \to X, F_1) P_{p_2}(Z_2 \to X, F_2)
$$
dans $A^p(Z_1 \cap Z_2 \to X)$.
\end{lemma}

\begin{proof}
Choisissons des morphismes propres $b_i : W_i \to \mathbf{P}^1_X$ et
$Q_i \in D(\mathcal{O}_{W_i})$, ainsi que des sous-schémas fermés
$T_i \subset W_{i, \infty}$, comme dans la construction des
classes de Chern localisées pour $F_i$, ou plus généralement comme dans le
lemme \ref{chow-lemma-independent-loc-chern-bQ}. Choisissons un diagramme
commutatif
$$
\xymatrix{
W \ar[d]^{g_1} \ar[rd]^b \ar[r]_{g_2} & W_2 \ar[d]^{b_2} \\
W_1 \ar[r]^{b_1} & \mathbf{P}^1_X
}
$$
où tous les morphismes sont propres et sont des isomorphismes au-dessus de
$\mathbf{A}^1_X$. Par exemple, nous pouvons prendre pour $W$ l'adhérence
du graphe de l'isomorphisme entre
$b_1^{-1}(\mathbf{A}^1_X)$ et $b_2^{-1}(\mathbf{A}^1_X)$.
Par le lemme \ref{chow-lemma-independent-loc-chern-bQ}, nous pouvons utiliser
$W$, $b = b_i \circ g_i$, $Lg_i^*Q_i$ et
$g_i^{-1}(T_i)$ pour construire les classes de Chern localisées
$c_p(Z_i \to X, F_i)$. Nous nous ramenons ainsi à la situation décrite
dans le paragraphe suivant.

\medskip\noindent
Supposons que nous ayons
\begin{enumerate}
\item un morphisme propre $b : W \to \mathbf{P}^1_X$ qui est un isomorphisme
au-dessus de $\mathbf{A}^1_X$,
\item $E_i \subset W_\infty$ est l'image inverse de $Z_i$,
\item des objets parfaits $Q_i \in D(\mathcal{O}_W)$ dont les classes de Chern
sont définies, tels que
\begin{enumerate}
\item la restriction de $Q_i$ à $b^{-1}(\mathbf{A}^1_X)$ soit
l'image inverse de $F_i$, et
\item il existe un sous-schéma fermé $T_i \subset W_\infty$ contenant
tous les points de $W_\infty$ situés au-dessus de $X \setminus Z_i$ tel que
$Q_i|_{T_i}$ soit nul.
\end{enumerate}
\end{enumerate}
Par le lemme \ref{chow-lemma-independent-loc-chern-bQ}, nous avons
$$
c_p(Z_i \to X, F_i) = c'_p(Q_i) =
(E_i \to Z_i)_* \circ c'_p(Q_i|_{E_i}) \circ C
$$
et
$$
P_p(Z_i \to X, F_i) = P'_p(Q_i) =
(E_i \to Z_i)_* \circ P'_p(Q_i|_{E_i}) \circ C
$$
pour $i = 1, 2$. Ensuite, nous observons que
$Q = Q_1 \otimes_{\mathcal{O}_W}^\mathbf{L} Q_2$
satisfait (3)(a) et (3)(b) pour $F_1 \otimes_{\mathcal{O}_X}^\mathbf{L} F_2$
et $T_1 \cup T_2$. Nous voyons donc que
$$
c_p(Z_1 \cap Z_2 \to X, F_1 \otimes_{\mathcal{O}_X}^\mathbf{L} F_2) =
(E_1 \cap E_2 \to Z_1 \cap Z_2)_* \circ
c'_p(Q|_{E_1 \cap E_2}) \circ C
$$
et
$$
P_p(Z_1 \cap Z_2 \to X, F_1 \otimes_{\mathcal{O}_X}^\mathbf{L} F_2) =
(E_1 \cap E_2 \to Z_1 \cap Z_2)_* \circ
P'_p(Q|_{E_1 \cap E_2}) \circ C
$$
par le même lemme. Par le lemme \ref{chow-lemma-silly-tensor-product},
les classes $c'_p(Q|_{E_1 \cap E_2})$ et $P'_p(Q|_{E_1 \cap E_2})$
peuvent être développées de la manière appropriée en fonction des classes
$c'_p(Q_i|_{E_i})$ et $P'_p(Q_i|_{E_i})$. Enfin, le
lemme \ref{chow-lemma-homomorphism-final}
nous dit que les polynômes en $c'_p(Q_i|_{E_i})$ et $P'_p(Q_i|_{E_i})$
coïncident avec les polynômes correspondants en
$c'_p(Q_i)$ et $P'_p(Q_i)$, comme désiré.
\end{proof}






\section{Éclatement à l'infini}
\label{chow-section-blowup-Z-first}

\noindent
Soit $X$ un schéma. Soit $Z \subset X$ un sous-schéma fermé découpé
par un faisceau quasi-cohérent d'idéaux de type fini. Notons $X' \to X$
l'éclatement de centre $Z$. Soit $b : W \to \mathbf{P}^1_X$
l'éclatement de centre $\infty(Z)$. Notons $E \subset W$ le diviseur
exceptionnel. Il existe un diagramme commutatif
$$
\xymatrix{
X' \ar[r] \ar[d] & W \ar[d]^b \\
X \ar[r]^\infty & \mathbf{P}^1_X
}
$$
dont les flèches horizontales sont des immersions fermées
(Diviseurs, lemme \ref{divisors-lemma-strict-transform}). Notons $E \subset W$
le diviseur exceptionnel et $W_\infty \subset W$ l'image inverse
de $(\mathbf{P}^1_X)_\infty$. Alors les assertions suivantes sont satisfaites :
\begin{enumerate}
\item $b$ est un isomorphisme au-dessus de
$\mathbf{A}^1_X \cup \mathbf{P}^1_{X \setminus Z}$,
\item $X'$ est un diviseur de Cartier effectif sur $W$,
\item $X' \cap E$ est le diviseur exceptionnel de $X' \to X$,
\item $W_\infty = X' + E$ comme diviseurs de Cartier effectifs sur $W$,
\item $E = \underline{\text{Proj}}_Z(\mathcal{C}_{Z/X, *}[S])$, où $S$
est une variable placée en degré $1$,
\item $X' \cap E = \underline{\text{Proj}}_Z(\mathcal{C}_{Z/X, *})$,
\item
\label{chow-item-cone-is-open}
$E \setminus X' = E \setminus (X' \cap E) =
\underline{\Spec}_Z(\mathcal{C}_{Z/X, *}) = C_ZX$,
\item
\label{chow-item-find-Z-in-blowup}
il existe une immersion fermée $\mathbf{P}^1_Z \to W$ dont la
composée avec $b$ est le morphisme d'inclusion
$\mathbf{P}^1_Z \to \mathbf{P}^1_X$ et dont le changement de base par $\infty$
est la composée $Z \to C_ZX \to E \to W_\infty$, où la première
flèche est le sommet du cône.
\end{enumerate}
Rappelons que $\mathcal{C}_{Z/X, *}$ est l'algèbre conormale de $Z$ dans $X$ ;
voir Diviseurs, définition \ref{divisors-definition-conormal-sheaf}, et
que $C_ZX$ est le cône normal de $Z$ dans $X$ ; voir
Diviseurs, définition \ref{divisors-definition-normal-cone}.

\medskip\noindent
Donnons maintenant la démonstration des assertions numérotées ci-dessus. Nous recommandons
vivement au lecteur d'étudier quelques exemples plutôt que de lire les
démonstrations.

\medskip\noindent
L'assertion (1) résulte de l'assertion correspondante de Diviseurs, lemme
\ref{divisors-lemma-blowing-up-gives-effective-Cartier-divisor}.
Observons que $E \subset W$ est un diviseur de Cartier effectif par le même lemme.

\medskip\noindent
Observons que $W_\infty$ est un diviseur de Cartier effectif par
Diviseurs, lemme \ref{divisors-lemma-blow-up-pullback-effective-Cartier}.
Puisque $E \subset W_\infty$, nous pouvons écrire $W_\infty = D + E$ pour un certain
diviseur de Cartier effectif $D$; voir
Diviseurs, lemme \ref{divisors-lemma-difference-effective-Cartier-divisors}.
Nous verrons ci-dessous que $D = X'$, ce qui démontrera (2) et (4).

\medskip\noindent
Puisque $X'$ est le transformé strict de l'immersion fermée
$\infty : X \to \mathbf{P}^1_X$ (voir ci-dessus), il s'ensuit que le diviseur
exceptionnel de $X' \to X$ est égal à l'intersection $X' \cap E$
(par exemple parce que tous deux sont découpés par l'image inverse du
faisceau d'idéaux de $Z$ sur $X'$). Cela démontre (3).

\medskip\noindent
L'intersection de $\infty(Z)$ avec $\mathbf{P}^1_Z$ est le diviseur de Cartier
effectif $(\mathbf{P}^1_Z)_\infty$ ; le transformé strict
de $\mathbf{P}^1_Z$ par l'éclatement $b$ s'envoie donc isomorphiquement sur
$\mathbf{P}^1_Z$ (voir Diviseurs, lemmes \ref{divisors-lemma-strict-transform}
et \ref{divisors-lemma-blow-up-effective-Cartier-divisor}).
Cela nous donne le morphisme $\mathbf{P}^1_Z \to W$ mentionné en (8).
C'est une immersion fermée puisque $b$ est séparé ; voir
Schémas, lemme \ref{schemes-lemma-section-immersion}.

\medskip\noindent
Supposons que $\Spec(A) \subset X$ soit un ouvert affine et que $Z \cap \Spec(A)$
corresponde à l'idéal de type fini $I \subset A$.
Un voisinage affine de $\infty(Z \cap \Spec(A))$ est un
espace affine sur $A$ de coordonnée $s = T_0/T_1$. Notons
$J = (I, s) \subset A[s]$ l'idéal engendré par $I$ et $s$.
Soit $B = A[s] \oplus J \oplus J^2 \oplus \ldots$ l'algèbre de Rees
de $(A[s], J)$. Observons que
$$
J^n =
I^n \oplus sI^{n - 1} \oplus s^2I^{n - 2} \ldots \oplus s^nA
\oplus s^{n + 1}A \oplus \ldots
$$
comme $A$-sous-module de $A[s]$ pour tout $n \geq 0$. Considérons le sous-schéma ouvert
$$
\text{Proj}(B) = \text{Proj}(A[s] \oplus J \oplus J^2 \oplus \ldots)
\subset W
$$
Enfin, notons $S$ l'élément $s \in J$ considéré comme un élément de degré $1$
de $B$.

\medskip\noindent
Puisque la formation de $\text{Proj}$ commute au changement de base
(Constructions, lemme \ref{constructions-lemma-base-change-map-proj}),
nous voyons que
$$
E = \text{Proj}(B \otimes_{A[s]} A/I) =
\text{Proj}((A/I \oplus I/I^2 \oplus I^2/I^3 \oplus \ldots)[S])
$$
La vérification de l'égalité $B \otimes_{A[s]} A/I = \bigoplus J^n/J^{n + 1}$
telle qu'elle a été donnée
résulte immédiatement de notre description des puissances $J^n$ ci-dessus.
Cela démontre (5), car l'algèbre conormale de $Z \cap \Spec(A)$
dans $\Spec(A)$ correspond à la $A$-algèbre graduée
$A/I \oplus I/I^2 \oplus I^2/I^3 \oplus \ldots$ par
Diviseurs, lemme \ref{divisors-lemma-affine-conormal-sheaf}.

\medskip\noindent
Rappelons que $\text{Proj}(B)$ est recouvert par les ouverts affines
$D_+(S)$ et $D_+(f^{(1)})$, pour $f \in I$, qui sont
les spectres des algèbres d'éclatement affines $A[s][\frac{J}{s}]$
et $A[s][\frac{J}{f}]$ ; voir
Diviseurs, lemme \ref{divisors-lemma-blowing-up-affine}, et
Algèbre, définition \ref{algebra-definition-blow-up}.
Nous allons décrire chacun de ces ouverts affines, ce qui achèvera la
démonstration.

\medskip\noindent
L'ouvert $D_+(S)$, c'est-à-dire le spectre de $A[s][\frac{J}{s}]$.
Il résulte de la description ci-dessus des puissances de $J$
que
$$
A[s][\textstyle{\frac{J}{s}}] = \sum s^{-n}I^n[s] \subset A[s, s^{-1}]
$$
L'élément $s$ est un non-diviseur de zéro dans cet anneau et définit aussi bien le diviseur
exceptionnel $E$ que $W_\infty$. Ainsi, $D \cap D_+(S) = \emptyset$.
Enfin, le quotient de $A[s][\frac{J}{s}]$ par $s$ est l'algèbre conormale
$$
A/I \oplus I/I^2 \oplus I^2/I^3 \oplus \ldots
$$
Cela démontre (7).

\medskip\noindent
L'ouvert $D_+(f^{(1)})$, c'est-à-dire le spectre de $A[s][\frac{J}{f}]$.
Il résulte de la description ci-dessus des puissances de $J$ que
$$
A[s][\textstyle{\frac{J}{f}}] =
A[\textstyle{\frac{I}{f}}][\textstyle{\frac{s}{f}}]
$$
où $\frac{s}{f}$ est une variable. L'élément $f$ est un non-diviseur de zéro
dans cet anneau, et son schéma des zéros définit le diviseur exceptionnel $E$.
Puisque $s$ définit $W_\infty$ et que $s = f \cdot \frac{s}{f}$,
nous concluons que $\frac{s}{f}$ définit
le diviseur $D$ construit ci-dessus. Nous voyons alors que
$$
D \cap D_+(f^{(1)}) = \Spec(A[\textstyle{\frac{I}{f}}])
$$
qui est l'ouvert correspondant de l'éclatement $X'$ au-dessus de $\Spec(A)$.
Plus précisément, le morphisme surjectif d'$A[s]$-algèbres graduées
$B \to A \oplus I \oplus I^2 \oplus \ldots$
vers l'algèbre de Rees de $(A, I)$ correspond à l'immersion
fermée $X' \to W$ au-dessus de $\Spec(A[s])$.
Cela démontre $D = X'$, comme désiré.

\medskip\noindent
Démontrons (6). Observons que le schéma des zéros de $\frac{s}{f}$
dans le paragraphe précédent est la restriction du schéma des zéros de $S$
à l'ouvert affine $D_+(f^{(1)})$. Nous voyons donc que $S = 0$ définit
$X' \cap E$ sur $E$. Ainsi, (6) résulte de (5).

\medskip\noindent
Enfin, nous devons démontrer la dernière partie de (8). C'est clair
car le morphisme $\mathbf{P}^1_Z \to W$ est décrit localement sur des ouverts affines
par la surjection
$$
B \to B \otimes_{A[s]} A/I =
(A/I \oplus I/I^2 \oplus I^2/I^3 \oplus \ldots)[S] \to
A/I[S]
$$
et l'identification $\text{Proj}(A/I[S]) = \Spec(A/I)$.
Certains détails sont omis.






\section{Homomorphismes de Gysin en codimension supérieure}
\label{chow-section-gysin-higher-codimension}

\noindent
Soit $(S, \delta)$ comme dans la situation \ref{chow-situation-setup}. Soit $X$ un schéma
localement de type fini sur $S$. Dans cette section, nous allons considérer des
triplets
$$
(Z \to X, \mathcal{N}, \sigma : \mathcal{N}^\vee \to \mathcal{C}_{Z/X})
$$
formés d'une immersion fermée $Z \to X$, d'un
$\mathcal{O}_Z$-module localement libre $\mathcal{N}$ et d'une surjection
$\sigma : \mathcal{N}^\vee \to \mathcal{C}_{Z/X}$ du dual
de $\mathcal{N}$ vers le faisceau conormal de $Z$ dans $X$ ; voir
Morphismes, section \ref{morphisms-section-conormal-sheaf}.
Nous dirons que
$\mathcal{N}$ est un {\it faisceau normal virtuel de $Z$ dans $X$}.

\begin{lemma}
\label{chow-lemma-pullback-virtual-normal-sheaf}
Soit $(S, \delta)$ comme dans la situation \ref{chow-situation-setup}. Soit
$$
\xymatrix{
Z' \ar[r] \ar[d]_g & X' \ar[d]^f \\
Z \ar[r] & X
}
$$
un diagramme cartésien de schémas localement de type fini sur $S$
dont les flèches horizontales sont des immersions fermées.
Si $\mathcal{N}$ est un faisceau normal virtuel pour $Z$ dans $X$, alors
$\mathcal{N}' = g^*\mathcal{N}$ est un faisceau normal virtuel pour
$Z'$ dans $X'$.
\end{lemma}

\begin{proof}
Cela résulte de la surjectivité de l'application
$g^*\mathcal{C}_{Z/X} \to \mathcal{C}_{Z'/X'}$ démontrée dans
Morphismes, lemme \ref{morphisms-lemma-conormal-functorial-flat}.
\end{proof}

\noindent
Soit $(S, \delta)$ comme dans la situation \ref{chow-situation-setup}. Soit $X$ un schéma
localement de type fini sur $S$. Soit $\mathcal{N}$ un fibré normal virtuel
pour une immersion fermée $Z \to X$. Dans cette situation, posons
$$
p : N = \underline{\Spec}_Z(\text{Sym}(\mathcal{N}^\vee)) \longrightarrow Z
$$
le fibré vectoriel sur $Z$
dont les sections correspondent aux sections de $\mathcal{N}$.
Dans cette situation, nous avons des immersions fermées canoniques
$$
C_ZX \longrightarrow N_ZX \longrightarrow N
$$
La première immersion fermée est donnée par Diviseurs, équation
(\ref{divisors-equation-normal-cone-in-normal-bundle}),
et la seconde immersion fermée correspond à la surjection
$\text{Sym}(\mathcal{N}^\vee) \to \text{Sym}(\mathcal{C}_{Z/X})$
induite par $\sigma$.
Soit
$$
b : W \longrightarrow \mathbf{P}^1_X
$$
l'éclatement en $\infty(Z)$ construit dans la
section \ref{chow-section-blowup-Z-first}. D'après le
lemme \ref{chow-lemma-gysin-at-infty},
nous avons une classe bivariante canonique dans
$$
C \in A^0(W_\infty \to X)
$$
Considérons l'immersion ouverte $j : C_ZX \to W_\infty$ de
(\ref{chow-item-cone-is-open}) et l'immersion fermée
$i : C_ZX \to N$ construite ci-dessus. D'après le lemme \ref{chow-lemma-vectorbundle},
pour tout $\alpha \in \CH_k(X)$, il existe un unique
$\beta \in \CH_*(Z)$ tel que
$$
i_*j^*(C \cap \alpha) = p^*\beta
$$
Posons $c(Z \to X, \mathcal{N}) \cap \alpha = \beta$.

\begin{lemma}
\label{chow-lemma-construction-gysin}
La construction ci-dessus définit une classe bivariante\footnote{La
notation $A^*(Z \to X)^\wedge$ est examinée dans la
remarque \ref{chow-remark-completion-bivariant}.
Si $X$ est quasi-compact, alors $A^*(Z \to X)^\wedge = A^*(Z \to X)$.}
$$
c(Z \to X, \mathcal{N}) \in A^*(Z \to X)^\wedge
$$
et, de plus, la construction est compatible au changement de base
comme dans le lemme \ref{chow-lemma-pullback-virtual-normal-sheaf}.
Si $\mathcal{N}$ est de rang constant $r$, alors
$c(Z \to X, \mathcal{N}) \in A^r(Z \to X)$.
\end{lemma}

\begin{proof}
Comme $i_* \circ j^* \circ C$ et $p^*$ sont tous deux des classes bivariantes
(voir les lemmes \ref{chow-lemma-flat-pullback-bivariant} et
\ref{chow-lemma-push-proper-bivariant}), nous pouvons utiliser l'égalité
$$
i_* \circ j^* \circ C  = p^* \circ c(Z \to X, \mathcal{N})
$$
(interprétée convenablement) pour définir $c(Z \to X, \mathcal{N})$
comme classe bivariante. Cela fonctionne parce que $p^*$ est toujours
bijectif sur les groupes de Chow d'après le lemme \ref{chow-lemma-vectorbundle}.

\medskip\noindent
Soient $X' \to X$, $Z' \to X'$ et $\mathcal{N}'$ comme dans le
lemme \ref{chow-lemma-pullback-virtual-normal-sheaf}. Écrivons
$c = c(Z \to X, \mathcal{N})$ et $c' = c(Z' \to X', \mathcal{N}')$.
La seconde assertion du lemme signifie que $c'$ est la restriction de $c$
comme dans la remarque \ref{chow-remark-restriction-bivariant}. Puisque nous affirmons que cela
vaut pour tout $X'/X$ localement de type fini, un argument formel
montre qu'il suffit de vérifier que $c' \cap \alpha' = c \cap \alpha'$
pour $\alpha' \in \CH_k(X')$.
Pour le voir, remarquons que nous avons un diagramme commutatif
$$
\xymatrix{
C_{Z'}X' \ar[d] \ar[r] &
W'_\infty \ar[d] \ar[r] &
W' \ar[d] \ar[r] &
\mathbf{P}^1_{X'} \ar[d] \\
C_ZX \ar[r] &
W_\infty \ar[r] &
W \ar[r] &
\mathbf{P}^1_X
}
$$
qui induit des immersions fermées :
$$
W' \to W \times_{\mathbf{P}^1_X} \mathbf{P}^1_{X'},\quad
W'_\infty \to W_\infty \times_X X',\quad
C_{Z'}X' \to C_ZX \times_Z Z'
$$
Pour obtenir $c \cap \alpha'$, nous utilisons la classe $C \cap \alpha'$
définie au moyen du morphisme
$W \times_{\mathbf{P}^1_X} \mathbf{P}^1_{X'} \to \mathbf{P}^1_{X'}$
dans le lemme \ref{chow-lemma-gysin-at-infty}.
Pour obtenir $c' \cap \alpha'$, d'autre part, nous utilisons la classe
$C' \cap \alpha'$ définie au moyen du morphisme $W' \to \mathbf{P}^1_{X'}$.
D'après le lemme \ref{chow-lemma-gysin-at-infty-independent}, l'image directe de
$C' \cap \alpha'$ par l'immersion fermée
$W'_\infty \to (W \times_{\mathbf{P}^1_X} \mathbf{P}^1_{X'})_\infty$
est égale à $C \cap \alpha'$. Il en va donc de même pour les images inverses
sur les ouverts
$$
C_{Z'}X' \subset W'_\infty,\quad
C_ZX \times_Z Z' \subset (W \times_{\mathbf{P}^1_X} \mathbf{P}^1_{X'})_\infty
$$
d'après le lemme \ref{chow-lemma-flat-pullback-proper-pushforward}.
Puisque nous avons un diagramme commutatif
$$
\xymatrix{
C_{Z'} X' \ar[d] \ar[r] & N' \ar@{=}[d] \\
C_ZX \times_Z Z' \ar[r] & N \times_Z Z'
}
$$
les images directes de ces classes donnent la même classe sur $N'$, ce qui
prouve que nous obtenons le même élément $c \cap \alpha' = c' \cap \alpha'$
dans $\CH_*(Z')$.
\end{proof}

\begin{lemma}
\label{chow-lemma-gysin-decompose}
Soit $(S, \delta)$ comme dans la situation \ref{chow-situation-setup}. Soit $X$ un schéma
localement de type fini sur $S$. Soit $\mathcal{N}$ un faisceau normal virtuel
pour un sous-schéma fermé $Z$ de $X$. Supposons que nous ayons une suite
exacte courte $0 \to \mathcal{N}' \to \mathcal{N} \to \mathcal{E} \to 0$
de $\mathcal{O}_Z$-modules localement libres de rang fini telle que la surjection donnée
$\sigma : \mathcal{N}^\vee \to \mathcal{C}_{Z/X}$ se factorise par une application
$\sigma' : (\mathcal{N}')^\vee \to \mathcal{C}_{Z/X}$.
Alors
$$
c(Z \to X, \mathcal{N}) = c_{top}(\mathcal{E}) \circ c(Z \to X, \mathcal{N}')
$$
comme classes bivariantes.
\end{lemma}

\begin{proof}
Notons $N' \to N$ l'immersion fermée de fibrés vectoriels correspondant
à la surjection $\mathcal{N}^\vee \to (\mathcal{N}')^\vee$. Nous avons alors
des immersions fermées
$$
C_ZX \to N' \to N
$$
Ainsi, la relation voulue entre les classes bivariantes résulte
immédiatement du lemme \ref{chow-lemma-easy-virtual-class}.
\end{proof}

\begin{lemma}
\label{chow-lemma-gysin-excess}
Soit $(S, \delta)$ comme dans la situation \ref{chow-situation-setup}. Considérons
un diagramme cartésien
$$
\xymatrix{
Z' \ar[r] \ar[d]_g & X' \ar[d]^f \\
Z \ar[r] & X
}
$$
de schémas localement de type fini sur $S$, dont les flèches horizontales
sont des immersions fermées. Soient $\mathcal{N}$, resp.\ $\mathcal{N}'$,
un faisceau normal virtuel pour $Z \subset X$, resp.\ $Z' \to X'$.
Supposons donnée une suite exacte courte
$0 \to \mathcal{N}' \to g^*\mathcal{N} \to \mathcal{E} \to 0$
de modules localement libres de rang fini sur $Z'$, telle que le diagramme
$$
\xymatrix{
g^*\mathcal{N}^\vee \ar[r] \ar[d] &
(\mathcal{N}')^\vee \ar[d] \\
g^*\mathcal{C}_{Z/X} \ar[r] &
\mathcal{C}_{Z'/X'}
}
$$
soit commutatif. Alors nous avons
$$
res(c(Z \to X, \mathcal{N})) =
c_{top}(\mathcal{E}) \circ c(Z' \to X', \mathcal{N}')
$$
dans $A^*(Z' \to X')^\wedge$.
\end{lemma}

\begin{proof}
D'après le lemme \ref{chow-lemma-construction-gysin}, nous avons
$res(c(Z \to X, \mathcal{N})) = c(Z' \to X', g^*\mathcal{N})$,
et l'égalité résulte du lemme \ref{chow-lemma-gysin-decompose}.
\end{proof}

\noindent
Soit $(S, \delta)$ comme dans la situation \ref{chow-situation-setup}. Soit $X$ un schéma
localement de type fini sur $S$. Soit $\mathcal{N}$ un faisceau normal virtuel
pour un sous-schéma fermé $Z$ de $X$. Soit $Y \to X$ un morphisme
localement de type fini. Supposons que $Z \times_X Y \to Y$ soit une
immersion fermée régulière, voir
Diviseurs, section \ref{divisors-section-regular-immersions}.
Dans ce cas, le faisceau conormal $\mathcal{C}_{Z \times_X Y/Y}$ est un
$\mathcal{O}_{Z \times_X Y}$-module localement libre de rang fini, et nous obtenons une suite
exacte courte
$$
0 \to \mathcal{E}^\vee \to
\mathcal{N}^\vee|_{Z \times_X Y} \to \mathcal{C}_{Z \times_X Y/Y} \to 0
$$
Le quotient $\mathcal{N}|_{Y \times_X Z} \to \mathcal{E}$ est appelé le
{\it faisceau normal d'excès} de la situation.

\begin{lemma}
\label{chow-lemma-gysin-fundamental}
Dans la situation décrite juste ci-dessus, supposons que $\dim_\delta(Y) = n$
et que $\mathcal{C}_{Y \times_X Z/Z}$ soit de rang constant $r$.
Alors
$$
c(Z \to X, \mathcal{N}) \cap [Y]_n =
c_{top}(\mathcal{E}) \cap [Z \times_X Y]_{n - r}
$$
dans $\CH_*(Z \times_X Y)$.
\end{lemma}

\begin{proof}
La classe bivariante $c_{top}(\mathcal{E}) \in A^*(Z \times_X Y)$ a été
définie dans la remarque \ref{chow-remark-top-chern-class}.
D'après le lemme \ref{chow-lemma-construction-gysin}, nous pouvons remplacer $X$ par $Y$.
Ainsi, nous pouvons supposer que $Z \to X$ est une immersion fermée régulière
de codimension $r$, que $\dim_\delta(X) = n$, et nous devons
montrer que $c(Z \to X, \mathcal{N}) \cap [X]_n =
c_{top}(\mathcal{E}) \cap [Z]_{n - r}$ dans $\CH_*(Z)$.
D'après le lemme \ref{chow-lemma-gysin-decompose}, nous pouvons même supposer que
$\mathcal{N}^\vee \to \mathcal{C}_{Z/X}$ est un isomorphisme.
Autrement dit, nous devons montrer que
$c(Z \to X, \mathcal{C}_{Z/X}^\vee) \cap [X]_n = [Z]_{n - r}$ dans $\CH_*(Z)$.

\medskip\noindent
Reprenons les étapes de la définition de
$c(Z \to X, \mathcal{C}_{Z/X}^\vee) \cap [X]_n$. Soit
$b : W \to \mathbf{P}^1_X$
l'éclatement de $\infty(Z)$. Nous devons d'abord calculer
$C \cap [X]_n$, où $C \in A^0(W_\infty \to X)$ est
la classe du lemme \ref{chow-lemma-gysin-at-infty}.
Pour ce faire, remarquons que $[W]_{n + 1}$
est un cycle sur $W$ dont la restriction à $\mathbf{A}^1_X$ est
égale à l'image inverse plate de $[X]_n$. Ainsi, $C \cap [X]_n$
est égal à $i_\infty^*[W]_{n + 1}$. Comme $W_\infty$ est un
diviseur de Cartier effectif sur $W$, nous avons
$i_\infty^*[W]_{n + 1} = [W_\infty]_n$, voir le lemme \ref{chow-lemma-easy-gysin}.
La restriction de cette classe à l'ouvert $C_ZX \subset W_\infty$
est bien sûr simplement $[C_ZX]_n$. Comme $Z \subset X$ est régulièrement
immergé, nous avons
$$
\mathcal{C}_{Z/X, *} = \text{Sym}(\mathcal{C}_{Z/X})
$$
comme $\mathcal{O}_Z$-algèbres graduées, voir
Diviseurs, lemme \ref{divisors-lemma-quasi-regular-immersion}.
Ainsi, $p : N = C_ZX \to Z$ est le morphisme structural du
fibré vectoriel associé au module localement libre de rang fini
$\mathcal{C}_{Z/X}$ de rang $r$. Il est alors clair que
$p^*[Z]_{n - r} = [C_ZX]_n$, et la démonstration est achevée.
\end{proof}

\begin{lemma}
\label{chow-lemma-gysin-easy}
Soit $(S, \delta)$ comme dans la situation \ref{chow-situation-setup}. Soit $X$ un schéma
localement de type fini sur $S$. Soit $\mathcal{N}$ un faisceau normal virtuel
pour un sous-schéma fermé $Z$ de $X$. Soit $Y \to X$ un morphisme
localement de type fini. Étant donnés des entiers $r$, $n$, supposons que
\begin{enumerate}
\item $\mathcal{N}$ soit localement libre de rang $r$,
\item toute composante irréductible de $Y$ soit de $\delta$-dimension $n$,
\item $\dim_\delta(Z \times_X Y) \leq n - r$, et
\item pour $\xi \in Z \times_X Y$ tel que $\delta(\xi) = n - r$,
l'anneau local $\mathcal{O}_{Y, \xi}$ soit de Cohen-Macaulay.
\end{enumerate}
Alors $c(Z \to X, \mathcal{N}) \cap [Y]_n = [Z \times_X Y]_{n - r}$
dans $\CH_{n - r}(Z \times_X Y)$.
\end{lemma}

\begin{proof}
L'énoncé a un sens puisque $Z \times_X Y$ est un sous-schéma fermé de $Y$.
Comme $\mathcal{N}$ est de rang $r$, nous savons que
$c(Z \to X, \mathcal{N}) \cap [Y]_n$ appartient à $\CH_{n - r}(Z \times_X Y)$.
Puisque $\dim_\delta(Z \cap Y) \leq n - r$, le groupe de Chow
$\CH_{n - r}(Z \times_X Y)$ est librement engendré par les
classes de cycles des composantes irréductibles $W \subset Z \times_X Y$
de $\delta$-dimension $n - r$. Soit $\xi \in W$ le point générique.
D'après l'hypothèse (2), nous voyons que $\dim(\mathcal{O}_{Y, \xi}) = r$.
D'autre part, comme $\mathcal{N}$ est de rang $r$ et que
$\mathcal{N}^\vee \to \mathcal{C}_{Z/X}$ est surjectif, nous voyons que
le faisceau d'idéaux de $Z$ est localement défini par $r$ équations.
Ainsi, le faisceau quasi-cohérent d'idéaux $\mathcal{I} \subset \mathcal{O}_Y$
de $Z \times_X Y$ dans $Y$ est localement engendré par $r$ éléments.
Comme $\mathcal{O}_{Y, \xi}$ est de Cohen-Macaulay de dimension $r$
et que $\mathcal{I}_\xi$ est un idéal de définition (puisque $\xi$ est
un point générique de $Z \times_X Y$), il s'ensuit que $\mathcal{I}_\xi$
est engendré par une suite régulière
(Algèbre, lemme \ref{algebra-lemma-reformulate-CM}).
D'après Diviseurs, lemme \ref{divisors-lemma-Noetherian-scheme-regular-ideal},
nous voyons que $\mathcal{I}$ est engendré par une suite régulière sur
un voisinage ouvert $V \subset Y$ de $\xi$. D'après notre description de
$\CH_{n - r}(Z \times_X Y)$, il suffit de montrer que
$c(Z \to X, \mathcal{N}) \cap [V]_n = [Z \times_X V]_{n - r}$
dans $\CH_{n - r}(Z \times_X V)$. Cela résulte du
lemme \ref{chow-lemma-gysin-fundamental},
car le faisceau normal d'excès est $0$ sur $V$.
\end{proof}

\begin{lemma}
\label{chow-lemma-gysin-agrees}
Soit $(S, \delta)$ comme dans la situation \ref{chow-situation-setup}. Soit $X$ un schéma
localement de type fini sur $S$. Soit $(\mathcal{L}, s, i : D \to X)$
un triplet comme dans la définition \ref{chow-definition-gysin-homomorphism}.
L'homomorphisme de Gysin $i^*$, considéré comme un élément de $A^1(D \to X)$
(voir le lemme \ref{chow-lemma-gysin-bivariant}), est égal à la classe bivariante
$c(D \to X, \mathcal{N}) \in A^1(D \to X)$
construite à l'aide de $\mathcal{N} = i^*\mathcal{L}$,
considéré comme un faisceau normal virtuel pour $D$ dans $X$.
\end{lemma}

\begin{proof}
Nous utiliserons le critère du lemme \ref{chow-lemma-bivariant-zero}.
Ainsi, nous pouvons supposer que $X$ est un schéma intègre et
nous devons montrer que $i^*[X]$ est égal à $c \cap [X]$.
Soit $n = \dim_\delta(X)$. Comme d'habitude, il y a deux cas.

\medskip\noindent
Si $X = D$, nous voyons que les deux classes sont représentées par
$c_1(\mathcal{N}) \cap [X]_n$. Voir le lemme \ref{chow-lemma-gysin-fundamental}
et la définition \ref{chow-definition-gysin-homomorphism}.

\medskip\noindent
Si $D \not = X$, alors $D \to X$ est un diviseur de Cartier effectif
et, en particulier, une immersion fermée régulière de codimension $1$.
De nouveau d'après le lemme \ref{chow-lemma-gysin-fundamental}, nous concluons que
$c(D \to X, \mathcal{N}) \cap [X]_n = [D]_{n - 1}$. Par définition, il en va
de même pour l'homomorphisme de Gysin, et nous concluons
une nouvelle fois.
\end{proof}

\begin{lemma}
\label{chow-lemma-gysin-commutes}
Soit $(S, \delta)$ comme dans la situation \ref{chow-situation-setup}. Soit $X$ un schéma
localement de type fini sur $S$. Soit $Z \subset X$ un sous-schéma fermé
muni d'un faisceau normal virtuel $\mathcal{N}$. Soit $Y \to X$ localement de
type fini et soit $c \in A^*(Y \to X)$. Alors $c$ et $c(Z \to X, \mathcal{N})$
commutent (remarque \ref{chow-remark-bivariant-commute}).
\end{lemma}

\begin{proof}
Pour le vérifier, nous pouvons utiliser le lemme \ref{chow-lemma-bivariant-zero}.
Ainsi, nous pouvons supposer que $X$ est un schéma intègre, et nous devons montrer
$c \cap c(Z \to X, \mathcal{N}) \cap [X] =
c(Z \to X, \mathcal{N}) \cap c \cap [X]$ dans $\CH_*(Z \times_X Y)$.

\medskip\noindent
Si $Z = X$, alors $c(Z \to X, \mathcal{N}) = c_{top}(\mathcal{N})$ d'après le
lemme \ref{chow-lemma-gysin-fundamental}, et cette classe commute
avec la classe bivariante $c$, voir le lemme \ref{chow-lemma-cap-commutative-chern}.

\medskip\noindent
Supposons que $Z$ ne soit pas égal à $X$. D'après le lemme \ref{chow-lemma-bivariant-zero},
il suffit même de prouver le résultat après avoir éclaté $X$ (le long d'un idéal non nul).
Éclatons $X$ le long du faisceau d'idéaux de $Z$. Cela nous ramène au cas
où $Z$ est un diviseur de Cartier effectif, voir
Diviseurs, lemme
\ref{divisors-lemma-blowing-up-gives-effective-Cartier-divisor},

\medskip\noindent
Si $Z$ est un diviseur de Cartier effectif, alors nous avons
$$
c(Z \to X, \mathcal{N}) =
c_{top}(\mathcal{E}) \circ i^*
$$
où $i^* \in A^1(Z \to X)$ est l'homomorphisme de Gysin
associé à $i : Z \to X$ (lemme \ref{chow-lemma-gysin-bivariant})
et $\mathcal{E}$ est le dual du noyau de
$\mathcal{N}^\vee \to \mathcal{C}_{Z/X}$, voir les
lemmes \ref{chow-lemma-gysin-decompose} et \ref{chow-lemma-gysin-agrees}.
Nous concluons alors parce que les classes de Chern sont dans le centre de
l'anneau bivariant (au sens fort formulé dans le
lemme \ref{chow-lemma-cap-commutative-chern}) et que $c$ commute
avec l'homomorphisme de Gysin $i^*$ par définition des classes bivariantes.
\end{proof}

\noindent
Soit $(S, \delta)$ comme dans la situation \ref{chow-situation-setup}. Soit $X$ un
schéma intègre localement de type fini sur $S$, de $\delta$-dimension $n$.
Soient $Z \subset Y \subset X$ des sous-schémas fermés qui sont tous deux des
diviseurs de Cartier effectifs dans $X$. Notons $o : Y \to C_Y X$ la section nulle du
cône normal en droites de $Y$ dans $X$. Comme $C_YX$ est un fibré en droites sur $Y$,
nous obtenons une classe bivariante $o^* \in A^1(Y \to C_YX)$, voir le
lemme \ref{chow-lemma-gysin-bivariant}.

\begin{lemma}
\label{chow-lemma-relation-normal-cones}
Avec les notations ci-dessus, nous avons
$$
o^*[C_ZX]_n = [C_Z Y]_{n - 1}
$$
dans $\CH_{n - 1}(Y \times_{o, C_Y X} C_ZX)$.
\end{lemma}

\begin{proof}
Notons $W \to \mathbf{P}^1_X$ l'éclatement de $\infty(Z)$ comme dans la
section \ref{chow-section-blowup-Z-first}.
De même, notons $W' \to \mathbf{P}^1_X$ l'éclatement de $\infty(Y)$.
Comme $\infty(Z) \subset \infty(Y)$, nous obtenons une inclusion opposée
des faisceaux d'idéaux, donc une application des algèbres graduées
définissant ces éclatements. Cela produit un morphisme rationnel de $W$
vers $W'$ qui admet en fait un représentant canonique
$$
W \supset U \longrightarrow W'
$$
Voir Constructions, lemme \ref{constructions-lemma-morphism-relative-proj}.
Un calcul local (omis) montre que $U$ contient au moins tous les points
de $W$ qui ne sont pas au-dessus de $\infty$, ainsi que le sous-schéma ouvert $C_Z X$ de la fibre
spéciale. Quitte à restreindre $U$, nous pouvons supposer que $U_\infty = C_Z X$ et que
$\mathbf{A}^1_X \subset U$. Un autre calcul local (omis)
montre que le morphisme $U_\infty \to W'_\infty$
induit le morphisme canonique $C_Z X \to C_Y X \subset W'_\infty$
de cônes normaux induit par l'inclusion des faisceaux d'idéaux
provenant de $Z \subset Y$. Notons $W'' \subset W$ la transformée stricte de
$\mathbf{P}^1_Y \subset \mathbf{P}^1_X$ dans $W$. Alors $W''$ est l'éclatement
de $\mathbf{P}^1_Y$ en $\infty(Z)$ d'après
Diviseurs, lemme \ref{divisors-lemma-strict-transform},
et donc $(W'' \cap U)_\infty = C_ZY$.

\medskip\noindent
Considérons le diviseur de Cartier effectif $i : \mathbf{P}^1_Y \to W'$
de (\ref{chow-item-find-Z-in-blowup}) et sa classe bivariante associée
$i^* \in A^1(\mathbf{P}^1_Y \to W')$ donnée par le lemme \ref{chow-lemma-gysin-bivariant}.
Notons de même $(i'_\infty)^* \in A^1(W'_\infty \to W')$
l'application de Gysin à l'infini. Remarquons que la restriction de $i'_\infty$
(remarque \ref{chow-remark-restriction-bivariant}) à $U$ est la restriction de
$i_\infty^* \in A^1(W_\infty \to W)$ à $U$. D'une part, nous avons
$$
(i'_\infty)^* i^* [U]_{n + 1} =
i_\infty^* i^* [U]_{n + 1} =
i_\infty^* [(W'' \cap U)_\infty]_{n + 1} =
[C_ZY]_n
$$
parce que $i_\infty^*$ annule toutes les classes à support au-dessus de $\infty$, parce que
$i^*[U]$ et $[W'']$ coïncident comme cycles sur $\mathbf{A}^1$, et parce que
$C_ZY$ est la fibre de $W'' \cap U$ au-dessus de $\infty$.
D'autre part, nous avons
$$
(i'_\infty)^* i^* [U]_{n + 1} =
i^* i_\infty^*[U]_{n + 1} =
i^* [U_\infty] =
o^*[C_YX]_n
$$
parce que $(i'_\infty)^*$ et $i^*$ commutent
(lemme \ref{chow-lemma-gysin-commutes-gysin})
et parce que la fibre de $i : \mathbf{P}^1_Y \to W'$ au-dessus de $\infty$
se factorise par $o : Y \to C_YX$ et l'immersion ouverte $C_YX \to W'_\infty$.
Le lemme en résulte.
\end{proof}

\begin{lemma}
\label{chow-lemma-gysin-composition}
Soit $(S, \delta)$ comme dans la situation \ref{chow-situation-setup}.
Soient $Z \subset Y \subset X$ des sous-schémas fermés d'un schéma localement
de type fini sur $S$.
Soit $\mathcal{N}$ un faisceau normal virtuel pour $Z \subset X$.
Soit $\mathcal{N}'$ un faisceau normal virtuel pour $Z \subset Y$.
Soit $\mathcal{N}''$ un faisceau normal virtuel pour $Y \subset X$.
Supposons qu'il existe un diagramme commutatif
$$
\xymatrix{
(\mathcal{N}'')^\vee|_Z \ar[r] \ar[d] &
\mathcal{N}^\vee \ar[r] \ar[d] &
(\mathcal{N}')^\vee \ar[d] \\
\mathcal{C}_{Y/X}|_Z \ar[r] &
\mathcal{C}_{Z/X} \ar[r] &
\mathcal{C}_{Z/Y}
}
$$
où la suite du bas provient de Compléments sur les morphismes, lemme
\ref{more-morphisms-lemma-transitivity-conormal}, et où la
suite du haut est une suite exacte courte. Alors
$$
c(Z \to X, \mathcal{N}) =
c(Z \to Y, \mathcal{N}') \circ c(Y \to X, \mathcal{N}'')
$$
dans $A^*(Z \to X)^\wedge$.
\end{lemma}

\begin{proof}
Remarquons que les hypothèses restent satisfaites après tout changement de base
par un morphisme $X' \to X$ localement de type fini (la suite
exacte courte de faisceaux normaux virtuels est localement scindée et
reste donc exacte après tout changement de base). Ainsi, pour vérifier
l'égalité des classes bivariantes, nous pouvons utiliser le lemme \ref{chow-lemma-bivariant-zero}.
Nous pouvons donc supposer que $X$ est un schéma intègre, et nous devons montrer
$c(Z \to X, \mathcal{N}) \cap [X] =
c(Z \to Y, \mathcal{N}') \cap c(Y \to X, \mathcal{N}'') \cap [X]$.

\medskip\noindent
Si $Y = X$, alors nous avons
\begin{align*}
c(Z \to Y, \mathcal{N}') \cap c(Y \to X, \mathcal{N}'') \cap [X]
& =
c(Z \to Y, \mathcal{N}') \cap c_{top}(\mathcal{N}'') \cap [Y] \\
& =
c_{top}(\mathcal{N}''|_Z) \cap c(Z \to Y, \mathcal{N}') \cap [Y] \\
& =
c(Z \to X, \mathcal{N}) \cap [X] 
\end{align*}
La première égalité résulte du lemme \ref{chow-lemma-gysin-decompose}.
La deuxième vient de ce que les classes de Chern commutent avec les classes bivariantes
(lemme \ref{chow-lemma-cap-commutative-chern}).
La troisième égalité résulte du lemme \ref{chow-lemma-gysin-decompose}.

\medskip\noindent
Supposons $Y \not = X$. D'après le lemme \ref{chow-lemma-bivariant-zero},
il suffit même de prouver le résultat après avoir éclaté $X$ en un idéal non nul.
Éclatons $X$ le long du produit du faisceau d'idéaux de $Y$ et du faisceau
d'idéaux de $Z$. Cela nous ramène au cas où $Y$ et $Z$ sont tous deux des
diviseurs de Cartier effectifs sur $X$, voir
Diviseurs, lemmes
\ref{divisors-lemma-blowing-up-gives-effective-Cartier-divisor} et
\ref{divisors-lemma-blowing-up-two-ideals}.

\medskip\noindent
Notons $\mathcal{N}'' \to \mathcal{E}$ la surjection de
$\mathcal{O}_Z$-modules localement libres de rang fini telle que
$0 \to \mathcal{E}^\vee \to (\mathcal{N}'')^\vee \to \mathcal{C}_{Y/X} \to 0$
soit une suite exacte courte. Alors $\mathcal{N} \to \mathcal{E}|_Z$
est également une surjection. Notons $\mathcal{N}_1$ le noyau localement libre de rang fini
de cette application et remarquons que $\mathcal{N}^\vee \to \mathcal{C}_{Z/X}$
se factorise par $\mathcal{N}_1$.
D'après le lemme \ref{chow-lemma-gysin-decompose}, nous avons
$$
c(Y \to X, \mathcal{N}'') = c_{top}(\mathcal{E}) \circ
c(Y \to X, \mathcal{C}_{Y/X}^\vee)
$$
et
$$
c(Z \to X, \mathcal{N}) = c_{top}(\mathcal{E}|_Z) \circ
c(Z \to X, \mathcal{N}_1)
$$
Comme les classes de Chern de fibrés commutent avec les classes bivariantes
(lemme \ref{chow-lemma-cap-commutative-chern}),
il suffit de prouver
$$
c(Z \to X, \mathcal{N}_1) =
c(Z \to Y, \mathcal{N}') \circ c(Y \to X, \mathcal{C}_{Y/X}^\vee)
$$
dans $A^*(Z \to X)$. Pour cela, nous pouvons supposer que $\mathcal{N}'' = \mathcal{C}_{Y/X}$.
Cela nous ramène au cas examiné dans l'alinéa suivant.

\medskip\noindent
Dans cet alinéa, $Z$ et $Y$ sont des diviseurs de Cartier effectifs sur $X$,
qui est intègre de dimension $n$, et nous avons $\mathcal{N}'' = \mathcal{C}_{Y/X}$.
Dans ce cas, $c(Y \to X, \mathcal{C}_{Y/X}^\vee) \cap [X] = [Y]_{n - 1}$ d'après le
lemme \ref{chow-lemma-gysin-fundamental}. Nous devons donc prouver que
$c(Z \to X, \mathcal{N}) \cap [X] = c(Z \to Y, \mathcal{N}') \cap [Y]_{n - 1}$.
Notons $N$ et $N'$ les fibrés vectoriels sur $Z$ associés à
$\mathcal{N}$ et $\mathcal{N}'$. Considérons le diagramme commutatif
$$
\xymatrix{
N' \ar[r]_i &
N \ar[r] &
(C_Y X) \times_Y Z \\
C_Z Y \ar[r] \ar[u] &
C_Z X \ar[u]
}
$$
de cônes et de fibrés vectoriels sur $Z$. Remarquons que $N'$ est un
diviseur de Cartier effectif relatif dans $N$ sur $Z$ et que
$$
\xymatrix{
N' \ar[d] \ar[r]_i & N \ar[d] \\
Z \ar[r]^-o & (C_Y X) \times_Y Z
}
$$
est cartésien, où $o$ est la section nulle du fibré en droites
$C_Y X$ sur $Y$. D'après le
lemme \ref{chow-lemma-relation-normal-cones}, nous avons $o^*[C_ZX]_n = [C_Z Y]_{n - 1}$
dans
$$
\CH_{n - 1}(Y \times_{o, C_Y X} C_ZX) =
\CH_{n - 1}(Z \times_{o, (C_Y X) \times_Y Z} C_ZX)
$$
Par la propriété cartésienne de
ce carré, il s'ensuit que
$$
i^*[C_ZX]_n = [C_Z Y]_{n - 1}
$$
dans $\CH_{n - 1}(N')$. Remarquons maintenant que
$\gamma = c(Z \to X, \mathcal{N}) \cap [X]$ et
$\gamma' = c(Z \to Y, \mathcal{N}') \cap [Y]_{n - 1}$
sont caractérisés par $p^*\gamma = [C_Z X]_n$ dans $\CH_n(N)$
et par $(p')^*\gamma' = [C_Z Y]_{n - 1}$ dans $\CH_{n - 1}(N')$.
La démonstration est donc achevée, puisque $i^* \circ p^* = (p')^*$ d'après le
lemme \ref{chow-lemma-relative-effective-cartier}.
\end{proof}

\begin{remark}[Variante pour les immersions]
\label{chow-remark-gysin-for-immersion}
Soit $(S, \delta)$ comme dans la situation \ref{chow-situation-setup}.
Soit $X$ un schéma localement de type fini sur $S$.
Soit $i : Z \to X$ une immersion de schémas.
Dans cette situation,
\begin{enumerate}
\item le faisceau conormal $\mathcal{C}_{Z/X}$
de $Z$ dans $X$ est défini
(Morphismes, définition \ref{morphisms-definition-conormal-sheaf}),
\item nous appelons un couple formé d'un $\mathcal{O}_Z$-module localement libre de rang fini
$\mathcal{N}$ et d'une surjection $\sigma : \mathcal{N}^\vee \to \mathcal{C}_{Z/X}$
un fibré normal virtuel pour l'immersion $Z \to X$,
\item choisissons un sous-schéma ouvert $U \subset X$ tel que $Z \to X$
se factorise par une immersion fermée $Z \to U$ et posons
$c(Z \to X, \mathcal{N}) = c(Z \to U, \mathcal{N}) \circ (U \to X)^*$.
\end{enumerate}
La classe bivariante $c(Z \to X, \mathcal{N})$ ne dépend pas du choix
du sous-schéma ouvert $U$. Tous les lemmes ont des analogues immédiats
pour cette construction légèrement plus générale. Nous omettons les détails.
\end{remark}







\section{Calcul de quelques classes}
\label{chow-section-calculate}

\noindent
Pour aller plus loin, nous devons calculer les valeurs de certaines des
classes construites ci-dessus.

\begin{lemma}
\label{chow-lemma-compute-koszul}
Soit $(S, \delta)$ comme dans la situation \ref{chow-situation-setup}. Soit $X$
un schéma localement de type fini sur $S$. Soit $\mathcal{E}$
un $\mathcal{O}_X$-module localement libre de rang $r$.
Alors
$$
\prod\nolimits_{n = 0, \ldots, r} c(\wedge^n \mathcal{E})^{(-1)^n} =
1 - (r - 1)! c_r(\mathcal{E}) + \ldots
$$
\end{lemma}

\begin{proof}
Par le principe de scindage, nous pouvons ramener ceci à un calcul dans
l'anneau de polynômes aux racines de Chern $x_1, \ldots, x_r$ de $\mathcal{E}$. Voir la
section \ref{chow-section-splitting-principle}. Remarquons que
$$
c(\wedge^n \mathcal{E}) =
\prod\nolimits_{1 \leq i_1 < \ldots < i_n \leq r}
(1 + x_{i_1} + \ldots + x_{i_n})
$$
Ainsi, le logarithme du membre de gauche de l'égalité du lemme est
$$
-
\sum\nolimits_{p \geq 1}
\sum\nolimits_{n = 0}^r
\sum\nolimits_{1 \leq i_1 < \ldots < i_n \leq r}
\frac{(-1)^{p + n}}{p}(x_{i_1} + \ldots + x_{i_n})^p
$$
Remarquons bien le signe moins devant. Cependant, nous avons
$$
\sum\nolimits_{p \geq 0}
\sum\nolimits_{n = 0}^r
\sum\nolimits_{1 \leq i_1 < \ldots < i_n \leq r}
\frac{(-1)^{p + n}}{p!}(x_{i_1} + \ldots + x_{i_n})^p
=
\prod (1 - e^{-x_i})
$$
Nous voyons donc que le premier terme non nul de notre classe de Chern
est de degré $r$ et égal à la valeur annoncée.
\end{proof}

\begin{lemma}
\label{chow-lemma-compute-section}
Soit $(S, \delta)$ comme dans la situation \ref{chow-situation-setup}. Soit $X$
un schéma localement de type fini sur $S$. Soit $\mathcal{C}$
un $\mathcal{O}_X$-module localement libre de rang $r$. Considérons les
morphismes
$$
X = \underline{\text{Proj}}_X(\mathcal{O}_X[T])
\xrightarrow{i}
E = \underline{\text{Proj}}_X(\text{Sym}^*(\mathcal{C})[T])
\xrightarrow{\pi}
X
$$
Alors $c_t(i_*\mathcal{O}_X) = 0$ pour $t = 1, \ldots, r - 1$ et, dans
$A^0(C \to E)$, nous avons
$$
p^* \circ \pi_* \circ c_r(i_*\mathcal{O}_X) = (-1)^{r - 1}(r - 1)! j^*
$$
où
$j : C \to E$ et $p : C \to X$ sont l'inclusion et le morphisme
structural du fibré vectoriel
$C = \underline{\Spec}(\text{Sym}^*(\mathcal{C}))$.
\end{lemma}

\begin{proof}
L'application canonique $\pi^*\mathcal{C} \to \mathcal{O}_E(1)$ s'annule
exactement le long de $i(X)$. Ainsi, le complexe de Koszul de l'application
$$
\pi^*\mathcal{C} \otimes \mathcal{O}_E(-1) \to \mathcal{O}_E
$$
est une résolution de $i_*\mathcal{O}_X$. En particulier, nous voyons que
$i_*\mathcal{O}_X$ est un objet parfait de $D(\mathcal{O}_E)$
dont les classes de Chern sont définies. L'annulation de $c_t(i_*\mathcal{O}_X)$
pour $t = 1, \ldots, t - 1$ résulte du lemme \ref{chow-lemma-compute-koszul}.
Ce lemme donne également
$$
c_r(i_*\mathcal{O}_X) = - (r - 1)!
c_r(\pi^*\mathcal{C} \otimes \mathcal{O}_E(-1))
$$
D'autre part, d'après le lemme \ref{chow-lemma-chern-classes-dual}, nous avons
$$
c_r(\pi^*\mathcal{C} \otimes \mathcal{O}_E(-1)) =
(-1)^r c_r(\pi^*\mathcal{C}^\vee \otimes \mathcal{O}_E(1))
$$
et $\pi^*\mathcal{C}^\vee \otimes \mathcal{O}_E(1)$ possède une section $s$
qui s'annule exactement le long de $i(X)$.

\medskip\noindent
Après avoir remplacé $X$ par un schéma localement de type fini sur $X$,
il suffit de prouver que les deux membres de l'égalité ont le
même effet sur un élément $\alpha \in \CH_*(E)$. Comme $C \to X$
est un fibré vectoriel, toute classe de cycle sur $C$ est de la forme $p^*\beta$
pour un certain $\beta \in \CH_*(X)$ (lemme \ref{chow-lemma-vectorbundle}).
Ainsi, d'après le lemme \ref{chow-lemma-restrict-to-open},
nous pouvons écrire $\alpha = \pi^*\beta + \gamma$, où $\gamma$
est à support dans $E \setminus C$. En utilisant les égalités ci-dessus,
il suffit de montrer que
$$
p^*(\pi_*(c_r(\pi^*\mathcal{C}^\vee \otimes \mathcal{O}_E(1)) \cap [W])) =
j^*[W]
$$
lorsque $W \subset E$ est un sous-schéma fermé intègre qui
soit (a) disjoint de $C$, soit (b) de la forme $W = \pi^{-1}Y$
pour un certain sous-schéma fermé intègre $Y \subset X$.
En utilisant la section $s$ et le lemme \ref{chow-lemma-top-chern-class}, nous trouvons
dans le cas (a) $c_r(\pi^*\mathcal{C}^\vee \otimes \mathcal{O}_E(1)) \cap [W] = 0$
et, dans le cas (b),
$c_r(\pi^*\mathcal{C}^\vee \otimes \mathcal{O}_E(1)) \cap [W] = [i(Y)]$.
Le résultat s'en déduit aisément ; nous omettons les détails.
\end{proof}

\begin{lemma}
\label{chow-lemma-agreement-with-loc-chern}
Soit $(S, \delta)$ comme dans la situation \ref{chow-situation-setup}. Soit $i : Z \to X$
une immersion fermée régulière de codimension $r$
entre des schémas localement de type fini sur $S$.
Soit $\mathcal{N} = \mathcal{C}_{Z/X}^\vee$ le faisceau normal. Si $X$
est quasi-compact (ou possède des composantes irréductibles quasi-compactes), alors
$c_t(Z \to X, i_*\mathcal{O}_Z) = 0$ pour $t = 1, \ldots, r - 1$ et
$$
c_r(Z \to X, i_*\mathcal{O}_Z) = (-1)^{r - 1} (r - 1)! c(Z \to X, \mathcal{N})
\quad\text{dans}\quad
A^r(Z \to X)
$$
où $c_t(Z \to X, i_*\mathcal{O}_Z)$
est la classe de Chern localisée
de la définition \ref{chow-definition-localized-chern}.
\end{lemma}

\begin{proof}
Pour tout $x \in Z$, nous pouvons choisir un voisinage ouvert affine
$\Spec(A) \subset X$ tel que $Z \cap \Spec(A) = V(f_1, \ldots, f_r)$,
où $f_1, \ldots, f_r \in A$ est une suite régulière.
Voir Diviseurs, définition \ref{divisors-definition-regular-immersion} et
lemme \ref{divisors-lemma-Noetherian-scheme-regular-ideal}.
Nous voyons alors que le complexe de Koszul de $f_1, \ldots, f_r$ est
une résolution de $A/(f_1, \ldots, f_r)$, par exemple d'après
Compléments d'algèbre, lemme \ref{more-algebra-lemma-regular-koszul-regular}.
Ainsi, $A/(f_1, \ldots, f_r)$ est parfait comme $A$-module.
Il s'ensuit que $F = i_*\mathcal{O}_Z$ est un objet parfait de
$D(\mathcal{O}_X)$ dont la restriction à $X \setminus Z$ est nulle.
L'hypothèse que $X$ est quasi-compact (ou possède des composantes
irréductibles quasi-compactes) signifie que les classes de Chern localisées
$c_t(Z \to X, i_*\mathcal{O}_Z)$ sont définies, voir
la situation \ref{chow-situation-loc-chern} et la
définition \ref{chow-definition-localized-chern}. En définitive,
nous concluons que l'énoncé a un sens.

\medskip\noindent
Notons $b : W \to \mathbf{P}^1_X$ l'éclatement en $\infty(Z)$
comme dans la section \ref{chow-section-blowup-Z-first}. D'après (\ref{chow-item-find-Z-in-blowup}),
nous avons une immersion fermée
$$
i' : \mathbf{P}^1_Z \longrightarrow W
$$
Nous affirmons que $Q = i'_*\mathcal{O}_{\mathbf{P}^1_Z}$
est un objet parfait de
$D(\mathcal{O}_W)$ et que $F$ et $Q$ satisfont aux hypothèses du
lemme \ref{chow-lemma-independent-loc-chern-bQ}.

\medskip\noindent
Supposons l'affirmation démontrée. La conclusion du lemme \ref{chow-lemma-independent-loc-chern-bQ}
est que nous avons
$$
c_p(Z \to X, F) = c'_p(Q) = (E \to Z)_* \circ c'_p(Q|_E) \circ C
$$
pour tout $p \geq 1$. Remarquons que $Q|_E$ est égal à l'image directe du
faisceau structural de $Z$ par le morphisme $Z \to E$ qui est le
changement de base de $i'$ par $\infty$.
Ainsi, l'annulation de $c_t(Z \to X, F)$ pour $1 \leq t \leq r - 1$
résulte du lemme \ref{chow-lemma-compute-section} appliqué à $E \to Z$.
Comme $\mathcal{C}_{Z/X} = \mathcal{N}^\vee$
est localement libre, la classe bivariante $c(Z \to X, \mathcal{N})$
est caractérisée par la relation
$$
j^* \circ C = p^* \circ c(Z \to X, \mathcal{N})
$$
où $j : C_ZX \to W_\infty$ et $p : C_ZX \to Z$ sont les applications données.
(Rappelons que $C \in A^0(W_\infty \to X)$ est la classe du
lemme \ref{chow-lemma-gysin-at-infty}.)
Ainsi, l'égalité affichée dans l'énoncé du lemme
résulte de l'égalité correspondante du lemme \ref{chow-lemma-compute-section}.

\medskip\noindent
Démonstration de l'affirmation. Soient $A$ et $f_1, \ldots, f_r$ comme ci-dessus.
Considérons l'ouvert affine $\Spec(A[s]) \subset \mathbf{P}^1_X$
comme dans la section \ref{chow-section-blowup-Z-first}. Rappelons que $s = 0$
définit $(\mathbf{P}^1_X)_\infty$ sur cet ouvert. Ainsi, au-dessus de
$\Spec(A[s])$, nous effectuons l'éclatement de l'idéal engendré par
la suite régulière $s, f_1, \ldots, f_r$. D'après Compléments d'algèbre, lemme
\ref{more-algebra-lemma-blowup-regular-sequence}, les $r + 1$
cartes affines sont des intersections complètes globales sur $A[s]$.
La carte correspondant à l'algèbre affine d'éclatement
$$
A[s][f_1/s, \ldots, f_r/s] = A[s, y_1, \ldots, y_r]/(sy_i - f_i)
$$
contient $i'(Z \cap \Spec(A))$ comme sous-schéma fermé défini par
$y_1, \ldots, y_r$. Comme $y_1, \ldots, y_r, sy_1 - f_1, \ldots, sy_r - f_r$
est une suite régulière dans l'anneau de polynômes $A[s, y_1, \ldots, y_r]$,
nous trouvons que $i'$ est une immersion régulière. Nous omettons quelques détails.
Comme ci-dessus, nous concluons que $Q = i'_*\mathcal{O}_{\mathbf{P}^1_Z}$
est un objet parfait de $D(\mathcal{O}_W)$. Toutes les
autres hypothèses sur $F$ et $Q$ du lemme \ref{chow-lemma-independent-loc-chern-bQ}
(et du lemme \ref{chow-lemma-localized-chern-pre}) se vérifient immédiatement.
\end{proof}

\begin{lemma}
\label{chow-lemma-actual-computation}
Dans la situation du lemme \ref{chow-lemma-agreement-with-loc-chern},
posons $\dim_\delta(X) = n$. Alors nous avons
\begin{enumerate}
\item $c_t(Z \to X, i_*\mathcal{O}_Z) \cap [X]_n = 0$ pour
$t = 1, \ldots, r - 1$,
\item $c_r(Z \to X, i_*\mathcal{O}_Z) \cap [X]_n =
(-1)^{r - 1}(r - 1)![Z]_{n - r}$,
\item $ch_t(Z \to X, i_*\mathcal{O}_Z) \cap [X]_n = 0$ pour
$t = 0, \ldots, r - 1$, et
\item $ch_r(Z \to X, i_*\mathcal{O}_Z) \cap [X]_n = [Z]_{n - r}$.
\end{enumerate}
\end{lemma}

\begin{proof}
Les parties (1) et (2) résultent immédiatement du
lemme \ref{chow-lemma-agreement-with-loc-chern}
combiné au lemme \ref{chow-lemma-gysin-fundamental}.
Nous en déduisons ensuite les parties (3) et (4) à l'aide de la relation
entre $ch_p = (1/p!)P_p$ et $c_p$ donnée dans le
lemme \ref{chow-lemma-loc-chern-character}. (En effet,
$(-1)^{r - 1}(r - 1)!ch_r = c_r$ pourvu que
$c_1 = c_2 = \ldots = c_{r - 1} = 0$.)
\end{proof}







\section{Un opérateur d'Adams}
\label{chow-section-adams}

\noindent
Nous effectuons le travail minimal nécessaire pour définir le deuxième opérateur d'Adams.
Soit $X$ un schéma. Rappelons que $\textit{Vect}(X)$ désigne la
catégorie des $\mathcal{O}_X$-modules localement libres de rang fini.
Rappelons en outre que nous avons construit le groupe $K$ d'indice zéro
$K_0(\textit{Vect}(X))$ associé à cette catégorie dans
Catégories dérivées de schémas, section \ref{perfect-section-K-groups}.
Enfin, $K_0(\textit{Vect}(X))$ est un anneau, voir
Catégories dérivées de schémas, remarque \ref{perfect-remark-K-ring}.

\begin{lemma}
\label{chow-lemma-second-adams-operator}
Soit $X$ un schéma. Il existe un homomorphisme d'anneaux
$$
\psi^2 :
K_0(\textit{Vect}(X))
\longrightarrow
K_0(\textit{Vect}(X))
$$
qui envoie $[\mathcal{L}]$ sur $[\mathcal{L}^{\otimes 2}]$
lorsque $\mathcal{L}$ est inversible et qui est compatible aux images inverses.
\end{lemma}

\begin{proof}
Soit $X$ un schéma.
Soit $\mathcal{E}$ un $\mathcal{O}_X$-module localement libre de rang fini.
Nous allons considérer l'élément
$$
\psi^2(\mathcal{E}) = [\text{Sym}^2(\mathcal{E})] - [\wedge^2(\mathcal{E})]
$$
de $K_0(\textit{Vect}(X))$.

\medskip\noindent
Soit $X$ un schéma et considérons une suite exacte courte
$$
0 \to \mathcal{E} \to \mathcal{F} \to \mathcal{G} \to 0
$$
de $\mathcal{O}_X$-modules localement libres de rang fini. Considérons-la
comme une filtration de $\mathcal{F}$ à $2$ crans. La filtration
induite sur $\text{Sym}^2(\mathcal{F})$ possède $3$ crans, de
gradués associés $\text{Sym}^2(\mathcal{E})$, $\mathcal{E} \otimes \mathcal{F}$
et $\text{Sym}^2(\mathcal{G})$. Ainsi
$$
[\text{Sym}^2(\mathcal{F})] =
[\text{Sym}^2(\mathcal{E})] +
[\mathcal{E} \otimes \mathcal{F}] +
[\text{Sym}^2(\mathcal{G})]
$$
Exactement de la même manière, on montre que
$$
[\wedge^2(\mathcal{F})] =
[\wedge^2(\mathcal{E})] +
[\mathcal{E} \otimes \mathcal{F}] +
[\wedge^2(\mathcal{G})]
$$
Nous voyons donc que
$\psi^2(\mathcal{F}) = \psi^2(\mathcal{E}) + \psi^2(\mathcal{G})$.
Nous concluons que nous obtenons une application additive bien définie
$\psi^2 : K_0(\textit{Vect}(X)) \to K_0(\textit{Vect}(X))$.

\medskip\noindent
Il est clair que cette application commute aux images inverses.

\medskip\noindent
Il reste à montrer que $\psi^2$ est un homomorphisme d'anneaux.
Soit $X$ un schéma et soient $\mathcal{E}$ et $\mathcal{F}$ des
$\mathcal{O}_X$-modules localement libres de rang fini.
Remarquons qu'il existe une suite exacte courte
$$
0 \to \wedge^2(\mathcal{E}) \otimes \wedge^2(\mathcal{F}) \to
\text{Sym}^2(\mathcal{E} \otimes \mathcal{F}) \to
\text{Sym}^2(\mathcal{E}) \otimes \text{Sym}^2(\mathcal{F}) \to 0
$$
où la première application envoie $(e \wedge e') \otimes (f \wedge f')$ sur
$(e \otimes f)(e' \otimes f') - (e' \otimes f)(e \otimes f')$ et où
la seconde application envoie $(e \otimes f) (e' \otimes f')$ sur $ee' \otimes ff'$.
De même, il existe une suite exacte courte
$$
0 \to \text{Sym}^2(\mathcal{E}) \otimes \wedge^2(\mathcal{F}) \to
\wedge^2(\mathcal{E} \otimes \mathcal{F}) \to
\wedge^2(\mathcal{E}) \otimes \text{Sym}^2(\mathcal{F}) \to 0
$$
où la première application envoie $e e' \otimes f \wedge f'$ sur
$(e \otimes f) \wedge (e' \otimes f') + (e' \otimes f) \wedge (e \otimes f')$
et où la seconde application envoie
$(e \otimes f) \wedge (e' \otimes f')$ sur
$(e \wedge e') \otimes (f f')$.
Comme ci-dessus, cela prouve que l'application $\psi^2$ est multiplicative.
Puisqu'il est clair que $\psi^2(1) = 1$, cela achève la démonstration.
\end{proof}

\begin{remark}
\label{chow-remark-adams-derived}
Soit $X$ un schéma tel que $2$ soit inversible sur $X$.
Alors l'opérateur d'Adams $\psi^2$ peut être défini sur le groupe $K$
$K_0(X) = K_0(D_{perf}(\mathcal{O}_X))$
(Catégories dérivées de schémas, définition \ref{perfect-definition-K-group})
d'une manière immédiate.
En effet, étant donné un complexe parfait $L$ sur $X$, nous obtenons une action
du groupe $\{\pm 1\}$ sur $L \otimes^\mathbf{L} L$ par permutation
des facteurs. Nous pouvons alors poser
$$
\psi^2(L) = [(L \otimes^\mathbf{L} L)^+] -
[(L \otimes^\mathbf{L} L)^-]
$$
où $(-)^+$ désigne la prise des invariants et $(-)^-$ celle des
anti-invariants (convenablement définis).
En utilisant l'exactitude de la prise des invariants et des anti-invariants, on peut
raisonner comme dans la démonstration du lemme \ref{chow-lemma-second-adams-operator}
pour montrer que cette définition est bonne.
Lorsque $2$ n'est pas inversible sur $X$, la situation est beaucoup plus
compliquée et une autre approche doit être employée.
\end{remark}

\begin{lemma}
\label{chow-lemma-minus-adams-operator}
Soit $X$ un schéma. Il existe un homomorphisme d'anneaux
$\psi^{-1} : K_0(\textit{Vect}(X)) \to K_0(\textit{Vect}(X))$
qui envoie $[\mathcal{E}]$ sur $[\mathcal{E}^\vee]$
lorsque $\mathcal{E}$ est localement libre de rang fini
et qui est compatible aux images inverses.
\end{lemma}

\begin{proof}
Il suffit de vérifier que la prise des duaux est compatible aux
suites exactes courtes et aux images inverses. C'est clair.
\end{proof}

\begin{remark}
\label{chow-remark-chern-classes-K}
Soit $(S, \delta)$ comme dans la situation \ref{chow-situation-setup}.
Soit $X$ localement de type fini sur $S$. La classe de Chern
définit une application canonique
$$
c : K_0(\textit{Vect}(X)) \longrightarrow \prod\nolimits_{i \geq 0} A^i(X)
$$
qui envoie un générateur $[\mathcal{E}]$ du membre de gauche sur
$c(\mathcal{E}) = 1 + c_1(\mathcal{E}) + c_2(\mathcal{E}) + \ldots$
et se prolonge multiplicativement. Ainsi, $-[\mathcal{E}]$ est envoyé sur
l'inverse formel $c(\mathcal{E})^{-1}$, ce qui explique le
produit infini du membre de droite. Cette application est bien définie d'après le
lemme \ref{chow-lemma-additivity-chern-classes}.
\end{remark}

\begin{remark}
\label{chow-remark-chern-character-K}
Soit $(S, \delta)$ comme dans la situation \ref{chow-situation-setup}.
Soit $X$ localement de type fini sur $S$. Le caractère de Chern
définit un homomorphisme d'anneaux canonique
$$
ch : K_0(\textit{Vect}(X)) \longrightarrow
\prod\nolimits_{i \geq 0} A^i(X) \otimes \mathbf{Q}
$$
qui envoie un générateur $[\mathcal{E}]$ du membre de gauche sur
$ch(\mathcal{E})$ et se prolonge additivement. Cette application est bien définie
d'après le lemme \ref{chow-lemma-chern-character-additive} et est un homomorphisme d'anneaux d'après le
lemme \ref{chow-lemma-chern-character-multiplicative}.
\end{remark}

\begin{lemma}
\label{chow-lemma-adams-and-chern}
Soit $(S, \delta)$ comme dans la situation \ref{chow-situation-setup}.
Soit $X$ localement de type fini sur $S$. Si $\psi^2$ est
comme dans le lemme \ref{chow-lemma-second-adams-operator}, et si $c$ et $ch$ sont comme dans les
remarques \ref{chow-remark-chern-classes-K} et \ref{chow-remark-chern-character-K},
alors nous avons $c_i(\psi^2(\alpha)) = 2^i c_i(\alpha)$ et
$ch_i(\psi^2(\alpha)) = 2^i ch_i(\alpha)$
pour tout $\alpha \in K_0(\textit{Vect}(X))$.
\end{lemma}

\begin{proof}
Remarquons que l'application $\prod_{i \geq 0} A^i(X) \to \prod_{i \geq 0} A^i(X)$
qui multiplie par $2^i$ sur $A^i(X)$ est un homomorphisme d'anneaux. Ainsi, puisque $\psi^2$
est également un homomorphisme d'anneaux, il suffit de prouver les formules pour des générateurs additifs
de $K_0(\textit{Vect}(X))$. Nous pouvons donc supposer que $\alpha = [\mathcal{E}]$
pour un certain $\mathcal{O}_X$-module localement libre de rang fini $\mathcal{E}$.
Par la construction des classes de Chern de $\mathcal{E}$, nous nous
ramenons immédiatement au cas où $\mathcal{E}$ est de rang constant $r$, voir la
remarque \ref{chow-remark-extend-to-finite-locally-free}.
Dans ce cas, nous pouvons choisir un morphisme projectif lisse $p : P \to X$
tel que la restriction $A^*(X) \to A^*(P)$ soit injective
et que $p^*\mathcal{E}$ possède une filtration finie dont les
gradués associés sont des $\mathcal{O}_P$-modules inversibles $\mathcal{L}_j$, voir le
lemme \ref{chow-lemma-splitting-principle}. Alors
$[p^*\mathcal{E}] = \sum [\mathcal{L}_j]$ et donc
$\psi^2([p^\mathcal{E}]) = \sum [\mathcal{L}_j^{\otimes 2}]$
par définition de $\psi^2$. En posant $x_j  = c_1(\mathcal{L}_j)$,
nous avons
$$
c(\alpha) = \prod (1 + x_j)
\quad\text{et}\quad
c(\psi^2(\alpha)) = \prod (1 + 2 x_j)
$$
dans $\prod A^i(P)$, et nous avons
$$
ch(\alpha) = \sum \exp(x_j)
\quad\text{et}\quad
ch(\psi^2(\alpha)) = \sum \exp(2 x_j)
$$
dans $\prod A^i(P)$. Le résultat voulu résulte de ces formules.
\end{proof}

\begin{remark}
\label{chow-remark-perf-Z-cohomology-K}
Soit $X$ un schéma localement noethérien.
Soit $Z \subset X$ un sous-schéma fermé. Considérons la sous-catégorie triangulée
strictement pleine et saturée
$$
D_{Z, perf}(\mathcal{O}_X) \subset D(\mathcal{O}_X)
$$
formée des complexes parfaits de $\mathcal{O}_X$-modules
dont les faisceaux de cohomologie ont leur support ensembliste dans $Z$.
Notons $\textit{Coh}_Z(X) \subset \textit{Coh}(X)$
la sous-catégorie de Serre des $\mathcal{O}_X$-modules cohérents dont le support
ensembliste est contenu dans $Z$. Remarquons que, pour
$E \in D_{Z, perf}(\mathcal{O}_X)$, localement pour la topologie de Zariski sur $X$,
seul un nombre fini des faisceaux de cohomologie $H^i(E)$ sont non nuls
(et ils ont tous leur support ensembliste dans $Z$).
Nous pouvons donc définir
$$
K_0(D_{Z, perf}(\mathcal{O}_X))
\longrightarrow
K_0(\textit{Coh}_Z(X)) = K'_0(Z)
$$
(égalité d'après le lemme \ref{chow-lemma-K-coherent-supported-on-closed}) par la règle
$$
E \longmapsto
[\bigoplus\nolimits_{i \in \mathbf{Z}} H^{2i}(E)] -
[\bigoplus\nolimits_{i \in \mathbf{Z}} H^{2i + 1}(E)]
$$
Cela fonctionne parce que, pour tout triangle distingué de
$D_{Z, perf}(\mathcal{O}_X)$, nous avons une suite exacte longue de
faisceaux de cohomologie.
\end{remark}

\begin{remark}
\label{chow-remark-perf-Z-regular}
Soient $X$, $Z$, $D_{Z, perf}(\mathcal{O}_X)$ comme dans la
remarque \ref{chow-remark-perf-Z-cohomology-K}. Supposons que $X$ soit régulier.
Il existe alors une application canonique
$$
K_0(\textit{Coh}(Z)) \longrightarrow K_0(D_{Z, perf}(\mathcal{O}_X))
$$
définie comme suit. Pour tout $\mathcal{O}_Z$-module cohérent
$\mathcal{F}$, notons $\mathcal{F}[0]$ l'objet de $D(\mathcal{O}_X)$
qui a $\mathcal{F}$ en degré $0$ et est nul dans les autres degrés.
Alors $\mathcal{F}[0]$ est un complexe parfait sur $X$ d'après
Catégories dérivées de schémas, lemme \ref{perfect-lemma-perfect-on-regular}.
Ainsi, $\mathcal{F}[0]$ est un objet de $D_{Z, perf}(\mathcal{O}_X)$.
D'autre part, étant donnée une suite exacte courte
$0 \to \mathcal{F} \to \mathcal{F}' \to \mathcal{F}'' \to 0$ de
$\mathcal{O}_Z$-modules cohérents, nous obtenons un triangle distingué
$\mathcal{F}[0] \to \mathcal{F}'[0] \to \mathcal{F}''[0] \to \mathcal{F}[1]$,
voir Catégories dérivées, section \ref{derived-section-canonical-delta-functor}.
Cela montre que nous obtenons une application
$K_0(\textit{Coh}(Z)) \to K_0(D_{Z, perf}(\mathcal{O}_X))$
en envoyant $[\mathcal{F}]$ sur $[\mathcal{F}[0]]$,
avec nos excuses pour cette affreuse notation.
\end{remark}

\begin{lemma}
\label{chow-lemma-perf-Z-regular}
Soit $X$ un schéma noethérien régulier.
Soit $Z \subset X$ un sous-schéma fermé. Les applications construites
dans les remarques \ref{chow-remark-perf-Z-cohomology-K} et
\ref{chow-remark-perf-Z-regular} sont inverses l'une de l'autre, et nous obtenons
$K'_0(Z) = K_0(D_{Z, perf}(\mathcal{O}_X))$.
\end{lemma}

\begin{proof}
Il est clair que la composée
$$
K_0(\textit{Coh}(Z)) \longrightarrow
K_0(D_{Z, perf}(\mathcal{O}_X)) \longrightarrow
K_0(\textit{Coh}(Z))
$$
est l'identité. Il suffit donc de montrer que la première flèche est
surjective. Soit $E$ un objet de $D_{Z, perf}(\mathcal{O}_X)$.
Rappelons que $D_{perf}(\mathcal{O}_X) = D^b_{\textit{Coh}}(\mathcal{O}_X)$
d'après Catégories dérivées de schémas, lemme
\ref{perfect-lemma-perfect-on-regular}.
Ainsi, les cohomologies $H^i(E)$ sont cohérentes, peuvent être considérées
comme des objets de $D_{Z, perf}(\mathcal{O}_X)$ et seul un nombre fini
d'entre elles sont non nulles. En utilisant les triangles distingués des troncations canoniques, le
lecteur voit que
$$
[E] = \sum (-1)^i[H^i(E)[0]]
$$
dans $K_0(D_{Z, perf}(\mathcal{O}_X))$. Il suffit alors de
montrer que $[\mathcal{F}[0]]$ est dans l'image de l'application
pour tout $\mathcal{O}_X$-module cohérent dont le support ensembliste
est contenu dans $Z$. Puisque nous pouvons trouver une filtration finie de
$\mathcal{F}$ dont les sous-quotients sont des $\mathcal{O}_Z$-modules,
la démonstration est achevée.
\end{proof}

\begin{remark}
\label{chow-remark-localized-chern-classes-K}
Soit $(S, \delta)$ comme dans la situation \ref{chow-situation-setup}.
Soit $X$ localement de type fini sur $S$.
Soit $Z \subset X$ un sous-schéma fermé et soit
$D_{Z, perf}(\mathcal{O}_X)$ comme dans la
remarque \ref{chow-remark-perf-Z-cohomology-K}.
Si $X$ est quasi-compact (ou, plus généralement, si les composantes
irréductibles de $X$ sont quasi-compactes), alors
les classes de Chern localisées définissent une application canonique
$$
c(Z \to X, -) : K_0(D_{Z, perf}(\mathcal{O}_X)) \longrightarrow
A^0(X) \times \prod\nolimits_{i \geq 1} A^i(Z \to X)
$$
qui envoie un générateur $[E]$ du membre de gauche sur
$$
c(Z \to X, E) = 1 + c_1(Z \to X, E) + c_2(Z \to X, E) + \ldots
$$
et se prolonge multiplicativement (le produit du membre de
droite étant celui de la remarque \ref{chow-remark-ring-loc-classes}).
La condition de quasi-compacité sur $X$ garantit que les
classes de Chern localisées sont définies (situation \ref{chow-situation-loc-chern} et
définition \ref{chow-definition-localized-chern})
et que ces classes de Chern localisées transforment les triangles distingués en
produits correspondants dans les anneaux de Chow bivariants
(lemme \ref{chow-lemma-additivity-loc-chern-c}).
\end{remark}

\begin{remark}
\label{chow-remark-localized-chern-character-K}
Soit $(S, \delta)$ comme dans la situation \ref{chow-situation-setup}.
Soit $X$ localement de type fini sur $S$.
Soit $Z \subset X$ un sous-schéma fermé et soit
$D_{Z, perf}(\mathcal{O}_X)$ comme dans la
remarque \ref{chow-remark-perf-Z-cohomology-K}.
Si les composantes irréductibles de $X$ sont quasi-compactes, alors
le caractère de Chern localisé définit une application canonique additive
et multiplicative
$$
ch(Z \to X, -) : K_0(D_{Z, perf}(\mathcal{O}_X)) \longrightarrow
\prod\nolimits_{i \geq 0} A^i(Z \to X) \otimes \mathbf{Q}
$$
qui envoie un générateur $[E]$ du membre de gauche sur
$ch(Z \to X, E)$ et se prolonge additivement.
En effet, la condition sur les composantes irréductibles de $X$ garantit que le
caractère de Chern localisé est défini (situation \ref{chow-situation-loc-chern} et
définition \ref{chow-definition-localized-chern})
et que ces caractères de Chern localisés
transforment les triangles distingués en les
sommes correspondantes dans les anneaux de Chow bivariants
(lemme \ref{chow-lemma-additivity-loc-chern-P}).
La multiplication sur
$K_0(D_{Z, perf}(X))$ est définie à l'aide du produit tensoriel dérivé
(Catégories dérivées de schémas, remarque \ref{perfect-remark-perf-Z}) ;
ainsi, $ch(Z \to X, \alpha \beta) = ch(Z \to X, \alpha) ch(Z \to X, \beta)$ d'après le
lemme \ref{chow-lemma-loc-chern-tensor-product}.
Si $X$ est quasi-compact, alors l'image de l'application $ch(Z \to X, -)$ est
contenue dans $A^*(Z \to X) \otimes \mathbf{Q}$ ; nous omettons
les détails.
\end{remark}

\begin{remark}
\label{chow-remark-chern-classes-agree}
Soit $(S, \delta)$ comme dans la situation \ref{chow-situation-setup}.
Soit $X$ localement de type fini sur $S$ et supposons que $X$
soit quasi-compact (ou, plus généralement, que les composantes irréductibles
de $X$ soient quasi-compactes). Avec $Z = X$ et les notations des
remarques \ref{chow-remark-localized-chern-classes-K} et
\ref{chow-remark-localized-chern-character-K},
nous avons $D_{Z, perf}(\mathcal{O}_X) = D_{perf}(\mathcal{O}_X)$,
et nous voyons que
$$
K_0(D_{Z, perf}(\mathcal{O}_X)) = K_0(D_{perf}(\mathcal{O}_X)) = K_0(X)
$$
voir
Catégories dérivées de schémas, définition \ref{perfect-definition-K-group}.
Nous obtenons donc
$$
c : K_0(X) \to \prod A^i(X)
\quad\text{et}\quad
ch : K_0(X) \to \prod A^i(X) \otimes \mathbf{Q}
$$
comme cas particulier des remarques \ref{chow-remark-localized-chern-classes-K} et
\ref{chow-remark-localized-chern-character-K}. Bien sûr, nous aurions aussi pu
utiliser directement la définition \ref{chow-definition-defined-on-perfect} et les
lemmes \ref{chow-lemma-additivity-on-perfect} et
\ref{chow-lemma-chern-classes-perfect-tensor-product} pour construire ces applications
(car on voit immédiatement que cela donne les mêmes classes).
Rappelons qu'il existe une application canonique $K_0(\textit{Vect}(X)) \to K_0(X)$
qui envoie un module localement libre de rang fini sur lui-même, considéré
comme un complexe parfait (placé en degré $0$), voir
Catégories dérivées de schémas, section \ref{perfect-section-K-groups}.
Alors le diagramme
$$
\xymatrix{
K_0((\textit{Vect}(X)) \ar[rd]_c \ar[rr] & &
K_0(D_{perf}(\mathcal{O}_X)) = K_0(X) \ar[ld]^c \\
& \prod A^i(X)
}
$$
est commutatif, où la flèche sud-est est celle construite dans la
remarque \ref{chow-remark-chern-classes-K}. De même, le diagramme
$$
\xymatrix{
K_0((\textit{Vect}(X)) \ar[rd]_{ch} \ar[rr] & &
K_0(D_{perf}(\mathcal{O}_X)) = K_0(X) \ar[ld]^{ch} \\
& \prod A^i(X) \otimes \mathbf{Q}
}
$$
est commutatif, où la flèche sud-est est celle construite dans la
remarque \ref{chow-remark-chern-character-K}.
\end{remark}











\section{Retour sur les groupes de Chow et les groupes K}
\label{chow-section-chow-and-K-II}

\noindent
Cette section fait suite à la section \ref{chow-section-chow-and-K}.
Soit $(S, \delta)$ comme dans la situation \ref{chow-situation-setup}.
Soit $X$ localement de type fini sur $S$. Le groupe K
$K'_0(X) = K_0(\textit{Coh}(X))$ des faisceaux cohérents sur $X$
possède une filtration croissante canonique
$$
F_kK'_0(X) =
\Im\Big(K_0(\textit{Coh}_{\leq k}(X)) \to K_0(\textit{Coh}(X)\Big)
$$
Celle-ci est appelée filtration par la dimension des supports. Remarquons que
$$
\text{gr}_k K'_0(X) \subset K'_0(X)/F_{k - 1}K'_0(X) =
K_0(\textit{Coh}(X)/\textit{Coh}_{\leq k - 1}(X))
$$
où l'égalité est valable
d'après Homologie, lemme \ref{homology-lemma-serre-subcategory-K-groups}.
La discussion de la remarque \ref{chow-remark-good-cases-K-A} montre
qu'il existe des applications canoniques
$$
\CH_k(X) \longrightarrow \text{gr}_k K'_0(X)
$$
définies en envoyant la classe d'un sous-schéma fermé intègre
$Z \subset X$ de $\delta$-dimension $k$ sur la classe de
$[\mathcal{O}_Z]$ dans le membre de droite.

\begin{proposition}
\label{chow-proposition-K-tensor-Q}
Soit $(S, \delta)$ comme dans la situation \ref{chow-situation-setup}. Supposons donnée une
immersion fermée $X \to Y$ de schémas localement de type fini sur $S$,
avec $Y$ régulier et quasi-compact. Alors la composée
$$
K'_0(X) \to
K_0(D_{X, perf}(\mathcal{O}_Y)) \to
A^*(X \to Y) \otimes \mathbf{Q} \to
\CH_*(X) \otimes \mathbf{Q}
$$
de l'application $\mathcal{F} \mapsto \mathcal{F}[0]$ de la
remarque \ref{chow-remark-perf-Z-regular}, de l'application $ch(X \to Y, -)$ de la
remarque \ref{chow-remark-localized-chern-character-K} et
de l'application $c \mapsto c \cap [Y]$ induit un isomorphisme
$$
K'_0(X) \otimes \mathbf{Q}
\longrightarrow
\CH_*(X) \otimes \mathbf{Q}
$$
qui dépend du choix de $Y$. De plus, l'application canonique
$$
\CH_k(X) \otimes \mathbf{Q}
\longrightarrow
\text{gr}_k K'_0(X) \otimes \mathbf{Q}
$$
(voir ci-dessus) est un isomorphisme d'espaces vectoriels sur $\mathbf{Q}$ pour tout
$k \in \mathbf{Z}$.
\end{proposition}

\begin{proof}
Comme $Y$ est régulier, la construction de la
remarque \ref{chow-remark-perf-Z-regular} s'applique.
Comme $Y$ est quasi-compact, la construction de la
remarque \ref{chow-remark-localized-chern-character-K} s'applique.
Le schéma $Y$ est localement équidimensionnel
(lemme \ref{chow-lemma-locally-equidimensional}) ; ainsi,
le « cycle fondamental » $[Y]$ est défini
comme un élément de $\CH_*(Y)$, voir la remarque \ref{chow-remark-fundamental-class}.
En combinant ceci avec l'application $\CH_k(X) \to \text{gr}_kK'_0(X)$
construite ci-dessus, nous voyons qu'il suffit de prouver :
\begin{enumerate}
\item si $\mathcal{F}$ est un $\mathcal{O}_X$-module cohérent
dont le support est de $\delta$-dimension $\leq k$, alors
la composée ci-dessus envoie $[\mathcal{F}]$ dans
$\bigoplus_{k' \leq k} \CH_{k'}(X) \otimes \mathbf{Q}$ ;
\item si $Z \subset X$ est un sous-schéma fermé intègre
de $\delta$-dimension $k$, alors la composée ci-dessus
envoie $[\mathcal{O}_Z]$ sur un élément dont la composante de degré $k$
est la classe de $[Z]$ dans $\CH_k(X) \otimes \mathbf{Q}$.
\end{enumerate}
En effet, si cela vaut, nos applications induisent des applications
$\text{gr}_kK'_0(X) \otimes \mathbf{Q} \to CH_k(X) \otimes \mathbf{Q}$
qui sont inverses des applications canoniques
$\CH_k(X) \otimes \mathbf{Q} \to \text{gr}_k K'_0(X) \otimes \mathbf{Q}$
données avant la proposition.

\medskip\noindent
Étant donné un $\mathcal{O}_X$-module cohérent $\mathcal{F}$,
la composée ci-dessus envoie $[\mathcal{F}]$ sur
$$
ch(X \to Y, \mathcal{F}[0]) \cap [Y] \in \CH_*(X) \otimes \mathbf{Q}
$$
Si $\mathcal{F}$ a son support ensembliste dans un sous-schéma fermé
$Z \subset X$, alors nous avons
$$
ch(X \to Y, \mathcal{F}[0]) = (Z \to X)_* \circ ch(Z \to Y, \mathcal{F}[0])
$$
d'après le lemme \ref{chow-lemma-loc-chern-shrink-Z}. Nous concluons que, dans ce
cas, nous aboutissons dans l'image de $\CH_*(Z) \to \CH_*(X)$. Nous obtenons donc
la condition (1).

\medskip\noindent
Soit $Z \subset X$ un sous-schéma fermé intègre de $\delta$-dimension $k$.
La composée ci-dessus envoie $[\mathcal{O}_Z]$ sur l'élément
$$
ch(X \to Y, \mathcal{O}_Z[0]) \cap [Y] =
(Z \to X)_* ch(Z \to Y, \mathcal{O}_Z[0]) \cap [Y]
$$
par le même argument que ci-dessus.
Il suffit donc de prouver que la composante de degré $k$ de
$ch(Z \to Y, \mathcal{O}_Z[0]) \cap [Y] \in
\CH_*(Z) \otimes \mathbf{Q}$ est $[Z]$.
Comme $\CH_k(Z) = \mathbf{Z}$, pour
le prouver, nous pouvons remplacer $Y$ par un voisinage ouvert du
point générique $\xi$ de $Z$. Comme l'idéal maximal de l'anneau
local régulier $\mathcal{O}_{X, \xi}$ est engendré par une
suite régulière (Algèbre, lemme \ref{algebra-lemma-regular-ring-CM}),
nous pouvons supposer que l'idéal de $Z$ est engendré par une suite régulière, voir
Diviseurs, lemme \ref{divisors-lemma-Noetherian-scheme-regular-ideal}.
Nous déduisons donc le résultat du lemme \ref{chow-lemma-actual-computation}.
\end{proof}







\section{Produits d'intersection rationnels sur les schémas réguliers}
\label{chow-section-intersection-regular}

\noindent
Nous montrerons que $\CH_*(X) \otimes \mathbf{Q}$ possède un produit
d'intersection si $X$ est noethérien, régulier, de dimension finie et à
diagonale affine. La construction repose sur le résultat suivant
(qui est un corollaire de la proposition de la section précédente).

\begin{lemma}
\label{chow-lemma-K-tensor-Q}
Soit $(S, \delta)$ comme dans la situation \ref{chow-situation-setup}.
Soit $X$ un schéma régulier quasi-compact de type fini sur $S$, à
diagonale affine, et tel que $\delta_{X/S} : X \to \mathbf{Z}$ soit bornée.
Alors la composée
$$
K_0(\textit{Vect}(X)) \otimes \mathbf{Q}
\longrightarrow
A^*(X) \otimes \mathbf{Q}
\longrightarrow
\CH_*(X) \otimes \mathbf{Q}
$$
de l'application $ch$ de la remarque \ref{chow-remark-chern-character-K} et
de l'application $c \mapsto c \cap [X]$ est un isomorphisme.
\end{lemma}

\begin{proof}
Nous avons $K'_0(X) = K_0(X) = K_0(\textit{Vect}(X))$ d'après
Catégories dérivées de schémas, lemmes \ref{perfect-lemma-Kprime-K},
\ref{perfect-lemma-regular-resolution-property} et
\ref{perfect-lemma-K-is-old-K}.
D'après la remarque \ref{chow-remark-chern-classes-agree},
la composée donnée coïncide avec l'application de la
proposition \ref{chow-proposition-K-tensor-Q} pour $X = Y$.
Le résultat découle donc de la proposition.
\end{proof}

\noindent
Soient $X, S, \delta$ comme dans le lemme \ref{chow-lemma-K-tensor-Q}.
Pour simplifier, travaillons avec des cycles d'une codimension donnée, voir la
section \ref{chow-section-cycles-codimension}.
Soit $[X]$ le cycle fondamental de $X$, voir la
remarque \ref{chow-remark-fundamental-class}.
Choisissons $\alpha \in CH^i(X)$ et
$\beta \in CH^j(X)$. D'après le lemme, nous pouvons trouver un unique
$\alpha' \in K_0(\textit{Vect}(X)) \otimes \mathbf{Q}$ tel que
$ch(\alpha') \cap [X]  = \alpha$.
Bien sûr, cela signifie que $ch_{i'}(\alpha') \cap [X] = 0$
si $i' \not = i$ et que $ch_i(\alpha') \cap [X] = \alpha$.
D'après le lemme \ref{chow-lemma-adams-and-chern}, nous voyons que
$\alpha'' = 2^{-i}\psi^2(\alpha')$ est une autre solution.
Par unicité, nous obtenons $\alpha'' = \alpha'$, et nous concluons
que $ch_{i'}(\alpha) = 0$ dans $A^{i'}(X) \otimes \mathbf{Q}$
pour $i' \not = i$. Nous pouvons alors définir
$$
\alpha \cdot \beta = ch(\alpha') \cap \beta =
ch_i(\alpha') \cap \beta
$$
dans $\CH^{i + j}(X) \otimes \mathbf{Q}$ grâce à la propriété de $\alpha'$
observée ci-dessus. C'est un accouplement symétrique : en effet, si nous choisissons
$\beta' \in K_0(\textit{Vect}(X)) \otimes \mathbf{Q}$ relevant
$\beta$, alors nous obtenons
$$
\alpha \cdot \beta = ch(\alpha') \cap \beta =
ch(\alpha') \cap ch(\beta') \cap [X]
$$
et nous savons que les classes de Chern commutent.
Le produit d'intersection est associatif pour la même raison :
$$
(\alpha \cdot \beta) \cdot \gamma =
ch(\alpha') \cap ch(\beta') \cap ch(\gamma') \cap [X]
$$
car nous savons que la composition des classes bivariantes est associative.
Une meilleure formulation est peut-être la suivante : il existe
un unique produit d'intersection commutatif et associatif sur
$\CH^*(X) \otimes \mathbf{Q}$, compatible à la graduation, tel que
l'isomorphisme
$K_0(\textit{Vect}(X)) \otimes \mathbf{Q} \to \CH^*(X) \otimes \mathbf{Q}$
soit un isomorphisme d'anneaux.






\section{Applications de Gysin pour les morphismes d'intersection complète locale}
\label{chow-section-koszul}

\noindent
Avant de lire cette section, nous suggérons au lecteur de se renseigner sur les
immersions régulières
(Diviseurs, section \ref{divisors-section-regular-immersions}) et les
morphismes d'intersection complète locale
(Compléments sur les morphismes, section \ref{more-morphisms-section-lci}).

\medskip\noindent
Soit $(S, \delta)$ comme dans la situation \ref{chow-situation-setup}.
Soit $i : X \to Y$ une
immersion régulière\footnote{Voir
Diviseurs, définition \ref{divisors-definition-regular-immersion}.
Remarquons que les immersions régulières sont la même chose que les
immersions Koszul-régulières ou quasi-régulières
pour les schémas localement noethériens, voir
Diviseurs, lemme \ref{divisors-lemma-regular-immersion-noetherian}.
Nous l'utiliserons sans autre mention dans cette section.}
de schémas localement de type fini sur $S$.
En particulier, le faisceau conormal $\mathcal{C}_{X/Y}$ est localement libre de rang fini
(voir Diviseurs, lemme \ref{divisors-lemma-quasi-regular-immersion}). Ainsi, le
faisceau normal
$$
\mathcal{N}_{X/Y} = \SheafHom_{\mathcal{O}_X}(\mathcal{C}_{X/Y}, \mathcal{O}_X)
$$
est lui aussi localement libre de rang fini, et nous avons une surjection
$\mathcal{N}_{X/Y}^\vee \to \mathcal{C}_{X/Y}$ (car un isomorphisme
est aussi une surjection).
La construction de la section \ref{chow-section-gysin-higher-codimension}
nous donne une classe bivariante canonique
$$
i^! = c(X \to Y, \mathcal{N}_{X/Y}) \in A^*(X \to Y)^\wedge
$$
Nous avons besoin de deux lemmes sur cette notion.

\begin{lemma}
\label{chow-lemma-composition-regular-immersion}
Soit $(S, \delta)$ comme dans la situation \ref{chow-situation-setup}.
Soient $i : X \to Y$ et $j : Y \to Z$ des immersions régulières
de schémas localement de type fini sur $S$. Alors
$j \circ i$ est une immersion régulière et
$(j \circ i)^! = i^! \circ j^!$.
\end{lemma}

\begin{proof}
La première assertion résulte de
Diviseurs, lemme \ref{divisors-lemma-composition-regular-immersion}.
D'après Diviseurs, lemme \ref{divisors-lemma-transitivity-conormal-quasi-regular},
il existe une suite exacte courte
$$
0 \to
i^*(\mathcal{C}_{Y/Z}) \to
\mathcal{C}_{X/Z} \to
\mathcal{C}_{X/Y} \to 0
$$
Le résultat découle donc du lemme plus général \ref{chow-lemma-gysin-composition}.
\end{proof}

\begin{lemma}
\label{chow-lemma-section-smooth}
Soit $(S, \delta)$ comme dans la situation \ref{chow-situation-setup}.
Soit $p : P \to X$ un morphisme lisse de schémas localement de type fini
sur $S$, et soit $s : X \to P$ une section. Alors $s$ est une
immersion régulière et $1 = s^! \circ p^*$ dans $A^*(X)^\wedge$,
où $p^* \in A^*(P \to X)^\wedge$ est la classe bivariante
du lemme \ref{chow-lemma-flat-pullback-bivariant}.
\end{lemma}

\begin{proof}
La première assertion est Diviseurs, lemme
\ref{divisors-lemma-section-smooth-regular-immersion}.
Il suffit de montrer que $s^! \cap p^*[Z] = [Z]$ dans
$\CH_*(X)$ pour tout sous-schéma fermé intègre $Z \subset X$,
car les hypothèses sont préservées par changement de base par $X' \to X$
localement de type fini. Après avoir remplacé $P$ par un voisinage ouvert
de $s(Z)$, nous pouvons supposer que $P \to X$ est lisse de dimension relative constante $r$.
Posons $\dim_\delta(Z) = n$. Alors toute composante irréductible de $p^{-1}(Z)$
est de dimension $r + n$, et $p^*[Z]$ est donné par $[p^{-1}(Z)]_{n + r}$.
Remarquons que $s(X) \cap p^{-1}(Z) = s(Z)$ schématiquement. Ainsi, d'après la
même référence que ci-dessus, $s(X) \cap p^{-1}(Z)$ est un sous-schéma fermé
régulièrement immergé dans $\overline{p}^{-1}(Z)$, de codimension $r$.
Nous concluons par le lemme \ref{chow-lemma-gysin-fundamental}.
\end{proof}

\noindent
Soit $(S, \delta)$ comme dans la situation \ref{chow-situation-setup}. Considérons un
diagramme commutatif
$$
\xymatrix{
X \ar[rd]_f \ar[rr]_i & & P \ar[ld]^g \\
& Y
}
$$
de schémas localement de type fini sur $S$, tel que $g$ soit lisse
et $i$ une immersion régulière. En combinant la classe bivariante
$i^!$ examinée ci-dessus à la classe bivariante $g^* \in A^*(P \to Y)^\wedge$
du lemme \ref{chow-lemma-flat-pullback-bivariant}, nous obtenons
$$
f^! = i^! \circ g^* \in A^*(X \to Y)
$$
Remarquons que le morphisme $f$ est un morphisme d'intersection complète locale, voir
Compléments sur les morphismes, définition \ref{more-morphisms-definition-lci}.
Réciproquement, si $f : X \to Y$ est un morphisme d'intersection complète locale
de schémas localement noethériens et $f = g \circ i$ avec $g$ lisse, alors
$i$ est une immersion régulière. Nous affirmons que notre construction de $f^!$
ne dépend que du morphisme $f$ et non du choix de la factorisation
$f = g \circ i$.

\begin{lemma}
\label{chow-lemma-lci-gysin-well-defined}
Soit $(S, \delta)$ comme dans la situation \ref{chow-situation-setup}.
Soit $f : X \to Y$ un morphisme d'intersection complète locale
de schémas localement de type fini sur $S$.
La classe bivariante $f^!$ est indépendante du choix de
la factorisation $f = g \circ i$ avec $g$ lisse (pourvu
qu'une telle factorisation existe).
\end{lemma}

\begin{proof}
Étant donnée une seconde factorisation $f = g' \circ i'$, nous pouvons
considérer le morphisme lisse $g'' : P \times_Y P' \to Y$, une
immersion $i'' : X \to P \times_Y P'$ et la factorisation
$f = g'' \circ i''$. Nous pouvons donc supposer que nous avons un diagramme
$$
\xymatrix{
& P' \ar[d]^p \ar[rd]^{g'} \\
X \ar[r]^i \ar[ru]^{i'} & P \ar[r]^g & Y
}
$$
où $p$ est un morphisme lisse. Alors $(g')^* = p^* \circ g^*$
(lemme \ref{chow-lemma-compose-flat-pullback}), et il suffit donc
de montrer que $i^! = (i')^! \circ p^*$
dans $A^*(X \to P)$. Considérons le diagramme commutatif
$$
\xymatrix{
& X \times_P P' \ar[d]^{\overline{p}} \ar[r]_j & P' \ar[d]^p \\
X \ar[ru]^s \ar[r]^1 & X \ar[r]^i & P
}
$$
où $s =(1, i')$. Alors $s$ et $j$ sont des immersions régulières
(d'après Diviseurs, lemme \ref{divisors-lemma-section-smooth-regular-immersion},
et Diviseurs, lemme \ref{divisors-lemma-flat-base-change-regular-immersion}),
et $i' = j \circ s$. D'après le lemme \ref{chow-lemma-composition-regular-immersion},
nous avons $(i')^! = s^! \circ j^!$.
Comme le carré est cartésien, la classe bivariante $j^!$
est la restriction (remarque \ref{chow-remark-restriction-bivariant})
de $i^!$ à $P'$, voir le lemme \ref{chow-lemma-construction-gysin}.
Comme les classes bivariantes commutent aux images inverses plates,
nous trouvons $j^! \circ p^* = \overline{p}^* \circ i^!$.
Il suffit donc de montrer que $s^! \circ \overline{p}^* = \text{id}$,
ce qui est fait dans le lemme \ref{chow-lemma-section-smooth}.
\end{proof}

\begin{definition}
\label{chow-definition-lci-gysin}
Soit $(S, \delta)$ comme dans la situation \ref{chow-situation-setup}.
Soit $f : X \to Y$ un morphisme d'intersection complète locale
de schémas localement de type fini sur $S$. Nous disons que
{\it l'application de Gysin pour $f$ existe} si nous pouvons écrire
$f = g \circ i$ avec $g$ lisse et $i$ une immersion.
Dans ce cas, nous définissons
{\it l'application de Gysin} $f^! = i^! \circ g^* \in A^*(X \to Y)$ comme ci-dessus.
\end{definition}

\noindent
Il résulte de la définition que, pour une immersion régulière,
ceci coïncide avec la construction antérieure et que, pour un morphisme
lisse, ceci coïncide avec l'image inverse plate. En fait, cette coïncidence
vaut pour tous les morphismes syntomiques.

\begin{lemma}
\label{chow-lemma-lci-gysin-flat}
Soit $(S, \delta)$ comme dans la situation \ref{chow-situation-setup}.
Soit $f : X \to Y$ un morphisme d'intersection complète locale
de schémas localement de type fini sur $S$. Si l'application de Gysin
existe pour $f$ et si $f$ est plat, alors $f^!$ est égal à la
classe bivariante du lemme \ref{chow-lemma-flat-pullback-bivariant}.
\end{lemma}

\begin{proof}
Choisissons une factorisation $f = g \circ i$, avec $i : X \to P$
une immersion et $g : P \to Y$ lisse. Remarquons que, pour
tout morphisme $Y' \to Y$ localement de type fini,
les changements de base $f'$, $g'$, $i'$ satisfont aux mêmes
hypothèses (voir Morphismes, lemmes \ref{morphisms-lemma-base-change-smooth}
et \ref{morphisms-lemma-base-change-syntomic}, et
Compléments sur les morphismes, lemme \ref{more-morphisms-lemma-flat-lci}).
Nous nous ramenons donc à prouver que $f^*[Y] = i^!(g^*[Y])$ lorsque $Y$
est intègre, voir le lemme \ref{chow-lemma-bivariant-zero}. Posons $n = \dim_\delta(Y)$.
Après avoir décomposé $X$ et $P$ en composantes connexes, nous
pouvons supposer que $f$ est plat de dimension relative $r$ et
que $g$ est lisse de dimension relative $t$.
Alors $f^*[Y] = [X]_{n + s}$ et $g^*[Y] = [P]_{n + t}$.
D'autre part, $i$ est une immersion régulière de codimension $t - s$.
Ainsi, $i^![P]_{n + t} = [X]_{n + s}$ (lemme \ref{chow-lemma-gysin-fundamental}),
et la démonstration est achevée.
\end{proof}

\begin{lemma}
\label{chow-lemma-lci-gysin-composition}
Soit $(S, \delta)$ comme dans la situation \ref{chow-situation-setup}.
Soient $f : X \to Y$ et $g : Y \to Z$ des morphismes d'intersection complète locale
de schémas localement de type fini sur $S$. Supposons que l'application de Gysin
existe pour $g \circ f$ et $g$. Alors l'application de Gysin existe pour $f$
et $(g \circ f)^! = f^! \circ g^!$.
\end{lemma}

\begin{proof}
Remarquons que $g \circ f$ est un morphisme d'intersection complète locale
d'après Compléments sur les morphismes, lemme \ref{more-morphisms-lemma-composition-lci},
et que l'énoncé du lemme a donc un sens.
Si $X \to P$ est une immersion de $X$ dans un schéma $P$ lisse sur $Z$,
alors $X \to P \times_Z Y$ est une immersion de $X$ dans un schéma lisse
sur $Y$. Cela prouve la première assertion du lemme.
Soit $Y \to P'$ une immersion de $Y$ dans un schéma $P'$ lisse sur $Z$.
Considérons le diagramme commutatif
$$
\xymatrix{
X \ar[r] \ar[d] &
P \times_Z Y \ar[r]_a \ar[ld]^p &
P \times_Z P' \ar[ld]^q \\
Y \ar[r]_b \ar[d] &
P' \ar[ld] \\
Z
}
$$
Ici, les flèches horizontales sont des immersions régulières, les flèches sud-ouest
sont lisses, et le carré est cartésien. Par conséquent,
$a^! \circ q^* = p^* \circ b^!$, puisque les classes bivariantes commutent
aux images inverses plates. En combinant ce fait avec les
lemmes \ref{chow-lemma-composition-regular-immersion} et
\ref{chow-lemma-compose-flat-pullback},
le lecteur constate que l'énoncé du lemme est vrai.
Nous omettons un petit détail.
\end{proof}

\begin{lemma}
\label{chow-lemma-lci-gysin-commutes}
Soit $(S, \delta)$ comme dans la situation \ref{chow-situation-setup}.
Considérons un diagramme commutatif
$$
\xymatrix{
X'' \ar[d] \ar[r] &
X' \ar[d] \ar[r] &
X \ar[d]^f \\
Y'' \ar[r] &
Y' \ar[r] &
Y
}
$$
de schémas localement de type fini sur $S$, dont les deux carrés sont cartésiens.
Supposons que $f : X \to Y$ soit un morphisme d'intersection complète locale
tel que l'application de Gysin existe pour $f$. Soit $c \in A^*(Y'' \to Y')$. Notons
$res(f^!) \in A^*(X' \to Y')$ la restriction de $f^!$ à $Y'$
(remarque \ref{chow-remark-restriction-bivariant}). Alors $c$ et $res(f^!)$ commutent
(remarque \ref{chow-remark-bivariant-commute}).
\end{lemma}

\begin{proof}
Choisissons une factorisation $f = g \circ i$ avec $g$ lisse et $i$ une immersion.
Comme $f^! = i^! \circ g^!$, il suffit de prouver le lemme pour $g^!$
(qui est donné par l'image inverse plate) et pour $i^!$. Le résultat pour l'image inverse plate
fait partie de la définition d'une classe bivariante. Le cas de $i^!$ découle
immédiatement du lemme \ref{chow-lemma-gysin-commutes}.
\end{proof}

\begin{lemma}
\label{chow-lemma-lci-gysin-easy}
Soit $(S, \delta)$ comme dans la situation \ref{chow-situation-setup}.
Considérons un diagramme cartésien
$$
\xymatrix{
X' \ar[d]_{f'} \ar[r] &
X \ar[d]^f \\
Y' \ar[r] &
Y
}
$$
de schémas localement de type fini sur $S$. Supposons que
\begin{enumerate}
\item $f$ soit un morphisme d'intersection complète locale et que
l'application de Gysin existe pour $f$,
\item $X$, $X'$, $Y$, $Y'$ vérifient les conditions équivalentes du
lemme \ref{chow-lemma-locally-equidimensional},
\item pour $x' \in X'$, d'images $x$, $y'$ et $y$
dans $X$, $Y'$ et $Y$, on ait $n_{x'} - n_{y'} = n_x - n_y$,
où $n_{x'}$, $n_x$, $n_{y'}$ et $n_y$ sont comme dans le lemme, et
\item pour tout point générique $\xi \in X'$, l'anneau local
$\mathcal{O}_{Y', f'(\xi)}$ soit de Cohen-Macaulay.
\end{enumerate}
Alors $f^![Y'] = [X']$, où $[Y']$ et $[X']$ sont comme dans la
remarque \ref{chow-remark-fundamental-class}.
\end{lemma}

\begin{proof}
Rappelons que $n_{x'}$ est la valeur commune de $\delta(\xi)$
lorsque $\xi$ parcourt les points génériques des composantes irréductibles
passant par $x'$. De plus, les fonctions
$x' \mapsto n_{x'}$, $x \mapsto n_x$, $y' \mapsto n_{y'}$ et
$y \mapsto n_y$ sont localement constantes. Soient $X'_n$, $X_n$, $Y'_n$
et $Y_n$ les sous-schémas ouverts et fermés de $X'$, $X$, $Y'$ et
$Y$ où la fonction vaut $n$. Rappelons que
$[X'] = \sum [X'_n]_n$ et $[Y'] = \sum [Y'_n]_n$.
Cela étant dit, il est clair que, pour démontrer le lemme, nous
pouvons remplacer $X'$ par l'une de ses composantes connexes
et $X$, $Y'$, $Y$ par la composante connexe dans laquelle
elle s'envoie. Nous savons alors que $X'$, $X$, $Y'$ et
$Y$ sont $\delta$-équidimensionnels, au sens où
chaque composante irréductible a la même $\delta$-dimension.
Notons $n'$, $n$, $m'$ et $m$ cette valeur commune
pour $X'$, $X$, $Y'$ et $Y$. La dernière hypothèse
signifie que $n' - m' = n - m$.

\medskip\noindent
Choisissons une factorisation $f = g \circ i$, où $i : X \to P$
est une immersion et $g : P \to Y$ est lisse. Comme $X$ est connexe,
nous voyons que la dimension relative de $P \to Y$ aux points de $i(X)$
est constante. Ainsi, après avoir remplacé $P$ par un voisinage ouvert
de $i(X)$, nous pouvons supposer que $P \to Y$ est de dimension relative constante
et que $i : X \to P$ est une immersion fermée.
Notons $g' : Y' \times_Y P \to Y'$ le changement de base de $g$ et
$i' : X' \to Y' \times_Y P$ le changement de base de $i$.
Il est clair que $g^*[Y] = [P]$ et $(g')^*[Y'] = [Y' \times_Y P]$.
Enfin, si $\xi' \in X'$ est un point générique, alors
$\mathcal{O}_{Y' \times_Y P, i'(\xi)}$ est de Cohen-Macaulay.
En effet, l'homomorphisme d'anneaux locaux
$\mathcal{O}_{Y', f'(\xi)} \to \mathcal{O}_{Y' \times_Y P, i'(\xi)}$
est plat à fibre régulière
(voir Algèbre, section \ref{algebra-section-smooth-overview}),
un anneau local régulier est de Cohen-Macaulay
(Algèbre, lemme \ref{algebra-lemma-regular-ring-CM}),
$\mathcal{O}_{Y', f'(\xi)}$ est de Cohen-Macaulay par l'hypothèse
(4), et le résultat voulu découle de
Algèbre, lemme \ref{algebra-lemma-CM-goes-up}.
Nous nous ramenons donc au cas traité dans le paragraphe suivant.

\medskip\noindent
Supposons que $f$ soit une immersion fermée régulière et que $X'$, $X$, $Y'$ et
$Y$ soient $\delta$-équidimensionnels, de $\delta$-dimensions
$n'$, $n$, $m'$ et $m$, avec $m' - n' = m - n$.
Dans ce cas, le résultat découle immédiatement du
lemme \ref{chow-lemma-gysin-easy}.
\end{proof}

\begin{remark}
\label{chow-remark-gysin-chern-classes}
Soit $(S, \delta)$ comme dans la situation \ref{chow-situation-setup}.
Soit $f : X \to Y$ un morphisme d'intersection complète locale
de schémas localement de type fini sur $S$. Supposons que l'application
de Gysin existe pour $f$. Alors
$f^! \circ c_i(\mathcal{E}) = c_i(f^*\mathcal{E}) \circ f^!$,
et de même pour le caractère de Chern, voir le
lemme \ref{chow-lemma-lci-gysin-commutes}.
Si $X$ et $Y$ vérifient les conditions équivalentes du
lemme \ref{chow-lemma-locally-equidimensional} et si $Y$ est de Cohen-Macaulay
(par exemple), alors $f^![Y] = [X]$ par le lemme \ref{chow-lemma-lci-gysin-easy}.
Dans ce cas, nous obtenons aussi
$f^!(c_i(\mathcal{E}) \cap [Y]) = c_i(f^*\mathcal{E}) \cap [X]$,
et de même pour le caractère de Chern.
\end{remark}

\begin{lemma}
\label{chow-lemma-compare-gysin-base-change}
Soit $(S, \delta)$ comme dans la situation \ref{chow-situation-setup}.
Considérons un carré cartésien
$$
\xymatrix{
X' \ar[d]_{f'} \ar[r]_{g'} &
X \ar[d]^f \\
Y' \ar[r]^g &
Y
}
$$
de schémas localement de type fini sur $S$. Supposons que
\begin{enumerate}
\item $f$ et $f'$ soient tous deux des morphismes d'intersection complète locale, et
\item l'application de Gysin existe pour $f$.
\end{enumerate}
Alors $\mathcal{C} = \Ker(H^{-1}((g')^*\NL_{X/Y}) \to H^{-1}(\NL_{X'/Y'}))$
est un $\mathcal{O}_{X'}$-module localement libre de rang fini, l'application de Gysin
existe pour $f'$, et nous avons
$$
res(f^!) = c_{top}(\mathcal{C}^\vee) \circ (f')^!
$$
dans $A^*(X' \to Y')$.
\end{lemma}

\begin{proof}
Le fait que $\mathcal{C}$ soit localement libre de rang fini découle immédiatement de
Compléments d'algèbre, lemme \ref{more-algebra-lemma-base-change-lci-bis}.
Choisissons une factorisation $f = h \circ i$ avec $h : P \to Y$ lisse et $i$
une immersion. Nous pouvons alors factoriser $f' = h' \circ i'$, où $h' : P' \to Y'$
et $i' : X' \to P'$ sont les changements de base correspondants. On a le diagramme
$$
\xymatrix{
X' \ar[r] \ar[d] &
P' \ar[r] \ar[d] &
Y' \ar[d] \\
X \ar[r]^i &
P \ar[r]^h &
Y
}
$$
En particulier, nous voyons que l'application de Gysin existe pour $f'$. Par
Compléments sur les morphismes, lemme \ref{more-morphisms-lemma-get-NL},
nous avons
$$
\NL_{X/Y} = \left( \mathcal{C}_{X/P} \to i^*\Omega_{P/Y} \right)
$$
où $\mathcal{C}_{X/P}$ est le faisceau conormal de l'immersion $i$.
De même pour la version avec primes. Nous avons
$(g')^*i^*\Omega_{P/Y} = (i')^*\Omega_{P'/Y'}$, car
$\Omega_{P/Y}$ devient $\Omega_{P'/Y'}$ par image inverse, d'après
Morphismes, lemme \ref{morphisms-lemma-base-change-differentials}.
Rappelons aussi que $(g')^*\mathcal{C}_{X/P} \to \mathcal{C}_{X'/P'}$
est surjectif, voir
Morphismes, lemme \ref{morphisms-lemma-conormal-functorial-flat}.
Nous en déduisons que le faisceau $\mathcal{C}$ est canoniquement
isomorphe au noyau de l'application
$(g')^*\mathcal{C}_{X/P} \to \mathcal{C}_{X'/P'}$
de modules localement libres de rang fini. Rappelons que $i^!$ est défini
à l'aide de $\mathcal{N} = \mathcal{C}_{Z/X}^\vee$, et de même
pour $(i')^!$. Ainsi, nous avons
$$
res(i^!) = c_{top}(\mathcal{C}^\vee) \circ (i')^!
$$
dans $A^*(X' \to P')$, par application du lemme \ref{chow-lemma-gysin-excess}.
Enfin, puisque $f^! = i^! \circ h^*$,
$(f')^! = (i')^! \circ (h')^*$ et $(h')^* = res(h^*)$, nous concluons.
\end{proof}

\begin{lemma}[Formule de l'éclatement]
\label{chow-lemma-blow-up-formula}
Soit $(S, \delta)$ comme dans la situation \ref{chow-situation-setup}.
Soit $i : Z \to X$ une immersion fermée régulière de schémas
localement de type fini sur $S$. Soit $b : X' \to X$
l'éclatement de centre $Z$. On a le diagramme
$$
\xymatrix{
E \ar[r]_j \ar[d]_\pi & X' \ar[d]^b \\
Z \ar[r]^i & X
}
$$
Supposons que l'application de Gysin existe pour $b$. Alors
$$
res(b^!) = c_{top}(\mathcal{F}^\vee) \circ \pi^*
$$
dans $A^*(E \to Z)$, où $\mathcal{F}$ est le noyau de l'application canonique
$\pi^*\mathcal{C}_{Z/X} \to \mathcal{C}_{E/X'}$.
\end{lemma}

\begin{proof}
Observons que le morphisme $b$ est un morphisme d'intersection complète locale
par Compléments d'algèbre, lemme \ref{more-algebra-lemma-blowup-regular-sequence},
et que l'énoncé a donc un sens. Comme $Z \to X$ est une
immersion régulière (et donc a fortiori quasi-régulière), nous voyons que $\mathcal{C}_{Z/X}$
est localement libre de rang fini et que l'application
$\text{Sym}^*(\mathcal{C}_{Z/X}) \to \mathcal{C}_{Z/X, *}$
est un isomorphisme, voir
Diviseurs, lemme \ref{divisors-lemma-quasi-regular-immersion}.
Comme $E = \text{Proj}(\mathcal{C}_{Z/X, *})$, nous concluons
que $E = \mathbf{P}(\mathcal{C}_{Z/X})$
est un fibré en espaces projectifs sur $Z$.
Ainsi $E \to Z$ est lisse et, en particulier, un morphisme d'intersection complète
locale. Le lemme \ref{chow-lemma-compare-gysin-base-change}
s'applique donc, et nous voyons que
$$
res(b^!) = c_{top}(\mathcal{C}^\vee) \circ \pi^!
$$
avec $\mathcal{C}$ comme dans l'énoncé de ce lemme.
Bien entendu, $\pi^* = \pi^!$ par le lemme \ref{chow-lemma-lci-gysin-flat}.
Il reste à montrer que $\mathcal{F}$ est égal au
noyau $\mathcal{C}$ de l'application
$H^{-1}(j^*\NL_{X'/X}) \to H^{-1}(\NL_{E/Z})$.

\medskip\noindent
Comme $E \to Z$ est lisse, nous avons $H^{-1}(\NL_{E/Z}) = 0$, voir
Compléments sur les morphismes, lemme \ref{more-morphisms-lemma-NL-smooth}.
Il suffit donc de montrer que $\mathcal{F}$ s'identifie
à $H^{-1}(j^*\NL_{X'/X})$. Par Compléments sur les morphismes, Lemmes
\ref{more-morphisms-lemma-get-triangle-NL} et
\ref{more-morphisms-lemma-NL-immersion}, nous avons une suite exacte
$$
0 \to H^{-1}(j^*\NL_{X'/X}) \to H^{-1}(\NL_{E/X}) \to
\mathcal{C}_{E/X'} \to \ldots
$$
Les mêmes lemmes appliqués à $E \to Z \to X$ donnent un isomorphisme
$\pi^*\mathcal{C}_{Z/X} = H^{-1}(\pi^*\NL_{Z/X}) \to H^{-1}(\NL_{E/X})$.
Nous pouvons donc conclure.
\end{proof}

\begin{lemma}
\label{chow-lemma-lci-gysin-product-regular}
Soit $(S, \delta)$ comme dans la situation \ref{chow-situation-setup}.
Soit $f : X \to Y$ un morphisme entre des schémas localement de type
fini sur $S$, tels que $X$ et $Y$ soient quasi-compacts,
réguliers, à diagonale affine et de dimension finie.
Alors $f$ est un morphisme d'intersection complète locale.
Supposons de plus que l'application de Gysin existe pour $f$. Alors
$$
f^!(\alpha \cdot \beta) = f^!\alpha \cdot f^!\beta
$$
dans $\CH^*(X) \otimes \mathbf{Q}$, où le produit d'intersection
est celui de la section \ref{chow-section-intersection-regular}.
\end{lemma}

\begin{proof}
Le premier énoncé découle de
Compléments sur les morphismes, Lemme
\ref{more-morphisms-lemma-morphism-regular-schemes-is-lci}.
Observons que $f^![Y] = [X]$, voir le lemme \ref{chow-lemma-lci-gysin-easy}.
Écrivons $\alpha = ch(\alpha') \cap [Y]$ et $\beta = ch(\beta') \cap [Y]$,
où $\alpha', \beta' \in K_0(\textit{Vect}(X)) \otimes \mathbf{Q}$,
comme dans la section \ref{chow-section-intersection-regular}.
En posant $c = ch(\alpha')$ et $c' = ch(\beta')$, nous obtenons
$\alpha \cdot \beta = c \cap c' \cap [Y]$ par construction.
D'après le lemme \ref{chow-lemma-lci-gysin-commutes}, nous savons que $f^!$
commute avec $c$ et $c'$. Par conséquent,
\begin{align*}
f^!(\alpha \cdot \beta)
& =
f^!(c \cap c' \cap [Y]) \\
& =
c \cap c' \cap f^![Y] \\
& =
c \cap c' \cap [X] \\
& =
(c \cap [X]) \cdot (c' \cap [X]) \\
& =
(c \cap f^![Y]) \cdot (c' \cap f^![Y]) \\
& =
f^!(\alpha) \cdot f^!(\beta)
\end{align*}
comme voulu.
\end{proof}

\begin{lemma}
\label{chow-lemma-projection-formula-regular}
Soit $(S, \delta)$ comme dans la situation \ref{chow-situation-setup}.
Soit $f : X \to Y$ un morphisme entre des schémas localement de type
fini sur $S$, tels que $X$ et $Y$ soient quasi-compacts,
réguliers, à diagonale affine et de dimension finie.
Alors $f$ est un morphisme d'intersection complète locale.
Supposons de plus que l'application de Gysin existe pour $f$
et que $f$ soit propre. Alors
$$
f_*(\alpha \cdot f^!\beta) = f_*\alpha \cdot \beta
$$
dans $\CH^*(Y) \otimes \mathbf{Q}$, où le produit d'intersection
est celui de la section \ref{chow-section-intersection-regular}.
\end{lemma}

\begin{proof}
Le premier énoncé découle de
Compléments sur les morphismes, Lemme
\ref{more-morphisms-lemma-morphism-regular-schemes-is-lci}.
Observons que $f^![Y] = [X]$, voir le lemme \ref{chow-lemma-lci-gysin-easy}.
Écrivons $\alpha = ch(\alpha') \cap [X]$ et $\beta = ch(\beta') \cap [Y]$,
avec $\alpha' \in K_0(\textit{Vect}(X)) \otimes \mathbf{Q}$ et
$\beta' \in K_0(\textit{Vect}(Y)) \otimes \mathbf{Q}$,
comme dans la section \ref{chow-section-intersection-regular}.
Posons $c = ch(\alpha')$ et $c' = ch(\beta')$. Nous avons
\begin{align*}
f_*(\alpha \cdot f^!\beta)
& =
f_*(c \cap f^!(c' \cap [Y]_e)) \\
& =
f_*(c \cap c' \cap f^![Y]_e) \\
& =
f_*(c \cap c' \cap [X]_d) \\
& =
f_*(c' \cap c \cap [X]_d) \\
& =
c' \cap f_*(c \cap [X]_d) \\
& =
\beta \cdot f_*(\alpha)
\end{align*}
La première égalité résulte de la construction du produit d'intersection.
D'après le lemme \ref{chow-lemma-lci-gysin-commutes}, nous savons que $f^!$
commute avec $c'$. Le fait que les classes de Chern soient dans le centre
de l'anneau bivariant justifie d'intervertir l'ordre dans lequel on intersecte
$[X]$ avec $c$ et $c'$. On peut commuter $c'$ avec $f_*$, car $c'$
est une classe bivariante. La dernière égalité résulte à nouveau de la construction
du produit d'intersection.
\end{proof}













\section{Applications de Gysin pour les diagonales}
\label{chow-section-gysin-for-diagonal}

\noindent
Soit $(S, \delta)$ comme dans la situation \ref{chow-situation-setup}. Soit $f : X \to Y$
un morphisme lisse de schémas localement de type fini sur $S$. Alors le
morphisme diagonal $\Delta : X \longrightarrow X \times_Y X$
est une immersion régulière, voir
Compléments sur les morphismes, lemme \ref{more-morphisms-lemma-smooth-diagonal-perfect}.
Nous avons donc l'application de Gysin
$$
\Delta^! \in A^*(X \to X \times_Y X)^\wedge
$$
construite à la section \ref{chow-section-koszul}. Si $X \to Y$ est de
dimension relative constante $d$, alors $\Delta^! \in A^d(X \to X \times_Y X)$.

\begin{lemma}
\label{chow-lemma-diagonal-identity}
Dans la situation ci-dessus, nous avons $\Delta^! \circ \text{pr}_i^! = 1$ dans $A^0(X)$.
\end{lemma}

\begin{proof}
Observons que les projections $\text{pr}_i : X \times_Y X \to X$ sont
lisses, de sorte que nous disposons aussi d'applications de Gysin pour ces projections.
Le lemme a donc un sens et est un cas particulier du
lemme \ref{chow-lemma-lci-gysin-composition}.
\end{proof}

\begin{proposition}
\label{chow-proposition-compute-bivariant}
\begin{reference}
\cite[Proposition 17.4.2]{F}
\end{reference}
Soit $(S, \delta)$ comme dans la situation \ref{chow-situation-setup}.
Soient $f : X \to Y$ et $g : Y \to Z$ des morphismes de schémas localement
de type fini sur $S$. Si $g$ est lisse de dimension relative $d$, alors
$A^p(X \to Y) = A^{p - d}(X \to Z)$.
\end{proposition}

\begin{proof}
Nous utiliserons le fait que les morphismes lisses sont des morphismes d'intersection complète
locale dont les applications de Gysin existent (voir la section \ref{chow-section-koszul}).
En particulier, nous avons $g^! \in A^{-d}(Y \to Z)$. Nous pouvons alors envoyer
$c \in A^p(X \to Y)$ sur $c \circ g^! \in A^{p - d}(X \to Z)$.

\medskip\noindent
Réciproquement, soit $c' \in A^{p - d}(X \to Z)$. Notons $res(c')$ la restriction
(remarque \ref{chow-remark-restriction-bivariant}) de $c'$ par le morphisme $Y \to Z$.
Puisque le diagramme
$$
\xymatrix{
X \times_Z Y \ar[r]_{\text{pr}_2} \ar[d]_{\text{pr}_1} & Y \ar[d]^g \\
X \ar[r]^f & Z
}
$$
est cartésien, nous obtenons $res(c') \in A^{p - d}(X \times_Z Y \to Y)$.
Soit $\Delta : Y \to Y \times_Z Y$ la diagonale, et notons
$res(\Delta^!)$ la restriction de $\Delta^!$
à $X \times_Z Y$ par le morphisme $X \times_Z Y \to Y \times_Z Y$.
Puisque le diagramme
$$
\xymatrix{
X \ar[r] \ar[d] & X \times_Z Y \ar[d] \\
Y \ar[r]^-\Delta & Y \times_Z Y
}
$$
est cartésien, nous voyons que $res(\Delta^!) \in A^d(X \to X \times_Z Y)$.
En composant ces deux restrictions, nous obtenons
$$
res(\Delta^!) \circ res(c') \in A^p(X \to Y)
$$
Nous avons ainsi construit des applications $A^p(X \to Y) \to A^{p - d}(X \to Z)$
et $A^{p - d}(X \to Z) \to A^p(X \to Y)$. Pour achever la démonstration, nous
allons montrer que ces applications sont réciproques l'une de l'autre.

\medskip\noindent
Commençons avec $c \in A^p(X \to Y)$. Considérons le diagramme
$$
\xymatrix{
X \ar[d] \ar[r] & Y \ar[d] \\
X \times_Z Y \ar[r] \ar[d]^{\text{pr}_1} &
Y \times_Z Y \ar[r]_{p_2} \ar[d]^{p_1} &
Y \ar[d]^g \\
X \ar[r]^f &
Y \ar[r]^g &
Z
}
$$
dont les carrés sont cartésiens. Les deux carrés inférieurs de ce diagramme
montrent que $res(c \circ g^!) = res(c) \cap p_2^!$, où, dans cette formule,
$res(c)$ désigne la restriction de $c$ par $p_1$. En considérant le carré supérieur
du diagramme et en utilisant le lemme \ref{chow-lemma-lci-gysin-commutes},
nous obtenons $c \circ \Delta^! = res(\Delta^!) \circ res(c)$.
Nous calculons
\begin{align*}
res(\Delta^!) \circ res(c \circ g^!)
& =
res(\Delta^!) \circ res(c) \circ p_2^! \\
& =
c \circ \Delta^! \circ p_2^! \\
& =
c
\end{align*}
La dernière égalité découle du lemme \ref{chow-lemma-diagonal-identity}.

\medskip\noindent
Réciproquement, partons de $c' \in A^{p - d}(X \to Z)$. En considérant
le rectangle inférieur du diagramme ci-dessus, nous obtenons
$res(c') \circ g^! = \text{pr}_1^! \circ c'$.
Nous calculons
\begin{align*}
res(\Delta^!) \circ res(c') \circ g^!
& =
res(\Delta^!) \circ \text{pr}_1^! \circ c' \\
& =
c'
\end{align*}
La dernière égalité résulte du fait que les deux carrés de gauche du
diagramme montrent que
$\text{id} = res(\Delta^! \circ p_1^!) = res(\Delta^!) \circ \text{pr}_1^!$.
Ceci achève la démonstration.
\end{proof}









\section{Produit extérieur}
\label{chow-section-exterior-product}

\noindent
Soit $k$ un corps. Dans cette section, nous travaillons sur $S = \Spec(k)$ avec
$\delta : S \to \mathbf{Z}$ définie en envoyant l'unique point sur $0$, voir
l'Exemple \ref{chow-example-field}.

\medskip\noindent
Considérons un carré cartésien
$$
\xymatrix{
X \times_k Y \ar[r] \ar[d] & Y \ar[d] \\
X \ar[r] & \Spec(k) = S
}
$$
de schémas localement de type fini sur $k$. Il existe alors une application canonique
$$
\times :
\CH_n(X) \otimes_{\mathbf{Z}} \CH_m(Y)
\longrightarrow
\CH_{n + m}(X \times_k Y)
$$
qui est uniquement déterminée par la règle suivante :
étant donnés des sous-schémas fermés intègres $X' \subset X$
et $Y' \subset Y$ de dimensions $n$ et $m$, nous avons
$$
[X'] \times [Y'] = [X' \times_k Y']_{n + m}
$$
dans $\CH_{n + m}(X \times_k Y)$.

\begin{lemma}
\label{chow-lemma-exterior-product-well-defined}
L'application
$\times : \CH_n(X) \otimes_{\mathbf{Z}} \CH_m(Y) \to \CH_{n + m}(X \times_k Y)$
est bien définie.
\end{lemma}

\begin{proof}
Une première remarque est que, si $\alpha = \sum n_i[X_i]$
et $\beta = \sum m_j[Y_j]$, avec $X_i \subset X$ et $Y_j \subset Y$
des familles localement finies de sous-schémas fermés intègres
de dimensions $n$ et $m$, alors
les $X_i \times_k Y_j$ forment une famille localement finie
de sous-schémas fermés de $X \times_k Y$, de
dimensions $n + m$, et nous pouvons effectivement considérer
$$
\alpha \times \beta = \sum n_i m_j [X_i \times_k Y_j]_{n + m}
$$
comme un $(n + m)$-cycle sur $X \times_k Y$. Nous obtenons ainsi une
application additive
$\times : Z_n(X) \otimes_{\mathbf{Z}} Z_m(Y) \to Z_{n + m}(X \times_k Y)$.
Il s'agit de montrer que
ce procédé est compatible avec l'équivalence rationnelle.

\medskip\noindent
Soit $i : X' \to X$ le morphisme d'inclusion d'un
sous-schéma fermé intègre de dimension $n$.
Alors l'image inverse plate par le morphisme $p' : X' \to \Spec(k)$ est un élément
$(p')^* \in A^{-n}(X' \to \Spec(k))$ d'après le
lemme \ref{chow-lemma-flat-pullback-bivariant},
et donc $c' = i_* \circ (p')^* \in A^{-n}(X \to \Spec(k))$ d'après le
lemme \ref{chow-lemma-push-proper-bivariant}.
Cela produit des applications
$$
c' \cap - : \CH_m(Y) \longrightarrow \CH_{m + n}(X \times_k Y)
$$
dont le lecteur vérifie aisément qu'elles envoient $[Y']$ sur $[X' \times_k Y']_{n + m}$
pour tout sous-schéma fermé intègre $Y' \subset Y$ de dimension
$m$. Par conséquent, la construction
$([X'], [Y']) \mapsto [X' \times_k Y']_{n + m}$
se factorise par l'équivalence rationnelle en la seconde variable, c'est-à-dire
donne une application bien définie
$Z_n(X) \otimes_{\mathbf{Z}} \CH_m(Y) \to \CH_{n + m}(X \times_k Y)$.
Par symétrie, il en est de même pour l'autre variable, et nous concluons.
\end{proof}

\begin{lemma}
\label{chow-lemma-chow-cohomology-towards-point}
Soit $k$ un corps. Soit $X$ un schéma localement de type fini sur $k$.
Nous avons alors une identification canonique
$$
A^p(X \to \Spec(k)) = \CH_{-p}(X)
$$
pour tout $p \in \mathbf{Z}$.
\end{lemma}

\begin{proof}
Considérons l'élément $[\Spec(k)] \in \CH_0(\Spec(k))$. Nous obtenons une application
$A^p(X \to \Spec(k)) \to \CH_{-p}(X)$ qui envoie $c$ sur $c \cap [\Spec(k)]$.

\medskip\noindent
Réciproquement, supposons que nous ayons $\alpha \in \CH_{-p}(X)$.
Nous pouvons alors définir $c_\alpha \in A^p(X \to \Spec(k))$ comme
suit : étant donnés $X' \to \Spec(k)$ et $\alpha' \in \CH_n(X')$,
nous posons
$$
c_\alpha \cap \alpha' = \alpha \times \alpha'
$$
dans $\CH_{n - p}(X \times_k X')$. Pour montrer qu'il s'agit d'une classe
bivariante, écrivons $\alpha = \sum n_i[X_i]$ comme dans la
définition \ref{chow-definition-cycles}. Considérons la composée
$$
\coprod X_i \xrightarrow{g} X \to \Spec(k)
$$
et notons $f : \coprod X_i \to \Spec(k)$ cette composée.
Alors $g$ est propre et $f$ est plat de dimension relative $-p$.
L'image inverse par $f$ est une classe bivariante
$f^* \in A^p(\coprod X_i \to \Spec(k))$ d'après le
lemme \ref{chow-lemma-flat-pullback-bivariant}.
Notons $\nu \in A^0(\coprod X_i)$ la classe bivariante
qui multiplie un cycle par $n_i$ sur la $i$-ième composante.
Ainsi, $\nu \circ f^* \in A^p(\coprod X_i \to X)$.
Enfin, nous avons une classe bivariante
$$
g_* \circ \nu \circ f^*
$$
par le lemme \ref{chow-lemma-push-proper-bivariant}. Le lecteur vérifie aisément
que $c_\alpha$ est égale à cette classe et qu'elle est donc
elle-même une classe bivariante.

\medskip\noindent
Pour achever la démonstration, nous devons montrer que les deux constructions
sont réciproques l'une de l'autre. Comme $c_\alpha \cap [\Spec(k)] = \alpha$,
c'est clair dans un sens. Pour l'autre sens, soit
$c \in A^p(X \to \Spec(k))$ et posons $\alpha = c \cap [\Spec(k)]$.
Il suffit de démontrer que
$$
c \cap [X'] = c_\alpha \cap [X']
$$
lorsque $X'$ est un schéma intègre localement de type fini sur $\Spec(k)$,
voir le lemme \ref{chow-lemma-bivariant-zero}. Or $p' : X' \to \Spec(k)$ est alors
plat de dimension relative $\dim(X')$, et donc
$[X'] = (p')^*[\Spec(k)]$. Ainsi, le fait que les classes bivariantes
$c$ et $c_\alpha$ coïncident sur $[\Spec(k)]$ implique qu'elles
coïncident après intersection avec $[X']$, ce qui achève la démonstration.
\end{proof}

\begin{lemma}
\label{chow-lemma-chow-cohomology-towards-point-commutes}
Soit $k$ un corps. Soit $X$ un schéma localement de type fini sur $k$.
Soit $c \in A^p(X \to \Spec(k))$. Soit $Y \to Z$ un morphisme de schémas
localement de type fini sur $k$. Soit $c' \in A^q(Y \to Z)$. Alors
$c \circ c' = c' \circ c$ dans $A^{p + q}(X \times_k Y \to Z)$.
\end{lemma}

\begin{proof}
Dans la démonstration du lemme \ref{chow-lemma-chow-cohomology-towards-point},
nous avons vu que $c$ est donnée par une combinaison d'images
directes propres, de multiplications par des entiers sur les composantes
connexes et d'images inverses plates. Comme $c'$ commute avec chacune de
ces opérations par définition des classes bivariantes, nous concluons.
Certains détails sont omis.
\end{proof}

\begin{remark}
\label{chow-remark-commuting-exterior}
La conséquence des lemmes \ref{chow-lemma-chow-cohomology-towards-point}
et \ref{chow-lemma-chow-cohomology-towards-point-commutes} est la suivante.
Soit $k$ un corps. Soit $X$ un schéma localement de type fini sur $k$.
Soit $\alpha \in \CH_*(X)$. Soit $Y \to Z$ un morphisme de schémas
localement de type fini sur $k$. Soit $c' \in A^q(Y \to Z)$. Alors
$$
\alpha \times (c' \cap \beta) = c' \cap (\alpha \times \beta)
$$
dans $\CH_*(X \times_k Y)$ pour tout $\beta \in \CH_*(Z)$. En effet, cela
s'obtient en prenant pour $c = c_\alpha \in A^*(X \to \Spec(k))$ la classe bivariante
correspondant à $\alpha$ ; voir la démonstration du
lemme \ref{chow-lemma-chow-cohomology-towards-point}.
\end{remark}

\begin{lemma}
\label{chow-lemma-exterior-product-associative}
Le produit extérieur est associatif. Plus précisément, soit $k$ un corps,
soient $X, Y, Z$ des schémas localement de type fini sur $k$, et soient
$\alpha \in \CH_*(X)$, $\beta \in \CH_*(Y)$, $\gamma \in \CH_*(Z)$.
Alors $(\alpha \times \beta) \times \gamma =
\alpha \times (\beta \times \gamma)$ dans $\CH_*(X \times_k Y \times_k Z)$.
\end{lemma}

\begin{proof}
Omis. Indication : associativité du produit fibré des schémas.
\end{proof}







\section{Produits d'intersection}
\label{chow-section-intersection-product}

\noindent
Soit $k$ un corps. Dans cette section, nous travaillons sur $S = \Spec(k)$ avec
$\delta : S \to \mathbf{Z}$ définie en envoyant l'unique point sur $0$ ; voir
l'exemple \ref{chow-example-field}.

\medskip\noindent
Soit $X$ un schéma lisse sur $k$. La classe bivariante $\Delta^!$
de la section \ref{chow-section-gysin-for-diagonal} nous permet de définir une sorte de
produit d'intersection sur les groupes de Chow des schémas localement de type fini sur $X$.
Supposons en effet que $Y \to X$ et $Z \to X$ soient des morphismes de schémas
localement de type fini. Observons alors que
$$
Y \times_X Z =  (Y \times_k Z) \times_{X \times_k X, \Delta} X
$$
Nous pouvons donc considérer la suite d'applications suivante
$$
\CH_n(Y) \otimes_\mathbf{Z} \CH_m(Z)
\xrightarrow{\times}
\CH_{n + m}(Y \times_k Z)
\xrightarrow{\Delta^!}
\CH_{n + m - *}(Y \times_X Z)
$$
Ici, la première flèche est le produit extérieur construit dans la
section \ref{chow-section-exterior-product}, et la seconde est l'homomorphisme de Gysin
pour la diagonale étudié dans la
section \ref{chow-section-gysin-for-diagonal}. Si $X$ est équidimensionnel
de dimension $d$, nous aboutissons à $\CH_{n + m - d}(Y \times_X Z)$ ;
en général, nous pouvons décomposer selon les parties situées au-dessus des sous-schémas
ouverts et fermés de $X$ sur lesquels $X$ est de dimension donnée.
Pour $\alpha \in \CH_*(Y)$ et $\beta \in \CH_*(Z)$, nous noterons
$$
\alpha \cdot \beta = \Delta^!(\alpha \times \beta)
\in \CH_*(Y \times_X Z)
$$
Dans le cas particulier où $X = Y = Z$, nous obtenons une multiplication
$$
\CH_*(X) \times \CH_*(X) \to \CH_*(X),\quad
(\alpha, \beta) \mapsto \alpha \cdot \beta
$$
appelée {\it produit d'intersection}. Nous observons que
ce produit est manifestement symétrique. Son associativité résulte du
lemme suivant.

\begin{lemma}
\label{chow-lemma-associative}
Le produit défini ci-dessus est associatif. Plus précisément, soit $k$ un corps,
soit $X$ lisse sur $k$,
soient $Y, Z, W$ des schémas localement de type fini sur $X$, et soient
$\alpha \in \CH_*(Y)$, $\beta \in \CH_*(Z)$, $\gamma \in \CH_*(W)$.
Alors $(\alpha \cdot \beta) \cdot \gamma =
\alpha \cdot (\beta \cdot \gamma)$ dans $\CH_*(Y \times_X Z \times_X W)$.
\end{lemma}

\begin{proof}
Par le lemme \ref{chow-lemma-exterior-product-associative}, nous avons
$(\alpha \times \beta) \times \gamma =
\alpha \times (\beta \times \gamma)$ dans $\CH_*(Y \times_k Z \times_k W)$.
Considérons les immersions fermées
$$
\Delta_{12} : X \times_k X \longrightarrow X \times_k X \times_k X,
\quad (x, x') \mapsto (x, x, x')
$$
et
$$
\Delta_{23} : X \times_k X \longrightarrow X \times_k X \times_k X,
\quad (x, x') \mapsto (x, x', x')
$$
Notons $\Delta_{12}^!$ et $\Delta_{23}^!$ les classes bivariantes
correspondantes ; observons que $\Delta_{12}^!$ est la restriction
(remarque \ref{chow-remark-restriction-bivariant}) de $\Delta^!$
à $X \times_k X \times_k X$ par l'application $\text{pr}_{12}$ et que
$\Delta_{23}^!$ est la restriction de $\Delta^!$
à $X \times_k X \times_k X$ par l'application $\text{pr}_{23}$.
Ainsi, la restriction de $\Delta_{12}^!$ par $\Delta_{23}$ est manifestement
$\Delta^!$, et la restriction de $\Delta_{23}^!$ par $\Delta_{12}$ est
également $\Delta^!$. Par le lemme \ref{chow-lemma-gysin-commutes}, nous avons donc
$$
\Delta^! \circ \Delta_{12}^! =
\Delta^! \circ \Delta_{23}^! 
$$
Nous pouvons maintenant démontrer le lemme par la suite d'égalités suivante :
\begin{align*}
(\alpha \cdot \beta) \cdot \gamma
& =
\Delta^!(\Delta^!(\alpha \times \beta) \times \gamma) \\
& =
\Delta^!(\Delta_{12}^!((\alpha \times \beta) \times \gamma)) \\
& =
\Delta^!(\Delta_{23}^!((\alpha \times \beta) \times \gamma)) \\
& =
\Delta^!(\Delta_{23}^!(\alpha \times (\beta \times \gamma)) \\
& =
\Delta^!(\alpha \times \Delta^!(\beta \times \gamma)) \\
& =
\alpha \cdot (\beta \cdot \gamma)
\end{align*}
Toutes les égalités découlent clairement de ce qui précède, sauf peut-être
la deuxième et l'avant-dernière. L'égalité
$\Delta_{23}^!(\alpha \times (\beta \times \gamma)) =
\alpha \times \Delta^!(\beta \times \gamma)$ résulte de la
remarque \ref{chow-remark-commuting-exterior}. Il en va de même pour la deuxième
égalité.
\end{proof}

\begin{lemma}
\label{chow-lemma-identify-chow-for-smooth}
Soit $k$ un corps. Soit $X$ un schéma lisse sur $k$, équidimensionnel
de dimension $d$. L'application
$$
A^p(X) \longrightarrow \CH_{d - p}(X),\quad
c \longmapsto c \cap [X]_d
$$
est un isomorphisme. Par cet isomorphisme, la composition des classes
bivariantes devient le produit d'intersection défini ci-dessus.
\end{lemma}

\begin{proof}
Notons $g : X \to \Spec(k)$ le morphisme structural.
L'application est la composée des isomorphismes
$$
A^p(X) \to A^{p - d}(X \to \Spec(k)) \to \CH_{d - p}(X)
$$
Le premier est l'isomorphisme $c \mapsto c \circ g^*$ de la
proposition \ref{chow-proposition-compute-bivariant},
et le second est l'isomorphisme $c \mapsto c \cap [\Spec(k)]$ du
lemme \ref{chow-lemma-chow-cohomology-towards-point}.
La démonstration du lemme \ref{chow-lemma-chow-cohomology-towards-point}
montre que l'inverse de la seconde flèche envoie $\alpha \in \CH_{d - p}(X)$
sur la classe bivariante $c_\alpha$ qui envoie $\beta \in \CH_*(Y)$,
pour $Y$ localement de type fini sur $k$,
sur $\alpha \times \beta$ dans $\CH_*(X \times_k Y)$. La démonstration de la
proposition \ref{chow-proposition-compute-bivariant} montre à son tour que l'inverse
de la première flèche envoie $c_\alpha$ sur la classe bivariante
qui envoie $\beta \in \CH_*(Y)$, pour $Y \to X$ localement de type fini,
sur $\Delta^!(\alpha \times \beta) = \alpha \cdot \beta$.
Le dernier énoncé du lemme en résulte.
\end{proof}

\begin{lemma}
\label{chow-lemma-lci-gysin-product}
Soit $k$ un corps. Soit $f : X \to Y$ un morphisme de schémas lisses
sur $k$. Alors l'homomorphisme de Gysin existe pour $f$ et
$f^!(\alpha \cdot \beta) = f^!\alpha \cdot f^!\beta$.
\end{lemma}

\begin{proof}
Observons que $X \to X \times_k Y$ est une immersion de $X$ dans un schéma
lisse sur $Y$. L'homomorphisme de Gysin existe donc pour $f$
(définition \ref{chow-definition-lci-gysin}).
Pour démontrer la formule, nous pouvons décomposer $X$ et $Y$ en leurs
composantes connexes ; nous pouvons donc supposer $X$ lisse sur $k$
et équidimensionnel de dimension $d$, et $Y$ lisse sur $k$
et équidimensionnel de dimension $e$. Observons que
$f^![Y]_e = [X]_d$ (voir par exemple le lemme \ref{chow-lemma-lci-gysin-easy}).
Écrivons $\alpha = c \cap [Y]_e$ et $\beta = c' \cap [Y]_e$,
de sorte que $\alpha \cdot \beta = c \cap c' \cap [Y]_e$ ;
voir le lemme \ref{chow-lemma-identify-chow-for-smooth}.
Par le lemme \ref{chow-lemma-lci-gysin-commutes}, nous savons que $f^!$
commute à la fois avec $c$ et $c'$. Par conséquent,
\begin{align*}
f^!(\alpha \cdot \beta)
& =
f^!(c \cap c' \cap [Y]_e) \\
& =
c \cap c' \cap f^![Y]_e \\
& =
c \cap c' \cap [X]_d \\
& =
(c \cap [X]_d) \cdot (c' \cap [X]_d) \\
& =
(c \cap f^![Y]_e) \cdot (c' \cap f^![Y]_e) \\
& =
f^!(\alpha) \cdot f^!(\beta)
\end{align*}
comme souhaité, où nous avons utilisé le Lemme \ref{chow-lemma-identify-chow-for-smooth}
pour $X$ également.

\medskip\noindent
On peut donner une autre démonstration en prouvant que
$(f \times f)^!(\alpha \times \beta) = f^!\alpha \times f^!\beta$
et en utilisant le lemme \ref{chow-lemma-lci-gysin-composition}.
\end{proof}

\begin{lemma}
\label{chow-lemma-projection-formula}
Soit $k$ un corps. Soit $f : X \to Y$ un morphisme propre de schémas lisses
sur $k$. Alors l'homomorphisme de Gysin existe pour $f$ et
$f_*(\alpha \cdot f^!\beta) = f_*\alpha \cdot \beta$.
\end{lemma}

\begin{proof}
Observons que $X \to X \times_k Y$ est une immersion de $X$ dans un schéma
lisse sur $Y$. L'homomorphisme de Gysin existe donc pour $f$
(définition \ref{chow-definition-lci-gysin}).
Pour démontrer la formule, nous pouvons décomposer $X$ et $Y$ en leurs
composantes connexes ; nous pouvons donc supposer $X$ lisse sur $k$
et équidimensionnel de dimension $d$, et $Y$ lisse sur $k$
et équidimensionnel de dimension $e$. Observons que
$f^![Y]_e = [X]_d$ (voir par exemple le lemme \ref{chow-lemma-lci-gysin-easy}).
Écrivons $\alpha = c \cap [X]_d$ et $\beta = c' \cap [Y]_e$ ;
voir le lemme \ref{chow-lemma-identify-chow-for-smooth}. Nous avons
\begin{align*}
f_*(\alpha \cdot f^!\beta)
& =
f_*(c \cap f^!(c' \cap [Y]_e)) \\
& =
f_*(c \cap c' \cap f^![Y]_e) \\
& =
f_*(c \cap c' \cap [X]_d) \\
& =
f_*(c' \cap c \cap [X]_d) \\
& =
c' \cap f_*(c \cap [X]_d) \\
& =
\beta \cdot f_*(\alpha)
\end{align*}
La première égalité résulte du lemme \ref{chow-lemma-identify-chow-for-smooth}
appliqué à $X$. Par le lemme \ref{chow-lemma-lci-gysin-commutes}, $f^!$
commute avec $c'$. La commutativité du produit
d'intersection justifie d'inverser l'ordre des intersections de $[X]_d$ avec $c$ et $c'$
(par le lemme). On peut commuter $c'$ avec $f_*$ puisque $c'$
est une classe bivariante. La dernière égalité résulte encore du lemme.
\end{proof}

\begin{lemma}
\label{chow-lemma-intersect-properly}
Soit $k$ un corps. Soit $X$ un schéma intègre lisse sur $k$.
Soient $Y, Z \subset X$ des sous-schémas fermés intègres. Posons
$d = \dim(Y) + \dim(Z) - \dim(X)$. Supposons que
\begin{enumerate}
\item $\dim(Y \cap Z) \leq d$, et
\item $\mathcal{O}_{Y, \xi}$ et $\mathcal{O}_{Z, \xi}$
soient de Cohen-Macaulay pour tout $\xi \in Y \cap Z$ tel que
$\delta(\xi) = d$.
\end{enumerate}
Alors $[Y] \cdot [Z] = [Y \cap Z]_d$ dans $\CH_d(X)$.
\end{lemma}

\begin{proof}
Rappelons que $[Y] \cdot [Z] = \Delta^!([Y \times Z])$, où
$\Delta^! = c(\Delta : X \to X \times X, \mathcal{T}_{X/k})$
est un homomorphisme de Gysin en codimension supérieure
(section \ref{chow-section-gysin-higher-codimension}), avec
$\mathcal{T}_{X/k} = \SheafHom(\Omega_{X/k}, \mathcal{O}_X)$
localement libre de rang $\dim(X)$. Nous avons l'égalité de schémas
$$
Y \cap Z = X \times_{\Delta, (X \times X)} (Y \times Z)
$$
et $\dim(Y \times Z) = \dim(Y) + \dim(Z)$ ; les conditions 
(1), (2) et (3) du lemme \ref{chow-lemma-gysin-easy} sont donc satisfaites.
Enfin, si $\xi \in Y \cap Z$, nous avons un homomorphisme local
plat
$$
\mathcal{O}_{Y, \xi} \longrightarrow
\mathcal{O}_{Y \times Z, \xi}
$$
dont la « fibre » est $\mathcal{O}_{Z, \xi}$. Il s'ensuit que si
$\mathcal{O}_{Y, \xi}$ et $\mathcal{O}_{Z, \xi}$
sont tous deux de Cohen-Macaulay, alors $\mathcal{O}_{Y \times Z, \xi}$ l'est aussi ; voir
Algèbre, lemme \ref{algebra-lemma-CM-goes-up}.
Nous voyons ainsi que toutes les hypothèses du
lemme \ref{chow-lemma-gysin-easy} sont satisfaites, et nous concluons.
\end{proof}

\begin{lemma}
\label{chow-lemma-intersect-regularly-embedded}
Soit $k$ un corps. Soit $X$ un schéma lisse sur $k$. Soit $i : Y \to X$
une immersion fermée régulière. Soit $\alpha \in \CH_*(X)$. Si $Y$ est
équidimensionnel de dimension $e$, alors
$\alpha \cdot [Y]_e = i_*(i^!(\alpha))$ dans $\CH_*(X)$.
\end{lemma}

\begin{proof}
Après avoir décomposé $X$ en composantes connexes, nous pouvons supposer, et supposons, $X$
équidimensionnel de dimension $d$. Écrivons $\alpha = c \cap [X]_n$
avec $x \in A^*(X)$ ; voir le lemme \ref{chow-lemma-identify-chow-for-smooth}. Alors
$$
i_*(i^!(\alpha)) = i_*(i^!(c \cap [X]_n)) =
i_*(c \cap i^![X]_n) = i_*(c \cap [Y]_e) =
c \cap i_*[Y]_e = \alpha \cdot [Y]_e
$$
La première égalité résulte du choix de $c$, puis la deuxième du
lemme \ref{chow-lemma-lci-gysin-commutes}. La troisième vient de ce que
$i^![X]_d = [Y]_e$ dans $\CH_*(Y)$ (lemme \ref{chow-lemma-lci-gysin-easy}).
La quatrième tient à ce que les classes bivariantes commutent à l'image directe propre.
La dernière égalité résulte du lemme \ref{chow-lemma-identify-chow-for-smooth}.
\end{proof}

\begin{lemma}
\label{chow-lemma-intersection-regular-smooth}
Soit $k$ un corps. Soit $X$ un schéma lisse sur $k$, qui est
quasi-compact et dont la diagonale est affine. Alors le produit
d'intersection sur $\CH^*(X)$ construit dans cette section coïncide,
après tensorisation par $\mathbf{Q}$, avec le produit d'intersection
construit dans la section \ref{chow-section-intersection-regular}.
\end{lemma}

\begin{proof}
Soient $\alpha \in \CH^i(X)$ et $\beta \in \CH^j(X)$. Écrivons
$\alpha = ch(\alpha') \cap [X]$ et $\beta = ch(\beta') \cap [X]$
$\alpha', \beta' \in K_0(\textit{Vect}(X)) \otimes \mathbf{Q}$
comme dans la section \ref{chow-section-intersection-regular}.
Posons $c = ch(\alpha')$ et $c' = ch(\beta')$.
Le produit d'intersection de la section \ref{chow-section-intersection-regular}
donne alors $c \cap c' \cap [X]$. C'est égal à $\alpha \cdot \beta$
par le lemme \ref{chow-lemma-identify-chow-for-smooth} (ou plutôt par sa
généralisation selon laquelle $A^i(X) \to \CH^i(X)$, $c \mapsto c \cap [X]$
est un isomorphisme pour tout schéma lisse $X$ sur $k$).
\end{proof}








\section{Produit extérieur sur les anneaux de Dedekind}
\label{chow-section-exterior-product-dim-1}

\noindent
Soit $S$ un schéma localement noethérien qui possède un recouvrement ouvert
par des spectres d'anneaux de Dedekind. Posons $\delta(s) = 0$ pour $s \in S$ fermé
et $\delta(s) = 1$ sinon. Alors $(S, \delta)$ est un cas particulier de notre
situation générale \ref{chow-situation-setup} ; voir
l'exemple \ref{chow-example-domain-dimension-1}.
Observons que $S$ est normal
(Algèbre, lemme \ref{algebra-lemma-characterize-Dedekind})
et est donc une réunion disjointe de schémas intègres normaux
(Propriétés, lemme \ref{properties-lemma-normal-locally-Noetherian}).
Ainsi, tous les arguments ci-dessous se ramènent au cas où $S$ est
irréductible. En revanche, nous permettons que $S$ soit non séparé (ainsi, $S$
pourrait être la droite affine avec $0$ dédoublé, par exemple).

\medskip\noindent
Considérons un carré cartésien
$$
\xymatrix{
X \times_S Y \ar[r] \ar[d] & Y \ar[d] \\
X \ar[r] & S
}
$$
de schémas localement de type fini sur $S$. Nous affirmons qu'il existe une application canonique
$$
\times :
\CH_n(X) \otimes_{\mathbf{Z}} \CH_m(Y)
\longrightarrow
\CH_{n + m - 1}(X \times_S Y)
$$
qui est uniquement déterminée par la règle suivante :
étant donnés des sous-schémas fermés intègres $X' \subset X$
et $Y' \subset Y$ de $\delta$-dimensions $n$ et $m$, nous posons
\begin{enumerate}
\item $[X'] \times [Y'] = [X' \times_S Y']_{n + m - 1}$ si
$X'$ ou $Y'$ domine une composante irréductible de $S$,
\item $[X'] \times [Y'] = 0$ si ni $X'$ ni $Y'$ ne domine de
composante irréductible de $S$.
\end{enumerate}

\begin{lemma}
\label{chow-lemma-exterior-product-well-defined-dim-1}
L'application $\times : \CH_n(X) \otimes_{\mathbf{Z}} \CH_m(Y) \to
\CH_{n + m - 1}(X \times_S Y)$ est bien définie.
\end{lemma}

\begin{proof}
Considérons les cycles de dimensions $n$ et $m$, $\alpha = \sum_{i \in I} n_i[X_i]$
et $\beta = \sum_{j \in J} m_j[Y_j]$, où $X_i \subset X$ et $Y_j \subset Y$
sont des familles localement finies de sous-schémas fermés intègres de
$\delta$-dimensions $n$ et $m$. Soit $K \subset I \times J$ l'ensemble
des couples $(i, j) \in I \times J$ tels que $X_i$ ou $Y_j$ domine
une composante irréductible de $S$.
Alors $\{X_i \times_S Y_j\}_{(i, j) \in K}$ est une famille localement finie
de sous-schémas fermés de $X \times_S Y$, de
$\delta$-dimension $n + m - 1$. Nous pouvons donc bien considérer
$$
\alpha \times \beta =
\sum\nolimits_{(i, j) \in K} n_i m_j [X_i \times_S Y_j]_{n + m - 1}
$$
comme un $(n + m - 1)$-cycle sur $X \times_S Y$. Nous obtenons ainsi une
application additive
$\times : Z_n(X) \otimes_{\mathbf{Z}} Z_m(Y) \to Z_{n + m}(X \times_S Y)$.
Le problème est de montrer que
ce procédé est compatible avec l'équivalence rationnelle.

\medskip\noindent
Soit $i : X' \to X$ le morphisme d'inclusion d'un sous-schéma fermé intègre
de $\delta$-dimension $n$ qui domine une composante irréductible
de $S$. Alors $p' : X' \to S$ est plat de dimension relative $n - 1$ ; voir
Compléments d'algèbre, lemme \ref{more-algebra-lemma-dedekind-torsion-free-flat}.
L'image inverse plate par $p'$ est donc un élément
$(p')^* \in A^{-n + 1}(X' \to S)$ par le
lemme \ref{chow-lemma-flat-pullback-bivariant},
et donc $c' = i_* \circ (p')^* \in A^{-n + 1}(X \to S)$ par le
lemme \ref{chow-lemma-push-proper-bivariant}.
Cela donne des applications
$$
c' \cap - : \CH_m(Y) \longrightarrow \CH_{m + n - 1}(X \times_S Y)
$$
qui envoient $[Y']$ sur $[X' \times_S Y']_{n + m - 1}$ pour tout
sous-schéma fermé intègre $Y' \subset Y$ de $\delta$-dimension $m$.

\medskip\noindent
Soit $i : X' \to X$ le morphisme d'inclusion d'un sous-schéma fermé intègre
de $\delta$-dimension $n$ tel que la composée $X' \to X \to S$ 
se factorise par un point fermé $s \in S$. Comme $s$ est un point fermé
du spectre d'un anneau de Dedekind, nous voyons que $s$ est un diviseur de
Cartier effectif sur $S$ dont le fibré normal est trivial. Notons
$c \in A^1(s \to S)$ l'homomorphisme de Gysin ; voir le
lemme \ref{chow-lemma-gysin-bivariant}. Le morphisme $p' : X' \to s$
est plat de dimension relative $n$. L'image inverse plate par $p'$
est donc un élément $(p')^* \in A^{-n}(X' \to S)$ par le
lemme \ref{chow-lemma-flat-pullback-bivariant}.
Ainsi,
$$
c' = i_* \circ (p')^* \circ c \in A^{-n}(X \to S)
$$
par le lemme \ref{chow-lemma-push-proper-bivariant}. Cela donne des applications
$$
c' \cap - : \CH_m(Y) \longrightarrow \CH_{m + n - 1}(X \times_S Y)
$$
qui, pour tout sous-schéma fermé intègre $Y' \subset Y$
de $\delta$-dimension $m$,
envoient $[Y']$ sur $[X' \times_S Y']_{n + m - 1}$ si $Y'$ domine
une composante irréductible de $S$, et sur $0$ sinon.

\medskip\noindent
Des deux paragraphes précédents, nous concluons que
la construction $([X'], [Y']) \mapsto [X' \times_S Y']_{n + m - 1}$
se factorise par l'équivalence rationnelle en la seconde variable, c'est-à-dire
qu'elle donne une application bien définie
$Z_n(X) \otimes_{\mathbf{Z}} \CH_m(Y) \to \CH_{n + m - 1}(X \times_S Y)$.
Par symétrie, il en va de même pour l'autre variable, et nous concluons.
\end{proof}

\begin{lemma}
\label{chow-lemma-chow-cohomology-towards-base-dim-1}
Soit $(S, \delta)$ comme ci-dessus. Soit $X$ un schéma localement de type fini
sur $S$. Nous avons alors une identification canonique
$$
A^p(X \to S) = \CH_{1 - p}(X)
$$
pour tout $p \in \mathbf{Z}$.
\end{lemma}

\begin{proof}
Considérons l'élément $[S]_1 \in \CH_1(S)$. Nous obtenons une application
$A^p(X \to S) \to \CH_{1 - p}(X)$ en envoyant $c$ sur $c \cap [S]_1$.

\medskip\noindent
Réciproquement, supposons donné $\alpha \in \CH_{1 - p}(X)$.
Nous pouvons alors définir $c_\alpha \in A^p(X \to S)$ de la manière
suivante : étant donnés $X' \to S$ et $\alpha' \in \CH_n(X')$,
nous posons
$$
c_\alpha \cap \alpha' = \alpha \times \alpha'
$$
dans $\CH_{n - p}(X \times_S X')$. Pour montrer qu'il s'agit d'une classe
bivariante, écrivons $\alpha = \sum_{i \in I} n_i[X_i]$ comme dans la
définition \ref{chow-definition-cycles}. En particulier, le morphisme
$$
g : \coprod\nolimits_{i \in I} X_i \longrightarrow X
$$
est propre. Fixons $i \in I$. Si $X_i$ domine une composante irréductible
de $S$, alors le morphisme structural $p_i : X_i \to S$ est plat et nous avons
$\xi_i = p_i^* \in A^p(X_i \to S)$. D'autre part, si $p_i$ se factorise
en $p'_i : X_i \to s_i$ suivi de l'inclusion $s_i \to S$
d'un point fermé, alors nous avons
$\xi_i = (p'_i)^* \circ c_i \in A^p(X_i \to S)$
où $c_i \in A^1(s_i \to S)$ est l'homomorphisme de Gysin et
$(p'_i)^*$ est l'image inverse plate. Observons que
$$
A^p(\coprod\nolimits_{i \in I} X_i \to S) =
\prod\nolimits_{i \in I} A^p(X_i \to S)
$$
Nous avons donc
$$
\xi = \sum n_i \xi_i \in A^p(\coprod\nolimits_{i \in I} X_i \to S)
$$
Enfin, puisque $g$ est propre, nous avons une classe bivariante
$$
g_* \circ \xi \in A^p(X \to S)
$$
par le lemme \ref{chow-lemma-push-proper-bivariant}. Le lecteur vérifie aisément
que $c_\alpha$ est égale à cette classe
(comparer avec la démonstration du
lemme \ref{chow-lemma-exterior-product-well-defined-dim-1})
et est donc elle-même une classe bivariante.

\medskip\noindent
Pour achever la démonstration, nous devons montrer que les deux constructions
sont réciproques l'une de l'autre. Comme $c_\alpha \cap [S]_1 = \alpha$,
c'est clair dans un sens. Pour l'autre sens, soit
$c \in A^p(X \to S)$ et posons $\alpha = c \cap [S]_1$.
Il suffit de démontrer que
$$
c \cap [X'] = c_\alpha \cap [X']
$$
lorsque $X'$ est un schéma intègre localement de type fini sur $S$ ;
voir le lemme \ref{chow-lemma-bivariant-zero}. Or, soit $p' : X' \to S$
est plat de dimension relative $\dim_\delta(X') - 1$, et donc
$[X'] = (p')^*[S]_1$, soit $X' \to S$ se factorise en $X' \to s \to S$,
et donc $[X'] = (p')^*(s \to S)^*[S]_1$. Ainsi, le fait que les
classes bivariantes $c$ et $c_\alpha$ coïncident sur $[S]_1$
implique qu'elles coïncident après intersection avec $[X']$ (puisque les classes bivariantes
commutent aux images inverses plates et aux homomorphismes de Gysin), ce qui achève la démonstration.
\end{proof}

\begin{lemma}
\label{chow-lemma-chow-cohomology-towards-base-dim-1-commutes}
Soit $(S, \delta)$ comme ci-dessus. Soit $X$ un schéma localement de type fini
sur $S$. Soit $c \in A^p(X \to S)$. Soit $Y \to Z$ un morphisme de schémas
localement de type fini sur $S$. Soit $c' \in A^q(Y \to Z)$. Alors
$c \circ c' = c' \circ c$ dans $A^{p + q}(X \times_S Y \to X \times_S Z)$.
\end{lemma}

\begin{proof}
Dans la démonstration du lemme \ref{chow-lemma-chow-cohomology-towards-base-dim-1},
nous avons vu que $c$ est donnée par une combinaison d'images
directes propres, de multiplications par des entiers sur les composantes
connexes, d'images inverses plates et d'homomorphismes de Gysin. Comme $c'$ commute avec chacune de
ces opérations par définition des classes bivariantes, nous concluons.
Certains détails sont omis.
\end{proof}

\begin{remark}
\label{chow-remark-commuting-exterior-dim-1}
La conséquence des lemmes \ref{chow-lemma-chow-cohomology-towards-base-dim-1}
et \ref{chow-lemma-chow-cohomology-towards-base-dim-1-commutes} est la suivante.
Soit $(S, \delta)$ comme ci-dessus. Soit $X$ un schéma localement de type fini
sur $S$.
Soit $\alpha \in \CH_*(X)$. Soit $Y \to Z$ un morphisme de schémas
localement de type fini sur $S$. Soit $c' \in A^q(Y \to Z)$. Alors
$$
\alpha \times (c' \cap \beta) = c' \cap (\alpha \times \beta)
$$
dans $\CH_*(X \times_S Y)$ pour tout $\beta \in \CH_*(Z)$. En effet, cela
s'obtient en prenant pour $c = c_\alpha \in A^*(X \to S)$ la classe bivariante
correspondant à $\alpha$ ; voir la démonstration du
lemme \ref{chow-lemma-chow-cohomology-towards-base-dim-1}.
\end{remark}

\begin{lemma}
\label{chow-lemma-exterior-product-associative-dim-1}
Le produit extérieur est associatif. Plus précisément, soit $(S, \delta)$
comme ci-dessus, soient $X, Y, Z$ des schémas localement de type fini sur $S$, et soient
$\alpha \in \CH_*(X)$, $\beta \in \CH_*(Y)$, $\gamma \in \CH_*(Z)$.
Alors $(\alpha \times \beta) \times \gamma =
\alpha \times (\beta \times \gamma)$ dans $\CH_*(X \times_S Y \times_S Z)$.
\end{lemma}

\begin{proof}
Omis. Indication : associativité du produit fibré des schémas.
\end{proof}















\section{Produits d'intersection sur les anneaux de Dedekind}
\label{chow-section-intersection-product-dim-1}

\noindent
Soit $S$ un schéma localement noethérien qui admet un recouvrement ouvert
par des spectres d'anneaux de Dedekind. Posons $\delta(s) = 0$ pour $s \in S$ fermé
et $\delta(s) = 1$ sinon. Alors $(S, \delta)$ est un cas particulier de notre
Situation générale \ref{chow-situation-setup} ; voir
l'Exemple \ref{chow-example-domain-dimension-1} et la discussion de la
section \ref{chow-section-exterior-product-dim-1}.

\medskip\noindent
Soit $X$ un schéma lisse sur $S$. La classe bivariante $\Delta^!$
de la section \ref{chow-section-gysin-for-diagonal} nous permet de définir une sorte de
produit d'intersection sur les groupes de Chow des schémas localement de type fini sur $X$.
Plus précisément, supposons que $Y \to X$ et $Z \to X$ soient des morphismes de schémas
qui sont localement de type fini. Observons alors que
$$
Y \times_X Z =  (Y \times_S Z) \times_{X \times_S X, \Delta} X
$$
Nous pouvons donc considérer la suite d'applications suivante
$$
\CH_n(Y) \otimes_\mathbf{Z} \CH_m(Y)
\xrightarrow{\times}
\CH_{n + m - 1}(Y \times_S Z)
\xrightarrow{\Delta^!}
\CH_{n + m - *}(Y \times_X Z)
$$
Ici, la première flèche est le produit extérieur construit à la
section \ref{chow-section-exterior-product-dim-1} et la seconde flèche
est l'homomorphisme de Gysin pour la diagonale étudié à la
section \ref{chow-section-gysin-for-diagonal}. Si $X$ est équidimensionnel
de dimension $d$, alors $X \to S$ est lisse de dimension relative $d - 1$
et nous aboutissons donc à $\CH_{n + m - d}(Y \times_X Z)$.
En général, nous pouvons décomposer selon les parties situées au-dessus des sous-schémas
ouverts et fermés de $X$ où $X$ a une dimension donnée.
Étant donnés $\alpha \in \CH_*(Y)$ et $\beta \in \CH_*(Z)$, nous noterons
$$
\alpha \cdot \beta = \Delta^!(\alpha \times \beta)
\in \CH_*(Y \times_X Z)
$$
Dans le cas particulier où $X = Y = Z$, nous obtenons une multiplication
$$
\CH_*(X) \times \CH_*(X) \to \CH_*(X),\quad
(\alpha, \beta) \mapsto \alpha \cdot \beta
$$
appelée le {\it produit d'intersection}. Nous observons que
ce produit est clairement symétrique. L'associativité résulte du
lemme suivant.

\begin{lemma}
\label{chow-lemma-associative-dim-1}
Le produit défini ci-dessus est associatif. Plus précisément, avec
$(S, \delta)$ comme ci-dessus, soit $X$ lisse sur $S$,
soient $Y, Z, W$ des schémas localement de type fini sur $X$, et soient
$\alpha \in \CH_*(Y)$, $\beta \in \CH_*(Z)$, $\gamma \in \CH_*(W)$.
Alors $(\alpha \cdot \beta) \cdot \gamma =
\alpha \cdot (\beta \cdot \gamma)$ dans $\CH_*(Y \times_X Z \times_X W)$.
\end{lemma}

\begin{proof}
Par le Lemme \ref{chow-lemma-exterior-product-associative-dim-1}, nous avons
$(\alpha \times \beta) \times \gamma =
\alpha \times (\beta \times \gamma)$ dans $\CH_*(Y \times_S Z \times_S W)$.
Considérons les immersions fermées
$$
\Delta_{12} : X \times_S X \longrightarrow X \times_S X \times_S X,
\quad (x, x') \mapsto (x, x, x')
$$
et
$$
\Delta_{23} : X \times_S X \longrightarrow X \times_S X \times_S X,
\quad (x, x') \mapsto (x, x', x')
$$
Notons $\Delta_{12}^!$ et $\Delta_{23}^!$ les classes bivariantes
correspondantes ; observons que $\Delta_{12}^!$ est la restriction
(Remarque \ref{chow-remark-restriction-bivariant}) de $\Delta^!$
à $X \times_S X \times_S X$ par l'application $\text{pr}_{12}$ et que
$\Delta_{23}^!$ est la restriction de $\Delta^!$
à $X \times_S X \times_S X$ par l'application $\text{pr}_{23}$.
Ainsi, la restriction de $\Delta_{12}^!$ par $\Delta_{23}$
est clairement $\Delta^!$, et la restriction de $\Delta_{23}^!$ par $\Delta_{12}$ est
également $\Delta^!$. Ainsi, par le Lemme \ref{chow-lemma-gysin-commutes}, nous avons
$$
\Delta^! \circ \Delta_{12}^! =
\Delta^! \circ \Delta_{23}^! 
$$
Nous pouvons maintenant démontrer le lemme par la suite d'égalités suivante :
\begin{align*}
(\alpha \cdot \beta) \cdot \gamma
& =
\Delta^!(\Delta^!(\alpha \times \beta) \times \gamma) \\
& =
\Delta^!(\Delta_{12}^!((\alpha \times \beta) \times \gamma)) \\
& =
\Delta^!(\Delta_{23}^!((\alpha \times \beta) \times \gamma)) \\
& =
\Delta^!(\Delta_{23}^!(\alpha \times (\beta \times \gamma)) \\
& =
\Delta^!(\alpha \times \Delta^!(\beta \times \gamma)) \\
& =
\alpha \cdot (\beta \cdot \gamma)
\end{align*}
Toutes les égalités ressortent clairement de ce qui précède, sauf peut-être
la deuxième et l'avant-dernière. L'égalité
$\Delta_{23}^!(\alpha \times (\beta \times \gamma)) =
\alpha \times \Delta^!(\beta \times \gamma)$ résulte de la
Remarque \ref{chow-remark-commuting-exterior}. Il en va de même pour la deuxième
égalité.
\end{proof}

\begin{lemma}
\label{chow-lemma-identify-chow-for-smooth-dim-1}
Soit $(S, \delta)$ comme ci-dessus. Soit $X$ un schéma lisse sur $S$,
équidimensionnel de dimension $d$. L'application
$$
A^p(X) \longrightarrow \CH_{d - p}(X),\quad
c \longmapsto c \cap [X]_d
$$
est un isomorphisme. Par cet isomorphisme, la composition des classes
bivariantes devient le produit d'intersection défini ci-dessus.
\end{lemma}

\begin{proof}
Notons $g : X \to S$ le morphisme structural.
L'application est la composée des isomorphismes
$$
A^p(X) \to A^{p - d + 1}(X \to S) \to \CH_{d - p}(X)
$$
Le premier est l'isomorphisme $c \mapsto c \circ g^*$ de la
Proposition \ref{chow-proposition-compute-bivariant},
et le second est l'isomorphisme $c \mapsto c \cap [S]_1$ du
Lemme \ref{chow-lemma-chow-cohomology-towards-base-dim-1}.
La démonstration du Lemme \ref{chow-lemma-chow-cohomology-towards-base-dim-1}
montre que l'inverse de la seconde flèche envoie $\alpha \in \CH_{d - p}(X)$
sur la classe bivariante $c_\alpha$ qui envoie $\beta \in \CH_*(Y)$,
pour $Y$ localement de type fini sur $k$,
sur $\alpha \times \beta$ dans $\CH_*(X \times_k Y)$. La démonstration de la
Proposition \ref{chow-proposition-compute-bivariant} montre à son tour que l'inverse
de la première flèche envoie $c_\alpha$ sur la classe bivariante
qui envoie $\beta \in \CH_*(Y)$, pour $Y \to X$ localement de type fini,
sur $\Delta^!(\alpha \times \beta) = \alpha \cdot \beta$.
Le dernier résultat du lemme en découle.
\end{proof}










\section{Classes de Todd}
\label{chow-section-todd-classes}

\noindent
Une dernière classe associée à un fibré vectoriel $\mathcal{E}$
de rang $r$ est sa {\it classe de Todd} $Todd(\mathcal{E})$.
En termes des racines de Chern $x_1, \ldots, x_r$, elle est
définie par
$$
Todd(\mathcal{E})
=
\prod\nolimits_{i = 1}^r
\frac{x_i}{1 - e^{-x_i}}
$$
En termes des classes de Chern $c_i = c_i(\mathcal{E})$,
nous avons
$$
Todd(\mathcal{E})
=
1
+
\frac{1}{2}c_1
+
\frac{1}{12}(c_1^2 + c_2)
+
\frac{1}{24}c_1c_2
+
\frac{1}{720}(-c_1^4 + 4c_1^2c_2 + 3c_2^2 + c_1c_3 - c_4)
+
\ldots
$$
Nous avons fait les remarques appropriées sur les dénominateurs
dans la section précédente. Pour toute suite exacte
donnée
$$
0
\to
{\mathcal E}_1
\to
{\mathcal E}
\to
{\mathcal E}_2
\to
0
$$
nous avons
$$
Todd({\mathcal E}) = Todd({\mathcal E}_1) Todd({\mathcal E}_2).
$$







\section{Grothendieck--Riemann--Roch}
\label{chow-section-grr}

\noindent
Soit $(S, \delta)$ comme dans la
Situation \ref{chow-situation-setup}.
Soient $X, Y$ localement de type fini sur $S$.
Soit $\mathcal{E}$ un faisceau localement libre de type fini
${\mathcal E}$ sur $X$, de rang $r$.
Soit $f : X \to Y$ un morphisme propre et lisse.
Supposons que les $R^if_*\mathcal{E}$ soient des faisceaux localement libres
sur $Y$, de rang fini.
Le théorème de Grothendieck--Riemann--Roch affirme dans ce
cas que
$$
f_*(Todd(T_{X/Y}) ch(\mathcal{E}))
=
\sum (-1)^i ch(R^if_*\mathcal{E})
$$
Ici,
$$
T_{X/Y} = \SheafHom_{\mathcal{O}_X}(\Omega_{X/Y}, \mathcal{O}_X)
$$
est le fibré tangent relatif de $X$ sur $Y$. Si $Y = \Spec(k)$,
où $k$ est un corps, nous pouvons reformuler ceci sous la forme
$$
\chi(X, \mathcal{E}) = \deg(Todd(T_{X/k}) ch(\mathcal{E}))
$$
Le théorème est plus général et devient plus facile à démontrer
lorsqu'il est formulé avec la bonne généralité. Nous y reviendrons
ailleurs (insérer ici une référence future).







\section{Changement de schéma de base}
\label{chow-section-change-base}

\noindent
Dans cette section, nous expliquons comment comparer les théories pour différents
schémas de base.

\begin{situation}
\label{chow-situation-setup-base-change}
Ici, $(S, \delta)$ et $(S', \delta')$ sont comme dans la
Situation \ref{chow-situation-setup}. En outre, $g : S' \to S$
est un morphisme plat de schémas et $c \in \mathbf{Z}$
est un entier tel que : pour tout $s \in S$ et tout $s' \in S'$
qui est un point générique d'une composante irréductible de $g^{-1}(\{s\})$, nous avons
$\delta'(s') = \delta(s) + c$.
\end{situation}

\noindent
Nous verrons que, pour un schéma $X$ localement de type fini sur $S$,
il existe une application bien définie $\CH_k(X) \to \CH_{k + c}(X \times_S S')$
de groupes de Chow qui commute (dans l'ensemble) avec les opérations
que nous avons définies dans ce chapitre.

\begin{lemma}
\label{chow-lemma-dimension-base-change}
Dans la Situation \ref{chow-situation-setup-base-change}, soit $X \to S$ localement
de type fini. Notons $X' \to S'$ le changement de base par $S' \to S$.
Si $X$ est intègre avec $\dim_\delta(X) = k$, alors
toute composante irréductible $Z'$ de $X'$ vérifie $\dim_{\delta'}(Z') = k + c$.
\end{lemma}

\begin{proof}
La projection $X' \to X$ est plate comme changement de base du morphisme plat
$S' \to S$ (Morphismes, Lemme \ref{morphisms-lemma-base-change-flat}).
Ainsi, tout point générique $x'$ d'une composante irréductible
de $X'$ s'envoie sur le point générique $x \in X$ (car les générisations
se relèvent le long de $X' \to X$ par
Morphismes, Lemme \ref{morphisms-lemma-generalizations-lift-flat}).
Soit $s \in S$ l'image de $x$.
Rappelons que le schéma $S'_s = S' \times_S s$
a le même espace topologique sous-jacent que $g^{-1}(\{s\})$
(Schémas, Lemme \ref{schemes-lemma-fibre-topological}).
Nous pouvons voir $x'$ comme un point du schéma $S'_s \times_s x$, lequel
est muni d'un monomorphisme $S'_s \times_s x \to S' \times_S X$.
Bien entendu, $x'$ est aussi un point générique d'une composante irréductible
de $S'_s \times_s x$.
En utilisant la platitude de $\Spec(\kappa(x)) \to \Spec(\kappa(s)) = s$
et en raisonnant comme ci-dessus, nous voyons que $x'$ s'envoie sur un point générique $s'$
d'une composante irréductible de $g^{-1}(\{s\})$. Ainsi,
$\delta'(s') = \delta(s) + c$ par hypothèse.
Nous avons $\dim_x(X_s) = \dim_{x'}(X_{s'})$ par
Morphismes, Lemme \ref{morphisms-lemma-dimension-fibre-after-base-change}.
Puisque $x$ est un point générique d'une composante irréductible $X_s$
(c'est un schéma irréductible, mais nous n'en avons pas besoin) et que
$x'$ est un point générique d'une composante irréductible de $X'_{s'}$, nous concluons
que $\text{trdeg}_{\kappa(s)}(\kappa(x)) = 
\text{trdeg}_{\kappa(s')}(\kappa(x'))$
par Morphismes, Lemme \ref{morphisms-lemma-dimension-fibre-at-a-point}.
Alors
$$
\delta_{X'/S'}(x') = \delta(s') + \text{trdeg}_{\kappa(s')}(\kappa(x')) =
\delta(s) + c + \text{trdeg}_{\kappa(s)}(\kappa(x)) = \delta_{X/S}(x) + c
$$
Ceci prouve ce que nous voulions par la Définition \ref{chow-definition-delta-dimension}.
\end{proof}

\noindent
Dans la Situation \ref{chow-situation-setup-base-change}, soit $X \to S$ localement
de type fini. Notons $X' \to S'$ le changement de base par $g : S' \to S$.
Il existe un unique homomorphisme
$$
g^* : Z_k(X) \longrightarrow Z_{k + c}(X')
$$
qui, étant donné un sous-schéma fermé intègre $Z \subset X$ de
$\delta$-dimension $k$, envoie $[Z]$ sur $[Z \times_S S']_{k + c}$.
Ceci a un sens par le Lemme \ref{chow-lemma-dimension-base-change}.

\begin{lemma}
\label{chow-lemma-pullback-coherent-base-change}
Dans la Situation \ref{chow-situation-setup-base-change}, soit $X \to S$
localement de type fini et soit $X' \to S$ le changement de base par $S' \to S$.
\begin{enumerate}
\item Soit $Z \subset X$ un sous-schéma fermé tel que
$\dim_\delta(Z) \leq k$ et de changement de base $Z' \subset X'$. Alors nous avons
$\dim_{\delta'}(Z')) \leq k + c$
et $[Z']_{k + c} = g^*[Z]_k$ dans $Z_{k + c}(X')$.
\item Soit $\mathcal{F}$ un faisceau cohérent sur $X$ tel que
$\dim_\delta(\text{Supp}(\mathcal{F})) \leq k$, et soit
$\mathcal{F}'$ son changement de base sur $X'$.
Alors nous avons $\dim_\delta(\text{Supp}(\mathcal{F}')) \leq k + c$
et $g^*[\mathcal{F}]_k = [\mathcal{F}']_{k + c}$
dans $Z_{k + c}(X')$.
\end{enumerate}
\end{lemma}

\begin{proof}
La démonstration est exactement la même que celle du
Lemme \ref{chow-lemma-pullback-coherent},
et nous suggérons au lecteur de la passer.

\medskip\noindent
Les assertions sur les dimensions résultent du
Lemme \ref{chow-lemma-dimension-base-change}.
La partie (1) résulte de la partie (2) par le Lemme \ref{chow-lemma-cycle-closed-coherent}
et par le fait que le changement de base du module cohérent $\mathcal{O}_Z$
est $\mathcal{O}_{Z'}$.

\medskip\noindent
Démonstration de (2). Puisque $X$, $X'$ sont localement noethériens, nous pouvons appliquer
Cohomologie des schémas, Lemme \ref{coherent-lemma-coherent-Noetherian}, pour voir
que $\mathcal{F}$ est de type fini, donc $\mathcal{F}'$ est
de type fini (Modules, Lemme \ref{modules-lemma-pullback-finite-type}),
et donc $\mathcal{F}'$ est cohérent
(de nouveau Cohomologie des schémas, Lemme \ref{coherent-lemma-coherent-Noetherian}).
Ainsi, le lemme a un sens. Soit $W \subset X$ un sous-schéma fermé intègre
de $\delta$-dimension $k$, et soit $W' \subset X'$ un sous-schéma
fermé intègre de $\delta'$-dimension $k + c$ s'envoyant dans $W$
par $X' \to X$. Nous devons montrer que le coefficient $n$ de
$[W']$ dans $g^*[\mathcal{F}]_k$ coïncide avec le coefficient
$m$ de $[W']$ dans $[\mathcal{F}']_{k + c}$. Soient $\xi \in W$ et
$\xi' \in W'$ les points génériques. Posons
$A = \mathcal{O}_{X, \xi}$, $B = \mathcal{O}_{X', \xi'}$
et $M = \mathcal{F}_\xi$ comme $A$-module. (Notons que
$M$ est de longueur finie par nos hypothèses de dimension, mais qu'en fait nous
n'avons pas besoin de le vérifier. Voir le
Lemme \ref{chow-lemma-length-finite}.)
Nous avons $\mathcal{F}'_{\xi'} = B \otimes_A M$.
Ainsi, nous voyons que
$$
n = \text{longueur}_B(B \otimes_A M)
\quad
\text{et}
\quad
m = \text{longueur}_A(M) \text{longueur}_B(B/\mathfrak m_AB)
$$
L'égalité résulte donc de
Algèbre, Lemme \ref{algebra-lemma-pullback-module}.
\end{proof}

\begin{lemma}
\label{chow-lemma-pullback-base-change}
Dans la Situation \ref{chow-situation-setup-base-change}, soit $X \to S$ localement
de type fini et soit $X' \to S'$ le changement de base par $S' \to S$.
L'application $g^* : Z_k(X) \to Z_{k + c}(X')$ ci-dessus passe au quotient par l'équivalence
rationnelle et donne une application
$$
g^* : \CH_k(X) \longrightarrow \CH_{k + c}(X')
$$
de groupes de Chow.
\end{lemma}

\begin{proof}
Supposons que $\alpha \in Z_k(X)$ soit un $k$-cycle rationnellement équivalent
à zéro. Par le Lemme \ref{chow-lemma-rational-equivalence-family},
il existe une famille localement finie de sous-schémas fermés intègres
$W_i \subset X \times \mathbf{P}^1$ de $\delta$-dimension $k$
qui ne sont pas contenus dans les diviseurs
$(X \times \mathbf{P}^1)_0$ ou $(X \times \mathbf{P}^1)_\infty$
de $X \times \mathbf{P}^1$, telle que
$\alpha = \sum ([(W_i)_0]_k - [(W_i)_\infty]_k)$.
Il suffit donc de prouver que, pour $W \subset X \times \mathbf{P}^1$
fermé intègre de $\delta$-dimension $k$ et non contenu dans les diviseurs
$(X \times \mathbf{P}^1)_0$ ou $(X \times \mathbf{P}^1)_\infty$
de $X \times \mathbf{P}^1$, nous avons
\begin{enumerate}
\item le changement de base $W' \subset X' \times \mathbf{P}^1$ satisfait aux
hypothèses du Lemme \ref{chow-lemma-closed-subscheme-cross-p1}, avec
$k$ remplacé par $k + c$, et
\item $g^*[W_0]_k = [(W')_0]_{k + c}$ et
$g^*[W_\infty]_k = [(W')_\infty]_{k + c}$.
\end{enumerate}
La partie (2) résulte immédiatement du
Lemme \ref{chow-lemma-pullback-coherent-base-change} et du fait que
$(W')_0$ est le changement de base de $W_0$ (par associativité des produits fibrés).
Pour la partie (1), l'assertion sur les dimensions résulte d'abord
du Lemme \ref{chow-lemma-dimension-base-change}.
Soit ensuite $w' \in (W')_0$, d'image $w \in W_0$
et $z \in \mathbf{P}^1_S$. Notons $t \in \mathcal{O}_{\mathbf{P}^1_S, z}$
l'équation usuelle de $0 : S \to \mathbf{P}^1_S$.
Puisque $\mathcal{O}_{W, w} \to \mathcal{O}_{W', w'}$ est plat
et puisque $t$ est un non-diviseur de zéro sur $\mathcal{O}_{W, w}$
(car $W$ est intègre et $W \not = W_0$), nous voyons que
$t$ est aussi un non-diviseur de zéro dans $\mathcal{O}_{W', w'}$. Ainsi,
$W'$ n'a aucun point associé situé sur $W'_0$.
\end{proof}

\begin{lemma}
\label{chow-lemma-pullback-base-change-pullback}
Dans la Situation \ref{chow-situation-setup-base-change}, soient $Y \to X \to S$ localement
de type fini et soit $Y' \to X' \to S'$ le changement de base par $S' \to S$.
Supposons $f : Y \to X$ plat de dimension relative $r$. Alors $f' : Y' \to X'$
est plat de dimension relative $r$ et les diagrammes
$$
\vcenter{
\xymatrix{
Z_{k + r}(Y) \ar[r]_{g^*} & Z_{k + c + r}(Y') \\
Z_k(X) \ar[r]^{g^*} \ar[u]^{f^*} & Z_{k + c}(X') \ar[u]_{(f')^*}
}
}
\quad\text{et}\quad
\vcenter{
\xymatrix{
\CH_{k + r}(Y) \ar[r]_{g^*} & \CH_{k + c + r}(Y') \\
\CH_k(X) \ar[r]^{g^*} \ar[u]^{f^*} & \CH_{k + c}(X') \ar[u]_{(f')^*}
}
}
$$
de groupes de cycles et de Chow commutent.
\end{lemma}

\begin{proof}
Il suffit de montrer que le premier diagramme commute. Pour le voir, soit
$Z \subset X$ un sous-schéma fermé intègre de $\delta$-dimension $k$
et notons $Z' \subset X'$ son changement de base. Par construction,
nous avons $g^*[Z] = [Z']_{k + c}$. Par le Lemme \ref{chow-lemma-pullback-coherent},
nous avons $(f')^*g^*[Z] = [Z' \times_{X'} Y']_{k + c + r}$.
Réciproquement, nous avons $f^*[Z] = [Z \times_X Y]_{k + r}$ par la
Définition \ref{chow-definition-flat-pullback}. Par le
Lemme \ref{chow-lemma-pullback-coherent-base-change},
nous avons $g^*f^*[Z] = [(Z \times_X Y)']_{k + r + c}$.
Puisque $(Z \times_X Y)' = Z' \times_{X'} Y'$ par
associativité du produit fibré, nous concluons.
\end{proof}

\begin{lemma}
\label{chow-lemma-pullback-base-change-pushforward}
Dans la Situation \ref{chow-situation-setup-base-change}, soient $Y \to X \to S$ localement
de type fini et soit $Y' \to X' \to S'$ le changement de base par $S' \to S$.
Supposons $f : Y \to X$ propre. Alors $f' : Y' \to X'$ est propre et les diagrammes
$$
\vcenter{
\xymatrix{
Z_k(Y) \ar[r]_{g^*} \ar[d]_{f_*} & Z_{k + c}(Y') \ar[d]^{f'_*} \\
Z_k(X) \ar[r]^{g^*} & Z_{k + c}(X')
}
}
\quad\text{et}\quad
\vcenter{
\xymatrix{
\CH_k(Y) \ar[r]_{g^*} \ar[d]_{f_*} & \CH_{k + c}(Y') \ar[d]^{f'_*} \\
\CH_k(X) \ar[r]^{g^*} & \CH_{k + c}(X')
}
}
$$
de groupes de cycles et de Chow commutent.
\end{lemma}

\begin{proof}
Il suffit de montrer que le premier diagramme commute. Pour le voir, soit
$Z \subset Y$ un sous-schéma fermé intègre de $\delta$-dimension $k$
et notons $Z' \subset X'$ son changement de base. Par construction,
nous avons $g^*[Z] = [Z']_{k + c}$. Par le Lemme \ref{chow-lemma-cycle-push-sheaf},
nous avons $(f')_*g^*[Z] = [f'_*\mathcal{O}_{Z'}]_{k + c}$.
Par le même lemme, nous avons $f_*[Z] = [f_*\mathcal{O}_Z]_k$. Par le
Lemme \ref{chow-lemma-pullback-coherent-base-change},
nous avons $g^*f_*[Z] = [(X' \to X)^*f_*\mathcal{O}_Z]_{k + r}$.
Il suffit donc de montrer que
$$
(X' \to X)^*f_*\mathcal{O}_Z \cong f'_*\mathcal{O}_{Z'}
$$
comme modules cohérents sur $X'$. Puisque $X' \to X$ est plat et que
$\mathcal{O}_{Z'} = (Y' \to Y)^*\mathcal{O}_Z$, ceci
résulte du changement de base plat, voir
Cohomologie des schémas, Lemme \ref{coherent-lemma-flat-base-change-cohomology}.
\end{proof}

\begin{lemma}
\label{chow-lemma-pullback-base-change-c1}
Dans la Situation \ref{chow-situation-setup-base-change}, soit $X \to S$ localement
de type fini et soit $X' \to S'$ le changement de base par $S' \to S$.
Soit $\mathcal{L}$ un $\mathcal{O}_X$-module inversible de
changement de base $\mathcal{L}'$ sur $X'$. Alors le
diagramme
$$
\xymatrix{
\CH_k(X) \ar[r]_{g^*} \ar[d]_{c_1(\mathcal{L}) \cap -} &
\CH_{k + c}(X') \ar[d]^{c_1(\mathcal{L}') \cap -} \\
\CH_{k - 1}(X) \ar[r]^{g^*} & \CH_{k + c - 1}(X')
}
$$
de groupes de Chow commute.
\end{lemma}

\begin{proof}
Soit $p : L \to X$ le fibré en droites associé à $\mathcal{L}$,
de section nulle $o : X \to L$. Pour $\alpha \in CH_k(X)$, nous
savons que $\beta = c_1(\mathcal{L}) \cap \alpha$
est l'unique élément de $\CH_{k - 1}(X)$ tel que
$o_*\alpha = - p^*\beta$, voir les Lemmes \ref{chow-lemma-linebundle} et
\ref{chow-lemma-linebundle-formulae}.
La même caractérisation vaut après image inverse. Le lemme résulte donc des
Lemmes \ref{chow-lemma-pullback-base-change-pullback} et
\ref{chow-lemma-pullback-base-change-pushforward}.
\end{proof}

\begin{lemma}
\label{chow-lemma-pullback-base-change-chern-classes}
Dans la Situation \ref{chow-situation-setup-base-change}, soit $X \to S$ localement
de type fini et soit $X' \to S'$ le changement de base par $S' \to S$.
Soit $\mathcal{E}$ un $\mathcal{O}_X$-module localement libre de type fini de
rang $r$, de changement de base $\mathcal{E}'$ sur $X'$. Alors le
diagramme
$$
\xymatrix{
\CH_k(X) \ar[r]_{g^*} \ar[d]_{c_i(\mathcal{E}) \cap -} &
\CH_{k + c}(X') \ar[d]^{c_i(\mathcal{E}') \cap -} \\
\CH_{k - i}(X) \ar[r]^{g^*} & \CH_{k + c - i}(X')
}
$$
de groupes de Chow commute pour tout $i$.
\end{lemma}

\begin{proof}
Posons $P = \mathbf{P}(\mathcal{E})$. Le changement de base $P'$ de $P$
est égal à $\mathbf{P}(\mathcal{E}')$. Puisque nous savons déjà que
l'image inverse plate et le cup-produit avec $c_1$ d'un module inversible
commutent au changement de base (Lemmes \ref{chow-lemma-pullback-base-change-pullback} et
\ref{chow-lemma-pullback-base-change-c1})
le lemme résulte de la caractérisation du cap-produit
avec $c_i(\mathcal{E})$ donnée dans le Lemme \ref{chow-lemma-determine-intersections}.
\end{proof}

\begin{lemma}
\label{chow-lemma-compose-base-change}
Soient $(S, \delta)$, $(S', \delta')$, $(S'', \delta'')$ comme dans la
Situation \ref{chow-situation-setup}. Soient $g : S' \to S$ et $g' : S'' \to S'$
des morphismes plats de schémas, et soient $c, c' \in \mathbf{Z}$
des entiers tels que $S, \delta, S', \delta', g, c$ et
$S', \delta', S'', g', c'$ soient comme dans la
Situation \ref{chow-situation-setup-base-change}.
Soit $X \to S$ localement de type fini et notons $X' \to S'$
et $X'' \to S''$ les changements de base par $S' \to S$ et $S'' \to S$.
Alors
\begin{enumerate}
\item $S, \delta, S'', \delta'', g \circ g', c + c'$ est comme dans la
Situation \ref{chow-situation-setup-base-change},
\item les applications $g^* : Z_k(X) \to Z_{k + c}(X')$ et
$(g')^* : Z_{k + c}(X') \to Z_{k + c + c'}(X'')$ se
composent pour donner l'application $(g \circ g')^* : Z_k(X) \to Z_{k + c + c'}(X'')$, et
\item les applications $g^* : \CH_k(X) \to \CH_{k + c}(X')$ et
$(g')^* : \CH_{k + c}(X') \to \CH_{k + c + c'}(X'')$ du
Lemme \ref{chow-lemma-pullback-base-change}
se composent pour donner l'application $(g \circ g')^* : \CH_k(X) \to \CH_{k + c + c'}(X'')$
du Lemme \ref{chow-lemma-pullback-base-change}.
\end{enumerate}
\end{lemma}

\begin{proof}
Soit $s \in S$, et soit $s'' \in S''$ un point générique d'une composante
irréductible de $(g \circ g')^{-1}(\{s\})$. Posons $s' = g'(s'')$.
Clairement, $s''$ est un point générique d'une composante irréductible de
$(g')^{-1}(\{s'\})$. De plus, puisque $g'$ est plat et que les générisations
se relèvent donc le long de $g'$ (Morphismes, Lemme \ref{morphisms-lemma-base-change-flat}),
nous voyons que $s'$ est aussi un point générique d'une composante irréductible
de $g^{-1}(\{s\})$. Ainsi, par hypothèse, $\delta'(s') = \delta(s) + c$
et $\delta''(s'') = \delta'(s') + c'$. Nous concluons que
$\delta''(s'') = \delta(s) + c + c'$ et que la première partie de cette
assertion est vraie.

\medskip\noindent
Pour la deuxième partie, soit $Z \subset X$ un sous-schéma fermé intègre
de $\delta$-dimension $k$. Notons $Z' \subset X'$ et $Z'' \subset X''$
les changements de base. Par définition, nous avons $g^*[Z] = [Z']_{k + c}$.
Par le Lemme \ref{chow-lemma-pullback-coherent-base-change}, nous avons
$(g')^*[Z']_{k + c} = [Z'']_{k + c + c'}$. Ceci prouve la dernière assertion.
\end{proof}

\begin{lemma}
\label{chow-lemma-chow-limit}
Dans la Situation \ref{chow-situation-setup-base-change}, supposons $c = 0$
et supposons que $S' = \lim_{i \in I} S_i$ soit une limite projective filtrante
de schémas $S_i$ affines sur $S$ tels que
\begin{enumerate}
\item avec $\delta_i$ égal à $S_i \to S \xrightarrow{\delta} \mathbf{Z}$,
le couple $(S_i, \delta_i)$ est comme dans la Situation \ref{chow-situation-setup},
\item $S_i, \delta_i, S, \delta, S \to S_i, c = 0$ est comme dans la
Situation \ref{chow-situation-setup-base-change},
\item $S_i, \delta_i, S_{i'}, \delta_{i'}, S_i \to S_{i'}, c = 0$ 
pour $i \geq i'$ est comme dans la Situation \ref{chow-situation-setup-base-change}.
\end{enumerate}
Alors, pour un schéma quasi-compact $X$ de type fini sur $S$,
de changements de base $X'$ et $X_i$ par $S' \to S$ et $S_i \to S$, nous avons
$Z_k(X') = \colim Z_k(X_i)$ et $\CH_k(X') = \colim \CH_k(X_i)$.
\end{lemma}

\begin{proof}
Par le résultat du Lemme \ref{chow-lemma-compose-base-change}, nous obtenons
un système de groupes de cycles $Z_k(X_i)$ et un système de
groupes de Chow $\CH_k(X_i)$ ainsi que des applications
$\colim Z_k(X_i) \to Z_k(X')$ et $\colim \CH_i(X_i) \to \CH_k(X')$.
Nous pouvons remplacer $S$ par un ouvert quasi-compact par lequel $X \to S$
se factorise ; nous pouvons donc supposer, et nous supposons, que tous les schémas intervenant dans
cette démonstration sont noethériens (et donc quasi-compacts et quasi-séparés).

\medskip\noindent
Montrons que l'application $\colim Z_k(X_i) \to Z_k(X')$ est surjective.
Soit en effet $Z' \subset X'$
un sous-schéma fermé intègre de $\delta'$-dimension $k$. Par
Limites, Lemme \ref{limits-lemma-descend-finite-presentation},
nous pouvons trouver un $i$ et un morphisme $Z_i \to X_i$ de présentation finie
dont le changement de base est $Z'$. En augmentant $i$, nous pouvons supposer que $Z_i$
est un sous-schéma fermé de $X_i$, voir
Limites, Lemme \ref{limits-lemma-descend-closed-immersion-finite-presentation}.
Alors $Z' \to X_i$ se factorise par $Z_i$ et nous pouvons remplacer $Z_i$
par l'image schématique de $Z' \to X_i$. Nous voyons ainsi
que nous pouvons supposer $Z_i$ fermé intègre dans $X_i$.
Par le Lemme \ref{chow-lemma-dimension-base-change}, nous concluons que
$\dim_{\delta_i}(Z_i) = \dim_{\delta'}(Z') = k$.
Ainsi, $Z_k(X_i) \to Z_k(X')$ envoie $[Z_i]$ sur $[Z']$ et
nous concluons à la surjectivité.

\medskip\noindent
Montrons que l'application $\colim Z_k(X_i) \to Z_k(X')$ est injective.
Soit $\alpha_i = \sum n_j[Z_j] \in Z_k(X_i)$ un cycle dont l'image
dans $Z_k(X')$ est nulle. Nous pouvons supposer, et supposons, $Z_j \not = Z_{j'}$ si
$j \not = j'$ et $n_j \not = 0$ pour tout $j$.
Notons $Z'_j \subset X'$ le changement de base de $Z_j$.
Par le Lemme \ref{chow-lemma-dimension-base-change}, chaque composante irréductible
de $Z'_j$ a pour $\delta'$-dimension $k$. De plus, comme $Z_j$ est irréductible
et $Z'_j \to Z_j$ est plat (comme changement de base de $S' \to S$),
nous voyons que $Z'_j \to Z_j$ est dominant.
Ainsi, si $Z'_j$ est non vide, une composante irréductible, disons
$Z'$, de $Z'_j$ domine $Z_j$. Il s'ensuit que $Z'$
ne peut être une composante irréductible de $Z'_{j'}$ pour $j' \not = j$.
Ainsi, si $Z'_j$ est non vide, nous voyons que $(S' \to S_i)^*\alpha_i
= \sum [Z'_j]_r$ est non nul (car le coefficient de $Z'$ serait non nul).
Nous voyons donc que $Z'_j = \emptyset$ pour tout $j$. Or, cela signifie
que le changement de base de $Z_j$ par une application de transition $S_{i'} \to S_i$
est vide par Limites, Lemme \ref{limits-lemma-limit-nonempty}.
Ainsi, $\alpha_i$ s'annule dans la colimite comme voulu.

\medskip\noindent
La surjectivité de $\colim Z_k(X_i) \to Z_k(X')$ implique que
$\colim \CH_k(X_i) \to \CH_k(X')$ est surjective.
Pour terminer la démonstration, montrons que cette application est injective.
Soit $\alpha_i \in \CH_k(X_i)$
un cycle dont l'image $\alpha' \in \CH_k(X')$ est nulle.
Il existe alors des sous-schémas fermés intègres
$W_l' \subset X'$, $l = 1, \ldots, r$, de $\delta"$-dimension $k + 1$
et des fonctions rationnelles non nulles $f'_l$ sur $W'_l$
tels que $\alpha' = \sum_{l = 1, \ldots, r} \text{div}_{W'_l}(f'_l)$.
En raisonnant comme ci-dessus, nous pouvons trouver un $i$ et des sous-schémas fermés intègres
$W_{i, l} \subset X_i$ de $\delta_i$-dimension $k + 1$
dont le changement de base est $W'_l$.
En augmentant $i$, nous pouvons supposer que nous avons des fonctions rationnelles
$f_{i, l}$ sur $W_{i, l}$. En effet, nous pouvons considérer $f'_l$ comme une
section du faisceau structural sur un ouvert non vide $U'_l \subset W'_l$ ;
nous pouvons descendre ces ouverts par Limites, Lemme \ref{limits-lemma-descend-opens},
et, en augmentant $i$, descendre $f'_l$ par
Limites, Lemme \ref{limits-lemma-descend-section}.
Nous affirmons que
$$
\alpha_i = \sum\nolimits_{l = 1, \ldots, r} \text{div}_{W_{i, l}}(f_{i, l})
$$
après avoir éventuellement augmenté $i$.

\medskip\noindent
Pour prouver l'affirmation, soit $Z'_{l, j} \subset W'_l$ une famille
finie de sous-schémas fermés intègres de $\delta'$-dimension $k$
tels que $f'_l$ soit une fonction régulière inversible en dehors de
$\bigcup_j Y'_{l, j}$. En augmentant $i$ (par les arguments ci-dessus),
nous pouvons supposer qu'il existe des sous-schémas fermés intègres $Z_{i, l, j} \subset W_i$
de $\delta_i$-dimension $k$ tels que $f_{i, l}$ soit une
fonction régulière inversible en dehors de $\bigcup_j Z_{i, l, j}$.
Nous pouvons alors écrire
$$
\text{div}_{W'_l}(f'_l) = \sum n_{l, j} [Z'_{l, j}]
$$
et
$$
\text{div}_{W_{i, l}}(f_{i, l}) = \sum n_{i, l, j} [Z_{i, l, j}]
$$
Pour prouver l'affirmation, il suffit de montrer que $n_{l, i} = n_{i, l, j}$.
En effet, ceci impliquera que $\beta_i =
\alpha_i - \sum\nolimits_{l = 1, \ldots, r} \text{div}_{W_{i, l}}(f_{i, l})$
est un cycle sur $X_i$ dont l'image inverse sur $X'$ est nulle comme cycle !
Il s'ensuit que l'image inverse de $\beta_i$ est nulle comme cycle sur $X_{i'}$
pour un certain $i' \geq i$, par un argument facile que nous omettons.

\medskip\noindent
Pour prouver l'égalité $n_{l, i} = n_{i, l, j}$, choisissons un
point générique $\xi' \in Z'_{l, j}$ et notons
$\xi \in Z_{i, l, j}$ son image, qui est également un point générique.
Alors, l'homomorphisme d'anneaux locaux
$$
\mathcal{O}_{W_{i, l}, \xi}
\longrightarrow
\mathcal{O}_{W'_l, \xi'}
$$
est plat car $W'_l \to W_{i, l}$ est le changement de base du morphisme
plat $S' \to S_i$. Nous avons aussi
$\mathfrak m_\xi \mathcal{O}_{W'_l, \xi'} = \mathfrak m_{\xi'}$
car l'image inverse de $Z_{i, l, j}$ est $Z'_{l, j}$ ! Ainsi, l'égalité de
$$
n_{l, j} = \text{ord}_{Z'_{l, j}}(f'_l) =
\text{ord}_{\mathcal{O}_{W'_l, \xi'}}(f'_l)
\quad\text{et}\quad
n_{i, l, j} = \text{ord}_{Z_{i, l, j}}(f_{i, l}) =
\text{ord}_{\mathcal{O}_{W_{i, l}, \xi}}(f_{i, l})
$$
résulte de Algèbre, Lemme \ref{algebra-lemma-pullback-module},
et de la construction de $\text{ord}$ dans
Algèbre, section \ref{algebra-section-orders-of-vanishing}.
\end{proof}












\section{Annexe A : autre approche du lemme clé}
\label{chow-section-appendix-A}

\noindent
Dans cette annexe, nous définissons d'abord les déterminants $\det_\kappa(M)$
des modules de longueur finie $M$ sur des anneaux locaux $(R, \mathfrak m, \kappa)$,
voir la Sous-section \ref{chow-subsection-determinants-finite-length}.
Le déterminant $\det_\kappa(M)$ est de dimension $1$ comme $\kappa$-espace vectoriel.
Nous l'utilisons dans la Sous-section \ref{chow-subsection-periodic-complexes-determinants}
pour définir le déterminant $\det_\kappa(M, \varphi, \psi) \in \kappa^*$
d'un complexe $(2, 1)$-périodique exact $(M, \varphi, \psi)$
avec $M$ de longueur finie. Dans la Sous-section \ref{chow-subsection-symbols},
nous utilisons ces déterminants pour construire un symbole modéré
$d_R(a, b) = \det_\kappa(R/ab, a, b)$ pour un couple de non-diviseurs de zéro
$a, b \in R$ lorsque $R$ est noethérien de dimension $1$.
Bien qu'il ne fasse aucun doute que
$$
d_R(a, b) = \partial_R(a, b)
$$
où $\partial_R$ est comme dans la section \ref{chow-section-tame-symbol},
nous n'avons pas (encore) ajouté la vérification. L'avantage du
symbole modéré construit dans cette annexe est qu'il s'étend
(par exemple) aux couples d'endomorphismes injectifs $\varphi, \psi$
d'un $R$-module de type fini $M$ de dimension $1$ tels que
$\varphi(\psi(M)) = \psi(\varphi(M))$. Dans la
Sous-section \ref{chow-subsection-length-determinant},
nous relions quotients de Herbrand et déterminants.
Une version facile à énoncer du résultat principal
(Proposition \ref{chow-proposition-length-determinant-periodic-complex})
est la formule
$$
-e_R(M, \varphi, \psi) =
\text{ord}_R(\det\nolimits_K(M_K, \varphi, \psi))
$$
lorsque $(M, \varphi, \psi)$ est un complexe $(2, 1)$-périodique
dont le quotient de Herbrand $e_R$ (Définition \ref{chow-definition-periodic-length})
est défini
sur un anneau local noethérien intègre de dimension $1$, noté $R$, de corps des fractions $K$.
Nous utilisons cette proposition pour donner une autre démonstration du lemme clé
(Lemme \ref{chow-lemma-milnor-gersten-low-degree})
pour le symbole modéré construit dans cette annexe, voir le
Lemme \ref{chow-lemma-secondary-ramification}.


\subsection{Déterminants de modules de longueur finie}
\label{chow-subsection-determinants-finite-length}

\noindent
Le contenu de cette section est apparenté à celui de
l'article \cite{determinant} et à celui de la
thèse \cite{Joe}.

\medskip\noindent
Soit $(R, \mathfrak m, \kappa)$ un anneau local. Soit
$\varphi : M \to M$ un endomorphisme $R$-linéaire d'un
$R$-module de longueur finie $M$. Dans Compléments sur l'algèbre, section
\ref{more-algebra-section-determinants-finite-length},
nous avons déjà défini le déterminant $\det_\kappa(\varphi)$
(ainsi que la trace et le polynôme caractéristique)
de $\varphi$ relativement à $\kappa$. Dans cette section, nous allons
construire un espace vectoriel canonique de dimension $1$ sur $\kappa$
$\det_\kappa(M)$ tel que
$\det_\kappa(\varphi : M \to M) : \det_\kappa(M) \to \det_\kappa(M)$
est égal à la multiplication par $\det_\kappa(\varphi)$.
Si $M$ est annulé par $\mathfrak m$, alors $M$ peut être considéré
comme un $\kappa$-espace vectoriel de dimension finie et nous avons alors
$\det_\kappa(M) = \wedge^n_\kappa(M)$, où $n = \dim_\kappa(M)$.
Notre construction généralisera ceci à tous les modules de longueur finie
sur $R$ et, si $R$ contient son corps résiduel, le déterminant
$\det_\kappa(M)$ sera donné par le déterminant usuel en un sens
approprié, voir la Remarque \ref{chow-remark-explain-determinant}.

\begin{definition}
\label{chow-definition-determinant}
Soit $R$ un anneau local d'idéal maximal $\mathfrak m$ et de
corps résiduel $\kappa$. Soit $M$ un $R$-module de longueur finie.
Posons $l = \text{longueur}_R(M)$.
\begin{enumerate}
\item Étant donnés des éléments $x_1, \ldots, x_r \in M$, nous notons
$\langle x_1, \ldots, x_r \rangle = Rx_1 + \ldots + Rx_r$ le
sous-$R$-module de $M$ engendré par $x_1, \ldots, x_r$.
\item Nous dirons qu'un $l$-uplet d'éléments
$(e_1, \ldots, e_l)$ de $M$ est {\it admissible} si
$\mathfrak m e_i \subset \langle e_1, \ldots, e_{i - 1} \rangle$
pour $i = 1, \ldots, l$.
\item Un {\it symbole} $[e_1, \ldots, e_l]$ signifiera que
$(e_1, \ldots, e_l)$ est un $l$-uplet admissible.
\item Une {\it relation admissible} entre symboles est l'une des suivantes :
\begin{enumerate}
\item si $(e_1, \ldots, e_l)$ est une suite admissible et si,
pour un certain $1 \leq a \leq l$, nous avons
$e_a \in \langle e_1, \ldots, e_{a - 1}\rangle$, alors
$[e_1, \ldots, e_l] = 0$,
\item si $(e_1, \ldots, e_l)$ est une suite admissible et si,
pour un certain $1 \leq a \leq l$, nous avons $e_a = \lambda e'_a + x$
avec $\lambda \in R^*$, et
$x \in \langle e_1, \ldots, e_{a - 1}\rangle$, alors
$$
[e_1, \ldots, e_l] =
\overline{\lambda} [e_1, \ldots, e_{a - 1}, e'_a, e_{a + 1}, \ldots, e_l]
$$
où $\overline{\lambda} \in \kappa^*$ est l'image de $\lambda$ dans
le corps résiduel, et
\item si $(e_1, \ldots, e_l)$ est une suite admissible et si
$\mathfrak m e_a \subset \langle e_1, \ldots, e_{a - 2}\rangle$, alors
$$
[e_1, \ldots, e_l] =
- [e_1, \ldots, e_{a - 2}, e_a, e_{a - 1}, e_{a + 1}, \ldots, e_l].
$$
\end{enumerate}
\item
Nous définissons le {\it déterminant du $R$-module de longueur finie $M$} par
$$
\det\nolimits_\kappa(M) =
\left\{
\frac{\kappa\text{-espace vectoriel engendré par les symboles}}
{\kappa\text{-combinaisons linéaires de relations admissibles}}
\right\}
$$
\end{enumerate}
\end{definition}

\noindent
Soulignons que l'on a toujours $l = \text{longueur}_R(M)$. Soulignons aussi qu'il
n'en résulte pas que le symbole $[e_1, \ldots, e_l]$ soit
additif en ses entrées (ce ne sera typiquement pas le cas).
Avant de pouvoir montrer que le déterminant $\det_\kappa(M)$ est effectivement
de dimension $1$, nous devons montrer qu'il est de dimension au plus $1$.

\begin{lemma}
\label{chow-lemma-dimension-at-most-one}
Avec les notations ci-dessus, nous avons $\dim_\kappa(\det_\kappa(M)) \leq 1$.
\end{lemma}

\begin{proof}
Fixons une suite admissible $(f_1, \ldots, f_l)$ de $M$ telle que
$$
\text{longueur}_R(\langle f_1, \ldots, f_i\rangle) = i
$$
pour $i = 1, \ldots, l$. Une telle suite admissible existe précisément parce que
$M$ est de longueur $l$. Nous allons montrer que tout élément de
$\det_\kappa(M)$ est un multiple par un élément de $\kappa$ du symbole
$[f_1, \ldots, f_l]$. Ceci prouvera le lemme.

\medskip\noindent
Soit $(e_1, \ldots, e_l)$ une suite admissible de $M$.
Il suffit de montrer que $[e_1, \ldots, e_l]$ est un multiple
de $[f_1, \ldots, f_l]$. Supposons d'abord que
$\langle e_1, \ldots, e_l\rangle \not = M$. Il existe alors
un $i \in [1, \ldots, l]$ tel que
$e_i \in \langle e_1, \ldots, e_{i - 1}\rangle$. Il résulte immédiatement
de la première relation admissible que
$[e_1, \ldots, e_n] = 0$ dans $\det_\kappa(M)$.
Nous pouvons donc supposer que $\langle e_1, \ldots, e_l\rangle = M$.
En particulier, il existe un plus petit indice $i \in \{1, \ldots, l\}$
tel que $f_1 \in \langle e_1, \ldots, e_i\rangle$. Ceci signifie
que $e_i = \lambda f_1 + x$, avec
$x \in \langle e_1, \ldots, e_{i - 1}\rangle$ et $\lambda \in R^*$.
Par la deuxième relation admissible, ceci signifie que
$[e_1, \ldots, e_l] =
\overline{\lambda}[e_1, \ldots, e_{i - 1}, f_1, e_{i + 1}, \ldots, e_l]$.
Notons que $\mathfrak m f_1 = 0$. Ainsi, en appliquant la troisième
relation admissible $i - 1$ fois, nous voyons que
$$
[e_1, \ldots, e_l] =
(-1)^{i - 1}\overline{\lambda}
[f_1, e_1, \ldots, e_{i - 1}, e_{i + 1}, \ldots, e_l].
$$
Notons que nous avons aussi
$ \langle f_1, e_1, \ldots, e_{i - 1}, e_{i + 1}, \ldots, e_l\rangle = M$.
Supposons par récurrence que nous ayons prouvé que notre symbole
initial est égal au produit d'un scalaire par
$$
[f_1, \ldots, f_j, e_{j + 1}, \ldots, e_l]
$$
pour une certaine suite admissible $(f_1, \ldots, f_j, e_{j + 1}, \ldots, e_l)$
dont les éléments engendrent $M$, c'est-à-dire \ avec
$\langle f_1, \ldots, f_j, e_{j + 1}, \ldots, e_l\rangle = M$.
Nous trouvons alors le plus petit $i$ tel que
$f_{j + 1} \in \langle f_1, \ldots, f_j, e_{j + 1}, \ldots, e_i\rangle$
et nous procédons comme ci-dessus pour voir que
$$
[f_1, \ldots, f_j, e_{j + 1}, \ldots, e_l]
=
(\text{scalaire}) [f_1, \ldots, f_j, f_{j + 1}, e_{j + 1},
\ldots, \hat{e_i}, \ldots, e_l]
$$
En continuant ainsi, nous obtenons le résultat désiré.
\end{proof}

\noindent
Avant de montrer que $\det_\kappa(M)$ est toujours de dimension $1$,
montrons qu'il coïncide avec la puissance extérieure maximale usuelle dans
le cas où le module est un espace vectoriel sur $\kappa$.

\begin{lemma}
\label{chow-lemma-compare-det}
Soit $R$ un anneau local d'idéal maximal $\mathfrak m$ et de
corps résiduel $\kappa$. Soit $M$ un $R$-module de longueur finie
annulé par $\mathfrak m$. Soit $l = \dim_\kappa(M)$.
Alors l'application
$$
\det\nolimits_\kappa(M) \longrightarrow \wedge^l_\kappa(M),
\quad
[e_1, \ldots, e_l] \longmapsto e_1 \wedge \ldots \wedge e_l
$$
est un isomorphisme.
\end{lemma}

\begin{proof}
Il est clair que la règle décrite dans le lemme donne une application $\kappa$-linéaire
puisque toutes les relations admissibles sont satisfaites par les
symboles usuels $e_1 \wedge \ldots \wedge e_l$. Cette application est aussi clairement
surjective. Puisque, par le Lemme \ref{chow-lemma-dimension-at-most-one}, le membre de gauche
est de dimension au plus un,
nous voyons que l'application est un isomorphisme.
\end{proof}

\begin{lemma}
\label{chow-lemma-determinant-dimension-one}
Soit $R$ un anneau local d'idéal maximal $\mathfrak m$ et de
corps résiduel $\kappa$. Soit $M$ un $R$-module de longueur finie.
Le déterminant $\det_\kappa(M)$ défini ci-dessus est un $\kappa$-espace vectoriel
de dimension $1$. Il est engendré par le symbole
$[f_1, \ldots, f_l]$ pour toute suite admissible telle
que $\langle f_1, \ldots f_l \rangle = M$.
\end{lemma}

\begin{proof}
Nous savons que $\det_\kappa(M)$ est de dimension au plus $1$, et en fait qu'il
est engendré par $[f_1, \ldots, f_l]$, par le
Lemme \ref{chow-lemma-dimension-at-most-one} et sa démonstration.
Nous allons montrer par récurrence sur $l = \text{longueur}(M)$
qu'il est non nul. Pour $l = 1$, ceci résulte du Lemme \ref{chow-lemma-compare-det}.
Choisissons un élément non nul $f \in M$
tel que $\mathfrak m f = 0$. Posons $\overline{M} = M /\langle f \rangle$,
et notons l'application quotient $x \mapsto \overline{x}$.
Nous allons définir une application surjective
$$
\psi : \det\nolimits_k(M) \to \det\nolimits_\kappa(\overline{M})
$$
qui prouvera le lemme puisque, par récurrence, le déterminant de
$\overline{M}$ est non nul.

\medskip\noindent
Nous définissons $\psi$ sur les symboles comme suit.
Soit $(e_1, \ldots, e_l)$ une suite admissible.
Si $f \not \in \langle e_1, \ldots, e_l \rangle$, alors
nous posons simplement $\psi([e_1, \ldots, e_l]) = 0$.
Si $f \in \langle e_1, \ldots, e_l \rangle$, alors nous choisissons
un $i$ minimal tel que $f \in \langle e_1, \ldots, e_i \rangle$.
Nous pouvons écrire $e_i = \lambda f + x$ pour une unité $\lambda \in R$
et $x \in \langle e_1, \ldots, e_{i - 1} \rangle$.
Dans ce cas, nous posons
$$
\psi([e_1, \ldots, e_l]) =
(-1)^i
\overline{\lambda}[\overline{e}_1, \ldots,
\overline{e}_{i - 1},
\overline{e}_{i + 1}, \ldots, \overline{e}_l].
$$
Notons que nous avons bien
$(\overline{e}_1, \ldots,
\overline{e}_{i - 1},
\overline{e}_{i + 1}, \ldots, \overline{e}_l)$
est une suite admissible dans $\overline{M}$, de sorte que ceci a un sens.
Montrons qu'étendre cette règle $\kappa$-linéairement aux
combinaisons linéaires de symboles conduit bien à une application sur les
déterminants. Pour cela, nous devons montrer que les relations
admissibles sont envoyées sur zéro.

\medskip\noindent
Relations de type (a). Supposons que $(e_1, \ldots, e_l)$ soit une
suite admissible et que, pour un certain $1 \leq a \leq l$, nous ayons
$e_a \in \langle e_1, \ldots, e_{a - 1}\rangle$.
Supposons que $f \in \langle e_1, \ldots, e_i\rangle$, avec $i$ minimal.
Alors $i \not = a$ et
$\overline{e}_a \in \langle \overline{e}_1, \ldots,
\hat{\overline{e}_i}, \ldots, \overline{e}_{a - 1}\rangle$ si $i < a$
ou
$\overline{e}_a \in \langle \overline{e}_1, \ldots,
\overline{e}_{a - 1}\rangle$ si $i > a$.
Ainsi, la même relation admissible pour $\det_\kappa(\overline{M})$ force
le symbole $[\overline{e}_1, \ldots,
\overline{e}_{i - 1},
\overline{e}_{i + 1}, \ldots, \overline{e}_l]$
à être nul comme voulu.

\medskip\noindent
Relations de type (b). Supposons que $(e_1, \ldots, e_l)$ soit une
suite admissible et que, pour un certain $1 \leq a \leq l$, nous ayons
$e_a = \lambda e'_a + x$, avec $\lambda \in R^*$, et
$x \in \langle e_1, \ldots, e_{a - 1}\rangle$.
Supposons que $f \in \langle e_1, \ldots, e_i\rangle$, avec $i$ minimal.
Écrivons $e_i = \mu f + y$, avec $y \in \langle e_1, \ldots, e_{i - 1}\rangle$.
Si $i < a$, alors l'égalité désirée est
$$
(-1)^i
\overline{\lambda}
[\overline{e}_1,
\ldots,
\overline{e}_{i - 1},
\overline{e}_{i + 1},
\ldots,
\overline{e}_l]
=
(-1)^i
\overline{\lambda}
[\overline{e}_1,
\ldots,
\overline{e}_{i - 1},
\overline{e}_{i + 1},
\ldots,
\overline{e}_{a - 1},
\overline{e}'_a,
\overline{e}_{a + 1},
\ldots,
\overline{e}_l]
$$
ce qui résulte de $\overline{e}_a = \lambda \overline{e}'_a + \overline{x}$
et de la relation admissible correspondante pour $\det_\kappa(\overline{M})$.
Si $i > a$, alors l'égalité désirée est
$$
(-1)^i
\overline{\lambda}
[\overline{e}_1,
\ldots,
\overline{e}_{i - 1},
\overline{e}_{i + 1},
\ldots,
\overline{e}_l]
=
(-1)^i
\overline{\lambda}
[\overline{e}_1,
\ldots,
\overline{e}_{a - 1},
\overline{e}'_a,
\overline{e}_{a + 1},
\ldots,
\overline{e}_{i - 1},
\overline{e}_{i + 1},
\ldots,
\overline{e}_l]
$$
ce qui résulte de $\overline{e}_a = \lambda \overline{e}'_a + \overline{x}$
et de la relation admissible correspondante pour $\det_\kappa(\overline{M})$.
Le cas intéressant est celui où $i = a$. Dans ce cas, nous avons
$e_a = \lambda e'_a + x = \mu f + y$. Donc aussi
$e'_a = \lambda^{-1}(\mu f + y - x)$. Ainsi, nous voyons que
$$
\psi([e_1, \ldots, e_l])
= (-1)^i \overline{\mu}
[\overline{e}_1, \ldots,
\overline{e}_{i - 1},
\overline{e}_{i + 1}, \ldots, \overline{e}_l]
=
\psi(
\overline{\lambda}
[e_1, \ldots, e_{a - 1}, e'_a, e_{a + 1}, \ldots, e_l]
)
$$
comme voulu.

\medskip\noindent
Relations de type (c). Supposons que $(e_1, \ldots, e_l)$
soit une suite admissible et que
$\mathfrak m e_a \subset \langle e_1, \ldots, e_{a - 2}\rangle$.
Supposons que $f \in \langle e_1, \ldots, e_i\rangle$, avec $i$ minimal.
Écrivons $e_i = \lambda f + x$, avec $x \in \langle e_1, \ldots, e_{i - 1}\rangle$.
Nous distinguons $4$ cas :

\medskip\noindent
Cas 1 : $i < a - 1$. L'égalité désirée est
\begin{align*}
& (-1)^i
\overline{\lambda}
[\overline{e}_1,
\ldots,
\overline{e}_{i - 1},
\overline{e}_{i + 1},
\ldots,
\overline{e}_l] \\
& =
(-1)^{i + 1}
\overline{\lambda}
[\overline{e}_1,
\ldots,
\overline{e}_{i - 1},
\overline{e}_{i + 1},
\ldots,
\overline{e}_{a - 2},
\overline{e}_a,
\overline{e}_{a - 1},
\overline{e}_{a + 1},
\ldots,
\overline{e}_l]
\end{align*}
ce qui résulte de la relation admissible de type (c) pour
$\det_\kappa(\overline{M})$.

\medskip\noindent
Cas 2 : $i > a$. L'égalité désirée est
\begin{align*}
& (-1)^i
\overline{\lambda}
[\overline{e}_1,
\ldots,
\overline{e}_{i - 1},
\overline{e}_{i + 1},
\ldots,
\overline{e}_l] \\
& =
(-1)^{i + 1}
\overline{\lambda}
[\overline{e}_1,
\ldots,
\overline{e}_{a - 2},
\overline{e}_a,
\overline{e}_{a - 1},
\overline{e}_{a + 1},
\ldots,
\overline{e}_{i - 1},
\overline{e}_{i + 1},
\ldots,
\overline{e}_l]
\end{align*}
ce qui résulte de la relation admissible de type (c) pour
$\det_\kappa(\overline{M})$.

\medskip\noindent
Cas 3 : $i = a$. Nous écrivons $e_a = \lambda f + \mu e_{a - 1} + y$,
avec $y \in \langle e_1, \ldots, e_{a - 2}\rangle$. Alors
$$
\psi([e_1, \ldots, e_l]) =
(-1)^a
\overline{\lambda}
[\overline{e}_1,
\ldots,
\overline{e}_{a - 1},
\overline{e}_{a + 1},
\ldots,
\overline{e}_l]
$$
par définition. Si $\overline{\mu}$ est non nul, alors nous avons
$e_{a - 1} = - \mu^{-1} \lambda f + \mu^{-1}e_a - \mu^{-1} y$
et nous obtenons
$$
\psi(-[e_1, \ldots, e_{a - 2}, e_a, e_{a - 1}, e_{a + 1}, \ldots, e_l]) =
(-1)^a
\overline{\mu^{-1}\lambda}
[\overline{e}_1,
\ldots,
\overline{e}_{a - 2},
\overline{e}_a,
\overline{e}_{a + 1},
\ldots,
\overline{e}_l]
$$
par définition. Puisque, dans $\overline{M}$, nous avons
$\overline{e}_a = \mu \overline{e}_{a - 1} + \overline{y}$, nous voyons
que les deux résultats sont égaux par la relation (a) pour $\det_\kappa(\overline{M})$.
Si, en revanche, $\overline{\mu}$ est nul, alors nous pouvons écrire
$e_a = \lambda f + y$, avec $y \in \langle e_1, \ldots, e_{a - 2}\rangle$,
et nous avons
$$
\psi(-[e_1, \ldots, e_{a - 2}, e_a, e_{a - 1}, e_{a + 1}, \ldots, e_l]) =
(-1)^a
\overline{\lambda}
[\overline{e}_1,
\ldots,
\overline{e}_{a - 1},
\overline{e}_{a + 1},
\ldots,
\overline{e}_l]
$$
ce qui est égal à $\psi([e_1, \ldots, e_l])$.

\medskip\noindent
Cas 4 : $i = a - 1$. Ici, nous avons
$$
\psi([e_1, \ldots, e_l]) =
(-1)^{a - 1}
\overline{\lambda}
[\overline{e}_1,
\ldots,
\overline{e}_{a - 2},
\overline{e}_a,
\ldots,
\overline{e}_l]
$$
par définition. Si $f \not \in \langle e_1, \ldots, e_{a - 2}, e_a \rangle$,
alors
$$
\psi(-[e_1, \ldots, e_{a - 2}, e_a, e_{a - 1}, e_{a + 1}, \ldots, e_l]) =
(-1)^{a + 1}\overline{\lambda}
[\overline{e}_1,
\ldots,
\overline{e}_{a - 2},
\overline{e}_a,
\ldots,
\overline{e}_l]
$$
Puisque $(-1)^{a - 1} = (-1)^{a + 1}$, les deux expressions sont les mêmes.
Enfin, supposons $f \in \langle e_1, \ldots, e_{a - 2}, e_a \rangle$.
Dans ce cas, nous voyons que $e_{a - 1} = \lambda f + x$, avec
$x \in \langle e_1, \ldots, e_{a - 2}\rangle$, et
$e_a = \mu f + y$, avec $y \in \langle e_1, \ldots, e_{a - 2}\rangle$,
pour des unités $\lambda, \mu \in R$.
Nous concluons que nous avons à la fois
$e_a \in \langle e_1, \ldots, e_{a - 1} \rangle$ et
$e_{a - 1} \in \langle e_1, \ldots, e_{a - 2}, e_a\rangle$.
Dans ce cas, une relation de type (a) s'applique à la fois à
$[e_1, \ldots, e_l]$ et
$[e_1, \ldots, e_{a - 2}, e_a, e_{a - 1}, e_{a + 1}, \ldots, e_l]$
et la compatibilité de $\psi$ avec celles-ci montrée ci-dessus fait voir que les deux
$$
\psi([e_1, \ldots, e_l])
\quad\text{et}\quad
\psi([e_1, \ldots, e_{a - 2}, e_a, e_{a - 1}, e_{a + 1}, \ldots, e_l])
$$
sont nuls, comme voulu.

\medskip\noindent
À ce stade, nous avons montré que $\psi$ est bien définie, et il ne reste
qu'à montrer qu'elle est surjective. Pour le voir, soit
$(\overline{f}_2, \ldots, \overline{f}_l)$ une suite admissible
dans $\overline{M}$. Nous pouvons choisir des relèvements $f_2, \ldots, f_l \in M$, et
alors $(f, f_2, \ldots, f_l)$ est une suite admissible dans $M$.
Puisque $\psi([f, f_2, \ldots, f_l]) = [f_2, \ldots, f_l]$, nous avons gagné.
\end{proof}

\noindent
Soit $R$ un anneau local d'idéal maximal $\mathfrak m$ et de
corps résiduel $\kappa$. Notons que si $\varphi : M \to N$ est un
isomorphisme de $R$-modules de longueur finie, alors nous obtenons un
isomorphisme
$$
\det\nolimits_\kappa(\varphi) :
\det\nolimits_\kappa(M)
\to
\det\nolimits_\kappa(N)
$$
simplement par la règle
$$
\det\nolimits_\kappa(\varphi)([e_1, \ldots, e_l])
=
[\varphi(e_1), \ldots, \varphi(e_l)]
$$
pour tout symbole $[e_1, \ldots, e_l]$ de $M$.
Ainsi, nous voyons que $\det\nolimits_\kappa$ est un foncteur
\begin{equation}
\label{chow-equation-functor}
\left\{
\begin{matrix}
\text{les }R\text{-modules de longueur finie}\\
\text{avec isomorphismes}
\end{matrix}
\right\}
\longrightarrow
\left\{
\begin{matrix}
1\text{ est la dimension des }\kappa\text{-espaces vectoriels}\\
\text{avec isomorphismes}
\end{matrix}
\right\}
\end{equation}
Ceci est typique d'un « foncteur déterminant »
(voir \cite{Knudsen}), tout comme la propriété d'additivité
suivante.

\begin{lemma}
\label{chow-lemma-det-exact-sequences}
Soit $(R, \mathfrak m, \kappa)$ un anneau local.
Pour toute suite exacte courte
$$
0 \to K \to L \to M \to 0
$$
de $R$-modules de longueur finie, il existe un isomorphisme canonique
$$
\gamma_{K \to L \to M} :
\det\nolimits_\kappa(K) \otimes_\kappa \det\nolimits_\kappa(M)
\longrightarrow
\det\nolimits_\kappa(L)
$$
défini sur les symboles non nuls par la règle
$$
[e_1, \ldots, e_k]
\otimes
[\overline{f}_1, \ldots, \overline{f}_m]
\longrightarrow
[e_1, \ldots, e_k, f_1, \ldots, f_m]
$$
ayant les propriétés suivantes :
\begin{enumerate}
\item Pour tout isomorphisme de suites exactes courtes, c'est-à-dire pour
tout diagramme commutatif
$$
\xymatrix{
0 \ar[r] &
K \ar[r] \ar[d]^u &
L \ar[r] \ar[d]^v &
M \ar[r] \ar[d]^w &
0 \\
0 \ar[r] &
K' \ar[r] &
L' \ar[r] &
M' \ar[r] &
0
}
$$
dont les lignes sont exactes courtes et où $u, v, w$ sont des isomorphismes, nous avons
$$
\gamma_{K' \to L' \to M'} \circ
(\det\nolimits_\kappa(u) \otimes \det\nolimits_\kappa(w))
=
\det\nolimits_\kappa(v) \circ
\gamma_{K \to L \to M},
$$
\item pour tout carré commutatif de $R$-modules de longueur finie
dont les lignes et les colonnes sont exactes
$$
\xymatrix{
& 0 \ar[d] & 0 \ar[d] & 0 \ar[d] & \\
0 \ar[r] & A \ar[r] \ar[d] & B \ar[r] \ar[d] & C \ar[r] \ar[d] & 0 \\
0 \ar[r] & D \ar[r] \ar[d] & E \ar[r] \ar[d] & F \ar[r] \ar[d] & 0 \\
0 \ar[r] & G \ar[r] \ar[d] & H \ar[r] \ar[d] & I \ar[r] \ar[d] & 0 \\
& 0  & 0  & 0  &
}
$$
le diagramme suivant est commutatif
$$
\xymatrix{
\det\nolimits_\kappa(A) \otimes
\det\nolimits_\kappa(C) \otimes
\det\nolimits_\kappa(G) \otimes
\det\nolimits_\kappa(I)
\ar[dd]_{\epsilon}
\ar[rrr]_-{\gamma_{A \to B \to C} \otimes \gamma_{G \to H \to I}}
& & &
\det\nolimits_\kappa(B) \otimes
\det\nolimits_\kappa(H)
\ar[d]^{\gamma_{B \to E \to H}}
\\
& & & \det\nolimits_\kappa(E)
\\
\det\nolimits_\kappa(A) \otimes
\det\nolimits_\kappa(G) \otimes
\det\nolimits_\kappa(C) \otimes
\det\nolimits_\kappa(I)
\ar[rrr]^-{\gamma_{A \to D \to G} \otimes \gamma_{C \to F \to I}}
& & &
\det\nolimits_\kappa(D) \otimes
\det\nolimits_\kappa(F)
\ar[u]_{\gamma_{D \to E \to F}}
}
$$
où $\epsilon$ est la permutation des facteurs du produit tensoriel
multipliée par $(-1)^{cg}$, avec $c = \text{longueur}_R(C)$ et $g = \text{longueur}_R(G)$,
et
\item l'application $\gamma_{K \to L \to M}$ coïncide avec l'isomorphisme usuel
si $0 \to K \to L \to M \to 0$ est en fait une suite exacte courte
de $\kappa$-espaces vectoriels.
\end{enumerate}
\end{lemma}

\begin{proof}
L'intérêt de prendre des symboles non nuls dans la description explicite
de l'application $\gamma_{K \to L \to M}$ tient simplement à ce que, si $(e_1, \ldots, e_l)$
est une suite admissible dans $K$, et si
$(\overline{f}_1, \ldots, \overline{f}_m)$ est une suite admissible dans
$M$, il n'est pas garanti que $(e_1, \ldots, e_l, f_1, \ldots, f_m)$
soit une suite admissible dans $L$ (où, bien sûr, $f_i \in L$ désigne
un relèvement de $\overline{f}_i$). Toutefois, si le symbole
$[e_1, \ldots, e_l]$ est non nul dans $\det_\kappa(K)$, alors
nécessairement $K = \langle e_1, \ldots, e_k\rangle$ (voir la
démonstration du Lemme \ref{chow-lemma-dimension-at-most-one}), et
dans ce cas $(e_1, \ldots, e_k, f_1, \ldots, f_m)$ est bien
une suite admissible.
De plus, par les relations admissibles de type (b) pour $\det_\kappa(L)$,
nous voyons que la valeur de $[e_1, \ldots, e_k, f_1, \ldots, f_m]$ dans
$\det_\kappa(L)$ est indépendante du choix des relèvements
$f_i$ dans ce cas également. Compte tenu de cette remarque, il est clair
qu'une relation admissible entre $e_1, \ldots, e_k$ dans $K$
se traduit par une relation admissible entre
$e_1, \ldots, e_k, f_1, \ldots, f_m$ dans $L$, et
de même pour une relation admissible entre les
$\overline{f}_1, \ldots, \overline{f}_m$.
Ainsi, $\gamma$ définit une application linéaire d'espaces vectoriels comme affirmé dans le lemme.

\medskip\noindent
Par le Lemme \ref{chow-lemma-determinant-dimension-one}, nous savons que
$\det_\kappa(L)$ est engendré par n'importe quel
symbole $[x_1, \ldots, x_{k + m}]$ tel que
$(x_1, \ldots, x_{k + m})$ soit une suite admissible
avec $L = \langle x_1, \ldots, x_{k + m}\rangle$. Il est donc
clair que l'application $\gamma_{K \to L \to M}$ est surjective et
qu'elle est donc un isomorphisme.

\medskip\noindent
La propriété (1) vaut car
\begin{eqnarray*}
& & \det\nolimits_\kappa(v)([e_1, \ldots, e_k, f_1, \ldots, f_m]) \\
& = &
[v(e_1), \ldots, v(e_k), v(f_1), \ldots, v(f_m)] \\
& = &
\gamma_{K' \to L' \to M'}([u(e_1), \ldots, u(e_k)]
\otimes [w(f_1), \ldots, w(f_m)]).
\end{eqnarray*}
La propriété (2) signifie qu'étant donnés un symbole
$[\alpha_1, \ldots, \alpha_a]$ engendrant $\det_\kappa(A)$,
un symbole $[\gamma_1, \ldots, \gamma_c]$ engendrant $\det_\kappa(C)$,
un symbole $[\zeta_1, \ldots, \zeta_g]$ engendrant $\det_\kappa(G)$, et
un symbole $[\iota_1, \ldots, \iota_i]$ engendrant $\det_\kappa(I)$,
nous avons
\begin{eqnarray*}
& & [\alpha_1, \ldots, \alpha_a, \tilde\gamma_1, \ldots, \tilde\gamma_c,
\tilde\zeta_1, \ldots, \tilde\zeta_g, \tilde\iota_1, \ldots, \tilde\iota_i] \\
& = &
(-1)^{cg} [\alpha_1, \ldots, \alpha_a, \tilde\zeta_1, \ldots, \tilde\zeta_g,
\tilde\gamma_1, \ldots, \tilde\gamma_c, \tilde\iota_1, \ldots, \tilde\iota_i]
\end{eqnarray*}
(pour des relèvements convenables $\tilde{x}$ dans $E$) dans $\det_\kappa(E)$.
Ceci vaut car nous pouvons utiliser les relations admissibles de type (c)
$cg$ fois dans l'ordre suivant : faire passer
$\tilde\zeta_1$ devant les éléments
$\tilde\gamma_c, \ldots, \tilde\gamma_1$
(ce qui est permis puisque $\mathfrak m\tilde\zeta_1 \subset A$),
puis faire passer $\tilde\zeta_2$ devant les éléments
$\tilde\gamma_c, \ldots, \tilde\gamma_1$
(ce qui est permis puisque $\mathfrak m\tilde\zeta_2 \subset A + R\tilde\zeta_1$),
et ainsi de suite.

\medskip\noindent
La partie (3) du lemme est évidente.
Ceci termine la démonstration.
\end{proof}

\noindent
Nous pouvons utiliser les applications $\gamma$ du lemme pour définir des applications
$\gamma$ plus générales comme suit. Supposons que $(R, \mathfrak m, \kappa)$ soit un
anneau local. Soit $M$ un $R$-module de longueur finie et supposons donnée
une filtration finie (voir
Homologie, Définition \ref{homology-definition-filtered})
$$
0 = F^m \subset F^{m - 1} \subset \ldots \subset F^{n + 1} \subset F^n = M
$$
alors il existe un isomorphisme canonique bien défini
$$
\gamma_{(M, F)} :
\det\nolimits_\kappa(F^{m - 1}/F^m) \otimes_\kappa \ldots \otimes_k 
\det\nolimits_\kappa(F^n/F^{n + 1})
\longrightarrow
\det\nolimits_\kappa(M)
$$
Pour le construire, nous utilisons les isomorphismes du Lemme \ref{chow-lemma-det-exact-sequences}
provenant des suites exactes courtes
$0 \to F^{i - 1}/F^i \to M/F^i \to M/F^{i - 1} \to 0$.
La partie (2) du Lemme \ref{chow-lemma-det-exact-sequences}, avec $G = 0$, montre
que nous obtenons le même isomorphisme si nous utilisons les suites exactes courtes
$0 \to F^i \to F^{i - 1} \to F^{i - 1}/F^i \to 0$.

\medskip\noindent
Voici un autre résultat typique sur les foncteurs déterminant.
Il n'est pas difficile à montrer. La partie délicate consiste habituellement à établir
l'existence d'un foncteur déterminant.

\begin{lemma}
\label{chow-lemma-uniqueness-det}
Soit $(R, \mathfrak m, \kappa)$ un anneau local quelconque.
Le foncteur
$$
\det\nolimits_\kappa :
\left\{
\begin{matrix}
\text{les }R\text{-modules de longueur finie} \\
\text{avec isomorphismes}
\end{matrix}
\right\}
\longrightarrow
\left\{
\begin{matrix}
1\text{ est la dimension des }\kappa\text{-espaces vectoriels} \\
\text{avec isomorphismes}
\end{matrix}
\right\}
$$
muni des applications $\gamma_{K \to L \to M}$ est caractérisé par
les propriétés suivantes
\begin{enumerate}
\item sa restriction à la sous-catégorie des modules annulés
par $\mathfrak m$ est isomorphe au foncteur déterminant usuel
(voir le Lemme \ref{chow-lemma-compare-det}), et
\item les points (1), (2) et (3) du Lemme \ref{chow-lemma-det-exact-sequences}
sont vérifiés.
\end{enumerate}
\end{lemma}

\begin{proof}
Omis.
\end{proof}

\begin{lemma}
\label{chow-lemma-determinant-quotient-ring}
Soit $(R', \mathfrak m') \to (R, \mathfrak m)$ un homomorphisme d'anneaux
locaux qui induit un isomorphisme sur les corps résiduels $\kappa$.
Alors, pour tout $R$-module de longueur finie, la restriction $M_{R'}$
est un $R'$-module de longueur finie et il existe un isomorphisme canonique
$$
\det\nolimits_{R, \kappa}(M)
\longrightarrow
\det\nolimits_{R', \kappa}(M_{R'})
$$
Cet isomorphisme est fonctoriel en $M$ et compatible avec les
isomorphismes $\gamma_{K \to L \to M}$ du Lemme \ref{chow-lemma-det-exact-sequences}
définis pour $\det_{R, \kappa}$ et $\det_{R', \kappa}$.
\end{lemma}

\begin{proof}
Si la longueur de $M$ comme $R$-module est $l$, alors la longueur
de $M$ comme $R'$-module (c'est-à-dire de $M_{R'}$) est aussi $l$, voir
Algèbre, Lemme \ref{algebra-lemma-pushdown-module}.
Notons qu'une suite admissible $x_1, \ldots, x_l$ de $M$
sur $R$ est une suite admissible de $M$ sur $R'$, puisque $\mathfrak m'$
s'envoie dans $\mathfrak m$.
L'isomorphisme est obtenu en envoyant le symbole
$[x_1, \ldots, x_l] \in \det\nolimits_{R, \kappa}(M)$
sur le symbole correspondant
$[x_1, \ldots, x_l] \in \det\nolimits_{R', \kappa}(M)$.
Il est immédiat de vérifier que ceci est fonctoriel pour les
isomorphismes et compatible avec les isomorphismes
$\gamma$ du Lemme \ref{chow-lemma-det-exact-sequences}.
\end{proof}

\begin{remark}
\label{chow-remark-explain-determinant}
Soit $(R, \mathfrak m, \kappa)$ un anneau local et supposons soit que
la caractéristique de $\kappa$ est nulle, soit qu'elle vaut $p$ et que $p R = 0$.
Soient $M_1, \ldots, M_n$ des $R$-modules de longueur finie.
Nous montrerons ci-dessous qu'il existe un
idéal $I \subset \mathfrak m$ annulant $M_i$ pour $i = 1, \ldots, n$
et une section $\sigma : \kappa \to R/I$ de la surjection canonique
$R/I \to \kappa$. La restriction $M_{i, \kappa}$ de $M_i$ par $\sigma$
est un $\kappa$-espace vectoriel de dimension $l_i = \text{longueur}_R(M_i)$ et,
en utilisant le Lemme \ref{chow-lemma-determinant-quotient-ring}, nous voyons que
$$
\det\nolimits_\kappa(M_i) = \wedge_\kappa^{l_i}(M_{i, \kappa})
$$
Ces isomorphismes sont compatibles avec les isomorphismes
$\gamma_{K \to M \to L}$ du Lemme \ref{chow-lemma-det-exact-sequences}
pour les suites exactes courtes de $R$-modules de longueur finie annulés
par $I$. Nous en concluons que vérifier une propriété de
$\det_\kappa$ se réduit souvent à vérifier les propriétés correspondantes
du déterminant usuel sur la catégorie des espaces vectoriels de dimension
finie.

\medskip\noindent
Pour $I$, nous pouvons prendre l'annulateur
(Algèbre, Définition \ref{algebra-definition-annihilator})
du module $M = \bigoplus M_i$. Dans ce cas, nous voyons que
$R/I \subset \text{End}_R(M)$ est donc de longueur finie.
Ainsi, $R/I$ est un anneau local artinien de corps résiduel $\kappa$.
Puisqu'un anneau local artinien est complet, nous voyons que $R/I$
admet un anneau de coefficients par le théorème de structure de Cohen
(Algèbre, Théorème \ref{algebra-theorem-cohen-structure-theorem}),
qui est un corps par notre hypothèse sur $R$.
\end{remark}

\noindent
Voici un cas où nous pouvons calculer le déterminant d'une application linéaire.
En fait, il n'y a rien de mystérieux dans aucun cas ; voir
l'Exemple \ref{chow-example-determinant-map} pour un exemple quelconque.

\begin{lemma}
\label{chow-lemma-times-u-determinant}
Soit $R$ un anneau local de corps résiduel $\kappa$.
Soit $u \in R^*$ une unité.
Soit $M$ un module de longueur finie sur $R$.
Notons $u_M : M \to M$ l'application de multiplication par $u$.
Alors
$$
\det\nolimits_\kappa(u_M) :
\det\nolimits_\kappa(M)
\longrightarrow
\det\nolimits_\kappa(M)
$$
est la multiplication par $\overline{u}^l$, où $l = \text{longueur}_R(M)$
et $\overline{u} \in \kappa^*$ est l'image de $u$.
\end{lemma}

\begin{proof}
Notons $f_M \in \kappa^*$ l'élément tel que
$\det\nolimits_\kappa(u_M) = f_M \text{id}_{\det\nolimits_\kappa(M)}$.
Supposons que $0 \to K \to L \to M \to 0$ soit une suite exacte
courte de $R$-modules de type fini. Nous voyons alors que
$u_k$, $u_L$, $u_M$ donnent un isomorphisme de suites exactes courtes.
Ainsi, par le Lemme \ref{chow-lemma-det-exact-sequences} (1), nous concluons que
$f_K f_M = f_L$.
Ceci signifie que, par récurrence sur la longueur, il suffit de prouver le
lemme dans le cas de longueur $1$, où il est trivial.
\end{proof}

\begin{example}
\label{chow-example-determinant-map}
Considérons l'anneau local $R = \mathbf{Z}_p$.
Posons $M = \mathbf{Z}_p/(p^2) \oplus \mathbf{Z}_p/(p^3)$.
Soit $u : M \to M$ l'application donnée par la matrice
$$
u =
\left(
\begin{matrix}
a & b \\
pc & d
\end{matrix}
\right)
$$
où $a, b, c, d \in \mathbf{Z}_p$, et $a, d \in \mathbf{Z}_p^*$.
Dans ce cas, $\det_\kappa(u)$ est égal à la multiplication par
$a^2d^3 \bmod p \in \mathbf{F}_p^*$. Ceci se voit facilement
en considérant l'effet de $u$ sur le symbole
$[p^2e, pe, pf, e, f]$, où $e = (0 , 1) \in M$ et
$f = (1, 0) \in M$.
\end{example}





\subsection{Complexes périodiques et déterminants}
\label{chow-subsection-periodic-complexes-determinants}

\noindent
Soit $R$ un anneau local de corps résiduel $\kappa$.
Soit $(M, \varphi, \psi)$ un complexe $(2, 1)$-périodique sur $R$.
Supposons que $M$ soit de longueur finie et que $(M, \varphi, \psi)$ soit
exact. Nous allons utiliser la construction du déterminant pour définir
un invariant de cette situation. Voir la
Sous-section \ref{chow-subsection-determinants-finite-length}.
Abrégeons
$K_\varphi = \Ker(\varphi)$,
$I_\varphi = \Im(\varphi)$,
$K_\psi = \Ker(\psi)$, et
$I_\psi = \Im(\psi)$.
Les suites exactes courtes
$$
0 \to K_\varphi \to M \to I_\varphi \to 0, \quad
0 \to K_\psi \to M \to I_\psi \to 0
$$
donnent des isomorphismes
$$
\gamma_\varphi :
\det\nolimits_\kappa(K_\varphi)
\otimes
\det\nolimits_\kappa(I_\varphi)
\longrightarrow
\det\nolimits_\kappa(M), \quad
\gamma_\psi :
\det\nolimits_\kappa(K_\psi)
\otimes
\det\nolimits_\kappa(I_\psi)
\longrightarrow
\det\nolimits_\kappa(M),
$$
voir le lemme \ref{chow-lemma-det-exact-sequences}.
D'autre part, l'exactitude du complexe donne les égalités
$K_\varphi = I_\psi$, et $K_\psi = I_\varphi$
et donc un isomorphisme
$$
\sigma :
\det\nolimits_\kappa(K_\varphi)
\otimes
\det\nolimits_\kappa(I_\varphi)
\longrightarrow
\det\nolimits_\kappa(K_\psi)
\otimes
\det\nolimits_\kappa(I_\psi)
$$
en permutant les facteurs. Avec ces notations, nous pouvons définir notre invariant.

\begin{definition}
\label{chow-definition-periodic-determinant}
Soit $R$ un anneau local de corps résiduel $\kappa$.
Soit $(M, \varphi, \psi)$ un complexe $(2, 1)$-périodique sur $R$.
Supposons que $M$ soit de longueur finie et que $(M, \varphi, \psi)$ soit
exact. Le {\it déterminant de $(M, \varphi, \psi)$} est
l'élément
$$
\det\nolimits_\kappa(M, \varphi, \psi) \in \kappa^*
$$
tel que la composée
$$
\det\nolimits_\kappa(M)
\xrightarrow{\gamma_\psi \circ \sigma \circ \gamma_\varphi^{-1}}
\det\nolimits_\kappa(M)
$$
soit la multiplication par
$(-1)^{\text{longueur}_R(I_\varphi)\text{longueur}_R(I_\psi)}
\det\nolimits_\kappa(M, \varphi, \psi)$.
\end{definition}

\begin{remark}
\label{chow-remark-more-elementary}
Voici une description plus concrète du déterminant
introduit ci-dessus. Soit $R$ un anneau local de corps résiduel $\kappa$.
Soit $(M, \varphi, \psi)$ un complexe $(2, 1)$-périodique sur $R$.
Supposons que $M$ soit de longueur finie et que $(M, \varphi, \psi)$ soit
exact. Abrégeons $I_\varphi = \Im(\varphi)$,
$I_\psi = \Im(\psi)$ comme ci-dessus.
Supposons que $\text{longueur}_R(I_\varphi) = a$ et que
$\text{longueur}_R(I_\psi) = b$, de sorte que $a + b = \text{longueur}_R(M)$
par exactitude. Choisissons des suites admissibles
$x_1, \ldots, x_a \in I_\varphi$ et $y_1, \ldots, y_b \in I_\psi$
telles que le symbole $[x_1, \ldots, x_a]$ engendre $\det_\kappa(I_\varphi)$
et que le symbole $[x_1, \ldots, x_b]$ engendre $\det_\kappa(I_\psi)$.
Choisissons $\tilde x_i \in M$ tels que $\varphi(\tilde x_i) = x_i$.
Choisissons $\tilde y_j \in M$ tels que $\psi(\tilde y_j) = y_j$.
Alors $\det_\kappa(M, \varphi, \psi)$ est caractérisé
par l'égalité
$$
[x_1, \ldots, x_a, \tilde y_1, \ldots, \tilde y_b]
=
(-1)^{ab} \det\nolimits_\kappa(M, \varphi, \psi)
[y_1, \ldots, y_b, \tilde x_1, \ldots, \tilde x_a]
$$
dans $\det_\kappa(M)$. Cela explique également le signe.
\end{remark}

\begin{lemma}
\label{chow-lemma-periodic-determinant-shift}
Soit $R$ un anneau local de corps résiduel $\kappa$.
Soit $(M, \varphi, \psi)$ un complexe $(2, 1)$-périodique sur $R$.
Supposons que $M$ soit de longueur finie et que $(M, \varphi, \psi)$ soit
exact. Alors
$$
\det\nolimits_\kappa(M, \varphi, \psi)
\det\nolimits_\kappa(M, \psi, \varphi)
= 1.
$$
\end{lemma}

\begin{proof}
La démonstration est omise.
\end{proof}

\begin{lemma}
\label{chow-lemma-periodic-determinant-sign}
Soit $R$ un anneau local de corps résiduel $\kappa$.
Soit $(M, \varphi, \varphi)$ un complexe $(2, 1)$-périodique sur $R$.
Supposons que $M$ soit de longueur finie et que $(M, \varphi, \varphi)$ soit
exact. Alors $\text{longueur}_R(M) = 2 \text{longueur}_R(\Im(\varphi))$
et
$$
\det\nolimits_\kappa(M, \varphi, \varphi)
=
(-1)^{\text{longueur}_R(\Im(\varphi))}
=
(-1)^{\frac{1}{2}\text{longueur}_R(M)}
$$
\end{lemma}

\begin{proof}
Cela résulte directement de la règle des signes dans les définitions.
\end{proof}

\begin{lemma}
\label{chow-lemma-periodic-determinant-easy-case}
Soit $R$ un anneau local de corps résiduel $\kappa$.
Soit $M$ un $R$-module de longueur finie.
\begin{enumerate}
\item si $\varphi : M \to M$ est un isomorphisme, alors
$\det_\kappa(M, \varphi, 0) = \det_\kappa(\varphi)$.
\item si $\psi : M \to M$ est un isomorphisme, alors
$\det_\kappa(M, 0, \psi) = \det_\kappa(\psi)^{-1}$.
\end{enumerate}
\end{lemma}

\begin{proof}
Démontrons (1). Posons $\psi = 0$. Avec les notations
précédant la définition \ref{chow-definition-periodic-determinant}, nous pouvons alors identifier
$K_\varphi = I_\psi = 0$, $I_\varphi = K_\psi = M$.
Sous ces identifications, l'application
$$
\gamma_\varphi :
\kappa \otimes \det\nolimits_\kappa(M)
=
\det\nolimits_\kappa(K_\varphi)
\otimes
\det\nolimits_\kappa(I_\varphi)
\longrightarrow
\det\nolimits_\kappa(M)
$$
s'identifie à $\det_\kappa(\varphi^{-1})$. D'autre part,
l'application $\gamma_\psi$ s'identifie à l'application identité. Par conséquent,
$\gamma_\psi \circ \sigma \circ \gamma_\varphi^{-1}$ est égale
à $\det_\kappa(\varphi)$ dans ce cas. D'où le résultat.
Nous omettons la démonstration de (2).
\end{proof}

\begin{lemma}
\label{chow-lemma-periodic-determinant}
Soit $R$ un anneau local de corps résiduel $\kappa$.
Supposons donnée une suite exacte courte de
complexes $(2, 1)$-périodiques
$$
0 \to (M_1, \varphi_1, \psi_1)
\to (M_2, \varphi_2, \psi_2)
\to (M_3, \varphi_3, \psi_3)
\to 0
$$
où tous les $M_i$ sont de longueur finie, et chaque $(M_1, \varphi_1, \psi_1)$ est exact.
Alors
$$
\det\nolimits_\kappa(M_2, \varphi_2, \psi_2) =
\det\nolimits_\kappa(M_1, \varphi_1, \psi_1)
\det\nolimits_\kappa(M_3, \varphi_3, \psi_3).
$$
. dans $\kappa^*$.
\end{lemma}

\begin{proof}
Abrégeons
$I_{\varphi, i} = \Im(\varphi_i)$,
$K_{\varphi, i} = \Ker(\varphi_i)$,
$I_{\psi, i} = \Im(\psi_i)$, et
$K_{\psi, i} = \Ker(\psi_i)$.
Remarquons que nous avons un carré commutatif
$$
\xymatrix{
& 0 \ar[d] & 0 \ar[d] & 0 \ar[d] & \\
0 \ar[r] &
K_{\varphi, 1} \ar[r] \ar[d] &
K_{\varphi, 2} \ar[r] \ar[d] &
K_{\varphi, 3} \ar[r] \ar[d] &
0 \\
0 \ar[r] &
M_1 \ar[r] \ar[d] &
M_2 \ar[r] \ar[d] &
M_3 \ar[r] \ar[d] &
0 \\
0 \ar[r] &
I_{\varphi, 1} \ar[r] \ar[d] &
I_{\varphi, 2} \ar[r] \ar[d] &
I_{\varphi, 3} \ar[r] \ar[d] &
0 \\
& 0  & 0  & 0  &
}
$$
de $R$-modules de longueur finie, à lignes et colonnes exactes.
La ligne supérieure est exacte puisqu'elle s'identifie à la
suite $I_{\psi, 1} \to I_{\psi, 2} \to I_{\psi, 3} \to 0$
des images, et de même pour la ligne inférieure. Il existe un diagramme analogue
faisant intervenir les modules $I_{\psi, i}$ et $K_{\psi, i}$.
Par définition, $\det_\kappa(M_2, \varphi_2, \psi_2)$
correspond, au signe près, à la composée des applications verticales de gauche
dans le diagramme suivant
$$
\xymatrix{
\det_\kappa(M_1) \otimes
\det_\kappa(M_3) \ar[r]^\gamma
\ar[d]^{\gamma^{-1} \otimes \gamma^{-1}} &
\det_\kappa(M_2)
\ar[d]^{\gamma^{-1}} \\
\det\nolimits_\kappa(K_{\varphi, 1})
\otimes
\det\nolimits_\kappa(I_{\varphi, 1})
\otimes
\det\nolimits_\kappa(K_{\varphi, 3})
\otimes
\det\nolimits_\kappa(I_{\varphi, 3})
\ar[d]^{\sigma \otimes \sigma}
\ar[r]^-{\gamma \otimes \gamma} &
\det\nolimits_\kappa(K_{\varphi, 2})
\otimes
\det\nolimits_\kappa(I_{\varphi, 2})
\ar[d]^\sigma
\\
\det\nolimits_\kappa(K_{\psi, 1})
\otimes
\det\nolimits_\kappa(I_{\psi, 1})
\otimes
\det\nolimits_\kappa(K_{\psi, 3})
\otimes
\det\nolimits_\kappa(I_{\psi, 3})
\ar[d]^{\gamma \otimes \gamma}
\ar[r]^-{\gamma \otimes \gamma}
&
\det\nolimits_\kappa(K_{\psi, 2})
\otimes
\det\nolimits_\kappa(I_{\psi, 2})
\ar[d]^\gamma \\
\det_\kappa(M_1)
\otimes
\det_\kappa(M_3) \ar[r]^\gamma
&
\det_\kappa(M_2)
}
$$
Les carrés supérieur et inférieur sont commutatifs au signe près
en appliquant le lemme \ref{chow-lemma-det-exact-sequences} (2).
Le carré du milieu est trivialement
commutatif (nous ne faisons que permuter les facteurs). Nous voyons donc
que
$\det\nolimits_\kappa(M_2, \varphi_2, \psi_2) =
\epsilon \det\nolimits_\kappa(M_1, \varphi_1, \psi_1)
\det\nolimits_\kappa(M_3, \varphi_3, \psi_3)
$
pour un certain signe $\epsilon$. Ce signe se calcule :
le rectangle extérieur du diagramme ci-dessus commute au signe près
\begin{eqnarray*}
\epsilon & = &
(-1)^{\text{longueur}(I_{\varphi, 1})\text{longueur}(K_{\varphi, 3})
+ \text{longueur}(I_{\psi, 1})\text{longueur}(K_{\psi, 3})} \\
& = &
(-1)^{\text{longueur}(I_{\varphi, 1})\text{longueur}(I_{\psi, 3})
+ \text{longueur}(I_{\psi, 1})\text{longueur}(I_{\varphi, 3})}
\end{eqnarray*}
(démonstration omise). Il en résulte aisément que les signes
conviennent également.
\end{proof}

\begin{example}
\label{chow-example-dual-numbers}
Soit $k$ un corps.
Considérons l'anneau $R = k[T]/(T^2)$ des nombres duaux sur $k$.
Notons $t$ la classe de $T$ dans $R$.
Soient $M = R$, $\varphi = ut$ et $\psi = vt$, où $u, v \in k^*$.
Dans ce cas, $\det_k(M)$ a pour générateur $e = [t, 1]$.
Nous identifions $I_\varphi = K_\varphi = I_\psi = K_\psi = (t)$.
Alors $\gamma_\varphi(t \otimes t) = u^{-1}[t, 1]$
(puisque $u^{-1} \in M$ est un relèvement de $t \in I_\varphi$)
et $\gamma_\psi(t \otimes t) = v^{-1}[t, 1]$ (pour la même raison).
Nous voyons donc que $\det_k(M, \varphi, \psi) = -u/v \in k^*$.
\end{example}

\begin{example}
\label{chow-example-Zp}
Soient $R = \mathbf{Z}_p$ et $M = \mathbf{Z}_p/(p^l)$.
Soient $\varphi = p^b u$ et $\varphi = p^a v$, où $a, b \geq 0$,
$a + b = l$ et $u, v \in \mathbf{Z}_p^*$.
Alors un calcul analogue à celui de l'exemple \ref{chow-example-dual-numbers}
montre que
\begin{eqnarray*}
\det\nolimits_{\mathbf{F}_p}(\mathbf{Z}_p/(p^l), p^bu, p^av) & = &
(-1)^{ab}u^a/v^b \bmod p \\
& = &
(-1)^{\text{ord}_p(\alpha)\text{ord}_p(\beta)}
\frac{\alpha^{\text{ord}_p(\beta)}}{\beta^{\text{ord}_p(\alpha)}} \bmod p
\end{eqnarray*}
où $\alpha = p^bu, \beta = p^av \in \mathbf{Z}_p$.
Voir le lemme \ref{chow-lemma-symbol-is-usual-tame-symbol}
pour un cas plus général (et une démonstration).
\end{example}

\begin{example}
\label{chow-example-generic-vector-space}
Soit $R = k$ un corps.
Soit $M = k^{\oplus a} \oplus k^{\oplus b}$, de dimension $l = a + b$.
Soient $\varphi$ et $\psi$ les matrices diagonales suivantes
$$
\varphi = \text{diag}(u_1, \ldots, u_a, 0, \ldots, 0),
\quad
\psi = \text{diag}(0, \ldots, 0, v_1, \ldots, v_b)
$$
où $u_i, v_j \in k^*$. Dans ce cas, nous avons
$$
\det\nolimits_k(M, \varphi, \psi)
=
\frac{u_1 \ldots u_a}{v_1 \ldots v_b}.
$$
On peut le voir par un calcul direct, ou en calculant le cas $l = 1$
et en utilisant l'additivité du lemme \ref{chow-lemma-periodic-determinant}.
\end{example}

\begin{example}
\label{chow-example-special-vector-space}
Soit $R = k$ un corps.
Soit $M = k^{\oplus a} \oplus k^{\oplus a}$, de dimension $l = 2a$.
Soient $\varphi$ et $\psi$ les matrices par blocs suivantes
$$
\varphi =
\left(
\begin{matrix}
0 & U \\
0 & 0
\end{matrix}
\right),
\quad
\psi =
\left(
\begin{matrix}
0 & V \\
0 & 0
\end{matrix}
\right),
$$
où $U, V \in \text{Mat}(a \times a, k)$ sont inversibles.
Dans ce cas, nous avons
$$
\det\nolimits_k(M, \varphi, \psi)
=
(-1)^a\frac{\det(U)}{\det(V)}.
$$
On peut le voir par un calcul direct.
Le cas $a = 1$ est analogue au calcul de
l'exemple \ref{chow-example-dual-numbers}.
\end{example}

\begin{example}
\label{chow-example-a-la-oort}
Soit $R = k$ un corps.
Soit $M = k^{\oplus 4}$.
Soit
$$
\varphi =
\left(
\begin{matrix}
  0 &   0 &   0 &   0 \\
u_1 &   0 &   0 &   0 \\
  0 &   0 &   0 &   0 \\
  0 &   0 & u_2 &   0
\end{matrix}
\right)
\quad
\varphi =
\left(
\begin{matrix}
  0 &   0 &   0 &   0 \\
  0 &   0 & v_2 &   0 \\
  0 &   0 &   0 &   0 \\
v_1 &   0 &   0 &   0
\end{matrix}
\right)
\quad
$$
où $u_1, u_2, v_1, v_2 \in k^*$.
Alors nous avons
$$
\det\nolimits_k(M, \varphi, \psi) = -\frac{u_1u_2}{v_1v_2}.
$$
\end{example}

\noindent
Nous en venons maintenant à l'analogue du fait que le déterminant
d'une composée d'endomorphismes linéaires est le produit des
déterminants. Pour éviter des formules très longues, nous
écrivons $I_\varphi = \Im(\varphi)$ et
$K_\varphi = \Ker(\varphi)$
pour toute application de $R$-modules $\varphi : M \to M$.
Nous notons également $\varphi\psi = \varphi \circ \psi$
pour une paire de morphismes $\varphi, \psi : M \to M$.

\begin{lemma}
\label{chow-lemma-multiplicativity-determinant}
Soit $R$ un anneau local de corps résiduel $\kappa$.
Soit $M$ un $R$-module de longueur finie.
Soient $\alpha, \beta, \gamma$ des endomorphismes de $M$.
Supposons que
\begin{enumerate}
\item $I_\alpha = K_{\beta\gamma}$, et de même pour toute permutation
de $\alpha, \beta, \gamma$,
\item $K_\alpha = I_{\beta\gamma}$, et de même pour toute permutation
de $\alpha, \beta, \gamma$.
\end{enumerate}
Alors
\begin{enumerate}
\item le triplet $(M, \alpha, \beta\gamma)$
est un complexe $(2, 1)$-périodique exact ;
\item le triplet $(I_\gamma, \alpha, \beta)$
est un complexe $(2, 1)$-périodique exact ;
\item le triplet $(M/K_\beta, \alpha, \gamma)$
est un complexe $(2, 1)$-périodique exact ;
\item nous avons
$$
\det\nolimits_\kappa(M, \alpha, \beta\gamma)
=
\det\nolimits_\kappa(I_\gamma, \alpha, \beta)
\det\nolimits_\kappa(M/K_\beta, \alpha, \gamma).
$$
\end{enumerate}
\end{lemma}

\begin{proof}
Il est clair que les hypothèses impliquent la partie (1) du lemme.

\medskip\noindent
Pour voir la partie (1), remarquons que les hypothèses impliquent que
$I_{\gamma\alpha} = I_{\alpha\gamma}$, et de même pour les noyaux
et pour toute autre paire de morphismes.
De plus, nous voyons que
$I_{\gamma\beta} =I_{\beta\gamma} = K_\alpha \subset I_\gamma$, et
de même pour toute autre paire. En particulier, nous obtenons une suite exacte courte
$$
0 \to I_{\beta\gamma} \to I_\gamma \xrightarrow{\alpha} I_{\alpha\gamma} \to 0
$$
et, de même, nous obtenons une suite exacte courte
$$
0 \to I_{\alpha\gamma} \to I_\gamma \xrightarrow{\beta} I_{\beta\gamma} \to 0.
$$
Cela démontre que $(I_\gamma, \alpha, \beta)$ est un complexe $(2, 1)$-périodique
exact. La partie (2) du lemme est donc vraie.

\medskip\noindent
Pour voir que $\alpha$ et $\gamma$ induisent des endomorphismes bien définis
de $M/K_\beta$, nous devons vérifier que $\alpha(K_\beta) \subset K_\beta$
et $\gamma(K_\beta) \subset K_\beta$. Cela est vrai car
$\alpha(K_\beta) = \alpha(I_{\gamma\alpha}) = I_{\alpha\gamma\alpha}
\subset I_{\alpha\gamma} = K_\beta$, et de même dans l'autre cas.
Le noyau de l'application $\alpha : M/K_\beta \to M/K_\beta$ est
$K_{\beta\alpha}/K_\beta = I_\gamma/K_\beta$. De même,
le noyau de $\gamma : M/K_\beta \to M/K_\beta$ est égal à
$I_\alpha/K_\beta$. Nous concluons donc que (3) est vraie.

\medskip\noindent
Introduisons $r = \text{longueur}_R(K_\alpha)$,
$s = \text{longueur}_R(K_\beta)$ et $t = \text{longueur}_R(K_\gamma)$.
D'après les suites exactes ci-dessus et nos hypothèses, nous avons
$\text{longueur}_R(I_\alpha) = s + t$, $\text{longueur}_R(I_\beta) = r + t$,
$\text{longueur}_R(I_\gamma) = r + s$, et
$\text{longueur}(M) = r + s + t$.
Choisissons
\begin{enumerate}
\item une suite admissible $x_1, \ldots, x_r \in K_\alpha$
engendrant $K_\alpha$,
\item une suite admissible $y_1, \ldots, y_s \in K_\beta$
engendrant $K_\beta$,
\item une suite admissible $z_1, \ldots, z_t \in K_\gamma$
engendrant $K_\gamma$,
\item des éléments $\tilde x_i \in M$ tels que $\beta\gamma\tilde x_i = x_i$,
\item des éléments $\tilde y_i \in M$ tels que $\alpha\gamma\tilde y_i = y_i$,
\item des éléments $\tilde z_i \in M$ tels que $\beta\alpha\tilde z_i = z_i$.
\end{enumerate}
Avec ces choix, la suite
$y_1, \ldots, y_s, \alpha\tilde z_1, \ldots, \alpha\tilde z_t$
est une suite admissible dans $I_\alpha$ qui l'engendre.
Ainsi, d'après la remarque \ref{chow-remark-more-elementary}, le déterminant
$D = \det_\kappa(M, \alpha, \beta\gamma)$ est
l'unique élément de $\kappa^*$ tel que
\begin{align*}
[y_1, \ldots, y_s,
\alpha\tilde z_1, \ldots, \alpha\tilde z_s,
\tilde x_1, \ldots, \tilde x_r] \\
= (-1)^{r(s + t)} D
[x_1, \ldots, x_r,
\gamma\tilde y_1, \ldots, \gamma\tilde y_s,
\tilde z_1, \ldots, \tilde z_t]
\end{align*}
D'après la même remarque, nous voyons que
$D_1 = \det_\kappa(M/K_\beta, \alpha, \gamma)$
est caractérisé par
$$
[y_1, \ldots, y_s,
\alpha\tilde z_1, \ldots, \alpha\tilde z_t,
\tilde x_1, \ldots, \tilde x_r]
=
(-1)^{rt} D_1
[y_1, \ldots, y_s,
\gamma\tilde x_1, \ldots, \gamma\tilde x_r,
\tilde z_1, \ldots, \tilde z_t]
$$
D'après la même remarque, nous voyons que
$D_2 = \det_\kappa(I_\gamma, \alpha, \beta)$ est caractérisé par
$$
[y_1, \ldots, y_s,
\gamma\tilde x_1, \ldots, \gamma\tilde x_r,
\tilde z_1, \ldots, \tilde z_t]
=
(-1)^{rs} D_2
[x_1, \ldots, x_r,
\gamma\tilde y_1, \ldots, \gamma\tilde y_s,
\tilde z_1, \ldots, \tilde z_t]
$$
En combinant les formules ci-dessus, nous voyons que $D = D_1 D_2$,
comme voulu.
\end{proof}


\begin{lemma}
\label{chow-lemma-tricky}
Soit $R$ un anneau local de corps résiduel $\kappa$.
Soit $\alpha : (M, \varphi, \psi) \to (M', \varphi', \psi')$
un morphisme de complexes $(2, 1)$-périodiques sur $R$.
Supposons que
\begin{enumerate}
\item $M$ et $M'$ soient de longueur finie ;
\item $(M, \varphi, \psi)$ et $(M', \varphi', \psi')$ soient exacts ;
\item les applications $\varphi$ et $\psi$ induisent l'application nulle sur
$K = \Ker(\alpha)$ ; et
\item les applications $\varphi$ et $\psi$ induisent l'application nulle sur
$Q = \Coker(\alpha)$.
\end{enumerate}
Notons $N = \alpha(M) \subset M'$. Nous obtenons deux suites exactes courtes
de complexes $(2, 1)$-périodiques
$$
\begin{matrix}
0 \to (N, \varphi', \psi') \to (M', \varphi', \psi') \to (Q, 0, 0) \to 0 \\
0 \to (K, 0, 0) \to (M, \varphi, \psi) \to (N, \varphi', \psi') \to 0
\end{matrix}
$$
qui induisent deux isomorphismes $\alpha_i : Q \to K$, $i = 0, 1$. Alors
$$
\det\nolimits_\kappa(M, \varphi, \psi)
=
\det\nolimits_\kappa(\alpha_0^{-1} \circ \alpha_1)
\det\nolimits_\kappa(M', \varphi', \psi')
$$
En particulier, si $\alpha_0 = \alpha_1$, alors
$\det\nolimits_\kappa(M, \varphi, \psi) =
\det\nolimits_\kappa(M', \varphi', \psi')$.
\end{lemma}

\begin{proof}
Il existe (au moins) deux façons de démontrer ce lemme. L'une consiste à produire un
énorme diagramme commutatif en utilisant les propriétés des déterminants.
L'autre consiste à employer la caractérisation des déterminants au moyen
de suites admissibles d'éléments. C'est la seconde approche que nous
utiliserons.

\medskip\noindent
Expliquons d'abord précisément ce que sont les applications $\alpha_i$.
À savoir, $\alpha_0$ est la composée
$$
\alpha_0 : Q = H^0(Q, 0, 0) \to H^1(N, \varphi', \psi') \to H^2(K, 0, 0) = K
$$
et $\alpha_1$ est la composée
$$
\alpha_1 : Q = H^1(Q, 0, 0) \to H^2(N, \varphi', \psi') \to H^3(K, 0, 0) = K
$$
provenant des morphismes de bord des suites exactes courtes de complexes
affichées dans le lemme. Le fait que les
complexes $(M, \varphi, \psi)$ et $(M', \varphi', \psi')$ soient exacts
implique que ces applications sont des isomorphismes.

\medskip\noindent
Nous utiliserons les notations $I_\varphi = \Im(\varphi)$,
$K_\varphi = \Ker(\varphi)$, et de même pour les autres applications.
L'exactitude pour $M$ et $M'$
signifie que $K_\varphi = I_\psi$, ainsi que trois égalités analogues.
Introduisons $k = \text{longueur}_R(K)$, $a = \text{longueur}_R(I_\varphi)$
et $b = \text{longueur}_R(I_\psi)$. Nous voyons alors que $\text{longueur}_R(M) = a + b$,
ainsi que $\text{longueur}_R(N) = a + b - k$, $\text{longueur}_R(Q) = k$
et $\text{longueur}_R(M') = a + b$. Les suites exactes ci-dessous montreront
qu'également $\text{longueur}_R(I_{\varphi'}) = a$ et
$\text{longueur}_R(I_{\psi'}) = b$.

\medskip\noindent
L'hypothèse $K \subset K_\varphi = I_\psi$ signifie que
$\varphi$ se factorise par $N$ pour donner une suite exacte
$$
0 \to \alpha(I_\psi) \to N \xrightarrow{\varphi\alpha^{-1}} I_\psi \to 0.
$$
Ici, $\varphi\alpha^{-1}(x') = y$ signifie que $x' = \alpha(x)$ et $y = \varphi(x)$.
De même, nous avons
$$
0 \to \alpha(I_\varphi) \to N \xrightarrow{\psi\alpha^{-1}} I_\varphi \to 0.
$$
L'hypothèse selon laquelle $\psi'$ induit l'application nulle sur
$Q$ signifie que $I_{\psi'} = K_{\varphi'} \subset N$.
Cela signifie que le quotient $\varphi'(N) \subset I_{\varphi'}$
s'identifie à $Q$. Remarquons que $\varphi'(N) = \alpha(I_\varphi)$.
Nous en concluons donc qu'il existe un isomorphisme
$$
\varphi' : Q \to I_{\varphi'}/\alpha(I_\varphi)
$$
décrit simplement par
$\varphi'(x' \bmod N) = \varphi'(x') \bmod \alpha(I_\varphi)$.
Exactement de la même manière, nous obtenons
$$
\psi' : Q \to I_{\psi'}/\alpha(I_\psi)
$$
Enfin, remarquons que $\alpha_0$ est la composée
$$
\xymatrix{
Q \ar[r]^-{\varphi'} &
I_{\varphi'}/\alpha(I_\varphi)
\ar[rrr]^-{\psi\alpha^{-1}|_{I_{\varphi'}/\alpha(I_\varphi)}} & & &
K
}
$$
et, de même,
$\alpha_1 = \varphi\alpha^{-1}|_{I_{\psi'}/\alpha(I_\psi)} \circ \psi'$.

\medskip\noindent
Pour abréger les formules ci-dessous, nous écrirons désormais $\alpha x$
au lieu de $\alpha(x)$. Il ne devrait en résulter aucune confusion puisque
toutes les applications sont désignées par des lettres grecques et les éléments par des lettres romaines.
Nous allons choisir
\begin{enumerate}
\item une suite admissible $z_1, \ldots, z_k \in K$
engendrant $K$,
\item des éléments $z'_i \in M$ tels que $\varphi z'_i = z_i$,
\item des éléments $z''_i \in M$ tels que $\psi z''_i = z_i$,
\item des éléments $x_{k + 1}, \ldots, x_a \in I_\varphi$ tels
que $z_1, \ldots, z_k, x_{k + 1}, \ldots, x_a$ soit une suite admissible
engendrant $I_\varphi$,
\item des éléments $\tilde x_i \in M$ tels que $\varphi \tilde x_i = x_i$,
\item des éléments $y_{k + 1}, \ldots, y_b \in I_\psi$ tels que
$z_1, \ldots, z_k, y_{k + 1}, \ldots, y_b$ soit une suite admissible
engendrant $I_\psi$,
\item des éléments $\tilde y_i \in M$ tels que $\psi \tilde y_i = y_i$, et
\item des éléments $w_1, \ldots, w_k \in M'$ tels que
$w_1 \bmod N, \ldots, w_k \bmod N$ forment une suite admissible
dans $Q$ qui engendre $Q$.
\end{enumerate}
D'après la remarque \ref{chow-remark-more-elementary}, l'élément
$D = \det_\kappa(M, \varphi, \psi) \in \kappa^*$ est
caractérisé par
\begin{eqnarray*}
& &
[z_1, \ldots, z_k,
x_{k + 1}, \ldots, x_a,
z''_1, \ldots, z''_k,
\tilde y_{k + 1}, \ldots, \tilde y_b] \\
& = &
(-1)^{ab} D
[z_1, \ldots, z_k,
y_{k + 1}, \ldots, y_b,
z'_1, \ldots, z'_k,
\tilde x_{k + 1}, \ldots, \tilde x_a]
\end{eqnarray*}
Remarquons que, d'après la discussion ci-dessus,
$\alpha x_{k + 1}, \ldots, \alpha x_a, \varphi w_1, \ldots, \varphi w_k$
est une suite admissible engendrant $I_{\varphi'}$, et que
$\alpha y_{k + 1}, \ldots, \alpha y_b, \psi w_1, \ldots, \psi w_k$
est une suite admissible engendrant $I_{\psi'}$.
Ainsi, d'après la remarque \ref{chow-remark-more-elementary}, l'élément
$D' = \det_\kappa(M', \varphi', \psi') \in \kappa^*$ est
caractérisé par
\begin{eqnarray*}
& &
[\alpha x_{k + 1}, \ldots, \alpha x_a,
\varphi' w_1, \ldots, \varphi' w_k,
\alpha \tilde y_{k + 1}, \ldots, \alpha \tilde y_b,
w_1, \ldots, w_k]
\\
& = &
(-1)^{ab} D'
[\alpha y_{k + 1}, \ldots, \alpha y_b,
\psi' w_1, \ldots, \psi' w_k,
\alpha \tilde x_{k + 1}, \ldots, \alpha \tilde x_a,
w_1, \ldots, w_k]
\end{eqnarray*}
Remarquons que, dans la première formule affichée, resp.\ la seconde,
les $k$ premières, resp.\ dernières, entrées des symboles des deux côtés
sont les mêmes. Ces formules sont donc réellement équivalentes aux
égalités
\begin{eqnarray*}
& &
[\alpha x_{k + 1}, \ldots, \alpha x_a,
\alpha z''_1, \ldots, \alpha z''_k,
\alpha \tilde y_{k + 1}, \ldots, \alpha \tilde y_b] \\
& = &
(-1)^{ab} D
[\alpha y_{k + 1}, \ldots, \alpha y_b,
\alpha z'_1, \ldots, \alpha z'_k,
\alpha \tilde x_{k + 1}, \ldots, \alpha \tilde x_a]
\end{eqnarray*}
et
\begin{eqnarray*}
& &
[\alpha x_{k + 1}, \ldots, \alpha x_a,
\varphi' w_1, \ldots, \varphi' w_k,
\alpha \tilde y_{k + 1}, \ldots, \alpha \tilde y_b]
\\
& = &
(-1)^{ab} D'
[\alpha y_{k + 1}, \ldots, \alpha y_b,
\psi' w_1, \ldots, \psi' w_k,
\alpha \tilde x_{k + 1}, \ldots, \alpha \tilde x_a]
\end{eqnarray*}
dans $\det_\kappa(N)$. Remarquons que
$\varphi' w_1, \ldots, \varphi' w_k$
et
$\alpha z''_1, \ldots, z''_k$
sont des suites admissibles engendrant le module
$I_{\varphi'}/\alpha(I_\varphi)$. Écrivons
$$
[\varphi' w_1, \ldots, \varphi' w_k]
= \lambda_0 [\alpha z''_1, \ldots, \alpha z''_k]
$$
dans $\det_\kappa(I_{\varphi'}/\alpha(I_\varphi))$
pour un certain $\lambda_0 \in \kappa^*$. De même,
écrivons
$$
[\psi' w_1, \ldots, \psi' w_k]
= \lambda_1 [\alpha z'_1, \ldots, \alpha z'_k]
$$
dans $\det_\kappa(I_{\psi'}/\alpha(I_\psi))$
pour un certain $\lambda_1 \in \kappa^*$. D'une part,
il est clair que
$$
\alpha_i([w_1, \ldots, w_k]) = \lambda_i[z_1, \ldots, z_k]
$$
pour $i = 0, 1$, d'après notre description de $\alpha_i$ ci-dessus,
ce qui signifie que
$$
\det\nolimits_\kappa(\alpha_0^{-1} \circ \alpha_1)
=
\lambda_1/\lambda_0
$$
et,
d'autre part, il est clair que
\begin{eqnarray*}
& &
\lambda_0 [\alpha x_{k + 1}, \ldots, \alpha x_a,
\alpha z''_1, \ldots, \alpha z''_k,
\alpha \tilde y_{k + 1}, \ldots, \alpha \tilde y_b] \\
& = &
[\alpha x_{k + 1}, \ldots, \alpha x_a,
\varphi' w_1, \ldots, \varphi' w_k,
\alpha \tilde y_{k + 1}, \ldots, \alpha \tilde y_b]
\end{eqnarray*}
et
\begin{eqnarray*}
& &
\lambda_1[\alpha y_{k + 1}, \ldots, \alpha y_b,
\alpha z'_1, \ldots, \alpha z'_k,
\alpha \tilde x_{k + 1}, \ldots, \alpha \tilde x_a] \\
& = &
[\alpha y_{k + 1}, \ldots, \alpha y_b,
\psi' w_1, \ldots, \psi' w_k,
\alpha \tilde x_{k + 1}, \ldots, \alpha \tilde x_a]
\end{eqnarray*}
ce qui implique $\lambda_0 D = \lambda_1 D'$. Le lemme en résulte.
\end{proof}








\subsection{Symboles}
\label{chow-subsection-symbols}

\noindent
Le bon degré de généralité pour cette construction est peut-être la
situation du lemme suivant.

\begin{lemma}
\label{chow-lemma-pre-symbol}
Soit $A$ un anneau local noethérien.
Soit $M$ un $A$-module fini de dimension $1$.
Supposons que $\varphi, \psi : M \to M$ soient deux applications injectives de
$A$-modules, et supposons que $\varphi(\psi(M)) = \psi(\varphi(M))$,
par exemple si $\varphi$ et $\psi$ commutent.
Alors $\text{longueur}_A(M/\varphi\psi M) < \infty$
et $(M/\varphi\psi M, \varphi, \psi)$ est un complexe
$(2, 1)$-périodique exact.
\end{lemma}

\begin{proof}
Soit $\mathfrak q$ un idéal premier minimal du support de $M$.
Alors $M_{\mathfrak q}$ est un $A_{\mathfrak q}$-module de longueur finie,
voir Algèbre, lemme \ref{algebra-lemma-support-point}.
Ainsi, $\varphi$ et $\psi$
induisent des isomorphismes $M_{\mathfrak q} \to M_{\mathfrak q}$.
Le support de $M/\varphi\psi M$ est donc $\{\mathfrak m_A\}$
et celui-ci est ainsi de longueur finie (voir le lemme cité ci-dessus).
Enfin, le noyau de $\varphi$ sur $M/\varphi\psi M$
est manifestement $\psi M/\varphi\psi M$ ; le noyau
de $\varphi$ est donc l'image de $\psi$ sur $M/\varphi\psi M$.
De même dans l'autre sens, puisque $M/\varphi\psi M = M/\psi\varphi M$
par hypothèse.
\end{proof}

\begin{lemma}
\label{chow-lemma-symbol-defined}
Soit $A$ un anneau local noethérien. Soient $a, b \in A$.
\begin{enumerate}
\item si $M$ est un $A$-module fini de dimension $1$
tel que $a, b$ soient des non-diviseurs de zéro dans $M$, alors
$\text{longueur}_A(M/abM) < \infty$ et
$(M/abM, a, b)$ est un complexe exact $(2, 1)$-périodique ;
\item si $a, b$ sont des non-diviseurs de zéro et $\dim(A) = 1$,
alors $\text{longueur}_A(A/(ab)) < \infty$ et
$(A/(ab), a, b)$ est un complexe exact $(2, 1)$-périodique.
\end{enumerate}
En particulier, dans ces cas,
$\det_\kappa(M/abM, a, b) \in \kappa^*$,
resp.\ $\det_\kappa(A/(ab), a, b) \in \kappa^*$
sont définis.
\end{lemma}

\begin{proof}
Cela résulte du lemme \ref{chow-lemma-pre-symbol}.
\end{proof}

\begin{definition}
\label{chow-definition-symbol-M}
Soit $A$ un anneau local noethérien de corps résiduel $\kappa$.
Soient $a, b \in A$.
Soit $M$ un $A$-module fini de dimension $1$
tel que $a, b$ soient des non-diviseurs de zéro dans $M$.
Nous définissons le {\it symbole associé à $M, a, b$}
comme l'élément
$$
d_M(a, b) =
\det\nolimits_\kappa(M/abM, a, b) \in \kappa^*
$$
\end{definition}

\begin{lemma}
\label{chow-lemma-multiplicativity-symbol}
Soit $A$ un anneau local noethérien.
Soient $a, b, c \in A$. Soit $M$ un $A$-module fini
tel que $\dim(\text{Supp}(M)) = 1$. Supposons que $a, b, c$ soient des non-diviseurs de zéro dans $M$.
Alors
$$
d_M(a, bc) = d_M(a, b) d_M(a, c)
$$
et $d_M(a, b)d_M(b, a) = 1$.
\end{lemma}

\begin{proof}
La première assertion résulte du lemme \ref{chow-lemma-multiplicativity-determinant}
appliqué à $M/abcM$ et aux endomorphismes $\alpha, \beta, \gamma$ donnés par
la multiplication par $a, b, c$.
La seconde provient du lemme \ref{chow-lemma-periodic-determinant-shift}.
\end{proof}

\begin{definition}
\label{chow-definition-tame-symbol}
Soit $A$ un anneau intègre local noethérien de dimension $1$
et de corps résiduel $\kappa$.
Soit $K$ le corps des fractions de $A$.
Nous définissons le {\it symbole modéré} de $A$ comme l'application
$$
K^* \times K^* \longrightarrow \kappa^*,
\quad
(x, y) \longmapsto d_A(x, y)
$$
où $d_A(x, y)$ est étendu à $K^* \times K^*$ par la multiplicativité du
lemme \ref{chow-lemma-multiplicativity-symbol}.
\end{definition}

\noindent
Il est clair que nous pouvons étendre plus généralement $d_M(-, -)$ à
certains anneaux de fractions de $A$ (même si $A$ n'est pas intègre).

\begin{lemma}
\label{chow-lemma-symbol-when-equal}
Soient $A$ un anneau local noethérien et $M$ un $A$-module fini de
dimension $1$. Soit $a \in A$ un non-diviseur de zéro dans $M$.
Alors $d_M(a, a) = (-1)^{\text{longueur}_A(M/aM)}$.
\end{lemma}

\begin{proof}
C'est une conséquence immédiate du lemme \ref{chow-lemma-periodic-determinant-sign}.
\end{proof}

\begin{lemma}
\label{chow-lemma-symbol-when-one-is-a-unit}
Soit $A$ un anneau local noethérien.
Soit $M$ un $A$-module fini de dimension $1$.
Soient $b \in A$ un non-diviseur de zéro dans $M$ et $u \in A^*$.
Alors
$$
d_M(u, b) = u^{\text{longueur}_A(M/bM)} \bmod \mathfrak m_A.
$$
En particulier, si $M = A$, alors
$d_A(u, b) = u^{\text{ord}_A(b)} \bmod \mathfrak m_A$.
\end{lemma}

\begin{proof}
Remarquons que, dans ce cas, $M/ubM = M/bM$, sur lequel la multiplication
par $b$ est nulle. Ainsi, $d_M(u, b) = \det_\kappa(u|_{M/bM})$
d'après le lemme \ref{chow-lemma-periodic-determinant-easy-case}. Le lemme
résulte alors du lemme \ref{chow-lemma-times-u-determinant}.
\end{proof}

\begin{lemma}
\label{chow-lemma-symbol-short-exact-sequence}
Soit $A$ un anneau local noethérien.
Soient $a, b \in A$.
Soit
$$
0 \to M \to M' \to M'' \to 0
$$
une suite exacte courte de $A$-modules de dimension $1$
tels que $a, b$ soient des non-diviseurs de zéro dans
les trois $A$-modules.
Alors
$$
d_{M'}(a, b) = d_M(a, b) d_{M''}(a, b)
$$
dans $\kappa^*$.
\end{lemma}

\begin{proof}
Il est facile de voir que cela donne une suite exacte courte
de complexes $(2, 1)$-périodiques exacts
$$
0 \to
(M/abM, a, b) \to
(M'/abM', a, b) \to
(M''/abM'', a, b) \to 0
$$
Le lemme résulte donc du lemme \ref{chow-lemma-periodic-determinant}.
\end{proof}

\begin{lemma}
\label{chow-lemma-symbol-compare-modules}
Soit $A$ un anneau local noethérien.
Soit $\alpha : M \to M'$ un homomorphisme de
$A$-modules finis de dimension $1$.
Soient $a, b \in A$. Supposons que
\begin{enumerate}
\item $a$ et $b$ soient des non-diviseurs de zéro dans $M$ et $M'$ ; et
\item $\dim(\Ker(\alpha)), \dim(\Coker(\alpha)) \leq 0$.
\end{enumerate}
Alors $d_M(a, b) = d_{M'}(a, b)$.
\end{lemma}

\begin{proof}
Si $a \in A^*$, l'égalité résulte de
l'égalité $\text{longueur}(M/bM) = \text{longueur}(M'/bM')$
et du lemme \ref{chow-lemma-symbol-when-one-is-a-unit}.
De même, si $b$ est une unité, le lemme est encore vrai
(par la symétrie du lemme \ref{chow-lemma-multiplicativity-symbol}).
Nous pouvons donc supposer que $a, b \in \mathfrak m_A$.
Cela implique en particulier que $\mathfrak m$ n'est pas
un idéal premier associé à $M$, et donc que $\alpha : M \to M'$
est injectif. Nous pouvons ainsi considérer $M$ comme un sous-module de $M'$.
Par hypothèse, $M'/M$ est un $A$-module fini de support
$\{\mathfrak m_A\}$, et il est donc de longueur finie.
Remarquons que, pour tout troisième module $M''$ tel que $M \subset M'' \subset M'$,
les applications $M \to M''$ et $M'' \to M'$ satisfont elles aussi les hypothèses du lemme.
Cela nous ramène, par récurrence sur la longueur de $M'/M$,
au cas où $\text{longueur}_A(M'/M) = 1$.
Enfin, dans ce cas, considérons l'application
$$
\overline{\alpha} : M/abM \longrightarrow M'/abM'.
$$
Par construction, le conoyau $Q$ de $\overline{\alpha}$ est de
longueur $1$. Puisque $a, b \in \mathfrak m_A$, ils agissent trivialement sur
$Q$. Il s'ensuit aussi que le noyau $K$ de $\overline{\alpha}$ est de
longueur $1$, et donc que $a$ et $b$ agissent également trivialement sur $K$.
Nous pouvons ainsi appliquer le lemme \ref{chow-lemma-tricky}. Il suffit donc de voir
que les deux applications $\alpha_i : Q \to K$ sont les mêmes.
En fait, elles sont toutes deux égales à l'application
$q = x' \bmod \Im(\overline{\alpha}) \mapsto abx' \in K$.
Nous omettons la vérification.
\end{proof}

\begin{lemma}
\label{chow-lemma-compute-symbol-M}
Soit $A$ un anneau local noethérien.
Soit $M$ un $A$-module fini tel que $\dim(\text{Supp}(M)) = 1$.
Soient $a, b \in A$ des non-diviseurs de zéro dans $M$.
Soient $\mathfrak q_1, \ldots, \mathfrak q_t$ les idéaux premiers minimaux
du support de $M$. Alors
$$
d_M(a, b)
=
\prod\nolimits_{i = 1, \ldots, t}
d_{A/\mathfrak q_i}(a, b)^{
\text{longueur}_{A_{\mathfrak q_i}}(M_{\mathfrak q_i})}
$$
comme éléments de $\kappa^*$.
\end{lemma}

\begin{proof}
Choisissons une filtration par des $A$-sous-modules
$$
0 = M_0 \subset M_1 \subset \ldots \subset M_n = M
$$
telle que chaque quotient $M_j/M_{j - 1}$ soit isomorphe
à $A/\mathfrak p_j$ pour un certain idéal premier $\mathfrak p_j$
de $A$. Voir Algèbre, lemme \ref{algebra-lemma-filter-Noetherian-module}.
Pour chaque $j$, nous avons soit $\mathfrak p_j = \mathfrak q_i$
pour un certain $i$, soit $\mathfrak p_j = \mathfrak m_A$. De plus,
pour $i$ fixé, le nombre de $j$ tels que
$\mathfrak p_j = \mathfrak q_i$ est égal à
$\text{longueur}_{A_{\mathfrak q_i}}(M_{\mathfrak q_i})$, d'après
Algèbre, lemme \ref{algebra-lemma-filter-minimal-primes-in-support}.
Ainsi, $d_{M_j}(a, b)$ est défini pour chaque $j$, et
$$
d_{M_j}(a, b)
=
\left\{
\begin{matrix}
d_{M_{j - 1}}(a, b) d_{A/\mathfrak q_i}(a, b) &
\text{si} & \mathfrak p_j = \mathfrak q_i \\
d_{M_{j - 1}}(a, b) & \text{si} & \mathfrak p_j = \mathfrak m_A
\end{matrix}
\right.
$$
d'après le lemme \ref{chow-lemma-symbol-short-exact-sequence} dans le premier cas
et le lemme \ref{chow-lemma-symbol-compare-modules} dans le second. D'où le lemme.
\end{proof}

\begin{lemma}
\label{chow-lemma-symbol-is-usual-tame-symbol}
Soit $A$ un anneau de valuation discrète de corps des fractions $K$.
Pour tous $x, y \in K$ non nuls, nous avons
$$
d_A(x, y)
=
(-1)^{\text{ord}_A(x)\text{ord}_A(y)}
\frac{x^{\text{ord}_A(y)}}{y^{\text{ord}_A(x)}} \bmod \mathfrak m_A,
$$
autrement dit, le symbole est égal au symbole modéré usuel.
\end{lemma}

\begin{proof}
Par multiplicativité, il suffit de le démontrer lorsque $x, y \in A$.
Soit $t \in A$ une uniformisante.
Écrivons $x = t^bu$ et $y = t^bv$ pour certains $a, b \geq 0$
et $u, v \in A^*$. Posons $l = a + b$. Alors
$t^{l - 1}, \ldots, t^b$ est une suite admissible dans
$(x)/(xy)$ et $t^{l - 1}, \ldots, t^a$ est une suite admissible
dans $(y)/(xy)$. Ainsi, d'après la remarque \ref{chow-remark-more-elementary},
nous voyons que $d_A(x, y)$ est caractérisé par l'équation
$$
[t^{l - 1}, \ldots, t^b, v^{-1}t^{b - 1}, \ldots, v^{-1}]
=
(-1)^{ab} d_A(x, y)
[t^{l - 1}, \ldots, t^a, u^{-1}t^{a - 1}, \ldots, u^{-1}].
$$
Ainsi, d'après les relations d'admissibilité pour les
symboles $[x_1, \ldots, x_l]$, nous voyons que
$$
d_A(x, y) = (-1)^{ab} u^a/v^b \bmod \mathfrak m_A
$$
comme voulu.
\end{proof}

\begin{lemma}
\label{chow-lemma-symbol-is-steinberg-prepare}
Soit $A$ un anneau local noethérien.
Soient $a, b \in A$.
Soit $M$ un $A$-module fini de dimension $1$ tel que, dans celui-ci,
chacun de $a$, $b$ et $b - a$ soit un non-diviseur de zéro.
Alors
$$
d_M(a, b - a)d_M(b, b) = d_M(b, b - a)d_M(a, b)
$$
dans $\kappa^*$.
\end{lemma}

\begin{proof}
D'après le lemme \ref{chow-lemma-compute-symbol-M}, il suffit de montrer la relation lorsque
$M = A/\mathfrak q$ pour un certain idéal premier $\mathfrak q \subset A$ tel que
$\dim(A/\mathfrak q) = 1$.

\medskip\noindent
Dans le cas $M = A/\mathfrak q$, nous pouvons remplacer $A$ par $A/\mathfrak q$ et
$a, b$ par leurs images dans $A/\mathfrak q$. Nous pouvons donc supposer que $A = M$
et que $A$ est un anneau intègre local noethérien de dimension $1$. La raison en est
que les corps résiduels $\kappa$ de $A$ et de $A/\mathfrak q$ sont
les mêmes, et que, pour tout $A/\mathfrak q$-module $M$, les déterminants
pris sur $A$ ou sur $A/\mathfrak q$ s'identifient canoniquement.
Voir le lemme \ref{chow-lemma-determinant-quotient-ring}.

\medskip\noindent
Il suffit de montrer la relation lorsque
$a, b$ appartiennent tous deux à l'idéal maximal. En effet, le cas où l'un
d'eux, ou les deux, est une unité résulte des lemmes \ref{chow-lemma-symbol-when-one-is-a-unit}
et \ref{chow-lemma-symbol-when-equal}.

\medskip\noindent
Choisissons une extension $A \subset A'$ et des factorisations
$a = ta'$, $b = tb'$ comme dans le
lemme \ref{chow-lemma-Noetherian-domain-dim-1-two-elements}.
Remarquons qu'aussi $b - a = t(b' - a')$ et que
$A' = (a', b') = (a', b' - a') = (b' - a', b')$.
Ici et dans la suite, nous considérons $A'$ comme un $A$-module, et
$a, b, a', b', t$ comme des endomorphismes du $A$-module $A'$.
Nous emploierons les notations $d^A_{A'}(a', b')$, etc., pour désigner
$$
d^A_{A'}(a', b')
=
\det\nolimits_\kappa(A'/a'b'A', a', b')
$$
qui est défini par le lemme \ref{chow-lemma-pre-symbol}. L'indice supérieur ${}^A$
sert à le distinguer du symbole déjà défini
$d_{A'}(a', b')$, qui est différent (par exemple parce qu'il prend ses valeurs
dans le corps résiduel de $A'$, lequel peut différer de $\kappa$).
D'après le lemme \ref{chow-lemma-symbol-compare-modules}, nous voyons que
$d_A(a, b) = d^A_{A'}(a, b)$,
et de même pour les autres combinaisons.
En utilisant ceci et la multiplicativité, nous voyons qu'il suffit de démontrer
$$
d^A_{A'}(a', b' - a') d^A_{A'}(b', b')
=
d^A_{A'}(b', b' - a') d^A_{A'}(a', b')
$$
Maintenant, puisque $(a', b') = A'$, et de même pour les autres idéaux, nous avons
$$
\begin{matrix}
A'/(a'(b' - a')) & \cong & A'/(a') \oplus A'/(b' - a') \\
A'/(b'(b' - a')) & \cong & A'/(b') \oplus A'/(b' - a') \\
A'/(a'b') & \cong & A'/(a') \oplus A'/(b')
\end{matrix}
$$
De plus, remarquons que la multiplication par $b' - a'$ sur
$A/(a')$ est égale à la multiplication par $b'$, et que
la multiplication par $b' - a'$ sur $A/(b')$ est égale à la multiplication par $-a'$.
En utilisant les lemmes
\ref{chow-lemma-periodic-determinant-easy-case} et
\ref{chow-lemma-periodic-determinant}
nous concluons que
$$
\begin{matrix}
d^A_{A'}(a', b' - a') & = &
\det\nolimits_\kappa(b'|_{A'/(a')})^{-1}
\det\nolimits_\kappa(a'|_{A'/(b' - a')}) \\
d^A_{A'}(b', b' - a') & = &
\det\nolimits_\kappa(-a'|_{A'/(b')})^{-1}
\det\nolimits_\kappa(b'|_{A'/(b' - a')}) \\
d^A_{A'}(a', b') & = &
\det\nolimits_\kappa(b'|_{A'/(a')})^{-1}
\det\nolimits_\kappa(a'|_{A'/(b')})
\end{matrix}
$$
Nous concluons donc que
$$
(-1)^{\text{longueur}_A(A'/(b'))}
d^A_{A'}(a', b' - a')
=
d^A_{A'}(b', b' - a') d^A_{A'}(a', b')
$$
le signe provenant du $-a'$ dans la deuxième égalité ci-dessus.
D'autre part, d'après le lemme \ref{chow-lemma-periodic-determinant-sign}, nous avons
$d^A_{A'}(b', b') = (-1)^{\text{longueur}_A(A'/(b'))}$, et le lemme
est démontré.
\end{proof}

\noindent
Le symbole modéré est un symbole de Steinberg.

\begin{lemma}
\label{chow-lemma-symbol-is-steinberg}
Soit $A$ un anneau intègre local noethérien de dimension $1$
et de corps des fractions $K$. Pour $x \in K \setminus \{0, 1\}$,
nous avons
$$
d_A(x, 1 -x) = 1
$$
\end{lemma}

\begin{proof}
Écrivons $x = a/b$ avec $a, b \in A$.
Puisque $1 - x = (b - a)/b$, l'hypothèse implique
que $b - a \not = 0$ aussi. Nous calculons donc
$$
d_A(x, 1 - x)
=
d_A(a, b - a)d_A(a, b)^{-1}d_A(b, b - a)^{-1}d_A(b, b)
$$
Nous devons ainsi montrer que
$d_A(a, b - a) d_A(b, b) = d_A(b, b - a) d_A(a, b)$.
C'est le lemme \ref{chow-lemma-symbol-is-steinberg-prepare}.
\end{proof}










\subsection{Longueurs et déterminants}
\label{chow-subsection-length-determinant}

\noindent
Dans cette section, nous utilisons le déterminant pour comparer des réseaux.
Le lemme clé est le suivant.

\begin{lemma}
\label{chow-lemma-key-lemma}
Soit $R$ un anneau local noethérien.
Soit $\mathfrak q \subset R$ un idéal premier tel que $\dim(R/\mathfrak q) = 1$.
Soit $\varphi : M \to N$ un homomorphisme de $R$-modules finis.
Supposons qu'il existe $x_1, \ldots, x_l \in M$ et $y_1, \ldots, y_l \in M$
ayant les propriétés suivantes :
\begin{enumerate}
\item $M = \langle x_1, \ldots, x_l\rangle$ ;
\item $\langle x_1, \ldots, x_i\rangle / \langle x_1, \ldots, x_{i - 1}\rangle
\cong R/\mathfrak q$ pour $i = 1, \ldots, l$ ;
\item $N = \langle y_1, \ldots, y_l\rangle$ ; et
\item $\langle y_1, \ldots, y_i\rangle / \langle y_1, \ldots, y_{i - 1}\rangle
\cong R/\mathfrak q$ pour $i = 1, \ldots, l$.
\end{enumerate}
Alors $\varphi$ est injectif si et seulement si $\varphi_{\mathfrak q}$ est un
isomorphisme, et dans ce cas nous avons
$$
\text{longueur}_R(\Coker(\varphi)) = \text{ord}_{R/\mathfrak q}(f)
$$
où $f \in \kappa(\mathfrak q)$ est l'élément tel que
$$
[\varphi(x_1), \ldots, \varphi(x_l)] = f [y_1, \ldots, y_l]
$$
dans $\det_{\kappa(\mathfrak q)}(N_{\mathfrak q})$.
\end{lemma}

\begin{proof}
Remarquons d'abord que le lemme est vrai dans le cas $l = 1$.
En effet, dans ce cas, $x_1$ est une base de $M$ sur $R/\mathfrak q$,
$y_1$ est une base de $N$ sur $R/\mathfrak q$, et nous avons
$\varphi(x_1) = fy_1$ pour un certain $f \in R$. Ainsi, $\varphi$ est injectif
si et seulement si $f \not \in \mathfrak q$. De plus,
$\Coker(\varphi) = R/(f, \mathfrak q)$, et donc le lemme
résulte de la définition de $\text{ord}_{R/q}(f)$
(voir Algèbre, définition \ref{algebra-definition-ord}).

\medskip\noindent
En fait, supposons plus généralement que $\varphi(x_i) = f_iy_i$ pour certains
$f_i \in R$, $f_i \not \in \mathfrak q$. Alors les applications induites
$$
\langle x_1, \ldots, x_i\rangle / \langle x_1, \ldots, x_{i - 1}\rangle
\longrightarrow
\langle y_1, \ldots, y_i\rangle / \langle y_1, \ldots, y_{i - 1}\rangle
$$
sont toutes injectives et ont des conoyaux isomorphes à
$R/(f_i, \mathfrak q)$. Nous voyons donc que
$$
\text{longueur}_R(\Coker(\varphi)) = \sum \text{ord}_{R/\mathfrak q}(f_i).
$$
D'autre part, il est clair que
$$
[\varphi(x_1), \ldots, \varphi(x_l)] = f_1 \ldots f_l [y_1, \ldots, y_l]
$$
dans ce cas, d'après la relation d'admissibilité (b) pour les symboles.
Nous voyons donc que le résultat est encore vrai dans ce cas.

\medskip\noindent
Nous démontrons le cas général par récurrence sur $l$. Supposons $l > 1$.
Soit $i \in \{1, \ldots, l\}$ minimal tel que
$\varphi(x_1) \in \langle y_1, \ldots, y_i\rangle$.
Nous allons raisonner par récurrence sur $i$.
Si $i = 1$, nous obtenons un diagramme commutatif
$$
\xymatrix{
0 \ar[r] &
\langle x_1 \rangle \ar[r] \ar[d] &
\langle x_1, \ldots, x_l \rangle \ar[r] \ar[d] &
\langle x_1, \ldots, x_l \rangle / \langle x_1 \rangle \ar[r] \ar[d] &
0 \\
0 \ar[r] &
\langle y_1 \rangle \ar[r] &
\langle y_1, \ldots, y_l \rangle \ar[r] &
\langle y_1, \ldots, y_l \rangle / \langle y_1 \rangle \ar[r] &
0
}
$$
et le lemme résulte du lemme du serpent et de la récurrence sur $l$.
Supposons maintenant que $i > 1$.
Écrivons $\varphi(x_1) = a_1 y_1 + \ldots + a_{i - 1} y_{i - 1} + a y_i$
avec $a_j, a \in R$ et $a \not \in \mathfrak q$ (sinon
$i$ ne serait pas minimal). Posons
$$
x'_j =
\left\{
\begin{matrix}
x_j & \text{si} & j = 1 \\
ax_j & \text{si} & j \geq 2
\end{matrix}
\right.
\quad\text{et}\quad
y'_j =
\left\{
\begin{matrix}
y_j & \text{si} & j < i \\
ay_j & \text{si} & j \geq i
\end{matrix}
\right.
$$
Soient $M' = \langle x'_1, \ldots, x'_l \rangle$ et
$N' = \langle y'_1, \ldots, y'_l \rangle$.
Puisque $\varphi(x'_1) = a_1 y'_1 + \ldots + a_{i - 1} y'_{i - 1} + y'_i$
par construction, et que pour $j > 1$ nous avons
$\varphi(x'_j) = a\varphi(x_i) \in \langle y'_1, \ldots, y'_l\rangle$,
nous obtenons un diagramme commutatif de $R$-modules et d'applications
$$
\xymatrix{
M' \ar[d] \ar[r]_{\varphi'} & N' \ar[d] \\
M \ar[r]^\varphi & N
}
$$
D'après le résultat du deuxième paragraphe de la démonstration, nous savons
que $\text{longueur}_R(M/M') = (l - 1)\text{ord}_{R/\mathfrak q}(a)$
et, de même,
$\text{longueur}_R(M/M') = (l - i + 1)\text{ord}_{R/\mathfrak q}(a)$.
Une chasse au diagramme montre alors que
$$
\text{longueur}_R(\Coker(\varphi')) =
\text{longueur}_R(\Coker(\varphi)) + i\ \text{ord}_{R/\mathfrak q}(a).
$$
D'autre part, il est clair qu'en écrivant
$$
[\varphi(x_1), \ldots, \varphi(x_l)] = f [y_1, \ldots, y_l],
\quad
[\varphi'(x'_1), \ldots, \varphi(x'_l)] = f' [y'_1, \ldots, y'_l]
$$
nous avons $f' = a^if$. Il suffit donc de démontrer le lemme dans le
cas où $\varphi(x_1) = a_1y_1 + \ldots a_{i - 1}y_{i - 1} + y_i$,
c'est-à-dire dans le cas où $a = 1$. Rappelons ensuite que
$$
[y_1, \ldots, y_l] = [y_1, \ldots, y_{i - 1},
a_1y_1 + \ldots a_{i - 1}y_{i - 1} + y_i, y_{i + 1}, \ldots, y_l]
$$
d'après les relations d'admissibilité pour les symboles. La suite
$y_1, \ldots, y_{i - 1},
a_1y_1 + \ldots + a_{i - 1}y_{i - 1} + y_i, y_{i + 1}, \ldots, y_l$
satisfait encore les conditions (3) et (4) du lemme.
Ainsi, nous pouvons en fait
supposer que $\varphi(x_1) = y_i$. Dans ce cas, remarquons que
$\mathfrak q x_1 = 0$, ce qui implique aussi $\mathfrak q y_i = 0$.
Nous avons
$$
[y_1, \ldots, y_l] =
- [y_1, \ldots, y_{i - 2}, y_i, y_{i - 1}, y_{i + 1}, \ldots, y_l]
$$
d'après la troisième des relations d'admissibilité qui définissent
$\det_{\kappa(\mathfrak q)}(N_{\mathfrak q})$. Nous pouvons donc
remplacer $y_1, \ldots, y_l$ par
la suite
$y'_1, \ldots, y'_l =
y_1, \ldots, y_{i - 2}, y_i, y_{i - 1}, y_{i + 1}, \ldots, y_l$
(qui satisfait elle aussi les conditions (3) et (4) du lemme).
Cela diminue manifestement l'invariant $i$ de $1$, et nous concluons par récurrence
sur $i$.
\end{proof}

\noindent
Pour utiliser le lemme précédent, montrons qu'il existe souvent des suites d'éléments
ayant les propriétés requises.

\begin{lemma}
\label{chow-lemma-good-sequence-exists}
Soit $R$ un anneau local noethérien.
Soit $\mathfrak q \subset R$ un idéal premier.
Soit $M$ un $R$-module fini tel que
$\mathfrak q$ soit l'un des idéaux premiers minimaux du support de $M$.
Alors il existe $x_1, \ldots, x_l \in M$ tels que
\begin{enumerate}
\item le support de $M / \langle x_1, \ldots, x_l\rangle$ ne contienne pas
$\mathfrak q$ ; et
\item $\langle x_1, \ldots, x_i\rangle / \langle x_1, \ldots, x_{i - 1}\rangle
\cong R/\mathfrak q$ pour $i = 1, \ldots, l$.
\end{enumerate}
De plus, dans ce cas, $l = \text{longueur}_{R_\mathfrak q}(M_\mathfrak q)$.
\end{lemma}

\begin{proof}
La condition selon laquelle $\mathfrak q$ est un idéal premier minimal du support
de $M$ implique que $l = \text{longueur}_{R_\mathfrak q}(M_\mathfrak q)$
est fini (voir Algèbre, lemme \ref{algebra-lemma-support-point}).
Nous pouvons donc trouver $y_1, \ldots, y_l \in M_{\mathfrak q}$
tels que
$\langle y_1, \ldots, y_i\rangle / \langle y_1, \ldots, y_{i - 1}\rangle
\cong \kappa(\mathfrak q)$ pour $i = 1, \ldots, l$.
Nous pouvons trouver $f_i \in R$, $f_i \not \in \mathfrak q$, tels que
$f_i y_i$ soit l'image d'un certain élément $z_i \in M$.
De plus, puisque $R$ est noethérien, nous pouvons écrire
$\mathfrak q = (g_1, \ldots, g_t)$ pour certains $g_j \in R$.
Par hypothèse, $g_j y_i \in \langle y_1, \ldots, y_{i - 1} \rangle$
dans le module $M_{\mathfrak q}$.
Grâce à notre choix de $z_i$, nous pouvons trouver d'autres éléments
$f_{ji} \in R$, $f_{ij} \not \in \mathfrak q$, tels que
$f_{ij} g_j z_i \in \langle z_1, \ldots, z_{i - 1} \rangle$
(égalité dans le module $M$).
Le lemme résulte du choix
$$
x_1 = f_{11}f_{12}\ldots f_{1t}z_1,
\quad
x_2 = f_{11}f_{12}\ldots f_{1t}f_{21}f_{22}\ldots f_{2t}z_2,
$$
et ainsi de suite. En effet, puisque tous les éléments $f_i, f_{ij}$ sont inversibles
dans $R_{\mathfrak q}$, nous avons encore
$R_{\mathfrak q}x_1 + \ldots + R_{\mathfrak q}x_i /
R_{\mathfrak q}x_1 + \ldots + R_{\mathfrak q}x_{i - 1}
\cong \kappa(\mathfrak q)$ pour $i = 1, \ldots, l$.
Par construction, $\mathfrak q x_i \in \langle x_1, \ldots, x_{i - 1}\rangle$.
Ainsi, $\langle x_1, \ldots, x_i\rangle / \langle x_1, \ldots, x_{i - 1}\rangle$
est un $R$-module engendré par un élément et annulé par $\mathfrak q$,
tel que sa localisation en $\mathfrak q$ donne un espace vectoriel de dimension $q$
sur $\kappa(\mathfrak q)$.
Il est donc isomorphe à $R/\mathfrak q$.
\end{proof}

\noindent
Voici le résultat principal de cette section.
Nous verrons ci-dessous les différentes
conséquences de cette proposition.
Le lecteur est invité à démontrer d'abord lui-même le plus facile
lemme \ref{chow-lemma-application-herbrand-quotient}.

\begin{proposition}
\label{chow-proposition-length-determinant-periodic-complex}
Soit $R$ un anneau local noethérien de corps résiduel $\kappa$.
Supposons que $(M, \varphi, \psi)$ soit un complexe $(2, 1)$-périodique
sur $R$. Supposons que
\begin{enumerate}
\item $M$ soit un $R$-module fini ;
\item les modules de cohomologie de $(M, \varphi, \psi)$ soient de longueur finie ; et
\item $\dim(\text{Supp}(M)) = 1$.
\end{enumerate}
Soient $\mathfrak q_i$, $i = 1, \ldots, t$, les idéaux
premiers minimaux du support de $M$. Alors nous avons\footnote{
Évidemment, nous pourrions supprimer le signe moins en redéfinissant
$\det_\kappa(M, \varphi, \psi)$ comme l'inverse de sa
valeur actuelle ; voir la définition \ref{chow-definition-periodic-determinant}.}
$$
- e_R(M, \varphi, \psi) =
\sum\nolimits_{i = 1, \ldots, t}
\text{ord}_{R/\mathfrak q_i}\left(
\det\nolimits_{\kappa(\mathfrak q_i)}
(M_{\mathfrak q_i}, \varphi_{\mathfrak q_i}, \psi_{\mathfrak q_i})
\right)
$$
\end{proposition}

\begin{proof}
Nous nous ramenons d'abord au cas $t = 1$ de la manière suivante.
Remarquons que
$\text{Supp}(M) = \{\mathfrak m, \mathfrak q_1, \ldots, \mathfrak q_t\}$,
où $\mathfrak m \subset R$ est l'idéal maximal.
Notons $M_i$ l'image de $M \to M_{\mathfrak q_i}$,
de sorte que $\text{Supp}(M_i) = \{\mathfrak m, \mathfrak q_i\}$.
L'application $\varphi$ (resp.\ $\psi$) induit une application de $R$-modules
$\varphi_i : M_i \to M_i$ (resp.\ $\psi_i : M_i \to M_i$).
Nous obtenons ainsi un morphisme de complexes $(2, 1)$-périodiques
$$
(M, \varphi, \psi) \longrightarrow
\bigoplus\nolimits_{i = 1, \ldots, t} (M_i, \varphi_i, \psi_i).
$$
Le noyau et le conoyau de cette application ont leur support contenu dans
$\{\mathfrak m\}$. Ainsi, d'après le lemme \ref{chow-lemma-compare-periodic-lengths},
nous avons
$$
e_R(M, \varphi, \psi) =
\sum\nolimits_{i = 1, \ldots, t}
e_R(M_i, \varphi_i, \psi_i)
$$
D'autre part, nous avons clairement $M_{\mathfrak q_i} = M_{i, \mathfrak q_i}$,
et donc les termes du membre de droite de la formule du
lemme sont égaux aux expressions
$$
\text{ord}_{R/\mathfrak q_i}\left(
\det\nolimits_{\kappa(\mathfrak q_i)}
(M_{i, \mathfrak q_i}, \varphi_{i, \mathfrak q_i}, \psi_{i, \mathfrak q_i})
\right)
$$
Autrement dit, si nous pouvons démontrer le lemme pour chacun des modules
$M_i$, alors le lemme est vrai. Cela nous ramène au cas $t = 1$.

\medskip\noindent
Supposons donné un complexe $(2, 1)$-périodique $(M, \varphi, \psi)$
sur un anneau local noethérien, où $M$ est un $R$-module fini,
$\text{Supp}(M) = \{\mathfrak m, \mathfrak q\}$ et les modules de cohomologie
sont de longueur finie. Dans ce cas, la démonstration
résulte du lemme \ref{chow-lemma-key-lemma} et d'une comptabilité soigneuse.
Notons
$K_\varphi = \Ker(\varphi)$,
$I_\varphi = \Im(\varphi)$,
$K_\psi = \Ker(\psi)$, et
$I_\psi = \Im(\psi)$.
Puisque $R$ est noethérien, ce sont tous des $R$-modules finis.
Posons
$$
a = \text{longueur}_{R_{\mathfrak q}}(I_{\varphi, \mathfrak q})
= \text{longueur}_{R_{\mathfrak q}}(K_{\psi, \mathfrak q}),
\quad
b = \text{longueur}_{R_{\mathfrak q}}(I_{\psi, \mathfrak q})
= \text{longueur}_{R_{\mathfrak q}}(K_{\varphi, \mathfrak q}).
$$
Ces égalités viennent de ce que le complexe devient exact après localisation en
$\mathfrak q$. Remarquons que $l = \text{longueur}_{R_{\mathfrak q}}(M_{\mathfrak q})$
est égal à $l = a + b$.

\medskip\noindent
Nous allons utiliser le lemme \ref{chow-lemma-good-sequence-exists}
pour choisir des suites d'éléments dans des $R$-modules finis
$N$ dont le support est contenu dans $\{\mathfrak m, \mathfrak q\}$.
Dans ce cas, $N_{\mathfrak q}$ est de longueur finie, disons $n \in \mathbf{N}$.
Appelons « bonne suite » une suite $w_1, \ldots, w_n \in N$
ayant les propriétés (1) et (2) du lemme \ref{chow-lemma-good-sequence-exists}.
Remarquons que le quotient
$N/\langle w_1, \ldots, w_n \rangle$ de $N$ par le sous-module engendré par
une bonne suite a son support (contenu) dans $\{\mathfrak m\}$
et est donc de longueur finie (Algèbre, lemme \ref{algebra-lemma-support-point}).
De plus, le symbole
$[w_1, \ldots, w_n] \in \det_{\kappa(\mathfrak q)}(N_{\mathfrak q})$
est un générateur ; voir le lemme \ref{chow-lemma-determinant-dimension-one}.

\medskip\noindent
Cela étant dit, choisissons de bonnes suites
$$
\begin{matrix}
x_1, \ldots, x_b & \text{dans} & K_\varphi, &
t_1, \ldots, t_a & \text{dans} & K_\psi, \\
y_1, \ldots, y_a & \text{dans} & I_\varphi \cap \langle t_1, \ldots t_a\rangle, &
s_1, \ldots, s_b & \text{dans} & I_\psi \cap \langle x_1, \ldots, x_b\rangle.
\end{matrix}
$$
Nous allons modifier légèrement nos choix comme suit.
Choisissons des relèvements $\tilde y_i \in M$ de $y_i \in I_\varphi$
et $\tilde s_i \in M$ de $s_i \in I_\psi$. Il se peut
que $\mathfrak q \tilde y_1 \subset \langle x_1, \ldots, x_b\rangle$ ne soit pas vérifié,
et il se peut que
$\mathfrak q \tilde s_1 \subset \langle t_1, \ldots, t_a\rangle$ ne le soit pas non plus.
Cependant, en utilisant le fait que $\mathfrak q$ est de type fini (comme dans la démonstration
du lemme \ref{chow-lemma-good-sequence-exists}), nous pouvons trouver
$d \in R$, $d \not \in \mathfrak q$, tel que
$\mathfrak q d\tilde y_1 \subset \langle x_1, \ldots, x_b\rangle$
et
$\mathfrak q d\tilde s_1 \subset \langle t_1, \ldots, t_a\rangle$.
Ainsi, après avoir remplacé $y_i$ par $dy_i$,
$\tilde y_i$ par $d\tilde y_i$, $s_i$ par $ds_i$ et $\tilde s_i$
par $d\tilde s_i$, nous voyons que nous pouvons aussi supposer que
$x_1, \ldots, x_b, \tilde y_1, \ldots, \tilde y_b$
et $t_1, \ldots, t_a, \tilde s_1, \ldots, \tilde s_b$
sont de bonnes suites dans $M$.

\medskip\noindent
Enfin, choisissons une bonne suite
$z_1, \ldots, z_l$ dans le $R$-module fini
$$
\langle
x_1, \ldots, x_b, \tilde y_1, \ldots, \tilde y_a
\rangle
\cap
\langle
t_1, \ldots, t_a, \tilde s_1, \ldots, \tilde s_b
\rangle.
$$
Remarquons que c'est aussi une bonne suite dans $M$.

\medskip\noindent
Puisque $I_{\varphi, \mathfrak q} = K_{\psi, \mathfrak q}$,
il existe un unique élément $h \in \kappa(\mathfrak q)$ tel que
$[y_1, \ldots, y_a] = h [t_1, \ldots, t_a]$
dans $\det_{\kappa(\mathfrak q)}(K_{\psi, \mathfrak q})$.
De même, comme $I_{\psi, \mathfrak q} = K_{\varphi, \mathfrak q}$,
il existe un unique élément $h \in \kappa(\mathfrak q)$ tel que
$[s_1, \ldots, s_b] = g [x_1, \ldots, x_b]$
dans $\det_{\kappa(\mathfrak q)}(K_{\varphi, \mathfrak q})$.
Nous pouvons également procéder ainsi avec les trois bonnes suites que nous avons
dans $M$. Au total, nous obtenons les identités suivantes :
\begin{align*}
[y_1, \ldots, y_a]
& =
h [t_1, \ldots, t_a] \\
[s_1, \ldots, s_b]
& =
g [x_1, \ldots, x_b] \\
[z_1, \ldots, z_l]
& =
f_\varphi [x_1, \ldots, x_b, \tilde y_1, \ldots, \tilde y_a] \\
[z_1, \ldots, z_l]
& =
f_\psi [t_1, \ldots, t_a, \tilde s_1, \ldots, \tilde s_b]
\end{align*}
pour certains $g, h, f_\varphi, f_\psi \in \kappa(\mathfrak q)$.

\medskip\noindent
Toutes ces
notations étant en place, calculons $\det_{\kappa(\mathfrak q)}(M, \varphi, \psi)$.
Considérons l'élément $[z_1, \ldots, z_l]$.
Sous l'application $\gamma_\psi \circ \sigma \circ \gamma_\varphi^{-1}$
de la définition \ref{chow-definition-periodic-determinant}, nous avons
\begin{eqnarray*}
[z_1, \ldots, z_l] & = &
f_\varphi [x_1, \ldots, x_b, \tilde y_1, \ldots, \tilde y_a] \\
& \mapsto & f_\varphi [x_1, \ldots, x_b] \otimes [y_1, \ldots, y_a] \\
& \mapsto &
f_\varphi h/g [t_1, \ldots, t_a] \otimes [s_1, \ldots, s_b] \\
& \mapsto &
f_\varphi h/g [t_1, \ldots, t_a, \tilde s_1, \ldots, \tilde s_b] \\
& = &
f_\varphi h/f_\psi g [z_1, \ldots, z_l]
\end{eqnarray*}
Cela signifie que
$\det_{\kappa(\mathfrak q)}
(M_{\mathfrak q}, \varphi_{\mathfrak q}, \psi_{\mathfrak q})$
est égal à $f_\varphi h/f_\psi g$ au signe près.

\medskip\noindent
Abrégeons les quantités suivantes :
\begin{eqnarray*}
k_\varphi & = & \text{longueur}_R(K_\varphi/\langle x_1, \ldots, x_b\rangle) \\
k_\psi    & = & \text{longueur}_R(K_\psi/\langle t_1, \ldots, t_a\rangle) \\
i_\varphi & = & \text{longueur}_R(I_\varphi/\langle y_1, \ldots, y_a\rangle) \\
i_\psi    & = & \text{longueur}_R(I_\psi/\langle s_1, \ldots, s_a\rangle) \\
m_\varphi & = & \text{longueur}_R(M/
\langle x_1, \ldots, x_b, \tilde y_1, \ldots, \tilde y_a\rangle) \\
m_\psi    & = & \text{longueur}_R(M/
\langle t_1, \ldots, t_a, \tilde s_1, \ldots, \tilde s_b\rangle) \\
\delta_\varphi & = & \text{longueur}_R(
\langle x_1, \ldots, x_b, \tilde y_1, \ldots, \tilde y_a\rangle
\langle z_1, \ldots, z_l\rangle) \\
\delta_\psi & = & \text{longueur}_R(
\langle t_1, \ldots, t_a, \tilde s_1, \ldots, \tilde s_b\rangle
\langle z_1, \ldots, z_l\rangle)
\end{eqnarray*}
En utilisant les suites exactes $0 \to K_\varphi \to M \to I_\varphi \to 0$,
nous obtenons $m_\varphi = k_\varphi + i_\varphi$. De même, nous avons
$m_\psi = k_\psi + i_\psi$. Nous avons
$\delta_\varphi + m_\varphi = \delta_\psi + m_\psi$, puisque cette quantité
est égale à la colongueur de $\langle z_1, \ldots, z_l \rangle$
dans $M$. Enfin, nous avons
$$
\delta_\varphi = \text{ord}_{R/\mathfrak q}(f_\varphi),
\quad
\delta_\psi = \text{ord}_{R/\mathfrak q}(f_\psi)
$$
d'après notre première application du lemme clé \ref{chow-lemma-key-lemma}.

\medskip\noindent
Calculons ensuite la multiplicité du complexe périodique :
\begin{eqnarray*}
e_R(M, \varphi, \psi) & = &
\text{longueur}_R(K_\varphi/I_\psi) - \text{longueur}_R(K_\psi/I_\varphi) \\
& = &
\text{longueur}_R(
\langle x_1, \ldots, x_b\rangle/
\langle s_1, \ldots, s_b\rangle)
+ k_\varphi - i_\psi \\
& & -
\text{longueur}_R(
\langle t_1, \ldots, t_a\rangle/
\langle y_1, \ldots, y_a\rangle)
- k_\psi + i_\varphi \\
& = &
\text{ord}_{R/\mathfrak q}(g/h) + k_\varphi - i_\psi - k_\psi + i_\varphi \\
& = &
\text{ord}_{R/\mathfrak q}(g/h) + m_\varphi - m_\psi \\
& = &
\text{ord}_{R/\mathfrak q}(g/h) + \delta_\psi - \delta_\varphi \\
& = &
\text{ord}_{R/\mathfrak q}(f_\psi g/f_\varphi h)
\end{eqnarray*}
où nous avons utilisé deux fois le lemme clé \ref{chow-lemma-key-lemma} dans la troisième égalité.
D'après notre calcul de $\det_{\kappa(\mathfrak q)}
(M_{\mathfrak q}, \varphi_{\mathfrak q}, \psi_{\mathfrak q})$,
cela démontre la proposition.
\end{proof}

\noindent
Dans la plupart des applications, le lemme suivant suffit.

\begin{lemma}
\label{chow-lemma-application-herbrand-quotient}
Soit $R$ un anneau local noethérien d'idéal maximal $\mathfrak m$.
Soient $M$ un $R$-module fini et $\psi : M \to M$ une
application de $R$-modules. Supposons que
\begin{enumerate}
\item $\Ker(\psi)$ et $\Coker(\psi)$ soient de longueur finie ; et
\item $\dim(\text{Supp}(M)) \leq 1$.
\end{enumerate}
Écrivons
$\text{Supp}(M) = \{\mathfrak m, \mathfrak q_1, \ldots, \mathfrak q_t\}$
et notons $f_i \in \kappa(\mathfrak q_i)^*$ l'élément tel que
$\det_{\kappa(\mathfrak q_i)}(\psi_{\mathfrak q_i}) :
\det_{\kappa(\mathfrak q_i)}(M_{\mathfrak q_i})
\to \det_{\kappa(\mathfrak q_i)}(M_{\mathfrak q_i})$
soit la multiplication par $f_i$. Alors
nous avons
$$
\text{longueur}_R(\Coker(\psi))
-
\text{longueur}_R(\Ker(\psi))
=
\sum\nolimits_{i = 1, \ldots, t}
\text{ord}_{R/\mathfrak q_i}(f_i).
$$
\end{lemma}

\begin{proof}
Rappelons que $H^0(M, 0, \psi) = \Coker(\psi)$ et
$H^1(M, 0, \psi) = \Ker(\psi)$ ; voir les remarques précédant
la définition \ref{chow-definition-periodic-length}.
Le lemme résulte de la combinaison de la
proposition \ref{chow-proposition-length-determinant-periodic-complex} avec le
lemme \ref{chow-lemma-periodic-determinant-easy-case}.

\medskip\noindent
Autre démonstration. Ramenons-nous au cas
$\text{Supp}(M) = \{\mathfrak m, \mathfrak q\}$
comme dans la démonstration de la
proposition \ref{chow-proposition-length-determinant-periodic-complex}.
Combinons alors directement les
lemmes \ref{chow-lemma-key-lemma} et
\ref{chow-lemma-good-sequence-exists}
pour démontrer ce cas particulier de la
proposition \ref{chow-proposition-length-determinant-periodic-complex}.
La comptabilité est bien moindre dans ce cas, et le lecteur est
invité à la mener. Les détails sont omis.
\end{proof}




\subsection{Application au lemme clé}
\label{chow-subsection-application-tame-symbol}

\noindent
Dans cette section, nous appliquons les résultats ci-dessus pour montrer
l'analogue du lemme clé (lemme \ref{chow-lemma-milnor-gersten-low-degree})
avec le symbole modéré $d_A$ construit ci-dessus. Voir la
remarque \ref{chow-remark-gersten-complex-milnor}
pour la relation avec la $K$-théorie de Milnor.

\begin{lemma}[Lemme clé]
\label{chow-lemma-secondary-ramification}
\begin{reference}
Lorsque $A$ est un anneau excellent, c'est \cite[Proposition 1]{Kato-Milnor-K}.
\end{reference}
Soit $A$ un anneau intègre local noethérien de dimension $2$ et de corps des fractions $K$.
Soient $f, g \in K^*$.
Soient $\mathfrak q_1, \ldots, \mathfrak q_t$ les idéaux premiers de hauteur
$1$, notés $\mathfrak q$, de $A$, tels que $f$ ou $g$ n'appartienne pas à
$A^*_{\mathfrak q}$.
Alors nous avons
$$
\sum\nolimits_{i = 1, \ldots, t}
\text{ord}_{A/\mathfrak q_i}(d_{A_{\mathfrak q_i}}(f, g))
=
0
$$
Nous pouvons aussi écrire ceci sous la forme
$$
\sum\nolimits_{\text{hauteur}(\mathfrak q) = 1}
\text{ord}_{A/\mathfrak q}(d_{A_{\mathfrak q}}(f, g))
=
0
$$
puisque, pour tout idéal premier $\mathfrak q$ de hauteur un
de $A$ tel que $f, g \in A^*_{\mathfrak q}$,
nous avons $d_{A_{\mathfrak q}}(f, g) = 1$ d'après le
lemme \ref{chow-lemma-symbol-when-one-is-a-unit}.
\end{lemma}

\begin{proof}
Puisque les symboles modérés $d_{A_{\mathfrak q}}(f, g)$ sont additifs
(lemme \ref{chow-lemma-multiplicativity-symbol}) et que les fonctions d'ordre
$\text{ord}_{A/\mathfrak q}$
sont additives (Algèbre, lemme \ref{algebra-lemma-ord-additive}),
il suffit de démontrer la formule lorsque $f = a \in A$ et
$g = b \in A$. Dans ce cas, nous voyons qu'il faut montrer
$$
\sum\nolimits_{\text{hauteur}(\mathfrak q) = 1}
\text{ord}_{A/\mathfrak q}(\det\nolimits_\kappa(A_{\mathfrak q}/(ab), a, b))
= 0
$$
D'après la proposition \ref{chow-proposition-length-determinant-periodic-complex},
cela équivaut à montrer que
$$
e_A(A/(ab), a, b) = 0.
$$
Puisque le complexe
$A/(ab) \xrightarrow{a} A/(ab) \xrightarrow{b} A/(ab) \xrightarrow{a} A/(ab)$
est exact, nous avons terminé.
\end{proof}





\section{Annexe B : Approches alternatives}
\label{chow-section-appendix-chow}

\noindent
Dans cette annexe, nous tentons d'abord de relier brièvement les résultats
du texte principal à la $K$-théorie des faisceaux cohérents. En
particulier, nous décrivons comment le produit par $c_1$ d'un
module inversible est lié au produit tensoriel par ce module inversible ; voir le
lemme \ref{chow-lemma-coherent-sheaf-cap-c1}.
Ces considérations sont manifestement très intéressantes et
méritent une présentation beaucoup plus détaillée et développée.


\subsection{Équivalence rationnelle et groupes de K-théorie}
\label{chow-subsection-rational-equivalence-K-groups}

\noindent
Cette section prolonge la section \ref{chow-section-chow-and-K}.
Le lemme suivant est motivé par
Homologie, lemme \ref{homology-lemma-serre-subcategory-K-groups}.

\begin{lemma}
\label{chow-lemma-maps-between-coherent-sheaves}
Soit $(S, \delta)$ comme dans la situation \ref{chow-situation-setup}.
Soit $X$ un schéma localement de type fini sur $S$.
Soit $\mathcal{F}$ un faisceau cohérent sur $X$.
Soit
$$
\xymatrix{
\ldots \ar[r] &
\mathcal{F} \ar[r]^\varphi &
\mathcal{F} \ar[r]^\psi &
\mathcal{F} \ar[r]^\varphi &
\mathcal{F} \ar[r] & \ldots
}
$$
un complexe comme dans Homologie, équation (\ref{homology-equation-cyclic-complex}).
Supposons que
\begin{enumerate}
\item $\dim_\delta(\text{Supp}(\mathcal{F})) \leq k + 1$ ;
\item $\dim_\delta(\text{Supp}(H^i(\mathcal{F}, \varphi, \psi))) \leq k$
pour $i = 0, 1$.
\end{enumerate}
Alors nous avons
$$
[H^0(\mathcal{F}, \varphi, \psi)]_k
\sim_{rat}
[H^1(\mathcal{F}, \varphi, \psi)]_k
$$
comme $k$-cycles sur $X$.
\end{lemma}

\begin{proof}
Soit $\{W_j\}_{j \in J}$ la famille des composantes irréductibles
de $\text{Supp}(\mathcal{F})$
de $\delta$-dimension $k + 1$. Remarquons que $\{W_j\}$
est une famille localement finie de fermés de
$X$, d'après le lemme \ref{chow-lemma-length-finite}.
Pour chaque $j$, soit $\xi_j \in W_j$ le point générique.
Posons
$$
f_j = \det\nolimits_{\kappa(\xi_j)}
(\mathcal{F}_{\xi_j}, \varphi_{\xi_j}, \psi_{\xi_j})
\in
R(W_j)^*.
$$
Voir la définition \ref{chow-definition-periodic-determinant} pour les notations.
Nous affirmons que
$$
- [H^0(\mathcal{F}, \varphi, \psi)]_k + [H^1(\mathcal{F}, \varphi, \psi)]_k
=
\sum (W_j \to X)_*\text{div}(f_j)
$$
Si nous démontrons ceci, le lemme en résulte.

\medskip\noindent
Soit $Z \subset X$ un sous-schéma fermé intègre de $\delta$-dimension $k$.
Pour démontrer l'égalité ci-dessus, il suffit de montrer que le coefficient $n$
de $[Z]$ dans
$
[H^0(\mathcal{F}, \varphi, \psi)]_k - [H^1(\mathcal{F}, \varphi, \psi)]_k
$
est égal au coefficient $m$ de $[Z]$ dans
$
\sum (W_j \to X)_*\text{div}(f_j)
$.
Soit $\xi \in Z$ le point générique.
Considérons l'anneau local $A = \mathcal{O}_{X, \xi}$.
Posons $M = \mathcal{F}_\xi$, considéré comme un $A$-module.
Notons $\varphi, \psi : M \to M$ les applications induites par $\varphi, \psi$ sur
le germe.
Par notre choix de $\xi \in Z$, nous avons $\delta(\xi) = k$,
et donc $\dim(\text{Supp}(M)) = 1$.
Enfin, les sous-schémas fermés intègres
$W_j$ passant par $\xi$ correspondent aux idéaux premiers minimaux
$\mathfrak q_i$ de $\text{Supp}(M)$.
Dans chaque cas, l'élément $f_j \in R(W_j)^*$ correspond à
l'élément $\det_{\kappa(\mathfrak q_i)}(M_{\mathfrak q_i}, \varphi, \psi)$
de $\kappa(\mathfrak q_i)^*$. Nous voyons donc que
$$
n = - e_A(M, \varphi, \psi)
$$
et
$$
m =
\sum
\text{ord}_{A/\mathfrak q_i}
(\det\nolimits_{\kappa(\mathfrak q_i)}(M_{\mathfrak q_i}, \varphi, \psi))
$$
Le résultat découle ainsi de la
proposition \ref{chow-proposition-length-determinant-periodic-complex}.
\end{proof}

\begin{lemma}
\label{chow-lemma-cycles-rational-equivalence-K-group}
Soit $(S, \delta)$ comme dans la situation \ref{chow-situation-setup}.
Soit $X$ un schéma localement de type fini sur $S$.
L'application
$$
\CH_k(X) \longrightarrow
K_0(\textit{Coh}_{\leq k + 1}(X)/\textit{Coh}_{\leq k - 1}(X))
$$
du lemme \ref{chow-lemma-from-chow-to-K} induit une bijection de
$\CH_k(X)$ sur l'image $B_k(X)$ de l'application
$$
K_0(\textit{Coh}_{\leq k}(X)/\textit{Coh}_{\leq k - 1}(X))
\longrightarrow
K_0(\textit{Coh}_{\leq k + 1}(X)/\textit{Coh}_{\leq k - 1}(X)).
$$
\end{lemma}

\begin{proof}
D'après le lemme \ref{chow-lemma-cycles-k-group}, nous avons
$Z_k(X) = K_0(\textit{Coh}_{\leq k}(X)/\textit{Coh}_{\leq k - 1}(X))$,
de manière compatible avec l'application du lemme \ref{chow-lemma-from-chow-to-K}.
Supposons donc donné un élément $[A] - [B]$ de
$K_0(\textit{Coh}_{\leq k}(X)/\textit{Coh}_{\leq k - 1}(X))$
qui s'envoie sur zéro dans $B_k(X)$, c'est-à-dire sur zéro dans
$K_0(\textit{Coh}_{\leq k + 1}(X)/\textit{Coh}_{\leq k - 1}(X))$.
Nous devons montrer que $[A] - [B]$ correspond à un cycle
rationnellement équivalent à zéro sur $X$.
Supposons que $[A] = [\mathcal{A}]$ et $[B] = [\mathcal{B}]$
pour certains faisceaux cohérents $\mathcal{A}, \mathcal{B}$ sur
$X$, à support de $\delta$-dimension $\leq k$.
L'hypothèse selon laquelle $[A] - [B]$ s'envoie sur zéro dans le groupe
$K_0(\textit{Coh}_{\leq k + 1}(X)/\textit{Coh}_{\leq k - 1}(X))$
signifie qu'il existe des faisceaux cohérents
$\mathcal{A}', \mathcal{B}'$ sur $X$, à support de
$\delta$-dimension $\leq k - 1$, tels que
$[\mathcal{A} \oplus \mathcal{A}'] - [\mathcal{B} \oplus \mathcal{B}']$
soit nul dans $K_0(\textit{Coh}_{k + 1}(X))$ (utiliser la partie (1) de
Homologie, lemme \ref{homology-lemma-serre-subcategory-K-groups}).
D'après la partie (2) de
Homologie, lemme \ref{homology-lemma-serre-subcategory-K-groups},
cela signifie qu'il existe un complexe $(2, 1)$-périodique
$(\mathcal{F}, \varphi, \psi)$ dans la catégorie $\textit{Coh}_{\leq k + 1}(X)$
tel que
$\mathcal{A} \oplus \mathcal{A}' = H^0(\mathcal{F}, \varphi, \psi)$
et $\mathcal{B} \oplus \mathcal{B}' = H^1(\mathcal{F}, \varphi, \psi)$.
D'après le lemme \ref{chow-lemma-maps-between-coherent-sheaves},
cela implique que
$$
[\mathcal{A} \oplus \mathcal{A}']_k
\sim_{rat}
[\mathcal{B} \oplus \mathcal{B}']_k
$$
Cela montre que $[A] - [B]$ s'envoie sur un cycle rationnellement
équivalent à zéro par l'application
$$
K_0(\textit{Coh}_{\leq k}(X)/\textit{Coh}_{\leq k - 1}(X))
\longrightarrow Z_k(X)
$$
du lemme \ref{chow-lemma-cycles-k-group}. C'est ce que nous
devions démontrer, et la démonstration est achevée.
\end{proof}














\subsection{Diviseurs de Cartier et groupes de K-théorie}
\label{chow-subsection-cartier-coherent}

\noindent
Dans cette section, nous décrivons comment l'intersection avec la
première classe de Chern d'un faisceau inversible $\mathcal{L}$
correspond au produit tensoriel par $\mathcal{L} - \mathcal{O}$
dans les groupes de $K$-théorie.

\begin{lemma}
\label{chow-lemma-no-embedded-points-modules}
Soit $A$ un anneau local noethérien.
Soit $M$ un $A$-module fini.
Soient $a, b \in A$.
Supposons que
\begin{enumerate}
\item $\dim(A) = 1$ ;
\item $a$ et $b$ soient tous deux des non-diviseurs de zéro dans $A$ ;
\item $A$ n'ait pas d'idéaux premiers immergés ;
\item $M$ n'ait pas d'idéaux premiers associés immergés ;
\item $\text{Supp}(M) = \Spec(A)$.
\end{enumerate}
Soit $I = \{x \in A \mid x(a/b) \in A\}$.
Soient $\mathfrak q_1, \ldots, \mathfrak q_t$ les idéaux premiers minimaux
de $A$. Alors $(a/b)IM \subset M$ et
$$
\text{longueur}_A(M/(a/b)IM)
-
\text{longueur}_A(M/IM)
=
\sum\nolimits_i
\text{longueur}_{A_{\mathfrak q_i}}(M_{\mathfrak q_i})
\text{ord}_{A/\mathfrak q_i}(a/b)
$$
\end{lemma}

\begin{proof}
Puisque $M$ n'a pas d'idéaux premiers associés immergés et que
le support de $M$ est $\Spec(A)$, nous voyons que
$\text{Ass}(M) = \{\mathfrak q_1, \ldots, \mathfrak q_t\}$.
Ainsi, $a$ et $b$ sont des non-diviseurs de zéro dans $M$. Remarquons que
\begin{align*}
& \text{longueur}_A(M/(a/b)IM) \\
& = \text{longueur}_A(bM/aIM) \\
& = \text{longueur}_A(M/aIM)
-
\text{longueur}_A(M/bM) \\
& = \text{longueur}_A(M/aM) + \text{longueur}_A(aM/aIM) - \text{longueur}_A(M/bM) \\
& = \text{longueur}_A(M/aM) + \text{longueur}_A(M/IM) - \text{longueur}_A(M/bM)
\end{align*}
car l'application injective $b : M \to bM$ envoie $(a/b)IM$ sur $aIM$,
et l'application injective $a : M \to aM$ envoie $IM$ sur $aIM$.
Le membre de gauche de l'équation du lemme est donc
égal à
$$
\text{longueur}_A(M/aM) - \text{longueur}_A(M/bM).
$$
En appliquant la deuxième formule du
lemme \ref{chow-lemma-additivity-divisors-restricted}
respectivement avec $x = a, b$,
et en utilisant Algèbre, définition \ref{algebra-definition-ord},
des fonctions $\text{ord}$, nous obtenons le résultat.
\end{proof}

\begin{lemma}
\label{chow-lemma-no-embedded-points}
Soit $(S, \delta)$ comme dans la situation \ref{chow-situation-setup}.
Soit $X$ un schéma localement de type fini sur $S$.
Soit $\mathcal{L}$ un $\mathcal{O}_X$-module inversible.
Soit $\mathcal{F}$ un $\mathcal{O}_X$-module cohérent.
Soit $s \in \Gamma(X, \mathcal{K}_X(\mathcal{L}))$ une
section méromorphe de $\mathcal{L}$.
Supposons que
\begin{enumerate}
\item $\dim_\delta(X) \leq k + 1$ ;
\item $X$ n'ait pas de points immergés ;
\item $\mathcal{F}$ n'ait pas de points associés immergés ;
\item le support de $\mathcal{F}$ soit $X$ ; et
\item la section $s$ soit méromorphe régulière.
\end{enumerate}
Dans cette situation, soit $\mathcal{I} \subset \mathcal{O}_X$
l'idéal des dénominateurs de $s$ ; voir
Diviseurs,
définition \ref{divisors-definition-regular-meromorphic-ideal-denominators}.
Nous avons alors ce qui suit :
\begin{enumerate}
\item il existe des suites exactes courtes
$$
\begin{matrix}
0 &
\to &
\mathcal{I}\mathcal{F} &
\xrightarrow{1} &
\mathcal{F} &
\to &
\mathcal{Q}_1 &
\to &
0 \\
0 &
\to &
\mathcal{I}\mathcal{F} &
\xrightarrow{s} &
\mathcal{F} \otimes_{\mathcal{O}_X} \mathcal{L} &
\to &
\mathcal{Q}_2 &
\to &
0
\end{matrix}
$$
\item les faisceaux cohérents $\mathcal{Q}_1$, $\mathcal{Q}_2$
sont à support de $\delta$-dimension $\leq k$ ;
\item la section $s$ se restreint en une section méromorphe régulière
$s_i$ sur toute composante irréductible $X_i$ de
$X$ de $\delta$-dimension $k + 1$ ; et
\item en écrivant $[\mathcal{F}]_{k + 1} = \sum m_i[X_i]$, nous avons
$$
[\mathcal{Q}_2]_k - [\mathcal{Q}_1]_k
=
\sum m_i(X_i \to X)_*\text{div}_{\mathcal{L}|_{X_i}}(s_i)
$$
dans $Z_k(X)$ ; en particulier,
$$
[\mathcal{Q}_2]_k - [\mathcal{Q}_1]_k
=
c_1(\mathcal{L}) \cap [\mathcal{F}]_{k + 1}
$$
dans $\CH_k(X)$.
\end{enumerate}
\end{lemma}

\begin{proof}
Rappelons, d'après Diviseurs, lemme \ref{divisors-lemma-make-maps-regular-section},
l'existence d'applications injectives
$1 : \mathcal{I}\mathcal{F} \to \mathcal{F}$ et
$s : \mathcal{I}\mathcal{F} \to \mathcal{F} \otimes_{\mathcal{O}_X}\mathcal{L}$
dont les conoyaux sont supportés par un fermé partout non dense $T$.
Notons $\mathcal{Q}_i$ leurs conoyaux comme dans le lemme.
Nous concluons que $\dim_\delta(\text{Supp}(\mathcal{Q}_i)) \leq k$.
D'après Diviseurs, lemmes \ref{divisors-lemma-pullback-meromorphic-sections-defined}
et \ref{divisors-lemma-meromorphic-sections-pullback}, les images réciproques $s_i$
sont définies et sont des sections méromorphes régulières de $\mathcal{L}|_{X_i}$.
L'égalité des cycles dans (4) implique l'égalité des classes de cycles
dans (4). Il reste donc seulement à montrer que
$$
[\mathcal{Q}_2]_k - [\mathcal{Q}_1]_k
=
\sum m_i(X_i \to X)_*\text{div}_{\mathcal{L}|_{X_i}}(s_i)
$$
est vraie dans $Z_k(X)$. Pour le voir, soit $Z \subset X$ un sous-schéma fermé intègre
de $\delta$-dimension $k$. Soit $\xi \in Z$ le point générique.
Soient $A = \mathcal{O}_{X, \xi}$ et $M = \mathcal{F}_\xi$.
De plus, choisissons un générateur $s_\xi \in \mathcal{L}_\xi$.
Nous pouvons alors écrire $s = (a/b) s_\xi$, où $a, b \in A$ sont des
non-diviseurs de zéro. Dans ce cas,
$I = \mathcal{I}_\xi = \{x \in A \mid x(a/b) \in A\}$.
Dans ce cas, le coefficient de $[Z]$ dans le membre de gauche est
$$
\text{longueur}_A(M/(a/b)IM) - \text{longueur}_A(M/IM)
$$
et le coefficient de $[Z]$ dans le membre de droite
est
$$
\sum
\text{longueur}_{A_{\mathfrak q_i}}(M_{\mathfrak q_i})
\text{ord}_{A/\mathfrak q_i}(a/b)
$$
où $\mathfrak q_1, \ldots, \mathfrak q_t$ sont les idéaux premiers minimaux
de l'anneau local de dimension $1$ noté $A$. Le résultat
découle donc du lemme \ref{chow-lemma-no-embedded-points-modules}.
\end{proof}

\begin{lemma}
\label{chow-lemma-coherent-sheaf-cap-c1}
Soit $(S, \delta)$ comme dans la situation \ref{chow-situation-setup}.
Soit $X$ un schéma localement de type fini sur $S$.
Soit $\mathcal{L}$ un $\mathcal{O}_X$-module inversible.
Soit $\mathcal{F}$ un $\mathcal{O}_X$-module cohérent.
Supposons $\dim_\delta(\text{Supp}(\mathcal{F})) \leq k + 1$.
Alors l'élément
$$
[\mathcal{F} \otimes_{\mathcal{O}_X} \mathcal{L}]
-
[\mathcal{F}]
\in
K_0(\textit{Coh}_{\leq k + 1}(X)/\textit{Coh}_{\leq k - 1}(X))
$$
appartient au sous-groupe $B_k(X)$ du
lemme \ref{chow-lemma-cycles-rational-equivalence-K-group} et s'envoie sur
l'élément $c_1(\mathcal{L}) \cap [\mathcal{F}]_{k + 1}$ par
l'application $B_k(X) \to \CH_k(X)$.
\end{lemma}

\begin{proof}
Soit
$$
0 \to \mathcal{K} \to \mathcal{F} \to \mathcal{F}' \to 0
$$
la suite exacte courte construite dans
Diviseurs, lemme \ref{divisors-lemma-remove-embedded-points}.
Cela signifie en particulier que $\mathcal{F}'$ n'a pas de points
associés immergés.
Puisque le support de $\mathcal{K}$ est partout non dense dans le
support de $\mathcal{F}$, nous voyons que
$\dim_\delta(\text{Supp}(\mathcal{K})) \leq k$. Nous pouvons
appliquer de nouveau
Diviseurs, lemme \ref{divisors-lemma-remove-embedded-points},
en partant de $\mathcal{K}$, pour obtenir une suite exacte courte
$$
0 \to \mathcal{K}'' \to \mathcal{K} \to \mathcal{K}' \to 0
$$
où maintenant $\dim_\delta(\text{Supp}(\mathcal{K}'')) < k$
et où $\mathcal{K}'$ n'a pas de points associés immergés.
Supposons que nous puissions démontrer le lemme pour les faisceaux cohérents
$\mathcal{F}'$ et $\mathcal{K}'$. Alors, d'après
les égalités
$$
[\mathcal{F}]_{k + 1}
=
[\mathcal{F}']_{k + 1}
+ [\mathcal{K}']_{k + 1}
+ [\mathcal{K}'']_{k + 1}
$$
(utiliser le lemme \ref{chow-lemma-additivity-sheaf-cycle}),
$$
[\mathcal{F} \otimes_{\mathcal{O}_X} \mathcal{L}]
-
[\mathcal{F}]
=
[\mathcal{F}' \otimes_{\mathcal{O}_X} \mathcal{L}]
-
[\mathcal{F}']
+
[\mathcal{K}' \otimes_{\mathcal{O}_X} \mathcal{L}]
-
[\mathcal{K}']
+
[\mathcal{K}'' \otimes_{\mathcal{O}_X} \mathcal{L}]
-
[\mathcal{K}'']
$$
(utiliser le fait que $\otimes \mathcal{L}$ est exact)
et l'annulation triviale de $[\mathcal{K}'']_{k + 1}$ et de
$[\mathcal{K}'' \otimes_{\mathcal{O}_X} \mathcal{L}]
- [\mathcal{K}'']$ dans
$K_0(\textit{Coh}_{\leq k + 1}(X)/\textit{Coh}_{\leq k - 1}(X))$,
le résultat est vrai
pour $\mathcal{F}$. Cela signifie que nous pouvons supposer que
le faisceau $\mathcal{F}$ n'a pas de points associés immergés.

\medskip\noindent
Supposons que $X$ et $\mathcal{F}$ soient comme dans le lemme, et supposons de plus
que $\mathcal{F}$ n'a pas de points associés immergés. Considérons le
faisceau d'idéaux $\mathcal{I} \subset \mathcal{O}_X$, le sous-schéma fermé
correspondant $i : Z \to X$ et le $\mathcal{O}_Z$-module cohérent
$\mathcal{G}$ construits dans
Diviseurs, lemme \ref{divisors-lemma-no-embedded-points-endos}.
Rappelons que $Z$ est un schéma localement noethérien sans points immergés,
que $\mathcal{G}$ est un faisceau cohérent sans points
associés immergés, que $\text{Supp}(\mathcal{G}) = Z$,
et que $i_*\mathcal{G} = \mathcal{F}$.
De plus, posons $\mathcal{N} = \mathcal{L}|_Z$.

\medskip\noindent
D'après Diviseurs, lemme \ref{divisors-lemma-regular-meromorphic-section-exists},
le faisceau inversible $\mathcal{N}$ possède une section méromorphe régulière $s$
sur $Z$. Notons $\mathcal{J} \subset \mathcal{O}_Z$ le faisceau
des dénominateurs de $s$. D'après le lemme \ref{chow-lemma-no-embedded-points},
il existe des suites exactes courtes
$$
\begin{matrix}
0 &
\to &
\mathcal{J}\mathcal{G} &
\xrightarrow{1} &
\mathcal{G} &
\to &
\mathcal{Q}_1 &
\to &
0 \\
0 &
\to &
\mathcal{J}\mathcal{G} &
\xrightarrow{s} &
\mathcal{G} \otimes_{\mathcal{O}_Z} \mathcal{N} &
\to &
\mathcal{Q}_2 &
\to &
0
\end{matrix}
$$
telles que $\dim_\delta(\text{Supp}(\mathcal{Q}_i)) \leq k$ et
telles que le cycle
$
[\mathcal{Q}_2]_k - [\mathcal{Q}_1]_k
$
soit un représentant de $c_1(\mathcal{N}) \cap [\mathcal{G}]_{k + 1}$.
Nous voyons (en utilisant le fait que
$i_*(\mathcal{G} \otimes \mathcal{N}) = \mathcal{F} \otimes \mathcal{L}$
par la formule de projection ; voir
Cohomologie, lemme \ref{cohomology-lemma-projection-formula})
que
$$
[\mathcal{F} \otimes_{\mathcal{O}_X} \mathcal{L}]
-
[\mathcal{F}]
=
[i_*\mathcal{Q}_2] - [i_*\mathcal{Q}_1]
$$
dans $K_0(\textit{Coh}_{\leq k + 1}(X)/\textit{Coh}_{\leq k - 1}(X))$.
Cela montre déjà que
$[\mathcal{F} \otimes_{\mathcal{O}_X} \mathcal{L}] - [\mathcal{F}]$
est un élément de $B_k(X)$. De plus, nous avons
\begin{eqnarray*}
[i_*\mathcal{Q}_2]_k - [i_*\mathcal{Q}_1]_k
& = &
i_*\left( [\mathcal{Q}_2]_k - [\mathcal{Q}_1]_k \right) \\
& = &
i_*\left(c_1(\mathcal{N}) \cap [\mathcal{G}]_{k + 1} \right) \\
& = &
c_1(\mathcal{L}) \cap i_*[\mathcal{G}]_{k + 1} \\
& = &
c_1(\mathcal{L}) \cap [\mathcal{F}]_{k + 1}
\end{eqnarray*}
d'après ce qui précède et les lemmes \ref{chow-lemma-pushforward-cap-c1}
et \ref{chow-lemma-cycle-push-sheaf}. Et cela coïncide avec
l'image de l'élément par $B_k(X) \to \CH_k(X)$, par définition.
Le lemme est donc démontré.
\end{proof}






































% OK, start here.
%

\chapter{Théorie des intersections}



\phantomsection
\label{intersection-section-phantom}



\section{Introduction}
\label{intersection-section-introduction}

\noindent
Dans ce chapitre, nous construisons le produit d'intersection sur les groupes de Chow
modulo l'équivalence rationnelle d'une variété projective non singulière sur un
corps algébriquement clos. Nos outils sont la formule des Tor de Serre
(voir \cite[chapitre V]{Serre_algebre_locale}), la réduction à la diagonale
et le lemme de déplacement.

\medskip\noindent
Nous rappelons d'abord les cycles et la construction de l'image directe propre et de
l'image inverse plate des cycles. Ensuite, nous introduisons l'équivalence rationnelle des cycles,
qui nous donne les groupes de Chow $\CH_*(X)$. L'image directe propre et l'image inverse plate
se factorisent par l'équivalence rationnelle et donnent des opérations sur les groupes de Chow.
Cela occupe les sections
\ref{intersection-section-cycles},
\ref{intersection-section-cycle-of-closed},
\ref{intersection-section-cycle-of-coherent-sheaf},
\ref{intersection-section-pushforward},
\ref{intersection-section-flat-pullback},
\ref{intersection-section-rational-equivalence},
\ref{intersection-section-alternative},
\ref{intersection-section-pushforward-and-rational-equivalence}, et
\ref{intersection-section-flat-pullback-and-rational-equivalence}.
Pour les démonstrations, nous renvoyons principalement au chapitre sur l'homologie de Chow,
où ces résultats ont été démontrés dans le cadre des
schémas localement de type fini sur une base noethérienne universellement caténaire ; voir
Homologie de Chow, section \ref{chow-section-setup} et suivantes.

\medskip\noindent
Puisque nous travaillons sur une variété projective non singulière $X$, toute composante irréductible
de l'intersection $V \cap W$ de deux sous-variétés fermées irréductibles
est de dimension au moins $\dim(V) + \dim(W) - \dim(X)$. Nous disons que $V$ et $W$
se coupent proprement si l'égalité vaut pour toute composante irréductible $Z$.
Dans ce cas, nous définissons la multiplicité d'intersection
$e_Z = e(X, V \cdot W, Z)$ par la formule
$$
e_Z = \sum\nolimits_i
(-1)^i
\text{longueur}_{\mathcal{O}_{X, Z}}
\text{Tor}_i^{\mathcal{O}_{X, Z}}(\mathcal{O}_{W, Z}, \mathcal{O}_{V, Z})
$$
Nous devons faire un peu d'algèbre commutative pour montrer que ces
multiplicités d'intersection concordent avec l'intuition dans les cas simples,
c'est-à-dire que, parfois,
$$
e_Z = \text{longueur}_{\mathcal{O}_{X, Z}} \mathcal{O}_{V \cap W, Z},
$$
autrement dit, seul $\text{Tor}_0$ contribue. Cela se produit lorsque
$V$ et $W$ sont de Cohen-Macaulay au point générique de $Z$, ou lorsque
$W$ est défini par une suite régulière dans $\mathcal{O}_{X, Z}$ qui
définit aussi une suite régulière sur $\mathcal{O}_{V, Z}$. Toutefois,
l'exemple \ref{intersection-example-naive-multiplicity-wrong} montre que les Tor supérieurs
sont nécessaires en général. De plus, il existe un lien
avec la multiplicité de Samuel. Ces questions sont étudiées dans les
sections
\ref{intersection-section-intersect-properly},
\ref{intersection-section-tor-formula},
\ref{intersection-section-multiplicities},
\ref{intersection-section-computing-intersection-multiplicities}, et
\ref{intersection-section-intersection-product}.

\medskip\noindent
La réduction à la diagonale affirme que nous pouvons couper
$V$ et $W$ en coupant $V \times W$ avec la diagonale dans $X \times X$.
Cet énoncé anodin, évident au niveau des intersections
schématiques, ramène l'intersection d'une paire quelconque
de sous-schémas fermés au cas où l'un des deux est localement défini
par une suite régulière. À la suite de Serre, nous l'utilisons pour obtenir la positivité
des multiplicités d'intersection. De plus, la réduction à la diagonale
conduit à l'additivité des multiplicités d'intersection, à l'associativité et à une
formule de projection. On trouvera cela dans les sections
\ref{intersection-section-exterior-product},
\ref{intersection-section-reduction-diagonal},
\ref{intersection-section-associative},
\ref{intersection-section-flat-pullback-and-intersection-products}, et
\ref{intersection-section-projection-formula-flat}.

\medskip\noindent
Enfin, nous abordons les lemmes de déplacement et leurs applications. Le lemme de
déplacement comporte deux parties. La première affirme qu'étant données des sous-variétés fermées
$$
Z \subset X \subset \mathbf{P}^N
$$
avec $X$ non singulière, nous pouvons trouver une sous-variété $C \subset \mathbf{P}^N$
qui coupe $X$ proprement et telle que
$$
C \cdot X = [Z] + \sum m_j [Z_j]
$$
et telle que les autres composantes $Z_j$ soient « plus générales » que $Z$.
La seconde partie affirme que l'on peut déplacer $C \subset \mathbf{P}^N$, le long
d'une courbe rationnelle, vers une sous-variété en position générale par rapport à
toute liste donnée de sous-variétés. Ensemble, ces résultats impliquent qu'il suffit
de définir le produit d'intersection des cycles sur $X$ qui se coupent
proprement, ce qui a été fait ci-dessus. Bien entendu, cela ne conduit à un produit d'intersection
sur $\CH_*(X)$ que si l'on montre, comme nous le faisons dans le texte, que ces produits
passent au quotient par l'équivalence rationnelle. Cette question et quelques applications sont étudiées
dans les sections
\ref{intersection-section-projection},
\ref{intersection-section-moving-lemma},
\ref{intersection-section-intersections-and-rational-equivalence},
\ref{intersection-section-chow-rings},
\ref{intersection-section-general-pullback}, et
\ref{intersection-section-pullback-cycles}.



\section{Conventions}
\label{intersection-section-conventions}

\noindent
Nous fixons un corps de base algébriquement clos $\mathbf{C}$, de caractéristique
quelconque. Tous les schémas et toutes les variétés sont définis sur $\mathbf{C}$, et tous les
morphismes sont des $\mathbf{C}$-morphismes. Une variété $X$ est
{\it non singulière} si $X$ est un schéma régulier (voir
Propriétés, définition \ref{properties-definition-regular}).
Dans notre cas, cela signifie que le morphisme $X \to \Spec(\mathbf{C})$
est lisse (voir
Variétés, lemme \ref{varieties-lemma-geometrically-regular-smooth}).


\section{Cycles}
\label{intersection-section-cycles}

\noindent
Soit $X$ une variété. Une {\it sous-variété fermée} de $X$ est un sous-schéma
fermé intègre $Z \subset X$. Un {\it $k$-cycle} sur $X$ est une somme formelle
finie $\sum n_i [Z_i]$, où chaque $Z_i$ est une sous-variété fermée
de dimension $k$. Chaque fois que nous employons la notation $\alpha = \sum n_i[Z_i]$
pour un $k$-cycle, nous supposons toujours que les sous-variétés $Z_i$ sont deux à deux
distinctes et que $n_i \not = 0$ pour tout $i$. Dans ce cas, le
{\it support} de $\alpha$ est le sous-ensemble fermé
$$
\text{Supp}(\alpha) = \bigcup Z_i \subset X
$$
de dimension $k$. Le groupe des $k$-cycles est noté $Z_k(X)$.
Voir Homologie de Chow, section \ref{chow-section-cycles}.


\section{Cycle associé à un sous-schéma fermé}
\label{intersection-section-cycle-of-closed}

\noindent
Supposons que $X$ soit une variété et que $Z \subset X$ soit un sous-schéma fermé
avec $\dim(Z) \leq k$. Soient $Z_i$ les composantes irréductibles de $Z$ de
dimension $k$, et soit $n_i$ la {\it multiplicité de $Z_i$ dans $Z$},
définie par
$$
n_i = \text{longueur}_{\mathcal{O}_{X, Z_i}} \mathcal{O}_{Z, Z_i}
$$
où $\mathcal{O}_{X, Z_i}$, resp.\ $\mathcal{O}_{Z, Z_i}$, est l'anneau
local de $X$, resp.\ de $Z$, au point générique de $Z_i$.
Nous définissons le $k$-cycle associé à $Z$ comme le $k$-cycle
$$
[Z]_k = \sum n_i [Z_i].
$$
Voir Homologie de Chow, section \ref{chow-section-cycle-of-closed-subscheme}.


\section{Cycle associé à un faisceau cohérent}
\label{intersection-section-cycle-of-coherent-sheaf}

\noindent
Supposons que $X$ soit une variété et que
$\mathcal{F}$ soit un $\mathcal{O}_X$-module cohérent tel que
$\dim(\text{Supp}(\mathcal{F})) \leq k$.
Soient $Z_i$ les composantes irréductibles de $\text{Supp}(\mathcal{F})$
de dimension $k$, et soit $n_i$ la
{\it multiplicité de $Z_i$ dans $\mathcal{F}$}, définie par
$$
n_i = \text{longueur}_{\mathcal{O}_{X, Z_i}} \mathcal{F}_{\xi_i}
$$
où $\mathcal{O}_{X, Z_i}$ est l'anneau
local de $X$ au point générique $\xi_i$ de $Z_i$,
et où $\mathcal{F}_{\xi_i}$ est le germe de $\mathcal{F}$ en ce point.
Nous définissons le $k$-cycle associé à $\mathcal{F}$ comme le $k$-cycle
$$
[\mathcal{F}]_k = \sum n_i [Z_i].
$$
Voir Homologie de Chow, section \ref{chow-section-cycle-of-coherent-sheaf}.
Notons que, si $Z \subset X$ est un sous-schéma fermé avec $\dim(Z) \leq k$, alors
$[Z]_k = [\mathcal{O}_Z]_k$ par définition.


\section{Image directe propre}
\label{intersection-section-pushforward}

\noindent
Supposons que $f : X \to Y$ soit un morphisme propre de variétés.
Soit $Z \subset X$ une sous-variété fermée de dimension $k$.
Nous définissons $f_*[Z]$ comme $0$ si $\dim(f(Z)) < k$,
et comme $d \cdot [f(Z)]$ si $\dim(f(Z)) = k$, où
$$
d = [\mathbf{C}(Z) : \mathbf{C}(f(Z))] = \deg(Z/f(Z))
$$
est le degré du morphisme dominant $Z \to f(Z)$ ; voir
Morphismes, définition \ref{morphisms-definition-degree}.
Soit $\alpha = \sum n_i [Z_i]$ un $k$-cycle sur $X$. L'
{\it image directe} de $\alpha$ est la somme $f_* \alpha = \sum n_i f_*[Z_i]$,
où chaque $f_*[Z_i]$ est défini comme ci-dessus. Cela définit un homomorphisme
$$
f_* : Z_k(X) \longrightarrow Z_k(Y)
$$
Voir Homologie de Chow, section \ref{chow-section-proper-pushforward}.

\begin{lemma}
\label{intersection-lemma-push-coherent}
\begin{reference}
Voir \cite[chapitre V]{Serre_algebre_locale}.
\end{reference}
Supposons que $f : X \to Y$ soit un morphisme propre de variétés.
Soit $\mathcal{F}$ un faisceau cohérent tel que
$\dim(\text{Supp}(\mathcal{F})) \leq k$ ; alors
$f_*[\mathcal{F}]_k = [f_*\mathcal{F}]_k$. En particulier, si
$Z \subset X$ est un sous-schéma fermé de dimension $\leq k$, alors
$f_*[Z]_k = [f_*\mathcal{O}_Z]_k$.
\end{lemma}

\begin{proof}
Voir Homologie de Chow, lemme \ref{chow-lemma-cycle-push-sheaf}.
\end{proof}

\begin{lemma}
\label{intersection-lemma-compose-pushforward}
Soient $f : X \to Y$ et $g : Y \to Z$ des morphismes propres de
variétés. Alors $g_* \circ f_* = (g \circ f)_*$ comme applications $Z_k(X) \to Z_k(Z)$.
\end{lemma}

\begin{proof}
Cas particulier de Homologie de Chow, lemme \ref{chow-lemma-compose-pushforward}.
\end{proof}



\section{Image inverse plate}
\label{intersection-section-flat-pullback}

\noindent
Supposons que $f : X \to Y$ soit un morphisme plat de variétés.
D'après Morphismes, lemme
\ref{morphisms-lemma-dimension-fibre-at-a-point-additive} 
toute fibre de $f$ est de dimension $r = \dim(X) - \dim(Y)$\footnote{Réciproquement,
si $f : X \to Y$ est un morphisme dominant de variétés, si
$X$ est de Cohen-Macaulay, si $Y$ est non singulière et si toutes les fibres ont
la même dimension $r$, alors $f$ est plat. Cela résulte de
Algèbre, lemme \ref{algebra-lemma-CM-over-regular-flat}, et de
Variétés, lemme \ref{varieties-lemma-dimension-fibres-locally-algebraic},
qui montre que $\dim(X) = \dim(Y) + r$.}.
Soit $Z \subset Y$ une sous-variété fermée de dimension $k$. Nous définissons
$f^*[Z]$ comme le $(k + r)$-cycle associé à l'image inverse
schématique : $f^*[Z] = [f^{-1}(Z)]_{k + r}$. Soit
$\alpha = \sum n_i [Z_i]$ un $k$-cycle sur $Y$. L'{\it image inverse} de
$\alpha$ est la somme $f^* \alpha = \sum n_i f^*[Z_i]$, où chaque $f^*[Z_i]$
est défini comme ci-dessus. Cela définit un homomorphisme
$$
f^* : Z_k(Y) \longrightarrow Z_{k + r}(X)
$$
Voir Homologie de Chow, section \ref{chow-section-flat-pullback}.

\begin{lemma}
\label{intersection-lemma-pullback}
Soit $f : X \to Y$ un morphisme plat de variétés. Posons $r = \dim(X) - \dim(Y)$.
Alors $f^*[\mathcal{F}]_k = [f^*\mathcal{F}]_{k + r}$
si $\mathcal{F}$ est un faisceau cohérent sur $Y$ et si la dimension du
support de $\mathcal{F}$ est au plus $k$.
\end{lemma}

\begin{proof}
Voir Homologie de Chow, lemme \ref{chow-lemma-pullback-coherent}.
\end{proof}

\begin{lemma}
\label{intersection-lemma-compose-flat-pullback}
Soient $f : X \to Y$ et $g : Y \to Z$ des morphismes plats de
variétés. Alors $g \circ f$ est plat et $f^* \circ g^* = (g \circ f)^*$
comme applications $Z_k(Z) \to Z_{k + \dim(X) - \dim(Z)}(X)$.
\end{lemma}

\begin{proof}
Cas particulier de Homologie de Chow, lemme \ref{chow-lemma-compose-flat-pullback}.
\end{proof}


\section{Équivalence rationnelle}
\label{intersection-section-rational-equivalence}

\noindent
Nous allons définir l'équivalence rationnelle d'une manière qui peut, au premier
abord, sembler différente de celle qui vous est familière ou de celle
de \cite[chapitre I]{F} ou de
Homologie de Chow, section \ref{chow-section-rational-equivalence}.
Cependant, nous montrerons dans la section \ref{intersection-section-alternative} que
les deux notions coïncident.

\medskip\noindent
Soit $X$ une variété. Soit $W \subset X \times \mathbf{P}^1$
une sous-variété fermée de dimension $k + 1$. Soient $a, b$ deux points fermés distincts
de $\mathbf{P}^1$. Supposons que $X \times a$, $X \times b$ et $W$
se coupent proprement :
$$
\dim (W \cap X \times a) \leq k,\quad
\dim (W \cap X \times b) \leq k.
$$
C'est le cas dès que $W \to \mathbf{P}^1$ est dominant, ou si $W$ est
contenu dans une fibre de la projection au-dessus d'un point fermé distinct de
$a$ et de $b$ (cas sans intérêt que nous écarterons). Dans cette
situation, la fibre schématique $W_a$ du morphisme
$W \to \mathbf{P}^1$ est égale à l'intersection schématique
$W \cap X \times a$ dans $X \times \mathbf{P}^1$. En identifiant $X \times a$
et $X \times b$ à $X$, nous pouvons considérer les fibres $W_a$ et $W_b$
comme des sous-schémas fermés de $X$ de dimension $\leq k$\footnote{Nous considérerons parfois
$W_a$ comme un sous-schéma fermé de $X \times \mathbf{P}^1$, et parfois
comme un sous-schéma fermé de $X$. Le contexte indiquera toujours clairement quel
point de vue est adopté.}. Un exemple fondamental d'
équivalence rationnelle est
$$
[W_a]_k \sim_{rat} [W_b]_k
$$
Les cycles $[W_a]_k$ et $[W_b]_k$ se calculent facilement en pratique
(une fois $W$ donné), car ils s'obtiennent par intersection propre avec
un diviseur de Cartier (nous le verrons dans la
section \ref{intersection-section-intersection-product}).
Comme le groupe des automorphismes de $\mathbf{P}^1$ est $2$-transitif, nous pouvons
envoyer la paire de points fermés $a, b$ sur toute paire de notre choix. Un choix
traditionnel consiste à prendre $a = 0$ et $b = \infty$.

\medskip\noindent
Plus généralement, soit $\alpha = \sum n_i [W_i]$ un $(k + 1)$-cycle sur
$X \times \mathbf{P}^1$. Soient $a_i, b_i$ des paires de points fermés distincts de
$\mathbf{P}^1$. Supposons que $X \times a_i$, $X \times b_i$ et $W_i$ se coupent
proprement ; autrement dit, chaque triplet $W_i, a_i, b_i$ vérifie la condition
étudiée ci-dessus. Un {\it cycle rationnellement équivalent à zéro} est tout cycle
de la forme
$$
\sum n_i([W_{i, a_i}]_k - [W_{i, b_i}]_k).
$$
Il s'agit bien d'un $k$-cycle. L'ensemble des $k$-cycles rationnellement
équivalents à zéro est un sous-groupe additif du groupe des $k$-cycles.
Nous disons que deux $k$-cycles sont {\it rationnellement équivalents}, et nous écrivons
$\alpha \sim_{rat} \alpha'$, si $\alpha - \alpha'$ est un cycle rationnellement
équivalent à zéro.

\medskip\noindent
Nous définissons
$$
\CH_k(X) = Z_k(X)/ \sim_{rat}
$$
comme le {\it groupe de Chow des $k$-cycles sur $X$}. Nous verrons dans le
lemme \ref{intersection-lemma-rational-equivalence}
que cela coïncide avec le groupe de Chow défini dans
Homologie de Chow, définition \ref{chow-definition-rational-equivalence}.


\section{Équivalence rationnelle et fonctions rationnelles}
\label{intersection-section-alternative}

\noindent
Soit $X$ une variété. Soit $W \subset X$ une sous-variété
de dimension $k + 1$. Soit $f \in \mathbf{C}(W)^*$ une fonction rationnelle non nulle
sur $W$. Pour toute sous-variété $Z \subset W$ de dimension $k$,
on peut définir l'ordre d'annulation $\text{ord}_{W, Z}(f)$ de $f$ le long de
$Z$. Si $f$ appartient à l'anneau local $\mathcal{O}_{W, Z}$,
on a alors
$$
\text{ord}_{W, Z}(f) =
\text{longueur}_{\mathcal{O}_{X, Z}} \mathcal{O}_{W, Z}/f\mathcal{O}_{W, Z}
$$
où $\mathcal{O}_{X, Z}$, resp.\ $\mathcal{O}_{W, Z}$, est l'anneau
local de $X$, resp.\ de $W$, au point générique de $Z$. En général, on
étend la définition par multiplicativité. Le {\it diviseur principal
associé à $f$} est
$$
\text{div}_W(f) = \sum \text{ord}_{W, Z}(f)[Z]
$$
dans $Z_k(W)$. Comme $W \subset X$ est une sous-variété fermée, nous pouvons considérer
$\text{div}_W(f)$ comme un cycle sur $X$.
Voir Homologie de Chow, section \ref{chow-section-principal-divisors}.

\begin{lemma}
\label{intersection-lemma-rational-equivalence}
Soit $X$ une variété. Soit $W \subset X$ une sous-variété
de dimension $k + 1$. Soit $f \in \mathbf{C}(W)^*$ une fonction rationnelle non nulle
sur $W$. Alors $\text{div}_W(f)$ est rationnellement équivalent à zéro sur
$X$. Réciproquement, ces diviseurs principaux engendrent le groupe abélien des
cycles rationnellement équivalents à zéro sur $X$.
\end{lemma}

\begin{proof}
La première assertion résulte de
Homologie de Chow, lemme \ref{chow-lemma-rational-function}.
Plus précisément, soit $W' \subset X \times \mathbf{P}^1$ l'adhérence
du graphe de $f$. Alors $\text{div}_W(f) = [W'_0]_k - [W'_\infty]$
dans $Z_k(W) \subset Z_k(X)$ ; voir la partie (6) de
Homologie de Chow, lemme \ref{chow-lemma-rational-function}.

\medskip\noindent
Pour la seconde assertion, soit $W' \subset X \times \mathbf{P}^1$ une sous-variété
fermée de dimension $k + 1$ qui domine $\mathbf{P}^1$.
Nous allons montrer que $[W'_0]_k - [W'_\infty]_k$ est un diviseur principal,
ce qui achèvera la démonstration. Soit $W \subset X$ l'image de $W'$
par la projection sur $X$. Alors $W \subset X$ est une sous-variété fermée
et $W' \to W$ est propre et dominant, à fibres de dimension $0$
ou $1$. Si $\dim(W) = k$, alors $W' = W \times \mathbf{P}^1$ et nous
voyons que $[W'_0]_k - [W'_\infty]_k = [W] - [W] = 0$.
Si $\dim(W) = k + 1$, alors
$W' \to W$ est génériquement fini\footnote{Si $W' \to W$ est birationnel,
le résultat résulte de
Homologie de Chow, lemme \ref{chow-lemma-rational-function}.
Notre tâche est de montrer que, même si $W' \to W$
est de degré $> 1$, l'équivalence rationnelle fondamentale
$[W'_0]_k \sim_{rat} [W'_\infty]_k$ provient d'un diviseur principal
sur une sous-variété de $X$.}. Notons $f$ la projection $W' \to \mathbf{P}^1$,
considérée comme un élément de $\mathbf{C}(W')^*$. Soit
$g = \text{Nm}(f) \in \mathbf{C}(W)^*$ la norme. D'après
Homologie de Chow, lemme \ref{chow-lemma-proper-pushforward-alteration},
nous avons
$$
\text{div}_W(g) = \text{pr}_{X, *}\text{div}_{W'}(f)
$$
Comme $\text{div}_{W'}(f) = [W'_0]_k - [W'_\infty]_k$
d'après Homologie de Chow, lemme \ref{chow-lemma-rational-function},
la démonstration est achevée.
\end{proof}


\section{Image directe propre et équivalence rationnelle}
\label{intersection-section-pushforward-and-rational-equivalence}

\noindent
Supposons que $f : X \to Y$ soit un morphisme propre de variétés.
Soit $\alpha \sim_{rat} 0$ un $k$-cycle sur
$X$ rationnellement équivalent à $0$. Alors l'{image directe}
de $\alpha$ est rationnellement équivalente à zéro :
$f_* \alpha \sim_{rat} 0$. Voir le chapitre I de \cite{F} ou
Homologie de Chow, lemme \ref{chow-lemma-proper-pushforward-rational-equivalence}.

\medskip\noindent
Nous obtenons donc un diagramme commutatif
$$
\xymatrix{
Z_k(X) \ar[r] \ar[d]_{f_*} & \CH_k(X) \ar[d]^{f_*} \\
Z_k(Y) \ar[r] & \CH_k(Y)
}
$$
de groupes de $k$-cycles.


\section{Image inverse plate et équivalence rationnelle}
\label{intersection-section-flat-pullback-and-rational-equivalence}

\noindent
Supposons que $f : X \to Y$ soit un morphisme plat de variétés.
Posons $r = \dim(X) - \dim(Y)$.
Soit $\alpha \sim_{rat} 0$ un $k$-cycle sur
$Y$ rationnellement équivalent à $0$. Alors l'image inverse
de $\alpha$ est rationnellement équivalente à zéro :
$f^* \alpha \sim_{rat} 0$. Voir le chapitre I de \cite{F} ou
Homologie de Chow, lemme \ref{chow-lemma-flat-pullback-rational-equivalence}.

\medskip\noindent
Nous obtenons donc un diagramme commutatif
$$
\xymatrix{
Z_{k + r}(X) \ar[r] & \CH_{k + r}(X) \\
Z_k(Y) \ar[r] \ar[u]^{f^*} & \CH_k(Y) \ar[u]_{f^*}
}
$$
de groupes de $k$-cycles.


\section{La suite exacte courte associée à un ouvert}
\label{intersection-section-ses}

\noindent
Soit $X$ une variété et soit $U \subset X$ une sous-variété ouverte.
Soit $X \setminus U = \bigcup Z_i$ sa décomposition en composantes
irréductibles\footnote{Comme, dans ce chapitre, nous ne considérons que les groupes de Chow
de variétés, nous ne pouvons pas prendre $Z_k(X \setminus U)$
et $\CH_k(X \setminus U)$ ; d'où l'approche qui emploie les variétés $Z_i$.}.
Alors, pour tout $k \geq 0$, il existe un diagramme commutatif
$$
\xymatrix{
\bigoplus Z_k(Z_i) \ar[r] \ar[d] &
Z_k(X) \ar[r] \ar[d] &
Z_k(U) \ar[d] \ar[r] &
0 \\
\bigoplus \CH_k(Z_i) \ar[r] &
\CH_k(X) \ar[r] &
\CH_k(U) \ar[r] &
0
}
$$
à lignes exactes. Les flèches verticales sont ici les applications canoniques de passage au quotient.
Les flèches horizontales de gauche sont données par l'image directe propre le long des
immersions fermées $Z_i \to X$. Les flèches horizontales de droite sont données par l'image
inverse plate le long de l'immersion ouverte $j : U \to X$. Comme nous avons vu que
ces applications se factorisent par l'équivalence rationnelle, les carrés
sont commutatifs. La ligne du haut est exacte simplement parce que toute sous-variété
de $X$ est soit contenue dans un des $Z_i$, soit a une intersection irréductible
avec $U$. La ligne du bas est exacte parce que tout diviseur principal
$\text{div}_W(f)$ sur $U$ est la restriction d'un diviseur principal sur $X$.
Plus précisément, si $W \subset U$ est une sous-variété fermée de dimension $(k + 1)$
et si $f \in \mathbf{C}(W)^*$, notons $\overline{W}$ l'adhérence de $W$
dans $X$. Alors $W \subset \overline{W}$ est une immersion ouverte, donc
$\mathbf{C}(W) = \mathbf{C}(\overline{W})$, et nous pouvons considérer $f$
comme une fonction rationnelle non constante sur $\overline{W}$. Alors, clairement,
$$
j^*\text{div}_{\overline{W}}(f) = \text{div}_W(f)
$$
dans $Z_k(X)$. L'exactitude de la ligne inférieure s'en déduit aisément.
Pour les détails, voir Homologie de Chow, lemme \ref{chow-lemma-restrict-to-open}.


\section{Intersections propres}
\label{intersection-section-intersect-properly}

\noindent
Commençons par quelques lemmes donnant des estimations de dimension.

\begin{lemma}
\label{intersection-lemma-dimension-product-varieties}
Soient $X$ et $Y$ des variétés. Alors $X \times Y$ est une variété et
$\dim(X \times Y) = \dim(X) + \dim(Y)$.
\end{lemma}

\begin{proof}
Le schéma $X \times Y = X \times_{\Spec(\mathbf{C})} Y$ est une variété d'après
Variétés, lemme \ref{varieties-lemma-product-varieties}.
L'assertion sur la dimension est
Variétés, lemme \ref{varieties-lemma-dimension-product-locally-algebraic}.
\end{proof}

\noindent
Rappelons qu'une immersion régulière $i : X \to Y$ de schémas
est une immersion fermée dont le
faisceau d'idéaux correspondant est localement engendré par une suite régulière ; voir
Diviseurs, section \ref{divisors-section-regular-immersions}.
De plus, le faisceau conormal $\mathcal{C}_{X/Y}$ est localement libre de type fini, de
rang égal à la longueur de la suite régulière. Nous dirons que $i$ est une
{\it immersion régulière de codimension $c$}
si $\mathcal{C}_{X/Y}$ est localement libre de rang $c$.

\medskip\noindent
Plus généralement, rappelons
(Morphismes, compléments, section \ref{more-morphisms-section-lci})
que $f : X \to Y$ est un morphisme d'intersection complète locale
si l'on peut recouvrir $X$ par des ouverts $U$ tels que $f|_U$ se factorise
comme suit :
$$
\xymatrix{
U \ar[rr]_i \ar[rd] & & \mathbf{A}^n_Y \ar[ld] \\
& Y
}
$$
où $i$ est une immersion Koszul-régulière (si $Y$ est localement noethérien,
cela revient à demander que $i$ soit une immersion régulière ; voir
Diviseurs, lemme \ref{divisors-lemma-regular-immersion-noetherian}).
Nous dirons que $f$ est un {\it morphisme d'intersection complète locale
de dimension relative $r$} si, pour toute factorisation comme ci-dessus,
l'immersion fermée $i$ a un faisceau conormal de rang $n - r$ (autrement
dit, si $i$ est une immersion Koszul-régulière de codimension $n - r$,
ce qui, dans le cas noethérien, signifie simplement qu'elle est une immersion régulière de
codimension $n - r$).

\begin{lemma}
\label{intersection-lemma-pullback-by-regular-immersion}
Soit $f : X \to Y$ un morphisme de variétés.
\begin{enumerate}
\item Si $Z \subset Y$ est une sous-variété de dimension $d$ et si $f$ est une
immersion régulière de codimension $c$, alors toute composante irréductible
de $f^{-1}(Z)$ est de dimension $\geq d - c$.
\item Si $Z \subset Y$ est une sous-variété de dimension $d$ et si
$f$ est un morphisme d'intersection complète locale de dimension relative $r$,
alors toute composante irréductible de $f^{-1}(Z)$ est de dimension $\geq d + r$.
\end{enumerate}
\end{lemma}

\begin{proof}
Démonstration de (1). Nous pouvons travailler localement ; nous pouvons donc supposer que
$Y = \Spec(A)$ et $X = V(f_1, \ldots, f_c)$, où $f_1, \ldots, f_c$
est une suite régulière dans $A$. Si $Z = \Spec(A/\mathfrak p)$, alors
nous voyons que $f^{-1}(Z) = \Spec(A/\mathfrak p + (f_1, \ldots, f_c))$.
Si $V$ est une composante irréductible de $f^{-1}(Z)$, nous pouvons
choisir un point fermé $v \in V$ qui n'appartient à aucune autre composante irréductible
de $f^{-1}(Z)$. Alors
$$
\dim(Z) = \dim \mathcal{O}_{Z, v}
\quad\text{et}\quad
\dim(V) = \dim \mathcal{O}_{V, v} = \dim \mathcal{O}_{Z, v}/(f_1, \ldots, f_c)
$$
La première égalité résulte par exemple de
Algèbre, lemme \ref{algebra-lemma-dimension-prime-polynomial-ring},
et la seconde de notre choix du point fermé.
Le résultat découle alors du fait que le quotient par un élément
de l'idéal maximal diminue la dimension d'au plus $1$ ; voir
Algèbre, lemme \ref{algebra-lemma-one-equation}.

\medskip\noindent
Démonstration de (2). Choisissons une factorisation comme dans la définition d'une
intersection complète locale et appliquons (1). Certains détails sont omis.
\end{proof}

\begin{lemma}
\label{intersection-lemma-diagonal-regular-immersion}
Soit $X$ une variété non singulière. Alors la diagonale
$\Delta : X \to X \times X$ est une immersion régulière de codimension $\dim(X)$.
\end{lemma}

\begin{proof}
En fait, toute immersion fermée entre variétés projectives non singulières
est une immersion régulière ; voir Diviseurs,
lemme \ref{divisors-lemma-immersion-smooth-into-smooth-regular-immersion}.
\end{proof}

\noindent
Le lemme suivant illustre le fonctionnement de la réduction à la diagonale.

\begin{lemma}
\label{intersection-lemma-intersect-in-smooth}
Soit $X$ une variété non singulière et soient $W,V \subset X$
des sous-variétés fermées telles que $\dim(W) = s$ et $\dim(V) = r$. Alors toute
composante irréductible $Z$ de $V \cap W$ est de dimension $\geq r + s - \dim(X)$.
\end{lemma}

\begin{proof}
Comme $V \cap W = \Delta^{-1}(V \times W)$ (au sens schématique),
nous concluons par les lemmes \ref{intersection-lemma-diagonal-regular-immersion} et
\ref{intersection-lemma-pullback-by-regular-immersion}.
\end{proof}

\noindent
Ce lemme suggère la définition suivante.

\begin{definition}
\label{intersection-definition-proper-intersection}
Soit $X$ une variété non singulière.
\begin{enumerate}
\item {\emergencystretch=1em Soient $W,V \subset X$ des sous-variétés fermées telles que
$\dim(W) = s$ et $\dim(V) = r$. Nous disons que $W$ et $V$
{\it se coupent proprement} si $\dim(V \cap W) \leq r + s - \dim(X)$.\par}
\item Soit $\alpha = \sum n_i [W_i]$ un $s$-cycle et soit
$\beta = \sum_j m_j [V_j]$ un $r$-cycle sur $X$. Nous disons
que $\alpha$ et $\beta$ {\it se coupent proprement} si
$W_i$ et $V_j$ se coupent proprement quels que soient $i$ et $j$.
\end{enumerate}
\end{definition}


\section{Multiplicités d'intersection au moyen de la formule des Tor}
\label{intersection-section-tor-formula}

\noindent
Un fait fondamental que nous utiliserons souvent est que, pour des faisceaux de
modules $\mathcal{F}$, $\mathcal{G}$ sur un espace annelé $(X, \mathcal{O}_X)$
et un point $x \in X$, on a
$$
\text{Tor}_p^{\mathcal{O}_X}(\mathcal{F}, \mathcal{G})_x =
\text{Tor}_p^{\mathcal{O}_{X, x}}(\mathcal{F}_x, \mathcal{G}_x)
$$
comme $\mathcal{O}_{X, x}$-modules. Cela se voit de plusieurs façons
à partir de notre construction des produits tensoriels dérivés dans
Cohomologie, section \ref{cohomology-section-flat}; cela résulte par exemple de
Cohomologie, lemme \ref{cohomology-lemma-check-K-flat-stalks}.
De plus, si $X$ est un schéma et si $\mathcal{F}$ et $\mathcal{G}$
sont quasi-cohérents, alors les modules
$\text{Tor}_p^{\mathcal{O}_X}(\mathcal{F}, \mathcal{G})$ sont eux aussi
quasi-cohérents, voir
Catégories dérivées des schémas, lemme
\ref{perfect-lemma-quasi-coherence-tensor-product}.
Le résultat suivant importe davantage pour ce qui nous occupe.

\begin{lemma}
\label{intersection-lemma-tensor-coherent}
Soit $X$ un schéma localement noethérien.
\begin{enumerate}
\item Si $\mathcal{F}$ et $\mathcal{G}$ sont des $\mathcal{O}_X$-modules cohérents,
alors $\text{Tor}_p^{\mathcal{O}_X}(\mathcal{F}, \mathcal{G})$ l'est aussi.
\item Si $L$ et $K$ appartiennent à $D^-_{\textit{Coh}}(\mathcal{O}_X)$, alors
il en est de même de $L \otimes_{\mathcal{O}_X}^\mathbf{L} K$.
\end{enumerate}
\end{lemma}

\begin{proof}
Expliquons comment démontrer (1) de manière plus élémentaire, et la partie (2)
à l'aide de la théorie générale développée précédemment.

\medskip\noindent
Démonstration de (1). Comme la formation de $\text{Tor}$ commute à la localisation,
nous pouvons supposer $X$ affine. Ainsi $X = \Spec(A)$ pour un anneau noethérien
$A$, et $\mathcal{F}$, $\mathcal{G}$ correspondent à des $A$-modules de type fini
$M$ et $N$ (Cohomologie des schémas, lemme
\ref{coherent-lemma-coherent-Noetherian}).
D'après Catégories dérivées des schémas, lemme
\ref{perfect-lemma-quasi-coherence-tensor-product}, nous pouvons
calculer les $\text{Tor}$ en calculant d'abord les $\text{Tor}$
de $M$ et $N$ sur $A$, puis en prenant le $\mathcal{O}_X$-module associé.
Comme les modules $\text{Tor}_p^A(M, N)$ sont de type fini d'après
Algèbre, lemme \ref{algebra-lemma-tor-noetherian},
nous concluons.

\medskip\noindent
D'après Catégories dérivées des schémas, lemme
\ref{perfect-lemma-identify-pseudo-coherent-noetherian},
l'hypothèse équivaut à demander que $L$ et $K$ soient
(localement) pseudo-cohérents. Alors $L \otimes_{\mathcal{O}_X}^\mathbf{L} K$
est pseudo-cohérent d'après
Cohomologie, lemme \ref{cohomology-lemma-tensor-pseudo-coherent}.
\end{proof}

\begin{lemma}
\label{intersection-lemma-compute-tor-nonsingular}
{\emergencystretch=2em Soit $X$ une variété non singulière.
Soient $\mathcal{F}$, $\mathcal{G}$ des $\mathcal{O}_X$-modules cohérents.
Le $\mathcal{O}_X$-module
$\text{Tor}_p^{\mathcal{O}_X}(\mathcal{F}, \mathcal{G})$
est cohérent, a pour germe en $x$
$\text{Tor}_p^{\mathcal{O}_{X, x}}(\mathcal{F}_x, \mathcal{G}_x)$,
a son support contenu dans
$\text{Supp}(\mathcal{F}) \cap \text{Supp}(\mathcal{G})$, et
n'est non nul que pour $p \in \{0, \ldots, \dim(X)\}$.\par}
\end{lemma}

\begin{proof}
Le résultat sur les fibres a été discuté ci-dessus et implique la condition
sur le support. Les $\text{Tor}$ sont cohérents d'après
le lemme \ref{intersection-lemma-tensor-coherent}. L'annulation des $\text{Tor}$
en degrés négatifs résulte immédiatement de la construction. L'annulation
de $\text{Tor}_p$ pour $p > \dim(X)$ se voit comme suit :
les anneaux locaux $\mathcal{O}_{X, x}$ sont réguliers
(puisque $X$ est non singulière) et de dimension $\leq \dim(X)$
(Algèbre, lemme \ref{algebra-lemma-dimension-prime-polynomial-ring});
ainsi $\mathcal{O}_{X, x}$ est de dimension globale finie $\leq \dim(X)$
(Algèbre, lemme \ref{algebra-lemma-finite-gl-dim-finite-dim-regular}),
ce qui implique que les groupes $\text{Tor}$ de modules s'annulent au-delà de la dimension
(Compléments d'algèbre, lemme \ref{more-algebra-lemma-finite-gl-dim-tor-dimension}).
\end{proof}

\noindent
Soit $X$ une variété non singulière, et soient $W, V \subset X$
des sous-variétés fermées telles que $\dim(W) = s$ et $\dim(V) = r$.
Supposons que $V$ et $W$ se coupent proprement.
Dans ce cas, le lemme \ref{intersection-lemma-intersect-in-smooth} nous dit que toutes les composantes
irréductibles de $V \cap W$ ont une dimension égale à $r + s - \dim(X)$.
Les faisceaux $\text{Tor}_j^{\mathcal{O}_X}(\mathcal{O}_W, \mathcal{O}_V)$ sont
cohérents, à support dans $V \cap W$, et nuls si $j < 0$ ou $j > \dim(X)$
(lemme \ref{intersection-lemma-compute-tor-nonsingular}).
Nous définissons le {\it produit d'intersection} par
$$
W \cdot V = \sum\nolimits_i (-1)^i
[\text{Tor}_i^{\mathcal{O}_X}(\mathcal{O}_W, \mathcal{O}_V)]_{r + s - \dim(X)}.
$$
Soulignons que ceci n'a de sens qu'en vertu de notre hypothèse que
$V$ et $W$ se coupent proprement. Ce fait rendra nécessaire un lemme de
déplacement afin de définir le produit d'intersection en général.

\medskip\noindent
Avec cette notation, le cycle $V \cdot W$ est une combinaison linéaire
formelle $\sum e_Z Z$ des composantes irréductibles $Z$
de l'intersection $V \cap W$. Les entiers $e_Z$ sont appelés
les {\it multiplicités d'intersection}
$$
e_Z = e(X, V \cdot W, Z) =
\sum\nolimits_i
(-1)^i
\text{longueur}_{\mathcal{O}_{X, Z}}
\text{Tor}_i^{\mathcal{O}_{X, Z}}(\mathcal{O}_{W, Z}, \mathcal{O}_{V, Z})
$$
où $\mathcal{O}_{X, Z}$, resp.\ $\mathcal{O}_{W, Z}$,
resp.\ $\mathcal{O}_{V, Z}$ désigne l'anneau local de $X$, resp.\ $W$,
resp.\ $V$ au point générique de $Z$.
Ces sommes alternées des longueurs des $\text{Tor}$ satisfont à de nombreuses bonnes
propriétés, comme nous le verrons plus loin.

\medskip\noindent
Dans le cas d'intersections transversales, le nombre d'intersection vaut $1$.

\begin{lemma}
\label{intersection-lemma-transversal}
Soit $X$ une variété non singulière. Soient $V, W \subset X$ des
sous-variétés fermées qui se coupent proprement. Soit $Z$ une composante irréductible
de $V \cap W$, et supposons que la multiplicité
(au sens de la section \ref{intersection-section-cycle-of-closed}) de $Z$
dans le sous-schéma fermé $V \cap W$ vaille $1$.
Alors $e(X, V \cdot W, Z) = 1$, et $V$ et $W$ sont lisses
en un point général de $Z$.
\end{lemma}

\begin{proof}
Soit $(A, \mathfrak m, \kappa) =
(\mathcal{O}_{X, \xi}, \mathfrak m_\xi, \kappa(\xi))$, où $\xi \in Z$
est le point générique. Alors $\dim(A) = \dim(X) - \dim(Z)$, voir
Variétés, lemme \ref{varieties-lemma-dimension-locally-algebraic}.
Soient $I, J \subset A$ les idéaux définissant les traces de $V$ et $W$
dans $\Spec(A)$. Posons $\overline{I} = I + \mathfrak m^2/\mathfrak m^2$.
Alors $\dim_\kappa \overline{I} \leq \dim(X) - \dim(V)$, avec égalité
si et seulement si $A/I$ est régulier (cela résulte du lemme cité
ci-dessus et de la définition des anneaux réguliers, voir
Algèbre, définition \ref{algebra-definition-regular-local}
et la discussion qui la précède). De même pour $\overline{J}$.
Si la multiplicité vaut $1$, alors
$\text{longueur}_A(A/I + J) = 1$, donc $I + J = \mathfrak m$, et donc
$\overline{I} + \overline{J} = \mathfrak m/\mathfrak m^2$.
On a alors égalité partout (puisque l'intersection est
propre). Nous trouvons donc $f_1, \ldots, f_a \in I$ et $g_1, \ldots g_b \in J$
tels que $\overline{f}_1, \ldots, \overline{g}_b$ soient une base
de $\mathfrak m/\mathfrak m^2$. Alors $f_1, \ldots, g_b$ est un
système régulier de paramètres et une suite régulière
(Algèbre, lemme \ref{algebra-lemma-regular-ring-CM}).
Le même lemme montre que $A/(f_1, \ldots, f_a)$ est un anneau local régulier
de dimension $\dim(X) - \dim(V)$ ; ainsi $A/(f_1, \ldots, f_a) \to A/I$
est un isomorphisme (si le noyau est non nul, alors la dimension
de $A/I$ est strictement plus petite, voir
Algèbre, lemmes \ref{algebra-lemma-regular-domain} et
\ref{algebra-lemma-one-equation}).
Nous concluons que $I = (f_1, \ldots, f_a)$ et $J = (g_1, \ldots, g_b)$
par symétrie. Ainsi, le complexe de Koszul $K_\bullet(A, f_1, \ldots, f_a)$
sur $f_1, \ldots, f_a$ est une résolution de $A/I$, voir
Compléments d'algèbre, lemme \ref{more-algebra-lemma-regular-koszul-regular}.
Par conséquent,
\begin{align*}
\text{Tor}_p^A(A/I, A/J)
& =
H_p(K_\bullet(A, f_1, \ldots, f_a) \otimes_A A/J) \\
& =
H_p(K_\bullet(A/J, f_1 \bmod J, \ldots, f_a \bmod J))
\end{align*}
Puisque nous avons vu ci-dessus que $f_1 \bmod J, \ldots, f_a \bmod J$ est
un système régulier de paramètres de l'anneau local régulier $A/J$,
nous concluons qu'il n'y a qu'un seul groupe de cohomologie, à savoir
$H_0 = A/(I + J) = \kappa$. Ceci achève la démonstration.
\end{proof}

\begin{example}
\label{intersection-example-naive-multiplicity-wrong}
Dans cet exemple, nous montrons qu'il est nécessaire d'utiliser les
Tor supérieurs dans la formule des multiplicités d'intersection ci-dessus.
Soit $X$ une variété non singulière de dimension $4$.
Soit $p \in X$ un point fermé. Soient $V, W \subset X$
des sous-variétés fermées de $X$. Supposons qu'il existe un
isomorphisme
$$
\mathcal{O}_{X, p}^\wedge \cong \mathbf{C}[[x, y, z, w]]
$$
tel que l'idéal de $V$ soit $(xz, xw, yz, yw)$ et l'idéal
de $W$ soit $(x - z, y - w)$. Un calcul montre alors que
$$
\text{longueur}\ \mathbf{C}[[x, y, z, w]]/
(xz, xw, yz, yw, x - z, y - w) = 3
$$
D'autre part, la multiplicité $e(X, V \cdot W, p) = 2$,
comme on le voit du fait qu'au niveau formel local, $V$ est la
réunion de deux plans lisses $x = y = 0$ et $z = w = 0$ en $p$,
dont chacun a une multiplicité d'intersection $1$ avec le plan
$x - z = y - w = 0$ (lemme \ref{intersection-lemma-transversal}). Pour obtenir un
véritable exemple, prenons
un morphisme général $f : \mathbf{P}^2 \to \mathbf{P}^4$ donné par
$5$ polynômes homogènes de degré $> 1$. L'image
$V \subset \mathbf{P}^4 = X$ aura des singularités du type
décrit ci-dessus, car il existera $p_1, p_2 \in \mathbf{P}^2$
tels que $f(p_1) = f(p_2)$. Pour trouver $W$, prenons un plan général passant
par un tel point.
\end{example}






\section{Multiplicités algébriques}
\label{intersection-section-multiplicities}

\noindent
Soit $(A, \mathfrak m, \kappa)$ un anneau local noethérien.
Soit $M$ un $A$-module de type fini, et soit $I \subset A$ un idéal
de définition (Algèbre, définition \ref{algebra-definition-ideal-definition}).
Rappelons que la fonction
$$
\chi_{I, M}(n) = \text{longueur}_A(M/I^nM) =
\sum\nolimits_{p = 0, \ldots, n - 1} \text{longueur}_A(I^pM/I^{p + 1}M)
$$
est un polynôme numérique
(Algèbre, proposition \ref{algebra-proposition-hilbert-function-polynomial}).
Le degré de ce polynôme est égal à $\dim(\text{Supp}(M))$ d'après
Algèbre, lemme \ref{algebra-lemma-support-dimension-d}.

\begin{definition}
\label{intersection-definition-multiplicity}
Dans la situation ci-dessus, supposons $d \geq \dim(\text{Supp}(M))$.
Dans ce cas, si $d > \dim(\text{Supp}(M))$, nous posons $e_I(M, d) = 0$,
et si $d = \dim(\text{Supp}(M))$, nous posons $e_I(M, d)$ égal à $d!$
fois le coefficient dominant du polynôme numérique $\chi_{I, M}$.
Ainsi, dans les deux cas, nous avons
$$
\chi_{I, M}(n) \sim e_I(M, d) \frac{n^d}{d!} + \text{termes d'ordre inférieur}
$$
La {\it multiplicité de $M$ pour l'idéal de définition $I$}
est $e_I(M) = e_I(M, \dim(\text{Supp}(M)))$.
\end{definition}

\noindent
Nous avons les propriétés suivantes pour ces multiplicités.

\begin{lemma}
\label{intersection-lemma-multiplicity-ses}
Soit $A$ un anneau local noethérien. Soit $I \subset A$ un idéal de
définition. Soit $0 \to M' \to M \to M'' \to 0$ une suite exacte courte
de $A$-modules de type fini. Soit $d \geq \dim(\text{Supp}(M))$. Alors
$$
e_I(M, d) = e_I(M', d) + e_I(M'', d)
$$
\end{lemma}

\begin{proof}
Cela résulte immédiatement des définitions et de
Algèbre, lemme \ref{algebra-lemma-hilbert-ses-chi}.
\end{proof}

\begin{lemma}
\label{intersection-lemma-multiplicity-as-a-sum}
Soit $A$ un anneau local noethérien. Soit $I \subset A$ un idéal de
définition. Soit $M$ un $A$-module de type fini. Soit $d \geq \dim(\text{Supp}(M))$.
Alors
$$
e_I(M, d) =
\sum \text{longueur}_{A_\mathfrak p}(M_\mathfrak p) e_I(A/\mathfrak p, d)
$$
où la somme porte sur les idéaux premiers $\mathfrak p \subset A$ tels que
$\dim(A/\mathfrak p) = d$.
\end{lemma}

\begin{proof}
Le membre de gauche et le membre de droite sont tous deux additifs dans les suites
exactes courtes de modules de dimension $\leq d$, voir
le lemme \ref{intersection-lemma-multiplicity-ses} et
Algèbre, lemme \ref{algebra-lemma-length-additive}.
Ainsi, d'après Algèbre, lemme \ref{algebra-lemma-filter-Noetherian-module},
il suffit de démontrer ceci lorsque $M = A/\mathfrak q$ pour un
idéal premier $\mathfrak q$ de $A$ tel que $\dim(A/\mathfrak q) \leq d$.
Ce cas est évident.
\end{proof}

\begin{lemma}
\label{intersection-lemma-leading-coefficient}
Soit $P$ un polynôme de degré $r$ et de coefficient dominant $a$.
Alors
$$
r! a = \sum\nolimits_{i = 0, \ldots, r} (-1)^i{r \choose i} P(t - i)
$$
pour tout $t$.
\end{lemma}

\begin{proof}
Notons $\Delta$ l'opérateur qui, à un polynôme $P$, associe
le polynôme $\Delta(P) = P(t) - P(t - 1)$. Nous affirmons que
$$
\Delta^r(P) = \sum\nolimits_{i = 0, \ldots, r} (-1)^i {r \choose i} P(t - i)
$$
C'est vrai pour $r = 0, 1$ par inspection. Supposons que ce soit vrai pour $r$.
Alors nous calculons
\begin{align*}
\Delta^{r + 1}(P)
& =
\sum\nolimits_{i = 0, \ldots, r} (-1)^i {r \choose i} \Delta(P)(t - i) \\
& =
\sum\nolimits_{n = -r, \ldots, 0} (-1)^i {r \choose i}
(P(t - i) - P(t - i - 1))
\end{align*}
L'affirmation résulte donc de l'égalité
$$
{r + 1 \choose i} = {r \choose i} + {r \choose i - 1}
$$
Le lemme résulte du fait que $\Delta(P)$ est de degré $r - 1$
et de coefficient dominant $ra$ si le degré de $P$ est $r$.
\end{proof}

\noindent
Un fait important est que l'on peut calculer la multiplicité à l'aide
du complexe de Koszul. Rappelons que, si $R$ est un anneau et si
$f_1, \ldots, f_r \in R$, alors $K_\bullet(f_1, \ldots, f_r)$
désigne le complexe de Koszul, voir
Compléments d'algèbre, section \ref{more-algebra-section-koszul}.

\begin{theorem}
\label{intersection-theorem-multiplicity-with-koszul}
\begin{reference}
\cite[théorème 1, partie B du chapitre IV]{Serre_algebre_locale}
\end{reference}
Soit $A$ un anneau local noethérien. Soit $I = (f_1, \ldots, f_r) \subset A$
un idéal de définition. Soit $M$ un $A$-module de type fini. Alors
$$
e_I(M, r) = \sum
(-1)^i\text{longueur}_A H_i(K_\bullet(f_1, \ldots, f_r) \otimes_A M)
$$
\end{theorem}

\begin{proof}
Transformons le complexe de Koszul $K_\bullet(f_1, \ldots, f_r)$ en un complexe de
cochaînes $K^\bullet$ en posant $K^n = K_{-n}(f_1, \ldots, f_r)$.
Alors $K^\bullet$ est placé en degrés $-r, \ldots, 0$ et
$H^i(K^\bullet \otimes_A M) = H_{-i}(K_\bullet(f_1, \ldots, f_r) \otimes_A M)$.
L'énoncé du théorème a un sens, car les modules
$H^i(K^\bullet \otimes M)$ sont annulés par $f_1, \ldots, f_r$
(Compléments d'algèbre, lemme \ref{more-algebra-lemma-homotopy-koszul})
et sont donc de longueur finie.
Définissons une filtration du complexe $K^\bullet$ en posant
$$
F^p(K^n \otimes_A M) =
I^{\max(0, p + n)}(K^n \otimes_A M),\quad p \in \mathbf{Z}
$$
Comme $f_i I^p \subset I^{p + 1}$, c'est une filtration par sous-complexes.
Nous avons ainsi un complexe filtré et obtenons une suite spectrale, voir
Homologie, section \ref{homology-section-filtered-complex}.
Nous avons
$$
E_0 = \bigoplus\nolimits_{p, q} E_0^{p, q} =
\bigoplus\nolimits_{p, q} \text{gr}^p(K^{p + q} \otimes_A M) =
\text{Gr}_I(K^\bullet \otimes_A M)
$$
Comme $K^n$ est libre de type fini, nous avons
$$
\text{Gr}_I(K^\bullet \otimes_A M) =
\text{Gr}_I(K^\bullet) \otimes_{\text{Gr}_I(A)} \text{Gr}_I(M)
$$
Remarquons que $\text{Gr}_I(K^\bullet)$ est le complexe de Koszul
sur $\text{Gr}_I(A)$ associé aux éléments
$\overline{f}_1, \ldots, \overline{f}_r \in I/I^2$.
Un calcul simple (omis)
montre que la différentielle $d_0$ sur $E_0$
coïncide avec la différentielle provenant du complexe de Koszul.
Comme $\text{Gr}_I(M)$ est un $\text{Gr}_I(A)$-module de type fini
et comme $\text{Gr}_I(A)$ est noethérien (en tant que quotient
de $A/I[x_1, \ldots, x_r]$ par $x_i \mapsto \overline{f}_i$), le
module de cohomologie $E_1 = \bigoplus E_1^{p, q}$
est un $\text{Gr}_I(A)$-module de type fini. Toutefois, comme ci-dessus,
$E_1$ est annulé par $\overline{f}_1, \ldots, \overline{f}_r$.
Nous concluons que $E_1$ est de longueur finie.
En particulier, nous trouvons que $\text{Gr}^p_F(K^\bullet \otimes M)$ est
acyclique pour $p \gg 0$.

\medskip\noindent
Vérifions ensuite que la suite spectrale ci-dessus converge
en utilisant Homologie, lemme \ref{homology-lemma-filtered-complex-ss-converges}.
Les égalités requises résultent aisément du lemme d'Artin-Rees
sous la forme énoncée dans Algèbre, lemme \ref{algebra-lemma-map-AR}.
Nous voyons ainsi que
\begin{align*}
\sum (-1)^i\text{longueur}_A(H^i(K^\bullet \otimes_A M))
& =
\sum (-1)^{p + q} \text{longueur}_A(E_\infty^{p, q}) \\
& =
\sum (-1)^{p + q} \text{longueur}_A(E_1^{p, q})
\end{align*}
car, comme nous l'avons vu ci-dessus, la longueur de $E_1$ est finie
(cela utilise bien sûr l'additivité des longueurs). Choisissons $t$ assez
grand pour que $\text{Gr}^p_F(K^\bullet \otimes M)$
soit acyclique pour $p \geq t$ (voir ci-dessus). En utilisant
de nouveau l'additivité, nous voyons que
$$
\sum (-1)^{p + q} \text{longueur}_A(E_1^{p, q}) =
\sum\nolimits_n \sum\nolimits_{p \leq t}
(-1)^n \text{longueur}_A(\text{gr}^p(K^n \otimes_A M))
$$
Ceci est égal à
$$
\sum\nolimits_{n = -r, \ldots, 0} (-1)^n{r \choose |n|} \chi_{I, M}(t + n)
$$
par le choix de la filtration ci-dessus et la définition de $\chi_{I, M}$ dans
Algèbre, section \ref{algebra-section-Noetherian-local}.
Le lemme résulte du lemme \ref{intersection-lemma-leading-coefficient}
et de la définition de $e_I(M, r)$.
\end{proof}

\begin{remark}[Généralisation triviale]
\label{intersection-remark-trivial-generalization}
Soit $(A, \mathfrak m, \kappa)$ un anneau local noethérien.
Soit $M$ un $A$-module de type fini. Soit $I \subset A$ un idéal.
Les conditions suivantes sont équivalentes :
\begin{enumerate}
\item $I' = I + \text{Ann}(M)$ est un idéal de définition
(Algèbre, définition \ref{algebra-definition-ideal-definition}),
\item l'image $\overline{I}$ de $I$ dans $\overline{A} = A/\text{Ann}(M)$
est un idéal de définition,
\item $\text{Supp}(M/IM) \subset \{\mathfrak m\}$,
\item $\dim(\text{Supp}(M/IM)) \leq 0$, et
\item $\text{longueur}_A(M/IM) < \infty$.
\end{enumerate}
Ceci résulte d'Algèbre, lemme \ref{algebra-lemma-support-point}
(détails omis). Si ces conditions sont satisfaites, on a $M/I^nM = M/(I')^nM$
pour tout $n$, et $M/I^nM = M/\overline{I}^nM$ pour tout $n$
si l'on considère $M$ comme un $\overline{A}$-module.
Nous pouvons donc définir
$$
\chi_{I, M}(n) = \text{longueur}_A(M/I^nM) =
\sum\nolimits_{p = 0, \ldots, n - 1} \text{longueur}_A(I^pM/I^{p + 1}M)
$$
et nous obtenons
$$
\chi_{I, M}(n) = \chi_{I', M}(n) = \chi_{\overline{I}, M}(n)
$$
pour tout $n$ par les égalités ci-dessus.
Tous les résultats d'Algèbre, section \ref{algebra-section-Noetherian-local},
ainsi que tous les résultats de la présente section, ont des analogues dans ce cadre.
En particulier, nous pouvons définir les multiplicités $e_I(M, d)$ pour
$d \geq \dim(\text{Supp}(M))$, et nous avons
$$
\chi_{I, M}(n) \sim e_I(M, d) \frac{n^d}{d!} + \text{termes d'ordre inférieur}
$$
comme dans le cas où $I$ est un idéal de définition.
\end{remark}



\section{Calcul des multiplicités d'intersection}
\label{intersection-section-computing-intersection-multiplicities}

\noindent
Dans cette section, nous étudions certains cas où les multiplicités d'intersection
peuvent se calculer par d'autres moyens. Voici un premier exemple.

\begin{lemma}
\label{intersection-lemma-intersection-multiplicity-CM}
Soit $X$ une variété non singulière, et soient $W, V \subset X$ des
sous-variétés fermées qui se coupent proprement. Soit $Z$ une composante irréductible
de $V \cap W$, de point générique $\xi$. Supposons que $\mathcal{O}_{W, \xi}$
et $\mathcal{O}_{V, \xi}$ soient de Cohen-Macaulay. Alors
$$
e(X, V \cdot W, Z) =
\text{longueur}_{\mathcal{O}_{X, \xi}}(\mathcal{O}_{V \cap W, \xi})
$$
où $V \cap W$ est l'intersection schématique.
En particulier, si $V$ et $W$ sont toutes deux de Cohen-Macaulay, alors
$V \cdot W = [V \cap W]_{\dim(V) + \dim(W) - \dim(X)}$.
\end{lemma}

\begin{proof}
Posons $A = \mathcal{O}_{X, \xi}$, $B = \mathcal{O}_{V, \xi}$, et
$C = \mathcal{O}_{W, \xi}$. D'après Auslander-Buchsbaum
(Algèbre, proposition \ref{algebra-proposition-Auslander-Buchsbaum}),
nous pouvons trouver une résolution finie $F_\bullet \to B$ par des modules libres
de type fini, de longueur
$$
\text{profondeur}(A) - \text{profondeur}(B) =
\dim(A) - \dim(B) = \dim(C)
$$
La première égalité vient de ce que $A$ et $B$ sont de Cohen-Macaulay, et la seconde
de ce que $V$ et $W$ se coupent proprement. Alors $F_\bullet \otimes_A C$ est un
complexe de modules libres de type fini représentant $B \otimes_A^\mathbf{L} C$;
ses modules de cohomologie ont donc leur support dans $\{\mathfrak m_A\}$.
D'après le lemme d'acyclicité (Algèbre, lemme \ref{algebra-lemma-acyclic}),
qui s'applique puisque $C$ est de Cohen-Macaulay,
nous concluons que $F_\bullet \otimes_A C$ n'a de
cohomologie non nulle qu'en degré $0$. Ceci achève la démonstration.
\end{proof}

\begin{lemma}
\label{intersection-lemma-one-ideal-ci}
Soit $A$ un anneau local noethérien. Soit $I = (f_1, \ldots, f_r)$ un idéal
engendré par une suite régulière. Soit $M$ un $A$-module de type fini. Supposons que
$\dim(\text{Supp}(M/IM)) = 0$. Alors
$$
e_I(M, r) = \sum (-1)^i\text{longueur}_A(\text{Tor}_i^A(A/I, M))
$$
Ici, $e_I(M, r)$ est défini comme dans la remarque \ref{intersection-remark-trivial-generalization}.
\end{lemma}

\begin{proof}
Puisque $f_1, \ldots, f_r$ est une suite régulière, le complexe de Koszul
$K_\bullet(f_1, \ldots, f_r)$ est une résolution de $A/I$ sur $A$, voir
Compléments d'algèbre, lemme
\ref{more-algebra-lemma-noetherian-finite-all-equivalent}.
Ainsi, le membre de droite est égal à
$$
\sum (-1)^i\text{longueur}_A H_i(K_\bullet(f_1, \ldots, f_r) \otimes_A M)
$$
Le résultat découle alors immédiatement du
théorème \ref{intersection-theorem-multiplicity-with-koszul} si $I$ est un idéal
de définition. En général, nous remplaçons $A$ par $\overline{A} = A/\text{Ann}(M)$
et $f_1, \ldots, f_r$ par $\overline{f}_1, \ldots, \overline{f}_r$,
ce qui est permis parce que
$$
K_\bullet(f_1, \ldots, f_r) \otimes_A M =
K_\bullet(\overline{f}_1, \ldots, \overline{f}_r) \otimes_{\overline{A}} M
$$
Puisque $e_I(M, r) = e_{\overline{I}}(M, r)$, où
$\overline{I} = (\overline{f}_1, \ldots, \overline{f}_r) \subset \overline{A}$
est un idéal de définition, le résultat découle aussi dans ce cas du
théorème \ref{intersection-theorem-multiplicity-with-koszul}.
\end{proof}

\begin{lemma}
\label{intersection-lemma-multiplicity-with-lci}
Soit $X$ une variété non singulière. Soient $W,V \subset X$ des
sous-variétés fermées qui se coupent proprement. Soit $Z$ une composante irréductible
de $V \cap W$, de point générique $\xi$.
Supposons que l'idéal de $V$ dans $\mathcal{O}_{X, \xi}$ soit défini par
une suite régulière $f_1, \ldots, f_c \in \mathcal{O}_{X, \xi}$.
Alors $e(X, V\cdot W, Z)$ est égal à $c!$ fois le coefficient dominant du
polynôme de Hilbert
$$
t \mapsto \text{longueur}_{\mathcal{O}_{X, \xi}}
\mathcal{O}_{W, \xi}/(f_1, \ldots, f_c)^t,\quad t \gg 0.
$$
En particulier, ce coefficient est $> 0$.
\end{lemma}

\begin{proof}
L'égalité
$$
e(X, V\cdot W, Z) = e_{(f_1, \ldots, f_c)}(\mathcal{O}_{W, \xi}, c)
$$
résulte du lemme plus général \ref{intersection-lemma-one-ideal-ci}.
Pour voir que $e_{(f_1, \ldots, f_c)}(\mathcal{O}_{W, \xi}, c)$ est
$> 0$, ou, de manière équivalente, que $e_{(f_1, \ldots, f_c)}(\mathcal{O}_{W, \xi}, c)$
est le coefficient dominant du polynôme de Hilbert,
il suffit de montrer que la
dimension de $\mathcal{O}_{W, \xi}$ vaut $c$, car le degré du
polynôme de Hilbert est égal à la dimension d'après
Algèbre, proposition \ref{algebra-proposition-dimension}.
Écrivons $\dim(V) = r$, $\dim(W) = s$ et $\dim(X) = n$. Alors
$\dim(Z) = r + s - n$ puisque l'intersection est propre. Ainsi,
le degré de transcendance de $\kappa(\xi)$ sur $\mathbf{C}$ est
$r + s - n$, voir Algèbre, lemme
\ref{algebra-lemma-dimension-prime-polynomial-ring}.
Nous avons $r + c = n$ parce que $V$ est défini par une suite régulière
dans un voisinage de $\xi$, voir
Diviseurs, lemme \ref{divisors-lemma-Noetherian-scheme-regular-ideal},
et alors le lemme \ref{intersection-lemma-pullback-by-regular-immersion}
s'applique (par exemple). Ainsi,
$$
\dim(\mathcal{O}_{W, \xi}) = s - (r + s - n) = s - ((n - c) + s - n) = c
$$
la première égalité résultant d'Algèbre, lemme
\ref{algebra-lemma-dimension-at-a-point-finite-type-field}.
\end{proof}

\begin{lemma}
\label{intersection-lemma-multiplicity-with-effective-Cartier-divisor}
Dans le lemme \ref{intersection-lemma-multiplicity-with-lci}, supposons que $c = 1$, c'est-à-dire que $V$
est un diviseur de Cartier effectif. Alors
$$
e(X, V \cdot W, Z) =
\text{longueur}_{\mathcal{O}_{X, \xi}}
(\mathcal{O}_{W, \xi}/f_1\mathcal{O}_{W, \xi}).
$$
\end{lemma}

\begin{proof}
Dans ce cas, l'image de $f_1$ dans $\mathcal{O}_{W, \xi}$ est non nulle par
propreté de l'intersection, donc est un non-diviseur de zéro. De plus,
$\mathcal{O}_{W, \xi}$ est un anneau local noethérien intègre de dimension $1$.
Ainsi,
$$
\text{longueur}_{\mathcal{O}_{X, \xi}}
(\mathcal{O}_{W, \xi}/f_1^t\mathcal{O}_{W, \xi}) =
t \text{longueur}_{\mathcal{O}_{X, \xi}}
(\mathcal{O}_{W, \xi}/f_1\mathcal{O}_{W, \xi})
$$
pour tout $t \geq 1$, voir Algèbre, lemme \ref{algebra-lemma-ord-additive}.
Ceci démontre le lemme.
\end{proof}

\begin{lemma}
\label{intersection-lemma-multiplicity-lci-CM}
Dans le lemme \ref{intersection-lemma-multiplicity-with-lci}, supposons que
l'anneau local $\mathcal{O}_{W, \xi}$ soit de Cohen-Macaulay. Alors nous
avons
$$
e(X, V \cdot W, Z) =
\text{longueur}_{\mathcal{O}_{X, \xi}} (\mathcal{O}_{W, \xi}/
f_1\mathcal{O}_{W, \xi} + \ldots + f_c\mathcal{O}_{W, \xi}).
$$
\end{lemma}

\begin{proof}
Ceci résulte immédiatement du lemme \ref{intersection-lemma-intersection-multiplicity-CM}.
On peut aussi le déduire du lemme \ref{intersection-lemma-multiplicity-with-lci}.
En effet, d'après Algèbre, lemme \ref{algebra-lemma-reformulate-CM},
nous voyons que $f_1, \ldots, f_c$ est une suite régulière dans
$\mathcal{O}_{W, \xi}$. Alors
Algèbre, lemme \ref{algebra-lemma-regular-quasi-regular} montre que
$f_1, \ldots, f_c$ est une suite quasi-régulière.
Ceci implique facilement que la longueur de
$\mathcal{O}_{W, \xi}/(f_1, \ldots, f_c)^t$ vaut
$$
{c + t \choose c}
\text{longueur}_{\mathcal{O}_{X, \xi}} (\mathcal{O}_{W, \xi}/
f_1\mathcal{O}_{W, \xi} + \ldots + f_c\mathcal{O}_{W, \xi}).
$$
L'examen du coefficient dominant permet de conclure.
\end{proof}


\section{Produit d'intersection à l'aide de la formule des Tor}
\label{intersection-section-intersection-product}

\noindent
Soit $X$ une variété non singulière. Soit
$\alpha = \sum n_i [W_i]$ un $r$-cycle et
$\beta = \sum_j m_j [V_j]$ un $s$-cycle sur $X$.
Supposons que $\alpha$ et $\beta$ se coupent proprement, voir
la définition \ref{intersection-definition-proper-intersection}.
Dans ce cas, nous posons
$$
\alpha \cdot \beta = \sum\nolimits_{i,j} n_i m_j W_i \cdot V_j.
$$
où $W_i \cdot V_j$ est défini comme à la section \ref{intersection-section-tor-formula}.
Si $\beta = [V]$, où $V$ est une sous-variété fermée de dimension $s$,
nous écrivons parfois $\alpha \cdot \beta = \alpha \cdot V$.

\begin{lemma}
\label{intersection-lemma-rational-equivalence-and-intersection}
Soit $X$ une variété non singulière. Soient $a, b \in \mathbf{P}^1$
deux points fermés distincts. Soit $k \geq 0$.
\begin{enumerate}
\item Si $W \subset X \times \mathbf{P}^1$ est une sous-variété fermée
de dimension $k + 1$ qui coupe $X \times a$ proprement, alors
\begin{enumerate}
\item $[W_a]_k = W \cdot X \times a$ comme cycles sur $X \times \mathbf{P}^1$, et
\item $[W_a]_k = \text{pr}_{X, *}(W \cdot X \times a)$ comme cycles sur $X$.
\end{enumerate}
\item Soit $\alpha$ un $(k + 1)$-cycle sur $X \times \mathbf{P}^1$
qui coupe $X \times a$ et $X \times b$ proprement. Alors
$pr_{X,*}( \alpha \cdot X \times a - \alpha \cdot X \times b)$
est rationnellement équivalent à zéro.
\item {\emergencystretch=1em Réciproquement, tout $k$-cycle rationnellement
équivalent à $0$ est de cette forme.\par}
\end{enumerate}
\end{lemma}

\begin{proof}
Observons d'abord que $X \times a$ est un diviseur de Cartier effectif de
$X \times \mathbf{P}^1$ et que $W_a$ est l'intersection schématique
de $W$ avec $X \times a$. Par conséquent, l'égalité de (1)(a) résulte
immédiatement des définitions et du calcul de la multiplicité d'intersection
dans le cas d'un diviseur de Cartier donné dans le
lemme \ref{intersection-lemma-multiplicity-with-effective-Cartier-divisor}.
L'assertion (1)(b) est vraie parce que $W_a \to X \times \mathbf{P}^1 \to X$ envoie
isomorphiquement sur son image, ce qui est la manière dont nous avons considéré $W_a$
comme sous-schéma fermé de $X$ à la section \ref{intersection-section-rational-equivalence}.
Les assertions (2) et (3) sont des conséquences formelles de (1) et des définitions.
\end{proof}

\noindent
Pour des intersections transverses de sous-schémas fermés, la multiplicité
d'intersection vaut $1$.

\begin{lemma}
\label{intersection-lemma-transversal-subschemes}
Soit $X$ une variété non singulière. Soient $r, s \geq 0$ et
$Y, Z \subset X$ des sous-schémas fermés tels que $\dim(Y) \leq r$ et
$\dim(Z) \leq s$. Supposons que $[Y]_r = \sum n_i[Y_i]$ et
$[Z]_s = \sum m_j[Z_j]$ se coupent proprement.
Soit $T$ une composante irréductible de $Y_{i_0} \cap Z_{j_0}$
pour certains $i_0$ et $j_0$, et supposons que la multiplicité
(au sens de la section \ref{intersection-section-cycle-of-closed}) de $T$
dans le sous-schéma fermé $Y \cap Z$ soit $1$.
Alors
\begin{enumerate}
\item le coefficient de $T$ dans $[Y]_r \cdot [Z]_s$ est $1$,
\item $Y$ et $Z$ sont non singuliers au point générique de $T$,
\item $n_{i_0} = 1$, $m_{j_0} = 1$, et
\item $T$ n'est contenu dans aucun $Y_i$ ni aucun $Z_j$ pour $i \not = i_0$ et
$j \not = j_0$.
\end{enumerate}
\end{lemma}

\begin{proof}
Posons $n = \dim(X)$, $a = n - r$, $b = n - s$. Remarquons que
$\dim(T) = r + s - n = n - a - b$ puisque, par hypothèse, les
intersections sont propres. Posons $(A, \mathfrak m, \kappa) =
(\mathcal{O}_{X, \xi}, \mathfrak m_\xi, \kappa(\xi))$, où $\xi \in T$
est le point générique. Alors $\dim(A) = a + b$, voir
Variétés, lemme \ref{varieties-lemma-dimension-locally-algebraic}.
Soient $I_0, I, J_0, J \subset A$ les idéaux qui définissent les traces de
$Y_{i_0}$, $Y$, $Z_{j_0}$, $Z$ dans $\Spec(A)$.
Alors $\dim(A/I) = \dim(A/I_0) = b$ et $\dim(A/J) = \dim(A/J_0) = a$
d'après la même référence. Posons $\overline{I} = I + \mathfrak m^2/\mathfrak m^2$.
Alors $I \subset I_0 \subset \mathfrak m$ et
$J \subset J_0 \subset \mathfrak m$, et $I + J = \mathfrak m$.
Le lemme \ref{intersection-lemma-transversal} et sa démonstration montrent que
$I_0 = (f_1, \ldots, f_a)$ et $J_0 = (g_1, \ldots, g_b)$
où $f_1, \ldots, g_b$ est un système régulier de paramètres
de l'anneau local régulier $A$. Comme $I + J = \mathfrak m$, l'application
$$
I \oplus J \to
\mathfrak m/\mathfrak m^2 =
\kappa f_1
\oplus \ldots \oplus
\kappa f_a
\oplus
\kappa g_1
\oplus \ldots \oplus
\kappa g_b
$$
est surjective. Nous en déduisons que l'on peut trouver
$f_1', \ldots, f_a' \in I$ et $g'_1, \ldots, g_b' \in J$
dont les classes résiduelles dans $\mathfrak m/\mathfrak m^2$ sont égales aux
classes résiduelles de $f_1, \ldots, f_a$ et de $g_1, \ldots, g_b$.
Alors $f'_1, \ldots, g'_b$ est un système régulier de paramètres de $A$.
Le lemme d'algèbre \ref{algebra-lemma-regular-ring-CM} montre que
$A/(f'_1, \ldots, f'_a)$ est un anneau local régulier de dimension $b$.
Ainsi, tout quotient non trivial de $A/(f'_1, \ldots, f'_a)$
est de dimension strictement inférieure
(Algèbre, lemmes \ref{algebra-lemma-regular-domain} et
\ref{algebra-lemma-one-equation}). Par conséquent, $I = (f'_1, \ldots, f'_a) = I_0$.
Par symétrie, $J = J_0$. Cela démontre (2), (3) et (4).
Enfin, le coefficient de $T$ dans $[Y]_r \cdot [Z]_s$
est le coefficient de $T$ dans $Y_{i_0} \cdot Z_{j_0}$, qui vaut
$1$ d'après le lemme \ref{intersection-lemma-transversal}.
\end{proof}



\section{Produit extérieur}
\label{intersection-section-exterior-product}

\noindent
Soient $X$ et $Y$ des variétés.
Soit $V$, resp.\ $W$, une sous-variété fermée de $X$, resp.\ de $Y$.
Le produit $V\times W$ est une sous-variété fermée de $X\times Y$
(lemme \ref{intersection-lemma-dimension-product-varieties}).
Pour un $k$-cycle $\alpha = \sum n_i [V_i]$ et un $l$-cycle
$\beta = \sum m_j [V_j]$ sur $Y$, nous définissons le
{\it produit extérieur} de $\alpha$ et $\beta$ comme le cycle
$\alpha \times \beta = \sum n_i m_j [V_i \times W_j]$.
Le produit extérieur définit une application $\mathbf{Z}$-linéaire
$$
Z_r(X) \otimes_\mathbf{Z} Z_s(Y) \longrightarrow Z_{r + s}(X \times Y)
$$
Montrons que le produit extérieur passe au quotient par l'équivalence rationnelle.

\begin{lemma}
\label{intersection-lemma-exterior-product-rational-equivalence}
Soient $X$ et $Y$ des variétés.
Soient $\alpha \in Z_r(X)$ et $\beta \in Z_s(Y)$.
Si $\alpha \sim_{rat} 0$ ou $\beta \sim_{rat} 0$, alors
$\alpha \times \beta \sim_{rat} 0$.
\end{lemma}

\begin{proof}
Par linéarité et symétrie en $X$ et $Y$, il suffit de le démontrer lorsque
$\alpha = [V]$ pour une sous-variété $V \subset X$ de dimension $r$, et
$\beta = [W_a]_s - [W_b]_s$ pour une sous-variété fermée
$W \subset Y \times \mathbf{P}^1$ de dimension $s + 1$ qui
coupe $Y \times a$ et $Y \times b$ proprement. Dans ce cas,
le lemme résulte de l'égalité
$$
[(V \times W)_a]_{r + s} = [V] \times [W_a]_s
$$
et de même en remplaçant $a$ par $b$. En effet, nous voyons alors que
$\alpha \times \beta = [(V \times W)_a]_{r + s} - [(V \times W)_b]_{r + s}$
comme voulu. Pour établir l'égalité affichée, remarquons l'égalité
$$
V \times W_a = (V \times W)_a
$$
de schémas. La projection $V \times W_a \to W_a$ induit une bijection
entre les composantes irréductibles (voir par exemple
Variétés, lemme \ref{varieties-lemma-bijection-irreducible-components}).
Soit $W' \subset W_a$ une composante irréductible de point générique $\zeta$.
Alors $V \times W'$ est la composante irréductible correspondante de
$V \times W_a$ (voir le lemme \ref{intersection-lemma-dimension-product-varieties}).
Soit $\xi$ le point générique de $V \times W'$. Nous devons montrer que
$$
\text{longueur}_{\mathcal{O}_{Y, \zeta}}(\mathcal{O}_{W_a, \zeta}) =
\text{longueur}_{\mathcal{O}_{X \times Y, \xi}}(
\mathcal{O}_{V \times W_a, \xi})
$$
Dans cette formule, nous pouvons remplacer
$\mathcal{O}_{Y, \zeta}$ par $\mathcal{O}_{W_a, \zeta}$ et
nous pouvons remplacer
$\mathcal{O}_{X \times Y, \zeta}$ par $\mathcal{O}_{V \times W_a, \zeta}$
(voir Algèbre, lemme \ref{algebra-lemma-length-independent}).
Puisque $\mathcal{O}_{W_a, \zeta} \to \mathcal{O}_{V \times W_a, \xi}$ est plat,
le lemme d'algèbre \ref{algebra-lemma-pullback-module} montre qu'il suffit
de vérifier que
$$
\text{longueur}_{\mathcal{O}_{V \times W_a, \xi}}(
\mathcal{O}_{V \times W_a, \xi}/
\mathfrak m_\zeta\mathcal{O}_{V \times W_a, \xi}) = 1
$$
C'est vrai, car le quotient de droite est l'anneau local
$\mathcal{O}_{V \times W', \xi}$ d'une variété en un point générique,
et est donc égal à $\kappa(\xi)$.
\end{proof}

\noindent
Nous en concluons que le produit extérieur définit un diagramme commutatif
$$
\xymatrix{
Z_r(X) \otimes_\mathbf{Z} Z_s(Y) \ar[r] \ar[d] &
Z_{r + s}(X \times Y) \ar[d] \\
\CH_r(X) \otimes_\mathbf{Z} \CH_s(Y) \ar[r] &
\CH_{r + s}(X \times Y)
}
$$
pour toute paire de variétés $X$ et $Y$. Pour les variétés non singulières,
nous pouvons considérer le produit extérieur comme un produit d'intersection
d'images inverses.

\begin{lemma}
\label{intersection-lemma-exterior-product}
Soient $X$ et $Y$ des variétés non singulières.
Soient $\alpha \in Z_r(X)$ et $\beta \in Z_s(Y)$.
Alors
\begin{enumerate}
\item $\text{pr}_Y^*(\beta) = [X] \times \beta$ et
$\text{pr}_X^*(\alpha) = \alpha \times [Y]$,
\item $\alpha \times [Y]$ et $[X]\times \beta$
se coupent proprement sur $X\times Y$, et
\item on a
$\alpha \times \beta =
(\alpha \times [Y])\cdot ([X]\times\beta) =
pr_Y^*(\alpha) \cdot pr_X^*(\beta)$
dans $Z_{r + s}(X \times Y)$.
\end{enumerate}
\end{lemma}

\begin{proof}
Par linéarité, nous pouvons supposer que $\alpha = [V]$ et $\beta = [W]$.
Alors (1) affirme que $\text{pr}_Y^{-1}(W) = X \times W$ et
$\text{pr}_X^{-1}(V) = V \times Y$. Cela est clair.
L'assertion (2) est vraie parce que $X \times W \cap V \times Y = V \times W$ et
$\dim(V \times W) = r + s$ d'après le lemme \ref{intersection-lemma-dimension-product-varieties}.

\medskip\noindent
Démonstration de (3).
Soit $\xi$ le point générique de $V \times W$.
Puisque la projection $X \times W \to W$ est lisse comme changement de base du morphisme
$X \to \Spec(\mathbf{C})$, nous voyons que $X \times W$ est non singulier
en tout point situé au-dessus du point générique de $W$, en particulier en $\xi$.
Il en va de même pour $V \times Y$. Ainsi, $\mathcal{O}_{X \times W, \xi}$
et $\mathcal{O}_{V \times Y, \xi}$ sont des anneaux locaux de Cohen-Macaulay,
et le lemme \ref{intersection-lemma-intersection-multiplicity-CM} s'applique.
Comme $V \times Y \cap X \times W = V \times W$ schématiquement,
cela achève la démonstration.
\end{proof}



\section{Réduction à la diagonale}
\label{intersection-section-reduction-diagonal}

\noindent
Soit $X$ une variété non singulière. Nous utiliserons $\Delta$
pour désigner soit le morphisme diagonal $\Delta : X \to X \times X$,
soit l'image $\Delta \subset X \times X$.
La réduction à la diagonale est l'énoncé selon lequel
les produits d'intersection sur $X$ se ramènent aux produits d'intersection
des produits extérieurs avec la diagonale dans $X \times X$.

\begin{lemma}
\label{intersection-lemma-tor-and-diagonal}
Soit $X$ une variété non singulière.
\begin{enumerate}
\item Si $\mathcal{F}$ et $\mathcal{G}$ sont des $\mathcal{O}_X$-modules cohérents,
alors il existe des isomorphismes canoniques
$$
\text{Tor}_i^{\mathcal{O}_{X \times X}}(\mathcal{O}_\Delta,
\text{pr}_1^*\mathcal{F} \otimes_{\mathcal{O}_{X \times X}}
\text{pr}_2^*\mathcal{G})
=
\Delta_*\text{Tor}_i^{\mathcal{O}_X}(\mathcal{F}, \mathcal{G})
$$
\item Si $K$ et $M$ appartiennent à $D_\QCoh(\mathcal{O}_X)$, alors
il existe un isomorphisme canonique
$$
L\Delta^* \left(
L\text{pr}_1^*K \otimes_{\mathcal{O}_{X \times X}}^\mathbf{L} L\text{pr}_2^*M
\right)
= K \otimes_{\mathcal{O}_X}^\mathbf{L} M
$$
dans $D_\QCoh(\mathcal{O}_X)$ et un isomorphisme canonique
$$
\mathcal{O}_\Delta \otimes_{\mathcal{O}_{X \times X}}^\mathbf{L}
L\text{pr}_1^*K \otimes_{\mathcal{O}_{X \times X}}^\mathbf{L} L\text{pr}_2^*M
= \Delta_*(K \otimes_{\mathcal{O}_X}^\mathbf{L} M)
$$
dans $D_\QCoh(X \times X)$.
\end{enumerate}
\end{lemma}

\begin{proof}
Expliquons comment démontrer (1) de manière plus élémentaire et (2) à l'aide
d'une théorie plus générale. Comme (2) implique (1), le lecteur peut omettre la démonstration
de (1).

\medskip\noindent
Démonstration de (1). Choisissons un ouvert affine $\Spec(A) \subset X$.
Alors $A$ est une $\mathbf{C}$-algèbre noethérienne et
$\mathcal{F}$, $\mathcal{G}$ correspondent aux $A$-modules de type fini
$M$ et $N$ (Cohomologie des schémas, lemme
\ref{coherent-lemma-coherent-Noetherian}).
D'après Catégories dérivées de schémas, lemme
\ref{perfect-lemma-quasi-coherence-tensor-product}, nous pouvons
calculer $\text{Tor}_i$ sur $\mathcal{O}_X$
en calculant d'abord les $\text{Tor}$
de $M$ et $N$ sur $A$, puis en prenant le $\mathcal{O}_X$-module associé.
Pour les $\text{Tor}_i$ sur $\mathcal{O}_{X \times X}$, nous calculons
les Tor de $A$ et $M \otimes_\mathbf{C} N$ sur $A \otimes_\mathbf{C} A$
puis prenons le $\mathcal{O}_{X \times X}$-module associé.
Ainsi, sur cet ouvert affine, nous devons démontrer que
$$
\text{Tor}_i^{A \otimes_\mathbf{C} A}(A, M \otimes_\mathbf{C} N) =
\text{Tor}_i^A(M, N)
$$
Pour le voir, choisissons des résolutions $F_\bullet \to M$ et $G_\bullet \to M$
par des $A$-modules libres de type fini
(Algèbre, lemme \ref{algebra-lemma-resolution-by-finite-free}).
Remarquons que $\text{Tot}(F_\bullet \otimes_\mathbf{C} G_\bullet)$
est une résolution de $M \otimes_\mathbf{C} N$ puisqu'il calcule les
groupes de Tor sur $\mathbf{C}$ ! Bien entendu, les termes de
$F_\bullet \otimes_\mathbf{C} G_\bullet$ sont des
$A \otimes_\mathbf{C} A$-modules libres de type fini. Ainsi, le membre de gauche
de l'égalité affichée est le module
$$
H_i(A \otimes_{A \otimes_\mathbf{C} A}
\text{Tot}(F_\bullet \otimes_\mathbf{C} G_\bullet))
$$
et le membre de droite est le module
$$
H_i(\text{Tot}(F_\bullet \otimes_A G_\bullet))
$$
Comme $A \otimes_{A \otimes_\mathbf{C} A} (F_p \otimes_\mathbf{C} G_q)
= F_p \otimes_A G_q$, nous voyons que ces modules sont égaux.
Cela définit un isomorphisme sur l'ouvert affine $\Spec(A) \times \Spec(A)$
(ce qui suffit pour l'application à l'égalité des nombres d'intersection).
Nous omettons la démonstration du recollement de ces isomorphismes.

\medskip\noindent
Démonstration de (2). La seconde assertion résulte de la première par la
formule de projection telle qu'énoncée dans
Catégories dérivées de schémas, lemme \ref{perfect-lemma-cohomology-base-change}.
Pour établir la première, représentons $K$ et $M$ par des complexes K-plats
$\mathcal{K}^\bullet$ et $\mathcal{M}^\bullet$.
Comme l'image inverse et le produit tensoriel préservent les complexes K-plats
(Cohomologie, lemmes \ref{cohomology-lemma-tensor-product-K-flat} et
\ref{cohomology-lemma-pullback-K-flat})
nous voyons qu'il suffit de montrer
$$
\Delta^*\text{Tot}(
\text{pr}_1^*\mathcal{K}^\bullet
\otimes_{\mathcal{O}_{X \times X}} \text{pr}_2^*\mathcal{M}^\bullet)
=
\text{Tot}(
\mathcal{K}^\bullet \otimes_{\mathcal{O}_X} \mathcal{M}^\bullet)
$$
Il suffit donc de voir qu'il existe des isomorphismes canoniques
$$
\Delta^*(\text{pr}_1^*\mathcal{K}
\otimes_{\mathcal{O}_{X \times X}} \text{pr}_2^*\mathcal{M})
\longrightarrow
\mathcal{K} \otimes_{\mathcal{O}_X} \mathcal{M}
$$
lorsque $\mathcal{K}$ et $\mathcal{M}$ sont des $\mathcal{O}_X$-modules
(pas nécessairement quasi-cohérents ni plats).
Nous omettons les détails.
\end{proof}

\begin{lemma}
\label{intersection-lemma-reduction-diagonal}
Soit $X$ une variété non singulière. Soient $\alpha$ et $\beta$,
respectivement un $r$-cycle et un $s$-cycle sur $X$. Supposons que $\alpha$ et $\beta$
se coupent proprement. Alors
\begin{enumerate}
\item $\alpha \times \beta$ et $[\Delta]$ se coupent proprement,
\item on a $\Delta_*(\alpha \cdot \beta) = [\Delta] \cdot \alpha\times\beta$
comme cycles sur $X \times X$,
\item si $X$ est propre, alors
$\text{pr}_{1, *}([\Delta] \cdot \alpha\times\beta) = \alpha\cdot\beta$,
où $pr_1 : X\times X \to X$ est la projection.
\end{enumerate}
\end{lemma}

\begin{proof}
Par linéarité, il suffit de le démontrer lorsque $\alpha = [V]$ et $\beta = [W]$
pour des sous-variétés fermées $V \subset X$ et $W \subset Y$ qui se coupent
proprement. Rappelons que $V \times W$ est une sous-variété fermée de dimension $r + s$.
Remarquons que, schématiquement, on a
$V \cap W = \Delta^{-1}(V \times W)$ ainsi que
$\Delta(V \cap W) = \Delta \cap V \times W$.
Cela démontre (1).

\medskip\noindent
Démonstration de (2). Soit $Z \subset V \cap W$ une composante irréductible
de point générique $\xi$. Nous devons montrer que le coefficient de
$Z$ dans $\alpha \cdot \beta$ est égal au coefficient de
$\Delta(Z)$ dans $[\Delta] \cdot \alpha \times \beta$. Le premier est donné
par l'entier
$$
\sum (-1)^i
\text{longueur}_{\mathcal{O}_{X, \xi}}
\text{Tor}_i^{\mathcal{O}_X}(\mathcal{O}_V, \mathcal{O}_W)_\xi
$$
et le second par l'entier
$$
\sum (-1)^i
\text{longueur}_{\mathcal{O}_{X \times Y, \Delta(\xi)}}
\text{Tor}_i^{\mathcal{O}_{X \times Y}}(
\mathcal{O}_\Delta, \mathcal{O}_{V \times W})_{\Delta(\xi)}
$$
Cependant, d'après le lemme \ref{intersection-lemma-tor-and-diagonal}, on a
$$
\text{Tor}_i^{\mathcal{O}_X}(\mathcal{O}_V, \mathcal{O}_W)_\xi \cong
\text{Tor}_i^{\mathcal{O}_{X \times Y}}(
\mathcal{O}_\Delta, \mathcal{O}_{V \times W})_{\Delta(\xi)}
$$
comme $\mathcal{O}_{X \times X, \Delta(\xi)}$-modules. L'égalité
des longueurs résulte alors, plus précisément, d'Algèbre,
lemme \ref{algebra-lemma-length-independent}.

\medskip\noindent
L'assertion (2) implique (3), car
$\text{pr}_{1, *} \circ \Delta_* = \text{id}$ d'après
le lemme \ref{intersection-lemma-compose-pushforward}.
\end{proof}

\begin{proposition}
\label{intersection-proposition-positivity}
\begin{reference}
C'est l'un des résultats principaux de \cite{Serre_algebre_locale}.
\end{reference}
Soit $X$ une variété non singulière. Soient $V \subset X$ et
$W \subset Y$ des sous-variétés fermées qui se coupent proprement.
Soit $Z \subset V \cap W$ une composante irréductible.
Alors $e(X, V \cdot W, Z) > 0$.
\end{proposition}

\begin{proof}
D'après le lemme \ref{intersection-lemma-reduction-diagonal}, on a
$$
e(X, V \cdot W, Z) = e(X \times X, \Delta \cdot V \times W, \Delta(Z))
$$
Puisque $\Delta : X \to X \times X$ est une immersion régulière
(voir le lemme \ref{intersection-lemma-diagonal-regular-immersion}), nous voyons que
$e(X \times X, \Delta \cdot V \times W, \Delta(Z))$ est un entier strictement
positif d'après le lemme \ref{intersection-lemma-multiplicity-with-lci}.
\end{proof}

\noindent
Le lemme suivant est un lemme clé du développement de la théorie
présenté dans ce chapitre. Essentiellement, ce lemme affirme que
les nombres d'intersection vérifient une propriété d'additivité convenable.

\begin{lemma}
\label{intersection-lemma-tor-sheaf}
\begin{reference}
\cite[chapitre V]{Serre_algebre_locale}
\end{reference}
Soit $X$ une variété non singulière. Soient $\mathcal{F}$ et
$\mathcal{G}$ des faisceaux cohérents sur $X$ tels que
$\dim(\text{Supp}(\mathcal{F})) \leq r$,
$\dim(\text{Supp}(\mathcal{G})) \leq s$, et
$\dim(\text{Supp}(\mathcal{F}) \cap \text{Supp}(\mathcal{G}) )
\leq r + s - \dim X$. Dans ce cas, $[\mathcal{F}]_r$ et $[\mathcal{G}]_s$
se coupent proprement et
$$
[\mathcal{F}]_r \cdot [\mathcal{G}]_s =
\sum (-1)^p
[\text{Tor}_p^{\mathcal{O}_X}(\mathcal{F}, \mathcal{G})]_{r + s - \dim(X)}.
$$
\end{lemma}

\begin{proof}
L'assertion selon laquelle $[\mathcal{F}]_r$ et $[\mathcal{G}]_s$ se coupent proprement
est immédiate. Comme nous démontrons une égalité de cycles, nous pouvons travailler
localement sur $X$. (Remarquons que la formation du produit d'intersection
des cycles, celle des faisceaux de $\text{Tor}$ et
celle du cycle associé à un faisceau cohérent commutent chacune à la
restriction aux sous-schémas ouverts.) Nous pouvons donc supposer, et supposons, $X$ affine.

\medskip\noindent
Posons
$$
RHS(\mathcal{F}, \mathcal{G}) = [\mathcal{F}]_r \cdot [\mathcal{G}]_s
\quad\text{et}\quad
LHS(\mathcal{F}, \mathcal{G}) = \sum (-1)^p
[\text{Tor}_p^{\mathcal{O}_X}(\mathcal{F}, \mathcal{G})]_{r + s - \dim(X)}
$$
Considérons une suite exacte
$$
0 \to \mathcal{F}_1 \to \mathcal{F}_2 \to \mathcal{F}_3 \to 0
$$
de faisceaux cohérents sur $X$ telle que
$\text{Supp}(\mathcal{F}_i) \subset \text{Supp}(\mathcal{F})$ ;
alors $LHS(\mathcal{F}_i, \mathcal{G})$ et
$RHS(\mathcal{F}_i, \mathcal{G})$ sont définis pour $i = 1, 2, 3$,
et l'on a
$$
RHS(\mathcal{F}_2, \mathcal{G}) =
RHS(\mathcal{F}_1, \mathcal{G}) + RHS(\mathcal{F}_3, \mathcal{G})
$$
et de même pour LHS. En effet, la condition sur les supports garantit que
tout est défini, et la suite exacte ainsi que l'additivité des longueurs
donnent
$$
[\mathcal{F}_2]_r = [\mathcal{F}_1]_r  + [\mathcal{F}_3]_r
$$
(Homologie de Chow, lemme \ref{chow-lemma-additivity-sheaf-cycle}),
ce qui entraîne l'additivité pour RHS. La suite exacte longue des $\text{Tor}$
$$
\ldots \to \text{Tor}_1(\mathcal{F}_3, \mathcal{G}) \to
\text{Tor}_0(\mathcal{F}_1, \mathcal{G}) \to
\text{Tor}_0(\mathcal{F}_2, \mathcal{G}) \to
\text{Tor}_0(\mathcal{F}_3, \mathcal{G}) \to 0
$$
et l'additivité des longueurs, comme précédemment, entraînent l'additivité pour LHS.

\medskip\noindent
{\emergencystretch=1em D'après Algèbre, lemme \ref{algebra-lemma-filter-Noetherian-module},
et puisque $X$ est affine, nous pouvons trouver une filtration de $\mathcal{F}$
dont les quotients gradués sont des faisceaux structuraux de sous-variétés fermées de
$\text{Supp}(\mathcal{F})$. L'additivité démontrée au paragraphe précédent
montre qu'il suffit de démontrer $LHS = RHS$ avec
$\mathcal{F}$ remplacé par $\mathcal{O}_V$, où
$V \subset \text{Supp}(\mathcal{F})$.
Par symétrie, nous pouvons faire de même pour $\mathcal{G}$.
Nous sommes ainsi ramenés à démontrer que\par}
$$
LHS(\mathcal{O}_V, \mathcal{O}_W) = RHS(\mathcal{O}_V, \mathcal{O}_W)
$$
où $W \subset \text{Supp}(\mathcal{G})$ est une sous-variété fermée.
Si $\dim(V) = r$ et $\dim(W) = s$, alors cette égalité est la
{\bf définition} de $V \cdot W$. En revanche, si
$\dim(V) < r$ ou $\dim(W) < s$, c'est-à-dire si $[V]_r = 0$ ou $[W]_s = 0$,
alors nous devons démontrer que $RHS(\mathcal{O}_V, \mathcal{O}_W) = 0$
\footnote{Le lecteur peut voir que ce n'est pas trivial en
prenant $r = s = 1$, $X$ une surface non singulière et $V = W$
un point fermé $x$ de $X$. Dans ce cas, il y a $3$
$\text{Tor}$ non nuls, de longueurs $1, 2, 1$ en $x$.}.

\medskip\noindent
Soit $Z \subset V \cap W$ une composante irréductible de dimension
$r + s - \dim(X)$. C'est la dimension maximale d'une composante,
et il suffit de montrer que le coefficient de $Z$ dans $RHS$ est nul.
Soit $\xi \in Z$ le point générique. Posons $A = \mathcal{O}_{X, \xi}$,
$B = \mathcal{O}_{X \times X, \Delta(\xi)}$, et
$C = \mathcal{O}_{V \times W, \Delta(\xi)}$.
D'après le lemme \ref{intersection-lemma-tor-and-diagonal}, on a
$$
\text{coefficient de }Z\text{ dans }
RHS(\mathcal{O}_V, \mathcal{O}_W) = 
\sum (-1)^i
\text{longueur}_B \text{Tor}_i^B(A, C)
$$
Comme $\dim(V) < r$ ou $\dim(W) < s$, on a $\dim(V \times W) < r + s$,
ce qui entraîne $\dim(C) < \dim(X)$ (on omet un petit détail). De plus, le
noyau $I$ de $B \to A$ est engendré par une suite régulière de
longueur $\dim(X)$ (lemme \ref{intersection-lemma-diagonal-regular-immersion}).
Il y a donc annulation d'après le lemme \ref{intersection-lemma-one-ideal-ci}, car
la fonction de Hilbert de $C$ relativement à $I$ est de degré $\dim(C) < n$
d'après Algèbre, proposition \ref{algebra-proposition-dimension}.
\end{proof}

\begin{remark}
\label{intersection-remark-Serre-conjectures}
Soit $(A, \mathfrak m, \kappa)$ un anneau local régulier.
Soient $M$ et $N$ des $A$-modules de type fini non nuls tels que $M \otimes_A N$
soit supporté par $\{\mathfrak m\}$. Alors
$$
\chi(M, N) = \sum (-1)^i \text{longueur}_A \text{Tor}_i^A(M, N)
$$
est fini. Posons $r = \dim(\text{Supp}(M))$ et $s = \dim(\text{Supp}(N))$.
Dans \cite{Serre_algebre_locale}, on démontre que $r + s \leq \dim(A)$
et les conjectures suivantes y sont formulées :
\begin{enumerate}
\item si $r + s < \dim(A)$, alors $\chi(M, N) = 0$, et
\item si $r + s = \dim(A)$, alors $\chi(M, N) > 0$.
\end{enumerate}
Les arguments qui démontrent le lemme \ref{intersection-lemma-tor-sheaf} et
la proposition \ref{intersection-proposition-positivity} peuvent être exploités
(comme dans le texte de Serre) pour montrer que (1) et (2) sont
vraies si $A$ contient un corps. À l'heure actuelle, la conjecture (1) est démontrée
en général, et l'on sait que $\chi(M, N) \geq 0$ en général (Gabber).
À notre connaissance, la positivité reste un problème ouvert.
\end{remark}



\section{Associativité des intersections}
\label{intersection-section-associative}

\noindent
Il est clair que les intersections propres définies ci-dessus sont commutatives.
À l'aide du lemme clé \ref{intersection-lemma-tor-sheaf}, nous pouvons démontrer que les produits
d'intersection (propres) sont associatifs.

\begin{lemma}
\label{intersection-lemma-associative}
Soit $X$ une variété non singulière. Soient $U, V, W$ des sous-variétés
fermées. Supposons que $U, V, W$ se coupent proprement deux à deux
et que $\dim(U \cap V \cap W) \leq \dim(U) + \dim(V) + \dim(W) - 2\dim(X)$.
Alors
$$
U \cdot (V \cdot W) = (U \cdot V) \cdot W
$$
comme cycles sur $X$.
\end{lemma}

\begin{proof}
Nous allons utiliser le lemme \ref{intersection-lemma-tor-sheaf} sans plus le mentionner.
Cela implique que
\begin{align*}
V \cdot W
& =
\sum (-1)^i [\text{Tor}_i(\mathcal{O}_V, \mathcal{O}_W)]_{b + c - n} \\
U \cdot (V \cdot W)
& =
\sum (-1)^{i + j}
[
\text{Tor}_j(\mathcal{O}_U, \text{Tor}_i(\mathcal{O}_V, \mathcal{O}_W))
]_{a + b + c - 2n} \\
U \cdot V
& =
\sum (-1)^i [\text{Tor}_i(\mathcal{O}_U, \mathcal{O}_V)]_{a + b - n} \\
(U \cdot V) \cdot W
& =
\sum (-1)^{i + j}
[
\text{Tor}_j(\text{Tor}_i(\mathcal{O}_U, \mathcal{O}_V), \mathcal{O}_W))
]_{a + b + c - 2n}
\end{align*}
où $\dim(U) = a$, $\dim(V) = b$, $\dim(W) = c$, $\dim(X) = n$.
Les hypothèses du lemme garantissent que les faisceaux cohérents
des formules ci-dessus vérifient les majorations requises pour les dimensions
des supports, de sorte que ces formules ont un sens. Considérons maintenant l'objet
$$
K =
\mathcal{O}_U \otimes^\mathbf{L}_{\mathcal{O}_X} \mathcal{O}_V
\otimes^\mathbf{L}_{\mathcal{O}_X} \mathcal{O}_W
$$
de la catégorie dérivée $D_{\textit{Coh}}(\mathcal{O}_X)$.
Nous affirmons que les expressions obtenues ci-dessus pour
$U \cdot (V \cdot W)$ et $(U \cdot V) \cdot W$
sont égales à
$$
\sum (-1)^k [H^k(K)]_{a + b + c - 2n}
$$
Cela démontrera le lemme. Par symétrie, il suffit de démontrer l'une
de ces égalités. Pour cela, représentons $\mathcal{O}_U$
et $\mathcal{O}_V \otimes_{\mathcal{O}_X}^\mathbf{L} \mathcal{O}_W$
par des complexes K-plats $M^\bullet$ et $L^\bullet$ et utilisons la
suite spectrale associée au complexe double
$M^\bullet \otimes_{\mathcal{O}_X} L^\bullet$ dans
Homologie, section \ref{homology-section-double-complex}.
Il s'agit d'une suite spectrale dont la page $E_2$ est
$$
E_2^{p, q} =
\text{Tor}_{-p}(\mathcal{O}_U, \text{Tor}_{-q}(\mathcal{O}_V, \mathcal{O}_W))
$$
et qui converge vers $H^{p + q}(K)$ (détails omis ; comparer à
Compléments d'algèbre, exemple \ref{more-algebra-example-tor}).
Puisque les longueurs sont additives dans les suites
exactes courtes, le résultat en découle.
\end{proof}




\section{Image inverse plate et produits d'intersection}
\label{intersection-section-flat-pullback-and-intersection-products}

\noindent
Brève discussion de l'interaction entre les intersections et
l'image inverse plate.

\begin{lemma}
\label{intersection-lemma-flat-pull-back-and-intersections-sheaves}
Soit $f : X \to Y$ un morphisme plat de variétés non singulières. Posons
$e = \dim(X) - \dim(Y)$. Soient $\mathcal{F}$ et $\mathcal{G}$ des faisceaux
cohérents sur $Y$ tels que $\dim(\text{Supp}(\mathcal{F})) \leq r$,
$\dim(\text{Supp}(\mathcal{G})) \leq s$ et
$\dim(\text{Supp}(\mathcal{F}) \cap \text{Supp}(\mathcal{G}) )
\leq r + s - \dim(Y)$. Alors les cycles
$[f^*\mathcal{F}]_{r + e}$ et $[f^*\mathcal{G}]_{s + e}$
se coupent proprement et
$$
f^*([\mathcal{F}]_r \cdot [\mathcal{G}]_s) =
[f^*\mathcal{F}]_{r + e} \cdot [f^*\mathcal{G}]_{s + e}
$$
\end{lemma}

\begin{proof}
Le fait que $[f^*\mathcal{F}]_{r + e}$ et $[f^*\mathcal{G}]_{s + e}$
se coupent proprement résulte immédiatement de l'hypothèse que $f$ est de
dimension relative $e$. D'après les
lemmes \ref{intersection-lemma-tor-sheaf} et \ref{intersection-lemma-pullback},
il suffit de montrer que
$$
f^*\text{Tor}_i^{\mathcal{O}_Y}(\mathcal{F}, \mathcal{G}) =
\text{Tor}_i^{\mathcal{O}_X}(f^*\mathcal{F}, f^*\mathcal{G})
$$
comme $\mathcal{O}_X$-modules. Cela résulte de
Cohomologie, lemme \ref{cohomology-lemma-pullback-tensor-product},
et du fait que $f^*$ est exact, donc $Lf^*\mathcal{F} = f^*\mathcal{F}$,
et de même pour $\mathcal{G}$.
\end{proof}

\begin{lemma}
\label{intersection-lemma-flat-pullback-and-intersections}
Soit $f : X \to Y$ un morphisme plat de variétés non singulières.
Soit $\alpha$ un $r$-cycle sur $Y$ et $\beta$ un $s$-cycle sur $Y$.
Supposons que $\alpha$ et $\beta$ se coupent proprement. Alors $f^*\alpha$
et $f^*\beta$ se coupent proprement et
$f^*( \alpha \cdot \beta ) = f^*\alpha \cdot f^*\beta$.
\end{lemma}

\begin{proof}
Par linéarité, nous pouvons supposer que $\alpha = [V]$ et $\beta = [W]$
pour des sous-variétés fermées $V, W \subset Y$ de dimensions $r, s$.
Disons que $f$ est de dimension relative $e$. Le lemme est alors un cas particulier du
lemme \ref{intersection-lemma-flat-pull-back-and-intersections-sheaves},
car $[V] = [\mathcal{O}_V]_r$, $[W] = [\mathcal{O}_W]_r$,
$f^*[V] = [f^{-1}(V)]_{r + e} = [f^*\mathcal{O}_V]_{r + e}$ et
$f^*[W] = [f^{-1}(W)]_{s + e} = [f^*\mathcal{O}_W]_{s + e}$.
\end{proof}


\section{Formule de projection pour les morphismes plats propres}
\label{intersection-section-projection-formula-flat}

\noindent
Brève discussion de la formule de projection pour les morphismes plats propres.

\begin{lemma}
\label{intersection-lemma-projection-formula-flat}
\begin{reference}
Voir \cite[chapitre V, C), section 7, formule (10)]{Serre_algebre_locale}
pour une formule plus générale.
\end{reference}
Soit $f : X \to Y$ un morphisme plat propre de variétés non singulières.
Posons $e = \dim(X) - \dim(Y)$. Soit $\alpha$ un $r$-cycle sur $X$ et soit
$\beta$ un $s$-cycle sur $Y$. Supposons que $\alpha$ et $f^*(\beta)$ se coupent
proprement. Alors $f_*(\alpha)$ et $\beta$ se coupent proprement et
$$
f_*(\alpha) \cdot \beta = f_*( \alpha \cdot f^*\beta)
$$
\end{lemma}

\begin{proof}
Par linéarité, on se ramène au cas où $\alpha = [V]$ et
$\beta = [W]$ pour des sous-variétés fermées $V \subset X$ et
$W \subset Y$ de dimensions $r$ et $s$. Alors $f^{-1}(W)$ est de
dimension pure $s + e$. Nous supposons que les cycles
$[V]$ et $f^*[W]$ se coupent proprement. Nous utiliserons sans
autre mention le fait que $V \cap f^{-1}(W) \to f(V) \cap W$
est surjectif.

\medskip\noindent
Soit $a$ la dimension de la fibre générique de $V \to f(V)$.
Si $a > 0$, alors $f_*[V] = 0$. En particulier, $f_*\alpha$ et $\beta$
se coupent proprement. Pour achever ce cas, il faut montrer que
$f_*([V] \cdot f^*[W]) = 0$. Cependant, puisque toute fibre de
$V \to f(V)$ est de dimension $\geq a$ (voir
Morphismes, lemme \ref{morphisms-lemma-openness-bounded-dimension-fibres}),
nous concluons que toute composante irréductible $Z$ de $V \cap f^{-1}(W)$
a des fibres de dimension $\geq a$ au-dessus de $f(Z)$. Cela donne
l'assertion voulue.

\medskip\noindent
Supposons que $V \to f(V)$ soit génériquement fini. Soit $Z \subset f(V) \cap W$
une composante irréductible. Soient $Z_i \subset V \cap f^{-1}(W)$,
$i = 1, \ldots, t$, les composantes irréductibles de $V \cap f^{-1}(W)$
qui dominent $Z$. Par hypothèse, chaque $Z_i$ est de dimension
$r + s + e - \dim(X) = r + s - \dim(Y)$. Par conséquent,
$\dim(Z) \leq r + s - \dim(Y)$. On voit donc que $f(V)$ et $W$
se coupent proprement, que $\dim(Z) = r + s - \dim(Y)$ et que chaque
$Z_i \to Z$ est génériquement fini. Il s'ensuit notamment que
$V \to f(V)$ a une fibre finie au point générique $\xi$ de $Z$.
Ainsi, $V \to Y$ est fini au-dessus d'un voisinage ouvert de $\xi$, voir
Cohomologie des schémas, lemme
\ref{coherent-lemma-proper-finite-fibre-finite-in-neighbourhood}.
En utilisant une formule de projection très générale pour les produits tensoriels dérivés, on obtient
$$
Rf_*(\mathcal{O}_V \otimes_{\mathcal{O}_X}^\mathbf{L} Lf^*\mathcal{O}_W) =
Rf_*\mathcal{O}_V \otimes_{\mathcal{O}_Y}^\mathbf{L} \mathcal{O}_W
$$
{\emergencystretch=1em voir Catégories dérivées de schémas, lemme
\ref{perfect-lemma-cohomology-base-change}.
Comme $f$ est plat, on a $Lf^*\mathcal{O}_W = f^*\mathcal{O}_W$.
Puisque $f|_V$ est fini au-dessus d'un voisinage ouvert de $\xi$, on a\par}
$$
(Rf_*\mathcal{F})_\xi = (f_*\mathcal{F})_\xi
$$
pour tout faisceau cohérent sur $X$ dont le support est contenu dans $V$
(voir Cohomologie des schémas, lemme
\ref{coherent-lemma-higher-direct-images-zero-finite-fibre}). Nous
en concluons que
\begin{equation}
\label{intersection-equation-stalks}
\left(
f_*\text{Tor}_i^{\mathcal{O}_X}(\mathcal{O}_V, f^*\mathcal{O}_W)
\right)_\xi =
\left(\text{Tor}_i^{\mathcal{O}_Y}(f_*\mathcal{O}_V, \mathcal{O}_W)\right)_\xi
\end{equation}
pour tout $i$. Puisque $f^*[W] = [f^*\mathcal{O}_W]_{s + e}$ d'après le
lemme \ref{intersection-lemma-pullback}, on a
$$
[V] \cdot f^*[W] =
\sum (-1)^i
[\text{Tor}_i^{\mathcal{O}_X}(\mathcal{O}_V,
f^*\mathcal{O}_W)]_{r + s - \dim(Y)}
$$
d'après le lemme \ref{intersection-lemma-tor-sheaf}. En appliquant le
lemme \ref{intersection-lemma-push-coherent},
on trouve
$$
f_*([V] \cdot f^*[W]) =
\sum (-1)^i
[f_*\text{Tor}_i^{\mathcal{O}_X}(\mathcal{O}_V,
f^*\mathcal{O}_W)]_{r + s - \dim(Y)}
$$
Comme $f_*[V] = [f_*\mathcal{O}_V]_r$ d'après le
lemme \ref{intersection-lemma-push-coherent}, on a
$$
[f_*V] \cdot [W] =
\sum (-1)^i
[\text{Tor}_i^{\mathcal{O}_X}(f_*\mathcal{O}_V,
\mathcal{O}_W)]_{r + s - \dim(Y)}
$$
à nouveau d'après le lemme \ref{intersection-lemma-tor-sheaf}.
En comparant la formule pour $f_*([V] \cdot f^*[W])$ à
la formule pour $f_*[V] \cdot [W]$ et en comparant, pour le
coefficient de $Z$, les longueurs des germes en $\xi$, on voit que
(\ref{intersection-equation-stalks}) achève la démonstration.
\end{proof}

\begin{lemma}
\label{intersection-lemma-transfer}
Soit $X \to P$ une immersion fermée de variétés non singulières.
Soit $C' \subset P \times \mathbf{P}^1$ une sous-variété fermée de dimension
$r + 1$. Supposons que
\begin{enumerate}
\item la fibre $C = C'_0$ soit de dimension $r$, c'est-à-dire que $C' \to \mathbf{P}^1$
soit dominant,
\item $C'$ coupe $X \times \mathbf{P}^1$ proprement,
\item $[C]_r$ coupe $X$ proprement.
\end{enumerate}
Alors, en posant $\alpha = [C]_r \cdot X$, considéré comme cycle sur $X$, et
$\beta = C' \cdot X \times \mathbf{P}^1$, considéré comme cycle sur
$X \times \mathbf{P}^1$, on a
$$
\alpha = \text{pr}_{X, *}(\beta \cdot X \times 0)
$$
comme cycles sur $X$, où $\text{pr}_X : X \times \mathbf{P}^1 \to X$ est la
projection.
\end{lemma}

\begin{proof}
Soit $\text{pr} : P \times \mathbf{P}^1 \to P$ la projection.
Puisque nous démontrons une égalité de cycles, il suffit de considérer
$\alpha$, resp.\ $\beta$, comme un cycle sur $P$, resp.\ $P \times \mathbf{P}^1$,
et de démontrer le résultat après image directe par $\text{pr}$.
Comme $\text{pr}^*X = X \times \mathbf{P}^1$ et que
$\text{pr}$ induit un isomorphisme de $C'_0$ sur $C$,
la formule de projection (lemme \ref{intersection-lemma-projection-formula-flat})
donne
$$
\text{pr}_*([C'_0]_r \cdot X \times \mathbf{P}^1) = [C]_r \cdot X = \alpha
$$
D'autre part, on a $[C'_0]_r = C' \cdot P \times 0$
comme cycles sur $P \times \mathbf{P}^1$,
d'après le lemme \ref{intersection-lemma-rational-equivalence-and-intersection}.
Par conséquent,
$$
[C'_0]_r \cdot X \times \mathbf{P}^1 =
(C' \cdot P \times 0) \cdot X \times \mathbf{P}^1 =
(C' \cdot X \times \mathbf{P}^1) \cdot P \times 0
$$
par associativité (lemme \ref{intersection-lemma-associative}) et commutativité du
produit d'intersection. Il reste à montrer que le produit d'intersection de
$C' \cdot X \times \mathbf{P}^1$ avec $P \times 0$ sur
$P \times \mathbf{P}^1$ est égal, en tant que cycle, au produit d'intersection de
$\beta$ avec $X \times 0$ sur $X \times \mathbf{P}^1$. Écrivons
$C' \cdot X \times \mathbf{P}^1 = \sum n_k[E_k]$, et donc
$\beta = \sum n_k[E_k]$, pour certaines sous-variétés
$E_k \subset X \times \mathbf{P}^1 \subset P \times \mathbf{P}^1$.
Les deux intersections sont alors égales à $\sum m_k[E_{k, 0}]$ d'après le
lemme \ref{intersection-lemma-rational-equivalence-and-intersection}, appliqué deux fois.
Cela achève la démonstration.
\end{proof}




\section{Projections}
\label{intersection-section-projection}

\noindent
Rappelons que nous travaillons sur un corps de base algébriquement clos fixé
$\mathbf{C}$. Si $V$ est un espace vectoriel de dimension finie sur $\mathbf{C}$,
nous posons
$$
\mathbf{P}(V) = \text{Proj}(\text{Sym}(V))
$$
où $\text{Sym}(V)$ est l'algèbre symétrique de $V$ sur $\mathbf{C}$.
Voir Constructions, exemple \ref{constructions-example-projective-space}.
La convention est choisie de telle sorte que
$V = \Gamma(\mathbf{P}(V), \mathcal{O}_{\mathbf{P}(V)}(1))$.
Bien entendu, on a $\mathbf{P}(V) \cong \mathbf{P}^n_{\mathbf{C}}$ si
$\dim(V) = n + 1$. Notons que $\mathbf{P}(V)$ est une variété projective
non singulière.

\medskip\noindent
Soit $p \in \mathbf{P}(V)$ un point fermé. Le point $p$ correspond à une
surjection $V \to L_p$ d'espaces vectoriels avec $\dim(L_p) = 1$, voir
Constructions, lemme \ref{constructions-lemma-apply}.
Notons $W_p = \Ker(V \to L_p)$.
La {\it projection de centre $p$} est le morphisme
$$
r_p : \mathbf{P}(V) \setminus \{p\} \longrightarrow \mathbf{P}(W_p)
$$
de Constructions, lemme \ref{constructions-lemma-morphism-proj}.
Voici un lemme préliminaire.

\begin{lemma}
\label{intersection-lemma-projection-generically-finite}
Soit $V$ un espace vectoriel de dimension $n + 1$.
Soit $X \subset \mathbf{P}(V)$ un sous-schéma fermé.
Si $X \not = \mathbf{P}(V)$, il existe un ouvert de Zariski non vide
$U \subset \mathbf{P}(V)$
tel que, pour tout point fermé $p \in U$, la restriction
de la projection $r_p$ définisse un morphisme fini
$r_p|_X : X \to \mathbf{P}(W_p)$.
\end{lemma}

\begin{proof}
Le lemme vaut, affirmons-nous, avec $U = \mathbf{P}(V) \setminus X$. Pour tout point
fermé $p$ de $U$, on obtient en effet un morphisme $r_p|_X : X \to \mathbf{P}(W_p)$.
Ce morphisme est propre, car $X$ est un schéma propre
(Morphismes, lemmes \ref{morphisms-lemma-locally-projective-proper} et
\ref{morphisms-lemma-image-proper-scheme-closed}). D'autre part, les
fibres de $r_p$ sont des droites affines, comme le montre un calcul direct.
Ainsi, les fibres de $r_p|X$ sont propres et affines, donc finies
(Morphismes, lemme \ref{morphisms-lemma-finite-proper}).
Enfin, un morphisme propre à fibres finies est fini
(Cohomologie des schémas, lemme \ref{coherent-lemma-characterize-finite}).
\end{proof}

\begin{lemma}
\label{intersection-lemma-projection-generically-immersion}
Soit $V$ un espace vectoriel de dimension $n + 1$.
Soit $X \subset \mathbf{P}(V)$ une sous-variété fermée.
Soit $x \in X$ un point non singulier.
\begin{enumerate}
\item Si $\dim(X) < n - 1$, il existe un ouvert de Zariski non vide
$U \subset \mathbf{P}(V)$ tel que, pour tout point fermé $p \in U$, le
morphisme $r_p|_X : X \to r_p(X)$ soit un
isomorphisme au-dessus d'un voisinage ouvert de $r_p(x)$.
\item Si $\dim(X) = n - 1$, il existe un ouvert de Zariski non vide
$U \subset \mathbf{P}(V)$ tel que, pour tout point fermé $p \in U$, le
morphisme $r_p|_X : X \to \mathbf{P}(W_p)$ soit \'etale en $x$.
\end{enumerate}
\end{lemma}

\begin{proof}
Démonstration de (1). Notons que si $x, y \in X$ ont la même image par
$r_p$, alors $p$ appartient à la droite $\overline{xy}$.
Considérons le schéma de type fini
$$
T = \{(y, p) \mid
y \in X \setminus \{x\},\ p \in \mathbf{P}(V),\ p \in \overline{xy}\}
$$
et les morphismes $T \to X$ et $T \to \mathbf{P}(V)$ donnés par
$(y, p) \mapsto y$ et $(y, p) \mapsto p$.
Puisque toute fibre de $T \to X$ est une droite, on voit que
la dimension de $T$ est $\dim(X) + 1 < \dim(\mathbf{P}(V))$.
Ainsi, $T \to \mathbf{P}(V)$ n'est pas surjectif. D'autre part,
considérons le schéma de type fini
$$
T' = \{p \mid
p \in \mathbf{P}(V) \setminus \{x\},
\ \overline{xp}\text{ tangente à }X\text{ en }x\}
$$
La dimension de $T'$ est alors $\dim(X) < \dim(\mathbf{P}(V))$.
Ainsi, le morphisme $T' \to \mathbf{P}(V)$ n'est pas non plus surjectif.
Soit $U \subset \mathbf{P}(V) \setminus X$ un ouvert non vide disjoint
de ces images ; un tel $U$ existe, car les images de $T$ et de $T'$
dans $\mathbf{P}(V)$ sont constructibles d'après
Morphismes, lemme \ref{morphisms-lemma-chevalley}.
Alors, pour $p \in U$ fermé, la projection
$r_p|_X : X \to \mathbf{P}(W_p)$ induit une application injective sur
l'espace tangent en $x$ et $r_p^{-1}(\{r_p(x)\}) = \{x\}$.
Cela signifie que $r_p$ est non ramifié en $x$
(Variétés, lemme \ref{varieties-lemma-injective-tangent-spaces-unramified}),
fini d'après le lemme \ref{intersection-lemma-projection-generically-finite},
et que $r_p^{-1}(\{r_p(x)\}) = \{x\}$ ; le lemme de Morphismes \'etales
\ref{etale-lemma-finite-unramified-one-point} s'applique donc, et
il existe un voisinage ouvert $R$ de $r_p(x)$
dans $\mathbf{P}(W_p)$ tel que $(r_p|_X)^{-1}(R) \to R$ soit une
immersion fermée, ce qui démontre (1).

\medskip\noindent
Démonstration de (2). Dans ce cas, on conclut encore que le morphisme
$T' \to \mathbf{P}(V)$ n'est pas surjectif.
En raisonnant comme ci-dessus, on conclut que, pour $U \subset \mathbf{P}(V)$
évitant $X$ et l'image de $T'$, la projection
$r_p|_X : X \to \mathbf{P}(W_p)$ est \'etale en $x$ et finie.
\end{proof}

\begin{lemma}
\label{intersection-lemma-projection-generically-immersion-refined}
Soit $V$ un espace vectoriel. Soient $Z \subset X \subset \mathbf{P}(V)$ des
sous-variétés fermées avec $Z \not = X \not = \mathbf{P}(V)$. Soit $x \in Z$ un point
non singulier sur $X$. Il existe alors un ouvert de Zariski non vide
$U \subset \mathbf{P}(V)$ tel que, pour tout point fermé $p \in U$, le
morphisme $r_p|_Z : Z \to r_p(Z)$ soit un isomorphisme au-dessus d'un voisinage ouvert
de $r_p(x)$.
\end{lemma}

\begin{proof}
En raisonnant comme dans la démonstration du lemme \ref{intersection-lemma-projection-generically-immersion},
on trouve un ouvert de Zariski non vide $U$ tel que, pour $p \in U$, l'application
$r_p|_X : X \to r_p(X)$ soit \'etale en $x$ et que
$r_p^{-1}(r_p(\{x\})) \cap Z = \{x\}$. Détails omis.
Posons $Z' = r_p(Z)$, considéré comme sous-variété fermée de $\mathbf{P}(W_p)$.
Alors le morphisme $\pi : Z \to Z'$ est fini
(lemme \ref{intersection-lemma-projection-generically-finite}), non ramifié en $x$
(petit détail omis) et $\pi^{-1}(\pi(\{x\})) = \{x\}$. Le résultat découle
du lemme de Morphismes \'etales \ref{etale-lemma-finite-unramified-one-point}.
\end{proof}

\begin{lemma}
\label{intersection-lemma-projection-injective}
Soit $V$ un espace vectoriel de dimension $n + 1$.
Soient $Y, Z \subset \mathbf{P}(V)$ des sous-variétés fermées.
Il existe un ouvert de Zariski non vide $U \subset \mathbf{P}(V)$
tel que, pour tout point fermé $p \in U$, on ait
$$
Y \cap r_p^{-1}(r_p(Z)) = (Y \cap Z) \cup E
$$
où $E \subset Y$ est fermé et
$\dim(E) \leq \dim(Y) + \dim(Z) + 1 - n$.
\end{lemma}

\begin{proof}
Posons $Y' = Y \setminus Y \cap Z$.
Soient $y \in Y'$, $z \in Z$ des points fermés tels que $r_p(y) = r_p(z)$.
Alors $p$ appartient à la droite $\overline{yz}$ passant par $y$ et $z$.
Considérons le schéma de type fini
$$
T = \{(y, z, p) \mid y \in Y', z \in Z, p \in \overline{yz}\}
$$
et le morphisme $T \to \mathbf{P}(V)$ donné par $(y, z, p) \mapsto p$.
Observons que $T$ est irréductible et que $\dim(T) = \dim(Y) + \dim(Z) + 1$.
La fibre générale de $T \to \mathbf{P}(V)$ est donc de dimension au plus
$\dim(Y) + \dim(Z) + 1 - n$ ; plus précisément, il existe un ouvert non vide
$U \subset \mathbf{P}(V) \setminus (Y \cup Z)$ au-dessus
duquel la fibre est de dimension au plus $\dim(Y) + \dim(Z) + 1 - n$
(Variétés, lemme \ref{varieties-lemma-dimension-fibres-locally-algebraic}).
Soit $p \in U$ un point fermé et soit $F \subset T$ la fibre
de $T \to \mathbf{P}(V)$ au-dessus de $p$. Alors
$$
(Y \cap r_p^{-1}(r_p(Z))) \setminus (Y \cap Z)
$$
est l'image de $F \to Y$, $(y, z, p) \mapsto y$. À nouveau d'après
Variétés, lemme \ref{varieties-lemma-dimension-fibres-locally-algebraic},
l'adhérence de l'image de $F \to Y$ est de dimension au plus
$\dim(Y) + \dim(Z) + 1 - n$.
\end{proof}

\begin{lemma}
\label{intersection-lemma-find-lines}
Soit $V$ un espace vectoriel. Soit $B \subset \mathbf{P}(V)$
une sous-variété fermée de codimension $\geq 2$.
Soit $p \in \mathbf{P}(V)$ un point fermé, $p \not \in B$.
Il existe alors une droite $\ell \subset \mathbf{P}(V)$
avec $p \in \ell$ et $\ell \cap B = \emptyset$. De plus, ces droites
recouvrent un ouvert de $\mathbf{P}(V)$.
\end{lemma}

\begin{proof}
Considérons l'image de $B$ par la projection
$r_p : \mathbf{P}(V) \to \mathbf{P}(W_p)$.
Puisque $\dim(W_p) = \dim(V) - 1$, on voit que $r_p(B)$
est de codimension $\geq 1$ dans $\mathbf{P}(W_p)$.
Pour tout $q \in \mathbf{P}(V)$ tel que $r_p(q) \not \in r_p(B)$,
on voit que la droite $\ell = \overline{pq}$ joignant $p$ et $q$ convient.
\end{proof}

\begin{lemma}
\label{intersection-lemma-doubly-transitive}
Soit $V$ un espace vectoriel. Soit $G = \text{PGL}(V)$.
Alors $G \times \mathbf{P}(V) \to \mathbf{P}(V)$ est
doublement transitif.
\end{lemma}

\begin{proof}
Omis. Indication : cela résulte du fait que $\text{GL}(V)$ agit de façon
doublement transitive sur les couples de vecteurs linéairement indépendants.
\end{proof}

\begin{lemma}
\label{intersection-lemma-determinant}
Soit $k$ un corps. Soit $n \geq 1$ un entier et soient
$x_{ij}, 1 \leq i, j \leq n$ des indéterminées. Alors
$$
\det
\left(
\begin{matrix}
x_{11} & x_{12} & \ldots & x_{1n} \\
x_{21} & \ldots & \ldots & \ldots \\
\ldots & \ldots & \ldots & \ldots \\
x_{n1} & \ldots & \ldots & x_{nn}
\end{matrix}
\right)
$$
est un élément irréductible de l'anneau de polynômes $k[x_{ij}]$.
\end{lemma}

\begin{proof}
Soit $V$ un espace vectoriel de dimension $n$. Géométriquement,
le lemme signifie que le lieu $C$ des applications linéaires non inversibles
$V \to V$ est irréductible (notons que, comme le déterminant est de degré
$1$ en chaque indéterminée, il est sans facteur carré). Soit $W$ un espace vectoriel de
dimension $n - 1$. Par algèbre linéaire élémentaire, le morphisme
$$
\Hom(W, V) \times \Hom(V, W) \longrightarrow \Hom(V, V),\quad
(\psi, \varphi) \longmapsto \psi \circ \varphi
$$
a pour image $C$. Puisque la source est irréductible, l'image l'est aussi.
Détails omis.
\end{proof}

\noindent
Soit $V$ un espace vectoriel de dimension $n + 1$. Posons $E = \text{End}(V)$.
Soit $E^\vee = \Hom(E, \mathbf{C})$ l'espace vectoriel dual.
Écrivons $\mathbf{P} = \mathbf{P}(E^\vee)$.
Il existe une application linéaire canonique
$$
V \longrightarrow V \otimes_\mathbf{C} E^\vee = \Hom(E, V)
$$
qui envoie $v \in V$ sur l'application $g \mapsto g(v)$ de $\Hom(E, V)$.
Rappelons que nous avons une application canonique
$E^\vee \to \Gamma(\mathbf{P}, \mathcal{O}_\mathbf{P}(1))$
qui est un isomorphisme. On obtient donc une application canonique
$$
\psi : V \otimes \mathcal{O}_\mathbf{P} \to V \otimes \mathcal{O}_\mathbf{P}(1)
$$
de faisceaux de modules sur $\mathbf{P}$ qui, sur les sections globales, redonne
l'application donnée. Rappelons qu'un fibré projectif $\mathbf{P}(\mathcal{E})$
est défini comme le Proj relatif de l'algèbre symétrique de $\mathcal{E}$, voir
Constructions, définition \ref{constructions-definition-projective-bundle}.
Nous allons étudier l'application rationnelle
entre $\mathbf{P}(V \otimes \mathcal{O}_\mathbf{P}(1))$ et
$\mathbf{P}(V \otimes \mathcal{O}_\mathbf{P})$ associée à $\psi$. D'après
Constructions, lemme \ref{constructions-lemma-relative-proj-base-change},
nous avons un isomorphisme canonique
$$
\mathbf{P}(V \otimes \mathcal{O}_\mathbf{P}) = \mathbf{P} \times \mathbf{P}(V)
$$
D'après Constructions, lemme \ref{constructions-lemma-twisting-and-proj},
on voit que
$$
\mathbf{P}(V \otimes \mathcal{O}_\mathbf{P}(1)) =
\mathbf{P}(V \otimes \mathcal{O}_\mathbf{P}) = \mathbf{P} \times \mathbf{P}(V)
$$
En combinant ceci avec
Constructions, lemme \ref{constructions-lemma-morphism-relative-proj},
on obtient
\begin{equation}
\label{intersection-equation-r-psi}
\mathbf{P} \times \mathbf{P}(V) \supset
U(\psi) \xrightarrow{r_\psi} \mathbf{P} \times \mathbf{P}(V)
\end{equation}
Pour mieux comprendre ceci, étudions ce qui se passe sur les fibres au-dessus de
$\mathbf{P}$. Soit $g \in E$ non nul. Il définit une application non nulle
$E^\vee \to \mathbf{C}$, donc un point $[g] \in \mathbf{P}$.
D'autre part, $g$ définit une application $\mathbf{C}$-linéaire $g : V \to V$.
Ainsi, d'après
Constructions, lemme \ref{constructions-lemma-morphism-proj},
on obtient une application
$$
\mathbf{P}(V) \supset U(g) \xrightarrow{r_g} \mathbf{P}(V)
$$
Nous utiliserons ci-dessous le fait que $U(g)$ est la fibre $U(\psi)_{[g]}$ et
que $r_g$ est la fibre de $r_\psi$ au-dessus du point $[g]$. Une autre observation
que nous utiliserons est que le complémentaire de $U(g)$ dans $\mathbf{P}(V)$ est
l'image de l'immersion fermée
$$
\mathbf{P}(\Coker(g)) \longrightarrow \mathbf{P}(V)
$$
et que l'image de $r_g$ est l'image de l'immersion fermée
$$
\mathbf{P}(\Im(g)) \longrightarrow \mathbf{P}(V)
$$

\begin{lemma}
\label{intersection-lemma-make-family}
Avec les notations ci-dessus, soient $X, Y$ des sous-variétés fermées de $\mathbf{P}(V)$
qui se coupent proprement, telles que $X \not = \mathbf{P}(V)$ et
$X \cap Y \not = \emptyset$. Pour une droite générale $\ell \subset \mathbf{P}$
avec $[\text{id}_V] \in \ell$, on a
\begin{enumerate}
\item $X \subset U_g$ pour tout $[g] \in \ell$,
\item $g(X)$ coupe $Y$ proprement pour tout $[g] \in \ell$.
\end{enumerate}
\end{lemma}

\begin{proof}
Soit $B \subset \mathbf{P}$ l'ensemble des points «~mauvais~», c'est-à-dire des
points $[g]$ qui ne vérifient pas (1) ou (2). Notons que
$[\text{id}_V] \not \in B$ par hypothèse. De plus, $B$ est fermé.
Il suffit donc de montrer que $\dim(B) \leq \dim(\mathbf{P}) - 2$
(lemme \ref{intersection-lemma-find-lines}).

\medskip\noindent
Considérons d'abord l'ouvert $G = \text{PGL}(V) \subset \mathbf{P}$
formé des points $[g]$ tels que $g : V \to V$ soit inversible.
Puisque $G$ agit doublement transitivement sur $\mathbf{P}(V)$
(lemme \ref{intersection-lemma-doubly-transitive}),
on voit que
$$
T = \{(x, y, [g]) \mid x \in X, y \in Y, [g] \in G, r_g(x) = y\}
$$
est une fibration localement triviale sur $X \times Y$ dont la fibre est égale
au stabilisateur d'un point dans $G$. Ainsi, $T$ est une variété.
Observons que la fibre de $T \to G$ au-dessus de $[g]$ est $r_g(X) \cap Y$.
Le morphisme $T \to G$ est surjectif, car toute transformée de $X$
rencontre $Y$ (notons que l'hypothèse selon laquelle $X$ et $Y$ se coupent
proprement et que $X \cap Y \not = \emptyset$ entraîne
$\dim(X) + \dim(Y) \geq \dim(\mathbf{P}(V))$, puis le
lemme de Variétés \ref{varieties-lemma-intersection-in-projective-space}
implique que toute transformée de $X$ rencontre $Y$).
Puisque les dimensions des fibres d'un morphisme dominant ne
sautent pas en codimension $1$
(Variétés, lemme \ref{varieties-lemma-dimension-fibres-locally-algebraic}),
on conclut que $B \cap G$ est de codimension $\geq 2$.

\medskip\noindent
Considérons ensuite le complémentaire $Z = \mathbf{P} \setminus G$.
C'est une variété irréductible, car le déterminant est un
polynôme irréductible (lemme \ref{intersection-lemma-determinant}).
Il suffit donc de montrer que $B$ ne contient pas le
point générique de $Z$. Pour un point général $[g] \in Z$,
le conoyau $V \to \Coker(g)$ est de dimension $1$, donc
$U(g)$ est le complémentaire d'un point. Comme $X \not = \mathbf{P}(V)$,
on voit que, pour $[g] \in Z$ général, on a $X \subset U(g)$.
De plus, le morphisme $r_g|_X : X \to r_g(X)$ est fini, donc
$\dim(r_g(X)) = \dim(X)$.
D'autre part, pour un tel $g$, l'image de $r_g$ est le
sous-espace fermé $H = \mathbf{P}(\Im(g)) \subset \mathbf{P}(V)$
qui est de codimension $1$.
Pour un point général de $Z$, on voit que $H \cap Y$ est de dimension $1$
de moins que $Y$ (comparer à
Variétés, lemme \ref{varieties-lemma-exact-sequence-induction}).
On voit ainsi qu'il faut montrer que $r_g(X)$ et $H \cap Y$
se coupent proprement dans $H$. Pour un choix fixé de $H$, on peut,
en composant $g$ à gauche avec un automorphisme, déplacer $r_g(X)$ par
un automorphisme arbitraire de $H = \mathbf{P}(\Im(g))$.
On peut donc raisonner comme ci-dessus pour conclure que l'intersection
de $H \cap Y$ avec $r_g(X)$ est propre pour $g$ général ayant
$H = \mathbf{P}(\Im(g))$. Certains détails sont omis.
\end{proof}







\section{Lemme de déplacement}
\label{intersection-section-moving-lemma}

\noindent
Le lemme de déplacement affirme qu'étant donnés un $r$-cycle $\alpha$ et un $s$-cycle
$\beta$, il existe $\alpha'$, $\alpha' \sim_{rat} \alpha$, tel que
$\alpha'$ et $\beta$ se coupent proprement (lemme \ref{intersection-lemma-moving-move}).
Voir \cite{Samuel}, \cite{ChevalleyI}, \cite{ChevalleyII}.
Le point clé est le lemme \ref{intersection-lemma-moving} ; le lecteur pourra trouver
ce lemme sous la forme énoncée dans
\cite[exemple 11.4.1]{F} et une démonstration dans \cite{Roberts}.

\begin{lemma}
\label{intersection-lemma-moving}
\begin{reference}
Voir \cite{Roberts}.
\end{reference}
Soit $X \subset \mathbf{P}^N$ une sous-variété fermée non singulière.
Posons $n = \dim(X)$ et $0 \leq d, d' < n$. Soit $Z \subset X$ une sous-variété
fermée de dimension $d$ et soient $T_i \subset X$, $i \in I$, une famille finie
de sous-variétés fermées de dimension $d'$. Il existe alors
une sous-variété $C \subset \mathbf{P}^N$ telle que $C$ coupe $X$
proprement et telle que
$$
C \cdot X = Z + \sum\nolimits_{j \in J} m_j Z_j
$$
où les $Z_j \subset X$ sont irréductibles de dimension $d$, distinctes de $Z$, et
$$
\dim(Z_j \cap T_i) \leq \dim(Z \cap T_i)
$$
avec inégalité stricte si $Z$ ne rencontre pas $T_i$ proprement dans $X$.
\end{lemma}

\begin{proof}
Écrivons $\mathbf{P}^N = \mathbf{P}(V_N)$, de sorte que $\dim(V_N) = N + 1$.
Nous allons choisir une suite de projections depuis des points
\begin{align*}
& r_N : 
\mathbf{P}(V_N) \setminus \{p_N\} \to \mathbf{P}(V_{N - 1}), \\
& r_{N - 1} :
\mathbf{P}(V_{N - 1}) \setminus \{p_{N - 1}\} \to \mathbf{P}(V_{N - 2}), \\
& \ldots,\\
& r_{n + 1} :
\mathbf{P}(V_{n + 1}) \setminus \{p_{n + 1}\} \to \mathbf{P}(V_n)
\end{align*}
comme à la section \ref{intersection-section-projection}. À chaque étape, nous choisirons
$p_N, p_{N - 1}, \ldots, p_{n + 1}$ dans un ouvert de Zariski convenable.
Choisissons un point fermé $x \in Z \subset X$. Pour tout $i$, choisissons des
points fermés $x_{it} \in T_i \cap Z$, au moins un dans chaque composante irréductible
de $T_i \cap Z$. En prenant la composée, nous obtenons
un morphisme
$$
\pi = (r_{n + 1} \circ \ldots \circ r_N)|_X :
X \longrightarrow \mathbf{P}(V_n)
$$
qui possède les propriétés suivantes :
\begin{enumerate}
\item $\pi$ est fini,
\item $\pi$ est \'etale en $x$ et en tous les $x_{it}$,
\item $\pi|_Z : Z \to \pi(Z)$ est un isomorphisme
au-dessus d'un voisinage ouvert de $\pi(x_{it})$,
\item $T_i \cap \pi^{-1}(\pi(Z)) = (T_i \cap Z) \cup E_i$, où
$E_i \subset T_i$ est fermé et
$\dim(E_i) \leq d + d' + 1 - (n + 1) = d + d' - n$.
\end{enumerate}
Il résulte immédiatement des
Lemmes \ref{intersection-lemma-projection-generically-finite},
\ref{intersection-lemma-projection-generically-immersion},
\ref{intersection-lemma-projection-generically-immersion-refined} et
\ref{intersection-lemma-projection-injective}, ainsi que d'une récurrence, que ce choix est possible ;
observons que la dernière projection part de $\mathbf{P}(V_{n + 1})$ et que
$\dim(V_{n + 1}) = n + 2$, ce qui explique l'inégalité de (4). Plus
précisément, soient $X_j, Z_j, T_{i, j} \subset \mathbf{P}(V_j)$ les images
de $X$, $Z$, $T_i$ pour $j = N, \ldots, n$. Alors
$T_{i, j + 1} \cap r_{j + 1}^{-1}(Z_j) = (T_{i, j + 1} \cap Z_{j + 1})
\cup E_{i, j + 1}$, où nous disposons d'une borne sur la dimension de $E_{i, j + 1}$
d'après le lemme \ref{intersection-lemma-projection-injective}. Notons alors que
$T_i \cap \pi^{-1}(\pi(Z))$ est la réunion de $T_i \cap Z$
et des images inverses des $E_{i, j}$. Comme les morphismes de transition
entre les $X_j$ sont finis, on obtient la borne de dimension voulue.

\medskip\noindent
Soit $C \subset \mathbf{P}(V_N)$ l'adhérence schématique de
$(r_{n + 1} \circ \ldots \circ r_N)^{-1}(\pi(Z))$. Puisque $\pi$
est \'etale au point $x$ de $Z$, le sous-schéma fermé
$C \cap X = \pi^{-1}(\pi(Z))$
contient $Z$ avec multiplicité $1$ (le calcul local est omis).
Le lemme \ref{intersection-lemma-transversal-subschemes} donne donc
$$
C \cdot X = [Z] + \sum m_j[Z_j]
$$
pour certaines sous-variétés $Z_j \subset X$ de dimension $d$. Notons que
$$
C \cap X = \pi^{-1}(\pi(Z))
$$
ensemblistement. Par conséquent,
$T_i \cap Z_j \subset T_i \cap \pi^{-1}(\pi(Z)) \subset (T_i \cap Z) \cup E_i$.
Pour toute composante irréductible de $T_i \cap Z_j$ contenue dans $E_i$, nous
avons la borne de dimension voulue. Enfin, soit $V$ une composante irréductible
de $T_i \cap Z_j$ contenue dans $T_i \cap Z$. Pour achever
la démonstration, il suffit de montrer que $V$ ne contient aucun des
points $x_{it}$, car alors $\dim(V) < \dim(Z \cap T_i)$.
Pour cela, il suffit de montrer que $x_{it} \not \in Z_j$
pour tous $i, t, j$.

\medskip\noindent
Posons $Z' = \pi(Z)$ et $Z'' = \pi^{-1}(Z')$, schématiquement. D'après
la condition (3), il existe un ouvert $U \subset \mathbf{P}(V_n)$ contenant
$\pi(x_{it})$ tel que $\pi^{-1}(U) \cap Z \to U \cap Z'$ soit un isomorphisme.
En particulier, $Z \to Z'$ est un isomorphisme local en $x_{it}$.
D'autre part, $Z'' \to Z'$ est \'etale en $x_{it}$ d'après la condition (2).
L'immersion fermée $Z \to Z''$ est donc \'etale en $x_{it}$
(Morphismes, lemme \ref{morphisms-lemma-etale-permanence}).
Ainsi $Z = Z''$ dans un voisinage de Zariski de $x_{it}$, ce qui démontre
l'assertion.
\end{proof}

\noindent
Le déplacement proprement dit s'effectue à l'aide du lemme suivant.

\begin{lemma}
\label{intersection-lemma-move}
Soit $C \subset \mathbf{P}^N$ une sous-variété fermée.
Soit $X \subset \mathbf{P}^N$ une sous-variété et soient $T_i \subset X$
une collection finie de sous-variétés fermées.
Supposons que $C$ et $X$ se coupent proprement.
Alors il existe une sous-variété fermée
$C' \subset \mathbf{P}^N \times \mathbf{P}^1$ telle que
\begin{enumerate}
\item $C' \to \mathbf{P}^1$ soit dominant,
\item $C'_0 = C$ schématiquement,
\item $C'$ et $X \times \mathbf{P}^1$ se coupent proprement,
\item $C'_\infty$ coupe proprement chacun des $T_i$ donnés.
\end{enumerate}
\end{lemma}

\begin{proof}
Si $C \cap X = \emptyset$, prenons la famille constante
$C' = C \times \mathbf{P}^1$. Nous pouvons donc supposer
$C \cap X \not = \emptyset$.

\medskip\noindent
Écrivons $\mathbf{P}^N = \mathbf{P}(V)$, de sorte que $\dim(V) = N + 1$. Soit
$E = \text{End}(V)$. Soit $E^\vee = \Hom(E, \mathbf{C})$. Posons
$\mathbf{P} = \mathbf{P}(E^\vee)$ comme dans le lemme \ref{intersection-lemma-make-family}.
Choisissons une droite générale $\ell \subset \mathbf{P}$ passant par $\text{id}_V$.
Définissons $C' \subset \ell \times \mathbf{P}(V)$ comme le
sous-schéma fermé de fibre $r_g(C)$ au-dessus de $[g] \in \ell$.
Plus précisément, $C'$ est l'image de
$$
\ell \times C \subset \mathbf{P} \times \mathbf{P}(V)
$$
par le morphisme (\ref{intersection-equation-r-psi}). D'après le lemme \ref{intersection-lemma-make-family},
cela a un sens, c'est-à-dire que $\ell \times C \subset U(\psi)$. Le morphisme
$\ell \times C \to C'$ est fini et $C'_{[g]} = r_g(C)$ ensemblistement
pour tout $[g] \in \ell$. Les parties (1) et (2) sont claires avec
$0 = [\text{id}_V] \in \ell$. La partie (3) résulte du fait
que $r_g(C)$ et $X$ se coupent proprement pour tout $[g] \in \ell$.
La partie (4) résulte du fait qu'un point général $\infty = [g] \in \ell$
est un point général de $\mathbf{P}$ et que, pour un tel point,
$r_g(C) \cap T$ est propre pour toute sous-variété fermée $T$ de $\mathbf{P}(V)$.
Les détails sont omis.
\end{proof}

\begin{lemma}
\label{intersection-lemma-moving-move}
Soit $X$ une variété projective non singulière. Soit $\alpha$ un
$r$-cycle et soit $\beta$ un $s$-cycle sur $X$. Alors il existe
un $r$-cycle $\alpha'$ tel que $\alpha' \sim_{rat} \alpha$ et
tel que $\alpha'$ et $\beta$ se coupent proprement.
\end{lemma}

\begin{proof}
Écrivons $\beta = \sum n_i[T_i]$ pour certaines sous-variétés $T_i \subset X$
de dimension $s$. Par linéarité, nous pouvons supposer que $\alpha = [Z]$ pour
une sous-variété fermée irréductible $Z \subset X$ de dimension $r$.
Nous démontrerons le lemme par récurrence sur le maximum $e$ des entiers
$$
\dim(Z \cap T_i)
$$
Le cas initial est $e = r + s - \dim(X)$. Dans ce cas, $Z$ rencontre
$\beta$ proprement et le lemme est trivial.

\medskip\noindent
Étape de récurrence. Supposons $e > r + s - \dim(X)$. Choisissons un plongement
$X \subset \mathbf{P}^N$ et appliquons le lemme \ref{intersection-lemma-moving} afin de trouver une
sous-variété fermée $C \subset \mathbf{P}^N$ telle que
$C \cdot X = [Z] + \sum m_j[Z_j]$ et telle que l'hypothèse de
récurrence s'applique à chaque $Z_j$. Appliquons ensuite le lemme \ref{intersection-lemma-move}
à $C$, $X$, $T_i$ pour trouver $C' \subset \mathbf{P}^N \times \mathbf{P}^1$.
Soit $\gamma = C' \cdot X \times \mathbf{P}^1$, considéré comme un cycle
sur $X \times \mathbf{P}^1$. D'après le lemme \ref{intersection-lemma-transfer}, nous avons
$$
[Z] + \sum m_j[Z_j] = \text{pr}_{X, *}(\gamma \cdot X \times 0)
$$
D'autre part, le cycle
$\gamma_\infty = \text{pr}_{X, *}(\gamma \cdot X \times \infty)$
est supporté sur $C'_\infty \cap X$ et rencontre donc $\beta$ proprement.
Ainsi $[Z] \sim_{rat} - \sum m_j[Z_j] + \gamma_\infty$
d'après le lemme \ref{intersection-lemma-rational-equivalence-and-intersection}. Comme, par
récurrence, chaque $[Z_j]$ est rationnellement équivalent à un cycle qui rencontre
$\beta$ proprement, cela achève la démonstration.
\end{proof}



\section{Produits d'intersection et équivalence rationnelle}
\label{intersection-section-intersections-and-rational-equivalence}

\noindent
Avec les définitions précédentes, montrons que le produit d'intersection
est bien défini modulo équivalence rationnelle. Traitons d'abord un
cas particulier.

\begin{lemma}
\label{intersection-lemma-well-defined-special-case}
Soit $X$ une variété non singulière. Soit
$W \subset X \times \mathbf{P}^1$ une sous-variété de dimension $(s + 1)$
dominant $\mathbf{P}^1$. Soient $W_a$, resp.\ $W_b$, les fibres de
$W \to \mathbf{P}^1$ au-dessus de $a$, resp.\ de $b$. Soit $V$ une sous-variété de dimension $r$
de $X$ telle que $V$ coupe proprement $W_a$ et
$W_b$. Alors $[V] \cdot [W_a]_r \sim_{rat} [V] \cdot [W_b]_r$.
\end{lemma}

\begin{proof}
Nous avons $[W_a]_r = \text{pr}_{X,*}(W \cdot X \times a)$ et de même pour
$[W_b]_r$, voir le lemme \ref{intersection-lemma-rational-equivalence-and-intersection}.
Nous nous ramenons donc à démontrer
$$
V \cdot \text{pr}_{X,*}( W \cdot X \times a) \sim_{rat} V \cdot
\text{pr}_{X,*}( W \cdot X\times b).
$$
En appliquant la formule de projection du
lemme \ref{intersection-lemma-projection-formula-flat}, on obtient
$$
V \cdot \text{pr}_{X,*}( W \cdot X\times a) =
\text{pr}_{X,*}(V \times \mathbf{P}^1 \cdot (W \cdot X\times a))
$$
et de même pour $b$. Nous nous ramenons donc à démontrer
$$
\text{pr}_{X,*}(V \times \mathbf{P}^1 \cdot (W \cdot X\times a))
\sim_{rat}
\text{pr}_{X,*}(V \times \mathbf{P}^1 \cdot (W \cdot X\times b))
$$
Si $V \times \mathbf{P}^1$ et $W$ se coupent proprement, alors
l'associativité des multiplicités d'intersection
(lemme \ref{intersection-lemma-associative})
donne $V \times \mathbf{P}^1 \cdot (W \cdot X\times a) =
(V \times \mathbf{P}^1 \cdot W) \cdot X \times a$,
et de même pour $b$. Nous nous ramenons donc à démontrer
$$
\text{pr}_{X,*}((V \times \mathbf{P}^1 \cdot W) \cdot X\times a)
\sim_{rat}
\text{pr}_{X,*}((V \times \mathbf{P}^1 \cdot W) \cdot X\times b)
$$
ce qui est vrai d'après le lemme \ref{intersection-lemma-rational-equivalence-and-intersection}.

\medskip\noindent
Le raisonnement précédent ne fonctionne pas tout à fait. L'obstacle est que
nous ne savons pas si $V \times \mathbf{P}^1$ et $W$ se coupent proprement.
Nous savons seulement que $V$ et $W_a$, ainsi que $V$ et $W_b$, se coupent proprement.
Soient $Z_i$, $i \in I$, les composantes irréductibles de
$V \times \mathbf{P}^1 \cap W$. Nous savons alors que
$\dim(Z_i) \geq r + 1 + s + 1 - n - 1 = r + s + 1 - n$, où $n = \dim(X)$, voir
le lemme \ref{intersection-lemma-intersect-in-smooth}. Comme nous avons supposé
que $V$ et $W_a$ se coupent proprement, on voit que
$\dim(Z_{i, a}) = r + s - n$ ou $Z_{i, a} = \emptyset$.
D'autre part, si $Z_{i, a} \not = \emptyset$, alors
$\dim(Z_{i, a}) \geq \dim(Z_i) - 1 = r + s - n$.
Il s'ensuit que $\dim(Z_i) = r + s + 1 - n$ si $Z_i$ rencontre $X \times a$
et, dans ce cas, $Z_i \to \mathbf{P}^1$ est surjectif.
Nous pouvons donc écrire $I = I' \amalg I''$, où $I'$ est l'ensemble des $i \in I$
tels que $Z_i \to \mathbf{P}^1$ soit surjectif, et $I''$ l'ensemble des
$i \in I$ tels que $Z_i$ soit au-dessus d'un point fermé $t_i \in \mathbf{P}^1$
avec $t_i \not = a$ et $t_i \not = b$. Considérons le cycle
$$
\gamma = \sum\nolimits_{i \in I'} e_i [Z_i]
$$
où nous posons
$$
e_i = \sum\nolimits_p (-1)^p
\text{longueur}_{\mathcal{O}_{X \times \mathbf{P}^1, Z_i}}
\text{Tor}_p^{\mathcal{O}_{X \times \mathbf{P}^1, Z_i}}(
\mathcal{O}_{V \times \mathbf{P}^1, Z_i}, \mathcal{O}_{W, Z_i})
$$
Nous montrerons que $\gamma$ peut remplacer
le produit d'intersection de $V \times \mathbf{P}^1$ et $W$.

\medskip\noindent
Nous le démontrerons à l'aide de l'associativité des produits d'intersection, exactement
comme ci-dessus. Soit $U = \mathbf{P}^1 \setminus \{t_i, i \in I''\}$.
Notons que $X \times a$ et $X \times b$ sont contenus dans $X \times U$.
Les sous-variétés
$$
V \times U,\quad W_U,\quad X \times a\quad\text{de}\quad X \times U
$$
se rencontrent deux à deux transversalement par notre choix de $U$ et, de plus,
$\dim(V \times U \cap W_U \cap X \times a) = \dim(V \cap W_a)$ a
la dimension attendue. On voit donc que
$$
V \times U \cdot (W_U \cdot X \times a) =
(V \times U \cdot W_U) \cdot X \times a
$$
comme cycles sur $X \times U$, d'après le lemme \ref{intersection-lemma-associative}.
Par construction, $\gamma$ se restreint au cycle $V \times U \cdot W_U$
sur $X \times U$. Trivialement,
$V \times \mathbf{P}^1 \cdot (W \times X \times a)$ se restreint
à $V \times U \cdot (W_U \cdot X \times a)$ sur $X \times U$.
Par conséquent,
$$
V \times \mathbf{P}^1 \cdot (W \cdot X \times a) =
\gamma \cdot X \times a
$$
comme cycles sur $X \times \mathbf{P}^1$ (car les deux membres
sont contenus dans $X \times U$ et sont égaux après restriction
à $X \times U$ d'après ce qui précède). Comme on a le même résultat pour $b$,
on conclut que
\begin{align*}
V \cdot [W_a]
& =
\text{pr}_{X,*}(V \times \mathbf{P}^1 \cdot (W \cdot X\times a)) \\
& =
\text{pr}_{X, *}(\gamma \cdot X \times a) \\
& \sim_{rat} 
\text{pr}_{X, *}(\gamma \cdot X \times b) \\
& =
\text{pr}_{X,*}(V \times \mathbf{P}^1 \cdot (W \cdot X\times b)) \\
& =
V \cdot [W_b]
\end{align*}
La première et la dernière égalité résultent du premier paragraphe de la démonstration,
la deuxième et l'avant-dernière ont été démontrées dans le présent paragraphe, et
l'équivalence centrale est donnée par le
lemme \ref{intersection-lemma-rational-equivalence-and-intersection}.
\end{proof}

\begin{theorem}
\label{intersection-theorem-well-defined}
Soit $X$ une variété projective non singulière. Soient $\alpha$, resp.\ $\beta$,
un $r$-cycle, resp.\ un $s$-cycle, sur $X$. Supposons que $\alpha$ et $\beta$
se coupent proprement, de sorte que $\alpha \cdot \beta$ soit défini. Enfin,
supposons que $\alpha \sim_{rat} 0$. Alors $\alpha \cdot \beta \sim_{rat} 0$.
\end{theorem}

\begin{proof}
Choisissons une immersion fermée $X \subset \mathbf{P}^N$.
Par linéarité, il suffit de démontrer le résultat lorsque $\beta = [Z]$ pour une
sous-variété fermée de dimension $s$, notée $Z \subset X$, qui rencontre $\alpha$
proprement. La condition $\alpha \sim_{rat} 0$ signifie qu'il
existe un nombre fini de sous-variétés fermées de dimension $(r + 1)$,
$W_i \subset X \times \mathbf{P}^1$, telles que
$$
\alpha = \sum [W_{i, a_i}]_r - [W_{i, b_i}]_r
$$
pour certains couples de points $a_i, b_i$ de $\mathbf{P}^1$.
Soient $W_{i, a_i}^t$ et $W_{i, b_i}^t$ les composantes irréductibles
de $W_{i, a_i}$ et $W_{i, b_i}$.
Nous raisonnons par récurrence sur le maximum $d$ des entiers
$$
\dim(Z \cap W_{i, a_i}^t),\quad \dim(Z \cap W_{i, b_i}^t)
$$
La principale difficulté dans la suite de la démonstration est que, bien que $Z$
rencontre $\alpha$ proprement, il se peut que $Z$ ne rencontre pas les
variétés « intermédiaires » $W_{i, a_i}^t$
et $W_{i, b_i}^t$ proprement, c'est-à-dire que l'on peut avoir $d >  r + s - \dim(X)$.

\medskip\noindent
Cas initial : $d = r + s - \dim(X)$. Dans ce cas, toutes les intersections de
$Z$ avec les $W_{i, a_i}^t$ et les $W_{i, b_i}^t$ sont propres et le
résultat voulu découle du lemme \ref{intersection-lemma-well-defined-special-case},
car celui-ci montre que
$[Z] \cdot [W_{i, a_i}]_r \sim_{rat} [Z] \cdot [W_{i, b_i}]_r$ pour
tout $i$.

\medskip\noindent
Étape de récurrence : $d >  r + s - \dim(X)$. Appliquons le
lemme \ref{intersection-lemma-moving} à $Z \subset X$ et à
la famille de sous-variétés $\{W_{i, a_i}^t, W_{i, b_i}^t\}$. On obtient alors une
sous-variété fermée $C \subset \mathbf{P}^N$ qui rencontre $X$
proprement et telle que
$$
C \cdot X = [Z] + \sum m_j [Z_j]
$$
et telle que
$$
\dim(Z_j \cap W_{i, a_i}^t) \leq \dim(Z \cap W_{i, a_i}^t),\quad
\dim(Z_j \cap W_{i, b_i}^t) \leq \dim(Z \cap W_{i, b_i}^t)
$$
avec inégalité stricte si le membre de droite est $> r + s - \dim(X)$.
Cela implique deux choses : (a) l'hypothèse de récurrence s'applique à
chaque $Z_j$, et (b) $C \cdot X$ et $\alpha$ se coupent proprement (car
$\alpha$ est une combinaison linéaire des $[W_{i, a_i}^t]$ et
$[W_{i, a_i}^t]$ qui rencontrent $Z$ proprement).
Choisissons ensuite $C' \subset \mathbf{P}^N \times \mathbf{P}^1$
comme dans le lemme \ref{intersection-lemma-move}, relativement à $C$, $X$ et
$W_{i, a_i}^t$, $W_{i, b_i}^t$.
Écrivons $C' \cdot X \times \mathbf{P}^1 = \sum n_k [E_k]$ pour
certaines sous-variétés $E_k \subset X \times \mathbf{P}^1$ de
dimension $s + 1$. Notons que $n_k > 0$ pour tout $k$, d'après la
proposition \ref{intersection-proposition-positivity}.
D'après le lemme \ref{intersection-lemma-transfer}, nous avons
$$
[Z] + \sum m_j [Z_j] = \sum n_k[E_{k, 0}]_s
$$
Comme $E_{k, 0} \subset C \cap X$, on voit que $[E_{k, 0}]_s$ et $\alpha$
se coupent proprement. D'autre part, le cycle
$$
\gamma = \sum n_k[E_{k, \infty}]_s
$$
est supporté sur $C'_\infty \cap X$ et
rencontre donc proprement chaque $W_{i, a_i}^t$, $W_{i, b_i}^t$.
Par le cas initial et la linéarité, on voit donc que
$$
\gamma \cdot \alpha \sim_{rat} 0
$$
Comme nous avons vu que $E_{k, 0}$ et $E_{k, \infty}$
coupent $\alpha$ proprement, le lemme \ref{intersection-lemma-well-defined-special-case},
appliqué à $E_k \subset X \times \mathbf{P}^1$ et $\alpha$, donne
$$
[E_{k, 0}] \cdot \alpha \sim_{rat} [E_{k, \infty}] \cdot \alpha
$$
En regroupant le tout, nous avons
\begin{align*}
[Z] \cdot \alpha
& =
(\sum n_k[E_{k, 0}]_r - \sum m_j[Z_j]) \cdot \alpha \\
& \sim_{rat}
\sum n_k [E_{k, 0}] \cdot \alpha \quad (\text{par l'hypothèse de récurrence})\\
& \sim_{rat}
\sum n_k [E_{k, \infty}] \cdot \alpha \quad (\text{par le lemme})\\
& =
\gamma \cdot \alpha \\
& \sim_{rat}
0 \quad (\text{par le cas initial})
\end{align*}
Cela achève la démonstration.
\end{proof}

\begin{remark}
\label{intersection-remark-quasi-projective}
Le lemme \ref{intersection-lemma-moving-move} et le
théorème \ref{intersection-theorem-well-defined}
valent aussi pour les variétés quasi-projectives non singulières, avec
la même démonstration. La seule modification consiste à démontrer la version suivante
du lemme de déplacement \ref{intersection-lemma-moving} : soit $X \subset \mathbf{P}^N$
une sous-variété fermée. Soient $n = \dim(X)$ et $0 \leq d, d' < n$. Soit
$X^{reg} \subset X$ l'ouvert des points non singuliers. Soit
$Z \subset X^{reg}$ une sous-variété fermée de dimension $d$ et soient
$T_i \subset X^{reg}$, $i \in I$, une collection finie de sous-variétés fermées
de dimension $d'$. Alors il existe une sous-variété $C \subset \mathbf{P}^N$
telle que $C$ rencontre $X$ proprement et que
$$
(C \cdot X)|_{X^{reg}} = Z + \sum\nolimits_{j \in J} m_j Z_j
$$
où les $Z_j \subset X^{reg}$ sont irréductibles de dimension $d$, distinctes
de $Z$, et
$$
\dim(Z_j \cap T_i) \leq \dim(Z \cap T_i)
$$
avec inégalité stricte si $Z$ ne rencontre pas $T_i$ proprement dans $X^{reg}$.
\end{remark}



\section{Anneaux de Chow}
\label{intersection-section-chow-rings}

\noindent
Soit $X$ une variété projective non singulière. Nous définissons le produit
d'intersection
$$
\CH_r(X) \times \CH_s(X) \longrightarrow \CH_{r + s - \dim(X)}(X),\quad
(\alpha, \beta) \longmapsto \alpha \cdot \beta
$$
comme suit. Soient $\alpha \in Z_r(X)$ et $\beta \in Z_s(X)$.
Si $\alpha$ et $\beta$ se coupent proprement, nous utilisons la
définition donnée à la section \ref{intersection-section-intersection-product}.
Sinon, choisissons $\alpha \sim_{rat} \alpha'$ comme dans le
lemme \ref{intersection-lemma-moving-move} et posons
$$
\alpha \cdot \beta =
\text{classe de }\alpha' \cdot \beta \in \CH_{r + s - \dim(X)}(X)
$$
Ceci est bien défini et passe au quotient par l'équivalence rationnelle d'après le
théorème \ref{intersection-theorem-well-defined}. Le produit d'intersection
sur $\CH_*(X)$ est commutatif (c'est clair), associatif
(lemme \ref{intersection-lemma-associative}) et possède une unité $[X] \in \CH_{\dim(X)}(X)$.

\medskip\noindent
Nous employons souvent $\CH^c(X) = \CH_{\dim X - c}(X)$ pour désigner le groupe de Chow des
cycles de codimension $c$, voir
Homologie de Chow, section \ref{chow-section-cycles-codimension}.
Le produit d'intersection définit un produit
$$
\CH^k(X) \times \CH^l(X) \longrightarrow \CH^{k + l}(X)
$$
qui est commutatif, associatif et possède une unité $1 = [X] \in \CH^0(X)$.


\section{Image inverse par un morphisme quelconque}
\label{intersection-section-general-pullback}

\noindent
Soit $f : X \to Y$ un morphisme de variétés projectives non singulières.
Nous définissons
$$
f^* : \CH_k(Y) \to \CH_{k+\dim X - \dim Y}(X)
$$
par la règle
$$
f^*(\alpha) = pr_{X, *}(\Gamma_f \cdot pr_Y^*(\alpha))
$$
où $\Gamma_f \subset X\times Y$ est le graphe de $f$. Notons qu'à ce
degré de généralité, cette application n'est définie que sur les classes de cycles, et non sur les cycles. Avec la
notation $\CH^*$ introduite à la section \ref{intersection-section-chow-rings},
on peut considérer l'image inverse comme une application
$$
f^* : \CH^*(Y) \to \CH^*(X)
$$
autrement dit, comme une application de groupes abéliens gradués.

\begin{lemma}
\label{intersection-lemma-pullback-and-intersection-product}
Soit $f : X \to Y$ un morphisme de variétés projectives non singulières.
L'application d'image inverse sur les groupes de Chow vérifie les propriétés suivantes :
\begin{enumerate}
\item $f^* : \CH^*(Y) \to \CH^*(X)$ est un morphisme d'anneaux,
\item $(g \circ f)^* = f^* \circ g^*$ pour un couple composable $f, g$,
\item la formule de projection vaut : $f_*(\alpha) \cdot \beta =
f_*( \alpha \cdot f^*\beta)$, et
\item si $f$ est plat, elle coïncide avec la définition précédente.
\end{enumerate}
\end{lemma}

\begin{proof}
Toutes ces assertions résultent immédiatement des résultats précédents.

\medskip\noindent
Pour (1), il suffit de montrer que
$\text{pr}_{X,*}( \Gamma_f \cdot \alpha \cdot \beta) =
\text{pr}_{X,*}(\Gamma_f \cdot \alpha) \cdot
\text{pr}_{X,*}(\Gamma_f \cdot \beta)$
pour des cycles $\alpha$, $\beta$ sur $X \times Y$. Si $\alpha$ est un cycle sur
$X \times Y$ qui rencontre $\Gamma_f$ proprement, alors il est facile
de voir que
$$
\Gamma_f \cdot \alpha =
\Gamma_f \cdot \text{pr}_X^*(\text{pr}_{X,*}(\Gamma_f \cdot \alpha))
$$
comme cycles, puisque $\Gamma_f$ est un graphe. On obtient ainsi la première
égalité dans
\begin{align*}
\text{pr}_{X,*}(\Gamma_f \cdot \alpha \cdot \beta)
& =
\text{pr}_{X,*}(
\Gamma_f \cdot
\text{pr}_X^*(\text{pr}_{X,*}(\Gamma_f \cdot \alpha)) \cdot \beta) \\
& =
\text{pr}_{X,*}(\text{pr}_X^*(\text{pr}_{X,*}(\Gamma_f \cdot \alpha))
\cdot (\Gamma_f \cdot \beta)) \\
& =
\text{pr}_{X,*}(\Gamma_f \cdot \alpha) \cdot
\text{pr}_{X,*}(\Gamma_f \cdot \beta)
\end{align*}
la dernière étape résultant de la formule de projection dans le cas plat
(lemme \ref{intersection-lemma-projection-formula-flat}).

\medskip\noindent
Si $g : Y \to Z$, la propriété (2) résulte formellement de l'observation suivante :
$$
\Gamma =
\text{pr}_{X \times Y}^*\Gamma_f \cdot
\text{pr}_{Y \times Z}^*\Gamma_g
$$
dans $Z_*(X \times Y \times Z)$, où $\Gamma = \{(x, f(x), g(f(x))\}$,
et s'envoie isomorphiquement sur $\Gamma_{g \circ f}$ dans $X \times Z$.
L'égalité résulte de l'égalité schématique et du
lemme \ref{intersection-lemma-transversal}.

\medskip\noindent
Pour (3), nous appliquons deux fois la formule de projection pour les applications plates :
\begin{align*}
f_*(\alpha \cdot pr_{X, *}(\Gamma_f \cdot pr_Y^*(\beta)))
& =
f_*(pr_{X, *}(pr_X^*\alpha \cdot \Gamma_f \cdot pr_Y^*(\beta))) \\
& =
pr_{Y, *}(pr_X^*\alpha \cdot \Gamma_f \cdot pr_Y^*(\beta))) \\
& =
pt_{Y, *}(pr_X^*\alpha \cdot \Gamma_f) \cdot \beta \\
& =
f_*(\alpha) \cdot \beta
\end{align*}
où, dans la dernière égalité, nous utilisons la remarque précédente sur les graphes.
Cela démontre (3).

\medskip\noindent
La propriété (4) repose sur l'identification du produit d'intersection
$\Gamma_f \cdot pr_Y^*\alpha$ lorsque $f$ est plat. Plus précisément, dans ce
cas, si $V \subset Y$ est une sous-variété fermée, tout point générique
$\xi$ du schéma $f^{-1}(V) \cong \Gamma_f \cap pr_Y^{-1}(V)$
est au-dessus du point générique de $V$. L'anneau local de
$pr_Y^{-1}(V) = X \times V$ en $\xi$ est donc de Cohen--Macaulay. Comme
$\Gamma_f \subset X \times Y$ est une immersion régulière (en tant que morphisme de
variétés projectives lisses), on obtient
$$
\Gamma_f \cdot pr_Y^*[V] = [\Gamma_f \cap pr_Y^{-1}(V)]_d
$$
où $d$ est la dimension de $\Gamma_f \cap pr_Y^{-1}(V)$, voir le
lemme \ref{intersection-lemma-multiplicity-lci-CM}. Comme $\Gamma_f \cap pr_Y^{-1}(V)$
s'envoie isomorphiquement sur $f^{-1}(V)$, on conclut.
\end{proof}


\section{Image inverse des cycles}
\label{intersection-section-pullback-cycles}

\noindent
Supposons que $X$ et $Y$ soient des variétés
projectives non singulières, et soit $f : X \to Y$ un morphisme.
Supposons que $Z \subset Y$ soit une sous-variété fermée. Notons $f^{-1}(Z)$
l'image inverse schématique, de sorte que le diagramme
$$
\xymatrix{
f^{-1}(Z) \ar[r] \ar[d] & Z \ar[d] \\
X \ar[r] & Y
}
$$
est cartésien. En particulier, $f^{-1}(Z) \subset X$
est un sous-schéma fermé de $X$. Dans ce cas, on a toujours
$$
\dim f^{-1}(Z) \geq \dim Z + \dim X - \dim Y.
$$
Si l'égalité a lieu dans la formule précédente, alors
$f^*[Z] = [f^{-1}(Z)]_{\dim Z + \dim X - \dim Y}$
pourvu que le schéma $Z$ soit de Cohen--Macaulay aux images
des points génériques de $f^{-1}(Z)$. Cela résulte de l'identification de
$f^{-1}(Z)$ avec l'intersection schématique de $\Gamma_f$
et $X \times Z$, puis du lemme \ref{intersection-lemma-multiplicity-lci-CM}.
Les détails sont analogues à la démonstration de la partie (4) du
lemme \ref{intersection-lemma-pullback-and-intersection-product} ci-dessus.









% OK, start here.
%

\chapter{Schémas de Picard des courbes}



\phantomsection
\label{pic-section-phantom}



\section{Introduction}
\label{pic-section-introduction}

\noindent
Dans ce chapitre, nous nous bornons à construire le schéma de Picard
d'une courbe projective non singulière sur un corps algébriquement clos.
On trouvera dans \cite{Kleiman-Picard} un exposé plus approfondi ainsi que
des indications historiques.

\medskip\noindent
Nous étudierons plus loin, dans le projet Champs, les foncteurs de Hilbert et de Quot
dans un cadre beaucoup plus général.


\section{Schéma de Hilbert des points}
\label{pic-section-hilbert-scheme-points}

\noindent
Soit $X \to S$ un morphisme de schémas. Soit $d \geq 0$ un entier.
Pour tout schéma $T$ sur $S$, posons
$$
\Hilbfunctor^d_{X/S}(T) =
\left\{
\begin{matrix}
Z \subset X_T\text{ sous-schéma fermé tel que }\\
Z \to T\text{ soit fini localement libre de degré }d
\end{matrix}
\right\}
$$
Si $T' \to T$ est un morphisme de schémas sur $S$ et si
$Z \in \Hilbfunctor^d_{X/S}(T)$, le changement de base
$Z_{T'} \subset X_{T'}$ est un élément de $\Hilbfunctor^d_{X/S}(T')$.
On obtient ainsi un foncteur
$$
\Hilbfunctor^d_{X/S} :
(\Sch/S)^{opp} \longrightarrow \textit{Ensembles},\quad
T \longrightarrow \Hilbfunctor^d_{X/S}(T)
$$
En général, $\Hilbfunctor^d_{X/S}$ est un espace algébrique
(insérer ici une référence ultérieure). Dans cette section, nous
montrerons que $\Hilbfunctor^d_{X/S}$ est représentable
par un schéma si tout ensemble fini de points d'une fibre de
$X \to S$ est contenu dans un ouvert affine.
Lorsque $\Hilbfunctor^d_{X/S}$ est représentable par un schéma, nous
notons souvent ce schéma $\underline{\Hilbfunctor}^d_{X/S}$.

\begin{lemma}
\label{pic-lemma-hilb-d-sheaf}
Soit $X \to S$ un morphisme de schémas. Le foncteur $\Hilbfunctor^d_{X/S}$
est un faisceau pour la topologie fpqc
(Topologies, définition \ref{topologies-definition-sheaf-property-fpqc}).
\end{lemma}

\begin{proof}
Soit $\{T_i \to T\}_{i \in I}$ un recouvrement fpqc de schémas sur $S$.
Posons $X_i = X_{T_i} = X \times_S T_i$.
Notons que $\{X_i \to X_T\}_{i \in I}$ est un recouvrement fpqc de
$X_T$ (Topologies, lemme \ref{topologies-lemma-fpqc})
et que $X_{T_i \times_T T_{i'}} = X_i \times_{X_T} X_{i'}$.
Soit $Z_i \in \Hilbfunctor^d_{X/S}(T_i)$ une famille
d'éléments telle que $Z_i$ et $Z_{i'}$ aient la même image dans
$\Hilbfunctor^d_{X/S}(T_i \times_T T_{i'})$. Par descente effective
des immersions fermées (Descente, lemme \ref{descent-lemma-closed-immersion}),
il existe une immersion fermée $Z \to X_T$ dont le changement de base par
$X_i \to X_T$ est $Z_i \to X_i$. Le morphisme $Z \to T$
a alors pour changement de base à $T_i$ le morphisme
$Z_i \to T_i$. Par conséquent, $Z \to T$ est fini localement libre de degré $d$
d'après Descente, lemme \ref{descent-lemma-descending-property-finite-locally-free}.
\end{proof}

\begin{lemma}
\label{pic-lemma-hilb-d-limit-preserving}
Soit $X \to S$ un morphisme de schémas. Si $X \to S$ est
de présentation finie, le foncteur $\Hilbfunctor^d_{X/S}$
commute aux limites projectives (Limites, remarque \ref{limits-remark-limit-preserving}).
\end{lemma}

\begin{proof}
Soit $T = \lim T_i$ une limite de schémas affines sur $S$. Il s'agit de montrer
que $\Hilbfunctor^d_{X/S}(T) = \colim \Hilbfunctor^d_{X/S}(T_i)$.
Si $Z \to X_T$ est un élément de $\Hilbfunctor^d_{X/S}(T)$,
le morphisme $Z \to T$ est de présentation finie. Ainsi, d'après
Limites, lemme \ref{limits-lemma-descend-finite-presentation},
il existe un indice $i$, un schéma $Z_i$ de présentation finie sur $T_i$,
et un morphisme $Z_i \to X_{T_i}$ sur $T_i$ dont le changement de base à $T$
est $Z \to X_T$. Appliquons Limites, lemme
\ref{limits-lemma-descend-closed-immersion-finite-presentation}
pour voir que l'on peut supposer que $Z_i \to X_{T_i}$ est une immersion fermée
quitte à augmenter $i$.
Appliquons Limites, lemme \ref{limits-lemma-descend-finite-locally-free}
pour voir que $Z_i \to T_i$ est fini localement libre de degré $d$
quitte à augmenter encore $i$.
Alors $Z_i \in \Hilbfunctor^d_{X/S}(T_i)$, comme voulu.
\end{proof}

\noindent
Soit $S$ un schéma. Soit $i : X \to Y$ une immersion fermée de schémas
sur $S$. On dispose alors d'une transformation de foncteurs
$$
\Hilbfunctor^d_{X/S} \longrightarrow \Hilbfunctor^d_{Y/S}
$$
qui associe à un élément $Z \in \Hilbfunctor^d_{X/S}(T)$
l'élément $i_T(Z) \subset Y_T$ de $\Hilbfunctor^d_{Y/S}$. Ici, $i_T : X_T \to Y_T$
est le changement de base de $i$.

\begin{lemma}
\label{pic-lemma-hilb-d-of-closed}
Soit $S$ un schéma. Soit $i : X \to Y$ une immersion fermée de schémas.
Si $\Hilbfunctor^d_{Y/S}$ est représentable par un schéma, il en est de même de
$\Hilbfunctor^d_{X/S}$, et le morphisme de schémas correspondant
$\underline{\Hilbfunctor}^d_{X/S} \to \underline{\Hilbfunctor}^d_{Y/S}$
est une immersion fermée.
\end{lemma}

\begin{proof}
Soit $T$ un schéma sur $S$ et soit $Z \in \Hilbfunctor^d_{Y/S}(T)$.
Montrons qu'il existe un sous-schéma fermé $T_X \subset T$ tel
qu'un morphisme de schémas $T' \to T$ se factorise par $T_X$ si
et seulement si $Z_{T'} \to Y_{T'}$ se factorise par $X_{T'}$.
En appliquant ceci à un schéma $T_{univ}$ représentant $\Hilbfunctor^d_{Y/S}$ et à
l'objet universel\footnote{Voir
Catégories, section \ref{categories-section-opposite}}
$Z_{univ} \in \Hilbfunctor^d_{Y/S}(T_{univ})$,
on obtient un sous-schéma fermé $T_{univ, X} \subset T_{univ}$ tel que
$Z_{univ, X} = Z_{univ} \times_{T_{univ}} T_{univ, X}$
soit un sous-schéma fermé de $X \times_S T_{univ, X}$, donc
définisse un élément de $\Hilbfunctor^d_{X/S}(T_{univ, X})$.
Un argument formel montre alors que $T_{univ, X}$ est un schéma
représentant $\Hilbfunctor^d_{X/S}$, d'objet universel $Z_{univ, X}$.

\medskip\noindent
Démontrons l'assertion. Considérons $Z' = X_T \times_{Y_T} Z$. Étant donné $T' \to T$,
on voit que $Z_{T'} \to Y_{T'}$ se factorise par $X_{T'}$ si et
seulement si $Z'_{T'} \to Z_{T'}$ est un isomorphisme. L'assertion résulte donc
du résultat très général
Compléments sur la platitude, lemme \ref{flat-lemma-Weil-restriction-closed-subschemes}.
Dans ce cas particulier, on peut toutefois la démontrer directement comme
suit : on se ramène d'abord au cas $T = \Spec(A)$ et $Z = \Spec(B)$.
En restreignant encore $T$, on peut supposer qu'il existe un isomorphisme
$\varphi : B \to A^{\oplus d}$ de $A$-modules. Alors $Z' = \Spec(B/J)$
pour un idéal $J \subset B$. Soit $g_\beta \in J$ une famille de
générateurs, et écrivons $\varphi(g_\beta) = (g_\beta^1, \ldots, g_\beta^d)$.
Il est alors clair que $T_X$ est donné par $\Spec(A/(g_\beta^j))$.
\end{proof}

\begin{lemma}
\label{pic-lemma-hilb-d-separated}
Soit $X \to S$ un morphisme de schémas. Si $X \to S$ est séparé et si
$\Hilbfunctor^d_{X/S}$ est représentable,
alors $\underline{\Hilbfunctor}^d_{X/S} \to S$ est séparé.
\end{lemma}

\begin{proof}
Dans cette démonstration, tous les produits sans indice sont pris sur $S$.
Posons $H = \underline{\Hilbfunctor}^d_{X/S}$ et soit
$Z \in \Hilbfunctor^d_{X/S}(H)$ l'objet universel.
Considérons les deux objets $Z_1, Z_2 \in \Hilbfunctor^d_{X/S}(H \times H)$
obtenus en prenant l'image réciproque de $Z$ par les deux projections $H \times H \to H$.
Alors $Z_1 = Z \times H \subset X_{H \times H}$ et $Z_2 = H \times Z
\subset X_{H \times H}$. Puisque $H$ représente le foncteur
$\Hilbfunctor^d_{X/S}$, le morphisme diagonal $\Delta : H \to H \times H$
a la propriété universelle suivante : un morphisme de schémas
$T \to H \times H$ se factorise par $\Delta$ si et seulement si
$Z_{1, T} = Z_{2, T}$ en tant qu'éléments de $\Hilbfunctor^d_{X/S}(T)$.
Posons $Z = Z_1 \times_{X_{H \times H}} Z_2$. On voit alors que
$T \to H \times H$ se factorise par $\Delta$ si et seulement si
les morphismes $Z_T \to Z_{1, T}$ et $Z_T \to Z_{2, T}$ sont des
isomorphismes. Le résultat très général
Compléments sur la platitude, lemme \ref{flat-lemma-Weil-restriction-closed-subschemes}
montre que $\Delta$ est une immersion fermée. Dans la démonstration du
lemme \ref{pic-lemma-hilb-d-of-closed},
le lecteur trouvera une autre démonstration, plus simple, du résultat nécessaire
dans notre cas particulier.
\end{proof}

\begin{lemma}
\label{pic-lemma-hilb-d-An}
Soit $X \to S$ un morphisme de schémas affines. Soit $d \geq 0$. Alors
$\Hilbfunctor^d_{X/S}$ est représentable.
\end{lemma}

\begin{proof}
Écrivons $S = \Spec(R)$. On peut choisir une immersion fermée de $X$
dans le spectre de $R[x_i; i \in I]$ pour un ensemble $I$ de cardinal
assez grand. Ainsi, d'après le lemme \ref{pic-lemma-hilb-d-of-closed},
on peut supposer que $X = \Spec(A)$, où $A = R[x_i; i \in I]$.
Nous utiliserons Schémas, lemme \ref{schemes-lemma-glue-functors} pour démontrer le
lemme dans ce cas.

\medskip\noindent
La condition (1) du lemme résulte du lemme \ref{pic-lemma-hilb-d-sheaf}.

\medskip\noindent
Pour toute partie $W \subset A$ de cardinal $d$, nous allons
construire un sous-foncteur $F_W$ de $\Hilbfunctor^d_{X/S}$.
(Il suffirait de considérer le cas où $W$ est formé de
monômes en les $x_i$, mais cela ne nous sera pas nécessaire.)
Nous dirons que $Z \in \Hilbfunctor^d_{X/S}(T)$ appartient à $F_W(T)$
si et seulement si l'application $\mathcal{O}_T$-linéaire
$$
\bigoplus\nolimits_{f \in W} \mathcal{O}_T
\longrightarrow
(Z \to T)_*\mathcal{O}_Z,\quad
(g_f) \longmapsto \sum g_f f|_Z
$$
est surjective (ou, ce qui revient au même, est un isomorphisme). Pour $f \in A$
et $Z \in \Hilbfunctor^d_{X/S}(T)$, on note ici $f|_Z$ l'image réciproque de $f$
par le morphisme $Z \to X_T \to X$.

\medskip\noindent
Ouverture, c'est-à-dire condition (2)(b) du lemme. Elle résulte de
Algèbre, lemme \ref{algebra-lemma-cokernel-flat}.

\medskip\noindent
Recouvrement, c'est-à-dire condition (2)(c) du lemme. Puisque
$$
A \otimes_R \mathcal{O}_T =
(X_T \to T)_*\mathcal{O}_{X_T} \to (Z \to T)_*\mathcal{O}_Z
$$
est surjective et que $(Z \to T)_*\mathcal{O}_Z$ est localement libre
de rang $d$, on peut, pour tout point $t \in T$, trouver une partie finie
$W \subset A$ de cardinal $d$ dont les images forment une base
de l'espace vectoriel de dimension $d$ sur $\kappa(t)$
$((Z \to T)_*\mathcal{O}_Z)_t \otimes_{\mathcal{O}_{T, t}} \kappa(t)$.
Le lemme de Nakayama donne un voisinage ouvert $V \subset T$
de $t$ tel que $Z_V \in F_W(V)$.

\medskip\noindent
Représentabilité, c'est-à-dire condition (2)(a) du lemme. Soit $W \subset A$
de cardinal $d$. Nous affirmons que $F_W$ est représentable par un schéma
affine sur $R$. Nous allons construire ce schéma affine, mais invitons
le lecteur à en chercher lui-même la construction. Numérotons
$f_1, \ldots, f_d$ les éléments de $W$. Nous construirons un élément
universel $Z_{univ} = \Spec(B_{univ})$ de $F_W$ sur $T_{univ} = \Spec(R_{univ})$
comme suit. Posons d'abord
$$
R_{univ} = R[c_{kl}^m, b^l, b_i^l]/\mathfrak a_{univ}
$$
pour un idéal $\mathfrak a_{univ}$ que nous décrirons ci-dessous.
Les indices $k, l, m, i$ parcourent tous ici $\{1, \ldots, d\}$.
Posons ensuite
$$
B_{univ} = R_{univ}e_1 \oplus R_{univ}e_2 \oplus \ldots \oplus R_{univ}e_d
$$
en tant que $R_{univ}$-module, muni de la structure d'algèbre donnée par
$$
e_ke_l = \sum c_{kl}^m e_m
$$
L'immersion fermée $Z_{univ} \to X_{T_{univ}}$ est donnée
par l'homomorphisme de $R_{univ}$-algèbres
$$
\Psi : A \otimes_R R_{univ} \longrightarrow B_{univ}
$$
envoyant $1 \otimes 1$ sur $\sum b^le_l$ et $x_i$ sur $\sum b_i^le_l$
(voir ci-dessous pour son existence). Quant à l'idéal $\mathfrak a_{univ}$,
il doit satisfaire aux prescriptions suivantes :
\begin{enumerate}
\item la multiplication dans $B_{univ}$ est commutative, c'est-à-dire que
l'on doit avoir $c_{lk}^m - c_{kl}^m \in \mathfrak a_{univ}$,
\item la multiplication dans $B_{univ}$ est associative, c'est-à-dire que
l'on doit avoir $c_{lk}^m c_{m n}^p - c_{lq}^p c_{kn}^q \in \mathfrak a_{univ}$,
\item $\sum b^le_l$ est un élément unité $1$ de $B_{univ}$,
autrement dit, on doit avoir $(\sum b^le_l)e_k = e_k$ pour tout $k$,
ce qui signifie que $\sum b^lc_{lk}^m - \delta_{km} \in \mathfrak a_{univ}$
(symbole de Kronecker).
\end{enumerate}
Puisque nous avons assuré que $B_{univ}$ est une
$R_{univ}$-algèbre commutative dont l'unité $1$ est $\sum b^le_l$, l'homomorphisme $\Psi$
existe. La dernière condition est la suivante :
\begin{enumerate}
\item[(4)] $f_l$ doit avoir pour image $e_l$ dans $B_{univ}$. Écrivons
$\Psi(f_l) - e_l \equiv \sum h_l^me_m$ pour des
$h_l^m \in R[c_{kl}^m, b^l, b_i^l]$ ;
on doit alors avoir $h_l^m \in \mathfrak a_{univ}$.
\end{enumerate}
En prenant pour $\mathfrak a_{univ} \subset R[c_{kl}^m, b^l, b_i^l]$ l'idéal engendré
par les éléments énumérés dans (1) -- (4),
il est clair que $F_W$ est représenté par $\Spec(R_{univ})$.
\end{proof}

\begin{proposition}
\label{pic-proposition-hilb-d-representable}
Soit $X \to S$ un morphisme de schémas. Soit $d \geq 0$. Supposons
que pour tout $(s, x_1, \ldots, x_d)$, où $s \in S$ et
$x_1, \ldots, x_d \in X_s$, il existe un ouvert affine $U \subset X$
tel que $x_1, \ldots, x_d \in U$. Alors $\Hilbfunctor^d_{X/S}$ est
représentable par un schéma.
\end{proposition}

\begin{proof}
En utilisant soit le recollement relatif (Constructions, section
\ref{constructions-section-relative-glueing}), soit
le point de vue fonctoriel
(Schémas, lemme \ref{schemes-lemma-glue-functors}),
on se ramène au cas où $S$ est affine. Nous omettons les détails.

\medskip\noindent
Supposons $S$ affine. Pour un ouvert affine $U \subset X$, notons
$F_U \subset \Hilbfunctor^d_{X/S}$ le sous-foncteur tel que,
pour tout schéma $T/S$, un élément $Z \in \Hilbfunctor^d_{X/S}(T)$
appartienne à $F_U(T)$ si et seulement si $Z \subset U_T$. Nous utiliserons
Schémas, lemme \ref{schemes-lemma-glue-functors}
et les sous-foncteurs $F_U$ pour conclure.

\medskip\noindent
La condition (1) est le lemme \ref{pic-lemma-hilb-d-sheaf}.

\medskip\noindent
La condition (2)(a) résulte de ce que $F_U = \Hilbfunctor^d_{U/S}$
et de sa représentabilité, donnée par le lemme \ref{pic-lemma-hilb-d-An}.
En effet, si $Z \in F_U(T)$, on peut regarder $Z$ comme un sous-schéma fermé
de $U_T$, fini localement libre de degré $d$ sur $T$, d'où
$Z \in \Hilbfunctor^d_{U/S}(T)$. Réciproquement, si $Z \in \Hilbfunctor^d_{U/S}(T)$,
alors $Z \to U_T \to X_T$ est une immersion fermée\footnote{C'est clair
si $X \to S$ est séparé, car dans ce cas Morphismes, lemme
\ref{morphisms-lemma-image-proper-scheme-closed}
montre que l'immersion $\varphi : Z \to X_T$ a une image fermée,
et est donc une immersion fermée d'après
Schémas, lemme \ref{schemes-lemma-immersion-when-closed}. Nous suggérons au
lecteur de passer la suite de cette note, car nous ne connaissons aucun exemple
où les hypothèses sur $X \to S$ soient satisfaites sans que $X \to S$ soit séparé.
Dans le cas général, soit $x \in X_T$ un point de l'adhérence de
$\varphi(Z)$. Il faut montrer que $x \in \varphi(Z)$. Soit $t \in T$ l'image
de $x$. L'hypothèse sur $X \to S$ permet de choisir un ouvert affine
$W \subset X_T$ contenant $x$ et $\varphi(Z_t)$. Alors $\varphi^{-1}(W)$
est un ouvert contenant toute la fibre $Z_t$ et, puisque $Z \to T$ est fermé,
on peut, en remplaçant $T$ par un voisinage ouvert de $t$, supposer que
$Z = \varphi^{-1}(W)$. Alors $\varphi(Z) \subset W$ est fermé d'après le
cas séparé (car $W \to T$ est séparé), d'où $x \in \varphi(Z)$.}
et l'on peut regarder $Z$ comme un élément de $F_U(T)$.

\medskip\noindent
Soit $Z \in \Hilbfunctor^d_{X/S}(T)$ pour un schéma $T$ sur $S$. Posons
$$
B = (Z \to T)\left((Z \to X_T \to X)^{-1}(X \setminus U)\right)
$$
C'est une partie fermée de $T$, et il est clair que sur l'ouvert
$T_{Z, U} = T \setminus B$, la restriction $Z_{T'}$ a son image dans $U_{T'}$.
D'autre part, pour tout $b \in B$, la fibre $Z_b$ n'a pas son image
dans $U$. Ainsi, étant donné un morphisme $T' \to T$, on
a $Z_{T'} \in F_U(T')$ $\Leftrightarrow$ $T' \to T$ se factorise par
l'ouvert $T_{Z, U}$. Cela démontre la condition (2)(b).

\medskip\noindent
La condition (2)(c) résulte de notre hypothèse sur $X/S$. Il suffit de
montrer ceci : si $T$ est le spectre d'un corps
et si $Z \subset X_T$ est un sous-schéma fermé, fini plat de degré
$d$ sur $T$, alors $Z \to X_T \to X$ se factorise par un ouvert affine
$U$ de $X$. C'est clair, car $Z$ a au plus $d$ points,
dont les images sont toutes dans la fibre de $X$ au-dessus du point image
de $T \to S$.
\end{proof}

\begin{remark}
\label{pic-remark-when-proposition-applies}
Soit $f : X \to S$ un morphisme de schémas. L'hypothèse de la
proposition \ref{pic-proposition-hilb-d-representable}, et par
suite sa conclusion, sont satisfaites dans chacun des cas suivants :
\begin{enumerate}
\item $X$ est quasi-affine,
\item $f$ est quasi-affine,
\item $f$ est quasi-projectif,
\item $f$ est localement projectif,
\item il existe un faisceau inversible ample sur $X$,
\item il existe un faisceau inversible $f$-ample sur $X$, et
\item il existe un faisceau inversible $f$-très ample sur $X$.
\end{enumerate}
En effet, dans chacun de ces cas, tout ensemble fini de points d'une
fibre $X_s$ est contenu dans un ouvert quasi-compact $U$ de $X$
qui est muni d'un faisceau inversible ample, est isomorphe
à un ouvert d'un schéma affine, ou est isomorphe à un ouvert
du $\text{Proj}$ d'un anneau gradué (dans chaque cas, cela résulte
directement des définitions). L'existence d'ouverts affines convenables
résulte donc de
Propriétés, lemme \ref{properties-lemma-ample-finite-set-in-affine}.
\end{remark}




\section{Espaces de modules de diviseurs sur les courbes lisses}
\label{pic-section-divisors}

\noindent
Pour un morphisme lisse $X \to S$ de dimension relative $1$, le foncteur
$\Hilbfunctor^d_{X/S}$ paramètre les diviseurs de Cartier effectifs relatifs
au sens de
Diviseurs, section \ref{divisors-section-effective-Cartier-morphisms}.

\begin{lemma}
\label{pic-lemma-divisors-on-curves}
Soit $X \to S$ un morphisme lisse de schémas de dimension relative $1$.
Soit $D \subset X$ un sous-schéma fermé. Considérons les conditions suivantes :
\begin{enumerate}
\item $D \to S$ est fini localement libre,
\item $D$ est un diviseur de Cartier effectif relatif sur $X/S$,
\item $D \to S$ est localement quasi-fini, plat et
localement de présentation finie, et
\item $D \to S$ est localement quasi-fini et plat.
\end{enumerate}
On a toujours les implications
$$
(1) \Rightarrow (2) \Leftrightarrow (3) \Rightarrow (4)
$$
Si $S$ est localement noethérien, la dernière flèche est une équivalence.
Si $X \to S$ est propre (et $S$ quelconque), la première flèche est
une équivalence.
\end{lemma}

\begin{proof}
Équivalence de (2) et (3). Elle résulte de
Diviseurs, lemme \ref{divisors-lemma-fibre-Cartier},
à condition de montrer l'équivalence de (2) et (3) lorsque
$S$ est le spectre d'un corps $k$. Soit $x \in X$ un point fermé.
Puisque $X$ est lisse de dimension relative $1$ sur $k$, on voit que
$\mathcal{O}_{X, x}$ est un anneau local régulier de dimension $1$
(voir Variétés, lemme \ref{varieties-lemma-smooth-regular}).
Ainsi, $\mathcal{O}_{X, x}$ est un anneau de valuation discrète
(Algèbre, lemme \ref{algebra-lemma-characterize-dvr}),
donc un anneau principal. Il en résulte que tout faisceau d'idéaux
$\mathcal{I} \subset \mathcal{O}_X$ non nul en tous
les points génériques de $X$ est inversible
(Diviseurs, lemme \ref{divisors-lemma-effective-Cartier-in-points}).
Autrement dit, tout sous-schéma fermé de $X$ ne contenant
aucun point générique est un diviseur de Cartier effectif.
Il s'ensuit que (2) et (3) sont équivalentes.

\medskip\noindent
Si $S$ est noethérien, tout morphisme localement quasi-fini
$D \to S$ est localement de présentation finie (Morphismes, lemme
\ref{morphisms-lemma-noetherian-finite-type-finite-presentation}),
d'où l'équivalence de (3) et (4).

\medskip\noindent
Si $X \to S$ est propre (et $S$ quelconque), alors $D \to S$ est
également propre. Un morphisme propre localement quasi-fini étant fini
(Compléments sur les morphismes, lemme \ref{more-morphisms-lemma-characterize-finite})
et un morphisme fini, plat et de présentation finie étant fini localement libre
(Morphismes, lemme \ref{morphisms-lemma-finite-flat}), on voit que
(1) et (2) sont équivalentes.
\end{proof}

\begin{lemma}
\label{pic-lemma-sum-divisors-on-curves}
Soit $X \to S$ un morphisme lisse de schémas de dimension relative $1$.
Soient $D_1, D_2 \subset X$ des sous-schémas fermés finis localement libres de
degrés $d_1$, $d_2$ sur $S$. Alors $D_1 + D_2$ est fini localement libre
de degré $d_1 + d_2$ sur $S$.
\end{lemma}

\begin{proof}
D'après le lemme \ref{pic-lemma-divisors-on-curves}, les sous-schémas $D_1$
et $D_2$ sont des diviseurs de Cartier effectifs relatifs sur $X/S$.
Ainsi, $D = D_1 + D_2$ est un diviseur de Cartier effectif relatif
sur $X/S$ d'après
Diviseurs, lemme \ref{divisors-lemma-sum-relative-effective-Cartier-divisor}.
Par conséquent, $D \to S$ est localement quasi-fini, plat et
localement de présentation finie d'après le
lemme \ref{pic-lemma-divisors-on-curves}.
En appliquant
Morphismes, lemme \ref{morphisms-lemma-image-universally-closed-separated}
au morphisme entier surjectif $D_1 \amalg D_2 \to D$,
on voit que $D \to S$ est séparé. Alors
Morphismes, lemme \ref{morphisms-lemma-image-proper-is-proper}
montre que $D \to S$ est propre.
Il en résulte que $D \to S$ est fini
(Compléments sur les morphismes, lemme \ref{more-morphisms-lemma-characterize-finite}),
puis que $D \to S$ est fini localement libre
(Morphismes, lemme \ref{morphisms-lemma-finite-flat}).
Il suffit donc de montrer que le degré de $D \to S$ est $d_1 + d_2$.
Pour cela, on peut effectuer un changement de base à une fibre de $X \to S$, donc
supposer que $S = \Spec(k)$ pour un corps $k$.
Dans ce cas, il existe un ensemble fini de points fermés
$x_1, \ldots, x_n \in X$ tel que les supports de $D_1$ et $D_2$
soient contenus dans $\{x_1, \ldots, x_n\}$.
Il existe en fait des éléments non diviseurs de zéro $f_{i, j} \in \mathcal{O}_{X, x_i}$
tels que
$$
D_1 = \coprod_{i = 1}^n \Spec(\mathcal{O}_{X, x_i}/(f_{i, 1}))
\quad\text{et}\quad
D_2 = \coprod_{i = 1}^n \Spec(\mathcal{O}_{X, x_i}/(f_{i, 2}))
$$
On voit alors que
$$
D = \coprod_{i = 1}^n \Spec(\mathcal{O}_{X, x_i}/(f_{i, 1}f_{i, 2}))
$$
On en déduit aisément que $D$ est de degré $d_1 + d_2$
sur $k$ (au besoin, utiliser Algèbre, lemme \ref{algebra-lemma-ord-additive}).
\end{proof}

\begin{lemma}
\label{pic-lemma-difference-divisors-on-curves}
Soit $X \to S$ un morphisme lisse de schémas de dimension relative $1$.
Soient $D_1, D_2 \subset X$ des sous-schémas fermés finis localement libres de
degrés $d_1$, $d_2$ sur $S$. Si $D_1 \subset D_2$ (comme sous-schémas fermés),
il existe un sous-schéma fermé $D \subset X$ fini localement libre de
degré $d_2 - d_1$ sur $S$ tel que $D_2 = D_1 + D$.
\end{lemma}

\begin{proof}
Cette démonstration est presque identique à celle du
lemme \ref{pic-lemma-sum-divisors-on-curves}.
D'après le lemme \ref{pic-lemma-divisors-on-curves}, les sous-schémas $D_1$
et $D_2$ sont des diviseurs de Cartier effectifs relatifs sur $X/S$.
D'après Diviseurs, lemme
\ref{divisors-lemma-difference-relative-effective-Cartier-divisor},
il existe un diviseur de Cartier effectif relatif $D \subset X$
tel que $D_2 = D_1 + D$. Ainsi, $D \to S$ est localement quasi-fini, plat et
localement de présentation finie d'après le
lemme \ref{pic-lemma-divisors-on-curves}.
Puisque $D$ est un sous-schéma fermé de $D_2$, on voit que
$D \to S$ est fini. Il s'ensuit que $D \to S$ est fini localement libre
(Morphismes, lemme \ref{morphisms-lemma-finite-flat}).
Il suffit donc de montrer que le degré de $D \to S$ est $d_2 - d_1$.
Cela résulte du lemme \ref{pic-lemma-sum-divisors-on-curves}.
\end{proof}

\noindent
Soit $X \to S$ un morphisme lisse de schémas de dimension relative $1$.
D'après le lemme \ref{pic-lemma-divisors-on-curves}, pour un schéma $T$ sur $S$ et
$D \in \Hilbfunctor^d_{X/S}(T)$, on peut regarder $D$ comme un diviseur de Cartier
effectif relatif sur $X_T/T$ tel que $D \to T$ soit fini
localement libre de degré $d$. Par conséquent, le
lemme \ref{pic-lemma-sum-divisors-on-curves} donne une transformation
de foncteurs
$$
\Hilbfunctor^{d_1}_{X/S} \times \Hilbfunctor^{d_2}_{X/S}
\longrightarrow
\Hilbfunctor^{d_1 + d_2}_{X/S},\quad
(D_1, D_2) \longmapsto D_1 + D_2
$$
Si $\Hilbfunctor^d_{X/S}$ est représentable pour tout degré $d$, cette
transformation de foncteurs correspond à un morphisme de schémas
$$
\underline{\Hilbfunctor}^{d_1}_{X/S}
\times_S
\underline{\Hilbfunctor}^{d_2}_{X/S}
\longrightarrow
\underline{\Hilbfunctor}^{d_1 + d_2}_{X/S}
$$
sur $S$. Notons que $\underline{\Hilbfunctor}^0_{X/S} = S$ et
$\underline{\Hilbfunctor}^1_{X/S} = X$.
Un cas particulier du morphisme précédent est le morphisme
$$
\underline{\Hilbfunctor}^d_{X/S} \times_S X
\longrightarrow
\underline{\Hilbfunctor}^{d + 1}_{X/S},\quad
(D, x) \longmapsto D + x
$$

\begin{lemma}
\label{pic-lemma-universal-object}
Soit $X \to S$ un morphisme lisse de schémas de dimension relative $1$
tel que les foncteurs $\Hilbfunctor^d_{X/S}$ soient représentables. Le morphisme
$\underline{\Hilbfunctor}^d_{X/S} \times_S X \to
\underline{\Hilbfunctor}^{d + 1}_{X/S}$
est fini localement libre de degré $d + 1$.
\end{lemma}

\begin{proof}
Soit $D_{univ} \subset X \times_S \underline{\Hilbfunctor}^{d + 1}_{X/S}$
l'objet universel. On a un diagramme commutatif
$$
\xymatrix{
\underline{\Hilbfunctor}^d_{X/S} \times_S X \ar[rr] \ar[rd] & &
D_{univ} \ar[ld] \ar@{^{(}->}[r] &
\underline{\Hilbfunctor}^{d + 1}_{X/S} \times_S X \\
& \underline{\Hilbfunctor}^{d + 1}_{X/S}
}
$$
où la flèche horizontale supérieure envoie $(D', x)$ sur $(D' + x, x)$.
Nous affirmons que ce morphisme est un isomorphisme,
ce qui démontre le lemme. En effet, étant donné un schéma $T$ sur $S$,
un $T$-point $\xi$ de $D_{univ}$ est donné par un couple $\xi = (D, x)$,
où $D \subset X_T$ est un sous-schéma fermé fini localement libre
de degré $d + 1$ sur $T$ et $x : T \to X$ est un morphisme dont
le graphe $x : T \to X_T$ se factorise par $D$. D'après le
lemme \ref{pic-lemma-difference-divisors-on-curves},
on peut alors écrire $D = D' + x$ pour un $D' \subset X_T$ fini localement
libre de degré $d$ sur $T$. L'application qui associe à $\xi = (D, x)$ le couple
$(D', x)$ est l'inverse cherché.
\end{proof}

\begin{lemma}
\label{pic-lemma-hilb-d-smooth}
Soit $X \to S$ un morphisme lisse de schémas de dimension relative $1$
tel que les foncteurs $\Hilbfunctor^d_{X/S}$ soient représentables. Les
schémas $\underline{\Hilbfunctor}^d_{X/S}$ sont lisses sur $S$ de
dimension relative $d$.
\end{lemma}

\begin{proof}
On a $\underline{\Hilbfunctor}^0_{X/S} = S$ et
$\underline{\Hilbfunctor}^1_{X/S} = X$ ; le résultat est donc vrai pour $d = 0, 1$.
Supposons le résultat établi pour $d$ ; alors
$\underline{\Hilbfunctor}^d_{X/S} \times_S X$ est lisse sur $S$
(Morphismes, lemme \ref{morphisms-lemma-base-change-smooth} et
\ref{morphisms-lemma-composition-smooth}). Puisque
$\underline{\Hilbfunctor}^d_{X/S} \times_S X \to
\underline{\Hilbfunctor}^{d + 1}_{X/S}$
est fini localement libre de degré $d + 1$ d'après le
lemme \ref{pic-lemma-universal-object},
le résultat découle de
Descente, lemme \ref{descent-lemma-smooth-permanence}.
Nous omettons la vérification de la valeur annoncée de la dimension
relative (on peut la faire en examinant les fibres, ou en
suivant les dimensions dans le raisonnement précédent).
\end{proof}

\noindent
Rassemblons les résultats obtenus jusqu'ici dans le cas
d'une courbe propre et lisse sur un corps.

\begin{proposition}
\label{pic-proposition-hilb-d}
Soit $X$ une courbe propre, lisse et géométriquement irréductible sur un corps $k$.
\begin{enumerate}
\item Les foncteurs $\Hilbfunctor^d_{X/k}$ sont représentables par des variétés
propres et lisses $\underline{\Hilbfunctor}^d_{X/k}$ de dimension
$d$ sur $k$.
\item Pour toute extension de corps $k'/k$, les points $k'$-rationnels
de $\underline{\Hilbfunctor}^d_{X/k}$ sont en correspondance biunivoque
avec les diviseurs de Cartier effectifs de degré $d$ sur $X_{k'}$.
\item Pour $d_1, d_2 \geq 0$, il existe un morphisme
$$
\underline{\Hilbfunctor}^{d_1}_{X/k}
\times_k
\underline{\Hilbfunctor}^{d_2}_{X/k}
\longrightarrow
\underline{\Hilbfunctor}^{d_1 + d_2}_{X/k}
$$
qui est fini localement libre de degré ${d_1 + d_2 \choose d_1}$.
\end{enumerate}
\end{proposition}

\begin{proof}
Les foncteurs $\Hilbfunctor^d_{X/k}$ sont représentables d'après la
proposition \ref{pic-proposition-hilb-d-representable}
(voir aussi la remarque \ref{pic-remark-when-proposition-applies})
et le fait que $X$ est projectif
(Variétés, lemme \ref{varieties-lemma-dim-1-proper-projective}).
Les schémas $\underline{\Hilbfunctor}^d_{X/k}$ sont séparés
sur $k$ d'après le lemme \ref{pic-lemma-hilb-d-separated}.
Les schémas $\underline{\Hilbfunctor}^d_{X/k}$ sont lisses
sur $k$ d'après le lemme \ref{pic-lemma-hilb-d-smooth}.
Partant de $X = \underline{\Hilbfunctor}^1_{X/k}$,
les morphismes du lemme \ref{pic-lemma-universal-object}
donnent par récurrence un morphisme
$$
X^d = X \times_k X \times_k \ldots \times_k X \longrightarrow
\underline{\Hilbfunctor}^d_{X/k},\quad
(x_1, \ldots, x_d) \longrightarrow x_1 + \ldots + x_d
$$
qui est fini localement libre de degré $d!$. Puisque $X$ est
propre sur $k$, il en est de même de $X^d$, donc
$\underline{\Hilbfunctor}^d_{X/k}$ est propre sur $k$ d'après
Morphismes, lemme \ref{morphisms-lemma-image-proper-is-proper}.
Puisque $X$ est géométriquement irréductible sur $k$, le produit
$X^d$ est irréductible
(Variétés, lemme \ref{varieties-lemma-bijection-irreducible-components}),
donc son image est irréductible (et même géométriquement irréductible).
Cela démontre (1). L'assertion (2) résulte des définitions. L'assertion (3) découle
du diagramme commutatif
$$
\xymatrix{
X^{d_1} \times_k X^{d_2} \ar[d] \ar@{=}[r] & X^{d_1 + d_2} \ar[d] \\
\underline{\Hilbfunctor}^{d_1}_{X/k}
\times_k
\underline{\Hilbfunctor}^{d_2}_{X/k}
\ar[r] &
\underline{\Hilbfunctor}^{d_1 + d_2}_{X/k}
}
$$
et de la multiplicativité des degrés des morphismes finis localement libres.
\end{proof}

\begin{remark}
\label{pic-remark-universal-object-hilb-d}
Soit $X$ une courbe propre, lisse et géométriquement irréductible sur un corps $k$,
comme dans la proposition \ref{pic-proposition-hilb-d}. Soit $d \geq 0$. Le sous-schéma
fermé universel est un diviseur effectif relatif
$$
D_{univ} \subset \underline{\Hilbfunctor}^{d + 1}_{X/k} \times_k X
$$
sur $\underline{\Hilbfunctor}^{d + 1}_{X/k}$ d'après le
lemme \ref{pic-lemma-divisors-on-curves}.
En fait, $D_{univ}$ est isomorphe, en tant que schéma, à
$\underline{\Hilbfunctor}^d_{X/k} \times_k X$, voir la démonstration du
lemme \ref{pic-lemma-universal-object}.
En particulier, $D_{univ}$ est un diviseur de Cartier effectif et
l'on obtient un module inversible
$\mathcal{O}(D_{univ})$. Si $[D] \in \underline{\Hilbfunctor}^{d + 1}_{X/k}$
désigne le point $k$-rationnel correspondant au diviseur de Cartier
effectif $D \subset X$ de degré $d + 1$, la restriction
de $\mathcal{O}(D_{univ})$ à la fibre $[D] \times X$ est
$\mathcal{O}_X(D)$.
\end{remark}


\section{Le foncteur de Picard}
\label{pic-section-picard-functor}

\noindent
Pour tout schéma $X$, on note $\Pic(X)$ l'ensemble des classes
d'isomorphisme de $\mathcal{O}_X$-modules inversibles.
Voir Modules, définition \ref{modules-definition-pic}.
Étant donné un morphisme $f : X \to Y$ de schémas, l'image réciproque définit
un homomorphisme de groupes $\Pic(Y) \to \Pic(X)$.
L'association
$X \leadsto \Pic(X)$ est un foncteur contravariant de la catégorie
des schémas dans celle des groupes abéliens. Ce foncteur n'est pas
représentable, mais une variante relative de cette
construction est parfois représentable.

\medskip\noindent
Définissons le foncteur de Picard d'un morphisme de schémas $f : X \to S$.
L'idée de la construction consiste à prendre pour ce foncteur le faisceau
$R^1f_*\mathbf{G}_m$, l'image directe supérieure étant calculée pour
la topologie fppf. En explicitant les définitions, on obtient la définition
plus directe suivante.

\begin{definition}
\label{pic-definition-picard-functor}
Soit $\Sch_{fppf}$ un grand site comme dans
Topologies, définition \ref{topologies-definition-big-small-fppf}.
Soit $f : X \to S$ un morphisme de ce site. Le {\it foncteur de Picard}
$\Picardfunctor_{X/S}$ est le faisceau fppf associé au foncteur
$$
(\Sch/S)_{fppf} \longrightarrow \textit{Ensembles},\quad
T \longmapsto \Pic(X_T)
$$
Si ce foncteur est représentable, on note
$\underline{\Picardfunctor}_{X/S}$ un schéma qui le représente.
\end{definition}

\noindent
Une remarque d'usage fréquent est que, si $T \in \Ob((\Sch/S)_{fppf})$, alors
$\Picardfunctor_{X_T/T}$ est la restriction de $\Picardfunctor_{X/S}$ à
$(\Sch/T)_{fppf}$.
Déterminer la valeur de $\Picardfunctor_{X/S}$ sur les schémas
$T$ sur $S$ n'est pas immédiat. Le lemme suivant facilite cette
tâche.

\begin{lemma}
\label{pic-lemma-flat-geometrically-connected-fibres}
Soit $f : X \to S$ comme dans la définition \ref{pic-definition-picard-functor}.
Si $\mathcal{O}_T \to f_{T, *}\mathcal{O}_{X_T}$ est un isomorphisme
pour tout $T \in \Ob((\Sch/S)_{fppf})$, alors
$$
0 \to \Pic(T) \to \Pic(X_T) \to \Picardfunctor_{X/S}(T)
$$
est une suite exacte pour tout $T$.
\end{lemma}

\begin{proof}
On peut remplacer $S$ par $T$ et $X$ par $X_T$, donc supposer $S = T$
pour simplifier les notations. Soit $\mathcal{N}$ un
$\mathcal{O}_S$-module inversible. Si $f^*\mathcal{N} \cong \mathcal{O}_X$, alors
$f_*f^*\mathcal{N} \cong f_*\mathcal{O}_X \cong \mathcal{O}_S$
par hypothèse. Puisque $\mathcal{N}$ est localement trivial, on voit que
l'application canonique $\mathcal{N} \to f_*f^*\mathcal{N}$ est localement
un isomorphisme (car $\mathcal{O}_S \to f_*f^*\mathcal{O}_S$
est un isomorphisme par hypothèse). On en conclut que
$\mathcal{N} \to f_*f^*\mathcal{N} \to \mathcal{O}_S$ est un isomorphisme,
donc que $\mathcal{N}$ est trivial. Cela démontre que la première flèche
est injective.

\medskip\noindent
Soit $\mathcal{L}$ un $\mathcal{O}_X$-module inversible dont la classe appartient
au noyau de $\Pic(X) \to \Picardfunctor_{X/S}(S)$. Il existe alors
un recouvrement fppf $\{S_i \to S\}$ tel que l'image réciproque de $\mathcal{L}$
soit le faisceau inversible trivial sur $X_{S_i}$. Choisissons une section
trivialisante $s_i$. Alors $\text{pr}_0^*s_i$ et $\text{pr}_1^*s_j$ sont deux
sections trivialisantes de $\mathcal{L}$ sur $X_{S_i \times_S S_j}$
et diffèrent donc par multiplication par une unité
$$
f_{ij} \in
\Gamma(X_{S_i \times_S S_j}, \mathcal{O}_{X_{S_i \times_S S_j}}^*) =
\Gamma(S_i \times_S S_j, \mathcal{O}_{S_i \times_S S_j}^*)
$$
(l'égalité résulte de notre hypothèse sur l'image directe des faisceaux structuraux).
Ces éléments satisfont évidemment à la condition de cocycle sur
$S_i \times_S S_j \times_S S_k$, et définissent donc une donnée de descente
de faisceaux inversibles pour le recouvrement fppf $\{S_i \to S\}$.
D'après Descente, proposition \ref{descent-proposition-fpqc-descent-quasi-coherent},
il existe un $\mathcal{O}_S$-module inversible $\mathcal{N}$
muni de trivialisations sur les $S_i$ dont la donnée de descente associée est
$\{f_{ij}\}$. Alors $f^*\mathcal{N} \cong \mathcal{L}$, car le
foncteur des données de descente vers les modules est pleinement fidèle (voir la proposition
citée ci-dessus).
\end{proof}

\begin{lemma}
\label{pic-lemma-flat-geometrically-connected-fibres-with-section}
Soit $f : X \to S$ comme dans la définition \ref{pic-definition-picard-functor}.
Supposons que $f$ admette une section $\sigma$ et que
$\mathcal{O}_T \to f_{T, *}\mathcal{O}_{X_T}$ soit un isomorphisme
pour tout $T \in \Ob((\Sch/S)_{fppf})$. Alors
$$
0 \to \Pic(T) \to \Pic(X_T) \to \Picardfunctor_{X/S}(T) \to 0
$$
est une suite exacte scindée, le scindage étant donné par
$\sigma_T^* : \Pic(X_T) \to \Pic(T)$.
\end{lemma}

\begin{proof}
Notons $K(T) = \Ker(\sigma_T^* : \Pic(X_T) \to \Pic(T))$.
Puisque $\sigma$ est une section de $f$, le groupe $\Pic(X_T)$ est la somme
directe de $\Pic(T)$ et de $K(T)$.
Le lemme \ref{pic-lemma-flat-geometrically-connected-fibres} montre donc que
$K(T) \subset \Picardfunctor_{X/S}(T)$ pour tout $T$. De plus, la construction
montre clairement que $\Picardfunctor_{X/S}$ est le faisceau associé
au préfaisceau $K$. Pour achever la démonstration, il suffit de montrer que
$K$ satisfait à la condition de faisceau pour les recouvrements fppf, ce que nous faisons
dans le paragraphe suivant.

\medskip\noindent
Soit $\{T_i \to T\}$ un recouvrement fppf. Soient $\mathcal{L}_i$ des
éléments de $K(T_i)$ dont les images dans $K(T_i \times_T T_j)$ sont égales
pour tous $i$ et $j$. Choisissons un isomorphisme
$\alpha_i : \mathcal{O}_{T_i} \to \sigma_{T_i}^*\mathcal{L}_i$
pour tout $i$. Choisissons un isomorphisme
$$
\varphi_{ij} :
\mathcal{L}_i|_{X_{T_i \times_T T_j}}
\longrightarrow
\mathcal{L}_j|_{X_{T_i \times_T T_j}}
$$
Si l'application
$$
\alpha_j^{-1}|_{T_i \times_T T_j} \circ
\sigma_{T_i \times_T T_j}^*\varphi_{ij} \circ
\alpha_i|_{T_i \times_T T_j} :
\mathcal{O}_{T_i \times_T T_j} \to \mathcal{O}_{T_i \times_T T_j}
$$
n'est pas la multiplication par $1$, mais par un élément $u_{ij}$, on peut multiplier
$\varphi_{ij}$ par $u_{ij}^{-1}$ pour y remédier. Ceci fait, considérons
l'endomorphisme
$$
\varphi_{ki}|_{X_{T_i \times_T T_j \times_T T_k}} \circ
\varphi_{jk}|_{X_{T_i \times_T T_j \times_T T_k}} \circ
\varphi_{ij}|_{X_{T_i \times_T T_j \times_T T_k}}
\quad\text{de}\quad
\mathcal{L}_i|_{X_{T_i \times_T T_j \times_T T_k}}
$$
qui est la multiplication par une fonction régulière $f_{ijk}$
sur le schéma $X_{T_i \times_T T_j \times_T T_k}$.
Par notre choix de $\varphi_{ij}$, l'image réciproque de
cette application par $\sigma$ est la multiplication par $1$. D'après
notre hypothèse sur les fonctions sur $X$, on voit que $f_{ijk} = 1$.
On obtient ainsi une donnée de descente pour le recouvrement fppf
$\{X_{T_i} \to X_T\}$. D'après
Descente, proposition \ref{descent-proposition-fpqc-descent-quasi-coherent},
il existe un $\mathcal{O}_{X_T}$-module inversible $\mathcal{L}$
et un isomorphisme $\alpha : \mathcal{O}_T \to \sigma_T^*\mathcal{L}$
dont l'image réciproque sur $X_{T_i}$ redonne $(\mathcal{L}_i, \alpha_i)$
(nous omettons un petit détail). Ainsi, $\mathcal{L}$ définit un objet
de $K(T)$, comme voulu.
\end{proof}



\section{Un critère de représentabilité}
\label{pic-section-representability}

\noindent
Pour démontrer que le foncteur de Picard est représentable, nous utiliserons le
critère suivant.

\begin{lemma}
\label{pic-lemma-criterion}
Soit $k$ un corps. Soit $G : (\Sch/k)^{opp} \to \textit{Groupes}$ un
foncteur. Avec la terminologie de
Schémas, définition \ref{schemes-definition-representable-by-open-immersions},
supposons que
\begin{enumerate}
\item $G$ soit un faisceau pour la topologie de Zariski,
\item il existe un sous-foncteur $F \subset G$ tel que
\begin{enumerate}
\item $F$ soit représentable,
\item $F \subset G$ soit représentable par une immersion ouverte,
\item pour toute extension de corps $K$ de $k$ et tout $g \in G(K)$,
il existe un $g' \in G(k)$ tel que $g'g \in F(K)$.
\end{enumerate}
\end{enumerate}
Alors $G$ est représentable par un schéma en groupes sur $k$.
\end{lemma}

\begin{proof}
Cela résulte de Schémas, lemme \ref{schemes-lemma-glue-functors}.
En effet, prenons $I = G(k)$ et, pour $i = g' \in I$, prenons pour $F_i \subset G$
le sous-foncteur qui associe à $T$ sur $k$ l'ensemble des éléments
$g \in G(T)$ tels que $g'g \in F(T)$. Alors $F_i \cong F$ par multiplication
par $g'$. L'application $F_i \to G$ est isomorphe à l'application $F \to G$
par multiplication par $g'$, donc est représentable par des immersions ouvertes.
Enfin, la famille $(F_i)_{i \in I}$ recouvre $G$ par l'hypothèse (2)(c).
Le lemme cité ci-dessus s'applique donc, ce qui achève la démonstration.
\end{proof}



\section{Le schéma de Picard d'une courbe}
\label{pic-section-picard-curve}

\noindent
Dans cette section, nous appliquerons le lemme \ref{pic-lemma-criterion} pour montrer que
$\Picardfunctor_{X/k}$ est représentable lorsque $k$ est un corps algébriquement
clos et $X$ une courbe projective lisse sur $k$. Pour cela,
nous utiliserons quelques résultats de cohomologie et de changement de base développés dans
le chapitre sur les catégories dérivées des schémas.

\begin{lemma}
\label{pic-lemma-check-conditions}
Soit $k$ un corps. Soit $X$ une courbe projective lisse sur $k$
possédant un point $k$-rationnel. Alors les hypothèses du
lemme \ref{pic-lemma-flat-geometrically-connected-fibres-with-section}
sont satisfaites.
\end{lemma}

\begin{proof}
L'expression « possède un point $k$-rationnel » signifie exactement que
le morphisme structural $f : X \to \Spec(k)$ admet une section, ce qui
vérifie la première condition.
D'après Variétés, lemme \ref{varieties-lemma-regular-functions-proper-variety},
le corps $k' = H^0(X, \mathcal{O}_X)$ est une extension de $k$.
Puisque $X$ possède un point $k$-rationnel, il existe un homomorphisme de $k$-algèbres
$k' \to k$, d'où $k' = k$.
Puisque $k$ est un corps, tout morphisme $T \to \Spec(k)$ est plat.
Par cohomologie et changement de base
(Cohomologie des schémas, lemme \ref{coherent-lemma-flat-base-change-cohomology}),
on voit donc que $\mathcal{O}_T \to f_{T, *}\mathcal{O}_{X_T}$ est un isomorphisme.
Cela achève la démonstration.
\end{proof}

\noindent
Soit $X$ une courbe projective lisse sur un corps $k$, munie d'un
point $k$-rationnel $\sigma$. Alors le foncteur
$$
\Picardfunctor_{X/k, \sigma} : (\Sch/k)^{opp} \longrightarrow \textit{Ab},\quad
T \longmapsto \Ker(\Pic(X_T) \xrightarrow{\sigma_T^*} \Pic(T))
$$
est isomorphe à $\Picardfunctor_{X/k}$ sur $(\Sch/k)_{fppf}$
d'après les lemmes \ref{pic-lemma-check-conditions} et
\ref{pic-lemma-flat-geometrically-connected-fibres-with-section}.
Il suffira donc de démontrer que $\Picardfunctor_{X/k, \sigma}$
est représentable. Nous emploierons la notation
« $\mathcal{L} \in \Picardfunctor_{X/k, \sigma}(T)$ » pour signifier que
$T$ est un schéma sur $k$ et que $\mathcal{L}$ est un
$\mathcal{O}_{X_T}$-module inversible dont la restriction à $T$ par $\sigma_T$
est isomorphe à $\mathcal{O}_T$.

\begin{lemma}
\label{pic-lemma-define-open}
Soit $k$ un corps. Soit $X$ une courbe projective lisse sur $k$
munie d'un point $k$-rationnel $\sigma$. Pour tout schéma $T$ sur $k$,
considérons la partie $F(T) \subset \Picardfunctor_{X/k, \sigma}(T)$ formée des
$\mathcal{L}$ tels que $Rf_{T, *}\mathcal{L}$ soit isomorphe à un
$\mathcal{O}_T$-module inversible placé en degré $0$. Alors
$F \subset \Picardfunctor_{X/k, \sigma}$ est un sous-foncteur, et l'inclusion est
représentable par des immersions ouvertes.
\end{lemma}

\begin{proof}
Cela résulte immédiatement de Catégories dérivées des schémas, lemme
\ref{perfect-lemma-open-where-cohomology-in-degree-i-rank-r-geometric},
appliqué avec $i = 0$ et $r = 1$, et de
Schémas, définition \ref{schemes-definition-representable-by-open-immersions}.
\end{proof}

\noindent
Pour poursuivre, il convient de poser la définition suivante.

\begin{definition}
\label{pic-definition-genus}
Soit $k$ un corps. Soit $X$ une courbe projective lisse géométriquement irréductible
sur $k$. Le {\it genre} de $X$ est $g = \dim_k H^1(X, \mathcal{O}_X)$.
\end{definition}

\begin{lemma}
\label{pic-lemma-open-representable}
Soit $k$ un corps. Soit $X$ une courbe projective lisse de genre $g$
sur $k$, munie d'un point $k$-rationnel $\sigma$. Le sous-foncteur ouvert $F$ défini
au lemme \ref{pic-lemma-define-open} est représentable par un sous-schéma ouvert de
$\underline{\Hilbfunctor}^g_{X/k}$.
\end{lemma}

\begin{proof}
Dans cette démonstration, les produits sans indice sont pris sur $\Spec(k)$.
D'après la proposition \ref{pic-proposition-hilb-d}, le schéma
$H = \underline{\Hilbfunctor}^g_{X/k}$ existe.
Considérons le diviseur universel $D_{univ} \subset H \times X$
et le faisceau inversible associé $\mathcal{O}(D_{univ})$, voir
la remarque \ref{pic-remark-universal-object-hilb-d}.
Nous le rigidifions par produit tensoriel à l'aide de l'image réciproque suivant
$\sigma_H : H \to H \times X$, pour obtenir
$$
\mathcal{L}_H =
\mathcal{O}(D_{univ})
\otimes_{\mathcal{O}_{H \times X}}
\text{pr}_H^*\sigma_H^*\mathcal{O}(D_{univ})^{\otimes -1}
\in
\Picardfunctor_{X/k, \sigma}(H)
$$
Par le lemme de Yoneda (Catégories, lemme \ref{categories-lemma-yoneda}),
le faisceau inversible $\mathcal{L}_H$ définit une transformation naturelle
$$
h_H \longrightarrow \Picardfunctor_{X/k, \sigma}
$$
Puisque $F$ est un sous-foncteur ouvert, il existe un plus grand ouvert
$W \subset H$ tel que $\mathcal{L}_H|_{W \times X}$ appartienne à
$F(W)$. Cet ouvert n'est autre que le
sous-schéma ouvert construit dans
Catégories dérivées des schémas, lemme
\ref{perfect-lemma-open-where-cohomology-in-degree-i-rank-r-geometric},
avec $i = 0$ et $r = 1$, pour le morphisme $H \times X \to H$ et le faisceau
$\mathcal{F} = \mathcal{O}(D_{univ})$. En appliquant à nouveau le lemme de Yoneda,
on obtient un diagramme commutatif
$$
\xymatrix{
h_W \ar[d] \ar[r] & F \ar[d] \\
h_H \ar[r] & \Picardfunctor_{X/k, \sigma}
}
$$
Pour achever la démonstration, nous montrerons que la flèche horizontale supérieure est un
isomorphisme.

\medskip\noindent
Soit $\mathcal{L} \in F(T) \subset \Picardfunctor_{X/k, \sigma}(T)$.
Soit $\mathcal{N}$ le $\mathcal{O}_T$-module inversible
tel que $Rf_{T, *}\mathcal{L} \cong \mathcal{N}[0]$.
L'application d'adjonction
$$
f_T^*\mathcal{N} \longrightarrow \mathcal{L}
\quad\text{correspond à une section }s\text{ de}\quad
\mathcal{L} \otimes f_T^*\mathcal{N}^{\otimes -1}
$$
sur $X_T$. Assertion : le schéma des zéros de $s$ est un diviseur de Cartier effectif
relatif $D$ sur $(T \times X)/T$, fini localement libre de degré $g$ sur $T$.

\medskip\noindent
Achevons la démonstration du lemme en admettant cette assertion.
Le diviseur $D$ définit un morphisme $m : T \to H$ tel que $D$ soit l'image réciproque de
$D_{univ}$. Alors
$$
(m \times \text{id}_X)^*\mathcal{O}(D_{univ}) \cong
\mathcal{O}_{T \times X}(D)
$$
Puisque $\mathcal{O}_{T \times X}(D) \cong
\mathcal{L} \otimes f_T^*\mathcal{N}^{\otimes -1}$,
on en conclut que $(m \times \text{id}_X)^*\mathcal{L}_H$ et $\mathcal{L}$
diffèrent par tensorisation avec l'image réciproque d'un module inversible sur $T$.
Cela montre que $m : T \to H$ se factorise par l'ouvert $W \subset H$ ci-dessus.
De plus, il s'ensuit que ces modules inversibles définissent, après rigidification
par image réciproque suivant $\sigma_T$ comme ci-dessus, le même élément de
$\Picardfunctor_{X/k, \sigma}(T)$. Une chasse au diagramme utilisant le lemme de Yoneda
montre que $m \in h_W(T)$ a pour image $\mathcal{L} \in F(T)$. Nous omettons
la vérification du fait que la règle $F(T) \to h_W(T)$,
$\mathcal{L} \mapsto m$, définit un inverse de la transformation
de foncteurs ci-dessus.

\medskip\noindent
Démonstration de l'assertion. Puisque $D$ est un sous-schéma fermé localement principal
de $T \times X$, il suffit de montrer que les fibres de $D$ sur $T$ sont des
diviseurs de Cartier effectifs, voir le lemme \ref{pic-lemma-divisors-on-curves} et
Diviseurs, lemme \ref{divisors-lemma-fibre-Cartier}. Puisque le calcul
de la cohomologie de $\mathcal{L}$ commute au changement de base
(Catégories dérivées des schémas, lemme
\ref{perfect-lemma-flat-proper-perfect-direct-image-general}),
on se ramène à $T = \Spec(K)$, où $K/k$ est une extension de corps.
Alors $\mathcal{L}$ est un faisceau inversible sur $X_K$ tel que
$H^0(X_K, \mathcal{L}) = K$ et $H^1(X_K, \mathcal{L}) = 0$. Ainsi,
$$
\deg(\mathcal{L}) = \chi(X_K, \mathcal{L}) - \chi(X_K, \mathcal{O}_{X_K})
= 1 - (1 - g) = g
$$
Voir Variétés, définition \ref{varieties-definition-degree-invertible-sheaf}.
Pour achever la démonstration, il faut montrer qu'une section non nulle de $\mathcal{L}$
définit un diviseur de Cartier effectif sur $X_K$.
C'est clair.
\end{proof}

\begin{lemma}
\label{pic-lemma-twist-with-general-divisor}
Soit $k$ un corps séparablement clos. Soit $X$ une courbe projective lisse
de genre $g$ sur $k$. Soit $K/k$ une extension de corps et soit
$\mathcal{L}$ un faisceau inversible sur $X_K$. Il existe alors un
faisceau inversible $\mathcal{L}_0$ sur $X$ tel que
$\dim_K H^0(X_K,
\mathcal{L} \otimes_{\mathcal{O}_{X_K}} \mathcal{L}_0|_{X_K}) = 1$ et
$\dim_K H^1(X_K,
\mathcal{L} \otimes_{\mathcal{O}_{X_K}} \mathcal{L}_0|_{X_K}) = 0$.
\end{lemma}

\begin{proof}
Cette démonstration est une variante de celle de
Variétés, lemme \ref{varieties-lemma-general-degree-g-line-bundle}.
Nous invitons le lecteur à lire d'abord cette dernière.

\medskip\noindent
Choisissons d'abord un faisceau inversible ample $\mathcal{L}_0$ et
remplaçons $\mathcal{L}$ par
$\mathcal{L} \otimes_{\mathcal{O}_{X_K}} \mathcal{L}_0^{\otimes n}|_{X_K}$
pour un $n \gg 0$. Nous pourrons ainsi supposer que
$H^0(X_K, \mathcal{L}) \not = 0$ et $H^1(X_K, \mathcal{L}) = 0$.
En effet, l'annulation résulte de Cohomologie des schémas, lemme
\ref{coherent-lemma-vanshing-gives-ample}, et la non-annulation de ce que
le degré du produit tensoriel est $\gg 0$.
Nous achèverons la démonstration par récurrence descendante sur
$t = \dim_K H^0(X_K, \mathcal{L})$. Le cas initial $t = 1$ est trivial.
Supposons $t > 1$.

\medskip\noindent
Notons que, pour un point $k$-rationnel $x$ de $X$, le point $x_K$
obtenu par changement de base est $K$-rationnel dans $X_K$. De plus, il existe une infinité de
points $k$-rationnels d'après Variétés, lemme
\ref{varieties-lemma-smooth-separable-closed-points-dense}. Par conséquent,
les points $x_K$ forment une famille Zariski-dense de points de $X_K$.

\medskip\noindent
Soit $s \in H^0(X_K, \mathcal{L})$ non nul. Le paragraphe précédent
montre qu'il existe un point $k$-rationnel
$x$ tel que $s$ ne s'annule pas en $x_K$. Soit $\mathcal{I}$
le faisceau d'idéaux de $i : x_K \to X_K$ comme dans
Variétés, lemme \ref{varieties-lemma-regular-point-on-curve}. Considérons la
suite exacte courte
$$
0 \to \mathcal{I} \otimes_{\mathcal{O}_{X_K}} \mathcal{L} \to
\mathcal{L} \to i_*i^*\mathcal{L} \to 0
$$
Notons que $H^0(X_K, i_*i^*\mathcal{L}) = H^0(x_K, i^*\mathcal{L})$
est de dimension $1$ sur $K$. Puisque $s$ ne s'annule pas en $x_K$, on en conclut que
$$
H^0(X_K, \mathcal{L}) \longrightarrow H^0(X_K, i_*i^*\mathcal{L})
$$
est surjective. Ainsi,
$\dim_K H^0(X_K, \mathcal{I} \otimes_{\mathcal{O}_{X_K}} \mathcal{L}) = t - 1$.
Enfin, la suite exacte longue de cohomologie montre aussi que
$H^1(X_K, \mathcal{I} \otimes_{\mathcal{O}_{X_K}} \mathcal{L}) = 0$,
ce qui achève l'étape de récurrence.
\end{proof}

\begin{proposition}
\label{pic-proposition-pic-curve}
Soit $k$ un corps séparablement clos. Soit $X$ une courbe projective lisse
sur $k$. Le foncteur de Picard $\Picardfunctor_{X/k}$ est représentable.
\end{proposition}

\begin{proof}
Puisque $k$ est séparablement clos, il existe un point $k$-rationnel $\sigma$
de $X$, voir
Variétés, lemme \ref{varieties-lemma-smooth-separable-closed-points-dense}.
Comme on l'a vu ci-dessus, il suffit de montrer que le foncteur
$\Picardfunctor_{X/k, \sigma}$ classifiant les modules inversibles triviaux le long de
$\sigma$ est représentable. Pour cela, nous vérifierons les conditions (1),
(2)(a), (2)(b) et (2)(c) du
lemme \ref{pic-lemma-criterion}.

\medskip\noindent
Le foncteur $\Picardfunctor_{X/k, \sigma}$ satisfait à la condition de faisceau
pour la topologie fppf, car il est isomorphe à $\Picardfunctor_{X/k}$.
Il serait plus exact de dire que nous avons démontré la condition de faisceau
pour $\Picardfunctor_{X/k, \sigma}$ dans la démonstration du
lemme \ref{pic-lemma-flat-geometrically-connected-fibres-with-section},
qui s'applique d'après le lemme \ref{pic-lemma-check-conditions}.
Cela démontre la condition (1).

\medskip\noindent
Prenons pour sous-foncteur $F$ celui du lemme \ref{pic-lemma-define-open}.
La condition (2)(b) en résulte.
La condition (2)(a) est le lemme \ref{pic-lemma-open-representable}.
La condition (2)(c) est le lemme \ref{pic-lemma-twist-with-general-divisor}.
\end{proof}

\noindent
La démonstration précédente donne en fait des renseignements supplémentaires, que nous
rassemblons ici.

\begin{lemma}
\label{pic-lemma-picard-pieces}
Soit $k$ un corps séparablement clos. Soit $X$ une courbe projective lisse
de genre $g$ sur $k$.
\begin{enumerate}
\item $\underline{\Picardfunctor}_{X/k}$ est une réunion disjointe de variétés
propres et lisses de dimension $g$, notées $\underline{\Picardfunctor}^d_{X/k}$,
\item les points à valeurs dans $k$ de $\underline{\Picardfunctor}^d_{X/k}$
correspondent aux $\mathcal{O}_X$-modules inversibles de degré $d$,
\item $\underline{\Picardfunctor}^0_{X/k}$
est un sous-schéma en groupes ouvert et fermé,
\item pour $d \geq 0$, il existe un morphisme canonique
$\gamma_d :
\underline{\Hilbfunctor}^d_{X/k} \to \underline{\Picardfunctor}^d_{X/k}$
\item les morphismes $\gamma_d$
sont surjectifs pour $d \geq g$ et lisses pour $d \geq 2g - 1$,
\item le morphisme
$\underline{\Hilbfunctor}^g_{X/k} \to \underline{\Picardfunctor}^g_{X/k}$
est birationnel.
\end{enumerate}
\end{lemma}

\begin{proof}
Choisissons un point $k$-rationnel $\sigma$ de $X$. Rappelons que $\Picardfunctor_{X/k}$
est isomorphe au foncteur $\Picardfunctor_{X/k, \sigma}$. D'après
Catégories dérivées des schémas, lemme
\ref{perfect-lemma-chi-locally-constant-geometric},
pour tout $d \in \mathbf{Z}$, il existe un sous-foncteur ouvert
$$
\Picardfunctor^d_{X/k, \sigma} \subset \Picardfunctor_{X/k, \sigma}
$$
dont la valeur sur un schéma $T$ sur $k$ est formée des
$\mathcal{L} \in \Picardfunctor_{X/k, \sigma}(T)$ tels que
$\chi(X_t, \mathcal{L}_t) = d + 1 - g$, et l'on a en outre
$$
\Picardfunctor_{X/k, \sigma} =
\coprod\nolimits_{d \in \mathbf{Z}} \Picardfunctor^d_{X/k, \sigma}
$$
en tant que faisceaux fppf. Le schéma $\underline{\Picardfunctor}_{X/k}$
(qui existe d'après la proposition \ref{pic-proposition-pic-curve})
admet donc une décomposition correspondante
$$
\underline{\Picardfunctor}_{X/k, \sigma} =
\coprod\nolimits_{d \in \mathbf{Z}} \underline{\Picardfunctor}^d_{X/k, \sigma}
$$
où les points de $\underline{\Picardfunctor}^d_{X/k, \sigma}$ correspondent
aux classes d'isomorphisme de modules inversibles de degré $d$ sur $X$.

\medskip\noindent
Fixons $d \geq 0$. Il existe un morphisme
$$
\gamma_d :
\underline{\Hilbfunctor}^d_{X/k}
\longrightarrow
\underline{\Picardfunctor}^d_{X/k}
$$
provenant du faisceau inversible $\mathcal{O}(D_{univ})$ sur
$\underline{\Hilbfunctor}^d_{X/k} \times_k X$
(remarque \ref{pic-remark-universal-object-hilb-d}) par le lemme de Yoneda
(Catégories, lemme \ref{categories-lemma-yoneda}).
Notre démonstration de la représentabilité du foncteur de Picard de $X/k$
dans la proposition \ref{pic-proposition-pic-curve} et le
lemme \ref{pic-lemma-open-representable} montre que $\gamma_g$
induit une immersion ouverte sur un ouvert non vide de
$\underline{\Hilbfunctor}^g_{X/k}$. De plus, la démonstration montre
que les translatés de cet ouvert par les points $k$-rationnels du
schéma en groupes $\underline{\Picardfunctor}_{X/k}$ définissent
un recouvrement ouvert. Puisque
$\underline{\Hilbfunctor}^g_{X/k}$ est lisse de dimension $g$
(proposition \ref{pic-proposition-hilb-d})
sur $k$, on en conclut que le
schéma en groupes $\underline{\Picardfunctor}_{X/k}$ est lisse de dimension $g$
sur $k$.

\medskip\noindent
D'après Groupoïdes, lemme \ref{groupoids-lemma-group-scheme-over-field-separated},
le schéma $\underline{\Picardfunctor}_{X/k}$ est séparé.
Ainsi, pour tout $d \geq 0$, l'image de $\gamma_d$
est une variété propre sur $k$
(Morphismes, lemme \ref{morphisms-lemma-scheme-theoretic-image-is-proper}).

\medskip\noindent
Soit $d \geq g$. Pour toute extension de corps $K/k$ et tout
$\mathcal{O}_{X_K}$-module inversible $\mathcal{L}$ de degré $d$, on a
$\chi(X_K, \mathcal{L}) = d + 1 - g > 0$. Ainsi, $\mathcal{L}$ admet une
section non nulle, et l'on en conclut que $\mathcal{L} = \mathcal{O}_{X_K}(D)$
pour un diviseur $D \subset X_K$ de degré $d$. Il s'ensuit que
$\gamma_d$ est surjectif.

\medskip\noindent
En combinant les faits précédents, on voit que
$\underline{\Picardfunctor}^d_{X/k}$ est propre pour $d \geq g$.
Cela achève la démonstration de (1), car nous savons maintenant que
$\underline{\Picardfunctor}^d_{X/k}$ est propre pour $d \geq g$, et
tous les $\underline{\Picardfunctor}^d_{X/k}$ sont alors propres par translation.

\medskip\noindent
Il reste à montrer que $\gamma_d$ est lisse pour $d \geq 2g - 1$.
Considérons un $\mathcal{O}_X$-module inversible $\mathcal{L}$ de degré
$d$. La fibre au-dessus du point correspondant à $\mathcal{L}$ est alors
$$
Z = \{D \subset X \mid \mathcal{O}_X(D) \cong \mathcal{L}\} \subset
\underline{\Hilbfunctor}^d_{X/k}
$$
munie de sa structure naturelle de schéma. Puisque tout isomorphisme
$\mathcal{O}_X(D) \to \mathcal{L}$ n'est déterminé qu'à
multiplication par un scalaire non nul près, la section canonique
$1 \in \mathcal{O}_X(D)$ a pour image une section
$s \in \Gamma(X, \mathcal{L})$ bien définie à multiplication
près par un scalaire non nul. On obtient ainsi un morphisme
$$
Z \longrightarrow
\text{Proj}(\text{Sym}(\Gamma(X, \mathcal{L})^*))
$$
(le dual intervient en raison de nos conventions). Ce morphisme est un isomorphisme,
car, étant donné une section de $\mathcal{L}$, on peut prendre le diviseur de Cartier
effectif associé ; autrement dit, on peut construire un inverse du
morphisme écrit ci-dessus. Nous omettons la formulation précise et la démonstration.
Puisque $\dim H^0(X, \mathcal{L}) = d + 1 - g$ pour tout
$\mathcal{L}$ de degré $d \geq 2g - 1$ d'après
Variétés, lemme \ref{varieties-lemma-vanishing-degree-2g-and-1-line-bundle},
on voit que $\text{Proj}(\text{Sym}(\Gamma(X, \mathcal{L})^*))
\cong \mathbf{P}^{d - g}_k$.
On en conclut que $\dim(Z) = \dim(\mathbf{P}^{d - g}_k) = d - g$.
Ainsi, les fibres du morphisme $\gamma_d$ sont toutes
de dimension égale à la différence des dimensions de
$\underline{\Hilbfunctor}^d_{X/k}$ et de $\underline{\Picardfunctor}^d_{X/k}$.
Il s'ensuit que $\gamma_d$ est plat, voir
Algèbre, lemme \ref{algebra-lemma-CM-over-regular-flat}.
Comme ses fibres sont en outre lisses, on en conclut que $\gamma_d$
est lisse d'après Morphismes, lemme \ref{morphisms-lemma-smooth-flat-smooth-fibres}.
\end{proof}





\section{Quelques remarques sur les groupes de Picard}
\label{pic-section-remarks-picard}

\noindent
Cette section poursuit la discussion de
Variétés, section \ref{varieties-section-picard-groups},
qui sera prolongée dans
Courbes algébriques, section \ref{curves-section-torsion-in-pic}.

\begin{lemma}
\label{pic-lemma-pic-descends}
Soit $k$ un corps. Soit $X$ un schéma quasi-compact et quasi-séparé
sur $k$ tel que $H^0(X, \mathcal{O}_X) = k$. Si $X$ possède un point $k$-rationnel,
alors, pour toute extension galoisienne $k'/k$, on a
$$
\Pic(X) = \Pic(X_{k'})^{\text{Gal}(k'/k)}
$$
De plus, l'action de $\text{Gal}(k'/k)$ sur $\Pic(X_{k'})$
est continue.
\end{lemma}

\begin{proof}
Puisque $\text{Gal}(k'/k) = \text{Aut}(k'/k)$, ce groupe opère (à droite)
sur $\Spec(k')$, donc opère (à droite) sur
$X_{k'} = X \times_{\Spec(k)} \Spec(k')$ ; comme $\Pic(-)$
est un foncteur contravariant, il opère (à gauche) sur $\Pic(X_{k'})$.
Si $k'/k$ est une extension galoisienne infinie, on écrit
$k' = \colim k'_\lambda$ comme limite inductive filtrante d'extensions galoisiennes
finies, voir Corps, lemme \ref{fields-lemma-infinite-galois-limit}.
Alors $X_{k'} = \lim X_{k'_\lambda}$
(comme dans Limites, section \ref{limits-section-limits})
et l'on obtient
$$
\Pic(X_{k'}) = \colim \Pic(X_{k'_\lambda})
$$
d'après Limites, lemme \ref{limits-lemma-descend-invertible-modules}.
De plus, les applications de transition de ce système de groupes abéliens
sont injectives d'après Variétés, lemme \ref{varieties-lemma-change-fields-pic}.
Il s'ensuit que tout élément de $\Pic(X_{k'})$ est fixé par
l'un des sous-groupes ouverts $\text{Gal}(k'/k'_\lambda)$, ce qui signifie exactement
que l'action est continue. L'injectivité des applications de transition
montre qu'il suffit de démontrer l'assertion sur les points fixes dans le
cas où $k'/k$ est une extension galoisienne finie.

\medskip\noindent
Supposons $k'/k$ galoisienne finie, de groupe de Galois $G = \text{Gal}(k'/k)$.
Soit $\mathcal{L}$ un élément de $\Pic(X_{k'})$ fixé par $G$.
Nous utiliserons la descente galoisienne (Descente, lemme \ref{descent-lemma-galois-descent})
pour montrer que $\mathcal{L}$ est l'image réciproque d'un faisceau inversible sur $X$.
Rappelons que $f_\sigma = \text{id}_X \times \Spec(\sigma) : X_{k'} \to X_{k'}$
et que $\sigma$ opère sur $\Pic(X_{k'})$ par image réciproque suivant $f_\sigma$.
Pour tout $\sigma \in G$, on peut donc choisir un isomorphisme
$\varphi_\sigma : \mathcal{L} \to f_\sigma^*\mathcal{L}$,
puisque $\mathcal{L}$ est fixé par l'action de $G$.
La difficulté est que l'on ne sait pas si l'on peut choisir
$\varphi_\sigma$ de manière à satisfaire à la condition de cocycle
$\varphi_{\sigma\tau} = f_\sigma^*\varphi_\tau \circ \varphi_\sigma$.
Pour voir que c'est possible, nous utilisons le fait que $X$ possède un point $k$-rationnel
$x \in X(k)$. Bien entendu, $x$ détermine de même un point $k'$-rationnel
$x' \in X_{k'}$ fixé par $f_\sigma$ pour tout $\sigma$.
Choisissons un élément non nul $s$ de la fibre de $\mathcal{L}$ en $x'$ ;
cette fibre est l'espace vectoriel de dimension $1$ sur $k' = \kappa(x')$
$$
\mathcal{L}_{x'} \otimes_{\mathcal{O}_{X_{k'}, x'}} \kappa(x').
$$
Alors $f_\sigma^*s$ est un élément non nul de la fibre de
$f_\sigma^*\mathcal{L}$ en $x'$. Puisqu'on peut multiplier $\varphi_\sigma$
par un élément de $(k')^*$, on peut supposer que $\varphi_\sigma$ envoie
$s$ sur $f_\sigma^*s$. On voit alors que $\varphi_{\sigma\tau}$ et
$f_\sigma^*\varphi_\tau \circ \varphi_\sigma$
envoient tous deux $s$ sur $f_{\sigma\tau}^*s = f_\tau^*f_\sigma^*s$.
Puisque $H^0(X_{k'}, \mathcal{O}_{X_{k'}}) = k'$, ces deux isomorphismes
sont nécessairement égaux (l'un est le produit de l'autre par une unité globale
et ils coïncident en $x'$), ce qui achève la démonstration.
\end{proof}

\begin{lemma}
\label{pic-lemma-torsion-descends}
Soit $k$ un corps de caractéristique $p > 0$. Soit $X$ un
schéma quasi-compact et quasi-séparé sur $k$ tel que
$H^0(X, \mathcal{O}_X) = k$. Soit $n$ un entier premier à $p$.
Alors l'application
$$
\Pic(X)[n] \longrightarrow \Pic(X_{k'})[n]
$$
est bijective pour toute extension purement inséparable $k'/k$.
\end{lemma}

\begin{proof}
Notons d'abord que l'application $\Pic(X) \to \Pic(X_{k'})$
est injective d'après Variétés, lemme \ref{varieties-lemma-change-fields-pic}.
Il faut donc montrer que l'application du lemme est surjective.
Soit $\mathcal{L}$ un $\mathcal{O}_{X_{k'}}$-module inversible
dont la classe est d'ordre divisant $n$ dans $\Pic(X_{k'})$.
Choisissons un isomorphisme
$\alpha : \mathcal{L}^{\otimes n} \to \mathcal{O}_{X_{k'}}$
de modules inversibles.
Nous allons montrer que le couple $(\mathcal{L}, \alpha)$ descend
à $X$.

\medskip\noindent
Posons $A = k' \otimes_k k'$. Puisque $k'/k$ est purement inséparable,
le noyau de l'application de multiplication $A \to k'$ est un
idéal localement nilpotent $I$ de $A$. Notons que
$$
X_A = X \times_{\Spec(k)} \Spec(A) = X_{k'} \times_X X_{k'}
$$
est muni de deux projections $\text{pr}_i : X_A \to X_{k'}$, $i = 0, 1$,
qui coïncident sur $A/I$. Les modules inversibles
$\mathcal{L}_i = \text{pr}_i^*\mathcal{L}$
coïncident donc sur le sous-schéma fermé $X_{A/I} = X_{k'}$.
Puisque $X_{A/I} \to X_A$ est un épaississement et que
les $\mathcal{L}_i$ sont de $n$-torsion, il
existe un isomorphisme $\varphi : \mathcal{L}_0 \to \mathcal{L}_1$ d'après
Compléments sur les morphismes, lemme \ref{more-morphisms-lemma-torsion-pic-thickening}.
On peut choisir $\varphi$ dont la réduction modulo $I$ soit l'identité. En effet,
$H^0(X, \mathcal{O}_X) = k$ entraîne
$H^0(X_{k'}, \mathcal{O}_{X_{k'}}) = k'$ d'après
Cohomologie des schémas, lemme \ref{coherent-lemma-flat-base-change-cohomology},
et $A \to k'$ est surjective ; on peut donc ajuster $\varphi$ en le multipliant par
une unité convenable de $A$. Considérons l'application
$$
\lambda : \mathcal{O}_{X_A}
\xrightarrow{\text{pr}_0^*\alpha^{-1}}
\mathcal{L}_0^{\otimes n}
\xrightarrow{\varphi^{\otimes n}}
\mathcal{L}_1^{\otimes n}
\xrightarrow{\text{pr}_1^*\alpha}
\mathcal{O}_{X_A}
$$
On peut regarder $\lambda$ comme un élément de $A$, car
$H^0(X_A, \mathcal{O}_{X_A}) = A$ (même référence que ci-dessus).
Puisque $\varphi$ se réduit à l'identité modulo $I$, on
voit que $\lambda = 1 \bmod I$. Il existe alors une unique
racine $n$-ième de $\lambda$ dans $1 + I$
(Algèbre, lemme \ref{algebra-lemma-lift-nth-roots})
et, en multipliant $\varphi$ par son inverse, on obtient $\lambda = 1$.
Nous affirmons que $(\mathcal{L}, \varphi)$ est une donnée de descente
pour le recouvrement fpqc $\{X_{k'} \to X\}$
(Descente, définition \ref{descent-definition-descent-datum-quasi-coherent}).
Si cette assertion est vraie, $\mathcal{L}$ est l'image réciproque d'un
$\mathcal{O}_X$-module inversible $\mathcal{N}$ d'après
Descente, proposition \ref{descent-proposition-fpqc-descent-quasi-coherent}.
L'injectivité de l'application entre les groupes de Picard montre que $\mathcal{N}$
est un élément de torsion de $\Pic(X)$ de même ordre que $\mathcal{L}$.

\medskip\noindent
Démonstration de l'assertion. Il faut vérifier que
$$
\text{pr}_{12}^*\varphi \circ
\text{pr}_{01}^*\varphi =
\text{pr}_{02}^*\varphi
\quad\text{sur}\quad
X_{k'} \times_X X_{k'} \times_X X_{k'} = X_{k' \otimes_k k' \otimes_k k'}
$$
Comme précédemment, le morphisme diagonal
$\Delta : X_{k'} \to X_{k' \otimes_k k' \otimes_k k'}$
est un épaississement. Les membres de gauche et de droite de l'égalité
sont des applications $a, b : p_0^*\mathcal{L} \to p_2^*\mathcal{L}$ compatibles avec
$p_0^*\alpha$ et $p_2^*\alpha$, où
$p_i : X_{k' \otimes_k k' \otimes_k k'} \to X_{k'}$
sont les morphismes de projection. Enfin, les images réciproques de $a, b$ par
$\Delta$ sont égales.
Localement sur des ouverts affines (dans des trivialisations locales), cela signifie que
$a, b$ sont données par multiplication par des fonctions inversibles
qui se réduisent à la même fonction modulo un idéal localement nilpotent
et dont les puissances $n$-ièmes sont égales. Il résulte alors de
Algèbre, lemme \ref{algebra-lemma-lift-nth-roots}
que ces fonctions sont égales.
\end{proof}









% OK, start here.
%

\chapter{Théories cohomologiques de Weil}



\phantomsection
\label{weil-section-phantom}




\section{Introduction}
\label{weil-section-introduction}

\noindent
Nous étudions dans ce chapitre les théories cohomologiques de Weil pour les schémas
projectifs lisses sur un corps de base. En bref, nous entendons par là une théorie
cohomologique $H^*$ munie d'une formule de K\"unneth, d'une dualité de Poincar\'e
et de classes de cycles (satisfaisant aux compatibilités convenables). Nous
avertissons le lecteur que la notion de « théorie cohomologique de Weil »
ne fait pas l'objet d'un accord universel dans la littérature.

\medskip\noindent
Avant d'aborder ce chapitre, le lecteur consultera
Catégories, section \ref{categories-section-monoidal}, et
Homologie, section \ref{homology-section-monoidal}, où
nous définissons les catégories monoïdales (symétriques) et développons
les rudiments de leur langage nécessaires au présent chapitre.
À l'aide de ce langage, nous construisons, à la
section \ref{weil-section-correspondences}, la catégorie graduée monoïdale
symétrique dont les objets sont les schémas projectifs lisses et
les morphismes les correspondances. À la section \ref{weil-section-chow-motives},
nous adjoignons les images des projecteurs et inversons le motif de Lefschetz
pour obtenir la catégorie karoubienne monoïdale symétrique $M_k$
des motifs de Chow. Cette catégorie est munie d'un foncteur contravariant
$$
h : \{\text{schémas projectifs lisses sur }k\} \longrightarrow M_k
$$
Comme nous le verrons ci-dessous, une propriété essentielle d'une théorie
cohomologique de Weil est de se factoriser par $h$.

\medskip\noindent
Tout d'abord, lorsque le corps de base est algébriquement clos, nous définissons
ce que nous appelons une « théorie cohomologique de Weil classique » ;
voir la section \ref{weil-section-axioms-classical}. Cette notion est
celle introduite dans \cite[section 1.2]{Kleiman-cycles} et
coïncide avec celle de \cite[page 65]{Kleiman-motives}.
En revanche, elle ne coïncide pas a priori avec celle introduite dans
\cite[page 10]{Kleiman-standard}, car l'auteur y ajoute deux axiomes
de type Lefschetz, et l'on ignore si toute théorie cohomologique de Weil
classique au sens du présent chapitre satisfait à ces axiomes.
À la fin de la section \ref{weil-section-axioms-classical}, nous montrons
qu'une théorie cohomologique de Weil classique est de la forme $H^* = G \circ h$,
où $G$ est un foncteur monoïdal symétrique de $M_k$ dans la catégorie
des espaces vectoriels gradués sur le corps de coefficients de $H^*$.

\medskip\noindent
À la section \ref{weil-section-cycles-nonclosed}, nous démontrons quelques lemmes
sur les groupes de cycles sur des corps non clos, qui serviront à étudier
les théories cohomologiques de Weil sur les schémas projectifs lisses
sur un corps quelconque.

\medskip\noindent
Les propriétés suivantes motivent notre choix d'axiomes pour une théorie
cohomologique de Weil $H^*$ sur un corps de base quelconque $k$ :
\begin{enumerate}
\item $H^* = G \circ h$ pour un foncteur monoïdal symétrique $G$ de $M_k$
dans la catégorie des espaces vectoriels gradués sur le corps de coefficients $F$ de $H^*$ ;
\item $G$ doit envoyer le motif de Tate (inverse du motif de Lefschetz)
sur un espace vectoriel de dimension $1$, noté $F(1)$ et placé en degré $-2$ ;
\item lorsque $k$ est algébriquement clos, on doit retrouver la notion
étudiée à la section \ref{weil-section-axioms-classical},
au choix près d'un élément d'une base de $F(1)$.
\end{enumerate}
Nous analysons d'abord les deux premières conditions à la section \ref{weil-section-axioms}.
Après avoir établi quelques résultats supplémentaires à la section \ref{weil-section-further},
nous ajoutons, à la section \ref{weil-section-old}, les axiomes nécessaires
pour obtenir la propriété (3).

\medskip\noindent
Dans la dernière section \ref{weil-section-c1}, nous exposons en détail
une autre approche des théories cohomologiques de Weil, fondée sur
une application de première classe de Chern au lieu des classes de cycles.
C'est cette approche qui se prêtera le mieux à la démonstration, dans
les chapitres ultérieurs, du fait que certaines théories cohomologiques
sont des théories cohomologiques de Weil ; voir
Cohomologie de de Rham, section \ref{derham-section-weil}.





\section{Conventions et notations}
\label{weil-section-conventions}

\noindent
Soit $F$ un corps. Dans ce chapitre, nous considérons la catégorie des espaces
vectoriels gradués sur $F$ comme une catégorie monoïdale symétrique $F$-linéaire,
munie de la contrainte d'associativité usuelle et d'une contrainte de commutativité
faisant intervenir des signes.
Voir Homologie, exemple \ref{homology-example-graded-vector-spaces}.

\medskip\noindent
Soit $R$ un anneau. Dans ce chapitre, une
{\it $R$-algèbre graduée commutative} $A$ est une
$R$-algèbre différentielle graduée commutative
(Algèbre différentielle graduée, définitions \ref{dga-definition-dga} et
\ref{dga-definition-cdga}) dont la différentielle est nulle. Ainsi, $A$
est un $R$-module muni d'une graduation
$A = \bigoplus_{n \in \mathbf{Z}} A^n$ par des
sous-$R$-modules. La multiplication $R$-bilinéaire
$$
A^n \times A^m \longrightarrow A^{n + m},\quad
\alpha \times \beta \longmapsto \alpha \cup \beta
$$
sera appelée le {\it cup-produit} dans ce chapitre.
La contrainte de commutativité s'écrit
$\alpha \cup \beta = (-1)^{nm} \beta \cup \alpha$ si
$\alpha \in A^n$ et $\beta \in A^m$. Enfin, il existe
un élément unité $1 \in A^0$ pour la multiplication, ou, ce qui revient au même,
une application additive et multiplicative $R \to A^0$, compatible avec
la structure de $R$-module de $A$.

\medskip\noindent
Soit $k$ un corps et soit $X$ un schéma de type fini sur $k$.
Les groupes de Chow $\CH_k(X)$ de $X$, formés des cycles de dimension $k$
modulo l'équivalence rationnelle, ont été définis dans
Homologie de Chow, définition \ref{chow-definition-rational-equivalence}.
Si $X$ est normal ou de Cohen-Macaulay, on peut aussi considérer
les groupes de Chow $\CH^p(X)$ de cycles de codimension $p$
(Homologie de Chow, section \ref{chow-section-cycles-codimension}) ;
alors $[X] \in \CH^0(X)$ désigne la « classe fondamentale » de $X$ ; voir
Homologie de Chow, remarque \ref{chow-remark-fundamental-class}.
Si $X$ est lisse et si $\alpha$ et $\beta$ sont des cycles sur $X$,
$\alpha \cdot \beta$ désigne le produit d'intersection de
$\alpha$ et $\beta$ ; voir
Homologie de Chow, section \ref{chow-section-intersection-product}.













\section{Correspondances}
\label{weil-section-correspondences}

\noindent
Soit $k$ un corps. Pour des schémas $X$ et $Y$ sur $k$, nous notons
$X \times Y$ le produit de $X$ et $Y$ dans la catégorie des schémas
sur $k$. Nous construisons dans cette section la catégorie graduée sur
$\mathbf{Q}$ dont les objets sont les schémas projectifs lisses sur $k$ et dont
les morphismes sont les correspondances.

\medskip\noindent
Soient $X$ et $Y$ des schémas projectifs lisses sur $k$.
Soit $X = \coprod X_d$ la décomposition de $X$ en
sous-schémas ouverts et fermés équidimensionnels
tels que $\dim(X_d) = d$. Nous définissons l'espace vectoriel sur $\mathbf{Q}$
{\it des correspondances de degré $r$ de $X$ dans $Y$}
par la formule :
$$
\text{Corr}^r(X, Y) =
\bigoplus\nolimits_d \CH^{d + r}(X_d \times Y) \otimes \mathbf{Q}
\subset
\CH^*(X \times Y) \otimes \mathbf{Q}
$$
Étant donnés $c \in \text{Corr}^r(X, Y)$ et $\beta \in \CH_j(Y) \otimes \mathbf{Q}$,
nous définissons l'{\it image réciproque} de $\beta$ par $c$ par la formule
$$
c^*(\beta) = \text{pr}_{1, *}(c \cdot \text{pr}_2^*\beta)
\quad\text{dans}\quad
\CH_{j - r}(X) \otimes \mathbf{Q}
$$
Cette expression a un sens, car $\text{pr}_2$ est plat de dimension relative
$d$ sur $X_d \times Y$ ; ainsi $\text{pr}_2^*\beta$ est un cycle de
dimension $d + j$ sur $X_d \times Y$, donc $c \cdot \text{pr}_2^*\beta$
est un cycle de dimension $j - r$ sur $X_d \times Y$, dont l'image directe
par le morphisme propre $\text{pr}_1$ est un cycle de même dimension.
De même, en passant à la graduation par la codimension,
pour $\alpha \in \CH^i(X) \otimes \mathbf{Q}$, nous définissons
l'{\it image directe} de $\alpha$ par $c$ par la formule
$$
c_*(\alpha) = \text{pr}_{2, *}(c \cdot \text{pr}_1^*\alpha)
\quad\text{dans}\quad
\CH^{i + r}(Y) \otimes \mathbf{Q}
$$
Cette expression a un sens, car $\text{pr}_1^*\alpha$ est un cycle de codimension
$i$ sur $X \times Y$ ; ainsi $c \cdot \text{pr}_1^*\alpha$ est un cycle
de codimension $i + d + r$ sur $X_d \times Y$, dont l'image directe
est un cycle de codimension $i + r$ sur $Y$.

\medskip\noindent
Étant donnés trois schémas projectifs lisses $X, Y, Z$ sur $k$, nous définissons une
{\it composition des correspondances}
$$
\text{Corr}^s(Y, Z)
\times
\text{Corr}^r(X, Y)
\longrightarrow
\text{Corr}^{r + s}(X, Z)
$$
par la règle
$$
(c', c)
\longmapsto
c' \circ c =
\text{pr}_{13, *}(\text{pr}_{12}^*c \cdot \text{pr}_{23}^*c')
$$
où $\text{pr}_{12} : X \times Y \times Z \to X \times Y$
est la projection, et de même pour $\text{pr}_{13}$ et $\text{pr}_{23}$.

\begin{lemma}
\label{weil-lemma-composition-correspondences}
Les correspondances possèdent les propriétés suivantes :
\begin{enumerate}
\item la composition des correspondances est $\mathbf{Q}$-bilinéaire
et associative ;
\item il existe un isomorphisme canonique
$$
\CH_{-r}(X) \otimes \mathbf{Q} = \text{Corr}^r(X, \Spec(k))
$$
par lequel l'image réciproque par les correspondances correspond à la composition ;
\item il existe un isomorphisme canonique
$$
\CH^r(X) \otimes \mathbf{Q} = \text{Corr}^r(\Spec(k), X)
$$
par lequel l'image directe par les correspondances correspond à la composition ;
\item la composition des correspondances est compatible avec l'image directe
et l'image réciproque des cycles.
\end{enumerate}
\end{lemma}

\begin{proof}
La bilinéarité résulte immédiatement de la linéarité de l'image directe
et de l'image réciproque, et de la bilinéarité du produit d'intersection.
Pour démontrer l'associativité, donnons-nous
$X, Y, Z, W$ ainsi que $c \in \text{Corr}(X, Y)$, $c' \in \text{Corr}(Y, Z)$ et
$c'' \in \text{Corr}(Z, W)$. On a alors
\begin{align*}
c'' \circ (c' \circ c)
& =
\text{pr}^{134}_{14, *}(
\text{pr}^{134, *}_{13}
\text{pr}^{123}_{13, *}(\text{pr}^{123, *}_{12}c \cdot
\text{pr}^{123, *}_{23}c')
\cdot \text{pr}^{134, *}_{34}c'') \\
& =
\text{pr}^{134}_{14, *}(
\text{pr}^{1234}_{134, *}
\text{pr}^{1234, *}_{123}(\text{pr}^{123, *}_{12}c \cdot
\text{pr}^{123, *}_{23}c')
\cdot \text{pr}^{134, *}_{34}c'') \\
& =
\text{pr}^{134}_{14, *}(
\text{pr}^{1234}_{134, *}
(\text{pr}^{1234, *}_{12}c \cdot
\text{pr}^{1234, *}_{23}c')
\cdot \text{pr}^{134, *}_{34}c'') \\
& =
\text{pr}^{134}_{14, *}
\text{pr}^{1234}_{134, *}
((\text{pr}^{1234, *}_{12}c \cdot
\text{pr}^{1234, *}_{23}c')
\cdot \text{pr}^{1234, *}_{34}c'') \\
& =
\text{pr}^{1234}_{14, *}(
(\text{pr}^{1234, *}_{12}c \cdot
\text{pr}^{1234, *}_{23}c') \cdot
\text{pr}^{1234, *}_{34}c'')
\end{align*}
Nous notons ici
$$
p^{1234}_{134} : X \times Y \times Z \times W \to X \times Z \times W
\quad\text{et}\quad
p^{134}_{14} : X \times Z \times W \to X \times W
$$
{\emergencystretch=1em les projections, et de même pour les autres indices.
La première égalité est la définition de la composition.
La deuxième résulte de l'égalité
$\text{pr}^{134, *}_{13} \text{pr}^{123}_{13, *} =
\text{pr}^{1234}_{134, *} \text{pr}^{1234, *}_{123}$,
d'après Homologie de Chow, lemme \ref{chow-lemma-flat-pullback-proper-pushforward}.
La troisième résulte de la compatibilité du produit d'intersection
avec le morphisme de Gysin de $p^{1234}_{123}$ (qui est donné par l'image réciproque plate) ; voir
Homologie de Chow, lemme \ref{chow-lemma-lci-gysin-product}.
La quatrième égalité résulte de la formule de projection pour
$p^{1234}_{134}$ ; voir Homologie de Chow, lemme \ref{chow-lemma-projection-formula}.
La quatrième égalité exprime la compatibilité de l'image directe propre
avec la composition ; voir Homologie de Chow, lemme \ref{chow-lemma-compose-pushforward}.
Le produit d'intersection étant associatif d'après
Homologie de Chow, lemme \ref{chow-lemma-associative},
cela achève la démonstration de l'associativité de la composition des correspondances.\par}

\medskip\noindent
Nous omettons les démonstrations de (2) et (3), qui consistent essentiellement
à vérifier avec soin les espaces où vivent les différents cycles et leurs (co)dimensions.

\medskip\noindent
L'assertion relative à l'image directe et à l'image réciproque des cycles
signifie que $(c' \circ c)^*(\alpha) = c^*((c')^*(\alpha))$ et
$(c' \circ c)_*(\alpha) = (c')_*(c_*(\alpha))$.
Elle résulte de la combinaison de (1), (2) et (3).
\end{proof}

\begin{example}
\label{weil-example-graph-correspondence}
Soit $f : Y \to X$ un morphisme de schémas projectifs lisses sur $k$.
Notons $\Gamma_f \subset X \times Y$ le graphe de $f$. Plus précisément,
$\Gamma_f$ est l'image de l'immersion fermée
$$
(f, \text{id}_Y) : Y \longrightarrow X \times Y
$$
Soit $X = \coprod X_d$ la décomposition de $X$ en ses
parties ouvertes et fermées $X_d$, équidimensionnelles de dimension $d$.
Alors $\Gamma_f \cap (X_d \times Y)$ est de codimension pure $d$. Par conséquent,
$[\Gamma_f] \in \CH^*(X \times Y) \otimes \mathbf{Q}$
appartient à $\text{Corr}^0(X \times Y)$ ; autrement dit, $[\Gamma_f]$
est une correspondance de degré $0$ de $X$ dans $Y$.
\end{example}

\begin{lemma}
\label{weil-lemma-category-correspondences}
Les schémas projectifs lisses sur $k$, munis des correspondances et
de leur composition définies ci-dessus, forment une catégorie graduée sur
$\mathbf{Q}$
(Algèbre différentielle graduée, définition \ref{dga-definition-graded-category}).
\end{lemma}

\begin{proof}
Tout résulte de la construction et du
lemme \ref{weil-lemma-composition-correspondences},
sauf l'existence des morphismes identiques.
Pour un schéma projectif lisse $X$, considérons
la classe $[\Delta]$ de la diagonale $\Delta \subset X \times X$
dans $\text{Corr}^0(X, X)$. Observons que $\Delta$
est le graphe de l'identité $\text{id}_X : X \to X$ ;
nous utiliserons ce fait ci-dessous.

\medskip\noindent
Pour démontrer que $[\Delta]$ joue le rôle d'identité, il faut établir
$[\Delta] \circ c = c$ et $c' \circ [\Delta] = c'$ pour toutes correspondances
$c \in \text{Corr}^r(Y, X)$ et $c' \in \text{Corr}^s(X, Y)$.
Dans le deuxième cas, il faut montrer que
$$
c' = \text{pr}_{13, *}(\text{pr}_{12}^*[\Delta] \cdot \text{pr}_{23}^*c')
$$,
où $\text{pr}_{12} : X \times X \times Y \to X \times X$ est la
projection, et de même pour $\text{pr}_{13}$ et $\text{pr}_{23}$.
On peut écrire $c' = \sum a_i [Z_i]$, pour certains sous-schémas fermés intègres
$Z_i \subset X \times Y$ et nombres rationnels $a_i$. Il suffit donc
manifestement de montrer que
$$
[Z] = \text{pr}_{13, *}(\text{pr}_{12}^*[\Delta] \cdot \text{pr}_{23}^*[Z])
$$
dans le groupe de Chow de $X \times Y$, pour tout sous-schéma fermé intègre $Z$
de $X \times Y$. En remplaçant $X$ et $Y$ par
la composante irréductible contenant l'image de $Z$ par chacune des deux projections,
on peut supposer $X$ et $Y$ également intègres. Il faut alors montrer
$$
[Z] = \text{pr}_{13, *}([\Delta \times Y] \cdot [X \times Z])
$$
Notons $Z' \subset X \times X \times Y$ l'image de $Z$ par le morphisme
$(\Delta, 1) : X \times Y \to X \times X \times Y$. Alors $Z'$
est un sous-schéma fermé de $X \times X \times Y$ isomorphe à $Z$, et
$Z' = \Delta \times Y \cap X \times Z$ au sens schématique.
D'après Homologie de Chow, lemme \ref{chow-lemma-intersect-properly}\footnote{Le
lecteur vérifiera que $\dim(Z') = \dim(\Delta \times Y) + \dim(X \times Z) -
\dim(X \times X \times Y)$ et que $Z'$ possède un unique point générique,
dont l'image est le point générique de $Z$ (où l'anneau local est de Cohen-Macaulay)
et un point de $X$ (où l'anneau local est de Cohen-Macaulay).
Toutes les hypothèses du lemme sont donc bien vérifiées.},
on en déduit
$$
[Z'] = [\Delta \times Y] \cdot [X \times Z]
$$
Puisque $Z'$ s'envoie aussi isomorphiquement sur $Z$ par $\text{pr}_{13}$,
on obtient la conclusion. La vérification de
$[\Delta] \circ c = c$ est analogue et nous l'omettons.
\end{proof}

\begin{lemma}
\label{weil-lemma-contravariant-functor}
Il existe un foncteur contravariant de la catégorie des schémas
projectifs lisses sur $k$ dans la catégorie des correspondances,
qui est l'identité sur les objets et envoie $f : Y \to X$ sur
l'élément $[\Gamma_f] \in \text{Corr}^0(X, Y)$.
\end{lemma}

\begin{proof}
Dans la démonstration du lemme \ref{weil-lemma-category-correspondences},
nous avons vu que cette construction envoie les identités sur
les identités. Pour achever la démonstration, il faut montrer que, si $g : Z \to Y$
est un autre morphisme de schémas projectifs lisses sur $k$, alors
$[\Gamma_g] \circ [\Gamma_f] = [\Gamma_{f \circ g}]$ dans
$\text{Corr}^0(X, Z)$. En raisonnant comme dans la démonstration du
lemme \ref{weil-lemma-category-correspondences}, on voit qu'il
suffit de montrer
$$
[\Gamma_{f \circ g}] =
\text{pr}_{13, *}([\Gamma_f \times Z] \cdot [X \times \Gamma_g])
$$
dans $\CH^*(X \times Z)$ lorsque $X$, $Y$ et $Z$ sont intègres.
Notons $Z' \subset X \times Y \times Z$ l'image de l'immersion fermée
$(f \circ g, g, 1) : Z \to X \times Y \times Z$.
Alors $Z' = \Gamma_f \times Z \cap X \times \Gamma_g$
au sens schématique, et l'on déduit de
Homologie de Chow, lemme \ref{chow-lemma-intersect-properly},
que
$$
[Z'] = [\Gamma_f \times Z] \cdot [X \times \Gamma_g]
$$
Comme il est clair que $\text{pr}_{13, *}([Z']) = [\Gamma_{f \circ g}]$,
la démonstration est terminée.
\end{proof}

\begin{remark}
\label{weil-remark-transpose}
Soient $X$ et $Y$ des schémas projectifs lisses sur $k$.
Supposons $X$ équidimensionnel de dimension $d$ et
$Y$ équidimensionnel de dimension $e$. L'isomorphisme
$X \times Y \to Y \times X$ qui échange les facteurs détermine alors
un isomorphisme
$$
\text{Corr}^r(X, Y) \longrightarrow \text{Corr}^{d - e + r}(Y, X),\quad
c \longmapsto c^t
$$,
appelé la {\it transposition}. Il agit aussi bien sur les cycles que sur leurs classes.
Un exemple parfois utile est la transposée
$[\Gamma_f]^t = [\Gamma_f^t]$ du graphe d'un morphisme $f : Y \to X$.
\end{remark}

\begin{lemma}
\label{weil-lemma-functor-and-cycles}
Soit $f : Y \to X$ un morphisme de schémas projectifs lisses sur $k$.
Soit $[\Gamma_f] \in \text{Corr}^0(X, Y)$ comme dans
l'exemple \ref{weil-example-graph-correspondence}. Alors :
\begin{enumerate}
\item l'image directe des cycles par la correspondance $[\Gamma_f]$
coïncide avec le morphisme de Gysin $f^! : \CH^*(X) \to \CH^*(Y)$ ;
\item l'image réciproque des cycles par la correspondance $[\Gamma_f]$
coïncide avec l'application image directe $f_* : \CH_*(Y) \to \CH_*(X)$ ;
\item si $X$ et $Y$ sont équidimensionnels de dimensions $d$ et $e$,
alors :
\begin{enumerate}
\item l'image directe des cycles par la correspondance
$[\Gamma_f^t]$ de la remarque \ref{weil-remark-transpose}
correspond à l'image directe des cycles par $f$ ;
\item l'image réciproque des cycles par la correspondance
$[\Gamma_f^t]$ de la remarque \ref{weil-remark-transpose}
correspond au morphisme de Gysin $f^!$.
\end{enumerate}
\end{enumerate}
\end{lemma}

\begin{proof}
Démontrons (1). Rappelons que
$[\Gamma_f]_*(\alpha) =
\text{pr}_{2, *}([\Gamma_f] \cdot \text{pr}_1^*\alpha)$.
On a
$$
[\Gamma_f] \cdot \text{pr}_1^*\alpha =
(f, 1)_*((f, 1)^! \text{pr}_1^*\alpha) =
(f, 1)_*((f, 1)^! \text{pr}_1^!\alpha) =
(f, 1)_*(f^!\alpha)
$$
La première égalité résulte de Homologie de Chow, lemme
\ref{chow-lemma-intersect-regularly-embedded}.
La deuxième résulte de
Homologie de Chow, lemme \ref{chow-lemma-lci-gysin-flat}.
La troisième résulte de $\text{pr}_1 \circ (f, 1) = f$ et de
Homologie de Chow, lemme \ref{chow-lemma-lci-gysin-composition}.
On conclut alors puisque
$\text{pr}_{2, *} \circ (f, 1)_* = 1_*$, d'après
Homologie de Chow, lemme \ref{chow-lemma-compose-pushforward}.

\medskip\noindent
Démontrons (2). Rappelons que $[\Gamma_f]_*(\beta) =
\text{pr}_{1, *}([\Gamma_f] \cdot \text{pr}_2^*\beta)$.
En raisonnant exactement comme ci-dessus, on a
$$
[\Gamma_f] \cdot \text{pr}_2^*\beta = (f, 1)_*\beta
$$
Le résultat s'en déduit comme précédemment.

\medskip\noindent
Démontrons (3). La démonstration est exactement analogue à la précédente.
\end{proof}

\begin{example}
\label{weil-example-decompose-P1}
Soit $X = \mathbf{P}^1_k$. On a alors
$$
\text{Corr}^0(X, X) = \CH^1(X \times X) \otimes \mathbf{Q} =
\CH_1(X \times X) \otimes \mathbf{Q}
$$
Choisissons un point $k$-rationnel $x \in X$ et
considérons les cycles $c_0 = [x \times X]$ et $c_2 = [X \times x]$.
Un calcul montre que $1 = [\Delta] = c_0 + c_2$ dans $\text{Corr}^0(X, X)$
et que l'on a les règles de composition suivantes :
$c_0 \circ c_0 = c_0$,
$c_0 \circ c_2 = 0$,
$c_2 \circ c_0 = 0$ et
$c_2 \circ c_2 = c_2$.
Autrement dit, $c_0$ et $c_2$ sont des idempotents orthogonaux dans
l'algèbre $\text{Corr}^0(X, X)$ ; en fait, on obtient
$$
\text{Corr}^0(X, X) = \mathbf{Q} \times \mathbf{Q}
$$
comme $\mathbf{Q}$-algèbre.
\end{example}

\noindent
La catégorie des correspondances est une catégorie monoïdale symétrique.
Pour des schémas projectifs lisses $X$ et $Y$ sur $k$, on pose
$X \otimes Y = X \times Y$. Pour quatre schémas projectifs lisses
$X, X', Y, Y'$ sur $k$, on définit un produit tensoriel
$$
\otimes :
\text{Corr}^r(X, Y) \times \text{Corr}^{r'}(X', Y')
\longrightarrow
\text{Corr}^{r + r'}(X \times X', Y \times Y')
$$
par la règle
$$
(c, c') \longmapsto
c \otimes c' = \text{pr}_{13}^*c \cdot \text{pr}_{24}^*c'
$$,
où $\text{pr}_{13} : X \times X' \times Y \times Y' \to X \times Y$
et $\text{pr}_{24} : X \times X' \times Y \times Y' \to X' \times Y'$
sont les projections. Comme contrainte d'associativité
$$
X \otimes (Y \otimes Z) = (X \otimes Y) \otimes Z
$$,
{\emergencystretch=1em on prend la contrainte d'associativité usuelle des produits de schémas.
La contrainte de commutativité est donnée par l'isomorphisme
$X \times Y \to Y \times X$ qui échange les facteurs.\par}

\begin{lemma}
\label{weil-lemma-tensor-product}
Le produit tensoriel des correspondances défini ci-dessus fait de la catégorie
des correspondances une catégorie monoïdale symétrique d'unité $\Spec(k)$.
\end{lemma}

\begin{proof}
Démonstration omise.
\end{proof}

\begin{lemma}
\label{weil-lemma-prep-dual}
Soit $f : Y \to X$ un morphisme de schémas projectifs lisses sur $k$.
Supposons $X$ et $Y$ équidimensionnels de dimensions $d$ et $e$.
Notons $a = [\Gamma_f] \in \text{Corr}^0(X, Y)$ et
$a^t = [\Gamma_f^t] \in \text{Corr}^{d - e}(Y, X)$.
Posons
$\eta_X = [\Gamma_{X \to X \times X}] \in \text{Corr}^0(X \times X, X)$,
$\eta_Y = [\Gamma_{Y \to Y \times Y}] \in \text{Corr}^0(Y \times Y, Y)$,
$[X] \in \text{Corr}^{-d}(X, \Spec(k))$ et
$[Y] \in \text{Corr}^{-e}(Y, \Spec(k))$. Le diagramme
$$
\xymatrix{
X \otimes Y \ar[r]_{a \otimes \text{id}} \ar[d]_{\text{id} \otimes a^t} &
Y \otimes Y \ar[r]_{\eta_Y} &
Y \ar[d]^{[Y]} \\
X \otimes X \ar[r]^{\eta_X} &
X \ar[r]^{[X]} &
\Spec(k)
}
$$
est commutatif dans la catégorie des correspondances.
\end{lemma}

\begin{proof}
{\emergencystretch=1em Rappelons que $\text{Corr}^r(W, \Spec(k)) = \CH_{-r}(W)$ pour tout
schéma projectif lisse $W$ sur $k$
et que, pour $c \in \text{Corr}^s(W', W)$, la composition
avec $c$ coïncide avec l'image réciproque par $c$, comme application
$\CH_{-r}(W) \to \CH_{-r - s}(W')$
(lemme \ref{weil-lemma-composition-correspondences}).
Enfin, le lemme \ref{weil-lemma-functor-and-cycles}
indique comment exprimer cela au moyen des images directes
et réciproques usuelles des cycles.
On a
$$
(a \otimes \text{id})^* \eta_Y^* [Y] =
(a \otimes \text{id})^* [\Delta_Y] =
(f \times \text{id})_*\Delta_Y = [\Gamma_f]
$$,
et, par l'autre chemin, on obtient
$$
(\text{id} \otimes a^t)^* \eta_X^* [X] =
(\text{id} \otimes a^t)^* [\Delta_X] =
(\text{id} \times f)^![\Delta_X] = [\Gamma_f]
$$
La dernière égalité résulte de
Homologie de Chow, lemme \ref{chow-lemma-lci-gysin-easy}.
Autrement dit, les deux chemins dans le diagramme donnent
l'élément de $\text{Corr}^d(X \times Y, \Spec(k))$
correspondant au cycle $\Gamma_f \subset X \times Y$.\par}
\end{proof}







\section{Motifs de Chow}
\label{weil-section-chow-motives}

\noindent
Fixons un corps de base $k$. Nous construisons dans cette section une catégorie
additive karoubienne $\mathbf{Q}$-linéaire $M_k$, munie
d'une structure monoïdale symétrique et d'un foncteur contravariant
$$
h : \{\text{schémas projectifs lisses sur }k\} \longrightarrow M_k
$$
qui transforme les produits en produits tensoriels et les réunions disjointes en sommes directes.
Notre construction sera caractérisée par le fait que $h$ se factorise par
la catégorie monoïdale symétrique dont les objets sont les variétés projectives lisses
et les morphismes les correspondances de degré $0$, de telle sorte que
l'image du projecteur $c_2$ sur $h(\mathbf{P}^1_k)$ de
l'exemple \ref{weil-example-decompose-P1} soit inversible dans $M_k$ ; voir
le lemme \ref{weil-lemma-characterize-motives}.
À la fin de cette section, nous montrerons que tout motif, c'est-à-dire
tout objet de $M_k$, admet un dual (à gauche) ; voir le lemme \ref{weil-lemma-dual-general}.

\medskip\noindent
Un {\it motif}, ou {\it motif de Chow}, sur $k$ sera un triplet
$(X, p, m)$, où
\begin{enumerate}
\item $X$ est un schéma projectif lisse sur $k$ ;
\item $p \in \text{Corr}^0(X, X)$ satisfait à $p \circ p = p$ ;
\item $m \in \mathbf{Z}$.
\end{enumerate}
Étant donné un second motif $(Y, q, n)$, nous appelons
{\it morphisme de motifs}, ou {\it morphisme de motifs de Chow},
un élément de
$$
\Hom((X, p, m), (Y, q, n)) =
q \circ \text{Corr}^{n - m}(X, Y) \circ p \subset \text{Corr}^{n - m}(X, Y)
$$
La composition des morphismes de motifs est définie au moyen
de la composition des correspondances définie ci-dessus.

\begin{lemma}
\label{weil-lemma-motives}
La catégorie $M_k$ dont les objets sont les motifs sur $k$ et dont les morphismes
sont les morphismes de motifs sur $k$ est une catégorie $\mathbf{Q}$-linéaire.
Il existe un foncteur contravariant
$$
h : \{\text{schémas projectifs lisses sur }k\} \longrightarrow M_k
$$
défini par $h(X) = (X, 1, 0)$ et $h(f) = [\Gamma_f]$.
\end{lemma}

\begin{proof}
Cela résulte immédiatement du lemme \ref{weil-lemma-contravariant-functor}.
\end{proof}

\begin{lemma}
\label{weil-lemma-Karoubian}
La catégorie $M_k$ est karoubienne.
\end{lemma}

\begin{proof}
Soit $M = (X, p, m)$ un motif et soit $a \in \Mor(M, M)$
un projecteur. Alors $a = a \circ a$ aussi bien dans $\Mor(M, M)$
que dans $\text{Corr}^0(X, X)$. Posons $N = (X, a, m)$.
Puisque $a = p \circ a \circ a$ dans $\text{Corr}^0(X, X)$,
on voit que $a : N \to M$ est un morphisme de $M_k$.
Supposons ensuite que $b : (Y, q, n) \to M$ soit un morphisme
tel que $(1 - a) \circ b = 0$. Alors $b = a \circ b$, et l'on a aussi
$b = b \circ q$. Par conséquent, $b$ est un morphisme $b : (Y, q, n) \to N$.
Le projecteur $1 - a$ admet donc un noyau, à savoir $N$,
et $M_k$ est karoubienne ; voir
Homologie, définition \ref{homology-definition-karoubian}.
\end{proof}

\noindent
Nous définissons un foncteur
$$
\otimes : M_k \times M_k \longrightarrow M_k
$$
Sur les objets, nous employons la formule
$$
(X, p, m) \otimes (Y, q, n) = (X \times Y, p \otimes q, m + n)
$$
Sur les morphismes, nous employons l'application
$$
\xymatrix{
\Mor((X, p, m), (Y, q, n)) \times
\Mor((X', p', m'), (Y', q', n')) \ar[d] \\
\Mor(
(X \times X', p \otimes p', m + m'),
(Y \times Y', q \otimes q', n + n'))
}
$$,
donnée par la règle $(a, a') \longmapsto a \otimes a'$, où
$\otimes$ est le produit des correspondances de la section \ref{weil-section-correspondences}.
Cela a un sens : par définition des morphismes de motifs,
on peut écrire $a = q \circ c \circ p$ et $a' = q' \circ c' \circ p'$,
avec $c \in \text{Corr}^{n - m}(X, Y)$ et
$c' \in \text{Corr}^{n' - m'}(X', Y')$ ;
on obtient alors
$$
a \otimes a' =
(q \circ c \circ p) \otimes (q' \circ c' \circ p') =
(q \otimes q') \circ (c \otimes c') \circ (p \otimes p')
$$,
qui est bien un morphisme de motifs de
$(X \times X', p \otimes p', m + m')$ dans
$(Y \times Y', q \otimes q', n + n')$.

\begin{lemma}
\label{weil-lemma-motives-monoidal}
La catégorie $M_k$, munie du produit tensoriel défini ci-dessus,
est monoïdale symétrique pour les contraintes évidentes d'associativité
et de commutativité, et pour l'unité $\mathbf{1} = (\Spec(k), 1, 0)$.
\end{lemma}

\begin{proof}
Cela résulte aisément du lemme \ref{weil-lemma-tensor-product}. Nous omettons les détails.
\end{proof}

\noindent
Les motifs $\mathbf{1}(n) = (\Spec(k), 1, n)$ seront utiles. Observons que
$$
\mathbf{1} = \mathbf{1}(0)
\quad\text{et}\quad
\mathbf{1}(n + m) = \mathbf{1}(n) \otimes \mathbf{1}(m)
$$
Le produit tensoriel par $\mathbf{1}(1)$ est donc une autoéquivalence de
la catégorie des motifs. Pour un motif $M$, nous écrivons parfois
$M(n) = M \otimes \mathbf{1}(n)$. Observons que, si $M = (X, p, m)$,
alors $M(n) = (X, p, m + n)$.

\begin{lemma}
\label{weil-lemma-inverse-h2}
Avec les notations de l'exemple \ref{weil-example-decompose-P1} :
\begin{enumerate}
\item
le motif $(X, c_0, 0)$ est isomorphe au motif
$\mathbf{1} = (\Spec(k), 1, 0)$ ;
\item
le motif $(X, c_2, 0)$ est isomorphe au motif
$\mathbf{1}(-1) = (\Spec(k), 1, -1)$.
\end{enumerate}
\end{lemma}

\begin{proof}
{\emergencystretch=2em Nous utiliserons le lemme \ref{weil-lemma-contravariant-functor} sans le mentionner de nouveau.
Le morphisme structural $X \to \Spec(k)$ donne une correspondance
$a \in \text{Corr}^0(\Spec(k), X)$. D'autre part, le point rationnel
$x$ est un morphisme $\Spec(k) \to X$, qui donne une correspondance
$b \in \text{Corr}^0(X, \Spec(k))$. On a $b \circ a = 1$ comme
correspondance sur $\Spec(k)$. La composée $a \circ b$ correspond
au graphe de la composée $X \to x \to X$, qui est
$c_0 = [x \times X]$. Ainsi, $a = a \circ b \circ a = c_0 \circ a$
et $b = b \circ a \circ b = b \circ c_0$.
En revenant aux définitions, on voit donc que
$a$ et $b$ sont des morphismes inverses l'un de l'autre
$a : (\Spec(k), 1, 0) \to (X, c_0, 0)$ et
$b : (X, c_0, 0) \to (\Spec(k), 1, 0)$.\par}

\medskip\noindent
Nous procédons exactement comme ci-dessus pour démontrer la seconde assertion.
Notons
$$
a' \in \text{Corr}^1(\Spec(k), X) = \CH^1(X)
$$
la classe du point $x$. Notons
$$
b' \in \text{Corr}^{-1}(X, \Spec(k)) = \CH_1(X)
$$
la classe de $[X]$. On a $b' \circ a' = 1$ comme correspondance
sur $\Spec(k)$, puisque $[x] \cdot [X] = [x]$ sur
$X = \Spec(k) \times X \times \Spec(k)$. Le calcul du
produit d'intersection $\text{pr}_{12}^*b' \cdot \text{pr}_{23}^*a'$
sur $X \times \Spec(k) \times X$ donne le cycle
$X \times \Spec(k) \times x$. Par conséquent,
la composée $a' \circ b'$ est égale à $c_2$ comme
correspondance sur $X$. Ainsi, $a' = a' \circ b \circ a' = c_2 \circ a'$
et $b' = b' \circ a' \circ b' = b' \circ c_2$. Rappelons que
$$
\Mor((\Spec(k), 1, -1), (X, c_2, 0)) =
c_2 \circ \text{Corr}^1(\Spec(k), X)
\subset
\text{Corr}^1(\Spec(k), X)
$$
et
$$
\Mor((X, c_2, 0), (\Spec(k), 1, -1)) =
\text{Corr}^{-1}(X, \Spec(k)) \circ c_2
\subset
\text{Corr}^{-1}(X, \Spec(k))
$$
On voit donc que $a'$ et $b'$ sont des morphismes inverses l'un de l'autre
$a' : (\Spec(k), 1, -1) \to (X, c_2, 0)$ et
$b' : (X, c_2, 0) \to (\Spec(k), 1, -1)$.
\end{proof}

\begin{remark}[Motifs de Lefschetz et de Tate]
\label{weil-remark-lefschetz-tate}
Soient $X = \mathbf{P}^1_k$ et $c_2$ comme dans l'exemple \ref{weil-example-decompose-P1}.
Dans la littérature, le motif $(X, c_2, 0)$ est parfois appelé
{\it motif de Lefschetz} et, selon les références, les notations
$L$, $\mathbf{L}$, $\mathbf{Q}(-1)$ ou $h^2(\mathbf{P}^1_k)$
peuvent le désigner. D'après le lemme \ref{weil-lemma-inverse-h2}, le motif de Lefschetz
est isomorphe à $\mathbf{1}(-1)$. Il est donc
inversible (Catégories, définition \ref{categories-definition-invertible}),
d'inverse
$\mathbf{1}(1)$. Le motif $\mathbf{1}(1)$ est parfois appelé
{\it motif de Tate} et, selon les références, les notations
$L^{-1}$, $\mathbf{L}^{-1}$, $\mathbf{T}$ ou $\mathbf{Q}(1)$
peuvent le désigner.
\end{remark}

\begin{lemma}
\label{weil-lemma-additive}
La catégorie $M_k$ est additive.
\end{lemma}

\begin{proof}
Soient $(Y, p, m)$ et $(Z, q, n)$ des motifs. Si $n = m$, une
somme directe est donnée par $(Y \amalg Z, p + q, m)$, avec des notations évidentes.
Nous omettons les détails.

\medskip\noindent
Supposons $n < m$. Soient $X$ et $c_2$ comme dans
l'exemple \ref{weil-example-decompose-P1}. Considérons alors
\begin{align*}
(Z, q, n)
& =
(Z, q, m) \otimes (\Spec(k), 1, -1) \otimes \ldots \otimes
(\Spec(k), 1, -1) \\
& \cong
(Z, q, m) \otimes (X, c_2, 0) \otimes \ldots \otimes (X, c_2, 0) \\
& \cong
(Z \times X^{m - n}, q \otimes c_2 \otimes \ldots \otimes c_2, m)
\end{align*},
où nous avons utilisé le lemme \ref{weil-lemma-inverse-h2}.
Nous sommes ainsi ramenés au cas traité dans le premier paragraphe.
\end{proof}

\begin{lemma}
\label{weil-lemma-decompose-P1}
Dans $M_k$, on a $h(\mathbf{P}^1_k) \cong \mathbf{1} \oplus \mathbf{1}(-1)$.
\end{lemma}

\begin{proof}
Cela résulte de l'exemple \ref{weil-example-decompose-P1} et du
lemme \ref{weil-lemma-inverse-h2}.
\end{proof}

\begin{lemma}
\label{weil-lemma-characterize-motives}
Soient $X$ et $c_2$ comme dans l'exemple \ref{weil-example-decompose-P1}.
Soit $\mathcal{C}$ une catégorie monoïdale symétrique karoubienne
$\mathbf{Q}$-linéaire. Tout foncteur $\mathbf{Q}$-linéaire
$$
F :
\left\{
\begin{matrix}
\text{schémas projectifs lisses sur }k\\
\text{les morphismes sont les correspondances de degré }0
\end{matrix}
\right\}
\longrightarrow
\mathcal{C}
$$
qui est monoïdal symétrique et tel que l'image de $F(c_2)$ dans
$F(X)$ soit un objet inversible se factorise de manière unique par un foncteur
$F : M_k \to \mathcal{C}$ monoïdal symétrique.
\end{lemma}

\begin{proof}
Notons $U$ l'objet inversible de $\mathcal{C}$ dont l'existence
est supposée dans l'énoncé du lemme. Nous prolongeons $F$ aux motifs en posant
$$
F(X, p, m) = \left(\text{l'image du
projecteur }F(p)\text{ dans }F(X)\right) \otimes U^{\otimes -m}
$$,
ce qui a un sens puisque $U$ est inversible et que $\mathcal{C}$
est karoubienne. Une propriété importante de ce choix est que
$F(X, c_2, 0) = U$. Observons que
\begin{align*}
F((X, p, m) \otimes (Y, q, n))
& =
F(X \times Y, p \otimes q, m + n) \\
& =
\left(\text{l'image de }F(p \otimes q)\text{ dans }F(X \times Y)\right)
\otimes U^{\otimes -m - n} \\
& =
F(X, p, m) \otimes F(Y, q, n)
\end{align*}
Notre règle est donc compatible avec les produits tensoriels
au niveau des objets (nous omettons les détails).

\medskip\noindent
Prolongeons ensuite $F$ aux morphismes de motifs. Supposons que
$$
a \in
\Hom((Y, p, m), (Z, q, n)) =
q \circ \text{Corr}^{n - m}(Y, Z) \circ p \subset \text{Corr}^{n - m}(Y, Z)
$$
soit un morphisme. Si $n = m$, alors $a$ est une correspondance de degré $0$
et l'on peut utiliser $F(a) : F(Y) \to F(Z)$ pour obtenir l'application voulue
$F(Y, p, m) \to F(Z, q, n)$. Si $n < m$, on a des identifications canoniques
\begin{align*}
s : F((Z, q, n))
& \to
F(Z, q, m) \otimes U^{m - n} \\
& \to
F(Z, q, m) \otimes F(X, c_2, 0) \otimes \ldots \otimes F(X, c_2, 0) \\
& \to
F((Z, q, m) \otimes (X, c_2, 0) \otimes \ldots \otimes (X, c_2, 0)) \\
& \to
F((Z \times X^{m - n}, q \otimes c_2 \otimes \ldots \otimes c_2, m))
\end{align*}
En effet, pour le premier isomorphisme, on utilise la définition de $F$ sur les motifs
donnée ci-dessus. Pour le deuxième, on utilise le choix de $U$. Pour le troisième,
on utilise la compatibilité de $F$ avec les produits tensoriels de motifs.
Le quatrième est la définition du produit tensoriel des motifs. D'autre part,
on a de même un isomorphisme
$$
\sigma : (Z, q, n) \to
(Z \times X^{m - n}, q \otimes c_2 \otimes \ldots \otimes c_2, m)
$$
(voir la démonstration du lemme \ref{weil-lemma-additive}). En composant $a$ avec cet isomorphisme,
on obtient
$$
\sigma \circ a \in
\Hom((Y, p, m),
(Z \times X^{m - n}, q \otimes c_2 \otimes \ldots \otimes c_2, m))
$$
En réunissant ces constructions, on obtient
$$
s^{-1} \circ F(\sigma \circ a) :
F(Y, p, m) \to
F(Z, q, n)
$$
Si $n > m$, on définit de même des isomorphismes
$$
t : F((Y, p, m)) \to
F((Y \times X^{n - m}, p \otimes c_2 \otimes \ldots \otimes c_2, n))
$$
et
$$
\tau : (Y, p, m)) \to
(Y \times X^{n - m}, p \otimes c_2 \otimes \ldots \otimes c_2, n)
$$,
et l'on pose $F(a) = F(a \circ \tau^{-1}) \circ t$.
Nous omettons la vérification que cette construction définit un foncteur
monoïdal symétrique.
\end{proof}

\begin{lemma}
\label{weil-lemma-dual}
Soit $X$ un schéma projectif lisse sur $k$, équidimensionnel
de dimension $d$. Alors $h(X)(d)$ est un dual à gauche de $h(X)$ dans $M_k$.
\end{lemma}

\begin{proof}
Nous utiliserons le lemme \ref{weil-lemma-composition-correspondences}
sans le mentionner de nouveau. On calcule
$$
\Hom(\mathbf{1}, h(X) \otimes h(X)(d)) =
\text{Corr}^d(\Spec(k), X \times X) = \CH^d(X \times X)
$$
On dispose ici de $\eta = [\Delta]$. D'autre part, on a
$$
\Hom(h(X)(d) \otimes h(X), \mathbf{1}) =
\text{Corr}^{-d}(X \times X, \Spec(k)) = \CH_d(X \times X)
$$,
où l'on dispose également de la classe $\epsilon = [\Delta]$
de la diagonale. La composition de la correspondance
$[\Delta] \otimes 1$ avec $1 \otimes [\Delta]$, dans un ordre ou dans l'autre,
est la correspondance $[\Delta] = 1$ dans $\text{Corr}^0(X, X)$ ; les diagrammes de
Catégories, définition \ref{categories-definition-dual},
sont donc commutatifs.
En effet, observons que
$$
[\Delta] \otimes 1 \in \text{Corr}^d(X, X \times X \times X) =
\CH^{2d}(X \times X \times X \times X)
$$
est donnée par la classe du cycle
$\text{pr}^{1234, -1}_{23}(\Delta) \cap \text{pr}^{1234, -1}_{14}(\Delta)$, avec
des notations évidentes. De même, la classe
$$
1 \otimes [\Delta] \in \text{Corr}^{-d}(X \times X \times X, X) =
\CH^{2d}(X \times X \times X \times X)
$$
est donnée par la classe du cycle
$\text{pr}^{1234, -1}_{23}(\Delta) \cap \text{pr}^{1234, -1}_{14}(\Delta)$.
La composée $(1 \otimes [\Delta]) \circ ([\Delta] \otimes 1)$
est, par définition, l'image directe par $\text{pr}^{12345}_{15, *}$
du produit d'intersection
$$
[\text{pr}^{12345, -1}_{23}(\Delta) \cap \text{pr}^{12345, -1}_{14}(\Delta)]
\cdot
[\text{pr}^{12345, -1}_{34}(\Delta) \cap \text{pr}^{12345, -1}_{15}(\Delta)]
=
[\text{petite diagonale dans } X^5]
$$,
qui est égale à $\Delta$, comme voulu. Nous omettons la démonstration
de la formule pour la composition dans l'autre ordre.
\end{proof}

\begin{lemma}
\label{weil-lemma-dual-general}
Tout objet de $M_k$ admet un dual à gauche.
\end{lemma}

\begin{proof}
Soit $M = (X, p, m)$ un objet de $M_k$. Alors $M$ est un facteur direct de
$(X, 1, m) = h(X)(m)$.
D'après Homologie, lemme \ref{homology-lemma-Karoubian-dual},
il suffit de montrer que
$h(X)(m) = h(X) \otimes \mathbf{1}(m)$ admet un dual.
Par construction, $\mathbf{1}(-m)$ est un dual à gauche de $\mathbf{1}(m)$.
Il suffit donc de montrer que $h(X)$ admet un dual à gauche ; voir
Catégories, lemme \ref{categories-lemma-tensor-dual}.
Soit $X = \coprod X_i$ la décomposition de $X$ en
composantes irréductibles. Alors $h(X) = \bigoplus h(X_i)$,
et il suffit de montrer que $h(X_i)$ admet un dual à gauche ; voir
Homologie, lemme \ref{homology-lemma-additive-dual}.
Cela résulte du lemme \ref{weil-lemma-dual}.
\end{proof}






\section{Groupes de Chow des motifs}
\label{weil-section-chow-groups-motives}

\noindent
Nous définissons les groupes de Chow d'un motif comme suit.

\begin{definition}
\label{weil-definition-chow-group-motives}
Soit $k$ un corps de base et soit $M = (X, p, m)$ un motif de Chow sur $k$.
Pour $i \in \mathbf{Z}$, nous définissons le {\it $i$-ième groupe de Chow de $M$}
par la formule
$$
\CH^i(M) = p\left(\CH^{i + m}(X) \otimes \mathbf{Q}\right)
$$
\end{definition}

\noindent
On a $\CH^i(h(X)) = \CH^i(X) \otimes \mathbf{Q}$
si $X$ est un schéma projectif lisse sur $k$.

\medskip\noindent
Observons que $\CH^i(-)$ est un foncteur de $M_k$ dans les espaces vectoriels sur $\mathbf{Q}$.
En effet, si $c : M \to N$ est un morphisme de motifs
$M = (X, p, m)$ et $N = (Y, q, n)$, alors $c$ est une correspondance de
degré $n - m$ de $X$ dans $Y$ ; l'image directe par $c$
(section \ref{weil-section-correspondences}) est donc une famille d'applications
$$
c_* :
\CH^{i + m}(X) \otimes \mathbf{Q}
\longrightarrow
\CH^{i + n}(Y) \otimes \mathbf{Q}
$$
Puisque $c = q \circ c \circ p$ par définition des morphismes de motifs,
on obtient bien
$$
c_* : \CH^i(M) \to \CH^i(N)
$$
pour tout $i \in \mathbf{Z}$. Cette construction est compatible avec la composition des
morphismes de motifs, d'après le lemme \ref{weil-lemma-composition-correspondences}.
Cette fonctorialité des groupes de Chow peut également se déduire
du lemme suivant.

\begin{lemma}
\label{weil-lemma-chow-groups-representable}
Soit $k$ un corps de base. Le foncteur $\CH^i(-)$ sur la catégorie
des motifs $M_k$ est représentable par $\mathbf{1}(-i)$ ; autrement dit,
on a
$$
\CH^i(M) = \Hom_{M_k}(\mathbf{1}(-i), M)
$$,
fonctoriellement en $M$ dans $M_k$.
\end{lemma}

\begin{proof}
Cela résulte immédiatement des définitions et du
lemme \ref{weil-lemma-composition-correspondences}.
\end{proof}

\noindent
Le lecteur peut prévoir que le
lemme \ref{weil-lemma-chow-groups-representable}, le lemme de Yoneda et
la dualité du lemme \ref{weil-lemma-dual} permettent d'obtenir le résultat suivant.

\begin{lemma}[Manin]
\label{weil-lemma-manin}
Soit $k$ un corps de base et soit $c : M \to N$ un morphisme de motifs.
Si, pour tout schéma projectif lisse $X$ sur $k$, l'application
$c \otimes 1 : M \otimes h(X) \to N \otimes h(X)$ induit un isomorphisme sur
les groupes de Chow, alors $c$ est un isomorphisme.
\end{lemma}

\begin{proof}
Tout objet $L$ de $M_k$ est un facteur direct de $h(X)(m)$ pour un schéma projectif lisse
$X$ sur $k$ et un entier $m \in \mathbf{Z}$ convenables. Observons que les groupes de Chow
de $M \otimes h(X)(m)$ sont les mêmes que ceux de $M \otimes h(X)$,
à un décalage des degrés près. Notre hypothèse entraîne donc
que $c \otimes 1 : M \otimes L \to N \otimes L$ induit un isomorphisme sur
les groupes de Chow pour tout objet $L$ de $M_k$. D'après le
lemme \ref{weil-lemma-chow-groups-representable},
l'application
$$
\Hom_{M_k}(\mathbf{1}, M \otimes L) \to
\Hom_{M_k}(\mathbf{1}, N \otimes L)
$$
est un isomorphisme pour tout $L$. Tout objet de $M_k$ admettant un dual à gauche
(lemme \ref{weil-lemma-dual-general}), il en résulte que
$$
\Hom_{M_k}(K, M) \to \Hom_{M_k}(K, N)
$$
est un isomorphisme pour tout objet $K$ de $M_k$ ; voir
Catégories, lemme \ref{categories-lemma-left-dual}.
On conclut par le lemme de Yoneda
(Catégories, lemme \ref{categories-lemma-yoneda}).
\end{proof}




\section{Formule du fibré projectif}
\label{weil-section-projective-space-bundle}

\noindent
Soit $k$ un corps de base et soit $X$ un schéma projectif lisse sur $k$.
Soit $\mathcal{E}$ un $\mathcal{O}_X$-module localement libre de rang $r$.
Nous adoptons la convention selon laquelle le {\it fibré projectif associé à
$\mathcal{E}$} est le morphisme
$$
\xymatrix{
P = \mathbf{P}(\mathcal{E}) =
\underline{\text{Proj}}_X(\text{Sym}^*(\mathcal{E}))
\ar[r]^-p
& X
}
$$
sur $X$, où $\mathcal{O}_P(1)$ est normalisé de telle sorte que
$p_*(\mathcal{O}_P(1)) = \mathcal{E}$. Rappelons que
$$
[\Gamma_p] \in \text{Corr}^0(X, P) \subset \CH^*(X \times P) \otimes \mathbf{Q}
$$
Voir l'exemple \ref{weil-example-graph-correspondence}.
Pour $i = 0, \ldots, r - 1$, considérons les correspondances
$$
c_i = c_1(\text{pr}_2^*\mathcal{O}_P(1))^i \cap [\Gamma_p]
\in \text{Corr}^i(X, P)
$$
Nous considérons désormais $c_i$ comme un morphisme $h(X)(-i) \to h(P)$.

\begin{lemma}[Formule du fibré projectif]
\label{weil-lemma-projective-space-bundle-formula}
Dans la situation ci-dessus, l'application
$$
\sum\nolimits_{i = 0, \ldots, r - 1} c_i :
\bigoplus\nolimits_{i = 0, \ldots, r - 1} h(X)(-i)
\longrightarrow
h(P)
$$
est un isomorphisme dans la catégorie des motifs.
\end{lemma}

\begin{proof}
D'après le lemme \ref{weil-lemma-manin}, il suffit de montrer que
notre application définit un isomorphisme sur les groupes de Chow des motifs
après produit par un schéma projectif lisse quelconque $Z$.
Observons que $P \times Z \to X \times Z$ est le fibré
projectif associé à l'image réciproque de $\mathcal{E}$ sur $X \times Z$.
L'assertion sur les groupes de Chow résulte donc de la formule
du fibré projectif donnée dans
Homologie de Chow, lemme \ref{chow-lemma-chow-ring-projective-bundle}.
En effet, l'image directe des cycles par $[\Gamma_p]$ est donnée par l'image réciproque
des cycles par $p$, d'après le lemme \ref{weil-lemma-functor-and-cycles} et
Homologie de Chow, lemme \ref{chow-lemma-lci-gysin-flat}. L'image directe
par $c_i$ envoie donc $\alpha$ sur $c_1(\mathcal{O}_P(1))^i \cap p^*\alpha$.
Nous omettons certains détails.
\end{proof}

\noindent
Dans la situation ci-dessus, pour $j = 0, \ldots, r - 1$, considérons
les correspondances
$$
c'_j = c_1(\text{pr}_1^*\mathcal{O}_P(1))^{r - 1 - j} \cap [\Gamma_p^t] \in
\text{Corr}^{-j}(P, X)
$$
Pour $i, j \in \{0, \ldots, r - 1\}$, on a
$$
c'_j \circ c_i =
\text{pr}_{13, *}\left(
c_1(\text{pr}_2^*\mathcal{O}_P(1))^{i + r - 1 - j} \cap
(\text{pr}_{12}^*[\Gamma_p] \cdot \text{pr}_{23}^*[\Gamma_p^t])
\right)
$$
Les cycles $\text{pr}_{12}^{-1}\Gamma_p$ et
$\text{pr}_{23}^{-1}\Gamma_p^t$ se coupent transversalement,
et leur intersection est l'image de
$(p, 1, p) : P \to X \times P \times X$.
Observons que les fibres de
$(p, p) = \text{pr}_{13} \circ (p, 1, p) : P \to X \times X$
sont de dimension $r - 1$. On en déduit immédiatement
$c'_j \circ c_i = 0$ pour $i + r - 1 - j < r - 1$, autrement dit lorsque $i < j$.
D'autre part, d'après la formule du fibré projectif
(Homologie de Chow, lemme \ref{chow-lemma-chow-ring-projective-bundle}),
le cycle $c_1(\mathcal{O}_P(1))^{r - 1} \cap [P]$ a pour image
$[X]$ dans $X$. Ainsi, pour $i = j$, l'image directe ci-dessus
donne la classe de la diagonale ; on a donc
$$
c'_i \circ c_i = 1 \in \text{Corr}^0(X, X)
$$
pour tout $i \in \{0, \ldots, r - 1\}$. La matrice
de la composée
$$
\bigoplus h(X)(-i)
\xrightarrow{\bigoplus c_i}
h(P)
\xrightarrow{\bigoplus c'_j}
\bigoplus h(X)(-j)
$$
est donc inversible (elle est triangulaire supérieure, avec des $1$ sur la diagonale).
La formule du fibré projectif
(lemme \ref{weil-lemma-projective-space-bundle-formula})
entraîne que la composée dans l'autre ordre est elle aussi
inversible, mais une démonstration directe semble un peu plus difficile.

\begin{lemma}
\label{weil-lemma-diagonal-projective-bundle}
Soit $p : P \to X$ comme dans le lemme \ref{weil-lemma-projective-space-bundle-formula}.
La classe $[\Delta_P]$ de la diagonale de $P$ dans $\CH^*(P \times P)$
s'écrit
$$
[\Delta_P] =
\left(\sum\nolimits_{i = 0, \ldots, r - 1}
{r - 1 \choose i} c_{r - 1 - i}(\text{pr}_1^*\mathcal{S}^\vee) \cap
c_1(\text{pr}_2^*\mathcal{O}_P(1))^i\right)
\cap
(p \times p)^*[\Delta_X]
$$,
où $\mathcal{S}$ est le noyau de la surjection canonique
$p^*\mathcal{E} \to \mathcal{O}_P(1)$.
\end{lemma}

\begin{proof}
Observons que $(p \times p)^*[\Delta_X] = [P \times_X P]$.
Puisque $\Delta_P \subset P \times_X P \subset P \times P$
et que le cap-produit par les classes de Chern commute à l'image directe propre
(Homologie de Chow, lemme \ref{chow-lemma-pushforward-cap-cj}),
il suffit de montrer que la classe de
$\Delta_P \subset P \times_X P$ dans $\CH^*(P \times_X P)$
est égale à
$$
\left(\sum\nolimits_{i = 0, \ldots, r - 1}
{r - 1 \choose i} c_{r - 1 - i}(q_1^*\mathcal{S}^\vee) \cap
c_1(q_2^*\mathcal{O}_P(1))^i\right)
\cap
[P \times_X P]
$$,
où $q_i : P \times_X P \to P$, pour $i = 1, 2$, sont les projections.
Posons $q = p \circ q_1 = p \circ q_2 : P \times_X P \to X$.
Considérons les applications
$$
q_1^*\mathcal{S} \otimes q_2^*\mathcal{O}_P(-1) \to
q^*\mathcal{E} \otimes q^*\mathcal{E}^\vee \to
\mathcal{O}_{P \times_X P}
$$,
où la dernière flèche est l'image réciproque par $q$ de l'application d'évaluation
$\mathcal{E} \otimes_{\mathcal{O}_X} \mathcal{E}^\vee \to \mathcal{O}_X$.
La source de la composée est un module localement libre de rang $r - 1$,
et un calcul local montre que cette application s'annule exactement le long de
$\Delta_P$. D'après Homologie de Chow, lemme \ref{chow-lemma-top-chern-class},
la classe $[\Delta_P]$ est la classe de Chern de degré maximal du dual
$$
q_1^*\mathcal{S}^\vee \otimes q_2^*\mathcal{O}_P(1)
$$
Le résultat voulu découle de Homologie de Chow, lemme
\ref{chow-lemma-chern-classes-E-tensor-L}.
\end{proof}








\section{Théories cohomologiques de Weil classiques}
\label{weil-section-axioms-classical}

\noindent
Nous définissons dans cette section ce que nous appellerons une théorie
cohomologique de Weil classique. Il s'agit exactement de ce qui est appelé
une théorie cohomologique de Weil dans \cite[section 1.2]{Kleiman-cycles}.

\medskip\noindent
Fixons un corps algébriquement clos $k$ (le corps de base).
Dans cette section, {\it variété} signifie variété sur $k$ ; voir
Variétés, section \ref{varieties-section-varieties}.
Fixons un corps $F$ de caractéristique $0$ (le corps de coefficients).
Une théorie cohomologique de Weil est définie par des données (D1), (D2) et (D3)
soumises aux axiomes (A), (B) et (C).

\medskip\noindent
Les données sont les suivantes :
\begin{enumerate}
\item[(D1)] Un foncteur contravariant $H^*$ de la catégorie
des variétés projectives lisses dans la catégorie des
$F$-algèbres graduées commutatives.
\item[(D2)] Pour toute variété projective lisse $X$,
un homomorphisme de groupes $\gamma : \CH^i(X) \to H^{2i}(X)$.
\item[(D3)] Pour toute variété projective lisse $X$ de dimension $d$,
une application $\int_X : H^{2d}(X) \to F$.
\end{enumerate}
Précisons le sens de ces données par quelques remarques
et introduisons la terminologie correspondante.

\medskip\noindent
Remarques sur (D1). Pour une variété projective lisse $X$,
on appelle $H^*(X)$ la {\it cohomologie} de $X$. Pour un morphisme
$f : X \to Y$ de variétés projectives lisses, on note
$f^* : H^*(Y) \to H^*(X)$ l'application $H^*(f)$ et on l'appelle
l'{\it application d'image réciproque}.

\medskip\noindent
Remarques sur (D2). L'application $\gamma$ est appelée l'{\it application classe de cycle}.
On appelle $\gamma(\alpha)$ la {\it classe de cohomologie} de $\alpha$.
Si $Z \subset Y \subset X$ sont des sous-schémas fermés, avec $Y$ et $X$
des variétés projectives lisses et $Z$ intègre, alors $[Z]$ peut
désigner la classe du cycle $[Z]$ dans $\CH^*(Y)$ ou dans $\CH^*(X)$.
Dans ce cas, la notation $\gamma([Z])$ est ambiguë, et son sens
doit être déduit du contexte.

\medskip\noindent
Remarques sur (D3). L'application $\int_X$ est parfois appelée
l'{\it application trace} et est parfois notée $\text{Tr}_X$.

\medskip\noindent
Le premier axiome est souvent appelé {\it dualité de Poincar\'e} :
\begin{enumerate}
\item[(A)] Soit $X$ une variété projective lisse de dimension $d$. Alors :
\begin{enumerate}
\item $\dim_F H^i(X) < \infty$ pour tout $i$ ;
\item $H^i(X) \times H^{2d - i}(X) \rightarrow H^{2d}(X) \rightarrow F$
est un accouplement parfait pour tout $i$, la dernière
application étant l'application trace $\int_X$ ;
\item $H^i(X) = 0$ sauf si $i \in [0, 2d]$ ;
\item $\int_X : H^{2d}(X) \to F$ est un isomorphisme.
\end{enumerate}
\end{enumerate}
Soit $f : X \to Y$ un morphisme de variétés projectives lisses,
avec $\dim(X) = d$ et $\dim(Y) = e$. À l'aide de la dualité de Poincar\'e,
on peut définir une {\it image directe}
$$
f_* : H^{2d - i}(X) \longrightarrow H^{2e - i}(Y)
$$
comme l'application contragrédiente de l'application linéaire $f^* : H^i(Y) \to H^i(X)$.
Explicitement, pour $a \in H^{2d - i}(X)$, l'élément $f_*a \in H^{2e - i}(Y)$
est caractérisé par
$$
\int_X f^*b \cup a = \int_Y b \cup f_*a
$$
pour tout $b \in H^i(Y)$.

\begin{lemma}
\label{weil-lemma-pushforward-classical}
Supposons que les données (D1) et (D3) satisfassent à (A). Pour $f : X \to Y$
un morphisme de variétés projectives lisses, on a
$f_*(f^*b \cup a) = b \cup f_*a$. Si $g : Y \to Z$ est un second morphisme
de variétés projectives lisses, alors $g_* \circ f_* = (g \circ f)_*$.
\end{lemma}

\begin{proof}
La première égalité résulte de
$$
\int_Y c \cup b \cup f_*a =
\int_X f^*c \cup f^*b \cup a =
\int_Y c \cup f_*(f^*b \cup a).
$$
La seconde résulte de
$$
\int_Z c \cup (g \circ f)_*a = \int_X (g \circ f)^*c \cup a =
\int_X f^* g^* c \cup a = \int_Y g^*c \cup f_*a = \int_Z c \cup g_*f_*a
$$
Cela achève la démonstration.
\end{proof}

\noindent
Le deuxième axiome exprime que $H^*$ respecte la structure monoïdale
donnée par les produits, au moyen de la {\it formule de K\"unneth} :
\begin{enumerate}
\item[(B)] Soient $X$ et $Y$ des variétés projectives lisses. L'application
$$
H^*(X) \otimes_F H^*(Y) \to H^*(X \times Y),\quad
a \otimes b \mapsto \text{pr}_1^*a \cup \text{pr}_2^*b
$$
est un isomorphisme.
\end{enumerate}
Le troisième axiome porte sur les applications classes de cycles :
\begin{enumerate}
\item[(C)] Les applications classes de cycles satisfont aux règles suivantes :
\begin{enumerate}
\item pour un morphisme $f : X \to Y$ de variétés projectives lisses,
on a $\gamma(f^!\beta) = f^*\gamma(\beta)$ pour $\beta \in \CH^*(Y)$ ;
\item pour un morphisme $f : X \to Y$ de variétés projectives lisses, on a
$\gamma(f_*\alpha) = f_*\gamma(\alpha)$ pour $\alpha \in \CH^*(X)$ ;
\item pour toute variété projective lisse $X$, on a
$\gamma(\alpha \cdot \beta) = \gamma(\alpha) \cup \gamma(\beta)$
pour $\alpha, \beta \in \CH^*(X)$ ;
\item $\int_{\Spec(k)} \gamma([\Spec(k)]) = 1$.
\end{enumerate}
\end{enumerate}

\begin{remark}
\label{weil-remark-replace-cup-product-classical}
Soit $X$ une variété projective lisse. On obtient des applications
$$
H^*(X) \otimes_F H^*(X) \longrightarrow H^*(X \times X)
\xrightarrow{\Delta^*} H^*(X)
$$,
où la première flèche est celle de l'axiome (B) et $\Delta^*$
est l'image réciproque par le morphisme diagonal $\Delta : X \to X \times X$.
La composée est le cup-produit, puisque l'image réciproque est un homomorphisme d'algèbres
et que $\text{pr}_i \circ \Delta = \text{id}$.
D'autre part, pour des cycles $\alpha, \beta$ sur $X$,
le produit d'intersection est défini par la formule
$$
\alpha \cdot \beta =
\Delta^!(\alpha \times \beta)
$$
Autrement dit, $\alpha \cdot \beta$ est l'image réciproque du
produit extérieur $\alpha \times \beta$ sur $X \times X$ par
la diagonale. Notons aussi que
$\alpha \times \beta = \text{pr}_1^*\alpha \cdot \text{pr}_2^*\beta$
dans $\CH^*(X \times X)$ (nous omettons la démonstration). Ainsi, sous l'axiome (C)(a),
l'axiome (C)(c) équivaut à la compatibilité de $\gamma$
avec le produit extérieur, au sens où
$\gamma(\alpha \times \beta)$ est égal à
$\text{pr}_1^*\gamma(\alpha) \cup \text{pr}_2^*\gamma(\beta)$.
C'est ainsi que l'axiome (C)(c) est formulé dans \cite{Kleiman-cycles}.
\end{remark}

\begin{definition}
\label{weil-definition-weil-cohomology-theory-classical}
Soit $k$ un corps algébriquement clos.
Soit $F$ un corps de caractéristique $0$.
Une {\it théorie cohomologique de Weil classique} sur $k$ à coefficients dans $F$
est définie par des données (D1), (D2) et (D3) satisfaisant
à la dualité de Poincar\'e, à la formule de K\"unneth et à la compatibilité
avec les classes de cycles, plus précisément aux axiomes (A), (B) et (C).
\end{definition}

\noindent
Établissons quelques conséquences élémentaires.

\begin{lemma}
\label{weil-lemma-degrees-cycles-classical}
Soit $H^*$ une théorie cohomologique de Weil classique
(définition \ref{weil-definition-weil-cohomology-theory-classical}).
Soit $X$ une variété projective lisse de dimension $d$. Le diagramme
$$
\xymatrix{
\CH^d(X) \ar[r]_-\gamma \ar@{=}[d] &
H^{2d}(X) \ar[d]^{\int_X} \\
\CH_0(X) \ar[r]^\deg & F
}
$$
est commutatif, où $\deg : \CH_0(X) \to \mathbf{Z}$ est le degré des
zéro-cycles étudié dans Homologie de Chow, section
\ref{chow-section-degree-zero-cycles}.
\end{lemma}

\begin{proof}
Le résultat vaut pour $\Spec(k)$ d'après l'axiome (C)(d). Soit $x : \Spec(k) \to X$
un point fermé de $X$. Alors $\gamma([x]) = x_*\gamma([\Spec(k)])$
dans $H^{2d}(X)$, d'après l'axiome (C)(b). Par conséquent, $\int_X \gamma([x]) = 1$
par définition de $x_*$.
\end{proof}

\begin{lemma}
\label{weil-lemma-trace-product-classical}
Soit $H^*$ une théorie cohomologique de Weil classique
(définition \ref{weil-definition-weil-cohomology-theory-classical}).
Soient $X$ et $Y$ des variétés projectives lisses.
Alors $\int_{X \times Y} = \int_X \otimes \int_Y$.
\end{lemma}

\begin{proof}
Posons $\dim(X) = d$ et $\dim(Y) = e$. D'après l'axiome (B), on a
$H^{2d + 2e}(X \times Y) = H^{2d}(X) \otimes H^{2e}(Y)$,
et cet espace est de dimension $1$ d'après l'axiome (A)(d).
D'après le lemme \ref{weil-lemma-degrees-cycles-classical},
cet espace vectoriel de dimension $1$ est engendré par
la classe $\gamma([x \times y])$ d'un point fermé $(x, y)$, et
$\int_{X \times Y} \gamma([x \times y]) = 1$.
Puisque $\gamma([x \times y]) = \gamma([x]) \otimes \gamma([y])$
d'après les axiomes (C)(a) et (C)(c), et puisque $\int_X \gamma([x]) = 1$ et
$\int_Y \gamma([y]) = 1$, on obtient la conclusion.
\end{proof}

\begin{lemma}
\label{weil-lemma-pr2star-classical}
Soit $H^*$ une théorie cohomologique de Weil classique
(définition \ref{weil-definition-weil-cohomology-theory-classical}).
Soient $X$ et $Y$ des variétés projectives lisses.
Alors $\text{pr}_{2, *} : H^*(X \times Y) \to H^*(Y)$
envoie $a \otimes b$ sur $(\int_X a) b$.
\end{lemma}

\begin{proof}
Cela équivaut au résultat du lemme \ref{weil-lemma-trace-product-classical}.
\end{proof}

\begin{lemma}
\label{weil-lemma-class-diagonal-classical}
Soit $H^*$ une théorie cohomologique de Weil classique
(définition \ref{weil-definition-weil-cohomology-theory-classical}).
Soit $X$ une variété projective lisse de dimension $d$.
Choisissons une base $e_{i, j}, j = 1, \ldots, \beta_i$ de $H^i(X)$ sur $F$.
À l'aide de la formule de K\"unneth, écrivons
$$
\gamma([\Delta]) =
\sum\nolimits_{i = 0, \ldots, 2d}
\sum\nolimits_j e_{i, j} \otimes e'_{2d - i , j}
\quad\text{dans}\quad
\bigoplus\nolimits_i H^i(X) \otimes_F H^{2d - i}(X)
$$
avec $e'_{2d - i, j} \in H^{2d - i}(X)$.
Alors $\int_X e_{i, j} \cup e'_{2d - i, j'} = (-1)^i\delta_{jj'}$.
\end{lemma}

\begin{proof}
Rappelons que $\Delta^* : H^*(X \times X) \to H^*(X)$ est l'application
donnée par le cup-produit $H^*(X) \otimes_F H^*(X) \to H^*(X)$ ; voir
la remarque \ref{weil-remark-replace-cup-product-classical}. D'autre part, on a
$\gamma([\Delta]) = \Delta_*\gamma([X]) = \Delta_*1$ d'après
l'axiome (C)(b) et le fait que $\gamma([X]) = 1$. En effet,
$[X] \cdot [X] = [X]$ ; ainsi, par l'axiome (C)(c), la classe de cohomologie
$\gamma([X])$ vaut $0$ ou $1$ dans l'algèbre de dimension $1$ sur $F$, $H^0(X)$ ;
nous avons aussi utilisé ici les axiomes (A)(d) et (A)(b).
Mais $\gamma([X])$ ne peut être nulle, car $[X] \cdot [x] = [x]$
pour un point fermé $x$ de $X$, et $\gamma([x])$
est non nulle d'après le lemme \ref{weil-lemma-degrees-cycles-classical}.
Par conséquent,
$$
\int_{X \times X} \gamma([\Delta]) \cup a \otimes b =
\int_{X \times X} \Delta_*1 \cup a \otimes b =
\int_X a \cup b
$$,
par définition de $\Delta_*$. D'autre part, on a
$$
\int_{X \times X} (\sum e_{i, j} \otimes e'_{2d -i , j}) \cup a \otimes b =
\sum (\int_X a \cup e_{i, j})(\int_X e'_{2d - i, j} \cup b)
$$,
d'après le lemme \ref{weil-lemma-trace-product-classical} ; notons que nous avons
effectué deux échanges, de sorte que le signe est $1$.
Ainsi, si l'on choisit $a$ de sorte que $\int_X a \cup e_{i, j} = 1$
et que tous les autres accouplements soient nuls, on en déduit
$\int_X e'_{2d - i, j} \cup b = \int_X a \cup b$ pour tout $b$, c'est-à-dire
$e'_{2d - i, j} = a$. Le lemme est démontré.
\end{proof}

\begin{lemma}
\label{weil-lemma-square-diagonal-classical}
Soit $H^*$ une théorie cohomologique de Weil classique
(définition \ref{weil-definition-weil-cohomology-theory-classical}).
Soit $X$ une variété projective lisse. On a
$$
\sum\nolimits_{i = 0, \ldots, 2\dim(X)} (-1)^i\dim_F H^i(X) =
\deg([\Delta] \cdot [\Delta]) = \deg(c_d(\mathcal{T}_X) \cap [X])
$$
\end{lemma}

\begin{proof}
Démontrons l'égalité de droite. On a
$[\Delta] \cdot [\Delta] = \Delta_*(\Delta^![\Delta])$
(Homologie de Chow, lemme \ref{chow-lemma-intersect-regularly-embedded}).
Puisque $\Delta_*$ préserve les degrés des $0$-cycles, il suffit de calculer
le degré de $\Delta^![\Delta]$. La classe $\Delta^![\Delta]$ est donnée
par le cap-produit de $[\Delta]$ avec
la classe de Chern de degré maximal du faisceau normal de $\Delta \subset X \times X$
(Homologie de Chow, lemme \ref{chow-lemma-gysin-fundamental}).
Puisque le faisceau conormal de $\Delta$ est $\Omega_{X/k}$
(Morphismes, lemme \ref{morphisms-lemma-differentials-diagonal}),
le faisceau normal est égal au faisceau tangent
$\mathcal{T}_X = \SheafHom_{\mathcal{O}_X}(\Omega_{X/k}, \mathcal{O}_X)$,
comme voulu.

\medskip\noindent
Démontrons l'égalité de gauche. D'après le lemme \ref{weil-lemma-degrees-cycles-classical}, on a
\begin{align*}
\deg([\Delta] \cdot [\Delta])
& =
\int_{X \times X} \gamma([\Delta]) \cup \gamma([\Delta]) \\
& =
\int_{X \times X} \Delta_*1 \cup \gamma([\Delta]) \\
& =
\int_{X \times X} \Delta_*(\Delta^*\gamma([\Delta])) \\
& =
\int_X \Delta^*\gamma([\Delta])
\end{align*}
Écrivons $\gamma([\Delta]) = \sum  e_{i, j} \otimes e'_{2d - i , j}$
comme dans le lemme \ref{weil-lemma-class-diagonal-classical}.
En rappelant que $\Delta^*$ est donné par le cup-produit, on obtient
$$
\int_X \sum\nolimits_{i, j} e_{i, j} \cup e'_{2d - i, j} =
\sum\nolimits_{i, j} \int_X e_{i, j} \cup e'_{2d - i, j} =
\sum\nolimits_{i, j} (-1)^i = \sum (-1)^i\beta_i
$$,
comme voulu.
\end{proof}




\noindent
Relions maintenant les théories cohomologiques de Weil classiques aux motifs.

\begin{lemma}
\label{weil-lemma-from-functor-to-weil-classical}
Soit $k$ un corps algébriquement clos et soit $F$ un corps de
caractéristique $0$. Considérons un foncteur $\mathbf{Q}$-linéaire
$$
G : M_k \longrightarrow \text{espaces vectoriels gradués sur }F\text{}
$$
qui est monoïdal symétrique et tel que $G(\mathbf{1}(1))$
ne soit non nul qu'en degré $-2$. On obtient alors des données (D1), (D2), (D3)
satisfaisant à tous les axiomes (A), (B), (C), sauf peut-être (A)(c) et (A)(d).
\end{lemma}

\begin{proof}
On obtient un foncteur contravariant de la catégorie des variétés
projectives lisses dans celle des espaces vectoriels gradués sur $F$
en posant $H^*(X) = G(h(X))$. Par hypothèse, on dispose d'un isomorphisme
canonique
$$
H^*(X \times Y) = G(h(X \times Y)) = G(h(X) \otimes h(Y)) =
G(h(X)) \otimes G(h(Y)) = H^*(X) \otimes H^*(Y)
$$,
compatible avec les images réciproques. L'image réciproque par la diagonale
$\Delta : X \to X \times X$ donne une application canonique
$$
H^*(X) \otimes H^*(X) = H^*(X \times X) \to H^*(X)
$$
d'espaces vectoriels gradués, compatible avec les images réciproques.
Cela définit une structure fonctorielle de $F$-algèbre graduée sur
$H^*(X)$. Puisque $\Delta$ commute à la contrainte de commutativité
$h(X) \otimes h(X) \to h(X) \otimes h(X)$ (qui échange les facteurs),
que $G$ est un foncteur monoïdal symétrique (donc compatible avec
les contraintes de commutativité), et compte tenu de notre convention dans
Homologie, exemple \ref{homology-example-graded-vector-spaces},
on en déduit que $H^*(X)$ est une algèbre graduée
commutative ! Nous obtenons ainsi la donnée (D1).

\medskip\noindent
Puisque $\mathbf{1}(1)$ est inversible dans la catégorie des motifs,
$G(\mathbf{1}(1))$ est inversible dans la catégorie des
espaces vectoriels gradués sur $F$. Ainsi, $\sum_i \dim_F G^i(\mathbf{1}(1)) = 1$.
Par hypothèse, le seul degré où cet objet est non nul est $-2$, et l'on peut
choisir un isomorphisme $F[2] \to G(\mathbf{1}(1))$ d'espaces vectoriels gradués sur $F$.
Ici et dans la suite, $F[n]$ désigne l'espace vectoriel gradué sur $F$ ayant
$F$ en degré $-n$ et zéro ailleurs. La compatibilité avec
les produits tensoriels donne, pour tout $n \in \mathbf{Z}$, un isomorphisme
$F[2n] \to G(\mathbf{1}(n))$ compatible avec les produits tensoriels.

\medskip\noindent
Soit $X$ une variété projective lisse. D'après
le lemme \ref{weil-lemma-composition-correspondences}, on a
$$
\CH^r(X) \otimes \mathbf{Q} = \text{Corr}^r(\Spec(k), X) =
\Hom(\mathbf{1}(-r), h(X))
$$
En appliquant le foncteur $G$, on obtient
$$
\gamma :
\CH^r(X) \otimes \mathbf{Q} \longrightarrow
\Hom(G(\mathbf{1}(-r)), H^*(X)) = H^{2r}(X)
$$
C'est la donnée (D2).

\medskip\noindent
Soit $X$ une variété projective lisse de dimension $d$. D'après
le lemme \ref{weil-lemma-composition-correspondences}, on a
$$
\Mor(h(X)(d), \mathbf{1}) = \Mor((X, 1, d), (\Spec(k), 1, 0)) =
\text{Corr}^{-d}(X, \Spec(k)) = \CH_d(X)
$$
La classe du cycle $[X]$ dans $\CH_d(X)$ définit donc un morphisme
$h(X)(d) \to \mathbf{1}$. En appliquant $G$, on obtient
$$
H^*(X) \otimes F[-2d] = G(h(X)(d)) \longrightarrow G(\mathbf{1}) = F
$$
Cette application est nulle sauf en degré $0$, où l'on obtient
$\int_X : H^{2d}(X) \to F$. C'est la donnée (D3).

\medskip\noindent
Soit $X$ une variété projective lisse de dimension $d$.
D'après le lemme \ref{weil-lemma-dual},
$h(X)(d)$ est un dual à gauche de $h(X)$. Par conséquent,
$G(h(X)(d)) = H^*(X) \otimes F[-2d]$ est un dual à gauche de
$H^*(X)$ dans la catégorie des espaces vectoriels gradués sur $F$.
D'après Homologie, lemme \ref{homology-lemma-left-dual-graded-vector-spaces},
on a $\sum_i \dim_F H^i(X) < \infty$, et
$\epsilon : h(X)(d) \otimes h(X) \to \mathbf{1}$ donne
des accouplements non dégénérés $H^{2d - i}(X) \otimes_F H^i(X) \to F$.
Dans la démonstration du lemme \ref{weil-lemma-dual}, nous avons vu que
$\epsilon$ est donné par $[\Delta]$ au moyen des identifications
$$
\Hom(h(X)(d) \otimes h(X), \mathbf{1}) =
\text{Corr}^{-d}(X \times X, \Spec(k)) =
\CH_d(X \times X)
$$
Ainsi, $\epsilon$ est la composée de $[X] : h(X)(d) \to \mathbf{1}$
et de $h(\Delta)(d) : h(X)(d) \otimes h(X) \to h(X)(d)$. Il s'ensuit
que les accouplements ci-dessus sont donnés par le cup-produit suivi de
$\int_X$. Cela démontre les parties (a) et (b) de l'axiome (A).

\medskip\noindent
L'axiome (B) résulte de la compatibilité de $G$ avec les structures
tensorielles, supposée par hypothèse, et de notre construction du cup-produit ci-dessus.

\medskip\noindent
Axiome (C). Notre construction de $\gamma$ prend un cycle $\alpha$ sur $X$,
l'interprète comme une correspondance $a$ de $\Spec(k)$ dans $X$ d'un certain degré,
puis lui applique $G$. Si $f : Y \to X$ est un morphisme de variétés projectives
lisses, alors $f^!\alpha$ est l'image directe (!) de $\alpha$
par la correspondance $[\Gamma_f]$ de $X$ dans $Y$ ; voir
le lemme \ref{weil-lemma-functor-and-cycles}. Ainsi,
$f^!\alpha$, considéré comme une correspondance de $\Spec(k)$ dans $Y$,
est égal à $a \circ [\Gamma_f]$ ; voir
le lemme \ref{weil-lemma-composition-correspondences}.
Puisque $G$ est un foncteur, on en déduit que
$\gamma$ est compatible avec les images réciproques, ce qui est l'axiome (C)(a).

\medskip\noindent
Soit $f : Y \to X$ un morphisme de variétés projectives lisses et
soit $\beta \in \CH^r(Y)$ un cycle sur $Y$. Il faut montrer que
$$
\int_Y \gamma(\beta) \cup f^*c = \int_X \gamma(f_*\beta) \cup c
$$
pour tout $c \in H^*(X)$. Soient $a, a^t, \eta_X, \eta_Y, [X], [Y]$
comme dans le lemme \ref{weil-lemma-prep-dual}.
Soit $b$ le cycle $\beta$ considéré comme une correspondance de $\Spec(k)$ dans $Y$
de degré $r$. Alors $f_*\beta$, considéré comme une correspondance de
$\Spec(k)$ dans $X$, est égal à $a^t \circ b$ ; voir
les lemmes \ref{weil-lemma-functor-and-cycles} et
\ref{weil-lemma-composition-correspondences}.
L'égalité affichée ci-dessus sera vérifiée si l'on montre que
$$
h(X) = \mathbf{1} \otimes h(X)
\xrightarrow{b \otimes 1}
h(Y)(r) \otimes h(X)
\xrightarrow{1 \otimes a}
h(Y)(r) \otimes h(Y)
\xrightarrow{\eta_Y}
h(Y)(r)
\xrightarrow{[Y]}
\mathbf{1}(r - e)
$$
est égal à
$$
h(X) = \mathbf{1} \otimes h(X)
\xrightarrow{a^t \circ b \otimes 1}
h(X)(r + d - e) \otimes h(X)
\xrightarrow{\eta_X}
h(X)(r + d - e)
\xrightarrow{[X]}
\mathbf{1}(r - e)
$$
Cela résulte immédiatement du lemme \ref{weil-lemma-prep-dual}.
L'axiome (C)(b) est donc vérifié.

\medskip\noindent
Pour démontrer l'axiome (C)(c), nous utilisons la discussion de
la remarque \ref{weil-remark-replace-cup-product-classical}.
Il suffit donc de démontrer que $\gamma$ est compatible avec
les produits extérieurs. Soient $X$ et $Y$ des variétés projectives lisses
et soient $\alpha$ et $\beta$ des cycles sur ces variétés. Notons
$a$ et $b$ les correspondances associées de $\Spec(k)$ dans
$X$ et $Y$. Alors $\alpha \times \beta$ correspond à
la correspondance $a \otimes b$ de $\Spec(k)$ dans $X \otimes Y = X \times Y$.
La propriété voulue résulte donc de la compatibilité de $G$
avec les structures tensorielles de part et d'autre.

\medskip\noindent
L'axiome (C)(d) résulte de ce que le cycle $[\Spec(k)]$
correspond au morphisme identique de $h(\Spec(k))$.
Cela achève la démonstration du lemme.
\end{proof}

\begin{lemma}
\label{weil-lemma-from-weil-to-functor-classical}
Soit $k$ un corps algébriquement clos et soit $F$ un corps de
caractéristique $0$. Soit $H^*$ une théorie cohomologique de Weil classique.
On peut alors construire un foncteur $\mathbf{Q}$-linéaire
$$
G : M_k \longrightarrow \text{espaces vectoriels gradués sur }F\text{}
$$
monoïdal symétrique tel que $H^*(X) = G(h(X))$.
\end{lemma}

\begin{proof}
D'après le lemme \ref{weil-lemma-characterize-motives}, il suffit de construire un foncteur
$G$ sur la catégorie des schémas projectifs lisses sur $k$,
dont les morphismes sont les correspondances de degré $0$, tel que
l'image de $G(c_2)$ dans $G(\mathbf{P}^1)$ soit un espace vectoriel gradué
inversible sur $F$.
Tout schéma projectif lisse étant canoniquement une réunion
disjointe de variétés projectives lisses, il suffit de construire
$G$ sur la catégorie dont les objets sont les variétés projectives lisses
et les morphismes les correspondances de degré $0$.
(Nous omettons certains détails.)

\medskip\noindent
Pour une variété projective lisse $X$, nous posons $G(X) = H^*(X)$.

\medskip\noindent
Pour une correspondance $c \in \text{Corr}^0(X, Y)$ entre variétés projectives
lisses, considérons l'application
$G(c) : G(X) = H^*(X) \to G(Y) = H^*(Y)$ donnée par la règle
$$
a \longmapsto
G(c)(a) = \text{pr}_{2, *}(\gamma(c) \cup \text{pr}_1^*a)
$$
Il est clair que $G(c)$ est additive en $c$, donc $\mathbf{Q}$-linéaire.
La compatibilité de $\gamma$ avec les images réciproques, les images directes
et les produits d'intersection, donnée par les axiomes (C)(a), (C)(b) et (C)(c),
montre que
$G(c' \circ c) = G(c') \circ G(c)$ si $c' \in \text{Corr}^0(Y, Z)$.
En effet, pour $a \in H^*(X)$, on a
\begin{align*}
(G(c') \circ G(c))(a)
& =
\text{pr}^{23}_{3, *}(\gamma(c') \cup
\text{pr}^{23, *}_2(\text{pr}^{12}_{2, *}(\gamma(c) \cup
\text{pr}^{12, *}_1a))) \\
& =
\text{pr}^{23}_{3, *}(\gamma(c') \cup
\text{pr}^{123}_{23, *}(\text{pr}^{123, *}_{12}(\gamma(c) \cup
\text{pr}^{12, *}_1 a))) \\
& =
\text{pr}^{23}_{3, *}
\text{pr}^{123}_{23, *}(
\text{pr}^{123, *}_{23}\gamma(c') \cup
\text{pr}^{123, *}_{12}\gamma(c) \cup
\text{pr}^{123, *}_1 a) \\
& =
\text{pr}^{23}_{3, *}
\text{pr}^{123}_{23, *}(
\gamma(\text{pr}^{123, *}_{23}c') \cup
\gamma(\text{pr}^{123, *}_{12}c) \cup
\text{pr}^{123, *}_1 a) \\
& =
\text{pr}^{13}_{3, *}
\text{pr}^{123}_{13, *}(
\gamma(\text{pr}^{123, *}_{23}c' \cdot \text{pr}^{123, *}_{12}c) \cup
\text{pr}^{123, *}_1 a) \\
& =
\text{pr}^{13}_{3, *}(
\gamma(\text{pr}^{123}_{13, *}(
\text{pr}^{123, *}_{23}c' \cdot \text{pr}^{123, *}_{12}c)) \cup
\text{pr}^{13, *}_1 a) \\
& =
G(c' \circ c)(a)
\end{align*},
avec des notations évidentes. La première égalité résulte des définitions.
La deuxième résulte de l'égalité
$\text{pr}^{23, *}_2 \circ \text{pr}^{12}_{2, *} =
\text{pr}^{123}_{23, *} \circ \text{pr}^{123, *}_{12}$,
qui découle immédiatement de la description de l'image directe
par les projections donnée dans le lemme \ref{weil-lemma-pr2star-classical}.
La troisième égalité résulte du lemme \ref{weil-lemma-pushforward-classical}
et du fait que $H^*$ est un foncteur.
La quatrième égalité résulte de l'axiome (C)(a) et du fait que
le morphisme de Gysin coïncide avec l'image réciproque plate pour les morphismes plats
(Homologie de Chow, lemme \ref{chow-lemma-lci-gysin-flat}).
La cinquième utilise l'axiome (C)(c), ainsi que
le lemme \ref{weil-lemma-pushforward-classical}, pour obtenir
$\text{pr}^{23}_{3, *} \circ \text{pr}^{123}_{23, *} =
\text{pr}^{13}_{3, *} \circ \text{pr}^{123}_{13, *}$.
La sixième utilise la formule de projection du
lemme \ref{weil-lemma-pushforward-classical}, ainsi que
l'axiome (C)(b), pour obtenir $
\text{pr}^{123}_{13, *}
\gamma(\text{pr}^{123, *}_{23}c' \cdot \text{pr}^{123, *}_{12}c) =
\gamma(\text{pr}^{123}_{13, *}(
\text{pr}^{123, *}_{23}c' \cdot \text{pr}^{123, *}_{12}c))$.
Enfin, la dernière égalité est la définition.

\medskip\noindent
Pour achever la démonstration que $G$ est un foncteur,
il faut montrer qu'il préserve les identités. Autrement dit, si
$1 = [\Delta] \in \text{Corr}^0(X, X)$ est l'identité
dans la catégorie des correspondances (voir
le lemme \ref{weil-lemma-category-correspondences} et sa démonstration),
il faut montrer que $G([\Delta]) = \text{id}$.
Cela résulte de la détermination de
$\gamma([\Delta])$ dans le lemme \ref{weil-lemma-class-diagonal-classical}
et du lemme \ref{weil-lemma-pr2star-classical}.
La construction de $G$ comme foncteur sur les
variétés projectives lisses et les correspondances de degré $0$ est ainsi achevée.

\medskip\noindent
Il résulte des axiomes (A)(c) et (A)(d) que
$G(\Spec(k)) = H^*(\Spec(k))$ est canoniquement isomorphe à $F$
comme $F$-algèbre.
L'axiome de K\"unneth (B) montre que notre foncteur est compatible avec les produits tensoriels.
Il est donc un foncteur monoïdal symétrique.

\medskip\noindent
Il reste à vérifier que l'image de $G(c_2)$ dans $G(\mathbf{P}^1)$
est un espace vectoriel gradué inversible sur $F$ (en particulier, nous ne savons pas encore
que $G$ se prolonge à $M_k$).
D'après l'axiome (A)(d), l'application $\int_{\mathbf{P}^1} : H^2(\mathbf{P}^1) \to F$
est un isomorphisme. D'après l'axiome (A)(b), on a $\dim_F H^0(\mathbf{P}^1) = 1$.
Le lemme \ref{weil-lemma-square-diagonal-classical} et l'axiome (A)(c)
donnent $2 - \dim_F H^1(\mathbf{P}^1) = c_1(T_{\mathbf{P}^1}) = 2$.
Ainsi, $H^1(\mathbf{P}^1) = 0$. Par conséquent,
$$
G(\mathbf{P}^1) = H^0(\mathbf{P}^1) \oplus H^2(\mathbf{P}^1)
$$
Rappelons que $1 = c_0 + c_2$ est une décomposition de l'identité
en somme d'idempotents orthogonaux dans
$\text{Corr}^0(\mathbf{P}^1, \mathbf{P}^1)$ ; voir
l'exemple \ref{weil-example-decompose-P1}. On a $c_0 = a \circ b$, où
$a \in \text{Corr}^0(\Spec(k), \mathbf{P}^1)$ et
$b \in \text{Corr}^0(\mathbf{P}^1, \Spec(k))$, et où
$b \circ a = 1$ dans $\text{Corr}^0(\Spec(k), \Spec(k))$ ; voir la démonstration du
lemme \ref{weil-lemma-inverse-h2}. Puisque $F = G(\Spec(k))$, la
fonctorialité entraîne que $G(c_0)$ est le projecteur sur le facteur direct
$H^0(\mathbf{P}^1) \subset G(\mathbf{P}^1)$. Par conséquent,
$G(c_2)$ est nécessairement la projection sur $H^2(\mathbf{P}^1)$,
ce qui achève la démonstration.
\end{proof}

\begin{proposition}
\label{weil-proposition-weil-cohomology-theory-classical}
Soit $k$ un corps algébriquement clos et soit $F$ un corps de
caractéristique $0$. Une théorie cohomologique de Weil classique revient à se donner
un foncteur $\mathbf{Q}$-linéaire
$$
G : M_k \longrightarrow \text{espaces vectoriels gradués sur }F\text{}
$$
monoïdal symétrique, muni d'un isomorphisme
$F[2] \to G(\mathbf{1}(1))$ d'espaces vectoriels gradués sur $F$, tel que,
de plus,
\begin{enumerate}
\item $G(h(X))$ soit concentré en degrés positifs ou nuls ;
\item $\dim_F G^0(h(X)) = 1$
\end{enumerate}
pour toute variété projective lisse $X$.
\end{proposition}

\begin{proof}
Étant donnés $G$ et $F[2] \to G(\mathbf{1}(1))$, on pose $H^*(X) = G(h(X))$
et l'on obtient des données (D1), (D2) et (D3) satisfaisant à tous les axiomes (A), (B) et (C),
sauf peut-être (A)(c) et (A)(d) ; voir
le lemme \ref{weil-lemma-from-functor-to-weil-classical} et sa démonstration.
Observons que les hypothèses (1) et (2) entraînent les axiomes (A)(c) et (A)(d),
compte tenu des axiomes (A)(a) et (A)(b) déjà établis.

\medskip\noindent
Réciproquement, à partir de $H^*$, on obtient un foncteur $G$ par la construction du
lemme \ref{weil-lemma-from-weil-to-functor-classical}.
Soient $X = \mathbf{P}^1, c_0, c_2$ comme dans l'exemple \ref{weil-example-decompose-P1}.
Nous avons construit un isomorphisme $1(-1) \to (X, c_2, 0)$ de motifs au
lemme \ref{weil-lemma-inverse-h2}. Dans la démonstration du
lemme \ref{weil-lemma-from-weil-to-functor-classical}, nous avons vu que
$G(1(-1)) = G(X, c_2, 0) = H^2(\mathbf{P}^1)[-2]$.
Ainsi, l'isomorphisme $\int_{\mathbf{P}^1} : H^2(\mathbf{P}^1) \to F$
de l'axiome (A)(d) donne un isomorphisme $G(1(-1)) \to F[-2]$,
qui détermine un isomorphisme $F[2] \to G(\mathbf{1}(1))$.
Enfin, puisque $G(h(X)) = H^*(X)$, les hypothèses (1) et (2)
résultent de l'axiome (A).
\end{proof}











\section{Cycles sur les corps non clos}
\label{weil-section-cycles-nonclosed}

\noindent
Voici quelques lemmes qui nous seront utiles pour étudier les motifs
sur des corps de base non algébriquement clos.

\begin{lemma}
\label{weil-lemma-generated-by-separable}
Soit $k$ un corps et soit $X$ un schéma projectif lisse sur $k$.
Alors $\CH_0(X)$ est engendré par les classes de points fermés dont les corps
résiduels sont séparables sur $k$.
\end{lemma}

\begin{proof}
Le lemme est immédiat si $k$ est de caractéristique $0$ ou s'il est parfait.
On peut donc supposer que $k$ est un corps infini de caractéristique $p > 0$.

\medskip\noindent
On peut supposer $X$ irréductible de dimension $d$.
Alors $k' = H^0(X, \mathcal{O}_X)$ est une extension finie séparable
de $k$ et $X$ est géométriquement intègre sur $k'$.
Voir Variétés, lemmes \ref{varieties-lemma-smooth-geometrically-normal},
\ref{varieties-lemma-proper-geometrically-reduced-global-sections} et
\ref{varieties-lemma-baby-stein}. Remplaçons donc $k$ par $k'$
et supposons désormais $X$ géométriquement intègre.

\medskip\noindent
Soit $x \in X$ un point fermé. Pour démontrer le lemme, nous allons montrer que
$[x] \in \CH_0(X)$ est rationnellement équivalent à une combinaison linéaire
à coefficients entiers de classes de points fermés dont les corps résiduels
sont séparables sur $k$. Choisissons un
$\mathcal{O}_X$-module inversible ample $\mathcal{L}$. Posons
$$
V = \{s \in H^0(X, \mathcal{L}) \mid s(x) = 0 \}
$$
En remplaçant $\mathcal{L}$ par une puissance, on peut supposer
que (a) $\mathcal{L}$ est très ample, (b) $V$ engendre
$\mathcal{L}$ sur $X \setminus x$, (c) le morphisme
$X \setminus x \to \mathbf{P}(V)$ est une immersion et (d)
l'application $V \to \mathfrak m_x\mathcal{L}_x/\mathfrak m_x^2\mathcal{L}_x$
est surjective ; voir
Morphismes, lemme \ref{morphisms-lemma-finite-type-ample-very-ample},
Variétés, lemme \ref{varieties-lemma-generate-over-complement}, et
Propriétés, proposition \ref{properties-proposition-characterize-ample}.
Considérons l'ensemble
$$
V^d \supset U =
\{
(s_1, \ldots, s_d) \in V^d \mid s_1, \ldots, s_d
\text{ engendrent }
\mathfrak m_x\mathcal{L}_x/\mathfrak m_x^2\mathcal{L}_x
\text{ sur }\kappa(x)
\}
$$
Puisque $\mathcal{O}_{X, x}$ est un anneau local régulier de dimension $d$,
on a $\dim_{\kappa(x)}(\mathfrak m_x/\mathfrak m_x^2) = d$ ;
ainsi, $U$ est un ouvert de Zariski non vide de $V^d$.
Pour $(s_1, \ldots, s_d) \in U$, posons $H_i = Z(s_i)$. Puisque
$s_1, \ldots, s_d$ engendrent $\mathfrak m_x\mathcal{L}_x$,
on a
$$
H_1 \cap \ldots \cap H_d = x \amalg Z
$$
au sens schématique, pour un certain sous-schéma fermé $Z \subset X$.
D'après Bertini (sous la forme de Variétés, lemme \ref{varieties-lemma-bertini}),
pour un élément général $s_1 \in V$, le schéma $H_1 \cap (X \setminus x)$
est lisse sur $k$, de dimension $d - 1$.
Une fois $s_1$ choisi, pour un élément général
$s_2 \in V$, le schéma $H_1 \cap H_2 \cap (X \setminus x)$
est lisse sur $k$, de dimension $d - 2$. Et ainsi de suite.
Il en résulte que, pour un choix suffisamment général de
$(s_1, \ldots, s_d) \in U$, le schéma $Z$ est \'etale sur $\Spec(k)$.
En particulier, $H_1 \cap \ldots \cap H_d$ est de dimension $0$,
d'où
$$
[H_1] \cdot \ldots \cdot [H_d] = [x] + [Z]
$$
dans $\CH_0(X)$, par applications successives de
Homologie de Chow, lemme \ref{chow-lemma-intersect-properly} (nous omettons les détails).
Cela achève la démonstration, car on obtient $[x] \sim_{rat} - [Z] + [Z']$,
où $Z' = H'_1 \cap \ldots \cap H'_d$ est une intersection complète générale
de lieux d'annulation de sections suffisamment générales
de $\mathcal{L}$, qui est \'etale sur $k$ par le même
raisonnement que précédemment.
\end{proof}

\begin{lemma}
\label{weil-lemma-chow-limit}
Soit $K/k$ une extension algébrique de corps et soit $X$ un schéma
de type fini sur $k$. Alors $\CH_i(X_K) = \colim \CH_i(X_{k'})$, où la
limite inductive porte sur les sous-extensions $K/k'/k$ telles que $k'/k$ soit finie.
\end{lemma}

\begin{proof}
C'est un cas particulier de
Homologie de Chow, lemme \ref{chow-lemma-chow-limit}.
\end{proof}

\begin{lemma}
\label{weil-lemma-divide-difference-points}
Soit $k$ un corps et soit $X$ un schéma projectif lisse
géométriquement irréductible sur $k$. Soient $x, x' \in X$ des points $k$-rationnels.
Soit $n$ un entier inversible dans $k$.
Il existe alors une extension finie séparable $k'/k$ telle que
l'image réciproque de $[x] - [x']$ sur $X_{k'}$
soit divisible par $n$ dans $\CH_0(X_{k'})$.
\end{lemma}

\begin{proof}
Soit $k'$ une clôture séparable de $k$. Supposons que l'on sache montrer
que l'image réciproque de $[x] - [x']$ sur $X_{k'}$ est divisible par $n$ dans
$\CH_0(X_{k'})$. On conclut alors par le lemme \ref{weil-lemma-chow-limit}.
On peut donc supposer désormais $k$ séparablement clos.

\medskip\noindent
Supposons $\dim(X) > 1$. Soit $\mathcal{L}$ un faisceau inversible ample sur $X$.
Posons
$$
V = \{s \in H^0(X, \mathcal{L}) \mid s(x) = 0\text{ et }s(x') = 0 \}
$$
En remplaçant $\mathcal{L}$ par une puissance, on voit que, pour
un élément général $v \in V$, le diviseur correspondant $H_v \subset X$ est lisse
en dehors de $x$ et de $x'$ ; voir
Variétés, lemmes \ref{varieties-lemma-generate-over-complement} et
\ref{varieties-lemma-bertini}. Pour trouver $v$, on utilise le fait que $k$ est infini,
car séparablement clos.
Si l'on choisit $s$ général, l'image de $s$ dans
$\mathfrak m_x\mathcal{L}_x/\mathfrak m_x^2\mathcal{L}_x$
est non nulle, ce qui entraîne que $H_v$ est lisse en $x$
(nous omettons les détails). Il en est de même pour $x'$. Ainsi, $H_v$ est lisse.
D'après Variétés, lemme \ref{varieties-lemma-connectedness-ample-divisor}
(appliqué après changement de base de tous les objets
à la clôture algébrique de $k$),
$H_v$ est géométriquement connexe.
Il suffit de démontrer le résultat pour
$[x] - [x']$ considéré comme un élément de $\CH_0(H_v)$.
On se ramène ainsi au cas d'une courbe.

\medskip\noindent
Supposons que $X$ soit une courbe. Alors $\mathcal{O}_X(x - x')$
définit un point $k$-rationnel $g$ de $J = \underline{\Pic}^0_{X/k}$ ; voir
Schémas de Picard des courbes, lemme \ref{pic-lemma-picard-pieces}.
Rappelons que $J$ est une variété propre et lisse sur $k$,
qui est aussi un schéma en groupes sur $k$ (même référence).
Par conséquent, $J$ est géométriquement intègre
(voir Variétés, lemmes \ref{varieties-lemma-geometrically-connected-criterion}
et \ref{varieties-lemma-smooth-geometrically-normal}).
Autrement dit, $J$ est une variété abélienne ; voir
Groupoïdes, définition \ref{groupoids-definition-abelian-variety}.
Ainsi, $[n] : J \to J$ est fini \'etale d'après
Groupoïdes, proposition \ref{groupoids-proposition-review-abelian-varieties}
(c'est ici que l'on utilise l'inversibilité de $n$ dans $k$).
Puisque $k$ est séparablement clos, on en déduit $g = [n](g')$
pour un certain $g' \in J(k)$. Si $\mathcal{L}$ est le module inversible de degré $0$
sur $X$ correspondant à $g'$, on en déduit
$\mathcal{O}_X(x - x') \cong \mathcal{L}^{\otimes n}$, comme voulu.
\end{proof}

\begin{lemma}
\label{weil-lemma-kernel-to-closure}
Soit $K/k$ une extension algébrique de corps.
Soit $X$ un schéma de type fini sur $k$.
Le noyau de l'application $\CH_i(X) \to \CH_i(X_K)$
construite au lemme \ref{weil-lemma-chow-limit}
est de torsion.
\end{lemma}

\begin{proof}
Il suffit manifestement de montrer que le noyau
de l'image réciproque plate $\CH_i(X) \to \CH_i(X_{k'})$
par $\pi : X_{k'} \to X$ est de torsion,
pour toute extension finie $k'/k$. Cela résulte de
$\pi_* \pi^* \alpha = [k' : k] \alpha$, d'après
Homologie de Chow, lemme \ref{chow-lemma-finite-flat}.
\end{proof}

\begin{lemma}[Voevodsky]
\label{weil-lemma-smash-nilpotence}
\begin{reference}
\cite{nilpotence}
\end{reference}
Soit $k$ un corps et soit $X$ un schéma projectif lisse
géométriquement irréductible sur $k$. Soient $x, x' \in X$ des points $k$-rationnels.
Pour $n$ assez grand, la classe du zéro-cycle
$$
([x] - [x']) \times \ldots \times ([x] - [x']) \in
\CH_0(X^n)
$$
est de torsion.
\end{lemma}

\begin{proof}
Si l'on sait le montrer après changement de base à la clôture algébrique de $k$,
le résultat s'ensuit sur $k$, car le noyau de l'image réciproque
est de torsion d'après le lemme \ref{weil-lemma-kernel-to-closure}.
On peut donc supposer désormais $k$ algébriquement clos.

\medskip\noindent
Le théorème de Bertini permet de choisir une courbe lisse $C \subset X$ passant par
$x$ et $x'$. Voir la démonstration du lemme \ref{weil-lemma-divide-difference-points}.
On peut donc supposer que $X$ est une courbe.

\medskip\noindent
Supposons que $X$ soit une courbe et que $k$ soit algébriquement clos.
Écrivons $S^n(X) = \underline{\Hilbfunctor}^n_{X/k}$, avec les notations de
Schémas de Picard des courbes, sections \ref{pic-section-hilbert-scheme-points}
et \ref{pic-section-divisors}. Il existe un morphisme canonique
$$
\pi : X^n \longrightarrow S^n(X)
$$
qui envoie le point $k$-rationnel $(x_1, \ldots, x_n)$ sur le point $k$-rationnel
correspondant au diviseur $[x_1] + \ldots + [x_n]$ sur $X$.
Le groupe symétrique $S_n$ agit fidèlement sur $X^n$.
Le morphisme $\pi$ est $S_n$-invariant, et les fibres de $\pi$ sont
des orbites sous $S_n$ (ensemblistement). Enfin, $\pi$ est fini plat de
degré $n!$ ; voir Schémas de Picard des courbes, lemme
\ref{pic-lemma-universal-object}.

\medskip\noindent
Soit $\alpha_n$ le zéro-cycle sur $X^n$ donné par la formule de
l'énoncé du lemme. Posons $\mathcal{L} = \mathcal{O}_X(x - x')$. Alors
$c_1(\mathcal{L}) \cap [X] = [x] - [x']$. Ainsi,
$$
\alpha_n = c_1(\mathcal{L}_1) \cap \ldots \cap c_1(\mathcal{L}_n) \cap [X^n]
$$,
où $\mathcal{L}_i = \text{pr}_i^*\mathcal{L}$ et où $\text{pr}_i : X^n \to X$
est la $i$-ième projection. D'après
Diviseurs, lemme \ref{divisors-lemma-finite-locally-free-has-norm}, ou
Diviseurs, lemme \ref{divisors-lemma-norm-in-normal-case},
il existe une norme pour $\pi$. Posons
$\mathcal{N} = \text{Norm}_\pi(\mathcal{L}_1)$ ;
voir Diviseurs, lemme \ref{divisors-lemma-norm-invertible}. Un calcul donne
$$
\pi^*\mathcal{N} =
(\mathcal{L}_1 \otimes \ldots \otimes \mathcal{L}_n)^{\otimes (n - 1)!}
$$
dans $\Pic(X^n)$. Nous omettons les détails ; indiquons que cela résulte
du fait que
$\text{Norm}_\pi : \pi_*\mathcal{O}_{X^n} \to \mathcal{O}_{S^n(X)}$,
composé avec l'application naturelle $\pi_*\mathcal{O}_{S^n(X)} \to \mathcal{O}_{X^n}$,
est égal au produit, pour tous les $\sigma \in S_n$, des actions de
$\sigma$ sur $\pi_*\mathcal{O}_{X^n}$. Considérons
$$
\beta_n = c_1(\mathcal{N})^n \cap [S^n(X)]
$$
dans $\CH_0(S^n(X))$. Observons que
$c_1(\mathcal{L}_i) \cap c_1(\mathcal{L}_i) = 0$,
car $\mathcal{L}_i$ provient d'une courbe par image réciproque ; voir
Homologie de Chow, lemme \ref{chow-lemma-vanish-above-dimension}. On obtient donc
\begin{align*}
\pi^*\beta_n
& =
((n - 1)!)^n
(\sum\nolimits_{i = 1, \ldots, n} c_1(\mathcal{L}_i))^n \cap [X^n] \\
& =
((n - 1)!)^n n^n 
c_1(\mathcal{L}_1) \cap \ldots \cap c_1(\mathcal{L}_n) \cap [X^n] \\
& =
(n!)^n \alpha_n
\end{align*}
Il suffit ainsi de montrer que $\beta_n$ est de torsion.

\medskip\noindent
Il existe un morphisme canonique
$$
f : S^n(X) \longrightarrow \underline{\Picardfunctor}^n_{X/k}
$$
Voir Schémas de Picard des courbes, lemme \ref{pic-lemma-picard-pieces}.
Pour $n \geq 2g - 1$, ce morphisme est un fibré projectif
(nous omettons les détails ; comparer avec la
démonstration de Schémas de Picard des courbes, lemme \ref{pic-lemma-picard-pieces}).
Le faisceau inversible $\mathcal{N}$ est trivial sur les
fibres de $f$, comme on le verra ci-dessous. Par la formule du fibré projectif
(Homologie de Chow, lemme \ref{chow-lemma-chow-ring-projective-bundle}),
on a donc $\mathcal{N} = f^*\mathcal{M}$ pour un certain module
inversible $\mathcal{M}$ sur $\underline{\Picardfunctor}^n_{X/k}$.
Il s'ensuit évidemment que
$$
c_1(\mathcal{N})^n = f^*(c_1(\mathcal{M})^n)
$$
est nul, puisque $n > g = \dim(\underline{\Picardfunctor}^n_{X/k})$
et que l'on peut appliquer Homologie de Chow, lemme \ref{chow-lemma-vanish-above-dimension},
comme précédemment.

\medskip\noindent
Il reste à montrer que $\mathcal{N}$ est trivial sur une fibre $F$
de $f$. Puisque les fibres de $f$ sont des espaces projectifs et que
$\Pic(\mathbf{P}^m_k) = \mathbf{Z}$
(Diviseurs, lemme \ref{divisors-lemma-Pic-projective-space-UFD}),
on pourrait le montrer en calculant le degré de $\mathcal{N}$
sur une droite contenue dans la fibre. Nous le démontrerons plutôt
en établissant que $\mathcal{N}$ est algébriquement
équivalent à zéro. Montrons d'abord qu'il existe un schéma connexe $T$ de type fini
sur $k$, un module inversible $\mathcal{L}'$ sur $T \times X$ et des points
$k$-rationnels $p, q \in T$ tels que
$\mathcal{M}_p \cong \mathcal{O}_X$ et $\mathcal{M}_q = \mathcal{L}$.
En effet, puisque $\mathcal{L} = \mathcal{O}_X(x - x')$, on peut prendre
$T = X$, $p = x'$, $q = x$ et
$\mathcal{L}' = \mathcal{O}_{X \times X}(\Delta)
\otimes \text{pr}_2^*\mathcal{O}_X(-x')$.
On note ensuite $\mathcal{L}'_i$, sur
$T \times X^n$, pour $i = 1, \ldots, n$,
l'image réciproque de $\mathcal{L}'$ par
$\text{id}_T \times \text{pr}_i : T \times X^n \to T \times X$.
Enfin, on pose
$\mathcal{N}' = \text{Norm}_{\text{id}_T \times \pi}(\mathcal{L}'_1)$
sur $T \times S^n(X)$.
Par construction, on a $\mathcal{N}'_p = \mathcal{O}_{S^n(X)}$
et $\mathcal{N}'_q = \mathcal{N}$.
Il en résulte que
$$
\mathcal{N}'|_{T \times F}
$$
est un module inversible sur $T \times F \cong T \times \mathbf{P}^m_k$,
dont la fibre en $p$ est le module inversible trivial et dont la fibre
en $q$ est $\mathcal{N}|_F$. Puisque la caractéristique d'Euler
du fibré trivial est $1$ et que cette caractéristique d'Euler
est localement constante en famille (Catégories dérivées de schémas,
lemme \ref{perfect-lemma-chi-locally-constant-geometric}),
on obtient $\chi(F, \mathcal{N}^{\otimes s}|_F) = 1$
pour tout $s \in \mathbf{Z}$. Cela n'est possible que si
$\mathcal{N}|_F \cong \mathcal{O}_F$ (voir
Cohomologie des schémas, lemme
\ref{coherent-lemma-cohomology-projective-space-over-ring}),
ce qui achève la démonstration. Nous omettons certains détails.
\end{proof}















\section{Théories cohomologiques de Weil, I}
\label{weil-section-axioms}

\noindent
Cette section est l'analogue de la section \ref{weil-section-axioms-classical}
sur un corps quelconque. Autrement dit, nous déterminons les données
et les axiomes qui correspondent aux foncteurs monoïdaux symétriques $G$
de la catégorie des motifs dans celle des espaces vectoriels gradués, tels que
$G(\mathbf{1}(1))$ soit placé en degré $-2$. À la section \ref{weil-section-old},
nous définirons une théorie cohomologique de Weil en ajoutant une seule
condition supplémentaire.

\medskip\noindent
Fixons un corps $k$ (le corps de base).
Fixons un corps $F$ de caractéristique $0$ (le corps de coefficients).
Les données sont les suivantes :
\begin{enumerate}
\item[(D0)] Un espace vectoriel de dimension $1$ sur $F$, noté $F(1)$.
\item[(D1)] Un foncteur contravariant $H^*$ de la catégorie
des schémas projectifs lisses sur $k$ dans la catégorie des
$F$-algèbres graduées commutatives.
\item[(D2)] Pour tout schéma projectif lisse $X$ sur $k$,
un homomorphisme de groupes $\gamma : \CH^i(X) \to H^{2i}(X)(i)$.
\item[(D3)] Pour tout schéma projectif lisse non vide $X$ sur $k$,
équidimensionnel de dimension $d$, une application
$\int_X : H^{2d}(X)(d) \to F$.
\end{enumerate}
Précisons le sens de ces données par quelques remarques
et introduisons la terminologie correspondante.

\medskip\noindent
Remarques sur (D0).
L'espace vectoriel $F(1)$ définit des {\it tordus de Tate} sur la catégorie
des espaces vectoriels sur $F$. En effet, pour $n \in \mathbf{Z}$, on pose
$F(n) = F(1)^{\otimes n}$ si $n \geq 0$, puis $F(-1) = \Hom_F(F(1), F)$,
et $F(n) = F(-1)^{\otimes - n}$ si $n < 0$. Comparer avec
Compléments d'algèbre, section \ref{more-algebra-section-picard}.
Pour un espace vectoriel sur $F$, noté $V$, on définit le {\it $n$-ième tordu de Tate}
$$
V(n) = V \otimes_F F(n)
$$
Nous emploierons des notations évidentes : par exemple, pour des espaces vectoriels sur $F$, notés $U$, $V$
et $W$, et une application linéaire $U \otimes_F V \to W$, on obtient une application
linéaire $U(n) \otimes_F V(m) \to W(n + m)$ pour $n, m \in \mathbf{Z}$.

\medskip\noindent
Remarques sur (D1).
Pour un schéma projectif lisse $X$ sur $k$, on appelle $H^*(X)$
la {\it cohomologie} de $X$. Pour un morphisme $f : X \to Y$
de schémas projectifs lisses sur $k$, on note $f^* : H^*(Y) \to H^*(X)$
l'application $H^*(f)$ et on l'appelle l'{\it application d'image réciproque}.

\medskip\noindent
Remarques sur (D2). L'application $\gamma$ est appelée l'{\it application classe de cycle}.
On appelle $\gamma(\alpha)$ la {\it classe de cohomologie} de $\alpha$.
Si $Z \subset Y \subset X$ sont des sous-schémas fermés, avec $Y$ et $X$
projectifs lisses sur $k$ et $Z$ intègre, alors $[Z]$ peut
désigner la classe du cycle $[Z]$ dans $\CH^*(Y)$ ou dans $\CH^*(X)$.
Dans ce cas, la notation $\gamma([Z])$ est ambiguë, et son sens
doit être déduit du contexte.

\medskip\noindent
Remarques sur (D3). L'application $\int_X$ est parfois appelée
l'{\it application trace} et est parfois notée $\text{Tr}_X$.

\medskip\noindent
Le premier axiome est souvent appelé {\it dualité de Poincar\'e} :
\begin{enumerate}
\item[(A)] Soit $X$ un schéma projectif lisse non vide sur $k$,
équidimensionnel de dimension $d$. Alors :
\begin{enumerate}
\item $\dim_F H^i(X) < \infty$ pour tout $i$ ;
\item $H^i(X) \times H^{2d - i}(X)(d) \rightarrow
H^{2d}(X)(d) \rightarrow F$
est un accouplement parfait pour tout $i$, la dernière
application étant l'application trace $\int_X$.
\end{enumerate}
\end{enumerate}
Soit $f : X \to Y$ un morphisme de schémas projectifs lisses non vides, avec $X$
équidimensionnel de dimension $d$ et $Y$ équidimensionnel de dimension $e$.
À l'aide de la dualité de Poincar\'e, on peut définir une {\it image directe}
$$
f_* : H^{2d - i}(X)(d) \longrightarrow H^{2e - i}(Y)(e)
$$
comme l'application contragrédiente de l'application linéaire $f^* : H^i(Y) \to H^i(X)$.
Explicitement, pour $a \in H^{2d - i}(X)(d)$, l'élément $f_*a \in H^{2e - i}(Y)(e)$
est caractérisé par
$$
\int_X f^*b \cup a = \int_Y b \cup f_*a
$$
pour tout $b \in H^i(Y)$.

\begin{lemma}
\label{weil-lemma-pushforward}
Supposons que les données (D0), (D1) et (D3) satisfassent à (A). Pour $f : X \to Y$
un morphisme de schémas projectifs lisses non vides et équidimensionnels sur $k$,
on a $f_*(f^*b \cup a) = b \cup f_*a$. Si $g : Y \to Z$ est un second morphisme,
avec $Z$ non vide, projectif lisse et équidimensionnel, alors
$g_* \circ f_* = (g \circ f)_*$.
\end{lemma}

\begin{proof}
La première égalité résulte de
$$
\int_Y c \cup b \cup f_*a =
\int_X f^*c \cup f^*b \cup a =
\int_Y c \cup f_*(f^*b \cup a).
$$
La seconde résulte de
$$
\int_Z c \cup (g \circ f)_*a = \int_X (g \circ f)^*c \cup a =
\int_X f^* g^* c \cup a = \int_Y g^*c \cup f_*a = \int_Z c \cup g_*f_*a
$$
Cela achève la démonstration.
\end{proof}

\noindent
Le deuxième axiome exprime que $H^*$ respecte la structure monoïdale
donnée par les produits, au moyen de la {\it formule de K\"unneth} :
\begin{enumerate}
\item[(B)] Soient $X$ et $Y$ des schémas projectifs lisses sur $k$.
\begin{enumerate}
\item L'application $H^*(X) \otimes_F H^*(Y) \to H^*(X \times Y)$,
$\alpha \otimes \beta \mapsto \text{pr}_1^*\alpha \cup \text{pr}_2^*\beta$,
est un isomorphisme ;
\item si $X$ et $Y$ sont non vides et équidimensionnels, alors
$\int_{X \times Y} = \int_X \otimes \int_Y$ par l'identification (a).
\end{enumerate}
\end{enumerate}
L'axiome (B)(b) permet de calculer les images directes par les projections.

\begin{lemma}
\label{weil-lemma-pr2star}
Supposons que les données (D0), (D1) et (D3) satisfassent à (A) et (B).
Soient $X$ et $Y$ des schémas projectifs lisses non vides sur $k$,
équidimensionnels de dimensions $d$ et $e$. Alors
$\text{pr}_{2, *} : H^*(X \times Y)(d + e) \to H^*(Y)(e)$ envoie
$a \otimes b$ sur $(\int_X a) b$.
\end{lemma}

\begin{proof}
Cela résulte des axiomes (B)(a) et (B)(b).
\end{proof}

\noindent
Le troisième axiome porte sur les applications classes de cycles :
\begin{enumerate}
\item[(C)] Les applications classes de cycles satisfont aux règles suivantes :
\begin{enumerate}
\item pour un morphisme $f : X \to Y$ de schémas projectifs lisses sur
$k$, on a $\gamma(f^!\beta) = f^*\gamma(\beta)$ pour $\beta \in \CH^*(Y)$ ;
\item pour un morphisme $f : X \to Y$ de schémas projectifs lisses non vides
et équidimensionnels sur $k$, on a
$\gamma(f_*\alpha) = f_*\gamma(\alpha)$ pour $\alpha \in \CH^*(X)$ ;
\item pour tout schéma projectif lisse $X$ sur $k$, on a
$\gamma(\alpha \cdot \beta) = \gamma(\alpha) \cup \gamma(\beta)$
pour $\alpha, \beta \in \CH^*(X)$ ;
\item $\int_{\Spec(k)} \gamma([\Spec(k)]) = 1$.
\end{enumerate}
\end{enumerate}
Explicitons l'axiome (C)(b). Soit $f : X \to Y$
comme dans (C)(b), avec $\dim(X) = d$ et $\dim(Y) = e$. L'image directe
sur les groupes de Chow donne alors
$$
f_* : \CH^{d - i}(X) = \CH_i(X) \to \CH_i(Y) = \CH^{e - i}(Y)
$$
Soit $\alpha \in \CH^{d - i}(X)$. D'une part, on a
$f_*\alpha \in \CH^{e - i}(Y)$, d'où
$\gamma(f_*\alpha) \in H^{2e - 2i}(Y)(e - i)$.
D'autre part, on a
$\gamma(\alpha) \in H^{2d - 2i}(X)(d - i)$, d'où
$f_*\gamma(\alpha) \in H^{2e - 2i}(Y)(e - i)$ également.
La condition $\gamma(f_*\alpha) = f_*\gamma(\alpha)$ a donc bien un sens.

\begin{remark}
\label{weil-remark-replace-cup-product}
Supposons que les données (D0), (D1), (D2) et (D3) satisfassent à (A), (B) et (C)(a).
Soit $X$ un schéma projectif lisse sur $k$. On obtient des applications
$$
H^*(X) \otimes_F H^*(X) \longrightarrow H^*(X \times X)
\xrightarrow{\Delta^*} H^*(X)
$$,
où la première flèche est celle de l'axiome (B) et $\Delta^*$
est l'image réciproque par le morphisme diagonal $\Delta : X \to X \times X$.
La composée est le cup-produit, puisque l'image réciproque est un homomorphisme d'algèbres
et que $\text{pr}_i \circ \Delta = \text{id}$.
D'autre part, pour des cycles $\alpha, \beta$ sur $X$,
le produit d'intersection est défini par la formule
$$
\alpha \cdot \beta =
\Delta^!(\alpha \times \beta)
$$
Autrement dit, $\alpha \cdot \beta$ est l'image réciproque du
produit extérieur $\alpha \times \beta$ sur $X \times X$ par
la diagonale. Notons aussi que
$\alpha \times \beta = \text{pr}_1^*\alpha \cdot \text{pr}_2^*\beta$
dans $\CH^*(X \times X)$ (nous omettons la démonstration). Ainsi, sous l'axiome (C)(a),
l'axiome (C)(c) équivaut à la compatibilité de $\gamma$
avec le produit extérieur, au sens où
$\gamma(\alpha \times \beta)$ est égal à
$\text{pr}_1^*\gamma(\alpha) \cup \text{pr}_2^*\gamma(\beta)$.
\end{remark}

\begin{lemma}
\label{weil-lemma-base}
Supposons que les données (D0), (D1), (D2) et (D3) satisfassent à (A), (B) et (C).
Alors $H^i(\Spec(k)) = 0$ pour $i \not = 0$, et il existe un
unique isomorphisme de $F$-algèbres $F = H^0(\Spec(k))$.
On a $\gamma([\Spec(k)]) = 1$ et $\int_{\Spec(k)} 1 = 1$.
\end{lemma}

\begin{proof}
D'après l'axiome (C)(d), $H^0(\Spec(k))$ est non nul et même
$\gamma([\Spec(k)])$ est non nul.
Puisque $\Spec(k) \times \Spec(k) = \Spec(k)$, on obtient
$$
H^*(\Spec(k)) \otimes_F H^*(\Spec(k)) = H^*(\Spec(k))
$$
par l'axiome (B)(a), ce qui entraîne, en considérant les dimensions, que seul
$H^0$ est non nul et qu'il est de dimension $1$. Ainsi,
$F = H^0(\Spec(k))$ par l'unique isomorphisme de $F$-algèbres
envoyant $1 \in F$ sur $1 \in H^0(\Spec(k))$.
Puisque $[\Spec(k)] \cdot [\Spec(k)] = [\Spec(k)]$ dans
l'anneau de Chow de $\Spec(k)$, on obtient
$\gamma([\Spec(k)) \cup \gamma([\Spec(k)]) = \gamma([\Spec(k)])$
par l'axiome (C)(c). Comme nous savons déjà que $\gamma([\Spec(k)])$ est non nul,
il est nécessairement égal à $1$.
Enfin, on a $\int_{\Spec(k)} 1 = 1$, puisque
$\int_{\Spec(k)} \gamma([\Spec(k)]) = 1$
d'après l'axiome (C)(d).
\end{proof}

\begin{lemma}
\label{weil-lemma-unit}
Supposons que les données (D0), (D1), (D2) et (D3) satisfassent à (A), (B) et (C).
Soit $X$ un schéma projectif lisse sur $k$.
Si $X = \emptyset$, alors $H^*(X) = 0$.
Si $X$ est non vide, alors $\gamma([X]) = 1$ et $1 \not = 0$ dans $H^0(X)$.
\end{lemma}

\begin{proof}
Supposons d'abord $X$ non vide.
Observons que $[X]$ est l'image réciproque de $[\Spec(k)]$ par le morphisme structural
$p : X \to \Spec(k)$. On obtient donc $\gamma([X]) = 1$ par l'axiome (C)(a)
et le lemme \ref{weil-lemma-base}. Soit $X' \subset X$ une composante irréductible.
Par fonctorialité, il suffit de montrer que $1 \not = 0$ dans $H^0(X')$.
On peut donc supposer $X$ irréductible, et en particulier
non vide et équidimensionnel, de dimension $d$.
Pour établir $1 \not = 0$, il suffit de montrer que $H^*(X)$ est non nul.

\medskip\noindent
Soit $x \in X$ un point fermé dont le corps résiduel $k'$
est séparable sur $k$ ; voir
Variétés, lemme \ref{varieties-lemma-smooth-separable-closed-points-dense}.
Soit $i : \Spec(k') \to X$ le morphisme d'inclusion.
Notons $p : X \to \Spec(k)$ le morphisme structural.
Observons que
$p_*i_*[\Spec(k')] = [k' : k][\Spec(k)]$ dans $\CH_0(\Spec(k))$.
En appliquant deux fois l'axiome (C)(b), puis le lemme \ref{weil-lemma-base},
on en déduit que
$$
p_*i_*\gamma([\Spec(k')]) = \gamma([k' : k][\Spec(k)]) = [k' : k]
\in F = H^0(\Spec(k))
$$
est non nul. Ainsi, $i_*\gamma([\Spec(k)]) \in H^{2d}(X)(d)$ est non nul,
car son image par $p_*$ est non nulle. Cela achève la démonstration
lorsque $X$ est non vide.

\medskip\noindent
Considérons enfin le cas du schéma vide. L'axiome (B)(a) donne
$H^*(\emptyset) \otimes H^*(\emptyset) = H^*(\emptyset)$,
d'où il résulte que $H^*(\emptyset)$ est soit nul, soit de dimension $1$
en degré $0$. L'axiome (B)(a) donne à nouveau
$H^*(\emptyset) \otimes H^*(X) = H^*(\emptyset)$ pour
tout schéma projectif lisse $X$ sur $k$. L'axiome (A)(b)
et la non-nullité de $H^0(X)$ établie ci-dessus
entraînent que $H^*(X)$ est non nul en au moins deux degrés
si $\dim(X) > 0$. Cela force alors $H^*(\emptyset)$ à être nul.
\end{proof}

\begin{lemma}
\label{weil-lemma-push-unit}
Supposons que les données (D0), (D1), (D2) et (D3) satisfassent à (A), (B) et (C).
Soit $i : X \to Y$ une immersion fermée de schémas projectifs lisses
non vides et équidimensionnels sur $k$. Alors
$\gamma([X]) = i_*1$ dans $H^{2c}(Y)(c)$, où $c = \dim(Y) - \dim(X)$.
\end{lemma}

\begin{proof}
En effet, $1 = \gamma([X])$ dans $H^0(X)$ d'après le lemme \ref{weil-lemma-unit},
et il suffit alors d'appliquer l'axiome (C)(b).
\end{proof}

\begin{lemma}
\label{weil-lemma-class-diagonal}
Supposons que les données (D0), (D1), (D2) et (D3) satisfassent à (A), (B) et (C).
Soit $X$ un schéma projectif lisse non vide sur $k$, équidimensionnel
de dimension $d$. Choisissons une base $e_{i, j}, j = 1, \ldots, \beta_i$ de
$H^i(X)$ sur $F$. À l'aide de la formule de K\"unneth, écrivons
$$
\gamma([\Delta]) =
\sum\nolimits_i
\sum\nolimits_j e_{i, j} \otimes e'_{2d - i , j}
\quad\text{dans}\quad
\bigoplus\nolimits_i H^i(X) \otimes_F H^{2d - i}(X)(d)
$$
avec $e'_{2d - i, j} \in H^{2d - i}(X)(d)$.
Alors $\int_X e_{i, j} \cup e'_{2d - i, j'} = (-1)^i\delta_{jj'}$.
\end{lemma}

\begin{proof}
Rappelons que $\Delta^* : H^*(X \times X) \to H^*(X)$ est l'application
donnée par le cup-produit $H^*(X) \otimes_F H^*(X) \to H^*(X)$ ; voir
la remarque \ref{weil-remark-replace-cup-product}. D'autre part,
rappelons que $\gamma([\Delta]) = \Delta_*1$ (lemme \ref{weil-lemma-push-unit}),
d'où
$$
\int_{X \times X} \gamma([\Delta]) \cup a \otimes b =
\int_{X \times X} \Delta_*1 \cup a \otimes b =
\int_X a \cup b
$$
d'après le lemme \ref{weil-lemma-pushforward}.
D'autre part, on a
$$
\int_{X \times X} (\sum e_{i, j} \otimes e'_{2d -i , j}) \cup a \otimes b =
\sum (\int_X a \cup e_{i, j})(\int_X e'_{2d - i, j} \cup b)
$$
par l'axiome (B)(b) ; notons que nous avons effectué deux échanges, de sorte que
chaque terme a le signe $1$. Ainsi, si l'on choisit $a$ de sorte que
$\int_X a \cup e_{i, j} = 1$ et que tous les autres accouplements soient nuls,
on obtient $\int_X e'_{2d - i, j} \cup b = \int_X a \cup b$
pour tout $b$, c'est-à-dire $e'_{2d - i, j} = a$. Le lemme est démontré.
\end{proof}

\begin{lemma}
\label{weil-lemma-cohomology-P1}
Supposons que les données (D0), (D1), (D2) et (D3) satisfassent à (A), (B) et (C).
Alors $H^*(\mathbf{P}^1_k)$ est de dimension $1$ en degrés $0$ et $2$
et nul en tout autre degré.
\end{lemma}

\begin{proof}
Soit $x \in \mathbf{P}^1_k$ un point $k$-rationnel. Observons que
$\Delta = \text{pr}_1^*x + \text{pr}_2^*x$ comme diviseurs sur
$\mathbf{P}^1_k \times \mathbf{P}^1_k$. L'axiome (C)(a)
et l'additivité de $\gamma$ donnent
$$
\gamma([\Delta]) =
\text{pr}_1^*\gamma([x]) +
\text{pr}_2^*\gamma([x]) =
\gamma([x]) \otimes 1 + 1 \otimes \gamma([x])
$$
dans $H^*(\mathbf{P}^1_k \times \mathbf{P}^1_k) =
H^*(\mathbf{P}^1_k) \otimes_F H^*(\mathbf{P}^1_k)$.
Or le lemme \ref{weil-lemma-class-diagonal}
montre que $\gamma([\Delta])$ ne peut s'écrire
comme somme de moins de $\sum \beta_i$ tenseurs purs,
où $\beta_i = \dim_F H^i(\mathbf{P}^1_k)$.
Ainsi, $\sum \beta_i \leq 2$.
D'après le lemme \ref{weil-lemma-unit}, on a $H^0(\mathbf{P}^1_k) \not = 0$.
La dualité de Poincar\'e, plus précisément l'axiome (A)(b),
donne $\beta_0 = \beta_2$. Le lemme s'ensuit.
\end{proof}

\begin{lemma}
\label{weil-lemma-weil-additive}
Supposons que les données (D0), (D1), (D2) et (D3) satisfassent à (A), (B) et (C).
Si $X$ et $Y$ sont des schémas projectifs lisses sur $k$, alors
$H^*(X \amalg Y) \to H^*(X) \times H^*(Y)$,
$a \mapsto (i^*a, j^*a)$, est un isomorphisme, où $i$ et $j$
sont les coprojections.
\end{lemma}

\begin{proof}
Si $X$ ou $Y$ est vide, cela résulte de
$H^*(\emptyset) = 0$, d'après le lemme \ref{weil-lemma-unit}.
On peut donc supposer $X$ et $Y$ tous deux non vides.

\medskip\noindent
Montrons d'abord que l'application est injective. Observons que
l'on peut trouver des morphismes $X' \to X$ et $Y' \to Y$
de schémas projectifs lisses, tels que $X'$ et $Y'$ soient
équidimensionnels de même dimension et que
$X' \to X$ et $Y' \to Y$ admettent chacun une section. En effet,
décomposons $X = \coprod X_d$ et $Y = \coprod Y_e$
en sous-schémas ouverts et fermés équidimensionnels de
dimensions $d$ et $e$. On prend alors
$X' = \coprod X_d \times \mathbf{P}^{n - d}$
et $Y' = \coprod Y_e \times \mathbf{P}^{n - e}$ pour un entier
$n$ assez grand. L'image réciproque par
$X' \amalg Y' \to X \amalg Y$ est donc injective,
car ce morphisme admet une section, et
il suffit de démontrer l'injectivité pour $X', Y'$,
ce que nous faisons au paragraphe suivant.

\medskip\noindent
Montrons que l'application est injective lorsque $X$ et $Y$ sont équidimensionnels
de même dimension $d$.
Observons que $[X \amalg Y] = [X] + [Y]$ dans $\CH^0(X \amalg Y)$
et que $[X]$ et $[Y]$ sont des idempotents orthogonaux dans $\CH^0(X \amalg Y)$.
Ainsi,
$$
1 = \gamma([X \amalg Y] = \gamma([X]) + \gamma([Y]) = i_*1 + j_*1
$$
est une décomposition en idempotents orthogonaux. Nous avons utilisé ici
les lemmes \ref{weil-lemma-unit} et \ref{weil-lemma-push-unit}, ainsi que l'axiome (C)(c).
On obtient alors
$$
a = a \cup 1 = a \cup i_*1 + a \cup j_*1 =
i_*(i^*a) + j_*(j^*a)
$$
par la formule de projection (lemme \ref{weil-lemma-pushforward}), d'où
l'injectivité.

\medskip\noindent
Montrons que l'application est surjective. Posons $e = \gamma([X])$ et $f = \gamma([Y])$,
considérés comme des éléments de $H^0(X \amalg Y)$. On a
$i^*e = 1$, $i^*f = 0$, $j^*e = 0$ et $j^*f = 1$ d'après l'axiome (C)(a).
Ainsi, si $i^* : H^*(X \amalg Y) \to  H^*(X)$
et $j^* : H^*(X \amalg Y) \to H^*(Y)$ sont surjectives,
$(i^*, j^*)$ l'est aussi. En effet, pour $a, a' \in H^*(X \amalg Y)$,
on a
$$
(i^*a, j^*a') = (i^*(a \cup e + a' \cup f), j^*(a \cup e + a' \cup f))
$$
Par symétrie, il suffit de montrer que $i^* : H^*(X \amalg Y) \to  H^*(X)$
est surjective. S'il existe un morphisme $Y \to X$, il existe un morphisme
$g : X \amalg Y \to X$ tel que $g \circ i = \text{id}_X$, et l'on conclut.
Pour achever la démonstration, observons que, pour montrer
la surjectivité de $i^*$, il suffit de le faire après produit tensoriel
par un espace vectoriel gradué non nul sur $F$. Ainsi, par l'axiome (B)(b)
et la non-nullité de la cohomologie (lemme \ref{weil-lemma-unit}),
il suffit de démontrer la surjectivité de $i^*$ après avoir remplacé
$X$ et $Y$ par $X \times \Spec(k')$ et $Y \times \Spec(k')$
pour une extension finie séparable $k'/k$ convenable.
Si l'on choisit $k'$ de sorte qu'il existe un point fermé
$x \in X$ tel que $\kappa(x) = k'$ (ce qui est possible d'après
Variétés, lemme \ref{varieties-lemma-smooth-separable-closed-points-dense}),
il existe alors un morphisme $Y \times \Spec(k') \to X \times \Spec(k')$,
ce qui achève la démonstration.
\end{proof}

\begin{lemma}
\label{weil-lemma-from-functor-to-weil}
Soit $k$ un corps et soit $F$ un corps de caractéristique $0$.
Supposons donné un foncteur $\mathbf{Q}$-linéaire
$$
G : M_k \longrightarrow \text{espaces vectoriels gradués sur }F\text{}
$$
qui est monoïdal symétrique et tel que $G(\mathbf{1}(1))$
ne soit non nul qu'en degré $-2$. On obtient alors des données (D0), (D1), (D2) et (D3)
satisfaisant à tous les axiomes (A), (B) et (C) ci-dessus.
\end{lemma}

\begin{proof}
Cette démonstration est la même que celle du
lemme \ref{weil-lemma-from-functor-to-weil-classical} ;
nous conseillons vivement au lecteur de lire plutôt la démonstration de ce lemme.

\medskip\noindent
On obtient un foncteur contravariant de la catégorie des schémas projectifs
lisses sur $k$ dans celle des espaces vectoriels gradués sur $F$
en posant $H^*(X) = G(h(X))$. Par hypothèse, on dispose d'un isomorphisme
canonique
$$
H^*(X \times Y) = G(h(X \times Y)) = G(h(X) \otimes h(Y)) =
G(h(X)) \otimes G(h(Y)) = H^*(X) \otimes H^*(Y)
$$,
compatible avec les images réciproques. L'image réciproque par la diagonale
$\Delta : X \to X \times X$ donne une application canonique
$$
H^*(X) \otimes H^*(X) = H^*(X \times X) \to H^*(X)
$$
d'espaces vectoriels gradués, compatible avec les images réciproques.
Cela définit une structure fonctorielle de $F$-algèbre graduée sur
$H^*(X)$. Puisque $\Delta$ commute à la contrainte de commutativité
$h(X) \otimes h(X) \to h(X) \otimes h(X)$ (qui échange les facteurs),
que $G$ est un foncteur monoïdal symétrique (donc compatible avec
les contraintes de commutativité), et compte tenu de notre convention dans
Homologie, exemple \ref{homology-example-graded-vector-spaces},
on en déduit que $H^*(X)$ est une algèbre graduée
commutative ! Nous obtenons ainsi la donnée (D1).

\medskip\noindent
Puisque $\mathbf{1}(1)$ est inversible dans la catégorie des motifs,
$G(\mathbf{1}(1))$ est inversible dans la catégorie des
espaces vectoriels gradués sur $F$. Ainsi, $\sum_i \dim_F G^i(\mathbf{1}(1)) = 1$.
Par hypothèse, le seul degré où cet objet est non nul est $-2$.
La donnée (D0) est l'espace vectoriel $F(1) = G^{-2}(\mathbf{1}(1))$.
Puisque $G$ est un foncteur monoïdal symétrique,
on a $F(n) = G^{-2n}(\mathbf{1}(n))$ pour tout $n \in \mathbf{Z}$.
Il s'ensuit que
$$
H^{2r}(X)(r) = G^{2r}(h(X)) \otimes G^{-2r}(\mathbf{1}(r)) =
G^0(h(X)(r))
$$,
formule dont nous ferons fréquemment usage ci-dessous.

\medskip\noindent
Soit $X$ un schéma projectif lisse sur $k$. D'après
le lemme \ref{weil-lemma-composition-correspondences}, on a
$$
\CH^r(X) \otimes \mathbf{Q} = \text{Corr}^r(\Spec(k), X) =
\Hom(\mathbf{1}(-r), h(X)) = \Hom(\mathbf{1}, h(X)(r))
$$
En appliquant le foncteur $G$, on obtient une application à valeurs dans
$\Hom(G(\mathbf{1}), G(h(X)(r)))$.
En prenant l'image de $1$ dans $G^0(\mathbf{1}) = F$
à valeurs dans $G^0(h(X)(r)) = H^{2r}(X)(r)$, on obtient
$$
\gamma :
\CH^r(X) \otimes \mathbf{Q} \longrightarrow H^{2r}(X)(r)
$$
C'est la donnée (D2).

\medskip\noindent
Soit $X$ un schéma projectif lisse non vide sur $k$,
équidimensionnel de dimension $d$. D'après
le lemme \ref{weil-lemma-composition-correspondences}, on a
$$
\Mor(h(X)(d), \mathbf{1}) = \Mor((X, 1, d), (\Spec(k), 1, 0)) =
\text{Corr}^{-d}(X, \Spec(k)) = \CH_d(X)
$$
La classe du cycle $[X]$ dans $\CH_d(X)$ définit donc un morphisme
$h(X)(d) \to \mathbf{1}$. En appliquant $G$ et en prenant les composantes de degré $0$,
on obtient
$$
H^{2d}(X)(d) = G^0(h(X)(d)) \longrightarrow G^0(\mathbf{1}) = F
$$
Cette application $\int_X : H^{2d}(X)(d) \to F$ est la donnée (D3).

\medskip\noindent
Soit $X$ un schéma projectif lisse sur $k$,
non vide et équidimensionnel de dimension $d$. D'après le lemme \ref{weil-lemma-dual},
$h(X)(d)$ est un dual à gauche de $h(X)$. Par conséquent,
$G(h(X)(d)) = H^*(X) \otimes_F F(d)[2d]$
est un dual à gauche de $H^*(X)$ dans la catégorie des espaces vectoriels gradués sur $F$.
Ici, $[n]$ désigne le foncteur de décalage des espaces vectoriels gradués.
D'après Homologie, lemme \ref{homology-lemma-left-dual-graded-vector-spaces},
on a $\sum_i \dim_F H^i(X) < \infty$, et
$\epsilon : h(X)(d) \otimes h(X) \to \mathbf{1}$ donne
des accouplements non dégénérés $H^{2d - i}(X)(d) \otimes_F H^i(X) \to F$.
Dans la démonstration du lemme \ref{weil-lemma-dual}, nous avons vu que
$\epsilon$ est donné par $[\Delta]$ au moyen des identifications
$$
\Hom(h(X)(d) \otimes h(X), \mathbf{1}) =
\text{Corr}^{-d}(X \times X, \Spec(k)) =
\CH_d(X \times X)
$$
Ainsi, $\epsilon$ est la composée de $[X] : h(X)(d) \to \mathbf{1}$
et de $h(\Delta)(d) : h(X)(d) \otimes h(X) \to h(X)(d)$. Il s'ensuit
que les accouplements ci-dessus sont donnés par le cup-produit suivi de
$\int_X$. Cela démontre l'axiome (A).

\medskip\noindent
L'axiome (B) résulte de la compatibilité de $G$ avec les structures
tensorielles, supposée par hypothèse, et de notre construction du cup-produit ci-dessus.

\medskip\noindent
Axiome (C). Notre construction de $\gamma$ prend un cycle $\alpha$ sur $X$,
l'interprète comme une correspondance $a$ de $\Spec(k)$ dans $X$ d'un certain degré,
puis lui applique $G$. Si $f : Y \to X$ est un morphisme de schémas projectifs lisses
non vides et équidimensionnels sur $k$, alors
$f^!\alpha$ est l'image directe (!) de $\alpha$
par la correspondance $[\Gamma_f]$ de $X$ dans $Y$ ; voir
le lemme \ref{weil-lemma-functor-and-cycles}. Ainsi,
$f^!\alpha$, considéré comme une correspondance de $\Spec(k)$ dans $Y$,
est égal à $a \circ [\Gamma_f]$ ; voir
le lemme \ref{weil-lemma-composition-correspondences}.
Puisque $G$ est un foncteur, on en déduit que
$\gamma$ est compatible avec les images réciproques, ce qui est l'axiome (C)(a).

\medskip\noindent
Soit $f : Y \to X$ un morphisme de schémas projectifs lisses
non vides et équidimensionnels sur $k$,
et soit $\beta \in \CH^r(Y)$ un cycle sur $Y$. Il faut montrer que
$$
\int_Y \gamma(\beta) \cup f^*c = \int_X \gamma(f_*\beta) \cup c
$$
pour tout $c \in H^*(X)$. Soient $a, a^t, \eta_X, \eta_Y, [X], [Y]$
comme dans le lemme \ref{weil-lemma-prep-dual}.
Soit $b$ le cycle $\beta$ considéré comme une correspondance de $\Spec(k)$ dans $Y$
de degré $r$. Alors $f_*\beta$, considéré comme une correspondance de
$\Spec(k)$ dans $X$, est égal à $a^t \circ b$ ; voir
les lemmes \ref{weil-lemma-functor-and-cycles} et
\ref{weil-lemma-composition-correspondences}.
L'égalité affichée ci-dessus sera vérifiée si l'on montre que
$$
h(X) = \mathbf{1} \otimes h(X)
\xrightarrow{b \otimes 1}
h(Y)(r) \otimes h(X)
\xrightarrow{1 \otimes a}
h(Y)(r) \otimes h(Y)
\xrightarrow{\eta_Y}
h(Y)(r)
\xrightarrow{[Y]}
\mathbf{1}(r - e)
$$
est égal à
$$
h(X) = \mathbf{1} \otimes h(X)
\xrightarrow{a^t \circ b \otimes 1}
h(X)(r + d - e) \otimes h(X)
\xrightarrow{\eta_X}
h(X)(r + d - e)
\xrightarrow{[X]}
\mathbf{1}(r - e)
$$
Cela résulte immédiatement du lemme \ref{weil-lemma-prep-dual}.
L'axiome (C)(b) est donc vérifié.

\medskip\noindent
Pour démontrer l'axiome (C)(c), nous utilisons la discussion de
la remarque \ref{weil-remark-replace-cup-product-classical}.
Il suffit donc de démontrer que $\gamma$ est compatible avec
les produits extérieurs. Soient $X$ et $Y$ des schémas projectifs lisses
non vides sur $k$, et soient $\alpha$ et $\beta$ des cycles sur ces schémas. Notons
$a$ et $b$ les correspondances associées de $\Spec(k)$ dans
$X$ et $Y$. Alors $\alpha \times \beta$ correspond à
la correspondance $a \otimes b$ de $\Spec(k)$ dans $X \otimes Y = X \times Y$.
La propriété voulue résulte donc de la compatibilité de $G$
avec les structures tensorielles de part et d'autre.

\medskip\noindent
L'axiome (C)(d) résulte de ce que le cycle $[\Spec(k)]$
correspond au morphisme identique de $h(\Spec(k))$.
Cela achève la démonstration du lemme.
\end{proof}

\begin{lemma}
\label{weil-lemma-from-weil-to-functor}
Soit $k$ un corps et soit $F$ un corps de caractéristique $0$. À partir de données
(D0), (D1), (D2) et (D3) satisfaisant à (A), (B) et (C),
on peut construire un foncteur $\mathbf{Q}$-linéaire
$$
G : M_k \longrightarrow \text{espaces vectoriels gradués sur }F\text{}
$$
monoïdal symétrique tel que $H^*(X) = G(h(X))$.
\end{lemma}

\begin{proof}
La démonstration de ce lemme est la même que celle du
lemme \ref{weil-lemma-from-weil-to-functor-classical} ;
nous conseillons vivement au lecteur de lire plutôt la démonstration de ce lemme.

\medskip\noindent
D'après le lemme \ref{weil-lemma-characterize-motives}, il suffit de construire un foncteur
$G$ sur la catégorie des schémas projectifs lisses sur $k$,
dont les morphismes sont les correspondances de degré $0$, tel que
l'image de $G(c_2)$ dans $G(\mathbf{P}^1_k)$ soit un espace vectoriel gradué
inversible sur $F$.

\medskip\noindent
Soit $X$ un schéma projectif lisse sur $k$. Il existe une décomposition
canonique
$$
X = \coprod\nolimits_{0 \leq d \le \dim(X)} X_d
$$
en sous-schémas ouverts et fermés tels que $X_d$ soit équidimensionnel
de dimension $d$. Le lemme \ref{weil-lemma-weil-additive} donne alors
$$
H^*(X) \longrightarrow \prod\nolimits_{0 \leq d \le \dim(X)} H^*(X_d)
$$
Si $Y$ est un second schéma projectif lisse sur $k$
et si l'on décompose de même $Y = \coprod Y_e$, alors
$$
\text{Corr}^0(X, Y) = \bigoplus \text{Corr}^0(X_d, Y_e)
$$
On a aussi $X \otimes Y = \coprod X_d \otimes Y_e$ dans
la catégorie des correspondances. Ces observations montrent
qu'il suffit de construire $G$ sur la catégorie dont les objets
sont les schémas projectifs lisses équidimensionnels sur $k$
et les morphismes les correspondances de degré $0$.
(Nous omettons certains détails.)

\medskip\noindent
Pour un schéma projectif lisse équidimensionnel
$X$ sur $k$, on pose $G(X) = H^*(X)$. Observons que $G(X) = 0$
si $X = \emptyset$ (lemme \ref{weil-lemma-unit}). Les applications
de source ou de but $G(\emptyset)$ sont donc nulles, et nous pouvons
supposer désormais nos schémas non vides dans la discussion ci-dessous.

\medskip\noindent
Pour une correspondance $c \in \text{Corr}^0(X, Y)$ entre
schémas projectifs lisses non vides et équidimensionnels sur $k$,
considérons l'application $G(c) : G(X) = H^*(X) \to G(Y) = H^*(Y)$
donnée par la règle
$$
a \longmapsto
G(c)(a) = \text{pr}_{2, *}(\gamma(c) \cup \text{pr}_1^*a)
$$
Il est clair que $G(c)$ est additive en $c$, donc $\mathbf{Q}$-linéaire.
La compatibilité de $\gamma$ avec les images réciproques, les images directes
et les produits d'intersection, donnée par les axiomes (C)(a), (C)(b) et (C)(c),
montre que
$G(c' \circ c) = G(c') \circ G(c)$ si $c' \in \text{Corr}^0(Y, Z)$.
En effet, pour $a \in H^*(X)$, on a
\begin{align*}
(G(c') \circ G(c))(a)
& =
\text{pr}^{23}_{3, *}(\gamma(c') \cup
\text{pr}^{23, *}_2(\text{pr}^{12}_{2, *}(\gamma(c) \cup
\text{pr}^{12, *}_1a))) \\
& =
\text{pr}^{23}_{3, *}(\gamma(c') \cup
\text{pr}^{123}_{23, *}(\text{pr}^{123, *}_{12}(\gamma(c) \cup
\text{pr}^{12, *}_1 a))) \\
& =
\text{pr}^{23}_{3, *}
\text{pr}^{123}_{23, *}(
\text{pr}^{123, *}_{23}\gamma(c') \cup
\text{pr}^{123, *}_{12}\gamma(c) \cup
\text{pr}^{123, *}_1 a) \\
& =
\text{pr}^{23}_{3, *}
\text{pr}^{123}_{23, *}(
\gamma(\text{pr}^{123, *}_{23}c') \cup
\gamma(\text{pr}^{123, *}_{12}c) \cup
\text{pr}^{123, *}_1 a) \\
& =
\text{pr}^{13}_{3, *}
\text{pr}^{123}_{13, *}(
\gamma(\text{pr}^{123, *}_{23}c' \cdot \text{pr}^{123, *}_{12}c) \cup
\text{pr}^{123, *}_1 a) \\
& =
\text{pr}^{13}_{3, *}(
\gamma(\text{pr}^{123}_{13, *}(\text{pr}^{123, *}_{23}c' \cdot
\text{pr}^{123, *}_{12}c)) \cup
\text{pr}^{13, *}_1 a) \\
& =
G(c' \circ c)(a)
\end{align*},
avec des notations évidentes. La première égalité résulte des définitions.
La deuxième résulte de l'égalité
$\text{pr}^{23, *}_2 \circ \text{pr}^{12}_{2, *} =
\text{pr}^{123}_{23, *} \circ \text{pr}^{123, *}_{12}$,
qui découle immédiatement de la description de l'image directe
par les projections donnée dans le lemme \ref{weil-lemma-pr2star}.
La troisième égalité résulte du lemme \ref{weil-lemma-pushforward}
et du fait que $H^*$ est un foncteur.
La quatrième égalité résulte de l'axiome (C)(a) et du fait que
le morphisme de Gysin coïncide avec l'image réciproque plate pour les morphismes plats
(Homologie de Chow, lemme \ref{chow-lemma-lci-gysin-flat}).
La cinquième utilise l'axiome (C)(c), ainsi que
le lemme \ref{weil-lemma-pushforward}, pour obtenir
$\text{pr}^{23}_{3, *} \circ \text{pr}^{123}_{23, *} =
\text{pr}^{13}_{3, *} \circ \text{pr}^{123}_{13, *}$.
La sixième utilise la formule de projection du
lemme \ref{weil-lemma-pushforward}, ainsi que
l'axiome (C)(b), pour obtenir $
\text{pr}^{123}_{13, *}
\gamma(\text{pr}^{123, *}_{23}c' \cdot \text{pr}^{123, *}_{12}c) =
\gamma(\text{pr}^{123}_{13, *}(
\text{pr}^{123, *}_{23}c' \cdot \text{pr}^{123, *}_{12}c))$.
Enfin, la dernière égalité est la définition.

\medskip\noindent
Pour achever la démonstration que $G$ est un foncteur,
il faut montrer qu'il préserve les identités. Autrement dit, si
$1 = [\Delta] \in \text{Corr}^0(X, X)$ est l'identité dans la catégorie
des correspondances (lemme \ref{weil-lemma-category-correspondences}),
il faut montrer que $G([\Delta]) = \text{id}$.
Cela résulte de la détermination
de $\gamma([\Delta])$ dans le lemme \ref{weil-lemma-class-diagonal}
et du lemme \ref{weil-lemma-pr2star}.
La construction de $G$ comme foncteur sur les
schémas projectifs lisses sur $k$ et les correspondances de degré $0$ est ainsi achevée.

\medskip\noindent
D'après le lemme \ref{weil-lemma-base},
$G(\Spec(k)) = H^*(\Spec(k))$ est canoniquement isomorphe à $F$
comme $F$-algèbre. L'axiome de K\"unneth (B)(a)
montre que notre foncteur est compatible avec les produits tensoriels.
Il est donc un foncteur monoïdal symétrique.

\medskip\noindent
Il reste à vérifier que l'image de $G(c_2)$ dans
$G(\mathbf{P}^1_k) = H^*(\mathbf{P}^1_k)$
est un espace vectoriel gradué inversible sur $F$ (en particulier, nous ne savons pas encore
que $G$ se prolonge à $M_k$). D'après le lemme \ref{weil-lemma-cohomology-P1},
la cohomologie n'est non nulle qu'en degrés $0$ et $2$,
où elle est de dimension $1$. L'égalité $1 = c_0 + c_2$ est une décomposition
de l'identité en somme d'idempotents orthogonaux dans
$\text{Corr}^0(\mathbf{P}^1_k, \mathbf{P}^1_k)$ ; voir
l'exemple \ref{weil-example-decompose-P1}. De plus, on a $c_0 = a \circ b$, où
$a \in \text{Corr}^0(\Spec(k), \mathbf{P}^1_k)$ et
$b \in \text{Corr}^0(\mathbf{P}^1_k, \Spec(k))$, et où
$b \circ a = 1$ dans $\text{Corr}^0(\Spec(k), \Spec(k))$ ; voir la démonstration du
lemme \ref{weil-lemma-inverse-h2}. Ainsi, $G(c_0)$ est le projecteur
sur la composante de degré $0$. Il s'ensuit que $G(c_2)$ est
le projecteur sur la composante de degré $2$, ce qui achève la démonstration.
\end{proof}

\begin{proposition}
\label{weil-proposition-weil-cohomology-theory}
Soit $k$ un corps et soit $F$ un corps de caractéristique $0$. Il existe
une correspondance biunivoque ($1$ à $1$) entre :
\begin{enumerate}
\item les données (D0), (D1), (D2) et (D3) satisfaisant à (A), (B) et (C) ;
\item les foncteurs monoïdaux symétriques $\mathbf{Q}$-linéaires
$$
G : M_k \longrightarrow \text{espaces vectoriels gradués sur }F\text{}
$$
tels que $G(\mathbf{1}(1))$ ne soit non nul qu'en degré $-2$.
\end{enumerate}
\end{proposition}

\begin{proof}
Étant donné $G$ comme dans (2), on pose $H^*(X) = G(h(X))$ et l'on obtient des données
(D0), (D1), (D2) et (D3) satisfaisant à (A), (B) et (C) ;
voir le lemme \ref{weil-lemma-from-functor-to-weil} et sa démonstration.

\medskip\noindent
Réciproquement, à partir de données (D0), (D1), (D2) et (D3)
satisfaisant à (A), (B) et (C), on obtient un foncteur $G$ comme dans (2)
par la construction de la démonstration du lemme \ref{weil-lemma-from-weil-to-functor}.

\medskip\noindent
Nous omettons la démonstration détaillée que ces constructions
sont inverses l'une de l'autre.
\end{proof}











\section{Propriétés complémentaires}
\label{weil-section-further}

\noindent
Nous démontrons dans cette section quelques conséquences supplémentaires
de données (D0), (D1), (D2) et (D3) satisfaisant à (A), (B) et (C),
comme à la section \ref{weil-section-axioms}.

\begin{lemma}
\label{weil-lemma-trace-disjoint-union}
Supposons que les données (D0), (D1), (D2) et (D3) satisfassent à (A), (B) et (C).
Soient $X, Y$ des schémas projectifs lisses non vides, tous deux équidimensionnels
de dimension $d$ sur $k$. Alors $\int_{X \amalg Y} = \int_X + \int_Y$.
\end{lemma}

\begin{proof}
Notons $i : X \to X \amalg Y$ et $j : Y \to X \amalg Y$ les coprojections.
D'après le lemme \ref{weil-lemma-weil-additive}, l'application
$(i^*, j^*) : H^*(X \amalg Y) \to H^*(X) \times H^*(Y)$ est un isomorphisme.
L'énoncé signifie que, par l'isomorphisme
$(i^*, j^*) : H^{2d}(X \amalg Y)(d) \to H^{2d}(X)(d) \oplus H^{2d}(Y)(d)$,
l'application $\int_X + \int_Y$ se transforme en $\int_{X \amalg Y}$.
En effet, on a
$$
\int_{X \amalg Y} a =
\int_{X \amalg Y} i_*(i^*a) + j_*(j^*a) =
\int_X i^*a + \int_Y j^*a
$$,
où l'égalité $a = i_*(i^*a) + j_*(j^*a)$ a été établie dans
la démonstration du lemme \ref{weil-lemma-weil-additive}.
\end{proof}

\begin{lemma}
\label{weil-lemma-dim-0}
Supposons que les données (D0), (D1), (D2) et (D3) satisfassent à (A), (B) et (C).
Soit $X$ un schéma projectif lisse de dimension zéro sur $k$.
Alors :
\begin{enumerate}
\item $H^i(X) = 0$ pour $i \not = 0$ ;
\item $H^0(X)$ est une algèbre finie séparable sur $F$ ;
\item $\dim_F H^0(X) = \deg(X \to \Spec(F))$ ;
\item $\int_X : H^0(X) \to F$ est l'application trace ;
\item $\gamma([X]) = 1$ ;
\item $\int_X \gamma([X]) = \deg(X \to \Spec(k))$.
\end{enumerate}
\end{lemma}

\begin{proof}
On peut écrire $X = \Spec(k')$, où $k'$ est une algèbre finie séparable
sur $k$. Observons que $\deg(X \to \Spec(k)) = [k' : k]$.
Choisissons une extension galoisienne finie $k''/k$ contenant chacun des
facteurs de $k'$. (Rappelons qu'une $k$-algèbre finie séparable est
un produit d'extensions finies séparables de $k$.)
Posons $\Sigma = \Hom_k(k', k'')$. On obtient
$$
k' \otimes_k k'' = \prod\nolimits_{\sigma \in \Sigma} k''
$$
En posant $Y = \Spec(k'')$, l'axiome (B)(a) et le lemme \ref{weil-lemma-weil-additive} donnent
$$
H^*(X) \otimes_F H^*(Y) =
\prod\nolimits_{\sigma \in \Sigma} H^*(Y)
$$
comme $F$-algèbres graduées commutatives. D'après le lemme \ref{weil-lemma-unit},
la $F$-algèbre $H^*(Y)$ est non nulle. En comparant les dimensions des deux membres
de l'égalité affichée, on en déduit que $H^*(X)$ est concentré en degré $0$
et que $\dim_F H^0(X) = [k' : k]$. En appliquant cela à $Y$, on obtient
$H^*(Y) = H^0(Y)$. Puisque
$$
H^0(X) \otimes_F H^0(Y) = H^0(Y) \times \ldots \times H^0(Y)
$$
comme $F$-algèbres, il s'ensuit que $H^0(X)$ est une $F$-algèbre séparable,
car cela se vérifie après le changement de base fidèlement plat
$F \to H^0(Y)$.

\medskip\noindent
L'isomorphisme affiché ci-dessus est donné par l'application
$$
H^0(X) \otimes_F H^0(Y) \longrightarrow
\prod\nolimits_{\sigma \in \Sigma} H^0(Y),\quad
a \otimes b \longmapsto \prod\nolimits_\sigma \Spec(\sigma)^*a \cup b
$$
Par cet isomorphisme, on a $\int_{X \times Y} = \sum_\sigma \int_Y$, d'après
le lemme \ref{weil-lemma-trace-disjoint-union}. Ainsi,
$$
\int_X a = \text{pr}_{1, *}(a \otimes 1) = \sum \Spec(\sigma)^*a
$$
dans $H^0(Y)$ ; la première égalité résulte du lemme \ref{weil-lemma-pr2star}
et la seconde de l'observation précédente. Choisissons une
clôture algébrique $\overline{F}$ et
un homomorphisme de $F$-algèbres $\tau : H^0(Y) \to \overline{F}$.
Par changement de base, l'isomorphisme ci-dessus devient
$$
H^0(X) \otimes_F \overline{F} \longrightarrow
\prod\nolimits_{\sigma \in \Sigma} \overline{F},\quad
a \otimes b \longmapsto \prod\nolimits_\sigma \tau(\Spec(\sigma)^*a) b
$$
Il s'ensuit que les applications $a \mapsto \tau(\Spec(\sigma)^*a)$ forment l'ensemble complet
des plongements de $H^0(X)$ dans $\overline{F}$. En appliquant $\tau$
à la formule de $\int_X a$ obtenue ci-dessus, on en déduit
que $\int_X$ est l'application trace.
D'après le lemme \ref{weil-lemma-unit}, on a $\gamma([X]) = 1$.
Enfin, on a $\int_X \gamma([X]) = \deg(X \to \Spec(k))$,
puisque $\gamma([X]) = 1$ et que la trace de $1$ est égale à $[k' : k]$.
\end{proof}

\begin{lemma}
\label{weil-lemma-degrees-cycles}
Supposons que les données (D0), (D1), (D2) et (D3) satisfassent à (A), (B) et (C).
Soit $X$ un schéma projectif lisse non vide,
équidimensionnel de dimension $d$ sur $k$. Le diagramme
$$
\xymatrix{
\CH^d(X) \ar[r]_-\gamma \ar@{=}[d] &
H^{2d}(X)(d) \ar[d]^{\int_X} \\
\CH_0(X) \ar[r]^\deg & F
}
$$
est commutatif, où $\deg : \CH_0(X) \to \mathbf{Z}$ est le degré des
zéro-cycles étudié dans Homologie de Chow, section
\ref{chow-section-degree-zero-cycles}.
\end{lemma}

\begin{proof}
Soit $x$ un point fermé de $X$ dont le corps résiduel est séparable
sur $k$. Considérons $x$ comme un schéma et notons $i : x \to X$ le morphisme d'inclusion.
Pour éviter toute confusion, notons $\gamma' : \CH_0(x) \to H^0(x)$ l'application classe de cycle
de $x$. On a alors
$$
\int_X \gamma([x]) = \int_X \gamma(i_*[x]) =
\int_X i_*\gamma'([x]) = \int_x \gamma'([x]) = \deg(x \to \Spec(k))
$$
La deuxième égalité est l'axiome (C)(b), et la troisième
la définition de $i_*$ en cohomologie. La dernière égalité résulte du
lemme \ref{weil-lemma-dim-0}. Cela démontre le lemme,
car $\CH_0(X)$ est engendré par les classes de points $x$ comme ci-dessus,
d'après le lemme \ref{weil-lemma-generated-by-separable}.
\end{proof}

\begin{lemma}
\label{weil-lemma-square-diagonal}
Supposons que les données (D0), (D1), (D2) et (D3) satisfassent à (A), (B) et (C).
Soit $X$ un schéma projectif lisse non vide sur $k$,
équidimensionnel de dimension $d$. On a
$$
\sum\nolimits_i (-1)^i\dim_F H^i(X) =
\deg(\Delta \cdot \Delta) = \deg(c_d(\mathcal{T}_{X/k}))
$$
\end{lemma}

\begin{proof}
Démontrons l'égalité de droite. On a
$[\Delta] \cdot [\Delta] = \Delta_*(\Delta^![\Delta])$
(Homologie de Chow, lemme \ref{chow-lemma-intersect-regularly-embedded}).
Puisque $\Delta_*$ préserve les degrés des $0$-cycles, il suffit de calculer
le degré de $\Delta^![\Delta]$. La classe $\Delta^![\Delta]$ est donnée
par le cap-produit de $[\Delta]$ avec
la classe de Chern de degré maximal du faisceau normal de $\Delta \subset X \times X$
(Homologie de Chow, lemme \ref{chow-lemma-gysin-fundamental}).
Puisque le faisceau conormal de $\Delta$ est $\Omega_{X/k}$
(Morphismes, lemme \ref{morphisms-lemma-differentials-diagonal}),
le faisceau normal est égal au faisceau tangent
$\mathcal{T}_{X/k} = \SheafHom_{\mathcal{O}_X}(\Omega_{X/k}, \mathcal{O}_X)$,
comme voulu.

\medskip\noindent
Démontrons l'égalité de gauche. D'après le lemme \ref{weil-lemma-degrees-cycles}, on a
\begin{align*}
\deg([\Delta] \cdot [\Delta])
& =
\int_{X \times X} \gamma([\Delta]) \cup \gamma([\Delta]) \\
& =
\int_{X \times X} \Delta_*1 \cup \gamma([\Delta]) \\
& =
\int_{X \times X} \Delta_*(\Delta^*\gamma([\Delta])) \\
& =
\int_X \Delta^*\gamma([\Delta])
\end{align*}
Nous avons utilisé les lemmes \ref{weil-lemma-push-unit} et
\ref{weil-lemma-pushforward}.
Écrivons $\gamma([\Delta]) = \sum  e_{i, j} \otimes e'_{2d - i , j}$
comme dans le lemme \ref{weil-lemma-class-diagonal}.
En rappelant que $\Delta^*$ est donné par le cup-produit
(remarque \ref{weil-remark-replace-cup-product}), on obtient
$$
\int_X \sum\nolimits_{i, j} e_{i, j} \cup e'_{2d - i, j} =
\sum\nolimits_{i, j} \int_X e_{i, j} \cup e'_{2d - i, j} =
\sum\nolimits_{i, j} (-1)^i = \sum (-1)^i\beta_i
$$,
comme voulu.
\end{proof}

\begin{lemma}
\label{weil-lemma-algebra-relations}
Soit $F$ un corps de caractéristique $0$.
Soient $F'$ et $F_i$, pour $i = 1, \ldots, r$,
des $F$-algèbres finies séparables. Soit $A$ une $F$-algèbre finie.
Soient $\sigma, \sigma' : A \to F'$ et $\sigma_i : A \to F_i$
des homomorphismes de $F$-algèbres. Supposons $\sigma$ et $\sigma'$ surjectifs.
S'il existe une relation
$$
\text{Tr}_{F'/F} \circ \sigma - \text{Tr}_{F'/F} \circ \sigma' =
n(\sum m_i \text{Tr}_{F_i/F} \circ \sigma_i)
$$,
où $n > 1$ et les $m_i$ sont des entiers, alors $\sigma = \sigma'$.
\end{lemma}

\begin{proof}
On peut écrire $A = \prod A_j$ comme un produit fini de
$F$-algèbres artiniennes locales $(A_j, \mathfrak m_j, \kappa_j)$ ; voir
Algèbre, lemme \ref{algebra-lemma-finite-dimensional-algebra} et
proposition \ref{algebra-proposition-dimension-zero-ring}.
Notons $A' = \prod \kappa_j$, où le produit porte sur les $j$
tels que $\kappa_j/k$ soit séparable. Chacune des applications
$\sigma, \sigma', \sigma_i$ se factorise alors par l'application
$A \to A'$. En remplaçant $A$ par $A'$, on peut supposer
que $A$ est une $F$-algèbre finie séparable.

\medskip\noindent
Choisissons une clôture algébrique $\overline{F}$. Posons
$\overline{A} = A \otimes_F \overline{F}$,
$\overline{F}' = F' \otimes_F \overline{F}$ et
$\overline{F}_i = F_i \otimes_F \overline{F}$.
Par changement de base de $\sigma$, $\sigma'$ et $\sigma_i$,
on obtient des homomorphismes de $\overline{F}$-algèbres $\overline{A} \to \overline{F}'$
et $\overline{A} \to \overline{F}_i$. De plus,
$\text{Tr}_{\overline{F}'/\overline{F}}$ est le changement de base
de $\text{Tr}_{F'/F}$, et de même pour
$\text{Tr}_{F_i/F}$. On peut donc remplacer
$F$ par $\overline{F}$ et se ramener au cas traité
dans le paragraphe suivant.

\medskip\noindent
Supposons $F$ algébriquement clos et $A$ une $F$-algèbre finie séparable.
Alors chacun de $A$, $F'$ et $F_i$ est un produit de copies de $F$.
Disons qu'un élément $e$ d'un produit
$F \times \ldots \times F$ de copies de $F$ est un idempotent minimal
s'il engendre l'un des facteurs, c'est-à-dire si
$e = (0, \ldots, 0, 1, 0, \ldots, 0)$. Soit $e \in A$ un idempotent minimal.
Puisque $\sigma$ et $\sigma'$
sont surjectifs, $\sigma(e)$ et $\sigma'(e)$ sont des idempotents
minimaux ou sont nuls. Si $\sigma \not = \sigma'$, on peut choisir
un idempotent minimal $e \in A$ tel que $\sigma(e) = 0$ et
$\sigma'(e) \not = 0$, ou inversement. Alors
$\text{Tr}_{F'/F}(\sigma(e)) = 0$ et
$\text{Tr}_{F'/F}(\sigma'(e)) = 1$, ou inversement.
D'autre part, $\sigma_i(e)$ est un idempotent,
donc $\text{Tr}_{F_i/F}(\sigma_i(e)) = r_i$ est un entier.
Il en résulte que
$$
-1 = \sum n m_i r_i = n (\sum m_i r_i)
\quad\text{ou}\quad
1 = \sum n m_i r_i = n (\sum m_i r_i)
$$,
ce qui est impossible.
\end{proof}

\begin{lemma}
\label{weil-lemma-relations-classes-points}
Supposons que les données (D0), (D1), (D2) et (D3) satisfassent à (A), (B) et (C).
Soit $k'/k$ une extension finie séparable.
Soit $X$ un schéma projectif lisse sur $k'$.
Soient $x, x' \in X$ des points $k'$-rationnels.
Si $\gamma(x) \not = \gamma(x')$, alors
$[x] - [x']$ n'est divisible par aucun entier $n > 1$ dans $\CH_0(X)$.
\end{lemma}

\begin{proof}
Si $x$ et $x'$ appartiennent à des composantes irréductibles distinctes de $X$,
le résultat est évident. On peut donc supposer $X$ irréductible de dimension $d$.
Supposons $[x] - [x']$ divisible par $n > 1$ dans $\CH_0(X)$.
On peut écrire $[x] - [x'] = n(\sum m_i [x_i])$ dans $\CH_0(X)$,
pour certains points fermés $x_i \in X$
dont les corps résiduels sont séparables sur $k$, d'après
le lemme \ref{weil-lemma-generated-by-separable}.
Alors
$$
\gamma([x]) - \gamma([x']) = n (\sum m_i \gamma([x_i]))
$$
dans $H^{2d}(X)(d)$. Notons $i^*, (i')^*, i_i^*$ les applications d'image réciproque
$H^0(X) \to H^0(x)$, $H^0(X) \to H^0(x')$ et $H^0(X) \to H^0(x_i)$.
Rappelons que $H^0(x)$ est une $F$-algèbre finie séparable
et que $\int_x : H^0(x) \to F$ est l'application trace
(lemme \ref{weil-lemma-dim-0}), que nous noterons $\text{Tr}_x$.
Il en est de même pour $x'$ et $x_i$. Par la dualité de Poincar\'e sous la forme
de l'axiome (A)(b), l'égalité ci-dessus est duale de
$$
\text{Tr}_x \circ i^* - \text{Tr}_{x'} \circ (i')^* =
n(\sum m_i \text{Tr}_{x_i} \circ i_i^*)
$$,
égalité dans $\Hom_F(H^0(X), F)$. Enfin, observons que
$i^*$ et $(i')^*$ sont surjectives, puisque $x$ et $x'$ sont des points $k'$-rationnels
et que les composées $H^0(\Spec(k')) \to H^0(X) \to H^0(x)$
et $H^0(\Spec(k')) \to H^0(X) \to H^0(x')$ sont donc des isomorphismes.
Le lemme \ref{weil-lemma-algebra-relations} donne $i^* = (i')^*$,
ce qui contredit l'hypothèse $\gamma([x]) \not = \gamma([x'])$.
\end{proof}

\begin{lemma}
\label{weil-lemma-classes-points}
Supposons que les données (D0), (D1), (D2) et (D3) satisfassent à (A), (B) et (C).
Soit $k'/k$ une extension finie séparable et soit $X$ un schéma projectif lisse
géométriquement irréductible sur $k'$, de dimension $d$.
Alors $\gamma : \CH_0(X) \to H^{2d}(X)(d)$ se factorise par
$\deg : \CH_0(X) \to \mathbf{Z}$.
\end{lemma}

\begin{proof}
D'après le lemme \ref{weil-lemma-generated-by-separable}, il suffit de montrer que,
pour des points fermés $x, x' \in X$ dont les corps résiduels sont séparables sur $k$,
on a $\deg(x') \gamma([x]) = \deg(x) \gamma([x'])$.

\medskip\noindent
Ramenons-nous d'abord au cas de points $k'$-rationnels. Soit $k''/k'$ une
extension galoisienne telle que $\kappa(x)$ et $\kappa(x')$ se plongent dans $k''$
sur $k$. Posons $Y = X \times_{\Spec(k')} \Spec(k'')$ et notons $p : Y \to X$
la projection. Par le choix de $k''/k'$, il existe un
point $k''$-rationnel $y$, resp.\ $y'$, de $Y$, d'image $x$, resp.\ $x'$.
Alors $p_*[y] = [k'' : \kappa(x)][x]$ et
$p_*[y'] = [k'' : \kappa(x')][x']$ dans $\CH_0(X)$.
Par la compatibilité avec les images directes de l'axiome (C)(b),
il suffit de démontrer $\gamma([y]) = \gamma([y'])$ dans $\CH^{2d}(Y)(d)$.
On est ainsi ramené à la discussion du paragraphe suivant.

\medskip\noindent
Supposons $x$ et $x'$ des points $k'$-rationnels. D'après
le lemme \ref{weil-lemma-divide-difference-points},
il existe une extension finie séparable de corps $k''/k'$
telle que l'image réciproque $[y] - [y']$
de la différence $[x] - [x']$ devienne divisible
par un entier $n > 1$ sur $Y = X \times_{\Spec(k')} \Spec(k'')$.
(Notons que $y, y' \in Y$ sont des points $k''$-rationnels.)
D'après le lemme \ref{weil-lemma-relations-classes-points}, on a
$\gamma([y]) = \gamma([y'])$ dans $H^{2d}(Y)(d)$.
Par la compatibilité avec l'image directe de l'axiome (C)(b),
on en déduit la même assertion pour $x$ et $x'$.
\end{proof}

\begin{lemma}
\label{weil-lemma-injective}
Supposons que les données (D0), (D1), (D2) et (D3) satisfassent à (A), (B) et (C). Soit
$f : X \to Y$ un morphisme dominant de schémas projectifs lisses irréductibles
sur $k$. Alors $H^*(Y) \to H^*(X)$ est injective.
\end{lemma}

\begin{proof}
Il existe un sous-schéma fermé intègre $Z \subset X$ de même
dimension que $Y$ et qui se surjecte sur $Y$. Ainsi, $f_*[Z] = m[Y]$ pour un certain $m > 0$.
Alors $f_* \gamma([Z]) = m \gamma([Y]) = m$ dans $H^*(Y)$, d'après
le lemme \ref{weil-lemma-unit}. Par la formule de projection
(lemme \ref{weil-lemma-pushforward}),
on a donc $f_*(f^*a \cup \gamma([Z])) = m a$, et l'on conclut.
\end{proof}

\begin{lemma}
\label{weil-lemma-otimes}
Supposons que les données (D0), (D1), (D2) et (D3) satisfassent à (A), (B) et (C). Soient
$k''/k'/k$ des algèbres finies séparables et soit $X$ un
schéma projectif lisse sur $k'$. Alors
$$
H^*(X) \otimes_{H^0(\Spec(k'))} H^0(\Spec(k'')) =
H^*(X \times_{\Spec(k')} \Spec(k''))
$$
\end{lemma}

\begin{proof}
Nous utiliserons les résultats du lemme \ref{weil-lemma-dim-0} sans le mentionner de nouveau.
Écrivons
$$
k' \otimes_k k'' = k'' \times l
$$,
pour une certaine $k'$-algèbre finie séparable $l$. Posons
$F' = H^0(\Spec(k'))$, $F'' = H^0(\Spec(k''))$ et $G = H^0(\Spec(l))$.
Puisque $\Spec(k') \times \Spec(k'') = \Spec(k'') \amalg \Spec(l)$,
l'axiome (B)(a) et le lemme \ref{weil-lemma-weil-additive}
donnent
$$
F' \otimes_F F'' = F'' \times G
$$
L'application de gauche à droite identifie $F''$ à $F' \otimes_{F'} F''$.
De même, on a
$$
H^*(X) \otimes_F F'' = H^*(X \times_{\Spec(k')} \Spec(k''))
\times H^*(X \times_{\Spec(k')} \Spec(l))
$$
comme modules sur $F' \otimes_F F'' = F'' \times G$. Le lemme est démontré.
\end{proof}






















\section{Théories cohomologiques de Weil, II}
\label{weil-section-old}

\noindent
Pour nous, une théorie cohomologique de Weil sera l'analogue d'une
théorie cohomologique de Weil classique (section \ref{weil-section-axioms-classical})
lorsque le corps de base $k$ n'est pas algébriquement clos.
À la section \ref{weil-section-axioms}, nous avons énoncé des axiomes garantissant
que notre théorie cohomologique provient d'un foncteur monoïdal symétrique
sur la catégorie des motifs sur $k$. Il manque encore à ces axiomes
la condition $H^i(X) = 0$ pour $i < 0$, ainsi
qu'une condition sur $H^{2d}(X)(d)$ pour $X$ équidimensionnel de dimension $d$,
correspondant aux axiomes classiques (A)(c) et (A)(d).
Montrons d'abord au lecteur qu'il est nécessaire d'imposer
de telles conditions.

\begin{example}
\label{weil-example-weird-weil}
Soient $k = \mathbf{C}$ et $F = \mathbf{C}$ tous deux égaux au corps
des nombres complexes. Pour $X$ projectif lisse sur $k$, notons
$H^{p, q}(X) = H^q(X, \Omega^p_{X/k})$. Soit $(H')^*$ le foncteur
qui envoie $X$ sur $(H')^*(X) = \bigoplus H^{p, q}(X)$, muni du
cup-produit usuel.
C'est une théorie cohomologique de Weil classique (insérer ici une référence ultérieure).
D'après la proposition \ref{weil-proposition-weil-cohomology-theory-classical},
on obtient un foncteur monoïdal symétrique $\mathbf{Q}$-linéaire $G'$ de $M_k$
dans la catégorie des espaces vectoriels gradués sur $F$. Bien entendu, dans ce cas,
pour tout $M$ dans $M_k$, la valeur $G'(M)$ est naturellement bigraduée, c'est-à-dire
que l'on a
$$
(G')(M) = \bigoplus (G')^{p, q}(M),\quad
(G')^n = \bigoplus\nolimits_{n = p + q} (G')^{p, q}(M)
$$,
où $(G')^{p, q}$ est placé en degré total $p + q$, comme indiqué.
Construisons maintenant un foncteur monoïdal symétrique $\mathbf{Q}$-linéaire
$G$ à valeurs dans les espaces vectoriels gradués sur $F$ en posant
$$
G^n(M) = \bigoplus\nolimits_{n = 3p - q} (G')^{p, q}(M)
$$
Nous omettons la vérification que cela définit un foncteur monoïdal
symétrique (un point technique est que le choix des nombres impairs
$3$ et $-1$ ci-dessus assure la compatibilité du foncteur $G$ avec
les contraintes de commutativité).
Observons que $G(\mathbf{1}(1))$ est toujours placé en degré $-2$ !
Ainsi, par le lemme \ref{weil-lemma-from-functor-to-weil-classical},
on obtient un foncteur $H^*$, des classes de cycles $\gamma$ et des applications traces
satisfaisant à tous les axiomes classiques (A), (B), (C), sauf peut-être
aux axiomes classiques (A)(a) et (A)(d).
Cependant, si $E$ est une courbe elliptique sur $k$, on trouve
$\dim H^{-1}(E) = 1$ : l'axiome (A)(a) est donc effectivement violé.
\end{example}

\begin{lemma}
\label{weil-lemma-H-0-separable}
Supposons que les données (D0), (D1), (D2) et (D3) satisfassent à (A), (B) et (C).
Soit $X$ un schéma projectif lisse sur $k$.
Posons $k' = \Gamma(X, \mathcal{O}_X)$. Les conditions suivantes sont équivalentes :
\begin{enumerate}
\item il existe un nombre fini de points fermés $x_1, \ldots, x_r \in X$,
dont les corps résiduels sont séparables sur $k$, tels que
$H^0(X) \to H^0(x_1) \oplus \ldots \oplus H^0(x_r)$ soit injective ;
\item l'application $H^0(\Spec(k')) \to H^0(X)$ est un isomorphisme.
\end{enumerate}
Si ces conditions sont vérifiées, $H^0(X)$ est une algèbre finie séparable sur $F$.
Si $X$ est équidimensionnel de dimension $d$, les conditions (1) et (2)
équivalent aussi à :
\begin{enumerate}
\item[(3)] les classes des points fermés engendrent $H^{2d}(X)(d)$
comme module sur $H^0(X)$.
\end{enumerate}
\end{lemma}

\begin{proof}
Observons que l'énoncé a un sens, car $k'$ est une algèbre finie séparable
sur $k$ (Variétés, lemme
\ref{varieties-lemma-proper-geometrically-reduced-global-sections}) ;
ainsi, $\Spec(k')$ est lisse et projectif sur $k$.
La compatibilité de $H^*$ avec les sommes directes
(lemmes \ref{weil-lemma-weil-additive} et \ref{weil-lemma-trace-disjoint-union})
montre qu'il suffit de démontrer le lemme lorsque $X$ est connexe.
On peut donc supposer $X$ irréductible, et il faut établir
l'équivalence de (1), (2) et (3). Posons $d = \dim(X)$.
Alors $k'$ est un corps, extension finie séparable
de $k$, et $X$ est géométriquement irréductible sur $k'$ ; voir
Variétés, lemmes
\ref{varieties-lemma-proper-geometrically-reduced-global-sections} et
\ref{varieties-lemma-baby-stein}.

\medskip\noindent
D'après le lemme \ref{weil-lemma-generated-by-separable}, on peut supposer que les points
fermés de (3) ont des corps résiduels séparables sur $k$.
Les axiomes (A)(a) et (A)(b) montrent que (1) et (3) sont équivalentes.

\medskip\noindent
Si (2) est vérifiée, choisissons un point fermé quelconque $x \in X$ dont le corps résiduel
est fini séparable sur $k'$. Alors
$H^0(\Spec(k')) = H^0(X) \to H^0(x)$ est injective, par exemple d'après
le lemme \ref{weil-lemma-injective}.

\medskip\noindent
Supposons vérifiées les conditions équivalentes (1) et (3). Choisissons
$x_1, \ldots, x_r \in X$ comme dans (1), puis une extension finie séparable
$k''/k'$. D'après le lemme \ref{weil-lemma-otimes}, on a
$$
H^0(X) \otimes_{H^0(\Spec(k'))} H^0(\Spec(k'')) =
H^0(X \times_{\Spec(k')} \Spec(k''))
$$
Pour montrer que
$H^0(\Spec(k')) \to H^0(X)$ est un isomorphisme,
on peut donc remplacer $k'$ par $k''$. On peut ainsi supposer que $x_1, \ldots, x_r$
sont des points $k'$-rationnels (chaque $x_i$ est remplacé par plusieurs
points, de sorte que $r$ augmente à cette étape). D'après le lemme \ref{weil-lemma-classes-points},
$\gamma(x_1) = \gamma(x_2) = \ldots = \gamma(x_r)$.
Par l'axiome (A)(b), toutes les applications $H^0(X) \to H^0(x_i)$
sont égales. Cela signifie que (2) est vérifiée.

\medskip\noindent
Enfin, le lemme \ref{weil-lemma-dim-0} entraîne
que $H^0(X)$ est une $F$-algèbre séparable si (1) est vérifiée.
\end{proof}

\begin{lemma}
\label{weil-lemma-negative-cohomology}
Supposons que les données (D0), (D1), (D2) et (D3) satisfassent à (A), (B) et (C).
S'il existe un schéma projectif lisse $Y$ sur $k$ tel que
$H^i(Y)$ soit non nul pour un certain $i < 0$, il existe un schéma projectif lisse
équidimensionnel $X$ sur $k$ pour lequel les conditions équivalentes
du lemme \ref{weil-lemma-H-0-separable} ne sont pas vérifiées pour $X$.
\end{lemma}

\begin{proof}
D'après le lemme \ref{weil-lemma-weil-additive}, on peut supposer $Y$ irréductible,
donc a fortiori équidimensionnel. Si $i$ est impair, en remplaçant
$Y$ par $Y \times Y$, on obtient un exemple où $Y$ est équidimensionnel
et où $i = -2l$ pour un certain $l > 0$. Posons $X = Y \times (\mathbf{P}^1_k)^l$.
L'axiome (B)(a) donne
$$
H^0(X) \supset H^0(Y) \oplus
H^i(Y) \otimes_F H^2(\mathbf{P}^1_k)^{\otimes_F l}
$$,
les deux facteurs directs étant non nuls. Il est donc clair que $H^0(X)$ ne peut être
isomorphe à $H^0$ du spectre de
$\Gamma(X, \mathcal{O}_X) = \Gamma(Y, \mathcal{O}_Y)$,
puisque ce dernier s'envoie dans le premier facteur direct.
\end{proof}

\noindent
Nous sommes ainsi conduits à poser enfin la définition suivante.

\begin{definition}
\label{weil-definition-weil-cohomology-theory}
Soit $k$ un corps et soit $F$ un corps de caractéristique $0$.
Une {\it théorie cohomologique de Weil} sur $k$ à coefficients dans $F$
est définie par des données (D0), (D1), (D2) et (D3) satisfaisant
à la dualité de Poincar\'e, à la formule de K\"unneth et à la compatibilité
avec les classes de cycles, plus précisément aux axiomes (A), (B) et (C)
de la section \ref{weil-section-axioms},
et telles que les conditions équivalentes (1) et (2) du
lemme \ref{weil-lemma-H-0-separable} soient vérifiées pour tout schéma projectif lisse $X$ sur $k$.
\end{definition}

\noindent
D'après le lemme \ref{weil-lemma-negative-cohomology}, cela signifie aussi que tous
les groupes de cohomologie de degré négatif sont nuls. En particulier, si $k$ est
algébriquement clos, une théorie cohomologique de Weil ci-dessus, munie d'un isomorphisme
$F \to F(1)$, revient à se donner une théorie cohomologique de Weil classique.

\begin{remark}
\label{weil-remark-betti-numbers-in-some-sense}
Soit $H^*$ une théorie cohomologique de Weil
(définition \ref{weil-definition-weil-cohomology-theory}).
Soit $X$ un schéma projectif lisse géométriquement irréductible
de dimension $d$ sur $k'$, où $k'/k$ est une extension finie séparable de corps.
Supposons
$$
H^0(\Spec(k')) = F_1 \times \ldots \times F_r
$$,
pour certains corps $F_i$. On peut alors écrire
$$
H^*(X) = \prod\nolimits_{i = 1, \ldots, r}
H^*(X) \otimes_{H^0(\Spec(k'))} F_i
$$
La dernière hypothèse de la définition \ref{weil-definition-weil-cohomology-theory}
signifie que $H^0(X)$ est libre de rang $1$ sur $\prod F_i$.
Autrement dit, chacun des facteurs
$H^0(X) \otimes_{H^0(\Spec(k'))} F_i$ est de dimension $1$ sur $F_i$.
La dualité de Poincar\'e donne la même assertion pour
la cohomologie en degré $2d$.
En revanche, on ignore s'il en est de même aux autres degrés.
Plus précisément, pour $0 < n < \dim(X)$, on ne sait pas si les entiers
$$
\dim_{F_i} H^n(X) \otimes_{H^0(\Spec(k'))} F_i
$$
sont indépendants de $i$ ! Cette question est étroitement liée à la question
ouverte suivante : étant donnés un corps de base algébriquement clos $\overline{k}$,
un corps de caractéristique zéro $F$, une théorie cohomologique de Weil classique
$H^*$ sur $\overline{k}$ de corps de coefficients $F$ et une variété projective lisse
$X$ sur $\overline{k}$, les nombres de Betti de $X$,
$$
\beta_i = \dim_F H^i(X)
$$,
sont-ils indépendants de $F$ et de la théorie cohomologique de Weil $H^*$ ?
\end{remark}

\begin{proposition}
\label{weil-proposition-weil-cohomology-theory-again}
Soit $k$ un corps et soit $F$ un corps de caractéristique $0$.
Une théorie cohomologique de Weil revient à se donner un foncteur
monoïdal symétrique $\mathbf{Q}$-linéaire
$$
G : M_k \longrightarrow \text{espaces vectoriels gradués sur }F\text{}
$$
tel que :
\begin{enumerate}
\item $G(\mathbf{1}(1))$ ne soit non nul qu'en degré $-2$ ;
\item pour tout schéma projectif lisse $X$ sur $k$, en posant
$k' = \Gamma(X, \mathcal{O}_X)$, l'homomorphisme
$G(h(\Spec(k'))) \to G(h(X))$ d'espaces vectoriels gradués sur $F$
soit un isomorphisme en degré $0$.
\end{enumerate}
\end{proposition}

\begin{proof}
C'est une conséquence immédiate de la proposition \ref{weil-proposition-weil-cohomology-theory}
et de la définition \ref{weil-definition-weil-cohomology-theory}. On pourrait bien entendu
remplacer (2) par la condition que $G(h(X)) \to \bigoplus G(h(x_i))$
soit injective en degré $0$ pour un choix de points fermés
$x_1, \ldots, x_r \in X$ dont les corps résiduels sont séparables sur $k$.
\end{proof}







\section{Classes de Chern}
\label{weil-section-chern}

\noindent
Nous expliquons dans cette section comment une première classe de Chern et
une formule du fibré projectif permettent d'obtenir toutes les classes de Chern.
On pourra consulter \cite{Grothendieck-chern}, bien que nos axiomes
soient légèrement différents.

\medskip\noindent
Soit $\mathcal{C}$ une catégorie de schémas possédant les propriétés suivantes :
\begin{enumerate}
\item Tout $X \in \Ob(\mathcal{C})$ est quasi-compact et quasi-séparé.
\item Si $X \in \Ob(\mathcal{C})$ et si $U \subset X$ est ouvert et fermé,
alors $U \to X$ est un morphisme de $\mathcal{C}$. Si $X' \to X$ est un morphisme
de $\mathcal{C}$ se factorisant par $U$, alors $X' \to U$ est un morphisme
de $\mathcal{C}$.
\item Si $X \in \Ob(\mathcal{C})$ et si $\mathcal{E}$ est un
$\mathcal{O}_X$-module localement libre de type fini, alors :
\begin{enumerate}
\item $p : \mathbf{P}(\mathcal{E}) \to X$ est un morphisme de $\mathcal{C}$ ;
\item pour un morphisme $f : X' \to X$ de $\mathcal{C}$,
le morphisme induit $\mathbf{P}(f^*\mathcal{E}) \to \mathbf{P}(\mathcal{E})$
est un morphisme de $\mathcal{C}$ ;
\item si $\mathcal{E} \to \mathcal{F}$ est une surjection sur un autre
$\mathcal{O}_X$-module localement libre de type fini, l'immersion fermée
$\mathbf{P}(\mathcal{F}) \to \mathbf{P}(\mathcal{E})$
est un morphisme de $\mathcal{C}$.
\end{enumerate}
\end{enumerate}
Supposons ensuite donné un foncteur contravariant $A$ de
la catégorie $\mathcal{C}$ dans la catégorie des algèbres graduées.
Ici, une algèbre graduée $A$ est une
$\mathbf{Z}$-algèbre unitaire associative, non nécessairement commutative, $A$, munie d'une graduation
$A = \bigoplus_{i \geq 0} A^i$. Pour un morphisme
$f : X' \to X$ de $\mathcal{C}$, nous notons $f^* : A(X) \to A(X')$
l'homomorphisme d'algèbres induit.
Nous noterons le produit de $a, b \in A(X)$ par $a \cup b$.
Enfin, nous supposons
donnée, pour tout objet $X$ de $\mathcal{C}$, une application additive
$$
c_1^A : \Pic(X) \longrightarrow A^1(X)
$$
Nous supposons vérifiés les axiomes suivants :
\begin{enumerate}
\item Pour $X \in \Ob(\mathcal{C})$ et $\mathcal{L} \in \Pic(X)$,
l'élément $c_1^A(\mathcal{L})$ appartient au centre de l'algèbre $A(X)$.
\item Si $X \in \Ob(\mathcal{C})$ et si $X = U \amalg V$, avec $U$ et $V$
ouverts et fermés, alors $A(X) = A(U) \times A(V)$ par les applications induites
$A(X) \to A(U)$ et $A(X) \to A(V)$.
\item Si $f : X' \to X$ est un morphisme de $\mathcal{C}$ et si $\mathcal{L}$
est un $\mathcal{O}_X$-module inversible, alors $f^*c_1^A(\mathcal{L}) =
c_1^A(f^*\mathcal{L})$.
\item Pour $X \in \Ob(\mathcal{C})$ et un $\mathcal{O}_X$-module localement libre
$\mathcal{E}$ de rang constant $r$, considérons le morphisme
$p : P = \mathbf{P}(\mathcal{E}) \to X$ de $\mathcal{C}$.
Alors l'application
$$
\bigoplus\nolimits_{i = 0, \ldots, r - 1} A(X)
\longrightarrow A(P),\quad
(a_0, \ldots, a_{r - 1}) \longmapsto
\sum c_1^A(\mathcal{O}_P(1))^i \cup p^*(a_i)
$$
est bijective.
\item Soit $X \in \Ob(\mathcal{C})$ et soit $\mathcal{E} \to \mathcal{F}$
une surjection de $\mathcal{O}_X$-modules localement libres de type fini,
de rangs $r + 1$ et $r$. Notons
$i : P' = \mathbf{P}(\mathcal{F}) \to \mathbf{P}(\mathcal{E}) = P$
le morphisme d'inclusion correspondant. C'est un morphisme de $\mathcal{C}$
qui présente $P'$ comme un diviseur de Cartier effectif sur $P$. Alors, pour
$a \in A(P)$ tel que $i^*a = 0$, on a
$a \cup c_1^A(\mathcal{O}_P(P')) = 0$.
\end{enumerate}
Pour énoncer notre résultat, rappelons que $\textit{Vect}(X)$ désigne
la catégorie (exacte) des $\mathcal{O}_X$-modules localement libres de type fini.
Dans Catégories dérivées de schémas, section \ref{perfect-section-K-groups},
nous avons défini le $K$-groupe de degré zéro
$K_0(\textit{Vect}(X))$ de cette catégorie.
Nous avons également vu que $K_0(\textit{Vect}(X))$ est un anneau ; voir
Catégories dérivées de schémas, remarque \ref{perfect-remark-K-ring}.

\begin{proposition}
\label{weil-proposition-chern-class}
Dans la situation ci-dessus, il existe une unique règle associant
à tout $X \in \Ob(\mathcal{C})$ une « classe de Chern totale »
$$
c^A : K_0(\textit{Vect}(X)) \longrightarrow  \prod\nolimits_{i \geq 0} A^i(X)
$$
possédant les propriétés suivantes :
\begin{enumerate}
\item Pour $X \in \Ob(\mathcal{C})$, on a
$c^A(\alpha + \beta) = c^A(\alpha) c^A(\beta)$
et $c^A(0) = 1$.
\item Si $f : X' \to X$ est un morphisme de $\mathcal{C}$, alors
$f^* \circ c^A =  c^A \circ f^*$.
\item Pour $X \in \Ob(\mathcal{C})$ et $\mathcal{L} \in \Pic(X)$,
on a $c^A([\mathcal{L}]) = 1 + c_1^A(\mathcal{L})$.
\end{enumerate}
\end{proposition}

\begin{proof}
Soit $X \in \Ob(\mathcal{C})$ et soit $\mathcal{E}$ un
$\mathcal{O}_X$-module localement libre de type fini. Montrons d'abord comment définir
un élément $c^A(\mathcal{E}) \in A(X)$.

\medskip\noindent
Soit d'abord $X = \bigcup X_r$ la décomposition en
sous-schémas ouverts et fermés tels que $\mathcal{E}|_{X_r}$ soit
de rang constant $r$. Puisque $X$ est quasi-compact, cette décomposition
est finie. Ainsi, $A(X) = \prod A(X_r)$. Il suffit donc de définir
$c^A(\mathcal{E})$ lorsque $\mathcal{E}$ est de rang constant $r$. Dans
ce cas, soit $p : P \to X$ le fibré projectif de $\mathcal{E}$.
On définit de manière unique des éléments $c_i^A(\mathcal{E}) \in A^i(X)$
pour $i \geq 0$, tels que $c_0^A(\mathcal{E}) = 1$ et que l'équation
\begin{equation}
\label{weil-equation-chern-classes}
\sum\nolimits_{i = 0}^r
(-1)^i c_1(\mathcal{O}_P(1))^i \cup p^*c^A_{r - i}(\mathcal{E})
= 0
\end{equation}
soit vérifiée. Comme d'habitude, on pose
$c^A(\mathcal{E}) = c_0^A(\mathcal{E}) + c_1^A(\mathcal{E}) + \ldots
+ c_r^A(\mathcal{E})$ dans $A(X)$.

\medskip\noindent
Si $\mathcal{E}$ est inversible, alors
$c^A(\mathcal{E}) = 1 + c_1^A(\mathcal{L})$.
Cela résulte immédiatement de la construction ci-dessus.

\medskip\noindent
Les éléments $c_i^A(\mathcal{E})$ appartiennent au centre de $A(X)$.
Pour le montrer, on peut supposer $\mathcal{E}$ de rang constant $r$.
Soit $p : P \to X$ le fibré projectif correspondant.
Si $a \in A(X)$, alors $p^*a \cup (-1)^r c_1(\mathcal{O}_P(1))^r =
(-1)^r c_1(\mathcal{O}_P(1))^r \cup p^*a$ ; la même propriété vaut donc
pour tous les autres termes de l'expression définissant $c_i^A(\mathcal{E})$,
ce qui permet de conclure.

\medskip\noindent
Si $f : X' \to X$ est un morphisme de $\mathcal{C}$, alors
$f^*c_i^A(\mathcal{E}) = c_i^A(f^*\mathcal{E})$.
Pour le montrer, on peut supposer $\mathcal{E}$ de rang constant $r$.
Soient $p : P \to X$ et $p' : P' \to X'$ les fibrés
projectifs associés à $\mathcal{E}$ et à $f^*\mathcal{E}$.
Le morphisme induit $g : P' \to P$ est un morphisme de $\mathcal{C}$.
L'image réciproque par $g$ de l'égalité définissant $c_i^A(\mathcal{E})$
est l'équation correspondante pour $f^*\mathcal{E}$, d'où la conclusion.

\medskip\noindent
Soit $X \in \Ob(\mathcal{C})$. Considérons une suite exacte courte
$$
0 \to \mathcal{L} \to \mathcal{E} \to \mathcal{F} \to 0
$$
de $\mathcal{O}_X$-modules localement libres de type fini, avec $\mathcal{L}$ inversible.
Alors
$$
c^A(\mathcal{E}) = c^A(\mathcal{L}) c^A(\mathcal{F})
$$
En effet, par la construction de $c^A_i$, on peut supposer $\mathcal{E}$
de rang constant $r + 1$ et $\mathcal{F}$ de rang constant $r$.
L'inclusion
$$
i : P' = \mathbf{P}(\mathcal{F}) \longrightarrow \mathbf{P}(\mathcal{E}) = P
$$
est un morphisme de $\mathcal{C}$ et est le schéma des zéros d'une section
régulière du module inversible
$\mathcal{L}^{\otimes -1} \otimes \mathcal{O}_P(1)$.
L'élément
$$
\sum\nolimits_{i = 0}^r (-1)^i c_1^A(\mathcal{O}_P(1))^i \cup
p^*c^A_i(\mathcal{F})
$$
a une image réciproque nulle sur $P'$ par définition. On obtient donc
$$
\left(c_1^A(\mathcal{O}_P(1)) - c_1^A(\mathcal{L})\right) \cup
\left(\sum\nolimits_{i = 0}^r (-1)^i c_1^A(\mathcal{O}_P(1))^i \cup
p^*c^A_i(\mathcal{F})\right) = 0
$$
dans $A^*(P)$, par l'hypothèse (5) sur notre cohomologie $A$.
Par définition de $c_1^A(\mathcal{E})$,
cela donne l'égalité voulue.

\medskip\noindent
Soit $X \in \Ob(\mathcal{C})$. Considérons une suite exacte courte
$$
0 \to \mathcal{E} \to \mathcal{F} \to \mathcal{G} \to 0
$$
de $\mathcal{O}_X$-modules localement libres de type fini. Alors
$$
c^A(\mathcal{F}) = c^A(\mathcal{E}) c^A(\mathcal{G})
$$
En effet, par la construction de $c^A_i$, on peut supposer
$\mathcal{E}$, $\mathcal{F}$ et $\mathcal{G}$
de rangs constants $r$, $s$ et $t$. On raisonne par récurrence sur $r$.
Le cas $r = 1$ a été traité ci-dessus. Si $r > 1$, il suffit de vérifier
l'égalité après image réciproque par le morphisme $\mathbf{P}(\mathcal{E}^\vee) \to X$.
On peut donc supposer qu'il existe un sous-module inversible
$\mathcal{L} \subset \mathcal{E}$ tel que
$\mathcal{E}' = \mathcal{E}/\mathcal{L}$ et
$\mathcal{F}' = \mathcal{E}/\mathcal{L}$ soient tous deux localement libres de type fini
(de rangs $s - 1$ et $t - 1$). Alors
$$
c^A(\mathcal{E}) = c^A(\mathcal{L}) c^A(\mathcal{E}')
\quad\text{et}\quad
c^A(\mathcal{F}) = c^A(\mathcal{L}) c^A(\mathcal{F}')
$$
Puisque l'on a la suite exacte courte
$$
0 \to \mathcal{E}' \to \mathcal{F}' \to \mathcal{G} \to 0
$$,
l'hypothèse de récurrence donne
$$
c^A(\mathcal{F}') = c^A(\mathcal{E}') c^A(\mathcal{G})
$$
Le résultat découle donc d'un calcul formel.

\medskip\noindent
Nous pouvons maintenant, pour $X \in \Ob(\mathcal{C})$,
définir $c^A : K_0(\textit{Vect}(X)) \to A(X)$.
On envoie le générateur $[\mathcal{E}]$ sur $c^A(\mathcal{E})$
et l'on prolonge multiplicativement. Ainsi, par exemple,
$c^A(-[\mathcal{E}]) = c^A(\mathcal{E})^{-1}$ est l'inverse
formel de $a^A([\mathcal{E}])$.
La multiplicativité sur les suites exactes courtes établie ci-dessus
assure que cette définition convient.

\medskip\noindent
Unicité. Soit $X \in \Ob(\mathcal{C})$ et soit $\mathcal{E}$
un $\mathcal{O}_X$-module localement libre de type fini. Nous voulons montrer
que les conditions (1), (2) et (3) du lemme déterminent de manière unique
$c^A([\mathcal{E}])$. On peut supposer $\mathcal{E}$
de rang constant $r$, ce qui utilise déjà (2). On raisonne alors par récurrence sur $r$.
Si $r = 1$, l'unicité résulte de (3).
Si $r > 1$, on passe, à l'aide de (2), à l'image réciproque sur le fibré projectif $p : P \to X$ ;
on peut ainsi supposer que l'on dispose d'une suite exacte courte
$0 \to \mathcal{E}' \to \mathcal{E} \to \mathcal{E}'' \to 0$,
où $\mathcal{E}'$ et $\mathcal{E}''$ sont de rangs inférieurs.
Par l'hypothèse de récurrence, $c^A(\mathcal{E}')$ et $c^A(\mathcal{E}'')$
sont déterminés de manière unique. L'unicité pour $\mathcal{E}$ résulte alors
de l'axiome (1).
\end{proof}

\begin{lemma}
\label{weil-lemma-splitting-principle}
Dans la situation ci-dessus, soit $X \in \Ob(\mathcal{C})$. Soit $\mathcal{E}_i$
une famille finie de $\mathcal{O}_X$-modules localement libres de rangs $r_i$.
Il existe un morphisme $p : P \to X$ dans $\mathcal{C}$ tel que :
\begin{enumerate}
\item $p^* : A(X) \to A(P)$ soit injective ;
\item chaque $p^*\mathcal{E}_i$ admette une filtration dont les quotients successifs
$\mathcal{L}_{i, 1}, \ldots, \mathcal{L}_{i, r_i}$
soient des $\mathcal{O}_P$-modules inversibles.
\end{enumerate}
\end{lemma}

\begin{proof}
On peut supposer $r_i \geq 1$ pour tout $i$. On démontre le lemme par récurrence
sur $\sum (r_i - 1)$. Si cet entier vaut $0$, alors $\mathcal{E}_i$
est inversible pour tout $i$, et l'on conclut en prenant $\pi = \text{id}_X$.
Sinon, choisissons un $i$ tel que $r_i > 1$ et considérons
le fibré projectif $p : P \to X$ associé à $\mathcal{E}_i$.
On a une suite exacte courte
$$
0 \to \mathcal{F} \to p^*\mathcal{E}_i \to \mathcal{O}_P(1) \to 0
$$
de $\mathcal{O}_P$-modules localement libres de type fini, de rangs $r_i - 1$,
$r_i$ et $1$. Observons que $p^* : A(X) \to A(P)$ est injective
par hypothèse. En appliquant l'hypothèse de récurrence aux
$\mathcal{O}_P$-modules localement libres de type fini $\mathcal{F}$ et $p^*\mathcal{E}_{i'}$
pour $i' \not = i$, on obtient un morphisme $p' : P' \to P$ possédant
les propriétés de l'énoncé. La composée
$p \circ p' : P' \to X$ convient alors.
\end{proof}

\begin{lemma}
\label{weil-lemma-chern-classes-E-tensor-L}
Soit $X \in \Ob(\mathcal{C})$. Soit $\mathcal{E}$ un
$\mathcal{O}_X$-module localement libre de type fini et soit $\mathcal{L}$ un
$\mathcal{O}_X$-module inversible. Alors
$$
c^A_i({\mathcal E} \otimes {\mathcal L})
=
\sum\nolimits_{j = 0}^i
\binom{r - i + j}{j} c^A_{i - j}({\mathcal E}) \cup c^A_1({\mathcal L})^j
$$
\end{lemma}

\begin{proof}
Par la construction de $c^A_i$, on peut supposer $\mathcal{E}$
de rang constant $r$. Soient $p : P \to X$ et $p' : P' \to X$
les fibrés projectifs associés à $\mathcal{E}$ et à
$\mathcal{E} \otimes \mathcal{L}$.
Il existe alors un isomorphisme $g : P \to P'$ tel que
$g^*\mathcal{O}_{P'}(1) = \mathcal{O}_P(1) \otimes p^*\mathcal{L}$.
Voir Constructions, lemme \ref{constructions-lemma-twisting-and-proj}.
Ainsi,
$$
g^*c_1^A(\mathcal{O}_{P'}(1)) =
c_1^A(\mathcal{O}_P(1)) + p^*c_1^A(\mathcal{L})
$$
L'égalité voulue en découle formellement, en utilisant la définition
des classes de Chern par l'équation (\ref{weil-equation-chern-classes}).
\end{proof}

\begin{proposition}
\label{weil-proposition-chern-character}
Dans la situation ci-dessus, supposons que $A(X)$ soit une $\mathbf{Q}$-algèbre pour tout
$X \in \Ob(\mathcal{C})$. Il existe alors une unique règle associant
à tout $X \in \Ob(\mathcal{C})$ un « caractère de Chern »
$$
ch^A : K_0(\textit{Vect}(X)) \longrightarrow
\prod\nolimits_{i \geq 0} A^i(X)
$$
possédant les propriétés suivantes :
\begin{enumerate}
\item $ch^A$ est un homomorphisme d'anneaux pour tout $X \in \Ob(\mathcal{C})$.
\item Si $f : X' \to X$ est un morphisme de $\mathcal{C}$, alors
$f^* \circ ch^A =  ch^A \circ f^*$.
\item Pour $X \in \Ob(\mathcal{C})$ et $\mathcal{L} \in \Pic(X)$,
on a $ch^A([\mathcal{L}]) = \exp(c_1^A(\mathcal{L}))$.
\end{enumerate}
\end{proposition}

\begin{proof}
Soit $X \in \Ob(\mathcal{C})$ et soit $\mathcal{E}$ un
$\mathcal{O}_X$-module localement libre de type fini. Montrons d'abord comment définir
le rang $r^A(\mathcal{E}) \in A^0(X)$. Soit $X = \bigcup X_r$
la décomposition en sous-schémas ouverts et fermés tels que
$\mathcal{E}|_{X_r}$ soit de rang constant $r$. Puisque $X$ est quasi-compact, cette
décomposition est finie ; écrivons-la $X = X_0 \amalg X_1 \amalg \ldots \amalg X_n$.
Alors $A(X) = A(X_0) \times A(X_1) \times \ldots \times A(X_n)$. On peut donc
définir $r^A(\mathcal{E}) = (0, 1, \ldots, n) \in A^0(X)$.

\medskip\noindent
Soient $P_p(c_1, \ldots, c_p)$ les polynômes construits dans
Homologie de Chow, exemple \ref{chow-example-power-sum}.
On peut alors définir
$$
ch^A(\mathcal{E}) = r^A(\mathcal{E}) +
\sum\nolimits_{i \geq 1} (1/i!)
P_i(c^A_1(\mathcal{E}), \ldots, c^A_i(\mathcal{E}))
\in \prod\nolimits_{i \geq 0} A^i(X)
$$,
où les $ci^A$ sont les classes de Chern de la
proposition \ref{weil-proposition-chern-class}.
Les propriétés (2) et (3) du lemme sont immédiatement vérifiées.

\medskip\noindent
Il reste à démontrer les trois assertions suivantes :
\begin{enumerate}
\item Si $0 \to \mathcal{E}_1 \to \mathcal{E} \to \mathcal{E}_2 \to 0$
est une suite exacte courte de $\mathcal{O}_X$-modules localement libres de type fini
sur $X \in \Ob(\mathcal{C})$, alors
$ch^A(\mathcal{E}) = ch^A(\mathcal{E}_1) + ch^A(\mathcal{E}_2)$.
\item Si $\mathcal{E}_1$ et $\mathcal{E}_2 \to 0$ sont des
$\mathcal{O}_X$-modules localement libres de type fini sur $X \in \Ob(\mathcal{C})$, alors
$ch^A(\mathcal{E}_1 \otimes \mathcal{E}_2) =
ch^A(\mathcal{E}_1) ch^A(\mathcal{E}_2)$.
\end{enumerate}
La première montrera que $ch^A$ se factorise par
$K_0(\textit{Vect}(X))$, et la première et la seconde ensemble
montreront que $ch^A$ est un homomorphisme d'anneaux.

\medskip\noindent
Pour démontrer ces assertions, on peut se ramener au cas où $\mathcal{E}_1$
et $\mathcal{E}_2$ sont de rangs constants $r_1$ et $r_2$. Les
égalités dans $A^0(X)$ sont alors immédiates. Pour les égalités en degrés
supérieurs, le lemme \ref{weil-lemma-splitting-principle} permet de
supposer que $\mathcal{E}_1$ et $\mathcal{E}_2$ admettent des filtrations
dont les composantes graduées sont des modules inversibles
$\mathcal{L}_{1, j}$, pour $j = 1, \ldots, r_1$, et
$\mathcal{L}_{2, j}$, pour $j = 1, \ldots, r_2$.
La multiplicativité des classes de Chern donne
$$
c_i^A(\mathcal{E}_1) =
s_i(c_1^A(\mathcal{L}_{1, 1}), \ldots, c_1^A(\mathcal{L}_{1, r_1}))
$$,
où $s_i$ est la $i$-ième fonction symétrique élémentaire, comme dans
Homologie de Chow, exemple \ref{chow-example-power-sum}.
Il en est de même pour $c_i^A(\mathcal{E}_2)$. Dans le cas (1), on obtient
$$
c_i^A(\mathcal{E}) =
s_i(c_1^A(\mathcal{L}_{1, 1}), \ldots, c_1^A(\mathcal{L}_{1, r_1}),
c_1^A(\mathcal{L}_{2, 1}), \ldots, c_1^A(\mathcal{L}_{2, r_2}))
$$,
et, dans le cas (2),
$$
c_i^A(\mathcal{E}_1 \otimes \mathcal{E}_2) =
s_i(c_1^A(\mathcal{L}_{1, 1}) + c_1^A(\mathcal{L}_{2, 1}),
\ldots, c_1^A(\mathcal{L}_{1, r_1}) + c_1^A(\mathcal{L}_{2, r_2}))
$$
Par définition des polynômes $P_i$, cela signifie que
$$
P_i(c^A_1(\mathcal{E}_1), \ldots, c^A_i(\mathcal{E}_1)) =
\sum\nolimits_{j = 1, \ldots, r_1} c_1^A(\mathcal{L}_{1, j})^i
$$,
et de même pour $\mathcal{E}_2$. Dans le cas (1), on a aussi
$$
P_i(c^A_1(\mathcal{E}), \ldots, c^A_i(\mathcal{E})) =
\sum\nolimits_{j = 1, \ldots, r_1} c_1^A(\mathcal{L}_{1, j})^i +
\sum\nolimits_{j = 1, \ldots, r_2} c_1^A(\mathcal{L}_{2, j})^i
$$
Dans le cas (2), on obtient de même
$$
P_i(c^A_1(\mathcal{E}_1 \otimes \mathcal{E}_2), \ldots,
c^A_i(\mathcal{E}_1 \otimes \mathcal{E}_2)) =
\sum\nolimits_{j = 1, \ldots, r_1}
\sum\nolimits_{j' = 1, \ldots, r_2}
(c_1^A(\mathcal{L}_{1, j}) + c_1^A(\mathcal{L}_{2, j'}))^i
$$
Les égalités voulues découlent donc d'identités élémentaires
entre polynômes symétriques.

\medskip\noindent
Nous omettons la démonstration de l'unicité.
\end{proof}

\begin{lemma}
\label{weil-lemma-adams-and-chern}
Dans la situation ci-dessus, soit $X \in \Ob(\mathcal{C})$.
Si $\psi^2$ est défini comme dans
Homologie de Chow, lemme \ref{chow-lemma-second-adams-operator},
et si $c^A$ et $ch^A$ sont définis comme dans les
propositions \ref{weil-proposition-chern-class} et
\ref{weil-proposition-chern-character},
alors on a $c^A_i(\psi^2(\alpha)) = 2^i c^A_i(\alpha)$ et
$ch^A_i(\psi^2(\alpha)) = 2^i ch^A_i(\alpha)$
pour tout $\alpha \in K_0(\textit{Vect}(X))$.
\end{lemma}

\begin{proof}
Observons que l'application $\prod_{i \geq 0} A^i(X) \to \prod_{i \geq 0} A^i(X)$,
qui multiplie par $2^i$ sur $A^i(X)$, est un homomorphisme d'anneaux. Puisque $\psi^2$
est également un homomorphisme d'anneaux, il suffit donc de démontrer les formules
sur des générateurs additifs de $K_0(\textit{Vect}(X))$. On peut ainsi supposer $\alpha = [\mathcal{E}]$
pour un $\mathcal{O}_X$-module localement libre de type fini $\mathcal{E}$.
La construction des classes de Chern de $\mathcal{E}$ ramène immédiatement
au cas où $\mathcal{E}$ est de rang constant $r$.
On peut alors choisir un morphisme projectif lisse $p : P \to X$
tel que l'application d'image réciproque $A^*(X) \to A^*(P)$ soit injective
et que $p^*\mathcal{E}$ admette une filtration finie dont
les composantes graduées soient des $\mathcal{O}_P$-modules inversibles $\mathcal{L}_j$ ; voir
le lemme \ref{weil-lemma-splitting-principle}. Alors
$[p^*\mathcal{E}] = \sum [\mathcal{L}_j]$, d'où
$\psi^2([p^*\mathcal{E}]) = \sum [\mathcal{L}_j^{\otimes 2}]$
par définition de $\psi^2$. En posant $x_j  = c^A_1(\mathcal{L}_j)$,
on a
$$
c^A(\alpha) = \prod (1 + x_j)
\quad\text{et}\quad
c^A(\psi^2(\alpha)) = \prod (1 + 2 x_j)
$$
dans $\prod A^i(P)$, et l'on a
$$
ch^A(\alpha) = \sum \exp(x_j)
\quad\text{et}\quad
ch^A(\psi^2(\alpha)) = \sum \exp(2 x_j)
$$
dans $\prod A^i(P)$. Le résultat voulu découle de ces formules.
\end{proof}







\section{Puissances extérieures et groupes K}
\label{weil-section-lambda-operations}

\noindent
Nous nous limitons aux constructions nécessaires pour définir les opérateurs lambda.
Soit $X$ un schéma. Rappelons que $\textit{Vect}(X)$ désigne
la catégorie des $\mathcal{O}_X$-modules localement libres de type fini.
Rappelons également que nous avons construit un groupe $K$ de degré zéro
$K_0(\textit{Vect}(X))$ associé à cette catégorie dans
Catégories dérivées de schémas, section \ref{perfect-section-K-groups}.
Enfin, $K_0(\textit{Vect}(X))$ est un anneau ; voir
Catégories dérivées de schémas, remarque \ref{perfect-remark-K-ring}.

\begin{lemma}
\label{weil-lemma-lambda-operations}
Soit $X$ un schéma. Il existe des applications
$$
\lambda^r : K_0(\textit{Vect}(X)) \longrightarrow K_0(\textit{Vect}(X))
$$
qui envoient $[\mathcal{E}]$ sur $[\wedge^r(\mathcal{E})]$
lorsque $\mathcal{E}$ est un $\mathcal{O}_X$-module localement libre de type fini,
et qui sont compatibles avec les images réciproques.
\end{lemma}

\begin{proof}
Considérons l'anneau $R = K_0(\textit{Vect}(X))[[t]]$, où $t$ est une
indéterminée. Pour un $\mathcal{O}_X$-module localement libre de type fini
$\mathcal{E}$, posons
$$
c(\mathcal{E}) = \sum\nolimits_{i = 0}^\infty [\wedge^i(\mathcal{E})] t^i
$$
dans $R$. Montrons que, pour une suite exacte courte
$$
0 \to \mathcal{E}' \to \mathcal{E} \to \mathcal{E}'' \to 0
$$
de $\mathcal{O}_X$-modules localement libres de type fini,
on a $c(\mathcal{E}) = c(\mathcal{E}') c(\mathcal{E}'')$.
Cette assertion entraîne que $c$ se prolonge en une application
$$
c :  K_0(\textit{Vect}(X)) \longrightarrow R
$$
qui transforme l'addition dans $K_0(\textit{Vect}(X))$ en la multiplication dans $R$.
En écrivant $c(\alpha) = \sum \lambda^i(\alpha) t^i$, on obtient les opérateurs
$\lambda^i$ cherchés.

\medskip\noindent
Pour démontrer l'assertion, considérons la suite exacte courte comme une
filtration de $\mathcal{E}$ à $2$ crans. On obtient une filtration
induite de $\wedge^r(\mathcal{E})$ à $r + 1$ crans, dont les
sous-quotients sont
$$
\wedge^r(\mathcal{E}'),
\wedge^{r - 1}(\mathcal{E}') \otimes \mathcal{E}'',
\wedge^{r - 2}(\mathcal{E}') \otimes \wedge^2(\mathcal{E}''), \ldots
\wedge^r(\mathcal{E}'')
$$
Ainsi, $[\wedge^r(\mathcal{E})]$ est égal à
$$
\sum\nolimits_{i = 0}^r
[\wedge^{r - i}(\mathcal{E}')] [\wedge^i(\mathcal{E}'')]
$$,
et le résultat s'en déduit aisément par un calcul algébrique élémentaire.
\end{proof}










\section{Théories cohomologiques de Weil, III}
\label{weil-section-c1}

\noindent
Soit $k$ un corps et soit $F$ un corps de caractéristique zéro.
Supposons que l'on dispose des données suivantes :
\begin{enumerate}
\item[(D0)] Un espace vectoriel de dimension $1$ sur $F$, noté $F(1)$.
\item[(D1)] Un foncteur contravariant $H^*(-)$ de la catégorie des schémas projectifs
lisses sur $k$ dans la catégorie des
$F$-algèbres graduées commutatives.
\item[(D2')] Pour tout schéma projectif lisse $X$ sur $k$, un homomorphisme
$c_1^H : \Pic(X) \to H^2(X)(1)$ de groupes abéliens.
\end{enumerate}
Nous employons la terminologie, les notations et les conventions relatives
à (D0) et (D1) exposées à la section \ref{weil-section-axioms}.
Pour un schéma projectif lisse $X$ sur $k$ et un
$\mathcal{O}_X$-module inversible $\mathcal{L}$, la classe de cohomologie
$c_1^H(\mathcal{L}) \in H^2(X)(1)$ de (D2')
est parfois appelée la {\it première classe de Chern de $\mathcal{L}$
en cohomologie}.

\medskip\noindent
Voici la liste des axiomes.
\begin{enumerate}
\item[(A1)] $H^*$ est compatible avec les sommes finies.
\item[(A2)] $c_1^H$ est compatible avec les images réciproques.
\item[(A3)] Soit $X$ un schéma projectif lisse sur $k$.
Soit $\mathcal{E}$ un $\mathcal{O}_X$-module localement libre de rang $r \geq 1$.
Considérons le morphisme $p : P = \mathbf{P}(\mathcal{E}) \to X$.
Alors l'application
$$
\bigoplus\nolimits_{i = 0, \ldots, r - 1} H^*(X)(-i)
\longrightarrow H^*(P),\quad
(a_0, \ldots, a_{r - 1}) \longmapsto
\sum c_1^H(\mathcal{O}_P(1))^i \cup p^*(a_i)
$$
est un isomorphisme d'espaces vectoriels sur $F$.
\item[(A4)] Soit $i : Y \to X$ l'inclusion d'un diviseur de Cartier
effectif sur $k$, avec $X$ et $Y$ lisses et projectifs
sur $k$. Pour $a \in H^*(X)$ tel que
$i^*a = 0$, on a $a \cup c_1^H(\mathcal{O}_X(Y)) = 0$.
\item[(A5)] $H^*$ est compatible avec les produits finis.
\item[(A6)] Soit $X$ un schéma projectif lisse non vide sur $k$,
équidimensionnel de dimension $d$. Il existe une application
$F$-linéaire $\lambda : H^{2d}(X)(d) \to F$ telle que
$(\text{id} \otimes \lambda) \gamma([\Delta]) =  1$ dans $H^*(X)$.
\item[(A7)] Si $b : X' \to X$ est l'éclatement d'un centre
lisse dans un schéma projectif lisse $X$ sur $k$\footnote{Alors $X'$
est également lisse et projectif sur $k$ ; voir
Compléments sur les morphismes, lemme \ref{more-morphisms-lemma-blowup}.}, alors
$b^* : H^*(X) \to H^*(X')$ est injective.
\item[(A8)] Si $X$ est un schéma projectif lisse sur $k$ et si
$k' = \Gamma(X, \mathcal{O}_X)$, alors l'application $H^0(\Spec(k')) \to H^0(X)$
est un isomorphisme.
\item[(A9)] Soit $X$ un schéma projectif lisse non vide sur $k$,
équidimensionnel de dimension $d$. Soit $i : Y \to X$ un diviseur de Cartier
effectif non vide, lisse sur $k$. Pour $a \in H^{2d - 2}(X)(d - 1)$,
on a $\lambda_Y(i^*(a)) = \lambda_X(a \cup c_1^H(\mathcal{O}_X(Y))$, où
$\lambda_Y$ et $\lambda_X$ sont celles de l'axiome (A6) pour $X$ et $Y$.
\end{enumerate}
Précisons le sens de chacun de ces axiomes.
Les axiomes (A3), (A4) et (A7) sont clairs tels qu'ils sont énoncés.

\medskip\noindent
À propos de (A1). Cela signifie que $H^*(\emptyset) = 0$ et que
$(i^*, j^*) : H^*(X \amalg Y) \to H^*(X) \times H^*(Y)$
est un isomorphisme, où $i$ et $j$ sont les coprojections.

\medskip\noindent
À propos de (A2). Cela signifie que, pour un morphisme $f : X \to Y$ de schémas projectifs
lisses sur $k$ et un $\mathcal{O}_Y$-module inversible $\mathcal{N}$,
on a $f^*c_1^H(\mathcal{L}) = c_1^H(f^*\mathcal{L})$.

\medskip\noindent
À propos de (A5). Cela signifie que $H^*(\Spec(k)) = F$ et que, pour $X$ et $Y$ projectifs
lisses sur $k$, l'application $H^*(X) \otimes_F H^*(Y) \to H^*(X \times Y)$,
$a \otimes b \mapsto p^*(a) \cup q^*(b)$, est un isomorphisme,
où $p$ et $q$ sont les projections.

\medskip\noindent
À propos de (A6). Soit $X$ un schéma projectif lisse non vide sur $k$,
équidimensionnel de dimension $d$. D'après le lemme \ref{weil-lemma-cycle-classes},
sous les axiomes (A1) -- (A4), on peut considérer la classe de la diagonale
$$
\gamma([\Delta]) \in
H^{2d}(X \times X)(d) = \bigoplus\nolimits_i H^i(X) \otimes_F H^{2d - i}(X)(d)
$$,
où la décomposition tensorielle provient de l'axiome (A5).
Étant donnée une application $F$-linéaire $\lambda : H^{2d}(X)(d) \to F$, on peut aussi considérer
$\lambda$ comme une application $F$-linéaire $\lambda : H^*(X)(d) \to F$, en la précomposant
avec la projection sur $H^{2d}(X)(d)$. Avec cette convention, l'axiome (A6)
a bien un sens.

\medskip\noindent
À propos de (A8). Soit $X$ un schéma projectif lisse sur $k$.
Alors $k' = \Gamma(X, \mathcal{O}_X)$ est une
$k$-algèbre finie séparable (Variétés, lemme
\ref{varieties-lemma-proper-geometrically-reduced-global-sections}) ;
ainsi, $\Spec(k')$ est lisse et projectif sur $k$.
On peut donc appliquer $H^*$ à $\Spec(k')$, et l'axiome (A8) a bien un sens.

\medskip\noindent
À propos de (A9). Nous verrons dans la remarque \ref{weil-remark-trace} que, sous
les axiomes (A1) -- (A7), l'application $\lambda$ de l'axiome (A6) est unique.

\begin{lemma}
\label{weil-lemma-chern-classes}
Supposons que les données (D0), (D1) et (D2') satisfassent aux axiomes (A1), (A2), (A3) et (A4).
Il existe une unique règle associant à tout schéma projectif lisse $X$ sur $k$
un homomorphisme d'anneaux
$$
ch^H :
K_0(\textit{Vect}(X))
\longrightarrow
\prod\nolimits_{i \geq 0} H^{2i}(X)(i)
$$,
compatible avec les images réciproques et tel que
$ch^H(\mathcal{L}) = \exp(c_1^H(\mathcal{L}))$
pour tout $\mathcal{O}_X$-module inversible $\mathcal{L}$.
\end{lemma}

\begin{proof}
Cela résulte immédiatement de la proposition \ref{weil-proposition-chern-character},
appliquée à la catégorie des schémas projectifs lisses sur $k$,
au foncteur $A : X \mapsto \bigoplus_{i \geq 0} H^{2i}(X)(i)$
et à l'application $c_1^H$.
\end{proof}

\begin{lemma}
\label{weil-lemma-cycle-classes}
Supposons que les données (D0), (D1) et (D2') satisfassent aux axiomes (A1), (A2), (A3) et (A4).
Il existe une unique règle associant à tout schéma projectif lisse $X$ sur $k$
un homomorphisme d'anneaux gradués
$$
\gamma : \CH^*(X) \longrightarrow \bigoplus\nolimits_{i \geq 0} H^{2i}(X)(i)
$$,
compatible avec les images réciproques et tel que $ch^H(\alpha) = \gamma(ch(\alpha))$
pour $\alpha$ dans $K_0(\textit{Vect}(X))$.
\end{lemma}

\begin{proof}
Rappelons que l'on dispose d'un isomorphisme
$$
K_0(\textit{Vect}(X)) \otimes \mathbf{Q}
\longrightarrow \CH^*(X) \otimes \mathbf{Q},\quad
\alpha \longmapsto ch(\alpha) \cap [X]
$$ ;
voir Homologie de Chow, lemme \ref{chow-lemma-K-tensor-Q}. C'est un isomorphisme
d'anneaux d'après Homologie de Chow, remarque \ref{chow-remark-chern-character-K}.
On définit $\gamma$ par la formule $\gamma(\alpha) = ch^H(\alpha')$,
où $ch^H$ est celui du lemme \ref{weil-lemma-chern-classes} et où
$\alpha' \in K_0(\textit{Vect}(X))$ est tel que
$ch(\alpha') \cap [X] = \alpha$ dans $\CH^*(X) \otimes \mathbf{Q}$.

\medskip\noindent
La construction $\alpha \mapsto \gamma(\alpha)$ est compatible
avec les images réciproques, car $ch^H$ et la formation des classes de Chern
le sont tous deux ; voir
le lemme \ref{weil-lemma-chern-classes} et
Homologie de Chow, remarque \ref{chow-remark-gysin-chern-classes}.

\medskip\noindent
Il reste à voir que $\gamma$ respecte la graduation.
Soit $\psi^2 : K_0(\textit{Vect}(X)) \to K_0(\textit{Vect}(X))$
le deuxième opérateur d'Adams ; voir Homologie de Chow,
lemme \ref{chow-lemma-second-adams-operator}.
Si $\alpha \in \CH^i(X)$ et si
$\alpha' \in K_0(\textit{Vect}(X)) \otimes \mathbf{Q}$
est l'unique élément tel que $ch(\alpha') \cap [X] = \alpha$,
nous avons vu dans
Homologie de Chow, section \ref{chow-section-intersection-regular},
que $\psi^2(\alpha') = 2^i \alpha'$.
On en déduit $ch^H(\alpha') \in H^{2i}(X)(i)$
par le lemme \ref{weil-lemma-adams-and-chern}, comme voulu.
\end{proof}

\begin{lemma}
\label{weil-lemma-divide-pullback-good-blowing-up}
Soit $b : X' \to X$ l'éclatement d'un schéma projectif lisse
sur $k$ le long d'un sous-schéma fermé lisse $Z \subset X$.
Considérons le diagramme
$$
\xymatrix{
E \ar[r]_j \ar[d]_\pi & X' \ar[d]^b \\
Z \ar[r]^i & X
}
$$
Supposons qu'il existe un élément de $K_0(X)$ dont la restriction à
$Z$ soit égale à la classe de $\mathcal{C}_{Z/X}$ dans $K_0(Z)$.
Supposons que toute composante irréductible de $Z$ soit de codimension $r$ dans $X$.
Il existe alors un cycle $\theta \in \CH^{r - 1}(X')$
tel que $b^![Z] = [E] \cdot \theta$ dans $\CH^r(X')$ et
$\pi_*j^!(\theta) = [Z]$ dans $\CH^r(Z)$.
\end{lemma}

\begin{proof}
Le schéma $X$ est lisse et projectif sur $k$, de sorte que
$K_0(X) = K_0(\textit{Vect}(X))$. Voir
Catégories dérivées de schémas, lemmes
\ref{perfect-lemma-resolution-property-ample} et
\ref{perfect-lemma-K-is-old-K}.
Soit $\alpha \in K_0(\text{Vect}(X))$ un élément
dont la restriction à $Z$ est $\mathcal{C}_{Z/X}$.
D'après Homologie de Chow, lemme \ref{chow-lemma-minus-adams-operator},
il existe un élément $\alpha^\vee$ dont la restriction est
$\mathcal{C}_{Z/X}^\vee$. Par la formule de l'éclatement
(Homologie de Chow, lemme \ref{chow-lemma-blow-up-formula}),
on a
$$
b^![Z] = b^!i_*[Z] = j_* res(b^!)([Z]) =
j_*(c_{r - 1}(\mathcal{F}^\vee) \cap \pi^*[Z]) =
j_*(c_{r - 1}(\mathcal{F}^\vee) \cap [E])
$$,
où $\mathcal{F}$ est le noyau de la surjection
$\pi^*\mathcal{C}_{Z/X} \to \mathcal{C}_{E/X'}$.
Observons que $b^*\alpha^\vee - [\mathcal{O}_{X'}(E)]$
est un élément de $K_0(\text{Vect}(X'))$
dont la restriction est $[\pi^*\mathcal{C}_{Z/X}^\vee] - [\mathcal{C}_{E/X'}^\vee] =
[\mathcal{F}^\vee]$ sur $E$. Puisque le cap-produit par les classes de Chern
commute à $j_*$, l'expression ci-dessus est égale à
$$
c_{r - 1}(b^*\alpha^\vee - [\mathcal{O}_{X'}(E)]) \cap [E]
$$
dans le groupe de Chow de $X'$. Ainsi, en posant
$$
\theta = c_{r - 1}(b^*\alpha^\vee - [\mathcal{O}_{X'}(E)]) \cap [X']
$$,
on obtient la première relation $\theta \cdot [E] = b^![Z]$,
par exemple d'après Homologie de Chow, lemme \ref{chow-lemma-identify-chow-for-smooth}.
Pour la seconde relation, observons que
$$
j^!\theta = j^!(c_{r - 1}(b^*\alpha^\vee - [\mathcal{O}_{X'}(E)]) \cap [X'])
= c_{r - 1}(\mathcal{F}^\vee) \cap j^![X'] =
c_{r - 1}(\mathcal{F}^\vee) \cap [E]
$$
dans les groupes de Chow de $E$. Pour montrer que son image par $\pi_*$ est égale à $[Z]$,
il suffit de montrer que le degré du cycle de codimension $r - 1$
$(-1)^{r - 1}c_{r - 1}(\mathcal{F}) \cap [E]$ sur les fibres de $\pi$ vaut $1$.
Nous omettons ce calcul.
\end{proof}

\begin{lemma}
\label{weil-lemma-A5-A6-imply}
Supposons que les données (D0), (D1) et (D2') satisfassent aux axiomes (A1) -- (A4)
et (A7). Soit $X$ un schéma projectif lisse sur $k$ et soit $Z \subset X$
un sous-schéma fermé lisse, tel que toute composante irréductible de $Z$
soit de codimension $r$ dans $X$. Supposons que la classe de
$\mathcal{C}_{Z/X}$ dans $K_0(Z)$ soit la restriction d'un élément de $K_0(X)$.
Si $a \in H^*(X)$ et si $a|_Z = 0$ dans $H^*(Z)$, alors
$\gamma([Z]) \cup a = 0$.
\end{lemma}

\begin{proof}
Soit $b : X' \to X$ l'éclatement. D'après (A7), il suffit de montrer
que
$$
b^*(\gamma([Z]) \cup a) = b^*\gamma([Z]) \cup b^*a = 0
$$
Par le lemme \ref{weil-lemma-divide-pullback-good-blowing-up}, on a
$$
b^*\gamma([Z]) = \gamma(b^*[Z]) =
\gamma([E] \cdot \theta) =
\gamma([E]) \cup \gamma(\theta)
$$
Puisque $b^*a$ se restreint à zéro sur $E$ et que
$\gamma([E]) = c^H_1(\mathcal{O}_{X'}(E))$, l'axiome (A4) donne le résultat voulu.
\end{proof}

\begin{lemma}
\label{weil-lemma-poincare-duality}
Supposons que les données (D0), (D1) et (D2') satisfassent aux axiomes (A1) -- (A7).
Alors l'axiome (A) de la section \ref{weil-section-axioms} est vérifié avec
$\int_X = \lambda$ comme dans l'axiome (A6).
\end{lemma}

\begin{proof}
Soit $X$ un schéma projectif lisse non vide sur $k$,
équidimensionnel de dimension $d$. Nous allons montrer que l'espace vectoriel gradué sur $F$
$H^*(X)(d)[2d]$ est un dual à gauche de $H^*(X)$. Cela donnera le résultat voulu par
Homologie, lemme \ref{homology-lemma-left-dual-graded-vector-spaces}.
Nous utiliserons l'axiome (A5), qui donne en particulier
$$
H^*(X \times X)(d) =
\bigoplus H^i(X) \otimes H^j(X)(d) =
\bigoplus H^i(X)(d) \otimes H^j(X)
$$
Définissons une application
$$
\eta : F \longrightarrow H^*(X \times X)(d)
$$
par multiplication par $\gamma([\Delta]) \in H^{2d}(X \times X)(d)$.
D'autre part, définissons une application
$$
\epsilon :
H^*(X \times X)(d) \longrightarrow H^*(X)(d) \xrightarrow{\lambda} F
$$
en prenant d'abord l'image réciproque $\Delta^*$ par le morphisme diagonal
$\Delta : X \to X \times X$, puis l'application $F$-linéaire
$\lambda : H^{2d}(X)(d) \to F$ de l'axiome (A6), précomposée
avec la projection $H^*(X)(d) \to H^{2d}(X)(d)$.
Pour montrer que $H^*(X)(d)$ est un dual à gauche de $H^*(X)$,
il faut montrer que la composée des applications
$$
\eta \otimes 1 :
H^*(X) \longrightarrow H^*(X \times X \times X)(d)
$$
et
$$
1 \otimes \epsilon : H^*(X \times X \times X)(d) \longrightarrow H^*(X)
$$
est l'identité. Pour $a \in H^*(X)$,
la composée envoie $a$ sur
$$
(1 \otimes \lambda)(\Delta_{23}^*(q_{12}^*\gamma([\Delta]) \cup q_3^*a)) =
(1 \otimes \lambda)(\gamma([\Delta]) \cup p_2^*a)
$$,
où $q_i : X \times X \times X \to X$ et
$q_{ij} : X \times X \times X \to X \times X$ sont les projections,
$\Delta_{23} : X \times X \to X \times X \times X$ est la diagonale et
$p_i : X \times X \to X$ sont les projections.
L'égalité résulte de $\Delta_{23}^*(q_{12}^*\gamma([\Delta]) =
\Delta_{23}^*\gamma([\Delta \times X]) = \gamma([\Delta])$
et de $\Delta_{23}^* q_3^*a = p_2^*a$.
Puisque $\gamma([\Delta]) \cup p_1^*a = \gamma([\Delta]) \cup p_2^*a$
(voir ci-dessous), l'expression ci-dessus se simplifie en
$$
(1 \otimes \lambda)(\gamma([\Delta]) \cup p_1^*a) = a
$$
par notre choix de $\lambda$, comme voulu. La seconde condition,
$(\epsilon \otimes 1) \circ (1 \otimes \eta) = \text{id}$,
de Catégories, définition \ref{categories-definition-dual},
se démontre exactement de la
même manière.

\medskip\noindent
Notons que $p_1^*a$ et $\text{pr}_2^*a$ se restreignent à la même
classe de cohomologie sur $\Delta \subset X \times X$. De plus, on
a $\mathcal{C}_{\Delta/X \times X} = \Omega^1_\Delta$, qui
est la restriction de $p_1^*\Omega^1_X$. Le
lemme \ref{weil-lemma-A5-A6-imply} donne donc
$\gamma([\Delta]) \cup p_1^*a = \gamma([\Delta]) \cup p_2^*a$,
ce qui achève la démonstration.
\end{proof}

\begin{remark}[Unicité des applications traces]
\label{weil-remark-trace}
Supposons que les données (D0), (D1) et (D2') satisfassent aux axiomes (A1) -- (A7).
Soit $X$ un schéma projectif lisse sur $k$, non vide
et équidimensionnel de dimension $d$. En combinant les arguments
des démonstrations du lemme \ref{weil-lemma-poincare-duality} et de
Homologie, lemme \ref{homology-lemma-left-dual-graded-vector-spaces},
on voit que
$$
\gamma([\Delta]) \in \bigoplus\nolimits_i H^i(X) \otimes H^{2d - i}(X)(d)
$$
définit une dualité parfaite entre $H^i(X)$ et $H^{2d - i}(X)(d)$
pour tout $i$.
En particulier, l'application linéaire $\int_X = \lambda : H^{2d}(X)(d) \to F$ de
l'axiome (A6) est unique ! Nous appellerons désormais l'application linéaire $\int_X$
l'application trace de $X$.
\end{remark}

\begin{lemma}
\label{weil-lemma-trace-product}
Supposons que les données (D0), (D1) et (D2') satisfassent aux axiomes (A1) -- (A7).
Alors l'axiome (B) de la section \ref{weil-section-axioms} est vérifié.
\end{lemma}

\begin{proof}
L'axiome (B)(a) résulte immédiatement de l'axiome (A5).
Soient $X$ et $Y$ des schémas projectifs lisses non vides sur $k$,
équidimensionnels de dimensions $d$ et $e$. Pour établir l'axiome (B)(b),
observons que la diagonale $\Delta_{X \times Y}$ de $X \times Y$
est le produit d'intersection des images réciproques des diagonales
$\Delta_X$ de $X$ et $\Delta_Y$ de $Y$ par les projections
$p : X \times Y \times X \times Y \to X \times X$ et
$q : X \times Y \times X \times Y \to Y \times Y$.
La compatibilité de $\gamma$ avec les produits d'intersection entraîne alors
que
$$
\gamma([\Delta_{X \times Y}]) \in
H^{2d + 2e}(X \times Y \times X \times Y)(d + e)
$$
est le cup-produit des images réciproques de $\gamma([\Delta_X])$
et de $\gamma([\Delta_Y])$ par $p$ et $q$. Écrivons
$$
\gamma([\Delta_{X \times Y}]) = \sum \eta_{X \times Y, i}
\text{ avec }
\eta_{X \times Y, i} \in
H^i(X \times Y) \otimes H^{2d + 2e - i}(X \times Y)(d + e)
$$,
et de même $\gamma([\Delta_X]) = \sum \eta_{X, i}$ et
$\gamma([\Delta_Y]) = \sum \eta_{Y, i}$. L'observation précédente
donne
$$
\eta_{X \times Y, 0} =
\sum\nolimits_{i \in \mathbf{Z}} p^*\eta_{X, i} \cup q^*\eta_{Y, -i}
$$
(Si notre théorie cohomologique s'annule en degrés négatifs, ce qui sera
le cas dans presque toutes les situations, seul le terme $i = 0$ contribue,
et $\eta_{X \times Y, 0}$ appartient à
$H^0(X) \otimes H^0(Y) \otimes H^{2d}(X)(d) \otimes H^{2e}(Y)(e)$, comme attendu ;
mais nous n'en avons pas besoin.) Puisque $\lambda_X : H^{2d}(X)(d) \to F$ et
$\lambda_Y : H^{2e}(Y)(e) \to F$ envoient $\eta_{X, 0}$ et $\eta_{Y, 0}$
sur $1$ dans $H^*(X)$ et $H^*(Y)$, on voit que $\lambda_X \otimes \lambda_Y$
envoie $\eta_{X \times Y, 0}$ sur $1$ dans
$H^*(X) \otimes H^*(Y) = H^*(X \times Y)$, ce qui achève la démonstration.
\end{proof}

\begin{lemma}
\label{weil-lemma-trace-base}
Supposons que les données (D0), (D1) et (D2') satisfassent aux axiomes (A1) -- (A7).
Alors l'axiome (C)(d) de la section \ref{weil-section-axioms} est vérifié.
\end{lemma}

\begin{proof}
On a $\gamma([\Spec(k)]) = 1 \in H^*(\Spec(k))$ par construction.
Puisque
$$
H^0(\Spec(k)) = F,\quad
H^0(\Spec(k) \times \Spec(k)) = H^0(\Spec(k)) \otimes_F H^0(\Spec(k))
$$,
l'application $\int_{\Spec(k)} = \lambda$ de l'axiome (A6) doit envoyer $1$ sur $1$,
car nous avons vu que
$\int_{\Spec(k) \times \Spec(k)} = \int_{\Spec(k)} \int_{\Spec(k)}$
au lemme \ref{weil-lemma-trace-product}.
\end{proof}

\noindent
Supposons que les données (D0), (D1) et (D2') satisfassent aux axiomes (A1) -- (A7).
On obtient alors les données (D0), (D1), (D2) et (D3) de la
section \ref{weil-section-axioms},
satisfaisant aux axiomes (A), (B), (C)(a), (C)(c) et (C)(d)
de la section \ref{weil-section-axioms} ; voir
les lemmes \ref{weil-lemma-poincare-duality}, \ref{weil-lemma-trace-product} et
\ref{weil-lemma-trace-base}.
De plus, on dispose des images directes $f_* : H^*(X) \to H^*(Y)$
construites à la section \ref{weil-section-axioms}. Le seul axiome de la
section \ref{weil-section-axioms}
qui ne soit pas encore établi est l'axiome (C)(b).

\begin{lemma}
\label{weil-lemma-ok-for-projective-bundle}
Supposons que les données (D0), (D1) et (D2') satisfassent aux axiomes (A1) -- (A7).
Soit $p : P \to X$ comme dans l'axiome (A3), avec $X$ non vide et équidimensionnel.
Alors $\gamma$ commute à l'image directe par $p$.
\end{lemma}

\begin{proof}
Il suffit de le démontrer sur des générateurs de $\CH_*(P)$.
Il suffit donc de le démontrer pour une classe de cycle de la
forme $\xi^i \cdot p^*\alpha$, où $0 \leq i \leq r - 1$
et $\alpha \in \CH_*(X)$. Notons que $p_*(\xi^i \cdot p^*\alpha) = 0$
si $i < r - 1$, et que $p_*(\xi^{r - 1} \cdot p^*\alpha) = \alpha$.
D'autre part, on a
$\gamma(\xi^i \cdot p^*\alpha) = c^i \cup p^*\gamma(\alpha)$,
et la formule de projection (lemme \ref{weil-lemma-pushforward})
donne
$$
p_*\gamma(\xi^i \cdot p^*\alpha) = p_*(c^i) \cup \gamma(\alpha)
$$
Il suffit ainsi de montrer que $p_*c^i = 0$ pour $i < r - 1$ et
que $p_*c^{r - 1} = 1$. De manière équivalente, il suffit de démontrer que
$\lambda_P : H^{2d + 2r - 2}(P)(d + r - 1) \to F$, définie par
les règles
$$
\lambda_P(c^i \cup p^*(a)) =
\left\{
\begin{matrix}
0 & \text{si} & i < r - 1 \\
\int_X(a) & \text{si} & i = r - 1
\end{matrix}
\right.
$$,
satisfait à la condition de l'axiome (A5). Cela résulte du
calcul de la classe de la diagonale de $P$ au
lemme \ref{weil-lemma-diagonal-projective-bundle}.
\end{proof}

\begin{lemma}
\label{weil-lemma-integrate-1}
Supposons que les données (D0), (D1) et (D2') satisfassent aux axiomes (A1) -- (A7).
Si $k'/k$ est une extension galoisienne, on a
$\int_{\Spec(k')} 1 = [k' : k]$.
\end{lemma}

\begin{proof}
Observons que
$$
\Spec(k') \times \Spec(k') =
\coprod\nolimits_{\sigma \in \text{Gal}(k'/k)}
(\Spec(\sigma) \times \text{id})^{-1} \Delta
$$
comme cycles sur $\Spec(k') \times \Spec(k')$.
Notre construction de $\gamma$ envoie toujours $[X]$ sur $1$ dans $H^0(X)$. Ainsi,
$1 \otimes 1 = 1 = \sum (\Spec(\sigma) \times \text{id})^*\gamma([\Delta])$.
Notons $\lambda : H^0(\Spec(k')) \to F$ l'application de
l'axiome (A6), c'est-à-dire que $(\text{id} \otimes \lambda)(\gamma(\Delta)) = 1$
dans $H^0(\Spec(k'))$. On obtient
\begin{align*}
\lambda(1) 1
& =
(\text{id} \otimes \lambda)(1 \otimes 1) \\
& =
(\text{id} \otimes \lambda)(
\sum (\Spec(\sigma) \times \text{id})^*\gamma([\Delta])) \\
& =
\sum (\Spec(\sigma) \times \text{id})^*(
(\text{id} \otimes \lambda)(\gamma([\Delta])) \\
& =
\sum (\Spec(\sigma) \times \text{id})^*(1) \\
& =
[k' : k]
\end{align*}
Puisque $\lambda$ est une autre notation pour $\int_{\Spec(k')}$
(remarque \ref{weil-remark-trace}), la démonstration est achevée.
\end{proof}

\begin{lemma}
\label{weil-lemma-enough}
Supposons que les données (D0), (D1) et (D2') satisfassent aux axiomes (A1) -- (A7).
Pour montrer que $\gamma$ commute à l'image directe, il suffit
de montrer que $i_*(1) = \gamma([Z])$ lorsque $i : Z \to X$ est une immersion
fermée de schémas projectifs lisses non vides et équidimensionnels sur $k$.
\end{lemma}

\begin{proof}
Nous utiliserons sans le mentionner de nouveau le fait que l'application
classe de cycle $\gamma$ construite au lemme \ref{weil-lemma-cycle-classes}
est compatible avec les produits d'intersection et les images réciproques,
et que nous avons déjà établi les axiomes
(A), (B), (C)(a), (C)(c) et (C)(d) de la section \ref{weil-section-axioms} ; voir
le lemme \ref{weil-lemma-poincare-duality},
la remarque \ref{weil-remark-trace} et
les lemmes \ref{weil-lemma-trace-product} et \ref{weil-lemma-trace-base}.
En particulier, on peut utiliser, par exemple, le lemme \ref{weil-lemma-pushforward}
pour voir que l'image directe sur $H^*$ est compatible avec la composition
et satisfait à la formule de projection.

\medskip\noindent
Soit $f : X \to Y$ un morphisme de schémas projectifs lisses
non vides et équidimensionnels sur $k$.
Nous cherchons à montrer $f_*\gamma(\alpha) = \gamma(f_*\alpha)$
pour toute classe de cycle $\alpha$ sur $X$.
On peut écrire $\alpha$ comme une combinaison $\mathbf{Q}$-linéaire de produits de
classes de Chern de $\mathcal{O}_X$-modules localement libres
(Homologie de Chow, lemme \ref{chow-lemma-K-tensor-Q}).
On peut donc supposer que $\alpha$ est un produit de classes de Chern de
$\mathcal{O}_X$-modules localement libres de type fini
$\mathcal{E}_1, \ldots, \mathcal{E}_r$.
Choisissons $p : P \to X$ comme dans le principe de scindage
(Homologie de Chow, lemme \ref{chow-lemma-splitting-principle}).
D'après Homologie de Chow, remarque \ref{chow-remark-the-proof-shows-more},
$p$ est une composée de fibrés projectifs et
$\alpha = p_*(\xi_1 \cap \ldots \cap \xi_d \cap \cdot p^*\alpha)$,
où les $\xi_i$ sont des premières classes de Chern de modules inversibles.
D'après le lemme \ref{weil-lemma-ok-for-projective-bundle},
$p_*$ commute aux classes de cycles.
Il suffit donc de démontrer la propriété pour la composée
$f \circ p$. Puisque $p^*\mathcal{E}_1, \ldots, p^*\mathcal{E}_r$
admettent des filtrations dont les quotients successifs sont des modules
inversibles, on est ramené au cas où $\alpha$ est
de la forme $\xi_1 \cap \ldots \cap \xi_t \cap [X]$,
pour certaines premières classes de Chern $\xi_i$ de modules inversibles $\mathcal{L}_i$.

\medskip\noindent
Supposons
$\alpha = c_1(\mathcal{L}_1) \cap \ldots \cap c_1(\mathcal{L}_t) \cap [X]$,
pour certains modules inversibles $\mathcal{L}_i$ sur $X$.
Soit $\mathcal{L}$ un $\mathcal{O}_X$-module inversible ample.
Pour $n \gg 0$, les $\mathcal{O}_X$-modules inversibles
$\mathcal{L}^{\otimes n}$ et
$\mathcal{L}_1 \otimes \mathcal{L}^{\otimes n}$ sont tous deux
très amples sur $X$ relativement à $k$ ; voir
Morphismes, lemme \ref{morphisms-lemma-invertible-add-enough-ample-very-ample}.
Puisque $c_1(\mathcal{L}_1) = c_1(\mathcal{L}_1 \otimes \mathcal{L}^{\otimes n})
- c_1(\mathcal{L}^{\otimes n})$, on est ramené au cas où
$\mathcal{L}_1$ est très ample. En répétant ce raisonnement avec $\mathcal{L}_i$
pour $i = 2, \ldots, t$, on se ramène au cas où $\mathcal{L}_i$
est très ample sur $X$ relativement à $k$ pour tout $i = 1, \ldots, t$.

\medskip\noindent
Supposons $k$ infini et $\alpha = 
c_1(\mathcal{L}_1) \cap \ldots \cap c_1(\mathcal{L}_t) \cap [X]$,
pour certains modules inversibles très amples $\mathcal{L}_i$ sur $X$ relativement à $k$.
Par Bertini, sous la forme de Variétés, lemme \ref{varieties-lemma-bertini},
on peut trouver successivement des sections régulières $s_i$ de $\mathcal{L}_i$
telles que les schémas $Z(s_1) \cap \ldots \cap Z(s_i)$
soient lisses sur $k$ et de codimension $i$ dans $X$.
Par la construction du cap-produit avec la première classe
d'un module inversible (qui remonte à
Homologie de Chow, définition \ref{chow-definition-divisor-invertible-sheaf}),
on est ainsi ramené au cas où $\alpha = [Z]$,
pour un sous-schéma fermé lisse non vide $Z \subset X$
équidimensionnel.

\medskip\noindent
Supposons $\alpha = [Z]$, où $Z \subset X$ est un sous-schéma fermé lisse.
Choisissons une immersion fermée $X \to \mathbf{P}^n$. On peut factoriser $f$ sous la forme
$$
X \to Y \times \mathbf{P}^n \to Y
$$
Puisque le résultat est connu pour le second morphisme par
le lemme \ref{weil-lemma-ok-for-projective-bundle},
il suffit de le démontrer lorsque
$\alpha = [Z]$, où $i : Z \to X$ est une immersion fermée
et $f$ une immersion fermée.
Alors $j = f \circ i$ est également une immersion fermée.
L'hypothèse appliquée à $i$ et à $j$ donne la conclusion.

\medskip\noindent
Il reste à démontrer le lemme lorsque $k$ est fini. Nous conseillons vivement au
lecteur de passer la fin de la démonstration. Tout ce qui précède reste valable,
sauf que nous ne savons pas choisir les sections
$s_i$ qui découpent $Z$ sur $k$, puisque $k$ est fini. Nous savons toutefois
que l'on peut trouver $s_i$ sur la clôture algébrique $\overline{k}$
de $k$, par le même lemme. Il existe donc
une extension finie $k'/k$ telle que nos sections $s_i$
soient définies sur $k'$. Notons $\pi : X_{k'} \to X$ la projection.
Les arguments ci-dessus donnent la conclusion voulue,
sous l'hypothèse du lemme,
pour le cycle $\pi^*\alpha$ et le morphisme
$f \circ \pi : X_{k'} \to Y$.
On a $\pi_*\pi^*\alpha = [k' : k] \alpha$ ; voir
Homologie de Chow, lemme \ref{chow-lemma-finite-flat}.
D'autre part, on a
$$
\pi_*\gamma(\pi^*\alpha) = \pi_*\pi^*\gamma(\alpha) =
\gamma(\alpha) \pi_*1
$$
par la formule de projection de notre théorie cohomologique. Observons
que $\pi$ est une projection (!) ; on a donc
$\pi_*(1) = \int_{\Spec(k')}(1) 1$ d'après
le lemme \ref{weil-lemma-pr2star}. Pour achever la démonstration dans
le cas d'un corps fini, il suffit donc de démontrer
$\int_{\Spec(k')}(1) = [k' : k]$, ce que nous faisons au
lemme \ref{weil-lemma-integrate-1}.
\end{proof}

\noindent
Dans les lemmes ci-dessous, nous utilisons les grassmanniennes définies et construites
dans Constructions, section \ref{constructions-section-grassmannian}.

\begin{lemma}
\label{weil-lemma-grassmanian}
Supposons que les données (D0), (D1) et (D2') vérifient les axiomes (A1) -- (A7).
Étant donnés des entiers $0 < l < n$ et un schéma projectif lisse
équidimensionnel non vide $X$ sur $k$, considérons la projection
$p : X \times \mathbf{G}(l, n) \to X$.
Alors $\gamma$ commute à l'image directe par $p$.
\end{lemma}

\begin{proof}
Si $l = 1$ ou $l = n - 1$, alors $p$ est un fibré projectif et
le résultat découle du lemme \ref{weil-lemma-ok-for-projective-bundle}.
En général, il existe un morphisme
$$
h : Y \to X \times \mathbf{G}(l, n)
$$
tel que $h$ et $p \circ h$ soient tous deux des composés de fibrés
en espaces projectifs. En effet, notons $\mathbf{G}(1, 2, \ldots, l; n)$
la variété de drapeaux partiels. Alors le morphisme
$$
\mathbf{G}(1, 2, \ldots, l; n) \to \mathbf{G}(l, n)
$$
est un composé de fibrés en espaces projectifs, et le morphisme
structural $\mathbf{G}(1, 2, \ldots, l; n) \to \Spec(k)$
est lui aussi de cette forme. On peut donc poser $Y = X \times \mathbf{G}(1, 2, \ldots, l; n)$.
Puisque tout cycle sur $X \times \mathbf{G}(l, n)$ est l'image directe
d'un cycle sur $Y$, le résultat pour $Y \to X$ et celui pour
$Y \to X \times \mathbf{G}(l, n)$ entraînent le résultat pour $p$.
\end{proof}

\begin{lemma}
\label{weil-lemma-enough-better}
Supposons que les données (D0), (D1) et (D2') vérifient les axiomes (A1) -- (A7).
Pour montrer que $\gamma$ commute à l'image directe, il suffit
de montrer que $i_*(1) = \gamma([Z])$ lorsque $i : Z \to X$ est une immersion
fermée de schémas projectifs lisses équidimensionnels non vides sur $k$
telle que la classe de $\mathcal{C}_{Z/X}$ dans $K_0(Z)$ soit
l'image réciproque d'une classe de $K_0(X)$.
\end{lemma}

\begin{proof}
D'après le lemme \ref{weil-lemma-enough}, il suffit de montrer que $i_*(1) = \gamma([Z])$
lorsque $i : Z \to X$ est une immersion fermée de schémas
projectifs lisses équidimensionnels non vides
sur $k$. Supposons $Z$ de codimension $r$ dans $X$.
Soit $\mathcal{L}$ un module inversible suffisamment ample sur $X$.
Choisissons $n > 0$ et une surjection
$$
\mathcal{O}_Z^{\oplus n} \to \mathcal{C}_{Z/X} \otimes \mathcal{L}|_Z
$$
Ceci donne un morphisme $g : Z \to \mathbf{G}(n - r, n)$
vers la grassmannienne sur $k$, voir
Constructions, section \ref{constructions-section-grassmannian}.
Considérons la composée
$$
Z \to X \times \mathbf{G}(n - r, n) \to X
$$
L'image directe par le second morphisme est compatible avec les classes
de cycles d'après le lemme \ref{weil-lemma-grassmanian}. Le faisceau conormal $\mathcal{C}$
de l'immersion fermée $Z \to X \times \mathbf{G}(n - r, n)$ s'insère dans
une suite exacte courte
$$
0 \to \mathcal{C}_{Z/X} \to \mathcal{C} \to
g^*\Omega_{\mathbf{G}(n - r, n)} \to 0
$$
Voir Compléments sur les morphismes, lemme
\ref{more-morphisms-lemma-two-unramified-morphisms-formally-smooth}.
Puisque $\mathcal{C}_{Z/X} \otimes \mathcal{L}|_Z$ est l'image
réciproque d'un faisceau localement libre de rang fini sur $\mathbf{G}(n - r, n)$,
on en déduit que la classe de $\mathcal{C}$ dans $K_0(Z)$
est l'image réciproque d'une classe de $K_0(X \times \mathbf{G}(n - r, n))$.
On a donc la propriété pour $Z \to X \times \mathbf{G}(n - r, n)$,
ce qui permet de conclure.
\end{proof}

\begin{lemma}
\label{weil-lemma-injective-H0}
Supposons que les données (D0), (D1) et (D2') vérifient les axiomes (A1) -- (A7).
Si $k''/k'/k$ sont des extensions finies séparables de corps, alors
$H^0(\Spec(k')) \to H^0(\Spec(k''))$ est injective.
\end{lemma}

\begin{proof}
On peut remplacer $k''$ par sa clôture normale sur $k$,
qui est galoisienne sur $k$, voir
Corps, lemme \ref{fields-lemma-normal-closure-galois}.
Alors $k''$ est aussi galoisien sur $k'$, voir
Corps, lemme \ref{fields-lemma-galois-goes-up}.
On en déduit un isomorphisme
$$
k' \otimes_k k'' \longrightarrow
\prod\nolimits_{\sigma \in \text{Gal}(k''/k')} k'',\quad
\eta \otimes \zeta \longmapsto (\eta \sigma(\zeta))_\sigma
$$
Il donne un isomorphisme
$\coprod_\sigma \Spec(k'') \to \Spec(k') \times \Spec(k'')$
qui induit en cohomologie l'isomorphisme
$$
H^*(\Spec(k')) \otimes_F H^*(\Spec(k''))
\to
\prod\nolimits_\sigma H^*(\Spec(k'')),\quad
a' \otimes a'' \longmapsto (\pi^*a' \cup \Spec(\sigma)^*a'')_\sigma
$$
où $\pi : \Spec(k'') \to \Spec(k')$ est le morphisme
correspondant à l'inclusion de $k'$ dans $k''$.
On conclut en prenant $a'' = 1$.
\end{proof}

\begin{lemma}
\label{weil-lemma-pushforward-blowup}
Supposons que les données (D0), (D1) et (D2') vérifient les axiomes (A1) -- (A8).
Soit $b : X' \to X$ l'éclatement d'un schéma projectif lisse $X$
sur $k$, non vide, équidimensionnel de dimension $d$,
le long d'un centre lisse partout non dense $Z$. Alors $b_*(1) = 1$.
\end{lemma}

\begin{proof}
On peut remplacer $X$ par une composante connexe de $X$ (certains détails
sont omis). On peut donc supposer $X$ connexe, donc irréductible.
Posons $k' = \Gamma(X, \mathcal{O}_X) = \Gamma(X', \mathcal{O}_{X'})$ ;
nous omettons la démonstration de cette égalité. Choisissons un point fermé $x' \in X'$
qui n'appartient pas au diviseur exceptionnel et dont le corps résiduel
$k''$ est séparable sur $k$ ; ceci est possible d'après
Variétés, lemme \ref{varieties-lemma-smooth-separable-closed-points-dense}.
Notons $x \in X$ son image (dont le corps résiduel est bien entendu
égal lui aussi à $k''$). Considérons le diagramme
$$
\xymatrix{
x' \times X' \ar[r] \ar[d] & X' \times X' \ar[d] \\
x \times X \ar[r] & X \times X
}
$$
L'image réciproque de la classe de la diagonale $\Delta = \Delta_X$ est la classe du
« point diagonal » $\delta_x : x \to x \times X$, et de même pour la classe de
la diagonale $\Delta'$. D'autre part, l'image réciproque du point diagonal $\delta_x$
par la flèche verticale de gauche est le point diagonal $\delta_{x'}$.
Écrivons $\gamma([\Delta]) = \sum \eta_i$ avec
$\eta_i \in H^i(X) \otimes H^{2d - i}(X)(d)$ et
$\gamma([\Delta']) = \sum \eta'_i$ avec
$\eta'_i \in H^i(X') \otimes H^{2d - i}(X')(d)$.
Les arguments précédents montrent que $\eta_0$ et $\eta'_0$ ont pour image la même
classe dans
$$
H^0(x') \otimes_F H^{2d}(X')(d)
$$
On a $H^0(\Spec(k')) = H^0(X) = H^0(X')$ d'après l'axiome (A8).
D'après le lemme \ref{weil-lemma-injective-H0}, cet espace commun s'injecte
dans $H^0(x')$. On en déduit que $\eta_0$ a pour image $\eta'_0$ par l'application
$$
H^0(X) \otimes_F H^{2d}(X)(d)
\longrightarrow
H^0(X') \otimes_F H^{2d}(X')(d)
$$
Cela signifie que $\int_X$ est égal à $\int_{X'}$ composé avec
l'application d'image réciproque. Le lemme est démontré.
\end{proof}

\begin{lemma}
\label{weil-lemma-done}
Supposons que les données (D0), (D1) et (D2') vérifient les axiomes (A1) -- (A8).
Alors l'application classe de cycle $\gamma$ commute à l'image directe.
\end{lemma}

\begin{proof}
Soit $i : Z \to X$ comme dans le lemme \ref{weil-lemma-enough-better}. Considérons
le diagramme
$$
\xymatrix{
E \ar[r]_j \ar[d]_\pi & X' \ar[d]^b \\
Z \ar[r]^i & X
}
$$
Soit $\theta \in \CH^{r - 1}(X')$ comme dans le
lemme \ref{weil-lemma-divide-pullback-good-blowing-up}.
Alors $\pi_*j^!\theta = [Z]$ dans $\CH_*(Z)$ entraîne que
$\pi_*\gamma(j^!\theta) = 1$ d'après le lemme \ref{weil-lemma-ok-for-projective-bundle},
car $\pi$ est un fibré en espaces projectifs.
On obtient donc
$$
i_*(1) = i_*(\pi_*(\gamma(j^!\theta))) =
b_*j_*(j^*\gamma(\theta)) =
b_*(j_*(1) \cup \gamma(\theta))
$$
On a $j_*(1) = \gamma([E])$ d'après (A9). Ceci est donc égal à
$$
b_*(\gamma([E]) \cup \gamma(\theta)) =
b_*(\gamma([E] \cdot \theta)) =
b_*(\gamma(b^*[Z])) =
b_*b^*\gamma([Z]) = b_*(1) \cup \gamma([Z])
$$
Comme $b_*(1) = 1$ d'après le lemme \ref{weil-lemma-pushforward-blowup}, la
démonstration est achevée.
\end{proof}

\begin{proposition}
\label{weil-proposition-get-weil}
Supposons que les données (D0), (D1) et (D2') vérifient les axiomes (A1) -- (A8).
On obtient alors une théorie cohomologique de Weil.
\end{proposition}

\begin{proof}
Les axiomes (A), (B), (C)(a), (C)(c) et (C)(d) de la
section \ref{weil-section-axioms} résultent des
lemmes \ref{weil-lemma-poincare-duality}, \ref{weil-lemma-trace-product} et
\ref{weil-lemma-trace-base}.
L'axiome (C)(b) résulte du
lemme \ref{weil-lemma-done}.
Enfin, la condition supplémentaire de la
définition \ref{weil-definition-weil-cohomology-theory}
est satisfaite, car elle coïncide avec notre axiome (A8).
\end{proof}

\noindent
Le lemme suivant est parfois utile pour montrer que l'on obtient une
théorie cohomologique de Weil sur un corps non algébriquement clos, en se
ramenant à un corps algébriquement clos.

\begin{lemma}
\label{weil-lemma-check-over-extension}
Soit $k'/k$ une extension de corps. Soit $F'/F$ une extension
de corps de caractéristique $0$. Supposons donnés
\begin{enumerate}
\item des données (D0), (D1), (D2') pour $k$ et $F$, notées
$F(1), H^*, c_1^H$,
\item des données (D0), (D1), (D2') pour $k'$ et $F'$, notées
$F'(1), (H')^*, c_1^{H'}$, et
\item un isomorphisme $F(1) \otimes_F F' \to F'(1)$, des isomorphismes fonctoriels
$H^*(X) \otimes_F F' \to (H')^*(X_{k'})$ sur la catégorie des schémas projectifs lisses
$X$ sur $k$, tels que les diagrammes
$$
\xymatrix{
\Pic(X) \ar[r]_{c_1^H} \ar[d] & H^2(X)(1) \ar[d] \\
\Pic(X_{k'}) \ar[r]^{c_1^{H'}} & (H')^2(X_{k'})(1)
}
$$
commutent.
\end{enumerate}
Dans ce cas, si $F'(1), (H')^*, c_1^{H'}$ vérifient les axiomes (A1) -- (A9),
il en est de même de $F(1), H^*, c_1^H$.
\end{lemma}

\begin{proof}
Nous vérifions les axiomes un à un.

\medskip\noindent
Axiome (A1). Il faut montrer que $H^*(\emptyset) = 0$ et que
$(i^*, j^*) : H^*(X \amalg Y) \to H^*(X) \times H^*(Y)$
est un isomorphisme, où $i$ et $j$ sont les coprojections.
Par fonctorialité des isomorphismes
$H^*(X) \otimes_F F' \to (H')^*(X_{k'})$, cela
résulte de l'algèbre linéaire : si $\varphi : V \to W$ est une application $F$-linéaire
d'espaces vectoriels sur $F$, alors $\varphi$ est un isomorphisme si et seulement si
$\varphi_{F'} : V \otimes_F F' \to W \otimes_F F'$ est un isomorphisme.

\medskip\noindent
Axiome (A2). Cela signifie qu'étant donnés un morphisme $f : X \to Y$ de schémas projectifs lisses
sur $k$ et un $\mathcal{O}_Y$-module inversible $\mathcal{N}$,
on a $f^*c_1^H(\mathcal{L}) = c_1^H(f^*\mathcal{L})$. Ceci résulte immédiatement
de la propriété correspondante pour $c_1^{H'}$, des diagrammes commutatifs
du lemme et du fait que l'application canonique
$V \to V \otimes_F F'$ est injective pour tout espace vectoriel sur $F$, noté $V$.

\medskip\noindent
Axiome (A3). Il découle du principe énoncé dans la démonstration de
l'axiome (A1) et de la compatibilité de $c_1^H$ et $c_1^{H'}$.

\medskip\noindent
Axiome (A4). Soit $i : Y \to X$ l'inclusion d'un diviseur de Cartier
effectif sur $k$, où $X$ et $Y$ sont tous deux projectifs et lisses
sur $k$. Pour $a \in H^*(X)$ tel que
$i^*a = 0$, il faut montrer que $a \cup c_1^H(\mathcal{O}_X(Y)) = 0$.
Notons $a' \in (H')^*(X_{k'})$ l'image de $a$.
L'hypothèse entraîne que $(i')^*a' = 0$, où $i' : Y_{k'} \to X_{k'}$
est le changement de base de $i$. On obtient donc
$a' \cup c_1^{H'}(\mathcal{O}_{X_{k'}}(Y_{k'})) = 0$ par l'axiome
pour $(H')^*$. Puisque $a' \cup c_1^{H'}(\mathcal{O}_{X_{k'}}(Y_{k'}))$
est l'image de $a \cup c_1^H(\mathcal{O}_X(Y))$, on conclut par
le principe énoncé dans la démonstration de l'axiome (A2).

\medskip\noindent
Axiome (A5). Cela signifie que $H^*(\Spec(k)) = F$ et que, pour $X$ et $Y$ projectifs
lisses sur $k$, l'application $H^*(X) \otimes_F H^*(Y) \to H^*(X \times Y)$,
$a \otimes b \mapsto p^*(a) \cup q^*(b)$ est un isomorphisme,
où $p$ et $q$ sont les projections. Ceci découle du principe
énoncé dans la démonstration de l'axiome (A1).

\medskip\noindent
Interrompons le cours de l'argumentation pour montrer que, pour tout
schéma projectif lisse $X$ sur $k$, le diagramme
$$
\xymatrix{
\CH^*(X) \ar[r]_-\gamma \ar[d]_{g^*} & \bigoplus H^{2i}(X)(i) \ar[d] \\
\CH^*(X_{k'}) \ar[r]^-{\gamma'} & \bigoplus (H')^{2i}(X_{k'})(i)
}
$$
commute. Remarquons que l'on dispose de $\gamma$, puisque l'on a déjà établi les axiomes
(A1) -- (A4) ; voir le lemme \ref{weil-lemma-cycle-classes}.
La flèche verticale de gauche est d'ailleurs celle qui est étudiée dans
Homologie de Chow, section \ref{chow-section-change-base},
pour le morphisme de schémas de base $g : \Spec(k') \to \Spec(k)$.
Plus précisément, c'est l'application donnée dans
Homologie de Chow, lemme \ref{chow-lemma-pullback-base-change}.
Prenons $\alpha \in \CH^*(X)$. Écrivons $\alpha = ch(\beta) \cap [X]$
dans $\CH^*(X) \otimes \mathbf{Q}$,
pour un certain $\beta \in K_0(\textit{Vect}(X)) \otimes \mathbf{Q}$,
de sorte que $\gamma(\alpha) = ch^{H}(\beta)$ ; telle est notre construction de $\gamma$.
Puisque l'application de Homologie de Chow, lemme \ref{chow-lemma-pullback-base-change},
est compatible avec le cap-produit par les classes de Chern d'après
Homologie de Chow, lemme \ref{chow-lemma-pullback-base-change-chern-classes},
on voit que $g^*\alpha = ch((X_{k'} \to X)^*\beta) \cap [X_{k'}]$.
Ainsi $\gamma'(g^*\alpha) = ch^{H'}((X_{k'} \to X)^*\beta)$.
Le diagramme sera donc commutatif si, pour tout
$\mathcal{O}_X$-module localement libre $\mathcal{E}$ de rang $r$ et tout $0 \leq i \leq r$,
l'élément $c_i^H(\mathcal{E})$ de $H^{2i}(X)(i)$
a pour image l'élément $c_i^{H'}(\mathcal{E}_{k'})$ de $(H')^{2i}(X_{k'})(i)$.
Comme on dispose de la formule du fibré en espaces projectifs aussi bien pour
$X$ que pour $X'$, on peut, pour le montrer, remplacer $X$ par un fibré en espaces projectifs
sur $X$ un nombre fini de fois. On peut donc supposer
que $\mathcal{E}$ possède une filtration dont les quotients gradués sont des
$\mathcal{O}_X$-modules inversibles
$\mathcal{L}_1, \ldots, \mathcal{L}_r$.
Voir Homologie de Chow, lemme \ref{chow-lemma-splitting-principle} et
remarque \ref{chow-remark-the-proof-shows-more}.
Alors $c^H_i(\mathcal{E}$ est le $i$-ième polynôme symétrique élémentaire
en $c^H_1(\mathcal{L}_1), \ldots, c^H_1(\mathcal{L}_r)$,
et l'on conclut par l'hypothèse de compatibilité des
premières classes de Chern.

\medskip\noindent
Axiome (A6). Supposons donnés des espaces vectoriels sur $F$,
$V$, $W$, un élément $v \in V$ et un tenseur $\xi \in V \otimes_F W$.
Notons $V' = V \otimes_F F'$, $W' = W \otimes_F F'$, et $v'$, $\xi'$
les images de $v$, $\xi$ dans $V'$, $V' \otimes_{F'} W'$. Le principe d'algèbre linéaire
utilisé pour démontrer l'axiome (A6) est le suivant :
il existe une application $F$-linéaire $\lambda : W \to F$ telle que
$(1 \otimes \lambda)\xi = v$ si et seulement s'il existe une application $F'$-linéaire
$\lambda' : W \otimes_F F' \to F'$ telle que
$(1 \otimes \lambda')\xi' = v'$.

\medskip\noindent
Soit $X$ un schéma projectif lisse équidimensionnel non vide
sur $k$, de dimension $d$. Notons $\gamma = \gamma([\Delta])$
dans $H^{2d}(X \times X)(d)$ (les produits fibrés sans indication de base seront pris sur $k$).
Remarquons, ou rappelons, que ceci a un sens, puisque l'on a déjà établi les axiomes (A1) -- (A4) ;
voir le lemme \ref{weil-lemma-cycle-classes}. On peut écrire
$$
\gamma = \sum \gamma_i, \quad
\gamma_i \in H^i(X) \otimes_F H^{2d - i}(X)(d)
$$
dans la décomposition de K\"unneth. De même, notons
$\gamma' = \gamma([\Delta']) = \sum \gamma'_i$
dans $(H')^{2d}(X_{k'} \times_{k'} X_{k'})(d)$.
D'après le principe d'algèbre linéaire rappelé ci-dessus, il suffit
de montrer que $\gamma_0$ a pour image $\gamma'_0$ dans
$(H')^0(X) \otimes_{F'} (H')^{2d}(X')(d)$.
La compatibilité des décompositions de K\"unneth
montre qu'il suffit de vérifier que $\gamma$ a pour image
$\gamma'$ dans
$$
(H')^{2d}(X_{k'} \times_{k'} X_{k'})(d) = (H')^{2d}((X \times X)_{k'})(d)
$$
Puisque $\Delta_{k'} = \Delta'$, ceci résulte de la discussion précédente.

\medskip\noindent
Axiome (A7). Il résulte du fait d'algèbre linéaire suivant : une
application linéaire $V \to W$ d'espaces vectoriels sur $F$ est injective
si et seulement si $V \otimes_F F' \to W \otimes_F F'$ est injective.

\medskip\noindent
Axiome (A8). Il résulte du fait d'algèbre linéaire utilisé dans
la démonstration de l'axiome (A1).

\medskip\noindent
Axiome (A9). Soit $X$ un schéma projectif lisse non vide sur $k$,
équidimensionnel de dimension $d$. Soit $i : Y \to X$ un diviseur de Cartier
effectif non vide, lisse sur $k$.
Soient $\lambda_Y$ et $\lambda_X$ comme dans l'axiome (A6) pour $X$ et $Y$.
Il faut montrer que, pour $a \in H^{2d - 2}(X)(d - 1)$,
on a $\lambda_Y(i^*(a)) = \lambda_X(a \cup c_1^H(\mathcal{O}_X(Y))$.
D'après la remarque \ref{weil-remark-trace},
on sait que $\lambda_X : H^{2d}(X)(d) \to F$ et
$\lambda_Y : H^{2d - 2}(Y)(d - 1)$ sont uniquement
déterminés par la condition de l'axiome (A6).
Cela étant, il résulte de la démonstration de l'axiome (A6) pour $H^*$ ci-dessus
que $\lambda_X \otimes \text{id}_{F'}$ correspond à $\lambda_{X_{k'}}$
par l'identification donnée $H^{2d}(X)(d) \otimes_F F' = H^{2d}(X_{k'})(d)$.
Ainsi, le fait que l'axiome (A9) soit vérifié pour $F'(1), (H')^*, c_1^{H'}$
entraîne qu'il l'est pour $F(1), H^*, c_1^H$, par une chasse au diagramme.
Ceci achève la démonstration du lemme.
\end{proof}



















% OK, start here.
%

\chapter{Modules adéquats}



\phantomsection
\label{adequate-section-phantom}


\section{Introduction}
\label{adequate-section-introduction}

\noindent
Pour tout schéma $X$, la catégorie $\QCoh(\mathcal{O}_X)$
des modules quasi-cohérents est abélienne et est une sous-catégorie
de Serre faible de la catégorie abélienne de tous les $\mathcal{O}_X$-modules.
Il en va de même pour la catégorie des modules quasi-cohérents sur
un espace algébrique $X$, considérée comme sous-catégorie de la catégorie
de tous les $\mathcal{O}_X$-modules sur le petit site étale de $X$.
De plus, pour un morphisme quasi-compact et quasi-séparé
$f : X \to Y$, l'image directe $f_*$ et les images directes supérieures
préservent la quasi-cohérence.

\medskip\noindent
Soit ensuite $X$ un schéma et soit $\mathcal{O}$ le faisceau structural
sur l'un des grands sites de $X$, par exemple le grand site fppf.
La catégorie des $\mathcal{O}$-modules quasi-cohérents est abélienne
(elle est en fait équivalente à la catégorie usuelle des
$\mathcal{O}_X$-modules quasi-cohérents sur le schéma $X$ mentionnée plus haut),
mais son plongement dans $\textit{Mod}(\mathcal{O})$ n'est pas exact.
Un exemple est l'application de modules quasi-cohérents
$$
\mathcal{O}_{\mathbf{A}^1_k}
\longrightarrow
\mathcal{O}_{\mathbf{A}^1_k}
$$
sur $\mathbf{A}^1_k = \Spec(k[x])$ donnée par la multiplication par $x$.
Dans la catégorie abélienne des faisceaux quasi-cohérents, cette application
est injective, tandis que, dans la catégorie abélienne de tous les
$\mathcal{O}$-modules sur le grand site de $\mathbf{A}^1_k$, elle possède
un noyau non trivial, comme on le voit en l'évaluant sur les sections au-dessus
de $\Spec(k[x]/(x)) = \Spec(k)$. De plus, pour un morphisme
quasi-compact et quasi-séparé $f : X \to Y$, le foncteur $f_{big, *}$
ne préserve pas la quasi-cohérence.

\medskip\noindent
Dans ce chapitre, nous introduisons la catégorie de ce que nous appellerons
les modules adéquats, étroitement apparentés aux modules quasi-cohérents,
qui « résout » les deux problèmes mentionnés ci-dessus. Une autre solution,
que nous mettrons en œuvre lorsque nous traiterons des modules quasi-cohérents
sur les champs algébriques, consiste à considérer les $\mathcal{O}$-modules
qui sont localement quasi-cohérents et satisfont la propriété de changement
de base plat. Voir Cohomologie des champs, section
\ref{stacks-cohomology-section-loc-qcoh-flat-base-change},
Cohomologie des champs, remarque
\ref{stacks-cohomology-remark-bousfield-colocalization}, et
Catégories dérivées des champs, section
\ref{stacks-perfect-section-derived}.






\section{Conventions}
\label{adequate-section-conventions}

\noindent
Dans ce chapitre, nous fixons
$\tau \in \{Zar, \etale, smooth, syntomic, fppf\}$
et un grand $\tau$-site $\Sch_\tau$ comme dans
Topologies, section \ref{topologies-section-procedure}.
Tous les schémas seront des objets de $\Sch_\tau$.
En particulier, pour tout schéma $S$, nous obtenons les sites
$(\textit{Aff}/S)_\tau \subset (\Sch/S)_\tau$.
Le faisceau structural $\mathcal{O}$ sur ces sites est défini par
la règle $\mathcal{O}(T) = \Gamma(T, \mathcal{O}_T)$.

\medskip\noindent
Tous les anneaux $A$ seront tels que $\Spec(A)$ soit (isomorphe à) un
objet de $\Sch_\tau$. Étant donné un anneau $A$, nous notons
$\textit{Alg}_A$ la catégorie des $A$-algèbres dont les objets sont les
$A$-algèbres $B$ de la forme $B = \Gamma(U, \mathcal{O}_U)$,
où $U$ est un objet affine de $\Sch_\tau$. Ainsi, pour tout
schéma affine $S = \Spec(A)$, le foncteur
$$
(\textit{Aff}/S)_\tau \longrightarrow \textit{Alg}_A,
\quad
U \longmapsto \mathcal{O}(U)
$$
est une équivalence.





\section{Foncteurs adéquats}
\label{adequate-section-quasi-coherent}

\noindent
Dans cette section, nous étudions un sujet étroitement lié aux
images directes des faisceaux quasi-cohérents. La majeure partie de ce qui suit
est tirée de l'article \cite{Jaffe}.

\begin{definition}
\label{adequate-definition-module-valued-functor}
Soit $A$ un anneau. Un {\it foncteur à valeurs dans les modules} est un foncteur
$F : \textit{Alg}_A \to \textit{Ab}$ tel que
\begin{enumerate}
\item pour tout objet $B$ de $\textit{Alg}_A$, le groupe
$F(B)$ soit muni d'une structure de $B$-module, et
\item pour tout morphisme $B \to B'$ de $\textit{Alg}_A$, l'application
$F(B) \to F(B')$ soit $B$-linéaire.
\end{enumerate}
Un {\it morphisme de foncteurs à valeurs dans les modules} est une transformation
de foncteurs $\varphi : F \to G$ telle que $F(B) \to G(B)$ soit $B$-linéaire
pour tout $B \in \Ob(\textit{Alg}_A)$.
\end{definition}

\noindent
Soit $S = \Spec(A)$ un schéma affine.
La catégorie des foncteurs à valeurs dans les modules sur $\textit{Alg}_A$
est équivalente à la catégorie
$\textit{PMod}((\textit{Aff}/S)_\tau, \mathcal{O})$
des préfaisceaux de $\mathcal{O}$-modules. L'équivalence est donnée
par la règle qui associe au foncteur à valeurs dans les modules $F$ le
préfaisceau $\mathcal{F}$ défini par
$\mathcal{F}(U) = F(\mathcal{O}(U))$.
Cela résulte immédiatement de l'équivalence
$(\textit{Aff}/S)_\tau \to \textit{Alg}_A$, $U \mapsto \mathcal{O}(U)$
donnée dans la section \ref{adequate-section-conventions}.
Le quasi-inverse pose $F(B) = \mathcal{F}(\Spec(B))$.

\medskip\noindent
Un cas particulier important de foncteur à valeurs dans les modules
s'obtient comme suit. Soit $M$ un $A$-module. Nous noterons alors
$\underline{M}$ le foncteur à valeurs dans les modules $B \mapsto M \otimes_A B$
(muni de la structure évidente de $B$-module). Remarquons que si $M \to N$
est une application de $A$-modules, il lui est associé un morphisme
$\underline{M} \to \underline{N}$ de foncteurs à valeurs dans les modules.
Réciproquement, tout morphisme de foncteurs à valeurs dans les modules
$\underline{M} \to \underline{N}$ provient d'une application de $A$-modules
$M \to N$, comme on le voit en l'évaluant en $B = A$. Autrement dit,
$\text{Mod}_A$ est une sous-catégorie pleine de la catégorie des foncteurs
à valeurs dans les modules sur $\textit{Alg}_A$.

\medskip\noindent
Étant donnée une application de $A$-modules $\varphi : M \to N$, on a
$\Coker(\underline{M} \to \underline{N}) =
\underline{Q}$ où $Q = \Coker(M \to N)$, car $\otimes$
est exact à droite. Mais il n'en va pas ainsi
du noyau en général : par exemple, une application injective de
$A$-modules ne reste pas nécessairement injective après changement de base.
La définition suivante est donc naturelle.

\begin{definition}
\label{adequate-definition-adequate-functor}
Soit $A$ un anneau. Un foncteur à valeurs dans les modules $F$ sur
$\textit{Alg}_A$ est dit
\begin{enumerate}
\item {\it adéquat} s'il existe une
application de $A$-modules $M \to N$ telle que $F$ soit isomorphe à
$\Ker(\underline{M} \to \underline{N})$ ;
\item {\it linéairement adéquat} si $F$ est isomorphe au
noyau d'une application $\underline{A^{\oplus n}} \to \underline{A^{\oplus m}}$.
\end{enumerate}
\end{definition}

\noindent
Remarquons que $F$ est adéquat si et seulement s'il existe une
suite exacte $0 \to F \to \underline{M} \to \underline{N}$, et que
$F$ est linéairement adéquat si et seulement s'il existe une suite exacte
$0 \to F \to \underline{A^{\oplus n}} \to \underline{A^{\oplus m}}$.

\medskip\noindent
Soit $A$ un anneau. Dans cette section, nous montrerons que la catégorie
des foncteurs adéquats sur $\textit{Alg}_A$ est abélienne
(lemmes \ref{adequate-lemma-cokernel-adequate} et \ref{adequate-lemma-kernel-adequate})
et possède un ensemble de générateurs
(lemme \ref{adequate-lemma-adequate-surjection-from-linear}).
Nous verrons aussi qu'elle est une sous-catégorie de Serre faible de la catégorie
de tous les foncteurs à valeurs dans les modules sur $\textit{Alg}_A$
(lemme \ref{adequate-lemma-extension-adequate})
et qu'elle admet des limites inductives arbitraires
(lemme \ref{adequate-lemma-colimit-adequate}).

\begin{lemma}
\label{adequate-lemma-adequate-finite-presentation}
Soit $A$ un anneau.
Soit $F$ un foncteur adéquat sur $\textit{Alg}_A$.
Si $B = \colim B_i$ est une limite inductive filtrante de
$A$-algèbres, alors $F(B) = \colim F(B_i)$.
\end{lemma}

\begin{proof}
Cela tient à ce que, pour tout $A$-module $M$, on a
$M \otimes_A B = \colim M \otimes_A B_i$ (voir
Algèbre, lemme \ref{algebra-lemma-tensor-products-commute-with-limits})
et à ce que les limites inductives filtrantes commutent aux suites exactes, voir
Algèbre, lemme \ref{algebra-lemma-directed-colimit-exact}.
\end{proof}

\begin{remark}
\label{adequate-remark-settheoretic}
Considérons la catégorie $\textit{Alg}_{fp, A}$ dont les objets sont les
$A$-algèbres $B$ de la forme $B = A[x_1, \ldots, x_n]/(f_1, \ldots, f_m)$
et dont les morphismes sont les applications de $A$-algèbres. Toute
$A$-algèbre $B$ est limite inductive filtrante de $A$-algèbres de présentation
finie, c'est-à-dire limite inductive filtrante d'objets de
$\textit{Alg}_{fp, A}$. D'après le
lemme \ref{adequate-lemma-adequate-finite-presentation},
nous en concluons que tout foncteur adéquat $F$ est déterminé par sa restriction
à $\textit{Alg}_{fp, A}$. Pour certaines questions, nous pouvons donc nous
restreindre aux foncteurs sur $\textit{Alg}_{fp, A}$. Par exemple, la catégorie
des foncteurs adéquats ne dépend pas du choix du grand $\tau$-site
fait à la
section \ref{adequate-section-conventions}.
\end{remark}

\begin{lemma}
\label{adequate-lemma-adequate-flat}
Soit $A$ un anneau.
Soit $F$ un foncteur adéquat sur $\textit{Alg}_A$.
Si $B \to B'$ est plat, alors $F(B) \otimes_B B' \to F(B')$
est un isomorphisme.
\end{lemma}

\begin{proof}
Choisissons une suite exacte $0 \to F \to \underline{M} \to \underline{N}$.
Elle donne le diagramme
$$
\xymatrix{
0 \ar[r] & F(B) \otimes_B B' \ar[r] \ar[d] &
(M \otimes_A B)\otimes_B B' \ar[r] \ar[d] &
(N \otimes_A B)\otimes_B B' \ar[d] \\
0 \ar[r] & F(B') \ar[r] &
M \otimes_A B' \ar[r] &
N \otimes_A B'
}
$$
dont les lignes sont exactes (celle du haut parce que $B \to B'$ est plat).
Puisque les deux flèches verticales de droite sont des isomorphismes, celle
de gauche en est également un.
\end{proof}

\begin{lemma}
\label{adequate-lemma-adequate-surjection-from-linear}
Soit $A$ un anneau.
Soit $F$ un foncteur adéquat sur $\textit{Alg}_A$. Il existe alors une
surjection $L \to F$, où $L$ est une somme directe de foncteurs linéairement adéquats.
\end{lemma}

\begin{proof}
Choisissons une suite exacte $0 \to F \to \underline{M} \to \underline{N}$,
où $\underline{M} \to \underline{N}$ est donnée par
$\varphi : M \to N$. D'après le
lemme \ref{adequate-lemma-adequate-finite-presentation},
il suffit de construire $L \to F$ de sorte que $L(B) \to F(B)$ soit surjective
pour toute $A$-algèbre de présentation finie $B$. Il suffit donc, étant donnés
une $A$-algèbre de présentation finie $B$ et un élément $\xi \in F(B)$,
de construire une application $L \to F$, avec $L$ linéairement adéquat, telle que $\xi$
appartienne à l'image de $L(B) \to F(B)$.
(En effet, les classes d'isomorphisme de tels couples $(B, \xi)$ forment un ensemble.)

\medskip\noindent
Pour cela, notons $\sum_{i = 1, \ldots, n} m_i \otimes b_i$ l'image de
$\xi$ dans $\underline{M}(B) = M \otimes_A B$. Nous savons que
$\sum \varphi(m_i) \otimes b_i = 0$ dans $N \otimes_A B$.
Comme $N$ est une limite inductive filtrante de $A$-modules de présentation finie,
on peut trouver un $A$-module de présentation finie $N'$ et un diagramme commutatif
de $A$-modules
$$
\xymatrix{
A^{\oplus n} \ar[r] \ar[d]_{m_1, \ldots, m_n} & N' \ar[d] \\
M \ar[r] & N
}
$$
tels que $(b_1, \ldots, b_n)$ ait une image nulle dans $N' \otimes_A B$.
Choisissons une présentation $A^{\oplus l} \to A^{\oplus k} \to N' \to 0$.
Choisissons un relèvement $A^{\oplus n} \to A^{\oplus k}$ de l'application
$A^{\oplus n} \to N'$ du diagramme. On voit alors qu'il existe
$(c_1, \ldots, c_l) \in B^{\oplus l}$ tel que
$(b_1, \ldots, b_n, c_1, \ldots, c_l)$ ait une image nulle dans $B^{\oplus k}$
par l'application $B^{\oplus n} \oplus B^{\oplus l} \to B^{\oplus k}$.
Considérons le diagramme commutatif
$$
\xymatrix{
A^{\oplus n} \oplus A^{\oplus l} \ar[r] \ar[d] & A^{\oplus k} \ar[d] \\
M \ar[r] & N
}
$$
où la flèche verticale de gauche est nulle sur le facteur direct $A^{\oplus l}$.
On voit alors que le noyau $L$ de $\underline{A^{\oplus n + l}}
\to \underline{A^{\oplus k}}$ convient, car l'élément
$(b_1, \ldots, b_n, c_1, \ldots, c_l) \in L(B)$ a pour image $\xi$.
\end{proof}

\noindent
Considérons une $A$-algèbre graduée $B = \bigoplus_{d \geq 0} B_d$. Il existe alors
deux homomorphismes de $A$-algèbres $p, a : B \to B[t, t^{-1}]$, à savoir
$p : b \mapsto b$ et $a : b \mapsto t^{\deg(b)} b$, où $b$ est homogène.
Si $F$ est un foncteur à valeurs dans les modules sur $\textit{Alg}_A$, posons
\begin{equation}
\label{adequate-equation-weight-k}
F(B)^{(k)} = \{\xi \in F(B) \mid t^k F(p)(\xi) = F(a)(\xi)\}.
\end{equation}
Pour les foncteurs qui se comportent bien relativement aux extensions plates
d'anneaux, cela donne une décomposition en somme directe. Cela revient au fait
que les représentations de $\mathbf{G}_m$ sont complètement réductibles.

\begin{lemma}
\label{adequate-lemma-flat-functor-split}
Soit $A$ un anneau.
Soit $F$ un foncteur à valeurs dans les modules sur $\textit{Alg}_A$.
Supposons que, pour $B \to B'$ plat, l'application
$F(B) \otimes_B B' \to F(B')$ soit un isomorphisme.
Soit $B$ une $A$-algèbre graduée. Alors
\begin{enumerate}
\item $F(B) = \bigoplus_{k \in \mathbf{Z}} F(B)^{(k)}$, et
\item l'application $B \to B_0 \to B$ induit une application $F(B) \to F(B)$
dont l'image est contenue dans $F(B)^{(0)}$.
\end{enumerate}
\end{lemma}

\begin{proof}
Soit $x \in F(B)$. L'application $p : B \to B[t, t^{-1}]$ est libre,
donc nous savons que
$$
F(B[t, t^{-1}]) =
\bigoplus\nolimits_{k \in \mathbf{Z}} F(p)(F(B)) \cdot t^k =
\bigoplus\nolimits_{k \in \mathbf{Z}} F(B) \cdot t^k
$$
comme indiqué, nous omettons $F(p)$ dans la suite de la démonstration.
Écrivons $F(a)(x) = \sum t^k x_k$, avec $x_k \in F(B)$.
Notons $\epsilon : B[t, t^{-1}] \to B$
l'homomorphisme de $B$-algèbres $t \mapsto 1$. Notons que les composés
$\epsilon \circ p, \epsilon \circ a : B \to B[t, t^{-1}] \to B$ sont
l'identité. On a donc
$$
x = F(\epsilon)(F(a)(x)) = F(\epsilon)(\sum t^k x_k) = \sum x_k.
$$
D'autre part, nous affirmons que $x_k \in F(B)^{(k)}$. Considérons en effet
le diagramme commutatif
$$
\xymatrix{
B \ar[r]_a \ar[d]_{a'} &
B[t, t^{-1}] \ar[d]^f \\
B[s, s^{-1}] \ar[r]^-g &
B[t, s, t^{-1}, s^{-1}]
}
$$
où $a'(b) = s^{\deg(b)}b$, $f(b) = b$, $f(t) = st$,
$g(b) = t^{\deg(b)}b$ et $g(s) = s$. Alors
$$
F(g)(F(a'))(x) = F(g)(\sum s^k x_k) =
\sum s^k F(a)(x_k)
$$
et, en suivant l'autre chemin, on obtient
$$
F(f)(F(a))(x) = F(f)(\sum t^k x_k) = \sum (st)^k x_k.
$$
{\emergencystretch=2em
Puisque $B \to B[s, t, s^{-1}, t^{-1}]$ est libre, on a
$F(B[t, s, t^{-1}, s^{-1}]) =
\bigoplus_{k, l \in \mathbf{Z}} F(B) \cdot t^ks^l$ ;
en comparant les coefficients dans les expressions ci-dessus, on trouve
$F(a)(x_k) = t^k x_k$, comme voulu.\par}

\medskip\noindent
Enfin, l'image de $F(B_0) \to F(B)$ est contenue dans $F(B)^{(0)}$
car $B_0 \to B \xrightarrow{a} B[t, t^{-1}]$ est égal à
$B_0 \to B \xrightarrow{p} B[t, t^{-1}]$.
\end{proof}

\noindent
Comme cas particulier du
lemme \ref{adequate-lemma-flat-functor-split},
notons que
$$
\underline{M}(B)^{(k)} = M \otimes_A B_k
$$
où $B_k$ est la composante de degré $k$ de la $A$-algèbre graduée $B$.

\begin{lemma}
\label{adequate-lemma-lift-map}
Soit $A$ un anneau. Étant donné un diagramme en traits pleins
$$
\xymatrix{
0 \ar[r] &
L \ar[d]_\varphi \ar[r] &
\underline{A^{\oplus n}} \ar[r] \ar@{..>}[ld] &
\underline{A^{\oplus m}} \\
& \underline{M}
}
$$
de foncteurs à valeurs dans les modules sur $\textit{Alg}_A$
dont la ligne est exacte, il existe une flèche pointillée rendant le diagramme commutatif.
\end{lemma}

\begin{proof}
Supposons que l'application $A^{\oplus n} \to A^{\oplus m}$ soit donnée par la
matrice de taille $m \times n$ notée $(a_{ij})$. Considérons l'anneau
$B = A[x_1, \ldots, x_n]/(\sum a_{ij}x_j)$. L'élément
$(x_1, \ldots, x_n) \in \underline{A^{\oplus n}}(B)$ a une image nulle dans
$\underline{A^{\oplus m}}(B)$ ; il est donc l'image d'un unique élément
$\xi \in L(B)$. Notons que $\xi$ possède la propriété universelle suivante :
pour toute $A$-algèbre $C$ et tout $\xi' \in L(C)$, il existe un homomorphisme de
$A$-algèbres $B \to C$ tel que $\xi$ ait pour image $\xi'$ par $L(B) \to L(C)$.

\medskip\noindent
Notons que $B$ est une $A$-algèbre graduée ; les
lemmes \ref{adequate-lemma-flat-functor-split} et \ref{adequate-lemma-adequate-flat}
permettent donc de décomposer les valeurs de nos foncteurs en $B$ en composantes graduées.
Notons que $\xi \in L(B)^{(1)}$, puisque $(x_1, \ldots, x_n)$ est un élément
de degré un de $\underline{A^{\oplus n}}(B)$. On voit donc que
$\varphi(\xi) \in \underline{M}(B)^{(1)} = M \otimes_A B_1$.
Puisque $B_1$ est engendré par $x_1, \ldots, x_n$ comme $A$-module,
on peut écrire $\varphi(\xi) = \sum m_i \otimes x_i$. Considérons l'application
$A^{\oplus n} \to M$ qui envoie le $i$-ième vecteur de base sur $m_i$.
Par construction, l'application associée
$\underline{A^{\oplus n}} \to \underline{M}$
envoie l'élément $\xi$ sur $\varphi(\xi)$. La propriété
universelle mentionnée ci-dessus entraîne que le diagramme commute.
\end{proof}

\begin{lemma}
\label{adequate-lemma-cokernel-into-module}
Soit $A$ un anneau.
Soit $\varphi : F \to \underline{M}$ une application de foncteurs à valeurs dans les modules
sur $\textit{Alg}_A$, avec $F$ adéquat.
Alors $\Coker(\varphi)$ est adéquat.
\end{lemma}

\begin{proof}
D'après le
lemme \ref{adequate-lemma-adequate-surjection-from-linear},
on peut supposer que $F = \bigoplus L_i$ est une somme directe de foncteurs
linéairement adéquats. Choisissons des suites exactes
$0 \to L_i \to \underline{A^{\oplus n_i}} \to \underline{A^{\oplus m_i}}$.
Pour chaque $i$, choisissons une application $A^{\oplus n_i} \to M$ comme dans le
lemme \ref{adequate-lemma-lift-map}.
Considérons le diagramme
$$
\xymatrix{
0 \ar[r] &
\bigoplus L_i  \ar[r] \ar[d] &
\bigoplus \underline{A^{\oplus n_i}} \ar[r] \ar[ld] &
\bigoplus \underline{A^{\oplus m_i}} \\
& \underline{M}
}
$$
Considérons les $A$-modules
$$
Q =
\Coker(\bigoplus A^{\oplus n_i} \to M \oplus \bigoplus A^{\oplus m_i})
\quad\text{et}\quad
P = \Coker(\bigoplus A^{\oplus n_i} \to \bigoplus A^{\oplus m_i}).
$$
On voit alors que $\Coker(\varphi)$ est isomorphe au
noyau de $\underline{Q} \to \underline{P}$.
\end{proof}

\begin{lemma}
\label{adequate-lemma-cokernel-adequate}
\begin{slogan}
Le conoyau d'une application de foncteurs adéquats sur la catégorie des algèbres
sur un anneau est adéquat.
\end{slogan}
Soit $A$ un anneau.
Soit $\varphi : F \to G$ une application de foncteurs adéquats sur $\textit{Alg}_A$.
Alors $\Coker(\varphi)$ est adéquat.
\end{lemma}

\begin{proof}
Choisissons une injection $G \to \underline{M}$.
On a alors une injection $G/F \to \underline{M}/F$. D'après le
lemme \ref{adequate-lemma-cokernel-into-module},
$\underline{M}/F$ est adéquat ; on peut donc trouver une injection
$\underline{M}/F \to \underline{N}$.
Par composition, on obtient une injection $G/F \to \underline{N}$. D'après le
lemme \ref{adequate-lemma-cokernel-into-module},
le conoyau de l'application induite $G \to \underline{N}$ est adéquat ;
on peut donc trouver une injection $\underline{N}/G \to \underline{K}$.
La suite $0 \to G/F \to \underline{N} \to \underline{K}$ est alors exacte,
ce qui conclut.
\end{proof}

\begin{lemma}
\label{adequate-lemma-kernel-adequate}
Soit $A$ un anneau.
Soit $\varphi : F \to G$ une application de foncteurs adéquats sur $\textit{Alg}_A$.
Alors $\Ker(\varphi)$ est adéquat.
\end{lemma}

\begin{proof}
Choisissons une injection $F \to \underline{M}$ et une injection
$G \to \underline{N}$. Notons $F \to \underline{M \oplus N}$
l'application diagonale, de sorte que
$$
\xymatrix{
F \ar[d] \ar[r] & G \ar[d] \\
\underline{M \oplus N} \ar[r] & \underline{N}
}
$$
commute. D'après le
lemme \ref{adequate-lemma-cokernel-adequate},
on peut trouver une application de modules $M \oplus N \to K$ telle que
$F$ soit le noyau de $\underline{M \oplus N} \to \underline{K}$.
Alors $\Ker(\varphi)$ est le noyau de
$\underline{M \oplus N} \to \underline{K \oplus N}$.
\end{proof}

\begin{lemma}
\label{adequate-lemma-colimit-adequate}
Soit $A$ un anneau.
Une somme directe arbitraire de foncteurs adéquats sur $\textit{Alg}_A$
est adéquate. Une limite inductive de foncteurs adéquats est adéquate.
\end{lemma}

\begin{proof}
L'assertion relative aux sommes directes est immédiate.
Une limite inductive quelconque peut s'écrire comme noyau d'une application entre
sommes directes ; voir
Catégories, lemme \ref{categories-lemma-colimits-coproducts-coequalizers}.
L'assertion résulte donc du
lemme \ref{adequate-lemma-kernel-adequate}.
\end{proof}

\begin{lemma}
\label{adequate-lemma-flat-linear-functor}
Soit $A$ un anneau.
Soient $F, G$ des foncteurs à valeurs dans les modules sur $\textit{Alg}_A$.
Soit $\varphi : F \to G$ une transformation de foncteurs. Supposons que
\begin{enumerate}
\item $\varphi$ soit additive,
\item pour toute $A$-algèbre $B$, tout $\xi \in F(B)$ et toute unité
$u \in B^*$, on ait $\varphi(u\xi) = u\varphi(\xi)$ dans $G(B)$, et
\item pour tout homomorphisme plat d'anneaux $B \to B'$, on ait
$G(B) \otimes_B B' = G(B')$.
\end{enumerate}
Alors $\varphi$ est un morphisme de foncteurs à valeurs dans les modules.
\end{lemma}

\begin{proof}
Soient $B$ une $A$-algèbre, $\xi \in F(B)$ et $b \in B$. Il faut montrer
que $\varphi(b \xi) = b \varphi(\xi)$. Considérons l'homomorphisme d'anneaux
$$
B \to B' = B[x, y, x^{-1}, y^{-1}]/(x + y - b).
$$
Cet homomorphisme est fidèlement plat, donc $G(B) \subset G(B')$. D'autre
part,
$$
\varphi(b\xi) = \varphi((x + y)\xi) =
\varphi(x\xi) + \varphi(y\xi) = x\varphi(\xi) + y\varphi(\xi)
= (x + y)\varphi(\xi) = b\varphi(\xi)
$$
car $x, y$ sont des unités de $B'$. Cela conclut.
\end{proof}

\begin{lemma}
\label{adequate-lemma-extension-adequate-key}
Soit $A$ un anneau.
Soit $0 \to \underline{M} \to G \to L \to 0$ une suite exacte courte
de foncteurs à valeurs dans les modules sur $\textit{Alg}_A$, avec $L$ linéairement adéquat.
Alors $G$ est adéquat.
\end{lemma}

\begin{proof}
Remarquons d'abord que, pour tout homomorphisme plat de $A$-algèbres
$B \to B'$, l'application $G(B) \otimes_B B' \to G(B')$ est un isomorphisme.
Cela vaut en effet pour $\underline{M}$ et $L$, voir le
lemme \ref{adequate-lemma-adequate-flat},
et en résulte pour $G$ par le lemme des cinq. En particulier, le
lemme \ref{adequate-lemma-flat-functor-split}
montre que $G(B) = \bigoplus_{k \in \mathbf{Z}} G(B)^{(k)}$
pour toute $A$-algèbre graduée $B$.

\medskip\noindent
Choisissons une suite exacte
$0 \to L \to \underline{A^{\oplus n}} \to \underline{A^{\oplus m}}$.
Supposons que l'application $A^{\oplus n} \to A^{\oplus m}$ soit donnée par la
matrice de taille $m \times n$ notée $(a_{ij})$. Considérons la $A$-algèbre graduée
$B = A[x_1, \ldots, x_n]/(\sum a_{ij}x_j)$. L'élément
$(x_1, \ldots, x_n) \in \underline{A^{\oplus n}}(B)$ a une image nulle dans
$\underline{A^{\oplus m}}(B)$ ; il est donc l'image d'un unique élément
$\xi \in L(B)$. Observons que $\xi \in L(B)^{(1)}$. L'application
$$
\Hom_A(B, C) \longrightarrow L(C), \quad
f \longmapsto L(f)(\xi)
$$
définit un isomorphisme de foncteurs. En effet, $f$ est
déterminé par les images $c_i = f(x_i) \in C$, qui doivent
satisfaire aux relations $\sum a_{ij}c_j = 0$. Or $L(C)$ est
l'ensemble des $n$-uplets $(c_1, \ldots, c_n)$ satisfaisant aux relations
$\sum a_{ij} c_j = 0$.

\medskip\noindent
Puisque la valeur de chacun des foncteurs $\underline{M}$, $G$, $L$
en $B$ est la somme directe de ses espaces de poids (d'après le lemme
mentionné ci-dessus), l'exactitude de $0 \to \underline{M} \to G \to L \to 0$ entraîne
celle de la suite $0 \to \underline{M}(B)^{(1)} \to G(B)^{(1)} \to L(B)^{(1)} \to 0$.
On peut donc choisir un élément $\theta \in G(B)^{(1)}$ ayant pour image
$\xi$.

\medskip\noindent
Considérons la $A$-algèbre graduée
$$
C = A[x_1, \ldots, x_n, y_1, \ldots, y_n]/
(\sum a_{ij}x_j, \sum a_{ij}y_j)
$$
Il existe trois homomorphismes de $A$-algèbres graduées $p_1, p_2, m : B \to C$
définis par les règles
$$
p_1(x_i) = x_i, \quad
p_1(x_i) = y_i, \quad
m(x_i) = x_i + y_i.
$$
Nous montrerons que l'élément
$$
\tau = G(m)(\theta) - G(p_1)(\theta) - G(p_2)(\theta) \in G(C)
$$
est nul. Tout d'abord, un calcul direct montre que $\tau$ a une image nulle dans $L(C)$.
Ainsi, $\tau$ est un élément de $\underline{M}(C)$.
De plus, puisque $m$, $p_1$, $p_2$ sont des homomorphismes d'algèbres graduées,
on voit que $\tau \in G(C)^{(1)}$ ; comme $\underline{M} \subset G$,
on en conclut que
$$
\tau \in \underline{M}(C)^{(1)} = M \otimes_A C_1.
$$
On peut écrire de manière unique
$\tau = \underline{M}(p_1)(\tau_1) + \underline{M}(p_2)(\tau_2)$, avec
$\tau_i \in M \otimes_A B_1 = \underline{M}(B)^{(1)}$, car
$C_1 = p_1(B_1) \oplus p_2(B_1)$.
Considérons l'homomorphisme d'anneaux $q_1 : C \to B$ défini par $x_i \mapsto x_i$ et
$y_i \mapsto 0$. Alors
$\underline{M}(q_1)(\tau) =
\underline{M}(q_1)(\underline{M}(p_1)(\tau_1) + \underline{M}(p_2)(\tau_2)) =
\tau_1$.
D'autre part, puisque
$q_1 \circ m = q_1 \circ p_1$, on voit que
$G(q_1)(\tau) = - G(q_1 \circ p_2)(\tau)$. Comme $q_1 \circ p_2$ se factorise en
$B \to A \to B$, on voit que $G(q_1 \circ p_2)(\tau)$ appartient à
$G(B)^{(0)}$ ; voir le
lemme \ref{adequate-lemma-flat-functor-split}.
Ainsi $\tau_1 = 0$, car cet élément appartient à
$G(B)^{(0)} \cap \underline{M}(B)^{(1)} \subset
G(B)^{(0)} \cap G(B)^{(1)} = 0$.
De même $\tau_2 = 0$, d'où $\tau = 0$.

\medskip\noindent
Puisque $\theta \in G(B)$, on obtient une transformation de foncteurs
$$
\psi : L(-) = \Hom_A(B, - ) \longrightarrow G(-)
$$
en envoyant $f : B \to C$ sur $G(f)(\theta)$. Puisque $\theta$ est un relèvement de
$\xi$, l'application $\psi$ est un inverse à droite de $G \to L$. En termes de
$\psi$, les assertions démontrées ci-dessus signifient ceci :
$\tau = 0$ signifie que $\psi$ est additive, et
$\theta \in G(B)^{(1)}$ entraîne que, pour toute $A$-algèbre $D$, on a
$\psi(ul) = u\psi(l)$ dans $G(D)$ pour $l \in L(D)$ et toute unité $u \in D^*$.
Cela entraîne que $\psi$ est un morphisme de foncteurs à valeurs dans les modules ; voir le
lemme \ref{adequate-lemma-flat-linear-functor}.
Il en résulte clairement que $G \cong \underline{M} \oplus L$, ce qui conclut.
\end{proof}

\begin{remark}
\label{adequate-remark-linearly-adequate}
Soit $A$ un anneau.
La démonstration du
lemme \ref{adequate-lemma-extension-adequate-key}
montre que toute extension $0 \to \underline{M} \to E \to L \to 0$
de foncteurs à valeurs dans les modules sur $\textit{Alg}_A$,
avec $L$ linéairement adéquat, est scindée. Elle n'utilise que les propriétés suivantes
du foncteur à valeurs dans les modules $F = \underline{M}$ :
\begin{enumerate}
\item $F(B) \otimes_B B' \to F(B')$ est un isomorphisme
pour un homomorphisme plat d'anneaux $B \to B'$, et
\item
$F(C)^{(1)} = F(p_1)(F(B)^{(1)}) \oplus F(p_2)(F(B)^{(1)})$,
où $B = A[x_1, \ldots, x_n]/(\sum a_{ij}x_j)$ et
$C = A[x_1, \ldots, x_n, y_1, \ldots, y_n]/
(\sum a_{ij}x_j, \sum a_{ij}y_j)$.
\end{enumerate}
Ces deux propriétés valent pour tout foncteur adéquat $F$ ; nous omettons les détails.
On voit donc que $L$ est un objet projectif de la catégorie abélienne des
foncteurs adéquats.
\end{remark}

\begin{lemma}
\label{adequate-lemma-extension-adequate}
Soit $A$ un anneau.
Soit $0 \to F \to G \to H \to 0$ une suite exacte courte de
foncteurs à valeurs dans les modules sur $\textit{Alg}_A$.
Si $F$ et $H$ sont adéquats, $G$ l'est aussi.
\end{lemma}

\begin{proof}
Choisissons une suite exacte $0 \to F \to \underline{M} \to \underline{N}$.
Si l'on peut montrer que $(\underline{M} \oplus G)/F$ est adéquat, alors
$G$ est le noyau de l'application de foncteurs adéquats
$(\underline{M} \oplus G)/F \to \underline{N}$, donc est
adéquat d'après le
lemme \ref{adequate-lemma-kernel-adequate}.
On peut donc supposer que $F = \underline{M}$.

\medskip\noindent
On peut choisir une surjection $L \to H$, où $L$ est une somme directe de
foncteurs linéairement adéquats ; voir le
lemme \ref{adequate-lemma-adequate-surjection-from-linear}.
Si l'on peut montrer que le produit fibré $G \times_H L$ est adéquat, alors
$G$ est le conoyau de l'application $\Ker(L \to H) \to G \times_H L$,
donc est adéquat d'après le
lemme \ref{adequate-lemma-cokernel-adequate}.
On peut donc supposer que $H = \bigoplus L_i$ est une somme directe de
foncteurs linéairement adéquats. D'après le
lemme \ref{adequate-lemma-extension-adequate-key},
chacun des produits fibrés $G \times_H L_i$ est adéquat. D'après le
lemme \ref{adequate-lemma-colimit-adequate},
on voit que $\bigoplus G \times_H L_i$ est adéquat.
Alors $G$ est le conoyau de
$$
\bigoplus\nolimits_{i \not = i'} F \longrightarrow
\bigoplus G \times_H L_i
$$
où $\xi$, dans le facteur direct $(i, i')$, a pour image
$(0, \ldots, 0, \xi, 0, \ldots, 0, -\xi, 0, \ldots, 0)$,
avec des composantes non nulles dans les facteurs directs $i$ et $i'$.
Ainsi $G$ est adéquat d'après le
lemme \ref{adequate-lemma-cokernel-adequate}.
\end{proof}

\begin{lemma}
\label{adequate-lemma-base-change-adequate}
Soit $A \to A'$ un homomorphisme d'anneaux. Si $F$ est un foncteur adéquat sur
$\textit{Alg}_A$, sa restriction $F'$ à
$\textit{Alg}_{A'}$ est également adéquate.
\end{lemma}

\begin{proof}
Choisissons une suite exacte $0 \to F \to \underline{M} \to \underline{N}$.
Alors $F'(B') = F(B') = \Ker(M \otimes_A B' \to N \otimes_A B')$.
Puisque $M \otimes_A B' = M \otimes_A A' \otimes_{A'} B'$, et de même
pour $N$, on voit que $F'$ est le noyau de
$\underline{M \otimes_A A'} \to \underline{N \otimes_A A'}$.
\end{proof}

\begin{lemma}
\label{adequate-lemma-pushforward-adequate}
Soit $A \to A'$ un homomorphisme d'anneaux. Si $F'$ est un foncteur adéquat sur
$\textit{Alg}_{A'}$, le foncteur à valeurs dans les modules
$F : B \mapsto F'(A' \otimes_A B)$ sur $\textit{Alg}_A$ est également adéquat.
\end{lemma}

\begin{proof}
Choisissons une suite exacte $0 \to F' \to \underline{M'} \to \underline{N'}$.
Alors
\begin{align*}
F(B) & = F'(A' \otimes_A B) \\
& = \Ker(M' \otimes_{A'} (
A' \otimes_A B) \to N' \otimes_{A'} (A' \otimes_A B)) \\
& = \Ker(M' \otimes_A B \to N' \otimes_A B)
\end{align*}
Ainsi $F$ est le noyau de
$\underline{M} \to \underline{N}$,
où $M = M'$ et $N = N'$ sont considérés comme $A$-modules.
\end{proof}

\begin{lemma}
\label{adequate-lemma-adequate-product}
Soit $A = A_1 \times \ldots \times A_n$ un produit d'anneaux.
Se donner un foncteur adéquat sur $A$ revient à se donner une suite
$F_1, \ldots, F_n$ de foncteurs adéquats $F_i$ sur $A_i$.
\end{lemma}

\begin{proof}
Cela résulte de ce qu'une $A$-algèbre $B$ est canoniquement un produit
$B_1 \times \ldots \times B_n$, et qu'il en va de même des $A$-modules.
La correspondance est donnée par $F(B) = \coprod F_i(B_i)$.
Nous omettons les détails.
\end{proof}

\begin{lemma}
\label{adequate-lemma-adequate-descent}
Soient $A \to A'$ un homomorphisme d'anneaux et $F$ un foncteur à valeurs dans les modules sur
$\textit{Alg}_A$ tel que
\begin{enumerate}
\item la restriction $F'$ de $F$ à la catégorie des $A'$-algèbres soit
adéquate, et
\item pour toute $A$-algèbre $B$, la suite
$$
0 \to F(B) \to F(B \otimes_A A') \to F(B \otimes_A A' \otimes_A A')
$$
soit exacte.
\end{enumerate}
Alors $F$ est adéquat.
\end{lemma}

\begin{proof}
Les foncteurs $B \to F(B \otimes_A A')$ et
$B \mapsto F(B \otimes_A A' \otimes_A A')$ sont adéquats ; voir les
lemmes \ref{adequate-lemma-pushforward-adequate} et
\ref{adequate-lemma-base-change-adequate}.
Ainsi $F$, noyau d'une application de foncteurs adéquats, est adéquat ; voir le
lemme \ref{adequate-lemma-kernel-adequate}.
\end{proof}






\section{Groupes Ext supérieurs des foncteurs adéquats}
\label{adequate-section-higher-ext}

\noindent
Soit $A$ un anneau. Dans le
lemme \ref{adequate-lemma-extension-adequate},
nous avons vu que toute extension de foncteurs adéquats dans la catégorie
des foncteurs à valeurs dans les modules sur $\textit{Alg}_A$ est adéquate. Dans cette
section, nous montrons que la même propriété vaut pour les groupes Ext supérieurs.

\begin{lemma}
\label{adequate-lemma-adjoint}
Soit $A$ un anneau.
Pour tout foncteur à valeurs dans les modules $F$ sur $\textit{Alg}_A$,
il existe un morphisme $Q(F) \to F$ de foncteurs à valeurs dans les modules sur
$\textit{Alg}_A$ tel que (1) $Q(F)$ soit adéquat et (2) pour tout
foncteur adéquat $G$, l'application $\Hom(G, Q(F)) \to \Hom(G, F)$
soit bijective.
\end{lemma}

\begin{proof}
Choisissons un ensemble $\{L_i\}_{i \in I}$ de foncteurs linéairement adéquats tel que
tout foncteur linéairement adéquat soit isomorphe à l'un des $L_i$.
C'est possible. Supposons que l'on puisse trouver $Q(F) \to F$ satisfaisant à (1) et à la condition
(2)' pour tout $i \in I$, l'application $\Hom(L_i, Q(F)) \to \Hom(L_i, F)$
est bijective. Alors (2) est vérifiée. En effet, en combinant les
lemmes \ref{adequate-lemma-adequate-surjection-from-linear} et
\ref{adequate-lemma-kernel-adequate},
on voit que tout foncteur adéquat $G$ s'insère dans une suite exacte
$$
K \to L \to G \to 0
$$
où $K$ et $L$ sont des sommes directes de foncteurs linéairement adéquats. Ainsi (2)'
entraîne que
$\Hom(L, Q(F)) \to \Hom(L, F)$
et
$\Hom(K, Q(F)) \to \Hom(K, F)$
sont bijectives, d'où la même assertion pour $G$.

\medskip\noindent
Considérons la catégorie $\mathcal{I}$ dont les objets sont les couples
$(i, \varphi)$, où $i \in I$ et où $\varphi : L_i \to F$ est un morphisme.
Un morphisme $(i, \varphi) \to (i', \varphi')$ est une application
$\psi : L_i \to L_{i'}$ telle que $\varphi' \circ \psi = \varphi$.
Posons
$$
Q(F) = \colim_{(i, \varphi) \in \Ob(\mathcal{I})} L_i
$$
Il existe une application naturelle $Q(F) \to F$ ; d'après le
lemme \ref{adequate-lemma-colimit-adequate},
ce foncteur est adéquat et, par construction, il possède la propriété (2)'.
\end{proof}

\begin{lemma}
\label{adequate-lemma-enough-injectives}
Soit $A$ un anneau. Notons $\mathcal{P}$ la catégorie des foncteurs à valeurs
dans les modules sur $\textit{Alg}_A$ et $\mathcal{A}$ la catégorie des foncteurs
adéquats sur $\textit{Alg}_A$. Notons $i : \mathcal{A} \to \mathcal{P}$
le foncteur d'inclusion. Notons $Q : \mathcal{P} \to \mathcal{A}$
la construction du lemme \ref{adequate-lemma-adjoint}.
Alors
\begin{enumerate}
\item $i$ est pleinement fidèle, exact, et son image est une sous-catégorie de Serre faible,
\item $\mathcal{P}$ possède suffisamment d'injectifs,
\item le foncteur $Q$ est adjoint à droite de $i$, donc exact à gauche,
\item $Q$ transforme les objets injectifs en objets injectifs,
\item $\mathcal{A}$ possède suffisamment d'injectifs.
\end{enumerate}
\end{lemma}

\begin{proof}
Ce lemme rassemble simplement des faits déjà établis.
L'assertion (1) résulte des définitions, de la caractérisation des
sous-catégories de Serre faibles (voir
Homologie, lemme \ref{homology-lemma-characterize-weak-serre-subcategory})
et des
lemmes \ref{adequate-lemma-cokernel-adequate}, \ref{adequate-lemma-kernel-adequate}
et \ref{adequate-lemma-extension-adequate}.
Rappelons que $\mathcal{P}$ est équivalente à la catégorie
$\textit{PMod}((\textit{Aff}/\Spec(A))_\tau, \mathcal{O})$.
D'où (2), d'après
Injectifs, proposition \ref{injectives-proposition-presheaves-modules}.
L'assertion (3) résulte du
lemme \ref{adequate-lemma-adjoint}
et de
Catégories, lemme \ref{categories-lemma-adjoint-exact}.
Les assertions (4) et (5) résultent de
Homologie, lemmes \ref{homology-lemma-adjoint-preserve-injectives} et
\ref{homology-lemma-adjoint-enough-injectives}.
\end{proof}

\noindent
Soit $A$ un anneau. Comme dans
Théorie des déformations formelles, section
\ref{formal-defos-section-tangent-spaces-functors},
étant donnés une $A$-algèbre $B$ et un $B$-module $N$, nous notons $B[N]$
la $R$-algèbre dont le $B$-module sous-jacent est $B \oplus N$ et dont la multiplication
est donnée par $(b, m)(b', m ') = (bb', bm' + b'm)$. Notons que cette construction
est fonctorielle par rapport au couple $(B, N)$, un morphisme $(B, N) \to (B', N')$
étant donné par un homomorphisme de $A$-algèbres $B \to B'$ et une application de $B$-modules
$N \to N'$. En un certain sens, le foncteur $TF$ sur les couples, défini dans le lemme
suivant, est l'espace tangent de $F$.
Ci-dessous, nous ne considérerons que des couples $(B, N)$ tels que
$B[N]$ soit un objet de $\textit{Alg}_A$.

\begin{lemma}
\label{adequate-lemma-tangent-functor}
Soit $A$ un anneau. Soit $F$ un foncteur à valeurs dans les modules.
Pour tout $B \in \Ob(\textit{Alg}_A)$ et tout $B$-module $N$,
il existe une décomposition canonique
$$
F(B[N]) = F(B) \oplus TF(B, N)
$$
caractérisée par les propriétés suivantes
\begin{enumerate}
\item $TF(B, N) = \Ker(F(B[N]) \to F(B))$,
\item il existe une structure de $B$-module sur $TF(B, N)$
compatible avec la structure de $B[N]$-module sur $F(B[N])$,
\item $TF$ est un foncteur défini sur la catégorie des couples $(B, N)$,
\item
\label{adequate-item-mult-map}
il existe des applications canoniques $N \otimes_B F(B) \to TF(B, N)$
induisant une transformation entre foncteurs définis sur la catégorie
des couples $(B, N)$,
\item $TF(B, 0) = 0$, et l'application $TF(B, N) \to TF(B, N')$ est
nulle lorsque $N \to N'$ est l'application nulle.
\end{enumerate}
\end{lemma}

\begin{proof}
Puisque $B \to B[N] \to B$ est l'identité, on voit que $F(B) \to F(B[N])$
est un facteur direct dont le supplémentaire est $TF(N, B)$, défini en (1).
Cette construction est fonctorielle par rapport au couple $(B, N)$, simplement parce que,
étant donné un morphisme de couples $(B, N) \to (B', N')$, on obtient un diagramme
commutatif
$$
\xymatrix{
B' \ar[r] & B'[N'] \ar[r] & B' \\
B \ar[r] \ar[u] & B[N] \ar[r] \ar[u] & B \ar[u]
}
$$
dans $\textit{Alg}_A$. La structure de $B$-module provient de la structure de $B[N]$-module
et de l'homomorphisme d'anneaux $B \to B[N]$. L'application de (4) est le
composé
$$
N \otimes_B F(B) \longrightarrow
B[N] \otimes_{B[N]} F(B[N]) \longrightarrow F(B[N])
$$
dont l'image est contenue dans $TF(B, N)$. (La première flèche utilise les inclusions
$N \to B[N]$ et $F(B) \to F(B[N])$, et la seconde est l'application
de multiplication.) Si $N = 0$, alors $B = B[N]$,
donc $TF(B, 0) = 0$. Si $N \to N'$ est nulle, elle se factorise en
$N \to 0 \to N'$ ; l'application induite est donc nulle puisque $TF(B, 0) = 0$.
\end{proof}

\noindent
Soit $A$ un anneau. Soit $M$ un $A$-module. Le foncteur à valeurs dans les modules
$\underline{M}$ a pour espace tangent $T\underline{M}$, défini par la règle
$T\underline{M}(B, N) = N \otimes_A M$. En particulier, pour $B$ donné, le
foncteur $N \mapsto T\underline{M}(B, N)$ est additif et exact à droite. Il se trouve
que cela vaut également pour les foncteurs injectifs à valeurs dans les modules.

\begin{lemma}
\label{adequate-lemma-tangent-injective}
Soit $A$ un anneau. Soit $I$ un objet injectif de la catégorie
des foncteurs à valeurs dans les modules. Alors, pour tout $B \in \Ob(\textit{Alg}_A)$
et toute suite exacte courte
$0 \to N_1 \to N \to N_2 \to 0$
de $B$-modules, la suite
$$
TI(B, N_1) \to TI(B, N) \to TI(B, N_2) \to 0
$$
est exacte.
\end{lemma}

\begin{proof}
Nous utiliserons les résultats du
lemme \ref{adequate-lemma-tangent-functor}
sans le mentionner à nouveau.
Notons $h : \textit{Alg}_A \to \textit{Ensembles}$ le foncteur donné par
$h(C) = \Mor_A(B[N], C)$. De même pour $h_1$ et $h_2$.
L'application $B[N] \to B[N_2]$ correspondant à la surjection $N \to N_2$
est surjective. Il lui correspond une application $h_2 \to h$ telle que
$h_2(C) \to h(C)$ soit injective pour toutes les $A$-algèbres $C$. D'autre
part, il existe deux applications $p, q : h \to h_1$, correspondant à
l'application nulle $N_1 \to N$ et à l'injection $N_1 \to N$. Notons que
$$
\xymatrix{
h_2 \ar[r] & h \ar@<1ex>[r] \ar@<-1ex>[r] & h_1
}
$$
est un diagramme de noyau de la double flèche. Notons $\mathcal{O}_h$ le foncteur à valeurs dans les modules
$C \mapsto \bigoplus_{h(C)} C$. De même pour $\mathcal{O}_{h_1}$ et
$\mathcal{O}_{h_2}$. Notons que
$$
\Hom_\mathcal{P}(\mathcal{O}_h, F) = F(B[N])
$$
où $\mathcal{P}$ est la catégorie des foncteurs à valeurs dans les modules sur
$\textit{Alg}_A$. Nous affirmons qu'il existe un diagramme de noyau de la double flèche
$$
\xymatrix{
\mathcal{O}_{h_2} \ar[r] &
\mathcal{O}_h \ar@<1ex>[r] \ar@<-1ex>[r] &
\mathcal{O}_{h_1}
}
$$
dans $\mathcal{P}$. En effet, supposons que $C \in \Ob(\textit{Alg}_A)$
et que $\xi = \sum_{i = 1, \ldots, n} c_i \cdot f_i$, où $c_i \in C$ et
$f_i : B[N] \to C$, soit un élément de
$\mathcal{O}_h(C)$. Si $p(\xi) = q(\xi)$, alors
on voit que
$$
\sum c_i \cdot f_i \circ z = \sum c_i \cdot f_i \circ y
$$
où $z, y : B[N_1] \to B[N]$ sont les applications $z : (b, m_1) \mapsto (b, 0)$
et $y : (b, m_1) \mapsto (b, m_1)$. Cela signifie que, pour tout $i$,
il existe un $j$ tel que $f_j \circ z = f_i \circ y$. Il en résulte clairement
que $f_i(N_1) = 0$, c'est-à-dire que $f_i$ se factorise par une unique application
$\overline{f}_i : B[N_2] \to C$. Ainsi $\xi$ est l'image de
$\overline{\xi} = \sum c_i \cdot \overline{f}_i$.
Puisque $I$ est injectif, il transforme ce diagramme de noyau de la double flèche
en un diagramme de conoyau de la double flèche
$$
\xymatrix{
I(B[N_1]) \ar@<1ex>[r] \ar@<-1ex>[r] &
I(B[N]) \ar[r] &
I(B[N_2])
}
$$
Ce diagramme est compatible avec les décompositions en somme directe
$I(B[N]) = I(B) \oplus TI(B, N)$ et $I(B[N_i]) = I(B) \oplus TI(B, N_i)$.
L'application nulle $N \to N_1$ induit l'application nulle $TI(B, N) \to TI(B, N_1)$.
On voit ainsi que la propriété de conoyau de la double flèche
ci-dessus signifie que l'on a une suite exacte
$TI(B, N_1) \to TI(B, N) \to TI(B, N_2) \to 0$,
comme voulu.
\end{proof}

\begin{lemma}
\label{adequate-lemma-exactness-implies}
Soit $A$ un anneau. Soit $F$ un foncteur à valeurs dans les modules
tel que, pour tout $B \in \Ob(\textit{Alg}_A)$, le
foncteur $TF(B, -)$ sur les $B$-modules transforme une suite exacte courte
de $B$-modules en une suite exacte à droite. Alors
\begin{enumerate}
\item $TF(B, N_1 \oplus N_2) = TF(B, N_1) \oplus TF(B, N_2)$,
\item il existe une seconde structure fonctorielle de $B$-module sur $TF(B, N)$,
définie en posant $x \cdot b = TF(B, b\cdot 1_N)(x)$ pour $x \in TF(B, N)$
et $b \in B$,
\item
\label{adequate-item-mult-map-linear}
l'application canonique $N \otimes_B F(B) \to TF(B, N)$ du
lemme \ref{adequate-lemma-tangent-functor}
est $B$-linéaire également pour la seconde structure de $B$-module,
\item
\label{adequate-item-tangent-right-exact}
étant donné un $B$-module de présentation finie $N$, il existe un
isomorphisme canonique $TF(B, B) \otimes_B N \to TF(B, N)$, où le produit
tensoriel utilise la seconde structure de $B$-module sur $TF(B, B)$.
\end{enumerate}
\end{lemma}

\begin{proof}
Nous utiliserons les résultats du
lemme \ref{adequate-lemma-tangent-functor}
sans le mentionner à nouveau.
Les applications $N_1 \to N_1 \oplus N_2$ et $N_2 \to N_1 \oplus N_2$ donnent
une application $TF(B, N_1) \oplus TF(B, N_2) \to TF(B, N_1 \oplus N_2)$,
qui est injective puisque les applications $N_1 \oplus N_2 \to N_1$ et
$N_1 \oplus N_2 \to N_2$ induisent un inverse.
Puisque $TF$ est exact à droite, on voit que
$TF(B, N_1) \to TF(B, N_1 \oplus N_2) \to TF(B, N_2) \to 0$ est exacte.
Ainsi $TF(B, N_1) \oplus TF(B, N_2) \to TF(B, N_1 \oplus N_2)$ est un
isomorphisme. Cela démontre (1).

\medskip\noindent
Pour (2), il suffit de montrer que
$x \cdot (b_1 + b_2) = x \cdot b_1 + x \cdot b_2$.
(L'associativité et l'additivité sont claires.) Pour cela, considérons
$$
N \xrightarrow{(b_1, b_2)} N \oplus N \xrightarrow{+} N
$$
et appliquons $TF(B, -)$.

\medskip\noindent
L'assertion (3) résulte immédiatement de ce que
$N \otimes_B F(B) \to TF(B, N)$ est fonctorielle par rapport au couple $(B, N)$.

\medskip\noindent
Supposons que $N$ soit un $B$-module de présentation finie. Choisissons une présentation
$B^{\oplus m} \to B^{\oplus n} \to N \to 0$. Elle donne une suite
exacte
$$
TF(B, B^{\oplus m}) \to TF(B, B^{\oplus n}) \to TF(B, N) \to 0
$$
par exactitude à droite de $TF(B, -)$. D'après (1), on peut écrire
$TF(B, B^{\oplus m}) = TF(B, B)^{\oplus m}$ et
$TF(B, B^{\oplus n}) = TF(B, B)^{\oplus n}$. Supposons ensuite que
$B^{\oplus m} \to B^{\oplus n}$ soit donnée par la matrice $T = (b_{ij})$.
Alors l'application induite $TF(B, B)^{\oplus m} \to TF(B, B)^{\oplus n}$
est donnée par la matrice de coefficients $TF(B, b_{ij} \cdot 1_B)$.
Ce fait, joint à l'exactitude à droite de $\otimes$, démontre (4).
\end{proof}

\begin{example}
\label{adequate-example-module-structure-different}
Soit $F$ un foncteur à valeurs dans les modules comme dans le
lemme \ref{adequate-lemma-exactness-implies}.
Les deux structures de module sur
$TF(B, N)$ ne coïncident pas toujours. En voici un exemple. Supposons que $A = \mathbf{F}_p$,
où $p$ est un nombre premier. Posons $F(B) = B$, mais avec la structure de $B$-module
donnée par $b \cdot x = b^px$. Alors $TF(B, N) = N$, avec la structure de $B$-module
donnée par $b \cdot x = b^px$ pour $x \in N$. Cependant, la seconde structure de $B$-module
est donnée par $x \cdot b = bx$. Notons que, dans ce cas, l'application canonique
$N \otimes_B F(B) \to TF(B, N)$ est nulle, car élever un élément
$n \in B[N]$ à la puissance $p$ donne zéro.
\end{example}

\noindent
Dans le lemme suivant, nous utiliserons fréquemment l'observation suivante :
si $0 \to F \to G \to H \to 0$ est une suite exacte de foncteurs à valeurs dans les modules
sur $\textit{Alg}_A$, alors, pour tout couple $(B, N)$, la
suite $0 \to TF(B, N) \to TG(B, N) \to TH(B, N) \to 0$ est exacte.
Cela résulte de ce que $0 \to F(B[N]) \to G(B[N]) \to H(B[N]) \to 0$
est exacte.

\begin{lemma}
\label{adequate-lemma-exactness-permanence}
Soit $A$ un anneau. Pour un foncteur à valeurs dans les modules $F$ sur $\textit{Alg}_A$,
disons que $(*)$ est vérifiée si, pour tout $B \in \Ob(\textit{Alg}_A)$, le
foncteur $TF(B, -)$ sur les $B$-modules transforme une suite exacte courte
de $B$-modules en une suite exacte à droite. Soit
$0 \to F \to G \to H \to 0$ une suite exacte courte de
foncteurs à valeurs dans les modules sur $\textit{Alg}_A$.
\begin{enumerate}
\item Si $(*)$ est vérifiée pour $F, G$, alors $(*)$ est vérifiée pour $H$.
\item Si $(*)$ est vérifiée pour $F, H$, alors $(*)$ est vérifiée pour $G$.
\item Si $H' \to H$ est un morphisme de foncteurs à valeurs dans les modules sur $\textit{Alg}_A$
et si $(*)$ est vérifiée pour $F$, $G$, $H$ et $H'$, alors $(*)$ est vérifiée pour
$G \times_H H'$.
\end{enumerate}
\end{lemma}

\begin{proof}
Soit $B$ donné. Soit $0 \to N_1 \to N_2 \to N_3 \to 0$ une suite exacte
courte de $B$-modules. L'assertion (1) résulte d'une chasse dans le
diagramme
$$
\xymatrix{
0 \ar[r] &
TF(B, N_1) \ar[r] \ar[d] &
TG(B, N_1) \ar[r] \ar[d] &
TH(B, N_1) \ar[r] \ar[d] & 0 \\
0 \ar[r] &
TF(B, N_2) \ar[r] \ar[d] &
TG(B, N_2) \ar[r] \ar[d] &
TH(B, N_2) \ar[r] \ar[d] & 0 \\
0 \ar[r] &
TF(B, N_3) \ar[r] \ar[d] &
TG(B, N_3) \ar[r] \ar[d] &
TH(B, N_3) \ar[r] & 0 \\
& 0 & 0
}
$$
dont les lignes horizontales et les colonnes faisant intervenir $TF$ et $TG$ sont exactes.
Pour démontrer (2), on effectue une chasse dans le diagramme
$$
\xymatrix{
0 \ar[r] &
TF(B, N_1) \ar[r] \ar[d] &
TG(B, N_1) \ar[r] \ar[d] &
TH(B, N_1) \ar[r] \ar[d] & 0 \\
0 \ar[r] &
TF(B, N_2) \ar[r] \ar[d] &
TG(B, N_2) \ar[r] \ar[d] &
TH(B, N_2) \ar[r] \ar[d] & 0 \\
0 \ar[r] &
TF(B, N_3) \ar[r] \ar[d] &
TG(B, N_3) \ar[r] &
TH(B, N_3) \ar[r] \ar[d] & 0 \\
& 0 & & 0
}
$$
dont les lignes horizontales et les colonnes faisant intervenir $TF$ et $TH$ sont exactes.
L'assertion (3) résulte de (2), puisque $G \times_H H'$ s'insère dans la suite exacte
$0 \to F \to G \times_H H' \to H' \to 0$.
\end{proof}

\noindent
L'essentiel du travail de cette section a été accompli pour démontrer le
résultat d'annulation fondamental suivant.

\begin{lemma}
\label{adequate-lemma-ext-group-zero-key}
Soit $A$ un anneau. Soient $M$, $P$ des $A$-modules, avec $P$ de présentation
finie. Alors
$\Ext^i_\mathcal{P}(\underline{P}, \underline{M}) = 0$
pour $i > 0$, où $\mathcal{P}$ est la catégorie des foncteurs à valeurs dans les modules
sur $\textit{Alg}_A$.
\end{lemma}

\begin{proof}
Choisissons une résolution injective $\underline{M} \to I^\bullet$ dans
$\mathcal{P}$ ; voir le
lemme \ref{adequate-lemma-enough-injectives}.
D'après
Catégories dérivées, lemme \ref{derived-lemma-compute-ext-resolutions},
tout élément de $\Ext^i_\mathcal{P}(\underline{P}, \underline{M})$
provient d'un morphisme $\varphi : \underline{P} \to I^i$ tel que
$d^i \circ \varphi = 0$. Nous montrerons que
l'extension de Yoneda
$$
E : 0 \to \underline{M} \to I^0 \to \ldots \to
I^{i - 1} \times_{\Ker(d^i)} \underline{P} \to \underline{P} \to 0
$$
de $\underline{P}$ par $\underline{M}$
associée à $\varphi$ est triviale, ce qui démontrera le lemme d'après
Catégories dérivées, lemme \ref{derived-lemma-yoneda-extension}.

\medskip\noindent
Pour un foncteur à valeurs dans les modules $F$ sur $\textit{Alg}_A$,
disons que $(*)$ est vérifiée si, pour tout $B \in \Ob(\textit{Alg}_A)$, le
foncteur $TF(B, -)$ sur les $B$-modules transforme une suite exacte courte
de $B$-modules en une suite exacte à droite.
Rappelons que les foncteurs à valeurs dans les modules $\underline{M}, I^n, \underline{P}$
possèdent chacun la propriété $(*)$ ; voir le
lemme \ref{adequate-lemma-tangent-injective}
et les remarques qui le précèdent.
En décomposant $0 \to \underline{M} \to I^\bullet$ en suites exactes
courtes, on trouve que chacun des foncteurs
$\Im(d^{n - 1}) = \Ker(d^n) \subset I^n$ possède la propriété $(*)$ d'après le
lemme \ref{adequate-lemma-exactness-permanence},
et que $I^{i - 1} \times_{\Ker(d^i)} \underline{P}$ possède également la propriété
$(*)$.

\medskip\noindent
On peut donc supposer que l'extension de Yoneda est donnée sous la forme
$$
E : 0 \to \underline{M} \to F_{i - 1} \to \ldots \to
F_0 \to \underline{P} \to 0
$$
où chacun des foncteurs à valeurs dans les modules $F_j$ possède la propriété $(*)$.
Posons $G_j(B) = TF_j(B, B)$, considéré comme $B$-module au moyen de la {\it seconde}
structure de $B$-module définie dans le
lemme \ref{adequate-lemma-exactness-implies}.
Puisque $TF_j$ est un foncteur sur les couples, on voit que $G_j$ est un foncteur
à valeurs dans les modules sur $\textit{Alg}_A$. De plus, puisque $E$ est une suite exacte,
la suite $G_{j + 1} \to G_j \to G_{j - 1}$ est exacte (voir la remarque
précédant le
lemme \ref{adequate-lemma-exactness-permanence}).
Observons que $T\underline{M}(B, B) = M \otimes_A B = \underline{M}(B)$
et que les deux structures de $B$-module y coïncident.
On obtient ainsi une extension de Yoneda
$$
E' :  0 \to \underline{M} \to G_{i - 1} \to \ldots \to
G_0 \to \underline{P} \to 0
$$
De plus, les applications canoniques
$$
F_j(B) = B \otimes_B F_j(B) \longrightarrow TF_j(B, B) = G_j(B)
$$
du
lemme \ref{adequate-lemma-tangent-functor} (\ref{adequate-item-mult-map})
sont $B$-linéaires d'après le
lemme \ref{adequate-lemma-exactness-implies} (\ref{adequate-item-mult-map-linear})
et fonctorielles en $B$. On a donc une application
$$
\xymatrix{
0 \ar[r] &
\underline{M} \ar[r] \ar[d]^1 &
F_{i - 1} \ar[r] \ar[d] &
\ldots \ar[r] &
F_0 \ar[r] \ar[d] &
\underline{P} \ar[r] \ar[d]^1 & 0 \\
0 \ar[r] &
\underline{M} \ar[r] &
G_{i - 1} \ar[r] &
\ldots \ar[r] &
G_0 \ar[r] &
\underline{P} \ar[r] & 0
}
$$
d'extensions de Yoneda. En particulier, on voit que $E$ et $E'$ ont la
même classe dans $\Ext^i_\mathcal{P}(\underline{P}, \underline{M})$,
d'après le lemme sur les Ext de Yoneda mentionné ci-dessus. Enfin, soit $N$ un
$A$-module de présentation finie. On voit alors que
$$
0 \to T\underline{M}(A, N) \to TF_{i - 1}(A, N) \to \ldots \to
TF_0(A, N) \to T\underline{P}(A, N) \to 0
$$
est exacte. D'après le
lemme \ref{adequate-lemma-exactness-implies} (\ref{adequate-item-tangent-right-exact}),
avec $B = A$, cela se traduit par l'exactitude de la suite de
$A$-modules
$$
0 \to M \otimes_A N \to G_{i - 1}(A) \otimes_A N \to \ldots \to
G_0(A) \otimes_A N \to P \otimes_A N \to 0
$$
Ainsi la suite de $A$-modules
$0 \to M \to G_{i - 1}(A) \to \ldots \to G_0(A) \to P \to 0$
est universellement exacte, au sens où elle reste exacte après tensorisation
par tout $A$-module de présentation finie $N$. Posons
$K = \Ker(G_0(A) \to P)$, de sorte que l'on a des suites exactes
$$
0 \to K \to G_0(A) \to P \to 0
\quad\text{et}\quad
G_2(A) \to G_1(A) \to K \to 0
$$
En tensorisant la seconde suite par $N$, on obtient
$K \otimes_A N = \Coker(G_2(A) \otimes_A N \to G_1(A) \otimes_A N)$.
L'exactitude de $G_2(A) \otimes_A N \to G_1(A) \otimes_A N \to G_0(A) \otimes_A N$
entraîne alors que $K \otimes_A N \to G_0(A) \otimes_A N$ est injective.
D'après
Algèbre, théorème \ref{algebra-theorem-universally-exact-criteria},
cela signifie que l'extension de $A$-modules $0 \to K \to G_0(A) \to P \to 0$
est exacte ; puisque $P$ est supposé de présentation finie, cela signifie
que la suite est scindée ; voir
Algèbre, lemme \ref{algebra-lemma-universally-exact-split}.
Tout scindage $P \to G_0(A)$ définit une application $\underline{P} \to G_0$
qui scinde la surjection $G_0 \to \underline{P}$. Ainsi
l'extension de Yoneda $E'$ est équivalente à l'extension de Yoneda triviale,
ce qui conclut.
\end{proof}

\begin{lemma}
\label{adequate-lemma-ext-group-zero}
Soit $A$ un anneau. Soit $M$ un $A$-module. Soit $L$ un foncteur linéairement
adéquat sur $\textit{Alg}_A$. Alors
$\Ext^i_\mathcal{P}(L, \underline{M}) = 0$
pour $i > 0$, où $\mathcal{P}$ est la catégorie des foncteurs à valeurs dans les modules
sur $\textit{Alg}_A$.
\end{lemma}

\begin{proof}
Puisque $L$ est linéairement adéquat, il existe une suite exacte
$$
0 \to L \to \underline{A^{\oplus m}} \to \underline{A^{\oplus n}} \to
\underline{P} \to 0
$$
Ici $P = \Coker(A^{\oplus m} \to A^{\oplus n})$ est le conoyau
de l'application de $A$-modules libres de rang fini donnée par la définition
des foncteurs linéairement adéquats. D'après le
lemme \ref{adequate-lemma-ext-group-zero-key},
on a l'annulation de
$\Ext^i_\mathcal{P}(\underline{P}, \underline{M})$
et de
$\Ext^i_\mathcal{P}(\underline{A}, \underline{M})$
pour $i > 0$.
Posons $K = \Ker(\underline{A^{\oplus n}} \to \underline{P})$.
La suite exacte longue des groupes Ext associée à la suite exacte
$0 \to K \to \underline{A^{\oplus n}} \to \underline{P} \to 0$
permet de conclure que
$\Ext^i_\mathcal{P}(K, \underline{M}) = 0$ pour $i > 0$.
En répétant l'argument avec la suite
$0 \to L \to \underline{A^{\oplus m}} \to K \to 0$,
on conclut.
\end{proof}

\begin{lemma}
\label{adequate-lemma-RQ-zero}
Avec les notations du
lemme \ref{adequate-lemma-enough-injectives},
on a $R^pQ(F) = 0$ pour tout $p > 0$ et tout foncteur adéquat $F$.
\end{lemma}

\begin{proof}
Choisissons une suite exacte $0 \to F \to \underline{M^0} \to \underline{M^1}$.
Posons $M^2 = \Coker(M^0 \to M^1)$, de sorte que
$0 \to F \to \underline{M^0} \to \underline{M^1}
\to \underline{M^2} \to 0$ est une résolution. D'après
Catégories dérivées, lemme \ref{derived-lemma-two-ss-complex-functor},
on obtient une suite spectrale
$$
R^pQ(\underline{M^q}) \Rightarrow R^{p + q}Q(F)
$$
Puisque $Q(\underline{M^q}) = \underline{M^q}$, il suffit de démontrer
$R^pQ(\underline{M}) = 0$, $p > 0$, pour tout $A$-module $M$.

\medskip\noindent
Choisissons une résolution injective $\underline{M} \to I^\bullet$ dans
la catégorie $\mathcal{P}$. Supposons que $R^iQ(\underline{M})$ soit non nul.
Alors $\Ker(Q(I^i) \to Q(I^{i + 1}))$ contient strictement
l'image de $Q(I^{i - 1}) \to Q(I^i)$. Ainsi, d'après le
lemme \ref{adequate-lemma-adequate-surjection-from-linear},
il existe un foncteur linéairement adéquat $L$ et une application
$\varphi : L \to Q(I^i)$ à valeurs dans le noyau de $Q(I^i) \to Q(I^{i + 1})$
qui ne se factorise pas par l'image de $Q(I^{i - 1}) \to Q(I^i)$.
Puisque $Q$ est adjoint à gauche du foncteur d'inclusion, l'application
$\varphi$ correspond à une application $\varphi' : L \to I^i$ ayant les mêmes propriétés.
Ainsi $\varphi'$ donne un élément non nul de
$\Ext^i_\mathcal{P}(L, \underline{M})$, en contradiction avec le
lemme \ref{adequate-lemma-ext-group-zero}.
\end{proof}














\section{Modules adéquats}
\label{adequate-section-adequate}

\noindent
Dans
Descente, section \ref{descent-section-quasi-coherent-sheaves},
nous avons vu que les modules quasi-cohérents sur un schéma $S$
s'identifient aux modules quasi-cohérents sur n'importe lequel des grands
sites $(\Sch/S)_\tau$ associés à $S$. Nous avons vu que cette
identification présente deux difficultés :
\begin{enumerate}
\item $\QCoh(\mathcal{O}_S) \to
\textit{Mod}((\Sch/S)_\tau, \mathcal{O})$,
$\mathcal{F} \mapsto \mathcal{F}^a$ n'est pas exact en général
(Descente, lemme \ref{descent-lemma-equivalence-quasi-coherent-limits}), et
\item étant donné un morphisme quasi-compact et quasi-séparé $f : X \to S$,
le foncteur $f_*$ ne préserve pas en général les faisceaux quasi-cohérents sur les
grands sites (Descente, proposition
\ref{descent-proposition-equivalence-quasi-coherent-functorial}).
\end{enumerate}
L'assertion (1) signifie que l'on ne peut pas définir une sous-catégorie triangulée
de $D(\mathcal{O})$ constituée des complexes dont les faisceaux de cohomologie
sont quasi-cohérents. L'assertion (2) signifie que $Rf_*\mathcal{F}$ n'est pas un
complexe à faisceaux de cohomologie quasi-cohérents, même lorsque $\mathcal{F}$
est quasi-cohérent et que $f$ est quasi-compact et quasi-séparé.
De plus, les exemples donnés dans les démonstrations de
Descente, lemme
\ref{descent-lemma-equivalence-quasi-coherent-limits},
et de
Descente, proposition
\ref{descent-proposition-equivalence-quasi-coherent-functorial},
ne sont pas de nature pathologique.

\medskip\noindent
Dans cette section, nous étudions une catégorie légèrement plus grande
de $\mathcal{O}$-modules sur $(\Sch/S)_\tau$, qui contient les
modules quasi-cohérents, est abélienne et est préservée par $f_*$ lorsque
$f$ est quasi-compact et quasi-séparé.
Pour cela, supposons que $S$ soit un schéma. Soit $\mathcal{F}$ un préfaisceau
de $\mathcal{O}$-modules sur $(\Sch/S)_\tau$.
Pour tout objet affine $U = \Spec(A)$ de $(\Sch/S)_\tau$,
on peut restreindre $\mathcal{F}$ à $(\textit{Aff}/U)_\tau$ pour obtenir
un préfaisceau de $\mathcal{O}$-modules sur ce site. Le foncteur correspondant
à valeurs dans les modules, voir la
section \ref{adequate-section-quasi-coherent},
sera noté
$$
F = F_{\mathcal{F}, A} :
\textit{Alg}_A \longrightarrow \textit{Ab},
\quad
B \longmapsto \mathcal{F}(\Spec(B))
$$
La correspondance $\mathcal{F} \mapsto F_{\mathcal{F}, A}$ est un foncteur exact
de catégories abéliennes.

\begin{definition}
\label{adequate-definition-adequate}
Un faisceau de $\mathcal{O}$-modules $\mathcal{F}$ sur $(\Sch/S)_\tau$ est
{\it adéquat} s'il existe un $\tau$-recouvrement
$\{\Spec(A_i) \to S\}_{i \in I}$ tel que $F_{\mathcal{F}, A_i}$ soit
adéquat pour tout $i \in I$.
\end{definition}

\noindent
Nous verrons ci-dessous que la catégorie des $\mathcal{O}$-modules adéquats
est indépendante de la topologie $\tau$ choisie.

\begin{lemma}
\label{adequate-lemma-adequate-local}
Soit $S$ un schéma. Soit $\mathcal{F}$ un $\mathcal{O}$-module adéquat sur
$(\Sch/S)_\tau$. Pour tout schéma affine $\Spec(A)$ au-dessus de $S$,
le foncteur $F_{\mathcal{F}, A}$ est adéquat.
\end{lemma}

\begin{proof}
Soit $\{\Spec(A_i) \to S\}_{i \in I}$ un $\tau$-recouvrement
tel que $F_{\mathcal{F}, A_i}$ soit adéquat pour tout $i \in I$.
On peut trouver un $\tau$-recouvrement affine standard
$\{\Spec(A'_j) \to \Spec(A)\}_{j = 1, \ldots, m}$
tel que $\Spec(A'_j) \to \Spec(A) \to S$ se factorise
par $\Spec(A_{i(j)})$ pour un certain $i(j) \in I$. On voit alors que
$F_{\mathcal{F}, A'_j}$ est la restriction de
$F_{\mathcal{F}, A_{i(j)}}$ à la catégorie des $A'_j$-algèbres.
Ainsi $F_{\mathcal{F}, A'_j}$ est adéquat d'après le
lemme \ref{adequate-lemma-base-change-adequate}.
D'après le
lemme \ref{adequate-lemma-adequate-product},
la suite
$F_{\mathcal{F}, A'_j}$ correspond à un foncteur « produit » adéquat
$F'$ sur $A' = A'_1 \times \ldots \times A'_m$. Comme $\mathcal{F}$ est un
faisceau (pour la topologie de Zariski), ce foncteur produit $F'$ est égal
à $F_{\mathcal{F}, A'}$, c'est-à-dire à la restriction de $F$ aux $A'$-algèbres.
Enfin, $\{\Spec(A') \to \Spec(A)\}$ est un $\tau$-recouvrement.
Il résulte du
lemme \ref{adequate-lemma-adequate-descent}
que $F_{\mathcal{F}, A}$ est adéquat.
\end{proof}

\begin{lemma}
\label{adequate-lemma-adequate-affine}
Soit $S = \Spec(A)$ un schéma affine. La catégorie des
$\mathcal{O}$-modules adéquats sur $(\Sch/S)_\tau$ est équivalente à la
catégorie des foncteurs adéquats à valeurs dans les modules sur $\textit{Alg}_A$.
\end{lemma}

\begin{proof}
Étant donné un module adéquat $\mathcal{F}$, le foncteur $F_{\mathcal{F}, A}$
est adéquat d'après le lemme \ref{adequate-lemma-adequate-local}.
Étant donné un foncteur adéquat $F$, choisissons une suite exacte
$0 \to F \to \underline{M} \to \underline{N}$ et considérons
le $\mathcal{O}$-module $\mathcal{F} = \Ker(M^a \to N^a)$, où
$M^a$ désigne le $\mathcal{O}$-module quasi-cohérent sur
$(\Sch/S)_\tau$ associé au faisceau quasi-cohérent
$\widetilde{M}$ sur $S$. Notons que $F = F_{\mathcal{F}, A}$ ; en particulier,
le module $\mathcal{F}$ est adéquat par définition.
Nous omettons la démonstration du fait que ces constructions définissent des équivalences
de catégories inverses l'une de l'autre.
\end{proof}

\begin{lemma}
\label{adequate-lemma-pullback-adequate}
Soit $f : T \to S$ un morphisme de schémas.
L'image inverse $f^*\mathcal{F}$ d'un $\mathcal{O}$-module adéquat
$\mathcal{F}$ sur $(\Sch/S)_\tau$ est un
$\mathcal{O}$-module adéquat sur $(\Sch/T)_\tau$.
\end{lemma}

\begin{proof}
Le foncteur d'image inverse
$f^* : \textit{Mod}((\Sch/S)_\tau, \mathcal{O}) \to
\textit{Mod}((\Sch/T)_\tau, \mathcal{O})$
est donné par restriction, c'est-à-dire que $f^*\mathcal{F}(V) = \mathcal{F}(V)$
pour tout schéma $V$ au-dessus de $T$. Ce lemme résulte donc immédiatement du
lemme \ref{adequate-lemma-adequate-local}
et de la définition.
\end{proof}

\noindent
Voici une caractérisation de la catégorie des $\mathcal{O}$-modules adéquats.
Pour en comprendre la portée, considérons une application $\mathcal{G} \to \mathcal{H}$
de $\mathcal{O}_S$-modules quasi-cohérents sur un schéma $S$.
Le conoyau de l'application associée $\mathcal{G}^a \to \mathcal{H}^a$
de $\mathcal{O}$-modules est quasi-cohérent, car il est égal à
$(\mathcal{H}/\mathcal{G})^a$. Mais le noyau de
$\mathcal{G}^a \to \mathcal{H}^a$ n'est pas quasi-cohérent
en général. Il est toutefois adéquat.

\begin{lemma}
\label{adequate-lemma-adequate-characterize}
Soit $S$ un schéma. Soit $\mathcal{F}$ un $\mathcal{O}$-module sur
$(\Sch/S)_\tau$. Les conditions suivantes sont équivalentes
\begin{enumerate}
\item $\mathcal{F}$ est adéquat,
\item il existe un recouvrement ouvert affine $S = \bigcup S_i$ et des
applications de $\mathcal{O}_{S_i}$-modules quasi-cohérents
$\mathcal{G}_i \to \mathcal{H}_i$
telles que $\mathcal{F}|_{(\Sch/S_i)_\tau}$ soit le
noyau de $\mathcal{G}_i^a \to \mathcal{H}_i^a$
\item il existe un $\tau$-recouvrement $\{S_i \to S\}_{i \in I}$ et des
applications de $\mathcal{O}_{S_i}$-modules quasi-cohérents
$\mathcal{G}_i \to \mathcal{H}_i$
telles que $\mathcal{F}|_{(\Sch/S_i)_\tau}$ soit le
noyau de $\mathcal{G}_i^a \to \mathcal{H}_i^a$,
\item il existe un $\tau$-recouvrement $\{f_i : S_i \to S\}_{i \in I}$
tel que chaque $f_i^*\mathcal{F}$ soit adéquat,
\item pour tout schéma affine $U$ au-dessus de $S$, la restriction
$\mathcal{F}|_{(\Sch/U)_\tau}$ est le noyau
d'une application $\mathcal{G}^a \to \mathcal{H}^a$ de
$\mathcal{O}_U$-modules quasi-cohérents.
\end{enumerate}
\end{lemma}

\begin{proof}
Soit $U = \Spec(A)$ un schéma affine au-dessus de $S$.
Posons $F = F_{\mathcal{F}, A}$. Par définition, le foncteur
$F$ est adéquat si et seulement s'il existe une application de $A$-modules
$M \to N$ telle que $F = \Ker(\underline{M} \to \underline{N})$.
En combinant cela avec les
lemmes \ref{adequate-lemma-adequate-local} et
\ref{adequate-lemma-adequate-affine},
on voit que (1) et (5) sont équivalentes.

\medskip\noindent
Il est clair que (5) entraîne (2) et que (2) entraîne (3).
Si (3) est vérifiée, on peut raffiner le recouvrement
$\{S_i \to S\}$ de sorte que chaque $S_i = \Spec(A_i)$ soit affine.
Les remarques préliminaires de la démonstration montrent alors que
$F_{\mathcal{F}, A_i}$ est adéquat. Ainsi $\mathcal{F}$
est adéquat par définition. Donc (3) entraîne (1).

\medskip\noindent
Enfin, (4) est équivalente à (1) : on utilise le
lemme \ref{adequate-lemma-pullback-adequate}
dans un sens, et le fait qu'un
composé de $\tau$-recouvrements est un $\tau$-recouvrement dans l'autre.
\end{proof}

\noindent
Tout comme pour les faisceaux quasi-cohérents, la catégorie des
modules adéquats est indépendante de la topologie.

\begin{lemma}
\label{adequate-lemma-adequate-fpqc}
Soit $\mathcal{F}$ un $\mathcal{O}$-module adéquat sur
$(\Sch/S)_\tau$. Pour tout morphisme plat surjectif
$\Spec(B) \to \Spec(A)$ de schémas affines au-dessus de $S$,
le complexe de {\v C}ech augmenté
$$
0 \to \mathcal{F}(\Spec(A)) \to
\mathcal{F}(\Spec(B)) \to
\mathcal{F}(\Spec(B \otimes_A B)) \to \ldots
$$
est exact. En particulier, $\mathcal{F}$ satisfait à la condition de faisceau
pour les recouvrements fpqc, et c'est un faisceau de $\mathcal{O}$-modules
sur $(\Sch/S)_{fppf}$.
\end{lemma}

\begin{proof}
Avec $A \to B$ comme dans le lemme, posons $F = F_{\mathcal{F}, A}$. Ce foncteur
est adéquat d'après le
lemme \ref{adequate-lemma-adequate-local}.
D'après le
lemme \ref{adequate-lemma-adequate-flat},
puisque $A \to B$, $A \to B \otimes_A B$, etc., sont plats, on voit que
$F(B) = F(A) \otimes_A B$,
$F(B \otimes_A B) = F(A) \otimes_A B \otimes_A B$, etc.
L'exactitude résulte de
Descente, lemme \ref{descent-lemma-ff-exact}.

\medskip\noindent
Ainsi $\mathcal{F}$ satisfait à la condition de faisceau pour les
$\tau$-recouvrements (en particulier les recouvrements de Zariski) et pour tout
recouvrement fidèlement plat d'un schéma affine par un schéma affine. En raisonnant
comme dans les démonstrations de Descente, lemme \ref{descent-lemma-standard-fpqc-covering},
et de
Descente, proposition \ref{descent-proposition-fpqc-descent-quasi-coherent},
on conclut que $\mathcal{F}$ satisfait à la condition de faisceau pour tous les
recouvrements fpqc (constitués d'objets de $(\Sch/S)_\tau$).
Nous omettons les détails.
\end{proof}

\noindent
Le lemme \ref{adequate-lemma-adequate-fpqc} montre en particulier que,
pour tout couple de topologies $\tau, \tau'$, les collections
des modules adéquats pour la topologie $\tau$ et pour la topologie $\tau'$
sont identiques (comme préfaisceaux de modules sur la catégorie sous-jacente $\Sch/S$).

\begin{definition}
\label{adequate-definition-category-adequate-modules}
Soit $S$ un schéma. La catégorie des $\mathcal{O}$-modules adéquats sur
$(\Sch/S)_\tau$ est notée {\it $\textit{Adeq}(\mathcal{O})$} ou
{\it $\textit{Adeq}((\Sch/S)_\tau, \mathcal{O})$}. Si l'on veut considérer
la catégorie abélienne des modules adéquats sans choisir de
topologie, on écrit simplement {\it $\textit{Adeq}(S)$}.
\end{definition}

\begin{lemma}
\label{adequate-lemma-same-cohomology-adequate}
Soit $S$ un schéma. Soit $\mathcal{F}$ un
$\mathcal{O}$-module adéquat sur $(\Sch/S)_\tau$.
\begin{enumerate}
\item La restriction $\mathcal{F}|_{S_{Zar}}$ est un
$\mathcal{O}_S$-module quasi-cohérent sur le schéma $S$.
\item La restriction $\mathcal{F}|_{S_\etale}$ est le
module quasi-cohérent associé à $\mathcal{F}|_{S_{Zar}}$.
\item Pour tout schéma affine $U$ au-dessus de $S$, on a $H^q(U, \mathcal{F}) = 0$
pour tout $q > 0$.
\item Il existe un isomorphisme canonique
$$
H^q(S, \mathcal{F}|_{S_{Zar}}) =
H^q((\Sch/S)_\tau, \mathcal{F}).
$$
\end{enumerate}
\end{lemma}

\begin{proof}
D'après le
lemme \ref{adequate-lemma-adequate-flat}
et le
lemme \ref{adequate-lemma-adequate-local},
pour tout morphisme plat de schémas affines $U \to V$ au-dessus de $S$,
on a
$\mathcal{F}(U) = \mathcal{F}(V) \otimes_{\mathcal{O}(V)} \mathcal{O}(U)$.
Cela vaut en particulier si $U \subset V \subset S$ sont des ouverts affines de
$S$ ; ainsi $\mathcal{F}|_{S_{Zar}}$ est quasi-cohérent.
On obtient donc (1).

\medskip\noindent
Soit $S' \to S$ un morphisme étale de schémas.
Pour un ouvert affine $U \subset S'$ dont l'image est contenue dans un ouvert affine
$V \subset S$, on voit alors que
$\mathcal{F}(U) = \mathcal{F}(V) \otimes_{\mathcal{O}(V)} \mathcal{O}(U)$,
car $U \to V$ est étale, donc plat. Par conséquent,
$\mathcal{F}|_{S'_{Zar}}$ est l'image inverse de $\mathcal{F}|_{S_{Zar}}$.
Cela démontre (2).

\medskip\noindent
Nous allons appliquer
Cohomologie sur les sites,
lemme \ref{sites-cohomology-lemma-cech-vanish-collection},
au site $(\Sch/S)_\tau$, avec
$\mathcal{B}$ l'ensemble des schémas affines au-dessus de $S$ et
$\text{Cov}$ l'ensemble des $\tau$-recouvrements affines standard.
L'hypothèse (3) du lemme est satisfaite d'après
Descente, lemme \ref{descent-lemma-standard-covering-Cech},
et le
lemme \ref{adequate-lemma-adequate-fpqc}
pour le cas d'un recouvrement par un seul schéma affine.
On en conclut que $H^p(U, \mathcal{F}) = 0$ pour tout
schéma affine $U$ au-dessus de $S$. Cela démontre (3).
Exactement comme dans la démonstration de
Descente, proposition \ref{descent-proposition-same-cohomology-quasi-coherent},
cela entraîne l'égalité des cohomologies (4).
\end{proof}

\begin{remark}
\label{adequate-remark-compare}
Soit $S$ un schéma. On dispose de foncteurs
$u : \QCoh(\mathcal{O}_S) \to \textit{Adeq}(\mathcal{O})$
et
$v : \textit{Adeq}(\mathcal{O}) \to \QCoh(\mathcal{O}_S)$.
Le foncteur $u : \mathcal{F} \mapsto \mathcal{F}^a$
consiste à prendre le $\mathcal{O}$-module associé, qui est
adéquat d'après le
lemme \ref{adequate-lemma-adequate-characterize}.
Réciproquement, le foncteur $v$ est donné par restriction,
$v : \mathcal{G} \mapsto \mathcal{G}|_{S_{Zar}}$ ; voir le
lemme \ref{adequate-lemma-same-cohomology-adequate}.
Puisque $\mathcal{F}^a$ peut être décrit comme l'image inverse de
$\mathcal{F}$ par un morphisme de topos annelés
$((\Sch/S)_\tau, \mathcal{O}) \to (S_{Zar}, \mathcal{O}_S)$, voir
Descente, remarque \ref{descent-remark-change-topologies-ringed-sites},
et puisque la restriction est l'image directe, on voit que $u$ et $v$
sont adjoints comme suit
$$
\SheafHom_{\mathcal{O}_S}(\mathcal{F}, v\mathcal{G})
=
\SheafHom_\mathcal{O}(u\mathcal{F}, \mathcal{G})
$$
où $\mathcal{O}$ désigne le faisceau structural sur le grand site.
La description montre immédiatement que le morphisme d'adjonction
$\mathcal{F} \to vu\mathcal{F}$ est un isomorphisme pour tout faisceau
quasi-cohérent.
\end{remark}

\begin{lemma}
\label{adequate-lemma-sheafification-adequate}
{\emergencystretch=1em
Soit $S$ un schéma. Soit $\mathcal{F}$ un préfaisceau de $\mathcal{O}$-modules
sur $(\Sch/S)_\tau$. Si, pour tout schéma affine
$\Spec(A)$ au-dessus de $S$, le foncteur $F_{\mathcal{F}, A}$ est
adéquat, alors le faisceau associé à $\mathcal{F}$ est un
$\mathcal{O}$-module adéquat.\par}
\end{lemma}

\begin{proof}
Soit $U = \Spec(A)$ un schéma affine au-dessus de $S$.
Posons $F = F_{\mathcal{F}, A}$.
Le faisceau associé est $\mathcal{F}^\# = (\mathcal{F}^+)^+$ ; voir
Sites, section \ref{sites-section-sheafification}.
Par construction,
$$
(\mathcal{F})^+(U) =
\colim_\mathcal{U} \check{H}^0(\mathcal{U}, \mathcal{F})
$$
où la limite inductive porte sur les recouvrements dans le site $(\Sch/S)_\tau$.
Puisque $U$ est affine, il suffit de prendre la limite inductive sur les
$\tau$-recouvrements affines standard
$\mathcal{U} = \{U_i \to U\}_{i \in I} =
\{\Spec(A_i) \to \Spec(A)\}_{i \in I}$ de $U$.
Puisque chaque $A \to A_i$ et chaque $A \to A_i \otimes_A A_j$ est plat, on voit que
$$
\check{H}^0(\mathcal{U}, \mathcal{F}) =
\Ker(\prod F(A) \otimes_A A_i \to \prod F(A) \otimes_A A_i \otimes_A A_j)
$$
d'après le
lemme \ref{adequate-lemma-adequate-flat}.
Puisque $A \to \prod A_i$ est fidèlement plat, ce module
est toujours canoniquement isomorphe à $F(A)$ d'après
Descente, lemme \ref{descent-lemma-ff-exact}.
Ainsi le préfaisceau $(\mathcal{F})^+$ a la même valeur que
$\mathcal{F}$ sur tous les schémas affines au-dessus de $S$. En répétant une fois
l'argument, on en déduit la même assertion pour $\mathcal{F}^\# = ((\mathcal{F})^+)^+$.
Ainsi $F_{\mathcal{F}, A} = F_{\mathcal{F}^\#, A}$, et l'on conclut
que $\mathcal{F}^\#$ est adéquat.
\end{proof}

\begin{lemma}
\label{adequate-lemma-abelian-adequate}
Soit $S$ un schéma.
\begin{enumerate}
\item La catégorie $\textit{Adeq}(\mathcal{O})$ est abélienne.
\item Le foncteur
$\textit{Adeq}(\mathcal{O}) \to
\textit{Mod}((\Sch/S)_\tau, \mathcal{O})$
est exact.
\item Si $0 \to \mathcal{F}_1 \to \mathcal{F}_2 \to \mathcal{F}_3 \to 0$
est une suite exacte courte de $\mathcal{O}$-modules et si
$\mathcal{F}_1$ et $\mathcal{F}_3$ sont adéquats, alors
$\mathcal{F}_2$ est adéquat.
\item La catégorie $\textit{Adeq}(\mathcal{O})$ admet des limites inductives, et
le foncteur
$\textit{Adeq}(\mathcal{O}) \to
\textit{Mod}((\Sch/S)_\tau, \mathcal{O})$
commute à ces limites.
\end{enumerate}
\end{lemma}

\begin{proof}
Soit $\varphi : \mathcal{F} \to \mathcal{G}$ une application de
$\mathcal{O}$-modules adéquats. Pour démontrer (1) et (2), il suffit de montrer que
$\mathcal{K} = \Ker(\varphi)$ et
$\mathcal{Q} = \Coker(\varphi)$, calculés dans
$\textit{Mod}((\Sch/S)_\tau, \mathcal{O})$, sont adéquats.
Soit $U = \Spec(A)$ un schéma affine au-dessus de $S$.
Posons $F = F_{\mathcal{F}, A}$ et $G = F_{\mathcal{G}, A}$.
D'après les
lemmes \ref{adequate-lemma-kernel-adequate} et
\ref{adequate-lemma-cokernel-adequate},
le noyau $K$ et le conoyau $Q$ de l'application induite
$F \to G$ sont des foncteurs adéquats.
Comme le noyau se calcule au niveau des préfaisceaux, on voit
que $K = F_{\mathcal{K}, A}$, et l'on conclut que $\mathcal{K}$ est adéquat.
Pour démontrer le résultat pour le conoyau, notons $\mathcal{Q}'$ le conoyau
préfaisceautique de $\varphi$. Alors $Q = F_{\mathcal{Q}', A}$ et
$\mathcal{Q} = (\mathcal{Q}')^\#$. Ainsi $\mathcal{Q}$
est adéquat d'après le
lemme \ref{adequate-lemma-sheafification-adequate}.

\medskip\noindent
Soit $0 \to \mathcal{F}_1 \to \mathcal{F}_2 \to \mathcal{F}_3 \to 0$
une suite exacte courte de $\mathcal{O}$-modules, avec
$\mathcal{F}_1$ et $\mathcal{F}_3$ adéquats.
Soit $U = \Spec(A)$ un schéma affine au-dessus de $S$.
Posons $F_i = F_{\mathcal{F}_i, A}$. La suite de foncteurs
$$
0 \to F_1 \to F_2 \to F_3 \to 0
$$
est exacte, car, pour $V = \Spec(B)$ affine au-dessus de $U$, on a
$H^1(V, \mathcal{F}_1) = 0$ d'après le
lemme \ref{adequate-lemma-same-cohomology-adequate}.
Puisque $F_1$ et $F_3$ sont des foncteurs adéquats d'après le
lemme \ref{adequate-lemma-adequate-local},
on voit que $F_2$ est adéquat d'après le
lemme \ref{adequate-lemma-extension-adequate}.
Ainsi $\mathcal{F}_2$ est adéquat.

\medskip\noindent
Soit $\mathcal{I} \to \textit{Adeq}(\mathcal{O})$, $i \mapsto \mathcal{F}_i$,
un diagramme. Notons $\mathcal{F} = \colim_i \mathcal{F}_i$
la limite inductive calculée dans
$\textit{Mod}((\Sch/S)_\tau, \mathcal{O})$.
Pour démontrer (4), il suffit de montrer que $\mathcal{F}$ est adéquat.
Soit $\mathcal{F}' = \colim_i \mathcal{F}_i$ la limite inductive calculée
dans les préfaisceaux de $\mathcal{O}$-modules. Alors
$\mathcal{F} = (\mathcal{F}')^\#$.
Soit $U = \Spec(A)$ un schéma affine au-dessus de $S$.
Posons $F_i = F_{\mathcal{F}_i, A}$. D'après le
lemme \ref{adequate-lemma-colimit-adequate},
le foncteur $\colim_i F_i = F_{\mathcal{F}', A}$ est adéquat.
Le lemme \ref{adequate-lemma-sheafification-adequate}
montre que $\mathcal{F}$ est adéquat.
\end{proof}

\noindent
Le lemme suivant affirme que l'image directe totale
$Rf_*\mathcal{F}$ d'un module adéquat par un morphisme quasi-compact et
quasi-séparé est un complexe dont les faisceaux de cohomologie
sont adéquats.

\begin{lemma}
\label{adequate-lemma-direct-image-adequate}
Soit $f : T \to S$ un morphisme quasi-compact et quasi-séparé
de schémas. Pour tout $\mathcal{O}_T$-module adéquat sur
$(\Sch/T)_\tau$, l'image directe
$f_*\mathcal{F}$ et les images directes supérieures $R^if_*\mathcal{F}$
sont des $\mathcal{O}_S$-modules adéquats sur $(\Sch/S)_\tau$.
\end{lemma}

\begin{proof}
Expliquons d'abord comment calculer les images directes supérieures.
Choisissons une résolution injective $\mathcal{F} \to \mathcal{I}^\bullet$.
Alors $R^if_*\mathcal{F}$ est le $i$-ième faisceau de cohomologie du
complexe $f_*\mathcal{I}^\bullet$.
Ainsi $R^if_*\mathcal{F}$ est le faisceau associé au préfaisceau
qui associe à un objet $U/S$ de $(\Sch/S)_\tau$
le module
\begin{align*}
\frac{\Ker(f_*\mathcal{I}^i(U) \to f_*\mathcal{I}^{i + 1}(U))}
{\Im(f_*\mathcal{I}^{i - 1}(U) \to f_*\mathcal{I}^i(U))}
& =
\frac{\Ker(\mathcal{I}^i(U \times_S T) \to
\mathcal{I}^{i + 1}(U \times_S T))}
{\Im(\mathcal{I}^{i - 1}(U \times_S T) \to \mathcal{I}^i(U \times_S T))}
\\
& =
H^i(U \times_S T, \mathcal{F}) \\
& = H^i((\Sch/U \times_S T)_\tau,
\mathcal{F}|_{(\Sch/U \times_S T)_\tau}) \\
& = H^i(U \times_S T, \mathcal{F}|_{(U \times_S T)_{Zar}})
\end{align*}
La première égalité résulte de
Topologies, lemme \ref{topologies-lemma-morphism-big-fppf}
(et de ses analogues pour les autres topologies),
la deuxième de la définition de la cohomologie de $\mathcal{F}$
au-dessus d'un objet de $(\Sch/T)_\tau$,
la troisième de
Cohomologie sur les sites, lemme \ref{sites-cohomology-lemma-cohomology-of-open},
et la dernière du
lemme \ref{adequate-lemma-same-cohomology-adequate}.
Ainsi, d'après le
lemme \ref{adequate-lemma-sheafification-adequate},
il suffit de démontrer l'assertion énoncée dans le paragraphe suivant.

\medskip\noindent
Soit $A$ un anneau. Soit $T$ un schéma quasi-compact et quasi-séparé
sur $A$. Soit $\mathcal{F}$ un $\mathcal{O}_T$-module adéquat sur
$(\Sch/T)_\tau$. Pour une $A$-algèbre $B$, posons
$T_B = T \times_{\Spec(A)} \Spec(B)$ et notons
$\mathcal{F}_B = \mathcal{F}|_{(T_B)_{Zar}}$ la restriction de
$\mathcal{F}$ au petit site de Zariski de $T_B$.
(Rappelons qu'il s'agit d'un faisceau quasi-cohérent « usuel » sur le schéma
$T_B$ ; voir le
lemme \ref{adequate-lemma-same-cohomology-adequate}.)
Assertion : le foncteur
$$
B \longmapsto H^q(T_B, \mathcal{F}_B)
$$
est adéquat. Nous démontrerons le lemme par la méthode usuelle
consistant à découper $T$ en morceaux.

\medskip\noindent
Cas I : $T$ est affine. Dans ce cas, les schémas $T_B$ sont tous affines
et $H^q(T_B, \mathcal{F}_B) = 0$ pour tout $q \geq 1$.
Le foncteur $B \mapsto H^0(T_B, \mathcal{F}_B)$ est adéquat d'après le
lemme \ref{adequate-lemma-pushforward-adequate}.

\medskip\noindent
Cas II : $T$ est séparé. Soit $n$ le nombre minimal d'ouverts affines nécessaires
pour recouvrir $T$. Raisonnons par récurrence sur $n$. Le cas initial est le cas I.
Choisissons un recouvrement ouvert affine $T = V_1 \cup \ldots \cup V_n$.
Posons $V = V_1 \cup \ldots \cup V_{n - 1}$ et $U = V_n$. Observons que
$$
U \cap V = (V_1 \cap V_n) \cup \ldots \cup (V_{n - 1} \cap V_n)
$$
est également une réunion de $n - 1$ ouverts affines, puisque $T$ est séparé ; voir
Schémas, lemme \ref{schemes-lemma-characterize-separated}.
Notons que, pour chaque $B$, les changements de base $U_B$, $V_B$ et
$(U \cap V)_B = U_B \cap V_B$ se comportent de la même manière. Ainsi,
pour chaque $B$, on a une suite exacte longue
$$
0 \to
H^0(T_B, \mathcal{F}_B) \to
H^0(U_B, \mathcal{F}_B) \oplus H^0(V_B, \mathcal{F}_B) \to
H^0((U \cap V)_B, \mathcal{F}_B) \to
H^1(T_B, \mathcal{F}_B) \to \ldots
$$
fonctorielle en $B$ ; voir
Cohomologie, lemme \ref{cohomology-lemma-mayer-vietoris}.
Par l'hypothèse de récurrence, les foncteurs
$B \mapsto H^q(U_B, \mathcal{F}_B)$,
$B \mapsto H^q(V_B, \mathcal{F}_B)$ et
$B \mapsto H^q((U \cap V)_B, \mathcal{F}_B)$
sont adéquats. Les
lemmes \ref{adequate-lemma-kernel-adequate} et
\ref{adequate-lemma-cokernel-adequate}
montrent que notre foncteur $B \mapsto H^q(T_B, \mathcal{F}_B)$ est le
terme médian d'une suite exacte courte dont les termes extrêmes sont adéquats.
L'assertion résulte donc du
lemme \ref{adequate-lemma-extension-adequate}.

\medskip\noindent
Cas III : cas général quasi-compact et quasi-séparé.
La démonstration se fait encore par récurrence sur le nombre $n$ d'ouverts affines
nécessaires pour recouvrir $T$. Le cas initial $n = 1$ est le cas I.
Choisissons un recouvrement ouvert affine $T = V_1 \cup \ldots \cup V_n$.
Posons $V = V_1 \cup \ldots \cup V_{n - 1}$ et $U = V_n$. Notons que,
puisque $T$ est quasi-séparé, $U \cap V$ est un ouvert quasi-compact d'un
schéma affine ; le cas II s'y applique donc. La suite de l'argument
se déroule exactement comme dans le paragraphe précédent, et nous
l'omettons.
\end{proof}







\section{Modules adéquats parasites}
\label{adequate-section-parasitic}

\noindent
Dans cette section, nous commençons à comparer les modules adéquats et les modules
quasi-cohérents sur un schéma $S$. Rappelons qu'il existe des foncteurs
$u : \QCoh(\mathcal{O}_S) \to \textit{Adeq}(\mathcal{O})$
et
$v : \textit{Adeq}(\mathcal{O}) \to \QCoh(\mathcal{O}_S)$
satisfaisant à l'adjonction
$$
\SheafHom_{\QCoh(\mathcal{O}_S)}(\mathcal{F}, v\mathcal{G})
=
\SheafHom_{\textit{Adeq}(\mathcal{O})}(u\mathcal{F}, \mathcal{G})
$$
et tels que $\mathcal{F} \to vu\mathcal{F}$ soit un isomorphisme pour
tout faisceau quasi-cohérent $\mathcal{F}$ ; voir la
remarque \ref{adequate-remark-compare}.
Ainsi $u$ est un plongement pleinement fidèle, et l'on peut identifier
$\QCoh(\mathcal{O}_S)$ à une sous-catégorie pleine de
$\textit{Adeq}(\mathcal{O})$.
Le foncteur $v$ est exact, mais $u$ n'est pas exact à gauche en général.
Le noyau de $v$ est la sous-catégorie des modules adéquats parasites.

\medskip\noindent
Dans Descente, définition \ref{descent-definition-parasitic},
nous donnons la définition d'un module parasite.
Pour les modules adéquats, cette notion ne dépend pas
de la topologie choisie.

\begin{lemma}
\label{adequate-lemma-parasitic-adequate}
Soit $S$ un schéma.
Soit $\mathcal{F}$ un $\mathcal{O}$-module adéquat sur
$(\Sch/S)_\tau$. Les conditions suivantes sont équivalentes :
\begin{enumerate}
\item $v\mathcal{F} = 0$,
\item $\mathcal{F}$ est parasite,
\item $\mathcal{F}$ est parasite pour la topologie $\tau$,
\item $\mathcal{F}(U) = 0$ pour tout ouvert $U \subset S$, et
\item il existe un recouvrement ouvert affine $S = \bigcup U_i$
tel que $\mathcal{F}(U_i) = 0$ pour tout $i$.
\end{enumerate}
\end{lemma}

\begin{proof}
Les implications (2) $\Rightarrow$ (3) $\Rightarrow$ (4) $\Rightarrow$ (5)
résultent immédiatement des définitions. Supposons (5). Supposons que
$S = \bigcup U_i$ soit un recouvrement ouvert affine tel que $\mathcal{F}(U_i) = 0$
pour tout $i$. Soit $V \to S$ un morphisme plat. Il existe un recouvrement ouvert
affine $V = \bigcup V_j$ tel que chaque $V_j$ ait son image dans un
$U_i$. Puisque le morphisme $V_j \to S$ est plat, $V_j \to U_i$ l'est aussi.
L'homomorphisme d'anneaux correspondant
$A_i = \mathcal{O}(U_i) \to \mathcal{O}(V_j) = B_j$ est donc plat. Ainsi, d'après le
lemme \ref{adequate-lemma-adequate-local}
et le
lemme \ref{adequate-lemma-adequate-flat},
on voit que $\mathcal{F}(U_i) \otimes_{A_i} B_j \to \mathcal{F}(V_j)$
est un isomorphisme. Ainsi $\mathcal{F}(V_j) = 0$. Puisque $\mathcal{F}$ est
un faisceau pour la topologie de Zariski, on conclut que $\mathcal{F}(V) = 0$.
On voit de cette manière que (5) entraîne (2).

\medskip\noindent
Cela démontre l'équivalence de (2), (3), (4) et (5).
Comme (1) est équivalente à (3) (voir la
remarque \ref{adequate-remark-compare}),
on conclut que les cinq conditions sont équivalentes.
\end{proof}

\noindent
Soit $S$ un schéma.
La sous-catégorie des modules adéquats parasites est une sous-catégorie de Serre de
$\textit{Adeq}(\mathcal{O})$. Le quotient est la catégorie des
modules quasi-cohérents.

\begin{lemma}
\label{adequate-lemma-adequate-by-parasitic}
Soit $S$ un schéma. La sous-catégorie
$\mathcal{C} \subset \textit{Adeq}(\mathcal{O})$ des modules adéquats
parasites est une sous-catégorie de Serre. De plus, le foncteur $v$ induit
une équivalence de catégories
$$
\textit{Adeq}(\mathcal{O}) / \mathcal{C} = \QCoh(\mathcal{O}_S).
$$
\end{lemma}

\begin{proof}
La catégorie $\mathcal{C}$ est le noyau du foncteur exact
$v : \textit{Adeq}(\mathcal{O}) \to \QCoh(\mathcal{O}_S)$ ; voir le
lemme \ref{adequate-lemma-parasitic-adequate}.
C'est donc une sous-catégorie de Serre d'après
Homologie, lemme \ref{homology-lemma-kernel-exact-functor}.
D'après
Homologie, lemme \ref{homology-lemma-serre-subcategory-is-kernel},
on obtient un foncteur exact induit
$\overline{v} :
\textit{Adeq}(\mathcal{O}) / \mathcal{C}
\to
\QCoh(\mathcal{O}_S)$.
Puisque $u$ est un inverse à droite de $v$, on voit immédiatement que
$\overline{v}$ est essentiellement surjectif.
On voit que $\overline{v}$ est fidèle d'après
Homologie, lemme \ref{homology-lemma-quotient-by-kernel-exact-functor}.
Puisque $u$ est un inverse à droite de $v$, on conclut enfin que
$\overline{v}$ est pleinement fidèle.
\end{proof}

\begin{lemma}
\label{adequate-lemma-direct-image-parasitic-adequate}
Soit $f : T \to S$ un morphisme quasi-compact et quasi-séparé
de schémas. Pour tout $\mathcal{O}_T$-module adéquat parasite sur
$(\Sch/T)_\tau$, l'image directe
$f_*\mathcal{F}$ et les images directes supérieures $R^if_*\mathcal{F}$
sont des $\mathcal{O}_S$-modules adéquats parasites sur $(\Sch/S)_\tau$.
\end{lemma}

\begin{proof}
Nous avons déjà vu dans le
lemme \ref{adequate-lemma-direct-image-adequate}
que ces images directes supérieures sont adéquates.
Il suffit donc de montrer que
$(R^if_*\mathcal{F})(U_i) = 0$ pour tout $\tau$-recouvrement
ouvert $\{U_i \to S\}$. Or $R^if_*\mathcal{F}$
est parasite d'après
Descente, lemme \ref{descent-lemma-direct-image-parasitic}.
\end{proof}













\section{Catégories dérivées de modules adéquats, I}
\label{adequate-section-comparison}


\noindent
Soit $S$ un schéma. Nous poursuivons la discussion commencée dans la
section \ref{adequate-section-parasitic}.
Le foncteur exact $v$ induit un foncteur
$$
D(\textit{Adeq}(\mathcal{O}))
\longrightarrow
D(\QCoh(\mathcal{O}_S))
$$
et de même pour les variantes bornées.

\begin{lemma}
\label{adequate-lemma-quotient-easy}
Soit $S$ un schéma. Notons
$\mathcal{C} \subset \textit{Adeq}(\mathcal{O})$ la
sous-catégorie pleine constituée des modules adéquats parasites.
Alors
$$
D(\textit{Adeq}(\mathcal{O}))/D_\mathcal{C}(\textit{Adeq}(\mathcal{O}))
= D(\QCoh(\mathcal{O}_S))
$$
et de même pour les variantes bornées.
\end{lemma}

\begin{proof}
{\emergencystretch=1em
Cela résulte immédiatement de
Catégories dérivées, lemme \ref{derived-lemma-quotient-by-serre-easy}.\par}
\end{proof}

\noindent
Cherchons ensuite une description en sens inverse en considérant
les foncteurs
$$
K^+(\QCoh(\mathcal{O}_S))
\longrightarrow
K^+(\textit{Adeq}(\mathcal{O}))
\longrightarrow
D^+(\textit{Adeq}(\mathcal{O}))
\longrightarrow
D^+(\QCoh(\mathcal{O}_S)).
$$
Dans certains cas, la catégorie dérivée des modules adéquats est une localisation
de la catégorie homotopique des complexes de modules quasi-cohérents par rapport
aux quasi-isomorphismes universels. Soit $S$ un schéma. Une application de complexes
$\varphi : \mathcal{F}^\bullet \to \mathcal{G}^\bullet$
de $\mathcal{O}_S$-modules quasi-cohérents est appelée
{\it quasi-isomorphisme universel} si, pour tout morphisme de schémas
$f : T \to S$, l'image inverse $f^*\varphi$ est un quasi-isomorphisme.

\begin{lemma}
\label{adequate-lemma-describe-Dplus-adequate}
Soit $U = \Spec(A)$ un schéma affine.
La catégorie dérivée bornée inférieurement
$D^+(\textit{Adeq}(\mathcal{O}))$ est la localisation
de $K^+(\QCoh(\mathcal{O}_U))$ par rapport au système multiplicatif des
quasi-isomorphismes universels.
\end{lemma}

\begin{proof}
{\emergencystretch=2em
Si $\varphi : \mathcal{F}^\bullet \to \mathcal{G}^\bullet$
est un morphisme de complexes de
$\mathcal{O}_U$-modules quasi-cohérents, alors $u\varphi : u\mathcal{F}^\bullet \to
u\mathcal{G}^\bullet$ est un quasi-isomorphisme si et seulement si $\varphi$ est
un quasi-isomorphisme universel. Ainsi la collection $S$
des quasi-isomorphismes universels est un système multiplicatif
saturé compatible avec la structure triangulée d'après
Catégories dérivées, lemme \ref{derived-lemma-triangle-functor-localize}.
Ainsi $S^{-1}K^+(\QCoh(\mathcal{O}_U))$ existe et est une
catégorie triangulée ; voir
Catégories dérivées, proposition
\ref{derived-proposition-construct-localization}.
On obtient un foncteur canonique
$can : S^{-1}K^+(\QCoh(\mathcal{O}_U)) \to
D^{+}(\textit{Adeq}(\mathcal{O}))$ d'après
Catégories dérivées, lemme \ref{derived-lemma-universal-property-localization}.
\par}

\medskip\noindent
Notons que, presque par définition, tout module adéquat sur $U$ se plonge
dans un faisceau quasi-cohérent ; voir le
lemme \ref{adequate-lemma-adequate-characterize}. Ainsi, d'après
Catégories dérivées, lemme \ref{derived-lemma-subcategory-right-resolution},
étant donné $\mathcal{F}^\bullet \in \Ob(K^+(\textit{Adeq}(\mathcal{O})))$,
il existe un quasi-isomorphisme
$\mathcal{F}^\bullet \to u\mathcal{G}^\bullet$,
où $\mathcal{G}^\bullet \in \Ob(K^+(\QCoh(\mathcal{O}_U)))$.
Cela démontre que $can$ est essentiellement surjectif.

\medskip\noindent
De même, supposons que $\mathcal{F}^\bullet$ et $\mathcal{G}^\bullet$
soient des complexes bornés inférieurement de $\mathcal{O}_U$-modules quasi-cohérents.
Un morphisme dans $D^+(\textit{Adeq}(\mathcal{O}))$ entre ces
complexes est constitué d'un couple $f : u\mathcal{F}^\bullet \to \mathcal{H}^\bullet$
et $s : u\mathcal{G}^\bullet \to \mathcal{H}^\bullet$, où $s$
est un quasi-isomorphisme. Choisissons un quasi-isomorphisme
$s' : \mathcal{H}^\bullet \to u\mathcal{E}^\bullet$. On voit alors que
$s' \circ f : \mathcal{F} \to \mathcal{E}^\bullet$ et le
quasi-isomorphisme universel
$s' \circ s : \mathcal{G}^\bullet \to \mathcal{E}^\bullet$ donnent
un morphisme dans $S^{-1}K^{+}(\QCoh(\mathcal{O}_U))$ ayant pour image
le morphisme donné. Cela démontre la plénitude.
La fidélité se démontre de manière analogue.
\end{proof}

\begin{lemma}
\label{adequate-lemma-right-adjoint-adequate}
Soit $U = \Spec(A)$ un schéma affine.
Le foncteur d'inclusion
$$
\textit{Adeq}(\mathcal{O}) \to
\textit{Mod}((\Sch/U)_\tau, \mathcal{O})
$$
admet un adjoint à droite $A$\footnote{C'est l'« adéquateur ».}.
De plus, le morphisme d'adjonction
$A(\mathcal{F}) \to \mathcal{F}$ est un isomorphisme pour tout
module adéquat $\mathcal{F}$.
\end{lemma}

\begin{proof}
D'après
Topologies, lemme \ref{topologies-lemma-affine-big-site-fppf}
(et de même pour les autres topologies),
on peut travailler avec les $\mathcal{O}$-modules sur $(\textit{Aff}/U)_\tau$.
Notons $\mathcal{P}$ la catégorie des foncteurs à valeurs
dans les modules sur $\textit{Alg}_A$ et $\mathcal{A}$ la catégorie des foncteurs adéquats
sur $\textit{Alg}_A$. Notons $i : \mathcal{A} \to \mathcal{P}$
le foncteur d'inclusion. Notons $Q : \mathcal{P} \to \mathcal{A}$
la construction du lemme \ref{adequate-lemma-adjoint}.
On a le diagramme commutatif
\begin{equation}
\label{adequate-equation-categories}
\vcenter{
\xymatrix{
\textit{Adeq}(\mathcal{O}) \ar[r]_-k \ar@{=}[d] &
\textit{Mod}((\textit{Aff}/U)_\tau, \mathcal{O}) \ar[r]_-j &
\textit{PMod}((\textit{Aff}/U)_\tau, \mathcal{O}) \ar@{=}[d] \\
\mathcal{A} \ar[rr]^-i & & \mathcal{P}
}
}
\end{equation}
L'égalité verticale de gauche est le
lemme \ref{adequate-lemma-adequate-affine},
et l'égalité verticale de droite a été expliquée dans la
section \ref{adequate-section-quasi-coherent}.
Définissons $A(\mathcal{F}) = Q(j(\mathcal{F}))$.
Puisque $j$ est pleinement fidèle, il en résulte immédiatement que $A$
est adjoint à droite du foncteur d'inclusion $k$. De plus, puisque
$k$ est lui aussi pleinement fidèle, la dernière assertion en résulte formellement.
\end{proof}

\noindent
Le foncteur $A$ est un adjoint à droite, donc est exact à gauche. Puisque le foncteur
d'inclusion est exact, voir le
lemme \ref{adequate-lemma-abelian-adequate},
on conclut que $A$ transforme les objets injectifs en objets injectifs et que
la catégorie $\textit{Adeq}(\mathcal{O})$ possède suffisamment d'injectifs ; voir
Homologie, lemme \ref{homology-lemma-adjoint-enough-injectives},
et
Injectifs, théorème \ref{injectives-theorem-sheaves-modules-injectives}.
Cela résulte également de l'équivalence dans
(\ref{adequate-equation-categories})
et du
lemme \ref{adequate-lemma-enough-injectives}.

\begin{lemma}
\label{adequate-lemma-RA-zero}
Soit $U = \Spec(A)$ un schéma affine.
Pour tout objet $\mathcal{F}$ de $\textit{Adeq}(\mathcal{O})$,
on a $R^pA(\mathcal{F}) = 0$ pour tout $p > 0$, où $A$ est
comme dans le
lemme \ref{adequate-lemma-right-adjoint-adequate}.
\end{lemma}

\begin{proof}
Avec les notations de la démonstration du
lemme \ref{adequate-lemma-right-adjoint-adequate},
choisissons une résolution injective $k(\mathcal{F}) \to \mathcal{I}^\bullet$
dans la catégorie des $\mathcal{O}$-modules sur $(\textit{Aff}/U)_\tau$.
D'après
Cohomologie sur les sites, lemme \ref{sites-cohomology-lemma-include-modules},
et le
lemme \ref{adequate-lemma-same-cohomology-adequate},
le complexe $j(\mathcal{I}^\bullet)$ est exact.
D'autre part, chaque $j(\mathcal{I}^n)$ est un objet injectif
de la catégorie des préfaisceaux de modules d'après
Cohomologie sur les sites, lemme
\ref{sites-cohomology-lemma-injective-module-injective-presheaf}.
Il en résulte que $R^pA(\mathcal{F}) = R^pQ(j(k(\mathcal{F})))$.
Le résultat découle donc du
lemme \ref{adequate-lemma-RQ-zero}.
\end{proof}

\noindent
Soit $S$ un schéma. D'après la discussion de la
section \ref{adequate-section-adequate},
le plongement
$\textit{Adeq}(\mathcal{O}) \subset
\textit{Mod}((\Sch/S)_\tau, \mathcal{O})$
présente $\textit{Adeq}(\mathcal{O})$ comme une sous-catégorie de Serre faible de
la catégorie de tous les $\mathcal{O}$-modules. Notons
$$
D_{\textit{Adeq}}(\mathcal{O}) \subset
D(\mathcal{O}) = D(\textit{Mod}((\Sch/S)_\tau, \mathcal{O}))
$$
la sous-catégorie triangulée des complexes dont les faisceaux de cohomologie
sont adéquats ; voir
Catégories dérivées, section \ref{derived-section-triangulated-sub}.
On obtient un foncteur canonique
$$
D(\textit{Adeq}(\mathcal{O}))
\longrightarrow
D_{\textit{Adeq}}(\mathcal{O})
$$
voir
Catégories dérivées, équation (\ref{derived-equation-compare}).

\begin{lemma}
\label{adequate-lemma-bounded-below}
Si $U = \Spec(A)$ est un schéma affine, alors la variante
bornée inférieurement
\begin{equation}
\label{adequate-equation-compare-bounded-adequate}
D^+(\textit{Adeq}(\mathcal{O}))
\longrightarrow
D^+_{\textit{Adeq}}(\mathcal{O})
\end{equation}
du foncteur ci-dessus est une équivalence.
\end{lemma}

\begin{proof}
Soit $A : \textit{Mod}(\mathcal{O}) \to \textit{Adeq}(\mathcal{O})$
l'adjoint à droite du foncteur d'inclusion construit dans le
lemme \ref{adequate-lemma-right-adjoint-adequate}.
Puisque $A$ est exact à gauche et que $\textit{Mod}(\mathcal{O})$
possède suffisamment d'injectifs, $A$ admet un foncteur dérivé à droite
$RA : D^+_{\textit{Adeq}}(\mathcal{O}) \to D^+(\textit{Adeq}(\mathcal{O}))$.
Nous affirmons que $RA$ est un quasi-inverse de
(\ref{adequate-equation-compare-bounded-adequate}).
Le fait essentiel est que, si $\mathcal{F}$ est un module adéquat,
le morphisme d'adjonction $\mathcal{F} \to RA(\mathcal{F})$ est un
quasi-isomorphisme d'après le lemme \ref{adequate-lemma-RA-zero}.

\medskip\noindent
Pour démontrer complètement le lemme, il suffit en effet de montrer que :
\begin{enumerate}
\item étant donné $\mathcal{F}^\bullet \in K^+(\textit{Adeq}(\mathcal{O}))$,
l'application canonique $\mathcal{F}^\bullet \to RA(\mathcal{F}^\bullet)$
est un quasi-isomorphisme, et
\item étant donné $\mathcal{G}^\bullet \in K^+(\textit{Mod}(\mathcal{O}))$,
l'application canonique $RA(\mathcal{G}^\bullet) \to \mathcal{G}^\bullet$
est un quasi-isomorphisme.
\end{enumerate}
Les assertions (1) et (2) résultent toutes deux du fait essentiel par un argument
de suite spectrale utilisant l'une des suites spectrales de
Catégories dérivées, lemme \ref{derived-lemma-two-ss-complex-functor}.
Nous omettons certains détails.
\end{proof}

\begin{lemma}
\label{adequate-lemma-ext-adequate}
Soit $U = \Spec(A)$ un schéma affine.
Soient $\mathcal{F}$ et $\mathcal{G}$ des $\mathcal{O}$-modules adéquats.
Pour tout $i \geq 0$, l'application naturelle
$$
\Ext^i_{\textit{Adeq}(\mathcal{O})}(\mathcal{F}, \mathcal{G})
\longrightarrow
\Ext^i_{\textit{Mod}(\mathcal{O})}(\mathcal{F}, \mathcal{G})
$$
est un isomorphisme.
\end{lemma}

\begin{proof}
Par définition, ces groupes Ext se calculent comme ensembles de morphismes dans
la catégorie dérivée. Cela résulte donc immédiatement du
lemme \ref{adequate-lemma-bounded-below}.
\end{proof}









\section{Extensions pures}
\label{adequate-section-pure}

\noindent
Nous voulons caractériser en termes algébriques les extensions de faisceaux
quasi-cohérents sur le grand site d'un schéma affine. Pour cela,
nous introduisons la notion suivante.

\begin{definition}
\label{adequate-definition-pure}
Soit $A$ un anneau.
\begin{enumerate}
\item Un $A$-module $P$ est dit {\it projectif pur}
si, pour toute suite universellement exacte
$0 \to K \to M \to N \to 0$ de $A$-modules, la suite
$0 \to \Hom_A(P, K) \to \Hom_A(P, M) \to \Hom_A(P, N) \to 0$
est exacte.
\item Un $A$-module $I$ est dit {\it injectif pur}
si, pour toute suite universellement exacte
$0 \to K \to M \to N \to 0$ de $A$-modules, la suite
$0 \to \Hom_A(N, I) \to \Hom_A(M, I) \to \Hom_A(K, I) \to 0$
est exacte.
\end{enumerate}
\end{definition}

\noindent
Caractérisons les modules projectifs purs.

\begin{lemma}
\label{adequate-lemma-pure-projective}
Soit $A$ un anneau.
\begin{enumerate}
\item Un module est projectif pur si et seulement s'il est
facteur direct d'une somme directe de $A$-modules de présentation finie.
\item Pour tout module $M$, il existe une suite universellement exacte
$0 \to N \to P \to M \to 0$, avec $P$ projectif pur.
\end{enumerate}
\end{lemma}

\begin{proof}
Notons d'abord qu'un $A$-module de présentation finie est projectif pur d'après
Algèbre, théorème \ref{algebra-theorem-universally-exact-criteria}.
Ainsi un facteur direct d'une somme directe de $A$-modules de présentation finie
est bien projectif pur. Soit $M$ un $A$-module quelconque.
Écrivons $M = \colim_{i \in I} P_i$ comme limite inductive filtrante de
$A$-modules de présentation finie. Considérons la suite
$$
0 \to N \to \bigoplus P_i \to M \to 0.
$$
Pour tout $A$-module de présentation finie $P$, l'application
$\Hom_A(P, \bigoplus P_i) \to \Hom_A(P, M)$
est surjective, car toute application $P \to M$ se factorise par un certain $P_i$.
Ainsi, d'après
Algèbre, théorème \ref{algebra-theorem-universally-exact-criteria},
cette suite est universellement exacte. Cela démontre (2).
Si maintenant $M$ est projectif pur, alors la suite est scindée et
l'on voit que $M$ est facteur direct de $\bigoplus P_i$.
\end{proof}

\noindent
Caractérisons les modules injectifs purs.

\begin{lemma}
\label{adequate-lemma-pure-injective}
{\emergencystretch=2em
Soit $A$ un anneau. Pour tout $A$-module $M$, posons
$M^\vee = \Hom_\mathbf{Z}(M, \mathbf{Q}/\mathbf{Z})$.
\par}
\begin{enumerate}
\item Pour tout $A$-module $M$, le $A$-module $M^\vee$ est injectif pur.
\item Un $A$-module $I$ est injectif pur si et seulement si l'application
$I \to (I^\vee)^\vee$ est scindée.
\item Pour tout module $M$, il existe une suite universellement exacte
$0 \to M \to I \to N \to 0$, avec $I$ injectif pur.
\end{enumerate}
\end{lemma}

\begin{proof}
Nous utiliserons les propriétés du foncteur $M \mapsto M^\vee$ données dans
Compléments d'algèbre, section \ref{more-algebra-section-injectives-modules},
sans le mentionner à nouveau. L'assertion (1) résulte de ce que
$\Hom_A(N, M^\vee) =
\Hom_\mathbf{Z}(N \otimes_A M, \mathbf{Q}/\mathbf{Z})$
et de ce que $\mathbf{Q}/\mathbf{Z}$ est injectif dans la catégorie des
groupes abéliens. Ainsi, si $I \to (I^\vee)^\vee$ est scindée,
$I$ est injectif pur. Nous affirmons que, pour tout $A$-module $M$,
l'application d'évaluation $ev : M \to (M^\vee)^\vee$ est universellement injective.
Pour le voir, notons que $ev^\vee : ((M^\vee)^\vee)^\vee \to M^\vee$
admet un inverse à droite, à savoir $ev' : M^\vee \to ((M^\vee)^\vee)^\vee$.
Alors, pour tout $A$-module $N$, l'application du foncteur exact et fidèle
${}^\vee$ au morphisme $N \otimes_A M \to N \otimes_A (M^\vee)^\vee$
fournit
$$
\Hom_A(N, ((M^\vee)^\vee)^\vee) =
\Big(N \otimes_A (M^\vee)^\vee\Big)^\vee
\to
\Big(N \otimes_A M\Big)^\vee =
\Hom_A(N, M^\vee)
$$
qui est surjective en raison de l'existence de l'inverse à droite. D'où l'assertion.
Celle-ci entraîne (3) et la nécessité de la condition de (2).
\end{proof}

\noindent
Avant de poursuivre, faisons l'observation suivante, que nous utiliserons
fréquemment dans la suite de cette section.

\begin{lemma}
\label{adequate-lemma-split-universally-exact-sequence}
Soit $A$ un anneau.
\begin{enumerate}
\item Soit $L \to M \to N$ une suite universellement exacte
de $A$-modules. Posons $K = \Im(M \to N)$.
Alors $K \to N$ est universellement injective.
\item Tout complexe universellement exact
peut être décomposé en suites exactes courtes universellement exactes.
\end{enumerate}
\end{lemma}

\begin{proof}
Démonstration de (1). Pour tout $A$-module $T$, la suite
$L \otimes_A T \to M \otimes_A T \to K \otimes_A T \to 0$ est exacte
par exactitude à droite de $\otimes$. Par hypothèse, la suite
$L \otimes_A T \to M \otimes_A T \to N \otimes_A T$ est exacte.
Ces deux faits montrent que $K \otimes_A T \to N \otimes_A T$ est injective.

\medskip\noindent
L'assertion (2) signifie ceci : supposons que $M^\bullet$ soit un complexe
universellement exact de $A$-modules. Posons $K^i = \Ker(d^i) \subset M^i$.
Alors les suites exactes courtes $0 \to K^i \to M^i \to K^{i + 1} \to 0$
sont universellement exactes. Cela résulte immédiatement de (1).
\end{proof}

\begin{definition}
\label{adequate-definition-pure-resolution}
Soit $A$ un anneau. Soit $M$ un $A$-module.
\begin{enumerate}
\item Une {\it résolution par des projectifs purs} $P_\bullet \to M$
est une suite universellement exacte
$$
\ldots \to P_1 \to P_0 \to M \to 0
$$
où chaque $P_i$ est projectif pur.
\item Une {\it résolution par des injectifs purs} $M \to I^\bullet$ est une suite
universellement exacte
$$
0 \to M \to I^0 \to I^1 \to \ldots
$$
où chaque $I^i$ est injectif pur.
\end{enumerate}
\end{definition}

\noindent
Ces résolutions satisfont aux propriétés d'unicité usuelles parmi la classe
de toutes les résolutions à gauche ou à droite universellement exactes.

\begin{lemma}
\label{adequate-lemma-pure-projective-resolutions}
Soit $A$ un anneau.
\begin{enumerate}
\item Tout $A$-module admet une résolution par des projectifs purs.
\end{enumerate}
Soit $M \to N$ une application de $A$-modules.
Soit $P_\bullet \to M$ une résolution par des projectifs purs et
soit $N_\bullet \to N$ une résolution universellement exacte.
\begin{enumerate}
\item[(2)] Il existe une application de complexes $P_\bullet \to N_\bullet$
induisant l'application donnée
$$
M = \Coker(P_1 \to P_0) \to \Coker(N_1 \to N_0) = N
$$
\item[(3)] deux applications $\alpha, \beta : P_\bullet \to N_\bullet$
induisant la même application $M \to N$ sont homotopes.
\end{enumerate}
\end{lemma}

\begin{proof}
L'assertion (1) résulte immédiatement du
lemme \ref{adequate-lemma-pure-projective}.
Avant de démontrer (2) et (3), notons que, d'après le
lemme \ref{adequate-lemma-split-universally-exact-sequence},
on peut décomposer le complexe universellement exact $N_\bullet \to N \to 0$
en suites exactes courtes universellement exactes $0 \to K_0 \to N_0 \to N \to 0$
et $0 \to K_i \to N_i \to K_{i - 1} \to 0$.

\medskip\noindent
Démonstration de (2). Puisque $P_0$ est projectif pur,
on peut trouver une application $P_0 \to N_0$ relevant l'application $P_0 \to M \to N$.
On obtient une application induite $P_1 \to F_0 \to N_0$ dont l'image est contenue dans $K_0$.
Puisque $P_1$ est projectif pur, on peut la relever
en une application $P_1 \to N_1$. Celle-ci induit à son tour une application
$P_2 \to P_1 \to N_1$ dont l'image dans
$N_0$ est nulle, c'est-à-dire dont l'image est contenue dans $K_1$. On peut donc la relever en une application
$P_2 \to N_2$. On répète le procédé.

\medskip\noindent
Démonstration de (3). Pour montrer que $\alpha, \beta$ sont homotopes, il suffit
de montrer que la différence $\gamma = \alpha - \beta$ est homotope
à zéro. Notons que l'image de $\gamma_0 : P_0 \to N_0$
est contenue dans $K_0$. On peut donc relever
$\gamma_0$ en une application $h_0 : P_0 \to N_1$. Considérons l'application
$\gamma_1' = \gamma_1 - h_0 \circ d_{P, 1} : P_1 \to N_1$.
D'après le choix de $h_0$, on voit que l'image de $\gamma_1'$
est contenue dans $K_1$. Puisque $P_1$ est projectif pur, on peut relever
$\gamma_1'$ en une application $h_1 : P_1 \to N_2$. À ce stade, on a
$\gamma_1 = h_0 \circ d_{F, 1} + d_{N, 2} \circ h_1$. On répète le procédé.
\end{proof}

\begin{lemma}
\label{adequate-lemma-pure-injective-resolutions}
Soit $A$ un anneau.
\begin{enumerate}
\item Tout $A$-module admet une résolution par des injectifs purs.
\end{enumerate}
Soit $M \to N$ une application de $A$-modules.
Soit $M \to M^\bullet$ une résolution universellement exacte et
soit $N \to I^\bullet$ une résolution par des injectifs purs.
\begin{enumerate}
\item[(2)] Il existe une application de complexes $M^\bullet \to I^\bullet$
induisant l'application donnée
$$
M = \Ker(M^0 \to M^1) \to \Ker(I^0 \to I^1) = N
$$
\item[(3)] deux applications $\alpha, \beta : M^\bullet \to I^\bullet$
induisant la même application $M \to N$ sont homotopes.
\end{enumerate}
\end{lemma}

\begin{proof}
Ce lemme est dual du
lemme \ref{adequate-lemma-pure-projective-resolutions}.
La démonstration est identique, sauf qu'il faut renverser toutes les flèches.
\end{proof}

\noindent
Les résultats précédents permettent de définir les groupes d'extensions pures
comme suit. Soient $A$ un anneau et $M$, $N$ des $A$-modules.
Choisissons une résolution par des injectifs purs $N \to I^\bullet$. D'après le
lemme \ref{adequate-lemma-pure-injective-resolutions},
le complexe
$$
\Hom_A(M, I^\bullet)
$$
est bien défini à homotopie près. Son $i$-ième module de cohomologie
est donc un invariant bien défini de $M$ et $N$.

\begin{definition}
\label{adequate-definition-pure-ext}
Soient $A$ un anneau et $M$, $N$ des $A$-modules.
Le $i$-ième {\it module d'extensions pures} $\text{Pext}^i_A(M, N)$
est le $i$-ième module de cohomologie du complexe
$\Hom_A(M, I^\bullet)$, où $I^\bullet$ est une résolution par des injectifs purs
de $N$.
\end{definition}

\noindent
Attention : il n'est pas vrai qu'une suite exacte de $A$-modules donne
lieu à une suite exacte longue de groupes d'extensions pures. (Il faut
pour cela une suite universellement exacte.)
Rassemblons quelques faits qui résultent immédiatement de ce qui précède.

\begin{lemma}
\label{adequate-lemma-facts-pext}
Soit $A$ un anneau.
\begin{enumerate}
\item $\text{Pext}^i_A(M, N) = 0$ pour $i > 0$ dès que $N$ est injectif pur,
\item $\text{Pext}^i_A(M, N) = 0$ pour $i > 0$ dès que $M$ est projectif pur,
en particulier si $M$ est un $A$-module de présentation finie,
\item {\emergencystretch=3em
$\text{Pext}^i_A(M, N)$ est aussi le $i$-ième module de cohomologie
du complexe $\Hom_A(P_\bullet, N)$, où $P_\bullet$
est une résolution par des projectifs purs de $M$.\par}
\end{enumerate}
\end{lemma}

\begin{proof}
Pour (3), considérons le complexe double
$$
A^{\bullet, \bullet} = \Hom_A(P_\bullet, I^\bullet)
$$
{\emergencystretch=3em
Chacune de ses lignes est exacte sauf en degré $0$, où sa cohomologie
est $\Hom_A(M, I^q)$. Chacune de ses colonnes est exacte sauf en degré $0$,
où sa cohomologie est $\Hom_A(P_p, N)$. Ainsi les deux suites spectrales
associées à ce complexe dans
Homologie, section \ref{homology-section-double-complex},
dégénèrent, ce qui donne l'égalité.\par}
\end{proof}





\section{Groupes Ext supérieurs des faisceaux quasi-cohérents sur le grand site}
\label{adequate-section-big}

\noindent
Il se trouve que le foncteur à valeurs dans les modules $\underline{I}$ associé à
un module injectif pur $I$ donne un objet injectif dans la
catégorie des foncteurs adéquats sur $\textit{Alg}_A$.
Attention : il n'est pas vrai qu'un module projectif pur donne
un objet projectif dans la catégorie des foncteurs adéquats. Nous disposons néanmoins
de nombreux objets projectifs, à savoir les foncteurs linéairement adéquats.

\begin{lemma}
\label{adequate-lemma-pure-injective-injective-adequate}
Soit $A$ un anneau.
Soit $\mathcal{A}$ la catégorie des foncteurs adéquats sur $\textit{Alg}_A$.
Les objets injectifs de $\mathcal{A}$ sont exactement les foncteurs
$\underline{I}$, où $I$ est un $A$-module injectif pur.
\end{lemma}

\begin{proof}
Soit $I$ un objet injectif de $\mathcal{A}$.
Choisissons un plongement $I \to \underline{M}$ pour un certain $A$-module $M$.
Comme $I$ est injectif, on voit que $\underline{M} = I \oplus F$ pour un certain
foncteur à valeurs dans les modules $F$. Alors $M = I(A) \oplus F(A)$, et il en résulte
que $I = \underline{I(A)}$. On voit donc que tout objet injectif
est de la forme $\underline{I}$ pour un certain $A$-module $I$.
Il est clair que le module $I$ doit être injectif pur,
puisque toute suite universellement exacte $0 \to M \to N \to L \to 0$
donne lieu à une suite exacte
$0 \to \underline{M} \to \underline{N} \to \underline{L} \to 0$
de $\mathcal{A}$.

\medskip\noindent
Enfin, supposons que $I$ soit un
$A$-module injectif pur. Choisissons un plongement $\underline{I} \to J$
dans un objet injectif de $\mathcal{A}$ (voir le
lemme \ref{adequate-lemma-enough-injectives}).
Nous avons vu ci-dessus que $J = \underline{I'}$
pour un certain $A$-module $I'$ injectif pur. Puisque
$\underline{I} \to \underline{I'}$ est injective,
l'application $I \to I'$ est universellement injective. Par l'hypothèse sur $I$,
elle est scindée. Ainsi $\underline{I}$ est facteur direct de $J = \underline{I'}$,
donc est un objet injectif de la catégorie $\mathcal{A}$.
\end{proof}

\noindent
Soit $U = \Spec(A)$ un schéma affine. Soit $M$ un $A$-module.
Nous utiliserons la notation $M^a$ pour désigner le faisceau quasi-cohérent
de $\mathcal{O}$-modules sur $(\Sch/U)_\tau$ associé
au faisceau quasi-cohérent $\widetilde{M}$ sur $U$.
Nous disposons maintenant de toutes les notations nécessaires pour formuler le lemme suivant.

\begin{lemma}
\label{adequate-lemma-big-ext}
Soit $U = \Spec(A)$ un schéma affine. Soient $M$, $N$ des $A$-modules.
Pour tout $i$, on a un isomorphisme canonique
$$
\Ext^i_{\textit{Mod}(\mathcal{O})}(M^a, N^a) = \text{Pext}^i_A(M, N)
$$
fonctoriel en $M$ et $N$.
\end{lemma}

\begin{proof}
Construisons une flèche canonique de droite à gauche. Si
$N \to I^\bullet$ est une résolution par des injectifs purs, alors
$M^a \to (I^\bullet)^a$ est un complexe exact de
$\mathcal{O}$-modules (adéquats). Ainsi tout élément de $\text{Pext}^i_A(M, N)$
donne lieu à une application $N^a \to M^a[i]$ dans $D(\mathcal{O})$, c'est-à-dire à
un élément du groupe de gauche.

\medskip\noindent
Pour démontrer que cette application est un isomorphisme, notons que l'on peut remplacer
$\Ext^i_{\textit{Mod}(\mathcal{O})}(M^a, N^a)$ par
$\Ext^i_{\textit{Adeq}(\mathcal{O})}(M^a, N^a)$ ; voir le
lemme \ref{adequate-lemma-ext-adequate}.
Soit $\mathcal{A}$ la catégorie des foncteurs adéquats
sur $\textit{Alg}_A$. Nous avons vu que $\mathcal{A}$ est
équivalente à $\textit{Adeq}(\mathcal{O})$ ; voir le
lemme \ref{adequate-lemma-adequate-affine}, ainsi que la démonstration du
lemme \ref{adequate-lemma-right-adjoint-adequate}.
Il suffit donc maintenant de démontrer que
$$
\Ext^i_\mathcal{A}(\underline{M}, \underline{N}) =
\text{Pext}^i_A(M, N)
$$
Or cela résulte du
lemme \ref{adequate-lemma-pure-injective-injective-adequate},
puisqu'une résolution par des injectifs purs $N \to I^\bullet$ correspond exactement
à une résolution injective de $\underline{N}$ dans $\mathcal{A}$.
\end{proof}







\section{Catégories dérivées de modules adéquats, II}
\label{adequate-section-derived-categories}

\noindent
Soit $S$ un schéma. Notons $\mathcal{O}_S$ le faisceau structural de $S$
et $\mathcal{O}$ le faisceau structural du grand site $(\Sch/S)_\tau$.
Dans
Descente, remarque \ref{descent-remark-change-topologies-ringed},
nous avons construit un morphisme de sites annelés
\begin{equation}
\label{adequate-equation-compare-big-small}
f :
((\Sch/S)_\tau, \mathcal{O})
\longrightarrow
(S_{Zar}, \mathcal{O}_S).
\end{equation}
Dans les sections précédentes, nous avons vu que le foncteur
$f_* : \textit{Mod}(\mathcal{O}) \to \textit{Mod}(\mathcal{O}_S)$
transforme les faisceaux adéquats en faisceaux quasi-cohérents,
induit un foncteur exact
$v : \textit{Adeq}(\mathcal{O}) \to \QCoh(\mathcal{O}_S)$, et
qu'en fait $f_* = v$ induit une équivalence
$\textit{Adeq}(\mathcal{O})/\mathcal{C} \to \QCoh(\mathcal{O}_S)$,
où $\mathcal{C}$ est la sous-catégorie des modules adéquats parasites.
De plus, le foncteur $f^*$ transforme les modules quasi-cohérents
en modules adéquats et induit un foncteur
$u : \QCoh(\mathcal{O}_S) \to \textit{Adeq}(\mathcal{O})$
qui est adjoint à gauche de $v$.

\medskip\noindent
Il existe une relation très analogue entre
$D_{\textit{Adeq}}(\mathcal{O})$ et $D_\QCoh(S)$.
Expliquons d'abord pourquoi la catégorie $D_{\textit{Adeq}}(\mathcal{O})$
est indépendante de la topologie choisie.

\begin{remark}
\label{adequate-remark-D-adeq-independence-topology}
{\emergencystretch=3em
Soit $S$ un schéma.
Soient $\tau, \tau' \in \{Zar, \etale, smooth, syntomic, fppf\}$.
Notons $\mathcal{O}_\tau$, resp.\ $\mathcal{O}_{\tau'}$,
le faisceau structural $\mathcal{O}$ considéré comme faisceau sur
$(\Sch/S)_\tau$, resp.\ $(\Sch/S)_{\tau'}$.
Alors $D_{\textit{Adeq}}(\mathcal{O}_\tau)$ et
$D_{\textit{Adeq}}(\mathcal{O}_{\tau'})$ sont canoniquement isomorphes.
Cela résulte de Cohomologie sur les sites, lemme
\ref{sites-cohomology-lemma-compare-topologies-derived-adequate-modules}.
Supposons en effet que $\tau$ soit plus fine que la topologie $\tau'$, posons
$\mathcal{C} = (\Sch/S)_{fppf}$ et notons $\mathcal{B}$ la collection
des schémas affines au-dessus de $S$. Nous avons vérifié ci-dessus les hypothèses (1) et (2).
L'hypothèse (3) est claire, et l'hypothèse (4) résulte du
lemme \ref{adequate-lemma-same-cohomology-adequate}.\par}
\end{remark}

\begin{remark}
\label{adequate-remark-D-adeq-and-D-QCoh}
Soit $S$ un schéma. Le morphisme $f$, voir
(\ref{adequate-equation-compare-big-small}), induit
des foncteurs adjoints
$Rf_* : D_{\textit{Adeq}}(\mathcal{O}) \to D_\QCoh(S)$
et
$Lf^* : D_\QCoh(S) \to D_{\textit{Adeq}}(\mathcal{O})$.
De plus, $Rf_* Lf^* \cong \text{id}_{D_\QCoh(S)}$.

\medskip\noindent
Esquissons la démonstration. D'après la
remarque \ref{adequate-remark-D-adeq-independence-topology},
on peut supposer que la topologie $\tau$ est la topologie de Zariski.
Nous utiliserons l'existence des foncteurs dérivés totaux non bornés
$Lf^*$ et $Rf_*$ sur les $\mathcal{O}$-modules et leur
adjonction ; voir
Cohomologie sur les sites, lemme \ref{sites-cohomology-lemma-adjoint}.
Dans ce cas, $f_*$ est simplement la restriction à la sous-catégorie
$S_{Zar}$ de $(\Sch/S)_{Zar}$. Il est donc clair que
$Rf_* = f_*$ induit
$Rf_* : D_{\textit{Adeq}}(\mathcal{O}) \to D_\QCoh(S)$.
Supposons que $\mathcal{G}^\bullet$ soit un objet de
$D_\QCoh(S)$. On peut choisir un système
$\mathcal{K}_1^\bullet \to \mathcal{K}_2^\bullet \to \ldots$
de complexes bornés supérieurement de $\mathcal{O}_S$-modules plats dont les
applications de transition sont des injections scindées terme à terme, et un diagramme
$$
\xymatrix{
\mathcal{K}_1^\bullet \ar[d] \ar[r] &
\mathcal{K}_2^\bullet \ar[d] \ar[r] & \ldots \\
\tau_{\leq 1}\mathcal{G}^\bullet \ar[r] &
\tau_{\leq 2}\mathcal{G}^\bullet \ar[r] & \ldots
}
$$
ayant les propriétés (1), (2), (3) énumérées dans
Catégories dérivées, lemme \ref{derived-lemma-special-direct-system},
où $\mathcal{P}$ est la collection des $\mathcal{O}_S$-modules plats.
Alors $Lf^*\mathcal{G}^\bullet$ se calcule par
$\colim f^*\mathcal{K}_n^\bullet$ ; voir
Cohomologie sur les sites, lemmes \ref{sites-cohomology-lemma-pullback-K-flat} et
\ref{sites-cohomology-lemma-derived-base-change}
(notons que nos sites ont suffisamment de points d'après
Cohomologie étale, lemme \ref{etale-cohomology-lemma-points-fppf}).
Il faut voir que $H^i(Lf^*\mathcal{G}^\bullet) =
\colim H^i(f^*\mathcal{K}_n^\bullet)$ est adéquat pour chaque $i$. D'après le
lemme \ref{adequate-lemma-abelian-adequate},
on conclut qu'il suffit de montrer que
chaque $H^i(f^*\mathcal{K}_n^\bullet)$ est adéquat.

\medskip\noindent
Le caractère adéquat de $H^i(f^*\mathcal{K}_n^\bullet)$ est local sur $S$ ;
on peut donc supposer que $S = \Spec(A)$ est affine. Puisque $S$ est affine,
$D_\QCoh(S) = D(\QCoh(\mathcal{O}_S))$ ; voir
la discussion dans
Catégories dérivées de schémas, section
\ref{perfect-section-derived-quasi-coherent}.
Il existe donc un quasi-isomorphisme
$\mathcal{F}^\bullet \to \mathcal{K}_n^\bullet$,
où $\mathcal{F}^\bullet$ est un complexe borné supérieurement de modules
quasi-cohérents plats.
Alors $f^*\mathcal{F}^\bullet \to f^*\mathcal{K}_n^\bullet$ est un
quasi-isomorphisme, et les faisceaux de cohomologie de
$f^*\mathcal{F}^\bullet$ sont adéquats.

\medskip\noindent
La dernière assertion,
$Rf_* Lf^* \cong \text{id}_{D_\QCoh(S)}$,
résulte de la description explicite des foncteurs ci-dessus.
(En termes simples : si $\mathcal{F}$ est quasi-cohérent et si $p > 0$, alors
$L_pf^*\mathcal{F}$ est un module adéquat parasite.)
\end{remark}


\begin{remark}
\label{adequate-remark-conclusion}
La remarque \ref{adequate-remark-D-adeq-and-D-QCoh}
ci-dessus entraîne que l'on a une équivalence de catégories dérivées
$$
D_{\textit{Adeq}}(\mathcal{O})/D_\mathcal{C}(\mathcal{O})
\longrightarrow
D_\QCoh(S)
$$
où $\mathcal{C}$ est la catégorie des modules adéquats parasites.
Il est en effet clair que $D_\mathcal{C}(\mathcal{O})$ est le noyau
de $Rf_*$, d'où un foncteur comme indiqué. Pour tout objet $X$ de
$D_{\textit{Adeq}}(\mathcal{O})$, l'application $Lf^*Rf_*X \to X$ a pour image
un quasi-isomorphisme dans $D_\QCoh(S)$ ; ainsi
$Lf^*Rf_*X \to X$ est un isomorphisme dans
$D_{\textit{Adeq}}(\mathcal{O})/D_\mathcal{C}(\mathcal{O})$.
Enfin, pour des objets $X, Y$ de $D_{\textit{Adeq}}(\mathcal{O})$,
l'application
$$
Rf_* :
\Hom_{D_{\textit{Adeq}}(\mathcal{O})/D_\mathcal{C}(\mathcal{O})}(X, Y)
\to
\Hom_{D_\QCoh(S)}(Rf_*X, Rf_*Y)
$$
est bijective, puisque $Lf^*$ en donne un inverse (d'après les remarques ci-dessus).
\end{remark}













% OK, start here.
%

\chapter{Complexes dualisants}



\phantomsection
\label{dualizing-section-phantom}


\section{Introduction}
\label{dualizing-section-introduction}

\noindent
Dans ce chapitre, nous étudions les complexes dualisants en algèbre commutative.
On pourra consulter \cite{RD}.

\medskip\noindent
Nous commençons par les surjections et injections essentielles,
les couvertures projectives, les enveloppes injectives et la dualité
sur les anneaux artiniens, puis étudions les enveloppes injectives
des corps résiduels, ce qui conduit rapidement à une démonstration
de la dualité de Matlis. Voir les sections \ref{dualizing-section-essential},
\ref{dualizing-section-injective-modules}, \ref{dualizing-section-projective-cover}, \ref{dualizing-section-injective-hull}, \ref{dualizing-section-artinian} et \ref{dualizing-section-hull-residue-field},
ainsi que la proposition \ref{dualizing-proposition-matlis}.

\medskip\noindent
Viennent ensuite trois sections consacrées à la cohomologie locale
dans une grande généralité ; voir les sections \ref{dualizing-section-bad-local-cohomology},
\ref{dualizing-section-local-cohomology} et \ref{dualizing-section-local-cohomology-noetherian}.
Nous en appliquons certains résultats à l'étude de la profondeur
dans la section \ref{dualizing-section-depth}. Une autre application montre comment,
étant donné un idéal de type fini $I$ d'un anneau $A$,
les objets « $I$-complets » et « de $I$-torsion »
de la catégorie dérivée de $A$ se correspondent par une équivalence ;
voir la section \ref{dualizing-section-torsion-and-complete}.
Pour approfondir la cohomologie locale, notamment le théorème de finitude
(qui repose sur la dualité locale -- voir ci-dessous), on se reportera à
Cohomologie locale, section \ref{local-cohomology-section-introduction}.

\medskip\noindent
L'essentiel de ce chapitre est consacré à la dualité pour un homomorphisme
d'anneaux et aux complexes dualisants. Voir les sections \ref{dualizing-section-trivial},
\ref{dualizing-section-base-change-trivial-duality}, \ref{dualizing-section-dualizing}, \ref{dualizing-section-dualizing-local}, \ref{dualizing-section-dimension-function}, \ref{dualizing-section-local-duality}, \ref{dualizing-section-dualizing-module},
\ref{dualizing-section-CM}, \ref{dualizing-section-gorenstein}, \ref{dualizing-section-ubiquity-dualizing} et \ref{dualizing-section-formal-fibres}.
La définition fondamentale est celle d'un complexe dualisant
$\omega_A^\bullet$ sur un anneau noethérien $A$ : c'est un objet
$\omega_A^\bullet \in D^{+}(A)$ dont les modules de cohomologie
$H^i(\omega_A^\bullet)$ sont des $A$-modules de type fini, qui est de dimension injective finie,
et tel que le morphisme
$$
A \longrightarrow R\Hom_A(\omega_A^\bullet, \omega_A^\bullet)
$$
soit un quasi-isomorphisme. Après avoir établi quelques propriétés
élémentaires des complexes dualisants, nous montrons qu'un complexe
dualisant définit une fonction de dimension. Nous démontrons ensuite
le théorème de dualité locale de Grothendieck. Après une brève étude
des modules dualisants et des anneaux de Cohen-Macaulay, nous introduisons
les anneaux de Gorenstein et montrons que beaucoup d'anneaux noethériens
usuels possèdent des complexes dualisants. Dans une dernière section,
nous appliquons ces résultats pour établir une bonne théorie des anneaux
locaux noethériens dont les fibres formelles sont de Gorenstein ou
localement des intersections complètes.

\medskip\noindent
Dans les dernières sections, nous décrivons une construction algébrique
des « foncteurs image inverse extraordinaire » utilisés en géométrie
algébrique, par exemple dans l'ouvrage \cite{RD}. Ce sujet est poursuivi
dans le chapitre sur la dualité pour les schémas. Voir
Dualité pour les schémas, section \ref{duality-section-introduction}.







\section{Surjections et injections essentielles}
\label{dualizing-section-essential}

\noindent
Nous travaillerons surtout dans des catégories de modules, mais autant
donner la définition dans le cas général.

\begin{definition}
\label{dualizing-definition-essential}
Soit $\mathcal{A}$ une catégorie abélienne.
\begin{enumerate}
\item Une injection $A \subset B$ de $\mathcal{A}$ est dite {\it essentielle},
ou l'on dit que $B$ est une {\it extension essentielle de} $A$,
si tout sous-objet non nul $B' \subset B$ a une intersection non nulle avec $A$.
\item Une surjection $f : A \to B$ de $\mathcal{A}$ est dite {\it essentielle}
si, pour tout sous-objet propre $A' \subset A$, on a $f(A') \not = B$.
\end{enumerate}
\end{definition}

\noindent
Voici quelques lemmes sur cette notion.

\begin{lemma}
\label{dualizing-lemma-essential}
Soit $\mathcal{A}$ une catégorie abélienne.
\begin{enumerate}
\item Si $A \subset B$ et $B \subset C$ sont des extensions essentielles,
alors $A \subset C$ est une extension essentielle.
\item Si $A \subset B$ est une extension essentielle et $C \subset B$
un sous-objet, alors $A \cap C \subset C$ est une extension essentielle.
\item Si $A \to B$ et $B \to C$ sont des surjections essentielles,
alors $A \to C$ est une surjection essentielle.
\item Étant donné une surjection essentielle $f : A \to B$ et une surjection
$A \to C$ de noyau $K$, le morphisme $C \to B/f(K)$ est une surjection essentielle.
\end{enumerate}
\end{lemma}

\begin{proof}
Omise.
\end{proof}

\begin{lemma}
\label{dualizing-lemma-union-essential-extensions}
Soit $R$ un anneau. Soit $M$ un $R$-module. Soit $E = \colim E_i$
une limite inductive filtrante de $R$-modules. Supposons donné
un système compatible d'injections essentielles $M \to E_i$ de $R$-modules.
Alors $M \to E$ est une injection essentielle.
\end{lemma}

\begin{proof}
Cela résulte immédiatement des définitions et de l'exactitude
des limites inductives filtrantes (Algèbre, lemme \ref{algebra-lemma-directed-colimit-exact}).
\end{proof}

\begin{lemma}
\label{dualizing-lemma-essential-extension}
Soit $R$ un anneau. Soient $M \subset N$ des $R$-modules.
Les conditions suivantes sont équivalentes :
\begin{enumerate}
\item $M \subset N$ est une extension essentielle ;
\item pour tout $x \in N$ non nul, il existe $f \in R$ tel que $fx \in M$
et $fx \not = 0$.
\end{enumerate}
\end{lemma}

\begin{proof}
Supposons (1) et soit $x \in N$ un élément non nul. D'après (1), on a
$Rx \cap M \not = 0$, ce qui entraîne (2).

\medskip\noindent
Supposons (2). Soit $N' \subset N$ un sous-module non nul. Choisissons $x \in N'$
non nul. D'après (2), il existe $f \in R$ tel que $fx \in M$ et $fx \not = 0$.
Ainsi $N' \cap M \not = 0$.
\end{proof}




\section{Modules injectifs}
\label{dualizing-section-injective-modules}

\noindent
Quelques résultats sur les modules injectifs sur des anneaux.

\begin{lemma}
\label{dualizing-lemma-product-injectives}
Soit $R$ un anneau. Tout produit de $R$-modules injectifs est injectif.
\end{lemma}

\begin{proof}
C'est un cas particulier de Homologie, lemme \ref{homology-lemma-product-injectives}.
\end{proof}

\begin{lemma}
\label{dualizing-lemma-injective-flat}
Soit $R \to S$ un homomorphisme plat d'anneaux. Si $E$ est un $S$-module injectif,
alors $E$ est injectif en tant que $R$-module.
\end{lemma}

\begin{proof}
Cela résulte de l'égalité $\Hom_R(M, E) = \Hom_S(M \otimes_R S, E)$
(Algèbre, lemme \ref{algebra-lemma-adjoint-tensor-restrict})
et de l'exactitude du produit tensoriel avec $S$.
\end{proof}

\begin{lemma}
\label{dualizing-lemma-injective-epimorphism}
Soit $R \to S$ un épimorphisme d'anneaux. Soit $E$ un $S$-module.
Si $E$ est injectif en tant que $R$-module, alors $E$ est un $S$-module injectif.
\end{lemma}

\begin{proof}
Cela résulte de l'égalité $\Hom_R(N, E) = \Hom_S(N, E)$ pour tout $S$-module $N$ ;
voir Algèbre, lemme \ref{algebra-lemma-epimorphism-modules}.
\end{proof}

\begin{lemma}
\label{dualizing-lemma-hom-injective}
Soit $R \to S$ un homomorphisme d'anneaux. Si $E$ est un $R$-module injectif,
alors $\Hom_R(S, E)$ est un $S$-module injectif.
\end{lemma}

\begin{proof}
Cela résulte de l'égalité $\Hom_S(N, \Hom_R(S, E)) = \Hom_R(N, E)$ donnée par
Algèbre, lemme \ref{algebra-lemma-adjoint-hom-restrict}.
\end{proof}

\begin{lemma}
\label{dualizing-lemma-essential-extensions-in-injective}
Soit $R$ un anneau. Soit $I$ un $R$-module injectif. Soit $E \subset I$
un sous-module. Les conditions suivantes sont équivalentes :
\begin{enumerate}
\item $E$ est injectif ;
\item pour tout $E \subset E' \subset I$ tel que $E \subset E'$ soit essentielle,
on a $E = E'$.
\end{enumerate}
En particulier, un $R$-module est injectif si et seulement si toute
extension essentielle est triviale.
\end{lemma}

\begin{proof}
La dernière assertion résulte de la première et du fait que la catégorie
des $R$-modules possède suffisamment d'injectifs
(Compléments d'algèbre, section \ref{more-algebra-section-injectives-modules}).

\medskip\noindent
Supposons (1). Soit $E \subset E' \subset I$ comme dans (2).
Alors le morphisme $\text{id}_E : E \to E$ se prolonge en un morphisme $\alpha : E' \to E$.
Le noyau de $\alpha$ est nécessairement nul, car son intersection avec
$E$ est nulle et $E'$ est une extension essentielle. Ainsi $E = E'$.

\medskip\noindent
Supposons (2). Soient $M \subset N$ des $R$-modules et $\varphi : M \to E$
un homomorphisme de $R$-modules. Pour prouver (1), il faut montrer que
$\varphi$ se prolonge en un morphisme $N \to E$. Considérons l'ensemble $\mathcal{S}$
des couples $(M', \varphi')$ où $M \subset M' \subset N$ et où $\varphi' : M' \to E$
est un homomorphisme de $R$-modules coïncidant avec $\varphi$ sur $M$.
Définissons un ordre sur $\mathcal{S}$ par la règle $(M', \varphi') \leq (M'', \varphi'')$
si et seulement si $M' \subset M''$ et $\varphi''|_{M'} = \varphi'$.
Il est clair que l'on peut prendre le maximum d'une partie totalement
ordonnée de $\mathcal{S}$. Le lemme de Zorn permet donc de supposer que $(M, \varphi)$
est un élément maximal.

\medskip\noindent
Choisissons un prolongement $\psi : N \to I$ de $\varphi$ composé avec l'inclusion $E \to I$.
C'est possible puisque $I$ est injectif.
Si $\psi(N) \subset E$, alors $\psi$ est le prolongement cherché.
Si $\psi(N)$ n'est pas contenu dans $E$, alors, d'après (2), l'inclusion
$E \subset E + \psi(N)$ n'est pas essentielle. On peut donc trouver un sous-module non nul
$K \subset E + \psi(N)$ dont l'intersection avec $E$ est $0$.
Cela signifie que $M' = \psi^{-1}(E + K)$ contient strictement $M$.
On peut donc prolonger $\varphi$ à $M'$ au moyen de
$$
M' \xrightarrow{\psi|_{M'}} E + K \to (E + K)/K = E
$$
Cela contredit la maximalité de $(M, \varphi)$.
\end{proof}

\begin{example}
\label{dualizing-example-reduced-ring-injective}
Soit $R$ un anneau réduit. Soit $\mathfrak p \subset R$ un idéal premier minimal ;
alors $K = R_\mathfrak p$ est un corps
(Algèbre, lemme \ref{algebra-lemma-minimal-prime-reduced-ring}).
Le module $K$ est un $R$-module injectif. En effet, on a
$\Hom_R(M, K) = \Hom_K(M_\mathfrak p, K)$ pour tout $R$-module $M$.
La localisation étant un foncteur exact, ainsi que le passage au dual
sur les $K$-espaces vectoriels, on en déduit que $\Hom_R(-, K)$
est un foncteur exact, c'est-à-dire que $K$ est un $R$-module injectif.
\end{example}

\begin{lemma}
\label{dualizing-lemma-sum-injective-modules}
Soit $R$ un anneau noethérien. Toute somme directe de modules injectifs
est injective.
\end{lemma}

\begin{proof}
Soit $E_i$ une famille de modules injectifs indexée par un ensemble $I$.
Posons $E = \bigoplus E_i$. Pour montrer que $E$ est injectif, utilisons
Injectifs, lemme \ref{injectives-lemma-criterion-baer}.
Soit donc $\varphi : J \to E$ un homomorphisme de modules d'un idéal de $R$
dans $E$. Puisque $J$ est un $R$-module de type fini
(car $R$ est noethérien), il existe un nombre fini d'éléments $i_1, \ldots, i_r \in I$
tels que $\varphi$ prenne ses valeurs dans $\bigoplus_{j = 1, \ldots, r} E_{i_j}$.
On peut alors prolonger $\varphi$ à valeurs dans $\bigoplus_{j = 1, \ldots, r} E_{i_j}$
en utilisant l'injectivité des modules $E_{i_j}$.
\end{proof}

\begin{lemma}
\label{dualizing-lemma-localization-injective-modules}
Soit $R$ un anneau noethérien. Soit $S \subset R$ une partie multiplicative.
Si $E$ est un $R$-module injectif, alors $S^{-1}E$ est un $S^{-1}R$-module injectif.
\end{lemma}

\begin{proof}
Puisque $R \to S^{-1}R$ est un épimorphisme d'anneaux, il suffit de montrer que
$S^{-1}E$ est injectif en tant que $R$-module ; voir le lemme \ref{dualizing-lemma-injective-epimorphism}.
Pour cela, utilisons Injectifs, lemme \ref{injectives-lemma-criterion-baer}.
Soit donc $I \subset R$ un idéal et soit $\varphi : I \to S^{-1} E$ un homomorphisme de $R$-modules.
Puisque $I$ est un $R$-module de présentation finie
(car $R$ est noethérien), il existe $f \in S$ et un homomorphisme
de $R$-modules $I \to E$ tels que $f\varphi$ soit le composé $I \to E \to S^{-1}E$
(Algèbre, lemme \ref{algebra-lemma-hom-from-finitely-presented}).
On peut alors prolonger $I \to E$ en un homomorphisme $R \to E$.
Le composé
$$
R \to E \to S^{-1}E \xrightarrow{f^{-1}} S^{-1}E
$$
est le prolongement cherché de $\varphi$ à $R$.
\end{proof}

\begin{lemma}
\label{dualizing-lemma-injective-module-divide}
Soit $R$ un anneau noethérien. Soit $I$ un $R$-module injectif.
\begin{enumerate}
\item Soit $f \in R$. Alors $E = \bigcup I[f^n] = I[f^\infty]$
est un sous-module injectif de $I$.
\item Soit $J \subset R$ un idéal. Le sous-module de torsion
par les puissances de $J$, noté $I[J^\infty]$, est un sous-module injectif de $I$.
\end{enumerate}
\end{lemma}

\begin{proof}
Utilisons le lemme \ref{dualizing-lemma-essential-extensions-in-injective} pour prouver (1).
Supposons que $E \subset E' \subset I$ et que $E'$ soit une extension essentielle de $E$.
Montrons que $E' = E$. Sinon, il existe $x \in E'$ tel que $x \not \in E$.
Posons $J = \{ a \in R \mid ax \in E\}$. Puisque $R$ est noethérien, on peut écrire $J = (g_1, \ldots, g_t)$
pour certains $g_i \in R$. Par définition, $E$ est l'ensemble des éléments
de $I$ annulés par des puissances de $f$ ; on peut donc choisir
des entiers $n_i$ tels que $f^{n_i}g_ix = 0$.
Posons $n = \mathrm{max}\{ n_i \}$. Alors $x' = f^n x$ est un élément de $E'$
qui n'appartient pas à $E$ et qui est annulé par $J$.
Posons $J' = \{ a \in R \mid ax' \in E \}$, de sorte que $J \subset J'$. Réciproquement, on a $a \in J'$
si et seulement si $ax' \in E$, ou encore si et seulement si $f^m a x' = 0$
pour un certain $m \geq 0$. Mais alors $f^m a x' = f^{m + n} a x$ entraîne $ax \in E$, c'est-à-dire $a \in J$.
Ainsi $J = J'$. On a donc $J = J' = \text{Ann}(x')$, d'où $Rx' \cap E  = 0$.
Par conséquent, $E'$ n'est pas une extension essentielle de $E$,
ce qui est contradictoire.

\medskip\noindent
Pour prouver (2), écrivons $J = (f_1, \ldots, f_t)$. Alors $I[J^\infty]$ est égal à
$$
(\ldots((I[f_1^\infty])[f_2^\infty])\ldots)[f_t^\infty]
$$
et le résultat découle de (1) par récurrence.
\end{proof}

\begin{lemma}
\label{dualizing-lemma-injective-dimension-over-polynomial-ring}
Soit $A$ un anneau noethérien. Soit $E$ un $A$-module injectif.
Alors $E \otimes_A A[x]$ est d'amplitude injective $[0, 1]$
en tant qu'objet de $D(A[x])$. En particulier, $E \otimes_A A[x]$
est de dimension injective finie en tant que $A[x]$-module.
\end{lemma}

\begin{proof}
Écrivons $E[x] = E \otimes_A A[x]$. Considérons la suite exacte courte de $A[x]$-modules
$$
0 \to E[x] \to \Hom_A(A[x], E[x]) \to \Hom_A(A[x], E[x]) \to 0
$$
où le premier morphisme envoie $p \in E[x]$ sur $f \mapsto fp$ et le second envoie
$\varphi$ sur $f \mapsto \varphi(xf) - x\varphi(f)$.
Le second morphisme est surjectif, car $\Hom_A(A[x], E[x]) = \prod_{n \geq 0} E[x]$ en tant que groupe abélien,
et le morphisme envoie $(e_n)$ sur $(e_{n + 1} - xe_n)$, ce qui définit une surjection.
En tant que $A$-module, on a $E[x] \cong \bigoplus_{n \geq 0} E$,
qui est injectif d'après le lemme \ref{dualizing-lemma-sum-injective-modules}.
Le $A[x]$-module $\Hom_A(A[x], E[x])$ est donc injectif d'après le lemme \ref{dualizing-lemma-hom-injective},
ce qui achève la démonstration.
\end{proof}



\section{Couvertures projectives}
\label{dualizing-section-projective-cover}

\noindent
Dans cette section, nous étudions brièvement les couvertures projectives.

\begin{definition}
\label{dualizing-definition-projective-cover}
Soit $R$ un anneau. Une surjection $P \to M$ de $R$-modules est appelée
{\it couverture projective}, ou parfois {\it enveloppe projective},
si $P$ est un $R$-module projectif et si $P \to M$ est une surjection essentielle.
\end{definition}

\noindent
Les couvertures projectives n'existent pas toujours. Par exemple, si $k$
est un corps et $R = k[x]$ l'anneau de polynômes sur $k$, le module $M = R/(x)$
n'a pas de couverture projective. En effet, pour toute surjection $f : P \to M$
avec $P$ projectif sur $R$, le sous-module propre $(x - 1)P$ se surjecte
sur $M$. Ainsi $f$ n'est pas essentielle.

\begin{lemma}
\label{dualizing-lemma-projective-cover-unique}
Soit $R$ un anneau et soit $M$ un $R$-module. Si une couverture projective
de $M$ existe, elle est unique à isomorphisme près.
\end{lemma}

\begin{proof}
Soient $P \to M$ et $P' \to M$ des couvertures projectives. Puisque $P$ est
un $R$-module projectif et que $P' \to M$ est surjective, il existe
un homomorphisme de $R$-modules $\alpha : P \to P'$ compatible avec les morphismes vers $M$.
Puisque $P' \to M$ est essentielle, $\alpha$ est surjectif.
Comme $P'$ est un $R$-module projectif, on peut choisir une décomposition
en somme directe $P = \Ker(\alpha) \oplus P'$. Puisque $P' \to M$ est surjective
et que $P \to M$ est essentielle, on en déduit que $\Ker(\alpha)$ est nul, comme voulu.
\end{proof}

\noindent
Voici un exemple où les couvertures projectives existent.

\begin{lemma}
\label{dualizing-lemma-projective-covers-local}
Soit $(R, \mathfrak m, \kappa)$ un anneau local. Tout $R$-module de type fini possède
une couverture projective.
\end{lemma}

\begin{proof}
Soit $M$ un $R$-module de type fini. Posons $r = \dim_\kappa(M/\mathfrak m M)$.
Choisissons $x_1, \ldots, x_r \in M$ dont les images forment une base de $M/\mathfrak m M$.
Considérons le morphisme $f : R^{\oplus r} \to M$. D'après le lemme de Nakayama, il est
surjectif (Algèbre, lemme \ref{algebra-lemma-NAK}). Si $N \subset R^{\oplus r}$ est un sous-module propre,
alors $N/\mathfrak m N \to \kappa^{\oplus r}$ n'est pas surjectif (encore par le lemme de Nakayama),
donc $N/\mathfrak m N \to M/\mathfrak m M$ n'est pas surjectif.
Ainsi $f$ est une surjection essentielle.
\end{proof}







\section{Enveloppes injectives}
\label{dualizing-section-injective-hull}

\noindent
Dans cette section, nous étudions brièvement les enveloppes injectives.

\begin{definition}
\label{dualizing-definition-injective-hull}
Soit $R$ un anneau. Une injection $M \to I$ de $R$-modules est appelée
{\it enveloppe injective} si $I$ est un $R$-module injectif
et si $M \to I$ est une injection essentielle.
\end{definition}

\noindent
Les enveloppes injectives existent toujours.

\begin{lemma}
\label{dualizing-lemma-injective-hull}
Soit $R$ un anneau. Tout $R$-module possède une enveloppe injective.
\end{lemma}

\begin{proof}
Soit $M$ un $R$-module. D'après Compléments d'algèbre, section \ref{more-algebra-section-injectives-modules},
la catégorie des $R$-modules possède suffisamment d'injectifs.
Choisissons une injection $M \to I$ où $I$ est un $R$-module injectif.
Considérons l'ensemble $\mathcal{S}$ des sous-modules $M \subset E \subset I$
tels que $E$ soit une extension essentielle de $M$.
Ordonnons $\mathcal{S}$ par inclusion. Si $\{E_\alpha\}$ est une partie totalement
ordonnée de $\mathcal{S}$, alors $\bigcup E_\alpha$ est aussi une extension essentielle
de $M$ (lemme \ref{dualizing-lemma-union-essential-extensions}).
Le lemme de Zorn fournit donc un élément maximal $E \in \mathcal{S}$.
Affirmons que $M \subset E$ est une enveloppe injective, c'est-à-dire
que $E$ est un $R$-module injectif. Cela résulte du lemme \ref{dualizing-lemma-essential-extensions-in-injective}.
\end{proof}

\begin{lemma}
\label{dualizing-lemma-injective-hull-unique}
Soit $R$ un anneau. Soient $M$, $N$ des $R$-modules et soient $M \to E$
et $N \to E'$ des enveloppes injectives. Alors :
\begin{enumerate}
\item pour tout homomorphisme de $R$-modules $\varphi : M \to N$, il existe un
homomorphisme de $R$-modules $\psi : E \to E'$ tel que le diagramme
$$
\xymatrix{
M \ar[r] \ar[d]_\varphi & E \ar[d]^\psi \\
N \ar[r] & E'
}
$$
soit commutatif ;
\item si $\varphi$ est injectif, alors $\psi$ est injectif ;
\item si $\varphi$ est une injection essentielle, alors $\psi$ est un isomorphisme ;
\item si $\varphi$ est un isomorphisme, alors $\psi$ est un isomorphisme ;
\item si $M \to I$ est un plongement de $M$ dans un $R$-module injectif,
alors il existe un isomorphisme $I \cong E \oplus I'$ compatible avec les plongements de $M$.
\end{enumerate}
En particulier, l'enveloppe injective $E$ de $M$ est unique à isomorphisme près.
\end{lemma}

\begin{proof}
L'assertion (1) résulte de ce que $E'$ est un $R$-module injectif.
L'assertion (2) résulte de ce que $\Ker(\psi) \cap M = 0$
et que $E$ est une extension essentielle de $M$.
Supposons que $\varphi$ soit une injection essentielle. Alors
$E \cong \psi(E) \subset E'$ d'après (2), d'où $E' = \psi(E) \oplus E''$ puisque $E$ est injectif.
Puisque $E'$ est une extension essentielle de $M$
(lemme \ref{dualizing-lemma-essential}), on obtient $E'' = 0$.
L'assertion (4) est un cas particulier de (3).
Supposons $M \to I$ comme dans (5).
Choisissons un morphisme $\alpha : E \to I$ prolongeant le morphisme $M \to I$.
Le même raisonnement montre que $\alpha$ est injectif.
Comme précédemment, $\alpha(E)$ est donc facteur direct de $I$.
Cela prouve (5).
\end{proof}

\begin{example}
\label{dualizing-example-injective-hull-domain}
Soit $R$ un anneau intègre de corps des fractions $K$.
Alors $R \subset K$ est une enveloppe injective de $R$.
En effet, l'exemple \ref{dualizing-example-reduced-ring-injective} montre que $K$ est un $R$-module injectif,
et le lemme \ref{dualizing-lemma-essential-extension} montre que $R \subset K$ est une extension essentielle.
\end{example}

\begin{definition}
\label{dualizing-definition-indecomposable}
Un objet $X$ d'une catégorie additive est dit {\it indécomposable}
s'il est non nul et si $X = Y \oplus Z$ entraîne que $Y = 0$ ou $Z = 0$.
\end{definition}

\begin{lemma}
\label{dualizing-lemma-indecomposable-injective}
Soit $R$ un anneau. Soit $E$ un $R$-module injectif indécomposable.
Alors :
\begin{enumerate}
\item $E$ est l'enveloppe injective de tout sous-module non nul de $E$ ;
\item l'intersection de deux sous-modules non nuls de $E$ est non nulle ;
\item $\text{End}_R(E)$ est un anneau local non commutatif dont l'idéal maximal
est formé des $\varphi : E \to E$ de noyau non nul ;
\item l'ensemble des diviseurs de zéro sur $E$ est un idéal premier $\mathfrak p$
de $R$, et $E$ est un $R_\mathfrak p$-module injectif.
\end{enumerate}
\end{lemma}

\begin{proof}
L'assertion (1) résulte du lemme \ref{dualizing-lemma-injective-hull-unique}.
L'assertion (2) découle de (1) et de la définition des enveloppes injectives.

\medskip\noindent
Preuve de (3). Posons $A = \text{End}_R(E)$ et $I = \{\varphi \in A \mid \Ker(\varphi) \not = 0\}$.
L'énoncé signifie que $I$ est un idéal bilatère et que tout
$\varphi \in A$ tel que $\varphi \not \in I$ est inversible.
Supposons que $\varphi$ et $\psi$ ne soient pas injectifs.
Alors $\Ker(\varphi) \cap \Ker(\psi)$ est non nul d'après (2).
Ainsi $\varphi + \psi \in I$. Il en résulte que $I$ est un idéal bilatère.
Si $\varphi \in A$ et $\varphi \not \in I$, alors $E \cong \varphi(E) \subset E$ est un sous-module injectif,
donc $E = \varphi(E)$ puisque $E$ est indécomposable.

\medskip\noindent
Preuve de (4). Considérons l'homomorphisme d'anneaux $R \to A$ et soit $\mathfrak p \subset R$
l'image réciproque de l'idéal maximal $I$. Il est alors clair
que $\mathfrak p$ est un idéal premier et que $R \to A$ se prolonge en $R_\mathfrak p \to A$.
Ainsi $E$ est un $R_\mathfrak p$-module.
Le lemme \ref{dualizing-lemma-injective-epimorphism} montre que $E$ est injectif en tant que $R_\mathfrak p$-module.
\end{proof}

\begin{lemma}
\label{dualizing-lemma-injective-hull-indecomposable}
Soit $\mathfrak p \subset R$ un idéal premier d'un anneau $R$.
Soit $E$ l'enveloppe injective de $R/\mathfrak p$. Alors :
\begin{enumerate}
\item $E$ est indécomposable ;
\item $E$ est l'enveloppe injective de $\kappa(\mathfrak p)$ ;
\item $E$ est l'enveloppe injective de $\kappa(\mathfrak p)$ sur l'anneau $R_\mathfrak p$.
\end{enumerate}
\end{lemma}

\begin{proof}
D'après le lemme \ref{dualizing-lemma-essential-extension}, l'inclusion $R/\mathfrak p \subset \kappa(\mathfrak p)$ est une extension essentielle.
Le lemme \ref{dualizing-lemma-injective-hull-unique} donne alors (2).
Pour $f \in R$ tel que $f \not \in \mathfrak p$, le morphisme $f : \kappa(\mathfrak p) \to \kappa(\mathfrak p)$ est un isomorphisme ;
le morphisme $f : E \to E$ est donc un isomorphisme, voir le lemme \ref{dualizing-lemma-injective-hull-unique}.
Ainsi $E$ est un $R_\mathfrak p$-module. Il est injectif en tant que
$R_\mathfrak p$-module d'après le lemme \ref{dualizing-lemma-injective-epimorphism}.
Enfin, soit $E' \subset E$ un sous-$R$-module injectif non nul.
Alors $J = (R/\mathfrak p) \cap E'$ est non nul. Quitte à remplacer $E'$ par un sous-module,
on peut supposer que $E'$ est l'enveloppe injective de $J$
(voir par exemple le lemme \ref{dualizing-lemma-injective-hull-unique}).
Observons que $R/\mathfrak p$ est une extension essentielle de $J$,
par exemple d'après le lemme \ref{dualizing-lemma-essential-extension}. Ainsi $E' \to E$
est un isomorphisme d'après le lemme \ref{dualizing-lemma-injective-hull-unique} (3).
Par conséquent, $E$ est indécomposable.
\end{proof}

\begin{lemma}
\label{dualizing-lemma-indecomposable-injective-noetherian}
Soit $R$ un anneau noethérien. Soit $E$ un $R$-module injectif
indécomposable. Il existe alors un idéal premier $\mathfrak p$ de $R$
tel que $E$ soit l'enveloppe injective de $\kappa(\mathfrak p)$.
\end{lemma}

\begin{proof}
Soit $\mathfrak p$ l'idéal premier obtenu dans le lemme \ref{dualizing-lemma-indecomposable-injective}.
Écrivons $\mathfrak p = (f_1, \ldots, f_r)$.
Choisissons un élément non nul $x \in \bigcap \Ker(f_i : E \to E)$, voir le lemme \ref{dualizing-lemma-indecomposable-injective}.
Alors $(R_\mathfrak p)x$ est un module isomorphe à $\kappa(\mathfrak p)$ contenu dans $E$.
On conclut par le lemme \ref{dualizing-lemma-indecomposable-injective}.
\end{proof}

\begin{proposition}[Structure des modules injectifs sur les anneaux noethériens]
\label{dualizing-proposition-structure-injectives-noetherian}
Soit $R$ un anneau noethérien.
Tout module injectif est somme directe de modules injectifs indécomposables.
Tout module injectif indécomposable est l'enveloppe injective
du corps résiduel en un idéal premier.
\end{proposition}

\begin{proof}
La seconde assertion est le lemme \ref{dualizing-lemma-indecomposable-injective-noetherian}.
Pour la première, soit $I$ un $R$-module injectif.
Construisons par récurrence transfinie des $I_\alpha \subset I$
pour des ordinaux $\alpha$, sommes directes de $R$-modules injectifs
indécomposables $E_{\beta + 1}$ pour $\beta < \alpha$.
Pour $\alpha = 0$, posons $I_0 = 0$. Supposons donné un ordinal $\alpha$
pour lequel $I_\alpha$ a été construit. Alors $I_\alpha$ est un
$R$-module injectif d'après le lemme \ref{dualizing-lemma-sum-injective-modules}.
Ainsi $I \cong I_\alpha \oplus I'$. Si $I' = 0$, c'est terminé.
Sinon, $I'$ possède un idéal premier associé d'après
Algèbre, lemme \ref{algebra-lemma-ass-zero}.
Le module $I'$ contient donc une copie de $R/\mathfrak p$ pour un idéal premier $\mathfrak p$.
Ainsi $I'$ contient un sous-module indécomposable $E$ d'après
les lemmes \ref{dualizing-lemma-injective-hull-unique} et \ref{dualizing-lemma-injective-hull-indecomposable}. Posons $I_{\alpha + 1} = I_\alpha \oplus E_\alpha$.
Si $\alpha$ est un ordinal limite et si $I_\beta$ a été construit
pour $\beta < \alpha$, posons $I_\alpha = \bigcup_{\beta < \alpha} I_\beta$.
Observons que $I_\alpha = \bigoplus_{\beta < \alpha} E_{\beta + 1}$.
Cela achève la démonstration.
\end{proof}



\section{Dualité sur les anneaux locaux artiniens}
\label{dualizing-section-artinian}

\noindent
Soit $(R, \mathfrak m, \kappa)$ un anneau local artinien.
Rappelons qu'alors $R$ est noethérien et que $R$ est de longueur finie
en tant que $R$-module. De plus, un $R$-module est de type fini
si et seulement s'il est de longueur finie. Nous utiliserons ces faits
sans autre mention dans cette section. Pour plus de détails, voir
Algèbre, sections \ref{algebra-section-length} et \ref{algebra-section-artinian},
ainsi que Algèbre, proposition \ref{algebra-proposition-dimension-zero-ring}.

\begin{lemma}
\label{dualizing-lemma-finite}
Soit $(R, \mathfrak m, \kappa)$ un anneau local artinien.
Soit $E$ une enveloppe injective de $\kappa$. Pour tout
$R$-module de type fini $M$, on a
$$
\text{longueur}_R(M) = \text{longueur}_R(\Hom_R(M, E))
$$
En particulier, l'enveloppe injective $E$ de $\kappa$ est un $R$-module de type fini.
\end{lemma}

\begin{proof}
Puisque $E$ est une extension essentielle de $\kappa$, on a
$\kappa = E[\mathfrak m]$, où $E[\mathfrak m]$ désigne la $\mathfrak m$-torsion de $E$
(avec les notations de Compléments d'algèbre, section \ref{more-algebra-section-torsion}).
Ainsi $\Hom_R(\kappa, E) \cong \kappa$ et l'égalité des longueurs est vraie pour $M = \kappa$.
Démontrons l'égalité affichée dans le lemme par récurrence sur
la longueur de $M$. Si $M$ est non nul, il existe une surjection
$M \to \kappa$ de noyau $M'$. Puisque le foncteur $M \mapsto \Hom_R(M, E)$
est exact, on obtient une suite exacte courte
$$
0 \to \Hom_R(\kappa, E) \to \Hom_R(M, E) \to \Hom_R(M', E) \to 0.
$$
L'additivité de la longueur pour cette suite et pour la suite $0 \to M' \to M \to \kappa \to 0$,
ainsi que l'égalité pour $M'$ (hypothèse de récurrence) et pour $\kappa$,
donnent l'égalité pour $M$.
La dernière assertion du lemme résulte de l'égalité $E = \Hom_R(R, E)$.
\end{proof}

\begin{lemma}
\label{dualizing-lemma-evaluate}
Soit $(R, \mathfrak m, \kappa)$ un anneau local artinien.
Soit $E$ une enveloppe injective de $\kappa$.
Pour tout $R$-module de type fini $M$, le morphisme d'évaluation
$$
M \longrightarrow \Hom_R(\Hom_R(M, E), E)
$$
est un isomorphisme. En particulier, $R = \Hom_R(E, E)$.
\end{lemma}

\begin{proof}
Observons que la flèche affichée est injective. En effet, si $x \in M$
est un élément non nul, il existe un morphisme non nul $Rx \to \kappa$ que l'on peut
prolonger en un morphisme $\varphi : M \to E$ ne s'annulant pas sur $x$.
Puisque la source et le but de la flèche ont même longueur d'après
le lemme \ref{dualizing-lemma-finite}, c'est un isomorphisme.
La dernière assertion s'obtient en prenant $M = R$.
\end{proof}

\noindent
Pour énoncer le lemme suivant, notons $\text{Mod}^{fg}_R$ la catégorie des
$R$-modules de type fini sur un anneau $R$.

\begin{lemma}
\label{dualizing-lemma-duality}
Soit $(R, \mathfrak m, \kappa)$ un anneau local artinien.
Soit $E$ une enveloppe injective de $\kappa$.
Le foncteur $D(-) = \Hom_R(-, E)$ induit une anti-équivalence exacte
$\text{Mod}^{fg}_R \to \text{Mod}^{fg}_R$, et $D \circ D \cong \text{id}$.
\end{lemma}

\begin{proof}
Nous avons vu que $D \circ D = \text{id}$ sur $\text{Mod}^{fg}_R$ dans le lemme \ref{dualizing-lemma-evaluate}.
Il en résulte immédiatement que $D$ est une anti-équivalence.
\end{proof}

\begin{lemma}
\label{dualizing-lemma-duality-torsion-cotorsion}
Gardons les hypothèses et les notations du lemme \ref{dualizing-lemma-duality}.
Soient $I \subset R$ un idéal et $M$ un $R$-module de type fini.
Alors
$$
D(M[I]) = D(M)/ID(M) \quad\text{et}\quad D(M/IM) = D(M)[I]
$$
\end{lemma}

\begin{proof}
Écrivons $I = (f_1, \ldots, f_t)$. Considérons le morphisme
$$
M^{\oplus t} \xrightarrow{f_1, \ldots, f_t} M
$$
de conoyau $M/IM$. En appliquant le foncteur exact $D$,
on en déduit que $D(M/IM)$ est $D(M)[I]$. L'autre cas se démontre de la même manière.
\end{proof}



\section{Enveloppe injective du corps résiduel}
\label{dualizing-section-hull-residue-field}

\noindent
Dans cette section, la plupart des résultats concernent les anneaux locaux noethériens.

\begin{lemma}
\label{dualizing-lemma-quotient}
Soit $R \to S$ un homomorphisme surjectif d'anneaux locaux, de noyau $I$.
Soit $E$ l'enveloppe injective du corps résiduel de $R$ sur $R$.
Alors $E[I]$ est l'enveloppe injective du corps résiduel de $S$ sur $S$.
\end{lemma}

\begin{proof}
Observons que $E[I] = \Hom_R(S, E)$ puisque $S = R/I$. Ainsi $E[I]$ est un $S$-module injectif
d'après le lemme \ref{dualizing-lemma-hom-injective}. Puisque $E$ est une extension essentielle
de $\kappa = R/\mathfrak m_R$, il en résulte que $E[I]$ est aussi une extension essentielle
de $\kappa$. D'où le résultat.
\end{proof}

\begin{lemma}
\label{dualizing-lemma-torsion-submodule-sum-injective-hulls}
Soit $(R, \mathfrak m, \kappa)$ un anneau local.
Soit $E$ l'enveloppe injective de $\kappa$.
Soit $M$ un module de torsion par les puissances de $\mathfrak m$ sur $R$,
tel que $n = \dim_\kappa(M[\mathfrak m]) < \infty$.
Alors $M$ est isomorphe à un sous-module de $E^{\oplus n}$.
\end{lemma}

\begin{proof}
Observons que $E^{\oplus n}$ est l'enveloppe injective de $\kappa^{\oplus n} = M[\mathfrak m]$.
Il existe donc un homomorphisme de $R$-modules
$M \to E^{\oplus n}$ qui est injectif sur $M[\mathfrak m]$.
Puisque $M$ est de torsion par les puissances de $\mathfrak m$,
l'inclusion $M[\mathfrak m] \subset M$ est une extension essentielle
(par exemple d'après le lemme \ref{dualizing-lemma-essential-extension}).
On en déduit que le noyau de $M \to E^{\oplus n}$ est nul.
\end{proof}

\begin{lemma}
\label{dualizing-lemma-union-artinian}
Soit $(R, \mathfrak m, \kappa)$ un anneau local noethérien.
Soit $E$ une enveloppe injective de $\kappa$ sur $R$.
Soit $E_n$ une enveloppe injective de $\kappa$ sur $R/\mathfrak m^n$.
Alors $E = \bigcup E_n$ et $E_n = E[\mathfrak m^n]$.
\end{lemma}

\begin{proof}
On a $E_n = E[\mathfrak m^n]$ d'après le lemme \ref{dualizing-lemma-quotient}.
On a $E = \bigcup E_n$ parce que $\bigcup E_n = E[\mathfrak m^\infty]$ est un sous-$R$-module injectif
qui contient $\kappa$ ; voir le lemme \ref{dualizing-lemma-injective-module-divide}.
\end{proof}

\noindent
Le lemme suivant montre que l'enveloppe injective du corps résiduel
d'un anneau local noethérien ne dépend que du complété.

\begin{lemma}
\label{dualizing-lemma-compare}
Soit $R \to S$ un homomorphisme local plat d'anneaux locaux noethériens
tel que $R/\mathfrak m_R \cong S/\mathfrak m_R S$.
Alors l'enveloppe injective du corps résiduel de $R$
est l'enveloppe injective du corps résiduel de $S$.
\end{lemma}

\begin{proof}
Notons que $\mathfrak m_RS = \mathfrak m_S$, puisque le quotient par le premier idéal est un corps.
Posons $\kappa = R/\mathfrak m_R = S/\mathfrak m_S$.
Soit $E_R$ l'enveloppe injective de $\kappa$ sur $R$.
Soit $E_S$ l'enveloppe injective de $\kappa$ sur $S$.
Observons que $E_S$ est un $R$-module injectif d'après le lemme \ref{dualizing-lemma-injective-flat}.
Choisissons un prolongement $E_R \to E_S$ de l'identification des corps résiduels.
Ce morphisme est un isomorphisme d'après le lemme \ref{dualizing-lemma-union-artinian},
car $R \to S$ induit un isomorphisme $R/\mathfrak m_R^n \to S/\mathfrak m_S^n$ pour tout $n$.
\end{proof}

\begin{lemma}
\label{dualizing-lemma-endos}
Soit $(R, \mathfrak m, \kappa)$ un anneau local noethérien.
Soit $E$ une enveloppe injective de $\kappa$ sur $R$.
Alors $\Hom_R(E, E)$ est canoniquement isomorphe au complété de $R$.
\end{lemma}

\begin{proof}
Écrivons $E = \bigcup E_n$ avec $E_n = E[\mathfrak m^n]$ comme dans le lemme \ref{dualizing-lemma-union-artinian}.
Tout endomorphisme de $E$ préserve cette filtration. Ainsi
$$
\Hom_R(E, E) = \lim \Hom_R(E_n, E_n)
$$
Le lemme en résulte, puisque
$\Hom_R(E_n, E_n) = \Hom_{R/\mathfrak m^n}(E_n, E_n) = R/\mathfrak m^n$
d'après le lemme \ref{dualizing-lemma-evaluate}.
\end{proof}

\begin{lemma}
\label{dualizing-lemma-injective-hull-has-dcc}
Soit $(R, \mathfrak m, \kappa)$ un anneau local noethérien.
Soit $E$ une enveloppe injective de $\kappa$ sur $R$.
Alors $E$ satisfait à la condition de chaîne descendante.
\end{lemma}

\begin{proof}
Si $E \supset M_1 \supset M_2 \supset \ldots$ est une suite de sous-modules, alors
$$
\Hom_R(E, E) \to \Hom_R(M_1, E) \to \Hom_R(M_2, E) \to \ldots
$$
est une suite de surjections. D'après le lemme \ref{dualizing-lemma-endos},
chacun de ces modules est un module sur le complété $R^\wedge = \Hom_R(E, E)$.
Puisque $R^\wedge$ est noethérien (Algèbre, lemme \ref{algebra-lemma-completion-Noetherian-Noetherian}),
la suite stationne : $\Hom_R(M_n, E) = \Hom_R(M_{n + 1}, E) = \ldots$.
Puisque $E$ est injectif, cela ne peut se produire que si $\Hom_R(M_n/M_{n + 1}, E)$ est nul.
Or, si $M_n/M_{n + 1}$ est non nul, il contient un élément non nul annulé par $\mathfrak m$,
car $E$ est de torsion par les puissances de $\mathfrak m$ d'après le lemme \ref{dualizing-lemma-union-artinian}.
Dans ce cas, $M_n/M_{n + 1}$ admet un morphisme non nul vers $E$,
ce qui contredit l'annulation supposée. Cela achève la démonstration.
\end{proof}

\begin{lemma}
\label{dualizing-lemma-describe-categories}
Soit $(R, \mathfrak m, \kappa)$ un anneau local noethérien.
Soit $E$ une enveloppe injective de $\kappa$.
\begin{enumerate}
\item Pour un $R$-module $M$, les conditions suivantes sont équivalentes :
\begin{enumerate}
\item $M$ satisfait à la condition de chaîne ascendante ;
\item $M$ est un $R$-module de type fini ;
\item il existe $n, m$ et une suite exacte
$R^{\oplus m} \to R^{\oplus n} \to M \to 0$.
\end{enumerate}
\item Pour un $R$-module $M$, les conditions suivantes sont équivalentes :
\begin{enumerate}
\item $M$ satisfait à la condition de chaîne descendante ;
\item $M$ est de torsion par les puissances de $\mathfrak m$ et
$\dim_\kappa(M[\mathfrak m]) < \infty$ ;
\item il existe $n, m$ et une suite exacte
$0 \to M \to E^{\oplus n} \to E^{\oplus m}$.
\end{enumerate}
\end{enumerate}
\end{lemma}

\begin{proof}
Nous omettons la démonstration de (1).

\medskip\noindent
Soit $M$ un $R$-module satisfaisant à la condition de chaîne descendante.
Soit $x \in M$. Alors $\mathfrak m^n x$ est une chaîne descendante de sous-modules,
donc elle stationne. Ainsi $\mathfrak m^nx = \mathfrak m^{n + 1}x$ pour un certain $n$.
Le lemme de Nakayama (Algèbre, lemme \ref{algebra-lemma-NAK}) entraîne alors $\mathfrak m^n x = 0$,
c'est-à-dire que $x$ est de torsion par les puissances de $\mathfrak m$.
Puisque $M[\mathfrak m]$ est un espace vectoriel sur $\kappa$, il doit être de dimension
finie pour satisfaire à la condition de chaîne descendante.

\medskip\noindent
Supposons que $M$ soit de torsion par les puissances de $\mathfrak m$
et que son sous-module de $\mathfrak m$-torsion $M[\mathfrak m]$ soit de dimension finie.
D'après le lemme \ref{dualizing-lemma-torsion-submodule-sum-injective-hulls}, on voit que $M$ est un sous-module
de $E^{\oplus n}$ pour un certain $n$.
Considérons le quotient $N = E^{\oplus n}/M$. D'après le lemme \ref{dualizing-lemma-injective-hull-has-dcc},
le module $E$ satisfait à la condition de chaîne descendante,
donc il en est de même de $E^{\oplus n}$ et de $N$.
Ainsi $N$ satisfait à (2)(a), ce qui entraîne que $N$ satisfait
à (2)(b) d'après le deuxième paragraphe de la démonstration.
En appliquant à nouveau le lemme \ref{dualizing-lemma-torsion-submodule-sum-injective-hulls}, on voit donc que $N$
est un sous-module de $E^{\oplus m}$ pour un certain $m$.
On a ainsi une suite exacte courte
$0 \to M \to E^{\oplus n} \to E^{\oplus m}$.

\medskip\noindent
Supposons donnée une suite exacte courte
$0 \to M \to E^{\oplus n} \to E^{\oplus m}$.
Puisque $E$ satisfait à la condition de chaîne descendante
d'après le lemme \ref{dualizing-lemma-injective-hull-has-dcc}, il en est de même de $M$.
\end{proof}

\begin{proposition}[Dualité de Matlis]
\label{dualizing-proposition-matlis}
Soit $(R, \mathfrak m, \kappa)$ un anneau local noethérien complet.
Soit $E$ l'enveloppe injective de $\kappa$ sur $R$.
Le foncteur $D(-) = \Hom_R(-, E)$ est une anti-équivalence
$$
\left\{
\begin{matrix}
R\text{-modules satisfaisant la} \\
\text{condition de chaîne descendante}
\end{matrix}
\right\}
\longleftrightarrow
\left\{
\begin{matrix}
R\text{-modules satisfaisant la} \\
\text{condition de chaîne ascendante}
\end{matrix}
\right\}
$$
et l'on a $D \circ D = \text{id}$ de chaque côté de l'équivalence.
\end{proposition}

\begin{proof}
D'après le lemme \ref{dualizing-lemma-endos}, on a $R = \Hom_R(E, E) = D(E)$.
Bien entendu, on a $E = \Hom_R(R, E) = D(R)$. Puisque $E$ est injectif,
le foncteur $D$ est exact. Le résultat découle alors immédiatement
de la description des catégories donnée dans le lemme \ref{dualizing-lemma-describe-categories}.
\end{proof}

\begin{remark}
\label{dualizing-remark-matlis}
Soit $(R, \mathfrak m, \kappa)$ un anneau local noethérien.
Soit $E$ une enveloppe injective de $\kappa$ sur $R$.
Voici un complément à la dualité de Matlis : si $N$ est un module
de torsion par les puissances de $\mathfrak m$ et si $M = \Hom_R(N, E)$ est un module de type fini
sur le complété de $R$, alors $N$ satisfait à la condition de chaîne
descendante. En effet, pour tous sous-modules $N'' \subset N' \subset N$ tels que $N'' \not = N'$,
il existe un plongement $\kappa \subset N''/N'$, donc un morphisme non nul $N' \to E$
annulant $N''$, que l'on peut prolonger en un morphisme $N \to E$ annulant $N''$.
Ainsi $N \supset N' \mapsto M' = \Hom_R(N/N', E) \subset M$ est une application préservant les inclusions des sous-modules
de $N$ vers les sous-modules de $M$, d'où la conclusion.
\end{remark}




















\section{Dérivation du foncteur de torsion}
\label{dualizing-section-bad-local-cohomology}

\noindent
Soient $A$ un anneau et $I \subset A$ un idéal de type fini
(si $I$ n'est pas de type fini, il faudrait peut-être utiliser
une autre définition). Posons $Z = V(I) \subset \Spec(A)$. Rappelons que la catégorie $I^\infty\text{-torsion}$
des modules de torsion par les puissances de $I$ ne dépend que
du fermé $Z$ et non du choix de l'idéal de type fini $I$ tel que $Z = V(I)$ ;
voir Compléments d'algèbre, lemme \ref{more-algebra-lemma-local-cohomology-closed}.
Dans cette section, nous considérerons le foncteur
$$
H^0_{I} : \text{Mod}_A \longrightarrow I^\infty\text{-torsion},\quad
M \longmapsto M[I^\infty] = \bigcup M[I^n]
$$
qui associe à $M$ son sous-module de torsion par les puissances de $I$.

\medskip\noindent
Soient $A$ un anneau et $I$ un idéal de type fini.
Notons que $I^\infty\text{-torsion}$ est une catégorie abélienne de Grothendieck
(les sommes directes existent, les limites inductives filtrantes sont
exactes et $\bigoplus A/I^n$ est un générateur d'après
Compléments d'algèbre, lemme \ref{more-algebra-lemma-I-power-torsion-presentation}).
La catégorie dérivée $D(I^\infty\text{-torsion})$ existe donc ; voir Injectifs, remarque \ref{injectives-remark-existence-D}.
Notre foncteur $H^0_I$ est exact à gauche et possède un prolongement dérivé
que nous noterons
$$
R\Gamma_I : D(A) \longrightarrow D(I^\infty\text{-torsion}).
$$
{\bf Attention :} ce foncteur ne mérite le nom de cohomologie
locale que si l'anneau $A$ est noethérien.
Les foncteurs $H^0_I$, $R\Gamma_I$ et les satellites $H^p_I$
ne dépendent que du fermé $Z \subset \Spec(A)$ et non du choix de l'idéal
de type fini $I$ tel que $V(I) = Z$.
Nous tenons cependant à utiliser l'indice $I$ pour les foncteurs
ci-dessus, car la notation $R\Gamma_Z$ sera réservée à un autre foncteur,
voir (\ref{dualizing-equation-local-cohomology}), qui ne coïncide avec le foncteur $R\Gamma_I$
(à notre connaissance) que si $A$ est noethérien
(voir le lemme \ref{dualizing-lemma-local-cohomology-noetherian}).

\begin{lemma}
\label{dualizing-lemma-adjoint}
Soient $A$ un anneau et $I \subset A$ un idéal de type fini.
Le foncteur $R\Gamma_I$ est adjoint à droite du foncteur
$D(I^\infty\text{-torsion}) \to D(A)$.
\end{lemma}

\begin{proof}
Cela résulte du fait que le passage au sous-module de torsion
par les puissances de $I$ est adjoint à droite du foncteur d'inclusion
$I^\infty\text{-torsion} \to \text{Mod}_A$. Voir Catégories dérivées, lemme \ref{derived-lemma-derived-adjoint-functors}.
\end{proof}

\begin{lemma}
\label{dualizing-lemma-local-cohomology-ext}
Soient $A$ un anneau et $I \subset A$ un idéal de type fini.
Pour tout objet $K$ de $D(A)$, on a
$$
R\Gamma_I(K) = \text{hocolim}\ R\Hom_A(A/I^n, K)
$$
dans $D(A)$ et
$$
R^q\Gamma_I(K) = \colim_n \Ext_A^q(A/I^n, K)
$$
en tant que modules, pour tout $q \in \mathbf{Z}$.
\end{lemma}

\begin{proof}
Soit $J^\bullet$ un complexe K-injectif représentant $K$. Alors
$$
R\Gamma_I(K) = J^\bullet[I^\infty] = \colim J^\bullet[I^n] =
\colim \Hom_A(A/I^n, J^\bullet)
$$
où la première égalité est la définition de $R\Gamma_I(K)$.
Catégories dérivées, lemme \ref{derived-lemma-colim-hocolim}, donne la première égalité
affichée dans l'énoncé. La seconde en résulte, puisque $H^q(\Hom_A(A/I^n, J^\bullet)) = \Ext^q_A(A/I^n, K)$
et puisque les limites inductives filtrantes sont exactes
dans la catégorie des groupes abéliens.
\end{proof}

\begin{lemma}
\label{dualizing-lemma-bad-local-cohomology-vanishes}
Soient $A$ un anneau et $I \subset A$ un idéal de type fini.
Soit $K^\bullet$ un complexe de $A$-modules tel que
$f : K^\bullet \to K^\bullet$ soit un isomorphisme pour un certain $f \in I$,
c'est-à-dire que $K^\bullet$ soit un complexe de $A_f$-modules. Alors
$R\Gamma_I(K^\bullet) = 0$.
\end{lemma}

\begin{proof}
En effet, les modules de cohomologie de $R\Gamma_I(K^\bullet)$ sont alors de torsion
par les puissances de $f$ et $f$ y agit par automorphismes.
Ces modules de cohomologie sont donc nuls, et l'objet est donc nul.
\end{proof}

\noindent
Soient $A$ un anneau et $I \subset A$ un idéal de type fini.
D'après Compléments d'algèbre, lemme \ref{more-algebra-lemma-I-power-torsion}, la catégorie des modules
de torsion par les puissances de $I$ est une sous-catégorie de Serre
de la catégorie de tous les $A$-modules. On a donc un foncteur
\begin{equation}
\label{dualizing-equation-compare-torsion}
D(I^\infty\text{-torsion}) \longrightarrow D_{I^\infty\text{-torsion}}(A)
\end{equation}
où le membre de droite est la sous-catégorie pleine de $D(A)$
formée des objets dont les cohomologies sont des modules de torsion
par les puissances de $I$. Le foncteur et la catégorie sont tous deux
étudiés dans Catégories dérivées, section \ref{derived-section-triangulated-sub}.

\begin{lemma}
\label{dualizing-lemma-not-equal}
Soient $A$ un anneau et $I$ un idéal de type fini.
Soient $M$ et $N$ des modules de torsion par les puissances de $I$.
\begin{enumerate}
\item $\Hom_{D(A)}(M, N) = \Hom_{D({I^\infty\text{-torsion}})}(M, N)$ ;
\item $\Ext^1_{D(A)}(M, N) =
\Ext^1_{D({I^\infty\text{-torsion}})}(M, N)$ ;
\item $\Ext^2_{D({I^\infty\text{-torsion}})}(M, N) \to
\Ext^2_{D(A)}(M, N)$ n'est pas surjectif en général ;
\item (\ref{dualizing-equation-compare-torsion}) n'est pas une équivalence en général.
\end{enumerate}
\end{lemma}

\begin{proof}
Les assertions (1) et (2) résultent immédiatement de ce que les modules
de torsion par les puissances de $I$ forment une sous-catégorie de Serre
de $\text{Mod}_A$. L'assertion (4) découle de (3).

\medskip\noindent
Pour (3), soit $A$ un anneau possédant un élément $f \in A$ tel que
$A[f]$ contienne un élément non nul $x$ annulé par $f$ et que
$A$ contienne des éléments $x_n$ tels que $f^nx_n = x$.
Un tel anneau $A$ existe, car on peut prendre
$$
A = \mathbf{Z}[f, x, x_n]/(fx, f^nx_n - x)
$$
Étant donné $A$, posons $I = (f)$. Alors la suite exacte
$$
0 \to A[f] \to A \xrightarrow{f} A \to A/fA \to 0
$$
définit un élément de $\Ext^2_A(A/fA, A[f])$. Affirmons que cet élément
ne provient pas d'un élément de
$\Ext^2_{D(f^\infty\text{-torsion})}(A/fA, A[f])$.
En effet, sinon, il existerait une suite exacte
$$
0 \to A[f] \to M \to N \to A/fA \to 0
$$
où $M$ et $N$ sont des modules de torsion par les puissances de $f$,
définissant la même classe de $2$-extension.
Puisque $A \to A$ est un complexe de modules libres et que les classes
de $2$-extensions sont les mêmes, on pourrait trouver un morphisme
$$
\xymatrix{
0 \ar[r] &
A[f] \ar[r] \ar[d] &
A \ar[r] \ar[d]_\varphi &
A \ar[r] \ar[d]_\psi &
A/fA \ar[r] \ar[d] & 0 \\
0 \ar[r] &
A[f] \ar[r] &
M \ar[r] &
N \ar[r] &
A/fA \ar[r] & 0
}
$$
(nous omettons certains détails). On pourrait alors remplacer $M$
par l'image de $\varphi$ et $N$ par l'image de $\psi$.
Alors $M$ serait un module cyclique, donc $f^n M = 0$ pour un certain $n$.
En considérant $\varphi(x_{n + 1})$, on obtient une contradiction avec le fait que
$f^{n + 1}x_n = x$ est non nul dans $A[f]$.
\end{proof}









\section{Cohomologie locale}
\label{dualizing-section-local-cohomology}

\noindent
Soient $A$ un anneau et $I \subset A$ un idéal de type fini.
Posons $Z = V(I) \subset \Spec(A)$. Nous construirons un foncteur
\begin{equation}
\label{dualizing-equation-local-cohomology}
R\Gamma_Z : D(A) \longrightarrow D_{I^\infty\text{-torsion}}(A).
\end{equation}
adjoint à droite du foncteur d'inclusion. Pour les notations,
voir la section \ref{dualizing-section-bad-local-cohomology}. Les modules de cohomologie de $R\Gamma_Z(K)$
sont les {\it groupes de cohomologie locale de $K$ relativement à $Z$}.
D'après le lemme \ref{dualizing-lemma-not-equal}, ce foncteur ne sera en général {\bf pas}
égal à $R\Gamma_I( - )$, même en les considérant comme foncteurs à valeurs dans $D(A)$.
Dans la section \ref{dualizing-section-local-cohomology-noetherian}, nous montrerons que, si $A$ est noethérien,
les deux coïncident.

\medskip\noindent
Nous poursuivrons l'étude de la cohomologie locale dans le chapitre
qui lui est consacré ; voir Cohomologie locale, section \ref{local-cohomology-section-introduction}.
Nous y montrerons notamment que $R\Gamma_Z$ calcule la cohomologie
à support dans $Z$ du complexe associé de faisceaux quasi-cohérents
sur $\Spec(A)$. Voir Cohomologie locale, lemme \ref{local-cohomology-lemma-local-cohomology-is-local-cohomology}.


\begin{lemma}
\label{dualizing-lemma-local-cohomology-adjoint}
Soient $A$ un anneau et $I \subset A$ un idéal de type fini.
Il existe un adjoint à droite $R\Gamma_Z$ (\ref{dualizing-equation-local-cohomology})
du foncteur d'inclusion $D_{I^\infty\text{-torsion}}(A) \to D(A)$.
En fait, si $I$ est engendré par $f_1, \ldots, f_r \in A$, on a
$$
R\Gamma_Z(K) =
(A \to \prod\nolimits_{i_0} A_{f_{i_0}} \to
\prod\nolimits_{i_0 < i_1} A_{f_{i_0}f_{i_1}}
\to \ldots \to A_{f_1\ldots f_r}) \otimes_A^\mathbf{L} K
$$
fonctoriellement en $K \in D(A)$.
\end{lemma}

\begin{proof}
Soit $I = (f_1, \ldots, f_r)$ un idéal.
Soit $K^\bullet$ un complexe de $A$-modules.
Il existe un morphisme canonique de complexes
$$
(A \to \prod\nolimits_{i_0} A_{f_{i_0}} \to
\prod\nolimits_{i_0 < i_1} A_{f_{i_0}f_{i_1}} \to
\ldots \to A_{f_1\ldots f_r}) \longrightarrow A.
$$
du complexe de {\v C}ech augmenté vers $A$.
En tensorisant par $K^\bullet$ et en prenant le complexe total associé,
on obtient un morphisme
$$
\text{Tot}\left(
K^\bullet \otimes_A
(A \to \prod\nolimits_{i_0} A_{f_{i_0}} \to
\prod\nolimits_{i_0 < i_1} A_{f_{i_0}f_{i_1}} \to
\ldots \to A_{f_1\ldots f_r})\right)
\longrightarrow
K^\bullet
$$
dans $D(A)$. Affirmons que les modules de cohomologie du complexe
de gauche sont de torsion par les puissances de $I$,
c'est-à-dire que le membre de gauche est un objet de $D_{I^\infty\text{-torsion}}(A)$.
En effet, on a
$$
(A \to \prod\nolimits_{i_0} A_{f_{i_0}} \to
\prod\nolimits_{i_0 < i_1} A_{f_{i_0}f_{i_1}} \to
\ldots \to A_{f_1\ldots f_r}) = \colim K(A, f_1^n, \ldots, f_r^n)
$$
d'après Compléments d'algèbre, lemme \ref{more-algebra-lemma-extended-alternating-Cech-is-colimit-koszul}.
De plus, la multiplication par $f_i^n$ sur le complexe
$K(A, f_1^n, \ldots, f_r^n)$ est homotope à zéro d'après Compléments d'algèbre, lemme \ref{more-algebra-lemma-homotopy-koszul}.
Puisque
$$
H^q\left( LHS \right) =
\colim H^q(\text{Tot}(K^\bullet \otimes_A K(A, f_1^n, \ldots, f_r^n)))
$$
on obtient l'assertion. D'autre part, si $K^\bullet$ est un objet de $D_{I^\infty\text{-torsion}}(A)$,
les complexes $K^\bullet \otimes_A A_{f_{i_0} \ldots f_{i_p}}$ ont une cohomologie nulle.
Dans ce cas, le morphisme $LHS \to K^\bullet$ est donc un isomorphisme dans $D(A)$.
La construction
$$
R\Gamma_Z(K^\bullet) =
\text{Tot}\left(
K^\bullet \otimes_A
(A \to \prod\nolimits_{i_0} A_{f_{i_0}} \to
\prod\nolimits_{i_0 < i_1} A_{f_{i_0}f_{i_1}} \to
\ldots \to A_{f_1\ldots f_r})\right)
$$
est fonctorielle en $K^\bullet$ et définit un foncteur exact
$D(A) \to D_{I^\infty\text{-torsion}}(A)$ entre catégories triangulées.
L'existence de la transformation naturelle $R\Gamma_Z \to \text{id}$ donnée ci-dessus,
et le fait qu'elle soit un isomorphisme sur les $K^\bullet$ appartenant
à la sous-catégorie, entraînent formellement que $R\Gamma_Z$
est l'adjoint à droite cherché.
\end{proof}

\begin{lemma}
\label{dualizing-lemma-local-cohomology-and-restriction}
Soient $A \to B$ un homomorphisme d'anneaux et $I \subset A$ un idéal de type fini.
Posons $J = IB$, puis $Z = V(I)$ et $Y = V(J)$. Alors
$$
R\Gamma_Z(M_A) = R\Gamma_Y(M)_A
$$
fonctoriellement en $M \in D(B)$. Ici $(-)_A$ désigne les restrictions
$D(B) \to D(A)$ et $D_{J^\infty\text{-torsion}}(B) \to D_{I^\infty\text{-torsion}}(A)$.
\end{lemma}

\begin{proof}
Cela résulte de l'unicité des adjoints, puisque $R\Gamma_Z((-)_A)$ et $R\Gamma_Y(-)_A$
sont tous deux adjoints à droite du foncteur $D_{I^\infty\text{-torsion}}(A) \to D(B)$,
$K \mapsto K \otimes_A^\mathbf{L} B$.
On peut aussi utiliser la description de $R\Gamma_Z$ et de $R\Gamma_Y$
par les complexes alternés de {\v C}ech
(lemme \ref{dualizing-lemma-local-cohomology-adjoint}).
En effet, si $I = (f_1, \ldots, f_r)$, alors $J$ est engendré par les images
$g_1, \ldots, g_r \in B$ de $f_1, \ldots, f_r$.
L'énoncé résulte alors de l'existence d'un isomorphisme canonique
\begin{align*}
& M_A \otimes_A (A \to \prod\nolimits_{i_0} A_{f_{i_0}} \to
\prod\nolimits_{i_0 < i_1} A_{f_{i_0}f_{i_1}}
\to \ldots \to A_{f_1\ldots f_r}) \\
& = 
M \otimes_B (B \to \prod\nolimits_{i_0} B_{g_{i_0}} \to
\prod\nolimits_{i_0 < i_1} B_{g_{i_0}g_{i_1}}
\to \ldots \to B_{g_1\ldots g_r})
\end{align*}
pour tout $B$-module $M$.
\end{proof}

\begin{lemma}
\label{dualizing-lemma-torsion-change-rings}
Soient $A \to B$ un homomorphisme d'anneaux et $I \subset A$ un idéal de type fini.
Posons $J = IB$. Soient $Z = V(I)$ et $Y = V(J)$. Alors
$$
R\Gamma_Z(K) \otimes_A^\mathbf{L} B = R\Gamma_Y(K \otimes_A^\mathbf{L} B)
$$
fonctoriellement en $K \in D(A)$.
\end{lemma}

\begin{proof}
Écrivons $I = (f_1, \ldots, f_r)$. Alors $J$ est engendré par les images
$g_1, \ldots, g_r \in B$ de $f_1, \ldots, f_r$. On a donc
$$
(A \to \prod A_{f_{i_0}} \to \ldots \to A_{f_1\ldots f_r}) \otimes_A B =
(B \to \prod B_{g_{i_0}} \to \ldots \to B_{g_1\ldots g_r})
$$
en tant que complexes de $B$-modules. Représentons $K$
par un complexe K-plat $K^\bullet$ de $A$-modules.
Puisque les complexes totaux associés à
$$
K^\bullet \otimes_A
(A \to \prod A_{f_{i_0}} \to \ldots \to A_{f_1\ldots f_r}) \otimes_A B
$$
et à
$$
K^\bullet \otimes_A B \otimes_B
(B \to \prod B_{g_{i_0}} \to \ldots \to B_{g_1\ldots g_r})
$$
représentent les membres de gauche et de droite de la formule
affichée dans le lemme (voir le lemme \ref{dualizing-lemma-local-cohomology-adjoint}), on conclut.
\end{proof}

\begin{lemma}
\label{dualizing-lemma-local-cohomology-vanishes}
Soient $A$ un anneau et $I \subset A$ un idéal de type fini.
Soit $K^\bullet$ un complexe de $A$-modules tel que
$f : K^\bullet \to K^\bullet$ soit un isomorphisme pour un certain $f \in I$,
c'est-à-dire que $K^\bullet$ soit un complexe de $A_f$-modules. Alors
$R\Gamma_Z(K^\bullet) = 0$.
\end{lemma}

\begin{proof}
En effet, les modules de cohomologie de $R\Gamma_Z(K^\bullet)$ sont alors de torsion
par les puissances de $f$ et $f$ y agit par automorphismes.
Ces modules de cohomologie sont donc nuls, et l'objet est donc nul.
\end{proof}

\begin{lemma}
\label{dualizing-lemma-torsion-tensor-product}
Soient $A$ un anneau et $I \subset A$ un idéal de type fini.
Pour $K, L \in D(A)$, on a
$$
R\Gamma_Z(K \otimes_A^\mathbf{L} L) =
K \otimes_A^\mathbf{L} R\Gamma_Z(L) =
R\Gamma_Z(K) \otimes_A^\mathbf{L} L =
R\Gamma_Z(K) \otimes_A^\mathbf{L} R\Gamma_Z(L)
$$
Si $K$ ou $L$ appartient à $D_{I^\infty\text{-torsion}}(A)$,
il en est de même de $K \otimes_A^\mathbf{L} L$.
\end{lemma}

\begin{proof}
D'après le lemme \ref{dualizing-lemma-local-cohomology-adjoint}, on sait que $R\Gamma_Z$ est donné par $C \otimes^\mathbf{L} -$
pour un certain $C \in D(A)$. Ainsi, pour $K, L \in D(A)$ quelconques, on a
$$
R\Gamma_Z(K \otimes_A^\mathbf{L} L) =
K \otimes^\mathbf{L} L \otimes_A^\mathbf{L} C =
K \otimes_A^\mathbf{L} R\Gamma_Z(L)
$$
Les autres égalités en découlent formellement.
Cela entraîne aussi la dernière assertion du lemme.
\end{proof}

\begin{lemma}
\label{dualizing-lemma-local-cohomology-ss}
Soient $A$ un anneau et $I, J \subset A$ des idéaux de type fini.
Posons $Z = V(I)$ et $Y = V(J)$. Alors $Z \cap Y = V(I + J)$
et $R\Gamma_Y \circ R\Gamma_Z = R\Gamma_{Y \cap Z}$ en tant que foncteurs
$D(A) \to D_{(I + J)^\infty\text{-torsion}}(A)$. Pour $K \in D^+(A)$,
il existe une suite spectrale
$$
E_2^{p, q} = H^p_Y(H^q_Z(K)) \Rightarrow H^{p + q}_{Y \cap Z}(K)
$$
comme dans Catégories dérivées, lemme \ref{derived-lemma-grothendieck-spectral-sequence}.
\end{lemma}

\begin{proof}
Le lemme comporte un léger abus de notation, car, à proprement parler,
on ne peut pas composer $R\Gamma_Y$ et $R\Gamma_Z$.
L'énoncé signifie simplement que l'on compose $R\Gamma_Z$ avec l'inclusion $D_{I^\infty\text{-torsion}}(A) \to D(A)$,
puis avec $R\Gamma_Y$. L'égalité $R\Gamma_Y \circ R\Gamma_Z = R\Gamma_{Y \cap Z}$ résulte alors de ce que
$$
D_{I^\infty\text{-torsion}}(A) \to D(A) \xrightarrow{R\Gamma_Y}
D_{(I + J)^\infty\text{-torsion}}(A)
$$
est adjoint à droite de l'inclusion
$D_{(I + J)^\infty\text{-torsion}}(A) \to D_{I^\infty\text{-torsion}}(A)$.
On peut aussi démontrer la formule à l'aide du lemme \ref{dualizing-lemma-local-cohomology-adjoint}
et du fait que le produit tensoriel des complexes de
{\v C}ech augmentés sur $f_1, \ldots, f_r$ et sur $g_1, \ldots, g_m$
est le complexe de {\v C}ech augmenté sur $f_1, \ldots, f_n. g_1, \ldots, g_m$.
La dernière assertion en résulte, avec le lemme cité.
\end{proof}

\noindent
Le lemme suivant est l'analogue de Compléments d'algèbre, lemme \ref{more-algebra-lemma-restriction-derived-complete-equivalence},
pour les complexes dont les cohomologies sont de torsion.

\begin{lemma}
\label{dualizing-lemma-torsion-flat-change-rings}
Soient $A \to B$ un homomorphisme plat d'anneaux et $I \subset A$ un idéal de type fini
tel que $A/I = B/IB$. Le changement de base et la restriction induisent alors
des équivalences quasi-inverses
$D_{I^\infty\text{-torsion}}(A) = D_{(IB)^\infty\text{-torsion}}(B)$.
\end{lemma}

\begin{proof}
Plus précisément, les foncteurs sont $K \mapsto K \otimes_A^\mathbf{L} B$
pour $K$ dans $D_{I^\infty\text{-torsion}}(A)$ et $M \mapsto M_A$
pour $M$ dans $D_{(IB)^\infty\text{-torsion}}(B)$. Cela fonctionne parce que $H^i(K \otimes_A^\mathbf{L} B) = H^i(K) \otimes_A B = H^i(K)$.
La première égalité résulte de la platitude de $A \to B$ et la seconde
de Compléments d'algèbre, lemme \ref{more-algebra-lemma-neighbourhood-isomorphism}.
\end{proof}

\noindent
Le lemme suivant a été démontré pour les $\Hom$ et les $\Ext^1$ de modules
dans Compléments d'algèbre, lemmes \ref{more-algebra-lemma-neighbourhood-equivalence} et \ref{more-algebra-lemma-neighbourhood-extensions}.

\begin{lemma}
\label{dualizing-lemma-neighbourhood-extensions}
Soient $A \to B$ un homomorphisme plat d'anneaux et $I \subset A$
un idéal de type fini tel que $A/I \to B/IB$ soit un isomorphisme.
Pour $K \in D_{I^\infty\text{-torsion}}(A)$ et $L \in D(A)$, le morphisme
$$
R\Hom_A(K, L) \longrightarrow R\Hom_B(K \otimes_A B, L \otimes_A B)
$$
est un quasi-isomorphisme. En particulier, si $M$, $N$
sont des $A$-modules et si $M$ est de torsion par les puissances de $I$,
le morphisme canonique
$$
\Ext^i_A(M, N)
\longrightarrow
\Ext^i_B(M \otimes_A B, N \otimes_A B)
$$
est un isomorphisme pour tout $i$.
\end{lemma}

\begin{proof}
Soient $Z = V(I) \subset \Spec(A)$ et $Y = V(IB) \subset \Spec(B)$.
Puisque les modules de cohomologie de $K$ sont de torsion par les puissances
de $I$, le morphisme canonique $R\Gamma_Z(L) \to L$ induit un isomorphisme
$$
R\Hom_A(K, R\Gamma_Z(L)) \to R\Hom_A(K, L)
$$
dans $D(A)$. De même, les modules de cohomologie de $K \otimes_A B$
sont de torsion par les puissances de $IB$ et l'on a un isomorphisme
$$
R\Hom_B(K \otimes_A B, R\Gamma_Y(L \otimes_A B)) \to 
R\Hom_B(K \otimes_A B, L \otimes_A B)
$$
dans $D(B)$.
D'après le lemme \ref{dualizing-lemma-torsion-change-rings}, on a $R\Gamma_Z(L) \otimes_A B = R\Gamma_Y(L \otimes_A B)$.
Il suffit donc de montrer que le morphisme
$$
R\Hom_A(K, R\Gamma_Z(L)) \to R\Hom_B(K \otimes_A B, R\Gamma_Z(L) \otimes_A B)
$$
est un quasi-isomorphisme, ce qui résulte du lemme \ref{dualizing-lemma-torsion-flat-change-rings}.
\end{proof}




\section{Cohomologie locale pour les anneaux noethériens}
\label{dualizing-section-local-cohomology-noetherian}

\noindent
Soient $A$ un anneau et $I \subset A$ un idéal de type fini.
Posons $Z = V(I) \subset \Spec(A)$. Rappelons que (\ref{dualizing-equation-compare-torsion}) est le foncteur
$$
D(I^\infty\text{-torsion}) \to D_{I^\infty\text{-torsion}}(A)
$$
En fait, il existe une transformation naturelle de foncteurs
\begin{equation}
\label{dualizing-equation-compare-torsion-functors}
(\ref{dualizing-equation-compare-torsion}) \circ R\Gamma_I(-)
\longrightarrow
R\Gamma_Z(-)
\end{equation}
En effet, étant donné un complexe de $A$-modules $K^\bullet$,
le morphisme canonique $R\Gamma_I(K^\bullet) \to K^\bullet$ dans $D(A)$ se factorise
(de manière unique) par $R\Gamma_Z(K^\bullet)$, puisque $R\Gamma_I(K^\bullet)$ a des modules
de cohomologie de torsion par les puissances de $I$
(voir le lemme \ref{dualizing-lemma-adjoint}).
En général, ce morphisme n'est pas un isomorphisme
(nous l'avons vu dans le lemme \ref{dualizing-lemma-not-equal}).

\begin{lemma}
\label{dualizing-lemma-local-cohomology-noetherian}
Soient $A$ un anneau noethérien et $I \subset A$ un idéal.
\begin{enumerate}
\item Le morphisme d'adjonction $R\Gamma_I(K) \to K$ est un isomorphisme
pour $K \in D_{I^\infty\text{-torsion}}(A)$ ;
\item le foncteur
(\ref{dualizing-equation-compare-torsion})
$D(I^\infty\text{-torsion}) \to D_{I^\infty\text{-torsion}}(A)$
est une équivalence ;
\item la transformation de foncteurs
(\ref{dualizing-equation-compare-torsion-functors}) est un isomorphisme,
autrement dit $R\Gamma_I(K) = R\Gamma_Z(K)$ pour $K \in D(A)$.
\end{enumerate}
\end{lemma}

\begin{proof}
Un argument formel, que nous omettons, montre qu'il suffit de prouver (1).

\medskip\noindent
Soit $M$ un module de torsion par les puissances de $I$ sur $A$.
Choisissons un plongement $M \to J$ dans un $A$-module injectif.
Alors $J[I^\infty]$ est un $A$-module injectif, voir le lemme \ref{dualizing-lemma-injective-module-divide},
et l'on obtient un plongement $M \to J[I^\infty]$.
Ainsi tout module de torsion par les puissances de $I$
possède une résolution injective $M \to J^\bullet$ dont les $J^n$
sont aussi de torsion par les puissances de $I$.
Il en résulte que $R\Gamma_I(M) = M$ (ce n'est pas une trivialité,
et ce n'est pas vrai en général si $A$ n'est pas noethérien).
Supposons ensuite que $K \in D_{I^\infty\text{-torsion}}^+(A)$. La suite spectrale
$$
R^q\Gamma_I(H^p(K)) \Rightarrow R^{p + q}\Gamma_I(K)
$$
(Catégories dérivées, lemme \ref{derived-lemma-two-ss-complex-functor}) converge,
et nous avons vu ci-dessus que seuls les termes avec $q = 0$
peuvent être non nuls. Ainsi $R\Gamma_I(K) \to K$ est un isomorphisme.

\medskip\noindent
Supposons que $K$ soit un objet quelconque de $D_{I^\infty\text{-torsion}}(A)$.
On a
$$
R^q\Gamma_I(K) = \colim \Ext^q_A(A/I^n, K)
$$
par le lemme \ref{dualizing-lemma-local-cohomology-ext}. Prenons $f_1, \ldots, f_r \in A$ engendrant $I$.
Soit $K_n^\bullet = K(A, f_1^n, \ldots, f_r^n)$ le complexe de Koszul placé en degrés $-r, \ldots, 0$.
Puisque les pro-objets $\{A/I^n\}$ et $\{K_n^\bullet\}$ dans $D(A)$ sont les mêmes
d'après Compléments d'algèbre, lemme \ref{more-algebra-lemma-sequence-Koszul-complexes}, on voit que
$$
R^q\Gamma_I(K) = \colim \Ext^q_A(K_n^\bullet, K)
$$
Choisissons un complexe quelconque $K^\bullet$ de $A$-modules représentant $K$.
Puisque $K_n^\bullet$ est un complexe fini de modules libres de type fini,
on voit que
$$
\Ext^q_A(K_n, K) =
H^q(\text{Tot}((K_n^\bullet)^\vee \otimes_A K^\bullet))
$$
où $(K_n^\bullet)^\vee$ est le dual du complexe $K_n^\bullet$.
Voir Compléments d'algèbre, lemme \ref{more-algebra-lemma-RHom-out-of-projective}.
Puisque $(K_n^\bullet)^\vee$ est un complexe de $A$-modules libres de type fini
placés en degrés $0, \ldots, r$, les termes du complexe $\text{Tot}((K_n^\bullet)^\vee \otimes_A K^\bullet)$
sont les mêmes que ceux du complexe $\text{Tot}((K_n^\bullet)^\vee \otimes_A \tau_{\geq q - r - 2} K^\bullet)$
en degrés $q - 1$ et supérieurs. On voit donc que
$$
\Ext^q_A(K_n, K) = \text{Ext}^q_A(K_n, \tau_{\geq q - r - 2}K)
$$
pour tout $n$. Il en résulte que
$$
R^q\Gamma_I(K) = R^q\Gamma_I(\tau_{\geq q - r - 2}K) =
H^q(\tau_{\geq q - r - 2}K) = H^q(K)
$$
Ainsi le morphisme $R\Gamma_I(K) \to K$ est un isomorphisme.
\end{proof}

\begin{lemma}
\label{dualizing-lemma-compute-local-cohomology-noetherian}
Soient $A$ un anneau noethérien et $I = (f_1, \ldots, f_r)$ un idéal de $A$.
Posons $Z = V(I) \subset \Spec(A)$. On a des isomorphismes canoniques
$$
R\Gamma_I(A) \to
(A \to \prod\nolimits_{i_0} A_{f_{i_0}} \to
\prod\nolimits_{i_0 < i_1} A_{f_{i_0}f_{i_1}} \to
\ldots \to A_{f_1\ldots f_r}) \to R\Gamma_Z(A)
$$
dans $D(A)$. Si $M$ est un $A$-module, on a de même
$$
R\Gamma_I(M) \cong
(M \to \prod\nolimits_{i_0} M_{f_{i_0}} \to
\prod\nolimits_{i_0 < i_1} M_{f_{i_0}f_{i_1}} \to
\ldots \to M_{f_1\ldots f_r}) \cong R\Gamma_Z(M)
$$
dans $D(A)$.
\end{lemma}

\begin{proof}
Cela résulte du lemme \ref{dualizing-lemma-local-cohomology-noetherian} et du calcul du foncteur $R\Gamma_Z$
donné dans le lemme \ref{dualizing-lemma-local-cohomology-adjoint}.
\end{proof}

\begin{lemma}
\label{dualizing-lemma-local-cohomology-change-rings}
Si $A \to B$ est un homomorphisme d'anneaux noethériens et $I \subset A$
un idéal, alors on a dans $D(B)$
$$
R\Gamma_I(A) \otimes_A^\mathbf{L} B =
R\Gamma_Z(A) \otimes_A^\mathbf{L} B =
R\Gamma_Y(B) = R\Gamma_{IB}(B)
$$
où $Y = V(IB) \subset \Spec(B)$.
\end{lemma}

\begin{proof}
Il suffit de combiner les lemmes \ref{dualizing-lemma-compute-local-cohomology-noetherian} et \ref{dualizing-lemma-torsion-change-rings}.
\end{proof}






\section{Profondeur}
\label{dualizing-section-depth}

\noindent
Dans cette section, nous revenons sur la notion de profondeur introduite
dans Algèbre, section \ref{algebra-section-depth}.

\begin{lemma}
\label{dualizing-lemma-depth}
Soient $A$ un anneau noethérien, $I \subset A$ un idéal
et $M$ un $A$-module de type fini tel que $IM \not = M$.
Alors les entiers suivants sont égaux :
\begin{enumerate}
\item $\text{profondeur}_I(M)$ ;
\item le plus petit entier $i$ tel que $\Ext_A^i(A/I, M)$ soit non nul ;
\item le plus petit entier $i$ tel que $H^i_I(M)$ soit non nul.
\end{enumerate}
De plus, on a $\Ext^i_A(N, M) = 0$ pour $i < \text{profondeur}_I(M)$
et pour tout $A$-module de type fini $N$ annulé par une puissance de $I$.
\end{lemma}

\begin{proof}
Montrons par récurrence sur $\text{profondeur}_I(M)$ que (1) et (2) sont égaux,
ce qui est permis par Algèbre, lemme \ref{algebra-lemma-depth-finite-noetherian}.

\medskip\noindent
Cas initial. Si $\text{profondeur}_I(M) = 0$, alors $I$ est contenu dans la réunion
des idéaux premiers associés de $M$
(Algèbre, lemme \ref{algebra-lemma-ass-zero-divisors}).
Par évitement des idéaux premiers (Algèbre, lemme \ref{algebra-lemma-silly}),
on voit que $I \subset \mathfrak p$ pour un idéal premier associé $\mathfrak p$.
Ainsi $\Hom_A(A/I, M)$ est non nul. L'égalité est donc vraie dans ce cas.

\medskip\noindent
Supposons que $\text{profondeur}_I(M) > 0$. Soit $f \in I$ un non-diviseur de zéro
sur $M$ tel que $\text{profondeur}_I(M/fM) = \text{profondeur}_I(M) - 1$.
Considérons la suite exacte courte
$$
0 \to M \to M \to M/fM \to 0
$$
et la suite exacte longue associée pour $\Ext^*_A(A/I, -)$.
Notons que $\Ext^i_A(A/I, M)$ est un $A/I$-module de type fini
(Algèbre, lemmes \ref{algebra-lemma-ext-noetherian} et \ref{algebra-lemma-annihilate-ext}). On obtient donc
$$
\Hom_A(A/I, M/fM) = \Ext^1_A(A/I, M)
$$
et des suites exactes courtes
$$
0 \to \Ext^i_A(A/I, M) \to \text{Ext}^i_A(A/I, M/fM) \to
\Ext^{i + 1}_A(A/I, M) \to 0
$$
D'où l'égalité de (1) et (2) par récurrence.

\medskip\noindent
Observons que $\text{profondeur}_I(M) = \text{profondeur}_{I^n}(M)$ pour tout $n \geq 1$,
par exemple d'après Algèbre, lemme \ref{algebra-lemma-regular-sequence-powers}.
L'égalité de (1) et (2) donne donc $\Ext^i_A(A/I^n, M) = 0$ pour tous $n$ et $i < \text{profondeur}_I(M)$.
Soit $N$ un $A$-module de type fini annulé par une puissance de $I$.
On peut alors choisir une suite exacte courte
$$
0 \to N' \to (A/I^n)^{\oplus m} \to N \to 0
$$
pour certains $n, m \geq 0$. Alors
$\Hom_A(N, M) \subset \Hom_A((A/I^n)^{\oplus m}, M)$
et
$\Ext^i_A(N, M) \subset \text{Ext}^{i - 1}_A(N', M)$
pour $i < \text{profondeur}_I(M)$. Un simple raisonnement par récurrence
établit donc la dernière assertion du lemme.

\medskip\noindent
Enfin, prouvons que (3) est égal à (1) et (2).
On a $H^p_I(M) = \colim \Ext^p_A(A/I^n, M)$ d'après le lemme \ref{dualizing-lemma-local-cohomology-ext}.
On voit donc que $H^i_I(M) = 0$ pour $i < \text{profondeur}_I(M)$.
Pour $i = \text{profondeur}_I(M)$, l'annulation de $\Ext_A^{i - 1}(I/I^n, M)$ montre que le morphisme
$\Ext_A^i(A/I, M) \to H_I^i(M)$ est injectif, ce qui prouve la non-annulation au degré voulu.
\end{proof}

\begin{lemma}
\label{dualizing-lemma-depth-in-ses}
Soit $A$ un anneau noethérien. Soit $0 \to N' \to N \to N'' \to 0$
une suite exacte courte de $A$-modules de type fini.
Soit $I \subset A$ un idéal.
\begin{enumerate}
\item
$\text{profondeur}_I(N) \geq \min\{\text{profondeur}_I(N'), \text{profondeur}_I(N'')\}$
\item
$\text{profondeur}_I(N'') \geq \min\{\text{profondeur}_I(N), \text{profondeur}_I(N') - 1\}$
\item
$\text{profondeur}_I(N') \geq \min\{\text{profondeur}_I(N), \text{profondeur}_I(N'') + 1\}$
\end{enumerate}
\end{lemma}

\begin{proof}
Supposons $IN \not = N$, $IN' \not = N'$ et $IN'' \not = N''$.
On peut alors utiliser la caractérisation de la profondeur
par les groupes Ext $\Ext^i(A/I, N)$, voir le lemme \ref{dualizing-lemma-depth},
ainsi que la suite exacte longue de cohomologie
$$
\begin{matrix}
0
\to \Hom_A(A/I, N')
\to \Hom_A(A/I, N)
\to \Hom_A(A/I, N'')
\\
\phantom{0\ }
\to \Ext^1_A(A/I, N')
\to \Ext^1_A(A/I, N)
\to \Ext^1_A(A/I, N'')
\to \ldots
\end{matrix}
$$
donnée par Algèbre, lemme \ref{algebra-lemma-long-exact-seq-ext}.
Cet argument fonctionne aussi si $IN = N$,
car alors $\Ext^i_A(A/I, N) = 0$ pour tout $i$.
De même dans le cas où $IN' \not = N'$ et/ou $IN'' \not = N''$.
\end{proof}

\begin{lemma}
\label{dualizing-lemma-depth-drops-by-one}
Soient $A$ un anneau noethérien, $I \subset A$ un idéal et $M$
un $A$-module de type fini tel que $IM \not = M$.
\begin{enumerate}
\item Si $x \in I$ est un non-diviseur de zéro sur $M$, alors
$\text{profondeur}_I(M/xM) = \text{profondeur}_I(M) - 1$.
\item Toute suite $M$-régulière $x_1, \ldots, x_r$ dans $I$
peut être prolongée en une suite $M$-régulière dans $I$
de longueur $\text{profondeur}_I(M)$.
\end{enumerate}
\end{lemma}

\begin{proof}
L'assertion (2) est une conséquence formelle de (1).
Soit $x \in I$ comme dans (1). La suite exacte courte $0 \to M \to M \to M/xM \to 0$
et le lemme \ref{dualizing-lemma-depth-in-ses} donnent $\text{profondeur}_I(M/xM) \geq \text{profondeur}_I(M) - 1$.
D'autre part, si $x_1, \ldots, x_r \in I$ est une suite régulière pour $M/xM$,
alors $x, x_1, \ldots, x_r$ est une suite régulière pour $M$. D'où (1).
\end{proof}

\begin{lemma}
\label{dualizing-lemma-depth-CM}
Soit $R$ un anneau local noethérien. Si $M$ est un
$R$-module de Cohen-Macaulay de type fini et $I \subset R$ un idéal non trivial, alors
$$
\text{profondeur}_I(M) = \dim(\text{Supp}(M)) - \dim(\text{Supp}(M/IM)).
$$
\end{lemma}

\begin{proof}
Nous allons le démontrer par récurrence sur $\text{profondeur}_I(M)$.

\medskip\noindent
Si $\text{profondeur}_I(M) = 0$, alors $I$ est contenu dans l'un des idéaux premiers
associés $\mathfrak p$ de $M$ (Algèbre, lemme \ref{algebra-lemma-ideal-nonzerodivisor}).
Alors $\mathfrak p \in \text{Supp}(M/IM)$, d'où
$\dim(\text{Supp}(M/IM)) \geq \dim(R/\mathfrak p) = \dim(\text{Supp}(M))$
où l'égalité résulte d'Algèbre, lemme \ref{algebra-lemma-CM-ass-minimal-support}.
Le lemme est donc vrai dans ce cas.

\medskip\noindent
Si $\text{profondeur}_I(M) > 0$, choisissons $x \in I$ qui soit un non-diviseur de zéro sur $M$.
Notons que $(M/xM)/I(M/xM) = M/IM$. D'autre part, on a
$\text{profondeur}_I(M/xM) = \text{profondeur}_I(M) - 1$
d'après le lemme \ref{dualizing-lemma-depth-drops-by-one}, et $\dim(\text{Supp}(M/xM)) = \dim(\text{Supp}(M)) -  1$
d'après Algèbre, lemme \ref{algebra-lemma-one-equation-module}.
Le résultat découle donc de l'hypothèse de récurrence.
\end{proof}

\begin{lemma}
\label{dualizing-lemma-depth-flat-CM}
Soit $R \to S$ un homomorphisme local plat d'anneaux locaux noethériens.
Notons $\mathfrak m \subset R$ l'idéal maximal.
Soit $I \subset S$ un idéal.
Si $S/\mathfrak mS$ est de Cohen-Macaulay, alors
$$
\text{profondeur}_I(S) \geq \dim(S/\mathfrak mS) - \dim(S/\mathfrak mS + I)
$$
\end{lemma}

\begin{proof}
D'après Algèbre, lemme \ref{algebra-lemma-grothendieck-regular-sequence}, toute suite dans $S$
dont l'image est une suite régulière dans $S/\mathfrak mS$
est une suite régulière dans $S$.
Il suffit donc de prouver le lemme lorsque $R$ est un corps.
C'est un cas particulier du lemme \ref{dualizing-lemma-depth-CM}.
\end{proof}

\begin{lemma}
\label{dualizing-lemma-divide-by-torsion}
Soient $A$ un anneau et $I \subset A$ un idéal de type fini.
Soit $M$ un $A$-module. Soit $Z = V(I)$.
Alors $H^0_I(M) = H^0_Z(M)$. Notons $N$ cette valeur commune
et posons $M' = M/N$. Alors :
\begin{enumerate}
\item $H^0_I(M') = 0$, $H^p_I(M) = H^p_I(M')$ et $H^p_I(N) = 0$
pour tout $p > 0$ ;
\item $H^0_Z(M') = 0$, $H^p_Z(M) = H^p_Z(M')$ et $H^p_Z(N) = 0$
pour tout $p > 0$.
\end{enumerate}
\end{lemma}

\begin{proof}
Par définition, $H^0_I(M) = M[I^\infty]$ est de torsion par les puissances de $I$.
D'après le lemme \ref{dualizing-lemma-local-cohomology-adjoint}, on voit que
$$
H^0_Z(M) = \Ker(M \longrightarrow M_{f_1} \times \ldots \times M_{f_r})
$$
si $I = (f_1, \ldots, f_r)$. Ainsi $H^0_I(M) \subset H^0_Z(M)$ et, réciproquement, si $x \in H^0_Z(M)$,
cet élément est annulé par un $f_i^{e_i}$ pour un certain $e_i \geq 1$,
donc par une puissance de $I$.
Cela prouve la première égalité ; de plus, $N$ est de torsion
par les puissances de $I$.
D'après le lemme \ref{dualizing-lemma-adjoint}, on voit que $R\Gamma_I(N) = N$.
D'après le lemme \ref{dualizing-lemma-local-cohomology-adjoint}, on voit que $R\Gamma_Z(N) = N$.
Cela prouve l'annulation de $H^p_I(N)$ et de $H^p_Z(N)$ en degrés strictement
positifs dans (1) et (2).
L'annulation de $H^0_I(M')$ et de $H^0_Z(M')$ résulte des remarques précédentes
et du fait que $M'$ est sans torsion par les puissances de $I$
d'après Compléments d'algèbre, lemme \ref{more-algebra-lemma-divide-by-torsion}.
L'égalité des cohomologies supérieures de $M$ et de $M'$
découle immédiatement de la suite exacte longue de cohomologie.
\end{proof}









\section{Modules de torsion et modules complets}
\label{dualizing-section-torsion-and-complete}

\noindent
Soient $A$ un anneau et $I$ un idéal de type fini.
On peut alors considérer la catégorie dérivée $D_{I^\infty\text{-torsion}}(A)$
des complexes dont les modules de cohomologie sont de torsion
par les puissances de $I$ (section \ref{dualizing-section-local-cohomology})
et la catégorie dérivée $D_{comp}(A, I)$ des complexes complets au sens dérivé
(Compléments d'algèbre, section \ref{more-algebra-section-derived-completion}).
Dans cette section, nous montrons que ces catégories sont équivalentes.
On trouvera un énoncé plus général dans \cite{Dwyer-Greenlees}.

\begin{lemma}
\label{dualizing-lemma-complete-and-local}
\begin{slogan}
Les résultats de cette nature sont parfois appelés dualité de Greenlees-May.
\end{slogan}
Soient $A$ un anneau et $I$ un idéal de type fini.
Soit $R\Gamma_Z$ comme dans le lemme \ref{dualizing-lemma-local-cohomology-adjoint}.
Notons ${\ }^\wedge$ la complétion dérivée, comme dans
Compléments d'algèbre, lemme \ref{more-algebra-lemma-derived-completion}.
Pour un objet $K$ de $D(A)$, on a
$$
R\Gamma_Z(K^\wedge) = R\Gamma_Z(K)
\quad\text{et}\quad
(R\Gamma_Z(K))^\wedge = K^\wedge
$$
dans $D(A)$.
\end{lemma}

\begin{proof}
Choisissons $f_1, \ldots, f_r \in A$ engendrant $I$. Rappelons que
$$
K^\wedge = R\Hom_A\left((A \to \prod A_{f_{i_0}}
\to \prod A_{f_{i_0i_1}} \to \ldots \to A_{f_1 \ldots f_r}), K\right)
$$
d'après Compléments d'algèbre, lemme \ref{more-algebra-lemma-derived-completion}.
Le cône $C = \text{Cone}(K \to K^\wedge)$ est donc donné par
$$
R\Hom_A\left((\prod A_{f_{i_0}}
\to \prod A_{f_{i_0i_1}} \to \ldots \to A_{f_1 \ldots f_r}), K\right)
$$
qui peut être représenté par un complexe muni d'une filtration finie
dont les quotients successifs sont isomorphes aux
$$
R\Hom_A(A_{f_{i_0} \ldots f_{i_p}}, K), \quad p > 0
$$
Ces complexes s'annulent lorsqu'on leur applique $R\Gamma_Z$,
voir le lemme \ref{dualizing-lemma-local-cohomology-vanishes}. En appliquant $R\Gamma_Z$
au triangle distingué $K \to K^\wedge \to C \to K[1]$,
on voit que la première formule du lemme est correcte.

\medskip\noindent
Rappelons que
$$
R\Gamma_Z(K) =
K \otimes^\mathbf{L} (A \to \prod A_{f_{i_0}}
\to \prod A_{f_{i_0i_1}} \to \ldots \to A_{f_1 \ldots f_r})
$$
d'après le lemme \ref{dualizing-lemma-local-cohomology-adjoint}.
Le cône $C = \text{Cone}(R\Gamma_Z(K) \to K)$ peut donc être représenté par un complexe muni
d'une filtration finie dont les quotients successifs sont isomorphes aux
$$
K \otimes_A A_{f_{i_0} \ldots f_{i_p}}, \quad p > 0
$$
Ces complexes s'annulent lorsqu'on leur applique ${\ }^\wedge$,
voir Compléments d'algèbre, lemme \ref{more-algebra-lemma-derived-completion-vanishes}.
En appliquant la complétion dérivée au triangle distingué
$R\Gamma_Z(K) \to K \to C \to R\Gamma_Z(K)[1]$
on voit que la seconde formule du lemme est correcte.
\end{proof}

\noindent
Le résultat suivant est un cas particulier d'un phénomène très général
concernant les sous-catégories admissibles d'une catégorie triangulée.

\begin{proposition}
\label{dualizing-proposition-torsion-complete}
\begin{reference}
C'est un cas particulier de \cite[théorème 1.1]{Porta-Liran-Yekutieli}.
\end{reference}
Soient $A$ un anneau et $I \subset A$ un idéal de type fini.
Les foncteurs $R\Gamma_Z$ et ${\ }^\wedge$ définissent des équivalences
quasi-inverses de catégories
$$
D_{I^\infty\text{-torsion}}(A) \leftrightarrow D_{comp}(A, I)
$$
\end{proposition}

\begin{proof}
Cela résulte immédiatement du lemme \ref{dualizing-lemma-complete-and-local}.
\end{proof}

\noindent
Le complément suivant à la proposition précédente précise
la correspondance sur les morphismes.

\begin{lemma}
\label{dualizing-lemma-compare-RHom}
Gardons les notations du lemme \ref{dualizing-lemma-complete-and-local}.
Pour des objets $K, L$ de $D(A)$, il existe un isomorphisme canonique
$$
R\Hom_A(K^\wedge, L^\wedge) \longrightarrow R\Hom_A(R\Gamma_Z(K), R\Gamma_Z(L))
$$
dans $D(A)$.
\end{lemma}

\begin{proof}
Écrivons $I = (f_1, \ldots, f_r)$. Notons $C = (A \to \prod A_{f_i} \to \ldots \to A_{f_1 \ldots f_r})$ le complexe alterné de
{\v C}ech. La complétion dérivée est alors donnée par
$R\Hom_A(C, -)$ (Compléments d'algèbre, lemme \ref{more-algebra-lemma-derived-completion})
et la cohomologie locale par $C \otimes^\mathbf{L} -$ (lemme \ref{dualizing-lemma-local-cohomology-adjoint}).
En combinant l'isomorphisme
$$
R\Hom_A(K \otimes^\mathbf{L} C, L \otimes^\mathbf{L} C) =
R\Hom_A(K, R\Hom_A(C,  L \otimes^\mathbf{L} C))
$$
(Compléments d'algèbre, lemme \ref{more-algebra-lemma-internal-hom})
et le morphisme
$$
L \to R\Hom_A(C,  L \otimes^\mathbf{L} C)
$$
(Compléments d'algèbre, lemme \ref{more-algebra-lemma-internal-hom-diagonal}), on obtient un morphisme
$$
\gamma :
R\Hom_A(K, L)
\longrightarrow
R\Hom_A(K \otimes^\mathbf{L} C, L \otimes^\mathbf{L} C)
$$
D'autre part, le membre de droite est complet au sens dérivé,
car il est égal à
$$
R\Hom_A(C, R\Hom_A(K, L \otimes^\mathbf{L} C)).
$$
Ainsi $\gamma$ se factorise par le complété dérivé de $R\Hom_A(K, L)$
en vertu de la propriété universelle de la complétion dérivée.
Or, la complétion dérivée passe à l'intérieur du $R\Hom_A$
d'après Compléments d'algèbre, lemme \ref{more-algebra-lemma-completion-RHom},
ce qui donne le morphisme cherché.

\medskip\noindent
Pour montrer que le morphisme du lemme est un isomorphisme,
on peut supposer que $K$ et $L$ sont complets au sens dérivé,
c'est-à-dire que $K = K^\wedge$ et $L = L^\wedge$. Il s'agit alors du morphisme
$$
\gamma : R\Hom_A(K, L) \longrightarrow R\Hom_A(R\Gamma_Z(K), R\Gamma_Z(L))
$$
D'après la proposition \ref{dualizing-proposition-torsion-complete}, les groupes de cohomologie
des membres de gauche et de droite coïncident.
Autrement dit, il faut vérifier que le morphisme $\gamma$
envoie un morphisme $\alpha : K \to L$ dans $D(A)$ sur le morphisme
$R\Gamma_Z(\alpha) : R\Gamma_Z(K) \to R\Gamma_Z(L)$.
Nous omettons cette vérification (indication : noter que $R\Gamma_Z(\alpha)$
est simplement le morphisme $\alpha \otimes \text{id}_C :
K \otimes^\mathbf{L} C
\to
L \otimes^\mathbf{L} C$, ce qui est presque identique
à la construction du morphisme dans Compléments d'algèbre, lemme \ref{more-algebra-lemma-internal-hom-diagonal}).
\end{proof}

\begin{lemma}
\label{dualizing-lemma-completion-local}
Soient $I$ et $J$ des idéaux d'un anneau noethérien $A$.
Soit $M$ un $A$-module de type fini. Posons $Z =V(J)$.
Considérons le complété $I$-adique dérivé $R\Gamma_Z(M)^\wedge$
de la cohomologie locale. Alors :
\begin{enumerate}
\item on a $R\Gamma_Z(M)^\wedge = R\lim R\Gamma_Z(M/I^nM)$ ;
\item on a des suites exactes courtes
$$
0 \to R^1\lim H^{i - 1}_Z(M/I^nM) \to H^i(R\Gamma_Z(M)^\wedge) \to
\lim H^i_Z(M/I^nM) \to 0
$$
\end{enumerate}
En particulier, $R\Gamma_Z(M)^\wedge$ a une cohomologie nulle en degrés négatifs.
\end{lemma}

\begin{proof}
Supposons que $J = (g_1, \ldots, g_m)$.
Alors $R\Gamma_Z(M)$ est calculé par le complexe
$$
M \to \prod M_{g_{j_0}} \to \prod M_{g_{j_0}g_{j_1}} \to
\ldots \to M_{g_1g_2\ldots g_m}
$$
d'après le lemme \ref{dualizing-lemma-local-cohomology-adjoint}.
D'après Compléments d'algèbre, lemme \ref{more-algebra-lemma-when-derived-completion-is-completion},
le complété $I$-adique dérivé de ce complexe
est donné par le complexe
$$
\lim M/I^nM \to \prod \lim (M/I^nM)_{g_{j_0}} \to
\ldots \to \lim (M/I^nM)_{g_1g_2\ldots g_m}
$$
des complétés usuels. Puisque $R\Gamma_Z(M/I^nM)$ est calculé par le complexe
$ M/I^nM \to \prod (M/I^nM)_{g_{j_0}} \to
\ldots \to (M/I^nM)_{g_1g_2\ldots g_m}$ et que les morphismes de transition entre ces complexes
sont surjectifs, on en déduit (1) par
Compléments d'algèbre, lemme \ref{more-algebra-lemma-compute-Rlim-modules}.
L'assertion (2) découle alors de Compléments d'algèbre, lemme \ref{more-algebra-lemma-break-long-exact-sequence-modules}.
\end{proof}

\begin{lemma}
\label{dualizing-lemma-completion-local-H0}
Sous les notations et hypothèses du lemme \ref{dualizing-lemma-completion-local},
supposons $A$ complet pour la topologie $I$-adique. Alors
$$
H^0(R\Gamma_Z(M)^\wedge) = \colim H^0_{V(J')}(M)
$$
où la limite inductive filtrante porte sur les $J' \subset J$ tels que
$V(J') \cap V(I) = V(J) \cap V(I)$.
\end{lemma}

\begin{proof}
Puisque $M$ est un $A$-module de type fini, $M$
est complet pour la topologie $I$-adique.
La démonstration du lemme \ref{dualizing-lemma-completion-local} montre que
$$
H^0(R\Gamma_Z(M)^\wedge) =
\Ker(M^\wedge \to \prod M_{g_j}^\wedge) =
\Ker(M \to \prod M_{g_j}^\wedge)
$$
où le membre de droite fait intervenir la complétion $I$-adique usuelle.
Le noyau $K_j$ de $M_{g_j} \to M_{g_j}^\wedge$ est $\bigcap I^n M_{g_j}$.
D'après Algèbre, lemme \ref{algebra-lemma-intersection-powers-ideal-module},
pour tout $\mathfrak p \in V(IA_{g_j})$, il existe $f \in A_{g_j}$ tel que $f \not \in \mathfrak p$ et $(K_j)_f = 0$.

\medskip\noindent
Soit $s \in H^0(R\Gamma_Z(M)^\wedge)$.
D'après ce qui précède, on peut considérer $s$ comme un élément de $M$.
L'intersection du support $Z'$ de $s$ avec $D(g_j)$
est disjointe de $D(g_j) \cap V(I)$ d'après les arguments ci-dessus.
Ainsi $Z'$ est un fermé de $\Spec(A)$ tel que $Z' \cap V(I) \subset V(J)$.
On a donc $Z' \cup V(J) = V(J')$ pour un idéal $J' \subset J$ tel que $V(J') \cap V(I) \subset V(J)$,
et l'on a $s \in H^0_{V(J')}(M)$.
Réciproquement, tout $s \in H^0_{V(J')}(M)$ avec $J' \subset J$ et $V(J') \cap V(I) \subset V(J)$
a une image nulle dans $M_{g_j}^\wedge$ pour tout $j$.
Cela prouve le lemme.
\end{proof}









\section{Dualité triviale pour un homomorphisme d'anneaux}
\label{dualizing-section-trivial}

\noindent
Soit $A \to B$ un homomorphisme d'anneaux. Considérons le foncteur
$$
\Hom_A(B, -) : \text{Mod}_A \longrightarrow \text{Mod}_B,\quad
M \longmapsto \Hom_A(B, M)
$$
Ce foncteur est exact à gauche et possède un prolongement dérivé
$R\Hom(B, -) : D(A) \to D(B)$.

\begin{lemma}
\label{dualizing-lemma-right-adjoint}
Soit $A \to B$ un homomorphisme d'anneaux. Le foncteur $R\Hom(B, -)$
construit ci-dessus est adjoint à droite du foncteur de restriction
$D(B) \to D(A)$.
\end{lemma}

\begin{proof}
Cela résulte de ce que la restriction et $\Hom_A(B, -)$ sont des foncteurs adjoints
d'après Algèbre, lemme \ref{algebra-lemma-adjoint-hom-restrict}.
Voir Catégories dérivées, lemme \ref{derived-lemma-derived-adjoint-functors}.
\end{proof}

\begin{lemma}
\label{dualizing-lemma-composition-right-adjoints}
Soient $A \to B \to C$ des homomorphismes d'anneaux. Alors
$R\Hom(C, -) \circ R\Hom(B, -) : D(A) \to D(C)$
est le foncteur $R\Hom(C, -) : D(A) \to D(C)$.
\end{lemma}

\begin{proof}
Cela résulte de l'unicité des adjoints à droite et du lemme \ref{dualizing-lemma-right-adjoint}.
\end{proof}

\begin{lemma}
\label{dualizing-lemma-RHom-ext}
Soit $\varphi : A \to B$ un homomorphisme d'anneaux. Pour $K$ dans $D(A)$, on a
$$
\varphi_*R\Hom(B, K) = R\Hom_A(B, K)
$$
où $\varphi_* : D(B) \to D(A)$ est la restriction. En particulier,
$R^q\Hom(B, K) = \Ext_A^q(B, K)$.
\end{lemma}

\begin{proof}
Choisissons un complexe K-injectif $I^\bullet$ représentant $K$.
Alors $R\Hom(B, K)$ est représenté par le complexe $\Hom_A(B, I^\bullet)$ de $B$-modules.
Puisque ce complexe, considéré comme complexe de $A$-modules,
représente $R\Hom_A(B, K)$, le lemme est vrai.
\end{proof}

\noindent
Soit $A$ un anneau noethérien. Nous noterons
$$
D_{\textit{Coh}}(A) \subset D(A)
$$
la sous-catégorie pleine formée des objets $K$ de $D(A)$
dont les modules de cohomologie sont tous des $A$-modules de type fini.
Cela a un sens d'après Catégories dérivées, section \ref{derived-section-triangulated-sub},
car, puisque $A$ est noethérien, la sous-catégorie des
$A$-modules de type fini est une sous-catégorie de Serre de $\text{Mod}_A$.

\begin{lemma}
\label{dualizing-lemma-exact-support-coherent}
Avec les notations ci-dessus, supposons que $A \to B$ soit un homomorphisme
fini d'anneaux noethériens. Alors $R\Hom(B, -)$ envoie $D^+_{\textit{Coh}}(A)$ dans $D^+_{\textit{Coh}}(B)$.
\end{lemma}

\begin{proof}
Il faut montrer que, si $K \in D^+(A)$ a des modules de cohomologie de type fini,
le complexe $R\Hom(B, K)$ a lui aussi des modules de cohomologie de type fini.
Cela résulte par exemple du lemme \ref{dualizing-lemma-RHom-ext}, si l'on montre que
les modules Ext $\Ext^i_A(B, K)$ sont des $A$-modules de type fini.
Puisque $K$ est borné inférieurement, on a une suite spectrale convergente
$$
\Ext^p_A(B, H^q(K)) \Rightarrow \text{Ext}^{p + q}_A(B, K)
$$
Cela achève la démonstration, car les modules $\Ext^p_A(B, H^q(K))$
sont de type fini d'après Algèbre, lemme \ref{algebra-lemma-ext-noetherian}.
\end{proof}

\begin{remark}
\label{dualizing-remark-exact-support}
Soient $A$ un anneau et $I \subset A$ un idéal. Posons $B = A/I$.
Dans ce cas, le foncteur $\Hom_A(B, -)$ est égal au foncteur
$$
\text{Mod}_A \longrightarrow \text{Mod}_B,\quad M \longmapsto M[I]
$$
qui associe à $M$ son sous-module de $I$-torsion.
\end{remark}

\begin{situation}
\label{dualizing-situation-resolution}
Soit $R \to A$ un homomorphisme d'anneaux.
Donnons une autre construction de $R\Hom(A, -)$, qui nous sera utile
dans la suite de ce chapitre.
Supposons donnée une algèbre différentielle graduée $(E, d)$
sur $R$ et un quasi-isomorphisme $E \to A$, où l'on considère $A$
comme une algèbre différentielle graduée sur $R$ à différentielle nulle.
On a alors des diagrammes commutatifs
$$
\vcenter{
\xymatrix{
D(E, \text{d}) \ar[rd] & &  D(A) \ar[ll] \ar[ld] \\
& D(R)
}
}
\quad\text{et}\quad
\vcenter{
\xymatrix{
D(E, \text{d}) \ar[rr]_{- \otimes_E^\mathbf{L} A} & &  D(A) \\
& D(R) \ar[lu]^{- \otimes_R^\mathbf{L} E} \ar[ru]_{- \otimes_R^\mathbf{L} A}
}
}
$$
où les flèches horizontales sont des équivalences de catégories
(Algèbre différentielle graduée, lemme \ref{dga-lemma-qis-equivalence}).
La commutativité du premier diagramme est claire.
Le second est commutatif parce que le premier l'est et que nos foncteurs
en sont les adjoints à gauche
(Algèbre différentielle graduée, exemple \ref{dga-example-map-hom-tensor}),
ou encore parce que l'on a $E \otimes^\mathbf{L}_E A = E \otimes_E A$ et que l'on peut appliquer
Algèbre différentielle graduée, lemme \ref{dga-lemma-compose-tensor-functors-general}.
\end{situation}

\begin{lemma}
\label{dualizing-lemma-RHom-dga}
Dans la situation \ref{dualizing-situation-resolution}, le foncteur $R\Hom(A, -)$ est égal au composé de
$R\Hom(E, -) : D(R) \to D(E, \text{d})$
et de l'équivalence $- \otimes^\mathbf{L}_E A : D(E, \text{d}) \to D(A)$.
\end{lemma}

\begin{proof}
Cela résulte de ce que $R\Hom(E, -)$ est adjoint à droite
de $- \otimes^\mathbf{L}_R E$ ; voir Algèbre différentielle graduée, lemme \ref{dga-lemma-tensor-hom-adjoint}.
Ce foncteur joue donc le même rôle que le foncteur $R\Hom(A, -)$
pour l'homomorphisme $R \to A$ (lemme \ref{dualizing-lemma-right-adjoint}).
Ces foncteurs doivent donc se correspondre par l'équivalence
$- \otimes^\mathbf{L}_E A : D(E, \text{d}) \to D(A)$.
\end{proof}

\begin{lemma}
\label{dualizing-lemma-RHom-is-tensor}
Dans la situation \ref{dualizing-situation-resolution}, supposons que
\begin{enumerate}
\item $E$, considéré comme objet de $D(R)$, soit compact ;
\item $N = \Hom^\bullet_R(E^\bullet, R)$ calcule $R\Hom(E, R)$.
\end{enumerate}
Alors $R\Hom(E, -) : D(R) \to D(E)$ est isomorphe à
$K \mapsto K \otimes_R^\mathbf{L} N$.
\end{lemma}

\begin{proof}
C'est un cas particulier d'Algèbre différentielle graduée, lemme \ref{dga-lemma-RHom-is-tensor}.
\end{proof}

\begin{lemma}
\label{dualizing-lemma-RHom-is-tensor-special}
Dans la situation \ref{dualizing-situation-resolution}, supposons que $A$ soit un $R$-module parfait.
Alors
$$
R\Hom(A, -) : D(R) \to D(A)
$$
est donné par $K \mapsto K \otimes_R^\mathbf{L} M$,
où $M = R\Hom(A, R) \in D(A)$.
\end{lemma}

\begin{proof}
Appliquons Algèbre à puissances divisées, lemme \ref{dpa-lemma-tate-resoluton-pseudo-coherent-ring-map},
pour choisir une résolution de Tate $(E, \text{d})$ de $A$ sur $R$.
Notons que $E^i = 0$ pour $i > 0$, que $E^0 = R[x_1, \ldots, x_n]$ est une algèbre de polynômes
et que $E^i$ est un $E^0$-module libre de type fini pour $i < 0$.
Il en résulte que $E$, considéré comme complexe de $R$-modules,
est un complexe borné supérieurement de $R$-modules libres.
Vérifions les hypothèses du lemme \ref{dualizing-lemma-RHom-is-tensor}.
La première est satisfaite puisque $A$ est parfait
(donc compact d'après Compléments d'algèbre, proposition \ref{more-algebra-proposition-perfect-is-compact}),
et la seconde d'après Compléments d'algèbre, lemme \ref{more-algebra-lemma-RHom-out-of-projective}.
Le lemme montre que $K \mapsto R\Hom(E, K)$ est isomorphe à $K \mapsto K \otimes_R^\mathbf{L} N$
pour un certain $E$-module différentiel gradué $N$.
Observons que
$$
(R \otimes_R E) \otimes_E^\mathbf{L} A = R \otimes_E E \otimes_E A
$$
dans $D(A)$. D'après Algèbre différentielle graduée, lemme \ref{dga-lemma-compose-tensor-functors-general-algebra},
on en déduit que le composé de $- \otimes_R^\mathbf{L} N$ et de $- \otimes_R^\mathbf{L} A$
est de la forme $- \otimes_R M$ pour un certain $M \in D(A)$.
Pour achever la démonstration, appliquons le lemme \ref{dualizing-lemma-RHom-dga}.
\end{proof}

\begin{lemma}
\label{dualizing-lemma-compute-for-effective-Cartier-algebraic}
Soit $R \to A$ un homomorphisme surjectif d'anneaux dont le noyau $I$
est un $R$-module inversible. Le foncteur
$R\Hom(A, -) : D(R) \to D(A)$
est isomorphe à $K \mapsto K \otimes_R^\mathbf{L} N[-1]$,
où $N$ est l'inverse du $A$-module inversible $I \otimes_R A$.
\end{lemma}

\begin{proof}
Puisque $A$ possède la résolution projective finie
$$
0 \to I \to R \to A \to 0
$$
on voit que $A$ est un $R$-module parfait.
D'après le lemme \ref{dualizing-lemma-RHom-is-tensor-special}, il suffit de prouver que $R\Hom(A, R)$
est représenté par $N[-1]$ dans $D(A)$.
Cela signifie que $R\Hom(A, R)$ possède un unique module de cohomologie non nul,
à savoir $N$ en degré $1$.
Puisque $\text{Mod}_A \to \text{Mod}_R$ est pleinement fidèle, il suffit de le prouver
après application du foncteur de restriction $i_* : D(A) \to D(R)$.
D'après le lemme \ref{dualizing-lemma-RHom-ext}, on a
$$
i_*R\Hom(A, R) = R\Hom_R(A, R)
$$
À l'aide de la résolution projective finie ci-dessus, on trouve que
ce dernier est représenté par le complexe $R \to I^{\otimes -1}$ avec $R$
en degré $0$. Le morphisme $R \to I^{\otimes -1}$ est injectif
et son conoyau est $N$.
\end{proof}





\section{Changement de base pour la dualité triviale}
\label{dualizing-section-base-change-trivial-duality}

\noindent
Dans cette section, considérons un carré cocartésien d'anneaux
$$
\xymatrix{
A \ar[r]_\alpha & A' \\
R \ar[u]^\varphi \ar[r]^\rho & R' \ar[u]_{\varphi'}
}
$$
Autrement dit, on a $A' = A \otimes_R R'$. Si $A$ et $R'$
sont {\bf Tor-indépendants sur} $R$,
il existe un morphisme canonique de changement de base
\begin{equation}
\label{dualizing-equation-base-change}
R\Hom(A, K) \otimes_A^\mathbf{L} A'
\longrightarrow
R\Hom(A', K \otimes_R^\mathbf{L} R')
\end{equation}
dans $D(A')$, fonctoriel en $K$ dans $D(R)$.
En effet, par l'adjonction du lemme \ref{dualizing-lemma-right-adjoint},
une telle flèche revient à se donner un morphisme
$$
\varphi'_*\left(R\Hom(A, K) \otimes_A^\mathbf{L} A'\right)
\longrightarrow
K \otimes_R^\mathbf{L} R'
$$
dans $D(R')$, où $\varphi'_* : D(A') \to D(R')$ est le foncteur de restriction.
On peut appliquer Compléments d'algèbre, lemme \ref{more-algebra-lemma-base-change-comparison},
au membre de gauche pour voir que cela revient à se donner un morphisme
$$
\varphi_*(R\Hom(A, K)) \otimes_R^\mathbf{L} R'
\longrightarrow
K \otimes_R^\mathbf{L} R'
$$
dans $D(R')$, où $\varphi_* : D(A) \to D(R)$ est le foncteur de restriction.
On peut choisir pour cela $can \otimes^\mathbf{L} \text{id}_{R'}$,
où $can : \varphi_*(R\Hom(A, K)) \to K$ est la counité de l'adjonction entre $R\Hom(A, -)$ et $\varphi_*$.

\begin{lemma}
\label{dualizing-lemma-check-base-change-is-iso}
Dans la situation ci-dessus, le morphisme (\ref{dualizing-equation-base-change})
est un isomorphisme si et seulement si le morphisme
$$
R\Hom_R(A, K) \otimes_R^\mathbf{L} R'
\longrightarrow
R\Hom_R(A, K \otimes_R^\mathbf{L} R')
$$
de Compléments d'algèbre, lemme \ref{more-algebra-lemma-internal-hom-diagonal-better}, est un isomorphisme.
\end{lemma}

\begin{proof}
Pour voir que le morphisme est un isomorphisme, il suffit de le vérifier
après application de $\varphi'_*$.
En appliquant le foncteur $\varphi'_*$ à (\ref{dualizing-equation-base-change})
et en utilisant l'égalité $A' = A \otimes_R^\mathbf{L} R'$, on obtient le morphisme de changement de base
$R\Hom_R(A, K) \otimes_R^\mathbf{L} R' \to
R\Hom_{R'}(A \otimes_R^\mathbf{L} R', K \otimes_R^\mathbf{L} R')$
pour le Hom dérivé de Compléments d'algèbre, équation (\ref{more-algebra-equation-base-change-RHom}).
En explicitant les membres de gauche et de droite comme dans la démonstration
de Compléments d'algèbre, lemme \ref{more-algebra-lemma-base-change-RHom}, et en utilisant notamment
Compléments d'algèbre, lemme \ref{more-algebra-lemma-upgrade-adjoint-tensor-RHom}, on obtient le résultat voulu.
\end{proof}

\begin{lemma}
\label{dualizing-lemma-flat-bc-surjection}
Soient $R \to A$ et $R \to R'$ des homomorphismes d'anneaux et $A' = A \otimes_R R'$.
Supposons que
\begin{enumerate}
\item $A$ soit pseudo-cohérent en tant que $R$-module ;
\item $R'$ soit de Tor-dimension finie en tant que $R$-module
(par exemple, $R \to R'$ est plat) ;
\item $A$ et $R'$ soient Tor-indépendants sur $R$.
\end{enumerate}
Alors (\ref{dualizing-equation-base-change}) est un isomorphisme pour $K \in D^+(R)$.
\end{lemma}

\begin{proof}
Cela résulte du lemme \ref{dualizing-lemma-check-base-change-is-iso} et de
Compléments d'algèbre, lemme \ref{more-algebra-lemma-internal-hom-evaluate-tensor-isomorphism} (4).
\end{proof}

\begin{lemma}
\label{dualizing-lemma-bc-surjection}
Soient $R \to A$ et $R \to R'$ des homomorphismes d'anneaux et $A' = A \otimes_R R'$.
Supposons que
\begin{enumerate}
\item $A$ soit parfait en tant que $R$-module ;
\item $A$ et $R'$ soient Tor-indépendants sur $R$.
\end{enumerate}
Alors (\ref{dualizing-equation-base-change}) est un isomorphisme pour tout $K \in D(R)$.
\end{lemma}

\begin{proof}
Cela résulte du lemme \ref{dualizing-lemma-check-base-change-is-iso} et de
Compléments d'algèbre, lemme \ref{more-algebra-lemma-internal-hom-evaluate-tensor-isomorphism} (1).
\end{proof}







\section{Complexes dualisants}
\label{dualizing-section-dualizing}

\noindent
Dans cette section, nous définissons les complexes dualisants pour les anneaux noethériens.

\begin{definition}
\label{dualizing-definition-dualizing}
Soit $A$ un anneau noethérien. Un {\it complexe dualisant}
est un complexe de $A$-modules $\omega_A^\bullet$ tel que
\begin{enumerate}
\item $\omega_A^\bullet$ soit de dimension injective finie ;
\item $H^i(\omega_A^\bullet)$ soit un $A$-module de type fini pour tout $i$ ;
\item $A \to R\Hom_A(\omega_A^\bullet, \omega_A^\bullet)$
soit un quasi-isomorphisme.
\end{enumerate}
\end{definition}

\noindent
Il faut un peu de temps pour se familiariser avec cette définition.
Il peut être utile de démontrer soi-même certains des lemmes suivants
avant d'en lire les démonstrations.

\begin{lemma}
\label{dualizing-lemma-finite-ext-into-bounded-injective}
Soit $A$ un anneau noethérien. Soient $K, L \in D_{\textit{Coh}}(A)$,
et supposons $L$ de dimension injective finie.
Alors $R\Hom_A(K, L)$ appartient à $D_{\textit{Coh}}(A)$.
\end{lemma}

\begin{proof}
Choisissons un entier $n$ et considérons le triangle distingué
$$
\tau_{\leq n}K \to K \to \tau_{\geq n + 1}K \to \tau_{\leq n}K[1]
$$
voir Catégories dérivées, remarque \ref{derived-remark-truncation-distinguished-triangle}.
Puisque $L$ est de dimension injective finie, $R\Hom_A(\tau_{\geq n + 1}K, L)$
a une cohomologie nulle en degrés $\geq c - n$ pour une constante $c$.
Ainsi, étant donné $i$, on voit que $\Ext^i_A(K, L) \to \Ext^i_A(\tau_{\leq n}K, L)$
est un isomorphisme pour un certain $n \gg - i$.
En appliquant Catégories dérivées de schémas, lemme \ref{perfect-lemma-coherent-internal-hom},
à $\tau_{\leq n}K$ et $L$, on en déduit que $\Ext^i_A(K, L)$
est un $A$-module de type fini pour tout $i$.
Ainsi $R\Hom_A(K, L)$ est bien un objet de $D_{\textit{Coh}}(A)$.
\end{proof}

\begin{lemma}
\label{dualizing-lemma-dualizing}
Soit $A$ un anneau noethérien. Si $\omega_A^\bullet$ est un complexe dualisant,
le foncteur
$$
D : K \longmapsto R\Hom_A(K, \omega_A^\bullet)
$$
est une anti-équivalence $D_{\textit{Coh}}(A) \to D_{\textit{Coh}}(A)$
qui échange $D^+_{\textit{Coh}}(A)$ et $D^-_{\textit{Coh}}(A)$
et induit une anti-équivalence
$D^b_{\textit{Coh}}(A) \to D^b_{\textit{Coh}}(A)$.
De plus, $D \circ D$ est isomorphe au foncteur identité.
\end{lemma}

\begin{proof}
Soit $K$ un objet de $D_{\textit{Coh}}(A)$. D'après le lemme \ref{dualizing-lemma-finite-ext-into-bounded-injective},
on voit que $R\Hom_A(K, \omega_A^\bullet)$ est un objet de $D_{\textit{Coh}}(A)$.
Compléments d'algèbre, lemme \ref{more-algebra-lemma-internal-hom-evaluate-isomorphism-technical},
et les hypothèses sur le complexe dualisant
donnent un isomorphisme canonique
$$
K = R\Hom_A(\omega_A^\bullet, \omega_A^\bullet) \otimes_A^\mathbf{L} K
\longrightarrow
R\Hom_A(R\Hom_A(K, \omega_A^\bullet), \omega_A^\bullet)
$$
Notre foncteur possède donc un quasi-inverse, ce qui achève la démonstration.
\end{proof}

\noindent
Soit $R$ un anneau. Rappelons qu'un objet $L$ de $D(R)$
est {\it inversible} s'il est inversible pour la structure
monoïdale symétrique sur $D(R)$ donnée par le produit tensoriel dérivé.
Dans Compléments d'algèbre, lemme \ref{more-algebra-lemma-invertible-derived}, nous avons vu que cela signifie
que $L$ est parfait, que $L = \bigoplus H^n(L)[-n]$, cette somme étant finie,
que chaque $H^n(L)$ est projectif de type fini
et qu'il existe un recouvrement ouvert $\Spec(R) = \bigcup D(f_i)$ tel que
$L \otimes_R R_{f_i} \cong R_{f_i}[-n_i]$ pour certains entiers $n_i$.

\begin{lemma}
\label{dualizing-lemma-equivalence-comes-from-invertible}
Soit $A$ un anneau noethérien. Soit $F : D^b_{\textit{Coh}}(A) \to D^b_{\textit{Coh}}(A)$
une équivalence de catégories $A$-linéaire.
Alors $F(A)$ est un objet inversible de $D(A)$.
\end{lemma}

\begin{proof}
Soit $\mathfrak m \subset A$ un idéal maximal de corps résiduel $\kappa$.
Considérons l'objet $F(\kappa)$. Puisque $\kappa = \Hom_{D(A)}(\kappa, \kappa)$,
tous les groupes de cohomologie de $F(\kappa)$ sont annulés par $\mathfrak m$.
On voit aussi que
$$
\Ext^i_A(\kappa, \kappa) = \text{Ext}^i_A(F(\kappa), F(\kappa))
= \Hom_{D(A)}(F(\kappa), F(\kappa)[i])
$$
est nul pour $i < 0$. Supposons $H^a(F(\kappa)) \not = 0$ et $H^b(F(\kappa)) \not = 0$
avec $a$ minimal et $b$ maximal
(en particulier $a \leq b$). Il existe alors un morphisme non nul
$$
F(\kappa) \to H^b(F(\kappa))[-b] \to H^a(F(\kappa))[-b]
\to F(\kappa)[a - b]
$$
dans $D(A)$ (non nul parce qu'il induit un morphisme non nul en cohomologie).
Cela prouve que $b = a$. On en déduit que $F(\kappa) = \kappa[-a]$.

\medskip\noindent
Soit $G$ un quasi-inverse du foncteur $F$.
En raisonnant comme ci-dessus, on trouve un entier $b$ tel que $G(\kappa) = \kappa[-b]$.
En composant, on obtient $a + b = 0$. Soit $E$ un $A$-module de type fini
annulé par une puissance de $\mathfrak m$. Par récurrence sur la longueur de $E$,
on trouve que $G(E) = E'[-b]$ pour un certain $A$-module de type fini $E'$
annulé par une puissance de $\mathfrak m$. Alors $E[-a] = F(E')$.
Considérons ensuite les groupes
$$
\Ext^i_A(A, E') = \text{Ext}^i_A(F(A), F(E')) =
\Hom_{D(A)}(F(A), E[-a + i])
$$
Le membre de gauche est non nul si et seulement si $i = 0$,
auquel cas on obtient $E'$. En appliquant cela à $E = E' = \kappa$
et en utilisant le lemme de Nakayama, on voit que $H^j(F(A))_\mathfrak m$
est nul pour $j > a$ et engendré par $1$ élément pour $j = a$.
D'autre part, si $H^j(F(A))_\mathfrak m$ n'est pas nul pour un certain $j < a$,
il existe un morphisme $F(A) \to E[-a + i]$ pour un certain $i < 0$
et un certain $E$ (Compléments d'algèbre, lemme \ref{more-algebra-lemma-detect-cohomology}),
ce qui est contradictoire.
Ainsi $F(A)_\mathfrak m = M[-a]$ pour un certain $A_\mathfrak m$-module $M$
engendré par $1$ élément. Cependant, puisque
$$
A_\mathfrak m = \Hom_{D(A)}(A, A)_\mathfrak m =
\Hom_{D(A)}(F(A), F(A))_\mathfrak m = \Hom_{A_\mathfrak m}(M, M)
$$
on voit que $M \cong A_\mathfrak m$. On en déduit qu'il existe un élément $f \in A$
tel que $f \not \in \mathfrak m$ et tel que $F(A)_f$ soit isomorphe à $A_f[-a]$.
Cela achève la démonstration.
\end{proof}

\begin{lemma}
\label{dualizing-lemma-dualizing-unique}
Soit $A$ un anneau noethérien. Si $\omega_A^\bullet$ et $(\omega'_A)^\bullet$
sont des complexes dualisants, alors $(\omega'_A)^\bullet$ est quasi-isomorphe à
$\omega_A^\bullet \otimes_A^\mathbf{L} L$
pour un certain objet inversible $L$ de $D(A)$.
\end{lemma}

\begin{proof}
D'après les lemmes \ref{dualizing-lemma-dualizing} et \ref{dualizing-lemma-equivalence-comes-from-invertible}, le foncteur $K \mapsto R\Hom_A(R\Hom_A(K, \omega_A^\bullet), (\omega_A')^\bullet)$
envoie $A$ sur un objet inversible $L$.
Autrement dit, il existe un isomorphisme
$$
L \longrightarrow R\Hom_A(\omega_A^\bullet, (\omega_A')^\bullet)
$$
Puisque $L$ est de Tor-dimension finie, on peut appliquer
Compléments d'algèbre, lemme \ref{more-algebra-lemma-internal-hom-evaluate-isomorphism-technical}, pour voir que
$$
R\Hom_A(\omega_A^\bullet, (\omega'_A)^\bullet) \otimes_A^\mathbf{L} K
\longrightarrow
R\Hom_A(R\Hom_A(K, \omega_A^\bullet), (\omega_A')^\bullet)
$$
est un isomorphisme pour $K$ dans $D^b_{\textit{Coh}}(A)$.
En particulier, en posant $K = \omega_A^\bullet$, on achève la démonstration.
\end{proof}

\begin{lemma}
\label{dualizing-lemma-dualizing-localize}
Soit $A$ un anneau noethérien. Soit $B = S^{-1}A$ un localisé.
Si $\omega_A^\bullet$ est un complexe dualisant, alors $\omega_A^\bullet \otimes_A B$
est un complexe dualisant pour $B$.
\end{lemma}

\begin{proof}
Soit $\omega_A^\bullet \to I^\bullet$ un quasi-isomorphisme, où $I^\bullet$ est un complexe borné d'injectifs.
Alors $S^{-1}I^\bullet$ est un complexe borné de $B = S^{-1}A$-modules injectifs
(lemme \ref{dualizing-lemma-localization-injective-modules}) représentant $\omega_A^\bullet \otimes_A B$.
Ainsi $\omega_A^\bullet \otimes_A B$ est de dimension injective finie.
Puisque $H^i(\omega_A^\bullet \otimes_A B) = H^i(\omega_A^\bullet) \otimes_A B$ par platitude de $A \to B$,
on voit que $\omega_A^\bullet \otimes_A B$ a des modules de cohomologie de type fini.
Enfin, le morphisme
$$
B \longrightarrow
R\Hom_A(\omega_A^\bullet \otimes_A B, \omega_A^\bullet \otimes_A B)
$$
est un quasi-isomorphisme, car la formation du Hom interne commute
au changement de base plat dans ce cas ; voir
Compléments d'algèbre, lemme \ref{more-algebra-lemma-base-change-RHom}.
\end{proof}

\begin{lemma}
\label{dualizing-lemma-dualizing-glue}
Soit $A$ un anneau noethérien. Soient $f_1, \ldots, f_n \in A$ engendrant l'idéal unité.
Si $\omega_A^\bullet$ est un complexe de $A$-modules tel que $(\omega_A^\bullet)_{f_i}$
soit un complexe dualisant pour $A_{f_i}$ pour tout $i$,
alors $\omega_A^\bullet$ est un complexe dualisant pour $A$.
\end{lemma}

\begin{proof}
Considérons le complexe double
$$
\prod\nolimits_{i_0} (\omega_A^\bullet)_{f_{i_0}}
\to
\prod\nolimits_{i_0 < i_1} (\omega_A^\bullet)_{f_{i_0}f_{i_1}}
\to \ldots
$$
Le complexe total associé est quasi-isomorphe à $\omega_A^\bullet$,
par exemple d'après Descente, remarque \ref{descent-remark-standard-covering},
ou Catégories dérivées de schémas, lemme \ref{perfect-lemma-alternating-cech-complex-complex-computes-cohomology}.
Par hypothèse, les complexes $(\omega_A^\bullet)_{f_i}$ sont de dimension injective finie
en tant que complexes de $A_{f_i}$-modules.
Cela entraîne que chacun des complexes $(\omega_A^\bullet)_{f_{i_0} \ldots f_{i_p}}$, $p > 0$,
est de dimension injective finie sur $A_{f_{i_0} \ldots f_{i_p}}$ ; voir le lemme \ref{dualizing-lemma-localization-injective-modules}.
Cela entraîne à son tour que chacun des complexes $(\omega_A^\bullet)_{f_{i_0} \ldots f_{i_p}}$, $p > 0$,
est de dimension injective finie sur $A$ ; voir le lemme \ref{dualizing-lemma-injective-flat}.
Ainsi $\omega_A^\bullet$ est de dimension injective finie comme complexe de $A$-modules
(car il peut être représenté par un complexe muni d'une filtration finie
dont les gradués sont de dimension injective finie).
Puisque $H^n(\omega_A^\bullet)_{f_i}$ est un $A_{f_i}$-module de type fini pour chaque $i$,
on voit que $H^i(\omega_A^\bullet)$ est un $A$-module de type fini ; voir Algèbre, lemme \ref{algebra-lemma-cover}.
Enfin, le changement de base dérivé du morphisme $A \to R\Hom_A(\omega_A^\bullet, \omega_A^\bullet)$ à $A_{f_i}$
est le morphisme $A_{f_i} \to R\Hom_A((\omega_A^\bullet)_{f_i}, (\omega_A^\bullet)_{f_i})$ d'après Compléments d'algèbre, lemme \ref{more-algebra-lemma-base-change-RHom}.
On en déduit donc que
$A \to R\Hom_A(\omega_A^\bullet, \omega_A^\bullet)$
est un isomorphisme, ce qui achève la démonstration.
\end{proof}

\begin{lemma}
\label{dualizing-lemma-dualizing-finite}
Soit $A \to B$ un homomorphisme fini d'anneaux noethériens.
Soit $\omega_A^\bullet$ un complexe dualisant.
Alors $R\Hom(B, \omega_A^\bullet)$ est un complexe dualisant pour $B$.
\end{lemma}

\begin{proof}
Soit $\omega_A^\bullet \to I^\bullet$ un quasi-isomorphisme, où $I^\bullet$ est un complexe borné d'injectifs.
Alors $\Hom_A(B, I^\bullet)$ est un complexe borné de $B$-modules injectifs
(lemme \ref{dualizing-lemma-hom-injective}) représentant $R\Hom(B, \omega_A^\bullet)$.
Ainsi $R\Hom(B, \omega_A^\bullet)$ est de dimension injective finie.
D'après le lemme \ref{dualizing-lemma-exact-support-coherent}, c'est un objet de $D_{\textit{Coh}}(B)$.
Enfin, on calcule
$$
\Hom_{D(B)}(R\Hom(B, \omega_A^\bullet), R\Hom(B, \omega_A^\bullet)) =
\Hom_{D(A)}(R\Hom(B, \omega_A^\bullet), \omega_A^\bullet) = B
$$
et, pour $n \not = 0$, on calcule
$$
\Hom_{D(B)}(R\Hom(B, \omega_A^\bullet), R\Hom(B, \omega_A^\bullet)[n]) =
\Hom_{D(A)}(R\Hom(B, \omega_A^\bullet), \omega_A^\bullet[n]) = 0
$$
ce qui prouve la dernière propriété d'un complexe dualisant.
Dans les égalités affichées, la première égalité résulte du lemme \ref{dualizing-lemma-right-adjoint}
et la seconde du lemme \ref{dualizing-lemma-dualizing}.
\end{proof}

\begin{lemma}
\label{dualizing-lemma-dualizing-quotient}
Soit $A \to B$ un homomorphisme surjectif d'anneaux noethériens.
Soit $\omega_A^\bullet$ un complexe dualisant.
Alors $R\Hom(B, \omega_A^\bullet)$ est un complexe dualisant pour $B$.
\end{lemma}

\begin{proof}
C'est un cas particulier du lemme \ref{dualizing-lemma-dualizing-finite}.
\end{proof}

\begin{lemma}
\label{dualizing-lemma-dualizing-polynomial-ring}
Soit $A$ un anneau noethérien. Si $\omega_A^\bullet$ est un complexe dualisant,
alors $\omega_A^\bullet \otimes_A A[x]$ est un complexe dualisant pour $A[x]$.
\end{lemma}

\begin{proof}
Posons $B = A[x]$ et $\omega_B^\bullet = \omega_A^\bullet \otimes_A B$.
Le lemme \ref{dualizing-lemma-injective-dimension-over-polynomial-ring} et Compléments d'algèbre, lemme \ref{more-algebra-lemma-finite-injective-dimension},
montrent que $\omega_B^\bullet$ est de dimension injective finie.
Puisque $H^i(\omega_B^\bullet) = H^i(\omega_A^\bullet) \otimes_A B$ par platitude de $A \to B$, on voit que $\omega_A^\bullet \otimes_A B$
a des modules de cohomologie de type fini.
Enfin, le morphisme
$$
B \longrightarrow R\Hom_B(\omega_B^\bullet, \omega_B^\bullet)
$$
est un quasi-isomorphisme, car la formation du Hom interne commute
au changement de base plat dans ce cas ; voir
Compléments d'algèbre, lemme \ref{more-algebra-lemma-base-change-RHom}.
\end{proof}

\begin{proposition}
\label{dualizing-proposition-dualizing-essentially-finite-type}
Soit $A$ un anneau noethérien possédant un complexe dualisant.
Alors toute $A$-algèbre essentiellement de type fini sur $A$
possède un complexe dualisant.
\end{proposition}

\begin{proof}
Cela résulte de la combinaison des lemmes \ref{dualizing-lemma-dualizing-localize},
\ref{dualizing-lemma-dualizing-quotient} et \ref{dualizing-lemma-dualizing-polynomial-ring}.
\end{proof}

\begin{lemma}
\label{dualizing-lemma-find-function}
Soit $A$ un anneau noethérien. Soit $\omega_A^\bullet$ un complexe dualisant.
Soit $\mathfrak m \subset A$ un idéal maximal et posons $\kappa = A/\mathfrak m$. Alors
$R\Hom_A(\kappa, \omega_A^\bullet) \cong \kappa[n]$ pour un certain $n \in \mathbf{Z}$.
\end{lemma}

\begin{proof}
Cela résulte de ce que $R\Hom_A(\kappa, \omega_A^\bullet)$ est un complexe dualisant sur $\kappa$
(lemme \ref{dualizing-lemma-dualizing-quotient}), de ce que les complexes dualisants sur $\kappa$
sont uniques à décalage près (lemme \ref{dualizing-lemma-dualizing-unique})
et de ce que $\kappa$ est un complexe dualisant sur $\kappa$.
\end{proof}




\section{Complexes dualisants sur les anneaux locaux}
\label{dualizing-section-dualizing-local}

\noindent
Dans cette section, $(A, \mathfrak m, \kappa)$ sera un anneau local noethérien
muni d'un complexe dualisant $\omega_A^\bullet$ tel que l'entier $n$
du lemme \ref{dualizing-lemma-find-function} soit nul. Plus précisément, supposons $R\Hom_A(\kappa, \omega_A^\bullet) = \kappa[0]$.
Nous dirons alors que le complexe dualisant est {\it normalisé}.
Observons qu'un complexe dualisant normalisé est unique à isomorphisme près
et que tout autre complexe dualisant pour $A$ est isomorphe
à un décalé d'un complexe normalisé (lemme \ref{dualizing-lemma-dualizing-unique}).

\begin{lemma}
\label{dualizing-lemma-normalized-finite}
Soit $(A, \mathfrak m, \kappa) \to (B, \mathfrak m', \kappa')$ un homomorphisme local fini d'anneaux locaux noethériens.
Soit $\omega_A^\bullet$ un complexe dualisant normalisé.
Alors $\omega_B^\bullet = R\Hom(B, \omega_A^\bullet)$ est un complexe dualisant normalisé pour $B$.
\end{lemma}

\begin{proof}
D'après le lemme \ref{dualizing-lemma-dualizing-finite}, le complexe $\omega_B^\bullet$ est dualisant pour $B$.
On a
$$
R\Hom_B(\kappa', \omega_B^\bullet) =
R\Hom_B(\kappa', R\Hom(B, \omega_A^\bullet)) =
R\Hom_A(\kappa', \omega_A^\bullet)
$$
d'après le lemme \ref{dualizing-lemma-right-adjoint}. Puisque $\kappa'$ est isomorphe à une somme directe
finie de copies de $\kappa$ comme $A$-module, et que $\omega_A^\bullet$ est normalisé,
ce complexe n'a de cohomologie qu'en degré $0$.
Ainsi $\omega_B^\bullet$ est lui aussi un complexe dualisant normalisé.
\end{proof}

\begin{lemma}
\label{dualizing-lemma-normalized-quotient}
Soit $(A, \mathfrak m, \kappa)$ un anneau local noethérien muni d'un complexe dualisant
normalisé $\omega_A^\bullet$. Soit $A \to B$ surjectif.
Alors $\omega_B^\bullet = R\Hom_A(B, \omega_A^\bullet)$ est un complexe dualisant normalisé pour $B$.
\end{lemma}

\begin{proof}
C'est un cas particulier du lemme \ref{dualizing-lemma-normalized-finite}.
\end{proof}

\begin{lemma}
\label{dualizing-lemma-equivalence-finite-length}
Soit $(A, \mathfrak m, \kappa)$ un anneau local noethérien. Soit $F$
une auto-équivalence $A$-linéaire de la catégorie des $A$-modules
de longueur finie. Alors $F$ est isomorphe au foncteur identité.
\end{lemma}

\begin{proof}
Puisque $\kappa$ est l'unique objet simple de la catégorie, on a $F(\kappa) \cong \kappa$.
Notre catégorie étant abélienne, $F$ est exact.
Ainsi $F(E)$ a la même longueur que $E$ pour tout module de longueur finie $E$.
Puisque $\Hom(E, \kappa) = \Hom(F(E), F(\kappa)) \cong \Hom(F(E), \kappa)$, le lemme de Nakayama montre que $E$ et $F(E)$
ont le même nombre de générateurs. Ainsi $F(A/\mathfrak m^n)$ est un $A$-module cyclique.
Choisissons un générateur $e \in F(A/\mathfrak m^n)$.
Puisque $F$ est $A$-linéaire, on a $\mathfrak m^n e = 0$.
Le morphisme $A/\mathfrak m^n \to F(A/\mathfrak m^n)$ est nécessairement un isomorphisme, les longueurs étant égales.
Choisissons un élément
$$
e \in \lim F(A/\mathfrak m^n)
$$
dont l'image soit un générateur pour tout $n$ (nous omettons un petit argument).
On obtient alors un système d'isomorphismes $A/\mathfrak m^n \to F(A/\mathfrak m^n)$
compatible avec tous les homomorphismes de $A$-modules $A/\mathfrak m^n \to A/\mathfrak m^{n'}$
(encore par $A$-linéarité de $F$).
Puisque tout module de longueur finie est conoyau d'un morphisme
entre sommes directes de modules cycliques, on obtient l'isomorphisme du lemme.
\end{proof}

\begin{lemma}
\label{dualizing-lemma-dualizing-finite-length}
Soit $(A, \mathfrak m, \kappa)$ un anneau local noethérien muni d'un complexe dualisant
normalisé $\omega_A^\bullet$. Soit $E$ une enveloppe injective de $\kappa$.
Il existe alors un isomorphisme fonctoriel
$$
R\Hom_A(N, \omega_A^\bullet) = \Hom_A(N, E)[0]
$$
pour $N$ parcourant les $A$-modules de longueur finie.
\end{lemma}

\begin{proof}
Par récurrence sur la longueur de $N$, on voit que $R\Hom_A(N, \omega_A^\bullet)$
est un module de longueur finie placé en degré $0$.
Ainsi $R\Hom_A(-, \omega_A^\bullet)$ induit une anti-équivalence de la catégorie
des modules de longueur finie. Il en est de même de $\Hom_A(-, E)$
d'après la proposition \ref{dualizing-proposition-matlis} ; par conséquent,
$$
N \longmapsto \Hom_A(R\Hom_A(N, \omega_A^\bullet), E)
$$
est une équivalence comme dans le lemme \ref{dualizing-lemma-equivalence-finite-length}.
Elle est donc isomorphe au foncteur identité.
Puisque $\Hom_A(-, E)$ appliqué deux fois est l'identité
(proposition \ref{dualizing-proposition-matlis}), on obtient l'énoncé du lemme.
\end{proof}

\begin{lemma}
\label{dualizing-lemma-sitting-in-degrees}
Soit $(A, \mathfrak m, \kappa)$ un anneau local noethérien muni d'un complexe dualisant normalisé $\omega_A^\bullet$.
Soit $M$ un $A$-module de type fini et soit $d = \dim(\text{Supp}(M))$. Alors :
\begin{enumerate}
\item si $\Ext^i_A(M, \omega_A^\bullet)$ est non nul, alors $i \in \{-d, \ldots, 0\}$ ;
\item la dimension du support de $\Ext^i_A(M, \omega_A^\bullet)$ est au plus $-i$ ;
\item $\text{profondeur}(M)$ est le plus petit entier $\delta \geq 0$ tel que $\Ext^{-\delta}_A(M, \omega_A^\bullet) \not = 0$.
\end{enumerate}
\end{lemma}

\begin{proof}
Démontrons cela par récurrence sur $d$.
Si $d = 0$, cela résulte du lemme \ref{dualizing-lemma-dualizing-finite-length} et de la dualité de Matlis
(proposition \ref{dualizing-proposition-matlis}), qui assure que $\Hom_A(M, E)$ est non nul si $M$ est non nul.

\medskip\noindent
Supposons le résultat vrai pour les modules dont le support est de dimension
$< d$, et supposons $M$ de profondeur $> 0$.
Choisissons $f \in \mathfrak m$ qui soit un non-diviseur de zéro sur $M$
et considérons la suite exacte courte
$$
0 \to M \to M \to M/fM \to 0
$$
Puisque $\dim(\text{Supp}(M/fM)) = d - 1$ (Algèbre, lemme \ref{algebra-lemma-one-equation-module}), on peut appliquer l'hypothèse de récurrence.
En écrivant $E^i = \Ext^i_A(M, \omega_A^\bullet)$ et $F^i = \Ext^i_A(M/fM, \omega_A^\bullet)$, on obtient une suite exacte longue
$$
\ldots \to F^i \to E^i \xrightarrow{f} E^i \to F^{i + 1} \to \ldots
$$
Par récurrence, $E^i/fE^i = 0$ pour $i + 1 \not \in \{-\dim(\text{Supp}(M/fM)), \ldots, -\text{profondeur}(M/fM)\}$.
Le lemme de Nakayama (Algèbre, lemme \ref{algebra-lemma-NAK})
et Algèbre, lemme \ref{algebra-lemma-depth-drops-by-one}, donnent $E^i = 0$ pour $i \not \in \{-\dim(\text{Supp}(M)), \ldots, -\text{profondeur}(M)\}$.
De plus, dans le cas limite $i = - \text{profondeur}(M)$, on en déduit que $E^i$
est non nul puisque $F^{i + 1}$ est non nul par récurrence.
Puisque $E^i/fE^i \subset F^{i + 1}$, on obtient
$$
\dim(\text{Supp}(F^{i + 1})) \geq \dim(\text{Supp}(E^i/fE^i))
\geq \dim(\text{Supp}(E^i)) - 1
$$
(voir le lemme utilisé ci-dessus), ce qui donne aussi la majoration
de la dimension dans (2).

\medskip\noindent
Si $M$ est de profondeur $0$ et si $d > 0$,
posons $N = M[\mathfrak m^\infty]$ et $M' = M/N$ (comparer au lemme \ref{dualizing-lemma-divide-by-torsion}).
Alors $M'$ est de profondeur $> 0$ et $\dim(\text{Supp}(M')) = d$.
On connaît donc le résultat pour $M'$ et, puisque
$R\Hom_A(N, \omega_A^\bullet) = \Hom_A(N, E)$
(lemme \ref{dualizing-lemma-dualizing-finite-length}), la suite exacte longue de cohomologie des $\Ext$
donne le résultat pour $M$.
\end{proof}

\begin{remark}
\label{dualizing-remark-vanishing-for-arbitrary-modules}
Soient $(A, \mathfrak m)$ et $\omega_A^\bullet$ comme dans le lemme \ref{dualizing-lemma-sitting-in-degrees}.
D'après Compléments d'algèbre, lemme \ref{more-algebra-lemma-injective-amplitude},
on voit que $\omega_A^\bullet$ a son amplitude injective dans $[-d, 0]$,
car l'assertion (3) de ce lemme s'applique.
En particulier, pour tout $A$-module $M$ (pas nécessairement de type fini),
on a $\Ext^i_A(M, \omega_A^\bullet) = 0$ pour $i \not \in \{-d, \ldots, 0\}$.
\end{remark}

\begin{lemma}
\label{dualizing-lemma-local-CM}
Soit $(A, \mathfrak m, \kappa)$ un anneau local noethérien muni d'un complexe dualisant normalisé $\omega_A^\bullet$.
Soit $M$ un $A$-module de type fini. Les conditions suivantes sont équivalentes :
\begin{enumerate}
\item $M$ est de Cohen-Macaulay ;
\item $\Ext^i_A(M, \omega_A^\bullet)$ est non nul pour au plus un $i$ ;
\item $\Ext^{-i}_A(M, \omega_A^\bullet)$ est nul pour $i \not = \dim(\text{Supp}(M))$.
\end{enumerate}
Notons $CM_d$ la catégorie des $A$-modules de Cohen-Macaulay
de type fini et de profondeur $d$.
Alors $M \mapsto \Ext^{-d}_A(M, \omega_A^\bullet)$ définit une anti-auto-équivalence de $CM_d$.
\end{lemma}

\begin{proof}
Nous utiliserons les résultats du lemme \ref{dualizing-lemma-sitting-in-degrees} sans autre mention.
Fixons un module de type fini $M$. Si $M$ est de Cohen-Macaulay,
seul $\Ext^{-d}_A(M, \omega_A^\bullet)$ peut être non nul, d'où (1) $\Rightarrow$ (3).
L'implication (3) $\Rightarrow$ (2) est immédiate.
Supposons (2) et notons $N = \Ext^{-\delta}_A(M, \omega_A^\bullet)$ le groupe $\Ext$ non nul, où $\delta = \text{profondeur}(M)$.
Alors, puisque
$$
M[0] = R\Hom_A(R\Hom_A(M, \omega_A^\bullet), \omega_A^\bullet) =
R\Hom_A(N[\delta], \omega_A^\bullet)
$$
(lemme \ref{dualizing-lemma-dualizing}), on en déduit que $M = \Ext_A^{-\delta}(N, \omega_A^\bullet)$.
Ainsi $\delta \geq \dim(\text{Supp}(M))$. Cependant, sachant aussi que $\delta \leq \dim(\text{Supp}(M))$
(Algèbre, lemme \ref{algebra-lemma-bound-depth}), on en déduit que $M$ est de Cohen-Macaulay.

\medskip\noindent
Pour prouver la dernière assertion, il suffit de montrer que $N = \Ext^{-d}_A(M, \omega_A^\bullet)$
appartient à $CM_d$ pour $M$ dans $CM_d$.
Nous avons vu ci-dessus que $M[0] = R\Hom_A(N[d], \omega_A^\bullet)$, ce qui prouve le résultat voulu
par l'équivalence de (1) et (3).
\end{proof}

\begin{lemma}
\label{dualizing-lemma-dualizing-artinian}
Soit $(A, \mathfrak m, \kappa)$ un anneau local noethérien muni d'un complexe dualisant normalisé $\omega_A^\bullet$.
Si $\dim(A) = 0$, alors $\omega_A^\bullet \cong E[0]$,
où $E$ est une enveloppe injective du corps résiduel.
\end{lemma}

\begin{proof}
Cela résulte immédiatement du lemme \ref{dualizing-lemma-dualizing-finite-length}.
\end{proof}

\begin{lemma}
\label{dualizing-lemma-divide-by-finite-length-ideal}
Soit $(A, \mathfrak m, \kappa)$ un anneau local noethérien muni d'un complexe dualisant normalisé.
Soit $I \subset \mathfrak m$ un idéal de longueur finie. Posons $B = A/I$.
Il existe alors un triangle distingué
$$
\omega_B^\bullet \to \omega_A^\bullet \to \Hom_A(I, E)[0] \to
\omega_B^\bullet[1]
$$
dans $D(A)$, où $E$ est une enveloppe injective de $\kappa$
et $\omega_B^\bullet$ un complexe dualisant normalisé pour $B$.
\end{lemma}

\begin{proof}
Utiliser la suite exacte courte $0 \to I \to A \to B \to 0$
et les lemmes \ref{dualizing-lemma-dualizing-finite-length} et \ref{dualizing-lemma-normalized-quotient}.
\end{proof}

\begin{lemma}
\label{dualizing-lemma-divide-by-nonzerodivisor}
Soit $(A, \mathfrak m, \kappa)$ un anneau local noethérien muni d'un complexe dualisant normalisé $\omega_A^\bullet$.
Soit $f \in \mathfrak m$ un non-diviseur de zéro. Posons $B = A/(f)$.
Il existe alors un triangle distingué
$$
\omega_B^\bullet \to \omega_A^\bullet \to \omega_A^\bullet \to
\omega_B^\bullet[1]
$$
dans $D(A)$, où $\omega_B^\bullet$ est un complexe dualisant normalisé pour $B$.
\end{lemma}

\begin{proof}
Utiliser la suite exacte courte $0 \to A \to A \to B \to 0$ et le lemme \ref{dualizing-lemma-normalized-quotient}.
\end{proof}

\begin{lemma}
\label{dualizing-lemma-nonvanishing-generically-local}
Soit $(A, \mathfrak m, \kappa)$ un anneau local noethérien muni d'un complexe dualisant normalisé $\omega_A^\bullet$.
Soit $\mathfrak p$ un idéal premier minimal de $A$ tel que $\dim(A/\mathfrak p) = e$.
Alors $H^i(\omega_A^\bullet)_\mathfrak p$ est non nul si et seulement si $i = -e$.
\end{lemma}

\begin{proof}
Puisque $A_\mathfrak p$ est de dimension zéro, il existe un entier $n > 0$
tel que $\mathfrak p^nA_\mathfrak p$ soit nul. Posons $B = A/\mathfrak p^n$ et $\omega_B^\bullet = R\Hom_A(B, \omega_A^\bullet)$.
Puisque $B_\mathfrak p = A_\mathfrak p$, on voit que
$$
(\omega_B^\bullet)_\mathfrak p =
R\Hom_A(B, \omega_A^\bullet) \otimes_A^\mathbf{L} A_\mathfrak p =
R\Hom_{A_\mathfrak p}(B_\mathfrak p, (\omega_A^\bullet)_\mathfrak p) =
(\omega_A^\bullet)_\mathfrak p
$$
La deuxième égalité résulte de Compléments d'algèbre, lemme \ref{more-algebra-lemma-base-change-RHom}.
D'après le lemme \ref{dualizing-lemma-normalized-quotient}, on peut remplacer $A$ par $B$.
On a alors $\dim(A) = e$. Ainsi $H^i(\omega_A^\bullet)_\mathfrak p$ ne peut être non nul que si $i = -e$
d'après le lemme \ref{dualizing-lemma-sitting-in-degrees} (1) et (2).
D'autre part, puisque $(\omega_A^\bullet)_\mathfrak p$ est un complexe dualisant
pour l'anneau non nul $A_\mathfrak p$ (lemme \ref{dualizing-lemma-dualizing-localize}),
le module restant doit être non nul.
\end{proof}





\section{Complexes dualisants et fonctions de dimension}
\label{dualizing-section-dimension-function}

\noindent
Nos résultats dans le cas local ont la conséquence suivante :
un anneau noethérien possédant un complexe dualisant est
universellement caténaire et de dimension finie.

\begin{lemma}
\label{dualizing-lemma-nonvanishing-generically}
Soit $A$ un anneau noethérien. Soit $\mathfrak p$ un idéal premier minimal de $A$.
Alors $H^i(\omega_A^\bullet)_\mathfrak p$ est non nul pour exactement un $i$.
\end{lemma}

\begin{proof}
Le complexe $\omega_A^\bullet \otimes_A A_\mathfrak p$ est un complexe dualisant pour $A_\mathfrak p$
(lemme \ref{dualizing-lemma-dualizing-localize}). La dimension de $A_\mathfrak p$ est zéro puisque $\mathfrak p$ est minimal.
Le résultat découle donc du lemme \ref{dualizing-lemma-dualizing-artinian}.
\end{proof}

\noindent
Soient $A$ un anneau noethérien et $\omega_A^\bullet$ un complexe dualisant.
Le lemme \ref{dualizing-lemma-find-function} permet de définir une fonction
$$
\delta = \delta_{\omega_A^\bullet} : \Spec(A) \longrightarrow \mathbf{Z}
$$
en associant à $\mathfrak p$ l'entier du lemme \ref{dualizing-lemma-find-function}
pour le complexe dualisant $(\omega_A^\bullet)_\mathfrak p$ sur $A_\mathfrak p$
(lemme \ref{dualizing-lemma-dualizing-localize}) et le corps résiduel $\kappa(\mathfrak p)$.
Plus précisément, $\delta(\mathfrak p)$ est l'unique entier tel que
$$
(\omega_A^\bullet)_\mathfrak p[-\delta(\mathfrak p)]
$$
soit un complexe dualisant normalisé sur l'anneau local noethérien $A_\mathfrak p$.

\begin{lemma}
\label{dualizing-lemma-quotient-function}
Soient $A$ un anneau noethérien et $\omega_A^\bullet$ un complexe dualisant.
Soit $A \to B$ un homomorphisme surjectif d'anneaux et soit $\omega_B^\bullet = R\Hom(B, \omega_A^\bullet)$
le complexe dualisant pour $B$ du lemme \ref{dualizing-lemma-dualizing-quotient}. Alors on a
$$
\delta_{\omega_B^\bullet} = \delta_{\omega_A^\bullet}|_{\Spec(B)}
$$
\end{lemma}

\begin{proof}
Cela résulte de la définition des fonctions et du lemme \ref{dualizing-lemma-normalized-quotient}.
\end{proof}

\begin{lemma}
\label{dualizing-lemma-dimension-function}
Soient $A$ un anneau noethérien et $\omega_A^\bullet$ un complexe dualisant.
La fonction $\delta = \delta_{\omega_A^\bullet}$ définie ci-dessus est une fonction de dimension
(Topologie, définition \ref{topology-definition-dimension-function}).
\end{lemma}

\begin{proof}
Soit $\mathfrak p \subset \mathfrak q$ une spécialisation immédiate. Il faut montrer que $\delta(\mathfrak p) = \delta(\mathfrak q) + 1$.
On peut remplacer $A$ par $A/\mathfrak p$, le complexe $\omega_A^\bullet$ par $\omega_{A/\mathfrak p}^\bullet = R\Hom(A/\mathfrak p, \omega_A^\bullet)$,
l'idéal premier $\mathfrak p$ par $(0)$ et l'idéal premier $\mathfrak q$ par $\mathfrak q/\mathfrak p$ ;
voir le lemme \ref{dualizing-lemma-quotient-function}. On peut donc supposer que $A$ est intègre,
que $\mathfrak p = (0)$ et que $\mathfrak q$ est un idéal premier de hauteur $1$.

\medskip\noindent
Alors $H^i(\omega_A^\bullet)_{(0)}$ est non nul pour exactement un $i$,
disons $i_0$, d'après le lemme \ref{dualizing-lemma-nonvanishing-generically}.
En fait, $i_0 = -\delta((0))$ parce que $(\omega_A^\bullet)_{(0)}[-\delta((0))]$ est un complexe dualisant normalisé
sur le corps $A_{(0)}$.

\medskip\noindent
D'autre part, $(\omega_A^\bullet)_\mathfrak q[-\delta(\mathfrak q)]$ est un complexe dualisant normalisé pour $A_\mathfrak q$.
D'après le lemme \ref{dualizing-lemma-nonvanishing-generically-local}, on voit que
$$
H^e((\omega_A^\bullet)_\mathfrak q[-\delta(\mathfrak q)])_{(0)} =
H^{e - \delta(\mathfrak q)}(\omega_A^\bullet)_{(0)}
$$
n'est non nul que pour $e = -\dim(A_\mathfrak q) = -1$.
On en déduit
$$
-\delta((0)) = -1 - \delta(\mathfrak q)
$$
comme voulu.
\end{proof}

\begin{lemma}
\label{dualizing-lemma-universally-catenary}
Soit $A$ un anneau noethérien possédant un complexe dualisant.
Alors $A$ est universellement caténaire et de dimension finie.
\end{lemma}

\begin{proof}
Puisque $\Spec(A)$ possède une fonction de dimension d'après le lemme \ref{dualizing-lemma-dimension-function},
il est caténaire ; voir Topologie, lemme \ref{topology-lemma-dimension-function-catenary}.
Ainsi $A$ est caténaire ; voir Algèbre, lemme \ref{algebra-lemma-catenary}.
La proposition \ref{dualizing-proposition-dualizing-essentially-finite-type} entraîne alors que $A$ est universellement caténaire.

\medskip\noindent
Puisque tout complexe dualisant $\omega_A^\bullet$ appartient à $D^b_{\textit{Coh}}(A)$,
les valeurs de la fonction $\delta_{\omega_A^\bullet}$ aux idéaux premiers minimaux
sont bornées d'après le lemme \ref{dualizing-lemma-nonvanishing-generically}.
D'autre part, pour un idéal maximal $\mathfrak m$ de corps résiduel $\kappa$,
l'entier $i = -\delta(\mathfrak m)$ est l'unique entier tel que $\Ext_A^i(\kappa, \omega_A^\bullet)$ soit non nul
(lemme \ref{dualizing-lemma-find-function}).
Puisque $\omega_A^\bullet$ est de dimension injective finie, ces valeurs sont aussi bornées.
La dimension de $A$ étant la valeur maximale de $\delta(\mathfrak p) - \delta(\mathfrak m)$,
où $\mathfrak p \subset \mathfrak m$ parcourt les couples formés d'un idéal premier minimal
et d'un idéal maximal, on voit que la dimension de $\Spec(A)$ est bornée.
\end{proof}

\begin{lemma}
\label{dualizing-lemma-depth-dualizing-module}
Soit $(A, \mathfrak m, \kappa)$ un anneau local noethérien muni d'un complexe dualisant normalisé $\omega_A^\bullet$.
Soient $d = \dim(A)$ et $\omega_A = H^{-d}(\omega_A^\bullet)$. Alors :
\begin{enumerate}
\item le support de $\omega_A$ est la réunion des composantes irréductibles
de $\Spec(A)$ de dimension $d$ ;
\item $\omega_A$ satisfait à $(S_2)$, voir Algèbre, définition \ref{algebra-definition-conditions}.
\end{enumerate}
\end{lemma}

\begin{proof}
Nous utiliserons le lemme \ref{dualizing-lemma-sitting-in-degrees} sans autre mention.
D'après le lemme \ref{dualizing-lemma-nonvanishing-generically-local}, le support de $\omega_A$ contient les composantes
irréductibles de dimension $d$. Soit $\mathfrak p \subset A$ un idéal premier.
D'après le lemme \ref{dualizing-lemma-dimension-function}, le complexe $(\omega_A^\bullet)_{\mathfrak p}[-\dim(A/\mathfrak p)]$ est un complexe dualisant
normalisé pour $A_\mathfrak p$. Ainsi, si $\dim(A/\mathfrak p) + \dim(A_\mathfrak p) < d$, alors $(\omega_A)_\mathfrak p = 0$.
Cela prouve que le support de $\omega_A$ est la réunion des composantes
irréductibles de dimension $d$, car le complémentaire de cette réunion
est exactement formé des idéaux premiers $\mathfrak p$ de $A$ pour lesquels
$\dim(A/\mathfrak p) + \dim(A_\mathfrak p) < d$, puisque $A$ est caténaire (lemme \ref{dualizing-lemma-universally-catenary}).
D'autre part, si $\dim(A/\mathfrak p) + \dim(A_\mathfrak p) = d$, alors
$$
(\omega_A)_\mathfrak p =
H^{-\dim(A_\mathfrak p)}\left(
(\omega_A^\bullet)_{\mathfrak p}[-\dim(A/\mathfrak p)] \right)
$$
Pour prouver que $\omega_A$ satisfait à $(S_2)$, il suffit donc de montrer
que la profondeur de $\omega_A$ est au moins $\min(\dim(A), 2)$.
Procédons par récurrence sur $\dim(A)$. Le cas $\dim(A) = 0$ est trivial.

\medskip\noindent
Supposons $\text{profondeur}(A) > 0$. Choisissons un non-diviseur de zéro $f \in \mathfrak m$
et posons $B = A/fA$. Alors $\dim(B) = \dim(A) - 1$ et l'on peut appliquer
l'hypothèse de récurrence à $B$.
D'après le lemme \ref{dualizing-lemma-divide-by-nonzerodivisor}, la multiplication par $f$ est injective sur $\omega_A$
et l'on obtient $\omega_A/f\omega_A \subset \omega_B$. La profondeur de $\omega_A$ est donc au moins $1$.
Si $\dim(A) > 1$, alors $\dim(B) > 0$ et $\omega_B$ est de profondeur $ > 0$.
Ainsi $\omega_A$ est de profondeur $> 1$, ce qui conclut dans ce cas.

\medskip\noindent
Supposons $\dim(A) > 0$ et $\text{profondeur}(A) = 0$. Soit $I = A[\mathfrak m^\infty]$ et posons $B = A/I$.
Alors $B$ est de profondeur $\geq 1$ et $\omega_A = \omega_B$ d'après le lemme \ref{dualizing-lemma-divide-by-finite-length-ideal}.
Puisque le résultat a été démontré ci-dessus pour $\omega_B$,
la démonstration est achevée.
\end{proof}





\section{Le théorème de dualité locale}
\label{dualizing-section-local-duality}

\noindent
Le résultat principal de cette section est dû à Grothendieck.

\begin{lemma}
\label{dualizing-lemma-local-cohomology-of-dualizing}
Soit $(A, \mathfrak m, \kappa)$ un anneau local noethérien.
Soit $\omega_A^\bullet$ un complexe dualisant normalisé.
Soit $Z = V(\mathfrak m) \subset \Spec(A)$.
Alors $E = R^0\Gamma_Z(\omega_A^\bullet)$ est une enveloppe injective de $\kappa$ et $R\Gamma_Z(\omega_A^\bullet) = E[0]$.
\end{lemma}

\begin{proof}
D'après le lemme \ref{dualizing-lemma-local-cohomology-noetherian}, on a $R\Gamma_{\mathfrak m} = R\Gamma_Z$. Ainsi
$$
R\Gamma_Z(\omega_A^\bullet) =
R\Gamma_{\mathfrak m}(\omega_A^\bullet) =
\text{hocolim}\ R\Hom_A(A/\mathfrak m^n, \omega_A^\bullet)
$$
d'après le lemme \ref{dualizing-lemma-local-cohomology-ext}. Soit $E'$ une enveloppe injective du corps résiduel.
D'après le lemme \ref{dualizing-lemma-dualizing-finite-length}, on peut trouver des isomorphismes
$$
R\Hom_A(A/\mathfrak m^n, \omega_A^\bullet) \cong \Hom_A(A/\mathfrak m^n, E')[0]
$$
compatibles avec les morphismes de transition. Puisque
$E' = \bigcup E'[\mathfrak m^n] = \colim \Hom_A(A/\mathfrak m^n, E')$
d'après le lemme \ref{dualizing-lemma-union-artinian}, on en déduit que $E \cong E'$
et que tous les autres groupes de cohomologie du complexe $R\Gamma_Z(\omega_A^\bullet)$ sont nuls.
\end{proof}

\begin{remark}
\label{dualizing-remark-specific-injective-hull}
Soit $(A, \mathfrak m, \kappa)$ un anneau local noethérien muni d'un complexe dualisant normalisé $\omega_A^\bullet$.
D'après le lemme \ref{dualizing-lemma-local-cohomology-of-dualizing} ci-dessus, $R\Gamma_Z(\omega_A^\bullet)$ est une enveloppe injective
du corps résiduel placée en degré $0$.
Cela donne en fait une « construction » ou une « réalisation »
de l'enveloppe injective un peu plus canonique que le choix
d'une enveloppe injective quelconque.
En effet, un complexe dualisant normalisé est unique à isomorphisme près,
et son groupe d'automorphismes est le groupe des unités de $A$,
tandis qu'une enveloppe injective de $\kappa$ est unique à isomorphisme près,
avec pour groupe d'automorphismes le groupe des unités du complété
$A^\wedge$ de $A$ relativement à $\mathfrak m$.
\end{remark}

\noindent
Voici le résultat principal de cette section.

\begin{theorem}
\label{dualizing-theorem-local-duality}
Soit $(A, \mathfrak m, \kappa)$ un anneau local noethérien.
Soit $\omega_A^\bullet$ un complexe dualisant normalisé.
Soit $E$ une enveloppe injective du corps résiduel.
Soit $Z = V(\mathfrak m) \subset \Spec(A)$.
Notons ${}^\wedge$ la complétion dérivée relativement à $\mathfrak m$.
Alors
$$
R\Hom_A(K, \omega_A^\bullet)^\wedge \cong R\Hom_A(R\Gamma_Z(K), E[0])
$$
pour $K$ dans $D(A)$.
\end{theorem}

\begin{proof}
Observons que $E[0] \cong R\Gamma_Z(\omega_A^\bullet)$ d'après le lemme \ref{dualizing-lemma-local-cohomology-of-dualizing}.
D'après Compléments d'algèbre, lemme \ref{more-algebra-lemma-completion-RHom},
la complétion dans le membre de gauche passe à l'intérieur.
Il faut donc prouver
$$
R\Hom_A(K^\wedge, (\omega_A^\bullet)^\wedge)
=
R\Hom_A(R\Gamma_Z(K), R\Gamma_Z(\omega_A^\bullet))
$$
Cela résulte de l'équivalence entre $D_{comp}(A, \mathfrak m)$ et $D_{\mathfrak m^\infty\text{-torsion}}(A)$
donnée dans la proposition \ref{dualizing-proposition-torsion-complete}.
Plus précisément, c'est un cas particulier du lemme \ref{dualizing-lemma-compare-RHom}.
\end{proof}

\noindent
Voici un cas particulier du théorème précédent.

\begin{lemma}
\label{dualizing-lemma-special-case-local-duality}
Soit $(A, \mathfrak m, \kappa)$ un anneau local noethérien.
Soit $\omega_A^\bullet$ un complexe dualisant normalisé.
Soit $E$ une enveloppe injective du corps résiduel.
Soit $K \in D_{\textit{Coh}}(A)$. Alors
$$
\Ext^{-i}_A(K, \omega_A^\bullet)^\wedge =
\Hom_A(H^i_{\mathfrak m}(K), E)
$$
où ${}^\wedge$ désigne la complétion $\mathfrak m$-adique.
\end{lemma}

\begin{proof}
D'après le lemme \ref{dualizing-lemma-dualizing}, on voit que $R\Hom_A(K, \omega_A^\bullet)$ est un objet de $D_{\textit{Coh}}(A)$.
Il en résulte que les modules de cohomologie du complété dérivé de $R\Hom_A(K, \omega_A^\bullet)$
sont les complétés usuels $\Ext^i_A(K, \omega_A^\bullet)^\wedge$ d'après
Compléments d'algèbre, lemme \ref{more-algebra-lemma-derived-completion-pseudo-coherent}.
D'autre part, on a $R\Gamma_{\mathfrak m} = R\Gamma_Z$ pour $Z = V(\mathfrak m)$ d'après le lemme \ref{dualizing-lemma-local-cohomology-noetherian}.
De plus, le foncteur $\Hom_A(-, E)$ est exact, donc se factorise par la cohomologie.
Le lemme est ainsi une conséquence du théorème \ref{dualizing-theorem-local-duality}.
\end{proof}





\section{Modules dualisants}
\label{dualizing-section-dualizing-module}

\noindent
Si $(A, \mathfrak m, \kappa)$ est un anneau local noethérien et $\omega_A^\bullet$ un complexe dualisant normalisé,
on dit que le module $\omega_A = H^{-\dim(A)}(\omega_A^\bullet)$, décrit dans le lemme \ref{dualizing-lemma-depth-dualizing-module},
est un {\it module dualisant} pour $A$.
Ce module est un module canonique de $A$.
Il semble généralement admis de définir un {\it module canonique}
pour un anneau local noethérien $(A, \mathfrak m, \kappa)$ comme un $A$-module de type fini
$K$ tel que
$$
\Hom_A(K, E) \cong H^{\dim(A)}_\mathfrak m(A)
$$
où $E$ est une enveloppe injective du corps résiduel.
Un module dualisant est canonique parce que
$$
\Hom_A(H^{\dim(A)}_\mathfrak m(A), E) = (\omega_A)^\wedge
$$
d'après le lemme \ref{dualizing-lemma-special-case-local-duality} ; en appliquant $\Hom_A(-, E)$, on obtient donc
\begin{align*}
\Hom_A(\omega_A, E)
& =
\Hom_A((\omega_A)^\wedge, E) \\
& =
\Hom_A(\Hom_A(H^{\dim(A)}_\mathfrak m(A), E), E) \\
& = H^{\dim(A)}_\mathfrak m(A)
\end{align*}
La première égalité vient de ce que $E$ est de torsion par les puissances
de $\mathfrak m$, la deuxième de ce qui précède et la troisième de la dualité de Matlis
(proposition \ref{dualizing-proposition-matlis}).
L'intérêt de la définition du module canonique donnée ci-dessus
est qu'elle a un sens même si $A$ ne possède pas de complexe dualisant.




\section{Anneaux de Cohen-Macaulay}
\label{dualizing-section-CM}

\noindent
Les modules et anneaux de Cohen-Macaulay ont été étudiés dans
Algèbre, sections \ref{algebra-section-CM} et \ref{algebra-section-CM-ring}.

\begin{lemma}
\label{dualizing-lemma-depth-in-terms-dualizing-complex}
Soit $(A, \mathfrak m, \kappa)$ un anneau local noethérien muni d'un complexe dualisant normalisé $\omega_A^\bullet$.
Alors $\text{profondeur}(A)$ est égale au plus petit entier $\delta \geq 0$ tel que $H^{-\delta}(\omega_A^\bullet) \not = 0$.
\end{lemma}

\begin{proof}
Cela résulte immédiatement du lemme \ref{dualizing-lemma-sitting-in-degrees}.
Voici deux autres façons de le voir.

\medskip\noindent
Première variante. Le lemme de Nakayama montre que $\delta$
est le plus petit entier tel que $\Hom_A(H^{-\delta}(\omega_A^\bullet), \kappa) \not = 0$.
Autrement dit, c'est le plus petit entier tel que $\Ext_A^{-\delta}(\omega_A^\bullet, \kappa)$ soit non nul.
Le lemme \ref{dualizing-lemma-dualizing} et le fait que $\omega_A^\bullet$ soit normalisé montrent qu'il s'agit
du plus petit entier tel que $\Ext_A^\delta(\kappa, A)$ soit non nul.
C'est la profondeur de $A$ d'après Algèbre, lemme \ref{algebra-lemma-depth-ext}.

\medskip\noindent
Deuxième variante. Le théorème de dualité locale
(sous la forme du lemme \ref{dualizing-lemma-special-case-local-duality}) montre que $\delta$ est le plus petit entier
tel que $H^\delta_\mathfrak m(A)$ soit non nul. C'est la profondeur de $A$
d'après le lemme \ref{dualizing-lemma-depth}.
\end{proof}

\begin{lemma}
\label{dualizing-lemma-apply-CM}
Soit $(A, \mathfrak m, \kappa)$ un anneau local noethérien muni d'un complexe dualisant normalisé $\omega_A^\bullet$
et d'un module dualisant $\omega_A = H^{-\dim(A)}(\omega_A^\bullet)$. Les conditions suivantes sont équivalentes :
\begin{enumerate}
\item $A$ est de Cohen-Macaulay ;
\item $\omega_A^\bullet$ est concentré en un seul degré ;
\item $\omega_A^\bullet = \omega_A[\dim(A)]$.
\end{enumerate}
Dans ce cas, $\omega_A$ est un module de Cohen-Macaulay maximal.
\end{lemma}

\begin{proof}
Cela résulte immédiatement du lemme \ref{dualizing-lemma-local-CM}.
\end{proof}

\begin{lemma}
\label{dualizing-lemma-has-dualizing-module-CM}
Soit $A$ un anneau noethérien. S'il existe un $A$-module de type fini $\omega_A$
tel que $\omega_A[0]$ soit un complexe dualisant, alors $A$ est de Cohen-Macaulay.
\end{lemma}

\begin{proof}
On peut remplacer $A$ par son localisé en un idéal premier
(lemme \ref{dualizing-lemma-dualizing-localize} et Algèbre, définition \ref{algebra-definition-ring-CM}).
Dans ce cas, le résultat découle immédiatement du lemme \ref{dualizing-lemma-apply-CM}.
\end{proof}

\begin{lemma}
\label{dualizing-lemma-CM-open}
Soit $A$ un anneau noethérien muni d'un complexe dualisant $\omega_A^\bullet$.
Soit $M$ un $A$-module de type fini. Alors
$$
U = \{\mathfrak p \in \Spec(A) \mid M_\mathfrak p\text{ est de Cohen-Macaulay}\}
$$
est un ouvert de $\Spec(A)$ dont l'intersection avec $\text{Supp}(M)$ est dense.
\end{lemma}

\begin{proof}
Si $\mathfrak p$ est un point générique de $\text{Supp}(M)$, alors $\text{profondeur}(M_\mathfrak p) = \dim(M_\mathfrak p) = 0$,
donc $\mathfrak p \in U$. Cela prouve la densité.
Si $\mathfrak p \in U$, on voit que
$$
R\Hom_A(M, \omega_A^\bullet)_\mathfrak p =
R\Hom_{A_\mathfrak p}(M_\mathfrak p, (\omega_A^\bullet)_\mathfrak p)
$$
a un unique module de cohomologie non nul, disons en degré $i_0$,
d'après le lemme \ref{dualizing-lemma-local-CM}.
Puisque $R\Hom_A(M, \omega_A^\bullet)$ n'a qu'un nombre fini de modules de cohomologie non nuls $H^i$,
et que chacun est un $A$-module de type fini,
il existe $f \in A$ tel que $f \not \in \mathfrak p$ et $(H^i)_f = 0$ pour $i \not = i_0$.
Alors $R\Hom_A(M, \omega_A^\bullet)_f$ a un unique module de cohomologie non nul ;
en reprenant les arguments précédents en sens inverse, on trouve $D(f) \subset U$.
\end{proof}

\begin{lemma}
\label{dualizing-lemma-CM}
Soit $A$ un anneau noethérien. Si $A$ possède un complexe dualisant $\omega_A^\bullet$, alors
$\{\mathfrak p \in \Spec(A) \mid A_\mathfrak p\text{ est de Cohen-Macaulay}\}$
est un ouvert dense de $\Spec(A)$.
\end{lemma}

\begin{proof}
Conséquence immédiate du lemme \ref{dualizing-lemma-CM-open} et des définitions.
\end{proof}






\section{Anneaux de Gorenstein}
\label{dualizing-section-gorenstein}

\noindent
Jusqu'ici, le seul complexe dualisant explicite que nous ayons rencontré
est $\kappa$ sur $\kappa$ pour un corps $\kappa$ ; voir la démonstration du lemme \ref{dualizing-lemma-find-function}.
D'après la proposition \ref{dualizing-proposition-dualizing-essentially-finite-type}, cela signifie que toute algèbre de type fini
sur un corps possède un complexe dualisant.
Cependant, il existe des anneaux noethériens (locaux) sans complexe dualisant.
En effet, nous avons vu qu'un anneau possédant un complexe dualisant
est universellement caténaire (lemme \ref{dualizing-lemma-universally-catenary}),
mais il existe des exemples d'anneaux locaux noethériens non caténaires ;
voir Exemples, section \ref{examples-section-non-catenary-Noetherian-local}.

\medskip\noindent
Néanmoins, beaucoup d'anneaux de la géométrie algébrique possèdent
des complexes dualisants simplement parce qu'ils sont quotients
d'anneaux de Gorenstein. Cette condition est en fait nécessaire et suffisante :
un anneau noethérien possède un complexe dualisant si et seulement s'il est
quotient d'un anneau de Gorenstein de dimension finie.
C'est la conjecture de Sharp (\cite{Sharp}), que l'on trouve dans la littérature
sous la forme \cite[corollaire 1.4]{Kawasaki}. Revenons à notre sujet et définissons les anneaux de Gorenstein.

\begin{definition}
\label{dualizing-definition-gorenstein}
Anneaux de Gorenstein.
\begin{enumerate}
\item Soit $A$ un anneau local noethérien. On dit que $A$ est
{\it de Gorenstein} si $A[0]$ est un complexe dualisant pour $A$.
\item Soit $A$ un anneau noethérien. On dit que $A$ est
{\it de Gorenstein} si $A_\mathfrak p$ est de Gorenstein
pour tout idéal premier $\mathfrak p$ de $A$.
\end{enumerate}
\end{definition}

\noindent
Cette définition a un sens, car, si $A[0]$ est un complexe dualisant
pour $A$, alors $S^{-1}A[0]$ est un complexe dualisant pour $S^{-1}A$
d'après le lemme \ref{dualizing-lemma-dualizing-localize}.
Un anneau noethérien de dimension finie est de Gorenstein
s'il est de dimension injective finie comme module sur lui-même ;
voir le lemme \ref{dualizing-lemma-characterize-gorenstein}.

\begin{lemma}
\label{dualizing-lemma-gorenstein-CM}
Un anneau de Gorenstein est de Cohen-Macaulay.
\end{lemma}

\begin{proof}
Cela résulte du lemme \ref{dualizing-lemma-apply-CM}.
\end{proof}

\noindent
Un anneau régulier est un exemple d'anneau de Gorenstein.

\begin{lemma}
\label{dualizing-lemma-regular-gorenstein}
Un anneau local régulier est de Gorenstein.
Un anneau régulier est de Gorenstein.
\end{lemma}

\begin{proof}
Soit $A$ un anneau régulier de dimension finie $d$.
Alors $A$ est de dimension globale finie $d$ ; voir Algèbre, lemme \ref{algebra-lemma-finite-gl-dim-finite-dim-regular}.
Ainsi $\Ext^{d + 1}_A(M, A) = 0$ pour tout $A$-module $M$ ; voir Algèbre, lemme \ref{algebra-lemma-projective-dimension-ext}.
Par conséquent, $A$ est de dimension injective finie comme $A$-module
d'après Compléments d'algèbre, lemme \ref{more-algebra-lemma-injective-amplitude}.
Il en résulte que $A[0]$ est un complexe dualisant, donc $A$
est de Gorenstein d'après la remarque qui suit la définition.
\end{proof}

\begin{lemma}
\label{dualizing-lemma-gorenstein}
Soit $A$ un anneau noethérien.
\begin{enumerate}
\item Si $A$ possède un complexe dualisant $\omega_A^\bullet$, alors
\begin{enumerate}
\item $A$ est de Gorenstein $\Leftrightarrow$ $\omega_A^\bullet$ est un objet inversible de $D(A)$ ;
\item $A_\mathfrak p$ est de Gorenstein $\Leftrightarrow$
$(\omega_A^\bullet)_\mathfrak p$ est un objet inversible de $D(A_\mathfrak p)$ ;
\item $\{\mathfrak p \in \Spec(A) \mid A_\mathfrak p\text{ est de Gorenstein}\}$ est un ouvert.
\end{enumerate}
\item Si $A$ est de Gorenstein, alors $A$ possède un complexe dualisant
si et seulement si $A[0]$ est un complexe dualisant.
\end{enumerate}
\end{lemma}

\begin{proof}
Pour les objets inversibles de $D(A)$, voir Compléments d'algèbre,
lemme \ref{more-algebra-lemma-invertible-derived} et la discussion de la section \ref{dualizing-section-dualizing}.

\medskip\noindent
D'après le lemme \ref{dualizing-lemma-dualizing-localize}, pour tout $\mathfrak p$, le complexe $(\omega_A^\bullet)_\mathfrak p$
est un complexe dualisant sur $A_\mathfrak p$.
La définition et l'unicité des complexes dualisants (lemme \ref{dualizing-lemma-dualizing-unique})
donnent (1)(b).

\medskip\noindent
Pour voir (1)(c), supposons $A_\mathfrak p$ de Gorenstein.
Soit $n_x$ l'unique entier tel que $H^{n_{x}}((\omega_A^\bullet)_\mathfrak p)$ soit non nul et isomorphe à $A_\mathfrak p$.
Puisque $\omega_A^\bullet$ appartient à $D^b_{\textit{Coh}}(A)$, il n'y a qu'un nombre fini
de $A$-modules de type fini non nuls $H^i(\omega_A^\bullet)$.
Il existe donc $f \in A$ tel que $f \not \in \mathfrak p$ et tel que seul $H^{n_x}((\omega_A^\bullet)_f)$
soit non nul et engendré par $1$ élément sur $A_f$.
Les complexes dualisants étant fidèles par définition, on en déduit que $A_f \cong H^{n_x}((\omega_A^\bullet)_f)$.
Ainsi $A_\mathfrak q$ est de Gorenstein pour tout $\mathfrak q \in D(f)$.
Cela prouve que l'ensemble de (1)(c) est ouvert.

\medskip\noindent
Preuve de (1)(a). L'implication $\Leftarrow$ résulte de (1)(b).
L'implication $\Rightarrow$ découle du paragraphe précédent,
où l'on a montré que, si $A_\mathfrak p$ est de Gorenstein,
il existe $f \in A$ tel que $f \not \in \mathfrak p$ et tel que le complexe $(\omega_A^\bullet)_f$
n'ait qu'un module de cohomologie non nul, qui est inversible.

\medskip\noindent
Si $A[0]$ est un complexe dualisant, alors $A$ est de Gorenstein d'après (1).
Réciproquement, (1) montre que $\omega_A^\bullet$ est localement isomorphe à un décalé de $A$.
Puisque la propriété d'être un complexe dualisant est locale
(lemme \ref{dualizing-lemma-dualizing-glue}), le résultat est clair.
\end{proof}

\begin{lemma}
\label{dualizing-lemma-gorenstein-ext}
Soit $(A, \mathfrak m, \kappa)$ un anneau local noethérien.
Alors $A$ est de Gorenstein si et seulement si $\Ext^i_A(\kappa, A)$ est nul pour $i \gg 0$.
\end{lemma}

\begin{proof}
Observons que $A[0]$ est un complexe dualisant pour $A$ si et seulement si
$A$ est de dimension injective finie comme $A$-module
(cela résulte immédiatement de la définition \ref{dualizing-definition-dualizing}).
Le lemme découle donc de Compléments d'algèbre, lemme \ref{more-algebra-lemma-finite-injective-dimension-Noetherian-local}.
\end{proof}

\begin{lemma}
\label{dualizing-lemma-characterize-gorenstein}
Soit $A$ un anneau noethérien. Les conditions suivantes sont équivalentes :
\begin{enumerate}
\item $A$ est de dimension injective finie comme module sur lui-même ;
\item $A$ est de Gorenstein et de dimension finie.
\end{enumerate}
Dans ce cas, $A[0]$ est un complexe dualisant.
\end{lemma}

\begin{proof}
Si $A$ est de dimension injective finie comme module sur $A$,
alors $A[0]$ est un complexe dualisant sur $A$
(voir la définition \ref{dualizing-definition-dualizing}). On en déduit que $A$ est de dimension finie
d'après le lemme \ref{dualizing-lemma-universally-catenary} et de Gorenstein d'après le lemme \ref{dualizing-lemma-gorenstein}.
Supposons $A$ de Gorenstein et $\dim(A) \leq d < \infty$.
Soit $I \subset A$ un idéal quelconque. Affirmons que $\Ext^i_A(A/I, A) = 0$ pour $i > d$ ;
cela prouve que $A$ est de dimension injective finie comme module
sur lui-même d'après Compléments d'algèbre, lemme \ref{more-algebra-lemma-injective-amplitude}.
La formation des groupes $\Ext$ commute à la localisation
(voir par exemple Compléments d'algèbre, lemme \ref{more-algebra-lemma-base-change-RHom}).
Pour prouver l'annulation, on peut donc supposer $A$ local de dimension $e < d$.
Puisque $A$ est de Gorenstein, $\omega_A^\bullet = A[e]$ est un complexe dualisant normalisé.
L'annulation annoncée découle alors, par exemple, du lemme \ref{dualizing-lemma-sitting-in-degrees}.
\end{proof}

\begin{lemma}
\label{dualizing-lemma-gorenstein-divide-by-nonzerodivisor}
Soit $(A, \mathfrak m, \kappa)$ un anneau local noethérien. Soit $f \in \mathfrak m$ un non-diviseur de zéro.
Posons $B = A/(f)$. Alors $A$ est de Gorenstein si et seulement si $B$ l'est.
\end{lemma}

\begin{proof}
Si $A$ est de Gorenstein, alors $B$ l'est d'après le lemme \ref{dualizing-lemma-divide-by-nonzerodivisor}.
Réciproquement, supposons $B$ de Gorenstein.
Alors $\Ext^i_B(\kappa, B)$ est nul pour $i \gg 0$ (lemme \ref{dualizing-lemma-gorenstein-ext}).
Rappelons que $R\Hom(B, -) : D(A) \to D(B)$ est adjoint à droite de la restriction (lemme \ref{dualizing-lemma-right-adjoint}).
Ainsi
$$
R\Hom_A(\kappa, A) = R\Hom_B(\kappa, R\Hom(B, A)) =
R\Hom_B(\kappa, B[1])
$$
La dernière égalité s'obtient par calcul direct ou par le lemme \ref{dualizing-lemma-compute-for-effective-Cartier-algebraic}.
On voit donc que $\Ext^i_A(\kappa, A)$ est nul pour $i \gg 0$,
et $A$ est de Gorenstein (lemme \ref{dualizing-lemma-gorenstein-ext}).
\end{proof}

\begin{lemma}
\label{dualizing-lemma-gorenstein-lci}
Si $A \to B$ est un homomorphisme d'anneaux localement d'intersection complète
et si $A$ est un anneau noethérien de Gorenstein,
alors $B$ est un anneau de Gorenstein.
\end{lemma}

\begin{proof}
D'après Compléments d'algèbre, définition \ref{more-algebra-definition-local-complete-intersection},
on peut écrire $B = A[x_1, \ldots, x_n]/I$, où $I$ est un idéal Koszul-régulier.
Observons qu'un anneau de polynômes sur un anneau de Gorenstein $A$
est de Gorenstein : se ramener au cas où $A$ est local,
puis utiliser les lemmes \ref{dualizing-lemma-dualizing-polynomial-ring} et \ref{dualizing-lemma-gorenstein}.
Un idéal Koszul-régulier est, par définition, localement engendré
par une suite Koszul-régulière ; voir Compléments d'algèbre, section \ref{more-algebra-section-ideals}.
En examinant les anneaux locaux de $A[x_1, \ldots, x_n]$, on voit qu'il suffit de montrer ceci :
si $R$ est un anneau local noethérien de Gorenstein et si $f_1, \ldots, f_c \in \mathfrak m_R$
est une suite Koszul-régulière, alors $R/(f_1, \ldots, f_c)$ est de Gorenstein.
Cela résulte du lemme \ref{dualizing-lemma-gorenstein-divide-by-nonzerodivisor} et du fait qu'une suite Koszul-régulière
dans $R$ est simplement une suite régulière
(Compléments d'algèbre, lemme \ref{more-algebra-lemma-noetherian-finite-all-equivalent}).
\end{proof}

\begin{lemma}
\label{dualizing-lemma-flat-under-gorenstein}
Soit $A \to B$ un homomorphisme local plat d'anneaux locaux noethériens.
Les conditions suivantes sont équivalentes :
\begin{enumerate}
\item $B$ est de Gorenstein ;
\item $A$ et $B/\mathfrak m_A B$ sont de Gorenstein.
\end{enumerate}
\end{lemma}

\begin{proof}
Nous utiliserons ci-dessous sans autre mention le fait qu'un anneau local
de Gorenstein est de dimension injective finie, ainsi que le lemme \ref{dualizing-lemma-gorenstein-ext}.
D'après Compléments d'algèbre, lemme \ref{more-algebra-lemma-pseudo-coherence-and-base-change-ext}, on a
$$
\Ext^i_A(\kappa_A, A) \otimes_A B =
\Ext^i_B(B/\mathfrak m_A B, B)
$$
pour tout $i$.

\medskip\noindent
Supposons (2). Puisque $R\Hom(B/\mathfrak m_A B, -) : D(B) \to D(B/\mathfrak m_A B)$ est adjoint à droite de la restriction
(lemme \ref{dualizing-lemma-right-adjoint}), on obtient
$$
R\Hom_B(\kappa_B, B) =
R\Hom_{B/\mathfrak m_A B}(\kappa_B, R\Hom(B/\mathfrak m_A B, B))
$$
Les modules de cohomologie de $R\Hom(B/\mathfrak m_A B, B)$ sont les modules $\Ext^i_B(B/\mathfrak m_A B, B) =
\Ext^i_A(\kappa_A, A) \otimes_A B$.
Puisque $A$ est de Gorenstein, seul un nombre fini d'entre eux est non nul
et chacun est isomorphe à une somme directe de copies de $B/\mathfrak m_A B$.
Puisque $B/\mathfrak m_A B$ est de Gorenstein, on en déduit que $R\Hom_B(B/\mathfrak m_B, B)$
n'a qu'un nombre fini de modules de cohomologie non nuls.
Ainsi $B$ est de Gorenstein.

\medskip\noindent
Supposons (1). Puisque $B$ est de dimension injective finie,
$\Ext^i_B(B/\mathfrak m_A B, B)$ est $0$ pour $i \gg 0$.
Puisque $A \to B$ est fidèlement plat, on en déduit que $\Ext^i_A(\kappa_A, A)$ est $0$ pour $i \gg 0$.
Ainsi $A$ est de Gorenstein. Cela entraîne que $\Ext^i_A(\kappa_A, A)$ est non nul
pour exactement un $i$, à savoir $i = \dim(A)$, et que $\Ext^{\dim(A)}_A(\kappa_A, A) \cong \kappa_A$
(voir les lemmes \ref{dualizing-lemma-normalized-finite}, \ref{dualizing-lemma-apply-CM} et \ref{dualizing-lemma-gorenstein-CM}).
On voit donc que $\Ext^i_B(B/\mathfrak m_A B, B)$ est nul sauf pour un $i$,
à savoir $i = \dim(A)$, et que $\Ext^{\dim(A)}_B(B/\mathfrak m_A B, B) \cong B/\mathfrak m_A B$.
Ainsi $B/\mathfrak m_A B$ est de Gorenstein d'après le lemme \ref{dualizing-lemma-normalized-finite}.
\end{proof}

\begin{lemma}
\label{dualizing-lemma-tor-injective-hull}
Soit $(A, \mathfrak m, \kappa)$ un anneau local noethérien de Gorenstein de dimension $d$.
Soit $E$ l'enveloppe injective de $\kappa$.
Alors $\text{Tor}_i^A(E, \kappa)$ est nul pour $i \not = d$ et $\text{Tor}_d^A(E, \kappa) = \kappa$.
\end{lemma}

\begin{proof}
Puisque $A$ est de Gorenstein, $\omega_A^\bullet = A[d]$ est un complexe dualisant normalisé pour $A$.
De plus, $E$ est l'unique module de cohomologie non nul de $R\Gamma_\mathfrak m(\omega_A^\bullet)$,
placé en degré $0$ ; voir le lemme \ref{dualizing-lemma-local-cohomology-of-dualizing}.
D'après le lemme \ref{dualizing-lemma-torsion-tensor-product}, on a
$$
E \otimes_A^\mathbf{L} \kappa =
R\Gamma_\mathfrak m(\omega_A^\bullet) \otimes_A^\mathbf{L} \kappa =
R\Gamma_\mathfrak m(\omega_A^\bullet \otimes_A^\mathbf{L} \kappa) =
R\Gamma_\mathfrak m(\kappa[d]) = \kappa[d]
$$
d'où le lemme.
\end{proof}




\section{L'abondance des complexes dualisants}
\label{dualizing-section-ubiquity-dualizing}

\noindent
Beaucoup d'anneaux noethériens possèdent des complexes dualisants.

\begin{lemma}
\label{dualizing-lemma-flat-unramified}
Soit $A \to B$ un homomorphisme local d'anneaux locaux noethériens.
Soit $\omega_A^\bullet$ un complexe dualisant normalisé.
Si $A \to B$ est plat et si $\mathfrak m_A B = \mathfrak m_B$,
alors $\omega_A^\bullet \otimes_A B$ est un complexe dualisant normalisé pour $B$.
\end{lemma}

\begin{proof}
Il est clair que $\omega_A^\bullet \otimes_A B$ appartient à $D^b_{\textit{Coh}}(B)$.
Soient $\kappa_A$ et $\kappa_B$ les corps résiduels de $A$ et de $B$.
D'après Compléments d'algèbre, lemme \ref{more-algebra-lemma-base-change-RHom}, on voit que
$$
R\Hom_B(\kappa_B, \omega_A^\bullet \otimes_A B) =
R\Hom_A(\kappa_A, \omega_A^\bullet) \otimes_A B =
\kappa_A[0] \otimes_A B = \kappa_B[0]
$$
Ainsi $\omega_A^\bullet \otimes_A B$ est de dimension injective finie d'après
Compléments d'algèbre, lemme \ref{more-algebra-lemma-finite-injective-dimension-Noetherian-local}.
Enfin, les mêmes arguments donnent
$$
R\Hom_B(\omega_A^\bullet \otimes_A B, \omega_A^\bullet \otimes_A B) =
R\Hom_A(\omega_A^\bullet, \omega_A^\bullet) \otimes_A B = A \otimes_A B = B
$$
comme voulu.
\end{proof}

\begin{lemma}
\label{dualizing-lemma-flat-iso-mod-I}
Soit $A \to B$ un homomorphisme plat d'anneaux noethériens.
Soit $I \subset A$ un idéal tel que $A/I = B/IB$ et tel que $IB$
soit contenu dans le radical de Jacobson de $B$.
Soit $\omega_A^\bullet$ un complexe dualisant.
Alors $\omega_A^\bullet \otimes_A B$ est un complexe dualisant pour $B$.
\end{lemma}

\begin{proof}
Il est clair que $\omega_A^\bullet \otimes_A B$ appartient à $D^b_{\textit{Coh}}(B)$.
D'après Compléments d'algèbre, lemme \ref{more-algebra-lemma-base-change-RHom}, on voit que
$$
R\Hom_B(K \otimes_A B, \omega_A^\bullet \otimes_A B) =
R\Hom_A(K, \omega_A^\bullet) \otimes_A B
$$
pour tout $K \in D^b_{\textit{Coh}}(A)$. Pour tout idéal $IB \subset J \subset B$,
il existe un unique idéal $I \subset J' \subset A$ tel que $A/J' \otimes_A B = B/J$.
Ainsi $\omega_A^\bullet \otimes_A B$ est de dimension injective finie
d'après Compléments d'algèbre, lemme \ref{more-algebra-lemma-finite-injective-dimension-Noetherian-radical}.
Enfin, on a aussi
$$
R\Hom_B(\omega_A^\bullet \otimes_A B, \omega_A^\bullet \otimes_A B) =
R\Hom_A(\omega_A^\bullet, \omega_A^\bullet) \otimes_A B = A \otimes_A B = B
$$
comme voulu.
\end{proof}

\begin{lemma}
\label{dualizing-lemma-completion-henselization-dualizing}
Soient $A$ un anneau noethérien et $I \subset A$ un idéal.
Soit $\omega_A^\bullet$ un complexe dualisant.
\begin{enumerate}
\item $\omega_A^\bullet \otimes_A A^h$ est un complexe dualisant sur l'hensélisé $(A^h, I^h)$ du couple $(A, I)$ ;
\item $\omega_A^\bullet \otimes_A A^\wedge$ est un complexe dualisant sur le complété $I$-adique $A^\wedge$ ;
\item si $A$ est local, alors $\omega_A^\bullet \otimes_A A^h$,
resp.\ $\omega_A^\bullet \otimes_A A^{sh}$, est un complexe dualisant
sur l'hensélisé, resp.\ l'hensélisé strict de $A$.
\end{enumerate}
\end{lemma}

\begin{proof}
Cela résulte immédiatement des lemmes \ref{dualizing-lemma-flat-unramified} et \ref{dualizing-lemma-flat-iso-mod-I}.
Voir Compléments d'algèbre, sections \ref{more-algebra-section-henselian-pairs}, \ref{more-algebra-section-permanence-completion} et \ref{more-algebra-section-permanence-henselization},
ainsi que Algèbre, sections \ref{algebra-section-completion} et \ref{algebra-section-completion-noetherian},
pour les complétions et les hensélisations.
\end{proof}

\begin{lemma}
\label{dualizing-lemma-ubiquity-dualizing}
Les anneaux des types suivants possèdent un complexe dualisant :
\begin{enumerate}
\item les corps ;
\item les anneaux locaux noethériens complets ;
\item $\mathbf{Z}$ ;
\item les anneaux de Dedekind ;
\item tout anneau obtenu à partir d'un des anneaux précédents
en prenant une algèbre essentiellement de type fini, un complété
pour la topologie définie par un idéal, un hensélisé ou un hensélisé strict.
\end{enumerate}
\end{lemma}

\begin{proof}
L'assertion (5) résulte de la proposition \ref{dualizing-proposition-dualizing-essentially-finite-type} et du lemme \ref{dualizing-lemma-completion-henselization-dualizing}.
D'après le lemme \ref{dualizing-lemma-regular-gorenstein}, un anneau local régulier possède un complexe dualisant.
Un anneau local noethérien complet est quotient d'un anneau local régulier
par le théorème de structure de Cohen (Algèbre, théorème \ref{algebra-theorem-cohen-structure-theorem}).
Soit $A$ un anneau de Dedekind. Tout idéal $I$ est alors un
$A$-module projectif de type fini (cela résulte d'Algèbre, lemme \ref{algebra-lemma-finite-projective},
et du fait que les anneaux locaux de $A$ sont des anneaux
de valuation discrète, donc principaux).
Ainsi tout $A$-module est de dimension injective finie au plus $1$
d'après Compléments d'algèbre, lemme \ref{more-algebra-lemma-injective-amplitude}.
Il en résulte aisément que $A[0]$ est un complexe dualisant.
\end{proof}





\section{Fibres formelles}
\label{dualizing-section-formal-fibres}

\noindent
Cette section poursuit Compléments d'algèbre, section \ref{more-algebra-section-properties-formal-fibres}.
Nous y avons établi une théorie assez satisfaisante des anneaux noethériens
$A$ dont les anneaux locaux ont des fibres formelles de Cohen-Macaulay.
Plus précisément, nous avons prouvé que
(1) il suffit de vérifier que les fibres formelles des localisés
aux idéaux maximaux sont de Cohen-Macaulay,
(2) la propriété se transmet aux anneaux de type fini sur $A$,
(3) les fibres de $A \to A^\wedge$ sont de Cohen-Macaulay pour tout complété $A^\wedge$ de $A$,
et (4) la propriété se transmet aux hensélisés de $A$.
Voir Compléments d'algèbre, lemme \ref{more-algebra-lemma-check-P-ring-maximal-ideals}, proposition \ref{more-algebra-proposition-finite-type-over-P-ring},
lemme \ref{more-algebra-lemma-map-P-ring-to-completion-P} et lemme \ref{more-algebra-lemma-henselization-pair-P-ring}.
Il en va de même pour les anneaux noethériens dont les anneaux locaux
ont des fibres formelles géométriquement réduites, géométriquement normales,
$(S_n)$ ou géométriquement $(R_n)$.
Nous verrons ici que les mêmes résultats valent pour les anneaux noethériens
dont les anneaux locaux ont des fibres formelles de Gorenstein
ou localement d'intersection complète.
Cela concerne ce chapitre, car un anneau noethérien possédant
un complexe dualisant en est un exemple.

\begin{lemma}
\label{dualizing-lemma-formal-fibres-gorenstein}
Les propriétés (A), (B), (C), (D) et (E) de
Compléments d'algèbre, section \ref{more-algebra-section-properties-formal-fibres},
sont satisfaites pour $P(k \to R) =$« $R$ est un anneau de Gorenstein ».
\end{lemma}

\begin{proof}
Puisque le résultat est déjà connu avec « Cohen-Macaulay » à la place
de « Gorenstein », on peut à chaque étape supposer l'anneau de Cohen-Macaulay.
Cela n'aide guère à la démonstration, mais peut être utile psychologiquement.

\medskip\noindent
Assertion (A). Soit $K/k$ une extension de corps de type fini.
Soit $R$ une $k$-algèbre de Gorenstein.
On peut trouver une intersection complète globale
$A = k[x_1, \ldots, x_n]/(f_1, \ldots, f_c)$
sur $k$ telle que $K$ soit isomorphe au corps des fractions de $A$ ;
voir Algèbre, lemme \ref{algebra-lemma-colimit-syntomic}.
Alors $R \to R \otimes_k A$ est une intersection complète globale relative.
Ainsi $R \otimes_k A$ est de Gorenstein d'après le lemme \ref{dualizing-lemma-gorenstein-lci}.
Il en est de même de $R \otimes_k K$, qui en est un localisé.

\medskip\noindent
Preuve de (B). C'est clair, car un anneau est de Gorenstein
si et seulement si tous ses anneaux locaux le sont.

\medskip\noindent
Assertion (C). Soient $A \to B \to C$ des homomorphismes plats d'anneaux noethériens.
Supposons que les fibres de $A \to B$ soient de Gorenstein et que $B \to C$ soit régulier.
Il faut montrer que les fibres de $A \to C$ sont de Gorenstein.
On peut clairement supposer que $A = k$ est un corps,
puis que $B \to C$ est un homomorphisme local régulier d'anneaux locaux noethériens.
Alors $B$ est de Gorenstein et $C/\mathfrak m_B C$ est régulier,
donc en particulier de Gorenstein (lemme \ref{dualizing-lemma-regular-gorenstein}).
Ainsi $C$ est de Gorenstein d'après le lemme \ref{dualizing-lemma-flat-under-gorenstein}.

\medskip\noindent
Assertion (D). Cela résulte du lemme \ref{dualizing-lemma-flat-under-gorenstein}.
L'assertion (E) est immédiate, car la condition ne fait pas intervenir le corps de base.
\end{proof}

\begin{lemma}
\label{dualizing-lemma-dualizing-gorenstein-formal-fibres}
Soit $A$ un anneau local noethérien. Si $A$ possède un complexe dualisant,
les fibres formelles de $A$ sont de Gorenstein.
\end{lemma}

\begin{proof}
Soit $\mathfrak p$ un idéal premier de $A$.
La fibre formelle de $A$ en $\mathfrak p$ est isomorphe à la fibre formelle
de $A/\mathfrak p$ en $(0)$. Le quotient $A/\mathfrak p$ possède un complexe dualisant
(lemme \ref{dualizing-lemma-dualizing-quotient}). Il suffit donc de vérifier l'énoncé lorsque $A$
est un anneau local intègre et $\mathfrak p = (0)$.
Soit $\omega_A^\bullet$ un complexe dualisant pour $A$.
Alors $\omega_A^\bullet \otimes_A A^\wedge$ est un complexe dualisant pour le complété $A^\wedge$
(lemme \ref{dualizing-lemma-flat-unramified}).
De plus, $\omega_A^\bullet \otimes_A K$ est un complexe dualisant pour le corps des fractions $K$
de $A$ (lemme \ref{dualizing-lemma-dualizing-localize}).
Ainsi $\omega_A^\bullet \otimes_A K$ est isomorphe à $K[n]$ pour un certain $n \in \mathbf{Z}$.
De même, on obtient pour la fibre formelle $A^\wedge \otimes_A K$ le complexe dualisant
$$
\omega_A^\bullet \otimes_A A^\wedge \otimes_{A^\wedge} (A^\wedge \otimes_A K) =
(\omega_A^\bullet \otimes_A K) \otimes_K (A^\wedge \otimes_A K) \cong
(A^\wedge \otimes_A K)[n]
$$
comme voulu.
\end{proof}

\noindent
Voici la vérification annoncée dans Algèbre à puissances divisées,
remarque \ref{dpa-remark-no-good-ci-map}.

\begin{lemma}
\label{dualizing-lemma-formal-fibres-lci}
Les propriétés (A), (B), (C), (D) et (E) de
Compléments d'algèbre, section \ref{more-algebra-section-properties-formal-fibres},
sont satisfaites pour $P(k \to R) =$« $R$ est localement une intersection complète ».
Voir Algèbre à puissances divisées, définition \ref{dpa-definition-lci}.
\end{lemma}

\begin{proof}
Assertion (A). Soit $K/k$ une extension de corps de type fini.
Soit $R$ une $k$-algèbre localement d'intersection complète.
On peut trouver une intersection complète globale
$A = k[x_1, \ldots, x_n]/(f_1, \ldots, f_c)$
sur $k$ telle que $K$ soit isomorphe au corps des fractions de $A$ ;
voir Algèbre, lemme \ref{algebra-lemma-colimit-syntomic}.
Alors $R \to R \otimes_k A$ est une intersection complète globale relative.
Il en résulte que $R \otimes_k A$ est localement d'intersection complète
d'après Algèbre à puissances divisées, lemme \ref{dpa-lemma-avramov}.

\medskip\noindent
Preuve de (B). C'est clair, car un anneau est localement d'intersection
complète si et seulement si tous ses anneaux locaux sont des intersections complètes.

\medskip\noindent
Assertion (C). Soient $A \to B \to C$ des homomorphismes plats d'anneaux noethériens.
Supposons que les fibres de $A \to B$ soient localement d'intersection complète
et que $B \to C$ soit régulier. Il faut montrer que les fibres de $A \to C$
sont localement d'intersection complète.
On peut clairement supposer que $A = k$ est un corps,
puis que $B \to C$ est un homomorphisme local régulier d'anneaux locaux noethériens.
Alors $B$ est une intersection complète et $C/\mathfrak m_B C$ est régulier,
donc en particulier une intersection complète par définition.
Ainsi $C$ est une intersection complète d'après
Algèbre à puissances divisées, lemme \ref{dpa-lemma-avramov}.

\medskip\noindent
Assertion (D). Les mêmes arguments que pour (C) s'appliquent,
en utilisant l'autre implication d'Algèbre à puissances divisées,
lemme \ref{dpa-lemma-avramov}.
L'assertion (E) est immédiate, car la condition ne fait pas intervenir le corps de base.
\end{proof}






\section{Construction algébrique de l'image inverse extraordinaire}
\label{dualizing-section-relative-dualizing-complex-algebraic}

\noindent
Pour un homomorphisme de type fini $R \to A$ d'anneaux noethériens,
nous construirons un foncteur $\varphi^! : D(R) \to D(A)$ bien défini à isomorphisme non unique près.
Comme nous le verrons dans Dualité pour les schémas, remarque \ref{duality-remark-local-calculation-shriek},
il coïncide à isomorphisme près avec les foncteurs image inverse
extraordinaire de la théorie de la dualité pour les schémas.
Pour motiver cette construction, signalons deux propriétés supplémentaires :
\begin{enumerate}
\item $\varphi^!$ envoie un complexe dualisant pour $R$, s'il en existe,
sur un complexe dualisant pour $A$ ;
\item $\omega_{A/R}^\bullet = \varphi^!(R)$ est une sorte de complexe dualisant relatif :
il appartient à $D^b_{\textit{Coh}}(A)$ et sa restriction aux fibres est un complexe
dualisant si $R \to A$ est plat.
\end{enumerate}
Ces énoncés sont les lemmes \ref{dualizing-lemma-shriek-dualizing-algebraic} et \ref{dualizing-lemma-relative-dualizing-algebraic}.

\medskip\noindent
Soit $\varphi : R \to A$ un homomorphisme de type fini d'anneaux noethériens.
Définissons un foncteur $\varphi^! : D(R) \to D(A)$ de la manière suivante :
\begin{enumerate}
\item si $\varphi : R \to A$ est surjectif, posons $\varphi^!(K) = R\Hom(A, K)$.
Nous utilisons ici le foncteur $R\Hom(A, -) : D(R) \to D(A)$ de la section \ref{dualizing-section-trivial} ;
\item en général, choisissons une surjection $\psi : P \to A$ avec $P = R[x_1, \ldots, x_n]$
et posons $\varphi^!(K) = \psi^!(K \otimes_R^\mathbf{L} P)[n]$.
Nous utilisons ici le foncteur $- \otimes_R^\mathbf{L} P : D(R) \to D(P)$
de Compléments d'algèbre, section \ref{more-algebra-section-derived-base-change}.
\end{enumerate}
Noter le décalage $[n]$ par le nombre de variables de l'anneau de polynômes.
Cette construction n'est {\bf pas} canonique,
et le foncteur $\varphi^!$ ne sera bien défini qu'à isomorphisme
de foncteurs non unique près\footnote{On peut rendre la construction canonique :
utiliser $\Omega^n_{P/R}[n]$ au lieu de $P[n]$ dans la construction,
puis dans le lemme \ref{dualizing-lemma-well-defined}.
Cela rend toutefois le contenu de cette section beaucoup plus compliqué.}.

\begin{lemma}
\label{dualizing-lemma-well-defined}
Soit $\varphi : R \to A$ un homomorphisme de type fini d'anneaux noethériens.
Le foncteur $\varphi^!$ est bien défini à isomorphisme près.
\end{lemma}

\begin{proof}
Supposons que $\psi_1 : P_1 = R[x_1, \ldots, x_n] \to A$ et $\psi_2 : P_2 = R[y_1, \ldots, y_m] \to A$ soient deux surjections
d'anneaux de polynômes sur $A$. On obtient un diagramme commutatif
$$
\xymatrix{
R[x_1, \ldots, x_n, y_1, \ldots, y_m]
\ar[d]^{x_i \mapsto g_i} \ar[rr]_-{y_j \mapsto f_j} & &
R[x_1, \ldots, x_n] \ar[d] \\
R[y_1, \ldots, y_m] \ar[rr] & & A
}
$$
où $f_j$ et $g_i$ sont choisis de façon que $\psi_1(f_j) = \psi_2(y_j)$ et $\psi_2(g_i) = \psi_1(x_i)$.
Par symétrie, il suffit de prouver que les foncteurs définis à l'aide de
$P \to A$ et de $P[y_1, \ldots, y_m] \to A$ sont isomorphes.
Par récurrence, on peut supposer $m = 1$.
On se ramène ainsi au cas traité dans le paragraphe suivant.

\medskip\noindent
Ici $\psi : P \to A$ est donné et $\chi : P[y] \to A$ induit $\psi$ sur $P$.
Écrivons $Q = P[y]$. Choisissons $g \in P$ tel que $\psi(g) = \chi(y)$.
Notons $\pi : Q \to P$ l'homomorphisme de $P$-algèbres tel que $\pi(y) = g$.
Alors $\chi = \psi \circ \pi$, donc $\chi^! = \psi^! \circ \pi^!$, puisque les deux sont adjoints
au foncteur de restriction $D(A) \to D(Q)$ d'après la section \ref{dualizing-section-trivial}.
Ainsi
$$
\chi^!\left(K \otimes_R^\mathbf{L} Q\right)[n + 1] =
\psi^!\left(\pi^!\left(K \otimes_R^\mathbf{L} Q\right)[1]\right)[n]
$$
Il suffit donc de montrer que $\pi^!(K \otimes_R^\mathbf{L} Q[1]) = K \otimes_R^\mathbf{L} P$.
Il suffit par conséquent de montrer que le foncteur $\pi^!(-) : D(Q) \to D(P)$
est isomorphe à $K \mapsto K \otimes_Q^\mathbf{L} P[-1]$, ce qui résulte du lemme \ref{dualizing-lemma-compute-for-effective-Cartier-algebraic}.
\end{proof}

\begin{lemma}
\label{dualizing-lemma-shriek-boundedness}
Soit $\varphi : R \to A$ un homomorphisme de type fini d'anneaux noethériens.
\begin{enumerate}
\item $\varphi^!$ envoie $D^+(R)$ dans $D^+(A)$ et $D^+_{\textit{Coh}}(R)$ dans $D^+_{\textit{Coh}}(A)$.
\item Si $\varphi$ est parfait, alors $\varphi^!$ envoie
$D^-(R)$ dans $D^-(A)$, $D^-_{\textit{Coh}}(R)$ dans $D^-_{\textit{Coh}}(A)$ et $D^b_{\textit{Coh}}(R)$ dans $D^b_{\textit{Coh}}(A)$.
\end{enumerate}
\end{lemma}

\begin{proof}
Choisissons une factorisation $R \to P \to A$ comme dans la définition de $\varphi^!$.
Le foncteur $- \otimes_R^\mathbf{L} : D(R) \to D(P)$ préserve les sous-catégories $D^+, D^+_{\textit{Coh}}, D^-, D^-_{\textit{Coh}}, D^b_{\textit{Coh}}$.
Le foncteur $R\Hom(A, -) : D(P) \to D(A)$ préserve $D^+$ et $D^+_{\textit{Coh}}$ d'après le lemme \ref{dualizing-lemma-exact-support-coherent}.
Si $R \to A$ est parfait, alors $A$ est parfait comme $P$-module ;
voir Compléments d'algèbre, lemme \ref{more-algebra-lemma-perfect-ring-map}.
Rappelons que la restriction de $R\Hom(A, K)$ à $D(P)$ est $R\Hom_P(A, K)$.
D'après Compléments d'algèbre, lemme \ref{more-algebra-lemma-dual-perfect-complex}, on a $R\Hom_P(A, K) = E \otimes_P^\mathbf{L} K$
pour un certain objet parfait $E \in D(P)$.
Puisque $E$ peut être représenté par un complexe fini de $P$-modules
projectifs de type fini, il est clair que $R\Hom_P(A, K)$ appartient à $D^-(P), D^-_{\textit{Coh}}(P), D^b_{\textit{Coh}}(P)$
dès que $K$ y appartient.
Puisque le foncteur de restriction $D(A) \to D(P)$ reflète l'appartenance
à ces sous-catégories, la démonstration est achevée.
\end{proof}

\begin{lemma}
\label{dualizing-lemma-shriek-dualizing-algebraic}
Soit $\varphi$ un homomorphisme de type fini d'anneaux noethériens.
Si $\omega_R^\bullet$ est un complexe dualisant pour $R$,
alors $\varphi^!(\omega_R^\bullet)$ est un complexe dualisant pour $A$.
\end{lemma}

\begin{proof}
Cela résulte des lemmes \ref{dualizing-lemma-dualizing-polynomial-ring} et \ref{dualizing-lemma-dualizing-quotient}.
\end{proof}

\begin{lemma}
\label{dualizing-lemma-flat-bc}
Soit $R \to R'$ un homomorphisme plat d'anneaux noethériens.
Soit $\varphi : R \to A$ un homomorphisme de type fini.
Soit $\varphi' : R' \to A' = A \otimes_R R'$ l'homomorphisme induit par $\varphi$.
On a alors des morphismes fonctoriels
$$
\varphi^!(K) \otimes_A^\mathbf{L} A' \longrightarrow
(\varphi')^!(K \otimes_R^\mathbf{L} R')
$$
pour $K$ dans $D(R)$, qui sont des isomorphismes pour $K \in D^+(R)$.
\end{lemma}

\begin{proof}
Choisissons une factorisation $R \to P \to A$ où $P$ est un anneau de polynômes sur $R$.
Elle donne par changement de base une factorisation correspondante $R' \to P' \to A'$.
Puisque l'on a $(K \otimes_R^\mathbf{L} P) \otimes_P^\mathbf{L} P' =
(K \otimes_R^\mathbf{L} R') \otimes_{R'}^\mathbf{L} P'$ d'après Compléments d'algèbre, lemme \ref{more-algebra-lemma-double-base-change},
il suffit de construire des morphismes
$$
R\Hom(A, K \otimes_R^\mathbf{L} P[n]) \otimes_A^\mathbf{L} A'
\longrightarrow
R\Hom(A', (K \otimes_R^\mathbf{L} P[n]) \otimes_P^\mathbf{L} P')
$$
fonctoriels en $K$. Pour cela, utilisons le morphisme (\ref{dualizing-equation-base-change})
construit dans la section \ref{dualizing-section-base-change-trivial-duality} pour $P, A, P', A'$.
Ce morphisme est un isomorphisme pour $K \in D^+(R)$ d'après le lemme \ref{dualizing-lemma-flat-bc-surjection}.
\end{proof}

\begin{lemma}
\label{dualizing-lemma-bc}
Soit $R \to R'$ un homomorphisme d'anneaux noethériens.
Soit $\varphi : R \to A$ un homomorphisme parfait
(Compléments d'algèbre, définition \ref{more-algebra-definition-pseudo-coherent-perfect})
tel que $R'$ et $A$ soient Tor-indépendants sur $R$.
Soit $\varphi' : R' \to A' = A \otimes_R R'$ l'homomorphisme induit par $\varphi$.
On a alors un isomorphisme fonctoriel
$$
\varphi^!(K) \otimes_A^\mathbf{L} A' =
(\varphi')^!(K \otimes_R^\mathbf{L} R')
$$
pour $K$ dans $D(R)$.
\end{lemma}

\begin{proof}
On peut choisir une factorisation $R \to P \to A$ où $P$ est un anneau de polynômes
sur $R$ tel que $A$ soit un $P$-module parfait ;
voir Compléments d'algèbre, lemme \ref{more-algebra-lemma-perfect-ring-map}.
Elle donne par changement de base une factorisation correspondante $R' \to P' \to A'$.
Puisque l'on a $(K \otimes_R^\mathbf{L} P) \otimes_P^\mathbf{L} P' =
(K \otimes_R^\mathbf{L} R') \otimes_{R'}^\mathbf{L} P'$ d'après Compléments d'algèbre, lemme \ref{more-algebra-lemma-double-base-change},
il suffit de construire des morphismes
$$
R\Hom(A, K \otimes_R^\mathbf{L} P[n]) \otimes_A^\mathbf{L} A'
\longrightarrow
R\Hom(A', (K \otimes_R^\mathbf{L} P[n]) \otimes_P^\mathbf{L} P')
$$
fonctoriels en $K$. On a
$$
A \otimes_P^\mathbf{L} P' = A \otimes_R^\mathbf{L} R' = A'
$$
La première égalité résulte de Compléments d'algèbre, lemme \ref{more-algebra-lemma-base-change-comparison},
appliqué à $R, R', P, P'$. La seconde vient de ce que $A$ et $R'$
sont Tor-indépendants sur $R$.
Ainsi $A$ et $P'$ sont Tor-indépendants sur $P$,
et l'on peut utiliser le morphisme (\ref{dualizing-equation-base-change}) construit dans la section \ref{dualizing-section-base-change-trivial-duality}
pour $P, A, P', A'$ afin d'obtenir la flèche cherchée.
D'après le lemme \ref{dualizing-lemma-bc-surjection}, il suffit pour conclure de prouver que $A$
est un $P$-module parfait, ce qui a été vu ci-dessus.
\end{proof}

\begin{lemma}
\label{dualizing-lemma-bc-flat}
Soit $R \to R'$ un homomorphisme d'anneaux noethériens.
Soit $\varphi : R \to A$ plat et de type fini.
Soit $\varphi' : R' \to A' = A \otimes_R R'$ l'homomorphisme induit par $\varphi$.
On a alors un isomorphisme fonctoriel
$$
\varphi^!(K) \otimes_A^\mathbf{L} A' =
(\varphi')^!(K \otimes_R^\mathbf{L} R')
$$
pour $K$ dans $D(R)$.
\end{lemma}

\begin{proof}
C'est un cas particulier du lemme \ref{dualizing-lemma-bc} d'après
Compléments d'algèbre, lemme \ref{more-algebra-lemma-flat-finite-presentation-perfect}.
\end{proof}

\begin{lemma}
\label{dualizing-lemma-composition-shriek-algebraic}
Soient $A \xrightarrow{a} B \xrightarrow{b} C$ des homomorphismes de type fini d'anneaux noethériens.
Il existe alors une transformation de foncteurs $b^! \circ a^! \to (b \circ a)^!$
qui est un isomorphisme sur $D^+(A)$.
\end{lemma}

\begin{proof}
Choisissons un anneau de polynômes $P = A[x_1, \ldots, x_n]$ sur $A$
et une surjection $P \to B$. Choisissons des éléments $c_1, \ldots, c_m \in C$
engendrant $C$ sur $B$. Posons $Q = P[y_1, \ldots, y_m]$ et notons $Q' = Q \otimes_P B = B[y_1, \ldots, y_m]$.
Soit $\chi : Q' \to C$ la surjection envoyant $y_j$ sur $c_j$.
On a le diagramme
$$
\xymatrix{
& Q \ar[r]_{\psi'} & Q' \ar[r]_\chi & C \\
A \ar[r] & P \ar[r]^\psi \ar[u] & B \ar[u]
}
$$
D'après le lemme \ref{dualizing-lemma-flat-bc-surjection}, pour $M \in D(P)$, on a une flèche $\psi^!(M) \otimes_B^\mathbf{L} Q' \to (\psi')^!(M \otimes_P^\mathbf{L} Q)$
qui est un isomorphisme dès que $M$ est borné inférieurement.
De plus, $\chi^! \circ (\psi')^! = (\chi \circ \psi')^!$ puisque les deux foncteurs sont adjoints
au foncteur de restriction $D(C) \to D(Q)$ d'après la section \ref{dualizing-section-trivial}.
On obtient alors
\begin{align*}
b^!(a^!(K))
& =
\chi^!(\psi^!(K \otimes_A^\mathbf{L} P)[n] \otimes_B^\mathbf{L} Q)[m] \\
& \to
\chi^!((\psi')^!(K \otimes_A^\mathbf{L} P \otimes_P^\mathbf{L} Q))[n + m] \\
& =
(\chi \circ \psi')^!(K\otimes_A^\mathbf{L} Q)[n + m] \\
& =
(b \circ a)^!(K)
\end{align*}
en utilisant, outre ce qui précède, Compléments d'algèbre, lemme \ref{more-algebra-lemma-double-base-change}.
\end{proof}

\begin{lemma}
\label{dualizing-lemma-upper-shriek-finite}
Soit $\varphi : R \to A$ un homomorphisme fini d'anneaux noethériens.
Alors $\varphi^!$ est isomorphe au foncteur $R\Hom(A, -) : D(R) \to D(A)$ de la section \ref{dualizing-section-trivial}.
\end{lemma}

\begin{proof}
Supposons $A$ engendré par $n > 1$ éléments sur $R$.
On peut factoriser $R \to A$ en deux homomorphismes finis d'anneaux,
avec à chaque étape un nombre de générateurs $< n$.
Les lemmes \ref{dualizing-lemma-composition-shriek-algebraic} et \ref{dualizing-lemma-composition-right-adjoints} montrent donc qu'il suffit de prouver
le lemme lorsque $A$ est engendré par un seul élément sur $R$.
Puisque $A$ est fini sur $R$, $A$ est quotient de $B = R[x]/(f)$,
où $f$ est un polynôme unitaire en $x$ (Algèbre, lemme \ref{algebra-lemma-finite-is-integral}).
En utilisant encore les lemmes de composition et l'accord des foncteurs
pour les surjections par définition, on se ramène à $R \to B = R[x]/(f)$.
Dans ce cas, le foncteur $\varphi^!$ est isomorphe à $K \mapsto K \otimes_R^\mathbf{L} B$ :
cela se prouve en utilisant le lemme \ref{dualizing-lemma-compute-for-effective-Cartier-algebraic} pour l'homomorphisme $R[x] \to B$
(noter que les décalages dans la définition de $\varphi^!$ et dans le lemme ont somme nulle).
Pour le foncteur $R\Hom(B, -) : D(R) \to D(B)$, le lemme \ref{dualizing-lemma-RHom-is-tensor-special} permet de se ramener
à montrer que $\Hom_R(B, R) \cong B$ comme $B$-modules.
Supposons $f$ de degré $d$.
Alors une $R$-base de $B$ est donnée par $1, x, \ldots, x^{d - 1}$.
Soit $\delta_i : B \to R$, pour $i = 0, \ldots, d - 1$, l'application $R$-linéaire
qui extrait le coefficient de $x^i$ dans cette base.
Alors $\delta_0, \ldots, \delta_{d - 1}$ est une base de $\Hom_R(B, R)$.
Enfin, pour $0 \leq i \leq d - 1$, un calcul donne
$$
x^i \delta_{d - 1} =
\delta_{d - 1 - i} + b_1 \delta_{d - i} + \ldots + b_i \delta_{d - 1}
$$
pour certains $c_1, \ldots, c_d \in R$\footnote{Si $f = x^d + a_1 x^{d - 1} + \ldots + a_d$, alors
$c_1 = -a_1$, $c_2 = a_1^2 - a_2$, $c_3 = -a_1^3 + 2a_1a_2 -a_3$, etc.}.
Ainsi $\Hom_R(B, R)$ est un $B$-module monogène de générateur $\delta_{d - 1}$.
En examinant les rangs, on en déduit que c'est un module libre
de rang $1$ sur $B$.
\end{proof}

\begin{lemma}
\label{dualizing-lemma-upper-shriek-localize}
Soit $R$ un anneau noethérien et soit $f \in R$.
Si $\varphi$ désigne l'homomorphisme $R \to R_f$, alors $\varphi^!$ est isomorphe à $- \otimes_R^\mathbf{L} R_f$.
Plus généralement, si $\varphi : R \to R'$ est un homomorphisme tel que
$\Spec(R') \to \Spec(R)$ soit une immersion ouverte, alors $\varphi^!$ est isomorphe à $- \otimes_R^\mathbf{L} R'$.
\end{lemma}

\begin{proof}
Choisissons la présentation $R \to R[x] \to R[x]/(fx - 1) = R_f$ et observons que $fx - 1$
est un non-diviseur de zéro dans $R[x]$.
On peut donc appliquer le lemme \ref{dualizing-lemma-compute-for-effective-Cartier-algebraic} pour calculer le foncteur $\varphi^!$.
Nous omettons les détails ; noter que les décalages dans la définition
de $\varphi^!$ et dans le lemme ont somme nulle.

\medskip\noindent
Dans le cas général, noter que $R' \otimes_R R' = R'$.
Le résultat découle donc des résultats de changement de base ci-dessus :
on peut utiliser soit le lemme \ref{dualizing-lemma-flat-bc}, soit le lemme \ref{dualizing-lemma-bc}.
\end{proof}

\begin{lemma}
\label{dualizing-lemma-upper-shriek-is-tensor-functor}
Soit $\varphi : R \to A$ un homomorphisme parfait d'anneaux noethériens
(par exemple, $\varphi$ est plat de type fini).
Alors $\varphi^!(K) = K \otimes_R^\mathbf{L} \varphi^!(R)$ pour $K \in D(R)$.
\end{lemma}

\begin{proof}
(L'assertion entre parenthèses résulte de Compléments d'algèbre, lemme \ref{more-algebra-lemma-flat-finite-presentation-perfect}.)
On peut choisir une factorisation $R \to P \to A$ où $P$ est un anneau de polynômes
en $n$ variables sur $R$ ; alors $A$ est un $P$-module parfait,
voir Compléments d'algèbre, lemme \ref{more-algebra-lemma-perfect-ring-map}.
Rappelons que $\varphi^!(K) = R\Hom(A, K \otimes_R^\mathbf{L} P[n])$.
Le résultat découle donc du lemme \ref{dualizing-lemma-RHom-is-tensor-special}
et de Compléments d'algèbre, lemme \ref{more-algebra-lemma-double-base-change}.
\end{proof}

\begin{lemma}
\label{dualizing-lemma-relative-dualizing-if-have-omega}
Soit $\varphi : A \to B$ un homomorphisme de type fini d'anneaux noethériens.
Soit $\omega_A^\bullet$ un complexe dualisant pour $A$. Posons $\omega_B^\bullet = \varphi^!(\omega_A^\bullet)$.
Notons $D_A(K) = R\Hom_A(K, \omega_A^\bullet)$ pour $K \in D_{\textit{Coh}}(A)$ et $D_B(L) = R\Hom_B(L, \omega_B^\bullet)$ pour $L \in D_{\textit{Coh}}(B)$.
Il existe alors un isomorphisme fonctoriel
$$
\varphi^!(K) = D_B(D_A(K) \otimes_A^\mathbf{L} B)
$$
pour $K \in D_{\textit{Coh}}(A)$.
\end{lemma}

\begin{proof}
Observons que $\omega_B^\bullet$ est un complexe dualisant pour $B$ d'après le lemme \ref{dualizing-lemma-shriek-dualizing-algebraic}.
Soient $A \to B \to C$ des homomorphismes de type fini d'anneaux noethériens.
Si le lemme vaut pour $A \to B$ et $B \to C$, il vaut pour $A \to C$.
Cela résulte du lemme \ref{dualizing-lemma-composition-shriek-algebraic} et du fait que $D_B \circ D_B \cong \text{id}$ d'après le lemme \ref{dualizing-lemma-dualizing}.
Il suffit donc de prouver le lemme lorsque $A \to B$ est une surjection
et lorsque $B$ est un anneau de polynômes sur $A$.

\medskip\noindent
Supposons $B = A[x_1, \ldots, x_n]$. Puisque $D_A \circ D_A \cong \text{id}$, il suffit de prouver
$D_B(K \otimes_A B) \cong D_A(K) \otimes_A B[n]$ pour $K$ dans $D_{\textit{Coh}}(A)$.
Choisissons un complexe borné $I^\bullet$ d'injectifs représentant $\omega_A^\bullet$.
Choisissons un quasi-isomorphisme $I^\bullet \otimes_A B \to J^\bullet$ où $J^\bullet$
est un complexe borné de $B$-modules.
Étant donné un complexe $K^\bullet$ de $A$-modules,
considérons le morphisme évident de complexes
$$
\Hom^\bullet(K^\bullet, I^\bullet) \otimes_A B[n]
\longrightarrow
\Hom^\bullet(K^\bullet \otimes_A B, J^\bullet[n])
$$
Le membre de gauche représente $D_A(K) \otimes_A B[n]$ et celui de droite représente $D_B(K \otimes_A B)$.
Il suffit donc de prouver que ce morphisme est un quasi-isomorphisme
si les modules de cohomologie de $K^\bullet$ sont des $A$-modules de type fini.
Observons que la cohomologie du complexe en degré $r$,
de chaque côté, ne dépend que d'un nombre fini des $K^i$.
On peut donc remplacer $K^\bullet$ par une troncature, c'est-à-dire supposer
que $K^\bullet$ représente un objet de $D^-_{\textit{Coh}}(A)$.
Alors $K^\bullet$ est quasi-isomorphe à un complexe borné supérieurement
de $A$-modules libres de type fini.
On peut donc supposer que $K^\bullet$ est un complexe borné supérieurement
de $A$-modules libres de type fini.
Il est alors facile de voir que le morphisme affiché est
un isomorphisme de complexes, ce qui conclut dans ce cas.

\medskip\noindent
Supposons $A \to B$ surjectif. Notons $i_* : D(B) \to D(A)$ le foncteur de restriction
et rappelons que $\varphi^!(-) = R\Hom(A, -)$ est adjoint à droite de $i_*$ (lemme \ref{dualizing-lemma-right-adjoint}).
Pour $F \in D(B)$, on a
\begin{align*}
\Hom_B(F, D_B(D_A(K) \otimes_A^\mathbf{L} B))
& =
\Hom_B((D_A(K) \otimes_A^\mathbf{L} B) \otimes_B^\mathbf{L} F,
\omega_B^\bullet) \\
& =
\Hom_A(D_A(K) \otimes_A^\mathbf{L} i_*F, \omega_A^\bullet) \\
& =
\Hom_A(i_*F, D_A(D_A(K))) \\
& =
\Hom_A(i_*F, K) \\
& =
\Hom_B(F, \varphi^!(K))
\end{align*}
La première égalité résulte de Compléments d'algèbre, lemme \ref{more-algebra-lemma-internal-hom},
et de la définition de $D_B$.
La deuxième résulte de l'adjonction mentionnée ci-dessus et de l'égalité
$i_*((D_A(K) \otimes_A^\mathbf{L} B) \otimes_B^\mathbf{L} F) =
D_A(K) \otimes_A^\mathbf{L} i_*F$
(Compléments d'algèbre, lemme \ref{more-algebra-lemma-derived-base-change}).
La troisième résulte de Compléments d'algèbre, lemme \ref{more-algebra-lemma-internal-hom}.
La quatrième vient de ce que $D_A \circ D_A = \text{id}$, et la dernière à nouveau de l'adjonction.
Le résultat découle donc du lemme de Yoneda.
\end{proof}






\section{Complexes dualisants relatifs dans le cas noethérien}
\label{dualizing-section-relative-dualizing-complexes-Noetherian}

\noindent
Soit $\varphi : R \to A$ un homomorphisme de type fini d'anneaux noethériens.
On définit le {\it complexe dualisant relatif de $A$ sur $R$}
comme l'objet
$$
\omega_{A/R}^\bullet = \varphi^!(R)
$$
de $D(A)$. Ici $\varphi^!$ est celui de la section \ref{dualizing-section-relative-dualizing-complex-algebraic}.
Les résultats de cette section montrent que $\omega_{A/R}^\bullet$
est bien défini à isomorphisme non unique près.

\begin{lemma}
\label{dualizing-lemma-base-change-relative-algebraic}
Soit $R \to R'$ un homomorphisme d'anneaux noethériens.
Soit $R \to A$ de type fini. Posons $A' = A \otimes_R R'$. Si
\begin{enumerate}
\item $R \to R'$ est plat, ou
\item $R \to A$ est plat, ou
\item $R \to A$ est parfait et $R'$ et $A$ sont Tor-indépendants sur $R$,
\end{enumerate}
alors il existe un isomorphisme
$\omega_{A/R}^\bullet \otimes_A^\mathbf{L} A' \to \omega^\bullet_{A'/R'}$
dans $D(A')$.
\end{lemma}

\begin{proof}
Cela résulte des lemmes \ref{dualizing-lemma-flat-bc}, \ref{dualizing-lemma-bc-flat} et \ref{dualizing-lemma-bc} et des définitions.
\end{proof}

\begin{lemma}
\label{dualizing-lemma-relative-dualizing-algebraic}
Soit $\varphi : R \to A$ un homomorphisme de type fini d'anneaux noethériens.
\begin{enumerate}
\item Si $\varphi$ est parfait, alors
\begin{enumerate}
\item $\omega_{A/R}^\bullet$ appartient à $D^b_{\textit{Coh}}(A)$ ;
\item $\omega_{A/R}^\bullet$ est de Tor-dimension finie sur $R$ ;
\item $A \to R\Hom_A(\omega_{A/R}^\bullet, \omega_{A/R}^\bullet)$ est un isomorphisme.
\end{enumerate}
\item Si $\varphi$ est plat, alors (1)(a), (1)(b) et (1)(c) sont satisfaites, et
\begin{enumerate}
\item $\omega_{A/R}^\bullet$ est $R$-parfait
(Compléments d'algèbre, définition \ref{more-algebra-definition-relatively-perfect}) ;
\item pour tout homomorphisme $R \to k$ vers un corps, le changement de base
$\omega_{A/R}^\bullet \otimes_A^\mathbf{L} (A \otimes_R k)$
est un complexe dualisant pour $A \otimes_R k$.
\end{enumerate}
\end{enumerate}
\end{lemma}

\begin{proof}
Supposons $R \to A$ parfait. Choisissons $R \to P \to A$ comme dans la définition de $\varphi^!$.
Alors $A$ est parfait comme $P$-module
(Compléments d'algèbre, lemme \ref{more-algebra-lemma-perfect-ring-map}).
Cela montre que $\omega_{A/R}^\bullet$ appartient à $D^b_{\textit{Coh}}(A)$ d'après le lemme \ref{dualizing-lemma-shriek-boundedness}, d'où (1)(a).
Pour montrer que $\omega_{A/R}^\bullet$ est de Tor-dimension finie comme complexe de $R$-modules,
observons que $\omega_{A/R}^\bullet = \varphi^!(R) = R\Hom(A, P)[n]$ a pour image $R\Hom_P(A, P)[n]$ dans $D(P)$.
Cet objet est parfait dans $D(P)$ (Compléments d'algèbre, lemme \ref{more-algebra-lemma-dual-perfect-complex}),
donc de Tor-dimension finie dans $D(R)$ puisque $R \to P$ est plat.
Cela prouve (1)(b). L'objet $R\Hom_A(\omega_{A/R}^\bullet, \omega_{A/R}^\bullet)$ de $D(A)$ a pour image dans $D(P)$
\begin{align*}
R\Hom_P(\omega_{A/R}^\bullet, R\Hom(A, P)[n])
& =
R\Hom_P(R\Hom_P(A, P)[n], P)[n] \\
& =
R\Hom_P(R\Hom_P(A, P), P)
\end{align*}
Cet objet est égal à $A$ d'après Compléments d'algèbre, lemme \ref{more-algebra-lemma-dual-perfect-complex},
déjà utilisé. Cela prouve (1)(c).

\medskip\noindent
Supposons $\varphi$ plat. Alors $R \to A$ est un homomorphisme parfait
(Compléments d'algèbre, lemme \ref{more-algebra-lemma-flat-finite-presentation-perfect}),
et (1)(a), (1)(b) et (1)(c) sont satisfaites.
Bien entendu, $\omega_{A/R}^\bullet$ est alors $R$-parfait d'après (1)(a), (1)(b) et les définitions.
Soit $R \to k$ comme dans (2)(b).
D'après le lemme \ref{dualizing-lemma-base-change-relative-algebraic}, il existe un isomorphisme
$$
\omega_{A/R}^\bullet \otimes_A^\mathbf{L} (A \otimes_R k) \cong
\omega^\bullet_{A \otimes_R k/k}
$$
et le membre de droite est un complexe dualisant d'après le lemme \ref{dualizing-lemma-shriek-dualizing-algebraic}.
Cela achève la démonstration.
\end{proof}

\begin{lemma}
\label{dualizing-lemma-base-change-dualizing-over-field}
Soit $K/k$ une extension de corps. Soit $A$ une $k$-algèbre de type fini.
Soit $A_K = A \otimes_k K$. Si $\omega_A^\bullet$ est un complexe dualisant pour $A$,
alors $\omega_A^\bullet \otimes_A A_K$ est un complexe dualisant pour $A_K$.
\end{lemma}

\begin{proof}
Par unicité des complexes dualisants, peu importe celui que l'on choisit
pour $A$ ; nous omettons la preuve détaillée.
Notons $\varphi : k \to A$ la structure d'algèbre.
On peut prendre $\omega_A^\bullet = \varphi^!(k[0])$ d'après le lemme \ref{dualizing-lemma-shriek-dualizing-algebraic}.
On conclut par le lemme \ref{dualizing-lemma-relative-dualizing-algebraic}.
\end{proof}

\begin{lemma}
\label{dualizing-lemma-lci-shriek}
Soit $\varphi : R \to A$ un homomorphisme d'anneaux noethériens localement d'intersection complète.
Alors $\omega_{A/R}^\bullet$ est un objet inversible de $D(A)$ et $\varphi^!(K) = K \otimes_R^\mathbf{L} \omega_{A/R}^\bullet$ pour tout $K \in D(R)$.
\end{lemma}

\begin{proof}
Rappelons qu'un homomorphisme localement d'intersection complète
est parfait d'après Compléments d'algèbre, lemme \ref{more-algebra-lemma-lci-perfect}.
La dernière assertion résulte donc du lemme \ref{dualizing-lemma-upper-shriek-is-tensor-functor}.
D'après Compléments d'algèbre, définition \ref{more-algebra-definition-local-complete-intersection},
on peut écrire $A = R[x_1, \ldots, x_n]/I$, où $I$ est un idéal Koszul-régulier.
La construction de $\varphi^!$ dans la section \ref{dualizing-section-relative-dualizing-complex-algebraic}
montre qu'il suffit de prouver le lemme lorsque $A = R/I$,
où $I \subset R$ est un idéal Koszul-régulier.
Vérifier que $\omega_{A/R}^\bullet$ est inversible dans $D(A)$ est local sur $\Spec(A)$
d'après Compléments d'algèbre, lemme \ref{more-algebra-lemma-invertible-derived}.
De plus, la formation de $\omega_{A/R}^\bullet$ commute à la localisation sur $R$
d'après le lemme \ref{dualizing-lemma-flat-bc}.
En combinant Compléments d'algèbre, définition \ref{more-algebra-definition-regular-ideal} et lemme \ref{more-algebra-lemma-noetherian-finite-all-equivalent},
avec Algèbre, lemme \ref{algebra-lemma-regular-sequence-in-neighbourhood}, on trouve $g_1, \ldots, g_r \in R$
engendrant l'idéal unité dans $A$ tels que $I_{g_j} \subset R_{g_j}$
soit engendré par une suite régulière.
On peut donc supposer $A = R/(f_1, \ldots, f_c)$, où $f_1, \ldots, f_c$ est une suite régulière dans $R$.
Considérons alors les homomorphismes d'anneaux
$$
R \to R/(f_1) \to R/(f_1, f_2) \to \ldots \to R/(f_1, \ldots, f_c) = A
$$
et utilisons le lemme \ref{dualizing-lemma-composition-shriek-algebraic} (ainsi que la dernière assertion déjà prouvée)
pour nous ramener à chaque étape.
Enfin, lorsque $A = R/(f)$ pour un non-diviseur de zéro $f$,
le lemme est vrai puisque $\varphi^!(R) = R\Hom(A, R)$ est inversible d'après le lemme \ref{dualizing-lemma-compute-for-effective-Cartier-algebraic}.
\end{proof}

\begin{lemma}
\label{dualizing-lemma-gorenstein-shriek}
Soit $\varphi : R \to A$ un homomorphisme plat de type fini d'anneaux noethériens.
Les conditions suivantes sont équivalentes :
\begin{enumerate}
\item les fibres $A \otimes_R \kappa(\mathfrak p)$ sont de Gorenstein pour tous les idéaux premiers $\mathfrak p \subset R$ ;
\item $\omega_{A/R}^\bullet$ est un objet inversible de $D(A)$,
voir Compléments d'algèbre, lemme \ref{more-algebra-lemma-invertible-derived}.
\end{enumerate}
\end{lemma}

\begin{proof}
Si (2) est satisfaite, les anneaux des fibres $A \otimes_R \kappa(\mathfrak p)$
ont des complexes dualisants inversibles, donc sont de Gorenstein.
Voir les lemmes \ref{dualizing-lemma-relative-dualizing-algebraic} et \ref{dualizing-lemma-gorenstein}.

\medskip\noindent
Réciproquement, supposons (1).
Observons que $\omega_{A/R}^\bullet$ appartient à $D^b_{\textit{Coh}}(A)$ d'après le lemme \ref{dualizing-lemma-shriek-boundedness}
(puisque les homomorphismes plats de type fini d'anneaux noethériens
sont parfaits, voir Compléments d'algèbre, lemme \ref{more-algebra-lemma-flat-finite-presentation-perfect}).
Prenons un idéal premier $\mathfrak q \subset A$ au-dessus de $\mathfrak p \subset R$. Alors
$$
\omega_{A/R}^\bullet \otimes_A^\mathbf{L} \kappa(\mathfrak q) =
\omega_{A/R}^\bullet \otimes_A^\mathbf{L}
(A \otimes_R \kappa(\mathfrak p))
\otimes_{(A \otimes_R \kappa(\mathfrak p))}^\mathbf{L}
\kappa(\mathfrak q)
$$
Les lemmes \ref{dualizing-lemma-relative-dualizing-algebraic} et \ref{dualizing-lemma-gorenstein} et l'hypothèse (1) montrent que ce complexe
a $1$ groupe de cohomologie non nul,
qui est un espace vectoriel de dimension $1$ sur $\kappa(\mathfrak q)$.
D'après Compléments d'algèbre, lemme \ref{more-algebra-lemma-lift-bounded-pseudo-coherent-to-perfect},
on en déduit que $(\omega_{A/R}^\bullet)_f$ est un objet inversible de $D(A_f)$
pour un certain $f \in A$ tel que $f \not \in \mathfrak q$.
Cela prouve (2).
\end{proof}

\noindent
Le lemme suivant permet de comprendre comment les fonctions de dimension
se transforment lorsqu'on passe à une algèbre de type fini sur un anneau noethérien.

\begin{lemma}
\label{dualizing-lemma-shriek-normalized}
Soit $\varphi : R \to A$ un homomorphisme de type fini d'anneaux noethériens.
Supposons $R$ local et soit $\mathfrak m \subset A$ un idéal maximal au-dessus
de l'idéal maximal de $R$.
Si $\omega_R^\bullet$ est un complexe dualisant normalisé pour $R$,
alors $\varphi^!(\omega_R^\bullet)_\mathfrak m$ est un complexe dualisant normalisé pour $A_\mathfrak m$.
\end{lemma}

\begin{proof}
Nous savons déjà que $\varphi^!(\omega_R^\bullet)$ est un complexe dualisant pour $A$ ; voir le lemme \ref{dualizing-lemma-shriek-dualizing-algebraic}.
Choisissons une factorisation $R \to P \to A$ avec $P = R[x_1, \ldots, x_n]$ comme dans la construction de $\varphi^!$.
Si l'on prouve le lemme pour $R \to P$ et l'idéal maximal $\mathfrak m'$ de $P$
correspondant à $\mathfrak m$, on obtient le résultat pour $R \to A$
en appliquant le lemme \ref{dualizing-lemma-normalized-quotient} à $P_{\mathfrak m'} \to A_\mathfrak m$,
ou le lemme \ref{dualizing-lemma-quotient-function} à $P \to A$.
Dans le cas $A = R[x_1, \ldots, x_n]$, on a $\dim(A_\mathfrak m) = \dim(R) + n$, par exemple d'après Algèbre, lemme \ref{algebra-lemma-dimension-base-fibre-equals-total}
(combiné avec Algèbre, lemme \ref{algebra-lemma-dim-affine-space} pour calculer la dimension de la fibre).
La normalisation de $\omega_R^\bullet$ signifie que $i = -\dim(R)$ est le plus petit indice
tel que $H^i(\omega_R^\bullet)$ soit non nul (cela résulte des lemmes \ref{dualizing-lemma-sitting-in-degrees} et \ref{dualizing-lemma-nonvanishing-generically-local}).
Alors $\varphi^!(\omega_R^\bullet)_\mathfrak m =
\omega_R^\bullet \otimes_R A_\mathfrak m[n]$ a son premier module de cohomologie non nul en degré $-\dim(R) - n$
et est donc le complexe dualisant normalisé pour $A_\mathfrak m$.
\end{proof}

\begin{lemma}
\label{dualizing-lemma-relative-dualizing-trivial-vanishing}
Soit $R \to A$ un homomorphisme de type fini d'anneaux noethériens.
Soit $\mathfrak q \subset A$ un idéal premier au-dessus de $\mathfrak p \subset R$. Alors
$$
H^i(\omega_{A/R}^\bullet)_\mathfrak q \not = 0
\Rightarrow - d \leq i
$$
où $d$ est la dimension de la fibre de $\Spec(A) \to \Spec(R)$ au-dessus de $\mathfrak p$ au point $\mathfrak q$.
\end{lemma}

\begin{proof}
Choisissons une factorisation $R \to P \to A$ avec $P = R[x_1, \ldots, x_n]$ comme dans la section \ref{dualizing-section-relative-dualizing-complex-algebraic},
de sorte que $\omega_{A/R}^\bullet = R\Hom(A, P)[n]$. Il faut montrer que $R\Hom(A, P)_\mathfrak q$
a une cohomologie nulle en degrés $< n - d$.
D'après le lemme \ref{dualizing-lemma-RHom-ext}, cela revient à montrer que $\Ext_P^i(P/I, P)_{\mathfrak r} = 0$ pour $i < n - d$,
où $\mathfrak r \subset P$ est l'idéal premier correspondant à $\mathfrak q$
et $I$ le noyau de $P \to A$.
On peut récrire le groupe considéré sous la forme $\Ext_{P_\mathfrak r}^i(P_\mathfrak r/IP_\mathfrak r, P_\mathfrak r)$
d'après Compléments d'algèbre, lemme \ref{more-algebra-lemma-pseudo-coherence-and-base-change-ext}.
Il faut donc montrer
$$
\text{profondeur}_{IP_\mathfrak r}(P_\mathfrak r) \geq n - d
$$
d'après le lemme \ref{dualizing-lemma-depth}. Le lemme \ref{dualizing-lemma-depth-flat-CM} donne
$$
\text{profondeur}_{IP_\mathfrak r}(P_\mathfrak r) \geq
\dim((P \otimes_R \kappa(\mathfrak p))_\mathfrak r) -
\dim((P/I \otimes_R \kappa(\mathfrak p))_\mathfrak r)
$$
Les membres de droite des deux inégalités coïncident
d'après Algèbre, lemme \ref{algebra-lemma-codimension}.
\end{proof}

\begin{lemma}
\label{dualizing-lemma-relative-dualizing-flat-vanishing}
Soit $R \to A$ un homomorphisme plat de type fini d'anneaux noethériens.
Soit $\mathfrak q \subset A$ un idéal premier au-dessus de $\mathfrak p \subset R$. Alors
$$
H^i(\omega_{A/R}^\bullet)_\mathfrak q \not = 0
\Rightarrow - d \leq i \leq 0
$$
où $d$ est la dimension de la fibre de $\Spec(A) \to \Spec(R)$ au-dessus de $\mathfrak p$ au point $\mathfrak q$.
Si toutes les fibres de $\Spec(A) \to \Spec(R)$ sont de dimension $\leq d$,
alors $\omega_{A/R}^\bullet$ a son amplitude de Tor dans $[-d, 0]$
en tant que complexe de $R$-modules.
\end{lemma}

\begin{proof}
La borne inférieure a été établie dans le lemme \ref{dualizing-lemma-relative-dualizing-trivial-vanishing}.
Choisissons une factorisation $R \to P \to A$ avec $P = R[x_1, \ldots, x_n]$ comme dans la section \ref{dualizing-section-relative-dualizing-complex-algebraic},
de sorte que $\omega_{A/R}^\bullet = R\Hom(A, P)[n]$.
La borne supérieure signifie que $\Ext^i_P(A, P)$ est nul pour $i > n$.
Cela résulte de Compléments d'algèbre, lemme \ref{more-algebra-lemma-perfect-over-polynomial-ring},
qui montre que $A$ est un $P$-module parfait
d'amplitude de Tor dans $[-n, 0]$.

\medskip\noindent
Preuve de la dernière assertion. Soit $R \to R'$ un homomorphisme d'anneaux noethériens.
Posons $A' = A \otimes_R R'$. Alors
$$
\omega_{A'/R'}^\bullet =
\omega_{A/R}^\bullet \otimes_A^\mathbf{L} A' =
\omega_{A/R}^\bullet \otimes_R^\mathbf{L} R'
$$
Le premier isomorphisme résulte du lemme \ref{dualizing-lemma-base-change-relative-algebraic} et le second,
dans $D(R')$, de Compléments d'algèbre, lemme \ref{more-algebra-lemma-base-change-comparison}.
D'après la première partie de la démonstration
(noter que les fibres de $\Spec(A') \to \Spec(R')$ sont de dimension $\leq d$),
on en déduit que $\omega_{A/R}^\bullet \otimes_R^\mathbf{L} R'$ n'a de cohomologie qu'en degrés $[-d, 0]$.
En prenant pour $R' = R \oplus M$ l'épaississement de carré nul de $R$
par un $R$-module de type fini $M$, on voit que $R\Hom(A, P) \otimes_R^\mathbf{L} M$
n'a de cohomologie que dans l'intervalle $[-d, 0]$
pour tout $R$-module de type fini $M$.
Puisque tout $R$-module est limite inductive filtrante de $R$-modules
de type fini et que les produits tensoriels commutent aux limites inductives,
on conclut.
\end{proof}

\begin{lemma}
\label{dualizing-lemma-relative-dualizing-CM-vanishing}
Soit $R \to A$ un homomorphisme de type fini d'anneaux noethériens.
Soit $\mathfrak p \subset R$ un idéal premier. Supposons que
\begin{enumerate}
\item $R_\mathfrak p$ soit de Cohen-Macaulay ;
\item pour tout idéal premier minimal $\mathfrak q \subset A$, on ait $\text{trdeg}_{\kappa(R \cap \mathfrak q)} \kappa(\mathfrak q) \leq r$.
\end{enumerate}
Alors
$$
H^i(\omega_{A/R}^\bullet)_\mathfrak p \not = 0 \Rightarrow - r \leq i
$$
et $H^{-r}(\omega_{A/R}^\bullet)_\mathfrak p$ satisfait à $(S_2)$ en tant que $A_\mathfrak p$-module.
\end{lemma}

\begin{proof}
On peut remplacer $R$ par $R_\mathfrak p$ d'après le lemme \ref{dualizing-lemma-base-change-relative-algebraic}.
On peut donc supposer que $R$ est un anneau local de Cohen-Macaulay,
et il faut établir les assertions du lemme pour les $A$-modules $H^i(\omega_{A/R}^\bullet)$.

\medskip\noindent
Soit $R^\wedge$ le complété de $R$.
L'homomorphisme $R \to R^\wedge$ est plat et $R^\wedge$ est de Cohen-Macaulay
(Compléments d'algèbre, lemme \ref{more-algebra-lemma-completion-CM}).
Observons que les idéaux premiers minimaux de $A \otimes_R R^\wedge$
sont au-dessus d'idéaux premiers minimaux de $A$
par platitude de $A \to A \otimes_R R^\wedge$ (et par le théorème de descente pour les morphismes plats,
voir Algèbre, lemme \ref{algebra-lemma-flat-going-down}).
La condition (2) vaut donc pour l'homomorphisme de type fini $R^\wedge \to A \otimes_R R^\wedge$
d'après Morphismes, lemme \ref{morphisms-lemma-dimension-fibre-after-base-change}.
En appliquant encore le lemme \ref{dualizing-lemma-base-change-relative-algebraic}, il suffit de prouver le lemme pour $R^\wedge \to A \otimes_R R^\wedge$.
De cette manière, en utilisant le lemme \ref{dualizing-lemma-ubiquity-dualizing},
on peut supposer que $R$ est un anneau local noethérien de Cohen-Macaulay
possédant un complexe dualisant $\omega_R^\bullet$.

\medskip\noindent
Soit $\mathfrak m \subset A$ un idéal maximal.
Il suffit de montrer les assertions du lemme pour $H^i(\omega_{A/R}^\bullet)_\mathfrak m$.
Si $\mathfrak m$ n'est pas au-dessus de l'idéal maximal de $R$,
on remplace $R$ par un localisé pour se ramener à ce cas
(nous omettons un petit détail).

\medskip\noindent
On peut supposer $\omega_R^\bullet$ normalisé. En posant $d = \dim(R)$,
on voit que $\omega_R^\bullet = \omega_R[d]$ pour un certain $R$-module $\omega_R$ ; voir le lemme \ref{dualizing-lemma-apply-CM}.
Posons $\omega_A^\bullet = \varphi^!(\omega_R^\bullet)$. D'après le lemme \ref{dualizing-lemma-relative-dualizing-if-have-omega}, on a
$$
\omega_{A/R}^\bullet =
R\Hom_A(\omega_R[d] \otimes_R^\mathbf{L} A, \omega_A^\bullet)
$$
La formule de dimension donne $\dim(A_\mathfrak m) \leq d + r$ ; voir Morphismes, lemme \ref{morphisms-lemma-dimension-formula-general},
et utiliser le fait que $\kappa(\mathfrak m)$ est fini sur le corps résiduel de $R$
par le Nullstellensatz de Hilbert.
D'après le lemme \ref{dualizing-lemma-shriek-normalized}, $(\omega_A^\bullet)_\mathfrak m$ est un complexe dualisant normalisé pour $A_\mathfrak m$.
Ainsi $H^i((\omega_A^\bullet)_\mathfrak m)$ n'est non nul que pour $-d - r \leq i \leq 0$ ; voir le lemme \ref{dualizing-lemma-sitting-in-degrees}.
Puisque $\omega_R[d] \otimes_R^\mathbf{L} A$ est placé en degrés $\leq -d$, on obtient l'annulation voulue.
Enfin, on voit aussi que
$$
H^{-r}(\omega_{A/R}^\bullet)_\mathfrak m =
\Hom_A(\omega_R \otimes_R A, H^{-d - r}(\omega_A^\bullet))_\mathfrak m
$$
Puisque $H^{-d - r}(\omega_A^\bullet)_\mathfrak m$ satisfait à $(S_2)$ d'après le lemme \ref{dualizing-lemma-depth-dualizing-module},
la dernière assertion résulte de Compléments d'algèbre, lemme \ref{more-algebra-lemma-hom-into-S2}.
\end{proof}



\section{Compléments sur les complexes dualisants}
\label{dualizing-section-more-dualizing}

\noindent
Quelques lemmes qui ne trouvent pas naturellement leur place ailleurs.

\begin{lemma}
\label{dualizing-lemma-descent}
Soit $A \to B$ un homomorphisme fidèlement plat d'anneaux noethériens.
Si $K \in D(A)$ et si $K \otimes_A^\mathbf{L} B$ est un complexe dualisant pour $B$,
alors $K$ est un complexe dualisant pour $A$.
\end{lemma}

\begin{proof}
Puisque $A \to B$ est plat, on a $H^i(K) \otimes_A B = H^i(K \otimes_A^\mathbf{L} B)$.
Puisque $K \otimes_A^\mathbf{L} B$ appartient à $D^b_{\textit{Coh}}(B)$, on trouve d'abord que $K$ appartient à $D^b(A)$,
puis que $H^i(K)$ est un $A$-module de type fini d'après Algèbre, lemme \ref{algebra-lemma-descend-properties-modules}.
Soit $M$ un $A$-module de type fini. Alors
$$
R\Hom_A(M, K) \otimes_A B = R\Hom_B(M \otimes_A B, K \otimes_A^\mathbf{L} B)
$$
d'après Compléments d'algèbre, lemme \ref{more-algebra-lemma-base-change-RHom}.
Puisque $K \otimes_A^\mathbf{L} B$ est de dimension injective finie,
disons d'amplitude injective dans $[a, b]$,
le membre de droite a une cohomologie nulle en degrés $> b$.
Puisque $A \to B$ est fidèlement plat, $R\Hom_A(M, K)$ a une cohomologie nulle en degrés $> b$.
Ainsi $K$ est de dimension injective finie
d'après Compléments d'algèbre, lemme \ref{more-algebra-lemma-injective-amplitude}.
Pour conclure, il faut montrer que le morphisme $A \to R\Hom_A(K, K)$ est un isomorphisme.
Utilisons encore Compléments d'algèbre, lemme \ref{more-algebra-lemma-base-change-RHom},
et le fait que
$B \to R\Hom_B(K \otimes_A^\mathbf{L} B, K \otimes_A^\mathbf{L} B)$
est un isomorphisme.
\end{proof}

\begin{lemma}
\label{dualizing-lemma-descent-ascent}
Soit $\varphi : A \to B$ un homomorphisme d'anneaux noethériens. Supposons que
\begin{enumerate}
\item $A \to B$ soit syntomique et induise une application surjective sur les spectres, ou
\item $A \to B$ soit fidèlement plat et localement d'intersection complète, ou
\item $A \to B$ soit fidèlement plat de type fini, à fibres de Gorenstein.
\end{enumerate}
Alors $K \in D(A)$ est un complexe dualisant pour $A$ si et seulement si
$K \otimes_A^\mathbf{L} B$ est un complexe dualisant pour $B$.
\end{lemma}

\begin{proof}
Observons que $A \to B$ satisfait à (1) si et seulement si $A \to B$ satisfait à (2)
d'après Compléments d'algèbre, lemme \ref{more-algebra-lemma-syntomic-lci}.
Dans les cas (2) et (3), le complexe dualisant relatif $\varphi^!(A) = \omega_{B/A}^\bullet$
est un objet inversible de $D(B)$ ; voir les lemmes \ref{dualizing-lemma-lci-shriek} et \ref{dualizing-lemma-gorenstein-shriek}.
De plus, on a $\varphi^!(K) = K \otimes_A^\mathbf{L} \omega_{B/A}^\bullet$ dans les deux cas ; voir le lemme \ref{dualizing-lemma-upper-shriek-is-tensor-functor} pour le cas (3).
Ainsi $\varphi^!(K)$ est identique à $K \otimes_A^\mathbf{L} B$ à tensorisation près
par un objet inversible de $D(B)$.
Par conséquent, $\varphi^!(K)$ est un complexe dualisant pour $B$
si et seulement si $K \otimes_A^\mathbf{L} B$ l'est
(la propriété d'être dualisant étant locale et invariante par décalage).
Si $K$ est dualisant pour $A$, alors $K \otimes_A^\mathbf{L} B$ est donc dualisant
pour $B$ d'après le lemme \ref{dualizing-lemma-shriek-dualizing-algebraic}.
Pour la descente de cette propriété, voir le lemme \ref{dualizing-lemma-descent}.
\end{proof}

\begin{lemma}
\label{dualizing-lemma-injective-hull-goes-up}
Soit $(A, \mathfrak m, \kappa) \to (B, \mathfrak n, l)$ un homomorphisme local plat d'anneaux noethériens tel que $\mathfrak n = \mathfrak m B$.
Si $E$ est l'enveloppe injective de $\kappa$,
alors $E \otimes_A B$ est l'enveloppe injective de $l$.
\end{lemma}

\begin{proof}
Écrivons $E = \bigcup E_n$ comme dans le lemme \ref{dualizing-lemma-union-artinian}.
Il suffit de montrer que $E_n \otimes_{A/\mathfrak m^n} B/\mathfrak n^n$ est l'enveloppe injective de $l$ sur $B/\mathfrak n$.
On se ramène ainsi au cas où $A$ et $B$ sont locaux artiniens.
Observons que $\text{longueur}_A(A) = \text{longueur}_B(B)$ et $\text{longueur}_A(E) = \text{longueur}_B(E \otimes_A B)$ d'après Algèbre, lemme \ref{algebra-lemma-pullback-module}.
D'après le lemme \ref{dualizing-lemma-finite}, on a $\text{longueur}_A(E) = \text{longueur}_A(A)$ et $\text{longueur}_B(E') = \text{longueur}_B(B)$,
où $E'$ est l'enveloppe injective de $l$ sur $B$.
On en déduit que $\text{longueur}_B(E') = \text{longueur}_B(E \otimes_A B)$. Observons que
$$
\dim_l((E \otimes_A B)[\mathfrak n]) =
\dim_l(E[\mathfrak m] \otimes_A B) =
\dim_\kappa(E[\mathfrak m]) = 1
$$
où l'on a utilisé la platitude de $A \to B$ et l'égalité $\mathfrak n = \mathfrak mB$.
Il existe donc un homomorphisme injectif de $B$-modules $E \otimes_A B \to E'$
d'après le lemme \ref{dualizing-lemma-torsion-submodule-sum-injective-hulls}.
C'est un isomorphisme par l'égalité des longueurs établie ci-dessus.
\end{proof}

\begin{lemma}
\label{dualizing-lemma-injective-goes-up}
Soit $\varphi : A \to B$ un homomorphisme plat d'anneaux noethériens tel que,
pour tout idéal premier $\mathfrak q \subset B$, on ait $\mathfrak p B_\mathfrak q = \mathfrak qB_\mathfrak q$,
où $\mathfrak p = \varphi^{-1}(\mathfrak q)$, par exemple si $\varphi$ est \'etale.
Si $I$ est un $A$-module injectif, alors $I \otimes_A B$ est un $B$-module injectif.
\end{lemma}

\begin{proof}
Les morphismes \'etales satisfont à l'hypothèse d'après Algèbre, lemme \ref{algebra-lemma-etale-at-prime}.
Le lemme \ref{dualizing-lemma-sum-injective-modules} et la proposition \ref{dualizing-proposition-structure-injectives-noetherian} permettent de supposer que $I$
est l'enveloppe injective de $\kappa(\mathfrak p)$ pour un idéal premier $\mathfrak p \subset A$.
Alors $I$ est un module sur $A_\mathfrak p$.
Il suffit de prouver que $I \otimes_A B = I \otimes_{A_\mathfrak p} B_\mathfrak p$ est injectif comme $B_\mathfrak p$-module ; voir le lemme \ref{dualizing-lemma-injective-flat}.
On peut donc supposer $(A, \mathfrak m, \kappa)$ local noethérien
et $I = E$ égal à l'enveloppe injective du corps résiduel $\kappa$.
Notre hypothèse entraîne que l'anneau noethérien $B/\mathfrak m B$
est produit de corps (nous omettons les détails).
Il n'y a donc qu'un nombre fini d'idéaux premiers $\mathfrak m_1, \ldots, \mathfrak m_n$ dans $B$
au-dessus de $\mathfrak m$, et ils sont tous maximaux.
Écrivons $E = \bigcup E_n$ comme dans le lemme \ref{dualizing-lemma-union-artinian}.
Alors $E \otimes_A B = \bigcup E_n \otimes_A B$, et $E_n \otimes_A B$ est un $B$-module de type fini
de support $\{\mathfrak m_1, \ldots, \mathfrak m_n\}$, donc se décompose en produit des localisés en les $\mathfrak m_i$.
Ainsi $E \otimes_A B = \prod (E \otimes_A B)_{\mathfrak m_i}$.
Puisque $(E \otimes_A B)_{\mathfrak m_i} = E \otimes_A B_{\mathfrak m_i}$ est l'enveloppe injective du corps résiduel de $\mathfrak m_i$
d'après le lemme \ref{dualizing-lemma-injective-hull-goes-up}, on conclut.
\end{proof}






\section{Complexes dualisants relatifs}
\label{dualizing-section-relative-dualizing-complexes}

\noindent
Pour un homomorphisme de type fini $\varphi : R \to A$ d'anneaux noethériens,
on dispose du complexe dualisant relatif $\omega_{A/R}^\bullet = \varphi^!(R)$ étudié dans la section \ref{dualizing-section-relative-dualizing-complexes-Noetherian}.
Si $R$ n'est pas noethérien, un complexe construit de la même manière
n'aura généralement pas de bonnes propriétés.
Nous donnons ici une définition du complexe dualisant relatif
pour les homomorphismes plats de présentation finie $R \to A$
d'anneaux non noethériens. Cette définition est choisie pour se globaliser
aux morphismes plats de présentation finie de schémas ;
voir Dualité pour les schémas, section \ref{duality-section-relative-dualizing-complexes}.
Nous montrerons que les complexes dualisants relatifs existent
dans le cadre de la définition, sont uniques à isomorphisme non canonique près,
et que l'on retrouve dans le cas noethérien le complexe de la section \ref{dualizing-section-relative-dualizing-complexes-Noetherian}.

\medskip\noindent
Le lecteur qui s'en tient au cas noethérien peut sans risque sauter cette section !

\begin{definition}
\label{dualizing-definition-relative-dualizing-complex}
Soit $R \to A$ un homomorphisme plat de présentation finie.
Un {\it complexe dualisant relatif} est un objet $K \in D(A)$ tel que
\begin{enumerate}
\item $K$ soit $R$-parfait (Compléments d'algèbre, définition \ref{more-algebra-definition-relatively-perfect}) ;
\item $R\Hom_{A \otimes_R A}(A, K \otimes_A^\mathbf{L} (A \otimes_R A))$ soit isomorphe à $A$.
\end{enumerate}
\end{definition}

\noindent
Pour comprendre cette définition, il peut être nécessaire de lire
et de comprendre certains des lemmes suivants.
Les lemmes \ref{dualizing-lemma-relative-dualizing-noetherian} et \ref{dualizing-lemma-uniqueness-relative-dualizing} montrent qu'elle est compatible
avec la définition de la section \ref{dualizing-section-relative-dualizing-complexes-Noetherian}.

\begin{lemma}
\label{dualizing-lemma-uniqueness-relative-dualizing}
Soit $R \to A$ un homomorphisme plat de présentation finie.
Deux complexes dualisants relatifs pour $R \to A$ sont isomorphes.
\end{lemma}

\begin{proof}
Soient $K$ et $L$ deux complexes dualisants relatifs pour $R \to A$.
Notons $K_1 = K \otimes_A^\mathbf{L} (A \otimes_R A)$ et $L_2 = (A \otimes_R A) \otimes_A^\mathbf{L} L$ les changements de base dérivés
par les première et seconde coprojections $A \to A \otimes_R A$.
Par symétrie, l'hypothèse sur $L_2$ entraîne que $R\Hom_{A \otimes_R A}(A, L_2)$ est isomorphe à $A$.
En appliquant deux fois Compléments d'algèbre, lemme \ref{more-algebra-lemma-internal-hom-evaluate-tensor-isomorphism} (3), on a
$$
A \otimes_{A \otimes_R A}^\mathbf{L} L_2 \cong
R\Hom_{A \otimes_R A}(A, K_1 \otimes_{A \otimes_R A}^\mathbf{L} L_2) \cong
A \otimes_{A \otimes_R A}^\mathbf{L} K_1
$$
L'application du foncteur de restriction $D(A \otimes_R A) \to D(A)$
pour l'une ou l'autre coprojection donne le résultat voulu.
\end{proof}

\begin{lemma}
\label{dualizing-lemma-relative-dualizing-noetherian}
Soit $\varphi : R \to A$ un homomorphisme plat de type fini d'anneaux noethériens.
Alors le complexe dualisant relatif $\omega_{A/R}^\bullet = \varphi^!(R)$ de la section \ref{dualizing-section-relative-dualizing-complexes-Noetherian}
est un complexe dualisant relatif au sens de la définition \ref{dualizing-definition-relative-dualizing-complex}.
\end{lemma}

\begin{proof}
Le lemme \ref{dualizing-lemma-relative-dualizing-algebraic} montre que $\varphi^!(R)$ est $R$-parfait.
Notons $\delta : A \otimes_R A \to A$ l'homomorphisme de multiplication et $p_1, p_2 : A \to A \otimes_R A$ les coprojections.
Alors
$$
\varphi^!(R) \otimes_A^\mathbf{L} (A \otimes_R A) =
\varphi^!(R) \otimes_{A, p_1}^\mathbf{L} (A \otimes_R A) =
p_2^!(A)
$$
d'après le lemme \ref{dualizing-lemma-flat-bc}. Rappelons que
$
R\Hom_{A \otimes_R A}(A, \varphi^!(R) \otimes_A^\mathbf{L} (A \otimes_R A))
$
est l'image de $\delta^!(\varphi^!(R) \otimes_A^\mathbf{L} (A \otimes_R A))$ par le foncteur de restriction $\delta_* : D(A) \to D(A \otimes_R A)$.
Utiliser la définition de $\delta^!$ dans la section \ref{dualizing-section-relative-dualizing-complex-algebraic}
et le lemme \ref{dualizing-lemma-RHom-ext}.
Puisque $\delta^!(p_2^!(A)) \cong A$ d'après le lemme \ref{dualizing-lemma-composition-shriek-algebraic}, on conclut.
\end{proof}

\begin{lemma}
\label{dualizing-lemma-base-change-relative-dualizing}
Soit $R \to A$ un homomorphisme plat de présentation finie. Alors :
\begin{enumerate}
\item il existe un complexe dualisant relatif $K$ dans $D(A)$ ;
\item pour tout homomorphisme $R \to R'$, si l'on pose $A' = A \otimes_R R'$
et $K' = K \otimes_A^\mathbf{L} A'$, alors $K'$ est un complexe dualisant relatif pour $R' \to A'$.
\end{enumerate}
De plus, si
$$
\xi : A \longrightarrow K \otimes_A^\mathbf{L} (A \otimes_R A)
$$
est un générateur du module cyclique
$\Hom_{D(A \otimes_R A)}(A, K \otimes_A^\mathbf{L} (A \otimes_R A))$
alors, dans (2), le changement de base dérivé de $\xi$ par
$A \otimes_R A \to A' \otimes_{R'} A'$ est un générateur du module cyclique
$\Hom_{D(A' \otimes_{R'} A')}(A',
K' \otimes_{A'}^\mathbf{L} (A' \otimes_{R'} A'))$
\end{lemma}

\begin{proof}
Ramenons-nous d'abord au cas noethérien.
D'après Algèbre, lemme \ref{algebra-lemma-flat-finite-presentation-limit-flat}, il existe une sous-$\mathbf{Z}$-algèbre de type fini
$R_0 \subset R$ et un homomorphisme plat de type fini $R_0 \to A_0$ tels que $A = A_0 \otimes_{R_0} R$.
D'après le lemme \ref{dualizing-lemma-relative-dualizing-noetherian}, il existe un complexe dualisant relatif $K_0 \in D(A_0)$.
Si l'on montre (2) pour $K_0$, alors $K_0 \otimes_{A_0}^\mathbf{L} A$ est un complexe dualisant pour $R \to A$
et satisfait lui aussi à (2) par transitivité du changement de base dérivé.
L'unicité des complexes dualisants relatifs (lemme \ref{dualizing-lemma-uniqueness-relative-dualizing})
montre alors que cela vaut pour tout complexe dualisant relatif.

\medskip\noindent
Supposons $R$ noethérien et soit $K$ un complexe dualisant relatif pour $R \to A$.
Étant donné un homomorphisme $R \to R'$, posons $A' = A \otimes_R R'$ et $K' = K \otimes_A^\mathbf{L} A'$.
Pour conclure, il faut montrer que $K'$ est un complexe dualisant relatif pour $R' \to A'$.
D'après Compléments d'algèbre, lemme \ref{more-algebra-lemma-base-change-relatively-perfect},
on voit que $K'$ est $R'$-parfait dans tous les cas.
D'après les lemmes \ref{dualizing-lemma-base-change-relative-algebraic} et \ref{dualizing-lemma-relative-dualizing-noetherian}, si $R'$ est noethérien,
alors $K'$ est un complexe dualisant relatif pour $R' \to A'$ dans les deux sens.
La transitivité du produit tensoriel dérivé donne $K \otimes_A^\mathbf{L} (A \otimes_R A)
\otimes_{A \otimes_R A}^\mathbf{L} (A' \otimes_{R'} A') =
K' \otimes_{A'}^\mathbf{L} (A' \otimes_{R'} A')$.
La platitude de $R \to A$ assure que $A \otimes_{A \otimes_R A}^\mathbf{L} (A' \otimes_{R'} A') = A'$ ;
en effet, $A \otimes_R A$ et $R'$ sont Tor-indépendants sur $R$,
ce qui permet d'appliquer Compléments d'algèbre, lemme \ref{more-algebra-lemma-base-change-comparison}.
Enfin, $A$ est pseudo-cohérent comme $A \otimes_R A$-module
d'après Compléments d'algèbre, lemme \ref{more-algebra-lemma-more-relative-pseudo-coherent-is-moot}.
Toutes les hypothèses de Compléments d'algèbre, lemme \ref{more-algebra-lemma-compute-RHom-relatively-perfect},
sont donc vérifiées.
On obtient un complexe borné inférieurement $E^\bullet$ de modules
plats sur $R$ et de présentation finie sur $A \otimes_R A$
tel que $E^\bullet \otimes_R R'$ représente $R\Hom_{A' \otimes_{R'} A'}(A',
K' \otimes_{A'}^\mathbf{L} (A' \otimes_{R'} A'))$,
ces identifications étant compatibles au changement de base dérivé.
Soit $n \in \mathbf{Z}$ tel que $n \not = 0$. Définissons $Q^n$ par la suite
$$
E^{n - 1} \to E^n \to Q^n \to 0
$$
Puisque $\kappa(\mathfrak p)$ est un anneau noethérien, on sait que $H^n(E^\bullet \otimes_R \kappa(\mathfrak p)) = 0$,
voir les remarques ci-dessus. Une chasse au diagramme montre alors que
$$
Q^n \otimes_R \kappa(\mathfrak p) \to E^{n + 1} \otimes_R \kappa(\mathfrak p)
$$
est injectif. Ainsi, pour un idéal premier $\mathfrak q$ de $A \otimes_R A$
au-dessus de $\mathfrak p$, le module $Q^n_\mathfrak q$ est $R_\mathfrak p$-plat
et $Q^n_\mathfrak p \to E^{n + 1}_\mathfrak q$ est $R_\mathfrak p$-universellement injectif ; voir Algèbre, lemme \ref{algebra-lemma-mod-injective}.
Puisque cela vaut pour tous les idéaux premiers,
$Q^n$ est $R$-plat et $Q^n \to E^{n + 1}$ est $R$-universellement injectif.
En particulier, $H^n(E^\bullet \otimes_R R') = 0$ pour tout homomorphisme $R \to R'$.
Soit $Z^0 = \Ker(E^0 \to E^1)$. La suite exacte $0 \to Z^0 \to E^0 \to E^1 \to Q^1 \to 0$ montre que $Z^0$ est $R$-plat et que
$Z^0 \otimes_R R' = \Ker(E^0 \otimes_R R' \to E^1 \otimes_R R')$
pour tout $R \to R'$. La suite exacte courte
$0 \to Q^{-1} \to Z^0 \to H^0(E^\bullet) \to 0$
donne alors
$$
H^0(E^\bullet \otimes_R R') = H^0(E^\bullet) \otimes_R R'
= A \otimes_R R' = A'
$$
comme voulu. Cette égalité donne aussi la dernière assertion du lemme.
\end{proof}

\begin{lemma}
\label{dualizing-lemma-relative-dualizing-RHom}
Soit $R \to A$ un homomorphisme plat de présentation finie.
Soit $K$ un complexe dualisant relatif.
Alors $A \to R\Hom_A(K, K)$ est un isomorphisme.
\end{lemma}

\begin{proof}
D'après Algèbre, lemme \ref{algebra-lemma-flat-finite-presentation-limit-flat}, il existe une sous-$\mathbf{Z}$-algèbre de type fini
$R_0 \subset R$ et un homomorphisme plat de type fini $R_0 \to A_0$ tels que $A = A_0 \otimes_{R_0} R$.
D'après les lemmes \ref{dualizing-lemma-uniqueness-relative-dualizing}, \ref{dualizing-lemma-relative-dualizing-noetherian} et \ref{dualizing-lemma-base-change-relative-dualizing},
il existe un complexe dualisant relatif $K_0 \in D(A_0)$
dont le changement de base dérivé est $K$.
On se ramène ainsi à la situation du paragraphe suivant.

\medskip\noindent
Supposons $R$ noethérien et soit $K$ un complexe dualisant relatif pour $R \to A$.
Étant donné un homomorphisme $R \to R'$, posons $A' = A \otimes_R R'$ et $K' = K \otimes_A^\mathbf{L} A'$.
Pour conclure, montrons que $R\Hom_{A'}(K', K') = A'$.
D'après le lemme \ref{dualizing-lemma-relative-dualizing-algebraic}, c'est vrai dès que $R'$ est noethérien.
Puisqu'un $R'$ quelconque est limite inductive filtrante
de $R$-algèbres noethériennes, le résultat découle de
Compléments d'algèbre, lemme \ref{more-algebra-lemma-colimit-relatively-perfect}.
\end{proof}

\begin{lemma}
\label{dualizing-lemma-relative-dualizing-composition}
Soient $R \to A \to B$ des homomorphismes plats de présentation finie.
Soient $K_{A/R}$ et $K_{B/A}$ des complexes dualisants relatifs pour $R \to A$ et $A \to B$.
Alors $K = K_{A/R} \otimes_A^\mathbf{L} K_{B/A}$ est un complexe dualisant relatif pour $R \to B$.
\end{lemma}

\begin{proof}
Utilisons la réduction au cas noethérien.
D'après Algèbre, lemme \ref{algebra-lemma-flat-finite-presentation-limit-flat}, il existe une sous-$\mathbf{Z}$-algèbre de type fini
$R_0 \subset R$ et un homomorphisme plat de type fini $R_0 \to A_0$ tels que $A = A_0 \otimes_{R_0} R$.
Quitte à agrandir $R_0$ et à remplacer $A_0$ en conséquence,
on peut supposer qu'il existe un homomorphisme plat de type fini $A_0 \to B_0$
tel que $B = B_0 \otimes_{R_0} R$ (utiliser le même lemme).
Si l'on prouve le lemme pour $R_0 \to A_0 \to B_0$, il en résulte
par les lemmes \ref{dualizing-lemma-uniqueness-relative-dualizing}, \ref{dualizing-lemma-relative-dualizing-noetherian} et \ref{dualizing-lemma-base-change-relative-dualizing}.
On se ramène ainsi à la situation du paragraphe suivant.

\medskip\noindent
Supposons $R$ noethérien et notons $\varphi : R \to A$ et $\psi : A \to B$
les homomorphismes donnés. Alors $K_{A/R} \cong \varphi^!(R)$ et $K_{B/A} \cong \psi^!(A)$,
voir les références ci-dessus.
On a alors
$$
K = K_{A/R} \otimes_A^\mathbf{L} K_{B/A} \cong
\varphi^!(R) \otimes_A^\mathbf{L} \psi^!(A) \cong
\psi^!(\varphi^!(R)) \cong (\psi \circ \varphi)^!(R)
$$
d'après les lemmes \ref{dualizing-lemma-upper-shriek-is-tensor-functor} et \ref{dualizing-lemma-composition-shriek-algebraic}.
Ainsi $K$ est un complexe dualisant relatif pour $R \to B$.
\end{proof}









% OK, start here.
%

\chapter{Dualité pour les schémas}



\phantomsection
\label{duality-section-phantom}


\section{Introduction}
\label{duality-section-introduction}

\noindent
Ce chapitre étudie la dualité relative pour les morphismes de schémas et le complexe dualisant sur un schéma. On pourra consulter \cite{RD}.

\medskip\noindent
Les complexes dualisants sur les anneaux noethériens ont été définis et étudiés dans Complexes dualisants, section \ref{dualizing-section-dualizing} et suivantes. Nous poursuivons ici cette étude en considérant les complexes dualisants sur les schémas, voir la section \ref{duality-section-dualizing-schemes}.

\medskip\noindent
L'essentiel de ce chapitre est consacré à l'étude de l'adjoint à droite de l'image directe dans le cadre des catégories dérivées de faisceaux de modules à faisceaux de cohomologie quasi-cohérents. Voir les sections \ref{duality-section-twisted-inverse-image}, \ref{duality-section-restriction-to-opens}, \ref{duality-section-base-change-map}, \ref{duality-section-base-change-II}, \ref{duality-section-trace}, \ref{duality-section-compare-with-pullback}, \ref{duality-section-sections-with-exact-support}, \ref{duality-section-duality-finite}, \ref{duality-section-perfect-proper}, \ref{duality-section-dualizing-Cartier} et \ref{duality-section-examples}. Nous suivons ici les articles \cite{Neeman-Grothendieck}, \cite{LN}, \cite{Lipman-notes} et \cite{Neeman-improvement}.

\medskip\noindent
Nous présentons les importants foncteurs image inverse extraordinaire $f^!$ pour les morphismes séparés de type fini entre schémas noethériens dans les sections \ref{duality-section-upper-shriek}, \ref{duality-section-upper-shriek-properties} et \ref{duality-section-base-change-shriek}, qui aboutissent à la synthèse de la section \ref{duality-section-duality}.

\medskip\noindent
Dans la section \ref{duality-section-glue}, nous exposons une autre théorie de la dualité et des complexes dualisants lorsque l'on travaille sur une base localement noethérienne fixée, munie d'un complexe dualisant (cette section correspond à une remarque du livre de Hartshorne).

\medskip\noindent
Les dernières sections présentent quelques applications.

\medskip\noindent
Ce chapitre se prolonge par celui consacré à la dualité sur les espaces algébriques, voir Dualité pour les espaces, section \ref{spaces-duality-section-introduction}.







\section{Complexes dualisants sur les schémas}
\label{duality-section-dualizing-schemes}

\noindent
Nous appelons complexe dualisant sur un schéma localement noethérien un complexe qui, localement pour la topologie des ouverts affines, provient d'un complexe dualisant sur l'anneau correspondant. Cette définition n'est pas tout à fait usuelle, mais elle coïncide avec toutes celles de la littérature pour les schémas noethériens de dimension finie.

\begin{lemma}
\label{duality-lemma-equivalent-definitions}
Soit $X$ un schéma localement noethérien. Soit $K$ un objet de $D(\mathcal{O}_X)$. Les conditions suivantes sont équivalentes :
\begin{enumerate}
\item Pour tout ouvert affine $U = \Spec(A) \subset X$, il existe un complexe dualisant $\omega_A^\bullet$ sur $A$ tel que $K|_U$ soit isomorphe à l'image de $\omega_A^\bullet$ par le foncteur $\widetilde{} : D(A) \to D(\mathcal{O}_U)$.
\item Il existe un recouvrement ouvert affine $X = \bigcup U_i$, $U_i = \Spec(A_i)$, tel que, pour tout $i$, il existe un complexe dualisant $\omega_i^\bullet$ sur $A_i$ tel que $K|_{U_i}$ soit isomorphe à l'image de $\omega_i^\bullet$ par le foncteur $\widetilde{} : D(A_i) \to D(\mathcal{O}_{U_i})$.
\end{enumerate}
\end{lemma}

\begin{proof}
Supposons (2) et soit $U = \Spec(A)$ un ouvert affine de $X$. Comme la condition (2) implique que $K$ appartient à $D_\QCoh(\mathcal{O}_X)$, il existe un objet $\omega_A^\bullet$ de $D(A)$ dont le complexe de faisceaux quasi-cohérents associé est isomorphe à $K|_U$, voir Catégories dérivées des schémas, lemme \ref{perfect-lemma-affine-compare-bounded}. Nous allons montrer que $\omega_A^\bullet$ est un complexe dualisant sur $A$, ce qui achèvera la démonstration.

\medskip\noindent
Puisque $X = \bigcup U_i$ est un recouvrement ouvert, on peut trouver un recouvrement par des ouverts principaux $U = D(f_1) \cup \ldots \cup D(f_m)$ tel que chaque $D(f_j)$ soit un ouvert principal de l'un des ouverts affines $U_i$, voir Schémas, lemme \ref{schemes-lemma-standard-open-two-affines}. Écrivons $D(f_j) = D(g_j)$, où $g_j \in A_{i_j}$. Alors $A_{f_j} \cong (A_{i_j})_{g_j}$ et l'on a
$$
(\omega_A^\bullet)_{f_j} \cong (\omega_i^\bullet)_{g_j}
$$
dans la catégorie dérivée, d'après Catégories dérivées des schémas, lemme \ref{perfect-lemma-affine-compare-bounded}. D'après Complexes dualisants, lemme \ref{dualizing-lemma-dualizing-localize}, le complexe $(\omega_A^\bullet)_{f_j}$ est dualisant sur $A_{f_j}$ pour $j = 1, \ldots, m$. Il en résulte que $\omega_A^\bullet$ est dualisant, d'après Complexes dualisants, lemme \ref{dualizing-lemma-dualizing-glue}.
\end{proof}

\begin{definition}
\label{duality-definition-dualizing-scheme}
Soit $X$ un schéma localement noethérien. Un objet $K$ de $D(\mathcal{O}_X)$ est appelé {\it complexe dualisant} si $K$ vérifie les conditions équivalentes du lemme \ref{duality-lemma-equivalent-definitions}.
\end{definition}

\noindent
Voir les remarques faites au début de cette section.

\begin{lemma}
\label{duality-lemma-affine-duality}
Soient $A$ un anneau noethérien et $X = \Spec(A)$. Soient $K, L$ des objets de $D(A)$. Si $K \in D_{\textit{Coh}}(A)$ et si $L$ est de dimension injective finie, alors
$$
R\SheafHom_{\mathcal{O}_X}(\widetilde{K}, \widetilde{L})
=
\widetilde{R\Hom_A(K, L)}
$$
dans $D(\mathcal{O}_X)$.
\end{lemma}

\begin{proof}
On peut supposer que $L$ est représenté par un complexe fini $I^\bullet$ de $A$-modules injectifs. Par récurrence sur la longueur de $I^\bullet$ et par compatibilité des constructions avec les triangles distingués, on se ramène au cas où $L = I[0]$, avec $I$ un $A$-module injectif. Dans ce cas, Catégories dérivées des schémas, lemme \ref{perfect-lemma-quasi-coherence-internal-hom}, affirme que le faisceau de cohomologie de degré $n$ de $R\SheafHom_{\mathcal{O}_X}(\widetilde{K}, \widetilde{L})$ est le faisceau associé au préfaisceau
$$
D(f) \longmapsto \Ext^n_{A_f}(K \otimes_A A_f, I \otimes_A A_f)
$$
Puisque $A$ est noethérien, le $A_f$-module $I \otimes_A A_f$ est injectif (Complexes dualisants, lemme \ref{dualizing-lemma-localization-injective-modules}). On obtient donc
\begin{align*}
\Ext^n_{A_f}(K \otimes_A A_f, I \otimes_A A_f)
& =
\Hom_{A_f}(H^{-n}(K \otimes_A A_f), I \otimes_A A_f) \\
& =
\Hom_{A_f}(H^{-n}(K) \otimes_A A_f, I \otimes_A A_f) \\
& =
\Hom_A(H^{-n}(K), I) \otimes_A A_f
\end{align*}
La dernière égalité résulte du fait que $H^{-n}(K)$ est un $A$-module de type fini, voir Algèbre, lemme \ref{algebra-lemma-hom-from-finitely-presented}. Cela montre que l'application canonique
$$
\widetilde{R\Hom_A(K, L)}
\longrightarrow
R\SheafHom_{\mathcal{O}_X}(\widetilde{K}, \widetilde{L})
$$
est un quasi-isomorphisme dans ce cas, ce qui achève la démonstration.
\end{proof}

\begin{lemma}
\label{duality-lemma-internal-hom-evaluate-isom}
Soit $X$ un schéma noethérien. Soient $K, L, M \in D_\QCoh(\mathcal{O}_X)$. Alors l'application
$$
R\SheafHom(L, M) \otimes_{\mathcal{O}_X}^\mathbf{L} K
\longrightarrow
R\SheafHom(R\SheafHom(K, L), M)
$$
de Cohomologie, lemme \ref{cohomology-lemma-internal-hom-evaluate}, est un isomorphisme dans les deux cas suivants :
\begin{enumerate}
\item $K \in D^-_{\textit{Coh}}(\mathcal{O}_X)$, $L \in D^+_{\textit{Coh}}(\mathcal{O}_X)$, et $M$ est localement sur les ouverts affines de dimension injective finie (voir la démonstration) ; ou
\item \begingroup\emergencystretch=12pt $K$ et $L$ appartiennent à $D_{\textit{Coh}}(\mathcal{O}_X)$, l'objet $R\SheafHom(L, M)$ est de Tor-dimension finie, et $L$ et $M$ sont localement sur les ouverts affines de dimension injective finie (en particulier, $L$ et $M$ sont bornés).\par\endgroup
\end{enumerate}
\end{lemma}

\begin{proof}
Démonstration de (1). On dit que $M$ est localement sur les ouverts affines de dimension injective finie si $X$ possède un recouvrement par des ouverts affines $U = \Spec(A)$ tels que l'objet de $D(A)$ correspondant à $M|_U$ (Catégories dérivées des schémas, lemme \ref{perfect-lemma-affine-compare-bounded}) soit de dimension injective finie\footnote{Cette condition ne dépend pas du choix du recouvrement ouvert affine du schéma noethérien $X$. Les détails sont omis.}. Pour démontrer le lemme, on peut remplacer $X$ par $U$, c'est-à-dire supposer que $X = \Spec(A)$ pour un anneau noethérien $A$. Observons que $R\SheafHom(K, L)$ appartient à $D^+_{\textit{Coh}}(\mathcal{O}_X)$ d'après Catégories dérivées des schémas, lemme \ref{perfect-lemma-coherent-internal-hom}. De plus, la formation des deux membres de la flèche commute au foncteur $D(A) \to D_\QCoh(\mathcal{O}_X)$ d'après le lemme \ref{duality-lemma-affine-duality} et Catégories dérivées des schémas, lemme \ref{perfect-lemma-quasi-coherence-internal-hom} (on utilise bien ici les hypothèses sur $K$, $L$, $M$ et ce que l'on vient de montrer au sujet de $R\SheafHom(K, L)$). Enfin, la flèche est un isomorphisme d'après Compléments d'algèbre, lemme \ref{more-algebra-lemma-internal-hom-evaluate-isomorphism}, (2).

\medskip\noindent
Démonstration de (2). On raisonne comme ci-dessus. La seule différence est que l'on obtient ici $R\SheafHom(K, L)$ dans $D_{\textit{Coh}}(\mathcal{O}_X)$, car on peut, localement sur les ouverts affines (ce que permet le lemme \ref{duality-lemma-affine-duality}), appliquer Complexes dualisants, lemme \ref{dualizing-lemma-finite-ext-into-bounded-injective}. On conclut alors par Compléments d'algèbre, lemme \ref{more-algebra-lemma-internal-hom-evaluate-isomorphism-technical}.
\end{proof}

\begin{lemma}
\label{duality-lemma-dualizing-schemes}
Soit $K$ un complexe dualisant sur un schéma localement noethérien $X$. Alors $K$ est un objet de $D_{\textit{Coh}}(\mathcal{O}_X)$ et $D = R\SheafHom_{\mathcal{O}_X}(-, K)$ induit une anti-équivalence
$$
D :
D_{\textit{Coh}}(\mathcal{O}_X)
\longrightarrow
D_{\textit{Coh}}(\mathcal{O}_X)
$$
munie d'un isomorphisme canonique $\text{id} \to D \circ D$. Si $X$ est quasi-compact, alors $D$ échange $D^+_{\textit{Coh}}(\mathcal{O}_X)$ et $D^-_{\textit{Coh}}(\mathcal{O}_X)$ et induit une anti-équivalence $D^b_{\textit{Coh}}(\mathcal{O}_X) \to D^b_{\textit{Coh}}(\mathcal{O}_X)$.
\end{lemma}

\begin{proof}
Soit $U \subset X$ un ouvert affine. Écrivons $U = \Spec(A)$ et soit $\omega_A^\bullet$ un complexe dualisant sur $A$ correspondant à $K|_U$ comme au lemme \ref{duality-lemma-equivalent-definitions}. D'après le lemme \ref{duality-lemma-affine-duality}, le diagramme
$$
\xymatrix{
D_{\textit{Coh}}(A) \ar[r] \ar[d]_{R\Hom_A(-, \omega_A^\bullet)} &
D_{\textit{Coh}}(\mathcal{O}_U) \ar[d]^{R\SheafHom_{\mathcal{O}_X}(-, K|_U)} \\
D_{\textit{Coh}}(A) \ar[r] &
D(\mathcal{O}_U)
}
$$
est commutatif. On en conclut que $D$ envoie $D_{\textit{Coh}}(\mathcal{O}_X)$ dans $D_{\textit{Coh}}(\mathcal{O}_X)$. De plus, l'application canonique
$$
L
\longrightarrow
R\SheafHom_{\mathcal{O}_X}(K, K) \otimes_{\mathcal{O}_X}^\mathbf{L} L
\longrightarrow
R\SheafHom_{\mathcal{O}_X}(R\SheafHom_{\mathcal{O}_X}(L, K), K)
$$
(où l'on utilise Cohomologie, lemme \ref{cohomology-lemma-internal-hom-evaluate}, pour la seconde flèche) est un isomorphisme pour tout $L$, car cela est vrai sur les ouverts affines d'après Complexes dualisants, lemme \ref{dualizing-lemma-dualizing}\footnote{On peut aussi commencer par montrer que $R\SheafHom_{\mathcal{O}_X}(K, K) = \mathcal{O}_X$ en travaillant localement sur les ouverts affines, puis utiliser le lemme \ref{duality-lemma-internal-hom-evaluate-isom}, (2), pour voir que l'application est un isomorphisme.}, et nous avons déjà vu que, sur les ouverts affines, on retrouve la construction algébrique. L'énoncé relatif aux propriétés de bornitude du foncteur $D$ dans le cas quasi-compact résulte également des énoncés correspondants de Complexes dualisants, lemme \ref{dualizing-lemma-dualizing}.
\end{proof}

\noindent
Soit $X$ un espace localement annelé. Rappelons qu'un objet $L$ de $D(\mathcal{O}_X)$ est {\it inversible} s'il est inversible pour la structure monoïdale symétrique sur $D(\mathcal{O}_X)$ définie par le produit tensoriel dérivé. Dans Cohomologie, lemme \ref{cohomology-lemma-invertible-derived}, nous avons vu que cela signifie que $L$ est parfait et qu'il existe un recouvrement ouvert $X = \bigcup U_i$ tel que $L|_{U_i} \cong \mathcal{O}_{U_i}[-n_i]$ pour certains entiers $n_i$. Dans ce cas, la fonction
$$
x \mapsto n_x,\quad
\text{où }n_x\text{ est l'unique entier tel que }
H^{n_x}(L_x) \not = 0
$$
est localement constante sur $X$. En particulier, on a $L = \bigoplus H^n(L)[-n]$, ce qui définit un complexe de $\mathcal{O}_X$-modules (à différentielles nulles) représentant $L$.

\begin{lemma}
\label{duality-lemma-dualizing-unique-schemes}
Soit $X$ un schéma localement noethérien. Si $K$ et $K'$ sont des complexes dualisants sur $X$, alors $K'$ est isomorphe à $K \otimes_{\mathcal{O}_X}^\mathbf{L} L$ pour un objet inversible $L$ de $D(\mathcal{O}_X)$.
\end{lemma}

\begin{proof}
Posons
$$
L = R\SheafHom_{\mathcal{O}_X}(K, K')
$$
C'est un objet inversible de $D(\mathcal{O}_X)$, car il en est ainsi localement sur les ouverts affines, voir Complexes dualisants, lemme \ref{dualizing-lemma-dualizing-unique} et sa démonstration. L'application d'évaluation $L \otimes_{\mathcal{O}_X}^\mathbf{L} K \to K'$ est un isomorphisme pour la même raison.
\end{proof}

\begin{lemma}
\label{duality-lemma-dimension-function-scheme}
Soit $X$ un schéma localement noethérien. Soit $\omega_X^\bullet$ un complexe dualisant sur $X$. Alors $X$ est universellement caténaire et la fonction $X \to \mathbf{Z}$ définie par
$$
x \longmapsto \delta(x)\text{ tel que }
\omega_{X, x}^\bullet[-\delta(x)]
\text{ soit un complexe dualisant normalisé sur }
\mathcal{O}_{X, x}
$$
est une fonction de dimension.
\end{lemma}

\begin{proof}
Cela résulte immédiatement du cas affine, Complexes dualisants, lemme \ref{dualizing-lemma-dimension-function}, et des définitions.
\end{proof}

\begin{lemma}
\label{duality-lemma-sitting-in-degrees}
Soit $X$ un schéma localement noethérien. Soit $\omega_X^\bullet$ un complexe dualisant sur $X$, de fonction de dimension associée $\delta$. Soit $\mathcal{F}$ un $\mathcal{O}_X$-module cohérent. Posons $\mathcal{E}^i = \SheafExt^{-i}_{\mathcal{O}_X}(\mathcal{F}, \omega_X^\bullet)$. Alors $\mathcal{E}^i$ est un $\mathcal{O}_X$-module cohérent et, pour $x \in X$, on a :
\begin{enumerate}
\item $\mathcal{E}^i_x$ ne peut être non nul que pour $\delta(x) \leq i \leq \delta(x) + \dim(\text{Supp}(\mathcal{F}_x))$ ;
\item $\dim(\text{Supp}(\mathcal{E}^{i + \delta(x)}_x)) \leq i$ ;
\item $\text{depth}(\mathcal{F}_x)$ est le plus petit entier $i \geq 0$ tel que $\mathcal{E}_x^{i + \delta(x)} \not = 0$ ;
\item on a $x \in \text{Supp}(\bigoplus_{j \leq i} \mathcal{E}^j)
\Leftrightarrow
\text{depth}_{\mathcal{O}_{X, x}}(\mathcal{F}_x) + \delta(x) \leq i$.
\end{enumerate}
\end{lemma}

\begin{proof}
Le lemme \ref{duality-lemma-dualizing-schemes} affirme que $\mathcal{E}^i$ est cohérent. En choisissant un voisinage affine de $x$ et en utilisant Catégories dérivées des schémas, lemme \ref{perfect-lemma-quasi-coherence-internal-hom}, ainsi que Compléments d'algèbre, lemme \ref{more-algebra-lemma-base-change-RHom}, (3), on obtient
$$
\mathcal{E}^i_x =
\SheafExt^{-i}_{\mathcal{O}_X}(\mathcal{F}, \omega_X^\bullet)_x =
\Ext^{-i}_{\mathcal{O}_{X, x}}(\mathcal{F}_x,
\omega_{X, x}^\bullet) =
\Ext^{\delta(x) - i}_{\mathcal{O}_{X, x}}(\mathcal{F}_x,
\omega_{X, x}^\bullet[-\delta(x)])
$$
La construction de $\delta$ au lemme \ref{duality-lemma-dimension-function-scheme} ramène ainsi les assertions (1), (2) et (3) à Complexes dualisants, lemme \ref{dualizing-lemma-sitting-in-degrees}. L'assertion (4) est une conséquence formelle de (3) et (1).
\end{proof}




\section{Adjoint à droite de l'image directe}
\label{duality-section-twisted-inverse-image}

\noindent
Pour cette section et les suivantes, on pourra consulter \cite{Neeman-Grothendieck}, \cite{LN}, \cite{Lipman-notes} et \cite{Neeman-improvement}.

\medskip\noindent
Soit $f : X \to Y$ un morphisme de schémas. Nous considérons dans cette section l'adjoint à droite du foncteur $Rf_* : D_\QCoh(\mathcal{O}_X) \to D_\QCoh(\mathcal{O}_Y)$. Dans la littérature, lorsqu'il existe, ce foncteur est parfois noté $f^{\times}$. Cette notation n'est pas universellement admise et nous ne l'utiliserons pas. Nous n'emploierons pas non plus la notation $f^!$ pour ce foncteur, car elle serait incompatible (pour des morphismes $f$ généraux) avec celle de \cite{RD}.

\begin{lemma}
\label{duality-lemma-twisted-inverse-image}
\begin{reference}
Cet énoncé est presque identique à \cite[Example 4.2]{Neeman-Grothendieck}.
\end{reference}
Soit $f : X \to Y$ un morphisme entre schémas quasi-séparés et quasi-compacts. Le foncteur $Rf_* : D_\QCoh(\mathcal{O}_X) \to D_\QCoh(\mathcal{O}_Y)$ admet un adjoint à droite.
\end{lemma}

\begin{proof}
Nous allons montrer l'existence d'un adjoint à droite en vérifiant les hypothèses de Catégories dérivées, proposition \ref{derived-proposition-brown}. Tout d'abord, la catégorie $D_\QCoh(\mathcal{O}_X)$ admet des sommes directes, voir Catégories dérivées des schémas, lemme \ref{perfect-lemma-quasi-coherence-direct-sums}. La catégorie $D_\QCoh(\mathcal{O}_X)$ est engendrée par des objets compacts d'après Catégories dérivées des schémas, théorème \ref{perfect-theorem-bondal-van-den-Bergh}. Puisque $X$ et $Y$ sont quasi-compacts et quasi-séparés, il en est de même de $f$, voir Schémas, lemmes \ref{schemes-lemma-compose-after-separated} et \ref{schemes-lemma-quasi-compact-permanence}. Le foncteur $Rf_*$ commute donc aux sommes directes, voir Catégories dérivées des schémas, lemme \ref{perfect-lemma-quasi-coherence-pushforward-direct-sums}. Cela achève la démonstration.
\end{proof}

\begin{example}
\label{duality-example-affine-twisted-inverse-image}
Soit $A \to B$ un homomorphisme d'anneaux. Soient $Y = \Spec(A)$ et $X = \Spec(B)$, et soit $f : X \to Y$ le morphisme correspondant à $A \to B$. Alors $Rf_* : D_\QCoh(\mathcal{O}_X) \to D_\QCoh(\mathcal{O}_Y)$ correspond à la restriction $D(B) \to D(A)$ par les équivalences $D(B) \to D_\QCoh(\mathcal{O}_X)$ et $D(A) \to D_\QCoh(\mathcal{O}_Y)$. L'adjoint à droite correspond donc au foncteur $K \longmapsto R\Hom(B, K)$ de Complexes dualisants, section \ref{dualizing-section-trivial}.
\end{example}

\begin{example}
\label{duality-example-does-not-preserve-coherent}
Si $f : X \to Y$ est un morphisme séparé de type fini entre schémas noethériens, l'adjoint à droite de $Rf_* : D_\QCoh(\mathcal{O}_X) \to D_\QCoh(\mathcal{O}_Y)$ n'envoie pas nécessairement $D_{\textit{Coh}}(\mathcal{O}_Y)$ dans $D_{\textit{Coh}}(\mathcal{O}_X)$. En effet, soit $k$ un corps et considérons le morphisme $f : \mathbf{A}^1_k \to \Spec(k)$. D'après l'exemple \ref{duality-example-affine-twisted-inverse-image}, cela revient à demander si $R\Hom(B, -)$ envoie $D_{\textit{Coh}}(A)$ dans $D_{\textit{Coh}}(B)$, où $A = k$ et $B = k[x]$. Ce n'est pas le cas, car
$$
R\Hom(k[x], k) = \left(\prod\nolimits_{n \geq 0} k\right)[0]
$$
n'est pas un $k[x]$-module de type fini. Ainsi, $a(\mathcal{O}_Y)$ n'a pas des faisceaux de cohomologie cohérents.
\end{example}

\begin{example}
\label{duality-example-does-not-preserve-bounded-above}
Si $f : X \to Y$ est un morphisme propre, ou même fini, entre schémas noethériens, l'adjoint à droite de $Rf_* : D_\QCoh(\mathcal{O}_X) \to D_\QCoh(\mathcal{O}_Y)$ n'envoie pas nécessairement $D_\QCoh^-(\mathcal{O}_Y)$ dans $D_\QCoh^-(\mathcal{O}_X)$. En effet, soit $k$ un corps, soit $k[\epsilon]$ l'anneau des nombres duaux sur $k$, et soient $X = \Spec(k)$ et $Y = \Spec(k[\epsilon])$. Alors $\Ext^i_{k[\epsilon]}(k, k)$ est non nul pour tout $i \geq 0$. Ainsi, $a(\mathcal{O}_Y)$ n'est pas borné supérieurement, d'après l'exemple \ref{duality-example-affine-twisted-inverse-image}.
\end{example}

\begin{lemma}
\label{duality-lemma-twisted-inverse-image-bounded-below}
Soit $f : X \to Y$ un morphisme de schémas quasi-compacts et quasi-séparés. Soit $a : D_\QCoh(\mathcal{O}_Y) \to D_\QCoh(\mathcal{O}_X)$ l'adjoint à droite de $Rf_*$ du lemme \ref{duality-lemma-twisted-inverse-image}. Alors $a$ envoie $D^+_\QCoh(\mathcal{O}_Y)$ dans $D^+_\QCoh(\mathcal{O}_X)$. Plus précisément, il existe un entier $N$ tel que $H^i(K) = 0$ pour $i \leq c$ implique $H^i(a(K)) = 0$ pour $i \leq c - N$.
\end{lemma}

\begin{proof}
D'après Catégories dérivées des schémas, lemme \ref{perfect-lemma-quasi-coherence-direct-image}, le foncteur $Rf_*$ est de dimension cohomologique finie. Autrement dit, il existe un entier $N$ tel que $H^i(Rf_*L) = 0$ pour $i \geq N + c$ si $H^i(L) = 0$ pour $i \geq c$. Supposons que $K \in D^+_\QCoh(\mathcal{O}_Y)$ vérifie $H^i(K) = 0$ pour $i \leq c$. Alors
$$
\Hom_{D(\mathcal{O}_X)}(\tau_{\leq c - N}a(K), a(K)) =
\Hom_{D(\mathcal{O}_Y)}(Rf_*\tau_{\leq c - N}a(K), K) = 0
$$
d'après ce qui précède. Cela implique clairement que $H^i(a(K)) = 0$ pour $i \leq c - N$.
\end{proof}

\noindent
Soit $f : X \to Y$ un morphisme de schémas quasi-séparés et quasi-compacts. Notons $a$ l'adjoint à droite de $Rf_* : D_\QCoh(\mathcal{O}_X) \to D_\QCoh(\mathcal{O}_Y)$. Pour tous $K \in D_\QCoh(\mathcal{O}_Y)$ et $L \in D_\QCoh(\mathcal{O}_X)$, on obtient une application canonique
\begin{equation}
\label{duality-equation-sheafy-trace}
Rf_*R\SheafHom_{\mathcal{O}_X}(L, a(K))
\longrightarrow
R\SheafHom_{\mathcal{O}_Y}(Rf_*L, K)
\end{equation}
Cette application est définie comme la composée
$$
Rf_*R\SheafHom_{\mathcal{O}_X}(L, a(K)) \to
R\SheafHom_{\mathcal{O}_Y}(Rf_*L, Rf_*a(K)) \to
R\SheafHom_{\mathcal{O}_Y}(Rf_*L, K)
$$
où la première flèche est celle de Cohomologie, remarque \ref{cohomology-remark-projection-formula-for-internal-hom}, et la seconde est la co-unité $Rf_*a(K) \to K$ de l'adjonction.

\begin{lemma}
\label{duality-lemma-iso-on-RSheafHom}
Soit $f : X \to Y$ un morphisme de schémas quasi-compacts et quasi-séparés. Soit $a$ l'adjoint à droite de $Rf_* : D_\QCoh(\mathcal{O}_X) \to D_\QCoh(\mathcal{O}_Y)$. Soient $L \in D_\QCoh(\mathcal{O}_X)$ et $K \in D_\QCoh(\mathcal{O}_Y)$. Alors l'application (\ref{duality-equation-sheafy-trace})
$$
Rf_*R\SheafHom_{\mathcal{O}_X}(L, a(K))
\longrightarrow
R\SheafHom_{\mathcal{O}_Y}(Rf_*L, K)
$$
\begingroup\emergencystretch=16pt
devient un isomorphisme après application du foncteur
$DQ_Y : D(\mathcal{O}_Y) \to D_\QCoh(\mathcal{O}_Y)$
étudié dans Catégories dérivées des schémas, section \ref{perfect-section-better-coherator}.\par\endgroup
\end{lemma}

\begin{proof}
L'énoncé a un sens puisque $DQ_Y$ existe d'après Catégories dérivées des schémas, lemme \ref{perfect-lemma-better-coherator}. Comme $DQ_Y$ est l'adjoint à droite du foncteur d'inclusion $D_\QCoh(\mathcal{O}_Y) \to D(\mathcal{O}_Y)$, il suffit de montrer que, pour tout $M \in D_\QCoh(\mathcal{O}_Y)$, l'application (\ref{duality-equation-sheafy-trace}) induit une bijection
$$
\Hom_Y(M, Rf_*R\SheafHom_{\mathcal{O}_X}(L, a(K)))
\longrightarrow
\Hom_Y(M, R\SheafHom_{\mathcal{O}_Y}(Rf_*L, K))
$$
Pour cela, on utilise la suite d'égalités suivante :
\begin{align*}
\Hom_Y(M, Rf_*R\SheafHom_{\mathcal{O}_X}(L, a(K)))
& =
\Hom_X(Lf^*M, R\SheafHom_{\mathcal{O}_X}(L, a(K))) \\
& =
\Hom_X(Lf^*M \otimes_{\mathcal{O}_X}^\mathbf{L} L, a(K)) \\
& =
\Hom_Y(Rf_*(Lf^*M \otimes_{\mathcal{O}_X}^\mathbf{L} L), K) \\
& =
\Hom_Y(M \otimes_{\mathcal{O}_Y}^\mathbf{L} Rf_*L, K) \\
& =
\Hom_Y(M, R\SheafHom_{\mathcal{O}_Y}(Rf_*L, K))
\end{align*}
La première égalité résulte de Cohomologie, lemme \ref{cohomology-lemma-adjoint}. La deuxième résulte de Cohomologie, lemme \ref{cohomology-lemma-internal-hom}. La troisième résulte de la construction de $a$. La quatrième résulte de Catégories dérivées des schémas, lemme \ref{perfect-lemma-cohomology-base-change} (c'est l'étape essentielle). La cinquième résulte de Cohomologie, lemme \ref{cohomology-lemma-internal-hom}.
\end{proof}

\begin{example}
\label{duality-example-iso-on-RSheafHom}
L'énoncé du lemme \ref{duality-lemma-iso-on-RSheafHom} est faux si l'on n'applique pas le « cohérateur » $DQ_Y$. Supposons en effet $Y = \Spec(R)$ et $X = \mathbf{A}^1_R$. Prenons $L = \mathcal{O}_X$ et $K = \mathcal{O}_Y$. Le membre de gauche de la flèche appartient à $D_\QCoh(\mathcal{O}_Y)$, tandis que celui de droite est isomorphe à $\prod_{n \geq 0} \mathcal{O}_Y$, qui n'est pas quasi-cohérent.
\end{example}

\begin{remark}
\label{duality-remark-iso-on-RSheafHom}
Dans la situation du lemme \ref{duality-lemma-iso-on-RSheafHom}, on a
$$
DQ_Y(Rf_*R\SheafHom_{\mathcal{O}_X}(L, a(K))) =
Rf_* DQ_X(R\SheafHom_{\mathcal{O}_X}(L, a(K)))
$$
\begingroup\emergencystretch=16pt
d'après Catégories dérivées des schémas, lemme \ref{perfect-lemma-pushforward-better-coherator}. Ainsi, si
\par\noindent\hfil
$R\SheafHom_{\mathcal{O}_X}(L, a(K)) \in D_\QCoh(\mathcal{O}_X)$,
\hfil\par\noindent
on peut « effacer » le $DQ_Y$ du membre de gauche de la flèche. D'autre part, si l'on sait que $R\SheafHom_{\mathcal{O}_Y}(Rf_*L, K) \in D_\QCoh(\mathcal{O}_Y)$, on peut « effacer » le $DQ_Y$ du membre de droite. Si les deux conditions sont satisfaites, on voit que (\ref{duality-equation-sheafy-trace}) est un isomorphisme. En combinant cela avec Catégories dérivées des schémas, lemme \ref{perfect-lemma-quasi-coherence-internal-hom}, on voit que $Rf_*R\SheafHom_{\mathcal{O}_X}(L, a(K)) \to
R\SheafHom_{\mathcal{O}_Y}(Rf_*L, K)$ est un isomorphisme si\par\endgroup
\begin{enumerate}
\item $L$ et $Rf_*L$ sont parfaits ; ou
\item $K$ est borné inférieurement et $L$ et $Rf_*L$ sont pseudo-cohérents.
\end{enumerate}
Pour (2), on utilise le fait que $a(K)$ est borné inférieurement si $K$ l'est, voir le lemme \ref{duality-lemma-twisted-inverse-image-bounded-below}.
\end{remark}

\begin{example}
\label{duality-example-iso-on-RSheafHom-noetherian}
\begingroup\emergencystretch=16pt
Soit $f : X \to Y$ un morphisme propre de schémas noethériens, et soient $L \in D^-_{\textit{Coh}}(X)$ et $K \in D^+_{\QCoh}(\mathcal{O}_Y)$. Alors l'application $Rf_*R\SheafHom_{\mathcal{O}_X}(L, a(K)) \to
R\SheafHom_{\mathcal{O}_Y}(Rf_*L, K)$ est un isomorphisme. En effet, les complexes $L$ et $Rf_*L$ sont pseudo-cohérents d'après Catégories dérivées des schémas, lemmes \ref{perfect-lemma-identify-pseudo-coherent-noetherian} et \ref{perfect-lemma-direct-image-coherent}, et l'on peut appliquer la remarque \ref{duality-remark-iso-on-RSheafHom}.\par\endgroup
\end{example}

\begin{lemma}
\label{duality-lemma-iso-global-hom}
Soit $f : X \to Y$ un morphisme de schémas quasi-séparés et quasi-compacts. Pour tous $L \in D_\QCoh(\mathcal{O}_X)$ et $K \in D_\QCoh(\mathcal{O}_Y)$, (\ref{duality-equation-sheafy-trace}) induit un isomorphisme $R\Hom_X(L, a(K)) \to R\Hom_Y(Rf_*L, K)$ entre les Hom dérivés globaux.
\end{lemma}

\begin{proof}
Par la construction de Cohomologie, section \ref{cohomology-section-global-RHom}, on a
$$
R\Hom_X(L, a(K)) =
R\Gamma(X, R\SheafHom_{\mathcal{O}_X}(L, a(K))) =
R\Gamma(Y, Rf_*R\SheafHom_{\mathcal{O}_X}(L, a(K)))
$$
et
$$
R\Hom_Y(Rf_*L, K) = R\Gamma(Y, R\SheafHom_{\mathcal{O}_Y}(Rf_*L, K))
$$
Le lemme est donc une conséquence du lemme \ref{duality-lemma-iso-on-RSheafHom}. En effet, une application $E \to E'$ dans $D(\mathcal{O}_Y)$ qui induit un isomorphisme $DQ_Y(E) \to DQ_Y(E')$ induit un quasi-isomorphisme $R\Gamma(Y, E) \to R\Gamma(Y, E')$. On a en effet $H^i(Y, E) = \Ext^i_Y(\mathcal{O}_Y, E) = \Hom(\mathcal{O}_Y[-i], E) =
\Hom(\mathcal{O}_Y[-i], DQ_Y(E))$, car $\mathcal{O}_Y[-i]$ appartient à $D_\QCoh(\mathcal{O}_Y)$ et $DQ_Y$ est l'adjoint à droite du foncteur d'inclusion $D_\QCoh(\mathcal{O}_Y) \to D(\mathcal{O}_Y)$.
\end{proof}







\section{Adjoint à droite de l'image directe et restriction aux ouverts}
\sectionmark{Adjoint à droite et ouverts}
\label{duality-section-restriction-to-opens}

\noindent
Nous étudions dans cette section dans quelle mesure l'adjoint à droite de l'image directe commute à la restriction aux sous-schémas ouverts. Il s'agit d'une question de changement de base ; commençons donc par l'aborder dans un cadre plus général.

\medskip\noindent
On souhaite souvent savoir si les adjoints à droite de l'image directe commutent au changement de base. Considérons donc un carré cartésien
\begin{equation}
\label{duality-equation-base-change}
\vcenter{
\xymatrix{
X' \ar[r]_{g'} \ar[d]_{f'} & X \ar[d]^f \\
Y' \ar[r]^g & Y
}
}
\end{equation}
de schémas quasi-compacts et quasi-séparés. Notons
$$
a  : D_\QCoh(\mathcal{O}_Y) \to D_\QCoh(\mathcal{O}_X)
\quad\text{et}\quad
a'  : D_\QCoh(\mathcal{O}_{Y'}) \to D_\QCoh(\mathcal{O}_{X'})
$$
les adjoints à droite de $Rf_*$ et $Rf'_*$ (lemme \ref{duality-lemma-twisted-inverse-image}). Considérons l'application de changement de base de Cohomologie, remarque \ref{cohomology-remark-base-change}. Elle induit une transformation de foncteurs
$$
Lg^* \circ Rf_* \longrightarrow Rf'_* \circ L(g')^*
$$
sur les catégories dérivées de faisceaux à cohomologie quasi-cohérente. On obtient donc une transformation entre les adjoints à droite, en sens opposé :
$$
a \circ Rg_* \longleftarrow Rg'_* \circ a'
$$

\begin{lemma}
\label{duality-lemma-flat-precompose-pus}
Dans le diagramme (\ref{duality-equation-base-change}), supposons que $g$ soit plat, ou plus généralement que $f$ et $g$ soient Tor-indépendants. Alors $a \circ Rg_* \leftarrow Rg'_* \circ a'$ est un isomorphisme.
\end{lemma}

\begin{proof}
Dans ce cas, l'application de changement de base
$Lg^* \circ Rf_* K \longrightarrow Rf'_* \circ L(g')^*K$
est un isomorphisme pour tout $K$ dans $D_\QCoh(\mathcal{O}_X)$ d'après Catégories dérivées des schémas, lemme \ref{perfect-lemma-compare-base-change}. La transformation correspondante entre foncteurs adjoints est donc également un isomorphisme.
\end{proof}

\noindent
Soit $f : X \to Y$ un morphisme de schémas quasi-compacts et quasi-séparés. Soit $V \subset Y$ un sous-schéma ouvert quasi-compact et posons $U = f^{-1}(V)$. On obtient un carré cartésien
$$
\xymatrix{
U \ar[r]_{j'} \ar[d]_{f|_U} & X \ar[d]^f \\
V \ar[r]^j & Y
}
$$
comme en (\ref{duality-equation-base-change}). D'après le lemme \ref{duality-lemma-flat-precompose-pus}, l'application $\xi : a \circ Rj_* \leftarrow Rj'_* \circ a'$ est un isomorphisme, où $a$ et $a'$ sont les adjoints à droite de $Rf_*$ et $R(f|_U)_*$. On obtient une transformation de foncteurs $D_\QCoh(\mathcal{O}_Y) \to D_\QCoh(\mathcal{O}_U)$
\begin{equation}
\label{duality-equation-sheafy}
(j')^* \circ a \to
(j')^* \circ a \circ Rj_* \circ j^* \xrightarrow{\xi^{-1}}
(j')^* \circ Rj'_* \circ a' \circ j^* \to a' \circ j^*
\end{equation}
où la première flèche provient de $\text{id} \to Rj_* \circ j^*$ et la dernière de l'isomorphisme $(j')^* \circ Rj'_* \to \text{id}$. En particulier, (\ref{duality-equation-sheafy}), évaluée en $K$, est un isomorphisme si et seulement si $a(K)|_U \to a(Rj_*(K|_V))|_U$ en est un.

\begin{example}
\label{duality-example-not-supported-on-inverse-image}
Il existe un morphisme fini $f : X \to Y$ de schémas noethériens tel que (\ref{duality-equation-sheafy}) ne soit pas un isomorphisme lorsqu'on l'évalue en un certain $K \in D_{\textit{Coh}}(\mathcal{O}_Y)$. En effet, prenons $X = \Spec(B) \to Y = \Spec(A)$ avec $A = k[x, \epsilon]$, où $k$ est un corps, et prenons $\epsilon^2 = 0$ et $B = k[x] = A/(\epsilon)$. Pour $n \in \mathbf{N}$, posons $M_n = A/(\epsilon, x^n)$. On observe que
$$
\Ext^i_A(B, M_n) = M_n,\quad i \geq 0
$$
car $B$ admet la résolution libre périodique $\ldots \to A \to A \to A$, dont les applications sont données par la multiplication par $\epsilon$. Considérons l'objet
$K = \bigoplus M_n[n] = \prod M_n[n]$
de $D_{\textit{Coh}}(A)$ (égalité dans $D(A)$ d'après Catégories dérivées, lemmes \ref{derived-lemma-direct-sums} et \ref{derived-lemma-products}). Alors $a(K)$ correspond à $R\Hom(B, K)$ d'après l'exemple \ref{duality-example-affine-twisted-inverse-image}, et
$$
H^0(R\Hom(B, K)) = \Ext^0_A(B, K) =
\prod\nolimits_{n \geq 1} \Ext^n_A(B. M_n) = 
\prod\nolimits_{n \geq 1} M_n
$$
d'après ce qui précède. Mais ce module possède des éléments qui ne sont annulés par aucune puissance de $x$, tandis que chaque élément de la cohomologie du complexe $K$ est annulé par une puissance de $x$. Autrement dit, l'application (\ref{duality-equation-sheafy}), avec $V = D(x)$ et $U = D(x)$, évaluée sur le complexe $K$, ne peut pas être un isomorphisme, car $(j')^*(a(K))$ est non nul et $a'(j^*K)$ est nul.
\end{example}

\begin{lemma}
\label{duality-lemma-when-sheafy}
Soit $f : X \to Y$ un morphisme de schémas quasi-compacts et quasi-séparés. Soit $a$ l'adjoint à droite de $Rf_* : D_\QCoh(\mathcal{O}_X) \to D_\QCoh(\mathcal{O}_Y)$. Soit $V \subset Y$ un ouvert quasi-compact d'image inverse $U \subset X$.
\begin{enumerate}
\item Pour tout $Q \in D_\QCoh^+(\mathcal{O}_Y)$ à support dans $Y \setminus V$, l'image $a(Q)$ est à support dans $X \setminus U$ si et seulement si (\ref{duality-equation-sheafy}) est un isomorphisme sur tous les $K$ de $D_\QCoh^+(\mathcal{O}_Y)$.
\item Pour tout $Q \in D_\QCoh(\mathcal{O}_Y)$ à support dans $Y \setminus V$, l'image $a(Q)$ est à support dans $X \setminus U$ si et seulement si (\ref{duality-equation-sheafy}) est un isomorphisme sur tous les $K$ de $D_\QCoh(\mathcal{O}_Y)$.
\item Si $a$ commute aux sommes directes, les conditions équivalentes de (1) impliquent celles de (2).
\end{enumerate}
\end{lemma}

\begin{proof}
Démonstration de (1). Soit $K \in D_\QCoh^+(\mathcal{O}_Y)$. Choisissons un triangle distingué
$$
K \to Rj_*K|_V \to Q \to K[1]
$$
On observe que $Q$ appartient à $D_\QCoh^+(\mathcal{O}_Y)$ (Catégories dérivées des schémas, lemme \ref{perfect-lemma-quasi-coherence-direct-image}) et est à support dans $Y \setminus V$ (Catégories dérivées des schémas, définition \ref{perfect-definition-supported-on}). En appliquant $a$, on obtient un triangle distingué
$$
a(K) \to a(Rj_*K|_V) \to a(Q) \to a(K)[1]
$$
sur $X$. Si $a(Q)$ est à support dans $X \setminus U$, la restriction à $U$ de l'application $a(K)|_U \to a(Rj_*K|_V)|_U$ est un isomorphisme, c'est-à-dire que (\ref{duality-equation-sheafy}) est un isomorphisme sur $K$. La réciproque est immédiate.

\medskip\noindent
La démonstration de (2) est exactement la même que celle de (1).

\medskip\noindent
Démonstration de (3). Supposons satisfaites les conditions équivalentes de (1). Posons $T = Y \setminus V$. Nous utiliserons les notations $D_{\QCoh, T}(\mathcal{O}_Y)$ et $D_{\QCoh, f^{-1}(T)}(\mathcal{O}_X)$ pour les complexes dont les faisceaux de cohomologie sont à support dans $T$ et $f^{-1}(T)$. Puisque $a$ commute aux sommes directes, la sous-catégorie triangulée, saturée et strictement pleine $\mathcal{D}$ dont les objets sont les
$$
\{Q \in D_{\QCoh, T}(\mathcal{O}_Y) \mid
a(Q) \in D_{\QCoh, f^{-1}(T)}(\mathcal{O}_X)\}
$$
est stable par sommes directes, donc par colimites dérivées. D'autre part, la catégorie $D_{\QCoh, T}(\mathcal{O}_Y)$ est engendrée par un objet parfait $E$ (voir Catégories dérivées des schémas, lemme \ref{perfect-lemma-generator-with-support}). Par hypothèse, on a $E \in \mathcal{D}$. D'après Catégories dérivées, lemme \ref{derived-lemma-write-as-colimit}, tout objet $Q$ de $D_{\QCoh, T}(\mathcal{O}_Y)$ est une colimite dérivée d'un système $Q_1 \to Q_2 \to Q_3 \to \ldots$ tel que les cônes des applications de transition soient des sommes directes de décalés de $E$. Par récurrence, on voit que $Q_n \in \mathcal{D}$ pour tout $n$, puis que $Q$ appartient à $\mathcal{D}$. Les conditions équivalentes de (2) sont donc satisfaites.
\end{proof}

\begin{lemma}
\label{duality-lemma-proper-noetherian}
Soit $Y$ un schéma quasi-compact et quasi-séparé. Soit $f : X \to Y$ un morphisme propre. Si\footnote{Cette démonstration s'applique aux morphismes de schémas quasi-compacts et quasi-séparés tels que $Rf_*P$ soit pseudo-cohérent pour tout $P$ parfait sur $X$. Un théorème de Kiehl \cite{Kiehl} montre aisément qu'il en est ainsi si $f$ est propre et pseudo-cohérent. C'est le degré de généralité approprié pour ce lemme et certains autres résultats de ce chapitre.}
\begin{enumerate}
\item $f$ est plat et de présentation finie ; ou
\item $Y$ est noethérien,
\end{enumerate}
alors les conditions équivalentes du lemme \ref{duality-lemma-when-sheafy}, (1), sont satisfaites pour tous les ouverts quasi-compacts $V$ de $Y$.
\end{lemma}

\begin{proof}
Soit $Q \in D^+_\QCoh(\mathcal{O}_Y)$ à support dans $Y \setminus V$. Pour obtenir une contradiction, supposons que $a(Q)$ ne soit pas à support dans $X \setminus U$. On peut alors trouver un complexe parfait $P_U$ sur $U$ et une application non nulle $P_U \to a(Q)|_U$ (cela résulte de Catégories dérivées des schémas, théorème \ref{perfect-theorem-bondal-van-den-Bergh}). En utilisant Catégories dérivées des schémas, lemme \ref{perfect-lemma-lift-map-from-perfect-complex-with-support}, on peut donc supposer qu'il existe un complexe parfait $P$ sur $X$ et une application $P \to a(Q)$ dont la restriction à $U$ est non nulle. Par définition de $a$, cette application est adjointe à une application $Rf_*P \to Q$.

\medskip\noindent
Le complexe $Rf_*P$ est pseudo-cohérent. Dans le cas (1), cela résulte de Catégories dérivées des schémas, lemme \ref{perfect-lemma-flat-proper-pseudo-coherent-direct-image-general}. Dans le cas (2), cela résulte de Catégories dérivées des schémas, lemmes \ref{perfect-lemma-direct-image-coherent} et \ref{perfect-lemma-identify-pseudo-coherent-noetherian}. On peut donc appliquer Catégories dérivées des schémas, lemme \ref{perfect-lemma-map-from-pseudo-coherent-to-complex-with-support}, pour obtenir une application $I \to \mathcal{O}_Y$ entre complexes parfaits, dont la restriction à $V$ est un isomorphisme et telle que la composée $I \otimes^\mathbf{L}_{\mathcal{O}_Y} Rf_*P \to Rf_*P \to Q$ soit nulle. D'après Catégories dérivées des schémas, lemme \ref{perfect-lemma-cohomology-base-change}, on a $I \otimes^\mathbf{L}_{\mathcal{O}_Y} Rf_*P =
Rf_*(Lf^*I \otimes^\mathbf{L}_{\mathcal{O}_X} P)$. On en conclut que la composée
$$
Lf^*I \otimes^\mathbf{L}_{\mathcal{O}_X} P \to P \to a(Q)
$$
est nulle. Cependant, sa restriction à $U$ est l'application $P|_U \to a(Q)|_U$, supposée non nulle. Cette contradiction achève la démonstration.
\end{proof}












\section{Adjoint à droite de l'image directe et changement de base, I}
\label{duality-section-base-change-map}

\begingroup\emergencystretch=16pt
\noindent
L'application (\ref{duality-equation-sheafy}) est un cas particulier d'une application de changement de base. Supposons en effet donné un diagramme cartésien
$$
\xymatrix{
X' \ar[r]_{g'} \ar[d]_{f'} & X \ar[d]^f \\
Y' \ar[r]^g & Y
}
$$
de schémas quasi-compacts et quasi-séparés, c'est-à-dire un diagramme comme en (\ref{duality-equation-base-change}). Supposons $f$ et $g$ {\bf Tor-indépendants}. On peut alors considérer le morphisme de foncteurs
$D_\QCoh(\mathcal{O}_Y) \to D_\QCoh(\mathcal{O}_{X'})$
défini par la composée
\begin{equation}
\label{duality-equation-base-change-map}
L(g')^* \circ a \to
L(g')^* \circ a \circ Rg_* \circ Lg^* \leftarrow
L(g')^* \circ Rg'_* \circ a' \circ Lg^* \to a' \circ Lg^*
\end{equation}
La première flèche provient de l'application d'adjonction $\text{id} \to Rg_* Lg^*$, et la dernière de l'application d'adjonction $L(g')^*Rg'_* \to \text{id}$. L'hypothèse de Tor-indépendance est nécessaire pour inverser la flèche du milieu, voir le lemme \ref{duality-lemma-flat-precompose-pus}. On peut aussi voir (\ref{duality-equation-base-change-map}), par adjonction entre $L(g')^*$ et $R(g')_*$, comme une transformation naturelle
$$
a \to a \circ Rg_* \circ Lg^* \leftarrow Rg'_* \circ a' \circ Lg^*
$$
où, là encore, la deuxième flèche est inversible. Si $M \in D_\QCoh(\mathcal{O}_X)$ et $K \in D_\QCoh(\mathcal{O}_Y)$, cette application est donnée sur les foncteurs de Yoneda par
\begin{align*}
\Hom_X(M, a(K))
& =
\Hom_Y(Rf_*M, K) \\
& \to
\Hom_Y(Rf_*M, Rg_* Lg^*K) \\
& =
\Hom_{Y'}(Lg^*Rf_*M, Lg^*K) \\
& \leftarrow
\Hom_{Y'}(Rf'_* L(g')^*M, Lg^*K) \\
& =
\Hom_{X'}(L(g')^*M, a'(Lg^*K)) \\
& =
\Hom_X(M, Rg'_*a'(Lg^*K))
\end{align*}
(où la flèche dirigée vers la gauche est inversible par le théorème de changement de base de Catégories dérivées des schémas, lemme \ref{perfect-lemma-compare-base-change}), ce qui explicite quelque peu la construction.\par\endgroup

\medskip\noindent
Nous montrons d'abord dans cette section que l'application de changement de base satisfait certaines compatibilités naturelles relatives à la juxtaposition de carrés, analogues à celles de Cohomologie, remarques \ref{cohomology-remark-compose-base-change} et \ref{cohomology-remark-compose-base-change-horizontal}, pour l'application usuelle de changement de base. On pourra omettre la suite de cette section en première lecture.

\begin{lemma}
\label{duality-lemma-compose-base-change-maps}
Considérons un diagramme commutatif
$$
\xymatrix{
X' \ar[r]_k \ar[d]_{f'} & X \ar[d]^f \\
Y' \ar[r]^l \ar[d]_{g'} & Y \ar[d]^g \\
Z' \ar[r]^m & Z
}
$$
de schémas quasi-compacts et quasi-séparés dont les deux carrés sont cartésiens et où $f$ et $l$, ainsi que $g$ et $m$, sont Tor-indépendants. Alors la composée des applications (\ref{duality-equation-base-change-map}) associées aux deux carrés est l'application de changement de base associée au rectangle extérieur (voir la démonstration pour un énoncé précis).
\end{lemma}

\begin{proof}
Les hypothèses impliquent que $g \circ f$ et $m$ sont Tor-indépendants (les détails sont omis), de sorte que l'énoncé a un sens. Dans cette démonstration, nous écrivons $k^*$ au lieu de $Lk^*$ et $f_*$ au lieu de $Rf_*$. Soient $a$, $b$ et $c$ les adjoints à droite du lemme \ref{duality-lemma-twisted-inverse-image} pour $f$, $g$ et $g \circ f$, et de même pour les versions munies d'un prime. La flèche correspondant au carré supérieur est la composée
$$
\gamma_{top} :
k^* \circ a \to k^* \circ a \circ l_* \circ l^*
\xleftarrow{\xi_{top}} k^* \circ k_* \circ a' \circ l^* \to a' \circ l^*
$$
où $\xi_{top} : k_* \circ a' \to a \circ l_*$ est un isomorphisme (donc inversible), qui est la flèche « duale » de l'application de changement de base $l^* \circ f_* \to f'_* \circ k^*$. Les flèches extérieures proviennent des applications canoniques $1 \to l_* \circ l^*$ et $k^* \circ k_* \to 1$. De même, pour le deuxième carré, on a
$$
\gamma_{bot} :
l^* \circ b \to l^* \circ b \circ m_* \circ m^*
\xleftarrow{\xi_{bot}} l^* \circ l_* \circ b' \circ m^* \to b' \circ m^*
$$
Pour le rectangle extérieur, on obtient
$$
\gamma_{rect} :
k^* \circ c \to k^* \circ c \circ m_* \circ m^*
\xleftarrow{\xi_{rect}} k^* \circ k_* \circ c' \circ m^* \to c' \circ m^*
$$
On a $(g \circ f)_* = g_* \circ f_*$, donc $c = a \circ b$, et de même $c' = a' \circ b'$. Le lemme affirme que $\gamma_{rect}$ est égale à la composée
$$
k^* \circ c = k^* \circ a \circ b \xrightarrow{\gamma_{top}}
a' \circ l^* \circ b \xrightarrow{\gamma_{bot}}
a' \circ b' \circ m^* = c' \circ m^*
$$
Pour le voir, considérons le diagramme suivant :
$$
\xymatrix{
& & k^* \circ a \circ b \ar[d] \ar[lldd] \\
& & k^* \circ a \circ l_* \circ l^* \circ b \ar[ld] \\
k^* \circ a \circ b \circ m_* \circ m^* \ar[r] &
k^* \circ a \circ l_* \circ l^* \circ b \circ m_* \circ m^* &
k^* \circ k_* \circ a' \circ l^* \circ b \ar[u]_{\xi_{top}} \ar[d] \ar[ld] \\
& k^*\circ k_* \circ a' \circ l^* \circ b \circ m_* \circ m^*
\ar[u]_{\xi_{top}} \ar[rd] &
a' \circ l^* \circ b \ar[d] \\
k^* \circ k_* \circ a' \circ b' \circ m^* \ar[uu]_{\xi_{rect}} \ar[ddrr] &
k^*\circ k_* \circ a' \circ l^* \circ l_* \circ b' \circ m^*
\ar[u]_{\xi_{bot}} \ar[l] \ar[dr] &
a' \circ l^* \circ b \circ m_* \circ m^* \\
& & a' \circ l^* \circ l_* \circ b' \circ m^* \ar[u]_{\xi_{bot}} \ar[d] \\
& & a' \circ b' \circ m^*
}
$$
En descendant par le côté droit, on obtient la composée ; en descendant par le côté gauche, on obtient $\gamma_{rect}$. Tous les quadrilatères de la partie droite de ce diagramme sont commutatifs d'après Catégories, lemme \ref{categories-lemma-properties-2-cat-cats}, ou, plus simplement, d'après la discussion précédant Catégories, définition \ref{categories-definition-horizontal-composition}. Il suffit donc de montrer que le diagramme
$$
\xymatrix{
a \circ l_* \circ l^* \circ b \circ m_* &
a \circ b \circ m_* \ar[l] \\
k_* \circ a' \circ l^* \circ b \circ m_* \ar[u]_{\xi_{top}} & \\
k_* \circ a' \circ l^* \circ l_* \circ b' \ar[u]_{\xi_{bot}} \ar[r] &
k_* \circ a' \circ b' \ar[uu]_{\xi_{rect}}
}
$$
devient commutatif lorsqu'on inverse les flèches $\xi_{top}$, $\xi_{bot}$ et $\xi_{rect}$ (ce qui ne revient pas à demander que le diagramme soit commutatif). Or le diagramme
$$
\xymatrix{
& a \circ l_* \circ l^* \circ b \circ m_* \\
a \circ l_* \circ l^* \circ l_* \circ b'
\ar[ru]^{\xi_{bot}} & &
k_* \circ a' \circ l^* \circ b \circ m_* \ar[ul]_{\xi_{top}} \\
& k_* \circ a' \circ l^* \circ l_* \circ b'
\ar[ul]^{\xi_{top}} \ar[ur]_{\xi_{bot}}
}
$$
est commutatif d'après Catégories, lemme \ref{categories-lemma-properties-2-cat-cats}. Puisque les diagrammes
$$
\vcenter{
\xymatrix{
a \circ l_* \circ l^* \circ b \circ m_* & a \circ b \circ m \ar[l] \\
a \circ l_* \circ l^* \circ l_* \circ b' \ar[u] &
a \circ l_* \circ b' \ar[l] \ar[u]
}
}
\quad\text{et}\quad
\vcenter{
\xymatrix{
a \circ l_* \circ l^* \circ l_* \circ b' \ar[r] & a \circ l_* \circ b' \\
k_* \circ a' \circ l^* \circ l_* \circ b' \ar[u] \ar[r] &
k_* \circ a' \circ b' \ar[u]
}
}
$$
sont commutatifs (voir les références citées) et que la composée de $l_* \to l_* \circ l^* \circ l_* \to l_*$ est l'identité, il suffit de montrer que
$$
k \circ a' \circ b' \xrightarrow{\xi_{bot}} a \circ l_* \circ b
\xrightarrow{\xi_{top}} a \circ b \circ m_*
$$
est égale à $\xi_{rect}$ (via les identifications $a \circ b = c$ et $a' \circ b' = c'$). C'est l'énoncé dual de Cohomologie, remarque \ref{cohomology-remark-compose-base-change}, ce qui achève la démonstration.
\end{proof}

\begin{lemma}
\label{duality-lemma-compose-base-change-maps-horizontal}
Considérons un diagramme commutatif
$$
\xymatrix{
X'' \ar[r]_{g'} \ar[d]_{f''} & X' \ar[r]_g \ar[d]_{f'} & X \ar[d]^f \\
Y'' \ar[r]^{h'} & Y' \ar[r]^h & Y
}
$$
de schémas quasi-compacts et quasi-séparés dont les deux carrés sont cartésiens et où $f$ et $h$, ainsi que $f'$ et $h'$, sont Tor-indépendants. Alors la composée des applications (\ref{duality-equation-base-change-map}) associées aux deux carrés est l'application de changement de base associée au rectangle extérieur (voir la démonstration pour un énoncé précis).
\end{lemma}

\begin{proof}
Les hypothèses impliquent que $f$ et $h \circ h'$ sont Tor-indépendants (les détails sont omis), de sorte que l'énoncé a un sens. Dans cette démonstration, nous écrivons $g^*$ au lieu de $Lg^*$ et $f_*$ au lieu de $Rf_*$. Soient $a$, $a'$ et $a''$ les adjoints à droite du lemme \ref{duality-lemma-twisted-inverse-image} pour $f$, $f'$ et $f''$. La flèche correspondant au carré de droite est la composée
$$
\gamma_{right} :
g^* \circ a \to g^* \circ a \circ h_* \circ h^*
\xleftarrow{\xi_{right}} g^* \circ g_* \circ a' \circ h^* \to a' \circ h^*
$$
où $\xi_{right} : g_* \circ a' \to a \circ h_*$ est un isomorphisme (donc inversible), qui est la flèche « duale » de l'application de changement de base $h^* \circ f_* \to f'_* \circ g^*$. Les flèches extérieures proviennent des applications canoniques $1 \to h_* \circ h^*$ et $g^* \circ g_* \to 1$. De même, pour le carré de gauche, on a
$$
\gamma_{left} :
(g')^* \circ a' \to (g')^* \circ a' \circ (h')_* \circ (h')^*
\xleftarrow{\xi_{left}}
(g')^* \circ (g')_* \circ a'' \circ (h')^* \to a'' \circ (h')^*
$$
Pour le rectangle extérieur, on obtient
$$
\gamma_{rect} :
k^* \circ a \to
k^* \circ a \circ m_* \circ m^* \xleftarrow{\xi_{rect}}
k^* \circ k_* \circ a'' \circ m^* \to
a'' \circ m^*
$$
où $k = g \circ g'$ et $m = h \circ h'$. On a $k^* = (g')^* \circ g^*$ et $m^* = (h')^* \circ h^*$. Le lemme affirme que $\gamma_{rect}$ est égale à la composée
$$
k^* \circ a =
(g')^* \circ g^* \circ a \xrightarrow{\gamma_{right}}
(g')^* \circ a' \circ h^* \xrightarrow{\gamma_{left}}
a'' \circ (h')^* \circ h^* = a'' \circ m^*
$$
Pour le voir, considérons le diagramme suivant :
$$
\xymatrix{
& (g')^* \circ g^* \circ a \ar[d] \ar[ddl] \\
& (g')^* \circ g^* \circ a \circ h_* \circ h^* \ar[ld] \\
(g')^* \circ g^* \circ a \circ h_* \circ (h')_* \circ (h')^* \circ h^* &
(g')^* \circ g^* \circ g_* \circ a' \circ h^*
\ar[u]_{\xi_{right}} \ar[d] \ar[ld] \\
(g')^* \circ g^* \circ g_* \circ a' \circ (h')_* \circ (h')^* \circ h^*
\ar[u]_{\xi_{right}} \ar[dr] &
(g')^* \circ a' \circ h^* \ar[d] \\
(g')^* \circ g^* \circ g_* \circ (g')_* \circ a'' \circ (h')^* \circ h^*
\ar[u]_{\xi_{left}} \ar[ddr] \ar[dr] &
(g')^* \circ a' \circ (h')_* \circ (h')^* \circ h^* \\
& (g')^*\circ (g')_* \circ a'' \circ (h')^* \circ h^*
\ar[u]_{\xi_{left}} \ar[d] \\
& a'' \circ (h')^* \circ h^*
}
$$
En descendant par le côté droit, on obtient la composée ; en descendant par le côté gauche, on obtient $\gamma_{rect}$. Tous les quadrilatères de la partie droite de ce diagramme sont commutatifs d'après Catégories, lemme \ref{categories-lemma-properties-2-cat-cats}, ou, plus simplement, d'après la discussion précédant Catégories, définition \ref{categories-definition-horizontal-composition}. Il suffit donc de montrer que
$$
g_* \circ (g')_* \circ a'' \xrightarrow{\xi_{left}}
g_* \circ a' \circ (h')_* \xrightarrow{\xi_{right}}
a \circ h_* \circ (h')_*
$$
est égale à $\xi_{rect}$. C'est l'énoncé dual de Cohomologie, remarque \ref{cohomology-remark-compose-base-change-horizontal}, ce qui achève la démonstration.
\end{proof}

\begin{remark}
\label{duality-remark-going-around}
Considérons un diagramme commutatif
$$
\xymatrix{
X'' \ar[r]_{k'} \ar[d]_{f''} & X' \ar[r]_k \ar[d]_{f'} & X \ar[d]^f \\
Y'' \ar[r]^{l'} \ar[d]_{g''} & Y' \ar[r]^l \ar[d]_{g'} & Y \ar[d]^g \\
Z'' \ar[r]^{m'} & Z' \ar[r]^m & Z
}
$$
de schémas quasi-compacts et quasi-séparés dont tous les carrés sont cartésiens, et où $(f, l)$, $(g, m)$, $(f', l')$, $(g', m')$ sont des couples de morphismes Tor-indépendants. Soient $a$, $a'$, $a''$, $b$, $b'$, $b''$ les adjoints à droite du lemme \ref{duality-lemma-twisted-inverse-image} pour $f$, $f'$, $f''$, $g$, $g'$, $g''$. Désignons les carrés du diagramme par $A$, $B$, $C$, $D$ comme suit :
$$
\begin{matrix}
A & B \\
C & D
\end{matrix}
$$
Les applications (\ref{duality-equation-base-change-map}) associées aux carrés sont alors (en utilisant $k^* = Lk^*$, etc.)
$$
\begin{matrix}
\gamma_A : (k')^* \circ a' \to a'' \circ (l')^* &
\gamma_B : k^* \circ a \to a' \circ l^* \\
\gamma_C : (l')^* \circ b' \to b'' \circ (m')^* &
\gamma_D : l^* \circ b \to b' \circ m^*
\end{matrix}
$$
Pour les rectangles $2 \times 1$ et $1 \times 2$, on dispose de quatre autres applications de changement de base
$$
\begin{matrix}
\gamma_{A + B} : (k \circ k')^* \circ a \to a'' \circ (l \circ l')^* \\
\gamma_{C + D} : (l \circ l')^* \circ b \to b'' \circ (m \circ m')^* \\
\gamma_{A + C} : (k')^* \circ (a' \circ b') \to (a'' \circ b'') \circ (m')^* \\
\gamma_{B + D} : k^* \circ (a \circ b) \to (a' \circ b') \circ m^*
\end{matrix}
$$
D'après le lemme \ref{duality-lemma-compose-base-change-maps-horizontal}, on a
$$
\gamma_{A + B} = \gamma_A \circ \gamma_B, \quad
\gamma_{C + D} = \gamma_C \circ \gamma_D
$$
et, d'après le lemme \ref{duality-lemma-compose-base-change-maps}, on a
$$
\gamma_{A + C} = \gamma_C \circ \gamma_A, \quad
\gamma_{B + D} = \gamma_D \circ \gamma_B
$$
Il serait ici plus correct d'écrire $\gamma_{A + B} = (\gamma_A \star \text{id}_{l^*}) \circ
(\text{id}_{(k')^*} \star \gamma_B)$ avec les notations de Catégories, section \ref{categories-section-formal-cat-cat}, et de même pour les autres expressions. Nous conservons toutefois l'abus de notation des démonstrations des lemmes \ref{duality-lemma-compose-base-change-maps} et \ref{duality-lemma-compose-base-change-maps-horizontal}, en omettant les produits $\star$ avec les identités : on peut reconstituer ceux-ci dès lors que la source et le but de la transformation sont connus. Cela étant, on obtient a priori deux transformations
$$
(k')^* \circ k^* \circ a \circ b
\longrightarrow
a'' \circ b'' \circ (m')^* \circ m^*
$$
à savoir
$$
\gamma_C \circ \gamma_A \circ \gamma_D \circ \gamma_B =
\gamma_{A + C} \circ \gamma_{B + D}
$$
et
$$
\gamma_C \circ \gamma_D \circ \gamma_A \circ \gamma_B =
\gamma_{C + D} \circ \gamma_{A + B}
$$
Cette remarque a pour objet de signaler que ces transformations sont égales. Pour le voir, il suffit de montrer que
$$
\xymatrix{
(k')^* \circ a' \circ l^* \circ b \ar[r]_{\gamma_D} \ar[d]_{\gamma_A} &
(k')^* \circ a' \circ b' \circ m^* \ar[d]^{\gamma_A} \\
a'' \circ (l')^* \circ l^* \circ b \ar[r]^{\gamma_D} &
a'' \circ (l')^* \circ b' \circ m^*
}
$$
est commutatif. Cela résulte de Catégories, lemme \ref{categories-lemma-properties-2-cat-cats}, ou, plus simplement, de la discussion précédant Catégories, définition \ref{categories-definition-horizontal-composition}.
\end{remark}







\section{Adjoint à droite de l'image directe et changement de base, II}
\label{duality-section-base-change-II}

\noindent
Nous montrons dans cette section que l'application de changement de base de la section \ref{duality-section-base-change-map} est un isomorphisme dans certains cas. Observons d'abord qu'il suffit de le vérifier sur des ouverts affines, pourvu que la formation de l'adjoint à droite de l'image directe commute à la restriction aux ouverts.

\begin{remark}
\label{duality-remark-check-over-affines}
Considérons un diagramme cartésien
$$
\xymatrix{
X' \ar[r]_{g'} \ar[d]_{f'} & X \ar[d]^f \\
Y' \ar[r]^g & Y
}
$$
de schémas quasi-compacts et quasi-séparés, avec $(g, f)$ Tor-indépendants. Soient $V \subset Y$ et $V' \subset Y'$ des ouverts affines tels que $g(V') \subset V$. Formons les diagrammes cartésiens
$$
\vcenter{
\xymatrix{
U \ar[r] \ar[d] & X \ar[d] \\
V \ar[r] & Y
}
}
\quad\text{et}\quad
\vcenter{
\xymatrix{
U' \ar[r] \ar[d] & X' \ar[d] \\
V' \ar[r] & Y'
}
}
$$
Supposons que (\ref{duality-equation-sheafy}), pour $K$ et le premier diagramme, et (\ref{duality-equation-sheafy}), pour $Lg^*K$ et le deuxième diagramme, soient des isomorphismes. Alors la restriction de l'application de changement de base (\ref{duality-equation-base-change-map})
$$
L(g')^*a(K) \longrightarrow a'(Lg^*K)
$$
à $U'$ est isomorphe à l'application de changement de base (\ref{duality-equation-base-change-map}) pour $K|_V$ et le diagramme cartésien
$$
\xymatrix{
U' \ar[r] \ar[d] & U \ar[d] \\
V' \ar[r] & V
}
$$
Cela résulte du fait que (\ref{duality-equation-sheafy}) est un cas particulier de l'application de changement de base (\ref{duality-equation-base-change-map}), et de la compatibilité de ces applications avec la juxtaposition horizontale de carrés, voir le lemme \ref{duality-lemma-compose-base-change-maps-horizontal}. Pour vérifier que la restriction de l'application de changement de base à $U'$ est un isomorphisme, il suffit donc de travailler avec le dernier diagramme.
\end{remark}

\begin{lemma}
\label{duality-lemma-more-base-change}
Dans le diagramme (\ref{duality-equation-base-change}), supposons que
\begin{enumerate}
\item $g : Y' \to Y$ soit un morphisme de schémas affines ;
\item $f : X \to Y$ soit propre ;
\item $f$ et $g$ soient Tor-indépendants.
\end{enumerate}
Alors l'application de changement de base (\ref{duality-equation-base-change-map}) induit un isomorphisme
$$
L(g')^*a(K) \longrightarrow a'(Lg^*K)
$$
dans les cas suivants :
\begin{enumerate}
\item pour tout $K \in D_\QCoh(\mathcal{O}_Y)$ si $f$ est plat de présentation finie ;
\item pour tout $K \in D_\QCoh(\mathcal{O}_Y)$ si $f$ est parfait et $Y$ noethérien ;
\item pour $K \in D_\QCoh^+(\mathcal{O}_Y)$ si $g$ est de Tor-dimension finie et $Y$ noethérien.
\end{enumerate}
\end{lemma}

\begin{proof}
Écrivons $Y = \Spec(A)$ et $Y' = \Spec(A')$. Étant obtenu par changement de base d'un morphisme affine, le morphisme $g'$ est affine. Soit $M$ un générateur parfait de $D_\QCoh(\mathcal{O}_X)$, voir Catégories dérivées des schémas, théorème \ref{perfect-theorem-bondal-van-den-Bergh}. Alors $L(g')^*M$ est un générateur de $D_\QCoh(\mathcal{O}_{X'})$, voir Catégories dérivées des schémas, remarque \ref{perfect-remark-pullback-generator}. Il suffit donc de montrer que (\ref{duality-equation-base-change-map}) induit un isomorphisme
\begin{equation}
\label{duality-equation-iso}
R\Hom_{X'}(L(g')^*M, L(g')^*a(K))
\longrightarrow
R\Hom_{X'}(L(g')^*M, a'(Lg^*K))
\end{equation}
entre les complexes Hom globaux, voir Cohomologie, section \ref{cohomology-section-global-RHom}, car cela impliquera que le cône de $L(g')^*a(K) \to a'(Lg^*K)$ est nul. La démonstration s'organise comme suit : nous montrerons d'abord que ces complexes Hom sont isomorphes, puis, dans la dernière partie, que l'isomorphisme est induit par (\ref{duality-equation-iso}).

\medskip\noindent
Le membre de gauche. Puisque $M$ est parfait, l'application canonique
$$
R\Hom_X(M, a(K)) \otimes^\mathbf{L}_A A'
\longrightarrow
R\Hom_{X'}(L(g')^*M, L(g')^*a(K))
$$
est un isomorphisme d'après Catégories dérivées des schémas, lemme \ref{perfect-lemma-affine-morphism-and-hom-out-of-perfect}. En la combinant avec l'isomorphisme
$R\Hom_Y(Rf_*M, K) = R\Hom_X(M, a(K))$
du lemme \ref{duality-lemma-iso-global-hom}, on obtient que le membre de gauche est égal à $R\Hom_Y(Rf_*M, K) \otimes^\mathbf{L}_A A'$.

\medskip\noindent
Le membre de droite. On utilise d'abord l'isomorphisme
$$
R\Hom_{X'}(L(g')^*M, a'(Lg^*K)) = R\Hom_{Y'}(Rf'_*L(g')^*M, Lg^*K)
$$
du lemme \ref{duality-lemma-iso-global-hom}. On utilise ensuite le fait que l'application de changement de base $Lg^*Rf_*M \to Rf'_*L(g')^*M$ est un isomorphisme d'après Catégories dérivées des schémas, lemme \ref{perfect-lemma-compare-base-change}. On peut donc réécrire ce membre sous la forme
\par\noindent\hfil
$R\Hom_{Y'}(Lg^*Rf_*M, Lg^*K)$.
\hfil\par\noindent
Puisque $Y$, $Y'$ sont affines et que $K$, $Rf_*M$ appartiennent à $D_\QCoh(\mathcal{O}_Y)$ (Catégories dérivées des schémas, lemme \ref{perfect-lemma-quasi-coherence-direct-image}), on dispose d'une application canonique
$$
\beta :
R\Hom_Y(Rf_*M, K) \otimes^\mathbf{L}_A A'
\longrightarrow
R\Hom_{Y'}(Lg^*Rf_*M, Lg^*K)
$$
dans $D(A')$. C'est la flèche de Compléments d'algèbre, équation (\ref{more-algebra-equation-base-change-RHom}), où l'on utilise Catégories dérivées des schémas, lemmes \ref{perfect-lemma-affine-compare-bounded} et \ref{perfect-lemma-quasi-coherence-internal-hom}, pour passer au cadre algébrique et en revenir.
\begin{enumerate}
\item Si $f$ est plat de présentation finie, le complexe $Rf_*M$ est parfait sur $Y$ d'après Catégories dérivées des schémas, lemme \ref{perfect-lemma-flat-proper-perfect-direct-image-general}, et $\beta$ est un isomorphisme d'après Compléments d'algèbre, lemme \ref{more-algebra-lemma-base-change-RHom}, (1).
\item Si $f$ est parfait et $Y$ noethérien, le complexe $Rf_*M$ est parfait sur $Y$ d'après Compléments sur les morphismes, lemme \ref{more-morphisms-lemma-perfect-proper-perfect-direct-image}, et $\beta$ est un isomorphisme comme précédemment.
\item Si $g$ est de Tor-dimension finie et $Y$ est noethérien, le complexe $Rf_*M$ est pseudo-cohérent sur $Y$ (Catégories dérivées des schémas, lemmes \ref{perfect-lemma-direct-image-coherent} et \ref{perfect-lemma-identify-pseudo-coherent-noetherian}), et $\beta$ est un isomorphisme d'après Compléments d'algèbre, lemme \ref{more-algebra-lemma-base-change-RHom}, (4).
\end{enumerate}
On obtient donc le même résultat qu'au paragraphe précédent.

\medskip\noindent
\begingroup\emergencystretch=2em
Dans la suite de la démonstration, nous montrons que les identifications des membres de gauche et de droite de (\ref{duality-equation-iso}) données aux deuxième et troisième paragraphes sont effectivement réalisées par (\ref{duality-equation-iso}). Pour alléger les formules, nous utiliserons $(-, -)_X = R\Hom_X(-, -)$, écrirons $- \otimes A'$ au lieu de $- \otimes_A^\mathbf{L} A'$, et poserons les abréviations $g^* = Lg^*$ et $f_* = Rf_*$. Considérons le diagramme commutatif suivant :\par\endgroup
\begingroup\small
$$
\xymatrix{
((g')^*M, (g')^*a(K))_{X'} \ar[d] &
(M, a(K))_X \otimes A' \ar[l]^-\alpha \ar[d] &
(f_*M, K)_Y \otimes A' \ar@{=}[l] \ar[d] \\
((g')^*M, (g')^*a(g_*g^*K))_{X'} &
(M, a(g_*g^*K))_X \otimes A' \ar[l]^-\alpha &
(f_*M, g_*g^*K)_Y \otimes A' \ar@{=}[l] \ar@/_4pc/[dd]_{\mu'} \\
((g')^*M, (g')^*g'_*a'(g^*K))_{X'} \ar[u] \ar[d] &
(M, g'_*a'(g^*K))_X \otimes A' \ar[u] \ar[l]^-\alpha \ar[ld]^\mu &
(f_*M, K) \otimes A' \ar[d]^\beta \\
((g')^*M, a'(g^*K))_{X'} &
(f'_*(g')^*M, g^*K)_{Y'} \ar@{=}[l] \ar[r] &
(g^*f_*M, g^*K)_{Y'}
}
$$
\endgroup
Les flèches notées $\alpha$ sont les applications de Catégories dérivées des schémas, lemme \ref{perfect-lemma-affine-morphism-and-hom-out-of-perfect}, pour le diagramme de sommets $X', X, Y', Y$. La partie supérieure du diagramme est commutative, car les flèches horizontales sont fonctorielles en leurs arguments. Les flèches verticales centrales proviennent de la transformation inversible $g'_* \circ a' \to a \circ g_*$ du lemme \ref{duality-lemma-flat-precompose-pus} ; le carré central est donc commutatif. La descente par le côté gauche est (\ref{duality-equation-iso}). Les flèches horizontales supérieures fournissent les identifications utilisées au deuxième paragraphe de la démonstration. Les flèches horizontales inférieures, y compris $\beta$, fournissent celles du troisième paragraphe. Étant donnés $E \in D(A)$, $E' \in D(A')$ et $c : E \to E'$ dans $D(A)$, nous noterons $\mu_c : E \otimes A' \to E'$ l'application induite par $c$ et par l'adjonction entre restriction et changement de base ; lorsque $c$ est clair, nous écrirons $\mu = \mu_c$, en omettant donc $c$. L'application $\mu$ du diagramme est de cette forme, avec $c$ donné par l'identification
$(M, g'_*a(g^*K))_X = ((g')^*M, a'(g^*K))_{X'}$
; le triangle faisant intervenir $\mu$ est commutatif d'après Catégories dérivées des schémas, remarque \ref{perfect-remark-multiplication-map}.

\medskip\noindent
Observons que
$$
\xymatrix{
(M, a(g_*g^*K))_X &
(f_*M, g_* g^*K)_Y \ar@{=}[l] &
(g^*f_*M, g^*K)_{Y'} \ar@{=}[l] \\
(M, g'_* a'(g^*K))_X \ar[u] &
((g')^*M, a'(g^*K))_{X'} \ar@{=}[l] &
(f'_*(g')^*M, g^*K)_{Y'} \ar@{=}[l] \ar[u]
}
$$
est commutatif par définition même de la transformation $g'_* \circ a' \to a \circ g_*$. En prenant $\mu'$ comme ci-dessus, correspondant à l'identification $(f_*M, g_*g^*K)_X = (g^*f_*M, g^*K)_{Y'}$, l'hexagone est lui aussi commutatif. Il suffit donc de montrer que $\beta$ est égale à la composée de
$(f_*M, K)_Y \otimes A' \to (f_*M, g_*g^*K)_X \otimes A'$
et de $\mu'$. Pour cela, il suffit de prouver que les deux applications induites $(f_*M, K)_Y \to (g^*f_*M, g^*K)_{Y'}$ sont égales. Autrement dit, il suffit de montrer que le diagramme
$$
\xymatrix{
R\Hom_A(E, K) \ar[rr]_{\text{induite par }\beta} \ar[rd] & &
R\Hom_{A'}(E \otimes_A^\mathbf{L} A', K \otimes_A^\mathbf{L} A') \\
& R\Hom_A(E, K \otimes_A^\mathbf{L} A') \ar[ru]
}
$$
est commutatif pour tout $E, K \in D(A)$. Comme c'est ainsi que $\beta$ est construit dans Compléments d'algèbre, section \ref{more-algebra-section-base-change-RHom}, la démonstration est achevée.
\end{proof}








\section{Adjoint à droite de l'image directe et morphismes trace}
\label{duality-section-trace}

\noindent
Soit $f : X \to Y$ un morphisme de schémas quasi-compacts et quasi-séparés. Soit $a : D_\QCoh(\mathcal{O}_Y) \to D_\QCoh(\mathcal{O}_X)$ l'adjoint à droite du lemme \ref{duality-lemma-twisted-inverse-image}. D'après Catégories, section \ref{categories-section-adjoint}, on obtient une transformation de foncteurs
$$
\text{Tr}_f : Rf_* \circ a \longrightarrow \text{id}
$$
Le morphisme correspondant $\text{Tr}_{f, K} : Rf_*a(K) \longrightarrow K$, pour $K \in D_\QCoh(\mathcal{O}_Y)$, est parfois appelé {\it morphisme trace}. Il se caractérise par le fait que la bijection
$$
\Hom_X(L, a(K)) \longrightarrow \Hom_Y(Rf_*L, K)
$$
pour $L \in D_\QCoh(\mathcal{O}_X)$, qui caractérise l'adjoint à droite, est donnée par
$$
\varphi \longmapsto \text{Tr}_{f, K} \circ Rf_*\varphi
$$
L'application (\ref{duality-equation-sheafy-trace})
$$
Rf_*R\SheafHom_{\mathcal{O}_X}(L, a(K))
\longrightarrow
R\SheafHom_{\mathcal{O}_Y}(Rf_*L, K)
$$
s'obtient par composition avec $\text{Tr}_{f, K}$. Tout morphisme trace considéré dans cette section sera un cas particulier de celui-ci. Avant d'aborder quelques cas particuliers, montrons que la formation de la trace commute au changement de base.

\begin{lemma}[Morphisme trace et changement de base]
\label{duality-lemma-trace-map-and-base-change}
Supposons donné un diagramme (\ref{duality-equation-base-change}) où $f$ et $g$ sont Tor-indépendants. Alors les applications $1 \star \text{Tr}_f : Lg^* \circ Rf_* \circ a \to Lg^*$ et $\text{Tr}_{f'} \star 1 : Rf'_* \circ a' \circ Lg^* \to Lg^*$ se correspondent par les applications de changement de base
$\beta : Lg^* \circ Rf_* \to Rf'_* \circ L(g')^*$
(Cohomologie, remarque \ref{cohomology-remark-base-change}) et $\alpha : L(g')^* \circ a \to a' \circ Lg^*$ (\ref{duality-equation-base-change-map}). Plus précisément, le diagramme
$$
\xymatrix{
Lg^* \circ Rf_* \circ a
\ar[d]_{\beta \star 1} \ar[r]_-{1 \star \text{Tr}_f} &
Lg^* \\
Rf'_* \circ L(g')^* \circ a \ar[r]^{1 \star \alpha} &
Rf'_* \circ a' \circ Lg^* \ar[u]_{\text{Tr}_{f'} \star 1}
}
$$
de transformations de foncteurs est commutatif.
\end{lemma}

\begin{proof}
Dans cette démonstration, nous écrivons $f_*$ pour $Rf_*$ et $g^*$ pour $Lg^*$, et nous omettons les produits $\star$ avec les identités : on peut reconstituer ceux-ci dès lors que la source et le but de la transformation sont connus. Rappelons que $\beta : g^* \circ f_* \to f'_* \circ (g')^*$ est un isomorphisme et que $\alpha$ est défini à l'aide de l'isomorphisme $\beta^\vee : g'_* \circ a' \to a \circ g_*$, adjoint de $\beta$, voir le lemme \ref{duality-lemma-flat-precompose-pus} et sa démonstration. Remarquons d'abord que la flèche horizontale supérieure du diagramme du lemme est égale à la composée
$$
g^* \circ f_* \circ a \to
g^* \circ f_* \circ a \circ g_* \circ g^* \to
g^* \circ g_* \circ g^* \to g^*
$$
où la première flèche est l'unité pour $(g^*, g_*)$, la deuxième est $\text{Tr}_f$ et la troisième est la co-unité pour $(g^*, g_*)$. Cela résulte simplement du fait que la composée $g^* \to g^* \circ g_* \circ g^* \to g^*$ de l'unité et de la co-unité est l'identité. Considérons le diagramme
\begingroup\small
$$
\xymatrix{
& g^* \circ f_* \circ a \ar[ld]_\beta \ar[d] \ar[r]_{\text{Tr}_f} & g^* \\
f'_* \circ (g')^* \circ a \ar[dr] &
g^* \circ f_* \circ a \circ g_* \circ g^* \ar[d]_\beta \ar[ru] &
g^* \circ f_* \circ g'_* \circ a' \circ g^* \ar[l]_{\beta^\vee} \ar[d]_\beta &
f'_* \circ a' \circ g^* \ar[lu]_{\text{Tr}_{f'}} \\
& f'_* \circ (g')^* \circ a \circ g_* \circ g^* &
f'_* \circ (g')^* \circ g'_* \circ a' \circ g^* \ar[ru] \ar[l]_{\beta^\vee}
}
$$
\endgroup
Les deux carrés y sont commutatifs d'après Catégories, lemme \ref{categories-lemma-properties-2-cat-cats}, ou, plus simplement, d'après la discussion précédant Catégories, définition \ref{categories-definition-horizontal-composition}. Le triangle est commutatif d'après ce qui précède. D'après Catégories, lemme \ref{categories-lemma-transformation-between-functors-and-adjoints}, le carré
$$
\xymatrix{
g^* \circ f_* \circ g'_* \circ a' \ar[d]_{\beta^\vee} \ar[r]_-\beta &
f'_* \circ (g')^* \circ g'_* \circ a' \ar[d] \\
g^* \circ f_* \circ a \circ g_* \ar[r] &
\text{id}
}
$$
est commutatif, ce qui implique la commutativité du pentagone dans le grand diagramme. Puisque $\beta$ et $\beta^\vee$ sont des isomorphismes et que le parcours extérieur du grand diagramme est, par définition, égal à $\text{Tr}_f \circ \alpha \circ \beta$, le lemme est démontré.
\end{proof}

\noindent
Soit $f : X \to Y$ un morphisme de schémas quasi-compacts et quasi-séparés. Soit $a : D_\QCoh(\mathcal{O}_Y) \to D_\QCoh(\mathcal{O}_X)$ l'adjoint à droite de $Rf_*$ du lemme \ref{duality-lemma-twisted-inverse-image}. D'après Catégories, section \ref{categories-section-adjoint}, on obtient une transformation de foncteurs
$$
\eta_f : \text{id} \to  a \circ Rf_*
$$
appelée unité de l'adjonction.

\begin{lemma}
\label{duality-lemma-unit-and-base-change}
\begingroup\emergencystretch=4em
Supposons donné un diagramme (\ref{duality-equation-base-change}) où $f$ et $g$ sont Tor-indépendants. Alors les applications $1 \star \eta_f : L(g')^* \to L(g')^* \circ a \circ Rf_*$ et $\eta_{f'} \star 1 : L(g')^* \to a' \circ Rf'_* \circ L(g')^*$ se correspondent par les applications de changement de base
$\beta : Lg^* \circ Rf_* \to Rf'_* \circ L(g')^*$
(Cohomologie, remarque \ref{cohomology-remark-base-change}) et $\alpha : L(g')^* \circ a \to a' \circ Lg^*$ (\ref{duality-equation-base-change-map}).\par\endgroup
Plus précisément, le diagramme
$$
\xymatrix{
L(g')^* \ar[r]_-{1 \star \eta_f} \ar[d]_{\eta_{f'} \star 1} &
L(g')^* \circ a \circ Rf_* \ar[d]^\alpha \\
a' \circ Rf'_* \circ L(g')^* &
a' \circ Lg^* \circ Rf_* \ar[l]_-\beta
}
$$
de transformations de foncteurs est commutatif.
\end{lemma}

\begin{proof}
Cette démonstration est duale de celle du lemme \ref{duality-lemma-trace-map-and-base-change}. Nous écrivons $f_*$ pour $Rf_*$ et $g^*$ pour $Lg^*$, et omettons les produits $\star$ avec les identités : on peut reconstituer ceux-ci dès lors que la source et le but de la transformation sont connus. Rappelons que $\beta : g^* \circ f_* \to f'_* \circ (g')^*$ est un isomorphisme et que $\alpha$ est défini à l'aide de l'isomorphisme $\beta^\vee : g'_* \circ a' \to a \circ g_*$, adjoint de $\beta$, voir le lemme \ref{duality-lemma-flat-precompose-pus} et sa démonstration. Remarquons d'abord que la flèche verticale de gauche du diagramme du lemme est égale à la composée
$$
(g')^* \to (g')^* \circ g'_* \circ (g')^* \to
(g')^* \circ g'_* \circ a' \circ f'_* \circ (g')^* \to
a' \circ f'_* \circ (g')^*
$$
où la première flèche est l'unité pour $((g')^*, g'_*)$, la deuxième est $\eta_{f'}$ et la troisième est la co-unité pour $((g')^*, g'_*)$. Cela résulte simplement du fait que la composée
$(g')^* \to (g')^* \circ (g')_* \circ (g')^* \to (g')^*$
de l'unité et de la co-unité est l'identité. Considérons le diagramme
$$
\xymatrix{
& (g')^* \circ a \circ f_* \ar[r] &
(g')^* \circ a \circ g_* \circ g^* \circ f_*
\ar[ld]_\beta \\
(g')^* \ar[ru]^{\eta_f} \ar[dd]_{\eta_{f'}} \ar[rd] &
(g')^* \circ a \circ g_* \circ f'_* \circ (g')^* &
(g')^* \circ g'_* \circ a' \circ g^* \circ f_*
\ar[u]_{\beta^\vee} \ar[ld]_\beta \ar[d] \\
& (g')^* \circ g'_* \circ a' \circ f'_* \circ (g')^*
\ar[ld] \ar[u]_{\beta^\vee} &
a' \circ g^* \circ f_* \ar[lld]^\beta \\
a' \circ f'_* \circ (g')^*
}
$$
Les deux carrés y sont commutatifs d'après Catégories, lemme \ref{categories-lemma-properties-2-cat-cats}, ou, plus simplement, d'après la discussion précédant Catégories, définition \ref{categories-definition-horizontal-composition}. Le triangle est commutatif d'après ce qui précède. Par l'énoncé dual de Catégories, lemme \ref{categories-lemma-transformation-between-functors-and-adjoints}, le carré
$$
\xymatrix{
\text{id} \ar[r] \ar[d] &
g'_* \circ a' \circ g^* \circ f_* \ar[d]^\beta \\
g'_* \circ a' \circ g^* \circ f_* \ar[r]^{\beta^\vee} &
a \circ g_* \circ f'_* \circ (g')^*
}
$$
est commutatif, ce qui implique la commutativité du pentagone dans le grand diagramme. Puisque $\beta$ et $\beta^\vee$ sont des isomorphismes et que le parcours extérieur du grand diagramme est, par définition, égal à $\beta \circ \alpha \circ \eta_f$, le lemme est démontré.
\end{proof}

\begin{example}
\label{duality-example-trace-affine}
Soit $A \to B$ un homomorphisme d'anneaux. Soient $Y = \Spec(A)$ et $X = \Spec(B)$, et soit $f : X \to Y$ le morphisme correspondant à $A \to B$. Comme on l'a vu dans l'exemple \ref{duality-example-affine-twisted-inverse-image}, l'adjoint à droite de
$Rf_* : D_\QCoh(\mathcal{O}_X) \to D_\QCoh(\mathcal{O}_Y)$
envoie un objet $K$ de $D(A) = D_\QCoh(\mathcal{O}_Y)$ sur $R\Hom(B, K)$ dans $D(B) = D_\QCoh(\mathcal{O}_X)$. Le morphisme trace est le morphisme
$$
\text{Tr}_{f, K} : R\Hom(B, K) \longrightarrow R\Hom(A, K) = K
$$
induite par l'application de $A$-modules $A \to B$.
\end{example}




\section{Adjoint à droite de l'image directe et image inverse}
\label{duality-section-compare-with-pullback}

\noindent
Soit $f : X \to Y$ un morphisme de schémas quasi-compacts et quasi-séparés. Soit $a$ l'adjoint à droite de l'image directe du lemme \ref{duality-lemma-twisted-inverse-image}. Pour $K, L \in D_\QCoh(\mathcal{O}_Y)$, on dispose d'une application canonique
$$
Lf^*K \otimes^\mathbf{L}_{\mathcal{O}_X} a(L)
\longrightarrow
a(K \otimes_{\mathcal{O}_Y}^\mathbf{L} L)
$$
Cette application est adjointe à une application
$$
Rf_*(Lf^*K \otimes^\mathbf{L}_{\mathcal{O}_X} a(L)) =
K \otimes^\mathbf{L}_{\mathcal{O}_Y} Rf_*(a(L))
\longrightarrow
K \otimes^\mathbf{L}_{\mathcal{O}_Y} L
$$
(égalité d'après Catégories dérivées des schémas, lemme \ref{perfect-lemma-cohomology-base-change}), pour laquelle on utilise le morphisme trace $Rf_*a(L) \to L$. Lorsque $L = \mathcal{O}_Y$, on obtient une application
\begin{equation}
\label{duality-equation-compare-with-pullback}
Lf^*K \otimes^\mathbf{L}_{\mathcal{O}_X} a(\mathcal{O}_Y) \longrightarrow a(K)
\end{equation}
fonctorielle en $K$ et compatible avec les triangles distingués.

\begin{lemma}
\label{duality-lemma-compare-with-pullback-perfect}
Soit $f : X \to Y$ un morphisme de schémas quasi-compacts et quasi-séparés. L'application
$Lf^*K \otimes^\mathbf{L}_{\mathcal{O}_X} a(L) \to
a(K \otimes_{\mathcal{O}_Y}^\mathbf{L} L)$
définie ci-dessus pour $K, L \in D_\QCoh(\mathcal{O}_Y)$ est un isomorphisme si $K$ est parfait. En particulier, (\ref{duality-equation-compare-with-pullback}) est un isomorphisme si $K$ est parfait.
\end{lemma}

\begin{proof}
Soit $K^\vee$ le « dual » de $K$, voir Cohomologie, lemme \ref{cohomology-lemma-dual-perfect-complex}. Pour $M \in D_\QCoh(\mathcal{O}_X)$, on a
\begin{align*}
\Hom_{D(\mathcal{O}_Y)}(Rf_*M, K \otimes^\mathbf{L}_{\mathcal{O}_Y} L)
& =
\Hom_{D(\mathcal{O}_Y)}(
Rf_*M \otimes^\mathbf{L}_{\mathcal{O}_Y} K^\vee, L) \\
& =
\Hom_{D(\mathcal{O}_X)}(
M \otimes^\mathbf{L}_{\mathcal{O}_X} Lf^*K^\vee, a(L)) \\
& =
\Hom_{D(\mathcal{O}_X)}(M,
Lf^*K \otimes^\mathbf{L}_{\mathcal{O}_X} a(L))
\end{align*}
La deuxième égalité résulte de la définition de $a$ et de la formule de projection (Cohomologie, lemme \ref{cohomology-lemma-projection-formula-perfect}), ou, plus généralement, de Catégories dérivées des schémas, lemme \ref{perfect-lemma-cohomology-base-change}. Le résultat découle donc du lemme de Yoneda.
\end{proof}

\begin{lemma}
\label{duality-lemma-restriction-compare-with-pullback}
\begingroup\emergencystretch=4em
Supposons donné un diagramme (\ref{duality-equation-base-change}) où $f$ et $g$ sont Tor-indépendants. Soit $K \in D_\QCoh(\mathcal{O}_Y)$.\par\endgroup
Le diagramme
$$
\xymatrix{
L(g')^*(Lf^*K \otimes^\mathbf{L}_{\mathcal{O}_X} a(\mathcal{O}_Y))
\ar[r] \ar[d] & L(g')^*a(K) \ar[d] \\
L(f')^*Lg^*K \otimes_{\mathcal{O}_{X'}}^\mathbf{L} a'(\mathcal{O}_{Y'})
\ar[r] & a'(Lg^*K)
}
$$
est commutatif, où les flèches horizontales sont les applications (\ref{duality-equation-compare-with-pullback}) pour $K$ et $Lg^*K$, et où les applications verticales sont construites à l'aide de Cohomologie, remarque \ref{cohomology-remark-base-change}, et de (\ref{duality-equation-base-change-map}).
\end{lemma}

\begin{proof}
Dans cette démonstration, nous écrirons $f_*$ pour $Rf_*$ et $f^*$ pour $Lf^*$, etc., ainsi que $\otimes$ pour $\otimes^\mathbf{L}_{\mathcal{O}_X}$, etc. Écrivons (\ref{duality-equation-compare-with-pullback}) comme la composée
\begin{align*}
f^*K \otimes a(\mathcal{O}_Y)
& \to
a(f_*(f^*K \otimes a(\mathcal{O}_Y))) \\
& \leftarrow
a(K \otimes f_*a(\mathcal{O}_K)) \\
& \to
a(K \otimes \mathcal{O}_Y) \\
& \to
a(K)
\end{align*}
La première flèche est l'unité $\eta_f$ ; la deuxième est $a$ appliqué à Cohomologie, équation (\ref{cohomology-equation-projection-formula-map}), qui est un isomorphisme d'après Catégories dérivées des schémas, lemme \ref{perfect-lemma-cohomology-base-change} ; la troisième est $a$ appliqué à $\text{id}_K \otimes \text{Tr}_f$ ; la quatrième est $a$ appliqué à l'isomorphisme $K \otimes \mathcal{O}_Y = K$. La démonstration du lemme consiste à montrer que chacune de ces applications donne un carré commutatif comme dans l'énoncé. Pour $\eta_f$ et $\text{Tr}_f$, cela résulte des lemmes \ref{duality-lemma-unit-and-base-change} et \ref{duality-lemma-trace-map-and-base-change}. Pour la flèche faisant intervenir Cohomologie, équation (\ref{cohomology-equation-projection-formula-map}), cela résulte de Cohomologie, remarque \ref{cohomology-remark-compatible-with-diagram}. Pour l'application de multiplication, c'est clair. Cela achève la démonstration.
\end{proof}

\begin{lemma}
\label{duality-lemma-compare-on-open}
Soit $f : X \to Y$ un morphisme propre de schémas noethériens. Soit $V \subset Y$ un ouvert tel que $f^{-1}(V) \to V$ soit un isomorphisme. Alors, pour $K \in D_\QCoh^+(\mathcal{O}_Y)$, l'application (\ref{duality-equation-compare-with-pullback}) se restreint à un isomorphisme au-dessus de $f^{-1}(V)$.
\end{lemma}

\begin{proof}
D'après le lemme \ref{duality-lemma-proper-noetherian}, l'application (\ref{duality-equation-sheafy}) est un isomorphisme pour les objets de $D_\QCoh^+(\mathcal{O}_Y)$. Le lemme \ref{duality-lemma-restriction-compare-with-pullback} montre donc que la restriction de (\ref{duality-equation-compare-with-pullback}), pour $K$, à $f^{-1}(V)$ est l'application (\ref{duality-equation-compare-with-pullback}) pour $K|_V$ et $f^{-1}(V) \to V$. Il suffit donc de montrer que l'application est un isomorphisme lorsque $f$ est le morphisme identité, ce qui est clair.
\end{proof}

\begin{lemma}
\label{duality-lemma-transitivity-compare-with-pullback}
Soient $f : X \to Y$ et $g : Y \to Z$ deux morphismes composables de schémas quasi-compacts et quasi-séparés, et posons $h = g \circ f$. Soient $a, b, c$ les adjoints du lemme \ref{duality-lemma-twisted-inverse-image} pour $f, g, h$. Pour tout $K \in D_\QCoh(\mathcal{O}_Z)$, le diagramme
$$
\xymatrix{
Lf^*(Lg^*K \otimes_{\mathcal{O}_Y}^\mathbf{L}
b(\mathcal{O}_Z)) \otimes_{\mathcal{O}_X}^\mathbf{L} a(\mathcal{O}_Y)
\ar@{=}[d] \ar[r] &
a(Lg^*K \otimes_{\mathcal{O}_Y}^\mathbf{L} b(\mathcal{O}_Z)) \ar[r] &
a(b(K)) \ar@{=}[d] \\
Lh^*K \otimes_{\mathcal{O}_X}^\mathbf{L} Lf^*b(\mathcal{O}_Z)
\otimes_{\mathcal{O}_X}^\mathbf{L} a(\mathcal{O}_Y) \ar[r] &
Lh^*K \otimes_{\mathcal{O}_X}^\mathbf{L} c(\mathcal{O}_Z) \ar[r] &
c(K)
}
$$
est commutatif, où les flèches sont celles de (\ref{duality-equation-compare-with-pullback}) et où l'on a utilisé $Lh^* = Lf^* \circ Lg^*$ et $c = a \circ b$.
\end{lemma}

\begin{proof}
Dans cette démonstration, nous écrirons $f_*$ pour $Rf_*$ et $f^*$ pour $Lf^*$, etc., ainsi que $\otimes$ pour $\otimes^\mathbf{L}_{\mathcal{O}_X}$, etc. La composée des flèches supérieures est adjointe à une application
$$
g_*f_*(f^*(g^*K \otimes b(\mathcal{O}_Z)) \otimes a(\mathcal{O}_Y)) \to K
$$
Le membre de gauche est égal à $K \otimes g_*f_*(f^*b(\mathcal{O}_Z) \otimes a(\mathcal{O}_Y))$ d'après Catégories dérivées des schémas, lemme \ref{perfect-lemma-cohomology-base-change}, et l'examen des définitions montre que l'application provient de l'application
$$
g_*f_*(f^*b(\mathcal{O}_Z) \otimes a(\mathcal{O}_Y))
\xleftarrow{g_*\epsilon}
g_*(b(\mathcal{O}_Z) \otimes f_*a(\mathcal{O}_Y)) \xrightarrow{g_*\alpha}
g_*(b(\mathcal{O}_Z)) \xrightarrow{\beta} \mathcal{O}_Z
$$
tensorisée par $\text{id}_K$. Ici, $\epsilon$ est l'isomorphisme de Catégories dérivées des schémas, lemme \ref{perfect-lemma-cohomology-base-change}, et $\beta$ provient de la co-unité $g_*b \to \text{id}$. De même, la composée des flèches horizontales inférieures est adjointe à $\text{id}_K$ tensorisé par la composée
$$
g_*f_*(f^*b(\mathcal{O}_Z) \otimes a(\mathcal{O}_Y)) \xrightarrow{g_*f_*\delta}
g_*f_*(ab(\mathcal{O}_Z)) \xrightarrow{g_*\gamma}
g_*(b(\mathcal{O}_Z)) \xrightarrow{\beta}
\mathcal{O}_Z
$$
où $\gamma$ provient de la co-unité $f_*a \to \text{id}$, et $\delta$ est l'application dont l'adjointe est la composée
$$
f_*(f^*b(\mathcal{O}_Z) \otimes a(\mathcal{O}_Y))
\xleftarrow{\epsilon}
b(\mathcal{O}_Z) \otimes f_*a(\mathcal{O}_Y) \xrightarrow{\alpha}
b(\mathcal{O}_Z)
$$
Les propriétés générales des foncteurs adjoints, des applications adjointes et des co-unités (voir Catégories, section \ref{categories-section-adjoint}) donnent $\gamma \circ f_*\delta = \alpha \circ \epsilon^{-1}$, comme voulu.
\end{proof}





\section{Adjoint à droite de l'image directe pour les immersions fermées}
\sectionmark{Adjoint à droite et immersions fermées}
\label{duality-section-sections-with-exact-support}

\noindent
Soit $i : (Z, \mathcal{O}_Z) \to (X, \mathcal{O}_X)$ un morphisme d'espaces annelés tel que $i$ soit un homéomorphisme sur une partie fermée et que $i^\sharp : \mathcal{O}_X \to i_*\mathcal{O}_Z$ soit surjective (par exemple, une immersion fermée de schémas). Soit $\mathcal{I} = \Ker(i^\sharp)$. Pour un faisceau de $\mathcal{O}_X$-modules $\mathcal{F}$, le faisceau
$$
\SheafHom_{\mathcal{O}_X}(i_*\mathcal{O}_Z, \mathcal{F})
$$
est un faisceau de $\mathcal{O}_X$-modules annulé par $\mathcal{I}$. D'après Modules, lemme \ref{modules-lemma-i-star-equivalence}, il existe donc un faisceau de $\mathcal{O}_Z$-modules, que nous noterons $\SheafHom(\mathcal{O}_Z, \mathcal{F})$, tel que
$$
i_*\SheafHom(\mathcal{O}_Z, \mathcal{F}) =
\SheafHom_{\mathcal{O}_X}(i_*\mathcal{O}_Z, \mathcal{F})
$$
comme $\mathcal{O}_X$-modules. Explicitons cet énoncé.

\begin{lemma}
\label{duality-lemma-compute-sheaf-with-exact-support}
Avec les notations précédentes, le foncteur $\SheafHom(\mathcal{O}_Z, -)$ est un adjoint à droite du foncteur $i_* : \textit{Mod}(\mathcal{O}_Z) \to \textit{Mod}(\mathcal{O}_X)$. Pour $V \subset Z$ ouvert, on a
$$
\Gamma(V, \SheafHom(\mathcal{O}_Z, \mathcal{F})) =
\{s \in \Gamma(U, \mathcal{F}) \mid \mathcal{I}s = 0\}
$$
où $U \subset X$ est un ouvert dont l'intersection avec $Z$ est $V$.
\end{lemma}

\begin{proof}
Soit $\mathcal{G}$ un faisceau de $\mathcal{O}_Z$-modules. Alors
$$
\Hom_{\mathcal{O}_X}(i_*\mathcal{G}, \mathcal{F}) =
\Hom_{i_*\mathcal{O}_Z}(i_*\mathcal{G},
\SheafHom_{\mathcal{O}_X}(i_*\mathcal{O}_Z, \mathcal{F})) =
\Hom_{\mathcal{O}_Z}(\mathcal{G}, \SheafHom(\mathcal{O}_Z, \mathcal{F}))
$$
La première égalité résulte de Modules, lemme \ref{modules-lemma-adjoint-tensor-restrict}, et la seconde de la pleine fidélité de $i_*$, voir Modules, lemme \ref{modules-lemma-i-star-equivalence}. La description des sections est laissée au lecteur.
\end{proof}

\noindent
Le foncteur
$$
\textit{Mod}(\mathcal{O}_X)
\longrightarrow
\textit{Mod}(\mathcal{O}_Z),
\quad
\mathcal{F} \longmapsto \SheafHom(\mathcal{O}_Z, \mathcal{F})
$$
est exact à gauche et admet un prolongement dérivé
$$
R\SheafHom(\mathcal{O}_Z, -) : D(\mathcal{O}_X) \to D(\mathcal{O}_Z).
$$

\begin{lemma}
\label{duality-lemma-sheaf-with-exact-support-adjoint}
Avec les notations précédentes, le foncteur $R\SheafHom(\mathcal{O}_Z, -)$ est l'adjoint à droite du foncteur $Ri_* : D(\mathcal{O}_Z) \to D(\mathcal{O}_X)$.
\end{lemma}

\begin{proof}
Cela résulte du fait que $i_*$ et $\SheafHom(\mathcal{O}_Z, -)$ sont des foncteurs adjoints d'après le lemme \ref{duality-lemma-compute-sheaf-with-exact-support}. Voir Catégories dérivées, lemme \ref{derived-lemma-derived-adjoint-functors}.
\end{proof}

\begin{lemma}
\label{duality-lemma-sheaf-with-exact-support-ext}
Avec les notations précédentes, on a
$$
Ri_*R\SheafHom(\mathcal{O}_Z, K) =
R\SheafHom_{\mathcal{O}_X}(i_*\mathcal{O}_Z, K)
$$
dans $D(\mathcal{O}_X)$ pour tout $K$ de $D(\mathcal{O}_X)$.
\end{lemma}

\begin{proof}
Cela résulte immédiatement de la construction du foncteur
\par\noindent\hfil
$R\SheafHom(\mathcal{O}_Z, -)$.
\hfil\par
\end{proof}

\begin{lemma}
\label{duality-lemma-sheaf-with-exact-support-internal-home}
Avec les notations précédentes, pour $M \in D(\mathcal{O}_Z)$, on a
$$
R\SheafHom_{\mathcal{O}_X}(Ri_*M, K) =
Ri_*R\SheafHom_{\mathcal{O}_Z}(M, R\SheafHom(\mathcal{O}_Z, K))
$$
dans $D(\mathcal{O}_Z)$ pour tout $K$ de $D(\mathcal{O}_X)$.
\end{lemma}

\begin{proof}
Cela résulte immédiatement de la construction du foncteur
\par\noindent\hfil
$R\SheafHom(\mathcal{O}_Z, -)$
\hfil\par\noindent
et du fait que, si $\mathcal{K}^\bullet$ est un complexe K-injectif de $\mathcal{O}_X$-modules, alors $\SheafHom(\mathcal{O}_Z, \mathcal{K}^\bullet)$ est un complexe K-injectif de $\mathcal{O}_Z$-modules, voir Catégories dérivées, lemme \ref{derived-lemma-adjoint-preserve-K-injectives}.
\end{proof}

\begin{lemma}
\label{duality-lemma-sheaf-with-exact-support-quasi-coherent}
Soit $i : Z \to X$ une immersion fermée pseudo-cohérente de schémas (toute immersion fermée convient si $X$ est localement noethérien). Alors
\begin{enumerate}
\item $R\SheafHom(\mathcal{O}_Z, -)$ envoie $D^+_\QCoh(\mathcal{O}_X)$ dans $D^+_\QCoh(\mathcal{O}_Z)$ ;
\item si $X = \Spec(A)$ et $Z = \Spec(B)$, le diagramme
$$
\xymatrix{
D^+(B) \ar[r] & D_\QCoh^+(\mathcal{O}_Z) \\
D^+(A) \ar[r] \ar[u]^{R\Hom(B, -)} &
D_\QCoh^+(\mathcal{O}_X) \ar[u]_{R\SheafHom(\mathcal{O}_Z, -)}
}
$$
est commutatif.
\end{enumerate}
\end{lemma}

\begin{proof}
Pour justifier la remarque entre parenthèses, si $X$ est localement noethérien, alors $i$ est pseudo-cohérent d'après Compléments sur les morphismes, lemme \ref{more-morphisms-lemma-Noetherian-pseudo-coherent}.

\medskip\noindent
Soit $K$ un objet de $D^+_\QCoh(\mathcal{O}_X)$. Pour démontrer (1), il suffit, d'après Morphismes, lemme \ref{morphisms-lemma-i-star-equivalence}, de montrer que $i_*$ appliqué à $H^n(R\SheafHom(\mathcal{O}_Z, K))$ donne un module quasi-cohérent sur $X$. D'après le lemme \ref{duality-lemma-sheaf-with-exact-support-ext}, cela revient à montrer que
\par\noindent\hfil
$R\SheafHom_{\mathcal{O}_X}(i_*\mathcal{O}_Z, K)$
\hfil\par\noindent
appartient à $D_\QCoh(\mathcal{O}_X)$. Puisque $i$ est pseudo-cohérent, le faisceau $\mathcal{O}_Z$ est un $\mathcal{O}_X$-module pseudo-cohérent. Le résultat découle donc de Catégories dérivées des schémas, lemme \ref{perfect-lemma-quasi-coherence-internal-hom}.

\medskip\noindent
Supposons $X = \Spec(A)$ et $Z = \Spec(B)$ comme en (2). Soit $I^\bullet$ un complexe borné inférieurement de $A$-modules injectifs représentant un objet $K$ de $D^+(A)$. On sait alors que $R\Hom(B, K) = \Hom_A(B, I^\bullet)$, considéré comme complexe de $B$-modules. Choisissons un quasi-isomorphisme
$$
\widetilde{I^\bullet} \longrightarrow \mathcal{I}^\bullet
$$
où $\mathcal{I}^\bullet$ est un complexe borné inférieurement de $\mathcal{O}_X$-modules injectifs. La description du foncteur $\SheafHom(\mathcal{O}_Z, -)$ au lemme \ref{duality-lemma-compute-sheaf-with-exact-support} donne une application
$$
\Hom_A(B, I^\bullet)
\longrightarrow
\Gamma(Z, \SheafHom(\mathcal{O}_Z, \mathcal{I}^\bullet))
$$
Observons que $\SheafHom(\mathcal{O}_Z, \mathcal{I}^\bullet)$ représente $R\SheafHom(\mathcal{O}_Z, \widetilde{K})$. En appliquant la propriété universelle du foncteur $\widetilde{\ }$, on obtient une application
$$
\widetilde{\Hom_A(B, I^\bullet)}
\longrightarrow
R\SheafHom(\mathcal{O}_Z, \widetilde{K})
$$
dans $D(\mathcal{O}_Z)$. On peut vérifier que cette application est un isomorphisme dans $D(\mathcal{O}_Z)$ après application de $i_*$. Or, après application de $i_*$, on obtient l'isomorphisme de Catégories dérivées des schémas, lemme \ref{perfect-lemma-quasi-coherence-internal-hom}, via l'identification du lemme \ref{duality-lemma-sheaf-with-exact-support-ext}.
\end{proof}

\begin{lemma}
\label{duality-lemma-sheaf-with-exact-support-coherent}
Soit $i : Z \to X$ une immersion fermée de schémas. Supposons $X$ localement noethérien. Alors $R\SheafHom(\mathcal{O}_Z, -)$ envoie $D^+_{\textit{Coh}}(\mathcal{O}_X)$ dans $D^+_{\textit{Coh}}(\mathcal{O}_Z)$.
\end{lemma}

\begin{proof}
La question est locale sur $X$ ; on peut donc supposer $X$ affine. Écrivons $X = \Spec(A)$ et $Z = \Spec(B)$, avec $A$ noethérien et $A \to B$ surjectif. On peut alors appliquer le lemme \ref{duality-lemma-sheaf-with-exact-support-quasi-coherent} pour traduire la question en algèbre. Le résultat algébrique correspondant est une conséquence de Complexes dualisants, lemme \ref{dualizing-lemma-exact-support-coherent}.
\end{proof}

\begin{lemma}
\label{duality-lemma-twisted-inverse-image-closed}
Soit $X$ un schéma quasi-compact et quasi-séparé. Soit $i : Z \to X$ une immersion fermée pseudo-cohérente (si $X$ est noethérien, toute immersion fermée est pseudo-cohérente). Soit $a : D_\QCoh(\mathcal{O}_X) \to D_\QCoh(\mathcal{O}_Z)$ l'adjoint à droite de $Ri_*$. Il existe alors un isomorphisme fonctoriel
$$
a(K) = R\SheafHom(\mathcal{O}_Z, K)
$$
pour $K \in D_\QCoh^+(\mathcal{O}_X)$.
\end{lemma}

\begin{proof}
(L'énoncé entre parenthèses résulte de Compléments sur les morphismes, lemme \ref{more-morphisms-lemma-Noetherian-pseudo-coherent}.) D'après le lemme \ref{duality-lemma-sheaf-with-exact-support-adjoint}, le foncteur $R\SheafHom(\mathcal{O}_Z, -)$ est un adjoint à droite de $Ri_* : D(\mathcal{O}_Z) \to D(\mathcal{O}_X)$. De plus, d'après les lemmes \ref{duality-lemma-sheaf-with-exact-support-quasi-coherent} et \ref{duality-lemma-twisted-inverse-image-bounded-below}, les foncteurs $R\SheafHom(\mathcal{O}_Z, -)$ et $a$ envoient tous deux $D_\QCoh^+(\mathcal{O}_X)$ dans $D_\QCoh^+(\mathcal{O}_Z)$. On obtient donc l'isomorphisme par unicité des foncteurs adjoints.
\end{proof}

\begin{example}
\label{duality-example-trace-closed-immersion}
Si $i : Z \to X$ est une immersion fermée de schémas noethériens, le diagramme
$$
\xymatrix{
i_*a(K) \ar[rr]_-{\text{Tr}_{i, K}} \ar@{=}[d] & &
K \ar@{=}[d] \\
i_*R\SheafHom(\mathcal{O}_Z, K) \ar@{=}[r] &
R\SheafHom_{\mathcal{O}_X}(i_*\mathcal{O}_Z, K)
\ar[r] & K
}
$$
est commutatif pour $K \in D_\QCoh^+(\mathcal{O}_X)$. Le signe d'égalité horizontal correspond ici au lemme \ref{duality-lemma-sheaf-with-exact-support-ext}, et la flèche horizontale inférieure est induite par l'application $\mathcal{O}_X \to i_*\mathcal{O}_Z$. La commutativité du diagramme est une conséquence du lemme \ref{duality-lemma-twisted-inverse-image-closed}.
\end{example}




\section{Adjoint à droite de l'image directe pour les immersions fermées et changement de base}
\label{duality-section-sections-with-exact-support-base-change}

\noindent
Considérons un diagramme cartésien de schémas
$$
\xymatrix{
Z' \ar[r]_{i'} \ar[d]_g & X' \ar[d]^f \\
Z \ar[r]^i & X
}
$$
où $i$ est une immersion fermée. Si $Z$ et $X'$ sont Tor-indépendants sur $X$, il existe une application canonique de changement de base
\begin{equation}
\label{duality-equation-base-change-exact-support}
Lg^*R\SheafHom(\mathcal{O}_Z, K)
\longrightarrow
R\SheafHom(\mathcal{O}_{Z'}, Lf^*K)
\end{equation}
dans $D(\mathcal{O}_{Z'})$, fonctorielle en $K$ dans $D(\mathcal{O}_X)$. En effet, par l'adjonction du lemme \ref{duality-lemma-sheaf-with-exact-support-adjoint}, une telle flèche revient à une application
$$
Ri'_*Lg^*R\SheafHom(\mathcal{O}_Z, K)
\longrightarrow
Lf^*K
$$
dans $D(\mathcal{O}_{X'})$. La Tor-indépendance donne $Ri'_* \circ Lg^* = Lf^* \circ Ri_*$ (voir Catégories dérivées des schémas, lemme \ref{perfect-lemma-compare-base-change-closed-immersion}). Cela revient donc à une application
$$
Lf^*Ri_*R\SheafHom(\mathcal{O}_Z, K)
\longrightarrow
Lf^*K
$$
Pour celle-ci, on peut prendre $Lf^*(can)$, où $can : Ri_* R\SheafHom(\mathcal{O}_Z, K) \to K$ est la co-unité de l'adjonction.

\begin{lemma}
\label{duality-lemma-check-base-change-is-iso}
Dans la situation précédente, l'application (\ref{duality-equation-base-change-exact-support}) est un isomorphisme si et seulement si l'application de changement de base
$$
Lf^*R\SheafHom_{\mathcal{O}_X}(\mathcal{O}_Z, K)
\longrightarrow
R\SheafHom_{\mathcal{O}_{X'}}(\mathcal{O}_{Z'}, Lf^*K)
$$
de Cohomologie, remarque \ref{cohomology-remark-prepare-fancy-base-change}, en est un.
\end{lemma}

\begin{proof}
L'énoncé a un sens, car $\mathcal{O}_{Z'} = Lf^*\mathcal{O}_Z$ par la Tor-indépendance supposée. Puisque $i'_*$ est exact et fidèle, il suffit de montrer que l'application (\ref{duality-equation-base-change-exact-support}) est un isomorphisme après application de $Ri'_*$. Puisque $Ri'_* \circ Lg^* = Lf^* \circ Ri_*$ par la Tor-indépendance supposée et Catégories dérivées des schémas, lemme \ref{perfect-lemma-compare-base-change-closed-immersion}, on obtient une application
$$
Lf^*Ri_*R\SheafHom(\mathcal{O}_Z, K)
\longrightarrow
Ri'_*R\SheafHom(\mathcal{O}_{Z'}, Lf^*K)
$$
dont la source et le but sont ceux de l'énoncé, d'après le lemme \ref{duality-lemma-sheaf-with-exact-support-ext}. Nous omettons de vérifier qu'il s'agit de la même application que celle construite dans Cohomologie, remarque \ref{cohomology-remark-prepare-fancy-base-change}.
\end{proof}

\begin{lemma}
\label{duality-lemma-flat-bc-sheaf-with-exact-support}
Dans la situation précédente, supposons $f$ plat et $i$ pseudo-cohérente. Alors (\ref{duality-equation-base-change-exact-support}) est un isomorphisme pour $K$ dans $D^+_\QCoh(\mathcal{O}_X)$.
\end{lemma}

\begin{proof}
Première démonstration. Pour montrer que cette application est un isomorphisme, on peut travailler localement. On peut donc supposer $X$, $X'$, $Z$, $Z'$ affines, correspondant respectivement aux anneaux $A$, $A'$, $B$, $B'$. Alors $B$ et $A'$ sont Tor-indépendants sur $A$. D'après le lemme \ref{duality-lemma-check-base-change-is-iso}, il suffit de vérifier que
$$
R\Hom_A(B, K) \otimes_A^\mathbf{L} A' =
R\Hom_{A'}(B', K \otimes_A^\mathbf{L} A')
$$
dans $D(A')$ pour tout $K \in D^+(A)$. On utilise ici Catégories dérivées des schémas, lemme \ref{perfect-lemma-quasi-coherence-internal-hom}, et le fait que $B$, resp.\ $B'$, soit pseudo-cohérent comme $A$-module, resp.\ $A'$-module, pour comparer les Hom dérivés au niveau des anneaux et des schémas. L'égalité affichée résulte de Compléments d'algèbre, lemme \ref{more-algebra-lemma-internal-hom-evaluate-tensor-isomorphism}, (3). Voir aussi la discussion dans Complexes dualisants, section \ref{dualizing-section-base-change-trivial-duality}.

\medskip\noindent
Deuxième démonstration\footnote{Cette démonstration montre qu'il suffit de supposer que $K$ appartient à $D^+(\mathcal{O}_X)$.}. Soit $z' \in Z'$ d'image $z \in Z$. Montrons d'abord que (\ref{duality-equation-base-change-exact-support}), sur les fibres en $z'$, induit l'application
$$
R\Hom(\mathcal{O}_{Z, z}, K_z)
\otimes_{\mathcal{O}_{Z, x}}^\mathbf{L} \mathcal{O}_{Z', z'}
\longrightarrow
R\Hom(\mathcal{O}_{Z', z'},
K_z \otimes_{\mathcal{O}_{X, z}}^\mathbf{L} \mathcal{O}_{X', z'})
$$
de Complexes dualisants, équation (\ref{dualizing-equation-base-change}). En effet, les constructions de ces applications sont identiques. On applique ensuite Complexes dualisants, lemme \ref{dualizing-lemma-flat-bc-surjection}.
\end{proof}

\begin{lemma}
\label{duality-lemma-sheaf-with-exact-support-tensor}
Soit $i : Z \to X$ une immersion fermée pseudo-cohérente de schémas. Soit $M \in D_\QCoh(\mathcal{O}_X)$ localement de Tor-amplitude dans $[a, \infty)$. Soit $K \in D_\QCoh^+(\mathcal{O}_X)$. Il existe alors un isomorphisme canonique
$$
R\SheafHom(\mathcal{O}_Z, K) \otimes_{\mathcal{O}_Z}^\mathbf{L} Li^*M =
R\SheafHom(\mathcal{O}_Z, K \otimes_{\mathcal{O}_X}^\mathbf{L} M)
$$
dans $D(\mathcal{O}_Z)$.
\end{lemma}

\begin{proof}
Une application du membre de gauche vers celui de droite revient à une application
$$
Ri_*R\SheafHom(\mathcal{O}_Z, K) \otimes_{\mathcal{O}_X}^\mathbf{L} M
\longrightarrow
K \otimes_{\mathcal{O}_X}^\mathbf{L} M
$$
d'après les lemmes \ref{duality-lemma-sheaf-with-exact-support-adjoint} et \ref{duality-lemma-sheaf-with-exact-support-ext}. Pour cette application, on prend la co-unité
\par\noindent\hfil
$Ri_*R\SheafHom(\mathcal{O}_Z, K) \to K$
\hfil\par\noindent
tensorisée par $\text{id}_M$. Pour voir qu'elle est un isomorphisme sous les hypothèses données, on traduit l'énoncé en algèbre à l'aide du lemme \ref{duality-lemma-sheaf-with-exact-support-quasi-coherent}, puis on utilise, par exemple, Compléments d'algèbre, lemme \ref{more-algebra-lemma-internal-hom-evaluate-tensor-isomorphism}, (3). Au lieu d'utiliser le lemme \ref{duality-lemma-sheaf-with-exact-support-quasi-coherent}, on peut considérer les fibres, comme dans la deuxième démonstration du lemme \ref{duality-lemma-flat-bc-sheaf-with-exact-support}.
\end{proof}



\section{Adjoint à droite de l'image directe pour les morphismes finis}
\sectionmark{Adjoint à droite et morphismes finis}
\label{duality-section-duality-finite}

\noindent
Si $i : Z \to X$ est une immersion fermée de schémas, il existe un adjoint à droite $\SheafHom(\mathcal{O}_Z, -)$ du foncteur $i_* : \textit{Mod}(\mathcal{O}_Z) \to \textit{Mod}(\mathcal{O}_X)$ dont le prolongement dérivé $R\SheafHom(\mathcal{O}_Z, -)$ est l'adjoint à droite de $Ri_* : D(\mathcal{O}_Z) \to D(\mathcal{O}_X)$. Voir la section \ref{duality-section-sections-with-exact-support}. Dans le cas d'un morphisme fini $f : Y \to X$, cette stratégie ne peut pas fonctionner, car le foncteur $f_* : \textit{Mod}(\mathcal{O}_Y) \to \textit{Mod}(\mathcal{O}_X)$ n'est pas exact en général et n'admet donc pas d'adjoint à droite. On peut le remplacer par le foncteur exact
$\textit{Mod}(f_*\mathcal{O}_Y) \to \textit{Mod}(\mathcal{O}_X)$
et considérer l'adjoint à droite correspondant ainsi que son prolongement dérivé.

\medskip\noindent
Soit $f : Y \to X$ un morphisme affine de schémas. Pour un faisceau de $\mathcal{O}_X$-modules $\mathcal{F}$, le faisceau
$$
\SheafHom_{\mathcal{O}_X}(f_*\mathcal{O}_Y, \mathcal{F})
$$
est un faisceau de $f_*\mathcal{O}_Y$-modules. On obtient un foncteur
$\textit{Mod}(\mathcal{O}_X) \to \textit{Mod}(f_*\mathcal{O}_Y)$
que nous noterons $\SheafHom(f_*\mathcal{O}_Y, -)$.

\begin{lemma}
\label{duality-lemma-compute-sheafhom-affine}
Avec les notations précédentes, le foncteur $\SheafHom(f_*\mathcal{O}_Y, -)$ est un adjoint à droite du foncteur de restriction $\textit{Mod}(f_*\mathcal{O}_Y) \to \textit{Mod}(\mathcal{O}_X)$. Pour un ouvert affine $U \subset X$, on a
$$
\Gamma(U, \SheafHom(f_*\mathcal{O}_Y, \mathcal{F})) =
\Hom_A(B, \mathcal{F}(U))
$$
où $A = \mathcal{O}_X(U)$ et $B = \mathcal{O}_Y(f^{-1}(U))$.
\end{lemma}

\begin{proof}
L'adjonction résulte de Modules, lemme \ref{modules-lemma-adjoint-tensor-restrict}. Puisque $f$ est affine, $f_*\mathcal{O}_Y$ est le faisceau quasi-cohérent correspondant à $B$ considéré comme $A$-module. La description des sections sur $U$ découle donc de Schémas, lemme \ref{schemes-lemma-compare-constructions}.
\end{proof}

\noindent
Le foncteur $\SheafHom(f_*\mathcal{O}_Y, -)$ est exact à gauche. Soit
$$
R\SheafHom(f_*\mathcal{O}_Y, -) :
D(\mathcal{O}_X)
\longrightarrow
D(f_*\mathcal{O}_Y)
$$
son prolongement dérivé.

\begin{lemma}
\label{duality-lemma-sheafhom-affine-adjoint}
Avec les notations précédentes, le foncteur $R\SheafHom(f_*\mathcal{O}_Y, -)$ est l'adjoint à droite du foncteur $D(f_*\mathcal{O}_Y) \to D(\mathcal{O}_X)$.
\end{lemma}

\begin{proof}
Cela résulte du lemme \ref{duality-lemma-compute-sheafhom-affine} et de Catégories dérivées, lemme \ref{derived-lemma-derived-adjoint-functors}.
\end{proof}

\begin{lemma}
\label{duality-lemma-sheafhom-affine-ext}
Avec les notations précédentes, la composée
$$
D(\mathcal{O}_X) \xrightarrow{R\SheafHom(f_*\mathcal{O}_Y, -)}
D(f_*\mathcal{O}_Y) \to D(\mathcal{O}_X)
$$
est le foncteur $K \mapsto R\SheafHom_{\mathcal{O}_X}(f_*\mathcal{O}_Y, K)$.
\end{lemma}

\begin{proof}
Cela résulte immédiatement de la construction.
\end{proof}

\begin{lemma}
\label{duality-lemma-finite-twisted}
Soit $f : Y \to X$ un morphisme fini pseudo-cohérent de schémas (un morphisme fini de schémas noethériens est pseudo-cohérent). Le foncteur $R\SheafHom(f_*\mathcal{O}_Y, -)$ envoie $D_\QCoh^+(\mathcal{O}_X)$ dans $D_\QCoh^+(f_*\mathcal{O}_Y)$. Si $X$ est quasi-compact et quasi-séparé, le diagramme
$$
\xymatrix{
D_\QCoh^+(\mathcal{O}_X) \ar[rr]_a \ar[rd]_{R\SheafHom(f_*\mathcal{O}_Y, -)}
& & D_\QCoh^+(\mathcal{O}_Y) \ar[ld]^\Phi \\
& D_\QCoh^+(f_*\mathcal{O}_Y)
}
$$
est commutatif, où $a$ est l'adjoint à droite du lemme \ref{duality-lemma-twisted-inverse-image} pour $f$ et où $\Phi$ est l'équivalence de Catégories dérivées des schémas, lemme \ref{perfect-lemma-affine-morphism-equivalence}.
\end{lemma}

\begin{proof}
\begingroup\emergencystretch=6em
(La remarque entre parenthèses résulte de Compléments sur les morphismes, lemme \ref{more-morphisms-lemma-Noetherian-pseudo-coherent}.) Puisque $f$ est pseudo-cohérent, le $\mathcal{O}_X$-module $f_*\mathcal{O}_Y$ est pseudo-cohérent, voir Compléments sur les morphismes, lemme \ref{more-morphisms-lemma-finite-pseudo-coherent}. Ainsi, $R\SheafHom(f_*\mathcal{O}_Y, -)$ envoie $D_\QCoh^+(\mathcal{O}_X)$ dans $D_\QCoh^+(f_*\mathcal{O}_Y)$, voir Catégories dérivées des schémas, lemme \ref{perfect-lemma-quasi-coherence-internal-hom}. Alors $\Phi \circ a$ et $R\SheafHom(f_*\mathcal{O}_Y, -)$ coïncident sur $D_\QCoh^+(\mathcal{O}_X)$, car ces deux foncteurs sont adjoints à droite du foncteur de restriction $D_\QCoh^+(f_*\mathcal{O}_Y) \to D_\QCoh^+(\mathcal{O}_X)$. Pour le voir, on utilise les lemmes \ref{duality-lemma-twisted-inverse-image-bounded-below} et \ref{duality-lemma-sheafhom-affine-adjoint}.\par\endgroup
\end{proof}

\begin{remark}
\label{duality-remark-trace-map-finite}
Si $f : Y \to X$ est un morphisme fini de schémas noethériens, le diagramme
$$
\xymatrix{
Rf_*a(K) \ar[r]_-{\text{Tr}_{f, K}} \ar@{=}[d] & K \ar@{=}[d] \\
R\SheafHom_{\mathcal{O}_X}(f_*\mathcal{O}_Y, K) \ar[r] & K
}
$$
est commutatif pour $K \in D_\QCoh^+(\mathcal{O}_X)$. Cela résulte du lemme \ref{duality-lemma-finite-twisted}. La flèche horizontale inférieure est induite par l'application $\mathcal{O}_X \to f_*\mathcal{O}_Y$, et la flèche horizontale supérieure est le morphisme trace de la section \ref{duality-section-trace}.
\end{remark}






\section{Adjoint à droite de l'image directe pour les morphismes propres et plats}
\label{duality-section-proper-flat}

\noindent
Pour les morphismes propres, plats et de présentation finie entre schémas quasi-compacts et quasi-séparés, l'adjoint à droite de l'image directe possède des propriétés remarquables.

\begin{lemma}
\label{duality-lemma-proper-flat}
Soit $Y$ un schéma quasi-compact et quasi-séparé. Soit $f : X \to Y$ un morphisme de schémas propre, plat et de présentation finie. Soit $a$ l'adjoint à droite de $Rf_* : D_\QCoh(\mathcal{O}_X) \to D_\QCoh(\mathcal{O}_Y)$ du lemme \ref{duality-lemma-twisted-inverse-image}. Alors $a$ commute aux sommes directes.
\end{lemma}

\begin{proof}
Soit $P$ un objet parfait de $D(\mathcal{O}_X)$. D'après Catégories dérivées des schémas, lemme \ref{perfect-lemma-flat-proper-perfect-direct-image-general}, le complexe $Rf_*P$ est parfait sur $Y$. Soit $K_i$ une famille d'objets de $D_\QCoh(\mathcal{O}_Y)$. Alors
\begin{align*}
\Hom_{D(\mathcal{O}_X)}(P, a(\bigoplus K_i))
& =
\Hom_{D(\mathcal{O}_Y)}(Rf_*P, \bigoplus K_i) \\
& =
\bigoplus \Hom_{D(\mathcal{O}_Y)}(Rf_*P, K_i) \\
& =
\bigoplus \Hom_{D(\mathcal{O}_X)}(P, a(K_i))
\end{align*}
car un objet parfait est compact (Catégories dérivées des schémas, proposition \ref{perfect-proposition-compact-is-perfect}). Puisque $D_\QCoh(\mathcal{O}_X)$ admet un générateur parfait (Catégories dérivées des schémas, théorème \ref{perfect-theorem-bondal-van-den-Bergh}), on en conclut que l'application $\bigoplus a(K_i) \to a(\bigoplus K_i)$ est un isomorphisme, c'est-à-dire que $a$ commute aux sommes directes.
\end{proof}

\begin{lemma}
\label{duality-lemma-proper-flat-relative}
Soit $Y$ un schéma quasi-compact et quasi-séparé. Soit $f : X \to Y$ un morphisme de schémas propre, plat et de présentation finie. Soit $a$ l'adjoint à droite de $Rf_* : D_\QCoh(\mathcal{O}_X) \to D_\QCoh(\mathcal{O}_Y)$ du lemme \ref{duality-lemma-twisted-inverse-image}. Alors
\begin{enumerate}
\item pour tout fermé $T \subset Y$, si $Q \in D_\QCoh(Y)$ est à support dans $T$, alors $a(Q)$ est à support dans $f^{-1}(T)$ ;
\item pour tout ouvert quasi-compact $V \subset Y$ et tout $K \in D_\QCoh(\mathcal{O}_Y)$, l'application (\ref{duality-equation-sheafy}) est un isomorphisme.
\end{enumerate}
\end{lemma}

\begin{proof}
Cela résulte des lemmes \ref{duality-lemma-when-sheafy}, \ref{duality-lemma-proper-noetherian} et \ref{duality-lemma-proper-flat}. À strictement parler, on n'obtient (1) que pour les $T$ dont le complémentaire est un ouvert quasi-compact ; le cas général s'en déduit néanmoins, car tout fermé est une intersection de tels fermés.
\end{proof}

\begin{lemma}
\label{duality-lemma-compare-with-pullback-flat-proper}
Soit $Y$ un schéma quasi-compact et quasi-séparé. Soit $f : X \to Y$ un morphisme de schémas propre, plat et de présentation finie. L'application (\ref{duality-equation-compare-with-pullback}) est un isomorphisme pour tout objet $K$ de $D_\QCoh(\mathcal{O}_Y)$.
\end{lemma}

\begin{proof}
D'après le lemme \ref{duality-lemma-proper-flat}, $a$ commute aux sommes directes. La collection des objets de $D_\QCoh(\mathcal{O}_Y)$ pour lesquels (\ref{duality-equation-compare-with-pullback}) est un isomorphisme est donc une sous-catégorie triangulée, saturée et strictement pleine de $D_\QCoh(\mathcal{O}_Y)$, qui est de plus stable par sommes directes. Puisque $D_\QCoh(\mathcal{O}_Y)$ est une catégorie de modules (Catégories dérivées des schémas, théorème \ref{perfect-theorem-DQCoh-is-Ddga}) engendrée par un unique objet parfait (Catégories dérivées des schémas, théorème \ref{perfect-theorem-bondal-van-den-Bergh}), on peut raisonner comme dans Compléments d'algèbre, remarque \ref{more-algebra-remark-P-resolution}, pour voir qu'il suffit de montrer que (\ref{duality-equation-compare-with-pullback}) est un isomorphisme pour un unique objet parfait. Or le résultat vaut pour les objets parfaits, voir le lemme \ref{duality-lemma-compare-with-pullback-perfect}.
\end{proof}

\noindent
Le lemme suivant montre que l'application de changement de base (\ref{duality-equation-base-change-map}) est un isomorphisme pour les morphismes propres, plats et de présentation finie. Nous verrons dans l'exemple \ref{duality-example-base-change-wrong} que cela ne reste pas vrai pour les morphismes parfaits et propres ; dans ce cas, une hypothèse de Tor-indépendance est nécessaire.

\begin{lemma}
\label{duality-lemma-proper-flat-base-change}
Soit $g : Y' \to Y$ un morphisme de schémas quasi-compacts et quasi-séparés. Soit $f : X \to Y$ un morphisme propre et plat de présentation finie. Alors l'application de changement de base (\ref{duality-equation-base-change-map}) est un isomorphisme pour tout $K \in D_\QCoh(\mathcal{O}_Y)$.
\end{lemma}

\begin{proof}
D'après le lemme \ref{duality-lemma-proper-flat-relative}, la formation des foncteurs $a$ et $a'$ commute à la restriction aux ouverts de $Y$ et $Y'$. On peut donc supposer que $Y' \to Y$ est un morphisme de schémas affines, voir la remarque \ref{duality-remark-check-over-affines}. Dans ce cas, l'énoncé résulte du lemme \ref{duality-lemma-more-base-change}.
\end{proof}

\begin{remark}
\label{duality-remark-relative-dualizing-complex}
Soit $Y$ un schéma quasi-compact et quasi-séparé. Soit $f : X \to Y$ un morphisme propre et plat de présentation finie. Soit $a$ l'adjoint du lemme \ref{duality-lemma-twisted-inverse-image} pour $f$. Dans cette situation, $\omega_{X/Y}^\bullet = a(\mathcal{O}_Y)$ est parfois appelé {\it complexe dualisant relatif}. D'après le lemme \ref{duality-lemma-compare-with-pullback-flat-proper}, il existe un isomorphisme fonctoriel
$a(K) = Lf^*K \otimes_{\mathcal{O}_X}^\mathbf{L} \omega_{X/Y}^\bullet$
pour $K \in D_\QCoh(\mathcal{O}_Y)$. De plus, le morphisme trace
$$
\text{Tr}_{f, \mathcal{O}_Y} : Rf_*\omega_{X/Y}^\bullet \to \mathcal{O}_Y
$$
de la section \ref{duality-section-trace} induit le morphisme trace pour tout $K$ dans $D_\QCoh(\mathcal{O}_Y)$. Plus précisément, on a le diagramme
$$
\xymatrix{
Rf_*a(K) \ar[rrr]_{\text{Tr}_{f, K}} \ar@{=}[d] & & &
K \ar@{=}[d] \\
Rf_*(Lf^*K \otimes_{\mathcal{O}_X}^\mathbf{L} \omega_{X/Y}^\bullet)
\ar@{=}[r] &
K \otimes_{\mathcal{O}_Y}^\mathbf{L} Rf_*\omega_{X/Y}^\bullet
\ar[rr]^-{\text{id}_K \otimes \text{Tr}_{f, \mathcal{O}_Y}} & & K
}
$$
où l'égalité en bas à droite est celle de Catégories dérivées des schémas, lemme \ref{perfect-lemma-cohomology-base-change}. Si $g : Y' \to Y$ est un morphisme de schémas quasi-compacts et quasi-séparés et si $X' = Y' \times_Y X$, le lemme \ref{duality-lemma-proper-flat-base-change} donne $\omega_{X'/Y'}^\bullet = L(g')^*\omega_{X/Y}^\bullet$, où $g' : X' \to X$ est la projection ; d'après le lemme \ref{duality-lemma-trace-map-and-base-change}, le morphisme trace
$$
\text{Tr}_{f', \mathcal{O}_{Y'}} :
Rf'_*\omega_{X'/Y'}^\bullet \to \mathcal{O}_{Y'}
$$
pour $f' : X' \to Y'$ est obtenue par changement de base de $\text{Tr}_{f, \mathcal{O}_Y}$ via l'isomorphisme de changement de base.
\end{remark}

\begin{remark}
\label{duality-remark-relative-dualizing-complex-relative-cup-product}
Soient $f : X \to Y$, $\omega^\bullet_{X/Y}$ et $\text{Tr}_{f, \mathcal{O}_Y}$ comme dans la remarque \ref{duality-remark-relative-dualizing-complex}. Soient $K$ et $M$ dans $D_\QCoh(\mathcal{O}_X)$, avec $M$ pseudo-cohérent (par exemple parfait). Supposons donnée une application
$K \otimes_{\mathcal{O}_X}^\mathbf{L} M \to \omega^\bullet_{X/Y}$
correspondant à un isomorphisme
$K \to R\SheafHom_{\mathcal{O}_X}(M, \omega^\bullet_{X/Y})$
via Cohomologie, équation (\ref{cohomology-equation-internal-hom}). Alors le cup-produit relatif (Cohomologie, remarque \ref{cohomology-remark-cup-product})
$$
Rf_*K \otimes_{\mathcal{O}_Y}^\mathbf{L} Rf_*M
\to
Rf_*(K \otimes_{\mathcal{O}_X}^\mathbf{L} M)
\to
Rf_*\omega^\bullet_{X/Y}
\xrightarrow{\text{Tr}_{f, \mathcal{O}_Y}}
\mathcal{O}_Y
$$
détermine un isomorphisme $Rf_*K \to R\SheafHom_{\mathcal{O}_Y}(Rf_*M, \mathcal{O}_Y)$. En effet, puisque $\omega^\bullet_{X/Y} = a(\mathcal{O}_Y)$, l'application canonique (\ref{duality-equation-sheafy-trace})
$$
Rf_*R\SheafHom_{\mathcal{O}_X}(M, \omega^\bullet_{X/Y}) \to
R\SheafHom_{\mathcal{O}_Y}(Rf_*M, \mathcal{O}_Y)
$$
est un isomorphisme d'après le lemme \ref{duality-lemma-iso-on-RSheafHom} et la remarque \ref{duality-remark-iso-on-RSheafHom}, et parce que $M$ et $Rf_*M$ sont pseudo-cohérents, voir Catégories dérivées des schémas, lemme \ref{perfect-lemma-flat-proper-pseudo-coherent-direct-image-general}. Pour voir que le cup-produit relatif induit cet isomorphisme, on utilise la commutativité du diagramme de Cohomologie, remarque \ref{cohomology-remark-relative-cup-and-composition}.
\end{remark}

\begin{lemma}
\label{duality-lemma-properties-relative-dualizing}
Soit $Y$ un schéma quasi-compact et quasi-séparé. Soit $f : X \to Y$ un morphisme de schémas propre, plat et de présentation finie, de complexe dualisant relatif $\omega_{X/Y}^\bullet$ (remarque \ref{duality-remark-relative-dualizing-complex}). Alors
\begin{enumerate}
\item $\omega_{X/Y}^\bullet$ est un objet $Y$-parfait de $D(\mathcal{O}_X)$ ;
\item les faisceaux de cohomologie de $Rf_*\omega_{X/Y}^\bullet$ sont nuls en degrés strictement positifs ;
\item $\mathcal{O}_X \to
R\SheafHom_{\mathcal{O}_X}(\omega_{X/Y}^\bullet, \omega_{X/Y}^\bullet)$ est un isomorphisme.
\end{enumerate}
\end{lemma}

\begin{proof}
Comme la formation de $\omega_{X/Y}^\bullet$ commute au changement de base (voir la remarque \ref{duality-remark-relative-dualizing-complex}), on peut supposer, et l'on suppose, $Y$ affine. Pour un objet parfait $E$ de $D(\mathcal{O}_X)$, on a
\begin{align*}
Rf_*(E \otimes_{\mathcal{O}_X}^\mathbf{L} \omega_{X/Y}^\bullet)
& =
Rf_*R\SheafHom_{\mathcal{O}_X}(E^\vee, \omega_{X/Y}^\bullet) \\
& =
R\SheafHom_{\mathcal{O}_Y}(Rf_*E^\vee, \mathcal{O}_Y) \\
& =
(Rf_*E^\vee)^\vee
\end{align*}
Pour la première égalité, voir Cohomologie, lemme \ref{cohomology-lemma-dual-perfect-complex}. Pour la deuxième, voir le lemme \ref{duality-lemma-iso-on-RSheafHom}, la remarque \ref{duality-remark-iso-on-RSheafHom} et Catégories dérivées des schémas, lemme \ref{perfect-lemma-flat-proper-perfect-direct-image-general}. La troisième égalité est la définition du dual. En particulier, ces références montrent aussi que le résultat est un objet parfait de $D(\mathcal{O}_Y)$. On en conclut que $\omega_{X/Y}^\bullet$ est $Y$-parfait d'après Compléments sur les morphismes, lemme \ref{more-morphisms-lemma-characterize-relatively-perfect}. Cela démontre (1).

\medskip\noindent
Soit $M$ un objet de $D_\QCoh(\mathcal{O}_Y)$. Alors
\begin{align*}
\Hom_Y(M, Rf_*\omega_{X/Y}^\bullet) & =
\Hom_X(Lf^*M, \omega_{X/Y}^\bullet) \\
& =
\Hom_Y(Rf_*Lf^*M, \mathcal{O}_Y) \\
& =
\Hom_Y(M \otimes_{\mathcal{O}_Y}^\mathbf{L} Rf_*\mathcal{O}_X, \mathcal{O}_Y)
\end{align*}
La première égalité résulte de Cohomologie, lemme \ref{cohomology-lemma-adjoint}. La deuxième résulte de la construction de $a$. La troisième résulte de Catégories dérivées des schémas, lemme \ref{perfect-lemma-cohomology-base-change}. Rappelons que $Rf_*\mathcal{O}_X$ est parfait de Tor-amplitude dans $[0, N]$ pour un certain $N$, voir Catégories dérivées des schémas, lemme \ref{perfect-lemma-flat-proper-perfect-direct-image-general}. On peut donc représenter $Rf_*\mathcal{O}_X$ par un complexe de modules projectifs de type fini situés en degrés $[0, N]$ (en utilisant Compléments d'algèbre, lemme \ref{more-algebra-lemma-perfect}, et le fait que $Y$ est affine). Par conséquent, si $M = \mathcal{O}_Y[-i]$ pour un certain $i > 0$, le dernier groupe est nul. Puisque $Y$ est affine, on en conclut que $H^i(Rf_*\omega_{X/Y}^\bullet) = 0$ pour $i > 0$. Cela démontre (2).

\medskip\noindent
Soit $E$ un objet parfait de $D_\QCoh(\mathcal{O}_X)$. On a alors
\begin{align*}
\Hom_X(E, R\SheafHom_{\mathcal{O}_X}(\omega_{X/Y}^\bullet, \omega_{X/Y}^\bullet)
& =
\Hom_X(E \otimes_{\mathcal{O}_X}^\mathbf{L} \omega_{X/Y}^\bullet,
\omega_{X/Y}^\bullet) \\
& =
\Hom_Y(Rf_*(E \otimes_{\mathcal{O}_X}^\mathbf{L} \omega_{X/Y}^\bullet),
\mathcal{O}_Y) \\
& =
\Hom_Y(Rf_*(R\SheafHom_{\mathcal{O}_X}(E^\vee, \omega_{X/Y}^\bullet)),
\mathcal{O}_Y) \\
& =
\Hom_Y(R\SheafHom_{\mathcal{O}_Y}(Rf_*E^\vee, \mathcal{O}_Y),
\mathcal{O}_Y) \\
& =
R\Gamma(Y, Rf_*E^\vee) \\
& =
\Hom_X(E, \mathcal{O}_X)
\end{align*}
La première égalité résulte de Cohomologie, lemme \ref{cohomology-lemma-internal-hom}. La deuxième est la définition de $\omega_{X/Y}^\bullet$. La troisième provient de la construction du complexe parfait dual $E^\vee$, voir Cohomologie, lemme \ref{cohomology-lemma-dual-perfect-complex}. La quatrième résulte de l'égalité
\par\noindent\hfil
$Rf_*R\SheafHom_{\mathcal{O}_X}(E^\vee, \omega_{X/Y}^\bullet) =
R\SheafHom_{\mathcal{O}_Y}(Rf_*E^\vee, \mathcal{O}_Y)$
\hfil\par\noindent
montrée au premier paragraphe de la démonstration. La cinquième résulte de la bidualité des complexes parfaits (Cohomologie, lemme \ref{cohomology-lemma-dual-perfect-complex}) et du fait que $Rf_*E$ est parfait d'après Catégories dérivées des schémas, lemme \ref{perfect-lemma-flat-proper-perfect-direct-image-general}. La dernière égalité est celle de Leray pour $f$. Cette suite d'égalités démontre essentiellement (3) par le lemme de Yoneda. En effet, l'objet
$R\SheafHom(\omega_{X/Y}^\bullet, \omega_{X/Y}^\bullet)$
appartient à $D_\QCoh(\mathcal{O}_X)$ d'après Catégories dérivées des schémas, lemme \ref{perfect-lemma-quasi-coherence-internal-hom}. En prenant $E = \mathcal{O}_X$ ci-dessus, on obtient une application
$\alpha : \mathcal{O}_X \to
R\SheafHom_{\mathcal{O}_X}(\omega_{X/Y}^\bullet, \omega_{X/Y}^\bullet)$
correspondant à $\text{id}_{\mathcal{O}_X} \in \Hom_X(\mathcal{O}_X, \mathcal{O}_X)$. Comme tous les isomorphismes ci-dessus sont fonctoriels en $E$, le cône de $\alpha$ est un objet $C$ de $D_\QCoh(\mathcal{O}_X)$ tel que $\Hom(E, C) = 0$ pour tout $E$ parfait. Puisque les objets parfaits engendrent la catégorie (Catégories dérivées des schémas, théorème \ref{perfect-theorem-bondal-van-den-Bergh}), on en conclut que $\alpha$ est un isomorphisme.
\end{proof}

\begin{lemma}[Rigidité]
\label{duality-lemma-van-den-bergh}
Soit $Y$ un schéma quasi-compact et quasi-séparé. Soit $f : X \to Y$ un morphisme propre et plat de présentation finie, de complexe dualisant relatif $\omega_{X/Y}^\bullet$ (remarque \ref{duality-remark-relative-dualizing-complex}). Il existe un isomorphisme canonique
\begin{equation}
\label{duality-equation-pre-rigid}
\mathcal{O}_X =
c(L\text{pr}_1^*\omega_{X/Y}^\bullet) =
c(L\text{pr}_2^*\omega_{X/Y}^\bullet)
\end{equation}
et un isomorphisme canonique
\begin{equation}
\label{duality-equation-rigid}
\omega_{X/Y}^\bullet =
c\left(L\text{pr}_1^*\omega_{X/Y}^\bullet
\otimes_{\mathcal{O}_{X \times_Y X}}^\mathbf{L}
L\text{pr}_2^*\omega_{X/Y}^\bullet\right)
\end{equation}
où $c$ est l'adjoint à droite du lemme \ref{duality-lemma-twisted-inverse-image} pour la diagonale $\Delta : X \to X \times_Y X$.
\end{lemma}

\begin{proof}
Soit $a$ l'adjoint à droite de $Rf_*$ du lemme \ref{duality-lemma-twisted-inverse-image}. Considérons le carré cartésien
$$
\xymatrix{
X \times_Y X \ar[r]_q \ar[d]_p & X \ar[d]_f \\
X \ar[r]^f & Y
}
$$
Soit $b$ l'adjoint à droite pour $p$ du lemme \ref{duality-lemma-twisted-inverse-image}. Alors
\begin{align*}
\omega_{X/Y}^\bullet
& =
c(b(\omega_{X/Y}^\bullet)) \\
& =
c(Lp^*\omega_{X/Y}^\bullet
\otimes_{\mathcal{O}_{X \times_Y X}}^\mathbf{L} b(\mathcal{O}_X)) \\
& =
c(Lp^*\omega_{X/Y}^\bullet
\otimes_{\mathcal{O}_{X \times_Y X}}^\mathbf{L}
Lq^*a(\mathcal{O}_Y)) \\
& =
c(Lp^*\omega_{X/Y}^\bullet
\otimes_{\mathcal{O}_{X \times_Y X}}^\mathbf{L}
Lq^*\omega_{X/Y}^\bullet)
\end{align*}
comme en (\ref{duality-equation-rigid}). Les justifications sont les suivantes :
\begin{enumerate}
\item La première égalité résulte de $\text{id} = c \circ b$, car $\text{id}_X = p \circ \Delta$.
\item La deuxième résulte du lemme \ref{duality-lemma-compare-with-pullback-flat-proper}.
\item La troisième résulte du lemme \ref{duality-lemma-proper-flat-base-change} et du fait que $\mathcal{O}_X = Lf^*\mathcal{O}_Y$.
\item La quatrième résulte de $\omega_{X/Y}^\bullet = a(\mathcal{O}_Y)$.
\end{enumerate}
L'équation (\ref{duality-equation-pre-rigid}) se démontre exactement de la même manière.
\end{proof}

\begin{remark}
\label{duality-remark-van-den-bergh}
Le lemme \ref{duality-lemma-van-den-bergh} signifie que notre complexe dualisant relatif est {\it rigide}, en un sens analogue à la notion introduite dans \cite{vdB-rigid}. En effet, le foncteur de droite dans (\ref{duality-equation-rigid}) étant « quadratique » en $\omega_{X/Y}^\bullet$, tandis que celui de gauche dans (\ref{duality-equation-rigid}) est « linéaire », cela détermine dans une certaine mesure le complexe $\omega_{X/Y}^\bullet$. Il existe une approche de la théorie de la dualité au moyen de complexes dualisants (relatifs) « rigides », voir par exemple \cite{Neeman-rigid}, \cite{Yekutieli-rigid} et \cite{Yekutieli-Zhang}. Nous y reviendrons à la section \ref{duality-section-relative-dualizing-complexes}.
\end{remark}





\section{Adjoint à droite de l'image directe pour les morphismes parfaits et propres}
\label{duality-section-perfect-proper}

\noindent
Le cadre général approprié à cette section serait celui des morphismes parfaits et propres entre schémas quasi-compacts et quasi-séparés, voir \cite{LN}.

\begin{lemma}
\label{duality-lemma-proper-flat-noetherian}
Soit $f : X \to Y$ un morphisme parfait et propre de schémas noethériens. Soit $a$ l'adjoint à droite de $Rf_* : D_\QCoh(\mathcal{O}_X) \to D_\QCoh(\mathcal{O}_Y)$ du lemme \ref{duality-lemma-twisted-inverse-image}. Alors $a$ commute aux sommes directes.
\end{lemma}

\begin{proof}
Soit $P$ un objet parfait de $D(\mathcal{O}_X)$. D'après Compléments sur les morphismes, lemme \ref{more-morphisms-lemma-perfect-proper-perfect-direct-image}, le complexe $Rf_*P$ est parfait sur $Y$. Soit $K_i$ une famille d'objets de $D_\QCoh(\mathcal{O}_Y)$. Alors
\begin{align*}
\Hom_{D(\mathcal{O}_X)}(P, a(\bigoplus K_i))
& =
\Hom_{D(\mathcal{O}_Y)}(Rf_*P, \bigoplus K_i) \\
& =
\bigoplus \Hom_{D(\mathcal{O}_Y)}(Rf_*P, K_i) \\
& =
\bigoplus \Hom_{D(\mathcal{O}_X)}(P, a(K_i))
\end{align*}
car un objet parfait est compact (Catégories dérivées des schémas, proposition \ref{perfect-proposition-compact-is-perfect}). Puisque $D_\QCoh(\mathcal{O}_X)$ admet un générateur parfait (Catégories dérivées des schémas, théorème \ref{perfect-theorem-bondal-van-den-Bergh}), on en conclut que l'application $\bigoplus a(K_i) \to a(\bigoplus K_i)$ est un isomorphisme, c'est-à-dire que $a$ commute aux sommes directes.
\end{proof}

\begin{lemma}
\label{duality-lemma-proper-flat-noetherian-relative}
Soit $f : X \to Y$ un morphisme parfait et propre de schémas noethériens. Soit $a$ l'adjoint à droite de $Rf_* : D_\QCoh(\mathcal{O}_X) \to D_\QCoh(\mathcal{O}_Y)$ du lemme \ref{duality-lemma-twisted-inverse-image}. Alors
\begin{enumerate}
\item pour tout fermé $T \subset Y$, si $Q \in D_\QCoh(Y)$ est à support dans $T$, alors $a(Q)$ est à support dans $f^{-1}(T)$ ;
\item pour tout ouvert $V \subset Y$ et tout $K \in D_\QCoh(\mathcal{O}_Y)$, l'application (\ref{duality-equation-sheafy}) est un isomorphisme, et
\end{enumerate}
\end{lemma}

\begin{proof}
Cela résulte des lemmes \ref{duality-lemma-when-sheafy}, \ref{duality-lemma-proper-noetherian} et \ref{duality-lemma-proper-flat-noetherian}.
\end{proof}

\begin{lemma}
\label{duality-lemma-compare-with-pullback-flat-proper-noetherian}
Soit $f : X \to Y$ un morphisme parfait et propre de schémas noethériens. L'application (\ref{duality-equation-compare-with-pullback}) est un isomorphisme pour tout objet $K$ de $D_\QCoh(\mathcal{O}_Y)$.
\end{lemma}

\begin{proof}
D'après le lemme \ref{duality-lemma-proper-flat-noetherian}, $a$ commute aux sommes directes. La collection des objets de $D_\QCoh(\mathcal{O}_Y)$ pour lesquels (\ref{duality-equation-compare-with-pullback}) est un isomorphisme est donc une sous-catégorie triangulée, saturée et strictement pleine de $D_\QCoh(\mathcal{O}_Y)$, qui est de plus stable par sommes directes. Puisque $D_\QCoh(\mathcal{O}_Y)$ est une catégorie de modules (Catégories dérivées des schémas, théorème \ref{perfect-theorem-DQCoh-is-Ddga}) engendrée par un unique objet parfait (Catégories dérivées des schémas, théorème \ref{perfect-theorem-bondal-van-den-Bergh}), on peut raisonner comme dans Compléments d'algèbre, remarque \ref{more-algebra-remark-P-resolution}, pour voir qu'il suffit de montrer que (\ref{duality-equation-compare-with-pullback}) est un isomorphisme pour un unique objet parfait. Or le résultat vaut pour les objets parfaits, voir le lemme \ref{duality-lemma-compare-with-pullback-perfect}.
\end{proof}

\begin{lemma}
\label{duality-lemma-proper-perfect-base-change}
Soit $f : X \to Y$ un morphisme parfait et propre de schémas noethériens. Soit $g : Y' \to Y$ un morphisme, avec $Y'$ noethérien. Si $X$ et $Y'$ sont Tor-indépendants sur $Y$, l'application de changement de base (\ref{duality-equation-base-change-map}) est un isomorphisme pour tout $K \in D_\QCoh(\mathcal{O}_Y)$.
\end{lemma}

\begin{proof}
D'après le lemme \ref{duality-lemma-proper-flat-noetherian-relative}, la formation des foncteurs $a$ et $a'$ commute à la restriction aux ouverts de $Y$ et $Y'$. On peut donc supposer que $Y' \to Y$ est un morphisme de schémas affines, voir la remarque \ref{duality-remark-check-over-affines}. Dans ce cas, l'énoncé résulte du lemme \ref{duality-lemma-more-base-change}.
\end{proof}



\section{Adjoint à droite de l'image directe pour les diviseurs de Cartier effectifs}
\label{duality-section-dualizing-Cartier}

\noindent
Soit $X$ un schéma et soit $i : D \to X$ l'inclusion d'un diviseur de Cartier effectif. Notons $\mathcal{N} = i^*\mathcal{O}_X(D)$ le faisceau normal de $i$, voir Morphismes, section \ref{morphisms-section-conormal-sheaf}, et Diviseurs, section \ref{divisors-section-effective-Cartier-divisors}. Rappelons que $R\SheafHom(\mathcal{O}_D, -)$ désigne l'adjoint à droite de $i_* : D(\mathcal{O}_D) \to D(\mathcal{O}_X)$ et vérifie $i_*R\SheafHom(\mathcal{O}_D, -) =
R\SheafHom_{\mathcal{O}_X}(i_*\mathcal{O}_D, -)$, voir la section \ref{duality-section-sections-with-exact-support}.

\begin{lemma}
\label{duality-lemma-compute-for-effective-Cartier}
Comme ci-dessus, soit $X$ un schéma et soit $D \subset X$ un diviseur de Cartier effectif. Il existe un isomorphisme canonique $R\SheafHom(\mathcal{O}_D, \mathcal{O}_X) = \mathcal{N}[-1]$ dans $D(\mathcal{O}_D)$.
\end{lemma}

\begin{proof}
Cela revient à dire que $R\SheafHom(\mathcal{O}_D, \mathcal{O}_X)$ possède un unique faisceau de cohomologie non nul, en degré $1$, et que ce faisceau est isomorphe à $\mathcal{N}$. Puisque $i_*$ est exact et pleinement fidèle, il suffit de montrer que $i_*R\SheafHom(\mathcal{O}_D, \mathcal{O}_X)$ est isomorphe à $i_*\mathcal{N}[-1]$. On a
$i_*R\SheafHom(\mathcal{O}_D, \mathcal{O}_X) =
R\SheafHom_{\mathcal{O}_X}(i_*\mathcal{O}_D, \mathcal{O}_X)$
d'après le lemme \ref{duality-lemma-sheaf-with-exact-support-ext}. On dispose d'une résolution
$$
0 \to \mathcal{I} \to \mathcal{O}_X \to i_*\mathcal{O}_D \to 0
$$
où $\mathcal{I}$ est le faisceau d'idéaux de $D$, qui permet d'effectuer le calcul. Puisque
\par\noindent\hfil
$R\SheafHom_{\mathcal{O}_X}(\mathcal{O}_X, \mathcal{O}_X) = \mathcal{O}_X$
\hfil\par\noindent
et $R\SheafHom_{\mathcal{O}_X}(\mathcal{I}, \mathcal{O}_X) = \mathcal{O}_X(D)$, par un calcul local, on voit que
$$
R\SheafHom_{\mathcal{O}_X}(i_*\mathcal{O}_D, \mathcal{O}_X) =
(\mathcal{O}_X \to \mathcal{O}_X(D))
$$
où le membre de droite comporte $\mathcal{O}_X$ en degré $0$ et $\mathcal{O}_X(D)$ en degré $1$. Le résultat découle de la suite exacte courte
$$
0 \to \mathcal{O}_X \to \mathcal{O}_X(D) \to i_*\mathcal{N} \to 0
$$
qui provient du fait que $D$ est le schéma des zéros de la section canonique de $\mathcal{O}_X(D)$, et du fait que $\mathcal{N} = i^*\mathcal{O}_X(D)$.
\end{proof}

\noindent
Pour tout objet $K$ de $D(\mathcal{O}_X)$, il existe une application canonique
\begin{equation}
\label{duality-equation-map-effective-Cartier}
Li^*K
\otimes_{\mathcal{O}_D}^\mathbf{L}
R\SheafHom(\mathcal{O}_D, \mathcal{O}_X)
\longrightarrow
R\SheafHom(\mathcal{O}_D, K)
\end{equation}
dans $D(\mathcal{O}_D)$, fonctorielle en $K$ et compatible avec les triangles distingués. Cette application est adjointe à une application
$$
i_*(Li^*K \otimes^\mathbf{L}_{\mathcal{O}_D}
R\SheafHom(\mathcal{O}_D, \mathcal{O}_X)) =
K \otimes^\mathbf{L}_{\mathcal{O}_X}
R\SheafHom_{\mathcal{O}_X}(i_*\mathcal{O}_D, \mathcal{O}_X)
\longrightarrow K
$$
où l'égalité est celle de Cohomologie, lemme \ref{cohomology-lemma-projection-formula-closed-immersion}, et où la flèche provient de l'application canonique
$R\SheafHom_{\mathcal{O}_X}(i_*\mathcal{O}_D, \mathcal{O}_X) \to \mathcal{O}_X$
induite par $\mathcal{O}_X \to i_*\mathcal{O}_D$.

\medskip\noindent
Si $K \in D_\QCoh(\mathcal{O}_X)$, alors (\ref{duality-equation-map-effective-Cartier}) est égale à (\ref{duality-equation-compare-with-pullback}) via l'identification $a(K) = R\SheafHom(\mathcal{O}_D, K)$ du lemme \ref{duality-lemma-twisted-inverse-image-closed}. Si $K \in D_\QCoh(\mathcal{O}_X)$ et si $X$ est noethérien, le lemme suivant est un cas particulier du lemme \ref{duality-lemma-compare-with-pullback-flat-proper-noetherian}.

\begin{lemma}
\label{duality-lemma-sheaf-with-exact-support-effective-Cartier}
Comme ci-dessus, soit $X$ un schéma et soit $D \subset X$ un diviseur de Cartier effectif. Alors (\ref{duality-equation-map-effective-Cartier}), combinée au lemme \ref{duality-lemma-compute-for-effective-Cartier}, définit un isomorphisme
$$
Li^*K \otimes_{\mathcal{O}_D}^\mathbf{L} \mathcal{N}[-1]
\longrightarrow
R\SheafHom(\mathcal{O}_D, K)
$$
fonctoriel en $K$ dans $D(\mathcal{O}_X)$.
\end{lemma}

\begin{proof}
Puisque $i_*$ est exact et pleinement fidèle sur les modules, il suffit, pour montrer que l'application est un isomorphisme, de le vérifier après application de $i_*$. Nous utiliserons sans autre mention les suites exactes courtes
$0 \to \mathcal{I} \to \mathcal{O}_X \to i_*\mathcal{O}_D \to 0$
et
$0 \to \mathcal{O}_X \to \mathcal{O}_X(D) \to i_*\mathcal{N} \to 0$
de la démonstration du lemme \ref{duality-lemma-compute-for-effective-Cartier}. D'après Cohomologie, lemme \ref{cohomology-lemma-projection-formula-closed-immersion}, utilisé pour définir l'application (\ref{duality-equation-map-effective-Cartier}), le membre de gauche devient
$$
K \otimes_{\mathcal{O}_X}^\mathbf{L} i_*\mathcal{N}[-1] =
K \otimes_{\mathcal{O}_X}^\mathbf{L} (\mathcal{O}_X \to \mathcal{O}_X(D))
$$
Le membre de droite devient
\begin{align*}
R\SheafHom_{\mathcal{O}_X}(i_*\mathcal{O}_D, K) & =
R\SheafHom_{\mathcal{O}_X}((\mathcal{I} \to \mathcal{O}_X), K) \\
& =
R\SheafHom_{\mathcal{O}_X}((\mathcal{I} \to \mathcal{O}_X), \mathcal{O}_X)
\otimes_{\mathcal{O}_X}^\mathbf{L} K
\end{align*}
la dernière égalité résultant de Cohomologie, lemme \ref{cohomology-lemma-dual-perfect-complex}. Puisque l'application provient de l'isomorphisme
$$
R\SheafHom_{\mathcal{O}_X}((\mathcal{I} \to \mathcal{O}_X), \mathcal{O}_X)
= (\mathcal{O}_X \to \mathcal{O}_X(D))
$$
le lemme est clair.
\end{proof}





\section{Exemples d'adjoints à droite de l'image directe}
\label{duality-section-examples}

\noindent
Dans cette section, nous calculons l'adjoint à droite de l'image directe dans quelques exemples. Les isomorphismes sont canoniques, mais seulement au sens le plus faible : nous ne démontrons ni n'affirmons leur compatibilité avec les diverses opérations telles que le changement de base et la composition des morphismes. Ces questions font l'objet d'une vaste littérature ; on pourra commencer par \cite{RD}, \cite{Conrad-GD} (ces références adoptent un autre point de départ pour la dualité, mais abordent la construction de représentants canoniques des complexes dualisants relatifs), puis consulter les travaux de Joseph Lipman et de ses collaborateurs.

\begin{lemma}
\label{duality-lemma-upper-shriek-P1}
Soit $Y$ un schéma quasi-compact et quasi-séparé. Soit $\mathcal{E}$ un $\mathcal{O}_Y$-module localement libre de rang fini $n + 1$, de déterminant $\mathcal{L} = \wedge^{n + 1}(\mathcal{E})$. Soit $f : X = \mathbf{P}(\mathcal{E}) \to Y$ la projection. Soit $a$ l'adjoint à droite de $Rf_* : D_\QCoh(\mathcal{O}_X) \to D_\QCoh(\mathcal{O}_Y)$ du lemme \ref{duality-lemma-twisted-inverse-image}. Il existe alors un isomorphisme
$$
c : f^*\mathcal{L}(-n - 1)[n] \longrightarrow a(\mathcal{O}_Y)
$$
En particulier, si $\mathcal{E} = \mathcal{O}_Y^{\oplus n + 1}$, alors $X = \mathbf{P}^n_Y$ et l'on obtient $a(\mathcal{O}_Y) = \mathcal{O}_X(-n - 1)[n]$.
\end{lemma}

\begin{proof}
Dans (la démonstration de) Cohomologie des schémas, lemme \ref{coherent-lemma-cohomology-projective-bundle}, nous avons construit un isomorphisme canonique
$$
R^nf_*(f^*\mathcal{L}(-n - 1)) \longrightarrow \mathcal{O}_Y
$$
De plus, $Rf_*(f^*\mathcal{L}(-n - 1))[n] = R^nf_*(f^*\mathcal{L}(-n - 1))$, c'est-à-dire que les autres images directes supérieures sont nulles. On obtient donc un isomorphisme
$$
Rf_*(f^*\mathcal{L}(-n - 1)[n]) \longrightarrow \mathcal{O}_Y
$$
Cet isomorphisme détermine $c$ comme dans l'énoncé, car $a$ est l'adjoint à droite de $Rf_*$. D'après le lemme \ref{duality-lemma-proper-noetherian}, la construction des $a$ est locale sur la base. En particulier, pour vérifier que $c$ est un isomorphisme, on peut travailler localement sur $Y$. Autrement dit, on peut supposer $Y$ affine et $\mathcal{E} = \mathcal{O}_Y^{\oplus n + 1}$. Dans ce cas, les faisceaux $\mathcal{O}_X, \mathcal{O}_X(-1), \ldots,
\mathcal{O}_X(-n)$ engendrent $D_\QCoh(X)$, voir Catégories dérivées des schémas, lemme \ref{perfect-lemma-generator-P1}. Il suffit donc de montrer que
$c : \mathcal{O}_X(-n - 1)[n] \to a(\mathcal{O}_Y)$
est transformée en un isomorphisme par les foncteurs
$$
F_{i, p}(-) = \Hom_{D(\mathcal{O}_X)}(\mathcal{O}_X(i), (-)[p])
$$
pour $i \in \{-n, \ldots, 0\}$ et $p \in \mathbf{Z}$. Pour $F_{0, p}$, c'est vrai par construction de la flèche $c$ ! Pour $i \in \{-n, \ldots, -1\}$, on a
$$
\Hom_{D(\mathcal{O}_X)}(\mathcal{O}_X(i), \mathcal{O}_X(-n - 1)[n + p]) =
H^p(X, \mathcal{O}_X(-n - 1 - i)) = 0
$$
par le calcul de la cohomologie de l'espace projectif (Cohomologie des schémas, lemme \ref{coherent-lemma-cohomology-projective-space-over-ring}), et l'on a
$$
\Hom_{D(\mathcal{O}_X)}(\mathcal{O}_X(i), a(\mathcal{O}_Y)[p]) =
\Hom_{D(\mathcal{O}_Y)}(Rf_*\mathcal{O}_X(i), \mathcal{O}_Y[p]) = 0
$$
car $Rf_*\mathcal{O}_X(i) = 0$ d'après le même lemme. La source et le but de $F_{i, p}(c)$ sont donc nuls et $F_{i, p}(c)$ est nécessairement un isomorphisme. Cela achève la démonstration.
\end{proof}

\begin{example}
\label{duality-example-base-change-wrong}
L'application de changement de base (\ref{duality-equation-base-change-map}) n'est pas nécessairement un isomorphisme lorsque $f$ est parfait et propre et que $g$ est parfait. Soit $k$ un corps. Soit $Y = \mathbf{A}^2_k$ et soit $f : X \to Y$ l'éclatement de $Y$ à l'origine. Notons $E \subset X$ le diviseur exceptionnel. On peut factoriser $f$ sous la forme
$$
X \xrightarrow{i} \mathbf{P}^1_Y \xrightarrow{p} Y
$$
On obtient ainsi une factorisation $a = c \circ b$, où $a$, $b$ et $c$ sont les adjoints à droite du lemme \ref{duality-lemma-twisted-inverse-image} de $Rf_*$, $Rp_*$ et $Ri_*$. Notons $\mathcal{O}(n)$ le faisceau structural tordu au sens de Serre sur $\mathbf{P}^1_Y$, et notons $\mathcal{O}_X(n)$ sa restriction à $X$. Remarquons que $X \subset \mathbf{P}^1_Y$ est défini par une équation de degré un, donc $\mathcal{O}(X) = \mathcal{O}(1)$. Le lemme \ref{duality-lemma-upper-shriek-P1} donne $b(\mathcal{O}_Y) = \mathcal{O}(-2)[1]$. D'après le lemme \ref{duality-lemma-twisted-inverse-image-closed}, on a
$$
a(\mathcal{O}_Y) = c(b(\mathcal{O}_Y)) =
c(\mathcal{O}(-2)[1]) =
R\SheafHom(\mathcal{O}_X, \mathcal{O}(-2)[1]) =
\mathcal{O}_X(-1)
$$
La dernière égalité résulte du lemme \ref{duality-lemma-sheaf-with-exact-support-effective-Cartier}. Soit $Y' = \Spec(k)$ l'origine de $Y$. La restriction de $a(\mathcal{O}_Y)$ à $X' = E = \mathbf{P}^1_k$ est un faisceau inversible de degré $-1$ placé en degré cohomologique $0$. Mais, d'autre part,
$a'(\mathcal{O}_{\Spec(k)}) = \mathcal{O}_E(-2)[1]$
est un faisceau inversible de degré $-2$ placé en degré cohomologique $-1$, donc différent. Dans cet exemple, l'hypothèse de Tor-indépendance du lemme \ref{duality-lemma-more-base-change} n'est pas satisfaite.
\end{example}

\begin{lemma}
\label{duality-lemma-ext}
Soit $Y$ un espace annelé. Soit $\mathcal{I} \subset \mathcal{O}_Y$ un faisceau d'idéaux. Posons $\mathcal{O}_X = \mathcal{O}_Y/\mathcal{I}$ et $\mathcal{N} =
\SheafHom_{\mathcal{O}_Y}(\mathcal{I}/\mathcal{I}^2, \mathcal{O}_X)$. Il existe un isomorphisme canonique $c : \mathcal{N} \to
\SheafExt^1_{\mathcal{O}_Y}(\mathcal{O}_X, \mathcal{O}_X)
$.
\end{lemma}

\begin{proof}
Considérons la suite exacte courte canonique
\begin{equation}
\label{duality-equation-second-order-thickening}
0 \to \mathcal{I}/\mathcal{I}^2 \to \mathcal{O}_Y/\mathcal{I}^2 \to
\mathcal{O}_X \to 0
\end{equation}
Soit $U \subset X$ un ouvert et soit $s \in \mathcal{N}(U)$. On peut former la somme amalgamée de (\ref{duality-equation-second-order-thickening}) suivant $s$ pour obtenir une extension $E_s$ de $\mathcal{O}_X|_U$ par $\mathcal{O}_X|_U$. Celle-ci définit à son tour une section $c(s)$ de
$\SheafExt^1_{\mathcal{O}_Y}(\mathcal{O}_X, \mathcal{O}_X)$
sur $U$. Voir Cohomologie, lemme \ref{cohomology-lemma-section-RHom-over-U}, et Catégories dérivées, lemme \ref{derived-lemma-ext-1}. Réciproquement, étant donnée une extension
$$
0 \to \mathcal{O}_X|_U \to \mathcal{E} \to \mathcal{O}_X|_U \to 0
$$
de $\mathcal{O}_U$-modules, on peut trouver un recouvrement ouvert $U = \bigcup U_i$ et des sections $e_i \in \mathcal{E}(U_i)$ d'image $1 \in \mathcal{O}_X(U_i)$. Alors $e_i$ définit une application $\mathcal{O}_Y|_{U_i} \to \mathcal{E}|_{U_i}$ dont le noyau contient $\mathcal{I}^2$. On voit ainsi que $\mathcal{E}|_{U_i}$ provient d'une somme amalgamée comme ci-dessus. Cela montre que $c$ est surjective. Nous omettons la démonstration de l'injectivité.
\end{proof}

\begin{lemma}
\label{duality-lemma-regular-ideal-ext}
Soit $Y$ un espace annelé. Soit $\mathcal{I} \subset \mathcal{O}_Y$ un faisceau d'idéaux. Posons $\mathcal{O}_X = \mathcal{O}_Y/\mathcal{I}$. Si $\mathcal{I}$ est Koszul-régulier (Diviseurs, définition \ref{divisors-definition-regular-ideal-sheaf}), la composition sur $R\SheafHom_{\mathcal{O}_Y}(\mathcal{O}_X, \mathcal{O}_X)$ définit des isomorphismes
$$
\wedge^i(\SheafExt^1_{\mathcal{O}_Y}(\mathcal{O}_X, \mathcal{O}_X))
\longrightarrow
\SheafExt^i_{\mathcal{O}_Y}(\mathcal{O}_X, \mathcal{O}_X)
$$
pour tout $i$.
\end{lemma}

\begin{proof}
Par composition, nous entendons l'application
$$
R\SheafHom_{\mathcal{O}_Y}(\mathcal{O}_X, \mathcal{O}_X)
\otimes_{\mathcal{O}_Y}^\mathbf{L}
R\SheafHom_{\mathcal{O}_Y}(\mathcal{O}_X, \mathcal{O}_X)
\longrightarrow
R\SheafHom_{\mathcal{O}_Y}(\mathcal{O}_X, \mathcal{O}_X)
$$
de Cohomologie, lemme \ref{cohomology-lemma-internal-hom-composition}. Elle induit des applications de multiplication
$$
\SheafExt^a_{\mathcal{O}_Y}(\mathcal{O}_X, \mathcal{O}_X)
\otimes_{\mathcal{O}_Y}
\SheafExt^b_{\mathcal{O}_Y}(\mathcal{O}_X, \mathcal{O}_X)
\longrightarrow
\SheafExt^{a + b}_{\mathcal{O}_Y}(\mathcal{O}_X, \mathcal{O}_X)
$$
À comparer avec Compléments d'algèbre, équation (\ref{more-algebra-equation-simple-tor-product}). Le lemme signifie que l'application induite
$$
\SheafExt^1_{\mathcal{O}_Y}(\mathcal{O}_X, \mathcal{O}_X)
\otimes \ldots \otimes
\SheafExt^1_{\mathcal{O}_Y}(\mathcal{O}_X, \mathcal{O}_X)
\longrightarrow
\SheafExt^i_{\mathcal{O}_Y}(\mathcal{O}_X, \mathcal{O}_X)
$$
se factorise par le produit extérieur et induit alors un isomorphisme. Pour le vérifier, on peut travailler localement sur $Y$. On peut donc supposer qu'il existe des sections globales $f_1, \ldots, f_r$ de $\mathcal{O}_Y$ qui engendrent $\mathcal{I}$ et forment une suite Koszul-régulière. Notons
$$
\mathcal{A} = \mathcal{O}_Y\langle \xi_1, \ldots, \xi_r\rangle
$$
le faisceau de $\mathcal{O}_Y$-algèbres différentielles graduées strictement commutatives qui est une algèbre de polynômes (à puissances divisées) sur les $\xi_1, \ldots, \xi_r$ en degré $-1$ sur $\mathcal{O}_Y$, de différentielle $\text{d}$ définie par $\text{d}\xi_i = f_i$. Notons $\mathcal{A}^\bullet$ le complexe sous-jacent de $\mathcal{O}_Y$-modules, qui est le complexe de Koszul mentionné ci-dessus. L'application canonique
$\mathcal{A}^\bullet \to \mathcal{O}_X$
est donc un quasi-isomorphisme. On obtient des quasi-isomorphismes
$$
R\SheafHom_{\mathcal{O}_Y}(\mathcal{O}_X, \mathcal{O}_X) \to
\SheafHom^\bullet(\mathcal{A}^\bullet, \mathcal{A}^\bullet) \to
\SheafHom^\bullet(\mathcal{A}^\bullet, \mathcal{O}_X)
$$
d'après Cohomologie, lemme \ref{cohomology-lemma-Rhom-strictly-perfect}. Les différentielles du dernier complexe sont nulles, d'où
$$
\SheafExt^i_{\mathcal{O}_Y}(\mathcal{O}_X, \mathcal{O}_X)
\cong \SheafHom_{\mathcal{O}_Y}(\mathcal{A}^{-i}, \mathcal{O}_X)
$$
Pour $j \in \{1, \ldots, r\}$, soit $\delta_j : \mathcal{A} \to \mathcal{A}$ la dérivation de degré $1$ telle que $\delta_j(\xi_i) = \delta_{ij}$ (symbole de Kronecker). Un calcul donne $\delta_j \circ \text{d} = - \text{d} \circ \delta_j$, ce qui fournit un morphisme de complexes
$$
\delta_j : \mathcal{A}^\bullet \to \mathcal{A}^\bullet[1].
$$
Ainsi, $\delta_j$ définit une section du faisceau $\SheafExt$ correspondant. Un autre calcul montre que les images de $\delta_1, \ldots, \delta_r$ forment une base de $\SheafHom_{\mathcal{O}_Y}(\mathcal{A}^{-1}, \mathcal{O}_X)$ sur $\mathcal{O}_X$. Puisqu'il est clair que $\delta_j \circ \delta_j = 0$ et $\delta_j \circ \delta_{j'} = - \delta_{j'} \circ \delta_j$ comme endomorphismes de $\mathcal{A}$, donc dans les faisceaux $\SheafExt$, on obtient la factorisation de notre application par la puissance extérieure. Pour vérifier que l'on obtient l'isomorphisme voulu, le lecteur vérifiera que les images des éléments
$$
\delta_{j_1} \circ \ldots \circ \delta_{j_i}
$$
pour $j_1 < \ldots < j_i$ forment une base du faisceau
$\SheafHom_{\mathcal{O}_Y}(\mathcal{A}^{-i}, \mathcal{O}_X)$
sur $\mathcal{O}_X$.
\end{proof}

\begin{lemma}
\label{duality-lemma-regular-immersion-ext}
Soit $Y$ un espace annelé. Soit $\mathcal{I} \subset \mathcal{O}_Y$ un faisceau d'idéaux. Posons $\mathcal{O}_X = \mathcal{O}_Y/\mathcal{I}$ et $\mathcal{N} =
\SheafHom_{\mathcal{O}_Y}(\mathcal{I}/\mathcal{I}^2, \mathcal{O}_X)$. Si $\mathcal{I}$ est Koszul-régulier (Diviseurs, définition \ref{divisors-definition-regular-ideal-sheaf}), alors
$$
R\SheafHom_{\mathcal{O}_Y}(\mathcal{O}_X, \mathcal{O}_Y) =
\wedge^r \mathcal{N}[-r]
$$
où $r : Y \to \{1, 2, 3, \ldots \}$ associe à $y$ le nombre minimal de générateurs de $\mathcal{I}$ nécessaires dans un voisinage de $y$.
\end{lemma}

\begin{proof}
Les lemmes \ref{duality-lemma-ext} et \ref{duality-lemma-regular-ideal-ext} donnent des isomorphismes
$\wedge^i\mathcal{N} \to
\SheafExt^i_{\mathcal{O}_Y}(\mathcal{O}_X, \mathcal{O}_X)$
pour $i \geq 0$. Il suffit donc de montrer que l'application $\mathcal{O}_Y \to \mathcal{O}_X$ induit un isomorphisme
$$
\SheafExt^r_{\mathcal{O}_Y}(\mathcal{O}_X, \mathcal{O}_Y)
\longrightarrow
\SheafExt^r_{\mathcal{O}_Y}(\mathcal{O}_X, \mathcal{O}_X)
$$
et que
$\SheafExt^i_{\mathcal{O}_Y}(\mathcal{O}_X, \mathcal{O}_Y)$
est nul pour $i \not = r$. Ces énoncés sont locaux sur $Y$. On peut donc supposer qu'il existe des sections globales $f_1, \ldots, f_r$ de $\mathcal{O}_Y$ qui engendrent $\mathcal{I}$ et forment une suite Koszul-régulière. Soit $\mathcal{A}^\bullet$ le complexe de Koszul sur $f_1, \ldots, f_r$ introduit dans la démonstration du lemme \ref{duality-lemma-regular-ideal-ext}. Alors
$$
R\SheafHom_{\mathcal{O}_Y}(\mathcal{O}_X, \mathcal{O}_Y) =
\SheafHom^\bullet(\mathcal{A}^\bullet, \mathcal{O}_Y)
$$
d'après Cohomologie, lemme \ref{cohomology-lemma-Rhom-strictly-perfect}. Notons $1 : \mathcal{A}^\bullet \to \mathcal{O}_Y$ l'application de $\mathcal{O}_Y$-algèbres différentielles graduées donnée par l'identité de $\mathcal{A}^0 = \mathcal{O}_Y \to \mathcal{O}_Y$ en degré $0$. Avec $\delta_j$ comme dans la démonstration du lemme \ref{duality-lemma-regular-ideal-ext}, on obtient un isomorphisme de $\mathcal{O}_Y$-modules gradués
$$
\mathcal{O}_Y\langle \delta_1, \ldots, \delta_r\rangle
\longrightarrow
\SheafHom^\bullet(\mathcal{A}^\bullet, \mathcal{O}_Y)
$$
en envoyant $\delta_{j_1} \ldots \delta_{j_i}$ sur $1 \circ \delta_{j_1} \circ \ldots \circ \delta_{j_i}$ en degré $i$. Par cet isomorphisme, la différentielle du membre de droite induit une différentielle $\text{d}$ sur celui de gauche. Nos règles de signes donnent $\text{d}(1) = - \sum f_j \delta_j$. Puisque $\delta_j : \mathcal{A}^\bullet \to \mathcal{A}^\bullet[1]$ est un morphisme de complexes, on en déduit que
$$
\text{d}(\delta_{j_1} \ldots \delta_{j_i}) =
(- \sum f_j \delta_j )\delta_{j_1} \ldots \delta_{j_i}
$$
Observons que $\text{d} = \sum f_j \delta_j$ sur l'algèbre différentielle graduée $\mathcal{A}$. L'application définie par la règle
$$
1 \circ \delta_{j_1} \ldots \delta_{j_i} \longmapsto
(\delta_{j_1} \circ \ldots \circ \delta_{j_i})(\xi_1 \ldots \xi_r)
$$
définit donc un isomorphisme de complexes
$$
\SheafHom^\bullet(\mathcal{A}^\bullet, \mathcal{O}_Y)
\longrightarrow \mathcal{A}^\bullet[-r]
$$
si $r$ est impair, et commute aux différentielles au signe près si $r$ est pair. Dans tous les cas, ces complexes ont des cohomologies isomorphes, ce qui donne l'annulation cherchée. L'isomorphisme sur la cohomologie en degré $r$ induit par l'application
$$
\SheafHom^\bullet(\mathcal{A}^\bullet, \mathcal{O}_Y)
\longrightarrow
\SheafHom^\bullet(\mathcal{A}^\bullet, \mathcal{O}_X)
$$
s'en déduit aussi directement. (On observe que notre choix de conventions pour les complexes de Koszul intervient effectivement dans la définition de l'isomorphisme $R\SheafHom_{\mathcal{O}_X}(\mathcal{O}_X, \mathcal{O}_Y) =
\wedge^r \mathcal{N}[-r]$.)
\end{proof}

\begin{lemma}
\label{duality-lemma-regular-immersion}
Soit $Y$ un schéma quasi-compact et quasi-séparé. Soit $i : X \to Y$ une immersion fermée Koszul-régulière. Soit $a$ l'adjoint à droite de $Ri_* : D_\QCoh(\mathcal{O}_X) \to D_\QCoh(\mathcal{O}_Y)$ du lemme \ref{duality-lemma-twisted-inverse-image}. Il existe alors un isomorphisme
$$
\wedge^r\mathcal{N}[-r] \longrightarrow a(\mathcal{O}_Y)
$$
où
$\mathcal{N} = \SheafHom_{\mathcal{O}_X}(\mathcal{C}_{X/Y}, \mathcal{O}_X)$
est le faisceau normal de $i$ (Morphismes, section \ref{morphisms-section-conormal-sheaf}), et où $r$ est son rang, considéré comme fonction localement constante sur $X$.
\end{lemma}

\begin{proof}
Rappelons, d'après les lemmes \ref{duality-lemma-twisted-inverse-image-closed} et \ref{duality-lemma-sheaf-with-exact-support-ext}, que $a(\mathcal{O}_Y)$ est un objet de $D_\QCoh(\mathcal{O}_X)$ dont l'image directe sur $Y$ est $R\SheafHom_{\mathcal{O}_Y}(i_*\mathcal{O}_X, \mathcal{O}_Y)$. Le résultat découle donc du lemme \ref{duality-lemma-regular-immersion-ext}.
\end{proof}

\begin{lemma}
\label{duality-lemma-smooth-proper}
Soit $S$ un schéma noethérien. Soit $f : X \to S$ un morphisme lisse et propre de dimension relative $d$. Soit $a$ l'adjoint à droite de $Rf_* : D_\QCoh(\mathcal{O}_X) \to D_\QCoh(\mathcal{O}_S)$ du lemme \ref{duality-lemma-twisted-inverse-image}. Il existe alors un isomorphisme
$$
\wedge^d \Omega_{X/S}[d] \longrightarrow a(\mathcal{O}_S)
$$
dans $D(\mathcal{O}_X)$.
\end{lemma}

\begin{proof}
Posons $\omega_{X/S}^\bullet = a(\mathcal{O}_S)$ comme dans la remarque \ref{duality-remark-relative-dualizing-complex}. Soit $c$ l'adjoint à droite du lemme \ref{duality-lemma-twisted-inverse-image} pour $\Delta : X \to X \times_S X$. Puisque $\Delta$ est la diagonale d'un morphisme lisse, c'est une immersion Koszul-régulière, voir Diviseurs, lemme \ref{divisors-lemma-immersion-smooth-into-smooth-regular-immersion}. En particulier, $\Delta$ est un morphisme parfait et propre (Compléments sur les morphismes, lemme \ref{more-morphisms-lemma-regular-immersion-perfect}), et l'on obtient
\begin{align*}
\mathcal{O}_X
& =
c(L\text{pr}_1^*\omega_{X/S}^\bullet) \\
& =
L\Delta^*(L\text{pr}_1^*\omega_{X/S}^\bullet)
\otimes_{\mathcal{O}_X}^\mathbf{L}
c(\mathcal{O}_{X \times_S X}) \\
& =
\omega_{X/S}^\bullet \otimes_{\mathcal{O}_X}^\mathbf{L}
c(\mathcal{O}_{X \times_S X}) \\
& =
\omega_{X/S}^\bullet
\otimes_{\mathcal{O}_X}^\mathbf{L}
\wedge^d(\mathcal{N}_\Delta)[-d]
\end{align*}
La première égalité est (\ref{duality-equation-pre-rigid}), car $\omega_{X/S}^\bullet = a(\mathcal{O}_S)$. La deuxième résulte du lemme \ref{duality-lemma-compare-with-pullback-flat-proper-noetherian}. La troisième résulte de $\text{pr}_1 \circ \Delta = \text{id}_X$. La quatrième résulte du lemme \ref{duality-lemma-regular-immersion}. Observons que $\wedge^d(\mathcal{N}_\Delta)$ est un $\mathcal{O}_X$-module inversible. Ainsi, $\wedge^d(\mathcal{N}_\Delta)[-d]$ est un objet inversible de $D(\mathcal{O}_X)$, et l'on en conclut que $a(\mathcal{O}_S) = \omega_{X/S}^\bullet = \wedge^d(\mathcal{C}_\Delta)[d]$. Puisque le faisceau conormal $\mathcal{C}_\Delta$ de $\Delta$ est $\Omega_{X/S}$ d'après Morphismes, lemme \ref{morphisms-lemma-differentials-diagonal}, la démonstration est achevée.
\end{proof}







\section{Foncteurs image inverse extraordinaire}
\label{duality-section-upper-shriek}

\noindent
Dans cette section, nous construisons, à l'aide de compactifications, les foncteurs $f^!$ pour les morphismes entre schémas séparés de type fini sur une base noethérienne fixée. Comme il est habituel en dualité cohérente, il faut vérifier la commutativité d'un certain nombre de diagrammes. Après avoir lu la construction, on pourra omettre les démonstrations des lemmes et passer à la section suivante, consacrée aux propriétés des foncteurs image inverse extraordinaire.

\begin{situation}
\label{duality-situation-shriek}
Ici, $S$ est un schéma noethérien et $\textit{FTS}_S$ est la catégorie dont
\begin{enumerate}
\item les objets sont les schémas $X$ sur $S$ dont le morphisme structural $X \to S$ est séparé et de type fini ;
\item les morphismes $f : X \to Y$ entre objets sont les morphismes de schémas sur $S$.
\end{enumerate}
\end{situation}

\noindent
Dans la situation \ref{duality-situation-shriek}, étant donné un morphisme $f : X \to Y$ de $\textit{FTS}_S$, nous définirons un foncteur exact
$$
f^! : D_\QCoh^+(\mathcal{O}_Y) \to D_\QCoh^+(\mathcal{O}_X)
$$
entre catégories triangulées. Choisissons pour cela une compactification $X \to \overline{X}$ sur $Y$, ce qui est possible d'après Compléments sur la platitude, théorème \ref{flat-theorem-nagata} et lemme \ref{flat-lemma-compactifyable}. Notons $\overline{f} : \overline{X} \to Y$ le morphisme structural. Soit
$\overline{a} : D_\QCoh(\mathcal{O}_Y) \to D_\QCoh(\mathcal{O}_{\overline{X}})$
l'adjoint à droite de $R\overline{f}_*$ construit au lemme \ref{duality-lemma-twisted-inverse-image}. On pose alors
$$
f^!K  = \overline{a}(K)|_X
$$
pour $K \in D_\QCoh^+(\mathcal{O}_Y)$. Le résultat est un objet de $D_\QCoh^+(\mathcal{O}_X)$ d'après le lemme \ref{duality-lemma-twisted-inverse-image-bounded-below}.

\begin{lemma}
\label{duality-lemma-shriek-well-defined}
Dans la situation \ref{duality-situation-shriek}, soit $f : X \to Y$ un morphisme de $\textit{FTS}_S$. Le foncteur $f^!$ est indépendant, à isomorphisme canonique près, du choix de la compactification.
\end{lemma}

\begin{proof}
La catégorie des compactifications de $X$ sur $Y$ est définie dans Compléments sur la platitude, section \ref{flat-section-compactify}. Elle est non vide d'après Compléments sur la platitude, théorème \ref{flat-theorem-nagata} et lemme \ref{flat-lemma-compactifyable}. À tout choix d'une compactification
$$
j : X \to \overline{X},\quad \overline{f} : \overline{X} \to Y
$$
la construction précédente associe le foncteur $j^* \circ \overline{a} :
D_\QCoh^+(\mathcal{O}_Y) \to D_\QCoh^+(\mathcal{O}_X)$, où $\overline{a}$ est l'adjoint à droite de $R\overline{f}_*$ construit au lemme \ref{duality-lemma-twisted-inverse-image}.

\medskip\noindent
Supposons donné un morphisme $g : \overline{X}_1 \to \overline{X}_2$ entre compactifications $j_i : X \to \overline{X}_i$ sur $Y$ tel que $g^{-1}(j_2(X)) = j_1(X)$\footnote{Cela peut être faux avec notre définition de compactification. Voir Compléments sur la platitude, section \ref{flat-section-compactify}.}. Soit $\overline{c}$ l'adjoint à droite du lemme \ref{duality-lemma-twisted-inverse-image} pour $g$. Alors $\overline{c} \circ \overline{a}_2 = \overline{a}_1$, car ces foncteurs sont adjoints de $R\overline{f}_{2, *} \circ Rg_* = R(\overline{f}_2 \circ g)_*$. Par (\ref{duality-equation-sheafy}), on dispose d'une transformation canonique
$$
j_1^* \circ \overline{c} \longrightarrow j_2^*
$$
de foncteurs
$D^+_\QCoh(\mathcal{O}_{\overline{X}_2}) \to D^+_\QCoh(\mathcal{O}_X)$
qui est un isomorphisme d'après le lemme \ref{duality-lemma-proper-noetherian}. La composée
$$
j_1^* \circ \overline{a}_1 \longrightarrow
j_1^* \circ \overline{c} \circ \overline{a}_2 \longrightarrow
j_2^* \circ \overline{a}_2
$$
est un isomorphisme de foncteurs que nous noterons $\alpha_g$.

\medskip\noindent
Considérons deux compactifications $j_i : X \to \overline{X}_i$, $i = 1, 2$ de $X$ sur $Y$. D'après Compléments sur la platitude, lemme \ref{flat-lemma-compactifications-cofiltered}, (b), on peut trouver une compactification $j : X \to \overline{X}$ d'image dense et des morphismes $g_i : \overline{X} \to \overline{X}_i$ de compactifications. D'après Compléments sur la platitude, lemme \ref{flat-lemma-compactifications-cofiltered}, (c), on a $g_i^{-1}(j_i(X)) = j(X)$. Le paragraphe précédent donne donc des isomorphismes
$$
\alpha_{g_i} :
j^* \circ \overline{a}
\longrightarrow
j_i^* \circ \overline{a}_i
$$
On obtient un isomorphisme
$$
\alpha_{g_2} \circ \alpha_{g_1}^{-1} :
j_1^* \circ \overline{a}_1 \to j_2^* \circ \overline{a}_2
$$
Pour achever la démonstration, il faut montrer que ces isomorphismes sont bien définis. Nous affirmons qu'il suffit de montrer que la composée d'isomorphismes construits au paragraphe précédent en est encore un (voir le paragraphe suivant pour l'énoncé précis). Le lecteur pourra le vérifier rapidement, mais nous détaillons aussi entièrement l'argument dans la suite de ce paragraphe. Considérons en effet un second choix de compactification $j' : X \to \overline{X}'$ d'image dense, avec des morphismes de compactifications $g'_i : \overline{X}' \to \overline{X}_i$. D'après Compléments sur la platitude, lemme \ref{flat-lemma-compactifications-cofiltered}, on peut trouver une compactification $j'' : X \to \overline{X}''$ d'image dense et des morphismes de compactifications $h : \overline{X}'' \to \overline{X}$ et $h' : \overline{X}'' \to \overline{X}'$. On peut même supposer $g_1 \circ h = g'_1 \circ h'$ et $g_2 \circ h = g'_2 \circ h'$. Le résultat du paragraphe suivant donne
$$
\alpha_{g_i} \circ \alpha_h = \alpha_{g_i \circ h} =
\alpha_{g'_i \circ h'} = \alpha_{g'_i} \circ \alpha_{h'}
$$
pour $i = 1, 2$. Puisque ce sont tous des isomorphismes de foncteurs, on en conclut que $\alpha_{g_2} \circ \alpha_{g_1}^{-1} =
\alpha_{g'_2} \circ \alpha_{g'_1}^{-1}$, comme voulu.

\medskip\noindent
Supposons données des compactifications $j_i : X \to \overline{X}_i$ pour $i = 1, 2, 3$, ainsi que des morphismes $g : \overline{X}_1 \to \overline{X}_2$ et $h : \overline{X}_2 \to \overline{X}_3$ de compactifications tels que $g^{-1}(j_2(X)) = j_1(X)$ et $h^{-1}(j_2(X)) = j_3(X)$. Soit $\overline{a}_i$ comme ci-dessus. L'affirmation précédente signifie que
$$
\alpha_g \circ \alpha_h = \alpha_{g \circ h} :
j_1^* \circ \overline{a}_1 \to j_3^* \circ \overline{a}_3
$$
Soit $\overline{c}$, resp.\ $\overline{d}$, l'adjoint à droite du lemme \ref{duality-lemma-twisted-inverse-image} pour $g$, resp.\ $h$. Alors $\overline{c} \circ \overline{a}_2 = \overline{a}_1$ et $\overline{d} \circ \overline{a}_3 = \overline{a}_2$, et l'on dispose de transformations canoniques
$$
j_1^* \circ \overline{c} \longrightarrow j_2^*
\quad\text{et}\quad
j_2^* \circ \overline{d} \longrightarrow j_3^*
$$
de foncteurs
$D^+_\QCoh(\mathcal{O}_{\overline{X}_2}) \to D^+_\QCoh(\mathcal{O}_X)$
et
$D^+_\QCoh(\mathcal{O}_{\overline{X}_3}) \to D^+_\QCoh(\mathcal{O}_X)$
pour les mêmes raisons que précédemment. Notons $\overline{e}$ l'adjoint à droite du lemme \ref{duality-lemma-twisted-inverse-image} pour $h \circ g$. On dispose d'une transformation canonique
$$
j_1^* \circ \overline{e} \longrightarrow j_3^*
$$
de foncteurs
$D^+_\QCoh(\mathcal{O}_{\overline{X}_3}) \to D^+_\QCoh(\mathcal{O}_X)$
donnée par (\ref{duality-equation-sheafy}). Explicitement, il faut montrer que la composée
$$
\alpha_h \circ \alpha_g :
j_1^* \circ \overline{a}_1 \to
j_1^* \circ \overline{c} \circ \overline{a}_2 \to
j_2^* \circ \overline{a}_2 \to
j_2^* \circ \overline{d} \circ \overline{a}_3 \to
j_3^* \circ \overline{a}_3
$$
est égale à la composée
$$
\alpha_{h \circ g} :
j_1^* \circ \overline{a}_1 \to
j_1^* \circ \overline{e} \circ \overline{a}_3 \to
j_3^* \circ \overline{a}_3
$$
On procède en deux étapes. Il faut d'abord montrer que le diagramme
$$
\xymatrix{
\overline{a}_1 \ar[r] \ar[d] & \overline{c} \circ \overline{a}_2 \ar[d] \\
\overline{e} \circ \overline{a}_3 \ar[r] &
\overline{c} \circ \overline{d} \circ \overline{a}_3
}
$$
est commutatif, où la flèche horizontale inférieure provient de l'identification $\overline{e} = \overline{c} \circ \overline{d}$. Cela résulte de la commutativité du diagramme correspondant de foncteurs image directe totale
$$
\xymatrix{
R\overline{f}_{1, *} \ar[r] \ar[d] & Rg_* \circ R\overline{f}_{2, *} \ar[d] \\
R(h \circ g)_* \circ R\overline{f}_{3, *} \ar[r] &
Rg_* \circ Rh_* \circ R\overline{f}_{3, *}
}
$$
(insérer ici une référence ultérieure). Il faut ensuite montrer que la composée
$$
j_1^* \circ \overline{c} \circ \overline{d} \to
j_2^* \circ \overline{d} \to j_3^*
$$
est égale à l'application
$$
j_1^* \circ \overline{e} \to j_3^*
$$
via l'identification $\overline{e} = \overline{c} \circ \overline{d}$. Cela a été démontré au lemme \ref{duality-lemma-compose-base-change-maps} (notons que, dans le cas présent, les morphismes $f', g'$ de ce lemme sont égaux à $\text{id}_X$).
\end{proof}

\begin{lemma}
\label{duality-lemma-upper-shriek-composition}
Dans la situation \ref{duality-situation-shriek}, soient $f : X \to Y$ et $g : Y \to Z$ deux morphismes composables de $\textit{FTS}_S$. Il existe alors un isomorphisme canonique $(g \circ f)^! \to f^! \circ g^!$.
\end{lemma}

\begin{proof}
Choisissons une compactification $i : Y \to \overline{Y}$ de $Y$ sur $Z$, puis une compactification $X \to \overline{X}$ de $X$ sur $\overline{Y}$. On utilise ici deux fois Compléments sur la platitude, théorème \ref{flat-theorem-nagata} et lemme \ref{flat-lemma-compactifyable}. Soit $\overline{a}$ l'adjoint à droite du lemme \ref{duality-lemma-twisted-inverse-image} pour $\overline{X} \to \overline{Y}$, et soit $\overline{b}$ l'adjoint à droite du lemme \ref{duality-lemma-twisted-inverse-image} pour $\overline{Y} \to Z$. Alors $\overline{a} \circ \overline{b}$ est l'adjoint à droite du lemme \ref{duality-lemma-twisted-inverse-image} pour la composée $\overline{X} \to Z$. Ainsi, $g^! = i^* \circ \overline{b}$ et $(g \circ f)^! = (X \to \overline{X})^* \circ \overline{a} \circ \overline{b}$. Soit $U$ l'image inverse de $Y$ dans $\overline{X}$, de sorte que l'on obtient le diagramme commutatif
$$
\xymatrix{
X \ar[r]_j \ar[d] & U \ar[dl] \ar[r]_{j'} & \overline{X} \ar[dl] \\
Y \ar[r]_i \ar[d] & \overline{Y} \ar[dl] \\
Z
}
$$
Soit $\overline{a}'$ l'adjoint à droite du lemme \ref{duality-lemma-twisted-inverse-image} pour $U \to Y$. Alors $f^! = j^* \circ \overline{a}'$. On obtient
$$
\gamma : (j')^* \circ \overline{a} \to \overline{a}' \circ i^*
$$
par (\ref{duality-equation-sheafy}), ce qui permet de définir
$$
(g \circ f)^! =
(j' \circ j)^* \circ \overline{a} \circ \overline{b} =
j^* \circ (j')^* \circ \overline{a} \circ \overline{b}
\to
j^* \circ \overline{a}' \circ i^* \circ \overline{b} =
f^! \circ g^!
$$
qui est un isomorphisme sur les objets de $D_\QCoh^+(\mathcal{O}_Z)$ d'après le lemme \ref{duality-lemma-proper-noetherian}. Pour achever la démonstration, montrons que cet isomorphisme ne dépend pas des choix effectués.

\medskip\noindent
Supposons donnés deux diagrammes
$$
\vcenter{
\xymatrix{
X \ar[r]_{j_1} \ar[d] & U_1 \ar[dl] \ar[r]_{j'_1} & \overline{X}_1 \ar[dl] \\
Y \ar[r]_{i_1} \ar[d] & \overline{Y}_1 \ar[dl] \\
Z
}
}
\quad\text{et}\quad
\vcenter{
\xymatrix{
X \ar[r]_{j_2} \ar[d] & U_2 \ar[dl] \ar[r]_{j'_2} & \overline{X}_2 \ar[dl] \\
Y \ar[r]_{i_2} \ar[d] & \overline{Y}_2 \ar[dl] \\
Z
}
}
$$
On peut d'abord choisir une compactification $i : Y \to \overline{Y}$ d'image dense de $Y$ sur $Z$ dominant à la fois $\overline{Y}_1$ et $\overline{Y}_2$, voir Compléments sur la platitude, lemme \ref{flat-lemma-compactifications-cofiltered}. D'après Compléments sur la platitude, lemme \ref{flat-lemma-right-multiplicative-system}, et Catégories, lemmes \ref{categories-lemma-morphisms-right-localization} et \ref{categories-lemma-equality-morphisms-right-localization}, on peut choisir une compactification $X \to \overline{X}$ d'image dense de $X$ sur $\overline{Y}$, munie de morphismes $\overline{X} \to \overline{X}_1$ et $\overline{X} \to \overline{X}_2$, telle que la composée $\overline{X} \to \overline{Y} \to \overline{Y}_1$ soit égale à la composée $\overline{X} \to \overline{X}_1 \to \overline{Y}_1$ et que la composée $\overline{X} \to \overline{Y} \to \overline{Y}_2$ soit égale à la composée $\overline{X} \to \overline{X}_2 \to \overline{Y}_2$. Il suffit donc de comparer les applications déterminées par nos diagrammes lorsqu'on dispose d'un diagramme commutatif de la forme
$$
\xymatrix{
X \ar[rr]_{j_1} \ar@{=}[d] & &
U_1 \ar[d] \ar[ddll] \ar[rr]_{j'_1} & &
\overline{X}_1 \ar[d] \ar[ddll] \\
X \ar'[r][rr]^-{j_2} \ar[d] & &
U_2 \ar'[dl][ddll] \ar'[r][rr]^-{j'_2} & &
\overline{X}_2 \ar[ddll] \\
Y \ar[rr]^{i_1} \ar@{=}[d] & & \overline{Y}_1 \ar[d] \\
Y \ar[rr]^{i_2} \ar[d] & & \overline{Y}_2 \ar[dll] \\
Z
}
$$
et que les compactifications $X \to \overline{X}_1$ et $Y \to \overline{Y}_2$ sont en outre d'image dense. On note $\overline{a}_i$, $\overline{a}'_i$, $\overline{c}$ et $\overline{c}'$ les adjoints à droite du lemme \ref{duality-lemma-twisted-inverse-image} pour $\overline{X}_i \to \overline{Y}_i$, $U_i \to Y$, $\overline{X}_1 \to \overline{X}_2$ et $U_1 \to U_2$. Chacun des carrés
$$
\xymatrix{
X \ar[r] \ar[d] \ar@{}[dr]|A & U_1 \ar[d] \\
X \ar[r] & U_2
}
\quad
\xymatrix{
U_2 \ar[r] \ar[d] \ar@{}[dr]|B & \overline{X}_2 \ar[d] \\
Y \ar[r] & \overline{Y}_2
}
\quad
\xymatrix{
U_1 \ar[r] \ar[d] \ar@{}[dr]|C & \overline{X}_1 \ar[d] \\
Y \ar[r] & \overline{Y}_1
}
\quad
\xymatrix{
Y \ar[r] \ar[d] \ar@{}[dr]|D & \overline{Y}_1 \ar[d] \\
Y \ar[r] & \overline{Y}_2
}
\quad
\xymatrix{
X \ar[r] \ar[d] \ar@{}[dr]|E & \overline{X}_1 \ar[d] \\
X \ar[r] & \overline{X}_2
}
$$
est cartésien (voir Compléments sur la platitude, lemme \ref{flat-lemma-compactifications-cofiltered}, (c), pour A, D, E ; pour B, C, rappeler que $U_i$ est l'image inverse de $Y$ par $\overline{X}_i \to \overline{Y}_i$), et donne donc une application de changement de base (\ref{duality-equation-sheafy}) comme suit :
$$
\begin{matrix}
\gamma_A : j_1^* \circ \overline{c}' \to j_2^* &
\gamma_B : (j_2')^* \circ \overline{a}_2 \to \overline{a}'_2 \circ i_2^* &
\gamma_C : (j_1')^* \circ \overline{a}_1 \to \overline{a}'_1 \circ i_1^* \\
\gamma_D : i_1^* \circ \overline{d} \to i_2^* &
\gamma_E : (j'_1 \circ j_1)^* \circ \overline{c} \to (j'_2 \circ j_2)^*
\end{matrix}
$$
Notons $f_1^! = j_1^* \circ \overline{a}'_1$, $f_2^! = j_2^* \circ \overline{a}'_2$, $g_1^! = i_1^* \circ \overline{b}_1$, $g_2^! = i_2^* \circ \overline{b}_2$, $(g \circ f)_1^! =
(j_1' \circ j_1)^* \circ \overline{a}_1 \circ \overline{b}_1$ et $(g \circ f)^!_2 =
(j_2' \circ j_2)^* \circ \overline{a}_2 \circ \overline{b}_2$. La construction du premier paragraphe de la démonstration et du lemme \ref{duality-lemma-shriek-well-defined} utilise
\begin{enumerate}
\item $\gamma_C$ pour l'application $(g \circ f)^!_1 \to f_1^! \circ g_1^!$ ;
\item $\gamma_B$ pour l'application $(g \circ f)^!_2 \to f_2^! \circ g_2^!$ ;
\item $\gamma_A$ pour l'application $f_1^! \to f_2^!$ ;
\item $\gamma_D$ pour l'application $g_1^! \to g_2^!$ ;
\item $\gamma_E$ pour l'application $(g \circ f)^!_1 \to (g \circ f)^!_2$.
\end{enumerate}
Il faut montrer que le diagramme
$$
\xymatrix{
(g \circ f)^!_1 \ar[r]_{\gamma_E} \ar[d]_{\gamma_C} &
(g \circ f)^!_2 \ar[d]_{\gamma_B} \\
f_1^! \circ g_1^! \ar[r]^{\gamma_A \circ \gamma_D} & f_2^! \circ g_2^!
}
$$
est commutatif. Nous utiliserons les lemmes \ref{duality-lemma-compose-base-change-maps} et \ref{duality-lemma-compose-base-change-maps-horizontal}, avec les notations (et abus de notation) de la remarque \ref{duality-remark-going-around} (en omettant notamment les produits $\star$ avec les transformations identités). On peut écrire $\gamma_E = \gamma_A \circ \gamma_F$, où
$$
\xymatrix{
U_1 \ar[r] \ar[d] \ar@{}[rd]|F & \overline{X}_1 \ar[d] \\
U_2 \ar[r] & \overline{X}_2
}
$$
On voit donc que
$$
\gamma_B \circ \gamma_E = \gamma_B \circ \gamma_A  \circ \gamma_F
= \gamma_A \circ \gamma_B \circ \gamma_F
$$
la dernière égalité résultant du fait que les deux carrés $A$ et $B$ ne se rencontrent qu'en un sommet (comme dans le dernier argument de la remarque \ref{duality-remark-going-around}). Il suffit donc de montrer que $\gamma_D \circ \gamma_C = \gamma_B \circ \gamma_F$. Puisque ces deux expressions sont égales à l'application (\ref{duality-equation-sheafy}) pour le carré
$$
\xymatrix{
U_1 \ar[r] \ar[d] & \overline{X}_1 \ar[d] \\
Y \ar[r] & \overline{Y}_2
}
$$
on conclut.
\end{proof}

\begin{lemma}
\label{duality-lemma-pseudo-functor}
Dans la situation \ref{duality-situation-shriek}, les constructions des lemmes \ref{duality-lemma-shriek-well-defined} et \ref{duality-lemma-upper-shriek-composition} définissent un pseudo-foncteur de la catégorie $\textit{FTS}_S$ vers la $2$-catégorie des catégories (voir Catégories, définition \ref{categories-definition-functor-into-2-category}).
\end{lemma}

\begin{proof}
Il faut montrer que, pour des morphismes $f : X \to Y$, $g : Y \to Z$, $h : Z \to T$, le diagramme
$$
\xymatrix{
(h \circ g \circ f)^! \ar[r]_{\gamma_{A + B}} \ar[d]_{\gamma_{B + C}} &
f^! \circ (h \circ g)^! \ar[d]^{\gamma_C} \\
(g \circ f)^! \circ h^! \ar[r]^{\gamma_A} & f^! \circ g^! \circ h^!
}
$$
est commutatif (voir ci-dessous pour la signification des $\gamma$). Pour cela, choisissons une compactification $\overline{Z}$ de $Z$ sur $T$, puis une compactification $\overline{Y}$ de $Y$ sur $\overline{Z}$, et enfin une compactification $\overline{X}$ de $X$ sur $\overline{Y}$. On utilise Compléments sur la platitude, théorème \ref{flat-theorem-nagata} et lemme \ref{flat-lemma-compactifyable}. Soit $W \subset \overline{Y}$ l'image inverse de $Z$ par $\overline{Y} \to \overline{Z}$, et soient $U \subset V \subset \overline{X}$ les images inverses de $Y \subset W$ par $\overline{X} \to \overline{Y}$. On obtient le diagramme suivant :
$$
\xymatrix{
X \ar[d]_f \ar[r] & U \ar[r] \ar[d] \ar@{}[dr]|A &
V \ar[d] \ar[r] \ar@{}[rd]|B & \overline{X} \ar[d] \\
Y \ar[d]_g \ar[r] & Y \ar[r] \ar[d] & W \ar[r] \ar[d] \ar@{}[rd]|C &
\overline{Y} \ar[d] \\
Z \ar[d]_h \ar[r] & Z \ar[d] \ar[r] & Z \ar[d] \ar[r] & \overline{Z} \ar[d] \\
T \ar[r] & T \ar[r] & T \ar[r] & T
}
$$
Sans introduire une profusion de notations, mais en raisonnant exactement comme dans la démonstration du lemme \ref{duality-lemma-upper-shriek-composition}, on voit que les applications du premier diagramme affiché utilisent les applications (\ref{duality-equation-sheafy}) associées aux rectangles $A + B$, $B + C$, $A$ et $C$, comme indiqué. D'après les lemmes \ref{duality-lemma-compose-base-change-maps} et \ref{duality-lemma-compose-base-change-maps-horizontal}, on a $\gamma_{A + B} = \gamma_A \circ \gamma_B$ et $\gamma_{B + C} = \gamma_C \circ \gamma_B$ ; l'égalité cherchée est donc vraie pourvu que $\gamma_A \circ \gamma_C = \gamma_C \circ \gamma_A$. Cela résulte du fait que les deux carrés $A$ et $C$ ne se rencontrent qu'en un sommet (comme dans le dernier argument de la remarque \ref{duality-remark-going-around}).
\end{proof}

\begin{lemma}
\label{duality-lemma-map-pullback-to-shriek-well-defined}
Dans la situation \ref{duality-situation-shriek}, soit $f : X \to Y$ un morphisme de $\textit{FTS}_S$. Il existe des applications canoniques
$$
\mu_{f, K} :
Lf^*K \otimes_{\mathcal{O}_X}^\mathbf{L} f^!\mathcal{O}_Y
\longrightarrow
f^!K
$$
fonctorielles en $K$ dans $D^+_\QCoh(\mathcal{O}_Y)$. Si $g : Y \to Z$ est un autre morphisme de $\textit{FTS}_S$, le diagramme
$$
\xymatrix{
Lf^*(Lg^*K \otimes_{\mathcal{O}_Y}^\mathbf{L} g^!\mathcal{O}_Z)
\otimes_{\mathcal{O}_X}^\mathbf{L} f^!\mathcal{O}_Y
\ar@{=}[d] \ar[r]_-{\mu_f} &
f^!(Lg^*K \otimes_{\mathcal{O}_Y}^\mathbf{L} g^!\mathcal{O}_Z)
\ar[r]_-{f^!\mu_g} &
f^!g^!K \ar@{=}[d] \\
Lf^*Lg^*K \otimes_{\mathcal{O}_X}^\mathbf{L} Lf^* g^!\mathcal{O}_Z
\otimes_{\mathcal{O}_X}^\mathbf{L} f^!\mathcal{O}_Y \ar[r]^-{\mu_f} &
Lf^*Lg^*K \otimes_{\mathcal{O}_X}^\mathbf{L} f^!g^!\mathcal{O}_Z
\ar[r]^-{\mu_{g \circ f}} & f^!g^!K
}
$$
est commutatif pour tout $K \in D^+_\QCoh(\mathcal{O}_Z)$.
\end{lemma}

\begin{proof}
Si $f$ est propre, alors $f^! = a$ et l'on peut utiliser (\ref{duality-equation-compare-with-pullback}) ; si $g$ est également propre, le lemme \ref{duality-lemma-transitivity-compare-with-pullback} démontre la commutativité du diagramme (dans un cadre plus général).

\medskip\noindent
Définissons l'application $\mu_{f, K}$. Choisissons une compactification $j : X \to \overline{X}$ de $X$ sur $Y$. Puisque $f^!$ est défini comme $j^* \circ \overline{a}$, on obtient $\mu_{f, K}$ en restreignant l'application (\ref{duality-equation-compare-with-pullback})
$$
L\overline{f}^*K \otimes_{\mathcal{O}_{\overline{X}}}^\mathbf{L}
\overline{a}(\mathcal{O}_Y)
\longrightarrow
\overline{a}(K)
$$
à $X$. Pour voir que cela ne dépend pas du choix de la compactification, on raisonne comme dans la démonstration du lemme \ref{duality-lemma-shriek-well-defined}. Nous recommandons au lecteur de lire d'abord cette démonstration.

\medskip\noindent
Supposons donné un morphisme $g : \overline{X}_1 \to \overline{X}_2$ entre compactifications $j_i : X \to \overline{X}_i$ sur $Y$ tel que $g^{-1}(j_2(X)) = j_1(X)$. Notons $\overline{c}$ l'adjoint à droite de l'image directe du lemme \ref{duality-lemma-twisted-inverse-image} pour le morphisme $g$. Les applications
$$
L\overline{f}_1^*K \otimes_{\mathcal{O}_{\overline{X}}}^\mathbf{L}
\overline{a}_1(\mathcal{O}_Y)
\longrightarrow
\overline{a}_1(K)
\quad\text{et}\quad
L\overline{f}_2^*K \otimes_{\mathcal{O}_{\overline{X}}}^\mathbf{L}
\overline{a}_2(\mathcal{O}_Y)
\longrightarrow
\overline{a}_2(K)
$$
s'insèrent dans le diagramme commutatif
$$
\xymatrix{
Lg^*(L\overline{f}_2^*K \otimes^\mathbf{L}
\overline{a}_2(\mathcal{O}_Y))
\otimes^\mathbf{L} \overline{c}(\mathcal{O}_{\overline{X}_2})
\ar@{=}[d] \ar[r]_-\sigma &
\overline{c}(L\overline{f}_2^*K \otimes^\mathbf{L}
\overline{a}_2(\mathcal{O}_Y)) \ar[r] &
\overline{c}(\overline{a}_2(K)) \ar@{=}[d] \\
L\overline{f}_1^*K \otimes^\mathbf{L} Lg^*\overline{a}_2(\mathcal{O}_Y)
\otimes^\mathbf{L} \overline{c}(\mathcal{O}_{\overline{X}_2})
\ar[r]^-{1 \otimes \tau} &
L\overline{f}_1^*K \otimes^\mathbf{L} \overline{a}_1(\mathcal{O}_Y) \ar[r] &
\overline{a}_1(K)
}
$$
d'après le lemme \ref{duality-lemma-transitivity-compare-with-pullback}. D'après le lemme \ref{duality-lemma-compare-on-open}, les applications $\sigma$ et $\tau$ se restreignent en des isomorphismes au-dessus de $X$. On peut en fait préciser davantage. Rappelons que la démonstration du lemme \ref{duality-lemma-shriek-well-defined} utilisait l'application (\ref{duality-equation-sheafy}) $\gamma : j_1^* \circ \overline{c} \to j_2^*$ pour construire l'isomorphisme $\alpha_g : j_1^* \circ \overline{a}_1 \to j_2^* \circ \overline{a}_2$. L'image inverse de l'application $\sigma$ par $j_1$ est l'identité de $j_2^*\left(L\overline{f}_2^*K \otimes^\mathbf{L}
\overline{a}_2(\mathcal{O}_Y)\right)$ si l'on identifie $j_1^*\overline{c}(\mathcal{O}_{\overline{X}_2})$ à $\mathcal{O}_X$ par $j_1^* \circ \overline{c} \to j_2^*$, voir le lemme \ref{duality-lemma-restriction-compare-with-pullback}. De même, l'image inverse de l'application $\tau : Lg^*\overline{a}_2(\mathcal{O}_Y)
\otimes^\mathbf{L} \overline{c}(\mathcal{O}_{\overline{X}_2}) \to
\overline{a}_1(\mathcal{O}_Y) = \overline{c}(\overline{a}_2(\mathcal{O}_Y))$ est l'identité de $j_2^*\overline{a}_2(\mathcal{O}_Y)$. En prenant l'image inverse par $j_1$ et en appliquant $\gamma$ partout où cela est possible, on obtient donc un diagramme commutatif
$$
\xymatrix{
j_2^*\left(L\overline{f}_2^*K \otimes^\mathbf{L}
\overline{a}_2(\mathcal{O}_Y)\right) \ar[r] \ar[d] &
j_2^*\overline{a}_2(K) \\
j_1^*L\overline{f}_1^*K \otimes^\mathbf{L} j_2^*\overline{a}_2(\mathcal{O}_Y) &
j_1^*(L\overline{f}_1^*K \otimes^\mathbf{L} \overline{a}_1(\mathcal{O}_Y))
\ar[r] \ar[l]_{1 \otimes \alpha_g} &
j_1^* \overline{a}_1(K) \ar[lu]_{\alpha_g}
}
$$
Sa commutativité signifie exactement que l'application $\mu_{f, K}$ construite à l'aide de la compactification $\overline{X}_1$ est égale à l'application $\mu_{f, K}$ construite à l'aide de la compactification $\overline{X}_2$, via l'identification $\alpha_g$ utilisée dans la démonstration du lemme \ref{duality-lemma-shriek-well-defined}. Des arguments catégoriques identiques à ceux de la démonstration du lemme \ref{duality-lemma-shriek-well-defined} montrent alors que $\mu_{f, K}$ est bien définie (un petit détail est omis).

\medskip\noindent
Cela étant, la commutativité du diagramme de notre lemme résulte de la construction de l'isomorphisme $(g \circ f)^! \to f^! \circ g^!$ (première partie de la démonstration du lemme \ref{duality-lemma-upper-shriek-composition}, utilisant $\overline{X} \to \overline{Y} \to Z$) et du résultat du lemme \ref{duality-lemma-transitivity-compare-with-pullback} pour $\overline{X} \to \overline{Y} \to Z$.
\end{proof}



\section{Propriétés des foncteurs image inverse extraordinaire}
\label{duality-section-upper-shriek-properties}

\noindent
Voici quelques propriétés des foncteurs image inverse extraordinaire.

\begin{lemma}
\label{duality-lemma-shriek-open-immersion}
Dans la situation \ref{duality-situation-shriek}, soit $Y$ un objet de $\textit{FTS}_S$ et soit $j : X \to Y$ une immersion ouverte. Il existe alors un isomorphisme canonique de foncteurs $j^! = j^*$.
\end{lemma}

\noindent
Pour un morphisme \'etale $f : X \to Y$ de $\textit{FTS}_S$, on a aussi $f^* \cong f^!$, voir le lemme \ref{duality-lemma-shriek-etale}.

\begin{proof}
Dans ce cas, on peut choisir $\overline{X} = Y$ comme compactification. L'adjoint à droite du lemme \ref{duality-lemma-twisted-inverse-image} pour $\text{id} : Y \to Y$ est alors le foncteur identité, d'où $j^! = j^*$ par définition.
\end{proof}

\begin{lemma}
\label{duality-lemma-restrict-before-or-after}
Dans la situation \ref{duality-situation-shriek}, soit
$$
\xymatrix{
U \ar[r]_j \ar[d]_g & X \ar[d]^f \\
V \ar[r]^{j'} & Y
}
$$
un diagramme commutatif de $\textit{FTS}_S$ où $j$ et $j'$ sont des immersions ouvertes. Alors $j^* \circ f^! = g^! \circ (j')^*$ comme foncteurs $D^+_\QCoh(\mathcal{O}_Y) \to D^+(\mathcal{O}_U)$.
\end{lemma}

\begin{proof}
Soit $h = f \circ j = j' \circ g$. D'après le lemme \ref{duality-lemma-upper-shriek-composition}, on a $h^! = j^! \circ f^! = g^! \circ (j')^!$. D'après le lemme \ref{duality-lemma-shriek-open-immersion}, on a $j^! = j^*$ et $(j')^! = (j')^*$.
\end{proof}

\begin{lemma}
\label{duality-lemma-shriek-affine-line}
Dans la situation \ref{duality-situation-shriek}, soit $Y$ un objet de $\textit{FTS}_S$ et soit $f : X = \mathbf{A}^1_Y \to Y$ la projection. Il existe alors un isomorphisme (non canonique) de foncteurs $f^!(-) \cong Lf^*(-) [1]$.
\end{lemma}

\begin{proof}
Puisque $X = \mathbf{A}^1_Y \subset \mathbf{P}^1_Y$ et puisque $\mathcal{O}_{\mathbf{P}^1_Y}(-2)|_X \cong \mathcal{O}_X$, cela résulte des lemmes \ref{duality-lemma-upper-shriek-P1} et \ref{duality-lemma-compare-with-pullback-flat-proper-noetherian}.
\end{proof}

\begin{lemma}
\label{duality-lemma-shriek-closed-immersion}
Dans la situation \ref{duality-situation-shriek}, soit $Y$ un objet de $\textit{FTS}_S$ et soit $i : X \to Y$ une immersion fermée. Il existe alors un isomorphisme canonique de foncteurs $i^!(-) = R\SheafHom(\mathcal{O}_X, -)$.
\end{lemma}

\begin{proof}
C'est une reformulation du lemme \ref{duality-lemma-twisted-inverse-image-closed}.
\end{proof}

\begin{remark}[Description locale de l'image inverse extraordinaire]
\label{duality-remark-local-calculation-shriek}
Dans la situation \ref{duality-situation-shriek}, soit $f : X \to Y$ un morphisme de $\textit{FTS}_S$. Les lemmes précédents permettent de calculer $f^!$ localement comme suit. Supposons donnés des ouverts affines
$$
\xymatrix{
U \ar[r]_j \ar[d]_g & X \ar[d]^f \\
V \ar[r]^i & Y
}
$$
Puisque $j^! \circ f^! = g^! \circ i^!$ (lemme \ref{duality-lemma-upper-shriek-composition}) et que $j^!$ et $i^!$ sont donnés par restriction (lemme \ref{duality-lemma-shriek-open-immersion}), on voit que
$$
(f^!E)|_U = g^!(E|_V)
$$
pour tout $E \in D^+_\QCoh(\mathcal{O}_X)$. Écrivons $U = \Spec(A)$ et $V = \Spec(R)$, et soit $\varphi : R \to A$ l'homomorphisme d'anneaux de type fini correspondant à $g$. Choisissons une présentation $A = P/I$ où $P = R[x_1, \ldots, x_n]$ est une algèbre de polynômes en $n$ variables sur $R$. Choisissons un objet $K \in D^+(R)$ correspondant à $E|_V$ (Catégories dérivées des schémas, lemme \ref{perfect-lemma-affine-compare-bounded}). Nous affirmons alors que $f^!E|_U$ correspond à
$$
\varphi^!(K) = R\Hom(A, K \otimes_R^\mathbf{L} P)[n]
$$
où $R\Hom(A, -) : D(P) \to D(A)$ est le foncteur de Complexes dualisants, section \ref{dualizing-section-trivial}, et où $\varphi^! : D(R) \to D(A)$ est celui de Complexes dualisants, section \ref{dualizing-section-relative-dualizing-complex-algebraic}. En effet, le choix de la présentation donne une factorisation
$$
U \rightarrow \mathbf{A}^n_V \to \mathbf{A}^{n - 1}_V \to \ldots \to
\mathbf{A}^1_V \to V
$$
En appliquant le lemme \ref{duality-lemma-shriek-affine-line} exactement $n$ fois, on voit que $(\mathbf{A}^n_V \to V)^!(E|_V)$ correspond à $K \otimes_R^\mathbf{L} P[n]$. D'après les lemmes \ref{duality-lemma-sheaf-with-exact-support-quasi-coherent} et \ref{duality-lemma-shriek-closed-immersion}, la dernière étape correspond à l'application de $R\Hom(A, -)$.
\end{remark}

\begin{lemma}
\label{duality-lemma-shriek-coherent}
Dans la situation \ref{duality-situation-shriek}, soit $f : X \to Y$ un morphisme de $\textit{FTS}_S$. Alors $f^!$ envoie $D_{\textit{Coh}}^+(\mathcal{O}_Y)$ dans $D_{\textit{Coh}}^+(\mathcal{O}_X)$.
\end{lemma}

\begin{proof}
La question est locale sur $X$ ; on peut donc supposer que $X$ et $Y$ sont affines. On peut alors factoriser $f : X \to Y$ sous la forme
$$
X \xrightarrow{i} \mathbf{A}^n_Y \to \mathbf{A}^{n - 1}_Y \to \ldots \to
\mathbf{A}^1_Y \to Y
$$
où $i$ est une immersion fermée. Le lemme résulte des lemmes \ref{duality-lemma-shriek-affine-line} et \ref{duality-lemma-sheaf-with-exact-support-coherent}, de Complexes dualisants, lemme \ref{dualizing-lemma-dualizing-polynomial-ring}, et d'une récurrence.
\end{proof}

\begin{lemma}
\label{duality-lemma-shriek-dualizing}
Dans la situation \ref{duality-situation-shriek}, soit $f : X \to Y$ un morphisme de $\textit{FTS}_S$. Si $K$ est un complexe dualisant sur $Y$, alors $f^!K$ est un complexe dualisant sur $X$.
\end{lemma}

\begin{proof}
La question est locale sur $X$ ; on peut donc supposer que $X$ et $Y$ sont affines. On peut alors factoriser $f : X \to Y$ sous la forme
$$
X \xrightarrow{i} \mathbf{A}^n_Y \to \mathbf{A}^{n - 1}_Y \to \ldots \to
\mathbf{A}^1_Y \to Y
$$
où $i$ est une immersion fermée. Le lemme \ref{duality-lemma-shriek-affine-line} et Complexes dualisants, lemme \ref{dualizing-lemma-dualizing-polynomial-ring}, montrent par récurrence que $p^!K$ est un complexe dualisant sur $\mathbf{A}^n_Y$, où $p : \mathbf{A}^n_Y \to Y$ est la projection. De même, Complexes dualisants, lemme \ref{dualizing-lemma-dualizing-quotient}, et les lemmes \ref{duality-lemma-sheaf-with-exact-support-quasi-coherent} et \ref{duality-lemma-shriek-closed-immersion} montrent que $i^!$ transforme les complexes dualisants en complexes dualisants.
\end{proof}

\begin{lemma}
\label{duality-lemma-shriek-via-duality}
Dans la situation \ref{duality-situation-shriek}, soit $f : X \to Y$ un morphisme de $\textit{FTS}_S$. Soit $K$ un complexe dualisant sur $Y$. Posons $D_Y(M) = R\SheafHom_{\mathcal{O}_Y}(M, K)$ pour $M \in D_{\textit{Coh}}(\mathcal{O}_Y)$ et $D_X(E) = R\SheafHom_{\mathcal{O}_X}(E, f^!K)$ pour $E \in D_{\textit{Coh}}(\mathcal{O}_X)$. Il existe alors un isomorphisme canonique
$$
f^!M \longrightarrow D_X(Lf^*D_Y(M))
$$
pour $M \in D_{\textit{Coh}}^+(\mathcal{O}_Y)$.
\end{lemma}

\begin{proof}
Choisissons une compactification $j : X \subset \overline{X}$ de $X$ sur $Y$ (Compléments sur la platitude, théorème \ref{flat-theorem-nagata} et lemme \ref{flat-lemma-compactifyable}). Soit $a$ l'adjoint à droite du lemme \ref{duality-lemma-twisted-inverse-image} pour $\overline{X} \to Y$. Posons
$D_{\overline{X}}(E) = R\SheafHom_{\mathcal{O}_{\overline{X}}}(E, a(K))$
pour $E \in D_{\textit{Coh}}(\mathcal{O}_{\overline{X}})$. Comme la formation de $R\SheafHom$ commute à la restriction aux ouverts et que $f^! = j^* \circ a$, il suffit de montrer qu'il existe un isomorphisme canonique
$$
a(M) \longrightarrow D_{\overline{X}}(L\overline{f}^*D_Y(M))
$$
pour $M \in D_{\textit{Coh}}(\mathcal{O}_Y)$. Pour $F \in D_\QCoh(\mathcal{O}_X)$, on a
\begin{align*}
\Hom_{\overline{X}}(
F, D_{\overline{X}}(L\overline{f}^*D_Y(M)))
& =
\Hom_{\overline{X}}(
F \otimes_{\mathcal{O}_X}^\mathbf{L} L\overline{f}^*D_Y(M), a(K)) \\
& =
\Hom_Y(
R\overline{f}_*(F \otimes_{\mathcal{O}_X}^\mathbf{L} L\overline{f}^*D_Y(M)),
K) \\
& =
\Hom_Y(
R\overline{f}_*(F) \otimes_{\mathcal{O}_Y}^\mathbf{L} D_Y(M),
K) \\
& =
\Hom_Y(
R\overline{f}_*(F), D_Y(D_Y(M))) \\
& =
\Hom_Y(R\overline{f}_*(F), M) \\
& = \Hom_{\overline{X}}(F, a(M))
\end{align*}
La première égalité résulte de Cohomologie, lemme \ref{cohomology-lemma-internal-hom}. La deuxième résulte de la définition de $a$. La troisième résulte de Catégories dérivées des schémas, lemme \ref{perfect-lemma-cohomology-base-change}. La quatrième résulte de Cohomologie, lemme \ref{cohomology-lemma-internal-hom}, et de la définition de $D_Y$. La cinquième résulte du lemme \ref{duality-lemma-dualizing-schemes}. La dernière résulte de la définition de $a$. Le lemme de Yoneda donne donc $a(M) = D_{\overline{X}}(L\overline{f}^*D_Y(M))$.
\end{proof}

\begin{lemma}
\label{duality-lemma-perfect-comparison-shriek}
Dans la situation \ref{duality-situation-shriek}, soit $f : X \to Y$ un morphisme de $\textit{FTS}_S$. Supposons $f$ parfait (par exemple plat). Alors
\begin{enumerate}
\item[(a)] $f^!\mathcal{O}_Y$ appartient à $D_{\textit{Coh}}^b(\mathcal{O}_X)$ ;
\item[(b)] $f^!\mathcal{O}_Y$ est de Tor-dimension finie dans $D(f^{-1}\mathcal{O}_Y)$ ;
\item[(c)] $\mathcal{O}_X \to
R\SheafHom_{\mathcal{O}_X}(f^!\mathcal{O}_Y, f^!\mathcal{O}_Y)$ est un isomorphisme ;
\item[(d)] $f^!$ envoie $D_{\textit{Coh}}^b(\mathcal{O}_Y)$ dans $D_{\textit{Coh}}^b(\mathcal{O}_X)$ ;
\item[(e)] l'application
$\mu_{f,  K} :
Lf^*K \otimes_{\mathcal{O}_X}^\mathbf{L} f^!\mathcal{O}_Y
\to
f^!K$
du lemme \ref{duality-lemma-map-pullback-to-shriek-well-defined} est un isomorphisme pour tout $K \in D_\QCoh^+(\mathcal{O}_Y)$.
\end{enumerate}
\end{lemma}

\begin{proof}
(Un morphisme plat de présentation finie est parfait, voir Compléments sur les morphismes, lemme \ref{more-morphisms-lemma-flat-finite-presentation-perfect}.) Les assertions (a), (b) et (c) sont locales sur $X$. On peut donc supposer $X$ et $Y$ affines. La remarque \ref{duality-remark-local-calculation-shriek} transforme alors (a), (b) et (c) en (1)(a), (1)(b) et (1)(c) de Complexes dualisants, lemme \ref{dualizing-lemma-relative-dualizing-algebraic}. (Utiliser Catégories dérivées des schémas, lemmes \ref{perfect-lemma-tor-dimension-rel-affine}, \ref{perfect-lemma-pseudo-coherent}, \ref{perfect-lemma-identify-pseudo-coherent-noetherian} \ref{perfect-lemma-quasi-coherence-internal-hom}, pour faire correspondre les énoncés.)

\medskip\noindent
Pour démontrer (d) et (e), commençons par quelques remarques préliminaires.
\begin{enumerate}
\item Nous savons déjà que $f^!$ envoie $D_{\textit{Coh}}^+(\mathcal{O}_Y)$ dans $D_{\textit{Coh}}^+(\mathcal{O}_X)$, voir le lemme \ref{duality-lemma-shriek-coherent}.
\item Si $f$ est une immersion ouverte, alors $f^! = f^*$, et (d) et (e) sont vraies, car on peut prendre $\overline{X} = Y$ dans la construction de $f^!$ et $\mu_f$, voir le lemme \ref{duality-lemma-shriek-open-immersion}.
\item Si $f$ est un morphisme parfait et propre, (e) est vraie d'après le lemme \ref{duality-lemma-compare-with-pullback-flat-proper-noetherian}.
\item S'il existe un recouvrement ouvert $X = \bigcup U_i$ tel que (d) soit vraie pour $U_i \to Y$, alors (d) est vraie pour $X \to Y$. Il en va de même pour (e). Cela tient au fait que la construction de $f^!$ et $\mu_f$ commute au passage aux sous-schémas ouverts.
\item Si $g : Y \to Z$ est un second morphisme parfait dans $\textit{FTS}_S$ et si (e) vaut pour $f$ et $g$, alors
$f^!g^!\mathcal{O}_Z =
Lf^*g^!\mathcal{O}_Z \otimes_{\mathcal{O}_X}^\mathbf{L} f^!\mathcal{O}_Y$
et (e) vaut pour $g \circ f$ par le diagramme commutatif du lemme \ref{duality-lemma-map-pullback-to-shriek-well-defined}.
\item Si (d) et (e) valent pour $f$ et $g$, elles valent pour $g \circ f$. En effet, $f^!g^!\mathcal{O}_Z$ est alors borné supérieurement (par le point précédent), $L(g \circ f)^*$ est de dimension cohomologique finie, et (d) résulte de (e), déjà démontrée ci-dessus.
\end{enumerate}
Ces remarques montrent qu'il suffit de démontrer (d) et (e) lorsque $X$ est affine. Choisissons une immersion $X \to \mathbf{A}^n_Y$ (Morphismes, lemme \ref{morphisms-lemma-quasi-affine-finite-type-over-S}), que l'on factorise en $X \to U \to \mathbf{A}^n_Y \to Y$, où $X \to U$ est une immersion fermée et $U \subset \mathbf{A}^n_Y$ est ouverte. Remarquons que $X \to U$ est une immersion fermée parfaite d'après Compléments sur les morphismes, lemme \ref{more-morphisms-lemma-perfect-permanence}. Il suffit donc de démontrer le lemme pour une immersion fermée parfaite et pour la projection $\mathbf{A}^n_Y \to Y$.

\medskip\noindent
Soit $f : X \to Y$ une immersion fermée parfaite. Nous savons déjà que (d) est vraie. Soit $K \in D^b_{\textit{Coh}}(\mathcal{O}_Y)$. Alors $f^!K = R\SheafHom(\mathcal{O}_X, K)$ (lemme \ref{duality-lemma-shriek-closed-immersion}) et $f_*f^!K = R\SheafHom_{\mathcal{O}_Y}(f_*\mathcal{O}_X, K)$. Puisque $f$ est parfait, le complexe $f_*\mathcal{O}_X$ est parfait ; ainsi, $R\SheafHom_{\mathcal{O}_Y}(f_*\mathcal{O}_X, K)$ est borné supérieurement. Cela démontre (d). Certains détails sont omis.

\medskip\noindent
Soit $f : \mathbf{A}^n_Y \to Y$ la projection. Alors (d) résulte d'applications successives du lemme \ref{duality-lemma-shriek-affine-line}. Enfin, (e) est vraie, car elle vaut pour $\mathbf{P}^n_Y \to Y$ (plat et propre) et parce que $\mathbf{A}^n_Y \subset \mathbf{P}^n_Y$ est un ouvert.
\end{proof}

\begin{lemma}
\label{duality-lemma-flat-shriek-relatively-perfect}
Dans la situation \ref{duality-situation-shriek}, soit $f : X \to Y$ un morphisme de $\textit{FTS}_S$. Si $f$ est plat, alors $f^!\mathcal{O}_Y$ est un objet $Y$-parfait de $D(\mathcal{O}_X)$, et
\par\noindent\hfil
$\mathcal{O}_X \to
R\SheafHom_{\mathcal{O}_X}(f^!\mathcal{O}_Y, f^!\mathcal{O}_Y)$
\hfil\par\noindent
est un isomorphisme.
\end{lemma}

\begin{proof}
Les deux assertions sont locales sur $X$. On peut donc supposer $X$ et $Y$ affines. La remarque \ref{duality-remark-local-calculation-shriek} transforme alors le lemme en un lemme algébrique, à savoir Complexes dualisants, lemme \ref{dualizing-lemma-relative-dualizing-algebraic}. (Utiliser Catégories dérivées des schémas, lemme \ref{perfect-lemma-affine-locally-rel-perfect}, pour passer d'un langage à l'autre.)
\end{proof}

\begin{lemma}
\label{duality-lemma-lci-shriek}
Dans la situation \ref{duality-situation-shriek}, soit $f : X \to Y$ un morphisme de $\textit{FTS}_S$. Supposons que $f : X \to Y$ soit un morphisme localement d'intersection complète. Alors
\begin{enumerate}
\item $f^!\mathcal{O}_Y$ est un objet inversible de $D(\mathcal{O}_X)$ ;
\item $f^!$ envoie les complexes parfaits sur des complexes parfaits.
\end{enumerate}
\end{lemma}

\begin{proof}
Rappelons qu'un morphisme localement d'intersection complète est parfait, voir Compléments sur les morphismes, lemme \ref{more-morphisms-lemma-lci-properties}. D'après le lemme \ref{duality-lemma-perfect-comparison-shriek}, il suffit de montrer que $f^!\mathcal{O}_Y$ est un objet inversible de $D(\mathcal{O}_X)$. Cette question est locale sur $X$ et $Y$. On peut donc supposer que $X \to Y$ se factorise en $X \to \mathbf{A}^n_Y \to Y$, où la première flèche est une immersion Koszul-régulière. Voir Compléments sur les morphismes, section \ref{more-morphisms-section-lci}. Le résultat vaut pour $\mathbf{A}^n_Y \to Y$ d'après le lemme \ref{duality-lemma-shriek-affine-line}. Il suffit donc de démontrer le lemme lorsque $f$ est une immersion Koszul-régulière. En travaillant encore localement, on se ramène au cas $X = \Spec(A)$ et $Y = \Spec(B)$, où $A = B/(f_1, \ldots, f_r)$ pour une suite régulière $f_1, \ldots, f_r \in B$ (utiliser le fait que, sur les anneaux locaux noethériens, les notions de suite Koszul-régulière et de suite régulière coïncident, voir Compléments d'algèbre, lemme \ref{more-algebra-lemma-noetherian-finite-all-equivalent}). Ainsi, $X \to Y$ est une composée
$$
X = X_r \to X_{r - 1} \to \ldots \to X_1 \to X_0 = Y
$$
où chaque flèche est l'inclusion d'un diviseur de Cartier effectif. On se ramène ainsi à l'inclusion d'un diviseur de Cartier effectif $i : D \to X$. Dans ce cas, $i^!\mathcal{O}_X = \mathcal{N}[1]$ d'après le lemme \ref{duality-lemma-compute-for-effective-Cartier}, ce qui achève la démonstration.
\end{proof}







\section{Changement de base pour l'image inverse extraordinaire}
\label{duality-section-base-change-shriek}

\noindent
Dans la situation \ref{duality-situation-shriek}, soit
$$
\xymatrix{
X' \ar[r]_{g'} \ar[d]_{f'} & X \ar[d]^f \\
Y' \ar[r]^g & Y
}
$$
un diagramme cartésien de $\textit{FTS}_S$ tel que $X$ et $Y'$ soient Tor-indépendants sur $Y$. Notre cadre actuel ne suffit pas à construire une application de changement de base $L(g')^* \circ f^! \to (f')^! \circ Lg^*$ dans cette généralité. En effet, on ne peut pas en général choisir une compactification $j : X \to \overline{X}$ sur $Y$ telle que $\overline{X}$ et $Y'$ soient Tor-indépendants sur $Y$ ; la construction de l'application de changement de base donnée à la section \ref{duality-section-base-change-map} ne s'applique donc pas\footnote{Le lecteur familier de la géométrie algébrique dérivée reconnaîtra qu'il ne s'agit pas d'une difficulté « réelle ». En prenant $\overline{X}'$ égal au produit fibré dérivé de $\overline{X}$ et $Y'$ sur $Y$, on peut raisonner exactement comme dans la démonstration du lemme \ref{duality-lemma-base-change-shriek-flat} pour définir cette application. Après tout, la Tor-indépendance de $X$ et $Y'$ garantit que $X'$ sera un sous-schéma ouvert du schéma dérivé $\overline{X}'$.}.

\medskip\noindent
On trouvera une solution partielle à la section \ref{duality-section-relative-dualizing-complexes}. En effet, si le morphisme $f$ est plat, on dispose d'une bonne notion de complexe dualisant relatif ; les lemmes \ref{duality-lemma-compactifyable-relative-dualizing} \ref{duality-lemma-base-change-relative-dualizing} et \ref{duality-lemma-perfect-comparison-shriek} permettent alors de construire un isomorphisme canonique de changement de base. Si nous devons l'utiliser, nous ajouterons des énoncés précis et leurs démonstrations plus loin dans ce chapitre.

\begin{lemma}
\label{duality-lemma-base-change-shriek-flat}
Dans la situation \ref{duality-situation-shriek}, soit
$$
\xymatrix{
X' \ar[r]_{g'} \ar[d]_{f'} & X \ar[d]^f \\
Y' \ar[r]^g & Y
}
$$
un diagramme cartésien de $\textit{FTS}_S$, avec $g$ plat. Il existe alors un isomorphisme $L(g')^* \circ f^! \to (f')^! \circ Lg^*$ sur $D_\QCoh^+(\mathcal{O}_Y)$.
\end{lemma}

\begin{proof}
Puisque $g$ est plat, pour tout choix de compactification $j : X \to \overline{X}$ de $X$ sur $Y$, le schéma $\overline{X}$ est Tor-indépendant de $Y'$. Notons $j' : X' \to \overline{X}'$ le changement de base de $j$ et $\overline{g}' : \overline{X}' \to \overline{X}$ la projection. On définit l'application de changement de base comme la composée
$$
L(g')^* \circ f^! = L(g')^* \circ j^* \circ a =
(j')^* \circ L(\overline{g}')^* \circ a \longrightarrow
(j')^* \circ a' \circ Lg^* = (f')^! \circ Lg^*
$$
où la flèche centrale est l'application de changement de base (\ref{duality-equation-base-change-map}), et où $a$ et $a'$ sont les adjoints à droite de l'image directe du lemme \ref{duality-lemma-twisted-inverse-image} pour $\overline{X} \to Y$ et $\overline{X}' \to Y'$. Cette construction ne dépend pas du choix de la compactification (nous formulerons et démontrerons un lemme précis si nous avons besoin de ce résultat).

\medskip\noindent
Pour achever la démonstration, il suffit de montrer que l'application de changement de base $L(g')^* \circ a \to a' \circ Lg^*$ est un isomorphisme sur $D_\QCoh^+(\mathcal{O}_Y)$. D'après le lemme \ref{duality-lemma-proper-noetherian}, la formation de $a$ et $a'$ commute à la restriction aux ouverts affines de $Y$ et $Y'$. D'après la remarque \ref{duality-remark-check-over-affines}, on peut donc supposer $Y$ et $Y'$ affines. Le résultat découle alors du lemme \ref{duality-lemma-more-base-change}.
\end{proof}

\begin{lemma}
\label{duality-lemma-shriek-etale}
Dans la situation \ref{duality-situation-shriek}, soit $f : X \to Y$ un morphisme \'etale de $\textit{FTS}_S$. Alors $f^! \cong f^*$ comme foncteurs sur $D^+_\QCoh(\mathcal{O}_Y)$.
\end{lemma}

\begin{proof}
Nous allons utiliser le fait qu'un morphisme \'etale est plat, syntomique et localement d'intersection complète (Morphismes, lemmes \ref{morphisms-lemma-etale-syntomic} et \ref{morphisms-lemma-etale-flat}, et Compléments sur les morphismes, lemme \ref{more-morphisms-lemma-flat-lci}). D'après le lemme \ref{duality-lemma-perfect-comparison-shriek}, il suffit de montrer que $f^!\mathcal{O}_Y = \mathcal{O}_X$. Le lemme \ref{duality-lemma-lci-shriek} affirme que $f^!\mathcal{O}_Y$ est un module inversible. Considérons le diagramme commutatif
$$
\xymatrix{
X \times_Y X \ar[r]_{p_2} \ar[d]_{p_1} & X \ar[d]^f \\
X \ar[r]^f & Y
}
$$
et la diagonale $\Delta : X \to X \times_Y X$. Puisque $\Delta$ est une immersion ouverte (Morphismes, lemmes \ref{morphisms-lemma-diagonal-unramified-morphism} et \ref{morphisms-lemma-etale-smooth-unramified}), le lemme \ref{duality-lemma-shriek-open-immersion} donne $\Delta^! = \Delta^*$. Le lemme \ref{duality-lemma-upper-shriek-composition} donne $\Delta^! \circ p_1^! \circ f^! = f^!$. En appliquant le lemme \ref{duality-lemma-base-change-shriek-flat} au diagramme, on obtient $p_1^!\mathcal{O}_X = p_2^*f^!\mathcal{O}_Y$. On en conclut que
$$
f^!\mathcal{O}_Y = \Delta^!p_1^!f^!\mathcal{O}_Y =
\Delta^*(p_1^*f^!\mathcal{O}_Y \otimes p_1^!\mathcal{O}_X) =
\Delta^*(p_2^*f^!\mathcal{O}_Y \otimes p_1^*f^!\mathcal{O}_Y) =
(f^!\mathcal{O}_Y)^{\otimes 2}
$$
où l'on utilise encore le lemme \ref{duality-lemma-perfect-comparison-shriek} à la deuxième étape. Ainsi, $f^!\mathcal{O}_Y = \mathcal{O}_X$, comme voulu.
\end{proof}


\noindent
Dans la suite de cette section, nous énonçons quelques résultats faciles à démontrer qui découleraient d'une bonne théorie de l'application de changement de base.

\begin{lemma}[Changement de base de remplacement]
\label{duality-lemma-base-change-locally}
Dans la situation \ref{duality-situation-shriek}, soit
$$
\xymatrix{
X' \ar[r]_{g'} \ar[d]_{f'} & X \ar[d]^f \\
Y' \ar[r]^g & Y
}
$$
un diagramme cartésien de $\textit{FTS}_S$. Soit $E \in D^+_\QCoh(\mathcal{O}_Y)$ un objet tel que $Lg^*E$ appartienne à $D^+(\mathcal{O}_Y)$. Si $f$ est plat, alors $L(g')^*f^!E$ et $(f')^!Lg^*E$ se restreignent en des objets isomorphes de $D(\mathcal{O}_{U'})$ pour tout ouvert affine $U' \subset X'$ dont les images sont contenues dans des ouverts affines de $Y$, $Y'$ et $X$.
\end{lemma}

\begin{proof}
Nos hypothèses permettent de se ramener immédiatement au cas où $X$, $Y$, $Y'$ et $X'$ sont affines. Écrivons $Y = \Spec(R)$, $Y' = \Spec(R')$, $X = \Spec(A)$ et $X' = \Spec(A')$. Alors $A' = A \otimes_R R'$. Soit $E$ correspondant à $K \in D^+(R)$. En notant $\varphi : R \to A$ et $\varphi' : R' \to A'$ les applications données, la remarque \ref{duality-remark-local-calculation-shriek} montre que $L(g')^*f^!E$ et $(f')^!Lg^*E$ correspondent à $\varphi^!(K) \otimes_A^\mathbf{L} A'$ et $(\varphi')^!(K \otimes_R^\mathbf{L} R')$, où $\varphi^!$ et $(\varphi')^!$ sont les foncteurs de Complexes dualisants, section \ref{dualizing-section-relative-dualizing-complex-algebraic}. Le résultat découle de Complexes dualisants, lemme \ref{dualizing-lemma-bc-flat}.
\end{proof}

\begin{lemma}
\label{duality-lemma-relative-dualizing-fibres}
Dans la situation \ref{duality-situation-shriek}, soit $f : X \to Y$ un morphisme de $\textit{FTS}_S$. Supposons $f$ plat. Posons $\omega_{X/Y}^\bullet = f^!\mathcal{O}_Y$ dans $D^b_{\textit{Coh}}(X)$. Soit $y \in Y$ et soit $h : X_y \to X$ la projection. Alors $Lh^*\omega_{X/Y}^\bullet$ est un complexe dualisant sur $X_y$.
\end{lemma}

\begin{proof}
Le complexe $\omega_{X/Y}^\bullet$ appartient à $D^b_{\textit{Coh}}$ d'après le lemme \ref{duality-lemma-perfect-comparison-shriek}. Être dualisant est une propriété locale. D'après le lemme \ref{duality-lemma-base-change-locally}, il suffit donc de montrer que $(X_y \to y)^!\mathcal{O}_y$ est un complexe dualisant sur $X_y$. Cela résulte du lemme \ref{duality-lemma-shriek-dualizing}.
\end{proof}








\section{Une théorie de la dualité}
\label{duality-section-duality}

\noindent
Nous explicitons dans cette section la théorie de la dualité que donnent les résultats très généraux précédents pour les schémas séparés de type fini sur un schéma de base noethérien fixé.

\medskip\noindent
Rappelons qu'un complexe dualisant sur un schéma noethérien $X$ est un objet de $D(\mathcal{O}_X)$ qui donne localement, sur les ouverts affines, un complexe dualisant sur les anneaux correspondants, voir la définition \ref{duality-definition-dualizing-scheme}.

\medskip\noindent
Étant donné un schéma noethérien $S$, notons $\textit{FTS}_S$ la catégorie des schémas séparés de type fini sur $S$. Alors :
\begin{enumerate}
\item les foncteurs $f^!$ font de $D_\QCoh^+$ un pseudo-foncteur sur $\textit{FTS}_S$ ;
\item si $f : X \to Y$ est un morphisme propre de $\textit{FTS}_S$, alors $f^!$ est la restriction de l'adjoint à droite de $Rf_* : D_\QCoh(\mathcal{O}_X) \to D_\QCoh(\mathcal{O}_Y)$ à $D_\QCoh^+(\mathcal{O}_Y)$, et il existe un isomorphisme canonique
$$
Rf_*R\SheafHom_{\mathcal{O}_X}(K, f^!M)
\to
R\SheafHom_{\mathcal{O}_Y}(Rf_*K, M)
$$
pour tous $K \in D_{\textit{Coh}}^-(\mathcal{O}_X)$ et $M \in D_\QCoh^+(\mathcal{O}_Y)$ ;
\item si un objet $X$ de $\textit{FTS}_S$ admet un complexe dualisant $\omega_X^\bullet$, le foncteur
$D_X = R\SheafHom_{\mathcal{O}_X}(-, \omega_X^\bullet)$
définit une involution de $D_{\textit{Coh}}(\mathcal{O}_X)$ échangeant $D_{\textit{Coh}}^+(\mathcal{O}_X)$ et $D_{\textit{Coh}}^-(\mathcal{O}_X)$ et préservant $D_{\textit{Coh}}^b(\mathcal{O}_X)$ ;
\item si $f : X \to Y$ est un morphisme de $\textit{FTS}_S$ et si $\omega_Y^\bullet$ est un complexe dualisant sur $Y$, alors
\begin{enumerate}
\item $\omega_X^\bullet = f^!\omega_Y^\bullet$ est un complexe dualisant sur $X$ ;
\item $f^!M = D_X(Lf^*D_Y(M))$ canoniquement pour $M \in D_{\textit{Coh}}^+(\mathcal{O}_Y)$ ;
\item si de plus $f$ est propre, alors
$$
Rf_*R\SheafHom_{\mathcal{O}_X}(K, \omega_X^\bullet) =
R\SheafHom_{\mathcal{O}_Y}(Rf_*K, \omega_Y^\bullet)
$$
pour $K$ dans $D^-_{\textit{Coh}}(\mathcal{O}_X)$ ;
\end{enumerate}
\item si $f : X \to Y$ est une immersion fermée dans $\textit{FTS}_S$, alors
\par\noindent\hfil
$f^!(-) = R\SheafHom(\mathcal{O}_X, -)$ ;
\hfil\par
\item si $f : Y \to X$ est un morphisme fini dans $\textit{FTS}_S$, alors
\par\noindent\hfil
$f_*f^!(-) = R\SheafHom_{\mathcal{O}_X}(f_*\mathcal{O}_Y, -)$ ;
\hfil\par
\item si $f : X \to Y$ est l'inclusion d'un diviseur de Cartier effectif dans un objet de $\textit{FTS}_S$, alors $f^!(-) = Lf^*(-) \otimes_{\mathcal{O}_X} f^*\mathcal{O}_Y(X)[-1]$ ;
\item si $f : X \to Y$ est une immersion Koszul-régulière de codimension $c$ dans un objet de $\textit{FTS}_S$, alors $f^!(-) \cong Lf^*(-) \otimes_{\mathcal{O}_X} \wedge^c\mathcal{N}[-c]$ ;
\item si $f : X \to Y$ est un morphisme lisse et propre de dimension relative $d$ dans $\textit{FTS}_S$, alors $f^!(-) \cong Lf^*(-)  \otimes_{\mathcal{O}_X} \Omega^d_{X/Y}[d]$.
\end{enumerate}
Cela résulte des lemmes \ref{duality-lemma-dualizing-schemes}, \ref{duality-lemma-iso-on-RSheafHom}, \ref{duality-lemma-twisted-inverse-image-closed}, \ref{duality-lemma-finite-twisted}, \ref{duality-lemma-sheaf-with-exact-support-effective-Cartier}, \ref{duality-lemma-regular-immersion}, \ref{duality-lemma-smooth-proper}, \ref{duality-lemma-upper-shriek-composition}, \ref{duality-lemma-pseudo-functor}, \ref{duality-lemma-shriek-closed-immersion}, \ref{duality-lemma-shriek-dualizing}, \ref{duality-lemma-shriek-via-duality} et \ref{duality-lemma-perfect-comparison-shriek}, ainsi que de l'exemple \ref{duality-example-iso-on-RSheafHom-noetherian}. Nous avons obtenu nos foncteurs par un procédé très abstrait qui repose en dernier ressort sur un théorème d'existence (Catégories dérivées, proposition \ref{derived-proposition-brown}). Cela signifie que nous n'avons, en général, aucune description explicite des foncteurs $f^!$. Cela peut parfois poser problème. En fait, il suffit pourtant souvent de connaître l'existence d'un complexe dualisant et l'isomorphisme de dualité pour déterminer $f^!$.






\section{Recollement des complexes dualisants}
\label{duality-section-glue}

\noindent
Nous allons maintenant utiliser le recollement des complexes dualisants pour obtenir une théorie valable pour tous les schémas de type fini sur $S$, étant donné un couple $(S, \omega_S^\bullet)$ comme dans la situation \ref{duality-situation-dualizing}. Cette construction est analogue à \cite[remarque on page 310]{RD}.

\begin{situation}
\label{duality-situation-dualizing}
Ici, $S$ est un schéma noethérien et $\omega_S^\bullet$ est un complexe dualisant.
\end{situation}

\noindent
Dans la situation \ref{duality-situation-dualizing}, soit $X$ un schéma de type fini sur $S$. Soit $\mathcal{U} : X = \bigcup_{i = 1, \ldots, n} U_i$ un recouvrement ouvert fini de $X$ par des objets de $\textit{FTS}_S$, voir la situation \ref{duality-situation-shriek}. Cela signifie simplement que les morphismes $U_i \to S$ sont séparés (ils sont déjà de type fini). Tout schéma affine de type fini sur $S$ est un objet de $\textit{FTS}_S$ d'après Schémas, lemme \ref{schemes-lemma-compose-after-separated} ; de tels recouvrements ouverts existent donc bien. Pour tout $i, j, k \in \{1, \ldots, n\}$, les morphismes $p_i : U_i \to S$, $p_{ij} : U_i \cap U_j \to S$ et $p_{ijk} : U_i \cap U_j \cap U_k \to S$ sont alors séparés, et chacun de ces schémas est un objet de $\textit{FTS}_S$. Un tel recouvrement fournit
\begin{enumerate}
\item $\omega_i^\bullet = p_i^!\omega_S^\bullet$, complexe dualisant sur $U_i$, voir la section \ref{duality-section-duality} ;
\item pour tout $i, j$, un isomorphisme canonique $\varphi_{ij} :
\omega_i^\bullet|_{U_i \cap U_j} \to \omega_j^\bullet|_{U_i \cap U_j}$ ;
\item
\label{duality-item-cocycle-glueing}
pour tout $i, j, k$, on a
$$
\varphi_{ik}|_{U_i \cap U_j \cap U_k} =
\varphi_{jk}|_{U_i \cap U_j \cap U_k} \circ
\varphi_{ij}|_{U_i \cap U_j \cap U_k}
$$
dans $D(\mathcal{O}_{U_i \cap U_j \cap U_k})$.
\end{enumerate}
Dans (2), on utilise le fait que $(U_i \cap U_j \to U_i)^!$ est donné par restriction (lemme \ref{duality-lemma-shriek-open-immersion}), ainsi que les isomorphismes canoniques
$$
(U_i \cap U_j \to U_i)^! \circ p_i^! = p_{ij}^! =
(U_i \cap U_j \to U_j)^! \circ p_j^!
$$
du lemme \ref{duality-lemma-upper-shriek-composition} ; pour obtenir (3), on utilise le fait que les foncteurs image inverse extraordinaire forment un pseudo-foncteur d'après le lemme \ref{duality-lemma-pseudo-functor}.

\medskip\noindent
Dans la situation qui vient d'être décrite, un {\it complexe dualisant normalisé relativement à $\omega_S^\bullet$ et $\mathcal{U}$} est un couple $(K, \alpha_i)$ où $K \in D(\mathcal{O}_X)$ et où les $\alpha_i : K|_{U_i} \to \omega_i^\bullet$ sont des isomorphismes tels que $\varphi_{ij}$ soit donné par $\alpha_j|_{U_i \cap U_j} \circ \alpha_i^{-1}|_{U_i \cap U_j}$. Puisqu'être un complexe dualisant sur un schéma est une propriété locale, les complexes dualisants normalisés relativement à $\omega_S^\bullet$ et $\mathcal{U}$ sont bien des complexes dualisants.

\begin{lemma}
\label{duality-lemma-good-dualizing-unique}
Dans la situation \ref{duality-situation-dualizing}, soit $X$ un schéma de type fini sur $S$ et soit $\mathcal{U}$ un recouvrement ouvert fini de $X$ par des schémas séparés sur $S$. S'il existe un complexe dualisant normalisé relativement à $\omega_S^\bullet$ et $\mathcal{U}$, il est unique à unique isomorphisme près.
\end{lemma}

\begin{proof}
Si $(K, \alpha_i)$ et $(K', \alpha_i')$ sont deux tels complexes, considérons $L = R\SheafHom_{\mathcal{O}_X}(K, K')$. D'après le lemme \ref{duality-lemma-dualizing-unique-schemes} et sa démonstration, c'est un objet inversible de $D(\mathcal{O}_X)$. À l'aide de $\alpha_i$ et $\alpha'_i$, on obtient un isomorphisme
$$
\alpha_i^t \otimes \alpha'_i :
L|_{U_i} \longrightarrow
R\SheafHom_{\mathcal{O}_X}(\omega_i^\bullet, \omega_i^\bullet) =
\mathcal{O}_{U_i}[0]
$$
Cela implique déjà que $L = H^0(L)[0]$ dans $D(\mathcal{O}_X)$. De plus, $H^0(L)$ est un faisceau inversible muni de trivialisations sur les ouverts $U_i$ de $X$. Enfin, le fait que $\alpha_j|_{U_i \cap U_j} \circ \alpha_i^{-1}|_{U_i \cap U_j}$ et $\alpha'_j|_{U_i \cap U_j} \circ (\alpha'_i)^{-1}|_{U_i \cap U_j}$ donnent tous deux $\varphi_{ij}$ implique que les applications de transition sont $1$, et l'on obtient un isomorphisme $H^0(L) = \mathcal{O}_X$.
\end{proof}

\begin{lemma}
\label{duality-lemma-good-dualizing-independence-covering}
Dans la situation \ref{duality-situation-dualizing}, soit $X$ un schéma de type fini sur $S$ et soient $\mathcal{U}$, $\mathcal{V}$ deux recouvrements ouverts finis de $X$ par des schémas séparés sur $S$. S'il existe un complexe dualisant normalisé relativement à $\omega_S^\bullet$ et $\mathcal{U}$, il en existe un normalisé relativement à $\omega_S^\bullet$ et $\mathcal{V}$, et ces complexes sont canoniquement isomorphes.
\end{lemma}

\begin{proof}
Il suffit de le démontrer lorsque $\mathcal{U}$ est donné par les ouverts $U_1, \ldots, U_n$ et $\mathcal{V}$ par les ouverts $U_1, \ldots, U_{n + m}$. On peut même supposer, et l'on suppose, $m = 1$. Pour passer d'un complexe dualisant $(K, \alpha_i)$ normalisé relativement à $\omega_S^\bullet$ et $\mathcal{V}$ à un complexe dualisant normalisé relativement à $\omega_S^\bullet$ et $\mathcal{U}$, il suffit d'oublier $\alpha_i$ pour $i = n + 1$. Réciproquement, soit $(K, \alpha_i)$ un complexe dualisant normalisé relativement à $\omega_S^\bullet$ et $\mathcal{U}$. Pour achever la démonstration, il faut construire une application $\alpha_{n + 1} : K|_{U_{n + 1}} \to \omega_{n + 1}^\bullet$ satisfaisant les conditions requises. Observons pour cela que $U_{n + 1} = \bigcup U_i \cap U_{n + 1}$ est un recouvrement ouvert. Il est clair que $(K|_{U_{n + 1}}, \alpha_i|_{U_i \cap U_{n + 1}})$ est un complexe dualisant normalisé relativement à $\omega_S^\bullet$ et au recouvrement $U_{n + 1} = \bigcup U_i \cap U_{n + 1}$. D'autre part, par la condition (\ref{duality-item-cocycle-glueing}), le couple
$(\omega_{n + 1}^\bullet|_{U_{n + 1}}, \varphi_{n + 1i})$
est un autre complexe dualisant normalisé relativement à $\omega_S^\bullet$ et au recouvrement $U_{n + 1} = \bigcup U_i \cap U_{n + 1}$. Le lemme \ref{duality-lemma-good-dualizing-unique} donne un unique isomorphisme
$$
\alpha_{n + 1} : K|_{U_{n + 1}} \longrightarrow \omega_{n + 1}^\bullet
$$
compatible avec les isomorphismes locaux donnés. Le lecteur pourra vérifier, à titre d'exercice, que cela signifie qu'il possède la propriété requise.
\end{proof}

\begin{lemma}
\label{duality-lemma-existence-good-dualizing}
Dans la situation \ref{duality-situation-dualizing}, soit $X$ un schéma de type fini sur $S$ et soit $\mathcal{U}$ un recouvrement ouvert fini de $X$ par des schémas séparés sur $S$. Il existe alors un complexe dualisant normalisé relativement à $\omega_S^\bullet$ et $\mathcal{U}$.
\end{lemma}

\begin{proof}
Écrivons $\mathcal{U} : X = \bigcup_{i = 1, \ldots, n} U_i$. Démontrons le lemme par récurrence sur $n$. Le cas initial $n = 1$ est immédiat. Supposons $n > 1$. Posons $X' = U_1 \cup \ldots \cup U_{n - 1}$ et soit $(K', \{\alpha'_i\}_{i = 1, \ldots, n - 1})$ un complexe dualisant normalisé relativement à $\omega_S^\bullet$ et $\mathcal{U}' : X' = \bigcup_{i = 1, \ldots, n - 1} U_i$. Il est clair que $(K'|_{X' \cap U_n}, \alpha'_i|_{U_i \cap U_n})$ est un complexe dualisant normalisé relativement à $\omega_S^\bullet$ et au recouvrement $X' \cap U_n = \bigcup_{i = 1, \ldots, n - 1} U_i \cap U_n$. D'autre part, par la condition (\ref{duality-item-cocycle-glueing}), le couple
$(\omega_n^\bullet|_{X' \cap U_n}, \varphi_{ni})$
est un autre complexe dualisant normalisé relativement à $\omega_S^\bullet$ et au recouvrement $X' \cap U_n = \bigcup_{i = 1, \ldots, n - 1} U_i \cap U_n$. Le lemme \ref{duality-lemma-good-dualizing-unique} donne un unique isomorphisme
$$
\epsilon : K'|_{X' \cap U_n} \longrightarrow \omega_i^\bullet|_{X' \cap U_n}
$$
compatible avec les isomorphismes locaux donnés. D'après Cohomologie, lemme \ref{cohomology-lemma-glue}, on obtient $K \in D(\mathcal{O}_X)$ muni d'isomorphismes $\beta : K|_{X'} \to K'$ et $\gamma : K|_{U_n} \to \omega_n^\bullet$ tels que $\epsilon = \gamma|_{X'\cap U_n} \circ \beta|_{X' \cap U_n}^{-1}$. On définit alors
$$
\alpha_i = \alpha'_i \circ \beta|_{U_i}, i = 1, \ldots, n - 1,
\text{ et }
\alpha_n = \gamma
$$
Il reste à vérifier que $\varphi_{ij}$ est donné par $\alpha_j|_{U_i \cap U_j} \circ \alpha_i^{-1}|_{U_i \cap U_j}$. Pour $i, j \leq n - 1$, cela résulte de la condition correspondante pour $\alpha_i'$. Pour $i = j = n$, c'est également clair. Si $i < j = n$, on obtient
$$
\alpha_n|_{U_i \cap U_n} \circ \alpha_i^{-1}|_{U_i \cap U_n} =
\gamma|_{U_i \cap U_n} \circ \beta^{-1}|_{U_i \cap U_n}
\circ (\alpha'_i)^{-1}|_{U_i \cap U_n} =
\epsilon|_{U_i \cap U_n} \circ (\alpha'_i)^{-1}|_{U_i \cap U_n}
$$
Cela est égal à $\alpha_{in}$ précisément parce que $\epsilon$ est l'unique application compatible avec $\alpha_i'$ et $\alpha_{ni}$.
\end{proof}

\noindent
Soit $(S, \omega_S^\bullet)$ comme dans la situation \ref{duality-situation-dualizing}. Les lemmes précédents montrent que, pour tout schéma $X$ de type fini sur $S$, il existe un couple $(K, \alpha_U)$, donné à unique isomorphisme près, constitué d'un objet $K \in D(\mathcal{O}_X)$ et d'isomorphismes $\alpha_U : K|_U \to \omega_U^\bullet$ pour tout sous-schéma ouvert $U \subset X$ séparé sur $S$. Ici, $\omega_U^\bullet = (U \to S)^!\omega_S^\bullet$ est un complexe dualisant sur $U$, voir la section \ref{duality-section-duality}. De plus, si $\mathcal{U} : X = \bigcup U_i$ est un recouvrement ouvert fini par des ouverts séparés sur $S$, alors $(K, \alpha_{U_i})$ est un complexe dualisant normalisé relativement à $\omega_S^\bullet$ et $\mathcal{U}$. En effet, l'unicité à unique isomorphisme près résulte du lemme \ref{duality-lemma-good-dualizing-unique}, l'existence pour un recouvrement ouvert du lemme \ref{duality-lemma-existence-good-dualizing}, et le fait que $K$ convienne alors pour tous les recouvrements ouverts du lemme \ref{duality-lemma-good-dualizing-independence-covering}.

\begin{definition}
\label{duality-definition-good-dualizing}
Soit $S$ un schéma noethérien et soit $\omega_S^\bullet$ un complexe dualisant sur $S$. Soit $X$ un schéma de type fini sur $S$. Le complexe $K$ construit ci-dessus est appelé {\it complexe dualisant normalisé relativement à $\omega_S^\bullet$} et est noté $\omega_X^\bullet$.
\end{definition}

\noindent
Comme le suggère la terminologie, un complexe dualisant normalisé relativement à $\omega_S^\bullet$ n'est pas seulement un objet de la catégorie dérivée de $X$ : il est muni des isomorphismes locaux décrits ci-dessus. Cela reste compatible avec la définition $\omega_X^\bullet = p^!\omega_S^\bullet$, où $p : X \to S$ est le morphisme structural, si $X$ est séparé sur $S$. Plus généralement, on dispose de la vérification de cohérence suivante.

\begin{lemma}
\label{duality-lemma-good-over-both}
Soit $(S, \omega_S^\bullet)$ comme dans la situation \ref{duality-situation-dualizing}. Soit $f : X \to Y$ un morphisme de schémas de type fini sur $S$. Soient $\omega_X^\bullet$ et $\omega_Y^\bullet$ des complexes dualisants normalisés relativement à $\omega_S^\bullet$. Alors $\omega_X^\bullet$ est un complexe dualisant normalisé relativement à $\omega_Y^\bullet$.
\end{lemma}

\begin{proof}
Il s'agit simplement de vérifier les compatibilités. Choisissons un recouvrement ouvert affine fini $\mathcal{V} : Y = \bigcup V_j$. Pour tout $j$, choisissons un recouvrement ouvert affine fini $f^{-1}(V_j) = U_{ji}$. Posons $\mathcal{U} : X = \bigcup U_{ji}$. Les schémas $V_j$ et $U_{ji}$ sont séparés sur $S$ ; on dispose donc des foncteurs image inverse extraordinaire pour $q_j : V_j \to S$, $p_{ji} : U_{ji} \to S$, $f_{ji} : U_{ji} \to V_j$ et $f_{ji}' : U_{ji} \to Y$. Soit $(L, \beta_j)$ un complexe dualisant normalisé relativement à $\omega_S^\bullet$ et $\mathcal{V}$. Soit $(K, \gamma_{ji})$ un complexe dualisant normalisé relativement à $\omega_S^\bullet$ et $\mathcal{U}$. (Autrement dit, $L = \omega_Y^\bullet$ et $K = \omega_X^\bullet$.) On peut définir
$$
\alpha_{ji} :
K|_{U_{ji}} \xrightarrow{\gamma_{ji}}
p_{ji}^!\omega_S^\bullet = f_{ji}^!q_j^!\omega_S^\bullet
\xrightarrow{f_{ji}^!\beta_j^{-1}} f_{ji}^!(L|_{V_j}) =
(f_{ji}')^!(L)
$$
Pour achever la démonstration, il faut montrer que
$\alpha_{ji}|_{U_{ji} \cap U_{j'i'}}
\circ \alpha_{j'i'}^{-1}|_{U_{ji} \cap U_{j'i'}}$
est l'isomorphisme canonique $(f_{ji}')^!(L)|_{U_{ji} \cap U_{j'i'}} \to
(f_{j'i'}')^!(L)|_{U_{ji} \cap U_{j'i'}}$. C'est formel ; nous omettons les détails.
\end{proof}

\begin{lemma}
\label{duality-lemma-open-immersion-good-dualizing-complex}
Soit $(S, \omega_S^\bullet)$ comme dans la situation \ref{duality-situation-dualizing}. Soit $j : X \to Y$ une immersion ouverte de schémas de type fini sur $S$. Soient $\omega_X^\bullet$ et $\omega_Y^\bullet$ des complexes dualisants normalisés relativement à $\omega_S^\bullet$. Il existe alors un isomorphisme canonique $\omega_X^\bullet = \omega_Y^\bullet|_X$.
\end{lemma}

\begin{proof}
Cela résulte immédiatement de la construction des complexes dualisants normalisés donnée juste avant la définition \ref{duality-definition-good-dualizing}.
\end{proof}

\begin{lemma}
\label{duality-lemma-proper-map-good-dualizing-complex}
Soit $(S, \omega_S^\bullet)$ comme dans la situation \ref{duality-situation-dualizing}. Soit $f : X \to Y$ un morphisme propre de schémas de type fini sur $S$. Soient $\omega_X^\bullet$ et $\omega_Y^\bullet$ des complexes dualisants normalisés relativement à $\omega_S^\bullet$. Soit $a$ l'adjoint à droite du lemme \ref{duality-lemma-twisted-inverse-image} pour $f$. Il existe alors un isomorphisme canonique $a(\omega_Y^\bullet) = \omega_X^\bullet$.
\end{lemma}

\begin{proof}
Soient $p : X \to S$ et $q : Y \to S$ les morphismes structuraux. Si $X$ et $Y$ sont séparés sur $S$, cela résulte de $\omega_X^\bullet = p^!\omega_S^\bullet$, $\omega_Y^\bullet = q^!\omega_S^\bullet$, $f^! = a$ et $f^! \circ q^! = p^!$ (lemme \ref{duality-lemma-upper-shriek-composition}). Dans le cas général, on utilise d'abord le lemme \ref{duality-lemma-good-over-both} pour se ramener au cas $Y = S$. Alors $X$ et $Y$ sont séparés sur $S$, et le résultat vient d'être établi.
\end{proof}

\noindent
Soit $(S, \omega_S^\bullet)$ comme dans la situation \ref{duality-situation-dualizing}. Pour un schéma $X$ de type fini sur $S$, notons $\omega_X^\bullet$ le complexe dualisant sur $X$ normalisé relativement à $\omega_S^\bullet$. Définissons $D_X(-) = R\SheafHom_{\mathcal{O}_X}(-, \omega_X^\bullet)$ comme au lemme \ref{duality-lemma-dualizing-schemes}. Soit $f : X \to Y$ un morphisme de schémas de type fini sur $S$. Définissons
$$
f_{new}^! = D_X \circ Lf^* \circ D_Y :
D_{\textit{Coh}}^+(\mathcal{O}_Y)
\to
D_{\textit{Coh}}^+(\mathcal{O}_X)
$$
Si $f : X \to Y$ et $g : Y \to Z$ sont des morphismes composables entre schémas de type fini sur $S$, définissons
\begin{align*}
(g \circ f)^!_{new} & = D_X \circ L(g \circ f)^* \circ D_Z \\
& = D_X \circ Lf^* \circ Lg^* \circ D_Z \\
& \to D_X \circ Lf^* \circ D_Y \circ D_Y \circ Lg^* \circ D_Z \\
& = f^!_{new} \circ g^!_{new}
\end{align*}
où la flèche est celle du lemme \ref{duality-lemma-dualizing-schemes}. Le lemme suivant rassemble les résultats obtenus.

\begin{lemma}
\label{duality-lemma-duality-bootstrap}
Soit $(S, \omega_S^\bullet)$ comme dans la situation \ref{duality-situation-dualizing}. Avec $f^!_{new}$ et $\omega_X^\bullet$ définis comme ci-dessus pour tous les schémas (et morphismes de schémas) de type fini sur $S$ :
\begin{enumerate}
\item les foncteurs $f^!_{new}$ et les flèches $(g \circ f)^!_{new} \to f^!_{new} \circ g^!_{new}$ font de $D_{\textit{Coh}}^+$ un pseudo-foncteur de la catégorie des schémas de type fini sur $S$ vers la $2$-catégorie des catégories ;
\item $\omega_X^\bullet = (X \to S)^!_{new} \omega_S^\bullet$ ;
\item le foncteur $D_X$ définit une involution de $D_{\textit{Coh}}(\mathcal{O}_X)$ échangeant $D_{\textit{Coh}}^+(\mathcal{O}_X)$ et $D_{\textit{Coh}}^-(\mathcal{O}_X)$ et préservant $D_{\textit{Coh}}^b(\mathcal{O}_X)$ ;
\item $\omega_X^\bullet = f^!_{new}\omega_Y^\bullet$ pour $f : X \to Y$ un morphisme de schémas de type fini sur $S$ ;
\item $f^!_{new}M = D_X(Lf^*D_Y(M))$ pour $M \in D_{\textit{Coh}}^+(\mathcal{O}_Y)$ ;
\item si de plus $f$ est propre, alors $f^!_{new}$ est isomorphe à la restriction de l'adjoint à droite de $Rf_* : D_\QCoh(\mathcal{O}_X) \to D_\QCoh(\mathcal{O}_Y)$ à $D_{\textit{Coh}}^+(\mathcal{O}_Y)$, et il existe un isomorphisme canonique
$$
Rf_*R\SheafHom_{\mathcal{O}_X}(K, f^!_{new}M)
\to
R\SheafHom_{\mathcal{O}_Y}(Rf_*K, M)
$$
pour $K \in D^-_{\textit{Coh}}(\mathcal{O}_X)$ et $M \in D_{\textit{Coh}}^+(\mathcal{O}_Y)$, et
$$
Rf_*R\SheafHom_{\mathcal{O}_X}(K, \omega_X^\bullet) =
R\SheafHom_{\mathcal{O}_Y}(Rf_*K, \omega_Y^\bullet)
$$
pour $K \in D^-_{\textit{Coh}}(\mathcal{O}_X)$, et
\end{enumerate}
\enlargethispage\smallskipamount
Si $X$ est séparé sur $S$, alors $\omega_X^\bullet$ est canoniquement isomorphe à $(X \to S)^!\omega_S^\bullet$ ; si $f$ est un morphisme entre schémas séparés sur $S$, il existe un isomorphisme canonique\footnote{Nous n'avons pas vérifié leur compatibilité avec les isomorphismes $(g \circ f)^! \to f^! \circ g^!$ et $(g \circ f)^!_{new} \to f^!_{new} \circ g^!_{new}$. Nous le ferons ici si nous en avons besoin ultérieurement.} $f_{new}^!K = f^!K$ pour $K$ dans $D_{\textit{Coh}}^+$.
\end{lemma}

\begin{proof}
Soient $f : X \to Y$, $g : Y \to Z$, $h : Z \to T$ des morphismes de schémas de type fini sur $S$. Il faut montrer que
$$
\xymatrix{
(h \circ g \circ f)^!_{new} \ar[r] \ar[d] &
f^!_{new} \circ (h \circ g)^!_{new} \ar[d] \\
(g \circ f)^!_{new} \circ h^!_{new} \ar[r] &
f^!_{new} \circ g^!_{new} \circ h^!_{new}
}
$$
est commutatif. Soient $\eta_Y : \text{id} \to D_Y^2$ et $\eta_Z : \text{id} \to D_Z^2$ les isomorphismes canoniques du lemme \ref{duality-lemma-dualizing-schemes}. En utilisant Catégories, lemme \ref{categories-lemma-properties-2-cat-cats}, un calcul (omis) montre que les deux flèches
$(h \circ g \circ f)^!_{new} \to f^!_{new} \circ g^!_{new} \circ h^!_{new}$
sont données par
$$
1 \star \eta_Y \star 1 \star \eta_Z \star 1 :
D_X \circ Lf^* \circ Lg^* \circ Lh^* \circ D_T
\longrightarrow
D_X \circ Lf^* \circ D_Y^2 \circ Lg^* \circ D_Z^2 \circ Lh^* \circ D_T
$$
Cela démontre (1). L'assertion (2) résulte immédiatement de la définition de $(X \to S)^!_{new}$ et du fait que $D_S(\omega_S^\bullet) = \mathcal{O}_S$. L'assertion (3) est le lemme \ref{duality-lemma-dualizing-schemes}. L'assertion (4) se démontre comme (2). L'assertion (5) est la définition de $f^!_{new}$.

\medskip\noindent
Démonstration de (6). Soit $a$ l'adjoint à droite du lemme \ref{duality-lemma-twisted-inverse-image} pour le morphisme propre $f : X \to Y$ de schémas de type fini sur $S$. La difficulté est que l'on ne sait pas si $X$ ou $Y$ est séparé sur $S$ (et ce sera faux en général) ; on ne peut donc pas appliquer directement le lemme \ref{duality-lemma-shriek-via-duality} à $f$ sur $S$. Pour contourner cette difficulté, on utilise l'identification canonique $\omega_X^\bullet = a(\omega_Y^\bullet)$ du lemme \ref{duality-lemma-proper-map-good-dualizing-complex}. Ainsi, $f^!_{new}$ est la restriction de $a$ à $D_{\textit{Coh}}^+(\mathcal{O}_Y)$ d'après le lemme \ref{duality-lemma-shriek-via-duality} appliqué à $f : X \to Y$ sur le schéma de base $Y$ ! Les égalités affichées résultent de l'exemple \ref{duality-example-iso-on-RSheafHom-noetherian}.

\medskip\noindent
Les dernières assertions résultent de la construction des complexes dualisants normalisés et du lemme \ref{duality-lemma-shriek-via-duality}, déjà utilisé.
\end{proof}

\begin{remark}
\label{duality-remark-independent-omega-S}
Soit $S$ un schéma noethérien admettant un complexe dualisant. Soit $f : X \to Y$ un morphisme de schémas de type fini sur $S$. Alors le foncteur
$$
f_{new}^! : D^+_{Coh}(\mathcal{O}_Y) \to D^+_{Coh}(\mathcal{O}_X)
$$
est indépendant, à isomorphisme canonique près, du choix du complexe dualisant $\omega_S^\bullet$. Esquissons la démonstration. Tout second complexe dualisant est de la forme $\omega_S^\bullet \otimes_{\mathcal{O}_S}^\mathbf{L} \mathcal{L}$, où $\mathcal{L}$ est un objet inversible de $D(\mathcal{O}_S)$, voir le lemme \ref{duality-lemma-dualizing-unique-schemes}. Pour tout morphisme séparé $p : U \to S$ de type fini, on a
$p^!(\omega_S^\bullet \otimes^\mathbf{L}_{\mathcal{O}_S} \mathcal{L}) =
p^!(\omega_S^\bullet) \otimes^\mathbf{L}_{\mathcal{O}_U} Lp^*\mathcal{L}$
d'après le lemme \ref{duality-lemma-compare-with-pullback-perfect}. Si $\omega_X^\bullet$ et $\omega_Y^\bullet$ sont les complexes dualisants normalisés relativement à $\omega_S^\bullet$, alors $\omega_X^\bullet \otimes_{\mathcal{O}_X}^\mathbf{L} La^*\mathcal{L}$ et $\omega_Y^\bullet \otimes_{\mathcal{O}_Y}^\mathbf{L} Lb^*\mathcal{L}$ sont donc les complexes dualisants normalisés relativement à $\omega_S^\bullet \otimes_{\mathcal{O}_S}^\mathbf{L} \mathcal{L}$ (où $a : X \to S$ et $b : Y \to S$ sont les morphismes structuraux). Le résultat s'en déduit, car
\begin{align*}
& R\SheafHom_{\mathcal{O}_X}(Lf^*R\SheafHom_{\mathcal{O}_Y}(K,
\omega_Y^\bullet \otimes_{\mathcal{O}_Y}^\mathbf{L} Lb^*\mathcal{L}),
\omega_X^\bullet \otimes_{\mathcal{O}_X}^\mathbf{L} La^*\mathcal{L}) \\
& = R\SheafHom_{\mathcal{O}_X}(Lf^*R(\SheafHom_{\mathcal{O}_Y}(K,
\omega_Y^\bullet) \otimes_{\mathcal{O}_Y}^\mathbf{L} Lb^*\mathcal{L}),
\omega_X^\bullet \otimes_{\mathcal{O}_X}^\mathbf{L} La^*\mathcal{L}) \\
& = R\SheafHom_{\mathcal{O}_X}(Lf^*R\SheafHom_{\mathcal{O}_Y}(K,
\omega_Y^\bullet) \otimes_{\mathcal{O}_X}^\mathbf{L} La^*\mathcal{L},
\omega_X^\bullet \otimes_{\mathcal{O}_X}^\mathbf{L} La^*\mathcal{L}) \\
& = R\SheafHom_{\mathcal{O}_X}(Lf^*R\SheafHom_{\mathcal{O}_Y}(K,
\omega_Y^\bullet), \omega_X^\bullet)
\end{align*}
pour $K \in D^+_{Coh}(\mathcal{O}_Y)$. La dernière égalité résulte de l'inversibilité de $La^*\mathcal{L}$ dans $D(\mathcal{O}_X)$.
\end{remark}


\begin{example}
\label{duality-example-trace-proper}
Soit $S$ un schéma noethérien et soit $\omega_S^\bullet$ un complexe dualisant. Soit $f : X \to Y$ un morphisme propre de schémas de type fini sur $S$. Soient $\omega_X^\bullet$ et $\omega_Y^\bullet$ des complexes dualisants normalisés relativement à $\omega_S^\bullet$. Dans cette situation, on a $a(\omega_Y^\bullet) = \omega_X^\bullet$ (lemme \ref{duality-lemma-proper-map-good-dualizing-complex}) ; le morphisme trace (section \ref{duality-section-trace}) est donc une flèche canonique
$$
\text{Tr}_f : Rf_*\omega_X^\bullet \longrightarrow \omega_Y^\bullet
$$
qui donne les isomorphismes (lemme \ref{duality-lemma-duality-bootstrap})
$$
\Hom_X(L, \omega_X^\bullet) = \Hom_Y(Rf_*L, \omega_Y^\bullet)
$$
et
$$
Rf_*R\SheafHom_{\mathcal{O}_X}(L, \omega_X^\bullet) =
R\SheafHom_{\mathcal{O}_Y}(Rf_*L, \omega_Y^\bullet)
$$
pour $L$ dans $D_\QCoh(\mathcal{O}_X)$.
\end{example}

\begin{remark}
\label{duality-remark-dualizing-finite}
Soit $S$ un schéma noethérien et soit $\omega_S^\bullet$ un complexe dualisant. Soit $f : X \to Y$ un morphisme fini entre schémas de type fini sur $S$. Soient $\omega_X^\bullet$ et $\omega_Y^\bullet$ des complexes dualisants normalisés relativement à $\omega_S^\bullet$. On a alors
$$
f_*\omega_X^\bullet = R\SheafHom(f_*\mathcal{O}_X, \omega_Y^\bullet)
$$
dans $D_\QCoh^+(f_*\mathcal{O}_X)$ d'après les lemmes \ref{duality-lemma-finite-twisted} et \ref{duality-lemma-proper-map-good-dualizing-complex}, et le morphisme trace de l'exemple \ref{duality-example-trace-proper} est le morphisme
$$
\text{Tr}_f : Rf_*\omega_X^\bullet = f_*\omega_X^\bullet =
R\SheafHom(f_*\mathcal{O}_X, \omega_Y^\bullet) \longrightarrow
\omega_Y^\bullet
$$
souvent appelée « évaluation en $1$ ».
\end{remark}

\begin{remark}
\label{duality-remark-relative-dualizing-complex-shriek}
Soit $f : X \to Y$ un morphisme plat et propre de schémas de type fini sur un couple $(S, \omega_S^\bullet)$ comme dans la situation \ref{duality-situation-dualizing}. Le complexe dualisant relatif (remarque \ref{duality-remark-relative-dualizing-complex}) est $\omega_{X/Y}^\bullet = a(\mathcal{O}_Y)$. Le lemme \ref{duality-lemma-proper-map-good-dualizing-complex} donne le premier isomorphisme canonique de
$$
\omega_X^\bullet = a(\omega_Y^\bullet) =
Lf^*\omega_Y^\bullet \otimes_{\mathcal{O}_X}^\mathbf{L} \omega_{X/Y}^\bullet
$$
dans $D(\mathcal{O}_X)$. Le second isomorphisme canonique résulte de la discussion de la remarque \ref{duality-remark-relative-dualizing-complex}.
\end{remark}




\section{Fonctions de dimension}
\label{duality-section-dimension-functions}

\noindent
Nous avons besoin de quelques précisions sur la variation des fonctions de dimension lorsqu'on passe à un schéma de type fini sur un autre.

\begin{lemma}
\label{duality-lemma-good-dualizing-normalized}
Soit $S$ un schéma noethérien et soit $\omega_S^\bullet$ un complexe dualisant. Soit $X$ un schéma de type fini sur $S$ et soit $\omega_X^\bullet$ le complexe dualisant normalisé relativement à $\omega_S^\bullet$. Si $x \in X$ est un point fermé au-dessus d'un point fermé $s$ de $S$, alors $\omega_{X, x}^\bullet$ est un complexe dualisant normalisé sur $\mathcal{O}_{X, x}$, pourvu que $\omega_{S, s}^\bullet$ soit un complexe dualisant normalisé sur $\mathcal{O}_{S, s}$.
\end{lemma}

\begin{proof}
On peut remplacer $X$ par un voisinage affine de $x$ ; on peut donc supposer, et l'on suppose, $f : X \to S$ séparé. Alors $\omega_X^\bullet = f^!\omega_S^\bullet$. Il faut montrer que
$R\Hom_{\mathcal{O}_{X, x}}(\kappa(x), \omega_{X, x}^\bullet)$
est concentré en degré $0$. Notons $i_x : x \to X$ le morphisme d'inclusion, qui est une immersion fermée puisque $x$ est un point fermé. Ainsi, $R\Hom_{\mathcal{O}_{X, x}}(\kappa(x), \omega_{X, x}^\bullet)$ représente $i_x^!\omega_X^\bullet$ d'après le lemme \ref{duality-lemma-shriek-closed-immersion}. Considérons le diagramme commutatif
$$
\xymatrix{
x \ar[r]_{i_x} \ar[d]_\pi & X \ar[d]^f \\
s \ar[r]^{i_s} & S
}
$$
D'après Morphismes, lemme \ref{morphisms-lemma-closed-point-fibre-locally-finite-type}, l'extension $\kappa(x)/\kappa(s)$ est finie ; ainsi, $\pi$ est un morphisme fini. On en conclut que
$$
i_x^!\omega_X^\bullet = i_x^! f^! \omega_S^\bullet =
\pi^! i_s^! \omega_S^\bullet
$$
Si $\omega_{S, s}^\bullet$ est un complexe dualisant normalisé sur $\mathcal{O}_{S, s}$, le même raisonnement donne donc $i_s^!\omega_S^\bullet = \kappa(s)[0]$. On a
$$
R\pi_*(\pi^!(\kappa(s)[0])) =
R\SheafHom_{\mathcal{O}_s}(R\pi_*(\kappa(x)[0]), \kappa(s)[0]) =
\widetilde{\Hom_{\kappa(s)}(\kappa(x), \kappa(s))}
$$
La première égalité résulte de l'exemple \ref{duality-example-iso-on-RSheafHom-noetherian} appliqué avec $L = \kappa(x)[0]$. La seconde résulte de l'exactitude de $\pi_*$. Ainsi, $\pi^!(\kappa(s)[0])$ est concentré en degré $0$, ce qui conclut.
\end{proof}

\begin{lemma}
\label{duality-lemma-good-dualizing-dimension-function}
Soit $S$ un schéma noethérien et soit $\omega_S^\bullet$ un complexe dualisant. Soit $f : X \to S$ de type fini et soit $\omega_X^\bullet$ le complexe dualisant normalisé relativement à $\omega_S^\bullet$. Pour tout $x \in X$, on a
$$
\delta_X(x) - \delta_S(f(x)) = \text{trdeg}_{\kappa(f(x))}(\kappa(x))
$$
où $\delta_S$, resp.\ $\delta_X$, est la fonction de dimension de $\omega_S^\bullet$, resp.\ $\omega_X^\bullet$, voir le lemme \ref{duality-lemma-dimension-function-scheme}.
\end{lemma}

\begin{proof}
On peut remplacer $X$ par un voisinage affine de $x$. On peut donc supposer, et l'on suppose, qu'il existe une compactification $X \subset \overline{X}$ sur $S$. On peut alors remplacer $X$ par $\overline{X}$ et supposer $X$ propre sur $S$. On peut aussi supposer $X$ connexe en remplaçant $X$ par la composante connexe de $X$ contenant $x$. Rappelons ensuite que $\delta_X$ et la fonction
$x \mapsto \delta_S(f(x)) + \text{trdeg}_{\kappa(f(x))}(\kappa(x))$
sont toutes deux des fonctions de dimension sur $X$, voir Morphismes, lemme \ref{morphisms-lemma-dimension-function-propagates} (et le fait que $S$ est universellement caténaire d'après le lemme \ref{duality-lemma-dimension-function-scheme}). D'après Topologie, lemme \ref{topology-lemma-dimension-function-unique}, leur différence est localement constante, donc constante puisque $X$ est connexe. Il suffit donc de montrer l'égalité en un point quelconque de $X$. D'après Propriétés, lemme \ref{properties-lemma-locally-Noetherian-closed-point}, le schéma $X$ possède un point fermé $x$. Puisque $X \to S$ est propre, l'image $s$ de $x$ est fermée dans $S$. On peut donc appliquer le lemme \ref{duality-lemma-good-dualizing-normalized} pour conclure.
\end{proof}

\begin{lemma}
\label{duality-lemma-shriek}
Dans la situation \ref{duality-situation-shriek}, soit $f : X \to Y$ un morphisme de $\textit{FTS}_S$. Soit $x \in X$ d'image $y \in Y$. Alors
$$
H^i(f^!\mathcal{O}_Y)_x \not = 0
\Rightarrow - \dim_x(X_y) \leq i.
$$
\end{lemma}

\begin{proof}
L'énoncé étant local sur $X$, on peut supposer que $X$ et $Y$ sont affines. Écrivons $X = \Spec(A)$ et $Y = \Spec(R)$. Alors $f^!\mathcal{O}_Y$ correspond au complexe dualisant relatif $\omega_{A/R}^\bullet$ de Complexes dualisants, section \ref{dualizing-section-relative-dualizing-complexes-Noetherian}, d'après la remarque \ref{duality-remark-local-calculation-shriek}. Le lemme résulte donc de Complexes dualisants, lemme \ref{dualizing-lemma-relative-dualizing-trivial-vanishing}.
\end{proof}

\begin{lemma}
\label{duality-lemma-flat-shriek}
Dans la situation \ref{duality-situation-shriek}, soit $f : X \to Y$ un morphisme de $\textit{FTS}_S$. Soit $x \in X$ d'image $y \in Y$. Si $f$ est plat, alors
$$
H^i(f^!\mathcal{O}_Y)_x \not = 0
\Rightarrow - \dim_x(X_y) \leq i \leq 0.
$$
En fait, si toutes les fibres de $f$ sont de dimension $\leq d$, alors $f^!\mathcal{O}_Y$ est de Tor-amplitude dans $[-d, 0]$ comme objet de $D(X, f^{-1}\mathcal{O}_Y)$.
\end{lemma}

\begin{proof}
En raisonnant exactement comme dans la démonstration du lemme \ref{duality-lemma-shriek}, cela résulte de Complexes dualisants, lemme \ref{dualizing-lemma-relative-dualizing-flat-vanishing}.
\end{proof}

\begin{lemma}
\label{duality-lemma-shriek-over-CM}
Dans la situation \ref{duality-situation-shriek}, soit $f : X \to Y$ un morphisme de $\textit{FTS}_S$. Soit $x \in X$ d'image $y \in Y$. Supposons que
\begin{enumerate}
\item $\mathcal{O}_{Y, y}$ soit de Cohen–Macaulay ;
\item $\text{trdeg}_{\kappa(f(\xi))}(\kappa(\xi)) \leq r$ pour tout point générique $\xi$ d'une composante irréductible de $X$ contenant $x$.
\end{enumerate}
Alors
$$
H^i(f^!\mathcal{O}_Y)_x \not = 0
\Rightarrow - r \leq i
$$
et la fibre $H^{-r}(f^!\mathcal{O}_Y)_x$ est $(S_2)$ comme $\mathcal{O}_{X, x}$-module.
\end{lemma}

\begin{proof}
Après avoir remplacé $X$ par un voisinage ouvert de $x$, on peut supposer que toute composante irréductible de $X$ passe par $x$. En raisonnant exactement comme dans la démonstration du lemme \ref{duality-lemma-shriek}, cela résulte alors de Complexes dualisants, lemme \ref{dualizing-lemma-relative-dualizing-CM-vanishing}.
\end{proof}

\begin{lemma}
\label{duality-lemma-flat-quasi-finite-shriek}
Dans la situation \ref{duality-situation-shriek}, soit $f : X \to Y$ un morphisme de $\textit{FTS}_S$. Si $f$ est plat et quasi-fini, alors
$$
f^!\mathcal{O}_Y = \omega_{X/Y}[0]
$$
pour un $\mathcal{O}_X$-module cohérent $\omega_{X/Y}$ plat sur $Y$.
\end{lemma}

\begin{proof}
Cela résulte du lemme \ref{duality-lemma-flat-shriek} et du fait que les faisceaux de cohomologie de $f^!\mathcal{O}_Y$ sont cohérents d'après le lemme \ref{duality-lemma-shriek-coherent}.
\end{proof}

\begin{lemma}
\label{duality-lemma-CM-shriek}
Dans la situation \ref{duality-situation-shriek}, soit $f : X \to Y$ un morphisme de $\textit{FTS}_S$. Si $f$ est de Cohen–Macaulay (Compléments sur les morphismes, définition \ref{more-morphisms-definition-CM}), alors
$$
f^!\mathcal{O}_Y = \omega_{X/Y}[d]
$$
pour un $\mathcal{O}_X$-module cohérent $\omega_{X/Y}$ plat sur $Y$, où $d$ est la fonction localement constante sur $X$ donnant la dimension relative de $X$ sur $Y$.
\end{lemma}

\begin{proof}
La dimension relative $d$ est bien définie et localement constante d'après Morphismes, lemme \ref{morphisms-lemma-flat-finite-presentation-CM-fibres-relative-dimension}. Les faisceaux de cohomologie de $f^!\mathcal{O}_Y$ sont cohérents d'après le lemme \ref{duality-lemma-shriek-coherent}. On obtiendra la platitude de $\omega_{X/Y}$ par le lemme \ref{duality-lemma-flat-shriek} si l'on montre que les autres faisceaux de cohomologie de $f^!\mathcal{O}_Y$ sont nuls.

\medskip\noindent
La question est locale sur $X$ ; on peut donc supposer $X$ et $Y$ affines et le morphisme de dimension relative $d$. Si $d = 0$, le résultat découle directement du lemme \ref{duality-lemma-flat-quasi-finite-shriek}. Si $d > 0$, on peut supposer qu'il existe une factorisation
$$
X \xrightarrow{g} \mathbf{A}^d_Y \xrightarrow{p} Y
$$
avec $g$ quasi-fini et plat, voir Compléments sur les morphismes, lemme \ref{more-morphisms-lemma-flat-finite-presentation-characterize-CM}. Alors $f^! = g^! \circ p^!$. Le lemme \ref{duality-lemma-shriek-affine-line} montre que $p^!\mathcal{O}_Y \cong \mathcal{O}_{\mathbf{A}^d_Y}[-d]$. On conclut par le cas $d = 0$.
\end{proof}

\begin{remark}
\label{duality-remark-the-same-is-true}
Soit $S$ un schéma noethérien muni d'un complexe dualisant $\omega_S^\bullet$. Dans ce cas, les lemmes \ref{duality-lemma-shriek}, \ref{duality-lemma-flat-shriek}, \ref{duality-lemma-flat-quasi-finite-shriek} et \ref{duality-lemma-CM-shriek} sont vrais pour tout morphisme $f : X \to Y$ de schémas de type fini sur $S$, en remplaçant toutefois $f^!$ par $f_{new}^!$. C'est clair, car chaque démonstration se ramène immédiatement au cas affine, où $f^! = f_{new}^!$ d'après le lemme \ref{duality-lemma-duality-bootstrap}.
\end{remark}




\section{Modules dualisants}
\label{duality-section-dualizing-module}


\noindent
Cette section prolonge Complexes dualisants, section \ref{dualizing-section-dualizing-module}.

\medskip\noindent
Soit $X$ un schéma noethérien et soit $\omega_X^\bullet$ un complexe dualisant. Soit $n \in \mathbf{Z}$ le plus petit entier tel que $H^n(\omega_X^\bullet)$ soit non nul. Autrement dit, $-n$ est la valeur maximale de la fonction de dimension associée à $\omega_X^\bullet$ (lemme \ref{duality-lemma-dimension-function-scheme}). On appelle parfois $H^n(\omega_X^\bullet)$ un {\it module dualisant} ou un {\it faisceau dualisant} sur $X$, et on le note alors souvent $\omega_X$. Nous écrirons « soit $\omega_X$ un module dualisant » pour désigner cette situation.

\medskip\noindent
\begingroup\emergencystretch=2em
L'utilisation des modules dualisants $\omega_X$ sur les schémas noethériens $X$ exige quelques précautions :\par\endgroup
\begin{enumerate}
\item l'entier $n$ peut changer lorsqu'on passe de $X$ à un ouvert $U$ de $X$, et l'on n'aura alors pas nécessairement $\omega_X|_U = \omega_U$ ;
\item le complexe dualisant n'est pas unique ; le module dualisant n'est unique qu'à tensorisation par un module inversible près.
\end{enumerate}
La seconde difficulté sera souvent sans conséquence, car nous travaillerons avec $X$ de type fini sur une base $S$ munie d'un complexe dualisant fixé $\omega_S^\bullet$, et $\omega_X^\bullet$ sera le complexe dualisant normalisé relativement à $\omega_S^\bullet$. La première difficulté ne se présente pas si $X$ est équidimensionnel ; plus précisément, si la fonction de dimension associée à $\omega_X^\bullet$ (lemme \ref{duality-lemma-dimension-function-scheme}) prend la même valeur entière en tout point générique de $X$.

\begin{example}
\label{duality-example-proper-over-local}
Supposons $S = \Spec(A)$, où $(A, \mathfrak m, \kappa)$ est un anneau local noethérien, et où $\omega_S^\bullet$ correspond à un complexe dualisant normalisé $\omega_A^\bullet$. Alors, si $f : X \to S$ est propre sur $S$ et si $\omega_X^\bullet = f^!\omega_S^\bullet$, le faisceau cohérent
$$
\omega_X = H^{-\dim(X)}(\omega_X^\bullet)
$$
est un module dualisant, souvent appelé le module dualisant de $X$ ($S$ et $\omega_S^\bullet$ étant sous-entendus). Nous verrons qu'il possède de bonnes propriétés.
\end{example}

\begin{example}
\label{duality-example-equidimensional-over-field}
Supposons que $X$ soit un schéma équidimensionnel de type fini sur un corps $k$. Il est alors usuel de prendre pour $\omega_X^\bullet$ le complexe dualisant normalisé relativement à $k[0]$, et d'appeler
$$
\omega_X = H^{-\dim(X)}(\omega_X^\bullet)
$$
le module dualisant de $X$. Si $X$ est séparé sur $k$, alors $\omega_X^\bullet = f^!\mathcal{O}_{\Spec(k)}$, où $f : X \to \Spec(k)$ est le morphisme structural, d'après le lemme \ref{duality-lemma-duality-bootstrap}. Si $X$ est propre sur $k$, c'est un cas particulier de l'exemple \ref{duality-example-proper-over-local}.
\end{example}

\begin{lemma}
\label{duality-lemma-dualizing-module}
Soit $X$ un schéma noethérien connexe et soit $\omega_X$ un module dualisant sur $X$. Le support de $\omega_X$ est la réunion des composantes irréductibles de dimension maximale pour toute fonction de dimension, et $\omega_X$ est un $\mathcal{O}_X$-module cohérent satisfaisant la propriété $(S_2)$.
\end{lemma}

\begin{proof}
Selon les conventions précédentes, il existe un complexe dualisant $\omega_X^\bullet$ dont $\omega_X$ est le premier faisceau de cohomologie non nul à partir de la gauche. Puisque $X$ est connexe, deux fonctions de dimension diffèrent d'une constante (Topologie, lemme \ref{topology-lemma-dimension-function-unique}). On peut donc utiliser la fonction de dimension associée à $\omega_X^\bullet$ (lemme \ref{duality-lemma-dimension-function-scheme}). Ces remarques faites, le lemme résulte de Complexes dualisants, lemme \ref{dualizing-lemma-depth-dualizing-module}, et des définitions (en particulier Cohomologie des schémas, définition \ref{coherent-definition-depth}).
\end{proof}

\begin{lemma}
\label{duality-lemma-vanishing-good-dualizing}
Soient $X/A$, avec $\omega_X^\bullet$, et $\omega_X$ comme dans l'exemple \ref{duality-example-proper-over-local}. Alors
\begin{enumerate}
\item $H^i(\omega_X^\bullet) \not = 0 \Rightarrow
i \in \{-\dim(X), \ldots, 0\}$ ;
\item la dimension du support de $H^i(\omega_X^\bullet)$ est au plus $-i$ ;
\item $\text{Supp}(\omega_X)$ est la réunion des composantes de dimension $\dim(X)$ ;
\item $\omega_X$ satisfait la propriété $(S_2)$.
\end{enumerate}
\end{lemma}

\begin{proof}
Soient $\delta_X$ et $\delta_S$ les fonctions de dimension associées à $\omega_X^\bullet$ et $\omega_S^\bullet$ comme au lemme \ref{duality-lemma-good-dualizing-dimension-function}. Puisque $X$ est propre sur $A$, tout sous-schéma fermé de $X$ contient un point fermé $x$ dont l'image est le point fermé $s \in S$, et $\delta_X(x) = \delta_S(s) = 0$. Ainsi, $\delta_X(\xi) = \dim(\overline{\{\xi\}})$ pour tout point $\xi \in X$. On peut donc vérifier chacune des assertions du lemme en considérant ce qui se passe sur $\Spec(\mathcal{O}_{X, x})$ ; le résultat découle alors de Complexes dualisants, lemmes \ref{dualizing-lemma-sitting-in-degrees} et \ref{dualizing-lemma-depth-dualizing-module}. Certains détails sont omis. Les deux dernières assertions peuvent aussi se déduire du lemme \ref{duality-lemma-dualizing-module}.
\end{proof}

\begin{lemma}
\label{duality-lemma-dualizing-module-proper-over-A}
Soit $X/A$, de module dualisant $\omega_X$, comme dans l'exemple \ref{duality-example-proper-over-local}. Soit $d = \dim(X_s)$ la dimension de la fibre fermée. Si $\dim(X) = d + \dim(A)$, alors le module dualisant $\omega_X$ représente le foncteur
$$
\mathcal{F} \longmapsto \Hom_A(H^d(X, \mathcal{F}), \omega_A)
$$
sur la catégorie des $\mathcal{O}_X$-modules cohérents.
\end{lemma}

\begin{proof}
On a
\begin{align*}
\Hom_X(\mathcal{F}, \omega_X)
& =
\Ext^{-\dim(X)}_X(\mathcal{F}, \omega_X^\bullet) \\
& =
\Hom_X(\mathcal{F}[\dim(X)], \omega_X^\bullet) \\
& =
\Hom_X(\mathcal{F}[\dim(X)], f^!(\omega_A^\bullet)) \\
& =
\Hom_S(Rf_*\mathcal{F}[\dim(X)], \omega_A^\bullet) \\
& =
\Hom_A(H^d(X, \mathcal{F}), \omega_A)
\end{align*}
La première égalité résulte de $H^i(\omega_X^\bullet) = 0$ pour $i < -\dim(X)$, voir le lemme \ref{duality-lemma-vanishing-good-dualizing} et Catégories dérivées, lemme \ref{derived-lemma-negative-exts}. La deuxième résulte de la définition des groupes Ext. La troisième est notre choix de $\omega_X^\bullet$. La quatrième résulte du fait que $f^!$ est l'adjoint à droite du lemme \ref{duality-lemma-twisted-inverse-image} pour $f$, voir la section \ref{duality-section-duality}. La dernière égalité résulte du fait que $R^if_*\mathcal{F}$ est nul pour $i > d$ (Cohomologie des schémas, lemme \ref{coherent-lemma-higher-direct-images-zero-above-dimension-fibre}), et que $H^j(\omega_A^\bullet)$ est nul pour $j < -\dim(A)$.
\end{proof}








\section{Schémas de Cohen--Macaulay}
\label{duality-section-CM}

\noindent
Cette section prolonge Complexes dualisants, section \ref{dualizing-section-CM}. La dualité prend une forme particulièrement simple pour les schémas de Cohen–Macaulay.

\begin{lemma}
\label{duality-lemma-dualizing-module-CM-scheme}
Soit $X$ un schéma localement noethérien muni d'un complexe dualisant $\omega_X^\bullet$.
\begin{enumerate}
\item $X$ est de Cohen–Macaulay $\Leftrightarrow$ $\omega_X^\bullet$ possède localement un unique faisceau de cohomologie non nul ;
\item $\mathcal{O}_{X, x}$ est de Cohen–Macaulay $\Leftrightarrow$ $\omega_{X, x}^\bullet$ possède une unique cohomologie non nulle ;
\item $U = \{x \in X \mid \mathcal{O}_{X, x}\text{ est de Cohen--Macaulay}\}$ est ouvert et de Cohen–Macaulay.
\end{enumerate}
Si $X$ est connexe et de Cohen–Macaulay, il existe un entier $n$ et un $\mathcal{O}_X$-module cohérent de Cohen–Macaulay $\omega_X$ tels que $\omega_X^\bullet = \omega_X[-n]$.
\end{lemma}

\begin{proof}
Par définition et d'après Complexes dualisants, lemme \ref{dualizing-lemma-dualizing-localize}, pour tout $x \in X$, le complexe $\omega_{X, x}^\bullet$ est un complexe dualisant sur $\mathcal{O}_{X, x}$. Complexes dualisants, lemme \ref{dualizing-lemma-apply-CM}, donne donc (2).

\medskip\noindent
Pour démontrer (3), supposons $\mathcal{O}_{X, x}$ de Cohen–Macaulay. Soit $n_x$ l'unique entier tel que $H^{n_{x}}(\omega_{X, x}^\bullet)$ soit non nul. Pour un voisinage affine $V \subset X$ de $x$, on a $\omega_X^\bullet|_V$ dans $D^b_{\textit{Coh}}(\mathcal{O}_V)$ ; il n'y a donc qu'un nombre fini de modules cohérents non nuls $H^i(\omega_X^\bullet)|_V$. Après avoir rétréci $V$, on peut ainsi supposer que seul $H^{n_x}$ est non nul, voir Modules, lemme \ref{modules-lemma-finite-type-stalk-zero}. On voit de cette façon que $\mathcal{O}_{X, v}$ est de Cohen–Macaulay pour tout $v \in V$. Cela montre que $U$ est ouvert et est un schéma de Cohen–Macaulay.

\medskip\noindent
Démonstration de (1). L'implication $\Leftarrow$ résulte de (2). L'implication $\Rightarrow$ résulte du paragraphe précédent, où l'on a montré que, si $\mathcal{O}_{X, x}$ est de Cohen–Macaulay, alors, dans un voisinage de $x$, le complexe $\omega_X^\bullet$ ne possède qu'un seul faisceau de cohomologie non nul.

\medskip\noindent
Supposons $X$ connexe et de Cohen–Macaulay. Ce qui précède montre que l'application $x \mapsto n_x$ est localement constante. Puisque $X$ est connexe, elle est constante, de valeur $n$, disons. En posant $\omega_X = H^n(\omega_X^\bullet)$, on obtient le lemme, car $\omega_X$ est de Cohen–Macaulay d'après Complexes dualisants, lemme \ref{dualizing-lemma-apply-CM} (et Cohomologie des schémas, définition \ref{coherent-definition-Cohen-Macaulay}).
\end{proof}

\begin{lemma}
\label{duality-lemma-has-dualizing-module-CM-scheme}
Soit $X$ un schéma localement noethérien. S'il existe un faisceau cohérent $\omega_X$ tel que $\omega_X[0]$ soit un complexe dualisant sur $X$, alors $X$ est un schéma de Cohen–Macaulay.
\end{lemma}

\begin{proof}
Cela résulte immédiatement de Complexes dualisants, lemme \ref{dualizing-lemma-has-dualizing-module-CM}, et de nos définitions.
\end{proof}

\begin{lemma}
\label{duality-lemma-affine-flat-Noetherian-CM}
Dans la situation \ref{duality-situation-shriek}, soit $f : X \to Y$ un morphisme de $\textit{FTS}_S$. Soit $x \in X$. Si $f$ est plat, les conditions suivantes sont équivalentes :
\begin{enumerate}
\item $f$ est de Cohen–Macaulay en $x$ ;
\item $f^!\mathcal{O}_Y$ possède un unique faisceau de cohomologie non nul dans un voisinage de $x$.
\end{enumerate}
\end{lemma}

\begin{proof}
Un sens du lemme résulte du lemme \ref{duality-lemma-CM-shriek}. Pour démontrer la réciproque, on peut supposer que $f^!\mathcal{O}_Y$ possède un unique faisceau de cohomologie non nul. Soit $y = f(x)$. Soient $\xi_1, \ldots, \xi_n \in X_y$ les points génériques de la fibre $X_y$ se spécialisant en $x$. Soient $d_1, \ldots, d_n$ les dimensions des composantes irréductibles correspondantes de $X_y$. Le morphisme $f : X \to Y$ est de Cohen–Macaulay en $\eta_i$ d'après Compléments sur les morphismes, lemme \ref{more-morphisms-lemma-flat-finite-presentation-CM-open}. Le lemme \ref{duality-lemma-CM-shriek} donne donc $d_1 = \ldots = d_n$. Si $d$ désigne la valeur commune, alors $d = \dim_x(X_y)$. Après avoir rétréci $X$, on peut supposer toutes les fibres de dimension au plus $d$ (Morphismes, lemme \ref{morphisms-lemma-openness-bounded-dimension-fibres}). L'unique faisceau de cohomologie non nul $\omega = H^{-d}(f^!\mathcal{O}_Y)$ est alors plat sur $Y$ d'après le lemme \ref{duality-lemma-flat-shriek}. Si $h : X_y \to X$ désigne le morphisme canonique, on a donc
$Lh^*(f^!\mathcal{O}_Y) = Lh^*(\omega[d]) = (h^*\omega)[d]$
d'après Catégories dérivées des schémas, lemme \ref{perfect-lemma-tor-independence-and-tor-amplitude}. Ainsi, $h^*\omega[d]$ est le complexe dualisant de $X_y$ d'après le lemme \ref{duality-lemma-relative-dualizing-fibres}. Par conséquent, $X_y$ est de Cohen–Macaulay d'après le lemme \ref{duality-lemma-dualizing-module-CM-scheme}. Cela montre que $f$ est de Cohen–Macaulay en $x$, comme voulu.
\end{proof}

\begin{remark}
\label{duality-remark-CM-morphism-compare-dualizing}
Dans la situation \ref{duality-situation-shriek}, soit $f : X \to Y$ un morphisme de $\textit{FTS}_S$. Supposons que $f$ soit un morphisme de Cohen–Macaulay de dimension relative $d$. Soit $\omega_{X/Y} = H^{-d}(f^!\mathcal{O}_Y)$ l'unique faisceau de cohomologie non nul de $f^!\mathcal{O}_Y$, voir le lemme \ref{duality-lemma-CM-shriek}. Il existe alors un isomorphisme canonique
$$
f^!K = Lf^*K \otimes_{\mathcal{O}_X}^\mathbf{L} \omega_{X/Y}[d]
$$
pour $K \in D^+_\QCoh(\mathcal{O}_Y)$, voir le lemme \ref{duality-lemma-perfect-comparison-shriek}. En particulier, si $S$ admet un complexe dualisant $\omega_S^\bullet$, si $\omega_Y^\bullet = (Y \to S)^!\omega_S^\bullet$ et si $\omega_X^\bullet = (X \to S)^!\omega_S^\bullet$, on a
$$
\omega_X^\bullet =
Lf^*\omega_Y^\bullet \otimes_{\mathcal{O}_X}^\mathbf{L} \omega_{X/Y}[d]
$$
Si de plus $X$ et $Y$ sont connexes et de Cohen–Macaulay, et si $\omega_Y$ et $\omega_X$ désignent les uniques faisceaux de cohomologie non nuls de $\omega_Y^\bullet$ et $\omega_X^\bullet$, on a donc
$$
\omega_X = f^*\omega_Y \otimes_{\mathcal{O}_X} \omega_{X/Y}.
$$
Des résultats analogues valent pour $X$ et $Y$ schémas quelconques de type fini sur $S$ (donc pas nécessairement séparés sur $S$), avec des complexes dualisants normalisés relativement à $\omega_S^\bullet$ comme à la section \ref{duality-section-glue}.
\end{remark}





\section{Schémas de Gorenstein}
\label{duality-section-gorenstein}

\noindent
Cette section prolonge Complexes dualisants, section \ref{dualizing-section-gorenstein}.

\begin{definition}
\label{duality-definition-gorenstein}
Soit $X$ un schéma. On dit que $X$ est {\it de Gorenstein} si $X$ est localement noethérien et si $\mathcal{O}_{X, x}$ est de Gorenstein pour tout $x \in X$.
\end{definition}

\noindent
Cette définition a un sens, car un anneau noethérien est dit de Gorenstein si et seulement si tous ses anneaux locaux sont de Gorenstein, voir Complexes dualisants, définition \ref{dualizing-definition-gorenstein}.

\begin{lemma}
\label{duality-lemma-gorenstein-CM}
Un schéma de Gorenstein est de Cohen–Macaulay.
\end{lemma}

\begin{proof}
En travaillant localement sur les ouverts affines, cela résulte du résultat algébrique correspondant, à savoir Complexes dualisants, lemme \ref{dualizing-lemma-gorenstein-CM}.
\end{proof}

\begin{lemma}
\label{duality-lemma-regular-gorenstein}
Un schéma régulier est de Gorenstein.
\end{lemma}

\begin{proof}
En travaillant localement sur les ouverts affines, cela résulte du résultat algébrique correspondant, à savoir Complexes dualisants, lemme \ref{dualizing-lemma-regular-gorenstein}.
\end{proof}

\begin{lemma}
\label{duality-lemma-gorenstein}
Soit $X$ un schéma localement noethérien.
\begin{enumerate}
\item Si $X$ admet un complexe dualisant $\omega_X^\bullet$, alors
\begin{enumerate}
\item $X$ est de Gorenstein $\Leftrightarrow$ $\omega_X^\bullet$ est un objet inversible de $D(\mathcal{O}_X)$ ;
\item $\mathcal{O}_{X, x}$ est de Gorenstein $\Leftrightarrow$ $\omega_{X, x}^\bullet$ est un objet inversible de $D(\mathcal{O}_{X, x})$ ;
\item $U = \{x \in X \mid \mathcal{O}_{X, x}\text{ est de Gorenstein}\}$ est un sous-schéma ouvert de Gorenstein.
\end{enumerate}
\item Si $X$ est de Gorenstein, alors $X$ admet un complexe dualisant si et seulement si $\mathcal{O}_X[0]$ est un complexe dualisant.
\end{enumerate}
\end{lemma}

\begin{proof}
En travaillant localement sur les ouverts affines, cela résulte du résultat algébrique correspondant, à savoir Complexes dualisants, lemme \ref{dualizing-lemma-gorenstein}.
\end{proof}

\begin{lemma}
\label{duality-lemma-gorenstein-lci}
Si $f : Y \to X$ est un morphisme localement d'intersection complète et si $X$ est un schéma de Gorenstein, alors $Y$ est de Gorenstein.
\end{lemma}

\begin{proof}
D'après Compléments sur les morphismes, lemme \ref{more-morphisms-lemma-affine-lci}, il suffit de démontrer l'énoncé correspondant pour les homomorphismes d'anneaux. C'est Complexes dualisants, lemme \ref{dualizing-lemma-gorenstein-lci}.
\end{proof}

\begin{lemma}
\label{duality-lemma-gorenstein-local-syntomic}
La propriété $\mathcal{P}(S) =$« $S$ est de Gorenstein » est locale pour la topologie syntomique.
\end{lemma}

\begin{proof}
Soit $\{S_i \to S\}$ un recouvrement syntomique. Le schéma $S$ est localement noethérien si et seulement si chaque $S_i$ est noethérien, voir Descente, lemme \ref{descent-lemma-Noetherian-local-fppf}. On peut donc supposer $S$ et $S_i$ localement noethériens. Si $S$ est de Gorenstein, chaque $S_i$ est de Gorenstein d'après le lemme \ref{duality-lemma-gorenstein-lci}. Réciproquement, si chaque $S_i$ est de Gorenstein, on peut, pour tout point $s \in S$, choisir $i$ et $t \in S_i$ d'image $s$. Alors $\mathcal{O}_{S, s} \to \mathcal{O}_{S_i, t}$ est un homomorphisme local plat d'anneaux, avec $\mathcal{O}_{S_i, t}$ de Gorenstein. Ainsi, $\mathcal{O}_{S, s}$ est de Gorenstein d'après Complexes dualisants, lemme \ref{dualizing-lemma-flat-under-gorenstein}.
\end{proof}






\section{Morphismes de Gorenstein}
\label{duality-section-gorenstein-morphisms}

\noindent
Cette section s'inscrit dans une série. Les sections correspondantes pour les morphismes normaux, réguliers et de Cohen–Macaulay se trouvent dans Compléments sur les morphismes, sections \ref{more-morphisms-section-normal}, \ref{more-morphisms-section-regular} et \ref{more-morphisms-section-CM}.

\medskip\noindent
Le lemme suivant montre qu'il n'y a pas lieu de définir les schémas géométriquement de Gorenstein, car ils coïncideraient avec les schémas de Gorenstein.

\begin{lemma}
\label{duality-lemma-gorenstein-base-change}
Soit $X$ un schéma localement noethérien sur le corps $k$. Soit $k'/k$ une extension de corps de type fini. Soit $x \in X$ un point, et soit $x' \in X_{k'}$ un point au-dessus de $x$. Alors
$$
\mathcal{O}_{X, x}\text{ est de Gorenstein}
\Leftrightarrow
\mathcal{O}_{X_{k'}, x'}\text{ est de Gorenstein}
$$
Si $X$ est localement de type fini sur $k$, le même énoncé vaut pour toute extension de corps $k'/k$.
\end{lemma}

\begin{proof}
Dans les deux cas, l'homomorphisme d'anneaux $\mathcal{O}_{X, x} \to \mathcal{O}_{X_{k'}, x'}$ est un homomorphisme local fidèlement plat d'anneaux locaux noethériens. Si $\mathcal{O}_{X_{k'}, x'}$ est de Gorenstein, il en est donc de même de $\mathcal{O}_{X, x}$ d'après Complexes dualisants, lemme \ref{dualizing-lemma-flat-under-gorenstein}. Pour remonter, on utilise également Complexes dualisants, lemme \ref{dualizing-lemma-flat-under-gorenstein}. Il faut donc montrer que
$$
\mathcal{O}_{X_{k'}, x'}/\mathfrak m_x \mathcal{O}_{X_{k'}, x'} =
\kappa(x) \otimes_k k'
$$
est de Gorenstein. Remarquons que, dans le premier cas, $k \to k'$ est de type fini, et que, dans le second, $k \to \kappa(x)$ est de type fini. Le résultat découle donc de la propriété (A) pour les anneaux de Gorenstein, voir Complexes dualisants, lemme \ref{dualizing-lemma-formal-fibres-gorenstein}.
\end{proof}

\noindent
Le lemme précédent garantit que la définition suivante des morphismes de Gorenstein est la bonne.

\begin{definition}
\label{duality-definition-gorenstein-morphism}
Soit $f : X \to Y$ un morphisme de schémas. Supposons que toutes les fibres $X_y$ soient des schémas localement noethériens.
\begin{enumerate}
\item Soient $x \in X$ et $y = f(x)$. On dit que $f$ est {\it de Gorenstein en $x$} si $f$ est plat en $x$ et si l'anneau local du schéma $X_y$ en $x$ est de Gorenstein.
\item On dit que $f$ est un {\it morphisme de Gorenstein} si $f$ est de Gorenstein en tout point de $X$.
\end{enumerate}
\end{definition}

\noindent
Voici une reformulation.

\begin{lemma}
\label{duality-lemma-gorenstein-morphism}
Soit $f : X \to Y$ un morphisme de schémas. Supposons toutes les fibres de $f$ localement noethériennes. Les conditions suivantes sont équivalentes :
\begin{enumerate}
\item $f$ est de Gorenstein ;
\item $f$ est plat et ses fibres sont des schémas de Gorenstein.
\end{enumerate}
\end{lemma}

\begin{proof}
Cela résulte directement des définitions.
\end{proof}

\begin{lemma}
\label{duality-lemma-gorenstein-CM-morphism}
Un morphisme de Gorenstein est de Cohen–Macaulay.
\end{lemma}

\begin{proof}
Cela résulte du lemme \ref{duality-lemma-gorenstein-CM} et des définitions.
\end{proof}

\begin{lemma}
\label{duality-lemma-lci-gorenstein}
Un morphisme syntomique est de Gorenstein. De manière équivalente, un morphisme plat localement d'intersection complète est de Gorenstein.
\end{lemma}

\begin{proof}
Rappelons qu'un morphisme syntomique est plat et que ses fibres sont localement des intersections complètes sur des corps, voir Morphismes, lemme \ref{morphisms-lemma-syntomic-flat-fibres}. Puisqu'une intersection complète locale sur un corps est un schéma de Gorenstein d'après le lemme \ref{duality-lemma-gorenstein-lci}, on conclut. Les propriétés « syntomique » et « plat et localement d'intersection complète » sont équivalentes d'après Compléments sur les morphismes, lemme \ref{more-morphisms-lemma-flat-lci}.
\end{proof}

\begin{lemma}
\label{duality-lemma-composition-gorenstein}
Soient $f : X \to Y$ et $g : Y \to Z$ des morphismes. Supposons que les fibres $X_y$, $Y_z$ et $X_z$ de $f$, $g$ et $g \circ f$ soient localement noethériennes.
\begin{enumerate}
\item Si $f$ est de Gorenstein en $x$ et si $g$ est de Gorenstein en $f(x)$, alors $g \circ f$ est de Gorenstein en $x$.
\item Si $f$ et $g$ sont de Gorenstein, alors $g \circ f$ est de Gorenstein.
\item Si $g \circ f$ est de Gorenstein en $x$ et si $f$ est plat en $x$, alors $f$ est de Gorenstein en $x$ et $g$ est de Gorenstein en $f(x)$.
\item Si $g \circ f$ est de Gorenstein et si $f$ est plat, alors $f$ est de Gorenstein et $g$ est de Gorenstein en tout point de l'image de $f$.
\end{enumerate}
\end{lemma}

\begin{proof}
Après traduction en algèbre, cela résulte de Complexes dualisants, lemme \ref{dualizing-lemma-flat-under-gorenstein}.
\end{proof}

\begin{lemma}
\label{duality-lemma-flat-morphism-from-gorenstein-scheme}
\begin{slogan}
Le caractère de Gorenstein de l'espace total d'une fibration plate entraîne le même caractère pour la base et les fibres
\end{slogan}
Soit $f : X \to Y$ un morphisme plat de schémas localement noethériens. Si $X$ est de Gorenstein, alors $f$ est de Gorenstein et $\mathcal{O}_{Y, f(x)}$ est de Gorenstein pour tout $x \in X$.
\end{lemma}

\begin{proof}
Après traduction en algèbre, cela résulte de Complexes dualisants, lemme \ref{dualizing-lemma-flat-under-gorenstein}.
\end{proof}

\begin{lemma}
\label{duality-lemma-base-change-gorenstein}
Soit $f : X \to Y$ un morphisme de schémas. Supposons que toutes les fibres $X_y$ soient des schémas localement noethériens. Soit $Y' \to Y$ localement de type fini. Soit $f' : X' \to Y'$ le changement de base de $f$. Soit $x' \in X'$ un point d'image $x \in X$.
\begin{enumerate}
\item Si $f$ est de Gorenstein en $x$, alors $f' : X' \to Y'$ est de Gorenstein en $x'$.
\item Si $f$ est plat en $x$ et si $f'$ est de Gorenstein en $x'$, alors $f$ est de Gorenstein en $x$.
\item Si $Y' \to Y$ est plat en $f'(x')$ et si $f'$ est de Gorenstein en $x'$, alors $f$ est de Gorenstein en $x$.
\end{enumerate}
\end{lemma}

\begin{proof}
Remarquons que l'hypothèse sur $Y' \to Y$ implique que, pour $y' \in Y'$ d'image $y \in Y$, l'extension de corps $\kappa(y')/\kappa(y)$ est de type fini. Par conséquent, toutes les fibres $X'_{y'} = (X_y)_{\kappa(y')}$ sont elles aussi localement noethériennes, voir Variétés, lemme \ref{varieties-lemma-locally-Noetherian-base-change}. Le lemme a donc un sens. Posons $y' = f'(x')$ et $y = f(x)$. On obtient ainsi le diagramme commutatif suivant d'anneaux locaux
$$
\xymatrix{
\mathcal{O}_{X', x'} & \mathcal{O}_{X, x} \ar[l] \\
\mathcal{O}_{Y', y'} \ar[u] & \mathcal{O}_{Y, y} \ar[l] \ar[u]
}
$$
dans lequel le coin supérieur gauche est un localisé du produit tensoriel des coins supérieur droit et inférieur gauche sur le coin inférieur droit.

\medskip\noindent
Supposons que $f$ soit de Gorenstein en $x$. La platitude de $\mathcal{O}_{Y, y} \to \mathcal{O}_{X, x}$ entraîne celle de $\mathcal{O}_{Y', y'} \to \mathcal{O}_{X', x'}$, voir Algèbre, lemme \ref{algebra-lemma-base-change-flat-up-down}. Le fait que $\mathcal{O}_{X, x}/\mathfrak m_y\mathcal{O}_{X, x}$ soit de Gorenstein implique que $\mathcal{O}_{X', x'}/\mathfrak m_{y'}\mathcal{O}_{X', x'}$ est de Gorenstein, voir le lemme \ref{duality-lemma-gorenstein-base-change}. On voit donc que $f'$ est de Gorenstein en $x'$.

\medskip\noindent
Supposons que $f$ soit plat en $x$ et que $f'$ soit de Gorenstein en $x'$. Le fait que $\mathcal{O}_{X', x'}/\mathfrak m_{y'}\mathcal{O}_{X', x'}$ soit de Gorenstein implique que $\mathcal{O}_{X, x}/\mathfrak m_y\mathcal{O}_{X, x}$ est de Gorenstein, voir le lemme \ref{duality-lemma-gorenstein-base-change}. On voit donc que $f$ est de Gorenstein en $x$.

\medskip\noindent
Supposons que $Y' \to Y$ soit plat en $y'$ et que $f'$ soit de Gorenstein en $x'$. La platitude de $\mathcal{O}_{Y', y'} \to \mathcal{O}_{X', x'}$ et de $\mathcal{O}_{Y, y} \to \mathcal{O}_{Y', y'}$ entraîne celle de $\mathcal{O}_{Y, y} \to \mathcal{O}_{X, x}$, voir Algèbre, lemme \ref{algebra-lemma-base-change-flat-up-down}. Le fait que $\mathcal{O}_{X', x'}/\mathfrak m_{y'}\mathcal{O}_{X', x'}$ soit de Gorenstein implique que $\mathcal{O}_{X, x}/\mathfrak m_y\mathcal{O}_{X, x}$ est de Gorenstein, voir le lemme \ref{duality-lemma-gorenstein-base-change}. On voit donc que $f$ est de Gorenstein en $x$.
\end{proof}

\begin{lemma}
\label{duality-lemma-flat-lft-base-change-gorenstein}
Soit $f : X \to Y$ un morphisme de schémas plat et localement de type fini. Alors la formation de l'ensemble
$\{x \in X \mid f\text{ est de Gorenstein en }x\}$
commute à tout changement de base.
\end{lemma}

\begin{proof}
L'hypothèse implique que toute fibre de $f$ est localement de type fini sur un corps, donc localement noethérienne, et il en est de même après tout changement de base. L'énoncé a donc un sens. En examinant les fibres, on se ramène au problème suivant : soit $X$ un schéma localement de type fini sur un corps $k$, soit $K/k$ une extension de corps, et soit $x_K \in X_K$ un point d'image $x \in X$. Problème : montrer que $\mathcal{O}_{X_K, x_K}$ est de Gorenstein si et seulement si $\mathcal{O}_{X, x}$ est de Gorenstein.

\medskip\noindent
On peut résoudre le problème avec un peu d'algèbre comme suit. Choisissons un ouvert affine $\Spec(A) \subset X$ contenant $x$. Écrivons que $x$ correspond à $\mathfrak p \subset A$. Avec $A_K = A \otimes_k K$, on voit que $\Spec(A_K) \subset X_K$ contient $x_K$. Écrivons que $x_K$ correspond à $\mathfrak p_K \subset A_K$. Soit $\omega_A^\bullet$ un complexe dualisant pour $A$. D'après Complexes dualisants, lemme \ref{dualizing-lemma-base-change-dualizing-over-field}, $\omega_A^\bullet \otimes_A A_K$ est un complexe dualisant pour $A_K$. On a terminé, car $A_\mathfrak p \to (A_K)_{\mathfrak p_K}$ est un homomorphisme local plat d'anneaux noethériens et, par conséquent, $(\omega_A^\bullet)_\mathfrak p$ est un objet inversible de $D(A_\mathfrak p)$ si et seulement si $(\omega_A^\bullet)_\mathfrak p \otimes_{A_\mathfrak p} (A_K)_{\mathfrak p_K}$ est un objet inversible de $D((A_K)_{\mathfrak p_K})$. Quelques détails sont omis ; indication : examiner les modules de cohomologie.
\end{proof}

\begin{lemma}
\label{duality-lemma-affine-flat-Noetherian-gorenstein}
Dans la situation \ref{duality-situation-shriek}, soit $f : X \to Y$ un morphisme de $\textit{FTS}_S$. Soit $x \in X$. Si $f$ est plat, les conditions suivantes sont équivalentes :
\begin{enumerate}
\item $f$ est de Gorenstein en $x$ ;
\item $f^!\mathcal{O}_Y$ est isomorphe à un objet inversible dans un voisinage de $x$.
\end{enumerate}
En particulier, l'ensemble des points où $f$ est de Gorenstein est ouvert dans $X$.
\end{lemma}

\begin{proof}
Posons $\omega^\bullet = f^!\mathcal{O}_Y$. D'après le lemme \ref{duality-lemma-relative-dualizing-fibres}, c'est un complexe borné à faisceaux de cohomologie cohérents dont la restriction dérivée $Lh^*\omega^\bullet$ à la fibre $X_y$ est un complexe dualisant sur $X_y$. Notons $i : x \to X_y$ l'inclusion d'un point. Les conditions suivantes sont alors équivalentes :
\begin{enumerate}
\item $f$ est de Gorenstein en $x$ ;
\item $\mathcal{O}_{X_y, x}$ est de Gorenstein ;
\item $Lh^*\omega^\bullet$ est inversible dans un voisinage de $x$ ;
\item $Li^* Lh^* \omega^\bullet$ possède exactement une cohomologie non nulle, de dimension $1$ sur $\kappa(x)$ ;
\item $L(h \circ i)^* \omega^\bullet$ possède exactement une cohomologie non nulle, de dimension $1$ sur $\kappa(x)$ ;
\item $\omega^\bullet$ est inversible dans un voisinage de $x$.
\end{enumerate}
L'équivalence de (1) et (2) résulte de la définition (puisque $f$ est plat). L'équivalence de (2) et (3) résulte du lemme \ref{duality-lemma-gorenstein}. L'équivalence de (3) et (4) résulte de Compléments sur l'algèbre, lemme \ref{more-algebra-lemma-lift-bounded-pseudo-coherent-to-perfect}. L'équivalence de (4) et (5) vaut parce que $Li^* Lh^* = L(h \circ i)^*$. L'équivalence de (5) et (6) résulte de Compléments sur l'algèbre, lemme \ref{more-algebra-lemma-lift-bounded-pseudo-coherent-to-perfect}. Le lemme est donc clair.
\end{proof}

\begin{lemma}
\label{duality-lemma-flat-finite-presentation-characterize-gorenstein}
Soit $f : X \to S$ un morphisme de schémas plat et localement de présentation finie. Soit $x \in X$ d'image $s \in S$. Posons $d = \dim_x(X_s)$. Les conditions suivantes sont équivalentes :
\begin{enumerate}
\item $f$ est de Gorenstein en $x$ ;
\item il existe un voisinage ouvert $U \subset X$ de $x$ et un morphisme localement quasi-fini $U \to \mathbf{A}^d_S$ sur $S$ qui est de Gorenstein en $x$ ;
\item il existe un voisinage ouvert $U \subset X$ de $x$ et un morphisme de Gorenstein localement quasi-fini $U \to \mathbf{A}^d_S$ sur $S$ ;
\item pour tout $S$-morphisme $g : U \to \mathbf{A}^d_S$ d'un voisinage ouvert $U \subset X$ de $x$, on a : $g$ est quasi-fini en $x$ $\Rightarrow$ $g$ est de Gorenstein en $x$.
\end{enumerate}
En particulier, l'ensemble des points où $f$ est de Gorenstein est ouvert dans $X$.
\end{lemma}

\begin{proof}
Choisissons des ouverts affines $U = \Spec(A) \subset X$ avec $x \in U$, et $V = \Spec(R) \subset S$ avec $f(U) \subset V$. Alors $R \to A$ est un homomorphisme d'anneaux plat de présentation finie. Soit $\mathfrak p \subset A$ l'idéal premier correspondant à $x$. Après avoir remplacé $A$ par une localisation principale, on peut supposer qu'il existe un homomorphisme quasi-fini $R[x_1, \ldots, x_d]  \to A$, voir Algèbre, lemme \ref{algebra-lemma-quasi-finite-over-polynomial-algebra}. Il existe donc au moins une paire $(U, g)$ formée d'un voisinage ouvert $U \subset X$ de $x$ et d'un morphisme localement\footnote{Si $S$ est quasi-séparé, alors $g$ sera quasi-fini.} quasi-fini $g : U \to \mathbf{A}^d_S$.

\medskip\noindent
Cela étant dit, le lemme se traduit par le problème algébrique suivant (traduction omise). Étant donnés $R \to A$ plat et de présentation finie, un idéal premier $\mathfrak p \subset A$ et $\varphi : R[x_1, \ldots, x_d] \to A$ quasi-fini en $\mathfrak p$, les conditions suivantes sont équivalentes :
\begin{enumerate}
\item[(a)] $\Spec(\varphi)$ est de Gorenstein en $\mathfrak p$ ;
\item[(b)] $\Spec(A) \to \Spec(R)$ est de Gorenstein en $\mathfrak p$ ;
\item[(c)] $\Spec(A) \to \Spec(R)$ est de Gorenstein dans un voisinage ouvert de $\mathfrak p$.
\end{enumerate}
Dans chaque cas, $R[x_1, \ldots, x_n] \to A$ est plat en $\mathfrak p$ ; l'ouverture de la platitude (Algèbre, théorème \ref{algebra-theorem-openness-flatness}) permet donc de supposer que $R[x_1, \ldots, x_n] \to A$ est plat (en remplaçant $A$ par une localisation principale convenable). D'après Algèbre, lemme \ref{algebra-lemma-flat-finite-presentation-limit-flat}, il existe $R_0 \subset R$ et $R_0[x_1, \ldots, x_n] \to A_0$ tels que $R_0$ soit de type fini sur $\mathbf{Z}$, que $R_0 \to A_0$ soit de type fini et que $R_0[x_1, \ldots, x_n] \to A_0$ soit plat. Remarquons que l'ensemble des points où un morphisme plat de type fini est de Gorenstein commute au changement de base d'après le lemme \ref{duality-lemma-base-change-gorenstein}. On se ramène ainsi au cas où $R$ est noethérien.

\medskip\noindent
On peut donc supposer $X$ et $S$ affines et disposer d'une factorisation de $f$ de la forme
$$
X \xrightarrow{g} \mathbf{A}^n_S \xrightarrow{p} S
$$
où $g$ est plat et quasi-fini et où $S$ est noethérien. Alors $X$ et $\mathbf{A}^n_S$ sont séparés sur $S$, et l'on a
$$
f^!\mathcal{O}_S = g^!p^!\mathcal{O}_S = g^!\mathcal{O}_{\mathbf{A}^n_S}[n]
$$
d'après les propriétés connues des foncteurs image inverse extraordinaire (lemmes \ref{duality-lemma-upper-shriek-composition} et \ref{duality-lemma-shriek-affine-line}). L'équivalence de (a), (b) et (c) résulte donc du lemme \ref{duality-lemma-affine-flat-Noetherian-gorenstein}.
\end{proof}

\begin{lemma}
\label{duality-lemma-gorenstein-local-source-and-target}
La propriété
$\mathcal{P}(f)=$« les fibres de $f$ sont localement noethériennes et $f$ est de Gorenstein »
est locale pour la topologie fppf sur le but et locale pour la topologie syntomique sur la source.
\end{lemma}

\begin{proof}
On a
$\mathcal{P}(f) =
\mathcal{P}_1(f) \wedge \mathcal{P}_2(f)$
où
$\mathcal{P}_1(f)=$« $f$ est plat » et
$\mathcal{P}_2(f)=$« les fibres de $f$ sont localement noethériennes et de Gorenstein ».
On sait que $\mathcal{P}_1$ est locale pour la topologie fppf sur la source et sur le but, voir Descente, lemmes \ref{descent-lemma-descending-property-flat} et \ref{descent-lemma-flat-fpqc-local-source}. Il reste donc à traiter $\mathcal{P}_2$.

\medskip\noindent
Soit $f : X \to Y$ un morphisme de schémas. Soit $\{\varphi_i : Y_i \to Y\}_{i \in I}$ un recouvrement fppf de $Y$. Notons $f_i : X_i \to Y_i$ le changement de base de $f$ par $\varphi_i$. Soit $i \in I$, et soit $y_i \in Y_i$ un point. Posons $y = \varphi_i(y_i)$. Remarquons que
$$
X_{i, y_i} = \Spec(\kappa(y_i)) \times_{\Spec(\kappa(y))} X_y.
$$
et que $\kappa(y_i)/\kappa(y)$ est une extension de corps de type fini. Ainsi, si $X_y$ est localement noethérien, alors $X_{i, y_i}$ est localement noethérien, voir Variétés, lemme \ref{varieties-lemma-locally-Noetherian-base-change}. Et si, de plus, $X_y$ est de Gorenstein, alors $X_{i, y_i}$ est de Gorenstein, voir le lemme \ref{duality-lemma-gorenstein-base-change}. Par conséquent, $\mathcal{P}_2$ est locale pour la topologie fppf sur le but.

\medskip\noindent
Soit $\{X_i \to X\}$ un recouvrement syntomique de $X$. Soit $y \in Y$. Dans ce cas, $\{X_{i, y} \to X_y\}$ est un recouvrement syntomique de la fibre. La localité de $\mathcal{P}_2$ pour la topologie syntomique sur la source résulte donc du lemme \ref{duality-lemma-gorenstein-local-syntomic}.
\end{proof}




\section{Compléments sur les complexes dualisants}
\label{duality-section-more-dualizing}

\noindent
Quelques lemmes qui ne trouvent pas très bien leur place ailleurs.

\begin{lemma}
\label{duality-lemma-descent-ascent}
Soit $f : X \to Y$ un morphisme de schémas localement noethériens. Supposons que
\begin{enumerate}
\item $f$ soit syntomique et surjectif ; ou que
\item $f$ soit un morphisme surjectif, plat et localement d'intersection complète ; ou que
\item $f$ soit un morphisme de Gorenstein surjectif de type fini.
\end{enumerate}
Alors $K \in D_\QCoh(\mathcal{O}_Y)$ est un complexe dualisant sur $Y$ si et seulement si $Lf^*K$ est un complexe dualisant sur $X$.
\end{lemma}

\begin{proof}
En prenant des ouverts affines et en utilisant Catégories dérivées de schémas, lemme \ref{perfect-lemma-affine-compare-bounded}, cela se traduit par Complexes dualisants, lemme \ref{dualizing-lemma-descent-ascent}.
\end{proof}








\section{Dualité pour les schémas propres sur des corps}
\label{duality-section-duality-proper-over-field}

\noindent
Dans cette section, nous développons les conséquences des résultats très généraux ci-dessus sur les complexes dualisants et la dualité pour les schémas propres sur des corps.

\begin{lemma}
\label{duality-lemma-duality-proper-over-field}
Soit $X$ un schéma propre sur un corps $k$. Il existe un complexe dualisant $\omega_X^\bullet$ possédant les propriétés suivantes :
\begin{enumerate}
\item $H^i(\omega_X^\bullet)$ n'est non nul que pour $i \in [-\dim(X), 0]$ ;
\item $\omega_X = H^{-\dim(X)}(\omega_X^\bullet)$ est un $(S_2)$-module cohérent dont le support est la réunion des composantes irréductibles de dimension $\dim(X)$ ;
\item la dimension du support de $H^i(\omega_X^\bullet)$ est au plus $-i$ ;
\item pour $x \in X$ fermé, le module
$H^i(\omega_{X, x}^\bullet) \oplus \ldots \oplus H^0(\omega_{X, x}^\bullet)$
est non nul si et seulement si $\text{depth}(\mathcal{O}_{X, x}) \leq -i$ ;
\item pour $K \in D_\QCoh(\mathcal{O}_X)$, il existe des isomorphismes fonctoriels\footnote{Cette propriété caractérise $\omega_X^\bullet$ dans $D_\QCoh(\mathcal{O}_X)$ à isomorphisme unique près par le lemme de Yoneda. Puisque $\omega_X^\bullet$ appartient à $D^b_{\textit{Coh}}(\mathcal{O}_X)$, il suffit en fait de considérer $K \in D^b_{\textit{Coh}}(\mathcal{O}_X)$.}
\begingroup\small
$$
\Ext^i_X(K, \omega_X^\bullet) = \Hom_k(H^{-i}(X, K), k)
$$
\endgroup
compatibles aux décalages et aux triangles distingués ;
\item il existe des isomorphismes fonctoriels
\par\noindent\hfil
$\Hom(\mathcal{F}, \omega_X) = \Hom_k(H^{\dim(X)}(X, \mathcal{F}), k)$
\hfil\par\noindent
pour $\mathcal{F}$ quasi-cohérent sur $X$ ;
\item si $X \to \Spec(k)$ est lisse de dimension relative $d$, alors $\omega_X^\bullet \cong \wedge^d\Omega_{X/k}[d]$ et $\omega_X \cong \wedge^d\Omega_{X/k}$.
\end{enumerate}
\end{lemma}

\begin{proof}
Notons $f : X \to \Spec(k)$ le morphisme structural. Soit $a$ l'adjoint à droite de l'image directe par ce morphisme, voir le lemme \ref{duality-lemma-twisted-inverse-image}. Considérons le complexe dualisant relatif
$$
\omega_X^\bullet = a(\mathcal{O}_{\Spec(k)})
$$
Comparer avec la remarque \ref{duality-remark-relative-dualizing-complex}. Puisque $f$ est propre, on a
\par\noindent\hfil
$f^!(\mathcal{O}_{\Spec(k)}) = a(\mathcal{O}_{\Spec(k)})$
\hfil\par\noindent
par définition, voir la section \ref{duality-section-upper-shriek}. En appliquant le lemme \ref{duality-lemma-shriek-dualizing}, on trouve que $\omega_X^\bullet$ est un complexe dualisant. En outre, on voit que $\omega_X^\bullet$ et $\omega_X$ sont comme dans les exemples \ref{duality-example-proper-over-local} et \ref{duality-example-equidimensional-over-field}.

\medskip\noindent
Les assertions (1), (2) et (3) résultent du lemme \ref{duality-lemma-vanishing-good-dualizing}.

\medskip\noindent
Pour un point fermé $x \in X$, on voit que $\omega_{X, x}^\bullet$ est un complexe dualisant normalisé sur $\mathcal{O}_{X, x}$, voir le lemme \ref{duality-lemma-good-dualizing-normalized}. L'assertion (4) résulte alors de Complexes dualisants, lemme \ref{dualizing-lemma-depth-in-terms-dualizing-complex}.

\medskip\noindent
L'assertion (5) vaut par construction, puisque $a$ est l'adjoint à droite de $Rf_* : D_\QCoh(\mathcal{O}_X) \to D(\mathcal{O}_{\Spec(k)}) = D(k)$, que l'on peut identifier à $K \mapsto R\Gamma(X, K)$. On utilise également que la catégorie dérivée $D(k)$ des $k$-modules est la même que la catégorie des $k$-espaces vectoriels gradués.

\medskip\noindent
L'assertion (6) résulte du lemme \ref{duality-lemma-dualizing-module-proper-over-A} pour $\mathcal{F}$ cohérent, et, en général, en explicitant (5) pour $K = \mathcal{F}[0]$ et $i = -\dim(X)$.

\medskip\noindent
L'assertion (7) résulte du lemme \ref{duality-lemma-smooth-proper}.
\end{proof}

\begin{remark}
\label{duality-remark-duality-proper-over-field}
Soient $k$, $X$ et $\omega_X^\bullet$ comme dans le lemme \ref{duality-lemma-duality-proper-over-field}. L'identité du complexe $\omega_X^\bullet$ correspond, via l'isomorphisme fonctoriel de l'assertion (5), à un morphisme
$$
t : H^0(X, \omega_X^\bullet) \longrightarrow k
$$
Pour un $K$ arbitraire dans $D_\QCoh(\mathcal{O}_X)$, l'identification
\par\noindent\hfil
$\Hom(K, \omega_X^\bullet)$ à $H^0(X, K)^\vee$
\hfil\par\noindent
dans l'assertion (5) correspond à l'accouplement
$$
\Hom_X(K, \omega_X^\bullet) \times H^0(X, K) \longrightarrow k,\quad
(\alpha, \beta) \longmapsto t(\alpha(\beta))
$$
Cela résulte de la fonctorialité des isomorphismes de (5). De même, pour tout $i \in \mathbf{Z}$, on obtient l'accouplement
$$
\Ext^i_X(K, \omega_X^\bullet) \times H^{-i}(X, K) \longrightarrow k,\quad
(\alpha, \beta) \longmapsto t(\alpha(\beta))
$$
Ici, on considère $\alpha$ comme un morphisme $K[-i] \to \omega_X^\bullet$ et $\beta$ comme un élément de $H^0(X, K[-i])$ afin de définir $\alpha(\beta)$. Remarquons que si $K$ est général, on sait seulement que cet accouplement est non dégénéré d'un côté : l'accouplement induit un isomorphisme de $\Hom_X(K, \omega_X^\bullet)$, resp.\ $\Ext^i_X(K, \omega_X^\bullet)$, sur le dual $k$-linéaire de $H^0(X, K)$, resp.\ $H^{-i}(X, K)$, mais en général pas réciproquement. Si $K$ appartient à $D^b_{\textit{Coh}}(\mathcal{O}_X)$, alors $\Hom_X(K, \omega_X^\bullet)$, $\Ext_X(K, \omega_X^\bullet)$, $H^0(X, K)$ et $H^i(X, K)$ sont des $k$-espaces vectoriels de dimension finie (d'après Catégories dérivées de schémas, lemmes \ref{perfect-lemma-coherent-internal-hom} et \ref{perfect-lemma-direct-image-coherent-bdd-below}), et les accouplements sont parfaits au sens usuel.
\end{remark}

\begin{remark}
\label{duality-remark-coherent-duality-proper-over-field}
Nous poursuivons la discussion de la remarque \ref{duality-remark-duality-proper-over-field} et conservons les mêmes notations $k$, $X$, $\omega_X^\bullet$ et $t$. Si $\mathcal{F}$ est un $\mathcal{O}_X$-module cohérent, on obtient des accouplements parfaits
$$
\langle -, - \rangle :
\Ext^i_X(\mathcal{F}, \omega_X^\bullet) \times H^{-i}(X,\mathcal{F})
\longrightarrow k,\quad
(\alpha, \beta) \longmapsto t(\alpha(\beta))
$$
d'espaces vectoriels de dimension finie sur $k$. Ces accouplements satisfont la fonctorialité (évidente) suivante : si $\varphi : \mathcal{F} \to \mathcal{G}$ est un homomorphisme de $\mathcal{O}_X$-modules cohérents, alors on a
$$
\langle \alpha \circ \varphi, \beta \rangle =
\langle \alpha, \varphi(\beta) \rangle
$$
pour $\alpha \in \Ext^i_X(\mathcal{G}, \omega_X^\bullet)$ et $\beta \in H^{-i}(X, \mathcal{F})$. Autrement dit, l'application $k$-linéaire $\Ext^i_X(\mathcal{G}, \omega_X^\bullet) \to
\Ext^i_X(\mathcal{F}, \omega_X^\bullet)$ induite par $\varphi$ est, via les accouplements, le dual $k$-linéaire de l'application $k$-linéaire $H^{-i}(X, \mathcal{F}) \to H^{-i}(X, \mathcal{G})$ induite par $\varphi$. Formulé ainsi, cela reste valable si $\varphi$ est un homomorphisme de $\mathcal{O}_X$-modules quasi-cohérents.
\end{remark}

\begin{lemma}
\label{duality-lemma-duality-proper-over-field-perfect}
Soient $k$, $X$ et $\omega_X^\bullet$ comme dans le lemme \ref{duality-lemma-duality-proper-over-field}. Soit $t : H^0(X, \omega_X^\bullet) \to k$ comme dans la remarque \ref{duality-remark-duality-proper-over-field}. Supposons $E \in D(\mathcal{O}_X)$ parfait. Alors les accouplements
$$
H^i(X, \omega_X^\bullet \otimes_{\mathcal{O}_X}^\mathbf{L} E^\vee)
\times
H^{-i}(X, E) \longrightarrow k, \quad
(\xi, \eta) \longmapsto
t((1_{\omega_X^\bullet} \otimes \epsilon)(\xi \cup \eta))
$$
sont parfaits pour tout $i$. Ici, $\cup$ désigne le cup-produit de Cohomologie, section \ref{cohomology-section-cup-product}, et $\epsilon : E^\vee \otimes_{\mathcal{O}_X}^\mathbf{L} E \to \mathcal{O}_X$ est comme dans Cohomologie, exemple \ref{cohomology-example-dual-derived}.
\end{lemma}

\begin{proof}
En remplaçant $E$ par $E[-i]$, on se ramène au cas $i = 0$. D'après Cohomologie, lemme \ref{cohomology-lemma-ext-composition-is-cup}, l'accouplement est le même que celui étudié dans la remarque \ref{duality-remark-duality-proper-over-field}, d'où le résultat par la discussion de cette remarque.
\end{proof}

\begin{lemma}
\label{duality-lemma-duality-proper-over-field-CM}
Soit $X$ un schéma propre sur un corps $k$, de Cohen–Macaulay et équidimensionnel de dimension $d$. Le module $\omega_X$ du lemme \ref{duality-lemma-duality-proper-over-field} possède les propriétés suivantes :
\begin{enumerate}
\item $\omega_X$ est un module dualisant sur $X$ (section \ref{duality-section-dualizing-module}) ;
\item $\omega_X$ est un module cohérent de Cohen–Macaulay dont le support est $X$ ;
\item il existe des isomorphismes fonctoriels
\par\noindent\hfil
$\Ext^i_X(K, \omega_X[d]) = \Hom_k(H^{-i}(X, K), k)$
\hfil\par\noindent
compatibles aux décalages et aux triangles distingués, pour $K \in D_\QCoh(X)$ ;
\item il existe des isomorphismes fonctoriels
$\Ext^{d - i}(\mathcal{F}, \omega_X) = \Hom_k(H^i(X, \mathcal{F}), k)$
pour $\mathcal{F}$ quasi-cohérent sur $X$.
\end{enumerate}
\end{lemma}

\begin{proof}
Il résulte clairement du lemme \ref{duality-lemma-duality-proper-over-field} que $\omega_X$ est un module dualisant (puisqu'il s'agit du faisceau de cohomologie non nul le plus à gauche d'un complexe dualisant). On a $\omega_X^\bullet = \omega_X[d]$, et $\omega_X$ est de Cohen–Macaulay puisque $X$ est de Cohen–Macaulay, voir le lemme \ref{duality-lemma-dualizing-module-CM-scheme}. Les autres assertions en résultent, combinées aux assertions correspondantes du lemme \ref{duality-lemma-duality-proper-over-field}.
\end{proof}

\begin{remark}
\label{duality-remark-rework-duality-locally-free-CM}
Soit $X$ un schéma propre de Cohen–Macaulay sur un corps $k$, équidimensionnel de dimension $d$. Soient $\omega_X^\bullet$ et $\omega_X$ comme dans le lemme \ref{duality-lemma-duality-proper-over-field}. D'après le lemme \ref{duality-lemma-duality-proper-over-field-CM}, on a $\omega_X^\bullet = \omega_X[d]$. Soit $t : H^d(X, \omega_X) \to k$ le morphisme de la remarque \ref{duality-remark-duality-proper-over-field}. Soit $\mathcal{E}$ un $\mathcal{O}_X$-module localement libre de type fini, de dual $\mathcal{E}^\vee$. On a alors des accouplements parfaits
$$
H^i(X, \omega_X \otimes_{\mathcal{O}_X} \mathcal{E}^\vee)
\times
H^{d - i}(X, \mathcal{E})
\longrightarrow
k,\quad
(\xi, \eta) \longmapsto t(1 \otimes \epsilon)(\xi \cup \eta))
$$
où $\cup$ est le cup-produit et $\epsilon : \mathcal{E}^\vee \otimes_{\mathcal{O}_X} \mathcal{E}
\to \mathcal{O}_X$ le morphisme d'évaluation. C'est un cas particulier du lemme \ref{duality-lemma-duality-proper-over-field-perfect}.
\end{remark}

\noindent
Voici une vérification de cohérence pour le complexe dualisant.

\begin{lemma}
\label{duality-lemma-sanity-check-duality}
Soit $X$ un schéma propre sur un corps $k$. Soient $\omega_X^\bullet$ et $\omega_X$ comme dans le lemme \ref{duality-lemma-duality-proper-over-field}.
\begin{enumerate}
\item Si $X \to \Spec(k)$ se factorise sous la forme $X \to \Spec(k') \to \Spec(k)$ pour un corps $k'$, alors $\omega_X^\bullet$ et $\omega_X$ sont comme dans le lemme \ref{duality-lemma-duality-proper-over-field} pour le morphisme $X \to \Spec(k')$.
\item Si $K/k$ est une extension de corps, alors les images inverses de $\omega_X^\bullet$ et $\omega_X$ sur le changement de base $X_K$ sont comme dans le lemme \ref{duality-lemma-duality-proper-over-field} pour le morphisme $X_K \to \Spec(K)$.
\end{enumerate}
\end{lemma}

\begin{proof}
Notons $f : X \to \Spec(k)$ le morphisme structural et $f' : X \to \Spec(k')$ la factorisation donnée. Dans la démonstration du lemme \ref{duality-lemma-duality-proper-over-field}, on a pris $\omega_X^\bullet = a(\mathcal{O}_{\Spec(k)})$, où $a$ est l'adjoint à droite du lemme \ref{duality-lemma-twisted-inverse-image} pour $f$. Il faut donc montrer
$a(\mathcal{O}_{\Spec(k)}) \cong a'(\mathcal{O}_{\Spec(k)})$
où $a'$ est l'adjoint à droite du lemme \ref{duality-lemma-twisted-inverse-image} pour $f'$. Puisque $k' \subset H^0(X, \mathcal{O}_X)$, on voit que $k'/k$ est une extension finie (Cohomologie des schémas, lemme \ref{coherent-lemma-proper-over-affine-cohomology-finite}). Par unicité des adjoints, on a $a = a' \circ b$, où $b$ est l'adjoint à droite du lemme \ref{duality-lemma-twisted-inverse-image} pour $g : \Spec(k') \to \Spec(k)$. On peut aussi l'exprimer ainsi : on a $f^! = (f')^! \circ g^!$. Il suffit donc de montrer que $\Hom_k(k', k) \cong k'$ comme $k'$-modules, voir l'exemple \ref{duality-example-affine-twisted-inverse-image}. Cela vaut parce que ce sont des $k'$-espaces vectoriels de même dimension (à savoir de dimension $1$).

\medskip\noindent
Démonstration de (2). Cela vaut parce que l'on dispose du changement de base pour $a$ d'après le lemme \ref{duality-lemma-more-base-change}. Voir la discussion de la remarque \ref{duality-remark-relative-dualizing-complex}.
\end{proof}








\section{Complexes dualisants relatifs}
\label{duality-section-relative-dualizing-complexes}

\noindent
Pour un morphisme propre, plat et de présentation finie, on dispose d'un complexe dualisant relatif rigide, voir la remarque \ref{duality-remark-relative-dualizing-complex} et le lemme \ref{duality-lemma-van-den-bergh}. Pour un morphisme séparé et de type fini $f : X \to Y$ de schémas noethériens, on peut considérer $f^!\mathcal{O}_Y$. Dans cette section, nous définissons les complexes dualisants relatifs pour les morphismes plats et localement de présentation finie (mais pas nécessairement quasi-séparés ni quasi-compacts) entre schémas (pas nécessairement localement noethériens). Nous montrons que ces complexes existent, sont uniques à isomorphisme unique près et coïncident avec ceux des cas mentionnés ci-dessus. Avant de lire cette section, prière de lire Complexes dualisants, section \ref{dualizing-section-relative-dualizing-complexes}.

\begin{definition}
\label{duality-definition-relative-dualizing-complex}
Soit $X \to S$ un morphisme de schémas plat et localement de présentation finie. Soit $W \subset X \times_S X$ un ouvert quelconque tel que la diagonale $\Delta_{X/S} : X \to X \times_S X$ se factorise par une immersion fermée $\Delta : X \to W$. Un {\it complexe dualisant relatif} est un couple $(K, \xi)$ formé d'un objet $K \in D(\mathcal{O}_X)$ et d'un morphisme
$$
\xi : \Delta_*\mathcal{O}_X \longrightarrow L\text{pr}_1^*K|_W
$$
dans $D(\mathcal{O}_W)$ tels que
\begin{enumerate}
\item $K$ soit $S$-parfait (Catégories dérivées de schémas, définition \ref{perfect-definition-relatively-perfect}), et que
\item $\xi$ définisse un isomorphisme de $\Delta_*\mathcal{O}_X$ sur
$R\SheafHom_{\mathcal{O}_W}(
\Delta_*\mathcal{O}_X, L\text{pr}_1^*K|_W)$.
\end{enumerate}
\end{definition}

\noindent
\begingroup\emergencystretch=4em
D'après le lemme \ref{duality-lemma-sheaf-with-exact-support-ext}, la condition (2) équivaut à l'existence d'un isomorphisme
$$
\mathcal{O}_X \longrightarrow
R\SheafHom(\mathcal{O}_X, L\text{pr}_1^*K|_W)
$$
dans $D(\mathcal{O}_X)$ dont l'image directe par $\Delta$ est égale à $\xi$. Puisque
\par\noindent\hfil
$R\SheafHom(\mathcal{O}_X, L\text{pr}_1^*K|_W)$
\hfil\par\noindent
est indépendant du choix de l'ouvert $W$, il en est de même de la catégorie des couples $(K, \xi)$. Si $X \to S$ est séparé, on peut choisir $W = X \times_S X$. Nous ramènerons nombre des arguments au cas des anneaux au moyen du lemme suivant.\par\endgroup

\begin{lemma}
\label{duality-lemma-relative-dualizing-complex-algebra}
Soit $X \to S$ un morphisme de schémas plat et localement de présentation finie. Soit $(K, \xi)$ un complexe dualisant relatif. Alors, pour tout diagramme commutatif
$$
\xymatrix{
\Spec(A) \ar[d] \ar[r] & X \ar[d] \\
\Spec(R) \ar[r] & S
}
$$
dont les flèches horizontales sont des immersions ouvertes, la restriction de $K$ à $\Spec(A)$ correspond, via Catégories dérivées de schémas, lemme \ref{perfect-lemma-affine-compare-bounded}, à un complexe dualisant relatif pour $R \to A$ au sens de Complexes dualisants, définition \ref{dualizing-definition-relative-dualizing-complex}.
\end{lemma}

\begin{proof}
Puisque la formation de $R\SheafHom$ commute aux restrictions aux ouverts, on peut aussi bien supposer $X = \Spec(A)$ et $S = \Spec(R)$. Remarquons que les objets relativement parfaits de $D(\mathcal{O}_X)$ sont pseudo-cohérents et appartiennent donc à $D_\QCoh(\mathcal{O}_X)$ (Catégories dérivées de schémas, lemme \ref{perfect-lemma-pseudo-coherent}). L'énoncé a donc un sens. Remarquons que les constructions de $\Delta_*$, $L\text{pr}_1^*$ et $R\SheafHom$ sont compatibles à ce qui se passe du côté algébrique d'après Catégories dérivées de schémas, lemmes \ref{perfect-lemma-quasi-coherence-pushforward}, \ref{perfect-lemma-quasi-coherence-pullback} et \ref{perfect-lemma-quasi-coherence-internal-hom}. Pour le dernier, on remarque que $L\text{pr}_1^*K$ est $S$-parfait (donc borné inférieurement) et que $\Delta_*\mathcal{O}_X$ est un objet pseudo-cohérent de $D(\mathcal{O}_W)$ ; traduit en algèbre, cela signifie que $A$ est pseudo-cohérent comme $A \otimes_R A$-module, ce qui résulte de Compléments sur l'algèbre, lemme \ref{more-algebra-lemma-more-relative-pseudo-coherent-is-moot}, appliqué à $R \to A \otimes_R A \to A$. On retrouve ainsi exactement les conditions de Complexes dualisants, définition \ref{dualizing-definition-relative-dualizing-complex}.
\end{proof}

\begin{lemma}
\label{duality-lemma-relative-dualizing-RHom}
Soit $X \to S$ un morphisme de schémas plat et localement de présentation finie. Soit $(K, \xi)$ un complexe dualisant relatif. Alors
$\mathcal{O}_X \to R\SheafHom_{\mathcal{O}_X}(K, K)$
est un isomorphisme.
\end{lemma}

\begin{proof}
En raisonnant localement sur les ouverts affines, le lemme \ref{duality-lemma-relative-dualizing-complex-algebra} ramène l'énoncé au cas algébrique, qui est Complexes dualisants, lemme \ref{dualizing-lemma-relative-dualizing-RHom}.
\end{proof}

\begin{lemma}
\label{duality-lemma-uniqueness-relative-dualizing}
Soit $X \to S$ un morphisme de schémas plat et localement de présentation finie. Si $(K, \xi)$ et $(L, \eta)$ sont deux complexes dualisants relatifs sur $X/S$, il existe un unique isomorphisme $K \to L$ qui envoie $\xi$ sur $\eta$.
\end{lemma}

\begin{proof}
Soit $U \subset X$ un ouvert affine dont l'image est contenue dans un ouvert affine de $S$. Il existe alors un isomorphisme $K|_U \to L|_U$ d'après le lemme \ref{duality-lemma-relative-dualizing-complex-algebra} et Complexes dualisants, lemme \ref{dualizing-lemma-uniqueness-relative-dualizing}. Le lecteur peut reprendre l'argument de ce lemme dans le cas des schémas pour obtenir une démonstration dans ce cas-ci. Nous utiliserons plutôt un argument de recollement.

\medskip\noindent
Supposons donné un isomorphisme $\alpha : K \to L$. Alors $\alpha(\xi) = u \eta$ pour une section inversible $u \in H^0(W, \Delta_*\mathcal{O}_X) = H^0(X, \mathcal{O}_X)$. (En effet, $\eta$ et $\alpha(\xi)$ sont tous deux des générateurs d'un $\Delta_*\mathcal{O}_X$-module inversible par hypothèse.) Après avoir remplacé $\alpha$ par $u^{-1}\alpha$, on voit donc que $\alpha(\xi) = \eta$. Puisque le groupe des automorphismes de $K$ est $H^0(X, \mathcal{O}_X^*)$ d'après le lemme \ref{duality-lemma-relative-dualizing-RHom}, il existe au plus un tel $\alpha$.

\medskip\noindent
Soit $\mathcal{B}$ la collection des ouverts affines de $X$ dont l'image est contenue dans un ouvert affine de $S$. Pour chaque $U \in \mathcal{B}$, les deux paragraphes précédents fournissent un unique isomorphisme $\alpha_U : K|_U \to L|_U$ envoyant $\xi$ sur $\eta$. Remarquons que $\text{Ext}^i(K|_U, K|_U) = 0$ pour $i < 0$ et pour tout ouvert $U$ de $X$ d'après le lemme \ref{duality-lemma-relative-dualizing-RHom}. En appliquant Cohomologie, lemme \ref{cohomology-lemma-vanishing-and-glueing}, à $\text{id} : X \to X$, on obtient un unique morphisme $\alpha : K \to L$ qui coïncide avec $\alpha_U$ pour tout $U \in \mathcal{B}$. Alors $\alpha$ envoie $\xi$ sur $\eta$ puisque c'est vrai localement.
\end{proof}

\begin{lemma}
\label{duality-lemma-existence-relative-dualizing}
Soit $X \to S$ un morphisme de schémas plat et localement de présentation finie. Il existe un complexe dualisant relatif $(K, \xi)$.
\end{lemma}

\begin{proof}
Soit $\mathcal{B}$ la collection des ouverts affines de $X$ dont l'image est contenue dans un ouvert affine de $S$. Pour chaque $U$, on dispose d'un complexe dualisant relatif $(K_U, \xi_U)$ pour $U$ sur $S$. En effet, choisissons un ouvert affine $V \subset S$ tel que $U \to X \to S$ se factorise par $V$. Écrivons $U = \Spec(A)$ et $V = \Spec(R)$. D'après Complexes dualisants, lemme \ref{dualizing-lemma-base-change-relative-dualizing}, il existe un complexe dualisant relatif $K_A \in D(A)$ pour $R \to A$. En remontant l'argument de la démonstration du lemme \ref{duality-lemma-relative-dualizing-complex-algebra}, celui-ci détermine un objet $V$-parfait $K_U \in D(\mathcal{O}_U)$ et un morphisme
$$
\xi : \Delta_*\mathcal{O}_U \to L\text{pr}_1^*K_U
$$
dans $D(\mathcal{O}_{U \times_V U})$. Comme être $V$-parfait équivaut à être $S$-parfait et comme $U \times_V U = U \times_S U$, on trouve que $(K_U, \xi_U)$ convient.

\medskip\noindent
Si $U' \subset U \subset X$ avec $U', U \in \mathcal{B}$, alors il existe un unique isomorphisme $\rho_{U'}^U : K_U|_{U'} \to K_{U'}$ dans $D(\mathcal{O}_{U'})$ envoyant $\xi_U|_{U' \times_S U'}$ sur $\xi_{U'}$ d'après le lemme \ref{duality-lemma-uniqueness-relative-dualizing} (remarquons que la restriction d'un complexe dualisant relatif à un ouvert est trivialement un complexe dualisant relatif). L'unicité garantit que
$\rho^U_{U''} = \rho^V_{U''} \circ \rho ^U_{U'}|_{U''}$
pour $U'' \subset U' \subset U$ dans $\mathcal{B}$. Remarquons que $\text{Ext}^i(K_U, K_U) = 0$ pour $i < 0$, pour $U \in \mathcal{B}$, d'après le lemme \ref{duality-lemma-relative-dualizing-RHom} appliqué à $U/S$ et $K_U$. Le lemme de recollement BBD (Cohomologie, théorème \ref{cohomology-theorem-glueing-bbd-general}) affirme donc qu'il existe une unique solution, à savoir un objet $K \in D(\mathcal{O}_X)$ et des isomorphismes $\rho_U : K|_U \to K_U$ tels que l'on ait $\rho^U_{U'} \circ \rho_U|_{U'} = \rho_{U'}$ pour tout $U' \subset U$, $U, U' \in \mathcal{B}$.

\medskip\noindent
Pour achever la démonstration, il faut construire le morphisme
$$
\xi : \Delta_*\mathcal{O}_X \longrightarrow L\text{pr}_1^*K|_W
$$
dans $D(\mathcal{O}_W)$ induisant un isomorphisme de $\Delta_*\mathcal{O}_X$ sur
\par\noindent\hfil
$R\SheafHom_{\mathcal{O}_W}(\Delta_*\mathcal{O}_X, L\text{pr}_1^*K|_W)$.
\hfil\par\noindent
Puisque l'on peut changer $W$, choisissons $W = \bigcup_{U \in \mathcal{B}} U \times_S U$. On peut utiliser $\rho_U$ pour obtenir des isomorphismes
$$
R\SheafHom_{\mathcal{O}_W}(
\Delta_*\mathcal{O}_X, L\text{pr}_1^*K|_W)|_{U \times_S U}
\xrightarrow{\rho_U}
R\SheafHom_{\mathcal{O}_{U \times_S U}}(
\Delta_*\mathcal{O}_U, L\text{pr}_1^*K_U)
$$
Comme $W$ est recouvert par les ouverts $U \times_S U$, on conclut que les faisceaux de cohomologie de
$R\SheafHom_{\mathcal{O}_W}(\Delta_*\mathcal{O}_X, L\text{pr}_1^*K|_W)$
sont nuls sauf en degré $0$. En outre, on obtient des isomorphismes
\begingroup\small
$$
H^0\left(U \times_S U, R\SheafHom_{\mathcal{O}_W}(\Delta_*\mathcal{O}_X,
L\text{pr}_1^*K|_W)\right)
\xrightarrow{\rho_U}
H^0\left((R\SheafHom_{\mathcal{O}_{U \times_S U}}(
\Delta_*\mathcal{O}_U, L\text{pr}_1^*K_U)\right)
$$
\endgroup
Soit $\tau_U$ dans le membre de gauche un élément envoyé sur $\xi_U$ par ce morphisme. Les compatibilités entre $\rho^U_{U'}$, $\xi_U$, $\xi_{U'}$, $\rho_U$ et $\rho_{U'}$ pour $U' \subset U \subset X$ ouvert de $U', U \in \mathcal{B}$ impliquent que $\tau_U|_{U' \times_S U'} = \tau_{U'}$. On obtient ainsi une section globale $\tau$ du $0$-ième faisceau de cohomologie $H^0(R\SheafHom_{\mathcal{O}_W}(\Delta_*\mathcal{O}_X, L\text{pr}_1^*K|_W))$. Puisque les autres faisceaux de cohomologie de
$R\SheafHom_{\mathcal{O}_W}(\Delta_*\mathcal{O}_X, L\text{pr}_1^*K|_W)$
sont nuls, cette section globale $\tau$ détermine comme voulu un morphisme $\xi$. Puisque la restriction de $\xi$ à $U \times_S U$ donne $\xi_U$, on voit qu'il satisfait la condition finale de la définition \ref{duality-definition-relative-dualizing-complex}.
\end{proof}

\begin{lemma}
\label{duality-lemma-base-change-relative-dualizing}
Considérons un carré cartésien
$$
\xymatrix{
X' \ar[d]_{f'} \ar[r]_{g'} & X \ar[d]^f \\
S' \ar[r]^g & S
}
$$
de schémas. Supposons $X \to S$ plat et localement de présentation finie. Soit $(K, \xi)$ un complexe dualisant relatif pour $f$. Posons $K' = L(g')^*K$. Soit $\xi'$ le changement de base dérivé de $\xi$ (voir la démonstration). Alors $(K', \xi')$ est un complexe dualisant relatif pour $f'$.
\end{lemma}

\begin{proof}
Considérons le carré cartésien
$$
\xymatrix{
X' \ar[d]_{\Delta_{X'/S'}} \ar[r] & X \ar[d]^{\Delta_{X/S}} \\
X' \times_{S'} X' \ar[r]^{g' \times g'} & X \times_S X
}
$$
Choisissons un ouvert $W \subset X \times_S X$ tel que $\Delta_{X/S}$ se factorise par une immersion fermée $\Delta : X \to W$. Choisissons un ouvert $W' \subset X' \times_{S'} X'$ tel que $\Delta_{X'/S'}$ se factorise par une immersion fermée $\Delta' : X \to W'$ et tel que $(g' \times g')(W') \subset W$. Notons encore $g' \times g' : W' \to W$ le morphisme induit. On a
$$
L(g' \times g')^*\Delta_*\mathcal{O}_X =
\Delta'_*\mathcal{O}_{X'}
\quad\text{et}\quad
L(g' \times g')^*L\text{pr}_1^*K|_W =
L\text{pr}_1^*K'|_{W'}
$$
La première égalité vaut parce que $X$ et $X' \times_{S'} X'$ sont Tor-indépendants sur $X \times_S X$ (voir par exemple Compléments sur les morphismes, lemme \ref{more-morphisms-lemma-case-of-tor-independence}). La seconde résulte de la transitivité de l'image inverse dérivée (Cohomologie, lemme \ref{cohomology-lemma-derived-pullback-composition}). Ainsi, $\xi' = L(g' \times g')^*\xi$ peut être considéré comme un morphisme
$$
\xi' : \Delta'_*\mathcal{O}_{X'} \longrightarrow L\text{pr}_1^*K'|_{W'}
$$
Cela étant dit, la démonstration du lemme est immédiate. Premièrement, $K'$ est $S'$-parfait d'après Catégories dérivées de schémas, lemme \ref{perfect-lemma-base-change-relatively-perfect}. Pour vérifier que $\xi'$ induit un isomorphisme de $\Delta'_*\mathcal{O}_{X'}$ sur
$R\SheafHom_{\mathcal{O}_{W'}}(
\Delta'_*\mathcal{O}_{X'}, L\text{pr}_1^*K'|_{W'})$
on peut travailler localement sur les ouverts affines. Le lemme \ref{duality-lemma-relative-dualizing-complex-algebra} ramène au résultat correspondant en algèbre, démontré dans Complexes dualisants, lemme \ref{dualizing-lemma-base-change-relative-dualizing}.
\end{proof}

\begin{lemma}
\label{duality-lemma-flat-proper-relative-dualizing}
Soit $S$ un schéma quasi-compact et quasi-séparé. Soit $f : X \to S$ un morphisme propre et plat de présentation finie. Le complexe dualisant relatif $\omega_{X/S}^\bullet$ de la remarque \ref{duality-remark-relative-dualizing-complex}, muni de (\ref{duality-equation-pre-rigid}), est un complexe dualisant relatif au sens de la définition \ref{duality-definition-relative-dualizing-complex}.
\end{lemma}

\begin{proof}
Dans le lemme \ref{duality-lemma-properties-relative-dualizing}, on a démontré que $\omega_{X/S}^\bullet$ est $S$-parfait. Soit $c$ l'adjoint à droite du lemme \ref{duality-lemma-twisted-inverse-image} pour la diagonale $\Delta : X \to X \times_S X$. On peut alors appliquer $\Delta_*$ à (\ref{duality-equation-pre-rigid}) pour obtenir un isomorphisme
$$
\Delta_*\mathcal{O}_X \to
\Delta_*(c(L\text{pr}_1^*\omega_{X/S}^\bullet)) =
R\SheafHom_{\mathcal{O}_{X \times_S X}}(
\Delta_*\mathcal{O}_X, L\text{pr}_1^*\omega_{X/S}^\bullet)
$$
L'égalité résulte des lemmes \ref{duality-lemma-twisted-inverse-image-closed} et \ref{duality-lemma-sheaf-with-exact-support-ext}. Cela achève la démonstration.
\end{proof}

\begin{remark}
\label{duality-remark-relative-dualizing-complex-bis}
Soit $X \to S$ un morphisme de schémas plat, propre et de présentation finie. D'après le lemme \ref{duality-lemma-existence-relative-dualizing}, il existe un complexe dualisant relatif $(\omega_{X/S}^\bullet, \xi)$ au sens de la définition \ref{duality-definition-relative-dualizing-complex}. Considérons un morphisme quelconque $g : S' \to S$, où $S'$ est quasi-compact et quasi-séparé (par exemple, un ouvert affine de $S$). Le lemme \ref{duality-lemma-base-change-relative-dualizing} montre que $(L(g')^*\omega_{X/S}^\bullet, L(g')^*\xi)$ est un complexe dualisant relatif pour le changement de base $f' : X' \to S'$ au sens de la définition \ref{duality-definition-relative-dualizing-complex}. Soit $\omega_{X'/S'}^\bullet$ le complexe dualisant relatif pour $X' \to S'$ au sens de la remarque \ref{duality-remark-relative-dualizing-complex}. En combinant les lemmes \ref{duality-lemma-flat-proper-relative-dualizing} et \ref{duality-lemma-uniqueness-relative-dualizing}, on voit qu'il existe un unique isomorphisme
$$
\omega_{X'/S'}^\bullet \longrightarrow L(g')^*\omega_{X/S}^\bullet
$$
compatible à (\ref{duality-equation-pre-rigid}) et à $L(g')^*\xi$. Ces isomorphismes sont compatibles aux morphismes entre schémas quasi-compacts et quasi-séparés sur $S$ et aux isomorphismes de changement de base du lemme \ref{duality-lemma-proper-flat-base-change} (si nous avons un jour besoin de cette compatibilité, nous l'énoncerons et la démontrerons soigneusement ici).
\end{remark}

\begin{lemma}
\label{duality-lemma-compactifyable-relative-dualizing}
Dans la situation \ref{duality-situation-shriek}, soit $f : X \to Y$ un morphisme de $\textit{FTS}_S$. Si $f$ est plat, alors $f^!\mathcal{O}_Y$ est (la première composante d')un complexe dualisant relatif pour $X$ sur $Y$ au sens de la définition \ref{duality-definition-relative-dualizing-complex}.
\end{lemma}

\begin{proof}
D'après le lemme \ref{duality-lemma-flat-shriek-relatively-perfect}, on sait que $f^!\mathcal{O}_Y$ est $Y$-parfait. Comme $f$ est séparé, la diagonale $\Delta : X \to X \times_Y X$ est une immersion fermée et $\Delta_*\Delta^!(-) =
R\SheafHom_{\mathcal{O}_{X \times_Y X}}(\mathcal{O}_X, -)$, voir les lemmes \ref{duality-lemma-twisted-inverse-image-closed} et \ref{duality-lemma-sheaf-with-exact-support-ext}. Pour achever la démonstration, il suffit donc de montrer
$\Delta^!(L\text{pr}_1^*f^!(\mathcal{O}_Y)) \cong \mathcal{O}_X$
où $\text{pr}_1 : X \times_Y X \to X$ est la première projection. On a
$$
\mathcal{O}_X = \Delta^! \text{pr}_1^!\mathcal{O}_X =
\Delta^! \text{pr}_1^! L\text{pr}_2^*\mathcal{O}_Y =
\Delta^!(L\text{pr}_1^* f^!\mathcal{O}_Y)
$$
où $\text{pr}_2 : X \times_Y X \to X$ est la seconde projection et où l'on a utilisé l'isomorphisme de changement de base $\text{pr}_1^! \circ L\text{pr}_2^* = L\text{pr}_1^* \circ f^!$ du lemme \ref{duality-lemma-base-change-shriek-flat}.
\end{proof}

\begin{lemma}
\label{duality-lemma-relative-dualizing-composition}
Soient $f : Y \to X$ et $X \to S$ des morphismes de schémas plats et de présentation finie. Soient $(K, \xi)$ et $(M, \eta)$ des complexes dualisants relatifs pour $X \to S$ et $Y \to X$. Posons $E = M \otimes_{\mathcal{O}_Y}^\mathbf{L} Lf^*K$. Alors $(E, \zeta)$ est un complexe dualisant relatif pour $Y \to S$ pour un $\zeta$ convenable.
\end{lemma}

\begin{proof}
En utilisant le lemme \ref{duality-lemma-relative-dualizing-complex-algebra} et la version algébrique du présent lemme (Complexes dualisants, lemme \ref{dualizing-lemma-relative-dualizing-composition}), on voit que $E$ est, localement sur les ouverts affines, la première composante d'un complexe dualisant relatif. En particulier, on voit que $E$ est $S$-parfait, puisque cela se vérifie localement sur les ouverts affines, voir Catégories dérivées de schémas, lemme \ref{perfect-lemma-affine-locally-rel-perfect}.

\medskip\noindent
Démontrons d'abord l'existence de $\zeta$ lorsque les morphismes $X \to S$ et $Y \to X$ sont séparés, de sorte que $\Delta_{X/S}$, $\Delta_{Y/X}$ et $\Delta_{Y/S}$ sont des immersions fermées. Considérons le diagramme suivant
$$
\xymatrix{
& & Y \ar@{=}[r] & Y \ar[d]^f \\
Y \ar[r]_{\Delta_{Y/X}} &
Y \times_X Y \ar[d]_m \ar[r]_\delta \ar[ru]_q &
Y \times_S Y \ar[d]^{f \times f} \ar[ru]_p & X\\
& X \ar[r]^{\Delta_{X/S}} & X \times_S X \ar[ru]_r
}
$$
où $p$, $q$ et $r$ sont les premières projections. D'après le lemme \ref{duality-lemma-sheaf-with-exact-support-internal-home}, on a
\begingroup\footnotesize
$$
R\SheafHom_{\mathcal{O}_{Y \times_S Y}}(
\Delta_{Y/S, *}\mathcal{O}_Y, Lp^*E) =
R\delta_*\left(R\SheafHom_{\mathcal{O}_{Y \times_X Y}}(
\Delta_{Y/X, *}\mathcal{O}_Y,
R\SheafHom(\mathcal{O}_{Y \times_X Y}, Lp^*E))\right)
$$
\endgroup
D'après le lemme \ref{duality-lemma-sheaf-with-exact-support-tensor}, on a
$$
R\SheafHom(\mathcal{O}_{Y \times_X Y}, Lp^*E) =
R\SheafHom(\mathcal{O}_{Y \times_X Y}, L(f \times f)^*Lr^*K)
\otimes_{\mathcal{O}_{Y \times_S Y}}^\mathbf{L} Lq^*M
$$
D'après le lemme \ref{duality-lemma-flat-bc-sheaf-with-exact-support}, on a
$$
R\SheafHom(\mathcal{O}_{Y \times_X Y}, L(f \times f)^*Lr^*K) =
Lm^*R\SheafHom(\mathcal{O}_X, Lr^*K)
$$
La dernière expression est isomorphe (via $\xi$) à $Lm^*\mathcal{O}_X = \mathcal{O}_{Y \times_X Y}$. Par conséquent, l'expression qui la précède est isomorphe à $Lq^*M$. Ainsi,
$$
R\SheafHom_{\mathcal{O}_{Y \times_S Y}}(
\Delta_{Y/S, *}\mathcal{O}_Y, Lp^*E) =
R\delta_*\left(R\SheafHom_{\mathcal{O}_{Y \times_X Y}}(
\Delta_{Y/X, *}\mathcal{O}_Y, Lq^*M)\right)
$$
Le terme entre parenthèses est isomorphe à $\Delta_{Y/X, *}*\mathcal{O}_X$ via $\eta$. Cela achève la démonstration dans le cas séparé.

\medskip\noindent
Dans le cas général, on choisit un ouvert $W \subset X \times_S X$ tel que $\Delta_{X/S}$ se factorise par une immersion fermée $\Delta : X \to W$, et un ouvert $V \subset Y \times_X Y$ tel que $\Delta_{Y/X}$ se factorise par une immersion fermée $\Delta' : Y \to V$. Enfin, choisissons un ouvert $W' \subset Y \times_S Y$ dont l'intersection avec $Y \times_X Y$ donne $V$ et dont l'image est contenue dans $W$. On considère alors le diagramme
$$
\xymatrix{
& & Y \ar@{=}[r] & Y \ar[d]^f \\
Y \ar[r]_{\Delta'} &
V \ar[d]_m \ar[r]_\delta \ar[ru]_q &
W' \ar[d]^{f \times f} \ar[ru]_p & X\\
& X \ar[r]^\Delta & W \ar[ru]_r
}
$$
et l'on utilise exactement le même argument que précédemment.
\end{proof}






\section{La classe fondamentale d'un morphisme localement d'intersection complète}
\label{duality-section-fundamental-class}

\noindent
Dans cette section, nous utiliserons les calculs effectués dans la section \ref{duality-section-examples}. Notre résultat souffrira donc du même type de défaut d'unicité que celui rencontré dans cette section.

\begin{lemma}
\label{duality-lemma-determinant}
Soit $X$ un espace localement annelé. Soit
$$
\mathcal{E}_1 \xrightarrow{\alpha} \mathcal{E}_0 \to \mathcal{F} \to 0
$$
une suite exacte courte de $\mathcal{O}_X$-modules. Supposons $\mathcal{E}_1$ et $\mathcal{E}_0$ localement libres de rangs $r_1, r_0$. Il existe alors un morphisme canonique
$$
\wedge^{r_0 - r_1}\mathcal{F}
\longrightarrow
\wedge^{r_1}(\mathcal{E}_1^\vee) \otimes \wedge^{r_0}\mathcal{E}_0
$$
qui est un isomorphisme sur la fibre en $x \in X$ si et seulement si $\mathcal{F}$ est localement libre de rang $r_0 - r_1$ dans un voisinage ouvert de $x$.
\end{lemma}

\begin{proof}
Si $r_1 > r_0$, alors $\wedge^{r_0 - r_1}\mathcal{F} = 0$ par convention, et l'unique morphisme ne peut être un isomorphisme. On peut donc supposer $r = r_0 - r_1 \geq 0$. Définissons le morphisme par la formule
$$
s_1 \wedge \ldots \wedge s_r \mapsto
t_1^\vee \wedge \ldots \wedge t_{r_1}^\vee \otimes
\alpha(t_1) \wedge \ldots \wedge \alpha(t_{r_1}) \wedge
\tilde s_1 \wedge \ldots \wedge \tilde s_r
$$
où $t_1, \ldots, t_{r_1}$ est une base locale de $\mathcal{E}_1$, où, en conséquence, $t_1^\vee, \ldots, t_{r_1}^\vee$ est la base duale de $\mathcal{E}_1^\vee$, et où $s'_i$ est un relèvement local de $s_i$ en une section de $\mathcal{E}_0$. Nous omettons la vérification que cette définition ne dépend pas des choix.

\medskip\noindent
Si $\mathcal{F}$ est localement libre de rang $r$, il est immédiat de vérifier que le morphisme est un isomorphisme. Réciproquement, supposons que le morphisme soit un isomorphisme sur les fibres en $x$. Alors $\wedge^r\mathcal{F}_x$ est inversible. Cela implique que $\mathcal{F}_x$ est engendré par au plus $r$ éléments. Cela n'est possible que si $\alpha$ est de rang $r$ modulo $\mathfrak m_x$, c'est-à-dire si $\alpha$ est de rang maximal modulo $\mathfrak m_x$. Il s'ensuit que $\alpha$ est de rang maximal dans un voisinage de $x$ et donc que $\mathcal{F}$ est localement libre de rang $r$ dans un voisinage, comme souhaité.
\end{proof}

\begin{lemma}
\label{duality-lemma-fundamental-class-lci}
Soit $Y$ un schéma noethérien. Soit $f : X \to Y$ un morphisme localement d'intersection complète qui se factorise en une immersion $X \to P$ suivie d'un morphisme propre et lisse $P \to Y$. Soit $r$ la fonction localement constante sur $X$ telle que $\omega_{X/Y} = H^{-r}(f^!\mathcal{O}_Y)$ soit l'unique faisceau de cohomologie non nul de $f^!\mathcal{O}_Y$, voir le lemme \ref{duality-lemma-lci-shriek}. Il existe alors un morphisme
$$
\wedge^r\Omega_{X/Y} \longrightarrow \omega_{X/Y}
$$
qui est un isomorphisme sur la fibre en un point $x$ si et seulement si $f$ est lisse en $x$.
\end{lemma}

\begin{proof}
L'hypothèse implique que $X$ est compactifiable sur $Y$ ; ainsi $f^!$ est défini, voir la section \ref{duality-section-upper-shriek}. Soit $j : W \to P$ un sous-schéma ouvert tel que $X \to P$ se factorise par une immersion fermée $i : X \to W$. En outre, on a $f^! = i^! \circ j^! \circ g^!$, où $g : P \to Y$ est le morphisme donné. On a $g^!\mathcal{O}_Y = \wedge^d\Omega_{P/Y}[d]$ d'après le lemme \ref{duality-lemma-smooth-proper}, où $d$ est la fonction localement constante donnant la dimension relative de $P$ sur $Y$. On a $j^! = j^*$. On a $i^!\mathcal{O}_W = \wedge^c\mathcal{N}[-c]$, où $c$ est la codimension de $X$ dans $W$ (une fonction localement constante sur $X$) et où $\mathcal{N}$ est le faisceau normal de l'immersion régulière au sens de Koszul $i$, voir le lemme \ref{duality-lemma-regular-immersion}. En combinant ce qui précède, on trouve
$$
f^!\mathcal{O}_Y =
\left(\wedge^c\mathcal{N} \otimes_{\mathcal{O}_X}
\wedge^d\Omega_{P/Y}|_X\right)[d - c]
$$
où l'on a également utilisé le lemme \ref{duality-lemma-perfect-comparison-shriek}. Ainsi, $r = d|_X - c$ comme fonctions localement constantes sur $X$. Le faisceau conormal de $X \to P$ est le module $\mathcal{I}/\mathcal{I}^2$, où $\mathcal{I} \subset \mathcal{O}_W$ est le faisceau d'idéaux de $i$, voir Morphismes, section \ref{morphisms-section-conormal-sheaf}. Considérons la suite exacte canonique
$$
\mathcal{I}/\mathcal{I}^2 \to
\Omega_{P/Y}|_X \to \Omega_{X/Y} \to 0
$$
de Morphismes, lemme \ref{morphisms-lemma-differentials-relative-immersion}. Notre morphisme s'obtient en appliquant le lemme \ref{duality-lemma-determinant}.

\medskip\noindent
Si $f$ est lisse en $x$, alors le morphisme est un isomorphisme par application du lemme \ref{duality-lemma-determinant} et parce que $\Omega_{X/Y}$ est localement libre en $x$ de rang $r$. Réciproquement, supposons que notre morphisme soit un isomorphisme sur les fibres en $x$. Le lemme montre alors que $\Omega_{X/Y}$ est libre de rang $r$ après avoir remplacé $X$ par un voisinage ouvert de $x$. D'autre part, on peut également supposer que $X = \Spec(A)$ et $Y = \Spec(R)$, où $A = R[x_1, \ldots, x_n]/(f_1, \ldots, f_m)$ et où $f_1, \ldots, f_m$ est une suite régulière au sens de Koszul (cela résulte de la définition des morphismes localement d'intersection complète). Cela implique clairement $r = n - m$. On conclut que le rang de la matrice des dérivées partielles $\partial f_j/\partial x_i$ dans le corps résiduel en $x$ est $m$. Après avoir réordonné les variables, on peut donc supposer que le déterminant de $(\partial f_j/\partial x_i)_{1 \leq i, j \leq m}$ est inversible dans un voisinage ouvert de $x$. Il s'ensuit que $R \to A$ est lisse en ce point, voir par exemple Algèbre, exemple \ref{algebra-example-make-standard-smooth}.
\end{proof}

\begin{lemma}
\label{duality-lemma-fundamental-class-almost-lci}
Soit $f : X \to Y$ un morphisme de schémas. Soit $r \geq 0$. Supposons que
\begin{enumerate}
\item $Y$ soit de Cohen–Macaulay (Propriétés, définition \ref{properties-definition-Cohen-Macaulay}) ;
\item $f$ se factorise sous la forme $X \to P \to Y$, où le premier morphisme est une immersion et le second est lisse et propre ;
\item si $x \in X$ et $\dim(\mathcal{O}_{X, x}) \leq 1$, alors $f$ soit de Koszul en $x$ (Compléments sur les morphismes, définition \ref{more-morphisms-definition-lci}) ; et
\item si $\xi$ est un point générique d'une composante irréductible de $X$, alors on ait
$\text{trdeg}_{\kappa(f(\xi))} \kappa(\xi) = r$.
\end{enumerate}
Alors, avec $\omega_{X/Y} = H^{-r}(f^!\mathcal{O}_Y)$, il existe un morphisme
$$
\wedge^r\Omega_{X/Y} \longrightarrow \omega_{X/Y}
$$
qui est un isomorphisme sur le lieu où $f$ est lisse.
\end{lemma}

\begin{proof}
Soit $U \subset X$ le sous-schéma ouvert au-dessus duquel $f$ est un morphisme localement d'intersection complète. Puisque $f$ est de dimension relative $r$ en tout point générique par l'hypothèse (4), on voit que la fonction localement constante du lemme \ref{duality-lemma-fundamental-class-lci} est constante de valeur $r$, et l'on obtient un morphisme
$$
\wedge^r\Omega_{X/Y}|_U = \wedge^r \Omega_{U/Y}
\longrightarrow
\omega_{U/Y} = \omega_{X/Y}|_U
$$
qui est un isomorphisme aux points lisses de $f$ (ce lieu est contenu dans $U$, puisqu'un morphisme lisse est localement d'intersection complète). D'après le lemme \ref{duality-lemma-shriek-over-CM} et l'hypothèse que $Y$ est de Cohen–Macaulay, le module $\omega_{X/Y}$ est $(S_2)$. Puisque $U$ contient tous les points de codimension $1$ d'après la condition (3), et en utilisant Diviseurs, lemme \ref{divisors-lemma-depth-2-hartog}, on voit que $j_*\omega_{U/Y} = \omega_{X/Y}$. Par conséquent, le morphisme sur $U$ s'étend à $X$, ce qui achève la démonstration.
\end{proof}






\section{Prolongement par zéro des modules cohérents}
\label{duality-section-extension-by-zero}

\noindent
Le contenu de cette section et des quelques sections suivantes se trouve dans l'appendice de Deligne de \cite{RD}.

\medskip\noindent
Dans cette section, $j : U \to X$ désignera une immersion ouverte de schémas noethériens. Nous allons considérer des systèmes projectifs $(K_n)$ dans $D^b_{\textit{Coh}}(\mathcal{O}_X)$ construits comme suit. Soit $\mathcal{F}^\bullet$ un complexe borné de $\mathcal{O}_X$-modules cohérents. Soit $\mathcal{I} \subset \mathcal{O}_X$ un faisceau quasi-cohérent d'idéaux tel que $V(\mathcal{I}) = X \setminus U$. On peut alors poser
$$
K_n = \mathcal{I}^n\mathcal{F}^\bullet
$$
Plus précisément, $K_n$ est l'objet de $D^b_{\textit{Coh}}(\mathcal{O}_X)$ représenté par le complexe dont le terme en degré $q$ est le sous-module cohérent $\mathcal{I}^n\mathcal{F}^q$ de $\mathcal{F}^q$. Remarquons que les morphismes $\ldots \to K_3 \to K_2 \to K_1$ induisent des isomorphismes après restriction à $U$. Appelons un tel système un {\it système de Deligne}.

\begin{lemma}
\label{duality-lemma-lift-map}
Soit $j : U \to X$ une immersion ouverte de schémas noethériens. Soit $(K_n)$ un système de Deligne, et notons $K \in D^b_{\textit{Coh}}(\mathcal{O}_U)$ la valeur du système constant $(K_n|_U)$. Soit $L$ un objet de $D^b_{\textit{Coh}}(\mathcal{O}_X)$. Alors $\colim \Hom_X(K_n, L) = \Hom_U(K, L|_U)$.
\end{lemma}

\begin{proof}
Soit $L \to M \to N \to L[1]$ un triangle distingué dans $D^b_{\textit{Coh}}(\mathcal{O}_X)$. On obtient alors un diagramme commutatif
\begingroup\scriptsize
$$
\xymatrix{
\ldots \ar[r] &
\colim \Hom_X(K_n, L) \ar[r] \ar[d] &
\colim \Hom_X(K_n, M) \ar[r] \ar[d] &
\colim \Hom_X(K_n, N) \ar[r] \ar[d] &
\ldots \\
\ldots \ar[r] &
\Hom_U(K, L|_U) \ar[r] &
\Hom_U(K, M|_U) \ar[r] &
\Hom_U(K, N|_U) \ar[r] &
\ldots
}
$$
\endgroup
dont les lignes sont exactes d'après Catégories dérivées, lemme \ref{derived-lemma-representable-homological}, et Algèbre, lemme \ref{algebra-lemma-directed-colimit-exact}. Ainsi, si l'énoncé du lemme vaut pour $N[-1]$, $L$, $N$ et $L[1]$, alors il vaut pour $M$ d'après le lemme des cinq. En utilisant les triangles distingués des troncatures canoniques de $L$ (voir Catégories dérivées, remarque \ref{derived-remark-truncation-distinguished-triangle}), on se ramène donc au cas où $L$ ne possède qu'un seul faisceau de cohomologie non nul.

\medskip\noindent
Choisissons un complexe borné $\mathcal{F}^\bullet$ de $\mathcal{O}_X$-modules cohérents et un idéal quasi-cohérent $\mathcal{I} \subset \mathcal{O}_X$ définissant $X \setminus U$, tels que $K_n$ soit représenté par $\mathcal{I}^n\mathcal{F}^\bullet$. Les troncatures « brutales » donnent des suites exactes courtes compatibles de complexes, scindées terme à terme,
$$
0 \to \sigma_{\geq a + 1} \mathcal{I}^n\mathcal{F}^\bullet \to
\mathcal{I}^n\mathcal{F}^\bullet \to
\sigma_{\leq a} \mathcal{I}^n\mathcal{F}^\bullet \to 0
$$
qui correspondent à leur tour à des systèmes compatibles de triangles distingués dans $D^b_{\textit{Coh}}(\mathcal{O}_X)$. En raisonnant comme ci-dessus, on se ramène au cas où $\mathcal{F}^\bullet$ ne possède qu'un seul terme non nul. On est ainsi ramené au cas examiné dans le paragraphe suivant.

\medskip\noindent
Étant donnés un $\mathcal{O}_X$-module cohérent $\mathcal{F}$ et un $\mathcal{O}_X$-module cohérent $\mathcal{G}$, il faut montrer que le morphisme canonique
$$
\colim \Ext^i_X(\mathcal{I}^n\mathcal{F}, \mathcal{G})
\longrightarrow
\Ext^i_U(\mathcal{F}|_U, \mathcal{G}|_U)
$$
est un isomorphisme pour tout $i \geq 0$. Pour $i = 0$, c'est Cohomologie des schémas, lemme \ref{coherent-lemma-homs-over-open}. Supposons $i > 0$.

\medskip\noindent
Injectivité. Soit $\xi \in \Ext^i_X(\mathcal{I}^n\mathcal{F}, \mathcal{G})$ un élément dont la restriction à $U$ est nulle. Il faut montrer qu'il existe un $m \geq n$ tel que la restriction de $\xi$ à $\mathcal{I}^m\mathcal{F} = \mathcal{I}^{m - n}\mathcal{I}^n\mathcal{F}$ soit nulle. Après avoir remplacé $\mathcal{F}$ par $\mathcal{I}^n\mathcal{F}$, on peut supposer $n = 0$, c'est-à-dire que l'on dispose de
$\xi \in \Ext^i_X(\mathcal{F}, \mathcal{G})$
dont la restriction à $U$ est nulle. D'après Catégories dérivées de schémas, proposition \ref{perfect-proposition-DCoh}, on a
$D^b_{\textit{Coh}}(\mathcal{O}_X) = D^b(\textit{Coh}(\mathcal{O}_X))$.
On peut donc calculer le groupe $\Ext$ dans la catégorie abélienne des $\mathcal{O}_X$-modules cohérents. Cela implique qu'il existe une surjection $\alpha : \mathcal{F}'' \to \mathcal{F}$ telle que $\xi \circ \alpha = 0$ (c'est ici que l'on utilise $i > 0$). Posons $\mathcal{F}' = \Ker(\alpha)$, de sorte que l'on ait une suite exacte courte
$$
0 \to \mathcal{F}' \to \mathcal{F}'' \to \mathcal{F} \to 0
$$
Il s'ensuit que $\xi$ est l'image d'un élément $\xi' \in \Ext^{i - 1}_X(\mathcal{F}', \mathcal{G})$ dont la restriction à $U$ appartient à l'image de $\Ext^{i - 1}_U(\mathcal{F}''|_U, \mathcal{G}|_U) \to
\Ext^{i - 1}_U(\mathcal{F}'|_U, \mathcal{G}|_U)$. Par Artin–Rees, les systèmes projectifs $(\mathcal{I}^n\mathcal{F}')$ et $(\mathcal{I}^n \mathcal{F}'' \cap \mathcal{F}')$ sont pro-isomorphes, voir Cohomologie des schémas, lemme \ref{coherent-lemma-Artin-Rees}. Puisque l'on dispose du système compatible de suites exactes courtes
$$
0 \to
\mathcal{F}' \cap \mathcal{I}^n\mathcal{F}'' \to
\mathcal{I}^n\mathcal{F}'' \to
\mathcal{I}^n\mathcal{F} \to 0
$$
on obtient un diagramme commutatif
\begingroup\small
$$
\xymatrix{
\colim \Ext^{i - 1}_X(\mathcal{I}^n\mathcal{F}'', \mathcal{G})
\ar[r] \ar[d] &
\colim \Ext^{i - 1}_X(\mathcal{F}' \cap \mathcal{I}^n\mathcal{F}'', \mathcal{G})
\ar[r] \ar[d] &
\colim \Ext^i_X(\mathcal{I}^n\mathcal{F}, \mathcal{G})
\ar[d] \\
\Ext^{i - 1}_U(\mathcal{F}''|_U, \mathcal{G}|_U) \ar[r] &
\Ext^{i - 1}_U(\mathcal{F}'|_U, \mathcal{G}|_U) \ar[r] &
\Ext^{i - 1}_U(\mathcal{F}|_U, \mathcal{G}|_U)
}
$$
\endgroup
à lignes exactes. Par récurrence sur $i$ et d'après l'observation ci-dessus sur les systèmes projectifs, on trouve que les deux flèches verticales de gauche sont des isomorphismes. Maintenant, $\xi$ fournit un élément du groupe supérieur droit qui est l'image de $\xi'$ dans le groupe supérieur médian, lequel s'envoie à son tour sur un élément du groupe inférieur médian provenant d'un élément du groupe inférieur gauche. On conclut que $\xi$ s'envoie sur zéro dans
$\Ext^i_X(\mathcal{I}^n\mathcal{F}, \mathcal{G})$
pour un certain $n$, comme souhaité.

\medskip\noindent
Surjectivité. Soit $\xi \in \Ext^i_U(\mathcal{F}|_U, \mathcal{G}|_U)$. En raisonnant comme ci-dessus et en utilisant $i > 0$, on peut trouver une surjection $\mathcal{H} \to \mathcal{F}|_U$ de $\mathcal{O}_U$-modules cohérents telle que $\xi$ s'envoie sur zéro dans $\Ext^i_U(\mathcal{H}, \mathcal{G}|_U)$. On peut alors trouver un morphisme $\varphi : \mathcal{F}'' \to \mathcal{F}$ de $\mathcal{O}_X$-modules cohérents dont la restriction à $U$ est $\mathcal{H} \to \mathcal{F}|_U$, voir Propriétés, lemme \ref{properties-lemma-extend-finite-presentation}. Remarquons que le lemme ne garantit pas que $\varphi$ soit surjectif, mais cela n'aura pas d'importance (on peut choisir un $\varphi$ surjectif au prix d'un peu de travail supplémentaire). Notons $\mathcal{F}' = \Ker(\varphi)$. La suite exacte courte
$$
0 \to \mathcal{F}'|_U \to \mathcal{F}''|_U \to \mathcal{F}|_U \to 0
$$
montre que $\xi$ est l'image de $\xi'$ dans $\Ext^{i - 1}_U(\mathcal{F}'|_U, \mathcal{G}|_U)$. Par récurrence sur $i$, on peut trouver un $n$ tel que $\xi'$ soit l'image d'un certain $\xi'_n$ dans $\Ext^{i - 1}_X(\mathcal{I}^n\mathcal{F}', \mathcal{G})$. Par Artin–Rees, on peut trouver un $m \geq n$ tel que $\mathcal{F}' \cap \mathcal{I}^m\mathcal{F}'' \subset
\mathcal{I}^n\mathcal{F}'$. En utilisant la suite exacte courte
$$
0 \to \mathcal{F}' \cap \mathcal{I}^m\mathcal{F}'' \to
\mathcal{I}^m\mathcal{F}'' \to \mathcal{I}^m\Im(\varphi) \to 0
$$
l'image de $\xi'_n$ dans
$\Ext^{i - 1}_X(\mathcal{F}' \cap \mathcal{I}^m\mathcal{F}'', \mathcal{G})$
est envoyée par le morphisme de bord sur un élément $\xi_m$ de $\Ext^i_X(\mathcal{I}^m\Im(\varphi), \mathcal{G})$ qui s'envoie sur $\xi$. Puisque $\Im(\varphi)$ et $\mathcal{F}$ coïncident sur $U$, on voit que $\mathcal{F}/\mathcal{I}^m\Im(\varphi)$ est à support dans $X \setminus U$. Il existe donc un $l \geq m$ tel que $\mathcal{I}^l\mathcal{F} \subset \mathcal{I}^m\Im(\varphi)$, voir Cohomologie des schémas, lemme \ref{coherent-lemma-power-ideal-kills-sheaf}. En prenant l'image de $\xi_m$ dans $\Ext^i_X(\mathcal{I}^l\mathcal{F}, \mathcal{G})$, on conclut.
\end{proof}

\begin{lemma}
\label{duality-lemma-lift-map-plus}
Le résultat du lemme \ref{duality-lemma-lift-map} vaut même pour $L \in D^+_{\textit{Coh}}(\mathcal{O}_X)$.
\end{lemma}

\begin{proof}
En effet, si $(K_n)$ est un système de Deligne, il existe un $b \in \mathbf{Z}$ tel que $H^i(K_n) = 0$ pour $i > b$. Alors $\Hom(K_n, L) = \Hom(K_n, \tau_{\leq b}L)$ et $\Hom(K, L) = \Hom(K, \tau_{\leq b}L)$. En utilisant le résultat du lemme pour $\tau_{\leq b}L$, on conclut.
\end{proof}

\begin{lemma}
\label{duality-lemma-extension-by-zero}
Soit $j : U \to X$ une immersion ouverte de schémas noethériens.
\begin{enumerate}
\item Soient $(K_n)$ et $(L_n)$ des systèmes de Deligne. Soient $K$ et $L$ les valeurs des systèmes constants $(K_n|_U)$ et $(L_n|_U)$. Étant donné un morphisme $\alpha : K \to L$ de $D(\mathcal{O}_U)$, il existe un unique morphisme de pro-systèmes $(K_n) \to (L_n)$ de $D^b_{\textit{Coh}}(\mathcal{O}_X)$ dont la restriction à $U$ est $\alpha$.
\item Étant donné $K \in D^b_{\textit{Coh}}(\mathcal{O}_U)$, il existe un système de Deligne $(K_n)$ tel que $(K_n|_U)$ soit constant de valeur $K$.
\item Le pro-objet $(K_n)$ de $D^b_{\textit{Coh}}(\mathcal{O}_X)$ de (2) est unique à isomorphisme unique près (comme pro-objet).
\end{enumerate}
\end{lemma}

\begin{proof}
L'assertion (1) résulte immédiatement du lemme \ref{duality-lemma-lift-map} et du fait que les morphismes entre pro-systèmes sont les mêmes que les morphismes entre les foncteurs qu'ils corereprésentent, voir Catégories, remarque \ref{categories-remark-pro-category-copresheaves}.

\medskip\noindent
Soit $K$ comme dans (2). On peut choisir $K' \in D^b_{\textit{Coh}}(\mathcal{O}_X)$ dont la restriction à $U$ est isomorphe à $K$, voir Catégories dérivées de schémas, lemme \ref{perfect-lemma-lift-coherent}. D'après Catégories dérivées de schémas, proposition \ref{perfect-proposition-DCoh}, on peut représenter $K'$ par un complexe borné $\mathcal{F}^\bullet$ de $\mathcal{O}_X$-modules cohérents. Choisissons un faisceau quasi-cohérent d'idéaux $\mathcal{I} \subset \mathcal{O}_X$ dont le lieu d'annulation est $X \setminus U$ (par exemple, choisissons $\mathcal{I}$ correspondant à la structure de sous-schéma induit réduit sur $X \setminus U$). On peut alors prendre pour $K_n$ l'objet représenté par le complexe $\mathcal{I}^n\mathcal{F}^\bullet$ comme dans l'introduction de cette section.

\medskip\noindent
L'assertion (3) résulte immédiatement des assertions (1) et (2).
\end{proof}

\begin{lemma}
\label{duality-lemma-extension-by-zero-triangle}
Soit $j : U \to X$ une immersion ouverte de schémas noethériens. Soit
$$
K \to L \to M \to K[1]
$$
un triangle distingué de $D^b_{\textit{Coh}}(\mathcal{O}_U)$. Il existe alors un système projectif de triangles distingués
$$
K_n \to L_n \to M_n \to K_n[1]
$$
dans $D^b_{\textit{Coh}}(\mathcal{O}_X)$ tel que $(K_n)$, $(L_n)$ et $(M_n)$ soient des systèmes de Deligne et que la restriction de ces triangles distingués à $U$ soit isomorphe au triangle distingué dont nous sommes partis.
\end{lemma}

\begin{proof}
Soit $(K_n)$ comme dans le lemme \ref{duality-lemma-extension-by-zero}, assertion (2). Choisissons un objet $L'$ de $D^b_{\textit{Coh}}(\mathcal{O}_X)$ dont la restriction à $U$ est $L$ (cela est possible d'après le lemme). D'après le lemme \ref{duality-lemma-lift-map}, on peut trouver un $n$ et un morphisme $K_n \to L'$ sur $X$ dont la restriction à $U$ est la flèche donnée $K \to L$. On conclut à l'existence d'un morphisme $K' \to L'$ de $D^b_{\textit{Coh}}(\mathcal{O}_X)$ dont la restriction à $U$ est la flèche donnée $K \to L$.

\medskip\noindent
D'après Catégories dérivées de schémas, proposition \ref{perfect-proposition-DCoh}, on peut trouver un morphisme $\alpha^\bullet : \mathcal{F}^\bullet \to
\mathcal{G}^\bullet$ de complexes bornés de $\mathcal{O}_X$-modules cohérents représentant $K' \to L'$. Choisissons un faisceau quasi-cohérent d'idéaux $\mathcal{I} \subset \mathcal{O}_X$ dont le lieu d'annulation est $X \setminus U$. Posons alors
$K_n = \mathcal{I}^n\mathcal{F}^\bullet$
et
$L_n = \mathcal{I}^n\mathcal{G}^\bullet$.
Remarquons que $\alpha^\bullet$ induit un morphisme de complexes $\alpha_n^\bullet : \mathcal{I}^n\mathcal{F}^\bullet \to
\mathcal{I}^n\mathcal{G}^\bullet$. La construction des cônes dans Catégories dérivées, section \ref{derived-section-cones}, montre clairement que
$$
C(\alpha_n)^\bullet = \mathcal{I}^nC(\alpha^\bullet)
$$
et l'on peut donc poser $M_n = C(\alpha_n)^\bullet$. On dispose en effet d'un système compatible de triangles distingués (voir la discussion de Catégories dérivées, section \ref{derived-section-canonical-delta-functor})
$$
K_n \to L_n \to M_n \to K_n[1]
$$
dont la restriction à $U$ est isomorphe au triangle distingué dont nous sommes partis, d'après l'axiome TR3 et Catégories dérivées, lemme \ref{derived-lemma-third-isomorphism-triangle}.
\end{proof}

\begin{remark}
\label{duality-remark-extension-by-zero}
Soit $j : U \to X$ une immersion ouverte de schémas noethériens. Envoyer $K \in D^b_{\textit{Coh}}(\mathcal{O}_U)$ sur un système de Deligne dont la restriction à $U$ est $K$ détermine un foncteur
$$
Rj_! :
D^b_{\textit{Coh}}(\mathcal{O}_U)
\longrightarrow
\text{Pro-}D^b_{\textit{Coh}}(\mathcal{O}_X)
$$
qui est « exact » d'après le lemme \ref{duality-lemma-extension-by-zero-triangle} et qui est « adjoint à gauche » du foncteur $j^* : D^b_{\textit{Coh}}(\mathcal{O}_X) \to
D^b_{\textit{Coh}}(\mathcal{O}_U)$ d'après le lemme \ref{duality-lemma-lift-map}.
\end{remark}

\begin{remark}
\label{duality-remark-extension-by-zero-linear-pro-system}
Soient $(A_n)$ et $(B_n)$ des systèmes projectifs d'une catégorie $\mathcal{C}$. Appelons pro-morphisme linéaire de $(A_n)$ vers $(B_n)$ la donnée d'une famille compatible de morphismes $\varphi_n : A_{cn + d} \to B_n$ pour tout $n \geq 1$, pour des entiers fixes $c, d \geq 1$. On dira que $(\varphi_n : A_{cn + d} \to B_n)$ et $(\psi_n : A_{c'n + d'} \to B_n)$ déterminent le même morphisme s'il existe $c'' \geq \max(c, c')$ et $d'' \geq \max(d, d')$ tels que les deux morphismes induits $A_{c'' n + d''} \to B_n$ soient égaux pour tout $n$. Il semble probable que les systèmes de Deligne $(K_n)$ de valeur donnée sur $U$ soient bien définis à pro-isomorphisme linéaire près. Si nous en avons un jour besoin, nous le formulerons et le démontrerons soigneusement ici.
\end{remark}

\begin{lemma}
\label{duality-lemma-deligne-system-2-out-of-3}
Soit $j : U \to X$ une immersion ouverte de schémas noethériens. Soit
$$
K_n \to L_n \to M_n \to K_n[1]
$$
un système projectif de triangles distingués dans $D^b_{\textit{Coh}}(\mathcal{O}_X)$. Si $(K_n)$ et $(M_n)$ sont pro-isomorphes à des systèmes de Deligne, alors il en est de même de $(L_n)$.
\end{lemma}

\begin{proof}
Remarquons que les systèmes $(K_n|_U)$ et $(M_n|_U)$ sont essentiellement constants puisqu'ils sont pro-isomorphes à des systèmes constants. Notons $K$ et $M$ leurs valeurs. D'après Catégories dérivées, lemme \ref{derived-lemma-essentially-constant-2-out-of-3}, le système projectif $L_n|_U$ est lui aussi essentiellement constant. Notons $L$ sa valeur. Soit $N \in D^b_{\textit{Coh}}(\mathcal{O}_X)$. Considérons le diagramme commutatif
\begingroup\scriptsize
$$
\xymatrix{
\ldots \ar[r] &
\colim \Hom_X(M_n, N) \ar[r] \ar[d] &
\colim \Hom_X(L_n, N) \ar[r] \ar[d] &
\colim \Hom_X(K_n, N) \ar[r] \ar[d] &
\ldots \\
\ldots \ar[r] &
\Hom_U(M, N|_U) \ar[r] &
\Hom_U(L, N|_U) \ar[r] &
\Hom_U(K, N|_U) \ar[r] &
\ldots
}
$$
\endgroup
D'après le lemme \ref{duality-lemma-lift-map} et le fait que des ind-systèmes isomorphes ont même colimite, les deux flèches verticales à droite et les deux à gauche de celle du milieu sont des isomorphismes. Le lemme des cinq montre que la flèche verticale du milieu est un isomorphisme. Maintenant, si $(L'_n)$ est un système de Deligne dont la restriction à $U$ a pour valeur constante $L$ (un tel système existe d'après le lemme \ref{duality-lemma-extension-by-zero}), alors on a aussi $\colim \Hom_X(L'_n, N) = \Hom_U(L, N|_U)$. Les pro-systèmes $(L_n)$ et $(L'_n)$ sont donc pro-isomorphes d'après Catégories, remarque \ref{categories-remark-pro-category-copresheaves}.
\end{proof}

\begin{lemma}
\label{duality-lemma-consequence-Artin-Rees-bis}
Soit $X$ un schéma noethérien. Soit $\mathcal{I} \subset \mathcal{O}_X$ un faisceau quasi-cohérent d'idéaux. Soit $\mathcal{F}^\bullet$ un complexe de $\mathcal{O}_X$-modules cohérents. Soit $p \in \mathbf{Z}$. Posons $\mathcal{H} = H^p(\mathcal{F}^\bullet)$ et $\mathcal{H}_n = H^p(\mathcal{I}^n\mathcal{F}^\bullet)$. Il existe alors des morphismes canoniques de $\mathcal{O}_X$-modules
$$
\ldots \to \mathcal{H}_3 \to \mathcal{H}_2 \to \mathcal{H}_1 \to \mathcal{H}
$$
Il existe un $c > 0$ tel que, pour $n \geq c$, l'image de $\mathcal{H}_n \to \mathcal{H}$ soit contenue dans $\mathcal{I}^{n - c}\mathcal{H}$, et il existe un morphisme canonique de $\mathcal{O}_X$-modules $\mathcal{I}^n\mathcal{H} \to \mathcal{H}_{n - c}$ tel que les composés
$$
\mathcal{I}^n \mathcal{H} \to \mathcal{H}_{n - c} \to
\mathcal{I}^{n - 2c}\mathcal{H}
\quad\text{et}\quad
\mathcal{H}_n \to \mathcal{I}^{n - c}\mathcal{H} \to \mathcal{H}_{n - 2c}
$$
soient les morphismes canoniques. En particulier, les systèmes projectifs $(\mathcal{H}_n)$ et $(\mathcal{I}^n\mathcal{H})$ sont isomorphes comme pro-objets de $\textit{Mod}(\mathcal{O}_X)$.
\end{lemma}

\begin{proof}
Si $X$ est affine, la traduction en algèbre donne Compléments sur l'algèbre, lemme \ref{more-algebra-lemma-consequence-Artin-Rees-bis}. Dans le cas général, on raisonne exactement comme dans la démonstration de ce lemme, en remplaçant la référence à Artin–Rees en algèbre par une référence à Cohomologie des schémas, lemme \ref{coherent-lemma-Artin-Rees}. Les détails sont omis.
\end{proof}

\begin{lemma}
\label{duality-lemma-characterize-extension-by-zero-algebra}
Soit $j : U \to X$ une immersion ouverte de schémas noethériens. Soient $a \leq b$ des entiers. Soit $(K_n)$ un système projectif de $D^b_{\textit{Coh}}(\mathcal{O}_X)$ tel que $H^i(K_n) = 0$ pour $i \not \in [a, b]$. Les conditions suivantes sont équivalentes :
\begin{enumerate}
\item $(K_n)$ est pro-isomorphe à un système de Deligne ;
\item pour tout $p \in \mathbf{Z}$, il existe un $\mathcal{O}_X$-module cohérent $\mathcal{F}$ tel que les pro-systèmes $(H^p(K_n))$ et $(\mathcal{I}^n\mathcal{F})$ soient pro-isomorphes.
\end{enumerate}
\end{lemma}

\begin{proof}
Supposons (1). Pour démontrer (2), on peut supposer que $(K_n)$ est un système de Deligne. Par définition, on peut choisir un complexe borné $\mathcal{F}^\bullet$ de $\mathcal{O}_X$-modules cohérents et un faisceau quasi-cohérent d'idéaux définissant $X \setminus U$, tels que $K_n$ soit représenté par $\mathcal{I}^n\mathcal{F}^\bullet$. Le résultat résulte donc du lemme \ref{duality-lemma-consequence-Artin-Rees-bis}.

\medskip\noindent
Supposons (2). Nous allons démontrer que $(K_n)$ est comme dans (1), par récurrence sur $b - a$. Si $a = b$, alors (1) vaut essentiellement par hypothèse. Si $a < b$, considérons le système compatible de triangles distingués
$$
\tau_{\leq a}K_n \to K_n \to \tau_{\geq a + 1}K_n \to (\tau_{\leq a}K_n)[1]
$$
Voir Catégories dérivées, remarque \ref{derived-remark-truncation-distinguished-triangle}. Par récurrence sur $b - a$, on sait que $\tau_{\leq a}K_n$ et $\tau_{\geq a + 1}K_n$ sont pro-isomorphes à des systèmes de Deligne. On conclut par le lemme \ref{duality-lemma-deligne-system-2-out-of-3}.
\end{proof}

\begin{lemma}
\label{duality-lemma-characterize-extension-by-zero}
Soit $j : U \to X$ une immersion ouverte de schémas noethériens. Soit $(K_n)$ un système projectif dans $D^b_{\textit{Coh}}(\mathcal{O}_X)$. Soit $X = W_1 \cup \ldots \cup W_r$ un recouvrement ouvert. Les conditions suivantes sont équivalentes :
\begin{enumerate}
\item $(K_n)$ est pro-isomorphe à un système de Deligne ;
\item pour chaque $i$, la restriction $(K_n|_{W_i})$ est pro-isomorphe à un système de Deligne relativement à l'immersion ouverte $U \cap W_i \to W_i$.
\end{enumerate}
\end{lemma}

\begin{proof}
On raisonne par récurrence sur $r$. Si $r = 1$, le résultat est clair. Supposons $r > 1$. Posons $V = W_1 \cup \ldots \cup W_{r - 1}$. Par récurrence, on voit que $(K_n|_V)$ est un système de Deligne. On est ainsi ramené à la discussion du paragraphe suivant.

\medskip\noindent
Supposons que $X = V \cup W$ soit un recouvrement ouvert et que $(K_n|_W)$ et $(K_n|_V)$ soient pro-isomorphes à des systèmes de Deligne. Il faut montrer que $(K_n)$ est pro-isomorphe à un système de Deligne. Remarquons que $(K_n|_{V \cap W})$ est pro-isomorphe à un système de Deligne (il résulte immédiatement de la construction des systèmes de Deligne que leur restriction aux ouverts les préserve). En particulier, les pro-systèmes
$(K_n|_{U \cap V})$,
$(K_n|_{U \cap W})$ et
$(K_n|_{U \cap V \cap W})$
sont essentiellement constants. Il résulte des triangles distingués de Cohomologie, lemme \ref{cohomology-lemma-exact-sequence-j-star}, et de Catégories dérivées, lemme \ref{derived-lemma-essentially-constant-2-out-of-3}, que $(K_n|_U)$ est essentiellement constant. Notons $K \in D^b_{\textit{Coh}}(\mathcal{O}_U)$ la valeur de ce système. Soit $L$ un objet de $D^b_{\textit{Coh}}(\mathcal{O}_X)$. Considérons le diagramme
\begingroup\tiny
$$
\xymatrix{
\colim \Ext^{-1}(K_n|_V, L|_V) \oplus
\colim \Ext^{-1}(K_n|_W, L|_W) \ar[r] \ar[d] &
\Ext^{-1}(K|_{U \cap V}, L|_{U \cap V}) \oplus
\Ext^{-1}(K|_{U \cap W}, L|_{U \cap W}) \ar[d] \\
\colim \Ext^{-1}(K_n|_{V \cap W}, L|_{V \cap W}) \ar[r] \ar[d] &
\Ext^{-1}(K|_{U \cap V \cap W}, L|_{U \cap V \cap W}) \ar[d] \\
\colim \Hom(K_n, L) \ar[d] \ar[r] &
\Hom(K|_U, L|_U) \ar[d] \\
\colim \Hom(K_n|_V, L|_V) \oplus \colim \Hom(K_n|_W, L|_W) \ar[r] \ar[d] &
\Hom(K|_{U \cap V}, L|_{U \cap V}) \oplus
\Hom(K|_{U \cap W}, L|_{U \cap W}) \ar[d] \\
\colim \Hom(K_n|_{V \cap W}, L|_{V \cap W}) \ar[r] &
\Hom(K|_{U \cap V \cap W}, L|_{U \cap V \cap W})
}
$$
\endgroup
Les suites verticales sont exactes d'après Cohomologie, lemme \ref{cohomology-lemma-mayer-vietoris-hom}, et parce que les colimites filtrantes sont exactes. Toutes les flèches horizontales, sauf celle du milieu, sont des isomorphismes d'après le lemme \ref{duality-lemma-lift-map} et le fait que des systèmes pro-isomorphes ont mêmes colimites. Par le lemme des cinq, celle du milieu est donc aussi un isomorphisme. Il s'ensuit que $(K_n)$ est pro-isomorphe à un système de Deligne pour $K$. En effet, si $(K'_n)$ est un système de Deligne dont la restriction à $U$ a pour valeur constante $K$ (un tel système existe d'après le lemme \ref{duality-lemma-extension-by-zero}), alors on a aussi $\colim \Hom_X(K'_n, L) = \Hom_U(K, L|_U)$. Les pro-systèmes $(K_n)$ et $(K'_n)$ sont donc pro-isomorphes d'après Catégories, remarque \ref{categories-remark-pro-category-copresheaves}.
\end{proof}

\begin{lemma}
\label{duality-lemma-tensoring-Deligne-system}
Soit $j : U \to X$ une immersion ouverte de schémas noethériens. Soit $\mathcal{I} \subset \mathcal{O}_X$ un faisceau quasi-cohérent d'idéaux tel que $V(\mathcal{I}) = X \setminus U$. Soit $K$ dans $D^b_{\textit{Coh}}(\mathcal{O}_X)$. Alors
$$
K \otimes_{\mathcal{O}_X}^\mathbf{L} \mathcal{I}^n
$$
est pro-isomorphe à un système de Deligne de valeur constante $K|_U$ sur $U$.
\end{lemma}

\begin{proof}
D'après le lemme \ref{duality-lemma-characterize-extension-by-zero}, la question est locale sur $X$. On peut donc supposer que $X$ est le spectre d'un anneau noethérien. Dans ce cas, l'énoncé résulte de la version algébrique, qui est Compléments sur l'algèbre, lemme \ref{more-algebra-lemma-tensoring-Deligne-system}.
\end{proof}







\section{Préliminaires à la cohomologie à support compact}
\label{duality-section-preliminaries-compactly-supported}

\noindent
Dans la situation \ref{duality-situation-shriek}, soit $f : X \to Y$ un morphisme de la catégorie $\textit{FTS}_S$. À l'aide des constructions de la section précédente, nous construirons un foncteur
$$
Rf_! :
D^b_{\textit{Coh}}(\mathcal{O}_X)
\longrightarrow
\text{Pro-}D^b_{\textit{Coh}}(\mathcal{O}_Y)
$$
qui se réduit au foncteur de la remarque \ref{duality-remark-extension-by-zero} si $f$ est une immersion ouverte et qui, en général, se construit au moyen d'une compactification de $f$. Avant cela, nous avons besoin des lemmes suivants pour démontrer que notre construction est bien définie.

\begin{lemma}
\label{duality-lemma-well-defined-pre}
Soit $f : X \to Y$ un morphisme propre de schémas noethériens. Soit $V \subset Y$ un sous-schéma ouvert et posons $U = f^{-1}(V)$. Diagramme
$$
\xymatrix{
U \ar[r]_j \ar[d]_g & X \ar[d]^f \\
V \ar[r]^{j'} & Y
}
$$
Il existe alors un isomorphisme canonique $Rj'_! \circ Rg_* \to Rf_* \circ Rj_!$ de foncteurs $D^b_{\textit{Coh}}(\mathcal{O}_U) \to
\text{Pro-}D^b_{\textit{Coh}}(\mathcal{O}_Y)$, où $Rj_!$ et $Rj'_!$ sont comme dans la remarque \ref{duality-remark-extension-by-zero}.
\end{lemma}

\begin{proof}[Première démonstration]
Soit $K$ un objet de $D^b_{\textit{Coh}}(\mathcal{O}_U)$. Soit $(K_n)$ un système de Deligne pour $U \to X$ dont la restriction à $U$ est constante de valeur $K$. Cela signifie bien sûr que $(K_n)$ représente $Rj_!K$ dans $\text{Pro-}D^b_{\textit{Coh}}(\mathcal{O}_X)$. Remarquons que $Rj'_!Rg_*K$ et $Rf_*Rj_!K$ se restreignent tous deux au pro-objet constant de valeur $Rg_*K$ sur $V$. C'est immédiat pour le premier et, pour le second, cela résulte de $(Rf_*K_n)|_V = Rg_*(K_n|_U) = Rg_*K$. Par l'unicité des systèmes de Deligne du lemme \ref{duality-lemma-extension-by-zero}, il suffit de montrer que $(Rf_*K_n)$ est pro-isomorphe à un système de Deligne. Le lemme cité montrera aussi que l'isomorphisme obtenu est fonctoriel.

\medskip\noindent
Démonstration que $(Rf_*K_n)$ est pro-isomorphe à un système de Deligne. Tout d'abord, remarquons que la question est indépendante du choix du système de Deligne $(K_n)$ correspondant à $K$ (par l'unicité précitée). D'après les lemmes \ref{duality-lemma-extension-by-zero-triangle} et \ref{duality-lemma-deligne-system-2-out-of-3}, si l'on dispose d'un triangle distingué
$$
K \to L \to M \to K[1]
$$
dans $D^b_{\textit{Coh}}(\mathcal{O}_U)$ et que le résultat vaut pour $K$ et $M$, alors il vaut pour $L$. En utilisant les triangles distingués des troncatures canoniques (Catégories dérivées, remarque \ref{derived-remark-truncation-distinguished-triangle}), on se ramène au problème étudié dans le paragraphe suivant.

\medskip\noindent
Soit $\mathcal{F}$ un $\mathcal{O}_X$-module cohérent. Soit $\mathcal{J} \subset \mathcal{O}_Y$ un faisceau quasi-cohérent d'idéaux définissant $Y \setminus V$. Notons $\mathcal{J}^n\mathcal{F}$ l'image de $f^*\mathcal{J}^n \otimes \mathcal{F} \to \mathcal{F}$. Il faut montrer que $(Rf_*(\mathcal{J}^n\mathcal{F}))$ est un système de Deligne. D'après le lemme \ref{duality-lemma-characterize-extension-by-zero}, la question est locale sur $Y$. On peut donc supposer $Y = \Spec(A)$ affine et que $\mathcal{J}$ correspond à un idéal $I \subset A$. D'après le lemme \ref{duality-lemma-characterize-extension-by-zero-algebra}, il suffit de montrer que le système projectif de modules de cohomologie $(H^p(X, I^n\mathcal{F}))$ est pro-isomorphe au système projectif $(I^n M)$ pour un certain $A$-module fini $M$. Cela est démontré dans Cohomologie des schémas, lemme \ref{coherent-lemma-cohomology-powers-ideal-application}.
\end{proof}

\begin{proof}[Seconde démonstration]
Soit $K$ un objet de $D^b_{\textit{Coh}}(\mathcal{O}_U)$. Soit $L$ un objet de $D^b_{\textit{Coh}}(\mathcal{O}_Y)$. Nous allons construire une bijection
$$
\Hom_{\text{Pro-}D^b_{\textit{Coh}}(\mathcal{O}_Y)}(Rj'_!Rg_*K, L)
\longrightarrow
\Hom_{\text{Pro-}D^b_{\textit{Coh}}(\mathcal{O}_Y)}(Rf_*Rj_!K, L)
$$
fonctorielle en $K$ et $L$. À $K$ fixé, cela déterminera un isomorphisme de pro-objets
$Rf_*Rj_!K \to Rj'_!Rg_*K$
d'après Catégories, remarque \ref{categories-remark-pro-category-copresheaves}, et, lorsque $K$ varie, cela déterminera un isomorphisme de foncteurs. Pour produire effectivement l'isomorphisme, on utilise la suite d'égalités fonctorielles
\begin{align*}
\Hom_{\text{Pro-}D^b_{\textit{Coh}}(\mathcal{O}_Y)}(Rj'_!Rg_*K, L)
& =
\Hom_V(Rg_*K, L|_V) \\
& =
\Hom_U(K, g^!(L|_V)) \\
& =
\Hom_U(K, f^!L|_U)) \\
& =
\Hom_{\text{Pro-}D^b_{\textit{Coh}}(\mathcal{O}_X)}(Rj_!K, f^!L) \\
& =
\Hom_{\text{Pro-}D^b_{\textit{Coh}}(\mathcal{O}_Y)}(Rf_*Rj_!K, L)
\end{align*}
La première égalité résulte du lemme \ref{duality-lemma-lift-map}. La deuxième vaut parce que $g$ est propre (comme changement de base de $f$ à $V$) et que $g^!$ est donc, par construction, l'adjoint à droite de l'image directe, voir la section \ref{duality-section-upper-shriek}. La troisième vaut puisque $g^!(L|_V) = f^!L|_U$ d'après le lemme \ref{duality-lemma-restrict-before-or-after}. Comme $f^!L$ appartient à $D^+_{\textit{Coh}}(\mathcal{O}_X)$ d'après le lemme \ref{duality-lemma-shriek-coherent}, la quatrième égalité résulte du lemme \ref{duality-lemma-lift-map-plus}. La cinquième vaut à nouveau parce que $f^!$ est l'adjoint à droite de $Rf_*$ puisque $f$ est propre.
\end{proof}

\begin{lemma}
\label{duality-lemma-well-defined}
Soit $j : U \to X$ une immersion ouverte de schémas noethériens. Soit $j' : U \to X'$ une compactification de $U$ sur $X$ (voir la démonstration), et notons $f : X' \to X$ le morphisme structural. Il existe alors un isomorphisme canonique $Rj_! \to Rf_* \circ R(j')_!$ de foncteurs $D^b_{\textit{Coh}}(\mathcal{O}_U) \to
\text{Pro-}D^b_{\textit{Coh}}(\mathcal{O}_X)$, où $Rj_!$ et $Rj'_!$ sont comme dans la remarque \ref{duality-remark-extension-by-zero}.
\end{lemma}

\begin{proof}
Dire que $X'$ est une compactification de $U$ sur $X$ signifie précisément que $f : X' \to X$ est propre, que $j'$ est une immersion ouverte et que $j = f \circ j'$. Voir Compléments sur la platitude, section \ref{flat-section-compactify}. Si $j'(U) = f^{-1}(j(U))$, le lemme résulte immédiatement du lemme \ref{duality-lemma-well-defined-pre}. Si $j'(U) \not = f^{-1}(j(U))$, notons $X'' \subset X'$ l'adhérence schématique de $j' : U \to X'$ et $j'' : U \to X''$ l'immersion ouverte correspondante. Diagramme
$$
\xymatrix{
& & X'' \ar[d]^{f'} \\
& & X' \ar[d]^f \\
U \ar[rr]^j \ar[rru]^{j'} \ar[rruu]^{j''} & & X
}
$$
D'après Compléments sur la platitude, lemme \ref{flat-lemma-compactifications-cofiltered}, assertion (c), et la discussion ci-dessus, on a des isomorphismes $Rf'_* \circ Rj''_! = Rj'_!$ et $R(f \circ f')_* \circ Rj''_! = Rj_!$. Puisque $R(f \circ f')_* = Rf_* \circ Rf'_*$, on conclut.
\end{proof}

\begin{remark}
\label{duality-remark-covariance-open-j-lower-shriek}
Soient $X \supset U \supset U'$ des sous-schémas ouverts d'un schéma noethérien $X$. Notons $j : U \to X$ et $j' : U' \to X$ les morphismes d'inclusion. Nous affirmons qu'il existe un morphisme canonique
$$
Rj'_!(K|_{U'}) \longrightarrow Rj_!K
$$
fonctoriel en $K$ dans $D^b_{\textit{Coh}}(\mathcal{O}_U)$. En effet, d'après le lemme \ref{duality-lemma-lift-map}, on dispose, pour tout $L$ dans $D^b_{\textit{Coh}}(\mathcal{O}_X)$, du morphisme
\begin{align*}
\Hom_{\text{Pro-}D^b_{\textit{Coh}}(\mathcal{O}_X)}(Rj_!K, L)
& =
\Hom_U(K, L|_U) \\
& \to
\Hom_{U'}(K|_{U'}, L|_{U'}) \\
& =
\Hom_{\text{Pro-}D^b_{\textit{Coh}}(\mathcal{O}_X)}(Rj'_!(K|_{U'}), L)
\end{align*}
fonctoriel en $L$ et $K'$. La fonctorialité en $L$ montre, d'après Catégories, remarque \ref{categories-remark-pro-category-copresheaves}, que l'on obtient un morphisme canonique $Rj'_!(K|_{U'}) \to Rj_!K$ qui est fonctoriel en $K$ par la fonctorialité en $K$ de la flèche ci-dessus.

\medskip\noindent
Voici une construction explicite de cette flèche. Supposons que $\mathcal{F}^\bullet$ soit un complexe borné de $\mathcal{O}_X$-modules cohérents dont la restriction à $U$ représente $K$ dans la catégorie dérivée. Nous avons vu dans la démonstration du lemme \ref{duality-lemma-extension-by-zero} qu'un tel complexe existe toujours. Soient $\mathcal{I}$, resp.\ $\mathcal{I}'$, des faisceaux quasi-cohérents d'idéaux sur $X$ tels que $V(\mathcal{I}) = X \setminus U$, resp.\ $V(\mathcal{I}') = X \setminus U'$. Après avoir remplacé $\mathcal{I}$ par $\mathcal{I} + \mathcal{I}'$, on peut supposer $\mathcal{I}' \subset \mathcal{I}$. Par construction, $Rj_!K$, resp.\ $Rj'_!(K|_{U'})$, est représenté par le système projectif $(K_n)$, resp.\ $(K'_n)$, de $D^b_{\textit{Coh}}(\mathcal{O}_X)$, avec
$$
K_n = \mathcal{I}^n\mathcal{F}^\bullet
\quad\text{resp.}\quad
K'_n = (\mathcal{I}')^n\mathcal{F}^\bullet
$$
Il est clair que le morphisme construit ci-dessus est donné par les morphismes $K'_n \to K_n$ provenant des inclusions $(\mathcal{I}')^n \subset \mathcal{I}^n$.
\end{remark}






\section{Cohomologie à support compact pour les modules cohérents}
\label{duality-section-compactly-supported}

\noindent
Dans la situation \ref{duality-situation-shriek}, étant donné un morphisme $f : X \to Y$ de $\textit{FTS}_S$, nous allons définir un foncteur
$$
Rf_! : D^b_{\textit{Coh}}(\mathcal{O}_X)
\longrightarrow
\text{Pro-}D^b_{\textit{Coh}}(\mathcal{O}_Y)
$$
Choisissons une compactification $j : X \to \overline{X}$ sur $Y$, ce qui est possible d'après Compléments sur la platitude, théorème \ref{flat-theorem-nagata}, et lemme \ref{flat-lemma-compactifyable}. Notons $\overline{f} : \overline{X} \to Y$ le morphisme structural. Posons alors
$$
Rf_!K  = R\overline{f}_* Rj_! K
$$
pour $K \in D^b_{\textit{Coh}}(\mathcal{O}_X)$, où $Rj_!$ est comme dans la remarque \ref{duality-remark-extension-by-zero}.

\begin{lemma}
\label{duality-lemma-lower-shriek-well-defined}
Le foncteur $Rf_!$ est, à isomorphisme près, indépendant du choix de la compactification.
\end{lemma}

\noindent
En fait, le foncteur $Rf_!$ sera caractérisé comme un « adjoint à gauche » de $f^!$, ce qui le déterminera à isomorphisme unique près.

\begin{proof}
Considérons la catégorie des compactifications de $X$ sur $Y$, qui est cofiltrante d'après Compléments sur la platitude, théorème \ref{flat-theorem-nagata}, et lemmes \ref{flat-lemma-compactifications-cofiltered} et \ref{flat-lemma-compactifyable}. À chaque choix d'une compactification
$$
j : X \to \overline{X},\quad \overline{f} : \overline{X} \to Y
$$
la construction ci-dessus associe le foncteur $R\overline{f}_* \circ Rj_!$. Supposons donné un morphisme $g : \overline{X}_1 \to \overline{X}_2$ entre des compactifications $j_i : X \to \overline{X}_i$ sur $Y$. On obtient alors un isomorphisme
$$
R\overline{f}_{2, *} \circ Rj_{2, !} =
R\overline{f}_{2, *} \circ Rg_* \circ j_{1, !} =
R\overline{f}_{1, *} \circ Rj_{1, !}
$$
en utilisant le lemme \ref{duality-lemma-well-defined} dans la première égalité. On voit ainsi que notre foncteur est indépendant, à isomorphisme près, du choix de la compactification.
\end{proof}

\begin{proposition}
\label{duality-proposition-duality-compactly-supported}
Dans la situation \ref{duality-situation-shriek}, soit $f : X \to Y$ un morphisme de $\textit{FTS}_S$. Les foncteurs $Rf_!$ et $f^!$ sont alors adjoints au sens suivant : pour tous $K \in D^b_{\textit{Coh}}(\mathcal{O}_X)$ et $L \in D^+_{\textit{Coh}}(\mathcal{O}_Y)$, on a
$$
\Hom_X(K, f^!L) =
\Hom_{\text{Pro-}D^+_{\textit{Coh}}(\mathcal{O}_Y)}(Rf_!K, L)
$$
bifonctoriellement en $K$ et $L$.
\end{proposition}

\begin{proof}
Choisissons une compactification $j : X \to \overline{X}$ sur $Y$ et notons $\overline{f} : \overline{X} \to Y$ le morphisme structural. On a alors
\begin{align*}
\Hom_X(K, f^!L)
& =
\Hom_X(K, j^*\overline{f}{}^!L) \\
& =
\Hom_{\text{Pro-}D^+_{\textit{Coh}}(\mathcal{O}_{\overline{X}})}
(Rj_!K, \overline{f}{}^!L) \\
& =
\Hom_{\text{Pro-}D^+_{\textit{Coh}}(\mathcal{O}_Y)}(Rf_*Rj_!K, L) \\
& =
\Hom_{\text{Pro-}D^+_{\textit{Coh}}(\mathcal{O}_Y)}(Rf_!K, L)
\end{align*}
La première égalité résulte immédiatement de la construction de $f^!$ dans la section \ref{duality-section-upper-shriek}. D'après le lemme \ref{duality-lemma-shriek-coherent}, on a $\overline{f}{}^!L$ dans $D^+_{\textit{Coh}}(\mathcal{O}_{\overline{X}})$ ; la deuxième égalité résulte donc du lemme \ref{duality-lemma-lift-map-plus}. Puisque $\overline{f}$ est propre, le foncteur $\overline{f}{}^!$ est, par construction, l'adjoint à droite de l'image directe. C'est ce qui donne la troisième égalité. La quatrième vaut parce que $Rf_! = Rf_* Rj_!$.
\end{proof}

\begin{lemma}
\label{duality-lemma-compactly-supported-triangle}
Dans la situation \ref{duality-situation-shriek}, soit $f : X \to Y$ un morphisme de $\textit{FTS}_S$. Soit
$$
K \to L \to M \to K[1]
$$
un triangle distingué de $D^b_{\textit{Coh}}(\mathcal{O}_X)$. Il existe alors un système projectif de triangles distingués
$$
K_n \to L_n \to M_n \to K_n[1]
$$
dans $D^b_{\textit{Coh}}(\mathcal{O}_Y)$ tel que les pro-systèmes $(K_n)$, $(L_n)$ et $(M_n)$ donnent $Rf_!K$, $Rf_!L$ et $Rf_!M$.
\end{lemma}

\begin{proof}
Choisissons une compactification $j : X \to \overline{X}$ sur $Y$ et notons $\overline{f} : \overline{X} \to Y$ le morphisme structural. Choisissons un système projectif de triangles distingués
$$
\overline{K}_n \to \overline{L}_n \to \overline{M}_n \to \overline{K}_n[1]
$$
dans $D^b_{\textit{Coh}}(\mathcal{O}_{\overline{X}})$ comme dans le lemme \ref{duality-lemma-extension-by-zero-triangle}, correspondant à l'immersion ouverte $j$ et au triangle distingué donné. Prenons $K_n = R\overline{f}_*\overline{K}_n$, et de même pour $L_n$ et $M_n$. Cela convient par la définition même de $Rf_!$.
\end{proof}

\begin{remark}
\label{duality-remark-compose-inverse-systems}
Soit $\mathcal{C}$ une catégorie. Supposons donné un système projectif
$$
\ldots \xrightarrow{\alpha_4} (M_{3, n}) \xrightarrow{\alpha_3} (M_{2, n})
\xrightarrow{\alpha_2} (M_{1, n})
$$
de systèmes projectifs dans la catégorie des pro-objets de $\mathcal{C}$. Autrement dit, les flèches $\alpha_i$ sont des morphismes de pro-objets. D'après Catégories, exemple \ref{categories-example-pro-morphism-inverse-systems}, on peut représenter chaque $\alpha_i$ par un couple $(m_i, a_i)$, où $m_i : \mathbf{N} \to \mathbf{N}$ est une fonction croissante et $a_{i, n} : M_{i, m_i(n)} \to M_{i - 1, n}$ est un morphisme de $\mathcal{C}$ rendant commutatifs les diagrammes
$$
\xymatrix{
\ldots \ar[r] &
M_{i, m_i(3)} \ar[d]^{a_{i, 3}} \ar[r] &
M_{i, m_i(2)} \ar[d]^{a_{i, 2}} \ar[r] &
M_{i, m_i(1)} \ar[d]^{a_{i, 1}} \\
\ldots \ar[r] &
M_{i - 1, 3} \ar[r] &
M_{i - 1, 2} \ar[r] &
M_{i - 1, 1}
}
$$
En remplaçant $m_i(n)$ par $\max(n, m_i(n))$ et en ajustant les morphismes $a_i(n)$ en conséquence (comme dans l'exemple cité), on peut supposer $m_i(n) \geq n$. Dans cette situation, considérons le système projectif
$$
\ldots \to
M_{4, m_4(m_3(m_2(4)))} \to
M_{3, m_3(m_2(3))} \to
M_{2, m_2(2)} \to
M_{1, 1}
$$
de terme général
$$
M_k = M_{k, m_k(m_{k - 1}(\ldots (m_2(k))\ldots))}
$$
Pour tout objet $N$ de $\mathcal{C}$, on a
$$
\colim_i \colim_n \Mor_\mathcal{C}(M_{i, n}, N) =
\colim_k \Mor_\mathcal{C}(M_k, N) 
$$
Nous omettons les détails. Autrement dit, on voit que le système projectif $(M_k)$ possède la propriété
$$
\colim_i \Mor_{\text{Pro-}\mathcal{C}}((M_{i, n}), N) =
\Mor_{\text{Pro-}\mathcal{C}}((M_k), N)
$$
Cette propriété détermine le système projectif $(M_k)$ à pro-isomorphisme près d'après la discussion de Catégories, remarque \ref{categories-remark-pro-category-copresheaves}. On peut ainsi transformer certains systèmes projectifs de $\text{Pro-}\mathcal{C}$ en pro-objets dont les catégories d'indices sont dénombrables.
\end{remark}

\begin{remark}
\label{duality-remark-composition-lower-shriek}
Dans la situation \ref{duality-situation-shriek}, soient $f : X \to Y$ et $g : Y \to Z$ des morphismes composables de $\textit{FTS}_S$. Définissons la composition
$$
Rg_! \circ Rf_! :
D^b_{\textit{Coh}}(\mathcal{O}_X)
\longrightarrow
\text{Pro-}D^b_{\textit{Coh}}(\mathcal{O}_Z)
$$
Par la construction même de $Rf_!$ pour $K$ dans $D^b_{\textit{Coh}}(\mathcal{O}_X)$, l'objet $Rf_!K$ obtenu est la classe de pro-isomorphisme d'un système projectif $(M_n)$ dans $D^b_{\textit{Coh}}(\mathcal{O}_Y)$. Ensuite, puisque $Rg_!$ se construit de la même manière, on voit que
$$
\ldots \to Rg_!M_3 \to Rg_!M_2 \to Rg_!M_1
$$
est un système projectif de $\text{Pro-}D^b_{\textit{Coh}}(\mathcal{O}_Y)$. D'après la discussion de la remarque \ref{duality-remark-compose-inverse-systems}, il existe une unique classe de pro-isomorphisme, que l'on notera $Rg_! Rf_! K$, représentée par des systèmes projectifs dans $D^b_{\textit{Coh}}(\mathcal{O}_Z)$ et caractérisée par l'égalité
$$
\Hom_{\text{Pro-}D^b_{\textit{Coh}}(\mathcal{O}_Z)}(Rg_!Rf_!K, L) =
\colim_n \Hom_{\text{Pro-}D^b_{\textit{Coh}}(\mathcal{O}_Z)}(Rg_!M_n, L)
$$
Nous omettons la discussion nécessaire pour voir que cette construction est fonctorielle en $K$, car cela résultera immédiatement du lemme suivant.
\end{remark}

\begin{lemma}
\label{duality-lemma-composition-lower-shriek}
Dans la situation \ref{duality-situation-shriek}, soient $f : X \to Y$ et $g : Y \to Z$ des morphismes composables de $\textit{FTS}_S$. Avec les notations de la remarque \ref{duality-remark-composition-lower-shriek}, on a
$Rg_! \circ Rf_! = R(g \circ f)_!$.
\end{lemma}

\begin{proof}
D'après la discussion de Catégories, remarque \ref{categories-remark-pro-category-copresheaves}, il suffit de montrer que l'on obtient le même résultat si l'on calcule $\Hom$ à valeurs dans $L$ dans la catégorie $D^b_{\textit{Coh}}(\mathcal{O}_Z)$. Pour ce faire, avec les notations de la remarque \ref{duality-remark-composition-lower-shriek}, on calcule comme suit :
\begin{align*}
\Hom_Z(Rg_!Rf_!K, L)
& =
\colim_n \Hom_Z(Rg_!M_n, L) \\
& =
\colim_n \Hom_Y(M_n, g^!L) \\
& =
\Hom_Y(Rf_!K, g^!L) \\
& =
\Hom_X(K, f^!g^!L) \\
& =
\Hom_X(K, (g \circ f)^!L) \\
& =
\Hom_Z(R(g \circ f)_!K, L)
\end{align*}
La première égalité est la définition de $Rg_!Rf_!K$. La deuxième est la proposition \ref{duality-proposition-duality-compactly-supported} pour $g$. La troisième exprime que $Rf_!K$ est donné par $(M_n)$. La quatrième est la proposition \ref{duality-proposition-duality-compactly-supported} pour $f$. La cinquième est le lemme \ref{duality-lemma-upper-shriek-composition}. La sixième est la proposition \ref{duality-proposition-duality-compactly-supported} pour $g \circ f$.
\end{proof}

\begin{remark}
\label{duality-remark-covariance-open-lower-shriek}
Dans la situation \ref{duality-situation-shriek}, soit $f : X \to Y$ un morphisme de $\textit{FTS}_S$, et soit $U \subset X$ un ouvert. Posons $g = f|_U : U \to Y$. Il existe alors un morphisme canonique
$$
Rg_!(K|_U) \longrightarrow Rf_!K
$$
fonctoriel en $K$ dans $D^b_{\textit{Coh}}(\mathcal{O}_X)$, que l'on peut définir d'au moins trois manières.
\begin{enumerate}
\item Notons $i : U \to X$ le morphisme d'inclusion. On a $Rg_! = Rf_! \circ Ri_!$ d'après le lemme \ref{duality-lemma-composition-lower-shriek}, et l'on peut appliquer $Rf_!$ au morphisme $Ri_!(K|_U) \to K$, qui est un cas particulier de la remarque \ref{duality-remark-covariance-open-j-lower-shriek}.
\item Choisissons une compactification $j : X \to \overline{X}$ de $X$ sur $Y$, de morphisme structural $\overline{f} : \overline{X} \to Y$. Posons $j' = j \circ i : U \to \overline{X}$. On peut utiliser $Rf_! = R\overline{f}_* \circ Rj_!$ et $Rg_! = R\overline{f}_* \circ Rj'_!$, puis appliquer $R\overline{f}_*$ au morphisme $Rj'_!(K|_U) \to Rj_!K$ de la remarque \ref{duality-remark-covariance-open-j-lower-shriek}.
\item On peut utiliser
\begin{align*}
\Hom_{\text{Pro-}D^b_{\textit{Coh}}(\mathcal{O}_Y)}(Rf_!K, L)
& =
\Hom_X(K, f^!L) \\
& \to
\Hom_U(K|_U, f^!L|_U) \\
& =
\Hom_U(K|_U, g^!L) \\
& =
\Hom_{\text{Pro-}D^b_{\textit{Coh}}(\mathcal{O}_Y)}(Rg_!(K|_U), L)
\end{align*}
fonctoriel en $L$ et $K$. Ici, on a appliqué deux fois la proposition \ref{duality-proposition-duality-compactly-supported}, ainsi que la construction des foncteurs image inverse extraordinaire, qui montre que $g^! = i^* \circ f^!$. La fonctorialité en $L$ montre, d'après Catégories, remarque \ref{categories-remark-pro-category-copresheaves}, que l'on obtient un morphisme canonique $Rg_!(K|_U) \to Rf_!K$ dans $\text{Pro-}D^b_{\textit{Coh}}(\mathcal{O}_Y)$, qui est fonctoriel en $K$ par la fonctorialité en $K$ de la flèche ci-dessus.
\end{enumerate}
Chacune de ces trois constructions donne la même flèche ; nous omettons les détails.
\end{remark}

\begin{remark}
\label{duality-remark-covariance-etale-lower-shriek}
Généralisons la covariance de la cohomologie à support compact donnée dans la remarque \ref{duality-remark-covariance-open-lower-shriek} aux morphismes \'etales. Dans la situation \ref{duality-situation-shriek}, supposons en effet donné un diagramme commutatif
$$
\xymatrix{
U \ar[rr]_h \ar[rd]_g & & X \ar[ld]^f \\
& Y
}
$$
de $\textit{FTS}_S$, avec $h$ \'etale. Il existe alors un morphisme canonique
$$
Rg_!(h^*K) \longrightarrow Rf_!K
$$
fonctoriel en $K$ dans $D^b_{\textit{Coh}}(\mathcal{O}_X)$. Nous définissons cette transformation au moyen de la suite de morphismes
\begin{align*}
\Hom_{\text{Pro-}D^b_{\textit{Coh}}(\mathcal{O}_Y)}(Rf_!K, L)
& =
\Hom_X(K, f^!L) \\
& \to
\Hom_U(h^*K, h^*(f^!L)) \\
& =
\Hom_U(h^*K, h^!f^!L) \\
& =
\Hom_U(h^*K, g^!L) \\
& =
\Hom_{\text{Pro-}D^b_{\textit{Coh}}(\mathcal{O}_Y)}(Rg_!(h^*K), L)
\end{align*}
fonctorielle en $L$ et $K$. Ici, on a appliqué deux fois la proposition \ref{duality-proposition-duality-compactly-supported}, utilisé l'égalité $h^* = h^!$ du lemme \ref{duality-lemma-shriek-etale} et l'égalité $h^! \circ f^! = g^!$ du lemme \ref{duality-lemma-upper-shriek-composition}. La fonctorialité en $L$ montre, d'après Catégories, remarque \ref{categories-remark-pro-category-copresheaves}, que l'on obtient un morphisme canonique $Rg_!(h^*K) \to Rf_!K$ dans $\text{Pro-}D^b_{\textit{Coh}}(\mathcal{O}_Y)$, qui est fonctoriel en $K$ par la fonctorialité en $K$ de la flèche ci-dessus.
\end{remark}

\begin{remark}
\label{duality-remark-covariance-lower-shriek}
Dans les remarques \ref{duality-remark-covariance-open-lower-shriek} et \ref{duality-remark-covariance-etale-lower-shriek}, nous avons vu que la construction de la cohomologie à support compact est covariante relativement aux immersions ouvertes et aux morphismes \'etales. En fait, le bon degré de généralité est le suivant : étant donné un diagramme commutatif
$$
\xymatrix{
U \ar[rr]_h \ar[rd]_g & & X \ar[ld]^f \\
& Y
}
$$
de $\textit{FTS}_S$, avec $h$ plat et quasi-fini, il existe une transformation canonique
$$
Rg_! \circ h^* \longrightarrow Rf_!
$$
Comme dans la remarque \ref{duality-remark-covariance-etale-lower-shriek}, ce morphisme peut se construire à l'aide d'une transformation de foncteurs $h^* \to h^!$ sur $D^+_{\textit{Coh}}(\mathcal{O}_X)$. Rappelons que $h^!K = h^*K \otimes \omega_{U/X}$, où $\omega_{U/X} = h^!\mathcal{O}_X$ est le faisceau dualisant relatif du morphisme plat quasi-fini $h$ (voir les lemmes \ref{duality-lemma-perfect-comparison-shriek} et \ref{duality-lemma-flat-quasi-finite-shriek}). Rappelons que $\omega_{U/X}$ est le même que le module dualisant relatif qui sera construit dans Discriminants, remarque \ref{discriminant-remark-relative-dualizing-for-quasi-finite}, d'après Discriminants, lemme \ref{discriminant-lemma-compare-dualizing}. On peut donc utiliser l'élément de trace
$\tau_{U/X} : \mathcal{O}_U \to \omega_{U/X}$
qui sera construit dans Discriminants, remarque \ref{discriminant-remark-relative-dualizing-for-flat-quasi-finite}, pour définir notre transformation. Si nous en avons un jour besoin, nous formulerons et démontrerons précisément le résultat ici.
\end{remark}







\section{Dualité pour la cohomologie à support compact}
\label{duality-section-duality-compactly-supported}

\noindent
Soit $k$ un corps. Soit $U$ un schéma séparé de type fini sur $k$. Soit $K$ un objet de $D^b_{\textit{Coh}}(\mathcal{O}_U)$. Définissons la cohomologie à support compact $H^i_c(U, K)$ de $K$ comme suit. Choisissons une immersion ouverte $j : U \to X$ dans un schéma propre sur $k$ et un système de Deligne $(K_n)$ pour $j : U \to X$ dont la restriction à $U$ est constante de valeur $K$. Posons alors
$$
H^i_c(U, K) = \lim H^i(X, K_n)
$$
On le considère comme un $k$-espace vectoriel topologique muni de la topologie limite (voir Compléments sur l'algèbre, section \ref{more-algebra-section-topological-ring}). Plusieurs remarques s'imposent.

\medskip\noindent
Premièrement, cette définition est indépendante du choix de $X$ et de $(K_n)$. En effet, si $p : U \to \Spec(k)$ désigne le morphisme structural, on sait déjà que $Rp_!K = (R\Gamma(X, K_n))$ est bien défini à pro-isomorphisme près dans $D(k)$ ; il en est donc de même de la limite qui définit $H^i_c(U, K)$.

\medskip\noindent
Deuxièmement, il pourrait paraître plus naturel d'utiliser l'expression
$$
H^i(R\lim R\Gamma(X, K_n)) = R\Gamma(X, R\lim K_n)
$$
mais cela donnerait le même résultat : comme les $k$-espaces vectoriels $H^j(X, K_n)$ sont de dimension finie, ces systèmes projectifs satisfont la condition de Mittag-Leffler, et les termes $R^1\lim$ de Cohomologie, lemme \ref{cohomology-lemma-RGamma-commutes-with-Rlim}, sont donc nuls.

\medskip\noindent
Si $U' \subset U$ est un sous-schéma ouvert, il existe un morphisme canonique
$$
H^i_c(U', K|_{U'}) \longrightarrow H^i_c(U, K)
$$
fonctoriel en $K$ dans $D^b_{\textit{Coh}}(\mathcal{O}_U)$. Voir par exemple la remarque \ref{duality-remark-covariance-open-lower-shriek}. En fait, en utilisant la remarque \ref{duality-remark-covariance-etale-lower-shriek}, on voit que, plus généralement, un tel morphisme existe pour un morphisme \'etale $U' \to U$ de schémas séparés de type fini sur $k$.

\medskip\noindent
Si $V$ est un $k$-espace vectoriel, on munit $\Hom_k(V, k)$ d'une topologie comme suit : écrivons $V = \bigcup V_i$ comme la réunion filtrante de ses sous-$k$-espaces vectoriels de dimension finie et utilisons la topologie limite sur $\Hom_k(V, k) = \lim \Hom_k(V_i, k)$. Si $\dim_k V < \infty$, la topologie de $\Hom_k(V, k)$ est discrète. Plus généralement, si $V = \colim_n V_n$ est écrit comme une colimite filtrante d'espaces vectoriels de dimension finie, alors $\Hom_k(V, k) = \lim \Hom_k(V_n, k)$ comme espaces vectoriels topologiques.

\begin{lemma}
\label{duality-lemma-duality-compact-support}
Soit $p : U \to \Spec(k)$ séparé et de type fini, où $k$ est un corps. Soit $\omega_{U/k}^\bullet = p^!\mathcal{O}_{\Spec(k)}$. Il existe des isomorphismes canoniques
$$
\Hom_k(H^i(U, K), k) =
H^{-i}_c(U, R\SheafHom_{\mathcal{O}_U}(K, \omega_{U/k}^\bullet))
$$
de $k$-espaces vectoriels topologiques, fonctoriels en $K$ dans $D^b_{\textit{Coh}}(\mathcal{O}_U)$.
\end{lemma}

\begin{proof}
Choisissons une compactification $j : U \to X$ sur $k$. Soit $\mathcal{I} \subset \mathcal{O}_X$ un faisceau quasi-cohérent d'idéaux tel que $V(\mathcal{I}) = X \setminus U$. D'après Catégories dérivées de schémas, proposition \ref{perfect-proposition-DCoh}, on peut choisir $M \in D^b_{\textit{Coh}}(\mathcal{O}_X)$ avec $K = M|_U$. On a
\begingroup\small
$$
H^i(U, K) =
\Ext^i_U(\mathcal{O}_U, M|_U) =
\colim \Ext^i_X(\mathcal{I}^n, M) =
\colim H^i(X, R\SheafHom_{\mathcal{O}_X}(\mathcal{I}^n, M))
$$
\endgroup
d'après le lemme \ref{duality-lemma-lift-map}. Puisque $\mathcal{I}^n$ est un $\mathcal{O}_X$-module cohérent, on a $\mathcal{I}^n$ dans $D^-_{\textit{Coh}}(\mathcal{O}_X)$ ; ainsi $R\SheafHom_{\mathcal{O}_X}(\mathcal{I}^n, M)$ appartient à $D^+_{\textit{Coh}}(\mathcal{O}_X)$ d'après Catégories dérivées de schémas, lemme \ref{perfect-lemma-coherent-internal-hom}.

\medskip\noindent
Soit $\omega_{X/k}^\bullet = q^!\mathcal{O}_{\Spec(k)}$, où $q : X \to \Spec(k)$ est le morphisme structural, voir la section \ref{duality-section-duality-proper-over-field}. On trouve que
\begin{align*}
\Hom_k(
&
H^i(X, R\SheafHom_{\mathcal{O}_X}(\mathcal{I}^n, M)), k) \\
& =
\Ext^{-i}_X(R\SheafHom_{\mathcal{O}_X}(\mathcal{I}^n, M),
\omega_{X/k}^\bullet) \\
& =
H^{-i}(X, R\SheafHom_{\mathcal{O}_X}(R\SheafHom_{\mathcal{O}_X}(
\mathcal{I}^n, M), \omega_{X/k}^\bullet))
\end{align*}
d'après le lemme \ref{duality-lemma-duality-proper-over-field}. D'après l'assertion (1) du lemme \ref{duality-lemma-internal-hom-evaluate-isom}, le morphisme canonique
$$
R\SheafHom_{\mathcal{O}_X}(M, \omega_{X/k}^\bullet)
\otimes_{\mathcal{O}_X}^\mathbf{L} \mathcal{I}^n
\longrightarrow
R\SheafHom_{\mathcal{O}_X}(R\SheafHom_{\mathcal{O}_X}(
\mathcal{I}^n, M), \omega_{X/k}^\bullet)
$$
est un isomorphisme. Remarquons que
$\omega^\bullet_{U/k} = \omega^\bullet_{X/k}|_U$
parce que $p^!$ se construit comme $q^!$ composé avec la restriction à $U$. Par conséquent, $R\SheafHom_{\mathcal{O}_X}(M, \omega_{X/k}^\bullet)$ est un objet de $D^b_{\textit{Coh}}(\mathcal{O}_X)$ dont la restriction est $R\SheafHom_{\mathcal{O}_U}(K, \omega_{U/k}^\bullet)$ sur $U$. D'après le lemme \ref{duality-lemma-tensoring-Deligne-system}, on conclut donc que
$$
\lim H^{-i}(X, R\SheafHom_{\mathcal{O}_X}(M, \omega_{X/k}^\bullet)
\otimes_{\mathcal{O}_X}^\mathbf{L} \mathcal{I}^n)
$$
est un représentant du membre de droite de l'égalité du lemme. En combinant tous les isomorphismes ainsi obtenus, on obtient l'isomorphisme du lemme.
\end{proof}

\begin{lemma}
\label{duality-lemma-duality-compact-support-restrict-open}
Avec les notations du lemme \ref{duality-lemma-duality-compact-support}, supposons que $U' \subset U$ soit un sous-schéma ouvert. Alors le diagramme
$$
\xymatrix{
\Hom_k(H^i(U, K), k) \ar[rr] & &
H^{-i}_c(U, R\SheafHom_{\mathcal{O}_U}(K, \omega_{U/k}^\bullet)) \\
\Hom_k(H^i(U', K|_{U'}), k) \ar[rr] \ar[u] & &
H^{-i}_c(U', R\SheafHom_{\mathcal{O}_{U'}}(K, \omega_{U'/k}^\bullet)) \ar[u]
}
$$
est commutatif. Ici, les flèches horizontales sont les isomorphismes du lemme \ref{duality-lemma-duality-compact-support}, la flèche verticale de gauche est la transposée du morphisme de restriction
\par\noindent\hfil
$H^i(U, K) \to H^i(U', K|_{U'})$,
\hfil\par\noindent
et la flèche verticale de droite est celle de la remarque \ref{duality-remark-covariance-open-lower-shriek} (voir la discussion précédant le lemme).
\end{lemma}

\begin{proof}
Nous recommandons vivement au lecteur de sauter cette démonstration. Choisissons $X$ et $M$ comme dans la démonstration du lemme \ref{duality-lemma-duality-compact-support}. Nous omettrons l'indice $\mathcal{O}_X$ de $R\SheafHom$ et $\otimes^\mathbf{L}$. Écrivons
$$
H^i(U, K) = \colim H^i(X, R\SheafHom(\mathcal{I}^n, M))
$$
et
$$
H^i(U', K|_{U'}) = \colim H^i(X, R\SheafHom((\mathcal{I}')^n, M))
$$
comme dans la démonstration du lemme \ref{duality-lemma-duality-compact-support}, où l'on choisit $\mathcal{I}' \subset \mathcal{I}$ comme dans la discussion de la remarque \ref{duality-remark-covariance-open-j-lower-shriek}, de sorte que le morphisme $H^i(U, K) \to H^i(U', K|_{U'})$ soit induit par les morphismes
$(\mathcal{I}')^n \to \mathcal{I}^n$.
De même, écrivons
$$
H^i_c(U, R\SheafHom(K, \omega_{U/k}^\bullet)) =
\lim H^i(X,  R\SheafHom(M, \omega_{X/k}^\bullet)
\otimes^\mathbf{L} \mathcal{I}^n)
$$
et
$$
H^i_c(U', R\SheafHom(K|_{U'}, \omega_{U'/k}^\bullet)) =
\lim H^i(X,  R\SheafHom(M, \omega_{X/k}^\bullet)
\otimes^\mathbf{L} (\mathcal{I}')^n)
$$
de sorte que la flèche $H^i_c(U', R\SheafHom(K|_{U'}, \omega_{U'/k}^\bullet))
\to H^i_c(U, R\SheafHom(K, \omega_{U/k}^\bullet))$ se déduise de façon analogue des morphismes $(\mathcal{I}')^n \to \mathcal{I}^n$. Les diagrammes
$$
\xymatrix{
R\SheafHom(M, \omega_{X/k}^\bullet)
\otimes^\mathbf{L} \mathcal{I}^n
\ar[rr] & &
R\SheafHom(R\SheafHom(\mathcal{I}^n, M), \omega_{X/k}^\bullet) \\
R\SheafHom(M, \omega_{X/k}^\bullet) \otimes^\mathbf{L} (\mathcal{I}')^n
\ar[rr] \ar[u] & &
R\SheafHom(R\SheafHom((\mathcal{I}')^n, M), \omega_{X/k}^\bullet) \ar[u]
}
$$
sont commutatifs parce que la construction des flèches horizontales dans Cohomologie, lemme \ref{cohomology-lemma-internal-hom-evaluate}, est fonctorielle dans les trois termes. On est ainsi finalement ramené à l'assertion que les diagrammes
\begingroup\small
$$
\xymatrix{
\Hom_k(H^i(X, R\SheafHom(\mathcal{I}^n, M)), k) \ar[r] &
H^{-i}(X, R\SheafHom(R\SheafHom(
\mathcal{I}^n, M), \omega_{X/k}^\bullet)) \\
\Hom_k(H^i(X, R\SheafHom((\mathcal{I}')^n, M)), k) \ar[r] \ar[u] &
H^{-i}(X, R\SheafHom(R\SheafHom(
(\mathcal{I}')^n, M), \omega_{X/k}^\bullet)) \ar[u]
}
$$
\endgroup
sont commutatifs. Cela vaut parce que l'isomorphisme de dualité
$$
\Hom_k(H^i(X, L), k) = \Ext^{-i}_X(L, \omega_{X/k}^\bullet) =
H^{-i}(X, R\SheafHom(L, \omega_{X/k}^\bullet))
$$
est fonctoriel en $L$ dans $D_\QCoh(\mathcal{O}_X)$.
\end{proof}

\begin{lemma}
\label{duality-lemma-h0-compactly-supported}
Soit $X$ un schéma propre sur un corps $k$. Soit $K \in D^b_{\textit{Coh}}(\mathcal{O}_X)$ avec $H^i(K) = 0$ pour $i < 0$. Posons $\mathcal{F} = H^0(K)$. Soit $Z \subset X$ un fermé dont le complémentaire est $U = X \setminus U$. Alors le sous-espace
$$
H^0_c(U, K|_U) \subset H^0(X, \mathcal{F})
$$
est constitué des sections globales de $\mathcal{F}$ qui s'annulent dans un voisinage ouvert de $Z$.
\end{lemma}

\begin{proof}
Considérons le morphisme
\par\noindent\hfil
$H^0_c(U, K|_U) \to H^0_X(X, K) = H^0(X, K) = H^0(X, \mathcal{F})$
\hfil\par\noindent
de la remarque \ref{duality-remark-covariance-open-lower-shriek}. Pour l'étudier, représentons $K$ par un complexe borné $\mathcal{F}^\bullet$ avec $\mathcal{F}^i = 0$ pour $i < 0$. On a alors, par définition,
$$
H^0_c(U, K|_U) = \lim H^0(X, \mathcal{I}^n\mathcal{F}^\bullet)
= \lim \Ker(
H^0(X, \mathcal{I}^n\mathcal{F}^0) \to
H^0(X, \mathcal{I}^n\mathcal{F}^1))
$$
Par Artin–Rees (Cohomologie des schémas, lemme \ref{coherent-lemma-Artin-Rees}), cela coïncide avec
\par\noindent\hfil
$\lim H^0(X, \mathcal{I}^n\mathcal{F})$.
\hfil\par\noindent
Ainsi, la flèche $H^0_c(U, K|_U) \to H^0(X, \mathcal{F})$ est injective, et son image est formée des sections globales de $\mathcal{F}$ qui sont contenues dans le sous-faisceau $\mathcal{I}^n\mathcal{F}$ pour tout $n$. Leur caractérisation comme les sections qui s'annulent dans un voisinage de $Z$ résulte du théorème d'intersection de Krull (Algèbre, lemme \ref{algebra-lemma-intersect-powers-ideal-module-zero}), en passant aux germes de $\mathcal{F}$. Voir la discussion d'Algèbre, remarque \ref{algebra-remark-intersection-powers-ideal}, pour le cas des fonctions.
\end{proof}




\section{Le théorème de Lichtenbaum}
\label{duality-section-lichtenbaum}

\noindent
Le théorème ci-dessous a été conjecturé par Lichtenbaum et démontré par Grothendieck (voir \cite{Hartshorne-local-cohomology}). Kleiman en donne une très belle démonstration dans \cite{Kleiman-Lichtenbaum}. Une généralisation du théorème au cas de la cohomologie à supports se trouve dans \cite{Lyubeznik-Lichtenbaum}. La partie la plus intéressante de l'argument figure dans la démonstration du lemme suivant.

\begin{lemma}
\label{duality-lemma-lichtenbaum}
Soit $U$ une variété. Soit $\mathcal{F}$ un $\mathcal{O}_U$-module cohérent. Si $H^d(U, \mathcal{F})$ est non nul, alors $\dim(U) \geq d$ et, en cas d'égalité, $U$ est propre.
\end{lemma}

\begin{proof}
Le théorème d'annulation de Grothendieck dans Cohomologie, proposition \ref{cohomology-proposition-vanishing-Noetherian}, permet de conclure que $\dim(U) \geq d$. Supposons $\dim(U) = d$. Choisissons une compactification $U \to X$ telle que $U$ soit dense dans $X$. (Cela est possible d'après Compléments sur la platitude, théorème \ref{flat-theorem-nagata}, et lemme \ref{flat-lemma-compactifyable}.) Après avoir remplacé $X$ par sa réduction, on trouve que $X$ est une variété propre de dimension $d$, et l'on voit que $U$ est propre si et seulement si $U = X$. Posons $Z = X \setminus U$. Nous montrerons que $H^d(U, \mathcal{F})$ est nul si $Z$ est non vide.

\medskip\noindent
Choisissons un $\mathcal{O}_X$-module cohérent $\mathcal{G}$ dont la restriction à $U$ est $\mathcal{F}$, voir Propriétés, lemme \ref{properties-lemma-lift-finite-presentation}. Notons $\omega_X^\bullet$ le complexe dualisant de $X$ comme dans la section \ref{duality-section-duality-proper-over-field}. Posons $\omega_U^\bullet = \omega_X^\bullet|_U$. Alors $H^d(U, \mathcal{F})$ est dual de
$$
H^{-d}_c(U, R\SheafHom_{\mathcal{O}_U}(\mathcal{F}, \omega_U^\bullet))
$$
d'après le lemme \ref{duality-lemma-duality-compact-support}. D'après le lemme \ref{duality-lemma-duality-proper-over-field}, on voit que les faisceaux de cohomologie de $\omega_X^\bullet$ sont nuls en degrés $< -d$ et que $H^{-d}(\omega_X^\bullet) = \omega_X$ est un $\mathcal{O}_X$-module cohérent qui est $(S_2)$ et dont le support est $X$. En particulier, $\omega_X$ est sans torsion, voir Diviseurs, lemme \ref{divisors-lemma-torsion-free-finite-noetherian-domain}. On voit donc que le faisceau de cohomologie
$$
H^{-d}(R\SheafHom_{\mathcal{O}_X}(\mathcal{G}, \omega_X^\bullet)) =
\SheafHom(\mathcal{G}, \omega_X)
$$
est sans torsion, voir Diviseurs, lemme \ref{divisors-lemma-hom-into-torsion-free}. Par conséquent, ce faisceau ne possède aucune section non nulle qui s'annule sur un ouvert non vide de $X$ (une telle section serait de torsion). Le lemme \ref{duality-lemma-h0-compactly-supported} montre donc que
$H^{-d}_c(U, R\SheafHom_{\mathcal{O}_U}(\mathcal{F}, \omega_U^\bullet))$
est nul, et ainsi $H^d(U, \mathcal{F})$ est nul, comme souhaité.
\end{proof}

\begin{theorem}
\label{duality-theorem-lichtenbaum}
Soit $X$ un schéma non vide, séparé et de type fini sur un corps $k$. Soit $d = \dim(X)$. Les conditions suivantes sont équivalentes :
\begin{enumerate}
\item $H^d(X, \mathcal{F}) = 0$ pour tout $\mathcal{O}_X$-module cohérent $\mathcal{F}$ sur $X$ ;
\item $H^d(X, \mathcal{F}) = 0$ pour tout $\mathcal{O}_X$-module quasi-cohérent $\mathcal{F}$ sur $X$ ; et
\item aucune composante irréductible $X' \subset X$ de dimension $d$ n'est propre sur $k$.
\end{enumerate}
\end{theorem}

\begin{proof}
Supposons qu'il existe une composante irréductible $X' \subset X$ (que l'on considère comme un sous-schéma fermé intègre) qui soit propre et de dimension $d$. Soit $\omega_{X'}$ un module dualisant de $X'$ sur $k$, voir le lemme \ref{duality-lemma-duality-proper-over-field}. Alors $H^d(X', \omega_{X'})$ est non nul puisqu'il est dual de $H^0(X', \mathcal{O}_{X'})$ d'après le lemme. On voit donc que $H^d(X, \omega_{X'}) = H^d(X', \omega_{X'})$ est non nul, et l'on conclut que (1) ne vaut pas. On voit ainsi que (1) implique (3).

\medskip\noindent
Démontrons que (3) implique (1). Soit $\mathcal{F}$ un $\mathcal{O}_X$-module cohérent tel que $H^d(X, \mathcal{F})$ soit non nul. Choisissons une filtration
$$
0 = \mathcal{F}_0 \subset \mathcal{F}_1 \subset
\ldots \subset \mathcal{F}_m = \mathcal{F}
$$
comme dans Cohomologie des schémas, lemme \ref{coherent-lemma-coherent-filter}. On obtient des suites exactes
$$
H^d(X, \mathcal{F}_i) \to H^d(X, \mathcal{F}_{i + 1}) \to
H^d(X, \mathcal{F}_{i + 1}/\mathcal{F}_i)
$$
Il existe donc un $i \in \{1, \ldots, m\}$ tel que $H^d(X, \mathcal{F}_{i + 1}/\mathcal{F}_i)$ soit non nul. Par le choix de la filtration, cela signifie qu'il existe un sous-schéma fermé intègre $Z \subset X$ et un faisceau cohérent non nul d'idéaux $\mathcal{I} \subset \mathcal{O}_Z$ tels que $H^d(Z, \mathcal{I})$ soit non nul. D'après le lemme \ref{duality-lemma-lichtenbaum}, on conclut que $\dim(Z) = d$ et que $Z$ est propre sur $k$, en contradiction avec (3). Ainsi, (3) implique (1).

\medskip\noindent
Enfin, montrons que (1) et (2) sont équivalentes pour tout schéma noethérien $X$. En effet, (2) implique trivialement (1). Réciproquement, supposons (1) et soit $\mathcal{F}$ un $\mathcal{O}_X$-module quasi-cohérent. On peut alors écrire $\mathcal{F} = \colim \mathcal{F}_i$ comme la colimite filtrante de ses sous-modules cohérents, voir Propriétés, lemme \ref{properties-lemma-quasi-coherent-colimit-finite-type}. On a alors $H^d(X, \mathcal{F}) = \colim H^d(X, \mathcal{F}_i) = 0$ d'après Cohomologie, lemme \ref{cohomology-lemma-quasi-separated-cohomology-colimit}. Ainsi, (2) est vraie.
\end{proof}










% OK, start here.
%

\chapter{Discriminants et différentes}



\phantomsection
\label{discriminant-section-phantom}


\section{Introduction}
\label{discriminant-section-introduction}

\noindent
Dans ce chapitre, nous étudions la différente et le discriminant
des morphismes localement quasi-finis de schémas.
On pourra consulter \cite{Kunz} pour une partie de ces résultats.

\medskip\noindent
Étant donné un morphisme quasi-fini $f : Y \to X$ de schémas noethériens,
on dispose d'un module dualisant relatif $\omega_{Y/X}$.
Dans la section \ref{discriminant-section-quasi-finite-dualizing},
nous construisons ce module directement, à l'aide
du théorème principal de Zariski et de méthodes de localisation étale.
Sa propriété essentielle est la suivante : étant donné un diagramme
$$
\xymatrix{
Y' \ar[d]_{f'} \ar[r]_{g'} & Y \ar[d]^f \\
X' \ar[r]^g & X
}
$$
où $g : X' \to X$ est plat, $Y' \subset X' \times_X Y$ est ouvert et
$f' : Y' \to X'$ est fini, il existe un isomorphisme canonique
$$
f'_*(g')^*\omega_{Y/X} =
\SheafHom_{\mathcal{O}_{X'}}(f'_*\mathcal{O}_{Y'}, \mathcal{O}_{X'})
$$
de faisceaux de $f'_*\mathcal{O}_{Y'}$-modules. Dans la
section \ref{discriminant-section-quasi-finite-traces}, nous montrons que,
si $f$ est plat, il existe une section globale canonique
$\tau_{Y/X} \in H^0(Y, \omega_{Y/X})$ telle que, pour tout diagramme
commutatif comme ci-dessus, l'isomorphisme envoie $(g')^*\tau_{Y/X}$ sur l'homomorphisme trace
de la section \ref{discriminant-section-discriminant}
associée au morphisme fini localement libre $f'$.
Dans la section \ref{discriminant-section-different},
nous définissons la différente d'un morphisme plat quasi-fini
de schémas noethériens comme l'annulateur du
conoyau de $\tau_{Y/X} : \mathcal{O}_X \to \omega_{Y/X}$.

\medskip\noindent
L'objectif principal de ce chapitre est de montrer que, pour
$f$ quasi-fini syntomique\footnote{Autrement dit plat et localement d'intersection complète.}, la
différente coïncide avec la différente de Kähler.
La différente de Kähler est l'idéal de Fitting d'ordre zéro de $\Omega_{Y/X}$, voir
la section \ref{discriminant-section-kahler-different}.
Cette égalité n'est pas évidente ; nous utilisons un argument élégant
dû à Tate, voir la section \ref{discriminant-section-formula-different}.
Nous étudions également, chemin faisant, la différente de Noether
et la différente de Dedekind.

\medskip\noindent
Ce n'est qu'à la fin de ce chapitre, voir
les sections \ref{discriminant-section-comparison} et \ref{discriminant-section-gorenstein-lci},
que nous établissons le lien avec les résultats plus avancés
sur la dualité des schémas.








\section{Modules dualisants des homomorphismes quasi-finis d'anneaux}
\label{discriminant-section-quasi-finite-dualizing}

\noindent
Soit $A \to B$ un homomorphisme quasi-fini d'anneaux noethériens. D'après
le théorème principal de Zariski
(Algèbre, lemme \ref{algebra-lemma-quasi-finite-open-integral-closure}),
il existe une factorisation $A \to B' \to B$ où
$A \to B'$ est fini et $B' \to B$ induit une immersion ouverte des spectres.
On pose
\begin{equation}
\label{discriminant-equation-dualizing}
\omega_{B/A} = \Hom_A(B', A) \otimes_{B'} B
\end{equation}
dans cette situation. On peut considérer ce module comme une sorte de module
dualisant relatif, voir les lemmes \ref{discriminant-lemma-compare-dualizing} et
\ref{discriminant-lemma-compare-dualizing-algebraic}.
Dans cette section, nous montrerons par des méthodes élémentaires d'algèbre commutative
que $\omega_{B/A}$ est indépendant du choix de la factorisation
et que la formation de $\omega_{B/A}$ commute au changement de base plat.
Pour établir l'indépendance à l'égard de la factorisation, nous comparons deux
factorisations données.

\begin{lemma}
\label{discriminant-lemma-dominate-factorizations}
Soit $A \to B$ un homomorphisme quasi-fini d'anneaux. Étant donné deux factorisations
$A \to B' \to B$ et $A \to B'' \to B$ où
$A \to B'$ et $A \to B''$ sont finis et $\Spec(B) \to \Spec(B')$
et $\Spec(B) \to \Spec(B'')$ sont des immersions ouvertes, il existe
une $A$-sous-algèbre $B''' \subset B$ finie sur $A$ telle que
$\Spec(B) \to \Spec(B''')$ soit une immersion ouverte et que $B' \to B$ et
$B'' \to B$ se factorisent par $B'''$.
\end{lemma}

\begin{proof}
Soit $B''' \subset B$ la $A$-sous-algèbre engendrée par les images
de $B' \to B$ et $B'' \to B$. Comme $B'$ et $B''$ sont chacun engendrés
par un nombre fini d'éléments entiers sur $A$, on voit que $B'''$ est
engendrée par un nombre fini d'éléments entiers sur $A$, et l'on en déduit
que $B'''$ est finie sur $A$
(Algèbre, lemme \ref{algebra-lemma-characterize-finite-in-terms-of-integral}).
Considérons les applications
$$
B = B' \otimes_{B'} B \to B''' \otimes_{B'} B \to B \otimes_{B'} B = B
$$
La dernière égalité tient au fait que $\Spec(B) \to \Spec(B')$ est une
immersion ouverte (donc un monomorphisme). La seconde flèche est injective
puisque $B' \to B$ est plat. Les deux flèches sont donc des isomorphismes.
Cela signifie que
$$
\xymatrix{
\Spec(B''') \ar[d] & \Spec(B) \ar[d] \ar[l] \\
\Spec(B') & \Spec(B) \ar[l]
}
$$
est cartésien. Comme le changement de base d'une immersion ouverte est une
immersion ouverte, on conclut.
\end{proof}

\begin{lemma}
\label{discriminant-lemma-dualizing-well-defined}
Le module (\ref{discriminant-equation-dualizing}) est bien défini, c'est-à-dire
indépendant du choix de la factorisation.
\end{lemma}

\begin{proof}
Soient $B', B'', B'''$ comme dans le lemme \ref{discriminant-lemma-dominate-factorizations}.
On obtient une application canonique
$$
\omega''' = \Hom_A(B''', A) \otimes_{B'''} B \longrightarrow
\Hom_A(B', A) \otimes_{B'} B = \omega'
$$
et une application analogue faisant intervenir $B''$. Il suffit de montrer que ces applications
sont des isomorphismes. Soit $g \in B'$ un élément tel que
$B'_g \to B_g$ soit un isomorphisme, de sorte que $B'_g \to (B''')_g \to B_g$
soient des isomorphismes. Il suffit de montrer que $(\omega''')_g \to \omega'_g$
est un isomorphisme. Le noyau et le conoyau de l'homomorphisme d'anneaux $B' \to B'''$
sont des $A$-modules de type fini et de torsion par rapport aux puissances de $g$.
Ils sont donc annulés par une puissance de $g$.
Le résultat s'en déduit aisément.
\end{proof}

\begin{lemma}
\label{discriminant-lemma-localize-dualizing}
Soit $A \to B$ un homomorphisme quasi-fini d'anneaux noethériens.
\begin{enumerate}
\item Si $A \to B$ se factorise en $A \to A_f \to B$ pour un $f \in A$,
alors $\omega_{B/A} = \omega_{B/A_f}$.
\item Si $g \in B$, alors $(\omega_{B/A})_g = \omega_{B_g/A}$.
\item Si $f \in A$, alors $\omega_{B_f/A_f} = (\omega_{B/A})_f$.
\end{enumerate}
\end{lemma}

\begin{proof}
Soit $A \to B' \to B$ une factorisation où $A \to B'$ est fini et
$\Spec(B) \to \Spec(B')$ est une immersion ouverte. Dans le cas (1), on peut utiliser
la factorisation $A_f \to B'_f \to B$ pour calculer $\omega_{B/A_f}$
et appliquer Algèbre, lemme \ref{algebra-lemma-hom-from-finitely-presented}.
Dans le cas (2), on utilise la factorisation $A \to B' \to B_g$ pour obtenir le résultat.
L'assertion (3) résulte de (1) et (2).
\end{proof}

\noindent
Soit $A \to B$ un homomorphisme quasi-fini d'anneaux noethériens, soit
$A \to A_1$ un homomorphisme quelconque d'anneaux noethériens, et posons
$B_1 = B \otimes_A A_1$. On obtient un diagramme cocartésien
$$
\xymatrix{
B \ar[r] & B_1 \\
A \ar[u] \ar[r] & A_1 \ar[u]
}
$$
Remarquons que $A_1 \to B_1$ est également quasi-fini (Algèbre, lemme
\ref{algebra-lemma-quasi-finite-base-change}).
Dans cette situation, nous allons définir une application canonique
$B$-linéaire de changement de base
\begin{equation}
\label{discriminant-equation-bc-dualizing}
\omega_{B/A} \longrightarrow \omega_{B_1/A_1}
\end{equation}
À cet effet, on choisit une factorisation $A \to B' \to B$ comme dans la construction
de $\omega_{B/A}$. Alors $B'_1 = B' \otimes_A A_1$ est fini sur $A_1$
et l'on peut utiliser la factorisation $A_1 \to B'_1 \to B_1$ pour construire
$\omega_{B_1/A_1}$. Il s'agit donc de construire une application
$$
\Hom_A(B', A) \otimes_{B'} B
\longrightarrow
\Hom_{A_1}(B' \otimes_A A_1, A_1) \otimes_{B'_1} B_1
$$
Il suffit ainsi de construire une application $B'$-linéaire
$\Hom_A(B', A) \to \Hom_{A_1}(B' \otimes_A A_1, A_1)$,
que l'on notera $\varphi \mapsto \varphi_1$.
Étant donné une application $A$-linéaire $\varphi : B' \to A$,
on prend pour $\varphi_1$ l'application définie par
$\varphi_1(b' \otimes a_1) = \varphi(b')a_1$.
Elle est manifestement $A_1$-linéaire, ce qui achève la construction.

\begin{lemma}
\label{discriminant-lemma-bc-map-dualizing}
L'application de changement de base (\ref{discriminant-equation-bc-dualizing})
est indépendante du choix de la
factorisation $A \to B' \to B$. Étant donné des homomorphismes d'anneaux $A \to A_1 \to A_2$,
la composée des applications de changement de base pour $A \to A_1$ et $A_1 \to A_2$
est l'application de changement de base pour $A \to A_2$.
\end{lemma}

\begin{proof}
Omise. Indication : raisonner exactement comme dans le
lemme \ref{discriminant-lemma-dualizing-well-defined},
en utilisant le lemme \ref{discriminant-lemma-dominate-factorizations}.
\end{proof}

\begin{lemma}
\label{discriminant-lemma-dualizing-flat-base-change}
Si $A \to A_1$ est plat, alors
l'application de changement de base (\ref{discriminant-equation-bc-dualizing}) induit un isomorphisme
$\omega_{B/A} \otimes_B B_1 \to \omega_{B_1/A_1}$.
\end{lemma}

\begin{proof}
Supposons $A \to A_1$ plat. Par construction de $\omega_{B/A}$, on peut
supposer que $A \to B$ est fini. Alors $\omega_{B/A} = \Hom_A(B, A)$ et
$\omega_{B_1/A_1} = \Hom_{A_1}(B_1, A_1)$. Puisque $B_1 = B \otimes_A A_1$,
le résultat découle de Compléments d'algèbre, lemme
\ref{more-algebra-lemma-pseudo-coherence-and-base-change-ext}.
\end{proof}

\begin{lemma}
\label{discriminant-lemma-dualizing-composition}
Soient $A \to B \to C$ des homomorphismes quasi-finis d'anneaux noethériens.
Il existe une application canonique
$\omega_{B/A} \otimes_B \omega_{C/B} \to \omega_{C/A}$.
\end{lemma}

\begin{proof}
Choisissons $A \to B' \to B$ où $A \to B'$ est fini et
$\Spec(B) \to \Spec(B')$ est une immersion ouverte. Alors
$B' \to C$ est également quasi-fini. Choisissons $B' \to C' \to C$
où $B' \to C'$ est fini et $\Spec(C) \to \Spec(C')$ est une
immersion ouverte. La source de la flèche est alors
$$
\Hom_A(B', A) \otimes_{B'} B \otimes_B
\Hom_B(B \otimes_{B'} C', B) \otimes_{B \otimes_{B'} C'} C
$$
qui est égale à
$$
\Hom_A(B', A) \otimes_{B'}
\Hom_{B'}(C', B) \otimes_{C'} C
$$
Ce module est bien muni d'une application canonique vers
$\Hom_A(C', A) \otimes_{C'} C = \omega_{C/A}$,
provenant de la composition
$\Hom_A(B', A) \times \Hom_{B'}(C', B) \to \Hom_A(C', A)$.
\end{proof}

\begin{lemma}
\label{discriminant-lemma-dualizing-product}
Soient $A \to B$ et $A \to C$ des homomorphismes quasi-finis d'anneaux noethériens.
Alors $\omega_{B \times C/A} = \omega_{B/A} \times \omega_{C/A}$
comme modules sur $B \times C$.
\end{lemma}

\begin{proof}
Choisissons des factorisations $A \to B' \to B$ et $A \to C' \to C$ telles que
$A \to B'$ et $A \to C'$ soient finis et que $\Spec(B) \to \Spec(B')$
et $\Spec(C) \to \Spec(C')$ soient des immersions ouvertes. Alors
$A \to B' \times C' \to B \times C$ est une factorisation analogue.
Le calcul de $\omega_{B \times C/A}$ à l'aide de cette factorisation
donne le lemme.
\end{proof}

\begin{lemma}
\label{discriminant-lemma-dualizing-associated-primes}
Soit $A \to B$ un homomorphisme quasi-fini d'anneaux noethériens.
Alors $\text{Ass}_B(\omega_{B/A})$ est l'ensemble des idéaux premiers de $B$
au-dessus d'idéaux premiers associés de $A$.
\end{lemma}

\begin{proof}
Choisissons une factorisation $A \to B' \to B$ où $A \to B'$ est fini et
$B' \to B$ induit une immersion ouverte des spectres. Comme
$\omega_{B/A} = \omega_{B'/A} \otimes_{B'} B$, il suffit
de démontrer l'assertion pour $\omega_{B'/A}$. On peut donc supposer $A \to B$
fini.

\medskip\noindent
Supposons $\mathfrak p \in \text{Ass}(A)$ et soit $\mathfrak q$ un idéal premier
de $B$ au-dessus de $\mathfrak p$. Soit $x \in A$ un élément dont
l'annulateur est $\mathfrak p$. Choisissons une application $\kappa(\mathfrak p)$-linéaire
non nulle $\lambda : \kappa(\mathfrak q) \to \kappa(\mathfrak p)$.
Puisque $A/\mathfrak p \subset B/\mathfrak q$ est une extension finie
d'anneaux, il existe $f \in A$, $f \not \in \mathfrak p$,
tel que $f\lambda$ envoie $B/\mathfrak q$ dans $A/\mathfrak p$.
On obtient donc une application $A$-linéaire non nulle
$$
B \to B/\mathfrak q \to A/\mathfrak p \to A,\quad
b \mapsto f\lambda(b)x
$$
Un calcul immédiat montre que cet élément de $\omega_{B/A}$
a pour annulateur $\mathfrak q$, d'où
$\mathfrak q \in \text{Ass}(\omega_{B/A})$.

\medskip\noindent
Réciproquement, supposons que $\mathfrak q \subset B$ soit un idéal premier
au-dessus d'un idéal premier $\mathfrak p \subset A$ qui n'est pas associé
à $A$. Il faut montrer que
$\mathfrak q \not \in \text{Ass}_B(\omega_{B/A})$.
En remplaçant $A$ par $A_\mathfrak p$ et $B$ par
$B_\mathfrak p$, on peut supposer que $\mathfrak p$ est un idéal maximal
de $A$. Ceci est permis par le lemme \ref{discriminant-lemma-dualizing-flat-base-change} et
Algèbre, lemme \ref{algebra-lemma-localize-ass}.
Il existe alors $f \in \mathfrak m$
qui est un non-diviseur de zéro de $A$.
Cet élément $f$ agit donc injectivement sur $\omega_{B/A}$,
et $\mathfrak q$ n'est donc pas un idéal premier associé à ce module.
\end{proof}

\begin{lemma}
\label{discriminant-lemma-dualizing-base-flat-flat}
Soit $A \to B$ un homomorphisme plat quasi-fini d'anneaux noethériens.
Alors $\omega_{B/A}$ est un $A$-module plat.
\end{lemma}

\begin{proof}
Soit $\mathfrak q \subset B$ un idéal premier au-dessus de $\mathfrak p \subset A$.
Nous allons montrer que la localisation $\omega_{B/A, \mathfrak q}$ est plate
sur $A_\mathfrak p$.
Cela suffit d'après Algèbre, lemme \ref{algebra-lemma-flat-localization}.
D'après
Algèbre, lemme \ref{algebra-lemma-etale-makes-quasi-finite-finite-one-prime},
on peut trouver un homomorphisme étale d'anneaux $A \to A'$ et un idéal
premier $\mathfrak p' \subset A'$ au-dessus de $\mathfrak p$
tels que $\kappa(\mathfrak p') = \kappa(\mathfrak p)$ et
que
$$
B' = B \otimes_A A' = C \times D
$$
avec $A' \to C$ fini, et tels que l'unique idéal premier $\mathfrak q'$
de $B \otimes_A A'$ au-dessus de $\mathfrak q$ et de $\mathfrak p'$
corresponde à un idéal premier de $C$. D'après le
lemme \ref{discriminant-lemma-dualizing-flat-base-change}
et Algèbre, lemme \ref{algebra-lemma-base-change-flat-up-down},
il suffit de montrer que $\omega_{B'/A', \mathfrak q'}$
est plat sur $A'_{\mathfrak p'}$.
Puisque $\omega_{B'/A'} = \omega_{C/A'} \times \omega_{D/A'}$
d'après le lemme \ref{discriminant-lemma-dualizing-product},
on est ramené au cas où $B$ est fini plat sur $A$.
Dans ce cas, $B$ est un $A$-module fini localement libre
et $\omega_{B/A} = \Hom_A(B, A)$ est le $A$-module dual,
fini localement libre.
\end{proof}

\begin{lemma}
\label{discriminant-lemma-dualizing-base-change-of-flat}
Si $A \to B$ est plat, l'application de changement de base (\ref{discriminant-equation-bc-dualizing})
induit un isomorphisme $\omega_{B/A} \otimes_B B_1 \to \omega_{B_1/A_1}$.
\end{lemma}

\begin{proof}
Si $A \to B$ est fini plat, $B$ est un $A$-module fini localement libre.
Dans ce cas, $\omega_{B/A} = \Hom_A(B, A)$ est le $A$-module dual,
fini localement libre, et sa formation commute
à tout changement de base, ce qui démontre le lemme dans ce cas.
Dans le paragraphe suivant, nous ramenons le cas général (quasi-fini plat)
au cas fini plat que l'on vient de traiter.

\medskip\noindent
Soit $\mathfrak q_1 \subset B_1$ un idéal premier. Nous allons montrer que la
l'application localisée en $\mathfrak q_1$ est un isomorphisme, ce qui
suffit d'après Algèbre, lemme \ref{algebra-lemma-characterize-zero-local}.
Soient $\mathfrak q \subset B$ et $\mathfrak p \subset A$ les idéaux
premiers au-dessous de $\mathfrak q_1$. D'après
Algèbre, lemme \ref{algebra-lemma-etale-makes-quasi-finite-finite-one-prime},
on peut trouver un homomorphisme étale d'anneaux $A \to A'$ et un idéal
premier $\mathfrak p' \subset A'$ au-dessus de $\mathfrak p$
tels que $\kappa(\mathfrak p') = \kappa(\mathfrak p)$ et
que
$$
B' = B \otimes_A A' = C \times D
$$
avec $A' \to C$ fini, et tels que l'unique idéal premier $\mathfrak q'$
de $B \otimes_A A'$ au-dessus de $\mathfrak q$ et de $\mathfrak p'$
corresponde à un idéal premier de $C$. Posons $A'_1 = A' \otimes_A A_1$ et
considérons les applications de changement de base
(\ref{discriminant-equation-bc-dualizing}) pour les homomorphismes d'anneaux
$A \to A' \to A'_1$ et $A \to A_1 \to A'_1$, dans le diagramme
$$
\xymatrix{
\omega_{B'/A'} \otimes_{B'} B'_1 \ar[r] & \omega_{B'_1/A'_1} \\
\omega_{B/A} \otimes_B B'_1 \ar[r] \ar[u] &
\omega_{B_1/A_1} \otimes_{B_1} B'_1 \ar[u]
}
$$
où $B' = B \otimes_A A'$, $B_1 = B \otimes_A A_1$, et
$B_1' = B \otimes_A (A' \otimes_A A_1)$. D'après le
lemme \ref{discriminant-lemma-bc-map-dualizing}, le diagramme commute. D'après le
lemme \ref{discriminant-lemma-dualizing-flat-base-change},
les flèches verticales sont des isomorphismes.
Comme $B_1 \to B'_1$ est étale, donc plat, il suffit
de montrer que la flèche horizontale supérieure est un isomorphisme après localisation
en un idéal premier $\mathfrak q'_1$ de $B'_1$ au-dessus de $\mathfrak q$
(un tel idéal premier existe ; appliquer
Algèbre, lemme \ref{algebra-lemma-local-flat-ff}).
On peut donc supposer que $B = C \times D$, avec $A \to C$
fini et $\mathfrak q$ correspondant à un idéal premier de $C$.
Dans ce cas, le module dualisant $\omega_{B/A}$ se décompose
de façon analogue (lemme \ref{discriminant-lemma-dualizing-product}),
ce qui ramène la question
au cas fini plat $A \to C$ traité ci-dessus.
\end{proof}

\begin{remark}
\label{discriminant-remark-relative-dualizing-for-quasi-finite}
Soit $f : Y \to X$ un morphisme localement quasi-fini de schémas localement
noethériens. Il résulte immédiatement du lemme \ref{discriminant-lemma-localize-dualizing}
qu'il existe un unique $\mathcal{O}_Y$-module cohérent
$\omega_{Y/X}$ sur $Y$ tel que, pour tout couple d'ouverts affines
$\Spec(B) = V \subset Y$, $\Spec(A) = U \subset X$ vérifiant $f(V) \subset U$,
on ait un isomorphisme canonique
$$
H^0(V, \omega_{Y/X}) = \omega_{B/A}
$$
et que ces isomorphismes soient compatibles aux applications de restriction.
\end{remark}

\begin{lemma}
\label{discriminant-lemma-compare-dualizing-algebraic}
Soit $A \to B$ un homomorphisme quasi-fini d'anneaux noethériens.
Soit $\omega_{B/A}^\bullet \in D(B)$ le complexe dualisant relatif
algébrique étudié dans Complexes dualisants, section
\ref{dualizing-section-relative-dualizing-complexes-Noetherian}.
Il existe alors un isomorphisme (non unique)
$\omega_{B/A} = H^0(\omega_{B/A}^\bullet)$.
\end{lemma}

\begin{proof}
Choisissons une factorisation $A \to B' \to B$
où $A \to B'$ est fini et $\Spec(B') \to \Spec(B)$
est une immersion ouverte. Alors
$\omega_{B/A}^\bullet = \omega_{B'/A}^\bullet \otimes_B^\mathbf{L} B'$
d'après Complexes dualisants, lemmes
\ref{dualizing-lemma-composition-shriek-algebraic} et
\ref{dualizing-lemma-upper-shriek-localize}, et
la définition de $\omega_{B/A}^\bullet$. Il suffit donc
de montrer qu'il existe un isomorphisme lorsque $A \to B$ est fini.
Dans ce cas, on peut utiliser
Complexes dualisants, lemme \ref{dualizing-lemma-upper-shriek-finite},
pour voir que $\omega_{B/A}^\bullet = R\Hom(B, A)$, d'où
$H^0(\omega^\bullet_{B/A}) = \Hom_A(B, A)$, ce qu'il fallait établir.
\end{proof}






\section{Discriminant d'un morphisme fini localement libre}
\label{discriminant-section-discriminant}

\noindent
Soient $X$ un schéma et $\mathcal{F}$ un $\mathcal{O}_X$-module fini localement
libre. Il existe alors un homomorphisme canonique, l'{\it homomorphisme trace},
$$
\text{Trace} :
\SheafHom_{\mathcal{O}_X}(\mathcal{F}, \mathcal{F})
\longrightarrow
\mathcal{O}_X
$$
Voir Exercices, exercice \ref{exercises-exercise-trace-det}. Cet homomorphisme vérifie
la propriété suivante : $\text{Trace}(\text{id})$ est la fonction localement constante
de $\mathcal{O}_X$ correspondant au rang de $\mathcal{F}$.

\medskip\noindent
Soit $\pi : X \to Y$ un morphisme fini localement
libre de schémas. Il existe un {\it homomorphisme trace de $\pi$} canonique,
qui est $\mathcal{O}_Y$-linéaire
$$
\text{Trace}_\pi : \pi_*\mathcal{O}_X \longrightarrow \mathcal{O}_Y
$$
envoyant une section locale $f$ de $\pi_*\mathcal{O}_X$ sur la
trace de la multiplication par $f$ sur $\pi_*\mathcal{O}_X$. Sur les
ouverts affines, on retrouve la construction de
Exercices, exercice \ref{exercises-exercise-trace-det-rings}.
La composée
$$
\mathcal{O}_Y \xrightarrow{\pi^\sharp} \pi_*\mathcal{O}_X
\xrightarrow{\text{Trace}_\pi} \mathcal{O}_Y
$$
est la multiplication par le degré de $\pi$ (qui est une fonction localement constante
sur $Y$). Par analogie avec
Corps, section \ref{fields-section-trace-pairing},
on peut définir la forme trace
$$
Q_\pi :
\pi_*\mathcal{O}_X \times \pi_*\mathcal{O}_X
\longrightarrow
\mathcal{O}_Y
$$
par $(f, g) \mapsto \text{Trace}_\pi(fg)$. On peut considérer
$Q_\pi$ comme une application linéaire
$\pi_*\mathcal{O}_X \to
\SheafHom_{\mathcal{O}_Y}(\pi_*\mathcal{O}_X, \mathcal{O}_Y)$
entre modules localement libres de même rang, et l'on obtient ainsi
un déterminant
$$
\det(Q_\pi) :
\wedge^{top}(\pi_*\mathcal{O}_X)
\longrightarrow
\wedge^{top}(\pi_*\mathcal{O}_X)^{\otimes -1}
$$
ou, en d'autres termes, une section globale
$$
\det(Q_\pi) \in \Gamma(Y, \wedge^{top}(\pi_*\mathcal{O}_X)^{\otimes -2})
$$
Le {\it discriminant de $\pi$} est, par définition, le sous-schéma
fermé $D_\pi \subset Y$ défini par cette section globale.
Il est clair que $D_\pi$ est un sous-schéma fermé localement principal de $Y$.

\begin{lemma}
\label{discriminant-lemma-discriminant}
Soit $\pi : X \to Y$ un morphisme fini localement
libre de schémas. Alors $\pi$ est étale si et seulement si son discriminant est vide.
\end{lemma}

\begin{proof}
D'après Morphismes, lemme \ref{morphisms-lemma-etale-flat-etale-fibres},
il suffit de vérifier que les fibres de $\pi$ sont étales.
Puisque la construction de la forme trace commute au changement de
base, on est ramené à la question suivante : soient $k$ un corps
et $A$ une $k$-algèbre de dimension finie. Montrer que
$A$ est étale sur $k$ si et seulement si la forme trace
$Q_{A/k} : A \times A \to k$, $(a, b) \mapsto \text{Trace}_{A/k}(ab)$,
est non dégénérée.

\medskip\noindent
Supposons $Q_{A/k}$ non dégénérée. Si $a \in A$ est nilpotent, alors
$ab$ est nilpotent pour tout $b \in A$, et $Q_{A/k}(a, -)$ est donc
identiquement nulle. Ainsi $A$ est réduit. On peut alors écrire
$A = K_1 \times \ldots \times K_n$ comme un produit où chaque $K_i$
est un corps (voir
Algèbre, lemmes \ref{algebra-lemma-finite-dimensional-algebra},
\ref{algebra-lemma-artinian-finite-length} et
\ref{algebra-lemma-minimal-prime-reduced-ring}).
Dans ce cas, l'espace
quadratique $(A, Q_{A/k})$ est la somme directe orthogonale des espaces
$(K_i, Q_{K_i/k})$. Il résulte de
Corps, lemme \ref{fields-lemma-separable-trace-pairing},
que chaque $K_i$ est séparable sur $k$. Cela signifie que $A$ est étale
sur $k$, d'après Algèbre, lemme \ref{algebra-lemma-etale-over-field}.
On démontre la réciproque en parcourant l'argument en sens inverse.
\end{proof}





\section{Traces des homomorphismes plats quasi-finis d'anneaux}
\label{discriminant-section-quasi-finite-traces}

\noindent
La trace dont il est question dans le titre de cette section est de nature tout
autre que celle étudiée dans
Dualité des schémas, section \ref{duality-section-trace}.
Il s'agit ici de la trace
étudiée dans Corps, section \ref{fields-section-trace-pairing},
et généralisée dans Exercices, exercices \ref{exercises-exercise-trace-det} et
\ref{exercises-exercise-trace-det-rings}.

\medskip\noindent
Soit $A \to B$ un homomorphisme fini plat d'anneaux noethériens. Alors $B$ est un
$A$-module fini plat, donc fini localement libre
(Algèbre, lemme \ref{algebra-lemma-finite-projective}).
Étant donné $b \in B$, on peut considérer la {\it trace} $\text{Trace}_{B/A}(b)$
de l'application $A$-linéaire $B \to B$ donnée par
la multiplication par $b$ sur $B$. D'après les références ci-dessus, cela définit
un homomorphisme $A$-linéaire $\text{Trace}_{B/A} : B \to A$.
Puisque $\omega_{B/A} = \Hom_A(B, A)$, car $A \to B$ est fini, on voit
que $\text{Trace}_{B/A} \in \omega_{B/A}$.

\medskip\noindent
Pour un homomorphisme plat quasi-fini d'anneaux quelconque, on définit la notion
de trace comme suit.

\begin{definition}
\label{discriminant-definition-trace-element}
Soit $A \to B$ un homomorphisme plat quasi-fini d'anneaux noethériens.
L'{\it élément trace} est l'unique\footnote{L'unicité
et l'existence seront justifiées dans les
lemmes \ref{discriminant-lemma-trace-unique} et \ref{discriminant-lemma-dualizing-tau}.}
élément
$\tau_{B/A} \in \omega_{B/A}$
ayant la propriété suivante : pour toute $A$-algèbre noethérienne $A_1$
telle que $B_1 = B \otimes_A A_1$ soit muni d'une
décomposition en produit $B_1 = C \times D$ avec $A_1 \to C$ fini,
l'image de $\tau_{B/A}$ dans $\omega_{C/A_1}$
est $\text{Trace}_{C/A_1}$.
On utilise ici l'application de changement de base (\ref{discriminant-equation-bc-dualizing}) et
le lemme \ref{discriminant-lemma-dualizing-product} pour obtenir
$\omega_{B/A} \to \omega_{B_1/A_1} \to \omega_{C/A_1}$.
\end{definition}

\noindent
Nous montrons d'abord l'unicité des éléments trace, puis
leur existence.

\begin{lemma}
\label{discriminant-lemma-trace-unique}
Soit $A \to B$ un homomorphisme plat quasi-fini d'anneaux noethériens.
Il existe alors au plus un élément trace dans $\omega_{B/A}$.
\end{lemma}

\begin{proof}
Soit $\mathfrak q \subset B$ un idéal premier au-dessus de l'idéal premier
$\mathfrak p \subset A$. D'après
Algèbre, lemme \ref{algebra-lemma-etale-makes-quasi-finite-finite-one-prime},
on peut trouver un homomorphisme étale d'anneaux $A \to A_1$ et un idéal
premier $\mathfrak p_1 \subset A_1$ au-dessus de $\mathfrak p$
tels que $\kappa(\mathfrak p_1) = \kappa(\mathfrak p)$ et
que
$$
B_1 = B \otimes_A A_1 = C \times D
$$
avec $A_1 \to C$ fini, et tels que l'unique idéal premier $\mathfrak q_1$
de $B \otimes_A A_1$ au-dessus de $\mathfrak q$ et de $\mathfrak p_1$
corresponde à un idéal premier de $C$. Remarquons que
$\omega_{C/A_1} = \omega_{B/A} \otimes_B C$
(combiner les lemmes \ref{discriminant-lemma-dualizing-flat-base-change} et
\ref{discriminant-lemma-dualizing-product}). Puisque la collection
des homomorphismes d'anneaux $B \to C$ ainsi obtenus forme une famille
conjointement injective d'homomorphismes plats et que l'image de $\tau_{B/A}$
dans $\omega_{C/A_1}$ est prescrite, l'unicité en résulte.
\end{proof}

\noindent
Vérifions la cohérence de cette définition.

\begin{lemma}
\label{discriminant-lemma-finite-flat-trace}
Soit $A \to B$ un homomorphisme fini plat d'anneaux noethériens.
Alors $\text{Trace}_{B/A} \in \omega_{B/A}$ est l'élément trace.
\end{lemma}

\begin{proof}
Supposons donnés $A \to A_1$, avec $A_1$ noethérien, et
une décomposition en produit $B \otimes_A A_1 = C \times D$ avec $A_1 \to C$
fini. Dans ce cas, $A_1 \to D$ est évidemment fini lui aussi.
Posons $B_1 = B \otimes_A A_1$.
Puisque la construction des traces commute au changement de base,
on voit que $\text{Trace}_{B/A}$ s'envoie sur $\text{Trace}_{B_1/A_1}$.
Il suffit donc de remarquer que
$\text{Trace}_{B_1/A_1} = (\text{Trace}_{C/A_1}, \text{Trace}_{D/A_1})$
par l'isomorphisme
$\omega_{B_1/A_1} = \omega_{C/A_1} \times \omega_{D/A_1}$
du lemme \ref{discriminant-lemma-dualizing-product}.
\end{proof}

\begin{lemma}
\label{discriminant-lemma-trace-base-change}
Soit $A \to B$ un homomorphisme plat quasi-fini d'anneaux noethériens.
Soit $\tau \in \omega_{B/A}$ un élément trace.
\begin{enumerate}
\item Si $A \to A_1$ est un homomorphisme où $A_1$ est noethérien, alors, en posant
$B_1 = A_1 \otimes_A B$, l'image de $\tau$ dans $\omega_{B_1/A_1}$ est un
élément trace.
\item Si $A = R_f$ pour un anneau $R$ et un $f \in R$, alors
$\tau$ est un élément trace dans $\omega_{B/R}$.
\item Si $g \in B$, l'image de $\tau$ dans $\omega_{B_g/A}$
est un élément trace.
\item Si $B = B_1 \times B_2$, alors $\tau$ s'envoie sur un élément trace
dans chacun des modules $\omega_{B_1/A}$ et $\omega_{B_2/A}$.
\end{enumerate}
\end{lemma}

\begin{proof}
L'assertion (1) est une conséquence formelle de la définition.

\medskip\noindent
L'assertion (2) a un sens car $\omega_{B/R} = \omega_{B/A}$
d'après le lemme \ref{discriminant-lemma-localize-dualizing}. Notons $\tau'$ l'élément
$\tau$ considéré dans $\omega_{B/R}$. Pour démontrer (2),
supposons donnés $R \to R_1$, avec $R_1$ noethérien, et une décomposition
en produit $B \otimes_R R_1 = C \times D$ où $R_1 \to C$ est fini.
En posant $A_1 = (R_1)_f$, on a $B \otimes_A A_1 = C \times D$.
Puisque $R_1 \to C$ est fini, $A_1 \to C$ l'est a fortiori.
On peut donc utiliser la propriété qui définit $\tau$ pour obtenir la propriété
correspondante de $\tau'$.

\medskip\noindent
L'assertion (3) a un sens car $\omega_{B_g/A} = (\omega_{B/A})_g$
d'après le lemme \ref{discriminant-lemma-localize-dualizing}. La démonstration est analogue à celle
de (2). Supposons donnés $A \to A_1$, avec $A_1$ noethérien, et
une décomposition en produit $B_g \otimes_A A_1 = C \times D$ où $A_1 \to C$
est fini. Posons $B_1 = B \otimes_A A_1$. Alors
$\Spec(C) \to \Spec(B_1)$ est une immersion ouverte puisque $B_g \otimes_A A_1 = (B_1)_g$,
et son image est fermée puisque $B_1 \to C$ est fini
(car $A_1 \to C$ est fini).
On a donc $B_1 = C \times D_1$ et $D = (D_1)_g$. On peut alors utiliser
la propriété qui définit $\tau$ pour obtenir la propriété correspondante
de l'image de $\tau$ dans $\omega_{B_g/A}$.

\medskip\noindent
L'assertion (4) a un sens car
$\omega_{B/A} = \omega_{B_1/A} \times \omega_{B_2/A}$ d'après le
lemme \ref{discriminant-lemma-dualizing-product}.
Supposons donnés $A \to A'$, avec $A'$ noethérien, et
une décomposition en produit $B \otimes_A A' = C \times D$ où $A' \to C$
est fini. On peut manifestement raffiner cette décomposition
en produit en $B \otimes_A A' = C_1 \times C_2 \times D_1 \times D_2$,
avec $A' \to C_i$ fini, de sorte que $B_i \otimes_A A' = C_i \times D_i$.
On peut alors utiliser la propriété qui définit $\tau$ pour obtenir la propriété
correspondante de l'image de $\tau$ dans $\omega_{B_i/A}$. On utilise ici le fait évident
que
$\text{Trace}_{C/A'} = (\text{Trace}_{C_1/A'}, \text{Trace}_{C_2/A'})$
pour la décomposition
$\omega_{C/A'} = \omega_{C_1/A'} \times \omega_{C_2/A'}$.
\end{proof}

\begin{lemma}
\label{discriminant-lemma-glue-trace}
Soit $A \to B$ un homomorphisme plat quasi-fini d'anneaux noethériens.
Soient $g_1, \ldots, g_m \in B$ des éléments engendrant l'idéal unité.
Soit $\tau \in \omega_{B/A}$ un élément dont l'image dans
$\omega_{B_{g_i}/A}$ est un élément trace pour $A \to B_{g_i}$.
Alors $\tau$ est un élément trace.
\end{lemma}

\begin{proof}
Supposons donnés $A \to A_1$, avec $A_1$ noethérien, et une décomposition
en produit $B \otimes_A A_1 = C \times D$ où $A_1 \to C$ est fini.
Il faut montrer que l'image de $\tau$ dans $\omega_{C/A_1}$ est
$\text{Trace}_{C/A_1}$. Remarquons que $g_1, \ldots, g_m$
engendrent l'idéal unité de $B_1 = B \otimes_A A_1$ et que
$\tau$ s'envoie sur un élément trace dans $\omega_{(B_1)_{g_i}/A_1}$
d'après le lemme \ref{discriminant-lemma-trace-base-change}. On peut donc remplacer
$A$ par $A_1$ et $B$ par $B_1$ pour se placer dans la situation décrite
au paragraphe suivant.

\medskip\noindent
On suppose ici que $B = C \times D$ avec $A \to C$ fini.
Soit $\tau_C$ l'image de $\tau$ dans $\omega_{C/A}$.
Il faut démontrer que $\tau_C = \text{Trace}_{C/A}$ dans $\omega_{C/A}$.
La compatibilité des éléments trace aux produits
(lemme \ref{discriminant-lemma-trace-base-change})
montre que $\tau_C$ s'envoie sur un élément trace dans $\omega_{C_{g_i}/A}$.
Ainsi, en remplaçant $B$ par $C$, on peut supposer $A \to B$
fini plat.

\medskip\noindent
Supposons $A \to B$ fini plat. Dans ce cas, $\text{Trace}_{B/A}$
est un élément trace d'après le lemme \ref{discriminant-lemma-finite-flat-trace}.
Par conséquent, $\text{Trace}_{B/A}$ s'envoie sur un élément trace dans
$\omega_{B_{g_i}/A}$ d'après le lemme \ref{discriminant-lemma-trace-base-change}.
L'unicité des éléments trace (lemme \ref{discriminant-lemma-trace-unique})
montre que $\text{Trace}_{B/A}$ et $\tau$ ont
les mêmes images dans $\omega_{B_{g_i}/A} = (\omega_{B/A})_{g_i}$.
Comme $g_1, \ldots, g_m$ engendrent l'idéal unité de $B$, l'application
$\omega_{B/A} \to \prod \omega_{B_{g_i}/A}$ est injective,
et l'on conclut que $\tau_C = \text{Trace}_{B/A}$, ce qu'il fallait établir.
\end{proof}

\begin{lemma}
\label{discriminant-lemma-dualizing-tau}
Soit $A \to B$ un homomorphisme plat quasi-fini d'anneaux noethériens.
Il existe un élément trace $\tau \in \omega_{B/A}$.
\end{lemma}

\begin{proof}
Choisissons une factorisation $A \to B' \to B$ où $A \to B'$ est fini et
$\Spec(B) \to \Spec(B')$ est une immersion ouverte. Soient $g_1, \ldots, g_n \in B'$
des éléments tels que $\Spec(B) = \bigcup D(g_i)$ comme ouverts de $\Spec(B')$.
Supposons que l'on puisse montrer l'existence d'éléments trace $\tau_i$ pour les
homomorphismes quasi-finis plats d'anneaux $A \to B_{g_i}$. Alors, pour tous $i, j$, les éléments
$\tau_i$ et $\tau_j$ s'envoient sur des éléments trace de $\omega_{B_{g_ig_j}/A}$
d'après le lemme \ref{discriminant-lemma-trace-base-change}. Par unicité des
éléments trace (lemme \ref{discriminant-lemma-trace-unique}), ils ont la même image.
La condition de faisceau pour le module quasi-cohérent associé à
$\omega_{B/A}$ (voir Algèbre, lemme \ref{algebra-lemma-cover-module})
fournit donc un élément $\tau \in \omega_{B/A}$.
Alors $\tau$ est un élément trace d'après le
lemme \ref{discriminant-lemma-glue-trace}.
On est ainsi ramené au cas traité au paragraphe suivant.

\medskip\noindent
Supposons $A \to B'$ fini et $g \in B'$ tel que $B = B'_g$ soit plat sur $A$.
Il s'agit de construire un élément trace dans
$\omega_{B/A} = \Hom_A(B', A) \otimes_{B'} B$.
Choisissons une résolution $F_1 \to F_0 \to B' \to 0$ de $B'$ par des
$A$-modules libres de type fini $F_0$ et $F_1$. On a alors une suite exacte
$$
0 \to \Hom_A(B', A) \to F_0^\vee \to F_1^\vee
$$
où $F_i^\vee = \Hom_A(F_i, A)$ est le module dual, libre de type fini.
De même, on a la suite exacte
$$
0 \to \Hom_A(B', B') \to F_0^\vee \otimes_A B' \to F_1^\vee \otimes_A B'
$$
L'idée de la construction de $\tau$ est d'utiliser le diagramme
$$
B' \xrightarrow{\mu} \Hom_A(B', B')
\leftarrow \Hom_A(B', A) \otimes_A B'
\xrightarrow{ev} A
$$
où la première flèche envoie $b' \in B'$ sur l'opérateur $A$-linéaire
de multiplication par $b'$, et la dernière flèche est l'application d'évaluation.
La difficulté est que la flèche du milieu, qui envoie $\lambda' \otimes b'$
sur l'application $b'' \mapsto \lambda'(b'')b'$, n'est pas un isomorphisme.
Si $B'$ est plat sur $A$, les suites exactes ci-dessus montrent que cette flèche
est un isomorphisme, et la composée de gauche à droite est la trace usuelle
$\text{Trace}_{B'/A}$. Dans le cas général, on considère
le diagramme
$$
\xymatrix{
& \Hom_A(B', A) \otimes_A B' \ar[r] \ar[d] &
\Hom_A(B', A) \otimes_A B'_g \ar[d] \\
B' \ar[r]_-\mu \ar@{..>}[rru] \ar@{..>}[ru]^\psi &
\Hom_A(B', B') \ar[r] &
\Ker(F_0^\vee \otimes_A B'_g \to F_1^\vee \otimes_A B'_g)
}
$$
La platitude de $A \to B'_g$ montre que la flèche verticale de droite est un
isomorphisme. On obtient donc la flèche pointillée sans étiquette.
Puisque $B'_g = \colim \frac{1}{g^n}B'$, que
les colimites commutent aux produits tensoriels,
et que $B'$ est un $A$-module de présentation finie,
on peut trouver $n \geq 0$ et une application $B'$-linéaire (pour la structure de $B'$-module à droite)
$\psi : B' \to \Hom_A(B', A) \otimes_A B'$
dont la composée avec la flèche verticale de gauche est $g^n\mu$.
En composant avec $ev$, on obtient un élément
$ev \circ \psi \in \Hom_A(B', A)$. On pose alors
$$
\tau = (ev \circ \psi) \otimes g^{-n} \in
\Hom_A(B', A) \otimes_{B'} B'_g = \omega_{B'_g/A} = \omega_{B/A}
$$
Nous omettons la vérification immédiate que cet élément ne dépend pas
du choix de $n$ et de $\psi$ ci-dessus.

\medskip\noindent
Montrons que $\tau$, construit au paragraphe précédent,
possède la propriété voulue dans un cas particulier. Supposons
$B' = C' \times D'$ et $g = (f, h)$, où $A \to C'$ est plat, $D'_h$ est plat, et
$f$ est inversible dans $C'$.
Il faut montrer que $\tau$ s'envoie sur $\text{Trace}_{C'/A}$ dans $\omega_{C'/A}$.
Dans ce cas, on choisit d'abord $n_D$ et
$\psi_D : D' \to \Hom_A(D', A) \otimes_A D'$ comme ci-dessus pour le couple
$(D', h)$, et l'on peut prendre pour
$\psi_C : C' \to \Hom_A(C', A) \otimes_A C' = \Hom_A(C', C')$
l'application envoyant $c' \in C'$ sur la multiplication par $c'$.
On prend alors $n = n_D$ et $\psi = (f^{n_D} \psi_C, \psi_D)$,
et la compatibilité voulue est claire car
$\text{Trace}_{C'/A} = ev \circ \psi_C$, comme on l'a remarqué ci-dessus.

\medskip\noindent
Pour établir la propriété voulue dans le cas général, supposons donnés
$A \to A_1$, avec $A_1$ noethérien, et une décomposition en produit
$B'_g \otimes_A A_1 = C \times D$ où $A_1 \to C$ est fini.
Posons $B'_1 = B' \otimes_A A_1$. Alors $\Spec(C) \to \Spec(B'_1)$
est une immersion ouverte puisque $B'_g \otimes_A A_1 = (B'_1)_g$, et
son image est fermée puisque $B'_1 \to C$ est fini (car $A_1 \to C$
est fini). Ainsi $B'_1 = C \times D'$ et $D'_g = D$.
On en déduit que $B'_1 = C \times D'$ et $g$, sur $A_1$,
sont dans la situation du paragraphe précédent.
Puisque la formation du diagramme affiché ci-dessus
commute au changement de base, la formation de $\tau$ commute
au changement de base $A \to A_1$ (détails omis ; utiliser la
résolution $F_1 \otimes_A A_1 \to F_0 \otimes_A A_1 \to B'_1 \to 0$
pour le voir). La compatibilité voulue découle donc du résultat
du paragraphe précédent.
\end{proof}

\begin{remark}
\label{discriminant-remark-relative-dualizing-for-flat-quasi-finite}
Soit $f : Y \to X$ un morphisme plat localement quasi-fini de schémas localement
noethériens. Soit $\omega_{Y/X}$ comme dans la
remarque \ref{discriminant-remark-relative-dualizing-for-quasi-finite}.
L'unicité, l'existence et la compatibilité à la
localisation des éléments trace
(lemmes \ref{discriminant-lemma-trace-unique}, \ref{discriminant-lemma-dualizing-tau} et
\ref{discriminant-lemma-trace-base-change})
montrent qu'il existe une section globale
$$
\tau_{Y/X} \in \Gamma(Y, \omega_{Y/X})
$$
telle que, pour tout couple d'ouverts affines
$\Spec(B) = V \subset Y$, $\Spec(A) = U \subset X$ vérifiant $f(V) \subset U$,
cet élément $\tau_{Y/X}$ s'envoie sur $\tau_{B/A}$ par
l'isomorphisme canonique
$H^0(V, \omega_{Y/X}) = \omega_{B/A}$.
\end{remark}

\begin{lemma}
\label{discriminant-lemma-tau-nonzero}
Soient $k$ un corps et $A$ une $k$-algèbre finie. Supposons $A$
locale, de corps résiduel $k'$. Les conditions suivantes sont équivalentes :
\begin{enumerate}
\item $\text{Trace}_{A/k}$ est non nulle ;
\item $\tau_{A/k} \in \omega_{A/k}$ est non nul ;
\item $k'/k$ est séparable et $\text{longueur}_A(A)$ est première
à la caractéristique de $k$.
\end{enumerate}
\end{lemma}

\begin{proof}
Les conditions (1) et (2) sont équivalentes d'après le lemme \ref{discriminant-lemma-finite-flat-trace}.
Soit $\mathfrak m \subset A$. Puisque $\dim_k(A) < \infty$, il est clair que
$A$ est de longueur finie sur $A$. Choisissons une filtration
$$
A = I_0 \supset \mathfrak m = I_1 \supset I_2 \supset \ldots I_n = 0
$$
par des idéaux tels que $I_i/I_{i + 1} \cong k'$ comme $A$-modules. Voir
Algèbre, lemme \ref{algebra-lemma-simple-pieces}, qui montre également que
$n = \text{longueur}_A(A)$. Si $a \in \mathfrak m$, alors $aI_i \subset I_{i + 1}$,
et l'on voit immédiatement que $\text{Trace}_{A/k}(a) = 0$.
Si $a \not \in \mathfrak m$ a pour image $\lambda \in k'$, alors
on obtient
$$
\text{Trace}_{A/k}(a) =
\sum\nolimits_{i = 0, \ldots, n - 1}
\text{Trace}_k(a : I_i/I_{i - 1} \to I_i/I_{i - 1}) =
n \text{Trace}_{k'/k}(\lambda)
$$
On achève la démonstration en appliquant
Corps, lemme \ref{fields-lemma-separable-trace-pairing}.
\end{proof}






\section{Morphismes finis}
\label{discriminant-section-finite-morphisms}

\noindent
Dans cette section, nous rassemblons quelques remarques sur les
constructions des sections précédentes dans le cas des morphismes finis.
Soit $f : Y \to X$ un morphisme fini de schémas localement noethériens.
Soit $\omega_{Y/X}$ comme dans la
remarque \ref{discriminant-remark-relative-dualizing-for-quasi-finite}.

\medskip\noindent
La première remarque est que
$$
f_*\omega_{Y/X} = \SheafHom_{\mathcal{O}_X}(f_*\mathcal{O}_Y, \mathcal{O}_X)
$$
comme faisceaux de $f_*\mathcal{O}_Y$-modules. Puisque $f$ est affine, cette
formule caractérise $\omega_{Y/X}$ de façon unique, voir
Morphismes, lemme \ref{morphisms-lemma-affine-equivalence-modules}.
La formule est vraie car, pour un ouvert affine $\Spec(A) = U \subset X$,
l'image inverse $V = f^{-1}(U)$ est le spectre d'une $A$-algèbre finie
$B$, d'où
$$
H^0(U, f_*\omega_{Y/X}) =
H^0(V, \omega_{Y/X}) =
\omega_{B/A} =
\Hom_A(B, A) =
H^0(U, \SheafHom_{\mathcal{O}_X}(f_*\mathcal{O}_Y, \mathcal{O}_X))
$$
par construction. On obtient en particulier une application canonique d'évaluation
$$
f_*\omega_{Y/X} \longrightarrow \mathcal{O}_X
$$
donnée par l'évaluation en $1$ lorsque l'on considère $f_*\omega_{Y/X}$
comme le faisceau $\SheafHom_{\mathcal{O}_X}(f_*\mathcal{O}_Y, \mathcal{O}_X)$.

\medskip\noindent
La deuxième remarque est que l'application d'évaluation fournit
des identifications canoniques
$$
\Hom_Y(\mathcal{F}, f^*\mathcal{G} \otimes_{\mathcal{O}_Y} \omega_{Y/X})
=
\Hom_X(f_*\mathcal{F}, \mathcal{G})
$$
fonctorielles par rapport au module quasi-cohérent $\mathcal{F}$ sur $Y$
et au module fini localement libre $\mathcal{G}$ sur $X$.
Si $\mathcal{G} = \mathcal{O}_X$, cela résulte immédiatement
de ce qui précède et de
Algèbre, lemme \ref{algebra-lemma-adjoint-hom-restrict}.
Pour $\mathcal{G}$ quelconque, on peut utiliser le même lemme et les
isomorphismes
$$
f_*(f^*\mathcal{G} \otimes_{\mathcal{O}_Y} \omega_{Y/X}) =
\mathcal{G} \otimes_{\mathcal{O}_X}
\SheafHom_{\mathcal{O}_X}(f_*\mathcal{O}_Y, \mathcal{O}_X) =
\SheafHom_{\mathcal{O}_X}(f_*\mathcal{O}_Y, \mathcal{G})
$$
de $f_*\mathcal{O}_Y$-modules, où la première égalité est la
formule de projection
(Cohomologie, lemme \ref{cohomology-lemma-projection-formula}).
On peut également démontrer la formule localement sur les ouverts affines par
un calcul direct.

\medskip\noindent
La troisième remarque est que, si $f$ est de plus plat, la
composée
$$
f_*\mathcal{O}_Y \xrightarrow{f_*\tau_{Y/X}} f_*\omega_{Y/X}
\longrightarrow \mathcal{O}_X
$$
est égale à l'homomorphisme trace $\text{Trace}_f$ étudié dans la
section \ref{discriminant-section-discriminant}. Cela se voit immédiatement
sur les ouverts affines.

\medskip\noindent
La quatrième remarque est que, si $f$ est plat et $X$ noethérien, on
obtient
$$
\Hom_Y(K, Lf^*M \otimes_{\mathcal{O}_Y} \omega_{Y/X})
=
\Hom_X(Rf_*K, M)
$$
pour tous $K$ dans $D_\QCoh(\mathcal{O}_Y)$ et $M$ dans $D_\QCoh(\mathcal{O}_X)$.
Cela résulte des résultats de
Dualité des schémas, section \ref{duality-section-proper-flat},
mais on peut le démontrer directement dans ce cas, comme suit.
Tout d'abord, si $X$ est affine, l'assertion découle de
Complexes dualisants, lemmes \ref{dualizing-lemma-right-adjoint} et
\ref{dualizing-lemma-RHom-is-tensor-special}\footnote{Ce lemme admet
une démonstration plus simple dans notre cas.}, et de
Catégories dérivées de schémas, lemme \ref{perfect-lemma-affine-compare-bounded}.
On peut ensuite utiliser le principe de récurrence
(Cohomologie des schémas, lemme \ref{coherent-lemma-induction-principle})
et Mayer-Vietoris
(sous la forme de Cohomologie, lemme \ref{cohomology-lemma-mayer-vietoris-hom})
pour achever la démonstration.







\section{La différente de Noether}
\label{discriminant-section-noether-different}

\noindent
La littérature contient de nombreuses notions de différente.
Nous en présentons quelques-unes dans cette section et les suivantes ; pour plus
de détails, nous renvoyons le lecteur à \cite{Kunz}.

\medskip\noindent
Soit $A \to B$ un homomorphisme d'anneaux. Notons
$$
\mu : B \otimes_A B \longrightarrow B,\quad
b \otimes b' \longmapsto bb'
$$
l'application de multiplication. Soit $I = \Ker(\mu)$. Il est clair que $I$ est
engendré par les éléments $b \otimes 1 - 1 \otimes b$ pour $b \in B$.
L'annulateur $J \subset B \otimes_A B$ de $I$ est donc canoniquement un $B$-module.
La {\it différente de Noether} de $B$ sur $A$ est
l'image de $J$ par l'application $\mu : B \otimes_A B \to B$. De façon équivalente,
la différente de Noether est l'image de l'application
$$
J = \Hom_{B \otimes_A B}(B, B \otimes_A B) \longrightarrow B,\quad
\varphi \longmapsto \mu(\varphi(1))
$$
Commençons par quelques lemmes indispensables.

\begin{lemma}
\label{discriminant-lemma-noether-different-product}
Soient $A \to B_i$, $i = 1, 2$, des homomorphismes d'anneaux. Posons $B = B_1 \times B_2$.
\begin{enumerate}
\item L'annulateur $J$ de $\Ker(B \otimes_A B \to B)$ est $J_1 \times J_2$,
où $J_i$ est l'annulateur de $\Ker(B_i \otimes_A B_i \to B_i)$.
\item La différente de Noether $\mathfrak{D}$ de $B$ sur $A$ est
$\mathfrak{D}_1 \times \mathfrak{D}_2$, où $\mathfrak{D}_i$ est
la différente de Noether de $B_i$ sur $A$.
\end{enumerate}
\end{lemma}

\begin{proof}
Omise.
\end{proof}

\begin{lemma}
\label{discriminant-lemma-noether-different-base-change}
Soit $A \to B$ un homomorphisme d'anneaux de type fini. Soit $A \to A'$ un homomorphisme plat d'anneaux.
Posons $B' = B \otimes_A A'$.
\begin{enumerate}
\item L'annulateur $J'$ de $\Ker(B' \otimes_{A'} B' \to B')$ est
$J \otimes_A A'$, où $J$ est l'annulateur de $\Ker(B \otimes_A B \to B)$.
\item La différente de Noether $\mathfrak{D}'$ de $B'$ sur $A'$ est
$\mathfrak{D}B'$, où $\mathfrak{D}$ est
la différente de Noether de $B$ sur $A$.
\end{enumerate}
\end{lemma}

\begin{proof}
Choisissons des générateurs $b_1, \ldots, b_n$ de $B$ comme $A$-algèbre.
Alors
$$
J = \Ker(B \otimes_A B \xrightarrow{b_i \otimes 1 - 1 \otimes b_i}
(B \otimes_A B)^{\oplus n})
$$
On voit donc que la formation de $J$ commute au changement de base plat.
Le résultat concernant la différente de Noether en découle immédiatement.
\end{proof}

\begin{lemma}
\label{discriminant-lemma-noether-different-localization}
Soient $A \to B' \to B$ des homomorphismes d'anneaux, avec $A \to B'$
de type fini et $B' \to B$ induisant une immersion ouverte des spectres.
\begin{enumerate}
\item L'annulateur $J$ de $\Ker(B \otimes_A B \to B)$ est
$J' \otimes_{B'} B$, où $J'$ est l'annulateur de
$\Ker(B' \otimes_A B' \to B')$.
\item La différente de Noether $\mathfrak{D}$ de $B$ sur $A$ est
$\mathfrak{D}'B$, où $\mathfrak{D}'$ est
la différente de Noether de $B'$ sur $A$.
\end{enumerate}
\end{lemma}

\begin{proof}
Posons $I = \Ker(B \otimes_A B \to B)$ et $I' = \Ker(B' \otimes_A B' \to B')$.
Comme $\Spec(B) \to \Spec(B')$ est une immersion ouverte, on a
$B = (B \otimes_A B) \otimes_{B' \otimes_A B'} B'$. Il en résulte
$I = I'(B \otimes_A B)$. Puisque $I'$ est de type fini et que
$B' \otimes_A B' \to B \otimes_A B$ est plat, on obtient
$J = J'(B \otimes_A B)$, voir
Algèbre, lemme \ref{algebra-lemma-annihilator-flat-base-change}.
Puisque la structure de $B' \otimes_A B'$-module de $J'$
se factorise par $B' \otimes_A B' \to B'$, on en déduit (1).
L'assertion (2) est une conséquence de (1).
\end{proof}

\begin{remark}
\label{discriminant-remark-construction-pairing}
Soit $A \to B$ un homomorphisme quasi-fini d'anneaux noethériens.
Soit $J$ l'annulateur de $\Ker(B \otimes_A B \to B)$.
Il existe un accouplement $B$-bilinéaire canonique
\begin{equation}
\label{discriminant-equation-pairing-noether}
\omega_{B/A} \times J \longrightarrow B
\end{equation}
défini comme suit. Choisissons une factorisation $A \to B' \to B$
où $A \to B'$ est fini et $B' \to B$ induit une immersion ouverte
des spectres. Soit $J'$ l'annulateur de $\Ker(B' \otimes_A B' \to B')$.
On définit d'abord
$$
\Hom_A(B', A) \times J' \longrightarrow B',\quad
(\lambda, \sum b_i \otimes c_i) \longmapsto \sum \lambda(b_i)c_i
$$
Cet accouplement est $B'$-bilinéaire précisément parce que, pour $\xi \in J'$ et $b \in B'$,
on a $(b \otimes 1)\xi = (1 \otimes b)\xi$. D'après le
lemme \ref{discriminant-lemma-noether-different-localization}
et puisque $\omega_{B/A} = \Hom_A(B', A) \otimes_{B'} B$,
on peut l'étendre en l'accouplement $B$-bilinéaire affiché ci-dessus.
\end{remark}

\begin{lemma}
\label{discriminant-lemma-noether-pairing-compatibilities}
Soit $A \to B$ un homomorphisme quasi-fini d'anneaux noethériens.
\begin{enumerate}
\item Si $A \to A'$ est un homomorphisme plat d'anneaux noethériens, alors
$$
\xymatrix{
\omega_{B/A} \times J \ar[r] \ar[d] & B \ar[d]  \\
\omega_{B'/A'} \times J' \ar[r] & B'
}
$$
est commutatif, avec les notations du
lemme \ref{discriminant-lemma-noether-different-base-change},
les flèches horizontales étant données par
(\ref{discriminant-equation-pairing-noether}).
\item Si $B = B_1 \times B_2$, alors
$$
\xymatrix{
\omega_{B/A} \times J \ar[r] \ar[d] & B \ar[d]  \\
\omega_{B_i/A} \times J_i \ar[r] & B_i
}
$$
est commutatif pour $i = 1, 2$, avec les notations du
lemme \ref{discriminant-lemma-noether-different-product},
les flèches horizontales étant données par
(\ref{discriminant-equation-pairing-noether}).
\end{enumerate}
\end{lemma}

\begin{proof}
En vertu de la construction de l'accouplement dans la
remarque \ref{discriminant-remark-construction-pairing},
(1) et (2) se ramènent au cas où $A \to B$ est fini.
L'assertion (1) résulte alors de ce que l'application de contraction
$\Hom_A(M, A) \otimes_A M \otimes_A M \to M$,
$\lambda \otimes m \otimes m' \mapsto \lambda(m)m'$,
commute au changement de base. Pour (2), utiliser le fait que
$J = J_1 \times J_2$ est contenu dans les facteurs
$B_1 \otimes_A B_1$ et $B_2 \otimes_A B_2$
de $B \otimes_A B$.
\end{proof}

\begin{lemma}
\label{discriminant-lemma-noether-pairing-flat-quasi-finite}
Soit $A \to B$ un homomorphisme plat quasi-fini d'anneaux noethériens.
L'accouplement de la remarque \ref{discriminant-remark-construction-pairing} induit un isomorphisme
$J \to \Hom_B(\omega_{B/A}, B)$.
\end{lemma}

\begin{proof}
Montrons d'abord l'assertion lorsque $A \to B$ est fini plat. Dans ce cas, on peut
localiser sur $A$ et supposer que $B$ est un $A$-module libre de type fini. Soit
$b_1, \ldots, b_n$ une base de $B$ comme $A$-module, et notons
$b_1^\vee, \ldots, b_n^\vee$ la base duale de $\omega_{B/A}$. Remarquons que
$\sum b_i \otimes c_i \in J$ s'envoie sur l'élément de $\Hom_B(\omega_{B/A}, B)$
qui envoie $b_i^\vee$ sur $c_i$. Supposons $\varphi : \omega_{B/A} \to B$
$B$-linéaire. Nous affirmons alors que $\xi = \sum b_i \otimes \varphi(b_i^\vee)$
est un élément de $J$. En effet, la $B$-linéarité de $\varphi$
entraîne précisément que $(b \otimes 1)\xi = (1 \otimes b)\xi$ pour tout $b \in B$.
Notre application possède donc un inverse et est un isomorphisme.

\medskip\noindent
Soit $\mathfrak q \subset B$ un idéal premier au-dessus de $\mathfrak p \subset A$.
Nous allons montrer que l'application localisée
$$
J_\mathfrak q
\longrightarrow
\Hom_B(\omega_B/A, B)_\mathfrak q
$$
est un isomorphisme.
Cela suffit d'après Algèbre, lemme \ref{algebra-lemma-characterize-zero-local}.
D'après
Algèbre, lemme \ref{algebra-lemma-etale-makes-quasi-finite-finite-one-prime},
on peut trouver un homomorphisme étale d'anneaux $A \to A'$ et un idéal
premier $\mathfrak p' \subset A'$ au-dessus de $\mathfrak p$
tels que $\kappa(\mathfrak p') = \kappa(\mathfrak p)$ et
que
$$
B' = B \otimes_A A' = C \times D
$$
avec $A' \to C$ fini, et tels que l'unique idéal premier $\mathfrak q'$
de $B \otimes_A A'$ au-dessus de $\mathfrak q$ et de $\mathfrak p'$
corresponde à un idéal premier de $C$. Soit $J'$ l'annulateur de
$\Ker(B' \otimes_{A'} B' \to B')$. D'après les
lemmes \ref{discriminant-lemma-dualizing-flat-base-change},
\ref{discriminant-lemma-noether-different-base-change} et
\ref{discriminant-lemma-noether-pairing-compatibilities},
l'application $J' \to \Hom_{B'}(\omega_{B'/A'}, B')$
s'obtient en appliquant le foncteur $- \otimes_B B'$
à l'application $J \to \Hom_B(\omega_{B/A}, B)$.
Puisque $B_\mathfrak q \to B'_{\mathfrak q'}$ est fidèlement plat,
il suffit d'établir le résultat pour $(A' \to B', \mathfrak q')$.
D'après les lemmes \ref{discriminant-lemma-dualizing-product},
\ref{discriminant-lemma-noether-different-product} et
\ref{discriminant-lemma-noether-pairing-compatibilities},
on est ramené au cas traité dans le premier
paragraphe de la démonstration.
\end{proof}

\begin{lemma}
\label{discriminant-lemma-noether-different-flat-quasi-finite}
Soit $A \to B$ un homomorphisme plat quasi-fini d'anneaux noethériens.
Le diagramme
$$
\xymatrix{
J \ar[rr] \ar[rd]_\mu & &
\Hom_B(\omega_{B/A}, B) \ar[ld]^{\varphi \mapsto \varphi(\tau_{B/A})} \\
& B
}
$$
commute, la flèche horizontale étant l'isomorphisme du
lemme \ref{discriminant-lemma-noether-pairing-flat-quasi-finite}.
La différente de Noether de $B$ sur $A$
est donc l'image de l'application $\Hom_B(\omega_{B/A}, B) \to B$.
\end{lemma}

\begin{proof}
Exactement comme dans la démonstration du lemme \ref{discriminant-lemma-noether-pairing-flat-quasi-finite},
on se ramène au cas d'un homomorphisme fini libre $A \to B$.
Dans ce cas, $\tau_{B/A} = \text{Trace}_{B/A}$.
Choisissons une base $b_1, \ldots, b_n$ de $B$ comme $A$-module.
Soit $\xi = \sum b_i \otimes c_i \in J$. Alors $\mu(\xi) = \sum b_i c_i$.
D'autre part, l'image de $\xi$ dans $\Hom_B(\omega_{B/A}, B)$
envoie $\text{Trace}_{B/A}$ sur $\sum \text{Trace}_{B/A}(b_i)c_i$.
Il faut donc montrer que
$$
\sum b_ic_i = \sum \text{Trace}_{B/A}(b_i)c_i
$$
lorsque $\xi = \sum b_i \otimes c_i \in J$. Écrivons $b_i b_j = \sum_k a_{ij}^k b_k$
avec $a_{ij}^k \in A$. Le membre de droite est alors
$\sum_{i, j} a_{ij}^j c_i$. D'autre part, $\xi \in J$ entraîne
$$
(b_j \otimes 1)(\sum\nolimits_i b_i \otimes c_i) =
(1 \otimes b_j)(\sum\nolimits_i b_i \otimes c_i)
$$
ce qui donne $b_j c_i = \sum_k a_{jk}^i c_k$. Le membre de gauche
est donc $\sum_{i, j} a_{ij}^i c_j$. Puisque $a_{ij}^k = a_{ji}^k$, on a bien l'égalité.
\end{proof}

\begin{lemma}
\label{discriminant-lemma-noether-different}
Soit $A \to B$ un homomorphisme d'anneaux de type fini. Soit $\mathfrak{D} \subset B$
la différente de Noether. Alors $V(\mathfrak{D})$ est l'ensemble des idéaux premiers
$\mathfrak q \subset B$ tels que $A \to B$ soit ramifié en $\mathfrak q$.
\end{lemma}

\begin{proof}
Supposons $A \to B$ non ramifié en $\mathfrak q$. En remplaçant
$B$ par $B_g$ pour un $g \in B$, $g \not \in \mathfrak q$, on peut
supposer $A \to B$ non ramifié (Algèbre, définition
\ref{algebra-definition-unramified} et
lemme \ref{discriminant-lemma-noether-different-localization}).
Dans ce cas, $\Omega_{B/A} = 0$. Ainsi, si $I = \Ker(B \otimes_A B \to B)$,
alors $I/I^2 = 0$ d'après
Algèbre, lemme \ref{algebra-lemma-differentials-diagonal}.
Puisque $A \to B$ est de type fini, $I$ est de type
fini. Le lemme de Nakayama
(Algèbre, lemme \ref{algebra-lemma-NAK})
assure donc l'existence d'un élément de la forme $1 + i$
annulant $I$. Il en résulte que $\mathfrak{D} = B$.

\medskip\noindent
Réciproquement, supposons $\mathfrak{D} \not \subset \mathfrak q$.
En remplaçant $B$ par une localisation principale comme ci-dessus,
on peut supposer $\mathfrak{D} = B$. Cela signifie qu'il existe un
élément de la forme $1 + i$ dans l'annulateur de $I$.
Réciproquement, cela entraîne que $I/I^2 = \Omega_{B/A}$ est nul,
ce qui permet de conclure.
\end{proof}







\section{La différente de Kähler}
\label{discriminant-section-kahler-different}

\noindent
Soit $A \to B$ un homomorphisme d'anneaux de type fini. La {\it différente de Kähler} est
l'idéal de Fitting d'ordre zéro de $\Omega_{B/A}$ comme $B$-module. On globalise cette
définition comme suit.

\begin{definition}
\label{discriminant-definition-kahler-different}
Soit $f : Y \to X$ un morphisme de schémas localement de type fini.
La {\it différente de Kähler} est l'idéal de Fitting d'ordre $0$ de $\Omega_{Y/X}$.
\end{definition}

\noindent
La différente de Kähler est un faisceau quasi-cohérent d'idéaux sur $Y$.

\begin{lemma}
\label{discriminant-lemma-base-change-kahler-different}
Considérons un diagramme cartésien de schémas
$$
\xymatrix{
Y' \ar[d]_{f'} \ar[r] & Y \ar[d]^f \\
X' \ar[r]^g & X
}
$$
où $f$ est localement de type fini. Soit $R \subset Y$, resp.\ $R' \subset Y'$,
le sous-schéma fermé défini par la différente de Kähler de $f$, resp.\ de $f'$.
Alors $Y' \to Y$ induit un isomorphisme $R' \to R \times_Y Y'$.
\end{lemma}

\begin{proof}
Cela tient au fait que $\Omega_{Y'/X'}$ est l'image inverse de $\Omega_{Y/X}$
(Morphismes, lemme \ref{morphisms-lemma-base-change-differentials}) ;
on peut alors appliquer
Compléments d'algèbre, lemme \ref{more-algebra-lemma-fitting-ideal-basics}.
\end{proof}

\begin{lemma}
\label{discriminant-lemma-kahler-different}
Soit $f : Y \to X$ un morphisme de schémas localement de type fini.
Soit $R \subset Y$ le sous-schéma fermé défini par
la différente de Kähler. Alors $R \subset Y$ est exactement
l'ensemble des points où $f$ est ramifié.
\end{lemma}

\begin{proof}
Il s'agit de l'énoncé de
Diviseurs, lemme \ref{divisors-lemma-zero-fitting-ideal-omega-unramified}.
\end{proof}

\begin{lemma}
\label{discriminant-lemma-kahler-different-complete-intersection}
Soient $A$ un anneau, $n \geq 1$ et
$f_1, \ldots, f_n \in A[x_1, \ldots, x_n]$.
Posons $B = A[x_1, \ldots, x_n]/(f_1, \ldots, f_n)$.
La différente de Kähler de $B$ sur $A$ est l'idéal
de $B$ engendré par $\det(\partial f_i/\partial x_j)$.
\end{lemma}

\begin{proof}
En effet, $\Omega_{B/A}$ admet la présentation
$$
\bigoplus\nolimits_{i = 1, \ldots, n} B f_i
\xrightarrow{\text{d}}
\bigoplus\nolimits_{j = 1, \ldots, n} B \text{d}x_j
\rightarrow \Omega_{B/A} \rightarrow 0
$$
d'après Algèbre, lemme \ref{algebra-lemma-differential-seq}.
\end{proof}



\section{La différente de Dedekind}
\label{discriminant-section-dedekind-different}

\noindent
Soit $A \to B$ un homomorphisme d'anneaux. On dit que {\it la différente de Dedekind est définie}
si $A$ est noethérien, $A \to B$ est fini,
tout non-diviseur de zéro de $A$ reste un non-diviseur de zéro dans $B$, et $K \to L$ est
étale, où $K = Q(A)$ et $L = B \otimes_A K$. Alors $K \subset L$ est
fini étale et
$$
\mathcal{L}_{B/A} = \{x \in L \mid \text{Trace}_{L/K}(bx) \in A
\text{ pour tout }b \in B\}
$$
est le module complémentaire de Dedekind. Dans cette situation, la
{\it différente de Dedekind} est
$$
\mathfrak{D}_{B/A} = \{x \in L \mid x\mathcal{L}_{B/A} \subset B\}
$$
considérée comme un $B$-sous-module de $L$.
D'après le lemme \ref{discriminant-lemma-dedekind-different-ideal}, la différente de Dedekind est un
idéal de $B$ si $A$ est normal ou si $B$ est plat sur $A$.

\begin{lemma}
\label{discriminant-lemma-dedekind-different-ideal}
Supposons que la différente de Dedekind de $A \to B$ soit définie. Considérons les assertions
\begin{enumerate}
\item $A \to B$ est plat ;
\item $A$ est un anneau normal ;
\item $\text{Trace}_{L/K}(B) \subset A$ ;
\item $1 \in \mathcal{L}_{B/A}$ ;
\item la différente de Dedekind $\mathfrak{D}_{B/A}$ est un idéal de $B$.
\end{enumerate}
On a alors (1) $\Rightarrow$ (3), (2) $\Rightarrow$ (3),
(3) $\Leftrightarrow$ (4) et (4) $\Rightarrow$ (5).
\end{lemma}

\begin{proof}
L'équivalence de (3) et (4) ainsi que
l'implication (4) $\Rightarrow$ (5) sont immédiates.

\medskip\noindent
Si $A \to B$ est plat, $\text{Trace}_{B/A} : B \to A$ est
définie et $\text{Trace}_{L/K}$ s'en déduit par changement de base. On a donc (3).

\medskip\noindent
Si $A$ est normal, $A$ est un produit fini d'anneaux intègres normaux,
et l'on se ramène donc au cas d'un anneau intègre normal. Alors $K$ est
le corps des fractions de $A$, et $L = \prod L_i$ est un produit fini
d'extensions finies séparables de $K$. On a
$\text{Trace}_{L/K}(b) = \sum \text{Trace}_{L_i/K}(b_i)$,
où $b_i \in L_i$ est l'image de $b$.
Puisque $b$ est entier sur $A$, car $B$ est fini sur $A$,
ces traces appartiennent à $A$. En effet, le
polynôme minimal de $b_i$ sur $K$ a ses coefficients dans $A$
(Algèbre, lemme \ref{algebra-lemma-minimal-polynomial-normal-domain}),
et $\text{Trace}_{L_i/K}(b_i)$ est un
multiple entier de l'un de ces coefficients
(Corps, lemme \ref{fields-lemma-trace-and-norm-from-minimal-polynomial}).
\end{proof}

\begin{lemma}
\label{discriminant-lemma-dedekind-complementary-module}
Si la différente de Dedekind de $A \to B$ est définie, alors
il existe un isomorphisme canonique
$\mathcal{L}_{B/A} \to \omega_{B/A}$.
\end{lemma}

\begin{proof}
Rappelons que $\omega_{B/A} = \Hom_A(B, A)$ puisque $A \to B$ est fini.
On envoie $x \in \mathcal{L}_{B/A}$ sur l'application
$b \mapsto \text{Trace}_{L/K}(bx)$.
Réciproquement, une application $A$-linéaire $\varphi : B \to A$
fournit une application $K$-linéaire $\varphi_K : L \to K$. Puisque $K \to L$ est fini
étale, la forme trace est non dégénérée
(lemme \ref{discriminant-lemma-discriminant}), et il existe donc $x \in L$ tel que
$\varphi_K(y) = \text{Trace}_{L/K}(xy)$ pour tout $y \in L$.
Alors $x \in \mathcal{L}_{B/A}$ s'envoie sur $\varphi$ dans $\omega_{B/A}$.
\end{proof}

\begin{lemma}
\label{discriminant-lemma-flat-dedekind-complementary-module-trace}
Si la différente de Dedekind de $A \to B$ est définie et si $A \to B$ est plat, alors
\begin{enumerate}
\item l'isomorphisme canonique $\mathcal{L}_{B/A} \to \omega_{B/A}$
envoie $1 \in \mathcal{L}_{B/A}$ sur l'élément trace
$\tau_{B/A} \in \omega_{B/A}$ ;
\item la différente de Dedekind est
$\mathfrak{D}_{B/A} = \{b \in B \mid b\omega_{B/A} \subset B\tau_{B/A}\}$.
\end{enumerate}
\end{lemma}

\begin{proof}
La première assertion
résulte de la démonstration du lemme \ref{discriminant-lemma-dedekind-different-ideal}
et du lemme \ref{discriminant-lemma-finite-flat-trace}.
La deuxième découle immédiatement de la première et des
définitions.
\end{proof}



\section{La différente}
\label{discriminant-section-different}

\noindent
La définition suivante est motivée par le fait qu'elle redonne la
différente de Dedekind dans le cas fini plat, comme nous le verrons ci-dessous.

\begin{definition}
\label{discriminant-definition-different}
Soit $f : Y \to X$ un morphisme plat localement quasi-fini de
schémas localement noethériens.
Soient $\omega_{Y/X}$ le module dualisant relatif et
$\tau_{Y/X} \in \Gamma(Y, \omega_{Y/X})$ l'élément trace
(remarques \ref{discriminant-remark-relative-dualizing-for-quasi-finite} et
\ref{discriminant-remark-relative-dualizing-for-flat-quasi-finite}).
L'annulateur de
$$
\Coker(\mathcal{O}_Y \xrightarrow{\tau_{Y/X}} \omega_{Y/X})
$$
est la {\it différente} de $Y/X$. C'est un idéal cohérent
$\mathfrak{D}_f \subset \mathcal{O}_Y$.
\end{definition}

\noindent
Nous généraliserons ceci dans la remarque \ref{discriminant-remark-different-generalization} ci-dessous.
Remarquons que $\mathfrak{D}_f$ est localement engendré par un élément si
$\omega_{Y/X}$ est un $\mathcal{O}_Y$-module inversible.
Énonçons d'abord l'égalité avec la différente de Dedekind.

\begin{lemma}
\label{discriminant-lemma-flat-agree-dedekind}
Soit $f : Y \to X$ un morphisme plat quasi-fini de schémas noethériens.
Soient $V = \Spec(B) \subset Y$, $U = \Spec(A) \subset X$
des sous-schémas ouverts affines tels que $f(V) \subset U$.
Si la différente de Dedekind de $A \to B$ est définie, alors
$$
\mathfrak{D}_f|_V = \widetilde{\mathfrak{D}_{B/A}}
$$
comme faisceaux cohérents d'idéaux sur $V$.
\end{lemma}

\begin{proof}
Cela résulte immédiatement des lemmes \ref{discriminant-lemma-dedekind-different-ideal} et
\ref{discriminant-lemma-flat-dedekind-complementary-module-trace}.
\end{proof}

\begin{lemma}
\label{discriminant-lemma-flat-gorenstein-agree-noether}
Soit $f : Y \to X$ un morphisme plat quasi-fini de schémas noethériens.
Soient $V = \Spec(B) \subset Y$, $U = \Spec(A) \subset X$
des sous-schémas ouverts affines tels que $f(V) \subset U$.
Si $\omega_{Y/X}|_V$ est inversible, c'est-à-dire si $\omega_{B/A}$
est un $B$-module inversible, alors
$$
\mathfrak{D}_f|_V = \widetilde{\mathfrak{D}}
$$
comme faisceaux cohérents d'idéaux sur $V$, où
$\mathfrak{D} \subset B$ est la différente de Noether de $B$ sur $A$.
\end{lemma}

\begin{proof}
Considérons l'application
$$
\SheafHom_{\mathcal{O}_Y}(\omega_{Y/X}, \mathcal{O}_Y)
\longrightarrow
\mathcal{O}_Y,\quad
\varphi \longmapsto \varphi(\tau_{Y/X})
$$
Son image correspond à la différente de Noether
sur les ouverts affines, voir le lemme \ref{discriminant-lemma-noether-different-flat-quasi-finite}.
Le résultat découle donc du fait élémentaire suivant : étant donné
un module inversible $\omega$ et une section globale $\tau$,
l'image de
$\tau : \SheafHom(\omega, \mathcal{O}) = \omega^{\otimes -1} \to \mathcal{O}$
est égale à l'annulateur de $\Coker(\tau : \mathcal{O} \to \omega)$.
\end{proof}

\begin{lemma}
\label{discriminant-lemma-base-change-different}
Considérons un diagramme cartésien de schémas noethériens
$$
\xymatrix{
Y' \ar[d]_{f'} \ar[r] & Y \ar[d]^f \\
X' \ar[r]^g & X
}
$$
où $f$ est plat et quasi-fini. Soit $R \subset Y$, resp.\ $R' \subset Y'$,
le sous-schéma fermé défini par la différente
$\mathfrak{D}_f$, resp.\ $\mathfrak{D}_{f'}$.
Alors $Y' \to Y$ induit une immersion fermée bijective $R' \to R \times_Y Y'$.
Si $g$ est plat ou si $\omega_{Y/X}$ est inversible, alors
$R' = R \times_Y Y'$.
\end{lemma}

\begin{proof}
On se ramène immédiatement au cas où $X$, $X'$, $Y$, $Y'$
sont affines. Autrement dit, on dispose d'un diagramme cocartésien d'anneaux
noethériens
$$
\xymatrix{
B' & B \ar[l] \\
A' \ar[u] & A \ar[l] \ar[u]
}
$$
où $A \to B$ est plat et quasi-fini. L'application de changement de base
$\omega_{B/A} \otimes_B B' \to \omega_{B'/A'}$ est un isomorphisme
(lemme \ref{discriminant-lemma-dualizing-base-change-of-flat}) et envoie
l'élément trace $\tau_{B/A}$ sur l'élément trace $\tau_{B'/A'}$
(lemme \ref{discriminant-lemma-trace-base-change}).
Le $B$-module de type fini $Q = \Coker(\tau_{B/A} : B \to \omega_{B/A})$
vérifie donc $Q \otimes_B B' = \Coker(\tau_{B'/A'} : B' \to \omega_{B'/A'})$.
Ainsi $\mathfrak{D}_{B/A}B' \subset \mathfrak{D}_{B'/A'}$, ce qui signifie
que l'on obtient l'immersion fermée $R' \to R \times_Y Y'$.
Puisque $R = \text{Supp}(Q)$ et $R' = \text{Supp}(Q \otimes_B B')$
(Algèbre, lemme \ref{algebra-lemma-support-closed}),
on voit que $R' \to R \times_Y Y'$ est bijective d'après
Algèbre, lemme \ref{algebra-lemma-support-base-change}.
L'égalité $\mathfrak{D}_{B/A}B' = \mathfrak{D}_{B'/A'}$ est vraie
si $B \to B'$ est plat, par exemple si $A \to A'$ est plat, voir
Algèbre, lemme \ref{algebra-lemma-annihilator-flat-base-change}.
Enfin, si $\omega_{B/A}$ est inversible, on peut localiser
et supposer $\omega_{B/A} = B \lambda$. En écrivant $\tau_{B/A} = b\lambda$,
on voit que $Q = B/bB$ et $\mathfrak{D}_{B/A} = bB$.
Le même raisonnement sur $B'$
donne $\mathfrak{D}_{B'/A'} = bB'$, ce qui démontre le lemme.
\end{proof}

\begin{lemma}
\label{discriminant-lemma-norm-different-in-discriminant}
Soit $f : Y \to X$ un morphisme fini plat de schémas noethériens.
Alors $\text{Norm}_f : f_*\mathcal{O}_Y \to \mathcal{O}_X$ envoie
$f_*\mathfrak{D}_f$ dans le faisceau d'idéaux du discriminant $D_f$.
\end{lemma}

\begin{proof}
L'application norme est construite dans
Diviseurs, lemme \ref{divisors-lemma-finite-locally-free-has-norm},
et le discriminant de $f$ dans la section \ref{discriminant-section-discriminant}.
La question est locale sur les ouverts affines ; on peut donc supposer $X = \Spec(A)$,
$Y = \Spec(B)$ et $f$ donné par un homomorphisme fini localement libre d'anneaux $A \to B$.
En localisant davantage, on peut supposer $B$ libre de type fini comme $A$-module.
Choisissons une base $b_1, \ldots, b_n \in B$ de $B$ comme $A$-module.
Notons $b_1^\vee, \ldots, b_n^\vee$ la base duale de
$\omega_{B/A} = \Hom_A(B, A)$ comme $A$-module.
Puisque la norme de $b$ est le déterminant de $b : B \to B$ comme
application $A$-linéaire, on a
$\text{Norm}_{B/A}(b) = \det(b_i^\vee(bb_j))$.
Le discriminant est le sous-schéma fermé principal de $\Spec(A)$
défini par $\det(\text{Trace}_{B/A}(b_ib_j))$.
Si $b \in \mathfrak{D}_{B/A}$, alors
il existe $c_i \in B$ tels que
$b \cdot b_i^\vee = c_i \cdot \text{Trace}_{B/A}$, où
le point désigne la structure de $B$-module sur $\omega_{B/A}$.
Écrivons $c_i = \sum a_{il} b_l$.
On a
\begin{align*}
\text{Norm}_{B/A}(b)
& =
\det(b_i^\vee(bb_j)) \\
& =
\det( (b \cdot b_i^\vee)(b_j)) \\
& =
\det((c_i \cdot \text{Trace}_{B/A})(b_j)) \\
& =
\det(\text{Trace}_{B/A}(c_ib_j)) \\
& =
\det(a_{il}) \det(\text{Trace}_{B/A}(b_l b_j))
\end{align*}
ce qui démontre le lemme.
\end{proof}

\begin{lemma}
\label{discriminant-lemma-different-ramification}
Soit $f : Y \to X$ un morphisme plat quasi-fini de schémas noethériens.
Le sous-schéma fermé $R \subset Y$ défini par la différente $\mathfrak{D}_f$
est exactement l'ensemble des points où $f$ n'est pas étale
(ou, de façon équivalente, est ramifié).
\end{lemma}

\begin{proof}
Puisque $f$ est de présentation finie et plat, il est étale
en un point si et seulement s'il est non ramifié en ce point. De plus, la
formation du lieu des points ramifiés commute au changement de base.
Voir Morphismes, section \ref{morphisms-section-etale}, et notamment
Morphismes, lemme \ref{morphisms-lemma-set-points-where-fibres-etale}.
D'après le lemme \ref{discriminant-lemma-base-change-different}, la formation de $R$ commute,
ensemblistement, au changement de base. Il suffit donc de démontrer le
lemme lorsque $X$ est le spectre d'un corps. D'autre part, la
construction de $(\omega_{Y/X}, \tau_{Y/X})$ est locale sur $Y$.
Puisque $Y$ est un espace discret fini (étant quasi-fini
sur un corps), on peut supposer que $Y$ possède un unique point.

\medskip\noindent
Écrivons $X = \Spec(k)$ et $Y = \Spec(B)$, où $k$ est un corps et $B$ une
$k$-algèbre locale finie. Si $Y \to X$ est étale, alors
$B$ est une extension finie séparable de $k$, et l'élément
trace $\text{Trace}_{B/k}$ est un élément de base de $\omega_{B/k}$
d'après Corps, lemme \ref{fields-lemma-separable-trace-pairing}.
Ainsi $\mathfrak{D}_{B/k} = B$ dans ce cas.
Réciproquement, si $\mathfrak{D}_{B/k} = B$, le
lemme \ref{discriminant-lemma-norm-different-in-discriminant}
et le fait que la norme de $1$ est $1$ montrent que le
discriminant est vide. Par conséquent,
$Y \to X$ est étale d'après le lemme \ref{discriminant-lemma-discriminant}.
\end{proof}

\begin{lemma}
\label{discriminant-lemma-norm-different-is-discriminant}
Soit $f : Y \to X$ un morphisme plat quasi-fini de schémas noethériens.
Soit $R \subset Y$ le sous-schéma fermé défini par $\mathfrak{D}_f$.
\begin{enumerate}
\item Si $\omega_{Y/X}$ est inversible,
alors $R$ est un sous-schéma fermé localement principal de $Y$.
\item Si $\omega_{Y/X}$ est inversible et $f$ fini, alors
la norme de $R$ est le discriminant $D_f$ de $f$.
\item Si $\omega_{Y/X}$ est inversible et $f$
est étale aux points associés de $Y$, alors $R$
est un diviseur de Cartier effectif et l'on a un
isomorphisme $\mathcal{O}_Y(R) = \omega_{Y/X}$.
\end{enumerate}
\end{lemma}

\begin{proof}
Preuve de (1). On peut travailler localement sur $Y$, donc supposer
$\omega_{Y/X}$ libre de rang $1$. Écrivons $\omega_{Y/X} = \mathcal{O}_Y\lambda$.
On peut alors écrire $\tau_{Y/X} = h \lambda$, et l'on voit que
$R$ est défini par $h$, c'est-à-dire que $R$ est localement principal.

\medskip\noindent
Preuve de (2). On peut supposer que $Y \to X$ est donné par un homomorphisme fini libre
d'anneaux $A \to B$ et que $\omega_{B/A}$ est libre de rang $1$ comme $B$-module.
Choisissons une $B$-base $\lambda$ de $\omega_{B/A}$ et écrivons
$\text{Trace}_{B/A} = b \cdot \lambda$ avec $b \in B$.
Alors $\mathfrak{D}_{B/A} = (b)$, et $D_f$ est défini par
$\det(\text{Trace}_{B/A}(b_ib_j))$, où $b_1, \ldots, b_n$ est une
base de $B$ comme $A$-module. Soit $b_1^\vee, \ldots, b_n^\vee$
la base duale.
En écrivant $b_i^\vee = c_i \cdot \lambda$, on voit que
$c_1, \ldots, c_n$ est également une base de $B$.
Ainsi, en posant $c_i = \sum a_{il}b_l$, on voit que $\det(a_{il})$
est inversible dans $A$. Il est clair que
$b \cdot b_i^\vee = c_i \cdot \text{Trace}_{B/A}$ ;
on conclut donc du calcul fait dans la démonstration du
lemme \ref{discriminant-lemma-norm-different-in-discriminant}
que $\text{Norm}_{B/A}(b)$ est le produit d'une unité par
$\det(\text{Trace}_{B/A}(b_ib_j))$.

\medskip\noindent
Preuve de (3). Avec les notations ci-dessus, le
lemme \ref{discriminant-lemma-different-ramification} et l'hypothèse montrent
que $h$ ne s'annule pas aux
points associés de $Y$, ce qui entraîne que $h$ est un non-diviseur de zéro.
L'isomorphisme canonique envoie $1$ sur $\tau_{Y/X}$, voir
Diviseurs, lemme \ref{divisors-lemma-characterize-OD}.
\end{proof}







\section{Morphismes quasi-finis syntomiques}
\label{discriminant-section-quasi-finite-syntomic}

\noindent
Cette section établit qu'un morphisme quasi-fini syntomique
possède un module dualisant relatif inversible.

\begin{lemma}
\label{discriminant-lemma-syntomic-quasi-finite}
Soit $f : Y \to X$ un morphisme de schémas. Les conditions suivantes sont équivalentes :
\begin{enumerate}
\item $f$ est localement quasi-fini et syntomique ;
\item $f$ est localement quasi-fini, plat et localement d'intersection
complète ;
\item $f$ est localement quasi-fini, plat, localement de présentation finie,
et les fibres de $f$ sont localement des intersections complètes ;
\item $f$ est localement quasi-fini et, pour tout $y \in Y$, il existe des
ouverts affines $y \in V = \Spec(B) \subset Y$, $U = \Spec(A) \subset X$
avec $f(V) \subset U$, un entier $n$ et
$h, f_1, \ldots, f_n \in A[x_1, \ldots, x_n]$ tels que
$B = A[x_1, \ldots, x_n, 1/h]/(f_1, \ldots, f_n)$ ;
\item pour tout $y \in Y$, il existe des ouverts affines
$y \in V = \Spec(B) \subset Y$, $U = \Spec(A) \subset X$
avec $f(V) \subset U$ tels que $A \to B$ soit une intersection complète globale
relative de la forme $B = A[x_1, \ldots, x_n]/(f_1, \ldots, f_n)$ ;
\item $f$ est localement quasi-fini, plat, localement de présentation finie,
et $\NL_{Y/X}$ est de Tor-amplitude contenue dans $[-1, 0]$ ;
\item $f$ est plat, localement de présentation finie,
$\NL_{Y/X}$ est parfait de rang $0$ et de Tor-amplitude contenue dans $[-1, 0]$.
\end{enumerate}
\end{lemma}

\begin{proof}
L'équivalence de (1) et (2) est donnée par
Compléments sur les morphismes, lemme \ref{more-morphisms-lemma-flat-lci}.
L'équivalence de (1) et (3) est donnée par
Morphismes, lemme \ref{morphisms-lemma-syntomic-flat-fibres}.

\medskip\noindent
Si $A \to B$ est comme dans (4), alors
$B = A[x, x_1, \ldots, x_n]/(xh - 1, f_1, \ldots, f_n)$
est une intersection complète globale relative, voir Algèbre, définition
\ref{algebra-definition-relative-global-complete-intersection}.
Ainsi (4) entraîne (5).
Il est clair que (5) entraîne (4).

\medskip\noindent
La condition (5) entraîne (1) : d'après
Algèbre, lemme \ref{algebra-lemma-relative-global-complete-intersection},
une intersection complète globale relative est syntomique, et
la définition d'une intersection complète globale relative
garantit qu'une telle intersection à
$n$ variables et $n$ équations est quasi-finie, voir
Algèbre, définition
\ref{algebra-definition-relative-global-complete-intersection} et
lemme \ref{algebra-lemma-isolated-point-fibre}.

\medskip\noindent
Algèbre, lemme \ref{algebra-lemma-syntomic}, ou
Morphismes, lemme \ref{morphisms-lemma-syntomic-locally-standard-syntomic},
montre que (1) entraîne (5).

\medskip\noindent
Compléments sur les morphismes, lemme \ref{more-morphisms-lemma-flat-fp-NL-lci}, montre que
(6) équivaut à (1). Si les conditions équivalentes (1) -- (6) sont satisfaites,
on voit que, localement sur les ouverts affines, $Y \to X$ est donné par une intersection complète
globale relative $B = A[x_1, \ldots, x_n]/(f_1, \ldots, f_n)$
ayant autant de variables que
d'équations. Cette présentation donne
$$
\NL_{B/A} =\left(
(f_1, \ldots, f_n)/(f_1, \ldots, f_n)^2
\longrightarrow
\bigoplus\nolimits_{i = 1, \ldots, n} B \text{d} x_i\right)
$$
D'après Algèbre, lemme
\ref{algebra-lemma-relative-global-complete-intersection-conormal},
le module $(f_1, \ldots, f_n)/(f_1, \ldots, f_n)^2$
est libre, de base formée des classes des éléments
$f_1, \ldots, f_n$. Ainsi $\NL_{B/A}$ est de rang $0$, de même que $\NL_{Y/X}$.
On voit ainsi que (1) -- (6) entraînent (7).

\medskip\noindent
Enfin, supposons (7). D'après
Compléments sur les morphismes, lemme \ref{more-morphisms-lemma-flat-fp-NL-lci},
le morphisme $f$ est syntomique. Sur des ouverts affines convenables,
$f$ est donc donné par une intersection complète globale relative
$A \to B = A[x_1, \ldots, x_n]/(f_1, \ldots, f_m)$, voir
Morphismes, lemme \ref{morphisms-lemma-syntomic-locally-standard-syntomic}.
Exactement comme ci-dessus, on voit que $\NL_{B/A}$ est un complexe parfait
de rang $n - m$. Ainsi $n = m$ et la condition (5) est satisfaite.
La démonstration est achevée.
\end{proof}

\begin{lemma}
\label{discriminant-lemma-characterize-invertible}
Inversibilité du module dualisant relatif.
\begin{enumerate}
\item Si $A \to B$ est un homomorphisme quasi-fini plat d'anneaux noethériens,
alors $\omega_{B/A}$ est un $B$-module inversible si et seulement si
$\omega_{B \otimes_A \kappa(\mathfrak p)/\kappa(\mathfrak p)}$
est un $B \otimes_A \kappa(\mathfrak p)$-module inversible
pour tout idéal premier $\mathfrak p \subset A$.
\item Si $Y \to X$ est un morphisme quasi-fini plat de
schémas noethériens, alors $\omega_{Y/X}$ est inversible
si et seulement si $\omega_{Y_x/x}$ est inversible pour tout $x \in X$.
\end{enumerate}
\end{lemma}

\begin{proof}
Preuve de (1). Comme $A \to B$ est plat, le module
$\omega_{B/A}$ est $A$-plat, voir le lemme \ref{discriminant-lemma-dualizing-base-flat-flat}.
Ainsi $\omega_{B/A}$ est un $B$-module inversible si et seulement si
$\omega_{B/A} \otimes_A \kappa(\mathfrak p)$
est un $B \otimes_A \kappa(\mathfrak p)$-module inversible pour
tout idéal premier $\mathfrak p \subset A$, voir Compléments sur les morphismes, lemme
\ref{more-morphisms-lemma-flat-and-free-at-point-fibre}.
Toujours par platitude de $A \to B$, la
formation de $\omega_{B/A}$ commute au changement de base, voir le
lemme \ref{discriminant-lemma-dualizing-base-change-of-flat}.
On voit donc que, sous l'hypothèse de platitude, l'inversibilité du module dualisant relatif
équivaut à l'inversibilité
du module dualisant relatif pour les homomorphismes
$\kappa(\mathfrak p) \to B \otimes_A \kappa(\mathfrak p)$.

\medskip\noindent
L'assertion (2) découle de (1) et de ce que, localement sur les ouverts affines,
les modules dualisants sont donnés par leurs analogues algébriques, voir la
remarque \ref{discriminant-remark-relative-dualizing-for-quasi-finite}.
\end{proof}

\begin{lemma}
\label{discriminant-lemma-dim-zero-global-complete-intersection-over-field}
Soit $k$ un corps. Soit $B = k[x_1, \ldots, x_n]/(f_1, \ldots, f_n)$
une intersection complète globale sur $k$, de dimension $0$.
Alors $\omega_{B/k}$ est inversible.
\end{lemma}

\begin{proof}
D'après le lemme de normalisation de Noether, voir
Algèbre, lemme \ref{algebra-lemma-Noether-normalization},
il existe une injection finie $k \to B$, c'est-à-dire
$\dim_k(B) < \infty$. Ainsi $\omega_{B/k} = \Hom_k(B, k)$
comme $B$-module.
D'après Complexes dualisants, lemme \ref{dualizing-lemma-dualizing-finite},
$R\Hom(B, k)$ est un complexe dualisant de $B$,
et Complexes dualisants, lemme \ref{dualizing-lemma-RHom-ext},
montre que $R\Hom(B, k)$ est égal à $\omega_{B/k}$
placé en degré $0$. Il suffit donc de montrer que
$B$ est de Gorenstein
(Complexes dualisants, lemme \ref{dualizing-lemma-gorenstein}).
C'est le cas d'après Complexes dualisants, lemme
\ref{dualizing-lemma-gorenstein-lci}.
\end{proof}

\begin{lemma}
\label{discriminant-lemma-dualizing-syntomic-quasi-finite}
Soit $f : Y \to X$ un morphisme de schémas localement noethériens. Si $f$
vérifie les conditions équivalentes du lemme \ref{discriminant-lemma-syntomic-quasi-finite},
alors $\omega_{Y/X}$ est un $\mathcal{O}_Y$-module inversible.
\end{lemma}

\begin{proof}
On peut supposer que $A \to B$ est une intersection complète globale
relative de la forme $B = A[x_1, \ldots, x_n]/(f_1, \ldots, f_n)$,
et il faut montrer que $\omega_{B/A}$ est inversible.
Cela résulte de la combinaison des lemmes \ref{discriminant-lemma-characterize-invertible} et
\ref{discriminant-lemma-dim-zero-global-complete-intersection-over-field}.
\end{proof}

\begin{example}
\label{discriminant-example-universal-quasi-finite-syntomic}
Soient $n \geq 1$ et $d \geq 1$ des entiers. Soit $T$ l'ensemble des
multi-indices $E = (e_1, \ldots, e_n)$ tels que $e_i \geq 0$ et
$\sum e_i \leq d$. Considérons l'anneau
$$
A = \mathbf{Z}[a_{i, E} ; 1 \leq i \leq n, E \in T]
$$
Dans $A[x_1, \ldots, x_n]$, considérons les éléments
$f_i = \sum_{E \in T} a_{i, E} x^E$, où $x^E = x_1^{e_1} \ldots x_n^{e_n}$
suivant la notation usuelle. Considérons la $A$-algèbre
$$
B = A[x_1, \ldots, x_n]/(f_1, \ldots, f_n)
$$
Notons $X_{n, d} = \Spec(A)$ et soit $Y_{n, d} \subset \Spec(B)$
le plus grand sous-schéma ouvert sur lequel la restriction du
morphisme $\Spec(B) \to \Spec(A) = X_{n, d}$ est quasi-finie, voir
Algèbre, lemme \ref{algebra-lemma-quasi-finite-open}.
\end{example}

\begin{lemma}
\label{discriminant-lemma-universal-quasi-finite-syntomic-etale}
Avec les notations de l'exemple \ref{discriminant-example-universal-quasi-finite-syntomic},
les schémas $X_{n, d}$ et $Y_{n, d}$ sont réguliers et irréductibles,
le morphisme $Y_{n, d} \to X_{n, d}$ est localement quasi-fini et
syntomique, et il existe un sous-schéma ouvert dense $V \subset Y_{n, d}$
tel que $Y_{n, d} \to X_{n, d}$ se restreigne en un morphisme étale
$V \to X_{n, d}$.
\end{lemma}

\begin{proof}
Le schéma $X_{n, d}$ est le spectre de l'anneau de polynômes $A$.
Ainsi $X_{n, d}$ est régulier et irréductible. Puisque l'on peut écrire
$$
f_i = a_{i, (0, \ldots, 0)} +
\sum\nolimits_{E \in T, E \not = (0, \ldots, 0)} a_{i, E} x^E
$$
on voit que l'anneau $B$ est isomorphe à l'anneau de polynômes
en $x_1, \ldots, x_n$ et en les éléments $a_{i, E}$ tels que
$E \not = (0, \ldots, 0)$. Ainsi $\Spec(B)$ est un schéma irréductible et
régulier, de même que l'ouvert $Y_{n, d}$. Le morphisme
$Y_{n, d} \to X_{n, d}$ est localement quasi-fini et syntomique d'après le
lemme \ref{discriminant-lemma-syntomic-quasi-finite}. Pour trouver $V$, il suffit
de trouver un point où $Y_{n, d} \to X_{n, d}$ est étale
(le lieu des points où un morphisme est étale est ouvert par
définition). Il suffit donc de trouver un point de $X_{n, d}$
où la fibre de $Y_{n, d} \to X_{n, d}$ est non vide et étale, voir
Morphismes, lemme \ref{morphisms-lemma-etale-at-point}. Choisissons
le point correspondant à l'homomorphisme d'anneaux $\chi : A \to \mathbf{Q}$
envoyant $f_i$ sur $x_i^d - 1$. Alors
$$
B \otimes_{A, \chi} \mathbf{Q} =
\mathbf{Q}[x_1, \ldots, x_n]/(x_1^d - 1, \ldots, x_n^d - 1)
$$
est une algèbre étale non nulle sur $\mathbf{Q}$.
\end{proof}

\begin{lemma}
\label{discriminant-lemma-locally-comes-from-universal}
Soit $f : Y \to X$ un morphisme de schémas. Si $f$ vérifie les conditions équivalentes
du lemme \ref{discriminant-lemma-syntomic-quasi-finite}, alors, pour tout
$y \in Y$, il existe $n, d$ et un diagramme commutatif
$$
\xymatrix{
Y \ar[d] &
V \ar[d] \ar[l] \ar[r] &
Y_{n, d} \ar[d] \\
X & U \ar[l] \ar[r] &
X_{n, d}
}
$$
où $U \subset X$ et $V \subset Y$ sont ouverts, avec $y \in V$,
où $Y_{n, d} \to X_{n, d}$
est comme dans l'exemple \ref{discriminant-example-universal-quasi-finite-syntomic}, et
où le carré de droite est cartésien.
\end{lemma}

\begin{proof}
D'après le lemme \ref{discriminant-lemma-syntomic-quasi-finite},
on peut choisir $U$ et $V$ affines, de sorte que
$U = \Spec(R)$ et $V = \Spec(S)$ avec
$S = R[y_1, \ldots, y_n]/(g_1, \ldots, g_n)$.
Avec les notations de l'exemple \ref{discriminant-example-universal-quasi-finite-syntomic},
si l'on choisit $d$ assez grand, on peut écrire chaque $g_i$ sous la forme
$g_i = \sum_{E \in T} g_{i, E}y^E$ avec $g_{i, E} \in R$.
L'application $A \to R$ envoyant $a_{i, E}$ sur $g_{i, E}$
et l'application $B \to S$ envoyant $x_i \to y_i$ donnent alors un diagramme cocartésien
d'anneaux
$$
\xymatrix{
S & B \ar[l] \\
R \ar[u] & A \ar[l] \ar[u]
}
$$
ce qui démontre le lemme.
\end{proof}











\section{Morphismes finis syntomiques}
\label{discriminant-section-finite-syntomic}

\noindent
Cette section est l'analogue de la section \ref{discriminant-section-quasi-finite-syntomic}
pour les morphismes finis syntomiques.

\begin{lemma}
\label{discriminant-lemma-syntomic-finite}
Soit $f : Y \to X$ un morphisme de schémas. Les conditions suivantes sont équivalentes :
\begin{enumerate}
\item $f$ est fini et syntomique ;
\item $f$ est fini, plat et localement d'intersection complète ;
\item $f$ est fini, plat, localement de présentation finie,
et les fibres de $f$ sont localement des intersections complètes ;
\item $f$ est fini et, pour tout $x \in X$, il existe un
ouvert affine $x \in U = \Spec(A) \subset X$, un entier $n$
et $f_1, \ldots, f_n \in A[x_1, \ldots, x_n]$ tels que
$f^{-1}(U)$ soit isomorphe au spectre de
$A[x_1, \ldots, x_n]/(f_1, \ldots, f_n)$ ;
\item $f$ est fini, plat, localement de présentation finie,
et $\NL_{X/Y}$ est de Tor-amplitude contenue dans $[-1, 0]$ ;
\item $f$ est fini, plat, localement de présentation finie, et
$\NL_{X/Y}$ est parfait de rang $0$ et de Tor-amplitude contenue dans $[-1, 0]$.
\end{enumerate}
\end{lemma}

\begin{proof}
L'équivalence de (1), (2), (3), (5) et (6)
ainsi que l'implication (4) $\Rightarrow$ (1) résultent immédiatement
du lemme \ref{discriminant-lemma-syntomic-quasi-finite}. Supposons les conditions équivalentes
(1), (2), (3), (5), (6) satisfaites.
Choisissons un point $x \in X$ et un voisinage ouvert affine $U = \Spec(A)$
de $x$ dans $X$, et supposons que $x$ corresponde à l'idéal premier
$\mathfrak p \subset A$. Écrivons $f^{-1}(U) = \Spec(B)$.
Écrivons $B = A[x_1, \ldots, x_n]/I$. Puisque $\NL_{B/A}$
est parfait de Tor-amplitude contenue dans $[-1, 0]$ d'après (6),
on voit que $I/I^2$ est un $B$-module fini localement libre
de rang $n$. Puisque $B_\mathfrak p$ est semi-local,
$(I/I^2)_\mathfrak p$ est libre de rang $n$, voir
Algèbre, lemme \ref{algebra-lemma-locally-free-semi-local-free}.
Ainsi, en remplaçant $A$ par sa localisation principale en
un élément n'appartenant pas à $\mathfrak p$, on peut supposer que $I/I^2$
est un $B$-module libre de rang $n$.
D'après Algèbre, lemme \ref{algebra-lemma-huber},
on peut donc trouver une présentation de $B$ sur $A$
ayant autant de variables que d'équations. Autrement dit,
on peut supposer $B = A[x_1, \ldots, x_n]/(f_1, \ldots, f_n)$.
Cela démontre (4).
\end{proof}

\begin{example}
\label{discriminant-example-universal-finite-syntomic}
Soit $d \geq 1$ un entier. Considérons des variables
$a_{ij}^l$ pour $1 \leq i, j, l \leq d$, et notons
$$
A_d = \mathbf{Z}[a_{ij}^k]/J
$$
où $J$ est l'idéal engendré par les éléments
$$
\left\{
\begin{matrix}
\sum_l a_{ij}^la_{lk}^m - \sum_l a_{il}^ma_{jk}^l & \forall i, j, k, m \\
a_{ij}^k - a_{ji}^k & \forall i, j, k \\
a_{i1}^j - \delta_{ij} & \forall i, j
\end{matrix}
\right.
$$
où $\delta_{ij}$ désigne le symbole de Kronecker.
On définit une $A_d$-algèbre $B_d$ comme suit : comme $A_d$-module, on pose
$$
B_d = A_d e_1 \oplus \ldots \oplus A_d e_d
$$
La structure d'algèbre est donnée par $A_d \to B_d$ envoyant $1$ sur $e_1$.
La multiplication sur $B_d$ est l'application $A_d$-bilinéaire
$$
m : B_d \times B_d \longrightarrow B_d, \quad
m(e_i, e_j) = \sum a_{ij}^k e_k
$$
On vérifie sans difficulté que les relations ci-dessus
imposent exactement une structure de $A_d$-algèbre.
Le morphisme
$$
\pi_d : Y_d = \Spec(B_d) \longrightarrow \Spec(A_d) = X_d
$$
est le morphisme fini libre « universel » de rang $d$.
\end{example}

\begin{lemma}
\label{discriminant-lemma-universal-finite-syntomic}
Avec les notations de l'exemple \ref{discriminant-example-universal-finite-syntomic},
il existe un sous-schéma ouvert $U_d \subset X_d$ possédant la propriété suivante :
un morphisme de schémas $X \to X_d$ se factorise par $U_d$ si et seulement
si $Y_d \times_{X_d} X \to X$ est syntomique.
\end{lemma}

\begin{proof}
Rappelons qu'être syntomique équivaut à être plat et
localement d'intersection complète, voir
Compléments sur les morphismes, lemme \ref{more-morphisms-lemma-flat-lci}.
L'ensemble $W_d \subset Y_d$ des points où $\pi_d$ est de Koszul
est ouvert dans $Y_d$, et sa formation commute à tout changement de base, voir
Compléments sur les morphismes, lemme \ref{more-morphisms-lemma-base-change-lci-fibres}.
Puisque $\pi_d$ est fini, donc fermé,
$Z = \pi_d(Y_d \setminus W_d)$ est fermé. Comme on a manifestement $U_d = X_d \setminus Z$
et que sa formation commute au changement de base, le lemme
en résulte.
\end{proof}

\begin{lemma}
\label{discriminant-lemma-universal-finite-syntomic-smooth}
Avec les notations de l'exemple \ref{discriminant-example-universal-finite-syntomic}
et $U_d$ comme dans le lemme \ref{discriminant-lemma-universal-finite-syntomic},
le schéma $U_d$ est lisse sur $\Spec(\mathbf{Z})$.
\end{lemma}

\begin{proof}
Utilisons Compléments sur les morphismes, lemme
\ref{more-morphisms-lemma-lifting-along-artinian-at-point},
pour montrer que $U_d \to \Spec(\mathbf{Z})$ est lisse.
Supposons donc que $\Spec(A) \to U_d$ soit un morphisme
et que $A' \to A$ soit une petite extension. Alors $B = A \otimes_{A_d} B_d$
est une $A$-algèbre finie libre, syntomique sur $A$
(par construction de $U_d$). D'après
Lissage des homomorphismes d'anneaux, proposition \ref{smoothing-proposition-lift-smooth},
il existe un homomorphisme syntomique d'anneaux $A' \to B'$ tel que
$B \cong B' \otimes_{A'} A$. Posons $e'_1 = 1 \in B'$. Pour $1 < i \leq d$,
choisissons des relèvements $e'_i \in B'$ des éléments
$1 \otimes e_i \in A \otimes_{A_d} B_d = B$. Alors $e'_1, \ldots, e'_d$
est une base de $B'$ sur $A'$ (voir par exemple Algèbre, lemme
\ref{algebra-lemma-local-artinian-basis-when-flat}).
On peut donc écrire $e'_i e'_j = \sum \alpha_{ij}^l e'_l$ pour d'uniques
éléments $\alpha_{ij}^l \in A'$ vérifiant les relations
$\sum_l \alpha_{ij}^l \alpha_{lk}^m = \sum_l \alpha_{il}^m \alpha _{jk}^l$
et $\alpha_{ij}^k = \alpha_{ji}^k$ et $\alpha_{i1}^j = \delta_{ij}$
dans $A'$. Cela détermine un morphisme $\Spec(A') \to X_d$ en
envoyant $a_{ij}^l \in A_d$ sur $\alpha_{ij}^l \in A'$. Ce morphisme
coïncide avec le morphisme donné $\Spec(A) \to U_d$. Puisque $\Spec(A')$
et $\Spec(A)$ ont le même espace topologique sous-jacent, on obtient
le relèvement voulu $\Spec(A') \to U_d$, et l'on
conclut que $U_d$ est lisse sur $\mathbf{Z}$.
\end{proof}

\begin{lemma}
\label{discriminant-lemma-universal-finite-syntomic-etale}
Avec les notations de l'exemple \ref{discriminant-example-universal-finite-syntomic},
considérons le sous-schéma ouvert $U'_d \subset X_d$ au-dessus duquel
$\pi_d$ est étale. Alors $U'_d$ est une partie dense de
l'ouvert $U_d$ du lemme \ref{discriminant-lemma-universal-finite-syntomic}.
\end{lemma}

\begin{proof}
En raisonnant exactement comme dans la démonstration du
lemme \ref{discriminant-lemma-universal-finite-syntomic}, à l'aide de
Morphismes, lemme \ref{morphisms-lemma-set-points-where-fibres-etale},
on obtient un plus grand ouvert $U'_d \subset X_d$ au-dessus duquel $\pi_d$ est
étale. De plus, un morphisme étale étant syntomique, on voit
que $U'_d \subset U_d$. Pour achever la démonstration, il faut montrer
que $U'_d \subset U_d$ est dense. Soit $u : \Spec(k) \to U_d$ un morphisme,
où $k$ est un corps. Soit $B = k \otimes_{A_d} B_d$ comme dans la
démonstration du lemme \ref{discriminant-lemma-universal-finite-syntomic-smooth}.
Nous allons montrer qu'il existe un anneau local intègre $A'$ de corps résiduel $k$
et une $A'$-algèbre finie syntomique $B'$ avec $B = k \otimes_{A'} B'$
dont la fibre générique est étale. Exactement comme au paragraphe précédent,
cela déterminera un morphisme $\Spec(A') \to U_d$ envoyant le
point générique dans $U'_d$ et le point fermé sur $u$, ce qui
achèvera la démonstration.

\medskip\noindent
D'après l'assertion (4) du lemme \ref{discriminant-lemma-syntomic-finite}, on peut choisir une présentation
$B = k[x_1, \ldots, x_n]/(f_1, \ldots, f_n)$.
Soit $d'$ le maximum des degrés totaux des polynômes $f_1, \ldots, f_n$.
Soit $Y_{n, d'} \to X_{n, d'}$ comme dans
l'exemple \ref{discriminant-example-universal-quasi-finite-syntomic}.
Par construction, il existe un morphisme $u' : \Spec(k) \to X_{n, d'}$
tel que
$$
\Spec(B) \cong Y_{n, d'} \times_{X_{n, d'}, u'} \Spec(k)
$$
Notons $A = \mathcal{O}_{X_{n, d'}, u'}^h$ le hensélisé de
l'anneau local de $X_{n, d'}$ à l'image de $u'$. On peut alors écrire
$$
Y_{n, d'} \times_{X_{n, d'}} \Spec(A) = Z \amalg W
$$
avec $Z \to \Spec(A)$ fini et $W \to \Spec(A)$ de fibre
fermée vide, voir
Algèbre, assertion (13) du lemme \ref{algebra-lemma-characterize-henselian},
ou la discussion dans Compléments sur les morphismes, section
\ref{more-morphisms-section-etale-localization}.
D'après le lemme \ref{discriminant-lemma-universal-quasi-finite-syntomic-etale},
l'anneau local $A$ est régulier (on utilise également ici
Compléments d'algèbre, lemme \ref{more-algebra-lemma-henselization-regular}),
et le morphisme $Z \to \Spec(A)$ est étale au-dessus du point générique de
$\Spec(A)$ (car celui-ci s'envoie sur le point générique de $X_{d, n'}$).
Par construction, $Z \times_{\Spec(A)} \Spec(k) \cong \Spec(B)$.
Cela donne le résultat voulu, à ceci près que l'application du
corps résiduel de $A$ vers $k$ n'est pas nécessairement un isomorphisme.
D'après Algèbre, lemme \ref{algebra-lemma-flat-local-given-residue-field},
il existe un homomorphisme local plat d'anneaux $A \to A'$ tel que le corps résiduel
de $A'$ soit $k$. Si $A'$ n'est pas intègre, on choisit un
idéal premier minimal $\mathfrak p \subset A'$ (situé au-dessus de
l'unique idéal premier minimal de $A$ par platitude) et l'on remplace
$A'$ par $A'/\mathfrak p$. On prend pour $B'$ l'unique $A'$-algèbre
telle que $Z \times_{\Spec(A)} \Spec(A') = \Spec(B')$.
La démonstration est achevée.
\end{proof}

\begin{remark}
\label{discriminant-remark-universal-finite-syntomic-smooth-top}
Soit $\pi_d : Y_d \to X_d$ comme dans
l'exemple \ref{discriminant-example-universal-finite-syntomic}.
Soit $U_d \subset X_d$ le plus grand ouvert au-dessus duquel
$V_d = \pi_d^{-1}(U_d)$ est fini syntomique, comme dans le
lemme \ref{discriminant-lemma-universal-finite-syntomic}.
Alors $V_d$ est également lisse sur $\mathbf{Z}$.
(Bien entendu, le morphisme $V_d \to U_d$ n'est pas lisse lorsque $d \geq 2$.)
En raisonnant comme dans la démonstration du lemme \ref{discriminant-lemma-universal-finite-syntomic-smooth},
on voit que cela correspond au problème de déformation
suivant : étant donné une petite extension $C' \to C$ et
une $C$-algèbre finie syntomique $B$ munie d'une section $B \to C$,
trouver une $C'$-algèbre finie syntomique $B'$ et une section $B' \to C'$
dont le produit tensoriel avec $C$ redonne $B \to C$.
D'après le lemme \ref{discriminant-lemma-syntomic-finite}, on peut écrire
$B = C[x_1, \ldots, x_n]/(f_1, \ldots, f_n)$ comme
une intersection complète globale relative.
Après un changement de coordonnées, on peut supposer
$x_1, \ldots, x_n$ dans le noyau de $B \to C$.
Les polynômes $f_i$ ont alors un terme constant nul.
Choisissons des relèvements quelconques $f'_i \in C'[x_1, \ldots, x_n]$ de $f_i$
ayant un terme constant nul. Alors
$B' = C'[x_1, \ldots, x_n]/(f'_1, \ldots, f'_n)$,
avec la section $B' \to C'$ envoyant $x_i$ sur zéro, convient.
\end{remark}

\begin{lemma}
\label{discriminant-lemma-locally-comes-from-universal-finite}
Soit $f : Y \to X$ un morphisme de schémas. Si $f$ vérifie les conditions équivalentes
du lemme \ref{discriminant-lemma-syntomic-finite}, alors, pour tout
$x \in X$, il existe un $d$ et un diagramme commutatif
$$
\xymatrix{
Y \ar[d] &
V \ar[d] \ar[l] \ar[r] &
V_d \ar[d] \ar[r] &
Y_d \ar[d]^{\pi_d}\\
X &
U \ar[l] \ar[r] &
U_d \ar[r] &
X_d
}
$$
possédant les propriétés suivantes :
\begin{enumerate}
\item $U \subset X$ est ouvert, $x \in U$ et $V = f^{-1}(U)$ ;
\item $\pi_d : Y_d \to X_d$ est comme dans
l'exemple \ref{discriminant-example-universal-finite-syntomic} ;
\item $U_d \subset X_d$ est comme dans le lemme \ref{discriminant-lemma-universal-finite-syntomic}
et $V_d = \pi_d^{-1}(U_d) \subset Y_d$ ;
\item le carré du milieu est cartésien.
\end{enumerate}
\end{lemma}

\begin{proof}
Choisissons un voisinage ouvert affine $U = \Spec(A) \subset X$ de $x$.
Écrivons $V = f^{-1}(U) = \Spec(B)$. Alors $B$ est un
$A$-module fini localement libre, et l'inclusion $A \subset B$ est localement un facteur direct.
En rétrécissant $U$, on peut donc choisir une base $1 = e_1, e_2, \ldots, e_d$
de $B$ comme $A$-module. Écrivons
$e_i e_j = \sum \alpha_{ij}^l e_l$ pour d'uniques
éléments $\alpha_{ij}^l \in A$ vérifiant les relations
$\sum_l \alpha_{ij}^l \alpha_{lk}^m = \sum_l \alpha_{il}^m \alpha _{jk}^l$
et $\alpha_{ij}^k = \alpha_{ji}^k$ et $\alpha_{i1}^j = \delta_{ij}$
dans $A$. Cela détermine un morphisme $\Spec(A) \to X_d$ en envoyant
$a_{ij}^l \in A_d$ sur $\alpha_{ij}^l \in A$. Par construction,
$V \cong \Spec(A) \times_{X_d} Y_d$. La définition de $U_d$
montre que $\Spec(A) \to X_d$ se factorise par $U_d$.
La démonstration est achevée.
\end{proof}










\section{Une formule pour la différente}
\label{discriminant-section-formula-different}

\noindent
Dans cette section, nous présentons les résultats dus à Tate dans \cite[appendice A]{Mazur-Roberts}.
Dans notre terminologie, ils montreront que la différente est
égale à la différente de Kähler pour un morphisme plat, quasi-fini,
localement d'intersection complète.
Calculons d'abord la différente de Noether dans un cas particulier.

\begin{lemma}
\label{discriminant-lemma-tate}
\begin{reference}
\cite[appendice]{Mazur-Roberts}
\end{reference}
Soit $A \to P$ un homomorphisme d'anneaux. Soit $f_1, \ldots, f_n \in P$ une
suite régulière au sens de Koszul. Supposons $B = P/(f_1, \ldots, f_n)$
plat sur $A$. Soit $g_1, \ldots, g_n \in P \otimes_A B$
une suite régulière au sens de Koszul engendrant le noyau de l'application de multiplication
$P \otimes_A B \to B$. Écrivons $f_i \otimes 1 = \sum g_{ij} g_j$.
Alors l'annulateur de $\Ker(B \otimes_A B \to B)$ est un idéal principal
engendré par l'image de $\det(g_{ij})$.
\end{lemma}

\begin{proof}
Le complexe de Koszul $K_\bullet = K(P, f_1, \ldots, f_n)$ est une résolution
de $B$ par des $P$-modules libres de type fini. Le complexe de Koszul
$M_\bullet = K(P \otimes_A B, g_1, \ldots, g_n)$ est une résolution
de $B$ par des $P \otimes_A B$-modules libres de type fini. Il existe un morphisme de
complexes
$$
K_\bullet \longrightarrow M_\bullet
$$
donné en degré $1$ par la matrice $(g_{ij})$ et
en degré $n$ par $\det(g_{ij})$. Voir
Compléments d'algèbre, lemme \ref{more-algebra-lemma-functorial}.
Comme $B$ est un $A$-module plat, on peut considérer $M_\bullet$ comme un complexe
de $P$-modules plats (via $P \to P \otimes_A B$, $p \mapsto p \otimes 1$).
On peut donc utiliser les deux complexes pour calculer $\text{Tor}_*^P(B, B)$,
et l'application affichée devient un quasi-isomorphisme après tensorisation
par $B$. Il est clair que $H_n(K_\bullet \otimes_P B) = B$.
D'autre part, $H_n(M_\bullet \otimes_P B)$ est le noyau de
$$
B \otimes_A B \xrightarrow{g_1, \ldots, g_n} (B \otimes_A B)^{\oplus n}
$$
Puisque $g_1, \ldots, g_n$ engendrent le noyau de $B \otimes_A B \to B$,
le lemme est démontré.
\end{proof}

\begin{lemma}
\label{discriminant-lemma-quasi-finite-complete-intersection}
Soient $A$ un anneau, $n \geq 1$ et
$h, f_1, \ldots, f_n \in A[x_1, \ldots, x_n]$.
Posons $B = A[x_1, \ldots, x_n, 1/h]/(f_1, \ldots, f_n)$.
Supposons que $B$ soit quasi-fini sur $A$.
Alors
\begin{enumerate}
\item $B$ est plat sur $A$ et $A \to B$ est une intersection complète locale
relative ;
\item l'annulateur $J$ de $I = \Ker(B \otimes_A B \to B)$
est libre de rang $1$ sur $B$ ;
\item la différente de Noether de $B$ sur $A$ est engendrée
par $\det(\partial f_i/\partial x_j)$ dans $B$.
\end{enumerate}
\end{lemma}

\begin{proof}
Remarquons que
$B = A[x, x_1, \ldots, x_n]/(xh - 1, f_1, \ldots, f_n)$
est une intersection complète globale relative sur $A$, voir
Algèbre, définition
\ref{algebra-definition-relative-global-complete-intersection}.
D'après Algèbre, lemme \ref{algebra-lemma-relative-global-complete-intersection},
on voit que $B$ est plat sur $A$.

\medskip\noindent
Posons $P' = A[x, x_1, \ldots, x_n]$ et
$P = P'/(xh - 1) = A[x_1, \ldots, x_n, 1/h]$.
On a alors $P' \to P \to B$.
D'après Compléments d'algèbre, lemme
\ref{more-algebra-lemma-relative-global-complete-intersection-koszul},
la suite $xh - 1, f_1, \ldots, f_n$ est régulière au sens de Koszul
dans $P'$. Puisque $xh - 1$ est une suite régulière au sens de Koszul de longueur
un dans $P'$ (par exemple d'après le même lemme), on en déduit que
$f_1, \ldots, f_n$ est une suite régulière au sens de Koszul dans $P$ d'après
Compléments d'algèbre, lemme \ref{more-algebra-lemma-truncate-koszul-regular}.

\medskip\noindent
Soit $g_i \in P \otimes_A B$ l'image de $x_i \otimes 1 - 1 \otimes x_i$.
Utilisons les notations abrégées $y_i = x_i \otimes 1$ et $z_i = 1 \otimes x_i$
dans $A[x_1, \ldots, x_n] \otimes_A A[x_1, \ldots, x_n]$,
de sorte que $g_i$ soit l'image de $y_i - z_i$. Pour un polynôme
$f \in A[x_1, \ldots, x_n]$, on écrit $f(y) = f \otimes 1$
et $f(z) = 1 \otimes f$ dans le produit tensoriel ci-dessus.
On a alors
$$
P \otimes_A B/(g_1, \ldots, g_n) =
\frac{A[y_1, \ldots, y_n, z_1, \ldots, z_n, \frac{1}{h(y)h(z)}]}
{(f_1(z), \ldots, f_n(z), y_1 - z_1, \ldots, y_n - z_n)}
$$
qui est manifestement isomorphe à $B$. Par les mêmes arguments
que ci-dessus, on trouve donc que $f_1(z), \ldots, f_n(z), y_1 - z_1, \ldots, y_n - z_n$
est une suite régulière au sens de Koszul dans
$A[y_1, \ldots, y_n, z_1, \ldots, z_n, \frac{1}{h(y)h(z)}]$.
La suite $f_1(z), \ldots, f_n(z)$ est régulière au sens de Koszul dans
$A[y_1, \ldots, y_n, z_1, \ldots, z_n, \frac{1}{h(y)h(z)}]$
par platitude de l'application
$$
P \longrightarrow A[y_1, \ldots, y_n, z_1, \ldots, z_n,
\textstyle{\frac{1}{h(y)h(z)}}],\quad x_i \longmapsto z_i
$$
et d'après Compléments d'algèbre, lemme
\ref{more-algebra-lemma-koszul-regular-flat-base-change}.
D'après Compléments d'algèbre, lemme \ref{more-algebra-lemma-truncate-koszul-regular},
on conclut que $g_1, \ldots, g_n$ est une suite régulière
dans $P \otimes_A B$.

\medskip\noindent
Nous avons maintenant vérifié toutes les hypothèses du lemme \ref{discriminant-lemma-tate}
ci-dessus, avec $P$, $f_1, \ldots, f_n$ et $g_i \in P \otimes_A B$ comme ci-dessus.
En particulier, l'annulateur $J$ de $I$ est libre, engendré par un
élément $\delta$ sur $B$.
Posons $f_{ij} = \partial f_i/\partial x_j \in A[x_1, \ldots, x_n]$.
Un calcul élémentaire montre que l'on peut écrire
$$
f_i(y) =
f_i(z_1 + g_1, \ldots, z_n + g_n) =
f_i(z) + \sum\nolimits_j f_{ij}(z) g_j +
\sum\nolimits_{j, j'} F_{ijj'}g_jg_{j'}
$$
avec $F_{ijj'} \in A[y_1, \ldots, y_n, z_1, \ldots, z_n]$.
En passant à l'image dans $P \otimes_A B$, les termes $f_i(z)$ deviennent
nuls et l'on obtient
$$
f_i \otimes 1 = \sum\nolimits_j
\left(1 \otimes f_{ij} + \sum\nolimits_{j'} F_{ijj'}g_{j'}\right)g_j
$$
On conclut donc du lemme \ref{discriminant-lemma-tate}
que $\delta = \det(g_{ij})$, avec
$g_{ij} = 1 \otimes f_{ij} + \sum_{j'} F_{ijj'}g_{j'}$.
Puisque $g_{j'}$ a une image nulle dans $B$, on en déduit
que l'image de $\det(\partial f_i/\partial x_j)$ dans $B$
engendre la différente de Noether de $B$ sur $A$.
\end{proof}

\begin{lemma}
\label{discriminant-lemma-different-syntomic-quasi-finite}
Soit $f : Y \to X$ un morphisme de schémas noethériens. Si $f$
vérifie les conditions équivalentes du lemme \ref{discriminant-lemma-syntomic-quasi-finite},
alors la différente $\mathfrak{D}_f$ de $f$ est la différente de Kähler
de $f$.
\end{lemma}

\begin{proof}
D'après les lemmes \ref{discriminant-lemma-flat-gorenstein-agree-noether} et
\ref{discriminant-lemma-dualizing-syntomic-quasi-finite},
la différente de $f$ coïncide, localement sur les ouverts affines, avec la
différente de Noether. Le lemme résulte alors du
calcul de la différente de Noether et de la différente de Kähler
sur les morceaux affines standards, effectué dans les
lemmes \ref{discriminant-lemma-kahler-different-complete-intersection} et
\ref{discriminant-lemma-quasi-finite-complete-intersection}.
\end{proof}

\begin{lemma}
\label{discriminant-lemma-different-quasi-finite-complete-intersection}
Soient $A$ un anneau, $n \geq 1$ et
$h, f_1, \ldots, f_n \in A[x_1, \ldots, x_n]$.
Posons $B = A[x_1, \ldots, x_n, 1/h]/(f_1, \ldots, f_n)$.
Supposons que $B$ soit quasi-fini sur $A$.
Il existe alors un isomorphisme $B \to \omega_{B/A}$
envoyant $\det(\partial f_i/\partial x_j)$ sur $\tau_{B/A}$.
\end{lemma}

\begin{proof}
Soit $J$ l'annulateur de $\Ker(B \otimes_A B \to B)$.
D'après le lemme \ref{discriminant-lemma-quasi-finite-complete-intersection},
l'homomorphisme $A \to B$ est plat et
$J$ est un $B$-module libre ayant un générateur $\xi$ dont l'image est
$\det(\partial f_i/\partial x_j)$ dans $B$.
Le lemme découle donc du
lemme \ref{discriminant-lemma-noether-different-flat-quasi-finite}
et du fait (lemme \ref{discriminant-lemma-dualizing-syntomic-quasi-finite})
que $\omega_{B/A}$ est un $B$-module inversible.
(Attention : il est nécessaire de montrer que $\omega_{B/A}$
est inversible, car un $B$-module de type fini $M$ tel
que $\Hom_B(M, B) \cong B$ n'est pas nécessairement libre.)
\end{proof}

\begin{example}
\label{discriminant-example-different-for-monogenic}
Soit $A$ un anneau noethérien. Soient $f, h \in A[x]$ tels que
$$
B = (A[x]/(f))_h = A[x, 1/h]/(f)
$$
soit quasi-fini sur $A$. Soit $f' \in A[x]$ la dérivée
de $f$ par rapport à $x$. L'idéal $\mathfrak{D} = (f') \subset B$
est la différente de Noether de $B$ sur $A$,
est la différente de Kähler de $B$ sur $A$, et
est l'idéal dont le faisceau quasi-cohérent d'idéaux associé est la
différente de $\Spec(B)$ sur $\Spec(A)$.
\end{example}

\begin{lemma}
\label{discriminant-lemma-discriminant-quasi-finite-morphism-smooth}
Soit $S$ un schéma noethérien. Soient $X$, $Y$ des schémas lisses
de dimension relative $n$ sur $S$. Soit $f : Y \to X$ un
morphisme localement quasi-fini sur $S$.
Alors $f$ est plat et le sous-schéma fermé $R \subset Y$
défini par la différente de $f$ est le sous-schéma fermé localement principal
défini par
$$
\wedge^n(\text{d}f) \in
\Gamma(Y,
(f^*\Omega^n_{X/S})^{\otimes -1} \otimes_{\mathcal{O}_Y} \Omega^n_{Y/S})
$$
Si $f$ est étale aux points associés de $Y$, alors $R$ est un
diviseur de Cartier effectif, et
$$
f^*\Omega^n_{X/S} \otimes_{\mathcal{O}_Y} \mathcal{O}(R) =
\Omega^n_{Y/S}
$$
comme faisceaux inversibles sur $Y$.
\end{lemma}

\begin{proof}
Pour montrer que $f$ est plat, il suffit de montrer que $Y_s \to X_s$
est plat pour tout $s \in S$ (Compléments sur les morphismes, lemme
\ref{more-morphisms-lemma-morphism-between-flat-Noetherian}).
La platitude de $Y_s \to X_s$ découle de
Algèbre, lemme \ref{algebra-lemma-CM-over-regular-flat}.
D'après Compléments sur les morphismes, lemme
\ref{more-morphisms-lemma-lci-permanence},
le morphisme $f$ est localement d'intersection complète.
L'assertion sur la différente découle donc de
l'assertion correspondante sur la différente de Kähler d'après le
lemme \ref{discriminant-lemma-different-syntomic-quasi-finite}.
Enfin, puisque l'on a la suite exacte
$$
f^*\Omega_{X/S} \xrightarrow{\text{d}f} \Omega_{Y/S} \to \Omega_{Y/X} \to 0
$$
d'après Morphismes, lemme \ref{morphisms-lemma-triangle-differentials},
et que $\Omega_{X/S}$ et $\Omega_{Y/S}$ sont finis localement libres
de rang $n$ (Morphismes, lemme
\ref{morphisms-lemma-smooth-omega-finite-locally-free}),
l'assertion concernant la différente de Kähler résulte immédiatement de la définition
de l'idéal de Fitting d'ordre zéro. Si $f$ est étale aux points associés
de $Y$, alors $\wedge^n\text{d}f$ ne s'annule pas aux
points associés de $Y$, ce qui entraîne que l'équation locale
de $R$ est un non-diviseur de zéro. Ainsi $R$ est un diviseur de Cartier effectif.
L'isomorphisme canonique envoie $1$ sur $\wedge^n\text{d}f$, voir
Diviseurs, lemme \ref{divisors-lemma-characterize-OD}.
\end{proof}






\section{L'application de Tate}
\label{discriminant-section-tate-map}

\noindent
Dans cette section, nous construisons un isomorphisme entre
le déterminant du complexe cotangent relatif et
le module dualisant relatif pour un morphisme localement quasi-fini
syntomique de schémas localement noethériens. Suivant
\cite[1.4.4]{Garel}, nous l'appelons l'application de Tate.
Notre méthode consiste à éviter autant que possible
les calculs locaux.

\medskip\noindent
Soit $Y \to X$ un morphisme localement quasi-fini syntomique de schémas.
Nous utiliserons dans cette section, sans le mentionner à nouveau, toutes les conditions
équivalentes données dans le lemme \ref{discriminant-lemma-syntomic-quasi-finite}
pour cette notion. On voit en particulier que $\NL_{Y/X}$ est un objet parfait
de $D(\mathcal{O}_Y)$ de Tor-amplitude contenue dans $[-1, 0]$. On dispose donc
d'un module inversible canonique
$\det(\NL_{Y/X})$ sur $Y$ et d'une section globale
$$
\delta(\NL_{Y/X}) \in \Gamma(Y, \det(\NL_{Y/X}))
$$
Voir Catégories dérivées de schémas, lemme
\ref{perfect-lemma-determinant-two-term-complexes}.
Supposons donné un diagramme commutatif de schémas
$$
\xymatrix{
Y' \ar[r]_b \ar[d] & Y \ar[d] \\
X' \ar[r] & X
}
$$
dont les flèches verticales sont localement quasi-finies syntomiques et qui
induit un isomorphisme de $Y'$ avec un ouvert de $X' \times_X Y$.
L'application canonique
$$
Lb^*\NL_{Y/X} \longrightarrow \NL_{Y'/X'}
$$
est alors un quasi-isomorphisme d'après
Compléments sur les morphismes, lemme \ref{more-morphisms-lemma-base-change-NL-flat}.
On obtient donc un isomorphisme canonique
$b^*\det(\NL_{Y/X}) \to \det(\NL_{Y'/X'})$ qui envoie la
section canonique $\delta(\NL_{Y/X})$ sur $\delta(\NL_{Y'/ X'})$, voir
Catégories dérivées de schémas, remarque \ref{perfect-remark-functorial-det}.

\begin{remark}
\label{discriminant-remark-local-description-delta}
Soit $Y \to X$ un morphisme localement quasi-fini syntomique de schémas.
Quelle est la forme locale du couple $(\det(\NL_{Y/X}), \delta(\NL_{Y/X}))$ ?
Choisissons des ouverts affines $V = \Spec(B) \subset Y$,
$U = \Spec(A) \subset X$ avec $f(V) \subset U$, un entier $n$ et
$f_1, \ldots, f_n \in A[x_1, \ldots, x_n]$ tels que
$B = A[x_1, \ldots, x_n]/(f_1, \ldots, f_n)$. Alors
$$
\NL_{B/A} = \left(
(f_1, \ldots, f_n)/(f_1, \ldots, f_n)^2
\longrightarrow
\bigoplus\nolimits_{i = 1, \ldots, n} B \text{d} x_i\right)
$$
et $(f_1, \ldots, f_n)/(f_1, \ldots, f_n)^2$ est libre de base
formée des classes $\overline{f}_i$. Voir la démonstration du
lemme \ref{discriminant-lemma-syntomic-quasi-finite}.
Ainsi $\det(L_{B/A})$ est libre sur le générateur
$$
\text{d}x_1 \wedge \ldots \wedge \text{d}x_n
\otimes
(\overline{f}_1 \wedge \ldots \wedge \overline{f}_n)^{\otimes -1}
$$
et la section $\delta(\NL_{B/A})$ est l'élément
$$
\delta(\NL_{B/A}) =
\det(\partial f_j/ \partial x_i) \cdot
\text{d}x_1 \wedge \ldots \wedge \text{d}x_n
\otimes
(\overline{f}_1 \wedge \ldots \wedge \overline{f}_n)^{\otimes -1}
$$
par définition.
\end{remark}

\noindent
Soit $Y \to X$ un morphisme localement quasi-fini syntomique de
schémas localement noethériens. D'après les
remarques \ref{discriminant-remark-relative-dualizing-for-quasi-finite} et
\ref{discriminant-remark-relative-dualizing-for-flat-quasi-finite}, on dispose
d'un $\mathcal{O}_Y$-module cohérent $\omega_{Y/X}$ et d'une section globale
canonique
$$
\tau_{Y/X} \in \Gamma(Y, \omega_{Y/X})
$$
qui redonne, localement sur les ouverts affines, le couple $\omega_{B/A}, \tau_{B/A}$.
D'après le lemme \ref{discriminant-lemma-dualizing-syntomic-quasi-finite}, le module
$\omega_{Y/X}$ est inversible. Supposons donné un diagramme commutatif de
schémas localement noethériens
$$
\xymatrix{
Y' \ar[r]_b \ar[d] & Y \ar[d] \\
X' \ar[r] & X
}
$$
dont les flèches verticales sont localement quasi-finies syntomiques et qui
induit un isomorphisme de $Y'$ avec un ouvert de $X' \times_X Y$.
Il existe alors une application canonique de changement de base
$$
b^*\omega_{Y/X} \longrightarrow \omega_{Y'/X'}
$$
qui est un isomorphisme
envoyant $\tau_{Y/X}$ sur $\tau_{Y'/X'}$. En effet, l'application de changement de base
dans le cadre affine est (\ref{discriminant-equation-bc-dualizing}) ; c'est un
isomorphisme d'après le lemme \ref{discriminant-lemma-dualizing-base-change-of-flat}, et elle
envoie $\tau_{Y/X}$ sur $\tau_{Y'/X'}$ d'après
l'assertion (1) du lemme \ref{discriminant-lemma-trace-base-change}.

\begin{proposition}
\label{discriminant-proposition-tate-map}
Il existe une unique règle qui associe à tout morphisme localement quasi-fini syntomique
de schémas localement noethériens $Y \to X$ un isomorphisme
$$
c_{Y/X} : \det(\NL_{Y/X}) \longrightarrow \omega_{Y/X}
$$
vérifiant les deux propriétés suivantes :
\begin{enumerate}
\item la section $\delta(\NL_{Y/X})$ s'envoie sur $\tau_{Y/X}$ ;
\item la règle est compatible avec la restriction aux ouverts et le
changement de base.
\end{enumerate}
\end{proposition}

\begin{proof}
Reformulons l'énoncé de la proposition. Considérons la catégorie
$\mathcal{C}$ dont les objets, notés $Y/X$, sont les morphismes localement quasi-finis syntomiques
$Y \to X$ de schémas localement noethériens, et dont les morphismes
$b/a : Y'/X' \to Y/X$ sont les diagrammes commutatifs
$$
\xymatrix{
Y' \ar[d] \ar[r]_b & Y \ar[d] \\
X' \ar[r]^a & X
}
$$
induisant un isomorphisme de $Y'$ avec un sous-schéma ouvert de
$X' \times_X Y$. La proposition signifie que, pour tout objet
$Y/X$ de $\mathcal{C}$, on dispose d'un isomorphisme
$c_{Y/X} : \det(\NL_{Y/X}) \to \omega_{Y/X}$
tel que $c_{Y/X}(\delta(\NL_{Y/X})) = \tau_{Y/X}$,
et que, pour tout morphisme $b/a : Y'/X' \to Y/X$ de $\mathcal{C}$, on a
$b^*c_{Y/X} = c_{Y'/X'}$ via les identifications
$b^*\det(\NL_{Y/X}) = \det(\NL_{Y'/X'})$ et
$b^*\omega_{Y/X} = \omega_{Y'/X'}$ décrites ci-dessus.

\medskip\noindent
Étant donné $Y/X$ dans $\mathcal{C}$ et $y \in Y$, on peut trouver
des ouverts affines $V \subset Y$ et $U \subset X$ avec
$y \in V$ et $f(V) \subset U$
tels qu'il existe un isomorphisme
$$
\det(\NL_{Y/X})|_V \longrightarrow \omega_{Y/X}|_V
$$
envoyant $\delta(\NL_{Y/X})|_V$ sur $\tau_{Y/X}|_V$. Cela résulte
du choix d'ouverts affines comme dans
l'assertion (5) du lemme \ref{discriminant-lemma-syntomic-quasi-finite}, de la description
locale affine de $\delta(\NL_{Y/X})$ dans la
remarque \ref{discriminant-remark-local-description-delta}, et du
lemme \ref{discriminant-lemma-different-quasi-finite-complete-intersection}.
Si l'annulateur de la section $\tau_{Y/X}$ est nul, ces
applications locales sont uniques et se recollent automatiquement. Ainsi, si l'annulateur
de $\tau_{Y/X}$ est nul, il existe un unique isomorphisme
$c_{Y/X} : \det(\NL_{Y/X}) \to \omega_{Y/X}$ tel que
$c_{Y/X}(\delta(\NL_{Y/X})) = \tau_{Y/X}$.
Si $b/a : Y'/X' \to Y/X$ est un morphisme de $\mathcal{C}$
et si l'annulateur de $\tau_{Y'/X'}$ est lui aussi nul,
alors $b^*c_{Y/X}$ est l'unique isomorphisme
$c_{Y'/X'} : \det(\NL_{Y'/X'}) \to \omega_{Y'/X'}$ tel que
$c_{Y'/X'}(\delta(\NL_{Y'/X'})) = \tau_{Y'/X'}$.
Cela découle formellement du fait que
$b^*\delta(\NL_{Y/X}) = \delta(\NL_{Y'/X'})$ et
$b^*\tau_{Y/X} = \tau_{Y'/X'}$.

\medskip\noindent
On peut résumer les résultats du paragraphe précédent comme suit.
Notons $\mathcal{C}_{nice} \subset \mathcal{C}$ la
sous-catégorie pleine des $Y/X$ tels que l'annulateur de
$\tau_{Y/X}$ soit nul. Nous avons alors résolu le problème
sur $\mathcal{C}_{nice}$. Pour $Y/X$ dans $\mathcal{C}_{nice}$,
nous continuons à noter $c_{Y/X}$ la solution que nous venons de trouver.

\medskip\noindent
Considérons des morphismes
$$
Y_1/X_1 \xleftarrow{b_1/a_1} Y/X \xrightarrow{b_2/a_2} Y_2/X_2
$$
dans $\mathcal{C}$ tels que $Y_1/X_1$ et $Y_2/X_2$ soient des objets
de $\mathcal{C}_{nice}$. {\bf Assertion.} $b_1^*c_{Y_1/X_1} = b_2^*c_{Y_2/X_2}$.
Nous montrerons d'abord que cette assertion entraîne la proposition,
puis nous la démontrerons.

\medskip\noindent
Soient $d, n \geq 1$ ; considérons le morphisme localement
quasi-fini syntomique $Y_{n, d} \to X_{n, d}$
construit dans l'exemple \ref{discriminant-example-universal-quasi-finite-syntomic}.
Alors $Y_{n, d}$ est un schéma irréductible régulier et le
morphisme $Y_{n, d} \to X_{n, d}$ est localement quasi-fini syntomique
et étale au-dessus d'un ouvert dense, voir le
lemme \ref{discriminant-lemma-universal-quasi-finite-syntomic-etale}.
Ainsi $\tau_{Y_{n, d}/X_{n, d}}$ est non nul, par exemple d'après le
lemme \ref{discriminant-lemma-different-ramification}. Or une section non nulle
d'un module inversible sur un schéma irréductible régulier
a un annulateur nul. Ainsi
$Y_{n, d}/X_{n, d}$ est un objet de $\mathcal{C}_{nice}$.

\medskip\noindent
Soit $Y/X$ un objet quelconque de $\mathcal{C}$. Soit $y \in Y$.
D'après le lemme \ref{discriminant-lemma-locally-comes-from-universal}, on peut trouver
$n, d \geq 1$ et des morphismes
$$
Y/X \leftarrow V/U \xrightarrow{b/a} Y_{n, d}/X_{n, d}
$$
de $\mathcal{C}$ tels que $V \subset Y$ et $U \subset X$ soient ouverts.
On peut donc prendre l'image inverse du morphisme canonique $c_{Y_{n, d}/X_{n, d}}$
construit ci-dessus par $b$ sur $V$. L'assertion garantit que ces isomorphismes
locaux se recollent ! On obtient ainsi un isomorphisme global bien défini
$c_{Y/X} : \det(\NL_{Y/X}) \to \omega_{Y/X}$ tel que
$c_{Y/X}(\delta(\NL_{Y/X})) = \tau_{Y/X}$.
Si $b/a : Y'/X' \to Y/X$ est un morphisme de $\mathcal{C}$,
l'assertion entraîne également que l'application construite de façon analogue
$c_{Y'/X'}$ est l'image inverse par $b$ de l'application construite localement
$c_{Y/X}$. Il reste donc à démontrer l'assertion.

\medskip\noindent
C'est ce que nous faisons dans la suite de la démonstration. On peut choisir un point
$y \in Y$ et montrer que les applications coïncident dans un voisinage ouvert de $y$.
On peut donc remplacer $Y_1$, $Y_2$ par des voisinages ouverts de
l'image de $y$ dans $Y_1$ et $Y_2$. On peut ainsi supposer que
$Y, X, Y_1, X_1, Y_2, X_2$ sont affines, disons les spectres
d'anneaux $B, A, B_1, A_1, B_2, A_2$. On a le diagramme
$$
\xymatrix{
B_1 \ar[r] & B & B_2 \ar[l] \\
A_1 \ar[u] \ar[r] & A \ar[u] & A_2 \ar[l] \ar[u]
}
$$
Par hypothèse, le spectre de $B$ est un ouvert affine
à la fois du spectre de $A \otimes_{A_1} B_1$ et de celui de $A \otimes_{A_2} B_2$.
En rétrécissant encore, on peut supposer qu'il existe des éléments
$g_i \in A \otimes_{A_i} B_i$ tels que nos applications donnent des isomorphismes
$(A \otimes_{A_i} B_i)_{g_i} = B$, voir
Propriétés, lemme \ref{properties-lemma-standard-open-two-affines}.
Soient $x_\alpha$, $y_\beta$ des collections assez grandes de
variables pour que l'on puisse choisir des surjections
$$
A'_1 = A_1[x_\alpha] \to A
\quad\text{et}\quad
A'_2 = A_2[y_\beta] \to A
$$
de $A_1$-algèbres et de $A_2$-algèbres. On peut alors choisir des relèvements
$h_i \in A'_i \otimes_{A_i} B_i$ de $g_i \in A \otimes_{A_i} B_i$,
et l'on considère le diagramme
$$
\xymatrix{
(A'_1 \otimes_{A_1} B_1)_{h_1} \ar[r] & B &
(A'_2 \otimes_{A_1} B_2)_{h_2} \ar[l] \\
A'_1 \ar[u] \ar[r] & A \ar[u] & A'_2 \ar[l] \ar[u]
}
$$
Par construction, les deux carrés sont cocartésiens. Considérons ensuite
l'homomorphisme d'anneaux
$$
A' = A'_1 \times_A A'_2 \longrightarrow
B' =
(A'_1 \otimes_{A_1} B_1)_{h_1} \times_B
(A'_2 \otimes_{A_1} B_2)_{h_2}
$$
D'après
Compléments d'algèbre, lemme \ref{more-algebra-lemma-module-over-fibre-product},
on a $A'_1 \otimes_{A'} B' = (A'_1 \otimes_{A_1} B_1)_{h_1}$
et $A'_2 \otimes_{A'} B' = (A'_2 \otimes_{A_2} B_2)_{h_2}$. En particulier,
les fibres du morphisme $Y' = \Spec(B') \to \Spec(A') = X'$ sont
des sous-schémas ouverts de changements de base des fibres des morphismes $Y_i \to X_i$,
donc finies et localement des intersections complètes
(voir le lemme \ref{discriminant-lemma-syntomic-quasi-finite}).
D'après Compléments d'algèbre, lemme
\ref{more-algebra-lemma-properties-algebras-over-fibre-product},
l'homomorphisme d'anneaux $A' \to B'$ est plat et de présentation finie.
D'après l'assertion (3) du lemme \ref{discriminant-lemma-syntomic-quasi-finite},
on conclut donc que $Y' \to X'$ est localement quasi-fini et syntomique.
On obtient ainsi le diagramme commutatif
$$
\xymatrix{
& & Y/X \ar[ld] \ar@/_2pc/[lldd]_{b_1/a_1} \ar[rd] \ar@/^2pc/[rrdd]^{b_2/a_2} \\
& Y'_1/X'_1 \ar[rd] \ar[ld] & & Y'_2/X'_2 \ar[ld] \ar[rd] \\
Y_1/X_1 & & Y'/X' & & Y_2/X_2
}
$$
où $Y'_i/X'_i$ correspond à $A'_i \to (A'_i \otimes_{A_i} B_i)_{h_i}$.

\medskip\noindent
Dans ce paragraphe, un passage à la limite remédie au fait que $A'$ et $A'_i$ ne sont pas
nécessairement noethériens ; le lecteur peut omettre ce paragraphe.
On peut trouver une $\mathbf{Z}$-sous-algèbre de type fini
$A'' \subset A'$ et un homomorphisme quasi-fini syntomique d'anneaux
$A'' \to B''$ tels que $B' = A' \otimes_{A''} B''$.
Cela résulte de Limites, lemmes
\ref{limits-lemma-descend-finite-presentation},
\ref{limits-lemma-descend-syntomic} et
\ref{limits-lemma-descend-quasi-finite}.
Les homomorphismes d'anneaux $A'' \to A'_1 = A_1[x_\alpha]$ se factorisent alors par une
sous-algèbre de polynômes à un nombre fini de variables $A_i[x_1, \ldots, x_n$, et l'homomorphisme
$A'' \to A'_2 = A_2[y_\beta]$ se factorise par $A_2[y_1, \ldots, y_m]$.
En agrandissant les ensembles de variables, on peut aussi supposer que
$h_1 \in A_1[x_1, \ldots, x_n] \otimes_{A_1} B_1$ et
$h_2 \in A_2[y_1, \ldots, y_m] \otimes_{A_2} B_2$.
En agrandissant encore les ensembles de variables, on peut supposer
que les applications $B'' \to (A'_i \otimes_{A_i} B_i)_{h_i}$
se factorisent par $(A_1[x_1, \ldots, x_n] \otimes_{A_1} B_1)_{h_1}$
et $(A_2[y_1, \ldots, y_m] \otimes_{A_2} B_2)_{h_2}$.
En remplaçant $Y'$, $X'$, $Y'_1$, $X'_1$, $Y'_2$, $X'_2$ par les spectres de
$B''$, $A''$, $(A_1[x_1, \ldots, x_n] \otimes_{A_1} B_1)_{h_1}$,
$A_1[x_1, \ldots, x_n]$, $(A_2[y_1, \ldots, y_m] \otimes_{A_2} B_2)_{h_2}$,
$A_2[y_1, \ldots, y_m]$, on obtient un diagramme comme ci-dessus
où tous les schémas sont noethériens.

\medskip\noindent
Remarquons que $Y'_i/X'_i$ est un objet de $\mathcal{C}_{nice}$, car il
s'obtient en prenant un sous-schéma ouvert du produit d'un espace affine
avec $Y_i/X_i$. En particulier, l'image inverse de
$c_{Y_i/X_i}$ par $(b_i/a_i)$ est la même que l'image inverse de
$c_{Y'_i/X'_i}$ par $Y/X \to Y'_i/X'_i$.
La démonstration serait alors achevée si $Y'/X'$ était un objet de $\mathcal{C}_{nice}$
(mais ce n'est probablement pas le cas). En effet, en prenant l'image inverse de $c_{Y'/X'}$
le long des deux côtés du carré, on obtiendrait la conclusion voulue.
Pour contourner le fait que $Y'/X'$ n'appartient pas à $\mathcal{C}_{nice}$,
remarquons que les arguments ci-dessus montrent qu'après avoir éventuellement rétréci tous
les schémas $X, Y, X'_1, Y'_1, X'_2, Y'_2, X', Y'$, on peut trouver
$n, d \geq 1$ et prolonger le diagramme comme suit :
$$
\xymatrix{
& Y/X \ar[ld] \ar[rd] \\
Y'_1/X'_1 \ar[rd] & & Y'_2/X'_2 \ar[ld] \\
& Y'/X' \ar[d] \\
& Y_{n, d}/X_{n, d}
}
$$
On peut alors utiliser l'argument déjà donné en prenant
l'image inverse de $c_{Y_{n, d}/X_{n, d}}$. La démonstration est achevée.
\end{proof}













\section{Une généralisation de la différente}
\label{discriminant-section-different-generalization}

\noindent
Dans cette section, nous généralisons la définition \ref{discriminant-definition-different}
pour englober tous les homomorphismes d'anneaux $A \to B$
pour lesquels la différente de Dedekind est définie et $1 \in \mathcal{L}_{B/A}$.
Expliquons d'abord la condition « $A \to B$ envoie les non-diviseurs de zéro sur
des non-diviseurs de zéro et induit un homomorphisme plat $Q(A) \to Q(A) \otimes_A B$ ».

\begin{lemma}
\label{discriminant-lemma-explain-condition}
Soit $A \to B$ un homomorphisme d'anneaux noethériens. Considérons les conditions
\begin{enumerate}
\item les non-diviseurs de zéro de $A$ s'envoient sur des non-diviseurs de zéro de $B$ ;
\item (1) est satisfaite et $Q(A) \to Q(A) \otimes_A B$ est plat ;
\item $A \to B_\mathfrak q$ est plat pour tout
$\mathfrak q \in \text{Ass}(B)$ ;
\item (3) est satisfaite et $A \to B_\mathfrak q$ est plat pour tout $\mathfrak q$
au-dessus d'un élément de $\text{Ass}(A)$.
\end{enumerate}
On a alors les implications suivantes :
$$
\xymatrix{
(1) & (2) \ar@{=>}[l] \ar@{=>}[d] \\
(3) \ar@{=>}[u] & (4) \ar@{=>}[l]
}
$$
Si $A \to B$ vérifie la propriété de montée, (2) et (4) sont équivalentes.
\end{lemma}

\begin{proof}
Les implications horizontales du diagramme sont triviales.
Soit $S \subset A$ l'ensemble des non-diviseurs de zéro, de sorte que
$Q(A) = S^{-1}A$ et $Q(A) \otimes_A B = S^{-1}B$. Rappelons que
$S = A \setminus \bigcup_{\mathfrak p \in \text{Ass}(A)} \mathfrak p$
d'après Algèbre, lemme \ref{algebra-lemma-ass-zero-divisors}.
Soit $\mathfrak q \subset B$ un idéal premier au-dessus de $\mathfrak p \subset A$.

\medskip\noindent
Supposons (2). Si $\mathfrak q \in \text{Ass}(B)$, alors
$\mathfrak q$ est constitué de diviseurs de zéro ; (1) entraîne donc
qu'il en est de même de $\mathfrak p$. Ainsi
$\mathfrak p$ correspond à un idéal premier de $S^{-1}A$.
Par conséquent, $A \to B_\mathfrak q$ est plat d'après l'hypothèse (2).
Si $\mathfrak q$ est au-dessus d'un idéal premier associé $\mathfrak p$
de $A$, alors on a certainement $\mathfrak p \in \Spec(S^{-1}A)$, et le
même argument s'applique.

\medskip\noindent
Supposons (3). Soit $f \in A$ un non-diviseur de zéro. Si $f$ était diviseur de zéro
sur $B$, alors $f$ appartiendrait à un idéal premier associé $\mathfrak q$
de $B$. Puisque $A \to B_\mathfrak q$ est plat par hypothèse, on en déduit que
$\mathfrak p$ est un idéal premier associé de $A$, d'après
Algèbre, lemme \ref{algebra-lemma-bourbaki}. Cela entraînerait que
$f$ est diviseur de zéro sur $A$, contradiction.

\medskip\noindent
Supposons (4) et la propriété de montée pour $A \to B$. Nous savons déjà que (1) est satisfaite.
Si $\mathfrak q$ correspond à un idéal premier de $S^{-1}B$, alors $\mathfrak p$
est contenu dans un idéal premier associé $\mathfrak p'$ de $A$. Par la propriété de montée,
il existe un idéal premier $\mathfrak q'$ contenant $\mathfrak q$ et situé
au-dessus de $\mathfrak p$. Alors $A \to B_{\mathfrak q'}$ est plat d'après
(4). Donc $A \to B_{\mathfrak q}$ est plat, en tant que localisation.
Ainsi $A \to S^{-1}B$ est plat, de même que $S^{-1}A \to S^{-1}B$, voir
Algèbre, lemme \ref{algebra-lemma-flat-localization}.
\end{proof}

\begin{remark}
\label{discriminant-remark-different-generalization}
On peut généraliser la définition \ref{discriminant-definition-different}.
Supposons que $f : Y \to X$ soit un morphisme quasi-fini de schémas noethériens
possédant les propriétés suivantes :
\begin{enumerate}
\item l'ouvert $V \subset Y$ où $f$ est plat contient
$\text{Ass}(\mathcal{O}_Y)$ et $f^{-1}(\text{Ass}(\mathcal{O}_X))$ ;
\item l'élément trace $\tau_{V/X}$ provient d'une section
$\tau \in \Gamma(Y, \omega_{Y/X})$.
\end{enumerate}
La condition (1) entraîne que $V$ contient les points associés de
$\omega_{Y/X}$, d'après le lemme \ref{discriminant-lemma-dualizing-associated-primes}.
En particulier, $\tau$ est unique s'il existe
(Diviseurs, lemme \ref{divisors-lemma-restriction-injective-open-contains-ass}).
Étant donné $\tau$, on peut définir la différente $\mathfrak{D}_f$ comme l'annulateur de
$\Coker(\tau : \mathcal{O}_Y \to \omega_{Y/X})$. Elle coïncide avec la
différente de Dedekind dans de nombreux cas (lemme \ref{discriminant-lemma-agree-dedekind}).
Cependant, pour les homomorphismes non plats entre anneaux non normaux, cette généralisation
ne mesure plus la ramification du morphisme, voir
l'exemple \ref{discriminant-example-no-different}.
\end{remark}

\begin{lemma}
\label{discriminant-lemma-agree-dedekind}
Supposons que la différente de Dedekind de $A \to B$ soit définie.
Posons $X = \Spec(A)$ et $Y = \Spec(B)$. La généralisation de la
remarque \ref{discriminant-remark-different-generalization}
s'applique au morphisme $f : Y \to X$ si et seulement si
$1 \in \mathcal{L}_{B/A}$ (par exemple si $A$ est normal, voir le
lemme \ref{discriminant-lemma-dedekind-different-ideal}).
Dans ce cas, $\mathfrak{D}_{B/A}$ est un idéal de $B$ et l'on a
$$
\mathfrak{D}_f = \widetilde{\mathfrak{D}_{B/A}}
$$
comme faisceaux cohérents d'idéaux sur $Y$.
\end{lemma}

\begin{proof}
La différente de Dedekind de $A \to B$ étant définie, on peut appliquer le
lemme \ref{discriminant-lemma-explain-condition} pour voir que
$Y \to X$ vérifie la condition (1) de la
remarque \ref{discriminant-remark-different-generalization}.
Rappelons qu'il existe un isomorphisme canonique
$c : \mathcal{L}_{B/A} \to \omega_{B/A}$, voir le
lemme \ref{discriminant-lemma-dedekind-complementary-module}.
Soient $K = Q(A)$ et $L = K \otimes_A B$ comme ci-dessus.
Par construction, l'application $c$ s'insère dans un diagramme commutatif
$$
\xymatrix{
\mathcal{L}_{B/A} \ar[r] \ar[d]_c & L \ar[d] \\
\omega_{B/A} \ar[r] & \Hom_K(L, K)
}
$$
où la flèche verticale de droite envoie $x \in L$ sur l'application
$y \mapsto \text{Trace}_{L/K}(xy)$, et la flèche horizontale
inférieure est l'application de changement de base (\ref{discriminant-equation-bc-dualizing}) pour $\omega_{B/A}$.
On peut factoriser l'application horizontale inférieure en
$$
\omega_{B/A} = \Gamma(Y, \omega_{Y/X})
\to \Gamma(V, \omega_{V/X}) \to \Hom_K(L, K)
$$
Puisque tous les points associés de $\omega_{V/X}$
s'envoient sur des idéaux premiers associés de $A$
(lemme \ref{discriminant-lemma-dualizing-associated-primes}),
la seconde application est injective.
L'élément $\tau_{V/X}$ s'envoie sur $\text{Trace}_{L/K}$ dans
$\Hom_K(L, K)$ par la définition même des éléments trace
(définition \ref{discriminant-definition-trace-element}).
Ainsi un $\tau$ comme dans la condition (2) de la
remarque \ref{discriminant-remark-different-generalization}
existe si et seulement si $1 \in \mathcal{L}_{B/A}$, et l'on a alors
$\tau = c(1)$. Dans ce cas, d'après le lemme \ref{discriminant-lemma-dedekind-different-ideal},
on voit que $\mathfrak{D}_{B/A} \subset B$.
Enfin, l'égalité de $\mathfrak{D}_f$ et de $\mathfrak{D}_{B/A}$
découle immédiatement des définitions et du fait que $\tau = c(1)$, établi ci-dessus.
\end{proof}

\begin{example}
\label{discriminant-example-no-different}
Soit $k$ un corps. Soient $A = k[x, y]/(xy)$ et $B = k[u, v]/(uv)$, et soit
$A \to B$ donné par $x \mapsto u^n$ et $y \mapsto v^m$, pour des
$n, m \in \mathbf{N}$ premiers à la caractéristique de $k$. Alors
$A_{x + y} \to B_{x + y}$ est (fini) étale, et nous sommes donc dans une situation
où la différente de Dedekind est définie. Un calcul donne
$$
\text{Trace}_{L/K}(1) = (nx + my)/(x + y),\quad
\text{Trace}_{L/K}(u^i) = 0,\quad \text{Trace}_{L/K}(v^j) = 0
$$
pour $1 \leq i < n$ et $1 \leq j < m$. On en déduit que
$1 \in \mathcal{L}_{B/A}$ si et seulement si $n = m$. De plus, un
calcul montre que, si $n = m$, alors $\mathcal{L}_{B/A} = B$,
et la différente de Dedekind est également $B$. Autrement dit,
la différente de la remarque \ref{discriminant-remark-different-generalization}
est définie pour $\Spec(B) \to \Spec(A)$
si et seulement si $n = m$, et, dans ce cas, elle est
l'idéal unité. On voit donc que, dans les cas non plats, la non-annulation
de la différente ne garantit pas que le morphisme soit étale ou non ramifié.
\end{example}







\section{Comparaison avec la théorie de la dualité}
\label{discriminant-section-comparison}

\noindent
Dans cette section, nous comparons les constructions algébriques élémentaires
ci-dessus aux constructions du chapitre sur la théorie de la dualité
des schémas.

\begin{lemma}
\label{discriminant-lemma-compare-dualizing}
Soit $f : Y \to X$ un morphisme quasi-fini séparé de schémas noethériens.
Pour tout couple d'ouverts affines $\Spec(B) = V \subset Y$,
$\Spec(A) = U \subset X$ avec $f(V) \subset U$, il existe un isomorphisme
$$
H^0(V, f^!\mathcal{O}_X) = \omega_{B/A}
$$
où $f^!$ est comme dans
Dualité des schémas, section \ref{duality-section-upper-shriek}.
Ces isomorphismes sont compatibles aux applications de restriction et définissent un
isomorphisme canonique $H^0(f^!\mathcal{O}_X) = \omega_{Y/X}$, avec
$\omega_{Y/X}$ comme dans la remarque \ref{discriminant-remark-relative-dualizing-for-quasi-finite}.
De même, si $f : Y \to X$ est un morphisme quasi-fini de schémas de
type fini sur une base noethérienne $S$ munie d'un complexe dualisant
$\omega_S^\bullet$, alors $H^0(f_{new}^!\mathcal{O}_X) = \omega_{Y/X}$.
\end{lemma}

\begin{proof}
Le théorème principal de Zariski permet de choisir une factorisation $f = f' \circ j$,
où $j : Y \to Y'$ est une immersion ouverte et $f' : Y' \to X$ un morphisme
fini, voir Compléments sur les morphismes, lemme
\ref{more-morphisms-lemma-quasi-finite-separated-pass-through-finite}.
Par la construction de
Dualité des schémas, lemme \ref{duality-lemma-shriek-well-defined}, on a
$f^! = j^* \circ a'$, où
$a' : D_\QCoh(\mathcal{O}_X) \to D_\QCoh(\mathcal{O}_{Y'})$
est l'adjoint à droite de $Rf'_*$ donné par
Dualité des schémas, lemme \ref{duality-lemma-twisted-inverse-image}.
D'après Dualité des schémas, lemme \ref{duality-lemma-finite-twisted},
on voit que
$\Phi(a'(\mathcal{O}_X)) = R\SheafHom(f'_*\mathcal{O}_{Y'}, \mathcal{O}_X)$ dans
$D_\QCoh^+(f'_*\mathcal{O}_{Y'})$. En particulier, $a'(\mathcal{O}_X)$ a des
faisceaux de cohomologie nuls en degrés $< 0$. Le faisceau de cohomologie de degré zéro
est déterminé par l'isomorphisme
$$
f'_*H^0(a'(\mathcal{O}_X)) =
\SheafHom_{\mathcal{O}_X}(f'_*\mathcal{O}_{Y'}, \mathcal{O}_X)
$$
de $f'_*\mathcal{O}_{Y'}$-modules via l'équivalence de
Morphismes, lemme \ref{morphisms-lemma-affine-equivalence-modules}.
En écrivant $(f')^{-1}U = V' = \Spec(B')$, on obtient
$$
H^0(V', a'(\mathcal{O}_X)) = \Hom_A(B', A).
$$
Puisque le faisceau de cohomologie de degré zéro de $a'(\mathcal{O}_X)$
est un module quasi-cohérent, on trouve que
sa restriction à $V$ est donnée par
$\omega_{B/A} = \Hom_A(B', A) \otimes_{B'} B$, comme annoncé.

\medskip\noindent
L'assertion sur les applications de restriction signifie que celles
du $\mathcal{O}_{Y'}$-module quasi-cohérent $H^0(a'(\mathcal{O}_X))$
pour les ouverts de $Y'$ coïncident avec les applications définies dans le
lemme \ref{discriminant-lemma-localize-dualizing}
pour les modules $\omega_{B/A}$ via les isomorphismes ci-dessus.
C'est clair.

\medskip\noindent
Soit $f : Y \to X$ un morphisme quasi-fini de schémas de type fini
sur une base noethérienne $S$ munie d'un complexe dualisant $\omega_S^\bullet$.
Considérons des ouverts $V \subset Y$ et $U \subset X$ avec $f(V) \subset U$,
et $V$ et $U$ séparés sur $S$. Notons $f|_V : V \to U$ la restriction
de $f$. D'après la discussion ci-dessus et
Dualité des schémas, lemme \ref{duality-lemma-duality-bootstrap},
on a des isomorphismes canoniques
$$
H^0(f_{new}^!\mathcal{O}_X)|_V = H^0((f|_V)^!\mathcal{O}_U) = \omega_{V/U} =
\omega_{Y/X}|_V
$$
Nous omettons la vérification que ces isomorphismes se recollent en un isomorphisme
global $H^0(f_{new}^!\mathcal{O}_X) \to \omega_{Y/X}$.
\end{proof}

\begin{lemma}
\label{discriminant-lemma-compare-trace}
Soit $f : Y \to X$ un morphisme fini plat de schémas noethériens.
L'application
$$
\text{Trace}_f : f_*\mathcal{O}_Y \longrightarrow \mathcal{O}_X
$$
de la section \ref{discriminant-section-discriminant}
correspond à une application $\mathcal{O}_Y \to f^!\mathcal{O}_X$ (voir la démonstration).
Notons $\tau_{Y/X} \in H^0(Y, f^!\mathcal{O}_X)$ l'image de $1$.
Via l'isomorphisme $H^0(f^!\mathcal{O}_X) = \omega_{X/Y}$ du
lemme \ref{discriminant-lemma-compare-dualizing},
cela coïncide avec la construction de la
remarque \ref{discriminant-remark-relative-dualizing-for-flat-quasi-finite}.
\end{lemma}

\begin{proof}
Le foncteur $f^!$ est défini dans
Dualité des schémas, section \ref{duality-section-upper-shriek}.
Puisque $f$ est fini (donc propre), $f^!$ est donné par
l'adjoint à droite de l'image directe par $f$. Dans
Dualité des schémas, section \ref{duality-section-duality-finite},
nous avons explicité cet adjoint. En particulier,
l'objet $f^!\mathcal{O}_X$ consiste en un seul
faisceau de cohomologie placé en degré $0$, et l'on a pour ce faisceau
$$
f_*f^!\mathcal{O}_X =
\SheafHom_{\mathcal{O}_X}(f_*\mathcal{O}_Y, \mathcal{O}_X)
$$
Pour le voir, on utilise aussi le fait que $f_*\mathcal{O}_Y$ est fini localement libre,
car $f$ est un morphisme fini plat de schémas noethériens,
et donc que tous les faisceaux Ext supérieurs sont nuls. Certains détails sont omis.
Finalement,
$$
\text{Trace}_f \in
\Hom_{\mathcal{O}_X}(f_*\mathcal{O}_Y, \mathcal{O}_X) =
\Gamma(X, f_*f^!\mathcal{O}_X) =
\Gamma(Y, f^!\mathcal{O}_X)
$$
D'autre part, on a $f^!\mathcal{O}_X = \omega_{Y/X}$
par l'identification du lemme \ref{discriminant-lemma-compare-dualizing}.
On dispose donc maintenant de deux éléments, à savoir $\text{Trace}_f$
et $\tau_{Y/X}$ de la
remarque \ref{discriminant-remark-relative-dualizing-for-flat-quasi-finite}, dans
$$
\Gamma(Y, f^!\mathcal{O}_X) = \Gamma(Y, \omega_{Y/X})
$$
et le lemme affirme qu'ils sont égaux.

\medskip\noindent
Soit $U = \Spec(A) \subset X$ un ouvert affine d'image inverse
$V = \Spec(B) \subset Y$. Puisque $f$ est fini,
$A \to B$ est fini et l'on a donc $\omega_{Y/X}(V) = \Hom_A(B,A)$
par construction ; cet isomorphisme coïncide avec l'identification
de $f_*f^!\mathcal{O}_Y$ avec
$\SheafHom_{\mathcal{O}_X}(f_*\mathcal{O}_Y, \mathcal{O}_X)$ décrite
ci-dessus. L'égalité de $\text{Trace}_f$ et de $\tau_{Y/X}$
découle donc du fait que $\tau_{B/A} = \text{Trace}_{B/A}$
d'après le lemme \ref{discriminant-lemma-finite-flat-trace}.
\end{proof}








\section{Morphismes quasi-finis de Gorenstein}
\label{discriminant-section-gorenstein-lci}

\noindent
Cette section traite des morphismes quasi-finis de Gorenstein.

\begin{lemma}
\label{discriminant-lemma-gorenstein-quasi-finite}
Soit $f : Y \to X$ un morphisme quasi-fini de schémas noethériens.
Les conditions suivantes sont équivalentes :
\begin{enumerate}
\item $f$ est de Gorenstein ;
\item $f$ est plat et les fibres de $f$ sont de Gorenstein ;
\item $f$ est plat et $\omega_{Y/X}$ est inversible
(remarque \ref{discriminant-remark-relative-dualizing-for-quasi-finite}) ;
\item pour tout $y \in Y$, il existe des ouverts affines
$y \in V = \Spec(B) \subset Y$, $U = \Spec(A) \subset X$
avec $f(V) \subset U$ tels que $A \to B$ soit plat
et $\omega_{B/A}$ soit un $B$-module inversible.
\end{enumerate}
\end{lemma}

\begin{proof}
Les conditions (1) et (2) sont équivalentes par définition. Les conditions (3) et (4)
sont équivalentes par la construction de $\omega_{Y/X}$ dans la
remarque \ref{discriminant-remark-relative-dualizing-for-quasi-finite}.
Il faut donc montrer que (1)-(2) équivaut à (3)-(4).

\medskip\noindent
Première démonstration. En travaillant localement sur les ouverts affines, on peut supposer $f$ séparé
et appliquer le lemme \ref{discriminant-lemma-compare-dualizing} pour voir que
$\omega_{Y/X}$ est le faisceau de cohomologie de degré zéro de $f^!\mathcal{O}_X$.
Sous chacune des deux hypothèses, $f$ est plat et quasi-fini ;
$f^!\mathcal{O}_X$ est donc isomorphe à $\omega_{Y/X}[0]$, voir
Dualité des schémas, lemme \ref{duality-lemma-flat-quasi-finite-shriek}. Ainsi
l'équivalence découle de
Dualité des schémas, lemme
\ref{duality-lemma-affine-flat-Noetherian-gorenstein}.

\medskip\noindent
Deuxième démonstration. D'après le lemme \ref{discriminant-lemma-characterize-invertible},
il suffit d'établir l'équivalence de
(2) et (3) lorsque $X$ est le spectre d'un corps $k$.
Alors $Y = \Spec(B)$, où $B$ est une $k$-algèbre finie.
Dans ce cas, $\omega_{B/A} = \omega_{B/k} = \Hom_k(B, k)$,
placé en degré $0$, est un complexe dualisant de $B$, voir
Complexes dualisants, lemme \ref{dualizing-lemma-dualizing-finite}.
L'équivalence découle donc de
Complexes dualisants, lemme \ref{dualizing-lemma-gorenstein}.
\end{proof}

\begin{remark}
\label{discriminant-remark-collect-results-qf-gorenstein}
Soit $f : Y \to X$ un morphisme quasi-fini de Gorenstein de schémas noethériens.
Soient $\mathfrak D_f \subset \mathcal{O}_Y$ la différente et
$R \subset Y$ le sous-schéma fermé défini par $\mathfrak D_f$.
On a alors
\begin{enumerate}
\item $\mathfrak D_f$ est un idéal localement principal ;
\item $R$ est un sous-schéma fermé localement principal ;
\item $\mathfrak D_f$ coïncide, localement sur les ouverts affines, avec la différente de Noether ;
\item la formation de $R$ commute au changement de base ;
\item si $f$ est fini, la norme de $R$ est le discriminant de $f$ ;
\item si $f$ est étale aux points associés de $Y$, alors
$R$ est un diviseur de Cartier effectif et $\omega_{Y/X} = \mathcal{O}_Y(R)$.
\end{enumerate}
Cela résulte des lemmes \ref{discriminant-lemma-flat-gorenstein-agree-noether},
\ref{discriminant-lemma-base-change-different} et
\ref{discriminant-lemma-norm-different-is-discriminant}.
\end{remark}

\begin{remark}
\label{discriminant-remark-collect-results-qf-gorenstein-two}
Soit $S$ un schéma noethérien muni d'un complexe dualisant
$\omega_S^\bullet$. Soit $f : Y \to X$ un morphisme quasi-fini de Gorenstein
de schémas compactifiables sur $S$. Supposons de plus
$Y$ et $X$ de Cohen-Macaulay et $f$ étale aux points génériques
de $Y$. On peut alors combiner
Dualité des schémas, remarque
\ref{duality-remark-CM-morphism-compare-dualizing} et la
remarque \ref{discriminant-remark-collect-results-qf-gorenstein}
pour obtenir un isomorphisme canonique
$$
\omega_Y = f^*\omega_X \otimes_{\mathcal{O}_Y} \omega_{Y/X} =
f^*\omega_X \otimes_{\mathcal{O}_Y} \mathcal{O}_Y(R)
$$
de $\mathcal{O}_Y$-modules. Si $f$ est en outre fini,
l'isomorphisme $\mathcal{O}_Y(R) = \omega_{Y/X}$ provient
de la section globale $\tau_{Y/X} \in H^0(Y, \omega_{Y/X})$
qui correspond par dualité à l'application
$\text{Trace}_f : f_*\mathcal{O}_Y \to \mathcal{O}_X$, voir le
lemme \ref{discriminant-lemma-compare-trace}.
\end{remark}









\bibliography{my}
\bibliographystyle{amsalpha}
\end{document}
